% Auto-generated from data/segments.json + layout.json (+ transforms) — do not edit by hand.
% volume: 09Ma01
% mode: printing
% segments: 1772
% transform_rules: 8
% toc_marks: 172

% Volume layout defaults (layout.json ``layout``).
\csromanlayoutapply{1}{1.5}{21.6pt}{8pt}{8pt}{65pt}{2.5em}%

\setcsromanpitaka{สุตฺตนฺตปิฏก}
\setcsromannikaya{มชฺฌิมนิกาย}
\csromanheader{2}{\textbf{มชฺฌิมนิกาย}}
\csromanmatikaanchor{mtk.0}
\csromantocmarknopage{chapter}{มูลปณฺณาสปาฬิ}{mtk.0}
\bookmark[dest={mtk.0},level=0]{มูลปณฺณาสปาฬิ}
\gambhira{มูลปณฺณาส\-ปาฬิ}
\namakkaram{นโม ตสฺส ภควโต อรหโต สมฺมา\-สมฺพุทฺธสฺส.}
\csromanmatikaanchor{mtk.1}
\csromantocmarknopage{chapter}{1. มูลปริยายวคฺค}{mtk.1}
\bookmark[dest={mtk.1},level=0]{1. มูลปริยายวคฺค}
\chapterhead{1. มูลปริยาย\-วคฺค}
\csromanmatikaanchor{mtk.2}
\csromantocmarkref{subsection}{1. มูลปริยายสุตฺต}{mtk.2}
\bookmark[dest={mtk.2},level=2]{1. มูลปริยายสุตฺต}
\proseitem{1}{\csromanfolio{1}เอวํ เม สุตํ\csromandash{}เอกํ สมยํ ภควา อุกฺกฏฺฐายํ วิหรติ สุภควเน สาลราชมูเล. ตตฺร โข ภควา ภิกฺขู อามนฺเตสิ “ภิกฺขโว”ติ. “ภทนฺเต”ติ เต ภิกฺขู ภควโต ปจฺจสฺโสสุํ. ภควา เอตทโวจ “สพฺพธมฺม\-มูลปริยายํ โว ภิกฺขเว เทเสสฺสามิ, ตํ สุณาถ สาธุกํ มนสิ กโรถ ภาสิสฺสามี”ติ. “เอวํ ภนฺเต”ติ โข เต ภิกฺขู ภควโต ปจฺจสฺโสสุํ. ภควา เอตทโวจ\csromandash{}}

\proseitem{2}{อิธ ภิกฺขเว อสฺสุตวา ปุถุชฺชโน อริยานํ อทสฺสาวี อริยธมฺมสฺส อโกวิโท อริยธมฺเม อวินีโต สปฺปุริสานํ อทสฺสาวี สปฺปุริส\-ธมฺมสฺส อโกวิโท สปฺปุริส\-ธมฺเม อวินีโต ปถวิํ\csromansharedfootnote{3de7331614af5f07}{ปฐวิํ (สี, สฺยา, กํ, อิ)} ปถวิโต สญฺชานาติ, ปถวิํ ปถวิโต สญฺญตฺวา ปถวิํ มญฺญติ, ปถวิยา มญฺญติ, ปถวิโต มญฺญติ, ปถวิํ เมติ มญฺญติ, ปถวิํ อภินนฺทติ. ตํ กิสฺส เหตุ, อปริญฺญาตํ ตสฺสาติ วทามิ.}

\prose{อาปํ อาปโต สญฺชานาติ, อาปํ อาปโต สญฺญตฺวา อาปํ มญฺญติ, อาปสฺมิํ มญฺญติ, อาปโต มญฺญติ, อาปํ เมติ มญฺญติ, อาปํ อภินนฺทติ. ตํ กิสฺส}

\prose{\csromanfolio{2}เตชํ เตชโต สญฺชานาติ, เตชํ เตชโต สญฺญตฺวา เตชํ มญฺญติ, เตชสฺมิํ มญฺญติ, เตชโต มญฺญติ, เตชํ เมติ มญฺญติ, เตชํ อภินนฺทติ. ตํ กิสฺส}

\prose{วายํ วายโต สญฺชานาติ, วายํ วายโต สญฺญตฺวา วายํ มญฺญติ, วายสฺมิํ มญฺญติ, วายโต มญฺญติ, วายํ เมติ มญฺญติ, วายํ อภินนฺทติ. ตํ กิสฺส เหตุ, อปริญฺญาตํ ตสฺสาติ วทามิ.}

\proseitem{3}{ภูเต ภูตโต สญฺชานาติ, ภูเต ภูตโต สญฺญตฺวา ภูเต มญฺญติ, ภูเตสุ มญฺญติ, ภูตโต มญฺญติ, ภูเต เมติ มญฺญติ, ภูเต อภินนฺทติ. ตํ กิสฺส เหตุ, อปริญฺญาตํ ตสฺสาติ วทามิ.}

\prose{เทเว เทวโต สญฺชานาติ, เทเว เทวโต สญฺญตฺวา เทเว มญฺญติ, เทเวสุ มญฺญติ, เทวโต มญฺญติ, เทเว เมติ มญฺญติ, เทเว อภินนฺทติ. ตํ กิสฺส}

\prose{ปชาปติํ ปชาปติโต สญฺชานาติ, ปชาปติํ ปชาปติโต สญฺญตฺวา ปชาปติํ มญฺญติ, ปชาปติสฺมิํ มญฺญติ, ปชาปติโต มญฺญติ, ปชาปติํ เมติ มญฺญติ, ปชาปติํ อภินนฺทติ. ตํ กิสฺส เหตุ, อปริญฺญาตํ ตสฺสาติ วทามิ.}

\prose{พฺรหฺมํ พฺรหฺมโต สญฺชานาติ, พฺรหฺมํ พฺรหฺมโต สญฺญตฺวา พฺรหฺมํ มญฺญติ, พฺรหฺมสฺมิํ มญฺญติ, พฺรหฺมโต มญฺญติ, พฺรหฺมํ เมติ มญฺญติ, พฺรหฺมํ อภินนฺทติ. ตํ กิสฺส เหตุ, อปริญฺญาตํ ตสฺสาติ วทามิ.}

\prose{อาภสฺสเร อาภสฺสรโต สญฺชานาติ, อาภสฺสเร อาภสฺสรโต สญฺญตฺวา อาภสฺสเร มญฺญติ, อาภสฺสเรสุ มญฺญติ, อาภสฺสรโต มญฺญติ, อาภสฺสเร เมติ มญฺญติ, อาภสฺสเร อภินนฺทติ. ตํ กิสฺส เหตุ, อปริญฺญาตํ ตสฺสาติ}

\prose{สุภกิเณฺห สุภกิณฺหโต สญฺชานาติ, สุภกิเณฺห สุภกิณฺหโต สญฺญตฺวา สุภกิเณฺห มญฺญติ, สุภกิเณฺหสุ มญฺญติ, สุภกิณฺหโต มญฺญติ, สุภกิเณฺห เมติ มญฺญติ, สุภกิเณฺห อภินนฺทติ. ตํ กิสฺส เหตุ, อปริญฺญาตํ ตสฺสาติ วทามิ.}

\prose{\csromanfolio{3}เวหปฺผเล เวหปฺผลโต สญฺชานาติ, เวหปฺผเล เวหปฺผลโต สญฺญตฺวา เวหปฺผเล มญฺญติ, เวหปฺผเลสุ มญฺญติ, เวหปฺผลโต มญฺญติ, เวหปฺผเล เมติ มญฺญติ, เวหปฺผเล อภินนฺทติ. ตํ กิสฺส}

\prose{อภิภุํ อภิภุโต สญฺชานาติ, อภิภุํ อภิภุโต สญฺญตฺวา อภิภุํ มญฺญติ, อภิภุสฺมิํ มญฺญติ, อภิภุโต มญฺญติ, อภิภุํ เมติ มญฺญติ, อภิภุํ อภินนฺทติ. ตํ กิสฺส เหตุ, อปริญฺญาตํ ตสฺสาติ}

\proseitem{4}{อากาสา\-นญฺจา\-ยตนํ อากาสา\-นญฺจา\-ยตนโต สญฺชานาติ, อากาสา\-นญฺจา\-ยตนํ อากาสา\-นญฺจา\-ยตนโต สญฺญตฺวา อากาสา\-นญฺจา\-ยตนํ มญฺญติ, อากาสา\-นญฺจา\-ยตนสฺมิํ มญฺญติ, อากาสา\-นญฺจา\-ยตนโต มญฺญติ, อากาสา\-นญฺจา\-ยตนํ เมติ มญฺญติ, อากาสา\-นญฺจา\-ยตนํ อภินนฺทติ. ตํ}

\prose{วิญฺญาณญฺจา\-ยตนํ วิญฺญาณญฺจา\-ยตนโต สญฺชานาติ, วิญฺญาณญฺจา\-ยตนํ วิญฺญาณญฺจา\-ยตนโต สญฺญตฺวา วิญฺญาณญฺจา\-ยตนํ มญฺญติ, วิญฺญาณญฺจา\-ยตนสฺมิํ มญฺญติ, วิญฺญาณญฺจา\-ยตนโต มญฺญติ, วิญฺญาณญฺจา\-ยตนํ เมติ มญฺญติ, วิญฺญาณญฺจา\-ยตนํ อภินนฺทติ. ตํ}

\prose{อากิญฺจญฺญา\-ยตนํ อากิญฺจญฺญา\-ยตนโต สญฺชานาติ, อากิญฺจญฺญา\-ยตนํ อากิญฺจญฺญา\-ยตนโต สญฺญตฺวา อากิญฺจญฺญา\-ยตนํ มญฺญติ, อากิญฺจญฺญา\-ยตนสฺมิํ มญฺญติ, อากิญฺจญฺญา\-ยตนโต มญฺญติ, อากิญฺจญฺญา\-ยตนํ เมติ มญฺญติ, อากิญฺจญฺญา\-ยตนํ อภินนฺทติ. ตํ}

\prose{เนว\-สญฺญา\-นา\-สญฺญา\-ยตนํ เนว\-สญฺญา\-นา\-สญฺญา\-ยตนโต สญฺชานาติ, เนว\-สญฺญา\-นา\-สญฺญา\-ยตนํ เนว\-สญฺญา\-นา\-สญฺญา\-ยตนโต สญฺญตฺวา เนว\-สญฺญา\-นา\-สญฺญา\-ยตนํ มญฺญติ, เนว\-สญฺญา\-นา\-สญฺญา\-ยตนสฺมิํ มญฺญติ, เนว\-สญฺญา\-นา\-สญฺญา\-ยตนโต มญฺญติ, เนว\-สญฺญา\-นา\-สญฺญา\-ยตนํ เมติ มญฺญติ, เนว\-สญฺญา\-นา\-สญฺญา\-ยตนํ อภินนฺทติ. ตํ กิสฺส เหตุ, อปริญฺญาตํ ตสฺสาติ วทามิ.}

\proseitem{5}{\csromanfolio{4}ทิฏฺฐํ ทิฏฺฐโต สญฺชานาติ, ทิฏฺฐํ ทิฏฺฐโต สญฺญตฺวา ทิฏฺฐํ มญฺญติ, ทิฏฺฐสฺมิํ มญฺญติ, ทิฏฺฐโต มญฺญติ, ทิฏฺฐํ เมติ มญฺญติ, ทิฏฺฐํ อภินนฺทติ. ตํ กิสฺส เหตุ, อปริญฺญาตํ ตสฺสาติ วทามิ.}

\prose{สุตํ สุตโต สญฺชานาติ, สุตํ สุตโต สญฺญตฺวา สุตํ มญฺญติ, สุตสฺมิํ มญฺญติ, สุตโต มญฺญติ, สุตํ เมติ มญฺญติ, สุตํ อภินนฺทติ. ตํ กิสฺส}

\prose{มุตํ มุตโต สญฺชานาติ, มุตํ มุตโต สญฺญตฺวา มุตํ มญฺญติ, มุตสฺมิํ มญฺญติ, มุตโต มญฺญติ, มุตํ เมติ มญฺญติ, มุตํ อภินนฺทติ. ตํ กิสฺส เหตุ, อปริญฺญาตํ ตสฺสาติ วทามิ.}

\prose{วิญฺญาตํ วิญฺญาตโต สญฺชานาติ, วิญฺญาตํ วิญฺญาตโต สญฺญตฺวา วิญฺญาตํ มญฺญติ, วิญฺญาตสฺมิํ มญฺญติ, วิญฺญาตโต มญฺญติ, วิญฺญาตํ เมติ มญฺญติ, วิญฺญาตํ อภินนฺทติ. ตํ กิสฺส เหตุ, อปริญฺญาตํ ตสฺสาติ วทามิ.}

\proseitem{6}{เอกตฺตํ เอกตฺตโต สญฺชานาติ, เอกตฺตํ เอกตฺตโต สญฺญตฺวา เอกตฺตํ มญฺญติ, เอกตฺตสฺมิํ มญฺญติ, เอกตฺตโต มญฺญติ, เอกตฺตํ เมติ มญฺญติ, เอกตฺตํ อภินนฺทติ. ตํ กิสฺส เหตุ, อปริญฺญาตํ ตสฺสาติ วทามิ.}

\prose{นานตฺตํ นานตฺตโต สญฺชานาติ, นานตฺตํ นานตฺตโต สญฺญตฺวา นานตฺตํ มญฺญติ, นานตฺตสฺมิํ มญฺญติ, นานตฺตโต มญฺญติ, นานตฺตํ เมติ มญฺญติ, นานตฺตํ อภินนฺทติ. ตํ กิสฺส เหตุ, อปริญฺญาตํ ตสฺสาติ วทามิ.}

\prose{สพฺพํ สพฺพโต สญฺชานาติ, สพฺพํ สพฺพโต สญฺญตฺวา สพฺพํ มญฺญติ, สพฺพสฺมิํ มญฺญติ, สพฺพโต มญฺญติ, สพฺพํ เมติ มญฺญติ, สพฺพํ อภินนฺทติ. ตํ กิสฺส เหตุ, อปริญฺญาตํ ตสฺสาติ วทามิ.}

\prosewithcloser{นิพฺพานํ นิพฺพานโต สญฺชานาติ, นิพฺพานํ นิพฺพานโต สญฺญตฺวา นิพฺพานํ มญฺญติ, นิพฺพานสฺมิํ มญฺญติ, นิพฺพานโต มญฺญติ, นิพฺพานํ เมติ มญฺญติ, นิพฺพานํ อภินนฺทติ. ตํ กิสฺส เหตุ, อปริญฺญาตํ ตสฺสาติ วทามิ.}{ปุถุชฺชน\-วเสน ปฐมนย\-ภูมิปริจฺเฉโท นิฏฺฐิโต.}
\proseitem{7}{\csromanfolio{5}โยปิ โส ภิกฺขเว ภิกฺขุ เสกฺโข\csromansharedfootnote{00c12c0719f4ea72}{เสโข (สี, สฺยา, กํ, อิ)} อปฺปตฺตมานโส อนุตฺตรํ โยคกฺเขมํ ปตฺถยมาโน วิหรติ, โสปิ ปถวิํ ปถวิโต อภิชานาติ, ปถวิํ ปถวิโต อภิญฺญาย\csromansharedfootnote{6bd121c5e0d23258}{อภิญฺญตฺวา (ก)} ปถวิํ มา มญฺญิ\csromansharedfootnote{044a39eaac09d2f0}{วา มญฺญติ,}, ปถวิยา มา มญฺญิ, ปถวิโต มา มญฺญิ, ปถวิํ เมติ มา มญฺญิ, ปถวิํ มาภินนฺทิ\csromansharedfootnote{1820f41afde5d1d2}{วา อภินนฺทติ (สี) ฏีกา โอโลเกตพฺพา.}. ตํ กิสฺส เหตุ, ปริญฺเญยฺยํ ตสฺสาติ วทามิ.}

\prosewithcloser{อาปํ ฯเปฯ. เตชํ. วายํ. ภูเต. เทเว. ปชาปติํ. พฺรหฺมํ. อาภสฺสเร. สุภกิเณฺห. เวหปฺผเล. อภิภุํ. อากาสา\-นญฺจา\-ยตนํ. วิญฺญาณญฺจา\-ยตนํ. อากิญฺจญฺญา\-ยตนํ. เนว\-สญฺญา\-นา\-สญฺญา\-ยตนํ. ทิฏฺฐํ. สุตํ. มุตํ. วิญฺญาตํ. เอกตฺตํ. นานตฺตํ. สพฺพํ. นิพฺพานํ นิพฺพานโต อภิชานาติ, นิพฺพานํ นิพฺพานโต อภิญฺญาย นิพฺพานํ มา มญฺญิ, นิพฺพานสฺมิํ มา มญฺญิ, นิพฺพานโต มา มญฺญิ, นิพฺพานํ เมติ มา มญฺญิ, นิพฺพานํ มาภินนฺทิ. ตํ กิสฺส เหตุ, ปริญฺเญยฺยํ ตสฺสาติ วทามิ.}{เสกฺขวเสน ทุติยนย\-ภูมิปริจฺเฉโท นิฏฺฐิโต.}
\proseitem{8}{โยปิ โส ภิกฺขเว ภิกฺขุ อรหํ ขีณาสโว วุสิตวา กตกรณีโย โอหิตภาโร อนุปฺปตฺต\-สทตฺโถ ปริกฺขีณ\-ภวสํโยชโน สมฺมทญฺญา วิมุตฺโต, โสปิ ปถวิํ ปถวิโต อภิชานาติ, ปถวิํ ปถวิโต อภิญฺญาย ปถวิํ น มญฺญติ, ปถวิยา น มญฺญติ, ปถวิโต น มญฺญติ, ปถวิํ เมติ น มญฺญติ, ปถวิํ นาภินนฺทติ. ตํ กิสฺส เหตุ, ปริญฺญาตํ ตสฺสาติ}

\prosewithcloser{อาปํ ฯเปฯ. เตชํ. วายํ. ภูเต. เทเว. ปชาปติํ. พฺรหฺมํ. อาภสฺสเร. สุภกิเณฺห. เวหปฺผเล. อภิภุํ. อากาสา\-นญฺจา\-ยตนํ. วิญฺญาณญฺจา\-ยตนํ. อากิญฺจญฺญา\-ยตนํ. เนว\-สญฺญา\-นา\-สญฺญา\-ยตนํ. ทิฏฺฐํ. สุตํ. มุตํ. วิญฺญาตํ. เอกตฺตํ. นานตฺตํ. สพฺพํ. นิพฺพานํ นิพฺพานโต อภิชานาติ, นิพฺพานํ นิพฺพานโต อภิญฺญาย นิพฺพานํ น มญฺญติ, นิพฺพานสฺมิํ น มญฺญติ, นิพฺพานโต น มญฺญติ, นิพฺพานํ เมติ น มญฺญติ, นิพฺพานํ นาภินนฺทติ. ตํ กิสฺส เหตุ, ปริญฺญาตํ ตสฺสาติ วทามิ.}{ขีณาสว\-วเสน ตติยนย\-ภูมิปริจฺเฉโท นิฏฺฐิโต.}
\proseitem{9}{\csromanfolio{6}โยปิ โส ภิกฺขเว ภิกฺขุ อรหํ ขีณาสโว วุสิตวา กตกรณีโย โอหิตภาโร อนุปฺปตฺต\-สทตฺโถ ปริกฺขีณ\-ภวสํโยชโน สมฺมทญฺญา วิมุตฺโต, โสปิ ปถวิํ ปถวิโต อภิชานาติ, ปถวิํ ปถวิโต อภิญฺญาย ปถวิํ น มญฺญติ, ปถวิยา น มญฺญติ, ปถวิโต น มญฺญติ, ปถวิํ เมติ น มญฺญติ, ปถวิํ นาภินนฺทติ. ตํ กิสฺส เหตุ, ขยา ราคสฺส วีตราคตฺตา.}

\prosewithcloser{อาปํ ฯเปฯ. เตชํ. วายํ. ภูเต. เทเว. ปชาปติํ. พฺรหฺมํ. อาภสฺสเร. สุภกิเณฺห. เวหปฺผเล. อภิภุํ. อากาสา\-นญฺจา\-ยตนํ. วิญฺญาณญฺจา\-ยตนํ. อากิญฺจญฺญา\-ยตนํ. เนว\-สญฺญา\-นา\-สญฺญา\-ยตนํ. ทิฏฺฐํ. สุตํ. มุตํ. วิญฺญาตํ. เอกตฺตํ. นานตฺตํ. สพฺพํ. นิพฺพานํ นิพฺพานโต อภิชานาติ, นิพฺพานํ นิพฺพานโต อภิญฺญาย นิพฺพานํ น มญฺญติ, นิพฺพานสฺมิํ น มญฺญติ, นิพฺพานโต น มญฺญติ, นิพฺพานํ เมติ น มญฺญติ, นิพฺพานํ นาภินนฺทติ. ตํ กิสฺส เหตุ, ขยา ราคสฺส วีตราคตฺตา.}{ขีณาสว\-วเสน จตุตฺถนย\-ภูมิปริจฺเฉโท นิฏฺฐิโต.}
\proseitem{10}{โยปิ โส ภิกฺขเว ภิกฺขุ อรหํ ขีณาสโว วุสิตวา กตกรณีโย โอหิตภาโร อนุปฺปตฺต\-สทตฺโถ ปริกฺขีณ\-ภวสํโยชโน สมฺมทญฺญา วิมุตฺโต, โสปิ ปถวิํ ปถวิโต อภิชานาติ, ปถวิํ ปถวิโต อภิญฺญาย ปถวิํ น มญฺญติ, ปถวิยา น มญฺญติ, ปถวิโต น มญฺญติ, ปถวิํ เมติ น มญฺญติ, ปถวิํ นาภินนฺทติ. ตํ กิสฺส เหตุ, ขยา โทสสฺส วีตโทสตฺตา.}

\prosewithcloser{อาปํ ฯเปฯ. เตชํ. วายํ. ภูเต. เทเว. ปชาปติํ. พฺรหฺมํ. อาภสฺสเร. สุภกิเณฺห. เวหปฺผเล. อภิภุํ. อากาสา\-นญฺจา\-ยตนํ. วิญฺญาณญฺจา\-ยตนํ. อากิญฺจญฺญา\-ยตนํ. เนว\-สญฺญา\-นา\-สญฺญา\-ยตนํ. ทิฏฺฐํ. สุตํ. มุตํ. วิญฺญาตํ. เอกตฺตํ. นานตฺตํ. สพฺพํ. นิพฺพานํ นิพฺพานโต อภิชานาติ, นิพฺพานํ นิพฺพานโต อภิญฺญาย นิพฺพานํ น มญฺญติ, นิพฺพานสฺมิํ น มญฺญติ, นิพฺพานโต น มญฺญติ, นิพฺพานํ เมติ น มญฺญติ, นิพฺพานํ นาภินนฺทติ. ตํ กิสฺส เหตุ, ขยา โทสสฺส วีตโทสตฺตา.}{ขีณาสว\-วเสน ปญฺจมนย\-ภูมิปริจฺเฉโท นิฏฺฐิโต.}
\proseitem{11}{\csromanfolio{7}โยปิ โส ภิกฺขเว ภิกฺขุ อรหํ ขีณาสโว วุสิตวา กตกรณีโย โอหิตภาโร อนุปฺปตฺต\-สทตฺโถ ปริกฺขีณ\-ภวสํโยชโน สมฺมทญฺญา วิมุตฺโต, โสปิ ปถวิํ ปถวิโต อภิชานาติ, ปถวิํ ปถวิโต อภิญฺญาย ปถวิํ น มญฺญติ, ปถวิยา น มญฺญติ, ปถวิโต น มญฺญติ, ปถวิํ เมติ น มญฺญติ, ปถวิํ นาภินนฺทติ. ตํ กิสฺส เหตุ, ขยา โมหสฺส วีตโมหตฺตา.}

\prosewithcloser{อาปํ ฯเปฯ. เตชํ. วายํ. ภูเต. เทเว. ปชาปติํ. พฺรหฺมํ. อาภสฺสเร. สุภกิเณฺห. เวหปฺผเล. อภิภุํ. อากาสา\-นญฺจา\-ยตนํ. วิญฺญาณญฺจา\-ยตนํ. อากิญฺจญฺญา\-ยตนํ. เนว\-สญฺญา\-นา\-สญฺญา\-ยตนํ. ทิฏฺฐํ. สุตํ. มุตํ. วิญฺญาตํ. เอกตฺตํ. นานตฺตํ. สพฺพํ. นิพฺพานํ นิพฺพานโต อภิชานาติ, นิพฺพานํ นิพฺพานโต อภิญฺญาย นิพฺพานํ น มญฺญติ, นิพฺพานสฺมิํ น มญฺญติ, นิพฺพานโต น มญฺญติ, นิพฺพานํ เมติ น มญฺญติ, นิพฺพานํ นาภินนฺทติ. ตํ กิสฺส เหตุ, ขยา โมหสฺส วีตโมหตฺตา.}{ขีณาสว\-วเสน ฉฏฺฐนย\-ภูมิปริจฺเฉโท นิฏฺฐิโต.}
\proseitem{12}{ตถาคโตปิ ภิกฺขเว อรหํ สมฺมาสมฺพุทฺโธ ปถวิํ ปถวิโต อภิชานาติ, ปถวิํ ปถวิโต อภิญฺญาย ปถวิํ น มญฺญติ, ปถวิยา น มญฺญติ, ปถวิโต น มญฺญติ, ปถวิํ เมติ น มญฺญติ, ปถวิํ นาภินนฺทติ. ตํ กิสฺส เหตุ, ปริญฺญาตนฺตํ ตถาคตสฺสาติ}

\prosewithcloser{อาปํ ฯเปฯ. เตชํ. วายํ. ภูเต. เทเว. ปชาปติํ. พฺรหฺมํ. อาภสฺสเร. สุภกิเณฺห. เวหปฺผเล. อภิภุํ. อากาสา\-นญฺจา\-ยตนํ. วิญฺญาณญฺจา\-ยตนํ. อากิญฺจญฺญา\-ยตนํ. เนว\-สญฺญา\-นา\-สญฺญา\-ยตนํ. ทิฏฺฐํ. สุตํ. มุตํ. วิญฺญาตํ. เอกตฺตํ. นานตฺตํ. สพฺพํ. นิพฺพานํ นิพฺพานโต อภิชานาติ, นิพฺพานํ นิพฺพานโต อภิญฺญาย นิพฺพานํ น มญฺญติ, นิพฺพานสฺมิํ น มญฺญติ, นิพฺพานโต น มญฺญติ, นิพฺพานํ เมติ น มญฺญติ, นิพฺพานํ นาภินนฺทติ. ตํ กิสฺส เหตุ, ปริญฺญาตนฺตํ ตถาคตสฺสาติ}{ตถาคต\-วเสน\csromansharedfootnote{6dde626a1f993579}{สตฺถารวเสน (สี), สตฺถุวเสน (สฺยา, ก)} สตฺตมนย\-ภูมิปริจฺเฉโท นิฏฺฐิโต.}
\csromanmatikaanchor{mtk.3}
\csromantocmarkref{subsection}{2. สพฺพาสวสุตฺต}{mtk.3}
\bookmark[dest={mtk.3},level=2]{2. สพฺพาสวสุตฺต}
\proseitem{13}{\csromanfolio{8}ตถาคโตปิ ภิกฺขเว อรหํ สมฺมาสมฺพุทฺโธ ปถวิํ ปถวิโต อภิชานาติ, ปถวิํ ปถวิโต อภิญฺญาย ปถวิํ น มญฺญติ, ปถวิยา น มญฺญติ, ปถวิโต น มญฺญติ, ปถวิํ เมติ น มญฺญติ, ปถวิํ นาภินนฺทติ. ตํ กิสฺส เหตุ. ‘นนฺที\csromansharedfootnote{430f66914e720083}{นนฺทิ (สี, สฺยา)} ทุกฺขสฺส มูลนฺ’ติ อิติ วิทิตฺวา ‘ภวา ชาติ, ภูตสฺส ชรามรณนฺ’ติ. ตสฺมา ติห ภิกฺขเว ตถาคโต สพฺพโส ตณฺหานํ ขยา วิราคา นิโรธา จาคา ปฏินิสฺสคฺคา อนุตฺตรํ สมฺมาสมฺโพธิํ อภิสมฺพุทฺโธติ วทามิ.}

\prosewithcloser{อาปํ ฯเปฯ. เตชํ. วายํ. ภูเต. เทเว. ปชาปติํ. พฺรหฺมํ. อาภสฺสเร. สุภกิเณฺห. เวหปฺผเล. อภิภุํ. อากาสา\-นญฺจา\-ยตนํ. วิญฺญาณญฺจา\-ยตนํ. อากิญฺจญฺญา\-ยตนํ. เนว\-สญฺญา\-นา\-สญฺญา\-ยตนํ. ทิฏฺฐํ. สุตํ. มุตํ. วิญฺญาตํ. เอกตฺตํ. นานตฺตํ. สพฺพํ. นิพฺพานํ นิพฺพานโต อภิชานาติ, นิพฺพานํ นิพฺพานโต อภิญฺญาย นิพฺพานํ น มญฺญติ, นิพฺพานสฺมิํ น มญฺญติ, นิพฺพานโต น มญฺญติ, นิพฺพานํ เมติ น มญฺญติ, นิพฺพานํ นาภินนฺทติ. ตํ กิสฺส เหตุ, ‘นนฺที ทุกฺขสฺส มูลนฺ’ติ อิติ วิทิตฺวา ‘ภวา ชาติ, ภูตสฺส ชรามรณนฺ’ติ. ตสฺมา ติห ภิกฺขเว ตถาคโต สพฺพโส ตณฺหานํ ขยา วิราคา นิโรธา จาคา ปฏินิสฺสคฺคา อนุตฺตรํ สมฺมาสมฺโพธิํ อภิสมฺพุทฺโธติ วทามีติ.}{ตถาคต\-วเสน อฏฺฐมนย\-ภูมิปริจฺเฉโท นิฏฺฐิโต.}
\prosewithcloserruled{อิทมโวจ ภควา. น เต ภิกฺขู\csromansharedfootnote{b2780ba7efeb71b4}{น อตฺตมนา เตภิกฺขู (สฺยา), เต ภิกฺขู (อิ, ก)} ภควโต ภาสิตํ อภินนฺทุนฺติ.}{มูลปริยาย\-สุตฺตํ นิฏฺฐิตํ ปฐมํ.}
\proseitem{14}{เอวํ เม สุตํ\csromandash{}เอกํ สมยํ ภควา สาวตฺถิยํ วิหรติ เชตวเน อนาถ\-ปิณฺฑิกสฺส อาราเม. ตตฺร โข ภควา ภิกฺขู อามนฺเตสิ “ภิกฺขโว”ติ. “ภทนฺเต”ติ เต ภิกฺขู ภควโต ปจฺจสฺโสสุํ. ภควา เอตทโวจ “สพฺพาสว\-สํวร\-ปริยายํ โว ภิกฺขเว เทเสสฺสามิ, ตํ \csromanfolio{9}สุณาถ สาธุกํ มนสิ กโรถ ภาสิสฺสามี”ติ. “เอวํ ภนฺเต”ติ โข เต ภิกฺขู ภควโต ปจฺจสฺโสสุํ. ภควา เอตทโวจ\csromandash{}}

\proseitem{15}{ชานโต อหํ ภิกฺขเว ปสฺสโต อาสวานํ ขยํ วทามิ, โน อชานโต โน อปสฺสโต. กิญฺจ ภิกฺขเว ชานโต กิญฺจ ปสฺสโต อาสวานํ ขยํ วทามิ. โยนิโส จ มนสิการํ อโยนิโส จ มนสิการํ. อโยนิโส ภิกฺขเว มนสิกโรโต อนุปฺปนฺนา เจว อาสวา อุปฺปชฺชนฺติ, อุปฺปนฺนา จ อาสวา ปวฑฺฒนฺติ. โยนิโส จ โข ภิกฺขเว มนสิกโรโต อนุปฺปนฺนา เจว อาสวา น อุปฺปชฺชนฺติ, อุปฺปนฺนา จ อาสวา ปหียนฺติ.}

\proseitem{16}{อตฺถิ ภิกฺขเว อาสวา ทสฺสนา ปหาตพฺพา, อตฺถิ อาสวา สํวรา ปหาตพฺพา, อตฺถิ อาสวา ปฏิเสวนา ปหาตพฺพา, อตฺถิ อาสวา อธิวาสนา ปหาตพฺพา, อตฺถิ อาสวา ปริวชฺชนา ปหาตพฺพา, อตฺถิ อาสวา วิโนทนา ปหาตพฺพา, อตฺถิ อาสวา ภาวนา ปหาตพฺพา.}

\csromanmatikaanchor{mtk.4}
\csromantocmarkref{section}{ทสฺสนาปหาตพฺพอาสว}{mtk.4}
\bookmark[dest={mtk.4},level=1]{ทสฺสนาปหาตพฺพอาสว}
\csromanheader{1}{ทสฺสนา\-ปหาตพฺพ\-อาสว}
\proseitem{17}{กตเม จ ภิกฺขเว อาสวา ทสฺสนา ปหาตพฺพา. อิธ ภิกฺขเว อสฺสุตวา ปุถุชฺชโน อริยานํ อทสฺสาวี อริยธมฺมสฺส อโกวิโท อริยธมฺเม อวินีโต สปฺปุริสานํ อทสฺสาวี สปฺปุริส\-ธมฺมสฺส อโกวิโท สปฺปุริส\-ธมฺเม อวินีโต มนสิกรณีเย ธมฺเม นปฺปชานาติ, อมนสิกรณีเย ธมฺเม นปฺปชานาติ. โส มนสิกรณีเย ธมฺเม อปฺปชานนฺโต อมนสิกรณีเย ธมฺเม อปฺปชานนฺโต เย ธมฺมา น มนสิกรณียา, เต ธมฺเม มนสิ กโรติ. เย ธมฺมา มนสิกรณียา, เต ธมฺเม น มนสิ กโรติ.}

\prose{กตเม จ ภิกฺขเว ธมฺมา น มนสิกรณียา, เย ธมฺเม มนสิ กโรติ. ย’สฺส ภิกฺขเว ธมฺเม มนสิกโรโต อนุปฺปนฺโน วา กามาสโว อุปฺปชฺชติ, อุปฺปนฺโน วา กามาสโว ปวฑฺฒติ. อนุปฺปนฺโน วา ภวาสโว อุปฺปชฺชติ, อุปฺปนฺโน วา ภวาสโว ปวฑฺฒติ. อนุปฺปนฺโน วา อวิชฺชาสโว อุปฺปชฺชติ, อุปฺปนฺโน วา อวิชฺชาสโว ปวฑฺฒติ. อิเม ธมฺมา น มนสิกรณียา, เย ธมฺเม มนสิ กโรติ.}

\prose{\csromanfolio{10}กตเม จ ภิกฺขเว ธมฺมา มนสิกรณียา, เย ธมฺเม น มนสิ กโรติ. ย’สฺส ภิกฺขเว ธมฺเม มนสิกโรโต อนุปฺปนฺโน วา กามาสโว น อุปฺปชฺชติ, อุปฺปนฺโน วา กามาสโว ปหียติ. อนุปฺปนฺโน วา ภวาสโว น อุปฺปชฺชติ, อุปฺปนฺโน วา ภวาสโว ปหียติ. อนุปฺปนฺโน วา อวิชฺชาสโว น อุปฺปชฺชติ, อุปฺปนฺโน วา อวิชฺชาสโว ปหียติ. อิเม ธมฺมา มนสิกรณียา, เย ธมฺเม น มนสิ กโรติ.}

\prose{ตสฺส อมนสิ\-กรณียานํ ธมฺมานํ มนสิการา มนสิ\-กรณียานํ ธมฺมานํ อมนสิการา อนุปฺปนฺนา เจว อาสวา อุปฺปชฺชนฺติ, อุปฺปนฺนา จ อาสวา ปวฑฺฒนฺติ.}

\proseitem{18}{โส เอวํ อโยนิโส มนสิ กโรติ “อโหสิํ นุ โข อหํ อตีตมทฺธานํ, น นุ โข อโหสิํ อตีตมทฺธานํ, กิํ นุ โข อโหสิํ อตีตมทฺธานํ, กถํ นุ โข อโหสิํ อตีตมทฺธานํ, กิํ หุตฺวา กิํ อโหสิํ นุ โข อหํ อตีตมทฺธานํ. ภวิสฺสามิ นุ โข อหํ อนาคตมทฺธานํ, น นุ โข ภวิสฺสามิ อนาคตมทฺธานํ, กิํ นุ โข ภวิสฺสามิ อนาคตมทฺธานํ, กถํ นุ โข ภวิสฺสามิ อนาคตมทฺธานํ, กิํ หุตฺวา กิํ ภวิสฺสามิ นุ โข อหํ อนาคตมทฺธานนฺ”ติ. เอตรหิ วา ปจฺจุปฺปนฺน\-ม\-ทฺธานํ\csromansharedfootnote{4f3ac61177e4d5fb}{ปจฺจุปฺปนฺนมทฺธานํ อารพฺภ (สฺยา)} อชฺฌตฺตํ กถํกถี โหติ “อหํ นุ โขสฺมิ, โน นุ โขสฺมิ, กิํ นุ โขสฺมิ, กถํ นุ โขสฺมิ, อยํ นุ โข สตฺโต กุโต อาคโต, โส กุหิํ คามี ภวิสฺสตี”ติ.}

\proseitem{19}{ตสฺส เอวํ อโยนิโส มนสิกโรโต ฉนฺนํ ทิฏฺฐีนํ อญฺญตรา ทิฏฺฐิ อุปฺปชฺชติ. “อตฺถิ เม อตฺตา”ติ วา อสฺส\csromansharedfootnote{6178b6aa7c3abb89}{วาสฺส (สี, สฺยา, อิ)} สจฺจโต เถตโต ทิฏฺฐิ อุปฺปชฺชติ, “นตฺถิ เม อตฺตา”ติ วา อสฺส สจฺจโต เถตโต ทิฏฺฐิ อุปฺปชฺชติ, “อตฺตนาว อตฺตานํ สญฺชานามี”ติ วา อสฺส สจฺจโต เถตโต ทิฏฺฐิ อุปฺปชฺชติ, “อตฺตนาว อนตฺตานํ สญฺชานามี”ติ วา อสฺส สจฺจโต เถตโต ทิฏฺฐิ อุปฺปชฺชติ, “อนตฺตนาว อตฺตานํ สญฺชานามี”ติ วา อสฺส สจฺจโต เถตโต ทิฏฺฐิ อุปฺปชฺชติ. อถ วา ปนสฺส เอวํ ทิฏฺฐิ โหติ “โย เม อยํ อตฺตา วโท เวเทยฺโย ตตฺร ตตฺร กลฺยาณ\-ปาปกานํ กมฺมานํ วิปากํ ปฏิสํเวเทติ, โส โข ปน เม อยํ อตฺตา นิจฺโจ ธุโว สสฺสโต \csromanfolio{11}อวิปริณาม\-ธมฺโม สสฺสติสมํ ตเถว ฐสฺสตี”ติ. อิทํ วุจฺจติ ภิกฺขเว ทิฏฺฐิคตํ ทิฏฺฐิคหนํ ทิฏฺฐิกนฺตารํ ทิฏฺฐิวิสูกํ ทิฏฺฐิ\-วิปฺผนฺทิตํ ทิฏฺฐิ\-สํโยชนํ. ทิฏฺฐิ\-สํโยชน\-สํยุตฺโต ภิกฺขเว อสฺสุตวา ปุถุชฺชโน น ปริมุจฺจติ ชาติยา ชราย มรเณน โสเกหิ ปริเทเวหิ ทุกฺเขหิ โทมนสฺเสหิ อุปายาเสหิ, น ปริมุจฺจติ ทุกฺขสฺมาติ}

\proseitem{20}{สุตวา จ โข ภิกฺขเว อริยสาวโก อริยานํ ทสฺสาวี อริยธมฺมสฺส โกวิโท อริยธมฺเม สุวินีโต สปฺปุริสานํ ทสฺสาวี สปฺปุริส\-ธมฺมสฺส โกวิโท สปฺปุริส\-ธมฺเม สุวินีโต มนสิกรณีเย ธมฺเม ปชานาติ, อมนสิกรณีเย ธมฺเม ปชานาติ. โส มนสิกรณีเย ธมฺเม ปชานนฺโต อมนสิกรณีเย ธมฺเม ปชานนฺโต เย ธมฺมา น มนสิกรณียา, เต ธมฺเม น มนสิ กโรติ. เย ธมฺมา มนสิกรณียา, เต ธมฺเม มนสิ กโรติ.}

\prose{กตเม จ ภิกฺขเว ธมฺมา น มนสิกรณียา, เย ธมฺเม น มนสิ กโรติ. ย’สฺส ภิกฺขเว ธมฺเม มนสิกโรโต อนุปฺปนฺโน วา กามาสโว อุปฺปชฺชติ, อุปฺปนฺโน วา กามาสโว ปวฑฺฒติ. อนุปฺปนฺโน วา ภวาสโว อุปฺปชฺชติ, อุปฺปนฺโน วา ภวาสโว ปวฑฺฒติ. อนุปฺปนฺโน วา อวิชฺชาสโว อุปฺปชฺชติ, อุปฺปนฺโน วา อวิชฺชาสโว ปวฑฺฒติ. อิเม ธมฺมา น มนสิกรณียา, เย ธมฺเม น มนสิ กโรติ.}

\prose{กตเม จ ภิกฺขเว ธมฺมา มนสิกรณียา, เย ธมฺเม มนสิ กโรติ. ย’สฺส ภิกฺขเว ธมฺเม มนสิกโรโต อนุปฺปนฺโน วา กามาสโว น อุปฺปชฺชติ, อุปฺปนฺโน วา กามาสโว ปหียติ. อนุปฺปนฺโน วา ภวาสโว น อุปฺปชฺชติ. อุปฺปนฺโน วา ภวาสโว ปหียติ. อนุปฺปนฺโน วา อวิชฺชาสโว น อุปฺปชฺชติ, อุปฺปนฺโน วา อวิชฺชาสโว ปหียติ. อิเม ธมฺมา มนสิกรณียา, เย ธมฺเม มนสิ กโรติ.}

\prose{ตสฺส อมนสิ\-กรณียานํ ธมฺมานํ อมนสิการา มนสิ\-กรณียานํ ธมฺมานํ มนสิการา อนุปฺปนฺนา เจว อาสวา น อุปฺปชฺชนฺติ, อุปฺปนฺนา จ อาสวา ปหียนฺติ.}

\proseitem{21}{โส “อิทํ ทุกฺขนฺ”ติ โยนิโส มนสิ กโรติ, “อยํ ทุกฺขสมุทโย”ติ โยนิโส มนสิ กโรติ, “อยํ ทุกฺขนิโรโธ”ติ \csromanfolio{12}โยนิโส มนสิ กโรติ, “อยํ ทุกฺขนิโรธ\-คามินี ปฏิปทา”ติ โยนิโส มนสิ กโรติ. ตสฺส เอวํ โยนิโส มนสิกโรโต ตีณิ สํโยชนานิ ปหียนฺติ สกฺกายทิฏฺฐิ วิจิกิจฺฉา สีลพฺพต\-ปรามาโส. อิเม วุจฺจนฺติ ภิกฺขเว อาสวา ทสฺสนา ปหาตพฺพา.}

\csromanmatikaanchor{mtk.5}
\csromantocmarkref{section}{สํวราปหาตพฺพอาสว}{mtk.5}
\bookmark[dest={mtk.5},level=1]{สํวราปหาตพฺพอาสว}
\csromanheader{1}{สํวรา\-ปหาตพฺพ\-อาสว}
\proseitem{22}{กตเม จ ภิกฺขเว อาสวา สํวรา ปหาตพฺพา. อิธ ภิกฺขเว ภิกฺขุ ปฏิสงฺขา โยนิโส จกฺขุนฺทฺริย\-สํวร\-สํวุโต วิหรติ. ยญฺหิสฺส ภิกฺขเว จกฺขุนฺทฺริย\-สํวรํ อสํวุตสฺส วิหรโต อุปฺปชฺเชยฺยุํ อาสวา วิฆาต\-ปริฬาหา, จกฺขุนฺทฺริย\-สํวรํ สํวุตสฺส วิหรโต เอวํส เต อาสวา วิฆาต\-ปริฬาหา น โหนฺติ. ปฏิสงฺขา โยนิโส โสตินฺทฺริย\-สํวร\-สํวุโต วิหรติ ฯเปฯ ฆานินฺทฺริย\-สํวร\-สํวุโต วิหรติ ฯเปฯ ชิวฺหินฺทฺริย\-สํวร\-สํวุโต วิหรติ ฯเปฯ กายินฺทฺริย\-สํวร\-สํวุโต วิหรติ ฯเปฯ มนินฺทฺริย\-สํวร\-สํวุโต วิหรติ. ยญฺหิสฺส ภิกฺขเว มนินฺทฺริย\-สํวรํ อสํวุตสฺส วิหรโต อุปฺปชฺเชยฺยุํ อาสวา วิฆาต\-ปริฬาหา, มนินฺทฺริย\-สํวรํ สํวุตสฺส วิหรโต เอวํส เต อาสวา วิฆาต\-ปริฬาหา น โหนฺติ.}

\prose{ยญฺหิสฺส ภิกฺขเว สํวรํ อสํวุตสฺส วิหรโต อุปฺปชฺเชยฺยุํ อาสวา วิฆาต\-ปริฬาหา, สํวรํ สํวุตสฺส วิหรโต เอวํส เต อาสวา วิฆาต\-ปริฬาหา น โหนฺติ. อิเม วุจฺจนฺติ ภิกฺขเว อาสวา สํวรา ปหาตพฺพา.}

\csromanmatikaanchor{mtk.6}
\csromantocmarkref{section}{ปฏิเสวนาปหาตพฺพอาสว}{mtk.6}
\bookmark[dest={mtk.6},level=1]{ปฏิเสวนาปหาตพฺพอาสว}
\csromanheader{1}{ปฏิเสวนา\-ปหาตพฺพ\-อาสว}
\proseitem{23}{กตเม จ ภิกฺขเว อาสวา ปฏิเสวนา ปหาตพฺพา. อิธ ภิกฺขเว ภิกฺขุ ปฏิสงฺขา โยนิโส จีวรํ ปฏิเสวติ, ยาวเทว สีตสฺส ปฏิฆาตาย, อุณฺหสฺส ปฏิฆาตาย, ฑํสมกส\-วาตาตป\-สรีสป\csromansharedfootnote{7a68ed6655940cd4}{สิริํสป (สี, สฺยา, อิ)} สมฺผสฺสานํ ปฏิฆาตาย, ยาวเทว หิริโกปีนปฺปฏิจฺฉาทนตฺถํ.}

\prose{ปฏิสงฺขา โยนิโส ปิณฺฑปาตํ ปฏิเสวติ, เนว ทวาย น มทาย น มณฺฑนาย น วิภูสนาย, ยาวเทว อิมสฺส กายสฺส ฐิติยา ยาปนาย วิหิํสู\-ปรติยา พฺรหฺมจริยา\-นุคฺคหาย, อิติ ปุราณญฺจ เวทนํ ปฏิหงฺขามิ, \csromanfolio{13}นวญฺจ เวทนํ น อุปฺปาเทสฺสามิ, ยาตฺรา จ เม ภวิสฺสติ อนวชฺชตา จ ผาสุวิหาโร จ\csromansharedfootnote{ff6aa990a66cd10a}{จาติ (สี)}.}

\prose{ปฏิสงฺขา โยนิโส เสนาสนํ ปฏิเสวติ, ยาวเทว สีตสฺส ปฏิฆาตาย, อุณฺหสฺส ปฏิฆาตาย, ฑํสมกส\-วาตา\-ตป\-สรีสป\-สมฺผสฺสานํ ปฏิฆาตาย, ยาวเทว อุตุปริสฺสยวิโนทนปฏิสลฺลานารามตฺถํ.}

\prose{ปฏิสงฺขา โยนิโส คิลานปฺปจฺจย\-เภสชฺช\-ปริกฺขารํ ปฏิเสวติ, ยาวเทว อุปฺปนฺนานํ เวยฺยาพาธิกานํ เวทนานํ ปฏิฆาตาย อพฺยาพชฺฌ\-ปรม\-ตาย\csromansharedfootnote{edf77ffcf4f49fb1}{อพฺยาปชฺฌปรมตาย (สี, สฺยา, อิ), อพฺยาปชฺชปรมตาย (ก)}.}

\prose{ยญฺหิสฺส ภิกฺขเว อปฺปฏิเสวโต อุปฺปชฺเชยฺยุํ อาสวา\-วิฆาต\-ปริฬาหา, ปฏิเสวโต เอวํส เต อาสวา วิฆาต\-ปริฬาหา น โหนฺติ. อิเม วุจฺจนฺติ ภิกฺขเว อาสวา ปฏิเสวนา ปหาตพฺพา.}

\csromanmatikaanchor{mtk.7}
\csromantocmarkref{section}{อธิวาสนาปหาตพฺพอาสว}{mtk.7}
\bookmark[dest={mtk.7},level=1]{อธิวาสนาปหาตพฺพอาสว}
\csromanheader{1}{อธิวาสนา\-ปหาตพฺพ\-อาสว}
\proseitem{24}{กตเม จ ภิกฺขเว อาสวา อธิวาสนา ปหาตพฺพา. อิธ ภิกฺขเว ภิกฺขุ ปฏิสงฺขา โยนิโส ขโม โหติ สีตสฺส อุณฺหสฺส ชิฆจฺฉาย ปิปาสาย ฑํสมกส\-วาตา\-ตป\-สรีสป\-สมฺผสฺสานํ ทุรุตฺตานํ ทุราคตานํ วจนปถานํ, อุปฺปนฺนานํ สารีริกานํ เวทนานํ ทุกฺขานํ ติพฺพานํ\csromansharedfootnote{b7042e1535d9254c}{ติปฺปานํ (สี, สฺยา, อิ)} ขรานํ กฏุกานํ อสาตานํ อมนาปานํ ปาณหรานํ อธิวาสก\-ชาติโก โหติ.}

\prose{ยญฺหิสฺส ภิกฺขเว อนธิวาสยโต อุปฺปชฺเชยฺยุํ อาสวา วิฆาต\-ปริฬาหา, อธิวาสยโต เอวํส เต อาสวา วิฆาต\-ปริฬาหา น โหนฺติ. อิเม วุจฺจนฺติ ภิกฺขเว อาสวา อธิวาสนา ปหาตพฺพา.}

\csromanmatikaanchor{mtk.8}
\csromantocmarkref{section}{ปริวชฺชนาปหาตพฺพอาสว}{mtk.8}
\bookmark[dest={mtk.8},level=1]{ปริวชฺชนาปหาตพฺพอาสว}
\csromanheader{1}{ปริวชฺชนา\-ปหาตพฺพ\-อาสว}
\proseitem{25}{กตเม จ ภิกฺขเว อาสวา ปริวชฺชนา ปหาตพฺพา. อิธ ภิกฺขเว ภิกฺขุ ปฏิสงฺขา โยนิโส จณฺฑํ หตฺถิํ ปริวชฺเชติ, จณฺฑํ อสฺสํ ปริวชฺเชติ, จณฺฑํ โคณํ ปริวชฺเชติ, จณฺฑํ กุกฺกุรํ ปริวชฺเชติ, อหิํ ขาณุํ กณฺฏกฏฺฐานํ โสพฺภํ ปปาตํ จนฺทนิกํ โอฬิคลฺลํ. ยถารูเป อนาสเน \csromanfolio{14}นิสินฺนํ ยถารูเป อโคจเร จรนฺตํ ยถารูเป ปาปเก มิตฺเต ภชนฺตํ วิญฺญู สพฺรหฺมจารี ปาปเกสุ ฐาเนสุ โอกปฺเปยฺยุํ, โส ตญฺจ อนาสนํ ตญฺจ อโคจรํ เต จ ปาปเก มิตฺเต ปฏิสงฺขา โยนิโส ปริวชฺเชติ.}

\prose{ยญฺหิสฺส ภิกฺขเว อปริวชฺชยโต อุปฺปชฺเชยฺยุํ อาสวา วิฆาต\-ปริฬาหา, ปริวชฺชยโต เอวํส เต อาสวา วิฆาต\-ปริฬาหา น โหนฺติ. อิเม วุจฺจนฺติ ภิกฺขเว อาสวา ปริวชฺชนา ปหาตพฺพา.}

\csromanmatikaanchor{mtk.9}
\csromantocmarkref{section}{วิโนทนาปหาตพฺพอาสว}{mtk.9}
\bookmark[dest={mtk.9},level=1]{วิโนทนาปหาตพฺพอาสว}
\csromanheader{1}{วิโนทนา\-ปหาตพฺพ\-อาสว}
\proseitem{26}{กตเม จ ภิกฺขเว อาสวา วิโนทนา ปหาตพฺพา. อิธ ภิกฺขเว ภิกฺขุ ปฏิสงฺขา โยนิโส อุปฺปนฺนํ กามวิตกฺกํ นาธิวาเสติ ปชหติ วิโนเทติ พฺยนฺตี กโรติ อนภาวํ คเมติ. อุปฺปนฺนํ พฺยาปาท\-วิตกฺกํ ฯเปฯ อุปฺปนฺนํ วิหิํสา\-วิตกฺกํ ฯเปฯ อุปฺปนฺนุปฺปนฺเน ปาปเก อกุสเล ธมฺเม นาธิวาเสติ ปชหติ วิโนเทติ พฺยนฺตี กโรติ อนภาวํ คเมติ.}

\prose{ยญฺหิสฺส ภิกฺขเว อวิโนทยโต อุปฺปชฺเชยฺยุํ อาสวา วิฆาต\-ปริฬาหา, วิโนทยโต เอวํส เต อาสวา วิฆาต\-ปริฬาหา น โหนฺติ. อิเม วุจฺจนฺติ ภิกฺขเว อาสวา วิโนทนา ปหาตพฺพา.}

\csromanmatikaanchor{mtk.10}
\csromantocmarkref{section}{ภาวนาปหาตพฺพอาสว}{mtk.10}
\bookmark[dest={mtk.10},level=1]{ภาวนาปหาตพฺพอาสว}
\csromanheader{1}{ภาวนา\-ปหาตพฺพ\-อาสว}
\proseitem{27}{กตเม จ ภิกฺขเว อาสวา ภาวนา ปหาตพฺพา. อิธ ภิกฺขเว ภิกฺขุ ปฏิสงฺขา โยนิโส สติสมฺโพชฺฌงฺคํ ภาเวติ วิเวกนิสฺสิตํ วิราคนิสฺสิตํ นิโรธ\-นิสฺสิตํ โวสฺสคฺค\-ปริณามิํ. ปฏิสงฺขา โยนิโส ธมฺมวิจย\-สมฺโพชฺฌงฺคํ ภาเวติ ฯเปฯ วีริย\-สมฺโพชฺฌงฺคํ ภาเวติ ฯเปฯ ปีติสมฺโพชฺฌงฺคํ ภาเวติ ฯเปฯ ปสฺสทฺธิ\-สมฺโพชฺฌงฺคํ ภาเวติ ฯเปฯ สมาธิ\-สมฺโพชฺฌงฺคํ ภาเวติ ฯเปฯ อุเปกฺขา\-สมฺโพชฺฌงฺคํ ภาเวติ วิเวกนิสฺสิตํ วิราคนิสฺสิตํ นิโรธ\-นิสฺสิตํ โวสฺสคฺค\-ปริณามิํ.}

\prose{ยญฺหิสฺส ภิกฺขเว อภาวยโต อุปฺปชฺเชยฺยุํ อาสวา วิฆาต\-ปริฬาหา, ภาวยโต เอวํส เต อาสวา วิฆาต\-ปริฬาหา น โหนฺติ. อิเม วุจฺจนฺติ ภิกฺขเว อาสวา ภาวนา ปหาตพฺพา.}

\csromanmatikaanchor{mtk.11}
\csromantocmarkref{subsection}{3. ธมฺมทายาทสุตฺต}{mtk.11}
\bookmark[dest={mtk.11},level=2]{3. ธมฺมทายาทสุตฺต}
\proseitemwithcloserruled{28}{ยโต โข ภิกฺขเว ภิกฺขุโน เย อาสวา ทสฺสนา ปหาตพฺพา, เต ทสฺสนา ปหีนา โหนฺติ. เย อาสวา สํวรา \csromanfolio{15}ปหาตพฺพา, เต สํวรา ปหีนา โหนฺติ. เย อาสวา ปฏิเสวนา ปหาตพฺพา, เต ปฏิเสวนา ปหีนา โหนฺติ. เย อาสวา อธิวาสนา ปหาตพฺพา, เต อธิวาสนา ปหีนา โหนฺติ. เย อาสวา ปริวชฺชนา ปหาตพฺพา, เต ปริวชฺชนา ปหีนา โหนฺติ. เย อาสวา วิโนทนา ปหาตพฺพา, เต วิโนทนา ปหีนา โหนฺติ. เย อาสวา ภาวนา ปหาตพฺพา, เต ภาวนา ปหีนา โหนฺติ. อยํ วุจฺจติ ภิกฺขเว ภิกฺขุ สพฺพาสว\-สํวร\-สํวุโต วิหรติ. อจฺเฉจฺฉิ\csromansharedfootnote{2d9309cfca9f7c87}{อจฺเฉชฺชิ (ก)} ตณฺหํ, วิวตฺตยิ\csromansharedfootnote{ad9184755c407031}{วาวตฺตยิ (สี, อิ)} สํโยชนํ, สมฺมา มานาภิสมยา อนฺตมกาสิ ทุกฺขสฺสาติ. อิทมโวจ ภควา. อตฺตมนา เต ภิกฺขู ภควโต ภาสิตํ อภินนฺทุนฺติ.}{สพฺพาสว\-สุตฺตํ นิฏฺฐิตํ ทุติยํ.}
\proseitem{29}{เอวํ เม สุตํ\csromandash{}เอกํ สมยํ ภควา สาวตฺถิยํ วิหรติ เชตวเน อนาถ\-ปิณฺฑิกสฺส อาราเม. ตตฺร โข ภควา ภิกฺขู อามนฺเตสิ “ภิกฺขโว”ติ. “ภทนฺเต”ติ เต ภิกฺขู ภควโต ปจฺจสฺโสสุํ. ภควา เอตทโวจ\csromandash{}}

\prose{ธมฺมทายาทา เม ภิกฺขเว ภวถ มา อามิสทายาทา. อตฺถิ เม ตุเมฺหสุ อนุกมฺปา “กินฺติ เม สาวกา ธมฺมทายาทา ภเวยฺยุํ, โน อามิสทายาทา”ติ. ตุเมฺห จ เม ภิกฺขเว อามิสทายาทา ภเวยฺยาถ, โน ธมฺมทายาทา, ตุเมฺหปิ เตน อาทิยา\csromansharedfootnote{5504600c8938febc}{อาทิสฺสา (สี, สฺยา, อิ)} ภเวยฺยาถ “อามิสทายาทา สตฺถุสาวกา วิหรนฺติ, โน ธมฺมทายาทา”ติ. อหมฺปิ เตน อาทิโย ภเวยฺยํ “อามิสทายาทา สตฺถุสาวกา วิหรนฺติ, โน ธมฺมทายาทา”ติ. ตุเมฺห จ เม ภิกฺขเว ธมฺมทายาทา ภเวยฺยาถ, โน อามิสทายาทา, ตุเมฺหปิ เตน น อาทิยา ภเวยฺยาถ “ธมฺมทายาทา สตฺถุสาวกา วิหรนฺติ, โน อามิสทายาทา”ติ. อหมฺปิ เตน น อาทิโย ภเวยฺยํ “ธมฺมทายาทา สตฺถุสาวกา วิหรนฺติ, โน อามิสทายาทา”ติ. ตสฺมาติห เม ภิกฺขเว ธมฺมทายาทา ภวถ, มา อามิสทายาทา. อตฺถิ \csromanfolio{16}เม ตุเมฺหสุ อนุกมฺปา “กินฺติ เม สาวกา ธมฺมทายาทา ภเวยฺยุํ, โน อามิสทายาทา”ติ.}

\proseitem{30}{อิธาหํ ภิกฺขเว ภุตฺตาวี อสฺสํ ปวาริโต ปริปุณฺโณ ปริโยสิโต สุหิโต ยาวทตฺโถ, สิยา จ เม ปิณฺฑปาโต อติเรกธมฺโม ฉฑฺฑนีย\-ธมฺโม\csromansharedfootnote{5ecfc8518b381d54}{ฉฑฺฑิยธมฺโม (สี, สฺยา, อิ)}. อถ เทฺว ภิกฺขู อาคจฺเฉยฺยุํ ชิฆจฺฉา\-ทุพฺพลฺย\csromansharedfootnote{20bbf10cea8c2178}{ชิฆจฺฉาทุพฺพลฺล (สี, อิ)} ปเรตา, ตฺยาหํ เอวํ วเทยฺยํ “อหํ โขมฺหิ ภิกฺขเว ภุตฺตาวี ปวาริโต ปริปุณฺโณ ปริโยสิโต สุหิโต ยาวทตฺโถ, อตฺถิ จ เม อยํ ปิณฺฑปาโต อติเรกธมฺโม ฉฑฺฑนีย\-ธมฺโม. สเจ อากงฺขถ ภุญฺชถ, โน เจ ตุเมฺห ภุญฺชิสฺสถ\csromansharedfootnote{e8759990cb7e6eb6}{สเจ ตุเมฺห น ภุญฺชิสฺสถ (สี, สฺยา, อิ)}, อิทานาหํ อปฺปหริเต วา ฉฑฺเฑสฺสามิ, อปฺปาณเก วา อุทเก โอปิลาเปสฺสามี”ติ. ตเตฺรกสฺส ภิกฺขุโน เอวมสฺส “ภควา โข ภุตฺตาวี ปวาริโต ปริปุณฺโณ ปริโยสิโต สุหิโต ยาวทตฺโถ, อตฺถิ จายํ ภควโต ปิณฺฑปาโต อติเรกธมฺโม ฉฑฺฑนีย\-ธมฺโม. สเจ มยํ น ภุญฺชิสฺสาม, อิทานิ ภควา อปฺปหริเต วา ฉฑฺเฑสฺสติ, อปฺปาณเก วา อุทเก โอปิลาเปสฺสติ. วุตฺตํ โข ปเนตํ ภควตา ‘ธมฺมทายาทา เม ภิกฺขเว ภวถ, มา อามิสทายาทา’ติ. อามิสญฺญตรํ โข ปเนตํ, ยทิทํ ปิณฺฑปาโต. ยํนูนาหํ อิมํ ปิณฺฑปาตํ อภุญฺชิตฺวา อิมินาว ชิฆจฺฉา\-ทุพฺพเลฺยน เอวํ อิมํ รตฺตินฺทิวํ\csromansharedfootnote{dd3fab0193bd20cb}{รตฺติทิวํ (ก)} วีตินาเมยฺยนฺ”ติ. โส ตํ ปิณฺฑปาตํ อภุญฺชิตฺวา เตเนว ชิฆจฺฉา\-ทุพฺพเลฺยน เอวํ ตํ รตฺตินฺทิวํ วีตินาเมยฺย. อถ ทุติยสฺส ภิกฺขุโน เอวมสฺส “ภควา โข ภุตฺตาวี ปวาริโต ปริปุณฺโณ ปริโยสิโต สุหิโต ยาวทตฺโถ, อตฺถิ จายํ ภควโต ปิณฺฑปาโต อติเรกธมฺโม ฉฑฺฑนีย\-ธมฺโม. สเจ มยํ น ภุญฺชิสฺสาม, อิทานิ ภควา อปฺปหริเต วา ฉฑฺเฑสฺสติ, อปฺปาณเก วา อุทเก โอปิลาเปสฺสติ. ยํนูนาหํ อิมํ ปิณฺฑปาตํ ภุญฺชิตฺวา ชิฆจฺฉา\-ทุพฺพลฺยํ ปฏิวิโนเทตฺวา\csromansharedfootnote{9c4b9e751ebbae48}{ปฏิวิเนตฺวา (สี, สฺยา, อิ)} เอวํ อิมํ รตฺตินฺทิวํ วีตินาเมยฺยนฺ”ติ. โส ตํ ปิณฺฑปาตํ ภุญฺชิตฺวา ชิฆจฺฉา\-ทุพฺพลฺยํ ปฏิวิโนเทตฺวา เอวํ ตํ รตฺตินฺทิวํ วีตินาเมยฺย. กิญฺจาปิ โส ภิกฺขเว ภิกฺขุ ตํ ปิณฺฑปาตํ ภุญฺชิตฺวา ชิฆจฺฉา\-ทุพฺพลฺยํ ปฏิวิโนเทตฺวา เอวํ ตํ รตฺตินฺทิวํ วีตินาเมยฺย, อถ โข อสุเยว เม ปุริโม ภิกฺขุ ปุชฺชตโร จ ปาสํสตโร จ. ตํ กิสฺส เหตุ, ตญฺหิ ตสฺส ภิกฺขเว ภิกฺขุโน \csromanfolio{17}ทีฆรตฺตํ อปฺปิจฺฉตาย สนฺตุฏฺฐิยา สลฺเลขาย สุภรตาย วีริยารมฺภาย สํวตฺติสฺสติ. ตสฺมาติห เม ภิกฺขเว ธมฺมทายาทา ภวถ, มา อามิสทายาทา. อตฺถิ เม ตุเมฺหสุ อนุกมฺปา “กินฺติ เม สาวกา ธมฺมทายาทา ภเวยฺยุํ, โน อามิสทายาทา”ติ.}

\prose{อิทมโวจ ภควา. อิทํ วตฺวาน\csromansharedfootnote{f786312130b697a4}{วตฺวา (สี, อิ) เอวมีทิเสสุ ฐาเนสุ.} สุคโต อุฏฺฐายาสนา วิหารํ ปาวิสิ.}

\proseitem{31}{ตตฺร โข อายสฺมา สาริปุตฺโต อจิร\-ปกฺกนฺตสฺส ภควโต ภิกฺขู อามนฺเตสิ “อาวุโส ภิกฺขเว”ติ. “อาวุโส”ติ โข เต ภิกฺขู อายสฺมโต สาริปุตฺตสฺส ปจฺจสฺโสสุํ. อายสฺมา สาริปุตฺโต เอตทโวจ\csromandash{}}

\prose{“กิตฺตาวตา นุ โข อาวุโส สตฺถุ ปวิวิตฺตสฺส วิหรโต สาวกา วิเวกํ นานุสิกฺขนฺติ, กิตฺตาวตา จ ปน สตฺถุ ปวิวิตฺตสฺส วิหรโต สาวกา วิเวกมนุสิกฺขนฺตี”ติ. ทูรโตปิ โข มยํ อาวุโส อาคจฺฉาม อายสฺมโต สาริปุตฺตสฺส สนฺติเก เอตสฺส ภาสิตสฺส อตฺถมญฺญาตุํ. สาธุ วตายสฺมนฺตํเยว สาริปุตฺตํ ปฏิภาตุ เอตสฺส ภาสิตสฺส อตฺโถ, อายสฺมโต สาริปุตฺตสฺส สุตฺวา ภิกฺขู ธาเรสฺสนฺตีติ. เตน หาวุโส สุณาถ สาธุกํ มนสิ กโรถ ภาสิสฺสามีติ. “เอวมาวุโส”ติ โข เต ภิกฺขู อายสฺมโต สาริปุตฺตสฺส ปจฺจสฺโสสุํ. อายสฺมา สาริปุตฺโต เอตทโวจ\csromandash{}}

\prose{กิตฺตาวตา นุ โข อาวุโส สตฺถุ ปวิวิตฺตสฺส วิหรโต สาวกา วิเวกํ นานุสิกฺขนฺติ. อิธาวุโส สตฺถุ ปวิวิตฺตสฺส วิหรโต สาวกา วิเวกํ นานุสิกฺขนฺติ, เยสญฺจ ธมฺมานํ สตฺถา ปหานมาห, เต จ ธมฺเม นปฺปชหนฺติ, พาหุลิกา\csromansharedfootnote{a8b7fc216b82e46b}{พาหุลฺลิกา (สฺยา)} จ โหนฺติ สาถลิกา, โอกฺกมเน ปุพฺพงฺคมา ปวิเวเก นิกฺขิตฺตธุรา. ตตฺราวุโส เถรา ภิกฺขู ตีหิ ฐาเนหิ คารยฺหา ภวนฺติ. “สตฺถุ ปวิวิตฺตสฺส วิหรโต สาวกา วิเวกํ นานุสิกฺขนฺตี”ติ อิมินา ปฐเมน ฐาเนน เถรา ภิกฺขู คารยฺหา ภวนฺติ. “เยสญฺจ ธมฺมานํ สตฺถา ปหานมาห, เต จ ธมฺเม นปฺปชหนฺตี”ติ อิมินา \csromanfolio{18}ทุติเยน ฐาเนน เถรา ภิกฺขู คารยฺหา ภวนฺติ. “พาหุลิกา จ สาถลิกา, โอกฺกมเน ปุพฺพงฺคมา ปวิเวเก นิกฺขิตฺตธุรา”ติ อิมินา ตติเยน ฐาเนน เถรา ภิกฺขู คารยฺหา ภวนฺติ. เถรา อาวุโส ภิกฺขู อิเมหิ ตีหิ ฐาเนหิ คารยฺหา ภวนฺติ. ตตฺราวุโส มชฺฌิมา ภิกฺขู ฯเปฯ นวา ภิกฺขู ตีหิ ฐาเนหิ คารยฺหา ภวนฺติ. “สตฺถุ ปวิวิตฺตสฺส วิหรโต สาวกา วิเวกํ นานุสิกฺขนฺตี”ติ อิมินา ปฐเมน ฐาเนน นวา ภิกฺขู คารยฺหา ภวนฺติ. “เยสญฺจ ธมฺมานํ สตฺถา ปหานมาห, เต จ ธมฺเม นปฺปชหนฺตี”ติ อิมินา ทุติเยน ฐาเนน นวา ภิกฺขู คารยฺหา ภวนฺติ. “พาหุลิกา จ โหนฺติ สาถลิกา, โอกฺกมเน ปุพฺพงฺคมา ปวิเวเก นิกฺขิตฺตธุรา”ติ อิมินา ตติเยน ฐาเนน นวา ภิกฺขู คารยฺหา ภวนฺติ. นวา อาวุโส ภิกฺขู อิเมหิ ตีหิ ฐาเนหิ คารยฺหา ภวนฺติ. เอตฺตาวตา โข อาวุโส สตฺถุ ปวิวิตฺตสฺส วิหรโต สาวกา วิเวกํ นานุสิกฺขนฺติ.}

\proseitem{32}{กิตฺตาวตา จ ปนาวุโส สตฺถุ ปวิวิตฺตสฺส วิหรโต สาวกา วิเวกมนุสิกฺขนฺติ. อิธาวุโส สตฺถุ ปวิวิตฺตสฺส วิหรโต สาวกา วิเวกมนุสิกฺขนฺติ, เยสญฺจ ธมฺมานํ สตฺถา ปหานมาห, เต จ ธมฺเม ปชหนฺติ, น จ พาหุลิกา โหนฺติ น สาถลิกา, โอกฺกมเน นิกฺขิตฺตธุรา ปวิเวเก ปุพฺพงฺคมา. ตตฺราวุโส เถรา ภิกฺขู ตีหิ ฐาเนหิ ปาสํสา ภวนฺติ. “สตฺถุ ปวิวิตฺตสฺส วิหรโต สาวกา วิเวกมนุสิกฺขนฺตี”ติ อิมินา ปฐเมน ฐาเนน เถรา ภิกฺขู ปาสํสา ภวนฺติ. “เยสญฺจ ธมฺมานํ สตฺถา ปหานมาห, เต จ ธมฺเม ปชหนฺตี”ติ อิมินา ทุติเยน ฐาเนน เถรา ภิกฺขู ปาสํสา ภวนฺติ. “น จ พาหุลิกา น สาถลิกา, โอกฺกมเน นิกฺขิตฺตธุรา ปวิเวเก ปุพฺพงฺคมา”ติ อิมินา ตติเยน ฐาเนน เถรา ภิกฺขู ปาสํสา ภวนฺติ. เถรา อาวุโส ภิกฺขู อิเมหิ ตีหิ ฐาเนหิ ปาสํสา ภวนฺติ. ตตฺราวุโส มชฺฌิมา ภิกฺขู ฯเปฯ นวา ภิกฺขู ตีหิ ฐาเนหิ ปาสํสา ภวนฺติ. “สตฺถุ ปวิวิตฺตสฺส วิหรโต สาวกา วิเวกมนุสิกฺขนฺตี”ติ อิมินา ปฐเมน ฐาเนน นวา ภิกฺขู ปาสํสา ภวนฺติ. “เยสญฺจ ธมฺมานํ สตฺถา ปหานมาห, เต จ ธมฺเม ปชหนฺตี”ติ อิมินา ทุติเยน ฐาเนน นวา ภิกฺขู ปาสํสา ภวนฺติ. “น จ พาหุลิกา น สาถลิกา, โอกฺกมเน นิกฺขิตฺตธุรา ปวิเวเก ปุพฺพงฺคมา”ติ อิมินา ตติเยน ฐาเนน นวา ภิกฺขู ปาสํสา ภวนฺติ. นวา อาวุโส ภิกฺขู \csromanfolio{19}อิเมหิ ตีหิ ฐาเนหิ ปาสํสา ภวนฺติ. เอตฺตาวตา โข อาวุโส สตฺถุ ปวิวิตฺตสฺส วิหรโต สาวกา วิเวกมนุสิกฺขนฺติ.}

\proseitem{33}{ตตฺราวุโส โลโภ จ ปาปโก, โทโส จ ปาปโก, โลภสฺส จ ปหานาย โทสสฺส จ ปหานาย อตฺถิ มชฺฌิมา ปฏิปทา จกฺขุกรณี ญาณกรณี อุปสมาย อภิญฺญาย สมฺโพธาย นิพฺพานาย สํวตฺตติ. กตมา จ สา อาวุโส มชฺฌิมา ปฏิปทา จกฺขุกรณี ญาณกรณี อุปสมาย อภิญฺญาย สมฺโพธาย นิพฺพานาย สํวตฺตติ. อยเมว อริโย อฏฺฐงฺคิโก มคฺโค, เสยฺยถิทํ\csromansharedfootnote{57d15287b23af282}{เสยฺยถีทํ (สี, สฺยา, อิ)}, สมฺมาทิฏฺฐิ สมฺมาสงฺกปฺโป สมฺมาวาจา สมฺมากมฺมนฺโต สมฺมาอาชีโว สมฺมาวายาโม สมฺมาสติ สมฺมาสมาธิ. อยํ โข สา อาวุโส มชฺฌิมา ปฏิปทา จกฺขุกรณี ญาณกรณี อุปสมาย อภิญฺญาย สมฺโพธาย นิพฺพานาย สํวตฺตติ.}

\prose{ตตฺราวุโส โกโธ จ ปาปโก, อุปนาโห จ ปาปโก ฯเปฯ มกฺโข จ ปาปโก, ปฬาโส จ ปาปโก. อิสฺสา จ ปาปิกา, มจฺเฉรญฺจ ปาปกํ. มายา จ ปาปิกา, สาเฐยฺยญฺจ ปาปกํ. ถมฺโภ จ ปาปโก, สารมฺโภ จ ปาปโก. มาโน จ ปาปโก, อติมาโน จ ปาปโก. มโท จ ปาปโก, ปมาโท จ ปาปโก, มทสฺส จ ปหานาย ปมาทสฺส จ ปหานาย อตฺถิ มชฺฌิมา ปฏิปทา จกฺขุกรณี ญาณกรณี อุปสมาย อภิญฺญาย สมฺโพธาย นิพฺพานาย สํวตฺตติ. กตมา จ สา อาวุโส มชฺฌิมา ปฏิปทา จกฺขุกรณี ญาณกรณี อุปสมาย อภิญฺญาย สมฺโพธาย นิพฺพานาย สํวตฺตติ. อยเมว อริโย อฏฺฐงฺคิโก มคฺโค, เสยฺยถิทํ, สมฺมาทิฏฺฐิ สมฺมาสงฺกปฺโป สมฺมาวาจา สมฺมากมฺมนฺโต สมฺมาอาชีโว สมฺมาวายาโม สมฺมาสติ สมฺมาสมาธิ. อยํ โข สา อาวุโส มชฺฌิมา ปฏิปทา จกฺขุกรณี ญาณกรณี อุปสมาย อภิญฺญาย สมฺโพธาย นิพฺพานาย สํวตฺตตีติ.}

\prosewithcloser{อิทมโวจายสฺมา สาริปุตฺโต. อตฺตมนา เต ภิกฺขู อายสฺมโต สาริปุตฺตสฺส ภาสิตํ อภินนฺทุนฺติ.}{ธมฺม\-ทายาท\-สุตฺตํ นิฏฺฐิตํ ตติยํ.}
\csromanmatikaanchor{mtk.12}
\csromantocmarkref{subsection}{4. ภยเภรวสุตฺต}{mtk.12}
\bookmark[dest={mtk.12},level=2]{4. ภยเภรวสุตฺต}
\proseitem{34}{\csromanfolio{20}เอวํ เม สุตํ\csromandash{}เอกํ สมยํ ภควา สาวตฺถิยํ วิหรติ เชตวเน อนาถ\-ปิณฺฑิกสฺส อาราเม. อถ โข ชาณุสฺโสณิ พฺราหฺมโณ เยน ภควา เตนุปสงฺกมิ, อุปสงฺกมิตฺวา ภควตา สทฺธิํ สมฺโมทิ, สมฺโมทนียํ กถํ สารณียํ\csromansharedfootnote{50a41b474ed73d9d}{สาราณียํ (สี, สฺยา, อิ)} วีติสาเรตฺวา เอกมนฺตํ นิสีทิ, เอกมนฺตํ นิสินฺโน โข ชาณุสฺโสณิ พฺราหฺมโณ ภควนฺตํ เอตทโวจ “เยเม โภ โคตม กุลปุตฺตา ภวนฺตํ โคตมํ อุทฺทิสฺส สทฺธา อคารสฺมา อนคาริยํ ปพฺพชิตา, ภวํ เตสํ โคตโม ปุพฺพงฺคโม, ภวํ เตสํ โคตโม พหุกาโร, ภวํ เตสํ โคตโม สมาทเปตา\csromansharedfootnote{251fb7f5126a923d}{สมาทาเปตา (?)}, โภโต จ ปน โคตมสฺส สา ชนตา ทิฏฺฐานุคติํ อาปชฺชตี”ติ. เอวเมตํ พฺราหฺมณ, เอวเมตํ พฺราหฺมณ, เย เต พฺราหฺมณ กุลปุตฺตา มมํ อุทฺทิสฺส สทฺธา อคารสฺมา อนคาริยํ ปพฺพชิตา, อหํ เตสํ ปุพฺพงฺคโม, อหํ เตสํ พหุกาโร, อหํ เตสํ สมาทเปตา, มม จ ปน สา ชนตา ทิฏฺฐานุคติํ อาปชฺชตีติ. ทุรภิสมฺภวานิ หิ โข โภ โคตม อรญฺญ\-วนปตฺถานิ ปนฺตานิ เสนาสนานิ, ทุกฺกรํ ปวิเวกํ, ทุรภิรมํ เอกตฺเต, หรนฺติ มญฺเญ มโน วนานิ สมาธิํ อลภมานสฺส ภิกฺขุโนติ. เอวเมตํ พฺราหฺมณ, เอวเมตํ พฺราหฺมณ, ทุรภิสมฺภวานิ หิ โข พฺราหฺมณ อรญฺญ\-วนปตฺถานิ ปนฺตานิ เสนาสนานิ, ทุกฺกรํ ปวิเวกํ, ทุรภิรมํ เอกตฺเต, หรนฺติ มญฺเญ มโน วนานิ สมาธิํ อลภมานสฺส ภิกฺขุโน.}

\proseitem{35}{มยฺหํปิ โข พฺราหฺมณ ปุพฺเพว สมฺโพธา อนภิสมฺพุทฺธสฺส โพธิสตฺตสฺเสว สโต เอตทโหสิ “ทุรภิสมฺภวานิ หิ โข อรญฺญ\-วนปตฺถานิ ปนฺตานิ เสนาสนานิ, ทุกฺกรํ ปวิเวกํ, ทุรภิรมํ เอกตฺเต, หรนฺติ มญฺเญ มโน วนานิ สมาธิํ อลภมานสฺส ภิกฺขุโน”ติ. ตสฺส มยฺหํ พฺราหฺมณ เอตทโหสิ “เย โข เกจิ สมณา วา พฺราหฺมณา วา อปริสุทฺธ\-กายกมฺมนฺตา อรญฺญ\-วนปตฺถานิ ปนฺตานิ เสนาสนานิ ปฏิเสวนฺติ, อปริสุทฺธกายกมฺมนฺต\-สนฺโทสเหตุ หเว เต โภนฺโต สมณพฺราหฺมณา อกุสลํ ภยเภรวํ อวฺหายนฺติ. น โข ปนาหํ อปริสุทฺธ\-กายกมฺมนฺโต อรญฺญ\-วนปตฺถานิ ปนฺตานิ เสนาสนานิ ปฏิเสวามิ, ปริสุทฺธกายกมฺมนฺโตหมสฺมิ. เย หิ โว อริยา ปริสุทฺธกาย \csromanfolio{21}กมฺมนฺตา อรญฺญ\-วนปตฺถานิ ปนฺตานิ เสนาสนานิ ปฏิเสวนฺติ, เตสมหํ อญฺญตโร”ติ. เอตมหํ พฺราหฺมณ ปริสุทฺธ\-กายกมฺมตํ อตฺตนิ สมฺปสฺสมาโน ภิยฺโย ปลฺโลมมาปาทิํ อรญฺเญ วิหาราย. (1)}

\proseitem{36}{ตสฺส มยฺหํ พฺราหฺมณ เอตทโหสิ “เย โข เกจิ สมณา วา พฺราหฺมณา วา อปริสุทฺธ\-วจีกมฺมนฺตา ฯเปฯ อปริสุทฺธ\-มโนกมฺมนฺตา ฯเปฯ อปริสุทฺธฺ อาชีวา อรญฺญ\-วนปตฺถานิ ปนฺตานิ เสนาสนานิ ปฏิเสวนฺติ, อปริสุทฺธาชีว\-สนฺโทสเหตุ หเว เต โภนฺโต สมณพฺราหฺมณา อกุสลํ ภยเภรวํ อวฺหายนฺติ. น โข ปนาหํ อปริสุทฺธา\-ชีโว อรญฺญ\-วนปตฺถานิ ปนฺตานิ เสนาสนานิ ปฏิเสวามิ, ปริสุทฺธาชีโวหมสฺมิ. เย หิ โว อริยา ปริสุทฺธาชีวา อรญฺญ\-วนปตฺถานิ ปนฺตานิ เสนาสนานิ ปฏิเสวนฺติ, เตสมหํ อญฺญตโร”ติ. เอตมหํ พฺราหฺมณ ปริสุทฺธา\-ชีวตํ อตฺตนิ สมฺปสฺสมาโน ภิยฺโย ปลฺโลมมาปาทิํ อรญฺเญ วิหาราย. (2-3-4)}

\proseitem{37}{ตสฺส มยฺหํ พฺราหฺมณ เอตทโหสิ “เย โข เกจิ สมณา วา พฺราหฺมณา วา อภิชฺฌาลู กาเมสุ ติพฺพสาราคา อรญฺญ\-วนปตฺถานิ ปนฺตานิ เสนาสนานิ ปฏิเสวนฺติ, อภิชฺฌาลุ\-กาเมสุ\-ติพฺพสาราค\-สนฺโทสเหตุ หเว เต โภนฺโต สมณพฺราหฺมณา อกุสลํ ภยเภรวํ อวฺหายนฺติ. น โข ปนาหํ อภิชฺฌาลุ กาเมสุ ติพฺพสาราโค อรญฺญ\-วนปตฺถานิ ปนฺตานิ เสนาสนานิ ปฏิเสวามิ, อนภิชฺฌาลูหมสฺมิ. เย หิ โว อริยา อนภิชฺฌาลู อรญฺญ\-วนปตฺถานิ ปนฺตานิ เสนาสนานิ ปฏิเสวนฺติ, เตสมหํ อญฺญตโร”ติ. เอตมหํ พฺราหฺมณ อนภิชฺฌาลุตํ อตฺตนิ สมฺปสฺสมาโน ภิยฺโย ปลฺโลมมาปาทิํ อรญฺเญ วิหาราย. (5)}

\proseitem{38}{ตสฺส มยฺหํ พฺราหฺมณ เอตทโหสิ “เย โข เกจิ สมณา วา พฺราหฺมณา วา พฺยาปนฺนจิตฺตา ปทุฏฺฐ\-มนสงฺกปฺปา อรญฺญ\-วนปตฺถานิ ปนฺตานิ เสนาสนานิ ปฏิเสวนฺติ, พฺยาปนฺนจิตฺต\-ปทุฏฺฐมนสงฺกปฺป\-สนฺโทสเหตุ หเว เต โภนฺโต สมณพฺราหฺมณา อกุสลํ ภยเภรวํ อวฺหายนฺติ. น โข ปนาหํ พฺยาปนฺนจิตฺโต ปทุฏฺฐ\-มนสงฺกปฺโป อรญฺญ\-วนปตฺถานิ ปนฺตานิ เสนาสนานิ ปฏิเสวามิ, เมตฺตจิตฺโตหมสฺมิ. เย หิ โว อริยา เมตฺตจิตฺตา อรญฺญ\-วนปตฺถานิ ปนฺตานิ เสนาสนานิ ปฏิเสวนฺติ, เตสมหํ อญฺญตโร”ติ. เอตมหํ พฺราหฺมณ เมตฺตจิตฺตตํ อตฺตนิ สมฺปสฺสมาโน ภิยฺโย ปลฺโลมมาปาทิํ อรญฺเญ วิหาราย. (6)}

\proseitem{39}{\csromanfolio{22}ตสฺส มยฺหํ พฺราหฺมณ เอตทโหสิ “เย โข เกจิ สมณา วา พฺราหฺมณา วา ถีนมิทฺธ\-ปริยุฏฺฐิตา อรญฺญ\-วนปตฺถานิ ปนฺตานิ เสนาสนานิ ปฏิเสวนฺติ, ถีนมิทฺธ\-ปริยุฏฺฐาน\-สนฺโทสเหตุ หเว เต โภนฺโต สมณพฺราหฺมณา อกุสลํ ภยเภรวํ อวฺหายนฺติ. น โข ปนาหํ ถีนมิทฺธ\-ปริยุฏฺฐิโต อรญฺญ\-วนปตฺถานิ ปนฺตานิ เสนาสนานิ ปฏิเสวามิ, วิคตถีนมิทฺโธหมสฺมิ. เย หิ โว อริยา วิคต\-ถีน\-มิทฺธา อรญฺญ\-วนปตฺถานิ ปนฺตานิ เสนาสนานิ ปฏิเสวนฺติ, เตสมหํ อญฺญตโร”ติ. เอตมหํ พฺราหฺมณ วิคต\-ถีนมิทฺธตํ อตฺตนิ สมฺปสฺสมาโน ภิยฺโย ปลฺโลมมาปาทิํ อรญฺเญ วิหาราย. (7)}

\proseitem{40}{ตสฺส มยฺหํ พฺราหฺมณ เอตทโหสิ “เย โข เกจิ สมณา วา พฺราหฺมณา วา อุทฺธตา อวูปสนฺต\-จิตฺตา อรญฺญ\-วนปตฺถานิ ปนฺตานิ เสนาสนานิ ปฏิเสวนฺติ, อุทฺธต\-อวูปสนฺตจิตฺต\-สนฺโทสเหตุ หเว เต โภนฺโต สมณพฺราหฺมณา อกุสลํ ภยเภรวํ อวฺหายนฺติ. น โข ปนาหํ อุทฺธโต อวูปสนฺต\-จิตฺโต อรญฺญ\-วนปตฺถานิ ปนฺตานิ เสนาสนานิ ปฏิเสวามิ, วูปสนฺตจิตฺโตหมสฺมิ. เย หิ โว อริยา วูปสนฺตจิตฺตา อรญฺญ\-วนปตฺถานิ ปนฺตานิ เสนาสนานิ ปฏิเสวนฺติ, เตสมหํ อญฺญตโร”ติ. เอตมหํ พฺราหฺมณ วูปสนฺต\-จิตฺตตํ อตฺตนิ สมฺปสฺสมาโน ภิยฺโย ปลฺโลมมาปาทิํ อรญฺเญ วิหาราย. (8)}

\proseitem{41}{ตสฺส มยฺหํ พฺราหฺมณ เอตทโหสิ “เย โข เกจิ สมณา วา พฺราหฺมณา วา กงฺขี วิจิกิจฺฉี อรญฺญ\-วนปตฺถานิ ปนฺตานิ เสนาสนานิ ปฏิเสวนฺติ, กงฺขิ\-วิจิกิจฺฉิ\-สนฺโทสเหตุ หเว เต โภนฺโต สมณพฺราหฺมณา อกุสลํ ภยเภรวํ อวฺหายนฺติ. น โข ปนาหํ กงฺขี วิจิกิจฺฉี อรญฺญ\-วนปตฺถานิ ปนฺตานิ เสนาสนานิ ปฏิเสวามิ, ติณฺณวิจิกิจฺโฉหมสฺมิ. เย หิ โว อริยา ติณฺณ\-วิจิกิจฺฉา อรญฺญ\-วนปตฺถานิ ปนฺตานิ เสนาสนานิ ปฏิ- เสวนฺติ, เตสมหํ อญฺญตโร”ติ. เอตมหํ พฺราหฺมณ ติณฺณวิจิกิจฺฉตํ อตฺตนิ สมฺปสฺสมาโน ภิยฺโย ปลฺโลมมาปาทิํ อรญฺเญ วิหาราย. (9)}

\proseitem{42}{ตสฺส มยฺหํ พฺราหฺมณ เอตทโหสิ “เย โข เกจิ สมณา วา พฺราหฺมณา วา อตฺตุกฺกํสกา ปรวมฺภี อรญฺญ\-วนปตฺถานิ ปนฺตานิ เสนาสนานิ ปฏิเสวนฺติ, อตฺตุกฺกํสน\-ปรวมฺภน\-สนฺโทสเหตุ หเว เต \csromanfolio{23}โภนฺโต สมณพฺราหฺมณา อกุสลํ ภยเภรวํ อวฺหายนฺติ. น โข ปนาหํ อตฺตุกฺกํสโก ปรวมฺภี อรญฺญ\-วนปตฺถานิ ปนฺตานิ เสนาสนานิ ปฏิเสวามิ, อนตฺตุกฺกํสโก อปรวมฺภีหมสฺมิ. เย หิ โว อริยา อนตฺตุกฺกํสกา อปรวมฺภี อรญฺญ\-วนปตฺถานิ ปนฺตานิ เสนาสนานิ ปฏิเสวนฺติ, เตสมหํ อญฺญตโร”ติ. เอตมหํ พฺราหฺมณ อนตฺตุกฺกํสกตํ อปรวมฺภิตํ อตฺตนิ สมฺปสฺสมาโน ภิยฺโย ปลฺโลมมาปาทิํ อรญฺเญ วิหาราย. (10)}

\proseitem{43}{ตสฺส มยฺหํ พฺราหฺมณ เอตทโหสิ “เย โข เกจิ สมณา วา พฺราหฺมณา วา ฉมฺภี ภีรุกชาติกา อรญฺญ\-วนปตฺถานิ ปนฺตานิ เสนาสนานิ ปฏิเสวนฺติ, ฉมฺภิ\-ภีรุกชาติก\-สนฺโทสเหตุ หเว เต โภนฺโต สมณพฺราหฺมณา อกุสลํ ภยเภรวํ อวฺหายนฺติ. น โข ปนาหํ ฉมฺภี ภีรุกชาติโก อรญฺญ\-วนปตฺถานิ ปนฺตานิ เสนาสนานิ ปฏิเสวามิ, วิคตโลมหํโสหมสฺมิ. เย หิ โว อริยา วิคต\-โลมหํสา อรญฺญ\-วนปตฺถานิ ปนฺตานิ เสนาสนานิ ปฏิเสวนฺติ, เตสมหํ อญฺญตโร”ติ. เอตมหํ พฺราหฺมณ วิคต\-โลม\-หํสตํ อตฺตนิ สมฺปสฺสมาโน ภิยฺโย ปลฺโลมมาปาทิํ อรญฺเญ วิหาราย. (11)}

\proseitem{44}{ตสฺส มยฺหํ พฺราหฺมณ เอตทโหสิ “เย โข เกจิ สมณา วา พฺราหฺมณา วา ลาภ\-สกฺการ\-สิโลกํ นิกามยมานา อรญฺญ\-วนปตฺถานิ ปนฺตานิ เสนาสนานิ ปฏิเสวนฺติ, ลาภ\-สกฺการ\-สิโลกนิกามน\csromansharedfootnote{2a6acf2d9eea7c80}{นิกามยมาน (สี, สฺยา)} สนฺโทสเหตุ หเว เต โภนฺโต สมณพฺราหฺมณา อกุสลํ ภยเภรวํ อวฺหายนฺติ. น โข ปนาหํ ลาภ\-สกฺการ\-สิโลกํ นิกามยมาโน อรญฺญ\-วนปตฺถานิ ปนฺตานิ เสนาสนานิ ปฏิเสวามิ, อปฺปิจฺโฉหมสฺมิ. เย หิ โว อริยา อปฺปิจฺฉา อรญฺญ\-วนปตฺถานิ ปนฺตานิ เสนาสนานิ ปฏิเสวนฺติ, เตสมหํ อญฺญตโร”ติ. เอตมหํ พฺราหฺมณ อปฺปิจฺฉตํ อตฺตนิ สมฺปสฺสมาโน ภิยฺโย ปลฺโลมมาปาทิํ อรญฺเญ วิหาราย. (12)}

\proseitem{45}{ตสฺส มยฺหํ พฺราหฺมณ เอตทโหสิ “เย โข เกจิ สมณา วา พฺราหฺมณา วา กุสีตา หีนวีริยา อรญฺญ\-วนปตฺถานิ ปนฺตานิ เสนาสนานิ \csromanfolio{24}ปฏิเสวนฺติ, กุสีต\-หีนวีริย\-สนฺโทสเหตุ หเว เต โภนฺโต สมณพฺราหฺมณา อกุสลํ ภยเภรวํ อวฺหายนฺติ. น โข ปนาหํ กุสีโต หีนวีริโย อรญฺญ\-วนปตฺถานิ ปนฺตานิ เสนาสนานิ ปฏิเสวามิ, อารทฺธวีริโยหมสฺมิ. เย หิ โว อริยา อารทฺธวีริยา อรญฺญ\-วนปตฺถานิ ปนฺตานิ เสนาสนานิ ปฏิเสวนฺติ, เตสมหํ อญฺญตโร”ติ. เอตมหํ พฺราหฺมณ อารทฺธ\-วีริยตํ อตฺตนิ สมฺปสฺสมาโน ภิยฺโย ปลฺโลมมาปาทิํ อรญฺเญ วิหาราย. (13)}

\proseitem{46}{ตสฺส มยฺหํ พฺราหฺมณ เอตทโหสิ “เย โข เกจิ สมณา วา พฺราหฺมณา วา มุฏฺฐสฺสตี อสมฺปชานา อรญฺญ\-วนปตฺถานิ ปนฺตานิ เสนาสนานิ ปฏิเสวนฺติ, มุฏฺฐสฺสติ\-อสมฺปชาน\-สนฺโทสเหตุ หเว เต โภนฺโต สมณพฺราหฺมณา อกุสลํ ภยเภรวํ อวฺหายนฺติ. น โข ปนาหํ มุฏฺฐสฺสติ อสมฺปชาโน อรญฺญ\-วนปตฺถานิ ปนฺตานิ เสนาสนานิ ปฏิเสวามิ, อุปฏฺฐิตสฺสติหมสฺมิ. เย หิ โว อริยา อุปฏฺฐิตสฺสตี อรญฺญ\-วนปตฺถานิ ปนฺตานิ เสนาสนานิ ปฏิเสวนฺติ, เตสมหํ อญฺญตโร”ติ. เอตมหํ พฺราหฺมณ อุปฏฺฐิต\-สฺสติตํ อตฺตนิ สมฺปสฺสมาโน ภิยฺโย ปลฺโลมมาปาทิํ อรญฺเญ วิหาราย. (14)}

\proseitem{47}{ตสฺส มยฺหํ พฺราหฺมณ เอตทโหสิ “เย โข เกจิ สมณา วา พฺราหฺมณา วา อสมาหิตา วิพฺภนฺตจิตฺตา อรญฺญ\-วนปตฺถานิ ปนฺตานิ เสนาสนานิ ปฏิเสวนฺติ, อสมาหิต\-วิพฺภนฺตจิตฺต\-สนฺโทสเหตุ หเว เต โภนฺโต สมณพฺราหฺมณา อกุสลํ ภยเภรวํ อวฺหายนฺติ. น โข ปนาหํ อสมาหิโต วิพฺภนฺตจิตฺโต อรญฺญ\-วนปตฺถานิ ปนฺตานิ เสนาสนานิ ปฏิเสวามิ, สมาธิสมฺปนฺโนหมสฺมิ. เย หิ โว อริยา สมาธิ\-สมฺปนฺนา อรญฺญ\-วนปตฺถานิ ปนฺตานิ เสนาสนานิ ปฏิเสวนฺติ, เตสมหํ อญฺญตโร”ติ. เอตมหํ พฺราหฺมณ สมาธิ\-สมฺปทํ อตฺตนิ สมฺปสฺสมาโน ภิยฺโย ปลฺโลมมาปาทิํ อรญฺเญ วิหาราย. (15)}

\proseitemwithcloser{48}{ตสฺส มยฺหํ พฺราหฺมณ เอตทโหสิ “เย โข เกจิ สมณา วา พฺราหฺมณา วา ทุปฺปญฺญา เอฬมูคา อรญฺญ\-วนปตฺถานิ ปนฺตานิ เสนาสนานิ ปฏิเสวนฺติ, ทุปฺปญฺญเอฬมูคสนฺโทสเหตุ หเว เต โภนฺโต สมณพฺราหฺมณา อกุสลํ ภยเภรวํ อวฺหายนฺติ. น โข ปนาหํ ทุปฺปญฺโญ \csromanfolio{25}เอฬมูโค อรญฺญ\-วนปตฺถานิ ปนฺตานิ เสนาสนานิ ปฏิเสวามิ, ปญฺญาสมฺปนฺโนหมสฺมิ. เย หิ โว อริยา ปญฺญาสมฺปนฺนา อรญฺญ\-วนปตฺถานิ ปนฺตานิ เสนาสนานิ ปฏิเสวนฺติ, เตสมหํ อญฺญตโร”ติ. เอตมหํ พฺราหฺมณ ปญฺญาสมฺปทํ อตฺตนิ สมฺปสฺสมาโน ภิยฺโย ปลฺโลมมาปาทิํ อรญฺเญ วิหาราย. (16)}{โสฬส\-ปริยายํ นิฏฺฐิตํ.}
\proseitem{49}{ตสฺส มยฺหํ พฺราหฺมณ เอตทโหสิ “ยํนูนาหํ ยา ตา รตฺติโย อภิญฺญาตา อภิลกฺขิตา จาตุทฺทสี ปญฺจทสี อฏฺฐมี จ ปกฺขสฺส, ตถารูปาสุ รตฺตีสุ ยานิ ตานิ อารามเจติยานิ วนเจติยานิ รุกฺขเจติยานิ ภิํสนกานิ สโลมหํสานิ, ตถารูเปสุ เสนาสเนสุ วิหเรยฺยํ ‘อปฺเปว นามาหํ ภยเภรวํ ปสฺเสยฺยนฺ’ติ”. โส โข อหํ พฺราหฺมณ อปเรน สมเยน ยา ตา รตฺติโย อภิญฺญาตา อภิลกฺขิตา จาตุทฺทสี ปญฺจทสี อฏฺฐมี จ ปกฺขสฺส, ตถารูปาสุ รตฺตีสุ ยานิ ตานิ อารามเจติยานิ วนเจติยานิ รุกฺขเจติยานิ ภิํสนกานิ สโลมหํสานิ, ตถารูเปสุ เสนาสเนสุ วิหรามิ. ตตฺถ จ เม พฺราหฺมณ วิหรโต มโค วา อาคจฺฉติ, โมโร วา กฏฺฐํ ปาเตติ, วาโต วา ปณฺณก สฏํ\csromansharedfootnote{3fe010249adc1f32}{ปณฺณสฏํ (สี, อิ)} เอเรติ. ตสฺส มยฺหํ พฺราหฺมณ เอตทโหสิ\csromansharedfootnote{190372115688e837}{ตสฺส มยฺหํ เอวํ โหติ (สี, สฺยา)} “เอตํนูน ตํ ภยเภรวํ อาคจฺฉตี”ติ. ตสฺส มยฺหํ พฺราหฺมณ เอตทโหสิ “กิํ นุ โข อหํ อญฺญทตฺถุ ภยปฏิกงฺขี\csromansharedfootnote{551da70598a48d0d}{ภยปาฏิกงฺขี (สี)} วิหรามิ. ยํนูนาหํ ยถาภูตํ ยถาภูตสฺส\csromansharedfootnote{6e29931573f9bf54}{ยถาภูตสฺส ยถาภูตสฺส (สี, สฺยา)} เม ตํ ภยเภรวํ อาคจฺฉติ, ตถาภูตํ ตถาภูโตว\csromansharedfootnote{dc1de1eb22eab767}{ยถาภูโต ยถาภูโตว (สี, สฺยา)} ตํ ภยเภรวํ ปฏิวิเนยฺยนฺ”ติ. ตสฺส มยฺหํ พฺราหฺมณ จงฺกมนฺตสฺส ตํ ภยเภรวํ อาคจฺฉติ, โส โข อหํ พฺราหฺมณ เนว ตาว ติฏฺฐามิ น นิสีทามิ น นิปชฺชามิ, ยาว จงฺกมนฺโตว ตํ ภยเภรวํ ปฏิวิเนมิ. ตสฺส มยฺหํ พฺราหฺมณ ฐิตสฺส ตํ ภยเภรวํ อาคจฺฉติ, โส โข อหํ พฺราหฺมณ เนว ตาว จงฺกมามิ น นิสีทามิ น นิปชฺชามิ, ยาว ฐิโตว ตํ ภยเภรวํ ปฏิวิเนมิ. ตสฺส มยฺหํ พฺราหฺมณ นิสินฺนสฺส ตํ ภยเภรวํ อาคจฺฉติ, โส โข อหํ พฺราหฺมณ เนว ตาว นิปชฺชามิ น \csromanfolio{26}ติฏฺฐามิ น จงฺกมามิ, ยาว นิสินฺโนว ตํ ภยเภรวํ ปฏิวิเนมิ. ตสฺส มยฺหํ พฺราหฺมณ นิปนฺนสฺส ตํ ภยเภรวํ อาคจฺฉติ, โส โข อหํ พฺราหฺมณ เนว ตาว นิสีทามิ น ติฏฺฐามิ น จงฺกมามิ, ยาว นิปนฺโนว ตํ ภยเภรวํ ปฏิวิเนมิ.}

\proseitem{50}{สนฺติ โข ปน พฺราหฺมณ เอเก สมณพฺราหฺมณา รตฺติํเยว สมานํ ทิวาติ สญฺชานนฺติ, ทิวาเยว สมานํ รตฺตีติ สญฺชานนฺติ. อิทมหํ เตสํ สมณ\-พฺราหฺมณานํ สมฺโมห\-วิหารสฺมิํ วทามิ. อหํ โข ปน พฺราหฺมณ รตฺติํ เยว สมานํ รตฺตีติ สญฺชานามิ, ทิวาเยว สมานํ ทิวาติ สญฺชานามิ. ยํ โข ตํ พฺราหฺมณ สมฺมา วทมาโน วเทยฺย “อสมฺโมหธมฺโม สตฺโต โลเก อุปฺปนฺโน พหุชนหิตาย พหุชน\-สุขาย โลกานุกมฺปาย อตฺถาย หิตาย สุขาย เทวมนุสฺสานนฺ”ติ, มเมว ตํ สมฺมา วทมาโน วเทยฺย “อสมฺโมหธมฺโม สตฺโต โลเก อุปฺปนฺโน พหุชนหิตาย พหุชน\-สุขาย โลกานุกมฺปาย อตฺถาย หิตาย สุขาย เทวมนุสฺสานนฺ”ติ.}

\proseitem{51}{อารทฺธํ โข ปน เม พฺราหฺมณ วีริยํ อโหสิ อสลฺลีนํ, อุปฏฺฐิตา สติ อสมฺมุฏฺฐา\csromansharedfootnote{7e2a70728e3ee633}{อปฺปมฺมุฏฺฐา (สฺยา)}, ปสฺสทฺโธ กาโย อสารทฺโธ, สมาหิตํ จิตฺตํ เอกคฺคํ. โส โข อหํ พฺราหฺมณ วิวิจฺเจว กาเมหิ วิวิจฺจ อกุสเลหิ ธมฺเมหิ สวิตกฺกํ สวิจารํ วิเวกชํ ปีติสุขํ ปฐมํ ฌานํ อุปสมฺปชฺช วิหาสิํ. วิตกฺก\-วิจารานํ วูปสมา อชฺฌตฺตํ สมฺปสาทนํ เจตโส เอโกทิภาวํ อวิตกฺกํ อวิจารํ สมาธิชํ ปีติสุขํ ทุติยํ ฌานํ อุปสมฺปชฺช วิหาสิํ. ปีติยา จ วิราคา อุเปกฺขโก จ วิหาสิํ สโต จ สมฺปชาโน, สุขญฺจ กาเยน ปฏิสํเวเทสิํ, ยํ ตํ อริยา อาจิกฺขนฺติ “อุเปกฺขโก สติมา สุขวิหารี”ติ ตติยํ ฌานํ อุปสมฺปชฺช วิหาสิํ. สุขสฺส จ ปหานา ทุกฺขสฺส จ ปหานา ปุพฺเพว โสมนสฺส\-โทมนสฺสานํ อตฺถงฺคมา อทุกฺขมสุขํ อุเปกฺขา\-สติ\-ปาริสุทฺธิํ จตุตฺถํ ฌานํ อุปสมฺปชฺช วิหาสิํ.}

\proseitem{52}{โส เอวํ สมาหิเต จิตฺเต ปริสุทฺเธ ปริโยทาเต อนงฺคเณ วิคตู\-ปกฺกิเลเส มุทุภูเต กมฺมนิเย ฐิเต อาเนญฺชปฺปตฺเต ปุพฺเพ\-นิวาสา\-นุสฺสติ\-ญาณาย จิตฺตํ อภินินฺนาเมสิํ, โส อเนกวิหิตํ \csromanfolio{27}ปุพฺเพนิวาสํ อนุสฺสรามิ, เสยฺยถิทํ, เอกมฺปิ ชาติํ เทฺวปิ ชาติโย ติสฺโสปิ ชาติโย จตสฺโสปิ ชาติโย ปญฺจปิ ชาติโย ทสปิ ชาติโย วีสมฺปิ ชาติโย ติํสมฺปิ ชาติโย จตฺตาลีสมฺปิ ชาติโย ปญฺญาสมฺปิ ชาติโย ชาติสตมฺปิ ชาติสหสฺสมฺปิ ชาติ\-สต\-สหสฺสมฺปิ อเนเกปิ สํวฏฺฏกปฺเป อเนเกปิ วิวฏฺฏกปฺเป อเนเกปิ สํวฏฺฏ\-วิวฏฺฏ\-กปฺเป, “อมุตฺราสิํ เอวํนาโม เอวํโคตฺโต เอวํวณฺโณ เอวมาหาโร เอวํ\-สุขทุกฺข\-ปฺปฏิสํเวที เอวมา\-ยุปริยนฺโต, โส ตโต จุโต อมุตฺร อุทปาทิํ, ตตฺราปาสิํ เอวํนาโม เอวํโคตฺโต เอวํวณฺโณ เอวมาหาโร เอวํ\-สุขทุกฺข\-ปฺปฏิสํเวที เอวมา\-ยุปริยนฺโต, โส ตโต จุโต อิธูปปนฺโน”ติ, อิติ สาการํ สฺเอาทฺเทสํ อเนกวิหิตํ ปุพฺเพนิวาสํ อนุสฺสรามิ. อยํ โข เม พฺราหฺมณ รตฺติยา ปฐเม ยาเม ปฐมา วิชฺชา อธิคตา, อวิชฺชา วิหตา, วิชฺชา อุปฺปนฺนา, ตโม วิหโต, อาโลโก อุปฺปนฺโน, ยถา ตํ อปฺปมตฺตสฺส อาตาปิโน ปหิตตฺตสฺส วิหรโต.}

\proseitem{53}{โส เอวํ สมาหิเต จิตฺเต ปริสุทฺเธ ปริโยทาเต อนงฺคเณ วิคตู\-ปกฺกิเลเส มุทุภูเต กมฺมนิเย ฐิเต อาเนญฺชปฺปตฺเต สตฺตานํ จุตูปปาตญาณย จิตฺตํ อภินินฺนาเมสิํ, โส ทิพฺเพน จกฺขุนา วิสุทฺเธน อติกฺกนฺต\-มานุสเกน สตฺเต ปสฺสามิ จวมาเน อุปปชฺชมาเน หีเน ปณีเต สุวณฺเณ ทุพฺพณฺเณ สุคเต ทุคฺคเต, ยถากมฺมูปเค สตฺเต ปชานามิ “อิเม วต โภนฺโต สตฺตา กาย\-ทุจฺจริเตน สมนฺนาคตา วจีทุจฺจริเตน สมนฺนาคตา มโน\-ทุจฺจริเตน สมนฺนาคตา อริยานํ อุปวาทกา มิจฺฉาทิฏฺฐิกา มิจฺฉาทิฏฺฐิ\-กมฺม\-สมาทานา, เต กายสฺส เภทา ปรํ มรณา อปายํ ทุคฺคติํ วินิปาตํ นิรยํ อุปปนฺนา. อิเม วา ปน โภนฺโต สตฺตา กายสุจริเตน สมนฺนาคตา วจีสุจริเตน สมนฺนาคตา มโนสุจริเตน สมนฺนาคตา อริยานํ อนุปวาทกา สมฺมาทิฏฺฐิกา สมฺมา\-ทิฏฺฐิ\-กมฺม\-สมาทานา, เต กายสฺส เภทา ปรํ มรณา สุคติํ สคฺคํ โลกํ อุปปนฺนา”ติ, อิติ ทิพฺเพน จกฺขุนา วิสุทฺเธน อติกฺกนฺต\-มานุสเกน สตฺเต ปสฺสามิ จวมาเน อุปปชฺชมาเน หีเน ปณีเต สุวณฺเณ ทุพฺพณฺเณ สุคเต ทุคฺคเต, ยถากมฺมูปเค สตฺเต ปชานามิ. อยํ โข เม พฺราหฺมณ รตฺติยา มชฺฌิเม ยาเม ทุติยา วิชฺชา อธิคตา, อวิชฺชา วิหตา, วิชฺชา อุปฺปนฺนา, ตโม วิหโต, อาโลโก อุปฺปนฺโน, ยถา ตํ อปฺปมตฺตสฺส อาตาปิโน ปหิตตฺตสฺส วิหรโต.}

\proseitem{54}{\csromanfolio{28}โส เอวํ สมาหิเต จิตฺเต ปริสุทฺเธ ปริโยทาเต อนงฺคเณ วิคตู\-ปกฺกิเลเส มุทุภูเต กมฺมนิเย ฐิเต อาเนญฺชปฺปตฺเต อาสวานํ ขยญาณาย จิตฺตํ อภินินฺนาเมสิํ, โส “อิทํ ทุกฺขนฺ”ติ ยถาภูตํ อพฺภญฺญาสิํ, “อยํ ทุกฺขสมุทโย”ติ ยถาภูตํ อพฺภญฺญาสิํ, “อยํ ทุกฺขนิโรโธ”ติ ยถาภูตํ อพฺภญฺญาสิํ, “อยํ ทุกฺขนิโรธ\-คามินี ปฏิปทา”ติ ยถาภูตํ อพฺภญฺญาสิํ. “อิเม อาสวา”ติ ยถาภูตํ อพฺภญฺญาสิํ, “อยํ อาสวสมุทโย”ติ ยถาภูตํ อพฺภญฺญาสิํ, “อยํ อาสวนิโรโธ”ติ ยถาภูตํ อพฺภญฺญาสิํ, “อยํ อาสว\-นิโรธ\-คามินี ปฏิปทา”ติ ยถาภูตํ อพฺภญฺญาสิํ. ตสฺส เม เอวํ ชานโต เอวํ ปสฺสโต กามาสวาปิ จิตฺตํ วิมุจฺจิตฺถ, ภวาสวาปิ จิตฺตํ วิมุจฺจิตฺถ, อวิชฺชาสวาปิ จิตฺตํ วิมุจฺจิตฺถ. วิมุตฺตสฺมิํ “วิมุตฺตมฺ”อิติ ญาณํ อโหสิ, “ขีณาชาติ, วุสิตํ พฺรหฺมจริยํ, กตํ กรณียํ, นาปรํ อิตฺถตฺตายา”ติ อพฺภญฺญาสิํ. อยํ โข เม พฺราหฺมณ รตฺติยา ปจฺฉิเม ยาเม ตติยา วิชฺชา อธิคตา, อวิชฺชา วิหตา, วิชฺชา อุปฺปนฺนา, ตโม วิหโต, อาโลโก อุปฺปนฺโน, ยถา ตํ อปฺปมตฺตสฺส อาตาปิโน ปหิตตฺตสฺส วิหรโต.}

\proseitem{55}{สิยา โข ปน เต พฺราหฺมณ เอวมสฺส “อชฺชาปิ นูน สมโณ โคตโม อวีตราโค อวีตโทโส อวีตโมโห, ตสฺมา อรญฺญ\-วนปตฺถานิ ปนฺตานิ เสนาสนานิ ปฏิเสวตี”ติ. น โข ปเนตํ พฺราหฺมณ เอวํ ทฏฺฐพฺพํ. เทฺว โข อหํ พฺราหฺมณ อตฺถวเส สมฺปสฺสมาโน อรญฺญ\-วนปตฺถานิ ปนฺตานิ เสนาสนานิ ปฏิเสวามิ, อตฺตโน จ ทิฏฺฐ\-ธมฺม\-สุข\-วิหารํ สมฺปสฺสมาโน ปจฺฉิมญฺจ ชนตํ อนุกมฺปมาโนติ.}

\proseitemwithcloser{56}{อนุกมฺปิตรูปา วตายํ โภตา โคตเมน ปจฺฉิมา ชนตา, ยถา ตํ อรหตา สมฺมา\-สมฺพุทฺเธน. อภิกฺกนฺตํ โภ โคตม, อภิกฺกนฺตํ โภ โคตม, เสยฺยถาปิ โภ โคตม นิกฺกุชฺชิตํ วา อุกฺกุชฺเชยฺย, ปฏิจฺฉนฺนํ วา วิวเรยฺย, มูฬฺหสฺส วา มคฺคํ อาจิกฺเขยฺย, อนฺธกาเร วา เตลปชฺโชตํ ธาเรยฺย “จกฺขุมนฺโต รูปานิ ทกฺขนฺตี”ติ, เอวเมวํ โภตา โคตเมน อเนก\-ปริยาเยน ธมฺโม ปกาสิโต, เอสาหํ ภวนฺตํ โคตมํ สรณํ คจฺฉามิ ธมฺมญฺจ ภิกฺขุสํฆญฺจ, อุปาสกํ มํ ภวํ โคตโม ธาเรตุ อชฺชตคฺเค ปาณุเปตํ สรณํ คตนฺติ.}{ภยเภรว\-สุตฺตํ นิฏฺฐิตํ จตุตฺถํ.}
\csromanmatikaanchor{mtk.13}
\csromantocmarkref{subsection}{5. อนงฺคณสุตฺต}{mtk.13}
\bookmark[dest={mtk.13},level=2]{5. อนงฺคณสุตฺต}
\csromanheader{2}{5. อนงฺคณสุตฺต \textbf{5. อนงฺคณสุตฺต}}
\proseitem{57}{\csromanfolio{29}เอวํ เม สุตํ\csromandash{}เอกํ สมยํ ภควา สาวตฺถิยํ วิหรติ เชตวเน อนาถ\-ปิณฺฑิกสฺส อาราเม. ตตฺร โข อายสฺมา สาริปุตฺโต ภิกฺขู อามนฺเตสิ “อาวุโส ภิกฺขเว”ติ. “อาวุโส”ติ โข เต ภิกฺขู อายสฺมโต สาริปุตฺตสฺส ปจฺจสฺโสสุํ. อายสฺมา สาริปุตฺโต เอตทโวจ\csromandash{}}

\prose{จตฺตาโรเม อาวุโส ปุคฺคลา สนฺโต สํวิชฺชมานา โลกสฺมิํ, กตเม จตฺตาโร, อิธาวุโส เอกจฺโจ ปุคฺคโล สางฺคโณว สมาโน “อตฺถิ เม อชฺฌตฺตํ องฺคณนฺ”ติ ยถาภูตํ นปฺปชานาติ, อิธ ปนาวุโส เอกจฺโจ ปุคฺคโล สางฺคโณว สมาโน “อตฺถิ เม อชฺฌตฺตํ องฺคณนฺ”ติ ยถาภูตํ ปชานาติ, อิธาวุโส เอกจฺโจ ปุคฺคโล อนงฺคโณว สมาโน “นตฺถิ เม อชฺฌตฺตํ องฺคณนฺ”ติ ยถาภูตํ นปฺปชานาติ, อิธ ปนาวุโส เอกจฺโจ ปุคฺคโล อนงฺคโณว สมาโน “นตฺถิ เม อชฺฌตฺตํ องฺคณนฺ”ติ ยถาภูตํ ปชานาติ. ตตฺราวุโส ยฺวายํ ปุคฺคโล สางฺคโณว สมาโน “อตฺถิ เม อชฺฌตฺตํ องฺคณนฺ”ติ ยถาภูตํ นปฺปชานาติ, อยํ อิเมสํ ทฺวินฺนํ ปุคฺคลานํ สางฺคณานํเยว สตํ หีนปุริโส อกฺขายติ. ตตฺราวุโส ยฺวายํ ปุคฺคโล สางฺคโณว สมาโน “อตฺถิ เม อชฺฌตฺตํ องฺคณนฺ”ติ ยถาภูตํ ปชานาติ, อยํ อิเมสํ ทฺวินฺนํ ปุคฺคลานํ สางฺคณานํเยว สตํ เสฏฺฐปุริโส อกฺขายติ. ตตฺราวุโส ยฺวายํ ปุคฺคโล อนงฺคโณว สมาโน “นตฺถิ เม อชฺฌตฺตํ องฺคณนฺ”ติ ยถาภูตํ นปฺปชานาติ, อยํ อิเมสํ ทฺวินฺนํ ปุคฺคลานํ อนงฺคณานํเยว สตํ หีนปุริโส อกฺขายติ. ตตฺราวุโส ยฺวายํ ปุคฺคโล อนงฺคโณว สมาโน “นตฺถิ เม อชฺฌตฺตํ องฺคณนฺ”ติ ยถาภูตํ ปชานาติ, อยํ อิเมสํ ทฺวินฺนํ ปุคฺคลานํ อนงฺคณานํเยว สตํ เสฏฺฐปุริโส อกฺขายตีติ.}

\proseitem{58}{เอวํ วุตฺเต อายสฺมา มหาโมคฺคลฺลาโน อายสฺมนฺตํ สาริปุตฺตํ เอตทโวจ\csromandash{}}

\prose{“โก นุ โข อาวุโส สาริปุตฺต เหตุ โก ปจฺจโย, เยนิเมสํ ทฺวินฺนํ ปุคฺคลานํ สางฺคณานํเยว สตํ เอโก หีนปุริโส อกฺขายติ, เอโก เสฏฺฐปุริโส อกฺขายติ. โก ปนาวุโส สาริปุตฺต เหตุ \csromanfolio{30}โก ปจฺจโย, เยนิเมสํ ทฺวินฺนํ ปุคฺคลานํ อนงฺคณานํเยว สตํ เอโก หีนปุริโส อกฺขายติ, เอโก เสฏฺฐปุริโส อกฺขายตี”ติ.}

\proseitem{59}{ตตฺราวุโส ยฺวายํ ปุคฺคโล สางฺคโณว สมาโน”อตฺถิ เม อชฺฌตฺตํ องฺคณนฺ”ติ ยถาภูตํ นปฺปชานาติ, ตสฺเสตํ ปาฏิกงฺขํ ‘น ฉนฺทํ ชเนสฺสติ น วายมิสฺสติ น วีริยํ อารภิสฺสติ ตสฺสงฺคณสฺส ปหานาย. โส สราโค สโทโส สโมโห สางฺคโณ สํกิลิฏฺฐ\-จิตฺโต กาลํ กริสฺสติ’. เสยฺยถาปิ อาวุโส กํสปาติ อาภตา อาปณา วา กมฺมารกุลา วา รเชน จ มเลน จ ปริโยนทฺธา, ตเมนํ สามิกา น เจว ปริภุญฺเชยฺยุํ น จ ปริโยทเปยฺยุํ\csromansharedfootnote{8405e6c03000ec98}{ปริโยทาเปยฺยุํ (?)}, รชาปเถ จ นํ นิกฺขิเปยฺยุํ. เอวญฺหิ สา อาวุโส กํสปาติ อปเรน สมเยน สํกิลิฏฺฐตรา อสฺส มลคฺคหิตาติ. เอวมาวุโสติ. เอวเมว โข อาวุโส ยฺวายํ ปุคฺคโล สางฺคโณว สมาโน “อตฺถิ เม อชฺฌตฺตํ องฺคณนฺ”ติ ยถาภูตํ นปฺปชานาติ, ตสฺเสตํ ปาฏิกงฺขํ ‘น ฉนฺทํ ชเนสฺสติ น วายมิสฺสติ น วีริยํ อารภิสฺสติ ตสฺสงฺคณสฺส ปหานาย. โส สราโค สโทโส สโมโห สางฺคโณ สํกิลิฏฺฐ\-จิตฺโต กาลํ กริสฺสติ’.}

\prose{ตตฺราวุโส ยฺวายํ ปุคฺคโล สางฺคโณว สมาโน “อตฺถิ เม อชฺฌตฺตํ องฺคณนฺ”ติ ยถาภูตํ ปชานาติ, ตสฺเสตํ ปาฏิกงฺขํ ‘ฉนฺทํ ชเนสฺสติ วายมิสฺสติ วีริยํ อารภิสฺสติ ตสฺสงฺคณสฺส ปหานาย. โส อราโค อโทโส อโมโห อนงฺคโณ อสํกิลิฏฺฐ\-จิตฺโต กาลํ กริสฺสติ’. เสยฺยถาปิ อาวุโส กํสปาติ อาภตา อาปณา วา กมฺมารกุลา วา รเชน จ มเลน จ ปริโยนทฺธา, ตเมนํ สามิกา ปริภุญฺเชยฺยุํ เจว ปริโยทเปยฺยุํ จ, น จ นํ รชาปเถ นิกฺขิเปยฺยุํ. เอวญฺหิ สา อาวุโส กํสปาติ อปเรน สมเยน ปริสุทฺธตรา อสฺส ปริโยทาตาติ. เอวมาวุโสติ. เอวเมว โข อาวุโส ยฺวายํ ปุคฺคโล สางฺคโณว สมาโน “อตฺถิ เม อชฺฌตฺตํ องฺคณนฺ”ติ ยถาภูตํ ปชานาติ, ตสฺเสตํ ปาฏิกงฺขํ ‘ฉนฺทํ ชเนสฺสติ วายมิสฺสติ วีริยํ อารภิสฺสติ ตสฺสงฺคณสฺส ปหานาย. โส อราโค อโทโส อโมโห อนงฺคโณ อสํกิลิฏฺฐ\-จิตฺโต กาลํ กริสฺสติ’.}

\prose{\csromanfolio{31}ตตฺราวุโส ยฺวายํ ปุคฺคโล อนงฺคโณว สมาโน “นตฺถิ เม อชฺฌตฺตํ องฺคณนฺ”ติ ยถาภูตํ นปฺปชานาติ, ตสฺเสตํ ปาฏิกงฺขํ ‘สุภนิมิตฺตํ มนสิ กริสฺสติ, ตสฺส สุภนิมิตฺตสฺส มนสิการา ราโค จิตฺตํ อนุทฺธํเสสฺสติ. โส สราโค สโทโส สโมโห สางฺคโณ สํกิลิฏฺฐ\-จิตฺโต กาลํ กริสฺสติ’. เสยฺยถาปิ อาวุโส กํสปาติ อาภตา อาปณา วา กมฺมารกุลา วา ปริสุทฺธา ปริโยทาตา, ตเมนํ สามิกา น เจว ปริภุญฺเชยฺยุํ น จ ปริโยทเปยฺยุํ, รชาปเถ จ นํ นิกฺขิเปยฺยุํ. เอวญฺหิ สา อาวุโส กํสปาติ อปเรน สมเยน สํกิลิฏฺฐตรา อสฺส มลคฺคหิตาติ. เอวมาวุโสติ. เอวเมว โข อาวุโส ยฺวายํ ปุคฺคโล อนงฺคโณว สมาโน “นตฺถิ เม อชฺฌตฺตํ องฺคณนฺ”ติ ยถาภูตํ นปฺปชานาติ, ตสฺเสตํ ปาฏิกงฺขํ ‘สุภนิมิตฺตํ มนสิ กริสฺสติ, ตสฺส สุภนิมิตฺตสฺส มนสิการา ราโค จิตฺตํ อนุทฺธํเสสฺสติ. โส สราโค สโทโส สโมโห สางฺคโณ สํกิลิฏฺฐ\-จิตฺโต กาลํ กริสฺสติ’.}

\prose{ตตฺราวุโส ยฺวายํ ปุคฺคโล อนงฺคโณว สมาโน ‘นตฺถิ เม อชฺฌตฺตํ องฺคณนฺ”ติ ยถาภูตํ ปชานาติ, ตสฺเสตํ ปาฏิกงฺขํ ‘สุภนิมิตฺตํ น มนสิ กริสฺสติ, ตสฺส สุภนิมิตฺตสฺส อมนสิการา ราโค จิตฺตํ นานุทฺธํเสสฺสติ. โส อราโค อโทโส อโมโห อนงฺคโณ อสํกิลิฏฺฐ\-จิตฺโต กาลํ กริสฺสติ’. เสยฺยถาปิ อาวุโส กํสปาติ อาภตา อาปณา วา กมฺมารกุลา วา ปริสุทฺธา ปริโยทาตา, ตเมนํ สามิกา ปริภุญฺเชยฺยุํ เจว ปริโยทเปยฺยุํ จ, น จ นํ รชาปเถ นิกฺขิเปยฺยุํ. เอวญฺหิ สา อาวุโส กํสปาติ อปเรน สมเยน ปริสุทฺธตรา อสฺส ปริโยทาตาติ. เอวมาวุโสติ. เอวเมว โข อาวุโส ยฺวายํ ปุคฺคโล อนงฺคโณว สมาโน “นตฺถิ เม อชฺฌตฺตํ องฺคณนฺ”ติ ยถาภูตํ ปชานาติ, ตสฺเสตํ ปาฏิกงฺขํ ‘สุภนิมิตฺตํ น มนสิ กริสฺสติ, ตสฺส สุภนิมิตฺตสฺส อมนสิการา ราโค จิตฺตํ นานุทฺธํเสสฺสติ. โส อราโค อโทโส อโมโห อนงฺคโณ อสํกิลิฏฺฐ\-จิตฺโต กาลํ กริสฺสติ’.}

\prose{อยํ โข อาวุโส โมคฺคลฺลาน เหตุ อยํ ปจฺจโย, เยนิเมสํ ทฺวินฺนํ ปุคฺคลานํ สางฺคณานํเยว สตํ เอโก หีนปุริโส อกฺขายติ, เอโก เสฏฺฐปุริโส อกฺขายติ. อยํ ปนาวุโส โมคฺคลฺลาน เหตุ}

\prose{\csromanfolio{32}อยํ ปจฺจโย, เยนิเมสํ ทฺวินฺนํ ปุคฺคลานํ อนงฺคณานํเยว สตํ เอโก หีนปุริโส อกฺขายติ, เอโก เสฏฺฐปุริโส อกฺขายตีติ.}

\proseitem{60}{“องฺคณํ องฺคณนฺ”ติ อาวุโส วุจฺจติ. กิสฺส นุ โข เอตํ อาวุโส อธิวจนํ, ยทิทํ องฺคณนฺติ. ปาปกานํ โข เอตํ อาวุโส อกุสลานํ อิจฺฉาวจรานํ อธิวจนํ, ยทิทํ องฺคณนฺติ.}

\prose{ฐานํ โข ปเนตํ อาวุโส วิชฺชติ, ยํ อิเธกจฺจสฺส ภิกฺขุโน เอวํ อิจฺฉา อุปฺปชฺเชยฺย “อาปตฺติํ จ วต อาปนฺโน อสฺสํ, น จ มํ ภิกฺขู ชาเนยฺยุํ อาปตฺติํ อาปนฺโน”ติ. ฐานํ โข ปเนตํ อาวุโส วิชฺชติ, ยํ ตํ ภิกฺขุํ ภิกฺขู ชาเนยฺยุํ “อาปตฺติํ อาปนฺโน”ติ. “ชานนฺติ มํ ภิกฺขู อาปตฺติํ อาปนฺโน”ติ อิติ โส กุปิโต โหติ อปฺปตีโต. โย เจว โข อาวุโส โกโป โย จ อปฺปจฺจโย, อุภยเมตํ องฺคณํ.}

\prose{ฐานํ โข ปเนตํ อาวุโส วิชฺชติ, ยํ อิเธกจฺจสฺส ภิกฺขุโน เอวํ อิจฺฉา อุปฺปชฺเชยฺย “อาปตฺติํ จ วต อาปนฺโน อสฺสํ, อนุรโห มํ ภิกฺขู โจเทยฺยุํ, โน สํฆมชฺเฌ”ติ. ฐานํ โข ปเนตํ อาวุโส วิชฺชติ, ยํ ตํ ภิกฺขุํ ภิกฺขู สํฆมชฺเฌ โจเทยฺยุํ, โน อนุรโห. “สํฆมชฺเฌ มํ ภิกฺขู โจเทนฺติ, โน อนุรโห”ติ อิติ โส กุปิโต โหติ อปฺปตีโต. โย เจว โข อาวุโส โกโป โย จ อปฺปจฺจโย, อุภยเมตํ องฺคณํ.}

\prose{ฐานํ โข ปเนตํ อาวุโส วิชฺชติ, ยํ อิเธกจฺจสฺส ภิกฺขุโน เอวํ อิจฺฉา อุปฺปชฺเชยฺย “อาปตฺติํ จ วต อาปนฺโน อสฺสํ, สปฺปฏิปุคฺคโล มํ โจเทยฺย, โน อปฺปฏิปุคฺคโล”ติ. ฐานํ โข ปเนตํ อาวุโส วิชฺชติ, ยํ ตํ ภิกฺขุํ อปฺปฏิปุคฺคโล โจเทยฺย, โน สปฺปฏิปุคฺคโล. “อปฺปฏิปุคฺคโล มํ โจเทติ, โน สปฺปฏิปุคฺคโล”ติ อิติ โส กุปิโต โหติ อปฺปตีโต. โย เจว โข อาวุโส โกโป โย จ อปฺปจฺจโย, อุภยเมตํ องฺคณํ.}

\prose{ฐานํ โข ปเนตํ อาวุโส วิชฺชติ, ยํ อิเธกจฺจสฺส ภิกฺขุโน เอวํ อิจฺฉา อุปฺปชฺเชยฺย “อโห วต มเมว สตฺถา ปฏิปุจฺฉิตฺวา ปฏิปุจฺฉิตฺวา ภิกฺขูนํ ธมฺมํ เทเสยฺย, น อญฺญํ ภิกฺขุํ สตฺถา ปฏิปุจฺฉิตฺวา ปฏิปุจฺฉิตฺวา ภิกฺขูนํ ธมฺมํ เทเสยฺยา”ติ. ฐานํ โข ปเนตํ อาวุโส วิชฺชติ, ยํ อญฺญํ ภิกฺขุํ สตฺถา \csromanfolio{33}ปฏิปุจฺฉิตฺวา ปฏิปุจฺฉิตฺวา ภิกฺขูนํ ธมฺมํ เทเสยฺย, น ตํ ภิกฺขุํ สตฺถา ปฏิปุจฺฉิตฺวา ปฏิปุจฺฉิตฺวา ภิกฺขูนํ ธมฺมํ เทเสยฺย. “อญฺญํ ภิกฺขุํ สตฺถา ปฏิปุจฺฉิตฺวา ปฏิปุจฺฉิตฺวา ภิกฺขูนํ ธมฺมํ เทเสติ, น มํ สตฺถา ปฏิปุจฺฉิตฺวา ปฏิปุจฺฉิตฺวา ภิกฺขูนํ ธมฺมํ เทเสตี”ติ อิติ โส กุปิโต โหติ อปฺปตีโต. โย เจว โข อาวุโส โกโป โย จ อปฺปจฺจโย, อุภยเมตํ องฺคณํ.}

\prose{ฐานํ โข ปเนตํ อาวุโส วิชฺชติ, ยํ อิเธกจฺจสฺส ภิกฺขุโน เอวํ อิจฺฉา อุปฺปชฺเชยฺย “อโห วต มเมว ภิกฺขู ปุรกฺขตฺวา ปุรกฺขตฺวา คามํ ภตฺตาย ปวิเสยฺยุํ, น อญฺญํ ภิกฺขุํ ภิกฺขู ปุรกฺขตฺวา ปุรกฺขตฺวา คามํ ภตฺตาย ปวิเสยฺยุนฺ”ติ. ฐานํ โข ปเนตํ อาวุโส วิชฺชติ, ยํ อญฺญํ ภิกฺขุํ ภิกฺขู ปุรกฺขตฺวา ปุรกฺขตฺวา คามํ ภตฺตาย ปวิเสยฺยุํ, น ตํ ภิกฺขุํ ภิกฺขู ปุรกฺขตฺวา ปุรกฺขตฺวา คามํ ภตฺตาย ปวิเสยฺยุํ. “อญฺญํ ภิกฺขุํ ภิกฺขู ปุรกฺขตฺวา ปุรกฺขตฺวา คามํ ภตฺตาย ปวิสนฺติ, น มํ ภิกฺขู ปุรกฺขตฺวา ปุรกฺขตฺวา คามํ ภตฺตาย ปวิสนฺตี”ติ อิติ โส กุปิโต โหติ อปฺปตีโต. โย เจว โข อาวุโส โกโป โย จ อปฺปจฺจโย, อุภยเมตํ องฺคณํ.}

\prose{ฐานํ โข ปเนตํ อาวุโส วิชฺชติ, ยํ อิเธกจฺจสฺส ภิกฺขุโน เอวํ อิจฺฉา อุปฺปชฺเชยฺย “อโห วต อหเมว ลเภยฺยํ ภตฺตคฺเค อคฺคาสนํ อคฺโคทกํ อคฺคปิณฺฑํ, น อญฺโญ ภิกฺขุ ลเภยฺย ภตฺตคฺเค อคฺคาสนํ อคฺโคทกํ อคฺคปิณฺฑนฺ”ติ. ฐานํ โข ปเนตํ อาวุโส วิชฺชติ, ยํ อญฺโญ ภิกฺขุ ลเภยฺย ภตฺตคฺเค อคฺคาสนํ อคฺโคทกํ อคฺคปิณฺฑํ, น โส ภิกฺขุ ลเภยฺย ภตฺตคฺเค อคฺคาสนํ อคฺโคทกํ อคฺคปิณฺฑํ. “อญฺโญ ภิกฺขุ ลภติ ภตฺตคฺเค อคฺคาสนํ อคฺโคทกํ อคฺคปิณฺฑํ, นาหํ ลภามิ ภตฺตคฺเค อคฺคาสนํ อคฺโคทกํ อคฺคปิณฺฑนฺ”ติ อิติ โส กุปิโต โหติ อปฺปตีโต. โย เจว โข อาวุโส โกโป โย จ อปฺปจฺจโย, อุภยเมตํ องฺคณํ.}

\prose{ฐานํ โข ปเนตํ อาวุโส วิชฺชติ, ยํ อิเธกจฺจสฺส ภิกฺขุโน เอวํ อิจฺฉา อุปฺปชฺเชยฺย “อโห วต อหเมว ภตฺตคฺเค ภุตฺตาวี อนุโมเทยฺยํ, น อญฺโญ ภิกฺขุ ภตฺตคฺเค ภุตฺตาวี อนุโมเทยฺยา”ติ. ฐานํ โข ปเนตํ อาวุโส วิชฺชติ, ยํ อญฺโญ ภิกฺขุ ภตฺตคฺเค ภุตฺตาวี อนุโมเทยฺย, น โส ภิกฺขุ ภตฺตคฺเค ภุตฺตาวี อนุโมเทยฺย. “อญฺโญ ภิกฺขุ ภตฺตคฺเค \csromanfolio{34}ภุตฺตาวี อนุโมทติ, นาหํ ภตฺตคฺเค ภุตฺตาวี อนุโมทามี”ติ อิติ โส กุปิโต โหติ อปฺปตีโต. โย เจว โข อาวุโส โกโป โย จ อปฺปจฺจโย, อุภยเมตํ องฺคณํ.}

\proseitem{5}{ฐานํ โข ปเนตํ อาวุโส วิชฺชติ, ยํ อิเธกจฺจสฺส ภิกฺขุโน เอวํ อิจฺฉา อุปฺปชฺเชยฺย “อโห วต อหเมว อารามคตานํ ภิกฺขูนํ ธมฺมํ เทเสยฺยํ, น อญฺโญ ภิกฺขุ อารามคตานํ ภิกฺขูนํ ธมฺมํ เทเสยฺยา”ติ. ฐานํ โข ปเนตํ อาวุโส วิชฺชติ, ยํ อญฺโญ ภิกฺขุ อารามคตานํ ภิกฺขูนํ ธมฺมํ เทเสยฺย, น โส ภิกฺขุ อารามคตานํ ภิกฺขูนํ ธมฺมํ เทเสยฺย. “อญฺโญ ภิกฺขุ อารามคตานํ ภิกฺขูนํ ธมฺมํ เทเสติ, นาหํ อารามคตานํ ภิกฺขูนํ ธมฺมํ เทเสมี”ติ อิติ โส กุปิโต โหติ อปฺปตีโต. โย เจว โข อาวุโส โกโป โย จ อปฺปจฺจโย, อุภยเมตํ องฺคณํ.}

\prose{ฐานํ โข ปเนตํ อาวุโส วิชฺชติ, ยํ อิเธกจฺจสฺส ภิกฺขุโน เอวํ อิจฺฉา อุปฺปชฺเชยฺย “อโห วต อหเมว อารามคตานํ ภิกฺขุนีนํ ธมฺมํ เทเสยฺยํ ฯเปฯ อุปาสกานํ ธมฺมํ เทเสยฺยํ ฯเปฯ อุปาสิกานํ ธมฺมํ เทเสยฺยํ, น อญฺโญ ภิกฺขุ อารามคตานํ อุปาสิกานํ ธมฺมํ เทเสยฺยา”ติ. ฐานํ โข ปเนตํ อาวุโส วิชฺชติ, ยํ อญฺโญ ภิกฺขุ อารามคตานํ อุปาสิกานํ ธมฺมํ เทเสยฺย, น โส ภิกฺขุ อารามคตานํ อุปาสิกานํ ธมฺมํ เทเสยฺย. “อญฺโญ ภิกฺขุ อารามคตานํ อุปาสิกานํ ธมฺมํ เทเสติ, นาหํ อารามคตานํ อุปาสิกานํ ธมฺมํ เทเสมี”ติ อิติ โส กุปิโต โหติ อปฺปตีโต. โย เจว โข อาวุโส โกโป โย จ อปฺปจฺจโย, อุภยเมตํ องฺคณํ.}

\prose{ฐานํ โข ปเนตํ อาวุโส วิชฺชติ, ยํ อิเธกจฺจสฺส ภิกฺขุโน เอวํ อิจฺฉา อุปฺปชฺเชยฺย “อโห วต มเมว ภิกฺขู สกฺกเรยฺยุํ ครุํ กเรยฺยุํ\csromansharedfootnote{1fa03322c836f9e8}{ครุกเรยฺยุํ (สี, สฺยา, อิ)} มาเนยฺยุํ ปูเชยฺยุํ, น อญฺญํ ภิกฺขุํ ภิกฺขู สกฺกเรยฺยุํ ครุํ กเรยฺยุํ มาเนยฺยุํ ปูเชยฺยุนฺ”ติ. ฐานํ โข ปเนตํ อาวุโส วิชฺชติ, ยํ อญฺญํ ภิกฺขุํ ภิกฺขู สกฺกเรยฺยุํ ครุํ กเรยฺยุํ มาเนยฺยุํ ปูเชยฺยุํ, น ตํ ภิกฺขุํ ภิกฺขู สกฺกเรยฺยุํ ครุํ กเรยฺยุํ มาเนยฺยุํ ปูเชยฺยุํ. “อญฺญํ ภิกฺขุํ ภิกฺขู สกฺกโรนฺติ ครุํ กโรนฺติ มาเนนฺติ ปูเชนฺติ, น มํ ภิกฺขู สกฺกโรนฺติ \csromanfolio{35}ครุํ กโรนฺติ มาเนนฺติ ปุเชนฺตี”ติ อิติ โส กุปิโต โหติ อปฺปตีโต. โย เจว โข อาวุโส โกโป โย จ อปฺปจฺจโย, อุภยเมตํ องฺคณํ.}

\prose{ฐานํ โข ปเนตํ อาวุโส วิชฺชติ, ยํ อิเธกจฺจสฺส ภิกฺขุโน เอวํ อิจฺฉา อุปฺปชฺเชยฺย “อโห วต มเมว ภิกฺขุนิโย ฯเปฯ อุปาสกา ฯเปฯ อุปาสิกา สกฺกเรยฺยุํ ครุํ กเรยฺยุํ มาเนยฺยุํ ปูเชยฺยุํ, น อญฺญํ ภิกฺขุํ อุปาสิกา สกฺกเรยฺยุํ ครุํ กเรยฺยุํ มาเนยฺยุํ ปูเชยฺยุนฺ”ติ. ฐานํ โข ปเนตํ อาวุโส วิชฺชติ, ยํ อญฺญํ ภิกฺขุํ อุปาสิกา สกฺกเรยฺยุํ ครุํ กเรยฺยุํ มาเนยฺยุํ ปูเชยฺยุํ, น ตํ ภิกฺขุํ อุปาสิกา สกฺกเรยฺยุํ ครุํ กเรยฺยุํ มาเนยฺยุํ ปูเชยฺยุํ. “อญฺญํ ภิกฺขุํ อุปาสิกา สกฺกโรนฺติ ครุํ กโรนฺติ มาเนนฺติ ปูเชนฺติ, น มํ อุปาสิกา สกฺกโรนฺติ ครุํ กโรนฺติ มาเนนฺติ ปูเชนฺตี”ติ อิติ โส กุปิโต โหติ อปฺปตีโต. โย เจว โข อาวุโส โกโป โย จ อปฺปจฺจโย, อุภยเมตํ องฺคณํ.}

\prose{ฐานํ โข ปเนตํ อาวุโส วิชฺชติ, ยํ อิเธกจฺจสฺส ภิกฺขุโน เอวํ อิจฺฉา อุปฺปชฺเชยฺย “อโห วต อหเมว ลาภี อสฺสํ ปณีตานํ จีวรานํ, น อญฺโญ ภิกฺขุ ลาภี อสฺส ปณีตานํ จีวรานนฺ”ติ. ฐานํ โข ปเนตํ อาวุโส วิชฺชติ, ยํ อญฺโญ ภิกฺขุ ลาภี อสฺส ปณีตานํ จีวรานํ, น โส ภิกฺขุ ลาภี อสฺส ปณีตานํ จีวรานํ. “อญฺโญ ภิกฺขุ ลาภี\csromansharedfootnote{455b5b17cac02bce}{ลาภี อสฺส (ก)} ปณีตานํ จีวรานํ, นาหํ ลาภี\csromansharedfootnote{a2a5890a62f62463}{ลาภี อสฺสํ (ก)} ปณีตานํ จีวรานนฺ”ติ อิติ โส กุปิโต โหติ อปฺปตีโต. โย เจว โข อาวุโส โกโป โย จ อปฺปจฺจโย, อุภยเมตํ องฺคณํ.}

\prose{ฐานํ โข ปเนตํ อาวุโส วิชฺชติ, ยํ อิเธกจฺจสฺส ภิกฺขุโน เอวํ อิจฺฉา อุปฺปชฺเชยฺย “อโห วต อหเมว ลาภี อสฺสํ ปณีตานํ ปิณฺฑปาตานํ ฯเปฯ ปณีตานํ เสนาสนานํ ฯเปฯ ปณีตานํ คิลานปฺปจฺจย\-เภสชฺช\-ปริกฺขารานํ, น อญฺโญ ภิกฺขุ ลาภี อสฺส ปณีตานํ คิลานปฺปจฺจยเภสชฺชปริกฺขารานนฺ”ติ. ฐานํ โข ปเนตํ อาวุโส วิชฺชติ, ยํ อญฺโญ ภิกฺขุ ลาภี อสฺส ปณีตานํ คิลานปฺปจฺจย\-เภสชฺช\-ปริกฺขารานํ, น โส ภิกฺขุ ลาภี อสฺส ปณีตานํ คิลานปฺปจฺจย\-เภสชฺช\-ปริกฺขารานํ. “อญฺโญ \csromanfolio{36}ภิกฺขุ ลาภี\csromansharedfootnote{455b5b17cac02bce}{ลาภี อสฺส (ก)} ปณีตานํ คิลานปฺปจฺจย\-เภสชฺช\-ปริกฺขารานํ, นาหํ ลาภี\csromansharedfootnote{a2a5890a62f62463}{ลาภี อสฺสํ (ก)} ปณีตานํ คิลานปฺปจฺจยเภสชฺชปริกฺขารานนฺ”ติ อิติ โส กุปิโต โหติ อปฺปตีโต. โย เจว โข อาวุโส โกโป โย จ อปฺปจฺจโย, อุภยเมตํ องฺคณํ. อิเมสํ โข เอตํ อาวุโส ปาปกานํ อกุสลานํ อิจฺฉาวจรานํ อธิวจนํ ยทิทํ องฺคณนฺติ.}

\proseitem{61}{ยสฺส กสฺสจิ อาวุโส ภิกฺขุโน อิเม ปาปกา อกุสลา อิจฺฉาวจรา อปฺปหีนา ทิสฺสนฺติ เจว สูยนฺติ จ, กิญฺจาปิ โส โหติ อารญฺญิโก ปนฺตเสนาสโน ปิณฺฑปาติโก สปทานจารี ปํสุกูลิโก ลูขจีวร\-ธโร, อถ โข นํ สพฺรหฺมจารี น เจว สกฺกโรนฺติ น ครุํ กโรนฺติ น มาเนนฺติ น ปูเชนฺติ. ตํ กิสฺส เหตุ, เต หิ ตสฺส อายสฺมโต ปาปกา อกุสลา อิจฺฉาวจรา อปฺปหีนา ทิสฺสนฺติ เจว สูยนฺติ จ. เสยฺยถาปิ อาวุโส กํสปาติ อาภตา อาปณา วา กมฺมารกุลา วา ปริสุทฺธา ปริโยทาตา, ตเมนํ สามิกา อหิกุณปํ วา กุกฺกุรกุณปํ วา มนุสฺสกุณปํ วา รจยิตฺวา อญฺญิสฺสา กํสปาติยา ปฏิกุชฺชิตฺวา อนฺตราปณํ ปฏิปชฺเชยฺยุํ, ตเมนํ ชโน ทิสฺวา เอวํ วเทยฺย “อมฺโภ กิเมวิทํ หรียติ ชญฺญชญฺญํ วิยา”ติ. ตเมนํ อุฏฺฐหิตฺวา อปาปุริตฺวา\csromansharedfootnote{d0291a320c4283de}{อวาปุริตฺวา (สี)} โอโลเกยฺย, ตสฺส สหทสฺสเนน อมนาปตา จ สณฺฐเหยฺย, ปาฏิกุลฺยตา\csromansharedfootnote{32e62137a0cbf256}{ปฏิกูลตา (ก), ปาฏิกูลฺยตา (สฺยา)} จ สณฺฐเหยฺย, เชคุจฺฉตา จ\csromansharedfootnote{39aaa7beecdf09a0}{เชคุจฺฉิตา จ (อิ, ก)} สณฺฐเหยฺย, ชิฆจฺฉิตานํปิ น โภตฺตุกมฺยตา อสฺส, ปเคว สุหิตานํ. เอวเมว โข อาวุโส ยสฺส กสฺสจิ ภิกฺขุโน อิเม ปาปกา อกุสลา อิจฺฉาวจรา อปฺปหีนา ทิสฺสนฺติ เจว สูยนฺติ จ, กิญฺจาปิ โส โหติ อารญฺญิโก ปนฺตเสนาสโน ปิณฺฑปาติโก สปทานจารี ปํสุกูลิโก ลูขจีวร\-ธโร, อถ โข นํ สพฺรหฺมจารี น เจว สกฺกโรนฺติ น ครุํ กโรนฺติ น มาเนนฺติ น ปูเชนฺติ. ตํ กิสฺส เหตุ, เต หิ ตสฺส อายสฺมโต ปาปกา อกุสลา อิจฺฉาวจรา อปฺปหีนา ทิสฺสนฺติ เจว สูยนฺติ จ.}

\proseitem{62}{ยสฺส กสฺสจิ อาวุโส ภิกฺขุโน อิเม ปาปกา อกุสลา อิจฺฉาวจรา ปหีนา ทิสฺสนฺติ เจว สูยนฺติ จ, กิญฺจาปิ โส โหติ \csromanfolio{37}คามนฺตวิหารี เนมนฺตนิโก คหปติจีวร\-ธโร, อถ โข นํ สพฺรหฺมจารี สกฺกโรนฺติ ครุํ กโรนฺติ มาเนนฺติ ปูเชนฺติ. ตํ กิสฺส เหตุ, เต หิ ตสฺส อายสฺมโต ปาปกา อกุสลา อิจฺฉาวจรา ปหีนา ทิสฺสนฺติ เจว สูยนฺติ จ. เสยฺยถาปิ อาวุโส กํสปาติ อาภตา อาปณา วา กมฺมารกุลา วา ปริสุทฺธา ปริโยทาตา, ตเมนํ สามิกา สาลีนํ โอทนํ วิจิตกาฬกํ\csromansharedfootnote{5084761a7cd8752c}{วิจินิตกาฬกํ (ก)} อเนกสูปํ อเนกพฺยญฺชนํ รจยิตฺวา อญฺญิสฺสา กํสปาติยา ปฏิกุชฺชิตฺวา อนฺตราปณํ ปฏิปชฺเชยฺยุํ, ตเมนํ ชโน ทิสฺวา เอวํ วเทยฺย “อมฺโภ กิเมวิทํ หรียติ ชญฺญชญฺญํ วิยา”ติ. ตเมนํ อุฏฺฐหิตฺวา อปาปุริตฺวา โอโลเกยฺย, ตสฺส สหทสฺสเนน มนาปตา จ สณฺฐเหยฺย, อปฺปาฏิกุลฺยตา จ สณฺฐเหยฺย, อเชคุจฺฉตา จ สณฺฐเหยฺย, สุหิตานํปิ โภตฺตุกมฺยตา อสฺส, ปเคว ชิฆจฺฉิตานํ. เอวเมว โข อาวุโส ยสฺส กสฺสจิ ภิกฺขุโน อิเม ปาปกา อกุสลา อิจฺฉาวจรา ปหีนา ทิสฺสนฺติ เจว สูยนฺติ จ, กิญฺจาปิ โส โหติ คามนฺตวิหารี เนมนฺตนิโก คหปติจีวร\-ธโร, อถ โข นํ สพฺรหฺมจารี สกฺกโรนฺติ ครุํ กโรนฺติ มาเนนฺติ ปูเชนฺติ. ตํ กิสฺส เหตุ, เต หิ ตสฺส อายสฺมโต ปาปกา อกุสลา อิจฺฉาวจรา ปหีนา ทิสฺสนฺติ เจว สูยนฺติ จาติ.}

\proseitem{63}{เอวํ วุตฺเต อายสฺมา มหาโมคฺคลฺลาโน อายสฺมนฺตํ สาริปุตฺตํ เอตทโวจ “อุปมา มํ อาวุโส สาริปุตฺต ปฏิภาตี”ติ. ปฏิภาตุ ตํ อาวุโส โมคฺคลฺลานาติ. เอกมิทาหํ อาวุโส สมยํ ราชคเห วิหรามิ คิริพฺพเช. อถ ขฺวาหํ อาวุโส ปุพฺพณฺหสมยํ นิวาเสตฺวา ปตฺต\-จีวรมา\-ทาย ราชคหํ ปิณฺฑาย ปาวิสิํ. เตน โข ปน สมเยน สมีติ ยานการปุตฺโต รถสฺส เนมิํ ตจฺฉติ. ตเมนํ ปณฺฑุปุตฺโต อาชีวโก ปุราณ\-ยานการ\-ปุตฺโต ปจฺจุปฏฺฐิโต โหติ. อถ โข อาวุโส ปณฺฑุปุตฺตสฺส อาชีวกสฺส ปุราณ\-ยานการ\-ปุตฺตสฺส เอวํ เจตโส ปริวิตกฺโก อุทปาทิ “อโห วตายํ สมีติ ยานการปุตฺโต อิมิสฺสา เนมิยา อิมญฺจ วงฺกํ อิมญฺจ ชิมฺหํ อิมญฺจ โทสํ ตจฺเฉยฺย, เอวายํ เนมิ อปคตวงฺกา อปคตชิมฺหา อปคตโทสา สุทฺธา อสฺส\csromansharedfootnote{66ac351fc7bfde71}{สุทฺธาสฺส (สี, อิ), สุทฺธา (ก)} สาเร ปติฏฺฐิตา”ติ. ยถา ยถา โข อาวุโส ปณฺฑุปุตฺตสฺส อาชีวกสฺส ปุราณ\-ยานการ\-ปุตฺตสฺส เจตโส ปริวิตกฺโก โหติ, ตถา \csromanfolio{38}ตถา สมีติ ยานการปุตฺโต ตสฺสา เนมิยา ตญฺจ วงฺกํ ตญฺจ ชิมฺหํ ตญฺจ โทสํ ตจฺฉติ. อถ โข อาวุโส ปณฺฑุปุตฺโต อาชีวโก ปุราณ\-ยานการ\-ปุตฺโต อตฺตมโน อตฺตมนวาจํ นิจฺฉาเรสิ “หทยา หทยํ มญฺเญ อญฺญาย ตจฺฉตี”ติ.}

\prose{เอวเมว โข อาวุโส เย เต ปุคฺคลา อสฺสทฺธา ชีวิกตฺถา น สทฺธา อคารสฺมา อนคาริยํ ปพฺพชิตา สฐา มายาวิโน เกตพิโน\csromansharedfootnote{84962f052b6df1f3}{เกฏุภิโน (พหูสุ)} อุทฺธตา อุนฺนฬา จปลา มุขรา วิกิณฺณวาจา อินฺทฺริเยสุ อคุตฺตทฺวารา โภชเน อมตฺตญฺญุโน ชาคริยํ อนนุยุตฺตา สามญฺเญ อนเปกฺขวนฺโต สิกฺขาย น ติพฺพคารวา พาหุลิกา สาถลิกา โอกฺกมเน ปุพฺพงฺคมา ปวิเวเก นิกฺขิตฺตธุรา กุสีตา หีนวีริยา มุฏฺฐสฺสตี อสมฺปชานา อสมาหิตา วิพฺภนฺตจิตฺตา ทุปฺปญฺญา เอฬมูคา, เตสํ อายสฺมา สาริปุตฺโต อิมินา ธมฺม\-ปริยาเยน หทยา หทยํ มญฺเญ อญฺญาย ตจฺฉติ.}

\prose{เย ปน เต กุลปุตฺตา สทฺธา อคารสฺมา อนคาริยํ ปพฺพชิตา อสฐา อมายาวิโน อเกตพิโน อนุทฺธตา อนุนฺนฬา อจปลา อมุขรา อวิกิณฺณวาจา อินฺทฺริเยสุ คุตฺตทฺวารา โภชเน มตฺตญฺญุโน ชาคริยํ อนุยุตฺตา สามญฺเญ อเปกฺขวนฺโต สิกฺขาย ติพฺพคารวา น พาหุลิกา น สาถลิกา โอกฺกมเน นิกฺขิตฺตธุรา ปวิเวเก ปุพฺพงฺคมา อารทฺธวีริยา ปหิตตฺตา อุปฏฺฐิตสฺสตี สมฺปชานา สมาหิตา เอกคฺคจิตฺตา ปญฺญวนฺโต อเนฬมูคา, เต อายสฺมโต สาริปุตฺตสฺส อิมํ ธมฺม\-ปริยายํ สุตฺวา ปิวนฺติ มญฺเญ ฆสนฺติ มญฺเญ วจสา เจว มนสา จ “สาธุ วต โภ สพฺรหฺมจารี อกุสลา วุฏฺฐาเปตฺวา กุสเล ปติฏฺฐาเปตี”ติ. เสยฺยถาปิ อาวุโส อิตฺถี วา ปุริโส วา ทหโร ยุวา มณฺฑนก\-ชาติโก สีสํนฺหาโต อุปฺปลมาลํ วา วสฺสิกมาลํ วา อติมุตฺตก\-มาลํ\csromansharedfootnote{7155413adc77967e}{อธิมุตฺตกมาลํ (สฺยา)} วา ลภิตฺวา อุโภหิ หตฺเถหิ ปฏิคฺคเหตฺวา อุตฺตมงฺเค สิรสฺมิํ ปติฏฺฐเปยฺย. เอวเมว โข อาวุโส เย เต กุลปุตฺตา สทฺธา อคารสฺมา อนคาริยํ ปพฺพชิตา อสฐา อมายาวิโน อเกตพิโน อนุทฺธตา อนุนฺนฬา อจปลา อมุขรา อวิกิณฺณวาจา อินฺทฺริเยสุ คุตฺตทฺวารา โภชเน มตฺตญฺญุโน ชาคริยํ อนุยุตฺตา สามญฺเญ อเปกฺขวนฺโต สิกฺขาย ติพฺพคารวา น พาหุลิกา น สาถลิกา โอกฺกมเน นิกฺขิตฺตธุรา ปวิเวเก ปุพฺพงฺคมา อารทฺธวีริยา ปหิตตฺตา}

\vspace{\tipitakaparskip}
\csromancenter{อากงฺเขยฺย\-สุตฺต อุปฏฺฐิตสฺสตี สมฺปชานา สมาหิตา เอกคฺคจิตฺตา ปญฺญวนฺโต อเนฬมูคา, เต อายสฺมโต สาริปุตฺตสฺส อิมํ ธมฺม\-ปริยายํ สุตฺวา ปิวนฺติ มญฺเญ ฆสนฺติ มญฺเญ วจสา เจว มนสา จ “สาธุ วต โภ สพฺรหฺมจารี อกุสลา วุฏฺฐาเปตฺวา กุสเล ปติฏฺฐาเปตี”ติ. อิติห เต อุโภ มหานาคา อญฺญมญฺญสฺส สุภาสิตํ สมนุ\-โมทิํสูติ.}
\nitthitamruled{อนงฺคณสุตฺตํ นิฏฺฐิตํ ปญฺจมํ.}
\csromanmatikaanchor{mtk.14}
\csromantocmarkref{subsection}{6. อากงฺเขยฺยสุตฺต}{mtk.14}
\bookmark[dest={mtk.14},level=2]{6. อากงฺเขยฺยสุตฺต}
\csromanheader{2}{6. อากงฺเขยฺย\-สุตฺต}
\proseitem{64}{\csromanfolio{39}เอวํ เม สุตํ\csromandash{}เอกํ สมยํ ภควา สาวตฺถิยํ วิหรติ เชตวเน อนาถ\-ปิณฺฑิกสฺส อาราเม. ตตฺร โข ภควา ภิกฺขู อามนฺเตสิ “ภิกฺขโว”ติ. “ภทนฺเต”ติ เต ภิกฺขู ภควโต ปจฺจสฺโสสุํ. ภควา เอตทโวจ\csromandash{}}

\prose{สมฺปนฺนสีลา ภิกฺขเว วิหรถ สมฺปนฺน\-ปาติโมกฺขา, ปาติโมกฺข\-สํวร\-สํวุตา วิหรถ อาจารโคจร\-สมฺปนฺนา, อณุมตฺเตสุ วชฺเชสุ ภยทสฺสาวิโน สมาทาย สิกฺขถ สิกฺขาปเทสุ.}

\proseitem{65}{อากงฺเขยฺย เจ ภิกฺขเว ภิกฺขุ “สพฺรหฺมจารีนํ ปิโย จ อสฺสํ มนาโป จ ครุ จ ภาวนีโย จา”ติ\csromansharedfootnote{8ec999580ef35f2e}{มนาโป ครุภาวนิโย จาติ (สี)}, สีเลเสฺววสฺส ปริปูรการี, อชฺฌตฺตํ เจโตสมถม\-นุยุตฺโต อนิรากต\-ชฺฌาโน วิปสฺสนาย สมนฺนาคโต พฺรูเหตา สุญฺญาคารานํ. (1)}

\prose{อากงฺเขยฺย เจ ภิกฺขเว ภิกฺขุ “ลาภี อสฺสํ จีวร\-ปิณฺฑปาต\-เสนาสน\-คิลานปฺปจฺจย\-เภสชฺช\-ปริกฺขารานนฺ”ติ, สีเลเสฺววสฺส ปริปูรการี, อชฺฌตฺตํ เจโตสมถมนุตฺโต อนิรากต\-ชฺฌาโน วิปสฺสนาย สมนฺนาคโต พฺรูเหตา สุญฺญาคารานํ. (2)}

\prose{อากงฺเขยฺย เจ ภิกฺขเว ภิกฺขุ “เยสาหํ จีวร\-ปิณฺฑปาต\-เสนาสน\-คิลานปฺปจฺจย\-เภสชฺช\-ปริกฺขารํ ปริภุญฺชามิ, เตสํ เต การา มหปฺผลา \csromanfolio{40}อสฺสุ มหานิสํสา”ติ, สีเลเสฺววสฺส ปริปูรการี, อชฺฌตฺตํ เจโตสมถม\-นุยุตฺโต อนิรากต\-ชฺฌาโน วิปสฺสนาย สมนฺนาคโต พฺรูเหตา สุญฺญาคารานํ. (3)}

\prose{อากงฺเขยฺย เจ ภิกฺขเว ภิกฺขุ “เย มํ\csromansharedfootnote{7b0abc85d1956527}{เย เม (สี, สฺยา)} ญาตี สาโลหิตา เปตา กาลงฺกตา\csromansharedfootnote{b0d0d79452bf051e}{กาลกตา (สี, สฺยา, อิ)} ปสนฺนจิตฺตา อนุสฺสรนฺติ, เตสํ ตํ มหปฺผลํ อสฺส มหานิสํสนฺ”ติ, สีเลเสฺววสฺส ปริปูรการี, อชฺฌตฺตํ เจโตสมถม\-นุยุตฺโต อนิรากต\-ชฺฌาโน วิปสฺสนาย สมนฺนาคโต พฺรูเหตา สุญฺญาคารานํ. (4)}

\proseitem{66}{อากงฺเขยฺย เจ ภิกฺขเว ภิกฺขุ “อรติรติสโห อสฺสํ, น จ มํ อรติ สเหยฺย, อุปฺปนฺนํ อรติํ อภิภุยฺย อภิภุยฺย วิหเรยฺยนฺ”ติ, สีเลเสฺววสฺส ปริปูรการี ฯเปฯ พฺรูเหตา สุญฺญาคารานํ. (5)}

\prose{อากงฺเขยฺย เจ ภิกฺขเว ภิกฺขุ “ภยเภรว\-สโห อสฺสํ, น จ มํ ภยเภรวํ สเหยฺย, อุปฺปนฺนํ ภยเภรวํ อภิภุยฺย อภิภุยฺย วิหเรยฺยนฺ”ติ, สีเลเสฺววสฺส ปริปูรการี ฯเปฯ พฺรูเหตา สุญฺญคารานํ. (6)}

\prose{อากงฺเขยฺย เจ ภิกฺขเว ภิกฺขุ “จตุนฺนํ ฌานานํ อาภิเจตสิกานํ ทิฏฺฐ\-ธมฺม\-สุข\-วิหารานํ นิกามลาภี อสฺสํ อกิจฺฉลาภี อกสิรลาภี”ติ, สีเลเสฺววสฺส ปริปูรการี ฯเปฯ พฺรูเหตา สุญฺญาคารานํ. (7)}

\prose{อากงฺเขยฺย เจ ภิกฺขเว ภิกฺขุ “เย เต สนฺตา วิโมกฺขา อติกฺกมฺม รูเป อารุปฺปา, เต กาเยน ผุสิตฺวา วิหเรยฺยนฺ”ติ, สีเลเสฺววสฺส ปริปูรการี ฯเปฯ พฺรูเหตา สุญฺญาคารานํ. (8)}

\proseitem{67}{อากงฺเขยฺย เจ ภิกฺขเว ภิกฺขุ “ติณฺณํ สํโยชนานํ ปริกฺขยา โสตาปนฺโน อสฺสํ อวินิปาต\-ธมฺโม นิยโต สมฺโพธิ\-ปรายโณ”ติ. สีเลเสฺววสฺส ปริปูรการี ฯเปฯ พฺรูเหตา สุญฺญาคารานํ. (9)}

\prose{อากงฺเขยฺย เจ ภิกฺขเว ภิกฺขุ “ติณฺณํ สํโยชนานํ ปริกฺขยา ราค\-โทส\-โมหานํ ตนุตฺตา สกทาคามี อสฺสํ สกิเทว อิมํ โลกํ อาคนฺตฺวา ทุกฺขสฺสนฺตํ กเรยฺยนฺ”ติ, สีเลเสฺววสฺส ปริปูรการี ฯเปฯ พฺรูเหตา สุญฺญาคารานํ. (10)}

\csromanheader{2}{6. อากงฺเขยฺย\-สุตฺต}
\prose{\csromanfolio{41}อากงฺเขยฺย เจ ภิกฺขเว ภิกฺขุ “ปญฺจนฺนํ โอรมฺภาคิยานํ สํโยชนานํ ปริกฺขยา โอปปาติโก อสฺสํ ตตฺถ ปรินิพฺพายี อนาวตฺติธมฺโม ตสฺมา โลกา”ติ, สีเลเสฺววสฺส ปริปูรการี ฯเปฯ พฺรูเหตา สุญฺญาคารานํ. (11)}

\proseitem{68}{อากงฺเขยฺย เจ ภิกฺขเว ภิกฺขุ “อเนกวิหิตํ อิทฺธิวิธํ ปจฺจนุภเวยฺยํ, เอโกปิ หุตฺวา พหุธา อสฺสํ, พหุธาปิ หุตฺวา เอโก อสฺสํ, อาวิภาวํ ติโรภาวํ ติโรกุฏฺฏํ ติโรปาการํ ติโรปพฺพตํ อสชฺชมาโน คจฺเฉยฺยํ เสยฺยถาปิ อากาเส, ปถวิยาปิ อุมฺมุชฺช\-นิมุชฺชํ กเรยฺยํ เสยฺยถาปิ อุทเก, อุทเกปิ อภิชฺชมาเน คจฺเฉยฺยํ เสยฺยถาปิ ปถวิยํ, อากาเสปิ ปลฺลงฺเกน กเมยฺยํ เสยฺยถาปิ ปกฺขี สกุโณ, อิเมปิ จนฺทิมสูริเย เอวํมหิทฺธิเก เอวํ\-มหา\-นุภาเว ปาณินา ปรามเสยฺยํ ปริมชฺเชยฺยํ, ยาว พฺรหฺมโลกาปิ กาเยน วสํ วตฺเตยฺยนฺ”ติ, สีเลเสฺววสฺส ปริปูรการี ฯเปฯ พฺรูเหตา สุญฺญาคารานํ. (12)}

\prose{อากงฺเขยฺย เจ ภิกฺขเว ภิกฺขุ “ทิพฺพาย โสตธาตุยา วิสุทฺธาย อติกฺกนฺตมานุสิกาย อุโภ สทฺเท สุเณยฺยํ ทิพฺเพ จ มานุเส จ เย ทูเร สนฺติเก จา”ติ, สีเลเสฺววสฺส ปริปูรการี ฯเปฯ พฺรูเหตา สุญฺญาคารานํ. (13)}

\prose{อากงฺเขยฺย เจ ภิกฺขเว ภิกฺขุ “ปรสตฺตานํ ปรปุคฺคลานํ เจตสา เจโต ปริจฺจ ปชาเนยฺยํ, สราคํ วา จิตฺตํ สราคํ จิตฺตนฺติ ปชาเนยฺยํ, วีตราคํ วา จิตฺตํ วีตราคํ จิตฺตนฺติ ปชาเนยฺยํ, สโทสํ วา จิตฺตํ สโทสํ จิตฺตนฺติ ปชาเนยฺยํ, วีตโทสํ วา จิตฺตํ วีตโทสํ จิตฺตนฺติ ปชาเนยฺยํ, สโมหํ วา จิตฺตํ สโมหํ จิตฺตนฺติ ปชาเนยฺยํ, วีตโมหํ วา จิตฺตํ วีตโมหํ จิตฺตนฺติ ปชาเนยฺยํ, สํขิตฺตํ วา จิตฺตํ สํขิตฺตํ จิตฺตนฺติ ปชาเนยฺยํ, วิกฺขิตฺตํ วา จิตฺตํ วิกฺขิตฺตํ จิตฺตนฺติ ปชาเนยฺยํ, มหคฺคตํ วา จิตฺตํ มหคฺคตํ จิตฺตนฺติ ปชาเนยฺยํ, อมหคฺคตํ วา จิตฺตํ อมหคฺคตํ จิตฺตนฺติ ปชาเนยฺยํ, สอุตฺตรํ วา จิตฺตํ สอุตฺตรํ จิตฺตนฺติ ปชาเนยฺยํ, อนุตฺตรํ วา จิตฺตํ อนุตฺตรํ จิตฺตนฺติ ปชาเนยฺยํ, สมาหิตํ วา จิตฺตํ สมาหิตํ จิตฺตนฺติ ปชาเนยฺยํ, อสมาหิตํ วา จิตฺตํ อสมาหิตํ จิตฺตนฺติ ปชาเนยฺยํ, วิมุตฺตํ วา จิตฺตํ วิมุตฺตํ จิตฺตนฺติ ปชาเนยฺยํ, อวิมุตฺตํ วา จิตฺตํ อวิมุตฺตํ จิตฺตนฺติ ปชาเนยฺยนฺ”ติ, สีเลเสฺววสฺส ปริปูรการี ฯเปฯ พฺรูเหตา สุญฺญาคารานํ. (14)}

\prose{\csromanfolio{42}อากงฺเขยฺย เจ ภิกฺขเว ภิกฺขุ “อเนกวิหิตํ ปุพฺเพนิวาสํ อนุสฺสเรยฺยํ. เสยฺยถิทํ, เอกมฺปิ ชาติํ เทฺวปิ ชาติโย ติสฺโสปิ ชาติโย จตสฺโสปิ ชาติโย ปญฺจปิ ชาติโย ทสปิ ชาติโย วีสมฺปิ ชาติโย ติํสมฺปิ ชาติโย จตฺตาลีสมฺปิ ชาติโย ปญฺญาสมฺปิ ชาติโย ชาติสตมฺปิ ชาติสหสฺสมฺปิ ชาติ\-สต\-สหสฺสมฺปิ อเนเกปิ สํวฏฺฏกปฺเป อเนเกปิ วิวฏฺฏกปฺเป อเนเกปิ สํวฏฺฏ\-วิวฏฺฏ\-กปฺเป ‘อมุตฺราสิํ เอวํนาโม เอวํโคตฺโต เอวํวณฺโณ เอวมาหาโร เอวํ\-สุขทุกฺข\-ปฺปฏิสํเวที เอวมา\-ยุปริยนฺโต. โส ตโต จุโต อมุตฺร อุทปาทิํ, ตตฺราปาสิํ เอวํนาโม เอวํโคตฺโต เอวํวณฺโณ เอวมาหาโร เอวํ\-สุขทุกฺข\-ปฺปฏิสํเวที เอวมา\-ยุปริยนฺโต, โส ตโต จุโต อิธูปปนฺโน’ติ, อิติ สาการํ สอุทฺเทสํ อเนกวิหิตํ ปุพฺเพนิวาสํ อนุสฺสเรยฺยนฺ”ติ, สีเลเสฺววสฺส ปริปูรการี ฯเปฯ พฺรูเหตา สุญฺญาคารานํ. (15)}

\prose{อากงฺเขยฺย เจ ภิกฺขเว ภิกฺขุ “ทิพฺเพน จกฺขุนา วิสุทฺเธน อติกฺกนฺต\-มานุสเกน สตฺเต ปสฺเสยฺยํ จวมาเน อุปปชฺชมาเน หีเน ปณีเต สุวณฺเณ ทุพฺพณฺเณ สุคเต ทุคฺคเต, ยถากมฺมูปเค สตฺเต ปชาเนยฺยํ, ‘อิเม วต โภนฺโต สตฺตา กาย\-ทุจฺจริเตน สมนฺนาคตา วจีทุจฺจริเตน สมนฺนาคตา มโน\-ทุจฺจริเตน สมนฺนาคตา อริยานํ อุปวาทกา มิจฺฉาทิฏฺฐิกา มิจฺฉาทิฏฺฐิ\-กมฺม\-สมาทานา, เต กายสฺส เภทา ปรํ มรณา อปายํ ทุคฺคติํ วินิปาตํ นิรยํ อุปปนฺนา. อิเม วา ปน โภนฺโต สตฺตา กายสุจริเตน สมนฺนาคตา วจีสุจริเตน สมนฺนาคตา มโนสุจริเตน สมนฺนาคตา อริยานํ อนุปวาทกา สมฺมาทิฏฺฐิกา สมฺมา\-ทิฏฺฐิ\-กมฺม\-สมาทานา, เต กายสฺส เภทา ปรํ มรณา สุคติํ สคฺคํ โลกํ อุปปนฺนา’ติ, อิติ ทิพฺเพน จกฺขุนา วิสุทฺเธน อติกฺกนฺต\-มานุสเกน สตฺเต ปสฺเสยฺยํ จวมาเน อุปปชฺชมาเน หีเน ปณีเต สุวณฺเณ ทุพฺพณฺเณ สุคเต ทุคฺคเต, ยถากมฺมูปเค สตฺเต ปชาเนยฺยนฺ”ติ, สีเลเสฺววสฺส ปริปูรการี, อชฺฌตฺตํ เจโตสมถม\-นุยุตฺโต อนิรากต\-ชฺฌาโน วิปสฺสนาย สมนฺนาคโต พฺรูเหตา สุญฺญาคารานํ. (16)}

\csromanmatikaanchor{mtk.15}
\csromantocmarkref{subsection}{7. วตฺถสุตฺต}{mtk.15}
\bookmark[dest={mtk.15},level=2]{7. วตฺถสุตฺต}
\proseitem{69}{อากงฺเขยฺย เจ ภิกฺขเว ภิกฺขุ “อาสวานํ ขยา อนาสวํ เจโตวิมุตฺติํ ปญฺญาวิมุตฺติํ ทิฏฺเฐวธมฺเม สยํ อภิญฺญา สจฺฉิกตฺวา \csromanfolio{43}อุปสมฺปชฺช วิหเรยฺยนฺ”ติ, สีเลเสฺววสฺส ปริปูรการี อชฺฌตฺตํ เจโตสมถม\-นุยุตฺโต อนิรากต\-ชฺฌาโน วิปสฺสนาย สมนฺนาคโต พฺรูเหตา สุญฺญาคารานํ. (17)}

\prosewithcloserruled{“สมฺปนฺนสีลา ภิกฺขเว วิหรถ สมฺปนฺน\-ปาติโมกฺขา, ปาติโมกฺข\-สํวร\-สํวุตา วิหรถ อาจารโคจร\-สมฺปนฺนา, อณุมตฺเตสุ วชฺเชสุ ภยทสฺสาวิโน สมาทาย สิกฺขถ สิกฺขาปเทสู”ติ, อิติ ยํ ตํ วุตฺตํ, อิทเมตํ ปฏิจฺจ วุตฺตนฺติ. อิทมโวจ ภควา. อตฺตมนา เต ภิกฺขู ภควโต ภาสิตํ อภินนฺทุนฺติ.}{อากงฺเขยฺย\-สุตฺตํ นิฏฺฐิตํ ฉฏฺฐํ.}
\proseitem{70}{เอวํ เม สุตํ\csromandash{}เอกํ สมยํ ภควา สาวตฺถิยํ วิหรติ เชตวเน อนาถ\-ปิณฺฑิกสฺส อาราเม. ตตฺร โข ภควา ภิกฺขู อามนฺเตสิ “ภิกฺขโว”ติ. “ภทนฺเต”ติ เต ภิกฺขู ภควโต ปจฺจสฺโสสุํ. ภควา เอตทโวจ\csromandash{}}

\prose{เสยฺยถาปิ ภิกฺขเว วตฺถํ สํกิลิฏฺฐํ มลคฺคหิตํ, ตเมนํ รชโก ยสฺมิํ ยสฺมิํ รงฺคชาเต อุปสํหเรยฺย, ยทิ นีลกาย ยทิ ปีตกาย ยทิ โลหิตกาย ยทิ มญฺชิฏฺฐกาย\csromansharedfootnote{69eebbfbf3531880}{มญฺเชฏฺฐกาย (สี, อิ), มญฺเชฏฺฐิกาย (สฺยา)}, ทุรตฺตวณฺณเม\-วสฺส, อปริสุทฺธวณฺณเม\-วสฺส. ตํ กิสฺส เหตุ, อปริสุทฺธตฺตา ภิกฺขเว วตฺถสฺส. เอวเมว โข ภิกฺขเว จิตฺเต สํกิลิฏฺเฐ ทุคฺคติ ปาฏิกงฺขา. เสยฺยถาปิ ภิกฺขเว วตฺถํ ปริสุทฺธํ ปริโยทาตํ, ตเมนํ รชโก ยสฺมิํ ยสฺมิํ รงฺคชาเต อุปสํหเรยฺย ยทิ นีลกาย ยทิ ปีตกาย ยทิ โลหิตกาย ยทิ มญฺชิฏฺฐกาย, สุรตฺตวณฺณเม\-วสฺส, ปริสุทฺธวณฺณเม\-วสฺส. ตํ กิสฺส เหตุ, ปริสุทฺธตฺตา ภิกฺขเว วตฺถสฺส. เอวเมว โข ภิกฺขเว จิตฺเต อสํกิลิฏฺเฐ สุคติ ปาฏิกงฺขา.}

\proseitem{71}{กตเม จ ภิกฺขเว จิตฺตสฺส อุปกฺกิเลสา. อภิชฺฌา\-วิสมโลโภ จิตฺตสฺส อุปกฺกิเลโส, พฺยาปาโท จิตฺตสฺส อุปกฺกิเลโส, โกโธ จิตฺตสฺส อุปกฺกิเลโส, อุปนาโห จิตฺตสฺส อุปกฺกิเลโส, \csromanfolio{44}มกฺโข จิตฺตสฺส อุปกฺกิเลโส, ปฬาโส จิตฺตสฺส อุปกฺกิเลโส, อิสฺสา จิตฺตสฺส อุปกฺกิเลโส, มจฺฉริยํ จิตฺตสฺส อุปกฺกิเลโส, มายา จิตฺตสฺส อุปกฺกิเลโส, สาเฐยฺยํ จิตฺตสฺส อุปกฺกิเลโส, ถมฺโภ จิตฺตสฺส อุปกฺกิเลโส, สารมฺโภ จิตฺตสฺส อุปกฺกิเลโส, มาโน จิตฺตสฺส อุปกฺกิเลโส, อติมาโน จิตฺตสฺส อุปกฺกิเลโส, มโท จิตฺตสฺส อุปกฺกิเลโส, ปมาโท จิตฺตสฺส อุปกฺกิเลโส.}

\proseitem{72}{ส โข โส ภิกฺขเว ภิกฺขุ อภิชฺฌา\-วิสมโลโภ จิตฺตสฺส อุปกฺกิเลโสติ อิติ วิทิตฺวา อภิชฺฌา\-วิสมโลภํ จิตฺตสฺส อุปกฺกิเลสํ ปชหติ. พฺยาปาโท จิตฺตสฺส อุปกฺกิเลโสติ อิติ วิทิตฺวา พฺยาปาทํ จิตฺตสฺส อุปกฺกิเลสํ ปชหติ. โกโธ จิตฺตสฺส อุปกฺกิเลโสติ อิติ วิทิตฺวา โกธํ จิตฺตสฺส อุปกฺกิเลสํ ปชหติ. อุปนาโห จิตฺตสฺส อุปกฺกิเลโสติ อิติ วิทิตฺวา อุปนาหํ จิตฺตสฺส อุปกฺกิเลสํ ปชหติ. มกฺโข จิตฺตสฺส อุปกฺกิเลโสติ อิติ วิทิตฺวา มกฺขํ จิตฺตสฺส อุปกฺกิเลสํ ปชหติ. ปฬาโส จิตฺตสฺส อุปกฺกิเลโสติ อิติ วิทิตฺวา ปฬาสํ จิตฺตสฺส อุปกฺกิเลสํ ปชหติ. อิสฺสา จิตฺตสฺส อุปกฺกิเลโสติ อิติ วิทิตฺวา อิสฺสํ จิตฺตสฺส อุปกฺกิเลสํ ปชหติ. มจฺฉริยํ จิตฺตสฺส อุปกฺกิเลโสติ อิติ วิทิตฺวา มจฺฉริยํ จิตฺตสฺส อุปกฺกิเลสํ ปชหติ. มายา จิตฺตสฺส อุปกฺกิเลโสติ อิติ วิทิตฺวา มายํ จิตฺตสฺส อุปกฺกิเลสํ ปชหติ. สาเฐยฺยํ จิตฺตสฺส อุปกฺกิเลโสติ อิติ วิทิตฺวา สาเฐยฺยํ จิตฺตสฺส อุปกฺกิเลสํ ปชหติ. ถมฺโภ จิตฺตสฺส อุปกฺกิเลโสติ อิติ วิทิตฺวา ถมฺภํ จิตฺตสฺส อุปกฺกิเลสํ ปชหติ. สารมฺโภ จิตฺตสฺส อุปกฺกิเลโสติ อิติ วิทิตฺวา สารมฺภํ จิตฺตสฺส อุปกฺกิเลสํ ปชหติ. มาโน จิตฺตสฺส อุปกฺกิเลโสติ อิติ วิทิตฺวา มานํ จิตฺตสฺส อุปกฺกิเลสํ ปชหติ. อติมาโน จิตฺตสฺส อุปกฺกิเลโสติ อิติ วิทิตฺวา อติมานํ จิตฺตสฺส อุปกฺกิเลสํ ปชหติ. มโท จิตฺตสฺส อุปกฺกิเลโสติ อิติ วิทิตฺวา มทํ จิตฺตสฺส อุปกฺกิเลสํ ปชหติ. ปมาโท จิตฺตสฺส อุปกฺกิเลโสติ อิติ วิทิตฺวา ปมาทํ จิตฺตสฺส อุปกฺกิเลสํ ปชหติ.}

\proseitem{73}{ยโต โข\csromansharedfootnote{f12fe225acbb9bc7}{ยโต จ โข (สี, สฺยา)} ภิกฺขเว ภิกฺขุโน อภิชฺฌา\-วิสมโลโภ จิตฺตสฺส อุปกฺกิเลโสติ อิติ วิทิตฺวา อภิชฺฌา\-วิสมโลโภ จิตฺตสฺส \csromanfolio{45}อุปกฺกิเลโส ปหีโน โหติ. พฺยาปาโท จิตฺตสฺส อุปกฺกิเลโสติ อิติ วิทิตฺวา พฺยาปาโท จิตฺตสฺส อุปกฺกิเลโส ปหีโน โหติ. โกโธ จิตฺตสฺส อุปกฺกิเลโสติ อิติ วิทิตฺวา โกโธ จิตฺตสฺส อุปกฺกิเลโส ปหีโน โหติ. อุปนาโห จิตฺตสฺส อุปกฺกิเลโสติ อิติ วิทิตฺวา อุปนาโห จิตฺตสฺส อุปกฺกิเลโส ปหีโน โหติ. มกฺโข จิตฺตสฺส อุปกฺกิเลโสติ อิติ วิทิตฺวา มกฺโข จิตฺตสฺส อุปกฺกิเลโส ปหีโน โหติ. ปฬาโส จิตฺตสฺส อุปกฺกิเลโสติ อิติ วิทิตฺวา ปฬาโส จิตฺตสฺส อุปกฺกิเลโส ปหีโน โหติ. อิสฺสา จิตฺตสฺส อุปกฺกิเลโสติ อิติ วิทิตฺวา อิสฺสา จิตฺตสฺส อุปกฺกิเลโส ปหีโน โหติ. มจฺฉริยํ จิตฺตสฺส อุปกฺกิเลโสติ อิติ วิทิตฺวา มจฺฉริยํ จิตฺตสฺส อุปกฺกิเลโส ปหีโน โหติ. มายา จิตฺตสฺส อุปกฺกิเลโสติ อิติ วิทิตฺวา มายา จิตฺตสฺส อุปกฺกิเลโส ปหีโน โหติ. สาเฐยฺยํ จิตฺตสฺส อุปกฺกิเลโสติ อิติ วิทิตฺวา สาเฐยฺยํ จิตฺตสฺส อุปกฺกิเลโส ปหีโน โหติ. ถมฺโภ จิตฺตสฺส อุปกฺกิเลโสติ อิติ วิทิตฺวา ถมฺโภ จิตฺตสฺส อุปกฺกิเลโส ปหีโน โหติ. สารมฺโภ จิตฺตสฺส อุปกฺกิเลโสติ อิติ วิทิตฺวา สารมฺโภ จิตฺตสฺส อุปกฺกิเลโส ปหีโน โหติ. มาโน จิตฺตสฺส อุปกฺกิเลโสติ อิติ วิทิตฺวา มาโน จิตฺตสฺส อุปกฺกิเลโส ปหีโน โหติ. อติมาโน จิตฺตสฺส อุปกฺกิเลโสติ อิติ วิทิตฺวา อติมาโน จิตฺตสฺส อุปกฺกิเลโส ปหีโน โหติ. มโท จิตฺตสฺส อุปกฺกิเลโสติ อิติ วิทิตฺวา มโท จิตฺตสฺส อุปกฺกิเลโส ปหีโน โหติ. ปมาโท จิตฺตสฺส อุปกฺกิเลโสติ อิติ วิทิตฺวา ปมาโท จิตฺตสฺส อุปกฺกิเลโส ปหีโน โหติ.}

\proseitem{74}{โส พุทฺเธ อเวจฺจ\-ปฺปสาเทน สมนฺนาคโต โหติ “อิติปิ โส ภควา อรหํ สมฺมาสมฺพุทฺโธ วิชฺชา\-จรณ\-สมฺปนฺโน สุคโต โลกวิทู อนุตฺตโร ปุริส\-ทมฺม\-สารถิ สตฺถา เทวมนุสฺสานํ พุทฺโธ ภควา”ติ. ธมฺเม อเวจฺจ\-ปฺปสาเทน สมนฺนาคโต โหติ “สฺวากฺขาโต ภควตา ธมฺโม สนฺทิฏฺฐิโก อกาลิโก เอหิปสฺสิโก โอปเนยฺยิโก ปจฺจตฺตํ เวทิตพฺโพ วิญฺญูหี”ติ. สํเฆ อเวจฺจ\-ปฺปสาเทน สมนฺนาคโต โหติ “สุปฺปฏิปนฺโน ภควโต สาวกสํโฆ, อุชุปฺปฏิปนฺโน ภควโต สาวกสํโฆ, ญายปฺปฏิปนฺโน ภควโต สาวกสํโฆ, \csromanfolio{46}สามีจิ\-ปฺปฏิปนฺโน ภควโต สาวกสํโฆ, ยทิทํ จตฺตาริ ปุริสยุคานิ อฏฺฐ ปุริสปุคฺคลา, เอส ภควโต สาวกสํโฆ อาหุเนยฺโย ปาหุเนยฺโย ทกฺขิเณยฺโย อญฺชลิกรณีโย อนุตฺตรํ ปุญฺญกฺเขตฺตํ โลกสฺสา”ติ.}

\proseitem{75}{ยโถธิ\csromansharedfootnote{b19edca1f64677c0}{ยโตธิ (อฏฺฐกถายํ ปาฐนฺตรํ)} โข ปนสฺส จตฺตํ โหติ วนฺตํ มุตฺตํ ปหีนํ ปฏินิสฺสฏฺฐํ. โส “พุทฺเธ อเวจฺจ\-ปฺปสาเทน สมนฺนาคโตมฺหี”ติ ลภติ อตฺถเวทํ, ลภติ ธมฺมเวทํ, ลภติ ธมฺมู\-ปสํหิตํ ปาโมชฺชํ, ปมุทิตสฺส ปีติ ชายติ, ปีติมนสฺส กาโย ปสฺสมฺภติ, ปสฺสทฺธกาโย สุขํ เวเทติ, สุขิโน จิตฺตํ สมาธิยติ. “ธมฺเม ฯเปฯ สํเฆ อเวจฺจ\-ปฺปสาเทน สมนฺนาคโตมฺหี”ติ ลภติ อตฺถเวทํ, ลภติ ธมฺมเวทํ, ลภติ ธมฺมู\-ปสํหิตํ ปาโมชฺชํ, ปมุทิตสฺส ปีติ ชายติ, ปีติมนสฺส กาโย ปสฺสมฺภติ, ปสฺสทฺธกาโย สุขํ เวเทติ, สุขิโน จิตฺตํ สมาธิยติ. “ยโถธิ โข ปน เม จตฺตํ วนฺตํ มุตฺตํ ปหีนํ ปฏินิสฺสฏฺฐนฺ”ติ ลภติ อตฺถเวทํ, ลภติ ธมฺมเวทํ, ลภติ ธมฺมู\-ปสํหิตํ ปาโมชฺชํ, ปมุทิตสฺส ปีติ ชายติ, ปีติมนสฺส กาโย ปสฺสมฺภติ, ปสฺสทฺธกาโย สุขํ เวเทติ, สุขิโน จิตฺตํ สมาธิยติ.}

\proseitem{76}{ส โข โส ภิกฺขเว ภิกฺขุ เอวํสีโล เอวํธมฺโม เอวํปญฺโญ สาลีนํ เจปิ ปิณฺฑปาตํ ภุญฺชติ วิจิตกาฬกํ อเนกสูปํ อเนกพฺยญฺชนํ, เนวสฺส ตํ โหติ อนฺตรายาย. เสยฺยถาปิ ภิกฺขเว วตฺถํ สํกิลิฏฺฐํ มลคฺคหิตํ อจฺโฉทกํ อาคมฺม ปริสุทฺธํ โหติ ปริโยทาตํ. อุกฺกามุขํ วา ปนาคมฺม ชาตรูปํ ปริสุทฺธํ โหติ ปริโยทาตํ. เอวเมว โข ภิกฺขเว ภิกฺขุ เอวํสีโล เอวํธมฺโม เอวํปญฺโญ สาลีนํ เจปิ ปิณฺฑปาตํ ภุญฺชติ วิจิตกาฬกํ อเนกสูปํ อเนกพฺยญฺชนํ, เนวสฺส ตํ โหติ อนฺตรายาย.}

\proseitem{77}{โส เมตฺตา\-สหคเตน เจตสา เอกํ ทิสํ ผริตฺวา วิหรติ. ตถา ทุติยํ. ตถา ตติยํ. ตถา จตุตฺถํ\csromansharedfootnote{7357c33c7ca002fb}{จตุตฺถิํ (สี, อิ)}. อิติ อุทฺธมโธ ติริยํ สพฺพธิ สพฺพตฺตตาย สพฺพาวนฺตํ โลกํ เมตฺตา\-สหคเตน เจตสา วิปุเลน มหคฺคเตน อปฺปมาเณน อเวเรน อพฺยาปชฺเชน ผริตฺวา วิหรติ. กรุณา\-สหคเตน เจตสา ฯเปฯ. มุทิตา\-สหคเตน เจตสา ฯเปฯ. \csromanfolio{47}อุเปกฺขา\-สหคเตน เจตสา เอกํ ทิสํ ผริตฺวา วิหรติ. ตถา ทุติยํ. ตถา ตติยํ. ตถา จตุตฺถํ. อิติ อุทฺธมโธ ติริยํ สพฺพธิ สพฺพตฺตตาย สพฺพาวนฺตํ โลกํ อุเปกฺขา\-สหคเตน เจตสา วิปุเลน มหคฺคเตน อปฺปมาเณน อเวเรน อพฺยาปชฺเชน ผริตฺวา วิหรติ.}

\proseitem{78}{โส “อตฺถิ อิทํ, อตฺถิ หีนํ, อตฺถิ ปณีตํ, อตฺถิ อิมสฺส สญฺญาคตสฺส อุตฺตริ นิสฺสรณนฺ”ติ ปชานาติ. ตสฺส เอวํ ชานโต เอวํ ปสฺสโต กามาสวาปิ จิตฺตํ วิมุจฺจติ, ภวาสวาปิ จิตฺตํ วิมุจฺจติ, อวิชฺชาสวาปิ จิตฺตํ วิมุจฺจติ. วิมุตฺตสฺมิํ “วิมุตฺตมฺ”อิติ ญาณํ โหติ, “ขีณา ชาติ, วุสิตํ พฺรหฺมจริยํ, กตํ กรณียํ, นาปรํ อิตฺถตฺตายา”ติ ปชานาติ. อยํ วุจฺจติ ภิกฺขเว ภิกฺขุ สินาโต อนฺตเรน สินาเนนาติ.}

\proseitem{79}{เตน โข ปน สมเยน สุนฺทริก\-ภารทฺวาโช พฺราหฺมโณ ภควโต อวิทูเร นิสินฺโน โหติ. อถ โข สุนฺทริก\-ภารทฺวาโช พฺราหฺมโณ ภควนฺตํ เอตทโวจ “คจฺฉติ ปน ภวํ โคตโม พาหุกํ นทิํ สินายิตุนฺ”ติ. กิํ พฺราหฺมณ พาหุกาย นทิยา, กิํ พาหุกา นที กริสฺสตีติ. โลกฺขสมฺมตา\csromansharedfootnote{62d9ca526fa3097e}{โลขฺยสมฺมตา (สี), โมกฺขสมฺมตา (อิ)} หิ โภ โคตม พาหุกา นที พหุชนสฺส, ปุญฺญสมฺมตา หิ โภ โคตม พาหุกา นที พหุชนสฺส. พาหุกาย ปน นทิยา พหุชโน ปาปกมฺมํ กตํ ปวาเหตีติ. อถ โข ภควา สุนฺทริก\-ภารทฺวาชํ พฺราหฺมณํ คาถาหิ อชฺฌภาสิ.}

\csromanmatikaanchor{mtk.16}
\csromantocmarkref{subsection}{8. สลฺเลขสุตฺต}{mtk.16}
\bookmark[dest={mtk.16},level=2]{8. สลฺเลขสุตฺต}
\setlength{\csromangathapagewidth}{0pt}
\setlength{\csromangathawakwidth}{0pt}
\csromangathameasuregroup{พาหุกํ อธิกกฺกญฺจ, \\ สรสฺสติํ ปยาคญฺจ, \\ นิจฺจํปิ พาโล ปกฺขนฺโท\textsuperscript{1}, \\ กิํ สุนฺทริกา กริสฺสติ, \\ เวริํ กตกิพฺพิสํ นรํ, \\ สุทฺธสฺส เว สทา ผคฺคุ, \\ สุทฺธสฺส สุจิกมฺมสฺส, \\ อิเธว สินาหิ พฺราหฺมณ, \\ สเจ มุสา น ภณสิ, \\ สเจ อทินฺนํ นาทิยสิ, \\ กิํ กาหสิ คยํ คนฺตฺวา,}{\csromangathabat{พาหุกํ อธิกกฺกญฺจ,}{คยํ สุนฺทริกํ มปิ\textsuperscript{1}.} \\ \csromangathabat{สรสฺสติํ ปยาคญฺจ,}{อโถ พาหุมติํ นทิํ.} \\ \csromangathabat{นิจฺจํปิ พาโล ปกฺขนฺโท\textsuperscript{1},}{กณฺหกมฺโม น สุชฺฌติ.} \\ \csromangathabat{กิํ สุนฺทริกา กริสฺสติ,}{กิํ ปยาคา\textsuperscript{1} กิํ พาหุกา นที.} \\ \csromangathabat{เวริํ กตกิพฺพิสํ นรํ,}{น หิ นํ โสธเย ปาปกมฺมินํ.} \\ \csromangathabat{สุทฺธสฺส เว สทา ผคฺคุ,}{สุทฺธสฺสุโปสโถ สทา.} \\ \csromangathabat{สุทฺธสฺส สุจิกมฺมสฺส,}{สทา สมฺปชฺชเต วตํ.} \\ \csromangathabat{อิเธว สินาหิ พฺราหฺมณ,}{สพฺพภูเตสุ กโรหิ เขมตํ.} \\ \csromangathabat{สเจ มุสา น ภณสิ,}{สเจ ปาณํ น หิํสสิ.} \\ \csromangathabat{สเจ อทินฺนํ นาทิยสิ,}{สทฺทหาโน อมจฺฉรี.} \\ \csromangathabat{กิํ กาหสิ คยํ คนฺตฺวา,}{อุทปาโนปิ เต คยาติ.}}
\csromangathasetleft{พาหุกํ อธิกกฺกญฺจ, \\ สรสฺสติํ ปยาคญฺจ, \\ นิจฺจํปิ พาโล ปกฺขนฺโท\textsuperscript{1}, \\ กิํ สุนฺทริกา กริสฺสติ, \\ เวริํ กตกิพฺพิสํ นรํ, \\ สุทฺธสฺส เว สทา ผคฺคุ, \\ สุทฺธสฺส สุจิกมฺมสฺส, \\ อิเธว สินาหิ พฺราหฺมณ, \\ สเจ มุสา น ภณสิ, \\ สเจ อทินฺนํ นาทิยสิ, \\ กิํ กาหสิ คยํ คนฺตฺวา,}
\csromangathagroup{\csromangathastanza{\csromangathabat{พาหุกํ อธิกกฺกญฺจ,}{คยํ สุนฺทริกํ มปิ\csromansharedfootnote{4c2f9096671bd79c}{สุนฺทริกามปิ (สี, สฺยา, อิ), สุนฺทริกํ มหิํ (อิติปิ)}.} \\ \csromangathabat{สรสฺสติํ ปยาคญฺจ,}{อโถ พาหุมติํ นทิํ.}}\csromangathastanza{\csromangathabat{นิจฺจํปิ พาโล ปกฺขนฺโท\csromansharedfootnote{7911997a0d110b1b}{ปกฺขนฺโน (สี, สฺยา, อิ)},}{กณฺหกมฺโม น สุชฺฌติ.} \\ \csromangathabat{กิํ สุนฺทริกา กริสฺสติ,}{กิํ ปยาคา\csromansharedfootnote{7911997a0d110b1b}{ปกฺขนฺโน (สี, สฺยา, อิ)} กิํ พาหุกา นที.}}\csromangathastanza{\csromangathabat{เวริํ กตกิพฺพิสํ นรํ,}{น หิ นํ โสธเย ปาปกมฺมินํ.} \\ \csromangathabat{สุทฺธสฺส เว สทา ผคฺคุ,}{สุทฺธสฺสุโปสโถ สทา.}}\csromangathastanza{\csromanfolio{48}\csromangathabat{สุทฺธสฺส สุจิกมฺมสฺส,}{สทา สมฺปชฺชเต วตํ.} \\ \csromangathabat{อิเธว สินาหิ พฺราหฺมณ,}{สพฺพภูเตสุ กโรหิ เขมตํ.}}\csromangathastanza{\csromangathabat{สเจ มุสา น ภณสิ,}{สเจ ปาณํ น หิํสสิ.} \\ \csromangathabat{สเจ อทินฺนํ นาทิยสิ,}{สทฺทหาโน อมจฺฉรี.} \\ \csromangathabat{กิํ กาหสิ คยํ คนฺตฺวา,}{อุทปาโนปิ เต คยาติ.}}}

\proseitemwithcloserruled{80}{เอวํ วุตฺเต สุนฺทริก\-ภารทฺวาโช พฺราหฺมโณ ภควนฺตํ เอตทโวจ “อภิกฺกนฺตํ โภ โคตม, อภิกฺกนฺตํ โภ โคตม, เสยฺยถาปิ โภ โคตม นิกฺกุชฺชิตํ วา อุกฺกุชฺเชยฺย, ปฏิจฺฉนฺนํ วา วิวเรยฺย, มูฬฺหสฺส วา มคฺคํ อาจิกฺเขยฺย, อนฺธกาเร วา เตลปชฺโชตํ ธาเรยฺย ‘จกฺขุมนฺโต รูปานิ ทกฺขนฺตี’ติ. เอวเมวํ โภตา โคตมเน อเนก\-ปริยาเยน ธมฺโม ปกาสิโต, เอสาหํ ภวนฺตํ โคตมํ สรณํ คจฺฉามิ ธมฺมญฺจ ภิกฺขุสํฆญฺจ. ลเภยฺยาหํ โภโต โคตมสฺส สนฺติเก ปพฺพชฺชํ, ลเภยฺยํอุปสมฺปทนฺ”ติ. อลตฺถ โข สุนฺทริก\-ภารทฺวาโช พฺราหฺมโณ ภควโต สนฺติเก ปพฺพชฺชํ, อลตฺถ อุปสมฺปทํ. อจิรู\-ปสมฺปนฺโน โข ปนายสฺมา ภารทฺวาโช เอโก วูปกฏฺโฐ อปฺปมตฺโต อาตาปี ปหิตตฺโต วิหรนฺโต นจิรสฺเสว ยสฺสตฺถาย กุลปุตฺตา สมฺมเทว อคารสฺมา อนคาริยํ ปพฺพชนฺติ, ตทนุตฺตรํ พฺรหฺมจริย\-ปริโยสานํ ทิฏฺเฐวธมฺเม สยํ อภิญฺญา สจฺฉิกตฺวา อุปสมฺปชฺช วิหาสิ, “ขีณา ชาติ, วุสิตํ พฺรหฺมจริยํ, กตํ กรณียํ, นาปรํ อิตฺถตฺตายา”ติ อพฺภญฺญาสิ. อญฺญตโร โข ปนายสฺมา ภารทฺวาโช อรหตํ อโหสีติ.}{วตฺถสุตฺตํ นิฏฺฐิตํ สตฺตมํ.}
\proseitem{81}{เอวํ เม สุตํ\csromandash{}เอกํ สมยํ ภควา สาวตฺถิยํ วิหรติ เชตวเน อนาถ\-ปิณฺฑิกสฺส อาราเม. อถ โข อายสฺมา มหาจุนฺโท สายนฺหสมยํ ปฏิสลฺลานา วุฏฺฐิโต เยน ภควา เตนุปสงฺกมิ, อุปสงฺกมิตฺวา ภควนฺตํ อภิวาเทตฺวา เอกมนฺตํ นิสีทิ, เอกมนฺตํ นิสินฺโน \csromanfolio{49}โข อายสฺมา มหาจุนฺโท ภควนฺตํ เอตทโวจ “ยา อิมา ภนฺเต อเนกวิหิตา ทิฏฺฐิโย โลเก อุปฺปชฺชนฺติ อตฺตวาท\-ปฏิสํยุตฺตา วา โลกวาท\-ปฏิสํยุตฺตา วา. อาทิเมว นุ โข ภนฺเต ภิกฺขุโน มนสิกโรโต เอวเมตาสํ ทิฏฺฐีนํ ปหานํ โหติ, เอวเมตาสํ ทิฏฺฐีนํ ปฏินิสฺสคฺโค โหตี”ติ.}

\proseitem{82}{ยา อิมา จุนฺท อเนกวิหิตา ทิฏฺฐิโย โลเก อุปฺปชฺชนฺติ อตฺตวาท\-ปฏิสํยุตฺตา วา โลกวาท\-ปฏิสํยุตฺตา วา. ยตฺถ เจตา ทิฏฺฐิโย อุปฺปชฺชนฺติ, ยตฺถ จ อนุเสนฺติ, ยตฺถ จ สมุทาจรนฺติ. ตํ “เนตํ มม, เนโสหมสฺมิ, น เม โส อตฺตา”ติ เอวเมตํ ยถาภูตํ สมฺมปฺปญฺญาย ปสฺสโต เอวเมตาสํ ทิฏฺฐีนํ ปหานํ โหติ, เอวเมตาสํ ทิฏฺฐีนํ ปฏินิสฺสคฺโค โหติ.}

\prose{ฐานํ โข ปเนตํ จุนฺท วิชฺชติ, ยํ อิเธกจฺโจ ภิกฺขุ วิวิจฺเจว กาเมหิ วิวิจฺจ อกุสเลหิ ธมฺเมหิ สวิตกฺกํ สวิจารํ วิเวกชํ ปีติสุขํ ปฐมํ ฌานํ อุปสมฺปชฺช วิหเรยฺย, ตสฺส เอวมสฺส “สลฺเลเขน วิหรามี”ติ. น โข ปเนเต จุนฺท อริยสฺส วินเย สลฺเลขา วุจฺจนฺติ, ทิฏฺฐ\-ธมฺม\-สุข\-วิหารา เอเต อริยสฺส วินเย วุจฺจนฺติ.}

\prose{ฐานํ โข ปเนตํ จุนฺท วิชฺชติ, ยํ อิเธกจฺโจ ภิกฺขุ วิตกฺก\-วิจารานํ วูปสมา อชฺฌตฺตํ สมฺปสาทนํ เจตโส เอโกทิภาวํ อวิตกฺกํ อวิจารํ สมาธิชํ ปีติสุขํ ทุติยํ ฌานํ อุปสมฺปชฺช วิหเรยฺย, ตสฺส เอวมสฺส “สลฺเลเขน วิหรามี”ติ. น โข ปเนเต จุนฺท อริยสฺส วินเย สลฺเลขา วุจฺจนฺติ, ทิฏฺฐ\-ธมฺม\-สุข\-วิหารา เอเต อริยสฺส วินเย วุจฺจนฺติ.}

\prose{ฐานํ โข ปเนตํ จุนฺท วิชฺชติ, ยํ อิเธกจฺโจ ภิกฺขุ ปีติยา จ วิราคา อุเปกฺขโก จ วิหเรยฺย, สโต จ สมฺปชาโน สุขญฺจ กาเยน ปฏิสํเวเทยฺย, ยํ ตํ อริยา อาจิกฺขนฺติ “อุเปกฺขโก สติมา สุขวิหารี”ติ, ตติยํ ฌานํ อุปสมฺปชฺช วิหเรยฺย, ตสฺส เอวมสฺส “สลฺเลเขน วิหรามี”ติ. น โข ปเนเต จุนฺท อริยสฺส วินเย สลฺเลขา วุจฺจนฺติ, ทิฏฺฐ\-ธมฺม\-สุข\-วิหารา เอเต อริยสฺส วินเย วุจฺจนฺติ.}

\prose{ฐานํ โข ปเนตํ จุนฺท วิชฺชติ, ยํ อิเธกจฺโจ ภิกฺขุ สุขสฺส จ ปหานา ทุกฺขสฺส จ ปหานา ปุพฺเพว โสมนสฺส\-โทมนสฺสานํ อตฺถงฺคมา อทุกฺขมสุขํ \csromanfolio{50}อุเปกฺขา\-สติ\-ปาริสุทฺธิํ จตุตฺถํ ฌานํ อุปสมฺปชฺช วิหเรยฺย, ตสฺส เอวมสฺส “สลฺเลเขน วิหรามี”ติ. น โข ปเนเต จุนฺท อริยสฺส วินเย สลฺเลขา วุจฺจนฺติ, ทิฏฺฐ\-ธมฺม\-สุข\-วิหารา เอเต อริยสฺส วินเย วุจฺจนฺติ.}

\prose{ฐานํ โข ปเนตํ จุนฺท วิชฺชติ, ยํ อิเธกจฺโจ ภิกฺขุ สพฺพโส รูปสญฺญานํ สมติกฺกมา ปฏิฆ\-สญฺญานํ อตฺถงฺคมา นานตฺต\-สญฺญานํ อมนสิการา “อนนฺโต อากาโส”ติ อากาสา\-นญฺจา\-ยตนํ อุปสมฺปชฺช วิหเรยฺย, ตสฺส เอวมสฺส “สลฺเลเขน วิหรามี”ติ. น โข ปเนเต จุนฺท อริยสฺส วินเย สลฺเลขา วุจฺจนฺติ, สนฺตา เอเต วิหารา อริยสฺส วินเย วุจฺจนฺติ.}

\prose{ฐานํ โข ปเนตํ จุนฺท วิชฺชติ, ยํ อิเธกจฺโจ ภิกฺขุ สพฺพโส อากาสา\-นญฺจา\-ยตนํ สมติกฺกมฺม “อนนฺตํ วิญฺญาณนฺ”ติ วิญฺญาณญฺจา\-ยตนํ อุปสมฺปชฺช วิหเรยฺย, ตสฺส เอวมสฺส “สลฺเลเขน วิหรามี”ติ. น โข ปเนเต จุนฺท อริยสฺส วินเย สลฺเลขา วุจฺจนฺติ, สนฺตา เอเต วิหารา อริยสฺส วินเย วุจฺจนฺติ.}

\prose{ฐานํ โข ปเนตํ จุนฺท วิชฺชติ, ยํ อิเธกจฺโจ ภิกฺขุ สพฺพโส วิญฺญาณญฺจา\-ยตนํ สมติกฺกมฺม “นตฺถิ กิญฺจี”ติ อากิญฺจญฺญา\-ยตนํ อุปสมฺปชฺช วิหเรยฺย, ตสฺส เอวมสฺส “สลฺเลเขน วิหรามี”ติ. น โข ปเนเต จุนฺท อริยสฺส วินเย สลฺเลขา วุจฺจนฺติ, สนฺตา เอเต วิหารา อริยสฺส วินเย วุจฺจนฺติ.}

\prose{ฐานํ โข ปเนตํ จุนฺท วิชฺชติ, ยํ อิเธกจฺโจ ภิกฺขุ สพฺพโส อากิญฺจญฺญา\-ยตนํ สมติกฺกมฺม เนว\-สญฺญา\-นา\-สญฺญา\-ยตนํ อุปสมฺปชฺช วิหเรยฺย, ตสฺส เอวมสฺส “สลฺเลเขน วิหรามี”ติ. น โข ปเนเต จุนฺท อริยสฺส วินเย สลฺเลขา วุจฺจนฺติ, สนฺตา เอเต วิหารา อริยสฺส วินเย วุจฺจนฺติ.}

\proseitem{83}{อิธ โข ปน โว จุนฺท สลฺเลโข กรณีโย. ปเร วิหิํสกา ภวิสฺสนฺติ, มยเมตฺถ อวิหิํสกา ภวิสฺสามาติ สลฺเลโข กรณีโย. ปเร ปาณาติปาตี ภวิสฺสนฺติ, มยเมตฺถ ปาณาติปาตา ปฏิวิรตา ภวิสฺสามาติ สลฺเลโข กรณีโย. ปเร อทินฺนาทายี ภวิสฺสนฺติ, มยเมตฺถ อทินฺนาทานา ปฏิวิรตา ภวิสฺสามาติ สลฺเลโข กรณีโย. \csromanfolio{51}ปเร อพฺรหฺมจารี ภวิสฺสนฺติ, มยเมตฺถ พฺรหฺมจารี ภวิสฺสามาติ สลฺเลโข กรณีโย. ปเร มุสาวาที ภวิสฺสนฺติ, มยเมตฺถ มุสาวาทา ปฏิวิรตา ภวิสฺสามาติ สลฺเลโข กรณีโย. ปเร ปิสุณวาจา\csromansharedfootnote{eb02e4ebadecba47}{ปิสุณา วาจา (สี, อิ)} ภวิสฺสนฺติ, มยเมตฺถ ปิสุณาย วาจาย ปฏิวิรตา ภวิสฺสามาติ สลฺเลโข กรณีโย. ปเร ผรุสวาจา\csromansharedfootnote{22839f23d1bb46c6}{ผรุสา วาจา (สี, อิ)} ภวิสฺสนฺติ, มยเมตฺถ ผรุสาย วาจาย ปฏิวิรตา ภวิสฺสามาติ สลฺเลโข กรณีโย. ปเร สมฺผปฺปลาปี ภวิสฺสนฺติ, มยเมตฺถ สมฺผปฺปลาปา ปฏิวิรตา ภวิสฺสามาติ สลฺเลโข กรณีโย. ปเร อภิชฺฌาลู ภวิสฺสนฺติ, มยเมตฺถ อนภิชฺฌาลู ภวิสฺสามาติ สลฺเลโข กรณีโย. ปเร พฺยาปนฺนจิตฺตา ภวิสฺสนฺติ, มยเมตฺถ อพฺยาปนฺนจิตฺตา ภวิสฺสามาติ สลฺเลโข กรณีโย. ปเร มิจฺฉาทิฏฺฐี ภวิสฺสนฺติ, มยเมตฺถ สมฺมาทิฏฺฐี ภวิสฺสมาติ สลฺเลโข กรณีโย. ปเร มิจฺฉาสงฺกปฺปา ภวิสฺสนฺติ, มยเมตฺถ สมฺมาสงฺกปฺปา ภวิสฺสามาติ สลฺเลโข กรณีโย. ปเร มิจฺฉาวาจา ภวิสฺสนฺติ, มยเมตฺถ สมฺมาวาจา ภวิสฺสามาติ สลฺเลโข กรณีโย. ปเร มิจฺฉากมฺมนฺตา ภวิสฺสนฺติ, มยเมตฺถ สมฺมากมฺมนฺตา ภวิสฺสามาติ สลฺเลโข กรณีโย. ปเร มิจฺฉาอาชีวา ภวิสฺสนฺติ, มยเมตฺถ สมฺมาอาชีวา ภวิสฺสามาติ สลฺเลโข กรณีโย. ปเร มิจฺฉาวายามา ภวิสฺสนฺติ, มยเมตฺถ สมฺมาวายามา ภวิสฺสามาติ สลฺเลโข กรณีโย. ปเร มิจฺฉาสตี ภวิสฺสนฺติ, มยเมตฺถ สมฺมาสตี ภวิสฺสามาติ สลฺเลโข กรณีโย. ปเร มิจฺฉาสมาธี ภวิสฺสนฺติ, มยเมตฺถ สมฺมาสมาธี ภวิสฺสามาติ สลฺเลโข กรณีโย. ปเร มิจฺฉาญาณี ภวิสฺสนฺติ, มยเมตฺถ สมฺมาญาณี ภวิสฺสามาติ สลฺเลโข กรณีโย. ปเร มิจฺฉาวิมุตฺตี ภวิสฺสนฺติ, มยเมตฺถ สมฺมาวิมุตฺตี ภวิสฺสามาติ สลฺเลโข กรณีโย.}

\prose{ปเร ถีนมิทฺธ\-ปริยุฏฺฐิตา ภวิสฺสนฺติ, มยเมตฺถ วิคต\-ถีน\-มิทฺธา ภวิสฺสามาติ สลฺเลโข กรณีโย. ปเร อุทฺธตา ภวิสฺสนฺติ, มยเมตฺถ อนุทฺธตา ภวิสฺสามาติ สลฺเลโข กรณีโย. ปเร วิจิกิจฺฉี\csromansharedfootnote{c0a2595c809f0833}{เวจิกิจฺฉี (สี, อิ, ก)} ภวิสฺสนฺติ, มยเมตฺถ ติณฺณ\-วิจิกิจฺฉา ภวิสฺสามาติ สลฺเลโข กรณีโย. ปเร โกธนา ภวิสฺสนฺติ, มยเมตฺถ อกฺโกธนา ภวิสฺสามาติ สลฺเลโข กรณีโย. ปเร อุปนาหี ภวิสฺสนฺติ, มยเมตฺถ อนุปนาหี ภวิสฺสามาติ สลฺเลโข กรณีโย. ปเร มกฺขี \csromanfolio{52}ภวิสฺสนฺติ, มยเมตฺถ อมกฺขี ภวิสฺสามาติ สลฺเลโข กรณีโย. ปเร ปฬาสี ภวิสฺสนฺติ, มยเมตฺถ อปฬาสี ภวิสฺสามาติ สลฺเลโข กรณีโย. ปเร อิสฺสุกี ภวิสฺสนฺติ, มยเมตฺถ อนิสฺสุกี ภวิสฺสามาติ สลฺเลโข กรณีโย. ปเร มจฺฉรี ภวิสฺสนฺติ, มยเมตฺถ อมจฺฉรี ภวิสฺสามาติ สลฺเลโข กรณีโย. ปเร สฐา ภวิสฺสนฺติ, มยเมตฺถ อสฐา ภวิสฺสามาติ สลฺเลโข กรณีโย. ปเร มายาวี ภวิสฺสนฺติ, มยเมตฺถ อมายาวี ภวิสฺสามาติ สลฺเลโข กรณีโย. ปเร ถทฺธา ภวิสฺสนฺติ, มยเมตฺถ อตฺถทฺธา ภวิสฺสามาติ สลฺเลโข กรณีโย. ปเร อติมานี ภวิสฺสนฺติ, มยเมตฺถ อนติมานี ภวิสฺสามาติ สลฺเลโข กรณีโย. ปเร ทุพฺพจา ภวิสฺสนฺติ, มยเมตฺถ สุวจา ภวิสฺสามาติ สลฺเลโข กรณีโย. ปเร ปาปมิตฺตา ภวิสฺสนฺติ, มยเมตฺถ กลฺยาณมิตฺตา ภวิสฺสมาติ สลฺเลโข กรณีโย. ปเร ปมตฺตา ภวิสฺสนฺติ, มยเมตฺถ อปฺปมตฺตา ภวิสฺสามาติ สลฺเลโข กรณีโย. ปเร อสฺสทฺธา ภวิสฺสนฺติ, มยเมตฺถ สทฺธา ภวิสฺสามาติ สลฺเลโข กรณีโย. ปเร อหิริกา ภวิสฺสนฺติ, มยเมตฺถ หิริมนา ภวิสฺสามาติ สลฺเลโข กรณีโย. ปเร อโนตฺตาปี\csromansharedfootnote{74a5b581c40eec1b}{อโนตฺตปฺปี (สฺยา, ก)} ภวิสฺสนฺติ, มยเมตฺถ โอตฺตาปี ภวิสฺสามาติ สลฺเลโข กรณีโย. ปเร อปฺปสฺสุตา ภวิสฺสนฺติ, มยเมตฺถ พหุสฺสุตา ภวิสฺสามาติ สลฺเลโข กรณีโย. ปเร กุสีตา ภวิสฺสนฺติ, มยเมตฺถ อารทฺธวีริยา ภวิสฺสามาติ สลฺเลโข กรณีโย. ปเร มุฏฺฐสฺสตี ภวิสฺสนฺติ, มยเมตฺถ อุปฏฺฐิตสฺสตี ภวิสฺสามาติ สลฺเลโข กรณีโย. ปเร ทุปฺปญฺญา ภวิสฺสนฺติ, มยเมตฺถ ปญฺญาสมฺปนฺนา ภวิสฺสามาติ สลฺเลโข กรณีโย. ปเร สนฺทิฏฺฐิ\-ปรามาสี อาธานคฺคาหี ทุปฺปฏินิสฺสคฺคี ภวิสฺสนฺติ, มยเมตฺถ อสนฺทิฏฺฐิ\-ปรามาสี อนาธานคฺคาหี สุปฺปฏินิสฺสคฺคี ภวิสฺสามาติ สลฺเลโข กรณีโย.}

\proseitem{84}{จิตฺตุปฺปาทมฺปิ โข อหํ จุนฺท กุสเลสุ ธมฺเมสุ พหุการํ\csromansharedfootnote{394a8154fe264c6c}{พหูปการํ (ก)} วทามิ, โก ปน วาโท กาเยน วาจาย อนุวิธียนาสุ. ตสฺมาติห จุนฺท ปเร วิหิํสกา ภวิสฺสนฺติ, มยเมตฺถ อวิหิํสกา ภวิสฺสามาติ จิตฺตํ อุปฺปาเทตพฺพํ. ปเร ปาณาติปาตี ภวิสฺสนฺติ, มยเมตฺถ ปาณาติปาตา \csromanfolio{53}ปฏิวิรตา ภวิสฺสามาติ จิตฺตํ อุปฺปาเทตพฺพํ ฯเปฯ. ปเร สนฺทิฏฺฐิ\-ปรามาสี อาธานคฺคาหี ทุปฺปฏินิสฺสคฺคี ภวิสฺสนฺติ, มยเมตฺถ อสนฺทิฏฺฐิ\-ปรามาสี อนาธานคฺคาหี สุปฺปฏินิสฺสคฺคี ภวิสฺสามาติ จิตฺตํ อุปฺปาเทตพฺพํ.}

\proseitem{85}{เสยฺยถาปิ จุนฺท วิสโม มคฺโค อสฺส, ตสฺส\csromansharedfootnote{fb97e67d1dad9f35}{มคฺโค ตสฺสาสฺส (สี, สฺยา, อิ)} อญฺโญ สโม มคฺโค ปริกฺกมนาย. เสยฺยถา วา ปน จุนฺท วิสมํ ติตฺถํ อสฺส, ตสฺส อญฺญํ สมํ ติตฺถํ ปริกฺกมนาย. เอวเมว โข จุนฺท วิหิํสกสฺส ปุริส\-ปุคฺคลสฺส อวิหิํสา โหติ ปริกฺกมนาย. ปาณาติปาติสฺส ปุริส\-ปุคฺคลสฺส ปาณาติปาตา เวรมณี โหติ ปริกฺกมนาย. อทินฺนาทายิสฺส ปุริส\-ปุคฺคลสฺส อทินฺนาทานา เวรมณี โหติ ปริกฺกมนาย. อพฺรหฺมจาริสฺส ปุริส\-ปุคฺคลสฺส อพฺรหฺมจริยา เวรมณี โหติ ปริกฺกมนาย. มุสาวาทิสฺส ปุริส\-ปุคฺคลสฺส มุสาวาทา เวรมณี โหติ ปริกฺกมนาย. ปิสุณวาจสฺส ปุริส\-ปุคฺคลสฺส ปิสุณาย วาจาย เวรมณี โหติ ปริกฺกมนาย. ผรุสวาจสฺส ปุริส\-ปุคฺคลสฺส ผรุสาย วาจาย เวรมณี โหติ ปริกฺกมนาย. สมฺผ\-ปฺปลาปิสฺส ปุริส\-ปุคฺคลสฺส สมฺผปฺปลาปา เวรมณี โหติ ปริกฺกมนาย. อภิชฺฌาลุสฺส ปุริส\-ปุคฺคลสฺส อนภิชฺฌา โหติ ปริกฺกมนาย. พฺยาปนฺน\-จิตฺตสฺส ปุริส\-ปุคฺคลสฺส อพฺยาปาโท โหติ ปริกฺกมนาย. มิจฺฉา\-ทิฏฺฐิสฺส ปุริส\-ปุคฺคลสฺส สมฺมาทิฏฺฐิ โหติ ปริกฺกมนาย. มิจฺฉา\-สงฺกปฺปสฺส ปุริส\-ปุคฺคลสฺส สมฺมาสงฺกปฺโป โหติ ปริกฺกมนาย. มิจฺฉาวาจสฺส ปุริส\-ปุคฺคลสฺส สมฺมาวาจา โหติ ปริกฺกมนาย. มิจฺฉา\-กมฺมนฺตสฺส ปุริส\-ปุคฺคลสฺส สมฺมากมฺมนฺโต โหติ ปริกฺกมนาย. มิจฺฉาอาชีวสฺส ปุริส\-ปุคฺคลสฺส สมฺมาอาชีโว โหติ ปริกฺกมนาย. มิจฺฉา\-วายามสฺส ปุริส\-ปุคฺคลสฺส สมฺมาวายาโม โหติ ปริกฺกมนาย. มิจฺฉาสติสฺส ปุริส\-ปุคฺคลสฺส สมฺมาสติ โหติ ปริกฺกมนาย. มิจฺฉา\-สมาธิสฺส ปุริส\-ปุคฺคลสฺส สมฺมาสมาธิ โหติ ปริกฺกมนาย. มิจฺฉาญาณิสฺส ปุริส\-ปุคฺคลสฺส สมฺมาญาณํ โหติ ปริกฺกมนาย. มิจฺฉา\-วิมุตฺติสฺส ปุริส\-ปุคฺคลสฺส สมฺมาวิมุตฺติ โหติ ปริกฺกมนาย.}

\prose{ถีนมิทฺธ\-ปริยุฏฺฐิตสฺส ปุริส\-ปุคฺคลสฺส วิคต\-ถีนมิทฺธตา โหติ ปริกฺกมนาย. อุทฺธตสฺส ปุริส\-ปุคฺคลสฺส อนุทฺธจฺจํ โหติ ปริกฺกมนาย. วิจิกิจฺฉิสฺส ปุริส\-ปุคฺคลสฺส ติณฺณวิจิกิจฺฉตา โหติ ปริกฺกมนาย. โกธนสฺส \csromanfolio{54}ปุริส\-ปุคฺคลสฺส อกฺโกโธ โหติ ปริกฺกมนาย. อุปนาหิสฺส ปุริส\-ปุคฺคลสฺส อนุปนาโห โหติ ปริกฺกมนาย. มกฺขิสฺส ปุริส\-ปุคฺคลสฺส อมกฺโข โหติ ปริกฺกมนาย. ปฬาสิสฺส ปุริส\-ปุคฺคลสฺส อปฬาโส โหติ ปริกฺกมนาย. อิสฺสุกิสฺส ปุริส\-ปุคฺคลสฺส อนิสฺสุกิตา โหติ ปริกฺกมนาย. มจฺฉริสฺส ปุริส\-ปุคฺคลสฺส อมจฺฉริยํ โหติ ปริกฺกมนาย. สฐสฺส ปุริส\-ปุคฺคลสฺส อสาเฐยฺยํ โหติ ปริกฺกมนาย. มายาวิสฺส ปุริส\-ปุคฺคลสฺส อมายา\csromansharedfootnote{a1a2174328e4ba1e}{อมายาวิตา (ก)} โหติ ปริกฺกมนาย. ถทฺธสฺส ปุริส\-ปุคฺคลสฺส อตฺถทฺธิยํ โหติ ปริกฺกมนาย. อติมานิสฺส ปุริส\-ปุคฺคลสฺส อนติมาโน โหติ ปริกฺกมนาย. ทุพฺพจสฺส ปุริส\-ปุคฺคลสฺส โสวจสฺสตา โหติ ปริกฺกมนาย. ปาปมิตฺตสฺส ปุริส\-ปุคฺคลสฺส กลฺยาณมิตฺตตา โหติ ปริกฺกมนาย. ปมตฺตสฺส ปุริส\-ปุคฺคลสฺส อปฺปมาโท โหติ ปริกฺกมนาย. อสฺสทฺธสฺส ปุริส\-ปุคฺคลสฺส สทฺธา โหติ ปริกฺกมนาย. อหิริกสฺส ปุริส\-ปุคฺคลสฺส หิรี โหติ ปริกฺกมนาย. อโนตฺตาปิสฺส ปุริส\-ปุคฺคลสฺส โอตฺตปฺปํ โหติ ปริกฺกมนาย. อปฺปสฺสุตสฺส ปุริส\-ปุคฺคลสฺส พาหุสจฺจํ โหติ ปริกฺกมนาย. กุสีตสฺส ปุริส\-ปุคฺคลสฺส วีริยารมฺโภ โหติ ปริกฺกมนาย. มุฏฺฐสฺสติสฺส ปุริส\-ปุคฺคลสฺส อุปฏฺฐิต\-สฺสติตา โหติ ปริกฺกมนาย. ทุปฺปญฺญสฺส ปุริส\-ปุคฺคลสฺส ปญฺญาสมฺปทา โหติ ปริกฺกมนาย. สนฺทิฏฺฐิ\-ปรามาสิ อาธานคฺคาหิ ทุปฺปฏินิสฺสคฺคิสฺส ปุริส\-ปุคฺคลสฺส อสนฺทิฏฺฐิ\-ปรามาสิ อนาธานคฺคาหิ สุปฺปฏินิสฺสคฺคิตา โหติ ปริกฺกมนาย.}

\proseitem{86}{เสยฺยถาปิ จุนฺท เยเกจิ อกุสลา ธมฺมา, สพฺเพ เต อโธภาค\-งฺคมนียา\csromansharedfootnote{e684b8c631ff4ae2}{อโธภาวงฺคมนียา (สี, สฺยา, อิ)}. เยเกจิ กุสลา ธมฺมา, สพฺเพ เต อุปริภาค\-งฺคมนียา\csromansharedfootnote{24814f0a32e53e0e}{อุปริภาวงฺคมนียา (สี, สฺยา, อิ)}. เอวเมว โข จุนฺท วิหิํสกสฺส ปุริส\-ปุคฺคลสฺส อวิหิํสา โหติ อุปริภาคาย\csromansharedfootnote{53c48f181aeee88c}{อุปริภาวาย (สี, สฺยา, ก)}. ปาณาติปาติสฺส ปุริส\-ปุคฺคลสฺส ปาณาติปาตา เวรมณี โหติ อุปริภาคาย ฯเปฯ สนฺทิฏฺฐิ\-ปรามาสิ อาธานคฺคาหิ ทุปฺปฏินิสฺสคฺคิสฺส ปุริส\-ปุคฺคลสฺส อสนฺทิฏฺฐิ\-ปรามาสิ อนาธานคฺคาหิ สุปฺปฏินิสฺสคฺคิตา โหติ อุปริภาคาย.}

\proseitem{87}{โส วต จุนฺท อตฺตนา ปลิป\-ปลิปนฺโน ปรํ ปลิป\-ปลิปนฺนํ อุทฺธริสฺสตีติ เนตํ ฐานํ วิชฺชติ. โส วต จุนฺท อตฺตนา อปลิป\-ปลิปนฺโน \csromanfolio{55}ปรํ ปลิป\-ปลิปนฺนํ อุทฺธริสฺสตีติ ฐานเมตํ วิชฺชติ. โส วต จุนฺท อตฺตนา อทนฺโต อวินีโต อปรินิพฺพุโต ปรํ ทเมสฺสติ วิเนสฺสติ ปรินิพฺพาเปสฺสตีติ เนตํ ฐานํ วิชฺชติ. โส วต จุนฺท อตฺตนา ทนฺโต วินีโต ปรินิพฺพุโต ปรํ ทเมสฺสติ วิเนสฺสติ ปรินิพฺพาเปสฺสตีติ ฐานเมตํ วิชฺชติ. เอวเมว โข จุนฺท วิหิํสกสฺส ปุริส\-ปุคฺคลสฺส อวิหิํสา โหติ ปรินิพฺพานาย. ปาณาติปาติสฺส ปุริส\-ปุคฺคลสฺส ปาณาติปาตา เวรมณี โหติ ปรินิพฺพานาย. อทินฺนาทายิสฺส ปุริส\-ปุคฺคลสฺส อทินฺนาทานา เวรมณี โหติ ปรินิพฺพานาย. อพฺรหฺมจาริสฺส ปุริส\-ปุคฺคลสฺส อพฺรหฺมจริยา เวรมณี โหติ ปรินิพฺพานาย. มุสาวาทิสฺส ปุริส\-ปุคฺคลสฺส มุสาวาทา เวรมณี โหติ ปรินิพฺพานาย. ปิสุณวาจสฺส ปุริส\-ปุคฺคลสฺส ปิสุณาย วาจาย เวรมณี โหติ ปรินิพฺพานาย. ผรุสวาจสฺส ปุริส\-ปุคฺคลสฺส ผรุสาย วาจาย เวรมณี โหติ ปรินิพฺพานาย. สมฺผ\-ปฺปลาปิสฺส ปุริส\-ปุคฺคลสฺส สมฺผปฺปลาปา เวรมณี โหติ ปรินิพฺพานาย. อภิชฺฌาลุสฺส ปุริส\-ปุคฺคลสฺส อนภิชฺฌา โหติ ปรินิพฺพานาย. พฺยาปนฺน\-จิตฺตสฺส ปุริส\-ปุคฺคลสฺส อพฺยาปาโท โหติ ปรินิพฺพานาย. มิจฺฉา\-ทิฏฺฐิสฺส ปุริส\-ปุคฺคลสฺส สมฺมาทิฏฺฐิ โหติ ปรินิพฺพานาย. มิจฺฉา\-สงฺกปฺปสฺส ปุริส\-ปุคฺคลสฺส สมฺมาสงฺกปฺโป โหติ ปรินิพฺพานาย. มิจฺฉาวาจสฺส ปุริส\-ปุคฺคลสฺส สมฺมาวาจา โหติ ปรินิพฺพานาย. มิจฺฉา\-กมฺมนฺตสฺส ปุริส\-ปุคฺคลสฺส สมฺมากมฺมนฺโต โหติ ปรินิพฺพานาย. มิจฺฉาอาชีวสฺส ปุริส\-ปุคฺคลสฺส สมฺมาอาชีโว โหติ ปรินิพฺพานาย. มิจฺฉา\-วายามสฺส ปุริส\-ปุคฺคลสฺส สมฺมาวายาโม โหติ ปรินิพฺพานาย. มิจฺฉาสติสฺส ปุริส\-ปุคฺคลสฺส สมฺมาสติ โหติ ปรินิพฺพานาย. มิจฺฉา\-สมาธิสฺส ปุริส\-ปุคฺคลสฺส สมฺมาสมาธิ โหติ ปรินิพฺพานาย. มิจฺฉาญาณิสฺส ปุริส\-ปุคฺคลสฺส สมฺมาญาณํ โหติ ปรินิพฺพานาย. มิจฺฉา\-วิมุตฺติสฺส ปุริส\-ปุคฺคลสฺส สมฺมาวิมุตฺติ โหติ ปรินิพฺพานาย.}

\prose{ถีนมิทฺธ\-ปริยุฏฺฐิตสฺส ปุริส\-ปุคฺคลสฺส วิคต\-ถีนมิทฺธตา โหติ ปรินิพฺพานาย. อุทฺธตสฺส ปุริส\-ปุคฺคลสฺส อนุทฺธจฺจํ โหติ ปรินิพฺพานาย. วิจิกิจฺฉิสฺส ปุริส\-ปุคฺคลสฺส ติณฺณวิจิกิจฺฉตา โหติ ปรินิพฺพานาย. โกธนสฺส ปุริส\-ปุคฺคลสฺส อกฺโกโธ โหติ ปรินิพฺพานาย. อุปนาหิสฺส ปุริส\-ปุคฺคลสฺส อนุปนาโห โหติ ปรินิพฺพานาย. มกฺขิสฺส ปุริส\-ปุคฺคลสฺส อมกฺโข โหติ ปรินิพฺพานาย. ปฬาสิสฺส ปุริส\-ปุคฺคลสฺส \csromanfolio{56}อปฬาโส โหติ ปรินิพฺพานาย. อิสฺสุกิสฺส ปุริส\-ปุคฺคลสฺส อนิสฺสุกิตา โหติ ปรินิพฺพานาย. มจฺฉริสฺส ปุริส\-ปุคฺคลสฺส อมจฺฉริยํ โหติ ปรินิพฺพานาย. สฐสฺส ปุริส\-ปุคฺคลสฺส อสาเฐยฺยํ โหติ ปรินิพฺพานาย. มายาวิสฺส ปุริส\-ปุคฺคลสฺส อมายา โหติ ปรินิพฺพานาย. ถทฺธสฺส ปุริส\-ปุคฺคลสฺส อตฺถทฺธิยํ โหติ ปรินิพฺพานาย. อติมานิสฺส ปุริส\-ปุคฺคลสฺส อนติมาโน โหติ ปรินิพฺพานาย. ทุพฺพจสฺส ปุริส\-ปุคฺคลสฺส โสวจสฺสตา โหติ ปรินิพฺพานาย. ปาปมิตฺตสฺส ปุริส\-ปุคฺคลสฺส กลฺยาณมิตฺตตา โหติ ปรินิพฺพานาย. ปมตฺตสฺส ปุริส\-ปุคฺคลสฺส อปฺปมาโท โหติ ปรินิพฺพานาย. อสฺสทฺธสฺส ปุริส\-ปุคฺคลสฺส สทฺธา โหติ ปรินิพฺพานาย. อหิริกสฺส ปุริส\-ปุคฺคลสฺส หิรี โหติ ปรินิพฺพานาย. อโนตฺตาปิสฺส ปุริส\-ปุคฺคลสฺส โอตฺตปฺปํ โหติ ปรินิพฺพานาย. อปฺปสฺสุตสฺส ปุริส\-ปุคฺคลสฺส พาหุสจฺจํ โหติ ปรินิพฺพานาย. กุสีตสฺส ปุริส\-ปุคฺคลสฺส วีริยารมฺโภ โหติ ปรินิพฺพานาย. มุฏฺฐสฺสติสฺส ปุริส\-ปุคฺคลสฺส อุปฏฺฐิต\-สฺสติตา โหติ ปรินิพฺพานาย. ทุปฺปญฺญสฺส ปุริส\-ปุคฺคลสฺส ปญฺญาสมฺปทา โหติ ปรินิพฺพานาย. สนฺทิฏฺฐิ\-ปรามาสิ อาธานคฺคาหิ ทุปฺปฏินิสฺสคฺคิสฺส ปุริส\-ปุคฺคลสฺส อสนฺทิฏฺฐิ\-ปรามาสิ อนาธานคฺคาหิ สุปฺปฏินิสฺสคฺคิตา โหติ ปรินิพฺพานาย.}

\proseitem{88}{อิติ โข จุนฺท เทสิโต มยา สลฺเลข\-ปริยาโย, เทสิโต จิตฺตุปฺปาท\-ปริยาโย, เทสิโต ปริกฺกมน\-ปริยาโย, เทสิโต อุปริภาค\-ปริยาโย, เทสิโต ปรินิพฺพาน\-ปริยาโย. ยํ โข จุนฺท สตฺถารา กรณียํ สาวกานํ หิเตสินา อนุกมฺปเกน อนุกมฺปํ อุปาทาย, กตํ โว ตํ มยา. เอตานิ จุนฺท รุกฺขมูลานิ, เอตานิ สุญฺญาคารานิ. ฌายถ จุนฺท มา ปมาทตฺถ, มา ปจฺฉา\-วิปฺปฏิสาริโน อหุวตฺถ. อยํ โข อมฺหากํ อนุสาสนีติ.}

\prose{อิทมโวจ ภควา. อตฺตมโน อายสฺมา มหาจุนฺโท ภควโต ภาสิตํ อภินนฺทีติ.}

\setlength{\csromangathapagewidth}{0pt}
\setlength{\csromangathawakwidth}{0pt}
\csromangathameasuregroup{จตุตฺตาลีสปทา วุตฺตา, \\ สลฺเลโข นาม สุตฺตนฺโต,}{\csromangathabat{จตุตฺตาลีสปทา วุตฺตา,}{สนฺธโย ปญฺจ เทสิตา.} \\ \csromangathabat{สลฺเลโข นาม สุตฺตนฺโต,}{คมฺภีโร สาครูปโมติ.}}
\csromangathasetleft{จตุตฺตาลีสปทา วุตฺตา, \\ สลฺเลโข นาม สุตฺตนฺโต,}
\csromangathagroupwithcloser{\csromangathastanza{\csromangathabat{จตุตฺตาลีสปทา วุตฺตา,}{สนฺธโย ปญฺจ เทสิตา.} \\ \csromangathabat{สลฺเลโข นาม สุตฺตนฺโต,}{คมฺภีโร สาครูปโมติ.}}}{สลฺเลขสุตฺตํ นิฏฺฐิตํ อฏฺฐมํ.}
\csromanmatikaanchor{mtk.17}
\csromantocmarkref{subsection}{9. สมฺมาทิฏฺฐิสุตฺต}{mtk.17}
\bookmark[dest={mtk.17},level=2]{9. สมฺมาทิฏฺฐิสุตฺต}
\csromanheader{2}{9. สมฺมาทิฏฺฐิ\-สุตฺต 9. สมฺมาทิฏฺฐิ\-สุตฺต}
\proseitem{89}{\csromanfolio{57}เอวํ เม สุตํ\csromandash{}เอกํ สมยํ ภควา สาวตฺถิยํ วิหรติ เชตวเน อนาถ\-ปิณฺฑิกสฺส อาราเม. ตตฺร โข อายสฺมา สาริปุตฺโต ภิกฺขู อามนฺเตสิ “อาวุโส ภิกฺขเว”ติ. “อาวุโส”ติ โข เต ภิกฺขู อายสฺมโต สาริปุตฺตสฺส ปจฺจสฺโสสุํ. อายสฺมา สาริปุตฺโต เอตทโวจ\csromandash{}}

\prose{“สมฺมาทิฏฺฐิ\csromansharedfootnote{b516d8861ad4b1c0}{สมฺมาทิฏฺฐี (สี, สฺยา)} สมฺมาทิฏฺฐีติ อาวุโส วุจฺจติ. กิตฺตาวตา นุ โข อาวุโส อริยสาวโก สมฺมาทิฏฺฐิ โหติ, อุชุคตาสฺส ทิฏฺฐิ, ธมฺเม อเวจฺจ\-ปฺปสาเทน สมนฺนาคโต, อาคโต อิมํ สทฺธมฺมนฺ”ติ.}

\prose{ทูรโตปิ โข มยํ อาวุโส อาคจฺเฉยฺยาม อายสฺมโต สาริปุตฺตสฺส สนฺติเก เอตสฺส ภาสิตสฺส อตฺถมญฺญาตุํ. สาธุ วตายสฺมนฺตํเยว สาริปุตฺตํ ปฏิภาตุ เอตสฺส ภาสิตสฺส อตฺโถ, อายสฺมโต สาริปุตฺตสฺส สุตฺวา ภิกฺขู ธาเรสฺสนฺตีติ. เตน หิ อาวุโส สุณาถ สาธุกํ มนสิกโรถ ภาสิสฺสามีติ. “เอวมาวุโส”ติ โข เต ภิกฺขู อายสฺมโต สาริปุตฺตสฺส ปจฺจสฺโสสุํ. อายสฺมา สาริปุตฺโต เอตทโวจ\csromandash{}}

\prose{“ยโต โข อาวุโส อริยสาวโก อกุสลญฺจ ปชานาติ, อกุสลมูลญฺจ ปชานาติ, กุสลญฺจ ปชานาติ, กุสลมูลญฺจ ปชานาติ. เอตฺตาวตาปิ โข อาวุโส อริยสาวโก สมฺมาทิฏฺฐิ โหติ. อุชุคตาสฺส ทิฏฺฐิ, ธมฺเม อเวจฺจ\-ปฺปสาเทน สมนฺนาคโต, อาคโต อิมํ สทฺธมฺมํ. กตมํ ปนาวุโส อกุสลํ, กตมํ อกุสลมูลํ, กตมํ กุสลํ, กตมํ กุสลมูลํ. ปาณาติปาโต โข อาวุโส อกุสลํ, อทินฺนาทานํ อกุสลํ, กาเมสุ\-มิจฺฉาจาโร อกุสลํ, มุสาวาโท อกุสลํ, ปิสุณา วาจา\csromansharedfootnote{6a383bc90811e18d}{ปิสุณวาจา (ก)} อกุสลํ, ผรุสา วาจา\csromansharedfootnote{99f9ecace336c802}{ผรุสวาจา (ก)} อกุสลํ, สมฺผปฺปลาโป อกุสลํ, อภิชฺฌา อกุสลํ, พฺยาปาโท อกุสลํ, มิจฺฉาทิฏฺฐิ อกุสลํ. อิทํ วุจฺจตาวุโส อกุสลํ. กตมญฺจาวุโส อกุสลมูลํ. โลโภ อกุสลมูลํ, โทโส อกุสลมูลํ, โมโห อกุสลมูลํ. อิทํ วุจฺจตาวุโส อกุสลมูลํ.}

\prose{\csromanfolio{58}กตมญฺจาวุโส กุสลํ. ปาณาติปาตา เวรมณี กุสลํ, อทินฺนาทานา เวรมณี กุสลํ, กาเมสุ\-มิจฺฉาจารา เวรมณี กุสลํ, มุสาวาทา เวรมณี กุสลํ, ปิสุณาย วาจาย เวรมณี กุสลํ, ผรุสาย วาจาย เวรมณี กุสลํ, สมฺผปฺปลาปา เวรมณี กุสลํ, อนภิชฺฌา กุสลํ, อพฺยาปาโท กุสลํ, สมฺมาทิฏฺฐิ กุสลํ. อิทํ วุจฺจตาวุโส กุสลํ. กตมญฺจาวุโส กุสลมูลํ. อโลโภ กุสลมูลํ, อโทโส กุสลมูลํ, อโมโห กุสลมูลํ. อิทํ วุจฺจตาวุโส กุสลมูลํ.}

\prose{ยโต โข อาวุโส อริยสาวโก เอวํ อกุสลํ ปชานาติ, เอวํ อกุสลมูลํ ปชานาติ, เอวํ กุสลํ ปชานาติ, เอวํ กุสลมูลํ ปชานาติ. โส สพฺพโส ราคานุสยํ ปหาย ปฏิฆานุสยํ ปฏิวิโนเทตฺวา อสฺมีติ ทิฏฺฐิ\-มานา\-นุสยํ สมูหนิตฺวา อวิชฺชํ ปหาย วิชฺชํ อุปฺปาเทตฺวา ทิฏฺเฐวธมฺเม ทุกฺขสฺส\-นฺตกโร โหติ. เอตฺตาวตาปิ โข อาวุโส อริยสาวโก สมฺมาทิฏฺฐิ โหติ, อุชุคตาสฺส ทิฏฺฐิ, ธมฺเม อเวจฺจ\-ปฺปสาเทน สมนฺนาคโต, อาคโต อิมํ สทฺธมฺมนฺ”ติ.}

\proseitem{90}{“สาธาวุโส”ติ โข เต ภิกฺขู อายสฺมโต สาริปุตฺตสฺส ภาสิตํ อภินนฺทิตฺวา อนุโมทิตฺวา อายสฺมนฺตํ สาริปุตฺตํ อุตฺตริ\csromansharedfootnote{40696c5953dcfaea}{อุตฺตริํ (สี, สฺยา, อิ)} ปญฺหํ อปุจฺฉุํ\csromansharedfootnote{05ba7053fbd095c0}{อปุจฺฉิํสุ (สฺยา)} “สิยา ปนาวุโส อญฺโญปิ ปริยาโย, ยถา อริยสาวโก สมฺมาทิฏฺฐิ โหติ, อุชุคตาสฺส ทิฏฺฐิ, ธมฺเม อเวจฺจ\-ปฺปสาเทน สมนฺนาคโต, อาคโต อิมํ สทฺธมฺมนฺ”ติ.}

\prose{สิยา อาวุโส. ยโต โข อาวุโส อริยสาวโก อาหารญฺจ ปชานาติ, อาหารสมุทยญฺจ ปชานาติ, อาหารนิโรธญฺจ ปชานาติ, อาหาร\-นิโรธ\-คามินิํ ปฏิปทญฺจ ปชานาติ. เอตฺตาวตาปิ โข อาวุโส อริยสาวโก สมฺมาทิฏฺฐิ โหติ, อุชุคตาสฺส ทิฏฺฐิ, ธมฺเม อเวจฺจ\-ปฺปสาเทน สมนฺนาคโต, อาคโต อิมํ สทฺธมฺมํ. กตโม ปนาวุโส อาหาโร, กตโม อาหารสมุทโย, กตโม อาหารนิโรโธ, กตมา อาหาร\-นิโรธ\-คามินี ปฏิปทา. จตฺตาโรเม อาวุโส อาหารา ภูตานํ วา สตฺตานํ ฐิติยา สมฺภเวสีนํ วา อนุคฺคหาย. กตเม จตฺตาโร, กพฬีกาโร อาหาโร \csromanfolio{59}โอฬาริโก วา สุขุโม วา, ผสฺโส ทุติโย, มโนสญฺเจตนา ตติยา, วิญฺญาณํ จตุตฺถํ. ตณฺหาสมุทยา อาหารสมุทโย, ตณฺหานิโรธา อาหารนิโรโธ. อยเมว อริโย อฏฺฐงฺคิโก มคฺโค อาหาร\-นิโรธ\-คามินี ปฏิปทา. เสยฺยถิทํ, สมฺมาทิฏฺฐิ สมฺมาสงฺกปฺโป สมฺมาวาจา สมฺมากมฺมนฺโต สมฺมาอาชีโว สมฺมาวายาโม สมฺมาสติ สมฺมาสมาธิ.}

\prose{ยโต โข อาวุโส อริยสาวโก เอวํ อาหารํ ปชานาติ, เอวํ อาหารสมุทยํ ปชานาติ, เอวํ อาหารนิโรธํ ปชานาติ, เอวํ อาหาร\-นิโรธ\-คามินิํ ปฏิปทํ ปชานาติ. โส สพฺพโส ราคานุสยํ ปหาย ปฏิฆานุสยํ ปฏิวิโนเทตฺวา อสฺมีติ ทิฏฺฐิ\-มานา\-นุสยํ สมูหนิตฺวา อวิชฺชํ ปหาย วิชฺชํ อุปฺปาเทตฺวา ทิฏฺเฐวธมฺเม ทุกฺขสฺส\-นฺตกโร โหติ. เอตฺตาวตาปิ โข อาวุโส อริยสาวโก สมฺมาทิฏฺฐิ โหติ, อุชุคตาสฺส ทิฏฺฐิ, ธมฺเม อเวจฺจ\-ปฺปสาเทน สมนฺนาคโต, อาคโต อิมํ สทฺธมฺมนฺติ.}

\proseitem{91}{“สาธาวุโส”ติ โข เต ภิกฺขู อายสฺมโต สาริปุตฺตสฺส ภาสิตํ อภินนฺทิตฺวา อนุโมทิตฺวา อายสฺมนฺตํ สาริปุตฺตํ อุตฺตริ ปญฺหํ อปุจฺฉุํ “สิยา ปนาวุโส อญฺโญปิ ปริยาโย, ยถา อริยสาวโก สมฺมาทิฏฺฐิ โหติ, อุชุคตาสฺส ทิฏฺฐิ, ธมฺเม อเวจฺจ\-ปฺปสาเทน สมนฺนาคโต, อาคโต อิมํ สทฺธมฺมนฺ”ติ.}

\prose{สิยา อาวุโส. ยโต โข อาวุโส อริยสาวโก ทุกฺขญฺจ ปชานาติ, ทุกฺขสมุทยญฺจ ปชานาติ, ทุกฺขนิโรธญฺจ ปชานาติ, ทุกฺขนิโรธ\-คามินิํ ปฏิปทญฺจ ปชานาติ. เอตฺตาวตาปิ โข อาวุโส อริยสาวโก สมฺมาทิฏฺฐิ โหติ, อุชุคตาสฺส ทิฏฺฐิ, ธมฺเม อเวจฺจ\-ปฺปสาเทน สมนฺนาคโต, อาคโต อิมํ สทฺธมฺมํ. กตมํ ปนาวุโส ทุกฺขํ, กตโม ทุกฺขสมุทโย, กตโม ทุกฺขนิโรโธ, กตมา ทุกฺขนิโรธ\-คามินี ปฏิปทา. ชาติปิ ทุกฺขา, ชราปิ ทุกฺขา, มรณมฺปิ ทุกฺขํ, โสก\-ปริเทว\-ทุกฺข\-โทมนสฺสุ\-ปายาสาปิ ทุกฺขา, อปฺปิเยหิ สมฺปโยโคปิ ทุกฺโข, ปิเยหิ วิปฺปโยโคปิ ทุกฺโข, ยมฺปิจฺฉํ น ลภติ ตมฺปิ ทุกฺขํ, สํขิตฺเตน ปญฺจุ\-ปาทาน\-กฺขนฺธา\csromansharedfootnote{feb51bbb692b3682}{ปญฺจุปาทานกฺขนฺธาปิ (ก)} ทุกฺขา. อิทํ วุจฺจตาวุโส ทุกฺขํ. กตโม จาวุโส ทุกฺขสมุทโย. ยายํ \csromanfolio{60}ตณฺหา โปโนพฺภวิกา\csromansharedfootnote{3459575e13e816f9}{โปโนภวิกา (สี, อิ)} นนฺที\-ราค\-สหคตา\csromansharedfootnote{6cb114d7ce9a643a}{นนฺทิราคสหคตา (สี, อิ)} ตตฺร\-ตตฺรา\-ภินนฺทินี. เสยฺยถิทํ, กามตณฺหา ภวตณฺหา วิภวตณฺหา. อยํ วุจฺจตาวุโส ทุกฺขสมุทโย. กตโม จาวุโส ทุกฺขนิโรโธ. โย ตสฺสาเยว ตณฺหาย อเสส\-วิราค\-นิโรโธ จาโค ปฏินิสฺสคฺโค มุตฺติ อนาลโย. อยํ วุจฺจตาวุโส ทุกฺขนิโรโธ. กตมา จาวุโส ทุกฺขนิโรธ\-คามินี ปฏิปทา. อยเมว อริโย อฏฺฐงฺคิโก มคฺโค. เสยฺยถิทํ, สมฺมาทิฏฺฐิ ฯเปฯ สมฺมาสมาธิ. อยํ วุจฺจตาวุโส ทุกฺขนิโรธ\-คามินี ปฏิปทา.}

\prose{ยโต โข อาวุโส อริยสาวโก เอวํ ทุกฺขํ ปชานาติ, เอวํ ทุกฺข\-สมุทยํ ปชานาติ, เอวํ ทุกฺขนิโรธํ ปชานาติ, เอวํ ทุกฺขนิโรธ\-คามินิํ ปฏิปทํ ปชานาติ. โส สพฺพโส ราคานุสยํ ปหาย ปฏิฆานุสยํ ปฏิวิโนเทตฺวา อสฺมีติ ทิฏฺฐิ\-มานา\-นุสยํ สมูหนิตฺวา อวิชฺชํ ปหาย วิชฺชํ อุปฺปาเทตฺวา ทิฏฺเฐวธมฺเม ทุกฺขสฺส\-นฺตกโร โหติ. เอตฺตาวตาปิ โข อาวุโส อริยสาวโก สมฺมาทิฏฺฐิ โหติ, อุชุคตาสฺส ทิฏฺฐิ, ธมฺเม อเวจฺจ\-ปฺปสาเทน สมนฺนาคโต, อาคโต อิมํ สทฺธมฺมนฺติ.}

\proseitem{92}{“สาธาวุโส”ติ โข เต ภิกฺขู อายสฺมโต สาริปุตฺตสฺส ภาสิตํ อภินนฺทิตฺวา อนุโมทิตฺวา อายสฺมนฺตํ สาริปุตฺตํ อุตฺตริ ปญฺหํ อปุจฺฉุํ “สิยา ปนาวุโส อญฺโญปิ ปริยาโย, ยถา อริยสาวโก สมฺมาทิฏฺฐิ โหติ, อุชุคตาสฺส ทิฏฺฐิ, ธมฺเม อเวจฺจ\-ปฺปสาเทน สมนฺนาคโต, อาคโต อิมํ สทฺธมฺมนฺ”ติ.}

\prose{สิยา อาวุโส. ยโต โข อาวุโส อริยสาวโก ชรามรณญฺจ ปชานาติ, ชรามรณสมุทยญฺจ ปชานาติ, ชรามรณนิโรธญฺจ ปชานาติ, ชรามรณนิโรธ\-คามินิํ ปฏิปทญฺจ ปชานาติ. เอตฺตาวตาปิ โข อาวุโส อริยสาวโก สมฺมาทิฏฺฐิ โหติ, อุชุคตาสฺส ทิฏฺฐิ, ธมฺเม อเวจฺจ\-ปฺปสาเทน สมนฺนาคโต, อาคโต อิมํ สทฺธมฺมํ. กตมํ ปนาวุโส ชรามรณํ, กตโม ชรามรณ\-สมุทโย, กตโม ชรามรณ\-นิโรโธ, กตมา ชรามรณนิโรธ\-คามินี ปฏิปทา. ยา เตสํ เตสํ สตฺตานํ ตมฺหิ ตมฺหิ สตฺตนิกาเย ชรา ชีรณตา ขณฺฑิจฺจํ ปาลิจฺจํ วลิตฺตจตา อายุโน สํหานิ อินฺทฺริยานํ ปริปาโก. อยํ วุจฺจตาวุโส ชรา. กตมญฺจาวุโส มรณํ. ยา\csromansharedfootnote{82cf5e8b312ff21e}{ยํ (อิ, ก), สติปฏฺฐานสุตฺเตปิ,} เตสํ เตสํ \csromanfolio{61}สตฺตานํ ตมฺหา ตมฺหา สตฺตนิกายา จุติ จวนตา เภโท อนฺตรธานํ มจฺจุ มรณํ กาลํกิริยา ขนฺธานํ เภโท กเฬวรสฺส นิกฺเขโป ชีวิตินฺทฺริยสฺสุ\-ปจฺเฉโท. อิทํ วุจฺจตาวุโส มรณํ. อิติ อยญฺจ ชรา อิทญฺจ มรณํ. อิทํ วุจฺจตาวุโส ชรามรณํ. ชาติสมุทยา ชรามรณ\-สมุทโย, ชาตินิโรธา ชรามรณ\-นิโรโธ. อยเมว อริโย อฏฺฐงฺคิโก มคฺโค ชรามรณนิโรธ\-คามินี ปฏิปทา. เสยฺยถิทํ, สมฺมาทิฏฺฐิ ฯเปฯ สมฺมาสมาธิ.}

\prose{ยโต โข อาวุโส อริยสาวโก เอวํ ชรามรณํ ปชานาติ, เอวํ ชรามรณ\-สมุทยํ ปชานาติ, เอวํ ชรามรณ\-นิโรธํ ปชานาติ, เอวํ ชรามรณนิโรธ\-คามินิํ ปฏิปทํ ปชานาติ. โส สพฺพโส ราคานุสยํ ปหาย ฯเปฯ ทุกฺขสฺส\-นฺตกโร โหติ. เอตฺตาวตาปิ โข อาวุโส อริยสาวโก สมฺมาทิฏฺฐิ โหติ, อุชุคตาสฺส ทิฏฺฐิ, ธมฺเม อเวจฺจ\-ปฺปสาเทน สมนฺนาคโต, อาคโต อิมํ สทฺธมฺมนฺติ.}

\proseitem{93}{“สาธาวุโส”ติ โข ฯเปฯ อปุจฺฉุํ\csromandash{}สิยา ปนาวุโส ฯเปฯ. สิยา อาวุโส. ยโต โข อาวุโส อริยสาวโก ชาติญฺจ ปชานาติ, ชาติสมุทยญฺจ ปชานาติ, ชาตินิโรธญฺจ ปชานาติ, ชาตินิโรธ\-คามินิํ ปฏิปทญฺจ ปชานาติ. เอตฺตาวตาปิ โข อาวุโส อริยสาวโก สมฺมาทิฏฺฐิ โหติ, อุชุคตาสฺส ทิฏฺฐิ, ธมฺเม อเวจฺจ\-ปฺปสาเทน สมนฺนาคโต, อาคโต อิมํ สทฺธมฺมํ. กตมา ปนาวุโส ชาติ, กตโม ชาติสมุทโย, กตโม ชาตินิโรโธ, กตมา ชาตินิโรธ\-คามินี ปฏิปทา. ยา เตสํ เตสํ สตฺตานํ ตมฺหิ ตมฺหิ สตฺตนิกาเย ชาติ สญฺชาติ โอกฺกนฺติ อภินิพฺพตฺติ ขนฺธานํ ปาตุภาโว อายตนานํ ปฏิลาโภ. อยํ วุจฺจตาวุโส ชาติ. ภวสมุทยา ชาติสมุทโย, ภวนิโรธา ชาตินิโรโธ. อยเมว อริโย อฏฺฐงฺคิโก มคฺโค ชาตินิโรธ\-คามินี ปฏิปทา. เสยฺยถิทํ, สมฺมาทิฏฺฐิ ฯเปฯ สมฺมาสมาธิ.}

\prose{ยโต โข อาวุโส อริยสาวโก เอวํ ชาติํ ปชานาติ, เอวํ ชาติสมุทยํ ปชานาติ, เอวํ ชาตินิโรธํ ปชานาติ, เอวํ ชาตินิโรธ\-คามินิํ ปฏิปทํ ปชานาติ. โส สพฺพโส ราคานุสยํ ปหาย ฯเปฯ ทุกฺขสฺส\-นฺตกโร โหติ. เอตฺตาวตาปิ โข อาวุโส อริยสาวโก \csromanfolio{62}สมฺมาทิฏฺฐิ โหติ, อุชุคตาสฺส ทิฏฺฐิ, ธมฺเม อเวจฺจ\-ปฺปสาเทน สมนฺนาคโต, อาคโต อิมํ สทฺธมฺมนฺติ.}

\proseitem{94}{“สาธาวุโส”ติ โข ฯเปฯ อปุจฺฉุํ\csromandash{}สิยา ปนาวุโส ฯเปฯ. สิยา อาวุโส. ยโต โข อาวุโส อริยสาวโก ภวญฺจ ปชานาติ, ภวสมุทยญฺจ ปชานาติ, ภวนิโรธญฺจ ปชานาติ, ภวนิโรธ\-คามินิํ ปฏิปทญฺจ ปชานาติ. เอตฺตาวตาปิ โข อาวุโส อริยสาวโก สมฺมาทิฏฺฐิ โหติ, อุชุคตาสฺส ทิฏฺฐิ, ธมฺเม อเวจฺจ\-ปฺปสาเทน สมนฺนาคโต, อาคโต อิมํ สทฺธมฺมํ. กตโม ปนาวุโส ภโว, กตโม ภวสมุทโย, กตโม ภวนิโรโธ, กตมา ภวนิโรธ\-คามินี ปฏิปทา. ตโยเม อาวุโส ภวา กามภโว รูปภโว อรูปภโว. อุปาทาน\-สมุทยา ภวสมุทโย, อุปาทานนิโรธา ภวนิโรโธ. อยเมว อริโย อฏฺฐงฺคิโก มคฺโค ภวนิโรธ\-คามินี ปฏิปทา. เสยฺยถิทํ, สมฺมาทิฏฺฐิ ฯเปฯ สมฺมาสมาธิ.}

\prose{ยโต โข อาวุโส อริยสาวโก เอวํ ภวํ ปชานาติ, เอวํ ภวสมุทยํ ปชานาติ, เอวํ ภวนิโรธํ ปชานาติ, เอวํ ภวนิโรธ\-คามินิํ ปฏิปทํ ปชานาติ. โส สพฺพโส ราคานุสยํ ปหาย ฯเปฯ ทุกฺขสฺส\-นฺตกโร โหติ. เอตฺตาวตาปิ โข อาวุโส อริยสาวโก สมฺมาทิฏฺฐิ โหติ, อุชุคตาสฺส ทิฏฺฐิ, ธมฺเม อเวจฺจ\-ปฺปสาเทน สมนฺนาคโต, อาคโต อิมํ สทฺธมฺมนฺติ.}

\proseitem{95}{“สาธาวุโส”ติ โข ฯเปฯ อปุจฺฉุํ\csromandash{}สิยา ปนาวุโส ฯเปฯ. สิยา อาวุโส. ยโต โข อาวุโส อริยสาวโก อุปาทานญฺจ ปชานาติ, อุปาทานสมุทยญฺจ ปชานาติ, อุปาทานนิโรธญฺจ ปชานาติ,อุปาทานนิโรธ\-คามินิํ ปฏิปทญฺจ ปชานาติ. เอตฺตาวตาปิ โข อาวุโส อริยสาวโก สมฺมาทิฏฺฐิ โหติ, อุชุคตาสฺส ทิฏฺฐิ, ธมฺเม อเวจฺจ\-ปฺปสาเทน สมนฺนาคโต, อาคโต อิมํ สทฺธมฺมํ. กตมํ ปนาวุโส อุปาทานํ, กตโม อุปาทาน\-สมุทโย, กตโม อุปาทานนิโรโธ, กตมา อุปาทานนิโรธ\-คามินี ปฏิปทา. จตฺตาริมานิ อาวุโส อุปาทานานิ กามุปาทานํ ทิฏฺฐุปาทานํ สีลพฺพตุ\-ปาทานํ อตฺตวาทุ\-ปาทานํ. ตณฺหาสมุทยา อุปาทาน\-สมุทโย, ตณฺหานิโรธา อุปาทานนิโรโธ. อยเมว \csromanfolio{63}อริโย อฏฺฐงฺคิโก มคฺโค อุปาทานนิโรธ\-คามินี ปฏิปทา. เสยฺยถิทํ, สมฺมาทิฏฺฐิ ฯเปฯ สมฺมาสมาธิ.}

\prose{ยโต โข อาวุโส อริยสาวโก เอวํ อุปาทานํ ปชานาติ, เอวํ อุปาทาน\-สมุทยํ ปชานาติ, เอวํ อุปาทาน\-นิโรธํ ปชานาติ, เอวํ อุปาทานนิโรธ\-คามินิํ ปฏิปทํ ปชานาติ. โส สพฺพโส ราคานุสยํ ปหาย ฯเปฯ ทุกฺขสฺส\-นฺตกโร โหติ. เอตฺตาวตาปิ โข อาวุโส อริยสาวโก สมฺมาทิฏฺฐิ โหติ, อุชุคตาสฺส ทิฏฺฐิ, ธมฺเม อเวจฺจ\-ปฺปสาเทน สมนฺนาคโต, อาคโต อิมํ สทฺธมฺมนฺติ.}

\proseitem{96}{“สาธาวุโส”ติ โข ฯเปฯ อปุจฺฉุํ\csromandash{}สิยา ปนาวุโส ฯเปฯ. สิยา อาวุโส. ยโต โข อาวุโส อริยสาวโก ตณฺหญฺจ ปชานาติ, ตณฺหาสมุทยญฺจ ปชานาติ, ตณฺหานิโรธญฺจ ปชานาติ, ตณฺหานิโรธ\-คามินิํ ปฏิปทญฺจ ปชานาติ. เอตฺตาวตาปิ โข อาวุโส อริยสาวโก สมฺมาทิฏฺฐิ โหติ, อุชุคตาสฺส ทิฏฺฐิ, ธมฺเม อเวจฺจ\-ปฺปสาเทน สมนฺนาคโต, อาคโต อิมํ สทฺธมฺมํ. กตมา ปนาวุโส ตณฺหา, กตโม ตณฺหาสมุทโย, กตโม ตณฺหานิโรโธ, กตมา ตณฺหานิโรธ\-คามินี ปฏิปทา. ฉยิเม อาวุโส ตณฺหากายา, รูปตณฺหา สทฺทตณฺหา คนฺธตณฺหา รสตณฺหา โผฏฺฐพฺพ\-ตณฺหา ธมฺมตณฺหา. เวทนาสมุทยา ตณฺหาสมุทโย, เวทนานิโรธา ตณฺหานิโรโธ. อยเมว อริโย อฏฺฐงฺคิโก มคฺโค ตณฺหานิโรธ\-คามินี ปฏิปทา. เสยฺยถิทํ, สมฺมาทิฏฺฐิ ฯเปฯ สมฺมาสมาธิ.}

\prose{ยโต โข อาวุโส อริยสาวโก เอวํ ตณฺหํ ปชานาติ, เอวํ ตณฺหาสมุทยํ ปชานาติ, เอวํ ตณฺหานิโรธํ ปชานาติ, เอวํ ตณฺหานิโรธ\-คามินิํ ปฏิปทํ ปชานาติ. โส สพฺพโส ราคานุสยํ ปหาย ฯเปฯ ทุกฺขสฺส\-นฺตกโร โหติ. เอตฺตาวตาปิ โข อาวุโส อริยสาวโก สมฺมาทิฏฺฐิ โหติ, อุชุคตาสฺส ทิฏฺฐิ, ธมฺเม อเวจฺจ\-ปฺปสาเทน สมนฺนาคโต, อาคโต อิมํ สทฺธมฺมนฺติ.}

\proseitem{97}{“สาธาวุโส”ติ โข ฯเปฯ อปุจฺฉุํ\csromandash{}สิยา ปนาวุโส ฯเปฯ. สิยา อาวุโส. ยโต โข อาวุโส อริยสาวโก เวทนญฺจ ปชานาติ, เวทนาสมุทยญฺจ ปชานาติ, เวทนานิโรธญฺจ ปชานาติ, เวทนา\-นิโรธ\-คามินิํ ปฏิปทญฺจ ปชานาติ. เอตฺตาวตาปิ โข อาวุโส \csromanfolio{64}อริยสาวโก สมฺมาทิฏฺฐิ โหติ, อุชุคตาสฺส ทิฏฺฐิ, ธมฺเม อเวจฺจ\-ปฺปสาเทน สมนฺนาคโต, อาคโต อิมํ สทฺธมฺมํ. กตมา ปนาวุโส เวทนา, กตโม เวทนาสมุทโย, กตโม เวทนานิโรโธ, กตมา เวทนา\-นิโรธ\-คามินี ปฏิปทา. ฉยิเม อาวุโส เวทนากายา, จกฺขุ\-สมฺผสฺสชา เวทนา โสต\-สมฺผสฺสชา เวทนา ฆาน\-สมฺผสฺสชา เวทนา ชิวฺหา\-สมฺผสฺสชา เวทนา กาย\-สมฺผสฺสชา เวทนา มโน\-สมฺผสฺสชา เวทนา. ผสฺสสมุทยา เวทนาสมุทโย, ผสฺสนิโรธา เวทนานิโรโธ. อยเมว อริโย อฏฺฐงฺคิโก มคฺโค เวทนา\-นิโรธ\-คามินี ปฏิปทา. เสยฺยถิทํ, สมฺมาทิฏฺฐิ ฯเปฯ สมฺมาสมาธิ.}

\prose{ยโต โข อาวุโส อริยสาวโก เอวํ เวทนํ ปชานาติ,เอวํ เวทนา\-สมุทยํ ปชานาติ, เอวํ เวทนานิโรธํ ปชานาติ, เอวํ เวทนา\-นิโรธ\-คามินิํ ปฏิปทํ ปชานาติ. โส สพฺพโส ราคานุสยํ ปหาย ฯเปฯ ทุกฺขสฺส\-นฺตกโร โหติ. เอตฺตาวตาปิ โข อาวุโส อริยสาวโก สมฺมาทิฏฺฐิ โหติ, อุชุคตาสฺส ทิฏฺฐิ, ธมฺเม อเวจฺจ\-ปฺปสาเทน สมนฺนาคโต, อาคโต อิมํ สทฺธมฺมนฺติ.}

\proseitem{98}{“สาธาวุโส”ติ โข ฯเปฯ อปุจฺฉุํ\csromandash{}สิยา ปนาวุโส ฯเปฯ. สิยา อาวุโส. ยโต โข อาวุโส อริยสาวโก ผสฺสญฺจ ปชานาติ, ผสฺสสมุทยญฺจ ปชานาติ, ผสฺสนิโรธญฺจ ปชานาติ, ผสฺสนิโรธ\-คามินิํ ปฏิปทญฺจ ปชานาติ. เอตฺตาวตาปิ โข อาวุโส อริยสาวโก สมฺมาทิฏฺฐิ โหติ, อุชุคตาสฺส ทิฏฺฐิ, ธมฺเม อเวจฺจ\-ปฺปสาเทน สมนฺนาคโต, อาคโต อิมํ สทฺธมฺมํ. กตโม ปนาวุโส ผสฺโส, กตโม ผสฺสสมุทโย, กตโม ผสฺสนิโรโธ, กตมา ผสฺสนิโรธ\-คามินี ปฏิปทา. ฉยิเม อาวุโส ผสฺสกายา, จกฺขุ\-สมฺผสฺโส โสตสมฺผสฺโส ฆานสมฺผสฺโส ชิวฺหาสมฺผสฺโส กายสมฺผสฺโส มโนสมฺผสฺโส. สฬายตน\-สมุทยา ผสฺสสมุทโย, สฬายตน\-นิโรธา ผสฺสนิโรโธ. อยเมว อริโย อฏฺฐงฺคิโก มคฺโค ผสฺสนิโรธ\-คามินี ปฏิปทา. เสยฺยถิทํ, สมฺมาทิฏฺฐิ ฯเปฯ สมฺมาสมาธิ.}

\prose{ยโต โข อาวุโส อริยสาวโก เอวํ ผสฺสํ ปชานาติ, เอวํ ผสฺส\-สมุทยํ ปชานาติ, เอวํ ผสฺสนิโรธํ ปชานาติ, เอวํ ผสฺสนิโรธ\-คามินิํ ปฏิปทํ ปชานาติ. โส สพฺพโส ราคานุสยํ ปหาย ฯเปฯ \csromanfolio{65}ทุกฺขสฺส\-นฺตกโร โหติ. เอตฺตาวตาปิ โข อาวุโส อริยสาวโก สมฺมาทิฏฺฐิ โหติ, อุชุคตาสฺส ทิฏฺฐิ, ธมฺเม อเวจฺจ\-ปฺปสาเทน สมนฺนาคโต, อาคโต อิมํ สทฺธมฺมนฺติ.}

\proseitem{99}{“สาธาวุโส”ติ โข ฯเปฯ อปุจฺฉุํ\csromandash{}สิยา ปนาวุโส ฯเปฯ. สิยา อาวุโส. ยโต โข อาวุโส อริยสาวโก สฬายตนญฺจ ปชานาติ, สฬายตนสมุทยญฺจ ปชานาติ, สฬายตนนิโรธญฺจ ปชานาติ, สฬายตนนิโรธ\-คามินิํ ปฏิปทญฺจ ปชานาติ. เอตฺตาวตาปิ โข อาวุโส อริยสาวโก สมฺมาทิฏฺฐิ โหติ, อุชุคตาสฺส ทิฏฺฐิ, ธมฺเม อเวจฺจ\-ปฺปสาเทน สมนฺนาคโต, อาคโต อิมํ สทฺธมฺมํ. กตมํ ปนาวุโส สฬายตนํ, กตโม สฬายตน\-สมุทโย, กตโม สฬายตน\-นิโรโธ, กตมา สฬายตนนิโรธ\-คามินี ปฏิปทา. ฉยิมานิ อาวุโส อายตนานิ, จกฺขายตนํ โสตายตนํ ฆานายตนํ ชิวฺหายตนํ กายายตนํ มนายตนํ. นามรูป\-สมุทยา สฬายตน\-สมุทโย, นามรูป\-นิโรธา สฬายตน\-นิโรโธ. อยเมว อริโย อฏฺฐงฺคิโก มคฺโค สฬายตนนิโรธ\-คามินี ปฏิปทา. เสยฺยถิทํ, สมฺมาทิฏฺฐิ ฯเปฯ สมฺมาสมาธิ.}

\prose{ยโต โข อาวุโส อริยสาวโก เอวํ สฬายตนํ ปชานาติ, เอวํ สฬายตน\-สมุทยํ ปชานาติ, เอวํ สฬายตน\-นิโรธํ ปชานาติ, เอวํ สฬายาตนนิโรธคามินิํ ปฏิปทํ ปชานาติ. โส สพฺพโส ราคานุสยํ ปหาย ฯเปฯ ทุกฺขสฺส\-นฺตกโร โหติ. เอตฺตาวตาปิ โข อาวุโส อริยสาวโก สมฺมาทิฏฺฐิ โหติ, อุชุคตาสฺส ทิฏฺฐิ, ธมฺเม อเวจฺจ\-ปฺปสาเทน สมนฺนาคโต, อาคโต อิมํ สทฺธมฺมนฺติ.}

\proseitem{100}{“สาธาวุโส”ติ โข ฯเปฯ อปุจฺฉุํ\csromandash{}สิยา ปนาวุโส ฯเปฯ. สิยา อาวุโส. ยโต โข อาวุโส อริยสาวโก นามรูปญฺจ ปชานาติ, นามรูปสมุทยญฺจ ปชานาติ, นามรูปนิโรธญฺจ ปชานาติ, นามรูปนิโรธ\-คามินิํ ปฏิปทญฺจ ปชานาติ. เอตฺตาวตาปิ โข อาวุโส อริยสาวโก สมฺมาทิฏฺฐิ โหติ, อุชุคตาสฺส ทิฏฺฐิ, ธมฺเม อเวจฺจ\-ปฺปสาเทน สมนฺนาคโต, อาคโต อิมํ สทฺธมฺมํ. กตมํ ปนาวุโส นามรูปํ, กตโม นามรูป\-สมุทโย, กตโม นามรูป\-นิโรโธ, กตมา นามรูปนิโรธ\-คามินี ปฏิปทา. เวทนา สญฺญา เจตนา ผสฺโส \csromanfolio{66}มนสิกาโร, อิทํ วุจฺจตาวุโส นามํ. จตฺตาริ จ มหาภูตานิ จตุนฺนญฺจ มหาภูตานํ อุปาทายรูปํ, อิทํ วุจฺจตาวุโส รูปํ. อิติ อิทญฺจ นามํ อิทญฺจ รูปํ, อิทํ วุจฺจตาวุโส นามรูปํ. วิญฺญาณ\-สมุทยา นามรูป\-สมุทโย, วิญฺญาณนิโรธา นามรูป\-นิโรโธ. อยเมว อริโย อฏฺฐงฺคิโก มคฺโค นามรูปนิโรธ\-คามินี ปฏิปทา. เสยฺยถิทํ, สมฺมาทิฏฺฐิ ฯเปฯ สมฺมาสมาธิ.}

\prose{ยโต โข อาวุโส อริยสาวโก เอวํ นามรูปํ ปชานาติ, เอวํ นามรูป\-สมุทยํ ปชานาติ, เอวํ นามรูป\-นิโรธํ ปชานาติ, เอวํ นามรูปนิโรธ\-คามินิํ ปฏิปทํ ปชานาติ. โส สพฺพโส ราคานุสยํ ปหาย ฯเปฯ ทุกฺขสฺส\-นฺตกโร โหติ. เอตฺตาวตาปิ โข อาวุโส อริยสาวโก สมฺมาทิฏฺฐิ โหติ, อุชุคตาสฺส ทิฏฺฐิ, ธมฺเม อเวจฺจ\-ปฺปสาเทน สมนฺนาคโต, อาคโต อิมํ สทฺธมฺมนฺติ.}

\proseitem{101}{“สาธาวุโส”ติ โข ฯเปฯ อปุจฺฉุํ \csromandash{} สิยา ปนาวุโส ฯเปฯ. สิยา อาวุโส. ยโต โข อาวุโส อริยสาวโก วิญฺญาณญฺจ ปชานาติ, วิญฺญาณสมุทยญฺจ ปชานาติ, วิญฺญาณนิโรธญฺจ ปชานาติ, วิญฺญาณ\-นิโรธ\-คามินิํ ปฏิปทญฺจ ปชานาติ. เอตฺตาวตาปิ โข อาวุโส อริยสาวโก สมฺมาทิฏฺฐิ โหติ, อุชุคตาสฺส ทิฏฺฐิ, ธมฺเม อเวจฺจ\-ปฺปสาเทน สมนฺนาคโต, อาคโต อิมํ สทฺธมฺมํ. กตมํ ปนาวุโส วิญฺญาณํ, กตโม วิญฺญาณ\-สมุทโย, กตโม วิญฺญณนิโรโธ, กตมา วิญฺญาณ\-นิโรธ\-คามินี ปฏิปทา. ฉยิเม อาวุโส วิญฺญาณกายา, จกฺขุวิญฺญาณํ โสตวิญฺญาณํ ฆานวิญฺญาณํ ชิวฺหาวิญฺญาณํ กายวิญฺญาณํ มโนวิญฺญาณํ. สงฺขาร\-สมุทยา วิญฺญาณ\-สมุทโย, สงฺขาร\-นิโรธา วิญฺญาณนิโรโธ. อยเมว อริโย อฏฺฐงฺคิโก มคฺโค วิญฺญาณ\-นิโรธ\-คามินี ปฏิปทา. เสยฺยถิทํ, สมฺมาทิฏฺฐิ ฯเปฯ สมฺมาสมาธิ.}

\prose{ยโต โข อาวุโส อริยสาวโก เอวํ วิญฺญาณํ ปชานาติ, เอวํ วิญฺญาณ\-สมุทยํ ปชานาติ, เอวํ วิญฺญาณ\-นิโรธํ ปชานาติ, เอวํ วิญฺญาณ\-นิโรธ\-คามินิํ ปฏิปทํ ปชานาติ. โส สพฺพโส ราคานุสยํ ปหาย ฯเปฯ ทุกฺขสฺส\-นฺตกโร โหติ. เอตฺตาวตาปิ โข อาวุโส อริยสาวโก สมฺมาทิฏฺฐิ โหติ, อุชุคตาสฺส ทิฏฺฐิ, ธมฺเม อเวจฺจ\-ปฺปสาเทน สมนฺนาคโต, อาคโต อิมํ สทฺธมฺมนฺติ.}

\proseitem{102}{\csromanfolio{67}“สาธาวุโส”ติ โข ฯเปฯ อปุจฺฉุํ\csromandash{}สิยา ปนาวุโส ฯเปฯ. สิยา อาวุโส. ยโต โข อาวุโส อริยสาวโก สงฺขาเร จ ปชานาติ, สงฺขารสมุทยญฺจ ปชานาติ, สงฺขาร\-นิโรธญฺจ ปชานาติ, สงฺขาร\-นิโรธ\-คามินิํ ปฏิปทญฺจ ปชานาติ. เอตฺตาวตาปิ โข อาวุโส อริยสาวโก สมฺมาทิฏฺฐิ โหติ, อุชุคตาสฺส ทิฏฺฐิ, ธมฺเม อเวจฺจ\-ปฺปสาเทน สมนฺนาคโต, อาคโต อิมํ สทฺธมฺมํ. กตเม ปนาวุโส สงฺขารา, กตโม สงฺขาร\-สมุทโย, กตโม สงฺขาร\-นิโรโธ, กตมา สงฺขาร\-นิโรธ\-คามินี ปฏิปทา. ตโยเม อาวุโส สงฺขารา, กายสงฺขาโร วจีสงฺขาโร จิตฺตสงฺขาโร. อวิชฺชาสมุทยา สงฺขาร\-สมุทโย, อวิชฺชานิโรธา สงฺขาร\-นิโรโธ. อยเมว อริโย อฏฺฐงฺคิโก มคฺโค สงฺขาร\-นิโรธ\-คามินี ปฏิปทา. เสยฺยถิทํ, สมฺมาทิฏฺฐิ ฯเปฯ สมฺมาสมาธิ.}

\prose{ยโต โข อาวุโส อริยสาวโก เอวํ สงฺขาเร ปชานาติ, เอวํ สงฺขาร\-สมุทยํ ปชานาติ, เอวํ สงฺขาร\-นิโรธํ ปชานาติ, เอวํ สงฺขาร\-นิโรธ\-คามินิํ ปฏิปทํ ปชานาติ. โส สพฺพโส ราคานุสยํ ปหาย ปฏิฆานุสยํ ปฏิวิโนเทตฺวา อสฺมีติ ทิฏฺฐิ\-มานา\-นุสยํ สมูหนิตฺวา อวิชฺชํ ปหาย วิชฺชํ อุปฺปาเทตฺวา ทิฏฺเฐว ธมฺเม ทุกฺขสฺส\-นฺตกโร โหติ. เอตฺตาวตาปิ โข อาวุโส อริยสาวโก สมฺมาทิฏฺฐิ โหติ, อุชุคตาสฺส ทิฏฺฐิ, ธมฺเม อเวจฺจ\-ปฺปสาเทน สมนฺนาคโต, อาคโต อิมํ สทฺธมฺมนฺติ.}

\proseitem{103}{“สาธาวุโส”ติ โข ฯเปฯ อปุจฺฉุํ\csromandash{}สิยา ปนาวุโส ฯเปฯ. สิยา อาวุโส. ยโต โข อาวุโส อริยสาวโก อวิชฺชญฺจ ปชานาติ, อวิชฺชาสมุทยญฺจ ปชานาติ, อวิชฺชานิโรธญฺจ ปชานาติ, อวิชฺชานิโรธ\-คามินิํ ปฏิปทญฺจ ปชานาติ. เอตฺตาวตาปิ โข อาวุโส อริยสาวโก สมฺมาทิฏฺฐิ โหติ, อุชุคตาสฺส ทิฏฺฐิ, ธมฺเม อเวจฺจ\-ปฺปสาเทน สมนฺนาคโต, อาคโต อิมํ สทฺธมฺมํ. กตมา ปนาวุโส อวิชฺชา, กตโม อวิชฺชาสมุทโย, กตโม อวิชฺชานิโรโธ, กตมา อวิชฺชานิโรธ\-คามินี ปฏิปทา. ยํ โข อาวุโส ทุกฺเข อญฺญาณํ, ทุกฺขสมุทเย อญฺญาณํ, ทุกฺขนิโรเธ อญฺญาณํ, ทุกฺขนิโรธ\-คามินิยา ปฏิปทาย อญฺญาณํ. อยํ วุจฺจตาวุโส อวิชฺชา. อาสวสมุทยา อวิชฺชาสมุทโย, อาสวนิโรธา อวิชฺชานิโรโธ. อยเมว อริโย อฏฺฐงฺคิโก มคฺโค อวิชฺชานิโรธ\-คามินี ปฏิปทา. เสยฺยถิทํ, สมฺมาทิฏฺฐิ ฯเปฯ สมฺมาสมาธิ.}

\prose{\csromanfolio{68}ยโต โข อาวุโส อริยสาวโก เอวํ อวิชฺชํ ปชานาติ, เอวํ อวิชฺชา\-สมุทยํ ปชานาติ, เอวํ อวิชฺชานิโรธํ ปชานาติ, เอวํ อวิชฺชานิโรธ\-คามินิํ ปฏิปทํ ปชานาติ. โส สพฺพโส ราคานุสยํ ปหาย ปฏิฆานุสยํ ปฏิวิโนเทตฺวา อสฺมีติ ทิฏฺฐิ\-มานา\-นุสยํ สมูหนิตฺวา อวิชฺชํ ปหาย วิชฺชํ อุปฺปาเทตฺวา ทิฏฺเฐว ธมฺเม ทุกฺขสฺส\-นฺตกโร โหติ. เอตฺตาวตาปิ โข อาวุโส อริยสาวโก สมฺมาทิฏฺฐิ โหติ, อุชุคตาสฺส ทิฏฺฐิ, ธมฺเม อเวจฺจ\-ปฺปสาเทน สมนฺนาคโต, อาคโต อิมํ สทฺธมฺมนฺติ.}

\proseitem{104}{“สาธาวุโส”ติ โข เต ภิกฺขู อายสฺมโต สาริปุตฺตสฺส ภาสิตํ อภินนฺทิตฺวา อนุโมทิตฺวา อายสฺมนฺตํ สาริปุตฺตํ อุตฺตริ ปญฺหํ อปุจฺฉุํ “สิยา ปนาวุโส อญฺโญปิ ปริยาโย, ยถา อริยสาวโก สมฺมาทิฏฺฐิ โหติ, อุชุคตาสฺส ทิฏฺฐิ, ธมฺเม อเวจฺจ\-ปฺปสาเทน สมนฺนาคโต, อาคโต อิมํ สทฺธมฺมนฺ”ติ.}

\prose{สิยา อาวุโส. ยโต โข อาวุโส อริยสาวโก อาสวญฺจ ปชานาติ, อาสวสมุทยญฺจ ปชานาติ, อาสวนิโรธญฺจ ปชานาติ, อาสว\-นิโรธ\-คามินิํ ปฏิปทญฺจ ปชานาติ. เอตฺตาวตาปิ โข อาวุโส อริยสาวโก สมฺมาทิฏฺฐิ โหติ, อุชุคตาสฺส ทิฏฺฐิ, ธมฺเม อเวจฺจ\-ปฺปสาเทน สมนฺนาคโต, อาคโต อิมํ สทฺธมฺมํ. กตโม ปนาวุโส อาสโว, กตโม อาสวสมุทโย, กตโม อาสวนิโรโธ, กตมา อาสว\-นิโรธ\-คามินี ปฏิปทาติ. ตโยเม อาวุโส อาสวา, กามาสโว ภวาสโว อวิชฺชาสโว. อวิชฺชาสมุทยา อาสวสมุทโย, อวิชฺชานิโรธา อาสวนิโรโธ. อยเมว อริโย อฏฺฐงฺคิโก มคฺโค อาสว\-นิโรธ\-คามินี ปฏิปทา. เสยฺยถิทํ, สมฺมาทิฏฺฐิ ฯเปฯ สมฺมาสมาธิ.}

\prose{ยโต โข อาวุโส อริยสาวโก เอวํ อาสวํ ปชานาติ, เอวํ อาสวสมุทยํ ปชานาติ, เอวํ อาสวนิโรธํ ปชานาติ, เอวํ อาสว\-นิโรธ\-คามินิํ ปฏิปทํ ปชานาติ. โส สพฺพโส ราคานุสยํ ปหาย ปฏิฆานุสยํ ปฏิวิโนเทตฺวา อสฺมีติ ทิฏฺฐิ\-มานา\-นุสยํ สมูหนิตฺวา อวิชฺชํ ปหาย วิชฺชํ อุปฺปาเทตฺวา ทิฏฺเฐว ธมฺเม ทุกฺขสฺส\-นฺตกโร โหติ. เอตฺตาวตาปิ โข อาวุโส อริยสาวโก สมฺมาทิฏฺฐิ \csromanfolio{69}โหติ, อุชุคตาสฺส ทิฏฺฐิ, ธมฺเม อเวจฺจ\-ปฺปสาเทน สมนฺนาคโต, อาคโต อิมํ สทฺธมฺมนฺติ.}

\prosewithcloser{อิทมโวจายสฺมา สาริปุตฺโต. อตฺตมนา เต ภิกฺขู อายสฺมโต สาริปุตฺตสฺส ภาสิตํ อภินนฺทุนฺติ.}{สมฺมาทิฏฺฐิ\-สุตฺตํ นิฏฺฐิตํ นวมํ\csromansharedfootnote{7a1a12fdcf49d1bf}{อิโต ปรํ เกสุจิ โปตฺถเกสุ อิมาปิ คาถาโย เอวํ ทิสฺสนฺติ\csromandash{} ทุกฺขํ ชรามรณํ อุปาทานํ, สฬายตนํ นามรูปํ. วิญฺญาณํ ยา สา ปเร, กตมา ปนาวุโส ปทานํ. กิํ ชาติ ตณฺหา จ เวทนา, อวิชฺชาย จตุกฺกนโย. จตฺตาริ ปเร กตมา, ปนาวุโส ปทานํ เกวลํ. อาหาโร จ ภโว ผสฺโส, สงฺขาโร อาสวปญฺจโม. ยาว ปญฺจ ปเร กตโม, ปนาวุโส ปทานํ กิํ. กตมนฺติ ฉพฺพิธา วุตฺตํ, กตมานิ จตุพฺพิธานิ. กตโม ปญฺจวิโธ วุตฺโต, สพฺเพสํ เอกสงฺขานํ ปญฺจนยปทานิ จาติ.}.}
\csromanmatikaanchor{mtk.18}
\csromantocmarkref{subsection}{10. มหาสติปฏฺฐานสุตฺต}{mtk.18}
\bookmark[dest={mtk.18},level=2]{10. มหาสติปฏฺฐานสุตฺต}
\proseitem{105}{\csromanfolio{70}เอวํ เม สุตํ\csromandash{}เอกํ สมยํ ภควา กุรูสุ วิหรติ กมฺมาสธมฺมํ นาม กุรูนํ นิคโม. ตตฺร โข ภควา ภิกฺขู อามนฺเตสิ “ภิกฺขโว”ติ. “ภทนฺเต”ติ เต ภิกฺขู ภควโต ปจฺจสฺโสสุํ. ภควา เอตทโวจ\csromandash{}}

\csromanmatikaanchor{mtk.19}
\csromantocmarkref{section}{อุทฺเทส}{mtk.19}
\bookmark[dest={mtk.19},level=1]{อุทฺเทส}
\csromanheader{1}{\textbf{อุทฺเทส}}
\proseitem{106}{เอกายโน อยํ ภิกฺขเว มคฺโค สตฺตานํ วิสุทฺธิยา โสก\-ปริเทวานํ\csromansharedfootnote{f357520afb87904e}{ปริทฺทวานํ (สี, อิ)} สมติกฺกมาย ทุกฺข\-โทมนสฺสานํ อตฺถงฺคมาย ญายสฺส อธิคมาย นิพฺพานสฺส สจฺฉิกิริยาย, ยทิทํ จตฺตาโร สติปฏฺฐานา.}

\prosewithcloserruled{กตเม จตฺตาโร. อิธ ภิกฺขเว ภิกฺขุ กาเย กายานุปสฺสี วิหรติ อาตาปี สมฺปชาโน สติมา วิเนยฺย โลเก อภิชฺฌา\-โทมนสฺสํ, เวทนาสุ เวทนานุปสฺสี วิหรติ อาตาปี สมฺปชาโน สติมา วิเนยฺย โลเก อภิชฺฌา\-โทมนสฺสํ, จิตฺเต จิตฺตานุปสฺสี วิหรติ อาตาปี สมฺปชาโน สติมา วิเนยฺย โลเก อภิชฺฌา\-โทมนสฺสํ, ธมฺเมสุ ธมฺมานุปสฺสี วิหรติ อาตาปี สมฺปชาโน สติมา วิเนยฺย โลเก อภิชฺฌา\-โทมนสฺสํ.}{อุทฺเทโส นิฏฺฐิโต.}
\csromanmatikaanchor{mtk.20}
\csromantocmarkref{section}{กายานุปสฺสนา อานาปานปพฺพ}{mtk.20}
\bookmark[dest={mtk.20},level=1]{กายานุปสฺสนา อานาปานปพฺพ}
\csromanheader{1}{\textbf{กายานุปสฺสนา อานาปานปพฺพ}}
\proseitem{107}{กถญฺจ ภิกฺขเว ภิกฺขุ กาเย กายานุปสฺสี วิหรติ. อิธ ภิกฺขเว ภิกฺขุ อรญฺญคโต วา รุกฺขมูลคโต วา สุญฺญาคารคโต วา นิสีทติ ปลฺลงฺกํ อาภุชิตฺวา อุชุํ กายํ ปณิธายปริมุขํ สติํ อุปฏฺฐเปตฺวา. โส สโตว อสฺสสติ, สโตว\csromansharedfootnote{441686471b68c19e}{สโต (สี, สฺยา)} ปสฺสสติ. ทีฆํ วา อสฺสสนฺโต “ทีฆํ อสฺสสามี”ติ ปชานาติ, ทีฆํ วา ปสฺสสนฺโต “ทีฆํ ปสฺสสามี”ติ ปชานาติ. รสฺสํ วา อสฺสสนฺโต “รสฺสํ อสฺสสามี”ติ ปชานาติ, รสฺสํ วา ปสฺสสนฺโต “รสฺสํ ปสฺสสามี”ติ ปชานาติ. “สพฺพ\-กาย\-ปฏิสํเวที อสฺสสิสฺสามี”ติ สิกฺขติ, “สพฺพ\-กาย\-ปฏิสํเวที ปสฺสสิสฺสามี”ติ \csromanfolio{71}สิกฺขติ. “ปสฺสมฺภยํ กายสงฺขารํ อสฺสสิสฺสามี”ติ สิกฺขติ, “ปสฺสมฺภยํ กายสงฺขารํ ปสฺสสิสฺสามี”ติ สิกฺขติ.}

\prosewithcloserruled{เสยฺยถาปิ ภิกฺขเว ทกฺโข ภมกาโร วา ภมการ\-นฺเตวาสี วา ทีฆํ วา อญฺฉนฺโต “ทีฆํ อญฺฉามี”ติ ปชานาติ, รสฺสํ วา อญฺฉนฺโต “รสฺสํ อญฺฉามี”ติ ปชานาติ. เอวเมว โข ภิกฺขเว ภิกฺขุ ทีฆํ วา อสฺสสนฺโต “ทีฆํ อสฺสสามี”ติ ปชานาติ, ทีฆํ วา ปสฺสสนฺโต “ทีฆํ ปสฺสสามี”ติ ปชานาติ. รสฺสํ วา อสฺสสนฺโต “รสฺสํ อสฺสสามี”ติ ปชานาติ, รสฺสํ วา ปสฺสสนฺโต “รสฺสํ ปสฺสสามี”ติ ปชานาติ. “สพฺพ\-กาย\-ปฏิสํเวที อสฺสสิสฺสามี”ติ สิกฺขติ, “สพฺพ\-กาย\-ปฏิสํเวที ปสฺสสิสฺสามี”ติ สิกฺขติ. “ปสฺสมฺภยํ กายสงฺขารํ อสฺสสิสฺสามี”ติ สิกฺขติ, “ปสฺสมฺภยํ กายสงฺขารํ ปสฺสสิสฺสามี”ติ สิกฺขติ. อิติ อชฺฌตฺตํ วา กาเย กายานุปสฺสี วิหรติ, พหิทฺธา วา กาเย กายานุปสฺสี วิหรติ, อชฺฌตฺต\-พหิทฺธา วา กาเย กายานุปสฺสี วิหรติ. สมุทยธมฺมา\-นุปสฺสี วา กายสฺมิํ วิหรติ, วยธมฺมา\-นุปสฺสี วา กายสฺมิํ วิหรติ, สมุทย\-วย\-ธมฺมานุปสฺสี วา กายสฺมิํ วิหรติ. “อตฺถิ กาโย”ติ วา ปนสฺส สติ ปจฺจุปฏฺฐิตา โหติ ยาวเทว ญาณมตฺตาย ปฏิสฺสติ\-มตฺตาย. อนิสฺสิโต จ วิหรติ, น จ กิญฺจิ โลเก อุปาทิยติ. เอวมฺปิ โข\csromansharedfootnote{ac84e5e12e7ee72c}{เอวมฺปิ (สี, สฺยา, อิ)} ภิกฺขเว ภิกฺขุ กาเย กายานุปสฺสี วิหรติ.}{อานาปานปพฺพํ นิฏฺฐิตํ.}
\csromanmatikaanchor{mtk.21}
\csromantocmarkref{section}{กายานุปสฺสนา อิริยาปถปพฺพ}{mtk.21}
\bookmark[dest={mtk.21},level=1]{กายานุปสฺสนา อิริยาปถปพฺพ}
\csromanheader{1}{กายานุปสฺสนา อิริยาปถ\-ปพฺพ}
\proseitemwithcloserruled{108}{ปุน จปรํ ภิกฺขเว ภิกฺขุ คจฺฉนฺโต วา “คจฺฉามี”ติ ปชานาติ, ฐิโต วา “ฐิโตมฺหี”ติ ปชานาติ, นิสินฺโน วา “นิสินฺโนมฺหี”ติ ปชานาติ, สยาโน วา “สยาโนมฺหี”ติ ปชานาติ. ยถา ยถา วา ปนสฺส กาโย ปณิหิโต โหติ, ตถา ตถา นํ ปชานาติ. อิติ อชฺฌตฺตํ วา กาเย กายานุปสฺสี วิหรติ, พหิทฺธา วา กาเย กายานุปสฺสี วิหรติ, อชฺฌตฺต\-พหิทฺธา วา กาเย กายานุปสฺสี วิหรติ. สมุทยธมฺมา\-นุปสฺสี วา กายสฺมิํ วิหรติ, วยธมฺมา\-นุปสฺสี วา กายสฺมิํ วิหรติ, สมุทย\-วย\-ธมฺมานุปสฺสี วา กายสฺมิํ วิหรติ. \csromanfolio{72}“อตฺถิ กาโย”ติ วา ปนสฺส สติ ปจฺจุปฏฺฐิตา โหติ ยาวเทว ญาณมตฺตาย ปฏิสฺสติ\-มตฺตาย. อนิสฺสิโต จ วิหรติ, น จ กิญฺจิ โลเก อุปาทิยติ. เอวมฺปิ โข ภิกฺขเว ภิกฺขุ กาเย กายานุปสฺสี วิหรติ.}{อิริยาปถ\-ปพฺพํ นิฏฺฐิตํ.}
\csromanmatikaanchor{mtk.22}
\csromantocmarkref{section}{กายานุปสฺสนา สมฺปชานปพฺพ}{mtk.22}
\bookmark[dest={mtk.22},level=1]{กายานุปสฺสนา สมฺปชานปพฺพ}
\csromanheader{1}{\textbf{กายานุปสฺสนา สมฺปชานปพฺพ}}
\proseitemwithcloserruled{109}{ปุน จปรํ ภิกฺขเว ภิกฺขุ อภิกฺกนฺเต ปฏิกฺกนฺเต สมฺปชานการี โหติ, อาโลกิเต วิโลกิเต สมฺปชานการี โหติ, สมิญฺชิเต ปสาริเต สมฺปชานการี โหติ, สํฆาฏิ\-ปตฺต\-จีวร\-ธารเณ สมฺปชานการี โหติ, อสิเต ปีเต ขายิเต สายิเต สมฺปชานการี โหติ, อุจฺจาร\-ปสฺสาว\-กมฺเม สมฺปชานการี โหติ, คเต ฐิเต นิสินฺเน สุตฺเต ชาคริเต ภาสิเต ตุณฺหีภาเว สมฺปชานการี โหติ. อิติ อชฺฌตฺตํ วา กาเย กายานุปสฺสี วิหรติ ฯเปฯ. เอวมฺปิ โข ภิกฺขเว ภิกฺขุ กาเย กายานุปสฺสี วิหรติ.}{สมฺปชาน\-ปพฺพํ นิฏฺฐิตํ.}
\csromanmatikaanchor{mtk.23}
\csromantocmarkref{section}{กายานุปสฺสนา ปฏิกูลมนสิการปพฺพ}{mtk.23}
\bookmark[dest={mtk.23},level=1]{กายานุปสฺสนา ปฏิกูลมนสิการปพฺพ}
\csromanheader{1}{กายานุปสฺสนา ปฏิกูล\-มนสิการ\-ปพฺพ}
\proseitem{110}{ปุน จปรํ ภิกฺขเว ภิกฺขุ อิมเมว กายํ อุทฺธํ ปาทตลา อโธ เกสมตฺถกา ตจปริยนฺตํ ปูรํ นานปฺปการสฺส อสุจิโน ปจฺจเวกฺขติ “อตฺถิ อิมสฺมิํ กาเย เกสา โลมา นขา ทนฺตา ตโจ, มํสํ นฺหารุ\csromansharedfootnote{04d6dcfeb695a177}{นหารุ (สี, สฺยา, อิ)} อฏฺฐิ อฏฺฐิมิญฺชํ วกฺกํ, หทยํ ยกนํ กิโลมกํ ปิหกํ ปปฺผาสํ, อนฺตํ อนฺตคุณํ อุทริยํ กรีสํ, ปิตฺตํ เสมฺหํ ปุพฺโพ โลหิตํ เสโท เมโท, อสฺสุ วสา เขโฬ สิงฺฆาณิกา ลสิกา มุตฺตนฺ”ติ\csromansharedfootnote{a482e95794c502aa}{มุตฺตํ มตฺถลุงฺคนฺติ (ก)}.}

\prosewithcloserruled{เสยฺยถาปิ ภิกฺขเว อุภโตมุขา ปุโตฬิ\csromansharedfootnote{bbfe134a5eb91b5d}{มูโตฬี (สี, สฺยา, อิ)} ปูรา นานาวิหิตสฺส ธญฺญสฺส. เสยฺยถิทํ, สาลีนํ วีหีนํ มุคฺคานํ มาสานํ ติลานํ ตณฺฑุลานํ. ตเมนํ จกฺขุมา ปุริโส มุญฺจิตฺวา ปจฺจเวกฺเขยฺย “อิเม สาลี อิเม วีหี \csromanfolio{73}อิเม มุคฺคา อิเม มาสา อิเม ติลา อิเม ตณฺฑุลา”ติ. เอวเมว โข ภิกฺขเว ภิกฺขู อิมเมว กายํ อุทฺธํ ปาทตลา อโธ เกสมตฺถกา ตจปริยนฺตํ ปูรํ นานปฺปการสฺส อสุจิโน ปจฺจเวกฺขติ “อตฺถิ อิมสฺมิํ กาเย เกสา โลมา ฯเปฯ มุตฺตนฺ”ติ. อิติ อชฺฌตฺตํ วา กาเย กายานุปสฺสี วิหรติ ฯเปฯ. เอวมฺปิ โข ภิกฺขเว ภิกฺขุ กาเย กายานุปสฺสี วิหรติ.}{ปฏิกูล\-มนสิการ\-ปพฺพํ นิฏฺฐิตํ.}
\csromanmatikaanchor{mtk.24}
\csromantocmarkref{section}{กายานุปสฺสนา ธาตุมนสิการปพฺพ}{mtk.24}
\bookmark[dest={mtk.24},level=1]{กายานุปสฺสนา ธาตุมนสิการปพฺพ}
\csromanheader{1}{กายานุปสฺสนา ธาตุ\-มนสิการ\-ปพฺพ}
\proseitem{111}{ปุน จปรํ ภิกฺขเว ภิกฺขุ อิมเมว กายํ ยถาฐิตํ ยถาปณิหิตํ ธาตุโส ปจฺจเวกฺขติ “อตฺถิ อิมสฺมิํ กาเย ปถวีธาตุ อาโปธาตุ เตโชธาตุ วาโยธาตู”ติ.}

\prosewithcloserruled{เสยฺยถาปิ ภิกฺขเว ทกฺโข โคฆาตโก วา โคฆาตก\-นฺเตวาสี วา คาวิํ วธิตฺวา จตุมหาปเถ\csromansharedfootnote{140db98a8d046582}{จาตุมฺมหาปเถ (สี, สฺยา, อิ)} พิลโส วิภชิตฺวา นิสินฺโน อสฺส. เอวเมว โข ภิกฺขเว ภิกฺขุ อิมเมว กายํ ยถาฐิตํ ยถาปณิหิตํ ธาตุโส ปจฺจเวกฺขติ “อตฺถิ อิมสฺมิํ กาเย ปถวีธาตุ อาโปธาตุ เตโชธาตุ วาโยธาตู”ติ. อิติ อชฺฌตฺตํ วา กาเย กายานุปสฺสี วิหรติ ฯเปฯ. เอวมฺปิ โข ภิกฺขเว ภิกฺขุ กาเย กายานุปสฺสี วิหรติ.}{ธาตุ\-มนสิการ\-ปพฺพํ นิฏฺฐิตํ.}
\csromanmatikaanchor{mtk.25}
\csromantocmarkref{section}{กายานุปสฺสนา นวสิวถิกปพฺพ}{mtk.25}
\bookmark[dest={mtk.25},level=1]{กายานุปสฺสนา นวสิวถิกปพฺพ}
\csromanheader{1}{กายานุปสฺสนา นวสิวถิก\-ปพฺพ}
\proseitem{112}{ปุน จปรํ ภิกฺขเว ภิกฺขุ เสยฺยถาปิ ปสฺเสยฺย สรีรํ สิวถิกาย ฉฑฺฑิตํ เอกาหมตํ วา ทฺวีหมตํ วา ตีหมตํ วา อุทฺธุมาตกํ วินีลกํ วิปุพฺพกชาตํ. โส อิมเมว กายํ อุปสํหรติ “อยมฺปิ โข กาโย เอวํธมฺโม เอวํภาวี เอวํอนตีโต”ติ\csromansharedfootnote{759d515d21c631aa}{เอตํ อนตีโตติ (สี, อิ)}. อิติ อชฺฌตฺตํ วา กาเย กายานุปสฺสี วิหรติ ฯเปฯ. เอวมฺปิ โข ภิกฺขเว ภิกฺขุ กาเย กายานุปสฺสี วิหรติ.}

\prose{\csromanfolio{74}ปุน จปรํ ภิกฺขเว ภิกฺขุ เสยฺยถาปิ ปสฺเสยฺย สรีรํ สิวถิกาย ฉฑฺฑิตํ กาเกหิ วา ขชฺชมานํ กุลเลหิ วา ขชฺชมานํ คิชฺเฌหิ วา ขชฺชมานํ กงฺเกหิ วา ขชฺชมานํ สุนเขหิ วา ขชฺชมานํ พฺยคฺเฆหิ วา ขชฺชมานํ ทีปีหิ วา ขชฺชมานํ สิงฺคาเลหิ วา\csromansharedfootnote{fb328f61f7dbfc13}{คิชฺเฌหิ วา ขชฺชมานํ, สุวาเนหิ วา ขชฺชมานํ, สิคาเลหิ วา (สฺยา, อิ)} ขชฺชมานํ วิวิเธหิ วา ปาณกชาเตหิ ขชฺชมานํ. โส อิมเมว กายํ อุปสํหรติ “อยมฺปิ โข กาโย เอวํธมฺโม เอวํภาวี เอวํอนตีโต”ติ. อิติ อชฺฌตฺตํ วา กาเย กายานุปสฺสี วิหรติ ฯเปฯ. เอวมฺปิ โข ภิกฺขเว ภิกฺขุ กาเย กายานุปสฺสี วิหรติ.}

\prose{ปุน จปรํ ภิกฺขเว ภิกฺขุ เสยฺยถาปิ ปสฺเสยฺย สรีรํ สิวถิกาย ฉฑฺฑิตํ อฏฺฐิก\-สงฺขลิกํ สมํสโลหิตํ นฺหารุ\-สมฺพนฺธํ ฯเปฯ อฏฺฐิก\-สงฺขลิกํ นิมํส\-โลหิต\-มกฺขิตํ นฺหารุ\-สมฺพนฺธํ ฯเปฯ อฏฺฐิก\-สงฺขลิกํ อปคต\-มํส\-โลหิตํ นฺหารุ\-สมฺพนฺธํ ฯเปฯ อฏฺฐิกานิ อปคต\-สมฺพนฺธานิ\csromansharedfootnote{0d6e59b13ed487cc}{อปคตนฺหารุสมฺพนฺธานิ (สฺยา)} ทิสา วิทิสาวิกฺขิตฺตานิ, อญฺเญน หตฺถฏฺฐิกํ อญฺเญน ปาทฏฺฐิกํ อญฺเญน โคปฺผก\-ฏฺฐิกํ\csromansharedfootnote{e2a8990ae2128cc2}{“อญฺเญน โคปฺผกฏฺฐิกนฺ”ติ อิทํ สีสฺยาอิโปตฺถเกสุ นตฺถิ. 4-4. อญฺเญน กฏฏฺฐิกํ อญฺเญน ปิฏฺฐฏฺฐิกํ อญฺเญน กณฺฑกฏฺฐิกํ อญฺเญน ผาสุกฏฺฐิกํ อญฺเญน อุรฏฺฐิกํ อญฺเญน อํสฏฺฐิกํ อญฺเญน พาหุฏฺฐิกํ (สฺยา)} อญฺเญน ชงฺฆฏฺฐิกํ อญฺเญน อูรุฏฺฐิกํ อญฺเญน กฏิฏฺฐิกํ4 อญฺเญน ผาสุกฏฺฐิกํ อญฺเญน ปิฏฺฐิฏฺฐิกํ อญฺเญน ขนฺธฏฺฐิกํ4 อญฺเญน คีวฏฺฐิกํ อญฺเญน หนุกฏฺฐิกํ อญฺเญน ทนฺตฏฺฐิกํ อญฺเญน สีสกฏาหํ. โส อิมเมว กายํ อุปสํหรติ “อยมฺปิ โข กาโย เอวํธมฺโม เอวํภาวี เอวํอนตีโต”ติ. อิติ อชฺฌตฺตํ วา กาเย กายานุปสฺสี วิหรติ ฯเปฯ. เอวมฺปิ โข ภิกฺขเว ภิกฺขุ กาเย กายานุปสฺสี วิหรติ.}

\prosewithcloser{ปุน จปรํ ภิกฺขเว ภิกฺขุ เสยฺยถาปิ ปสฺเสยฺย สรีรํ สิวถิกาย ฉฑฺฑิตํ อฏฺฐิกานิ เสตานิ สงฺข\-วณฺณ\-ปฏิภาคานิ\csromansharedfootnote{7a3ac8a606240a91}{สงฺขวณฺณูปนิภานิ (สี, สฺยา, อิ)} ฯเปฯ อฏฺฐิกานิ ปุญฺชกิตานิ เตโรวสฺสิกานิ ฯเปฯ อฏฺฐิกานิ ปูตีนิ จุณฺณกชาตานิ. โส อิมเมว กายํ อุปสํหรติ “อยมฺปิ โข กาโย เอวํธมฺโม เอวํภาวี เอวํอนตีโต”ติ. อิติ อชฺฌตฺตํ วา กาเย กายานุปสฺสี วิหรติ, พหิทฺธา วา กาเย กายานุปสฺสี วิหรติ, อชฺฌตฺต\-พหิทฺธา วา กาเย \csromanfolio{75}กายานุปสฺสี วิหรติ. สมุทยธมฺมา\-นุปสฺสี วา กายสฺมิํ วิหรติ, วยธมฺมา\-นุปสฺสี วา กายสฺมิํ วิหรติ, สมุทย\-วย\-ธมฺมานุปสฺสี วา กายสฺมิํ วิหรติ. “อตฺถิ กาโย”ติ วา ปนสฺส สติ ปจฺจุปฏฺฐิตา โหติ ยาวเทว ญาณมตฺตาย ปฏิสฺสติ\-มตฺตาย. อนิสฺสิโต จ วิหรติ, น จ กิญฺจิ โลเก อุปาทิยติ. เอวมฺปิ โข ภิกฺขเว ภิกฺขุ กาเย กายานุปสฺสี วิหรติ.}{นวสิวถิก\-ปพฺพํ นิฏฺฐิตํ.}
\nitthitamruled{จุทฺทส\-กายานุปสฺสนา นิฏฺฐิตา.}
\csromanmatikaanchor{mtk.26}
\csromantocmarkref{section}{เวทนานุปสฺสนา}{mtk.26}
\bookmark[dest={mtk.26},level=1]{เวทนานุปสฺสนา}
\csromanheader{1}{เวทนา\-นุปสฺสนา}
\proseitemwithcloserruled{113}{กถญฺจ ภิกฺขเว ภิกฺขุ เวทนาสุ เวทนานุปสฺสี วิหรติ. อิธ ภิกฺขเว ภิกฺขุ สุขํ วา\csromansharedfootnote{0cbb29fcd0dde78c}{สุขํ, ทุกฺขํ, อทุกฺขมสุขํ (สี, สฺยา, อิ, ก)} เวทนํ เวทยมาโน “สุขํ เวทนํ เวทยามี”ติ ปชานาติ. ทุกฺขํ วา\csromansharedfootnote{0cbb29fcd0dde78c}{สุขํ, ทุกฺขํ, อทุกฺขมสุขํ (สี, สฺยา, อิ, ก)} เวทนํ เวทยมาโน “ทุกฺขํ เวทนํ เวทยามี”ติ ปชานาติ. อทุกฺขมสุขํ วา\csromansharedfootnote{0cbb29fcd0dde78c}{สุขํ, ทุกฺขํ, อทุกฺขมสุขํ (สี, สฺยา, อิ, ก)} เวทนํ เวทยมาโน “อทุกฺขมสุขํ เวทนํ เวทยามี”ติ ปชานาติ. สามิสํ วา สุขํ เวทนํ เวทยมาโน “สามิสํ สุขํ เวทนํ เวทยามี”ติ ปชานาติ, นิรามิสํ วา สุขํ เวทนํ เวทยมาโน “นิรามิสํ สุขํ เวทนํ เวทยามี”ติ ปชานาติ. สามิสํ วา ทุกฺขํ เวทนํ เวทยมาโน “สามิสํ ทุกฺขํ เวทนํ เวทยามี”ติ ปชานาติ, นิรามิสํ วา ทุกฺขํ เวทนํ เวทยมาโน “นิรามิสํ ทุกฺขํ เวทนํ เวทยามี”ติ ปชานาติ. สามิสํ วา อทุกฺขมสุขํ เวทนํ เวทยมาโน สามิสํ อทุกฺขมสุขํ เวทนํ เวทยามี”ติ ปชานาติ, นิรามิสํ วา อทุกฺขมสุขํ เวทนํ เวทยมาโน “นิรามิสํ อทุกฺขมสุขํ เวทนํ เวทยามี”ติ ปชานาติ. อิติ อชฺฌตฺตํ วา เวทนาสุ เวทนานุปสฺสี วิหรติ, พหิทฺธา วา เวทนาสุ เวทนานุปสฺสี วิหรติ, อชฺฌตฺต\-พหิทฺธา วา เวทนาสุ เวทนานุปสฺสี วิหรติ. สมุทยธมฺมา\-นุปสฺสี วา เวทนาสุ วิหรติ, วยธมฺมา\-นุปสฺสี วา เวทนาสุ วิหรติ, สมุทย\-วย\-ธมฺมานุปสฺสี วา เวทนาสุ วิหรติ. “อตฺถิ เวทนา”ติ วา ปนสฺส สติ ปจฺจุปฏฺฐิตา โหติ ยาวเทว ญาณมตฺตาย ปฏิสฺสติ\-มตฺตาย. อนิสฺสิโต จ \csromanfolio{76}วิหรติ, น จ กิญฺจิ โลเก อุปาทิยติ. เอวมฺปิ โข ภิกฺขเว ภิกฺขุ เวทนาสุ เวทนานุปสฺสี วิหรติ.}{เวทนา\-นุปสฺสนา นิฏฺฐิตา.}
\csromanmatikaanchor{mtk.27}
\csromantocmarkref{section}{จิตฺตานุปสฺสนา}{mtk.27}
\bookmark[dest={mtk.27},level=1]{จิตฺตานุปสฺสนา}
\csromanheader{1}{\textbf{จิตฺตานุปสฺสนา}}
\proseitemwithcloserruled{114}{กถญฺจ ภิกฺขเว ภิกฺขุ จิตฺเต จิตฺตานุปสฺสี วิหรติ. อิธ ภิกฺขเว ภิกฺขุ สราคํ วา จิตฺตํ “สราคํ จิตฺตนฺ”ติ ปชานาติ, วีตราคํ วา จิตฺตํ “วีตราคํ จิตฺตนฺ”ติ ปชานาติ. สโทสํ วา จิตฺตํ “สโทสํ จิตฺตนฺ”ติ ปชานาติ, วีตโทสํ วา จิตฺตํ “วีตโทสํ จิตฺตนฺ”ติ ปชานาติ. สโมหํ วา จิตฺตํ “สโมหํ จิตฺตนฺ”ติ ปชานาติ, วีตโมหํ วา จิตฺตํ “วีตโมหํ จิตฺตนฺ”ติ ปชานาติ. สํขิตฺตํ วา จิตฺตํ “สํขิตฺตํ จิตฺตนฺ”ติ ปชานาติ, วิกฺขิตฺตํ วา จิตฺตํ “วิกฺขิตฺตํ จิตฺตํ”ติ ปชานาติ. มหคฺคตํ วา จิตฺตํ “มหคฺคตํ จิตฺตนฺ”ติ ปชานาติ, อมหคฺคตํ วา จิตฺตํ “อมหคฺคตํ จิตฺตํ”ติ ปชานาติ. สอุตฺตรํ วา จิตฺตํ “สอุตฺตรํ จิตฺตนฺ”ติ ปชานาติ, อนุตฺตรํ วา จิตฺตํ “อนุตฺตรํ จิตฺตนฺ”ติ ปชานาติ. สมาหิตํ วา จิตฺตํ “สมาหิตํ จิตฺตนฺ”ติ ปชานาติ, อสมาหิตํ วา จิตฺตํ “อสมาหิตํ จิตฺตนฺ”ติ ปชานาติ. วิมุตฺตํ วา จิตฺตํ “วิมุตฺตํ จิตฺตนฺ”ติ ปชานาติ, อวิมุตฺตํ วา จิตฺตํ “อวิมุตฺตํ จิตฺตนฺ”ติ ปชานาติ. อิติ อชฺฌตฺตํ วา จิตฺเต จิตฺตานุปสฺสี วิหรติ, พหิทฺธา วา จิตฺเต จิตฺตานุปสฺสี วิหรติ, อชฺฌตฺต\-พหิทฺธา วา จิตฺเต จิตฺตานุปสฺสี วิหรติ. สมุทยธมฺมา\-นุปสฺสี วา จิตฺตสฺมิํ วิหรติ, วยธมฺมา\-นุปสฺสี วา จิตฺตสฺมิํ วิหรติ, สมุทย\-วย\-ธมฺมานุปสฺสี วา จิตฺตสฺมิํ วิหรติ. “อตฺถิ จิตฺตนฺ”ติ วา ปนสฺส สติ ปจฺจุปฏฺฐิตา โหติ ยาวเทว ญาณมตฺตาย ปฏิสฺสติ\-มตฺตาย. อนิสฺสิโต จ วิหรติ, น จ กิญฺจิ โลเก อุปาทิยติ. เอวมฺปิ โข ภิกฺขเว ภิกฺขุ จิตฺเต จิตฺตานุปสฺสี วิหรติ.}{\textbf{จิตฺตานุปสฺสนา} นิฏฺฐิตา.}
\csromanmatikaanchor{mtk.28}
\csromantocmarkref{section}{ธมฺมานุปสฺสนา นีวรณปพฺพ}{mtk.28}
\bookmark[dest={mtk.28},level=1]{ธมฺมานุปสฺสนา นีวรณปพฺพ}
\csromanheader{1}{ธมฺมา\-นุปสฺสนา นีวรณปพฺพ}
\proseitem{115}{กถญฺจ ภิกฺขเว ภิกฺขุ ธมฺเมสุ ธมฺมานุปสฺสี วิหรติ. อิธ ภิกฺขเว ภิกฺขุ ธมฺเมสุ ธมฺมานุปสฺสี วิหรติ ปญฺจสุ นีวรเณสุ. กถญฺจ ปน ภิกฺขเว ภิกฺขุ ธมฺเมสุ ธมฺมานุปสฺสี วิหรติ ปญฺจสุ นีวรเณสุ.}

\prose{\csromanfolio{77}อิธ ภิกฺขเว ภิกฺขุ สนฺตํ วา อชฺฌตฺตํ กามจฺฉนฺทํ “อตฺถิ เม อชฺฌตฺตํ กามจฺฉนฺโท”ติ ปชานาติ, อสนฺตํ วา อชฺฌตฺตํ กามจฺฉนฺทํ “นตฺถิ เม อชฺฌตฺตํ กามจฺฉนฺโท”ติ ปชานาติ, ยถา จ อนุปฺปนฺนสฺส กามจฺฉนฺทสฺส อุปฺปาโท โหติ ตญฺจ ปชานาติ, ยถา จ อุปฺปนฺนสฺส กามจฺฉนฺทสฺส ปหานํ โหติ ตญฺจ ปชานาติ, ยถา จ ปหีนสฺส กามจฺฉนฺทสฺส อายติํ อนุปฺปาโท โหติ ตญฺจ ปชานาติ.}

\prose{สนฺตํ วา อชฺฌตฺตํ พฺยาปาทํ “อตฺถิ เม อชฺฌตฺตํ พฺยาปาโท”ติ ปชานาติ, อสนฺตํ วา อชฺฌตฺตํ พฺยาปาทํ “นตฺถิ เม อชฺฌตฺตํ พฺยาปาโท”ติ ปชานาติ, ยถา จ อนุปฺปนฺนสฺส พฺยาปาทสฺส อุปฺปาโท โหติ ตญฺจ ปชานาติ, ยถา จ อุปฺปนฺนสฺส พฺยาปาทสฺส ปหานํ โหติ ตญฺจ ปชานาติ, ยถา จ ปหีนสฺส พฺยาปาทสฺส อายติํ อนุปฺปาโท โหติ ตญฺจ ปชานาติ.}

\prose{สนฺตํ วา อชฺฌตฺตํ ถินมิทฺธํ “อตฺถิ เม อชฺฌตฺตํ ถินมิทฺธนฺ”ติ ปชานาติ, อสนฺตํ วา อชฺฌตฺตํ ถินมิทฺธํ “นตฺถิ เม อชฺฌตฺตํ ถินมิทฺธนฺ”ติ ปชานาติ, ยถา จ อนุปฺปนฺนสฺส ถินมิทฺธสฺส อุปฺปาโท โหติ ตญฺจ ปชานาติ, ยถา จ อุปฺปนฺนสฺส ถินมิทฺธสฺส ปหานํ โหติ ตญฺจ ปชานาติ, ยถา จ ปหีนสฺส ถินมิทฺธสฺส อายติํ อนุปฺปาโท โหติ ตญฺจ ปชานาติ.}

\prose{สนฺตํ วา อชฺฌตฺตํ อุทฺธจฺจ\-กุกฺกุจฺจํ “อตฺถิ เม อชฺฌตฺตํ อุทฺธจฺจกุกฺกุจฺจนฺ”ติ ปชานาติ, อสนฺตํ วา อชฺฌตฺตํ อุทฺธจฺจ\-กุกฺกุจฺจํ “นตฺถิ เม อชฺฌตฺตํ อุทฺธจฺจกุกฺกุจฺจนฺ”ติ ปชานาติ, ยถา จ อนุปฺปนฺนสฺส อุทฺธจฺจ\-กุกฺกุจฺจสฺส อุปฺปาโท โหติ ตญฺจ ปชานาติ, ยถา จ อุปฺปนฺนสฺส อุทฺธจฺจ\-กุกฺกุจฺจสฺส ปหานํ โหติ ตญฺจ ปชานาติ, ยถา จ ปหีนสฺส อุทฺธจฺจ\-กุกฺกุจฺจสฺส อายติํ อนุปฺปาโท โหติ ตญฺจ ปชานาติ.}

\prose{สนฺตํ วา อชฺฌตฺตํ วิจิกิจฺฉํ “อตฺถิ เม อชฺฌตฺตํ วิจิกิจฺฉา”ติ ปชานาติ, อสนฺตํ วา อชฺฌตฺตํ วิจิกิจฺฉํ “นตฺถิ เม อชฺฌตฺตํ วิจิกิจฺฉา”ติ ปชานาติ, ยถา จ อนุปฺปนฺนาย วิจิกิจฺฉาย อุปฺปาโท โหติ ตญฺจ ปชานาติ, ยถา จ อุปฺปนฺนาย วิจิกิจฺฉาย ปหานํ โหติ ตญฺจ ปชานาติ, ยถา จ ปหีนาย วิจิกิจฺฉาย อายติํ อนุปฺปาโท โหติ ตญฺจ ปชานาติ.}

\prose{อิติ อชฺฌตฺตํ วา ธมฺเมสุ ธมฺมานุปสฺสี วิหรติ, พหิทฺธา วา ธมฺเมสุ ธมฺมานุปสฺสี วิหรติ, อชฺฌตฺต\-พหิทฺธา วา ธมฺเมสุ ธมฺมานุปสฺสี วิหรติ. สมุทยธมฺมา\-นุปสฺสี วา ธมฺเมสุ วิหรติ, วยธมฺมา\-นุปสฺสี วา ธมฺเมสุ}

\prosewithcloserruled{\csromanfolio{78}วิหรติ, สมุทย\-วย\-ธมฺมานุปสฺสี วา ธมฺเมสุ วิหรติ. “อตฺถิ ธมฺมา”ติ วา ปนสฺส สติ ปจฺจุปฏฺฐิตา โหติ ยาวเทว ญาณมตฺตาย ปฏิสฺสติ\-มตฺตาย. อนิสฺสิโต จ วิหรติ, น จ กิญฺจิ โลเก อุปาทิยติ. เอวมฺปิ โข ภิกฺขเว ภิกฺขุ ธมฺเมสุ ธมฺมานุปสฺสี วิหรติ ปญฺจสุ นีวรเณสุ.}{นีวรณปพฺพํ นิฏฺฐิตํ.}
\csromanmatikaanchor{mtk.29}
\csromantocmarkref{section}{ธมฺมานุปสฺสนา ขนฺธปพฺพ}{mtk.29}
\bookmark[dest={mtk.29},level=1]{ธมฺมานุปสฺสนา ขนฺธปพฺพ}
\csromanheader{1}{ธมฺมา\-นุปสฺสนา ขนฺธปพฺพ}
\proseitemwithcloserruled{116}{ปุน จปรํ ภิกฺขเว ภิกฺขุ ธมฺเมสุ ธมฺมานุปสฺสี วิหรติ ปญฺจสุ อุปาทาน\-กฺขนฺเธสุ. กถญฺจ ปน ภิกฺขเว ภิกฺขุ ธมฺเมสุ ธมฺมานุปสฺสี วิหรติ ปญฺจสุ อุปาทาน\-กฺขนฺเธสุ. อิธ ภิกฺขเว ภิกฺขุ อิติ รูปํ, อิติ รูปสฺส สมุทโย, อิติ รูปสฺส อตฺถงฺคโม. อิติ เวทนา, อิติ เวทนาย สมุทโย, อิติ เวทนาย อตฺถงฺคโม. อิติ สญฺญา, อิติ สญฺญาย สมุทโย, อิติ สญฺญาย อตฺถงฺคโม. อิติ สงฺขารา, อิติ สงฺขารานํ สมุทโย, อิติ สงฺขารานํ อตฺถงฺคโม. อิติ วิญฺญาณํ, อิติ วิญฺญาณสฺส สมุทโย, อิติ วิญฺญาณสฺส อตฺถงฺคโมติ. อิติ อชฺฌตฺตํ วา ธมฺเมสุ ธมฺมานุปสฺสี วิหรติ พหิทฺธา วา ธมฺเมสุ ธมฺมานุปสฺสี วิหรติ, อชฺฌตฺต\-พหิทฺธา วา ธมฺเมสุ ธมฺมานุปสฺสี วิหรติ. สมุทยธมฺมา\-นุปสฺสี วา ธมฺเมสุ วิหรติ, วยธมฺมา\-นุปสฺสี วา ธมฺเมสุ วิหรติ, สมุทย\-วย\-ธมฺมานุปสฺสี วา ธมฺเมสุ วิหรติ. “อตฺถิ ธมฺมา”ติ วา ปนสฺส สติ ปจฺจุปฏฺฐิตา โหติ ยาวเทว ญาณมตฺตาย ปฏิสฺสติ\-มตฺตาย. อนิสฺสิโต จ วิหรติ, น จ กิญฺจิ โลเก อุปาทิยติ. เอวมฺปิ โข ภิกฺขเว ภิกฺขุ ธมฺเมสุ ธมฺมานุปสฺสี วิหรติ ปญฺจสุ อุปาทาน\-กฺขนฺเธสุ.}{ขนฺธปพฺพํ นิฏฺฐิตํ.}
\csromanmatikaanchor{mtk.30}
\csromantocmarkref{section}{ธมฺมานุปสฺสนา อายตนปพฺพ}{mtk.30}
\bookmark[dest={mtk.30},level=1]{ธมฺมานุปสฺสนา อายตนปพฺพ}
\csromanheader{1}{ธมฺมา\-นุปสฺสนา อายตนปพฺพ}
\proseitem{117}{ปุน จปรํ ภิกฺขเว ภิกฺขุ ธมฺเมสุ ธมฺมานุปสฺสี วิหรติ ฉสุ อชฺฌตฺติก\-พาหิเรสุ อายตเนสุ. กถญฺจ ปน ภิกฺขเว ภิกฺขุ ธมฺเมสุ ธมฺมานุปสฺสี วิหรติ ฉสุ อชฺฌตฺติก\-พาหิเรสุ อายตเนสุ.}

\prose{\csromanfolio{79}อิธ ภิกฺขเว ภิกฺขุ จกฺขุญฺจ ปชานาติ, รูเป จ ปชานาติ, ยญฺจ ตทุภยํ ปฏิจฺจ อุปฺปชฺชติ สํโยชนํ ตญฺจ ปชานาติ, ยถา จ อนุปฺปนฺนสฺส สํโยชนสฺส อุปฺปาโท โหติ ตญฺจ ปชานาติ, ยถา จ อุปฺปนฺนสฺส สํโยชนสฺส ปหานํ โหติ ตญฺจ ปชานาติ, ยถา จ ปหีนสฺส สํโยชนสฺส อายติํ อนุปฺปาโท โหติ ตญฺจ ปชานาติ.}

\prose{โสตญฺจ ปชานาติ, สทฺเท จ ปชานาติ, ยญฺจ ตทุภยํ ปฏิจฺจ อุปฺปชฺชติ สํโยชนํ ตญฺจ ปชานาติ, ยถา จ อนุปฺปนฺนสฺส สํโยชนสฺส อุปฺปาโท โหติ ตญฺจ ปชานาติ, ยถา จ อุปฺปนฺนสฺส สํโยชนสฺส ปหานํ โหติ ตญฺจ ปชานาติ, ยถา จ ปหีนสฺส สํโยชนสฺส อายติํ อนุปฺปาโท โหติ ตญฺจ ปชานาติ.}

\prose{ฆานญฺจ ปชานาติ, คนฺเธ จ ปชานาติ, ยญฺจ ตทุภยํ ปฏิจฺจ อุปฺปชฺชติ สํโยชนํ ตญฺจ ปชานาติ, ยถา จ อนุปฺปนฺนสฺส สํโยชนสฺส อุปฺปาโท โหติ ตญฺจ ปชานาติ, ยถา จ อุปฺปนฺนสฺส สํโยชนสฺส ปหานํ โหติ ตญฺจ ปชานาติ, ยถา จ ปหีนสฺส สํโยชนสฺส อายติํ อนุปฺปาโท โหติ ตญฺจ ปชานาติ.}

\prose{ชิวฺหญฺจ ปชานาติ, รเส จ ปชานาติ, ยญฺจ ตทุภยํ ปฏิจฺจ อุปฺปชฺชติ สํโยชนํ ตญฺจ ปชานาติ, ยถา จ อนุปฺปนฺนสฺส สํโยชนสฺส อุปฺปาโท โหติ ตญฺจ ปชานาติ, ยถา จ อุปฺปนฺนสฺส สํโยชนสฺส ปหานํ โหติ ตญฺจ ปชานาติ, ยถา จ ปหีนสฺส สํโยชนสฺส อายติํ อนุปฺปาโท โหติ ตญฺจ ปชานาติ.}

\prose{กายญฺจ ปชานาติ, โผฏฺฐพฺเพ จ ปชานาติ, ยญฺจ ตทุภยํ ปฏิจฺจ อุปฺปชฺชติ สํโยชนํ ตญฺจ ปชานาติ, ยถา จ อนุปฺปนฺนสฺส สํโยชนสฺส อุปฺปาโท โหติ ตญฺจ ปชานาติ, ยถา จ อุปฺปนฺนสฺส สํโยชนสฺส ปหานํ โหติ ตญฺจ ปชานาติ, ยถา จ ปหีนสฺส สํโยชนสฺส อายติํ อนุปฺปาโท โหติ ตญฺจ ปชานาติ.}

\prose{มนญฺจ ปชานาติ, ธมฺเม จ ปชานาติ, ยญฺจ ตทุภยํ ปฏิจฺจ อุปฺปชฺชติ สํโยชนํ ตญฺจ ปชานาติ, ยถา จ อนุปฺปนฺนสฺส สํโยชนสฺส อุปฺปโท โหติ ตญฺจ ปชานาติ, ยถา จ อุปฺปนฺนสฺส สํโยชนสฺส ปหานํ โหติ ตญฺจ ปชานาติ, ยถา จ ปหีนสฺส สํโยชนสฺส อายติํ อนุปฺปาโท โหติ ตญฺจ ปชานาติ.}

\prosewithcloserruled{\csromanfolio{80}อิติ อชฺฌตฺตํ วา ธมฺเมสุ ธมฺมานุปสฺสี วิหรติ, พหิทฺธา วา ธมฺเมสุ ธมฺมานุปสฺสี วิหรติ, อชฺฌตฺต\-พหิทฺธา วา ธมฺเมสุ ธมฺมานุปสฺสี วิหรติ. สมุทยธมฺมา\-นุปสฺสี วา ธมฺเมสุ วิหรติ, วยธมฺมา\-นุปสฺสี วา ธมฺเมสุ วิหรติ, สมุทย\-วย\-ธมฺมานุปสฺสี วา ธมฺเมสุ วิหรติ. “อตฺถิ ธมฺมา”ติ วา ปนสฺส สติ ปจฺจุปฏฺฐิตา โหติ ยาวเทว ญาณมตฺตาย ปฏิสฺสติ\-มตฺตาย. อนิสฺสิโต จ วิหรติ, น จ กิญฺจิ โลเก อุปาทิยติ. เอวมฺปิ โข ภิกฺขเว ภิกฺขุ ธมฺเมสุ ธมฺมานุปสฺสี วิหรติ ฉสุ อชฺฌตฺติก\-พาหิเรสุ อายตเนสุ.}{อายตนปพฺพํ นิฏฺฐิตํ.}
\csromanmatikaanchor{mtk.31}
\csromantocmarkref{section}{ธมฺมานุปสฺสนา โพชฺฌงฺคปพฺพ}{mtk.31}
\bookmark[dest={mtk.31},level=1]{ธมฺมานุปสฺสนา โพชฺฌงฺคปพฺพ}
\csromanheader{1}{ธมฺมา\-นุปสฺสนา โพชฺฌงฺคปพฺพ}
\proseitem{118}{ปุน จปรํ ภิกฺขเว ภิกฺขุ ธมฺเมสุ ธมฺมานุปสฺสี วิหรติ สตฺตสุ โพชฺฌงฺเคสุ. กถญฺจ ปน ภิกฺขเว ภิกฺขุ ธมฺเมสุ ธมฺมานุปสฺสี วิหรติ สตฺตสุ โพชฺฌงฺเคสุ. อิธ ภิกฺขเว ภิกฺขุ สนฺตํ วา อชฺฌตฺตํ สติสมฺโพชฺฌงฺคํ “อตฺถิ เม อชฺฌตฺตํ สติสมฺโพชฺฌงฺโค”ติ ปชานาติ, อสนฺตํ วา อชฺฌตฺตํ สติสมฺโพชฺฌงฺคํ “นตฺถิ เม อชฺฌตฺตํ สติสมฺโพชฺฌงฺโค”ติ ปชานาติ, ยถา จ อนุปฺปนฺนสฺส สติสมฺโพชฺฌงฺคสฺส อุปฺปาโท โหติ ตญฺจ ปชานาติ, ยถา จ อุปฺปนฺนสฺส สติสมฺโพชฺฌงฺคสฺส ภาวนาย ปาริปูรี โหติ ตญฺจ ปชานาติ.}

\prose{สนฺตํ วา อชฺฌตฺตํ ธมฺมวิจย\-สมฺโพชฺฌงฺคํ “อตฺถิ เม อชฺฌตฺตํ ธมฺมวิจย\-สมฺโพชฺฌงฺโค”ติ ปชานาติ, อสนฺตํ วา อชฺฌตฺตํ ธมฺมวิจย\-สมฺโพชฺฌงฺคํ “นตฺถิ เม อชฺฌตฺตํ ธมฺมวิจย\-สมฺโพชฺฌงฺโค”ติ ปชานาติ, ยถา จ อนุปฺปนฺนสฺส ธมฺมวิจย\-สมฺโพชฺฌงฺคสฺส อุปฺปาโท โหติ ตญฺจ ปชานาติ, ยถา จ อุปฺปนฺนสฺส ธมฺมวิจย\-สมฺโพชฺฌงฺคสฺส ภาวนาย ปาริปูรี โหติ ตญฺจ ปชานาติ.}

\prose{สนฺตํ วา อชฺฌตฺตํ วีริย\-สมฺโพชฺฌงฺคํ “อตฺถิ เม อชฺฌตฺตํ วีริย\-สมฺโพชฺฌงฺโค”ติ ปชานาติ, อสนฺตํ วา อชฺฌตฺตํ วีริย\-สมฺโพชฺฌงฺคํ “นตฺถิ เม อชฺฌตฺตํ วีริย\-สมฺโพชฺฌงฺโค”ติ ปชานาติ, ยถา จ อนุปฺปนฺนสฺส วีริย\-สมฺโพชฺฌงฺคสฺส อุปฺปาโท โหติ ตญฺจ ปชานาติ, ยถา จ อุปฺปนฺนสฺส วีริย\-สมฺโพชฺฌงฺคสฺส ภาวนาย ปาริปูรี โหติ ตญฺจ ปชานาติ.}

\prose{\csromanfolio{81}สนฺตํ วา อชฺฌตฺตํ ปีติสมฺโพชฺฌงฺคํ “อตฺถิ เม อชฺฌตฺตํ ปีติสมฺโพชฺฌงฺโค”ติ ปชานาติ, อสนฺตํ วา อชฺฌตฺตํ ปีติสมฺโพชฺฌงฺคํ “นตฺถิ เม อชฺฌตฺตํ ปีติสมฺโพชฺฌงฺโค”ติ ปชานาติ, ยถา จ อนุปฺปนฺนสฺส ปีติสมฺโพชฺฌงฺคสฺส อุปฺปาโท โหติ ตญฺจ ปชานาติ, ยถา จ อุปฺปนฺนสฺส ปีติสมฺโพชฺฌงฺคสฺส ภาวนาย ปาริปูรี โหติ ตญฺจ ปชานาติ.}

\prose{สนฺตํ วา อชฺฌตฺตํ ปสฺสทฺธิ\-สมฺโพชฺฌงฺคํ “อตฺถิ เม อชฺฌตฺตํ ปสฺสทฺธิ\-สมฺโพชฺฌงฺโค”ติ ปชานาติ, อสนฺตํ วา อชฺฌตฺตํ ปสฺสทฺธิ\-สมฺโพชฺฌงฺคํ “นตฺถิ เม อชฺฌตฺตํ ปสฺสทฺธิ\-สมฺโพชฺฌงฺโค”ติ ปชานาติ, ยถา จ อนุปฺปนฺนสฺส ปสฺสทฺธิ\-สมฺโพชฺฌงฺคสฺส อุปฺปาโท โหติ ตญฺจ ปชานาติ, ยถา จ อุปฺปนฺนสฺส ปสฺสทฺธิ\-สมฺโพชฺฌงฺคสฺส ภาวนาย ปาริปูรี โหติ ตญฺจ ปชานาติ.}

\prose{สนฺตํ วา อชฺฌตฺตํ สมาธิ\-สมฺโพชฺฌงฺคํ “อตฺถิ เม อชฺฌตฺตํ สมาธิ\-สมฺโพชฺฌงฺโค”ติ ปชานาติ, อสนฺตํ วา อชฺฌตฺตํ สมาธิ\-สมฺโพชฺฌงฺคํ “นตฺถิ เม อชฺฌตฺตํ สมาธิ\-สมฺโพชฺฌงฺโค”ติ ปชานาติ, ยถา จ อนุปฺปนฺนสฺส สมาธิ\-สมฺโพชฺฌงฺคสฺส อุปฺปาโท โหติ ตญฺจ ปชานาติ, ยถา จ อุปฺปนฺนสฺส สมาธิ\-สมฺโพชฺฌงฺคสฺส ภาวนาย ปาริปูรี โหติ ตญฺจ ปชานาติ.}

\prose{สนฺตํ วา อชฺฌตฺตํ อุเปกฺขา\-สมฺโพชฺฌงฺคํ “อตฺถิ เม อชฺฌตฺตํ อุเปกฺขา\-สมฺโพชฺฌงฺโค”ติ ปชานาติ, อสนฺตํ วา อชฺฌตฺตํ อุเปกฺขา\-สมฺโพชฺฌงฺคํ “นตฺถิ เม อชฺฌตฺตํ อุเปกฺขา\-สมฺโพชฺฌงฺโค”ติ ปชานาติ, ยถา จ อนุปฺปนฺนสฺส อุเปกฺขา\-สมฺโพชฺฌงฺคสฺส อุปฺปาโท โหติ ตญฺจ ปชานาติ, ยถา จ อุปฺปนฺนสฺส อุเปกฺขา\-สมฺโพชฺฌงฺคสฺส ภาวนา ปาริปูรี โหติ ตญฺจ ปชานาติ.}

\prosewithcloser{อิติ อชฺฌตฺตํ วา ธมฺเมสุ ธมฺมานุปสฺสี วิหรติ,พหิทฺธา วา ธมฺเมสุ ธมฺมานุปสฺสี วิหรติ, อชฺฌตฺต\-พหิทฺธา วา ธมฺเมสุ ธมฺมานุปสฺสี วิหรติ. สมุทยธมฺมา\-นุปสฺสี วา ธมฺเมสุ วิหรติ, วยธมฺมา\-นุปสฺสี วา ธมฺเมสุ วิหรติ, สมุทย\-วย\-ธมฺมานุปสฺสี วา ธมฺเมสุ วิหรติ. “อตฺถิ ธมฺมา”ติ วา ปนสฺส สติ ปจฺจุปฏฺฐิตา โหติ ยาวเทว ญาณมตฺตาย ปฏิสฺสติ\-มตฺตาย. อนิสฺสิโต จ วิหรติ, น จ กิญฺจิ โลเก อุปาทิยติ. เอวมฺปิ โข ภิกฺขเว ภิกฺขุ ธมฺเมสุ ธมฺมานุปสฺสี วิหรติ สตฺตสุ โพชฺฌงฺเคสุ.}{โพชฺฌงฺค\-ปพฺพํ นิฏฺฐิตํ\csromansharedfootnote{41d70ce1b8b762d9}{โพชฺฌงฺคปพฺพํ นิฏฺฐิตํ. ปฐมภาณวารํ (สฺยา)}.}
\csromanmatikaanchor{mtk.32}
\csromantocmarkref{section}{ธมฺมานุปสฺสนา สจฺจปพฺพ}{mtk.32}
\bookmark[dest={mtk.32},level=1]{ธมฺมานุปสฺสนา สจฺจปพฺพ}
\csromanheader{1}{ธมฺมา\-นุปสฺสนา สจฺจปพฺพ}
\proseitemwithcloserruled{119}{\csromanfolio{82}ปุน จปรํ ภิกฺขเว ภิกฺขุ ธมฺเมสุ ธมฺมานุปสฺสี วิหรติ จตูสุ อริยสจฺเจสุ. กถญฺจ ปน ภิกฺขเว ภิกฺขุ ธมฺเมสุ ธมฺมานุปสฺสี วิหรติ จตูสุ อริยสจฺเจสุ. อิธ ภิกฺขเว ภิกฺขุ “อิทํ ทุกฺขนฺ”ติ ยถาภูตํ ปชานาติ, “อยํ ทุกฺขสมุทโย”ติ ยถาภูตํ ปชานาติ, “อยํ ทุกฺขนิโรโธ”ติ ยถาภูตํ ปชานาติ, “อยํ ทุกฺขนิโรธ\-คามินี ปฏิปทา”ติ ยถาภูตํ ปชานาติ.}{ปฐม\-ภาณวาโร นิฏฺฐิโต.}
\csromanmatikaanchor{mtk.33}
\csromantocmarkref{section}{ทุกฺขสจฺจนิทฺเทส}{mtk.33}
\bookmark[dest={mtk.33},level=1]{ทุกฺขสจฺจนิทฺเทส}
\csromanheader{1}{ทุกฺขสจฺจ\-นิทฺเทส}
\proseitem{120}{กตมญฺจ ภิกฺขเว ทุกฺขํ อริยสจฺจํ. ชาติปิ ทุกฺขา, ชราปิ ทุกฺขา, มรณมฺปิ ทุกฺขํ, โสก\-ปริเทว\-ทุกฺข\-โทมนสฺสุ\-ปายาสาปิ ทุกฺขา, อปฺปิเยหิ สมฺปโยโคปิ ทุกฺโข, ปิเยหิ วิปฺปโยโคปิ ทุกฺโข\csromansharedfootnote{4c3602a2ea9e78b0}{“อปฺปิเยหิ ฯเปฯ วิปฺปโยโคปิ ทุกฺโข”ติ ปาโฐ เจว ตํนิทฺเทโส จ สีอิโปตฺถเกสุ น ทิสฺสติ, สุมงฺคลวิลาสินิยํปิ ตํสํวณฺณนา นตฺถิ.}, ยมฺปิจฺฉํ น ลภติ ตมฺปิ ทุกฺขํ, สํขิตฺเตน ปญฺจุ\-ปาทาน\-กฺขนฺธา\csromansharedfootnote{feb51bbb692b3682}{ปญฺจุปาทานกฺขนฺธาปิ (ก)} ทุกฺขา.}

\proseitem{121}{กตมา จ ภิกฺขเว ชาติ. ยา เตสํ เตสํ สตฺตานํ ตมฺหิ ตมฺหิ สตฺตนิกาเย ชาติ สญฺชาติ โอกฺกนฺติ อภินิพฺพตฺติ ขนฺธานํ ปาตุภาโว อายตนานํ ปฏิลาโภ. อยํ วุจฺจติ ภิกฺขเว ชาติ.}

\proseitem{122}{กตมา จ ภิกฺขเว ชรา. ยา เตสํ เตสํ สตฺตานํ ตมฺหิ ตมฺหิ สตฺตนิกาเย ชรา ชีรณตา ขณฺฑิจฺจํ ปาลิจฺจํ วลิตฺตจตา อายุโน สํหานิ อินฺทฺริยานํ ปริปาโก. อยํ วุจฺจติ ภิกฺขเว ชรา.}

\proseitem{123}{กตมญฺจ ภิกฺขเว มรณํ. ยํ\csromansharedfootnote{44882fa2c3546924}{สุมงฺคลวิลาสินี โอโลเกตพฺพา.} เตสํ เตสํ สตฺตานํ ตมฺหา ตมฺหา สตฺตนิกายา จุติ จวนตา เภโท อนฺตรธานํ มจฺจุ มรณํ กาลํกิริยา ขนฺธานํ เภโท กเฬวรสฺส นิกฺเขโป ชีวิตินฺทฺริยสฺสุ\-ปจฺเฉโท. อิทํ วุจฺจติ ภิกฺขเว มรณํ.}

\proseitem{124}{\csromanfolio{83}กตโม จ ภิกฺขเว โสโก. โย โข ภิกฺขเว อญฺญตร\-ญฺญตเรน พฺยสเนน สมนฺนาคตสฺส อญฺญตร\-ญฺญตเรน ทุกฺขธมฺเมน ผุฏฺฐสฺส โสโก โสจนา โสจิตตฺตํ อนฺโตโสโก อนฺโตปริโสโก. อยํ วุจฺจติ ภิกฺขเว โสโก.}

\proseitem{125}{กตโม จ ภิกฺขเว ปริเทโว. โย โข ภิกฺขเว อญฺญตร\-ญฺญตเรน พฺยสเนน สมนฺนาคตสฺส อญฺญตร\-ญฺญตเรน ทุกฺขธมฺเมน ผุฏฺฐสฺส อาเทโว ปริเทโว อาเทวนา ปริเทวนา อาเทวิตตฺตํ ปริเทวิตตฺตํ. อยํ วุจฺจติ ภิกฺขเว ปริเทโว.}

\proseitem{126}{กตมญฺจ ภิกฺขเว ทุกฺขํ. ยํ โข ภิกฺขเว กายิกํ ทุกฺขํ กายิกํ อสาตํ กาย\-สมฺผสฺสชํ ทุกฺขํ อสาตํ เวทยิตํ. อิทํ วุจฺจติ ภิกฺขเว ทุกฺขํ.}

\proseitem{127}{กตมญฺจ ภิกฺขเว โทมนสฺสํ. ยํ โข ภิกฺขเว เจตสิกํ ทุกฺขํ เจตสิกํ อสาตํ มโน\-สมฺผสฺสชํ ทุกฺขํ อสาตํ เวทยิตํ. อิทํ วุจฺจติ ภิกฺขเว โทมนสฺสํ.}

\proseitem{128}{กตโม จ ภิกฺขเว อุปายาโส. โย โข ภิกฺขเว อญฺญตร\-ญฺญตเรน พฺยสเนน สมนฺนาคตสฺส อญฺญตร\-ญฺญตเรน ทุกฺขธมฺเมน ผุฏฺฐสฺส อายาโส อุปายาโส อายาสิตตฺตํ อุปายาสิตตฺตํ. อยํ วุจฺจติ ภิกฺขเว อุปายาโส.}

\proseitem{129}{กตโม จ ภิกฺขเว อปฺปิเยหิ สมฺปโยโค ทุกฺโข. อิธ ยสฺส เต โหนฺติ อนิฏฺฐา อกนฺตา อมนาปา รูปา สทฺทา คนฺธา รสา โผฏฺฐพฺพา ธมฺมา, เย วา ปนสฺส เต โหนฺติ อนตฺถกามา อหิตกามา อผาสุกกามา อโยคกฺเขม\-กามา, ยา เตหิ สทฺธิํ สงฺคติ สมาคโม สโมธานํ มิสฺสีภาโว. อยํ วุจฺจติ ภิกฺขเว อปฺปิเยหิ สมฺปโยโค ทุกฺโข.}

\proseitem{130}{กตโม จ ภิกฺขเว ปิเยหิ วิปฺปโยโค ทุกฺโข. อิธ ยสฺส เต โหนฺติ อิฏฺฐา กนฺตา มนาปา รูปา สทฺทา คนฺธา รสา โผฏฺฐพฺพา ธมฺมา, เย วา ปนสฺส เต โหนฺติ อตฺถกามา หิตกามา ผาสุกกามา โยคกฺเขมกามา มาตา วา ปิตา วา ภาตา วา ภคินี วา มิตฺตา วา อมจฺจา วา ญาติสาโลหิตา วา, ยา เตหิ สทฺธิํ อสงฺคติ \csromanfolio{84}อสมาคโม อสโมธานํ อมิสฺสีภาโว. อยํ วุจฺจติ ภิกฺขเว ปิเยหิ วิปฺปโยโค ทุกฺโข.}

\proseitem{131}{กตมญฺจ ภิกฺขเว ยมฺปิจฺฉํ น ลภติ ตมฺปิ ทุกฺขํ. ชาติธมฺมานํ ภิกฺขเว สตฺตานํ เอวํ อิจฺฉา อุปฺปชฺชติ “อโห วต มยํ น ชาติธมฺมา อสฺสาม, น จ วต โน ชาติ อาคจฺเฉยฺยา”ติ. น โข ปเนตํ อิจฺฉาย ปตฺตพฺพํ, อิทมฺปิ ยมฺปิจฺฉํ น ลภติ ตมฺปิ ทุกฺขํ. ชราธมฺมานํ ภิกฺขเว สตฺตานํ เอวํ อิจฺฉา อุปฺปชฺชติ “อโห วต มยํ น ชราธมฺมา อสฺสาม, น จ วต โน ชรา อาคจฺเฉยฺยา”ติ. น โข ปเนตํ อิจฺฉาย ปตฺตพฺพํ, อิทมฺปิ ยมฺปิจฺฉํ น ลภติ ตมฺปิ ทุกฺขํ. พฺยาธิ\-ธมฺมานํ ภิกฺขเว สตฺตานํ เอวํ อิจฺฉา อุปฺปชฺชติ “อโห วต มยํ น พฺยาธิธมฺมา อสฺสาม, น จ วต โน พฺยาธิ อาคจฺเฉยฺยา”ติ. น โข ปเนตํ อิจฺฉาย ปตฺตพฺพํ, อิทมฺปิ ยมฺปิจฺฉํ น ลภติ ตมฺปิ ทุกฺขํ. มรณ\-ธมฺมานํ ภิกฺขเว สตฺตานํ เอวํ อิจฺฉา อุปฺปชฺชติ “อโห วต มยํ น มรณธมฺมา อสฺสาม, น จ วต โน มรณํ อาคจฺเฉยฺยา”ติ. น โข ปเนตํ อิจฺฉาย ปตฺตพฺพํ, อิทมฺปิ ยมฺปิจฺฉํ น ลภติ ตมฺปิ ทุกฺขํ. โสก\-ปริเทว\-ทุกฺข\-โทมนสฺสุ\-ปายาส\-ธมฺมานํ ภิกฺขเว สตฺตานํ เอวํ อิจฺฉา อุปฺปชฺชติ “อโห วต มยํ น โสก\-ปริเทว\-ทุกฺข\-โทมนสฺสุ\-ปายาส\-ธมฺมา อสฺสาม, น จ วต โน โสก\-ปริเทว\-ทุกฺข\-โทมนสฺสุ\-ปายาสา อาคจฺเฉยฺยุนฺ”ติ. น โข ปเนตํ อิจฺฉาย ปตฺตพฺพํ, อิทมฺปิ ยมฺปิจฺฉํ น ลภติ ตมฺปิ ทุกฺขํ.}

\proseitem{132}{กตเม จ ภิกฺขเว สํขิตฺเตน ปญฺจุ\-ปาทาน\-กฺขนฺธา ทุกฺขา. เสยฺยถิทํ, รูปุปาทาน\-กฺขนฺโธ, เวทนุ\-ปาทาน\-กฺขนฺโธ, สญฺญุ\-ปาทาน\-กฺขนฺโธ, สงฺขารุ\-ปาทาน\-กฺขนฺโธ, วิญฺญาณุ\-ปาทาน\-กฺขนฺโธ. อิเม วุจฺจนฺติ ภิกฺขเว สํขิตฺเตน ปญฺจุ\-ปาทาน\-กฺขนฺธา ทุกฺขา. อิทํ วุจฺจติ ภิกฺขเว ทุกฺขํ อริยสจฺจํ.}

\csromanmatikaanchor{mtk.34}
\csromantocmarkref{section}{สมุทยสจฺจนิทฺเทส}{mtk.34}
\bookmark[dest={mtk.34},level=1]{สมุทยสจฺจนิทฺเทส}
\csromanheader{1}{สมุทยสจฺจ\-นิทฺเทส}
\proseitem{133}{กตมญฺจ ภิกฺขเว ทุกฺข\-สมุทยํ\csromansharedfootnote{7268e5b3db3909f2}{ทุกฺขสมุทโย (สฺยา)} อริยสจฺจํ. ยายํ ตณฺหา โปโนพฺภวิกา\csromansharedfootnote{3459575e13e816f9}{โปโนภวิกา (สี, อิ)} นนฺที\-ราค\-สหคตา\csromansharedfootnote{25ed4f7992388687}{นนฺทิราคสหคตา (สี, สฺยา, อิ)} ตตฺร\-ตตฺรา\-ภินนฺทินี. เสยฺยถิทํ, กามตณฺหา ภวตณฺหา วิภวตณฺหา.}

\prose{\csromanfolio{85}สา โข ปเนสา ภิกฺขเว ตณฺหา กตฺถ อุปฺปชฺชมานา อุปฺปชฺชติ, กตฺถ นิวิสมานา นิวิสติ. ยํ โลเก ปิยรูปํ สาตรูปํ, เอตฺเถสา ตณฺหา อุปฺปชฺชมานา อุปฺปชฺชติ, เอตฺถ นิวิสมานา นิวิสติ.}

\prose{กิญฺจ โลเก ปิยรูปํ สาตรูปํ. จกฺขุ โลเก ปิยรูปํ สาตรูปํ, เอตฺเถสา ตณฺหา อุปฺปชฺชมานา อุปฺปชฺชติ, เอตฺถ นิวิสมานา นิวิสติ. โสตํ โลเก ฯเปฯ. ฆานํ โลเก. ชิวฺหา โลเก. กาโย โลเก. มโน โลเก ปิยรูปํ สาตรูปํ, เอตฺเถสา ตณฺหา อุปฺปชฺชมานา อุปฺปชฺชติ, เอตฺถ นิวิสมานา นิวิสติ.}

\prose{รูปา โลเก. สทฺทา โลเก. คนฺธา โลเก. รสา โลเก. โผฏฺฐพฺพา โลเก. ธมฺมา โลเก ปิยรูปํ สาตรูปํ, เอตฺเถสา ตณฺหา อุปฺปชฺชมานา อุปฺปชฺชติ, เอตฺถ นิวิสมานา นิวิสติ.}

\prose{จกฺขุวิญฺญาณํ โลเก. โสตวิญฺญาณํ โลเก. ฆานวิญฺญาณํ โลเก. ชิวฺหาวิญฺญาณํ โลเก. กายวิญฺญาณํ โลเก. มโนวิญฺญาณํ โลเก ปิยรูปํ สาตรูปํ, เอตฺเถสา ตณฺหา อุปฺปชฺชมานา อุปฺปชฺชติ, เอตฺถ นิวิสมานา นิวิสติ.}

\prose{จกฺขุ\-สมฺผสฺโส โลเก. โสตสมฺผสฺโส โลเก. ฆานสมฺผสฺโส โลเก. ชิวฺหาสมฺผสฺโส โลเก. กายสมฺผสฺโส โลเก. มโนสมฺผสฺโส โลเก ปิยรูปํ สาตรูปํ, เอตฺเถสา ตณฺหา อุปฺปชฺชมานา อุปฺปชฺชติ, เอตฺถ นิวิสมานา นิวิสติ.}

\prose{จกฺขุ\-สมฺผสฺสชา เวทนา โลเก. โสต\-สมฺผสฺสชา เวทนา โลเก. ฆาน\-สมฺผสฺสชา เวทนา โลเก. ชิวฺหา\-สมฺผสฺสชา เวทนา โลเก. กาย\-สมฺผสฺสชา เวทนา โลเก. มโน\-สมฺผสฺสชา เวทนา โลเก ปิยรูปํ สาตรูปํ, เอตฺเถสา ตณฺหา อุปฺปชฺชมานา อุปฺปชฺชติ, เอตฺถ นิวิสมานา นิวิสติ.}

\prose{รูปสญฺญา โลเก. สทฺทสญฺญา โลเก. คนฺธสญฺญา โลเก. รสสญฺญา โลเก. โผฏฺฐพฺพ\-สญฺญา โลเก. ธมฺมสญฺญา โลเก ปิยรูปํ สาตรูปํ, เอตฺเถสา ตณฺหา อุปฺปชฺชมานา อุปฺปชฺชติ, เอตฺถ นิวิสมานา นิวิสติ.}

\prose{\csromanfolio{86}รูปสญฺเจตนา โลเก. สทฺทญฺเจตนา โลเก. คนฺธ\-สญฺเจตนา โลเก. รสสญฺเจตนา โลเก. โผฏฺฐพฺพ\-สญฺเจตนา โลเก. ธมฺม\-สญฺเจตนา โลเก ปิยรูปํ สาตรูปํ, เอตฺเถสา ตณฺหา อุปฺปชฺชมานา อุปฺปชฺชติ, เอตฺถ นิวิสมานา นิวิสติ.}

\prose{รูปตณฺหา โลเก. สทฺทตณฺหา โลเก. คนฺธตณฺหา โลเก. รสตณฺหา โลเก. โผฏฺฐพฺพ\-ตณฺหา โลเก. ธมฺมตณฺหา โลเก ปิยรูปํ สาตรูปํ, เอตฺเถสา ตณฺหา อุปฺปชฺชมานา อุปฺปชฺชติ, เอตฺถ นิวิสมานา นิวิสติ.}

\prose{รูปวิตกฺโก โลเก. สทฺทวิตกฺโก โลเก. คนฺธวิตกฺโก โลเก. รสวิตกฺโก โลเก. โผฏฺฐพฺพ\-วิตกฺโก โลเก. ธมฺมวิตกฺโก โลเก ปิยรูปํ สาตรูปํ, เอตฺเถสา ตณฺหา อุปฺปชฺชมานา อุปฺปชฺชติ, เอตฺถ นิวิสมานา นิวิสติ.}

\prose{รูปวิจาโร โลเก. สทฺทวิจาโร โลเก. คนฺธวิจาโร โลเก. รสวิจาโร โลเก. โผฏฺฐพฺพ\-วิจาโร โลเก. ธมฺมวิจาโร โลเก ปิยรูปํ สาตรูปํ, เอตฺเถสา ตณฺหา อุปฺปชฺชมานา อุปฺปชฺชติ, เอตฺถ นิวิสมานา นิวิสติ. อิทํ วุจฺจติ ภิกฺขเว ทุกฺข\-สมุทยํ อริยสจฺจํ.}

\csromanmatikaanchor{mtk.35}
\csromantocmarkref{section}{นิโรธสจฺจนิทฺเทส}{mtk.35}
\bookmark[dest={mtk.35},level=1]{นิโรธสจฺจนิทฺเทส}
\csromanheader{1}{นิโรธสจฺจ\-นิทฺเทส}
\proseitem{134}{กตมญฺจ ภิกฺขเว ทุกฺขนิโรธํ\csromansharedfootnote{c6e4e2ea3d0b5954}{ทุกฺขนิโรโธ (สฺยา)} อริยสจฺจํ. โย ตสฺสาเยว ตณฺหาย อเสส\-วิราค\-นิโรโธ จาโค ปฏินิสฺสคฺโค มุตฺติ อนาลโย.}

\prose{สา โข ปเนสา ภิกฺขเว ตณฺหา กตฺถ ปหียมานา ปหียติ, กตฺถ นิรุชฺฌมานา นิรุชฺฌติ. ยํ โลเก ปิยรูปํ สาตรูปํ, เอตฺเถสา ตณฺหา ปหียมานา ปหียติ, เอตฺถ นิรุชฺฌมานา นิรุชฺฌติ.}

\prose{กิญฺจ โลเก ปิยรูปํ สาตรูปํ. จกฺขุ โลเก ปิยรูปํ สาตรูปํ, เอตฺเถสา ตณฺหา ปหียมานา ปหียติ, เอตฺถ นิรุชฺฌมานา นิรุชฺฌติ. โสตํ โลเก ฯเปฯ. ฆานํ โลเก. ชิวฺหา โลเก. กาโย โลเก. \csromanfolio{87}มโน โลเก ปิยรูปํ สาตรูปํ, เอตฺเถสา ตณฺหา ปหียมานา ปหียติ, เอตฺถ นิรุชฺฌมานา นิรุชฺฌติ.}

\prose{รูปา โลเก. สทฺทา โลเก. คนฺธา โลเก. รสา โลเก. โผฏฺฐพฺพา โลเก. ธมฺมา โลเก ปิยรูปํ สาตรูปํ, เอตฺเถสา ตณฺหา ปหียมานา ปหียติ, เอตฺถ นิรุชฺฌมานา นิรุชฺฌติ.}

\prose{จกฺขุวิญฺญาณํ โลเก. โสตวิญฺญาณํ โลเก. ฆานวิญฺญาณํ โลเก. ชิวฺหาวิญฺญาณํ โลเก. กายวิญฺญาณํ โลเก. มโนวิญฺญาณํ โลเก ปิยรูปํ สาตรูปํ, เอตฺเถสา ตณฺหา ปหียมานา ปหียติ, เอตฺถ นิรุชฺฌมานา นิรุชฺฌติ.}

\prose{จกฺขุ\-สมฺผสฺโส โลเก. โสตสมฺผสฺโส โลเก. ฆานสมฺผสฺโส โลเก. ชิวฺหาสมฺผสฺโส โลเก. กายสมฺผสฺโส โลเก. มโนสมฺผสฺโส โลเก ปิยรูปํ สาตรูปํ, เอตฺเถสา ตณฺหา ปหียมานา ปหียติ, เอตฺถ นิรุชฺฌมานา นิรุชฺฌติ.}

\prose{จกฺขุ\-สมฺผสฺสชา เวทนา โลเก. โสต\-สมฺผสฺสชา เวทนา โลเก. ฆาน\-สมฺผสฺสชา เวทนา โลเก. ชิวฺหา\-สมฺผสฺสชา เวทนา โลเก. กาย\-สมฺผสฺสชา เวทนา โลเก. มโน\-สมฺผสฺสชา เวทนา โลเก ปิยรูปํ สาตรูปํ, เอตฺเถสา ตณฺหา ปหียมานา ปหียติ, เอตฺถ นิรุชฺฌมานา นิรุชฺฌติ.}

\prose{รูปสญฺญา โลเก. สทฺทสญฺญา โลเก. คนฺธสญฺญา โลเก. รสสญฺญา โลเก. โผฏฺฐพฺพ\-สญฺญา โลเก. ธมฺมสญฺญา โลเก ปิยรูปํ สาตรูปํ, เอตฺเถสา ตณฺหา ปหียมานา ปหียติ, เอตฺถ นิรุชฺฌมานา นิรุชฺฌติ.}

\prose{รูปสญฺเจตนา โลเก. สทฺทสญฺเจตนา โลเก. คนฺธ\-สญฺเจตนา โลเก. รสสญฺเจตนา โลเก. โผฏฺฐพฺพ\-สญฺเจตนา โลเก. ธมฺม\-สญฺเจตนา โลเก ปิยรูปํ สาตรูปํ, เอตฺเถสา ตณฺหา ปหียมานา ปหียติ, เอตฺถ นิรุชฺฌมานา นิรุชฺฌติ.}

\prose{รูปตณฺหา โลเก. สทฺทตณฺหา โลเก. คนฺธตณฺหา โลเก. รสตณฺหา โลเก. โผฏฺฐพฺพ\-ตณฺหา โลเก. ธมฺมตณฺหา โลเก \csromanfolio{88}ปิยรูปํ สาตรูปํ, เอตฺเถสา ตณฺหา ปหียมานา ปหียติ, เอตฺถ นิรุชฺฌมานา นิรุชฺฌติ.}

\proseitem{10}{รูปวิตกฺโก โลเก. สทฺทวิตกฺโก โลเก. คนฺธวิตกฺโก โลเก. รสวิตกฺโก โลเก. โผฏฺฐพฺพ\-วิตกฺโก โลเก. ธมฺมวิตกฺโก โลเก ปิยรูปํ สาตรูปํ, เอตฺเถสา ตณฺหา ปหียมานา ปหียติ, เอตฺถ นิรุชฺฌมานา นิรุชฺฌติ.}

\prose{รูปวิจาโร โลเก. สทฺทวิจาโร โลเก. คนฺธวิจาโร โลเก. รสวิจาโร โลเก. โผฏฺฐพฺพ\-วิจาโร โลเก. ธมฺมวิจาโร โลเก ปิยรูปํ สาตรูปํ, เอตฺเถสา ตณฺหา ปหียมานา ปหียติ, เอตฺถ นิรุชฺฌมานา นิรุชฺฌติ. อิทํ วุจฺจติ ภิกฺขเว ทุกฺขนิโรธํ อริยสจฺจํ.}

\csromanmatikaanchor{mtk.36}
\csromantocmarkref{section}{มคฺคสจฺจนิทฺเทส}{mtk.36}
\bookmark[dest={mtk.36},level=1]{มคฺคสจฺจนิทฺเทส}
\csromanheader{1}{มคฺคสจฺจ\-นิทฺเทส}
\proseitem{135}{กตมญฺจ ภิกฺขเว ทุกฺขนิโรธ\-คามินี ปฏิปทา อริยสจฺจํ. อยเมว อริโย อฏฺฐงฺคิโก มคฺโค. เสยฺยถิทํ, สมฺมาทิฏฺฐิ สมฺมาสงฺกปฺโป สมฺมาวาจา สมฺมากมฺมนฺโต สมฺมาอาชีโว สมฺมาวายาโม สมฺมาสติ สมฺมาสมาธิ.}

\prose{กตมา จ ภิกฺขเว สมฺมาทิฏฺฐิ. ยํ โข ภิกฺขเว ทุกฺเข ญาณํ, ทุกฺขสมุทเย ญาณํ, ทุกฺขนิโรเธ ญาณํ, ทุกฺขนิโรธ\-คามินิยา ปฏิปทาย ญาณํ. อยํ วุจฺจติ ภิกฺขเว สมฺมาทิฏฺฐิ.}

\prose{กตโม จ ภิกฺขเว สมฺมาสงฺกปฺโป. เนกฺขมฺม\-สงฺกปฺโป อพฺยาปาทสงฺกปฺโป อวิหิํสาสงฺกปฺโป. อยํ วุจฺจติ ภิกฺขเว สมฺมาสงฺกปฺโป.}

\prose{กตมา จ ภิกฺขเว สมฺมาวาจา. มุสาวาทา เวรมณี\csromansharedfootnote{f96664fc1fbeb341}{เวรมณิ (ก)} ปิสุณาย วาจาย เวรมณี ผรุสาย วาจาย เวรมณี สมฺผปฺปลาปา เวรมณี. อยํ วุจฺจติ ภิกฺขเว สมฺมาวาจา.}

\prose{กตโม จ ภิกฺขเว สมฺมากมฺมนฺโต. ปาณาติปาตา เวรมณี อทินฺนาทานา เวรมณี กาเมสุ\-มิจฺฉาจารา เวรมณี. อยํ วุจฺจติ ภิกฺขเว สมฺมากมฺมนฺโต.}

\prose{\csromanfolio{89}กตโม จ ภิกฺขเว สมฺมาอาชีโว. อิธ ภิกฺขเว อริยสาวโก มิจฺฉา- อาชีวํ ปหาย สมฺมาอาชีเวน ชีวิตํ กปฺเปติ. อยํ วุจฺจติ ภิกฺขเว สมฺมาอาชีโว.}

\prose{กตโม จ ภิกฺขเว สมฺมาวายาโม. อิธ ภิกฺขเว ภิกฺขุ อนุปฺปนฺนานํ ปาปกานํ อกุสลานํ ธมฺมานํ อนุปฺปาทาย ฉนฺทํ ชเนติ วายมติ วีริยํ อารภติ จิตฺตํ ปคฺคณฺหาติ ปทหติ, อุปฺปนฺนานํ ปาปกานํ อกุสลานํ ธมฺมานํ ปหานาย ฉนฺทํ ชเนติ วายมติ วีริยํ อารภติ จิตฺตํ ปคฺคณฺหาติ ปทหติ, อนุปฺปนฺนานํ กุสลานํ ธมฺมานํ อุปฺปาทาย ฉนฺทํ ชเนติ วายมติ วีริยํ อารภติ จิตฺตํ ปคฺคณฺหาติ ปทหติ, อุปฺปนฺนานํ กุสลานํ ธมฺมานํ ฐิติยา อสมฺโมสาย ภิยฺโยภาวาย เวปุลฺลาย ภาวนาย ปาริปูริยา ฉนฺทํ ชเนติ วายมติ วีริยํ อารภติ จิตฺตํ ปคฺคณฺหาติ ปทหติ. อยํ วุจฺจติ ภิกฺขเว สมฺมาวายาโม.}

\prose{กตมา จ ภิกฺขเว สมฺมาสติ. อิธ ภิกฺขเว ภิกฺขุ กาเย กายานุปสฺสี วิหรติ อาตาปี สมฺปชาโน สติมา วิเนยฺย โลเก อภิชฺฌา\-โทมนสฺสํ, เวทนาสุ เวทนานุปสฺสี วิหรติ อาตาปีสมฺปชาโน สติมา วิเนยฺย โลเก อภิชฺฌา\-โทมนสฺสํ, จิตฺเต จิตฺตานุปสฺสี วิหรติ อาตาปี สมฺปชาโน สติมา วิเนยฺย โลเก อภิชฺฌา\-โทมนสฺสํ, ธมฺเมสุ ธมฺมานุปสฺสี วิหรติ อาตาปี สมฺปชาโน สติมา วิเนยฺย โลเก อภิชฺฌา\-โทมนสฺสํ. อยํ วุจฺจติ ภิกฺขเว สมฺมาสติ.}

\prose{กตโม จ ภิกฺขเว สมฺมาสมาธิ. อิธ ภิกฺขเว ภิกฺขุ วิวิจฺเจว กาเมหิ วิวิจฺจ อกุสเลหิ ธมฺเมหิ สวิตกฺกํ สวิจารํ วิเวกชํ ปีติสุขํ ปฐมํ ฌานํ อุปสมฺปชฺช วิหรติ. วิตกฺก\-วิจารานํ วูปสมา อชฺฌตฺตํ สมฺปสาทนํ เจตโส เอโกทิภาวํ อวิตกฺกํ อวิจารํ สมาธิชํ ปีติสุขํ ทุติยํ ฌานํ อุปสมฺปชฺช วิหรติ. ปีติยา จ วิราคา อุเปกฺขโก จ วิหรติ, สโต จ สมฺปชาโน, สุขญฺจ กาเยน ปฏิสํเวเทติ, ยํ ตํ อริยา อาจิกฺขนฺติ “อุเปกฺขโก สติมา สุขวิหารี”ติ, ตติยํ ฌานํ อุปสมฺปชฺช วิหรติ. สุขสฺส จ ปหานา ทุกฺขสฺส จ ปหานา ปุพฺเพว โสมนสฺส\-โทมนสฺสานํ อตฺถงฺคมา อทุกฺขมสุขํ อุเปกฺขา\-สติ\-ปาริสุทฺธิํ จตุตฺถํ ฌานํ อุปสมฺปชฺช วิหรติ. อยํ วุจฺจติ ภิกฺขเว สมฺมาสมาธิ. อิทํ วุจฺจติ ภิกฺขเว ทุกฺขนิโรธ\-คามินี ปฏิปทา อริยสจฺจํ.}

\proseitemwithcloser{136}{\csromanfolio{90}อิติ อชฺฌตฺตํ วา ธมฺเมสุ ธมฺมานุปสฺสี วิหรติ, พหิทฺธา วา ธมฺเมสุ ธมฺมานุปสฺสี วิหรติ, อชฺฌตฺต\-พหิทฺธา วา ธมฺเมสุ ธมฺมานุปสฺสี วิหรติ. สมุทยธมฺมา\-นุปสฺสี วา ธมฺเมสุ วิหรติ, วยธมฺมา\-นุปสฺสี วา ธมฺเมสุ วิหรติ, สมุทย\-วย\-ธมฺมานุปสฺสี วา ธมฺเมสุ วิหรติ. “อตฺถิ ธมฺมา”ติ วา ปนสฺส สติ ปจฺจุปฏฺฐิตา โหติ ยาวเทว ญาณมตฺตาย ปฏิสฺสติ\-มตฺตาย. อนิสฺสิโต จ วิหรติ, น จ กิญฺจิ โลเก อุปาทิยติ. เอวมฺปิ โข ภิกฺขเว ภิกฺขุ ธมฺเมสุ ธมฺมานุปสฺสี วิหรติ จตูสุ อริยสจฺเจสุ.}{สจฺจปพฺพํ นิฏฺฐิตํ.}
\nitthitam{ธมฺมา\-นุปสฺสนา นิฏฺฐิตา.}
\proseitem{137}{โย หิ โกจิ ภิกฺขเว อิเม จตฺตาโร สติปฏฺฐาเน เอวํ ภาเวยฺย สตฺต วสฺสานิ, ตสฺส ทฺวินฺนํ ผลานํ อญฺญตรํ ผลํ ปาฏิกงฺขํ ทิฏฺเฐว ธมฺเม อญฺญา สติ วา อุปาทิเสเส อนาคามิตา.}

\prose{ติฏฺฐนฺตุ ภิกฺขเว สตฺต วสฺสานิ. โย หิ โกจิ ภิกฺขเว อิเม จตฺตาโร สติปฏฺฐาเน เอวํ ภาเวยฺย ฉ วสฺสานิ ฯเปฯ ปญฺจ วสฺสานิ. จตฺตาริ วสฺสานิ. ตีณิ วสฺสานิ. เทฺว วสฺสานิ. เอกํ วสฺสํ.}

\prose{ติฏฺฐตุ ภิกฺขเว เอกํ วสฺสํ. โย หิ โกจิ ภิกฺขเว อิเม จตฺตาโร สติปฏฺฐาเน เอวํ ภาเวยฺย สตฺต มาสานิ, ตสฺส ทฺวินฺนํ ผลานํ อญฺญตรํ ผลํ ปาฏิกงฺขํ ทิฏฺเฐว ธมฺเม อญฺญา สติ วา อุปาทิเสเส อนาคามิตา.}

\prose{ติฏฺฐนฺตุ ภิกฺขเว สตฺต มาสานิ. โย หิ โกจิ ภิกฺขเว อิเม จตฺตาโร สติปฏฺฐาเน เอวํ ภาเวยฺย ฉ มาสานิ ฯเปฯ ปญฺจ มาสานิ. จตฺตาริ มาสานิ. ตีณิ มาสานิ. เทฺว มาสานิ. เอกํ มาสํ อฑฺฒมาสํ.}

\prose{ติฏฺฐตุ ภิกฺขเว อฑฺฒมาโส. โย หิ โกจิ ภิกฺขเว อิเม จตฺตาโร สติปฏฺฐาเน เอวํ ภาเวยฺย สตฺตาหํ, ตสฺส ทฺวินฺนํ ผลานํ อญฺญตรํ ผลํ ปาฏิกงฺขํ ทิฏฺเฐว ธมฺเม อญฺญา สติ วา อุปาทิเสเส อนาคามิตาติ.}

\csromanmatikaanchor{mtk.37}
\csromantocmarkref{section}{อุทฺทานคาถา}{mtk.37}
\bookmark[dest={mtk.37},level=1]{อุทฺทานคาถา}
\proseitem{138}{\csromanfolio{91}“เอกายโน อยํ ภิกฺขเว มคฺโค สตฺตานํ วิสุทฺธิยา โสก\-ปริเทวานํ สมติกฺกมาย ทุกฺข\-โทมนสฺสานํ อตฺถงฺคมาย ญายสฺส อธิคมาย นิพฺพานสฺส สจฺฉิกิริยาย, ยทิทํ จตฺตาโร สติปฏฺฐานา”ติ อิติ ยํ ตํ วุตฺตํ, อิทเมตํ ปฏิจฺจ วุตฺตนฺติ.}

\prosewithcloser{อิทมโวจ ภควา. อตฺตมนา เต ภิกฺขู ภควโต ภาสิตํ อภินนฺทุนฺติ.}{มหา\-สติปฏฺฐาน\-สุตฺตํ นิฏฺฐิตํ ทสมํ.}
\nitthitammid{มูลปริยาย\-วคฺโค นิฏฺฐิโต ปฐโม.}
\csromancenter{\textbf{ตสฺสุทฺทานํ}\csromansharedfootnote{731b477bf1f64ec1}{อิโต ปรํ เกสุจิ โปตฺถเกสุ อิมาปิ คาถาโย เอวํ ทิสฺสนฺติ\csromandash{} อชรํ อมตํ อมตาธิคมํ, ผลมคฺคนิทสฺสนํ ทุกฺขนุทํ. สหิตตฺตํ มหารสหสฺสกรํ, ภูตมิติ สารํ วิวิธํ สุณาถ. ตฬากํ วสุปูริตํ ฆมฺมปเถ, ติวิธคฺคิปิเฬสิตนิพฺพาปนํ. พฺยาธิปนุทนโอสธโย, ปจฺฉิมสุตฺตปวรา ฐปิตา. มธุมนฺทวรสามทานํ, ขิฑฺฑารติ ชนนิมนุสงฺฆาตํ. ตถา สุตฺเต เวยฺยากรณา ฐปิตา, สกฺยปุตฺตานมภิทมนตฺถาย. ปญฺญาสํ จ ทิยฑฺฒสตํ, เทฺว จ เวยฺยากรณํ อปเร จ. เตวนามคตํ จ อนุปุพฺพํ, เอกมนา นิสาเมถ มุทคฺคํ.}}
\setlength{\csromangathapagewidth}{0pt}
\setlength{\csromangathawakwidth}{0pt}
\csromangathameasuregroup{มูลสุสํวรธมฺมทายาทา, \\ สลฺเลขสมฺมาทิฏฺฐิสติปฏฺฐํ,}{\csromangathabat{มูลสุสํวรธมฺมทายาทา,}{เภรวานงฺคณากงฺเขยฺยวตฺถํ.} \\ \csromangathabat{สลฺเลขสมฺมาทิฏฺฐิสติปฏฺฐํ,}{วคฺควโร อสโม สุสมตฺโต.}}
\csromangathasetleft{มูลสุสํวรธมฺมทายาทา, \\ สลฺเลขสมฺมาทิฏฺฐิสติปฏฺฐํ,}
\csromangathagroup{\csromangathastanza{\csromangathabat{มูลสุสํวรธมฺมทายาทา,}{เภรวานงฺคณากงฺเขยฺยวตฺถํ.} \\ \csromangathabat{สลฺเลขสมฺมาทิฏฺฐิสติปฏฺฐํ,}{วคฺควโร อสโม สุสมตฺโต.}}}

\csromanmatikaanchor{mtk.38}
\csromantocmarknopage{chapter}{2. สีหนาทวคฺค}{mtk.38}
\bookmark[dest={mtk.38},level=0]{2. สีหนาทวคฺค}
\chapterheadpageread{2. \textbf{สีหนาทวคฺค}}
\csromanmatikaanchor{mtk.39}
\csromantocmarkref{subsection}{1. จูฬสีหนาทสุตฺต}{mtk.39}
\bookmark[dest={mtk.39},level=2]{1. จูฬสีหนาทสุตฺต}
\csromanheader{2}{1. จูฬสีหนาท\-สุตฺต}
\proseitem{139}{\csromanfolio{92}เอวํ เม สุตํ\csromandash{}เอกํ สมยํ ภควา สาวตฺถิยํ วิหรติ เชตวเน อนาถ\-ปิณฺฑิกสฺส อาราเม. ตตฺร โข ภควา ภิกฺขู อามนฺเตสิ ภิกฺขโวติ. “ภทนฺเต”ติ เต ภิกฺขู ภควโต ปจฺจสฺโสสุํ. ภควา เอตทโวจ\csromandash{}}

\prose{“อิเธว ภิกฺขเว สมโณ, อิธ ทุติโย สมโณ, อิธ ตติโย สมโณ, อิธ จตุตฺโถ สมโณ, สุญฺญา ปรปฺปวาทา สมเณภิ อญฺเญหี”ติ\csromansharedfootnote{1e49c3a5c2e31b4d}{สมเณหิ อญฺเญติ (สี, อิ, ก) เอตฺถ อญฺเญหีติ สกาย ปฏิญฺญาย สจฺจาภิญฺเญหีติ อตฺโถ เวทิตพฺโพ.}, เอวเมตํ\csromansharedfootnote{9cd2a5d40e3418e7}{เอวเมว (สฺยา, ก)} ภิกฺขเว สมฺมา สีหนาทํ นทถ.}

\proseitem{140}{ฐานํ โข ปเนตํ ภิกฺขเว วิชฺชติ, ยํ อญฺญติตฺถิยา ปริพฺพาชกา เอวํ วเทยฺยุํ “โก ปนา\-ยสฺมนฺตานํ อสฺสาโส กิํ พลํ, เยน ตุเมฺห อายสฺมนฺโต เอวํ วเทถ ‘อิเธว สมโณ, อิธ ทุติโย สมโณ, อิธ ตติโย สมโณ, อิธ จตุตฺโถ สมโณ, สุญฺญา ปรปฺปวาทา สมเณภิ อญฺเญหี’ติ”. เอวํวาทิโน ภิกฺขเว อญฺญติตฺถิยา ปริพฺพาชกา เอวมสฺสุ วจนียา “อตฺถิ โข โน อาวุโส เตน ภควตา ชานตา ปสฺสตา อรหตา สมฺมา\-สมฺพุทฺเธน จตฺตาโร ธมฺมา อกฺขาตา, เย มยํ อตฺตนิ สมฺปสฺสมานา เอวํ วเทม ‘อิเธว สมโณ, อิธ ทุติโย สมโณ, อิธ ตติโย สมโณ, อิธ จตุตฺโถ สมโณ, สุญฺญา ปรปฺปวาทา สมเณภิ อญฺเญหี”ติ. กตเม จตฺตาโร. อตฺถิ โข โน อาวุโส สตฺถริ ปสาโท, อตฺถิ ธมฺเม ปสาโท, อตฺถิ สีเลสุ ปริปูรการิตา, สหธมฺมิกา โข ปน ปิยา มนาปา คหฏฺฐา เจว ปพฺพชิตา จ. อิเม โข โน อาวุโส เตน ภควตา ชานตา ปสฺสตา อรหตา สมฺมา\-สมฺพุทฺเธน จตฺตาโร ธมฺมา อกฺขาตา, เย มยํ อตฺตนิ สมฺปสฺสมานา เอวํ วเทม ‘อิเธว สมโณ, อิธ ทุติโย สมโณ, อิธ ตติโย สมโณ, อิธ จตุตฺโถ สมโณ, สุญฺญา ปรปฺปวาทา สมเณภิ อญฺเญหี”ติ.}

\proseitem{141}{\csromanfolio{93}ฐานํ โข ปเนตํ ภิกฺขเว วิชฺชติ, ยํ อญฺญติตฺถิยา ปริพฺพาชกา เอวํ วเทยฺยุํ “อมฺหากมฺปิ โข อาวุโส อตฺถิ สตฺถริ ปสาโท, โย อมฺหากํ สตฺถา. อมฺหากมฺปิ อตฺถิ ธมฺเม ปสาโท, โย อมฺหากํ ธมฺโม. มยมฺปิ สีเลสุ ปริปูรการิโน, ยานิ อมฺหากํ สีลานิ. อมฺหากมฺปิ สหธมฺมิกา ปิยา มนาปา คหฏฺฐา เจว ปพฺพชิตา จ. อิธ โน อาวุโส โก วิเสโส โก อธิปฺปยาโส\csromansharedfootnote{b22ff14bb1ef04b6}{อธิปฺปาโย (กสี, สฺยา, อิ), อธิปฺปโยโค (ก)} กิํ นานากรณํ, ยทิทํ ตุมฺหากญฺเจว อมฺหากญฺจา”ติ.}

\prose{เอวํวาทิโน ภิกฺขเว อญฺญติตฺถิยา ปริพฺพาชกา เอวมสฺสุ วจนียา “กิํ ปนาวุโส เอกา นิฏฺฐา, อุทาหุ ปุถุ นิฏฺฐา”ติ. สมฺมา พฺยากรมานา ภิกฺขเว อญฺญติตฺถิยา ปริพฺพาชกา เอวํ พฺยากเรยฺยุํ “เอกาวุโส นิฏฺฐา, น ปุถุ นิฏฺฐา”ติ.}

\prose{สา ปนาวุโส นิฏฺฐา สราคสฺส, อุทาหุ วีตราคสฺสาติ. สมฺมา พฺยากรมานา ภิกฺขเว อญฺญติตฺถิยา ปริพฺพาชกา เอวํ พฺยากเรยฺยุํ “วีตราคสฺสา\-วุโส สา นิฏฺฐา, น สา นิฏฺฐา สราคสฺสา”ติ.}

\prose{สา ปนาวุโส นิฏฺฐา สโทสสฺส, อุทาหุ วีตโทสสฺสาติ. สมฺมา พฺยากรมานา ภิกฺขเว อญฺญติตฺถิยา ปริพฺพาชกา เอวํ พฺยากเรยฺยุํ “วีตโทสสฺสา\-วุโส สา นิฏฺฐา, น สา นิฏฺฐา สโทสสฺสา”ติ.}

\prose{สา ปนาวุโส นิฏฺฐา สโมหสฺส, อุทาหุ วีตโมหสฺสาติ. สมฺมา พฺยากรมานา ภิกฺขเว อญฺญติตฺถิยา ปริพฺพาชกา เอวํ พฺยากเรยฺยุํ “วีตโมหสฺสา\-วุโส สา นิฏฺฐา, น สา นิฏฺฐา สโมหสฺสา”ติ.}

\prose{สา ปนาวุโส นิฏฺฐา สตณฺหสฺส, อุทาหุ วีตตณฺหสฺสาติ. สมฺมา พฺยากรมานา ภิกฺขเว อญฺญติตฺถิยา ปริพฺพาชกา เอวํ พฺยากเรยฺยุํ “วีตตณฺหสฺสา\-วุโส สา นิฏฺฐา, น สา นิฏฺฐา สตณฺหสฺสา”ติ.}

\prose{สา ปนาวุโส นิฏฺฐา สอุปาทานสฺส อุทาหุ อนุปาทานสฺสาติ. สมฺมา พฺยากรมานา ภิกฺขเว อญฺญติตฺถิยา ปริพฺพาชกา เอวํ พฺยากเรยฺยุํ “อนุปาทานสฺสา\-วุโส สา นิฏฺฐา, น สา นิฏฺฐา สอุปาทานสฺสา”ติ.}

\prose{\csromanfolio{94}สา ปนาวุโส นิฏฺฐา วิทฺทสุโน, อุทาหุ อวิทฺทสุโนติ. สมฺมา พฺยากรมานา ภิกฺขเว อญฺญติตฺถิยา ปริพฺพาชกา เอวํ พฺยากเรยฺยุํ “วิทฺทสุโน อาวุโส สา นิฏฺฐา, น สา นิฏฺฐา อวิทฺทสุโน”ติ.}

\prose{สา ปนาวุโส นิฏฺฐา อนุรุทฺธ\-ปฺปฏิวิรุทฺธสฺส, อุทาหุ อนนุรุทฺธ- อปฺปฏิวิรุทฺธสฺสาติ. สมฺมา พฺยากรมานา ภิกฺขเว อญฺญติตฺถิยา ปริพฺพาชกา เอวํ พฺยากเรยฺยุํ “อนนุรุทฺธอปฺปฏิวิรุทฺธสฺสา\-วุโส สา นิฏฺฐา, น สา นิฏฺฐา อนุรุทฺธ\-ปฺปฏิวิรุทฺธสฺสา”ติ.}

\prose{สา ปนาวุโส นิฏฺฐา ปปญฺจารามสฺส ปปญฺจรติโน, อุทาหุ นิปฺปปญฺจา\-รามสฺส นิปฺปปญฺจ\-รติโนติ. สมฺมา พฺยากรมานา ภิกฺขเว อญฺญติตฺถิยา ปริพฺพาชกา เอวํ พฺยากเรยฺยุํ “นิปฺปปญฺจา\-รามสฺสา\-วุโส สา นิฏฺฐา นิปฺปปญฺจ\-รติโน, น สา นิฏฺฐา ปปญฺจารามสฺส ปปญฺจรติโน”ติ.}

\proseitem{142}{เทฺวมา ภิกฺขเว ทิฏฺฐิโย ภวทิฏฺฐิ จ วิภวทิฏฺฐิ จ. เย หิ เกจิ ภิกฺขเว สมณา วา พฺราหฺมณา วา ภวทิฏฺฐิํ อลฺลีนา ภวทิฏฺฐิํ อุปคตา ภวทิฏฺฐิํ อชฺโฌสิตา, วิภว\-ทิฏฺฐิยา เต ปฏิวิรุทฺธา. เย หิ เกจิ ภิกฺขเว สมณา วา พฺราหฺมณา วา วิภวทิฏฺฐิํ อลฺลีนา วิภวทิฏฺฐิํ อุปคตา วิภวทิฏฺฐิํ อชฺโฌสิตา, ภวทิฏฺฐิยา เต ปฏิวิรุทฺธา. เย หิ เกจิ ภิกฺขเว สมณา วา พฺราหฺมณา วา อิมาสํ ทฺวินฺนํ ทิฏฺฐีนํ สมุทยญฺจ อตฺถงฺคมญฺจ อสฺสาทญฺจ อาทีนวญฺจ นิสฺสรณญฺจ ยถาภูตํ นปฺปชานนฺติ, เต สราคา, เต สโทสา, เต สโมหา, เต สตณฺหา, เต ส- อุปาทานา, เต อวิทฺทสุโน, เต อนุรุทฺธ\-ปฺปฏิวิรุทฺธา, เต ปปญฺจารามา ปปญฺจรติโน, เต น ปริมุจฺจนฺติ ชาติยา ชราย มรเณน โสเกหิ ปริเทเวหิ ทุกฺเขหิ โทมนสฺเสหิ อุปายาเสหิ, น ปริมุจฺจนฺติ ทุกฺขสฺมาติ วทามิ. เย จ โข เกจิ ภิกฺขเว สมณา วา พฺราหฺมณา วา อิมาสํ ทฺวินฺนํ ทิฏฺฐีนํ สมุทยญฺจ อตฺถงฺคมญฺจ อสฺสาทญฺจ อาทีนวญฺจ นิสฺสรณญฺจ ยถาภูตํ ปชานนฺติ, เต วีตราคา, เต วีตโทสา, เต วีตโมหา, เต วีตตณฺหา, เต อนุปาทานา, เต วิทฺทสุโน, เต อนนุรุทฺธ\-อปฺปฏิวิรุทฺธา, เต นิปฺปปญฺจารามา นิปฺปปญฺจ\-รติโน, เต ปริมุจฺจนฺติ ชาติยา ชราย มรเณน โสเกหิ ปริเทเวหิ ทุกฺเขหิ โทมนสฺเสหิ อุปายาเสหิ, ปริมุจฺจนฺติ ทุกฺขสฺมาติ วทามิ.}

\proseitem{143}{\csromanfolio{95}จตฺตาริมานิ ภิกฺขเว อุปาทานานิ. กตมานิ จตฺตาริ. กามุปาทานํ ทิฏฺฐุปาทานํ สีลพฺพตุ\-ปาทานํ อตฺตวาทุ\-ปาทานํ. สนฺติ ภิกฺขเว เอเก สมณพฺราหฺมณา สพฺพุ\-ปาทาน\-ปริญฺญา\-วาทา ปฏิชานมานา, เต น สมฺมา สพฺพุ\-ปาทาน\-ปริญฺญํ ปญฺญเปนฺติ, กามุปาทานสฺส ปริญฺญํ ปญฺญเปนฺติ, น ทิฏฺฐุ\-ปาทานสฺส ปริญฺญํ ปญฺญเปนฺติ, น สีลพฺพตุ\-ปาทานสฺส ปริญฺญํ ปญฺญเปนฺติ, น อตฺตวาทุ\-ปาทานสฺส ปริญฺญํ ปญฺญเปนฺติ. ตํ กิสฺส เหตุ, อิมานิ หิ เต โภนฺโต สมณพฺราหฺมณา ตีณิ ฐานานิ ยถาภูตํ นปฺปชานนฺติ, ตสฺมา เต โภนฺโต สมณพฺราหฺมณา สพฺพุ\-ปาทาน\-ปริญฺญา\-วาทา ปฏิชานมานา, เต น สมฺมา\csromansharedfootnote{3fa21c6d6bf0ef92}{ปฏิชานมานา น สมฺมา (?)} สพฺพุ\-ปาทาน\-ปริญฺญํ ปญฺญเปนฺติ, กามุปาทานสฺส ปริญฺญํ ปญฺญเปนฺติ, น ทิฏฺฐุ\-ปาทานสฺส ปริญฺญํ ปญฺญเปนฺติ, น สีลพฺพตุ\-ปาทานสฺส ปริญฺญํ ปญฺญเปนฺติ, น อตฺตวาทุ\-ปาทานสฺส ปริญฺญํ ปญฺญเปนฺติ.}

\prose{สนฺติ ภิกฺขเว เอเก สมณพฺราหฺมณา สพฺพุ\-ปาทาน\-ปริญฺญา\-วาทา ปฏิชานมานา, เต น สมฺมา สพฺพุ\-ปาทาน\-ปริญฺญํ ปญฺญเปนฺติ, กามุปาทานสฺส ปริญฺญํ ปญฺญเปนฺติ, ทิฏฺฐุ\-ปาทานสฺส ปริญฺญํ ปญฺญเปนฺติ, น สีลพฺพตุ\-ปาทานสฺส ปริญฺญํ ปญฺญเปนฺติ, น อตฺตวาทุ\-ปาทานสฺส ปริญฺญํ ปญฺญเปนฺติ. ตํ กิสฺส เหตุ, อิมานิ หิ เต โภนฺโต สมณพฺราหฺมณา เทฺว ฐานานิ ยถาภูตํ นปฺปชานนฺติ, ตสฺมา เต โภนฺโต สมณพฺราหฺมณา สพฺพุ\-ปาทาน\-ปริญฺญา\-วาทา ปฏิชานมานา, เต น สมฺมา\csromansharedfootnote{3fa21c6d6bf0ef92}{ปฏิชานมานา น สมฺมา (?)} สพฺพุ\-ปาทาน\-ปริญฺญํ ปญฺญเปนฺติ, กามุปาทานสฺส ปริญฺญํ ปญฺญเปนฺติ, ทิฏฺฐุ\-ปาทานสฺส ปริญฺญํ ปญฺญเปนฺติ, น สีลพฺพตุ\-ปาทานสฺส ปริญฺญํ ปญฺญเปนฺติ, น อตฺตวาทุ\-ปาทานสฺส ปริญฺญํ ปญฺญเปนฺติ.}

\prose{สนฺติ ภิกฺขเว เอเก สมณพฺราหฺมณา สพฺพุ\-ปาทาน\-ปริญฺญา\-วาทา ปฏิชานมานา, เต น สมฺมา สพฺพุ\-ปาทาน\-ปริญฺญํ ปญฺญเปนฺติ, กามุปาทานสฺส ปริญฺญํ ปญฺญเปนฺติ, ทิฏฺฐุ\-ปาทานสฺส ปริญฺญํ ปญฺญเปนฺติ, สีลพฺพตุ\-ปาทานสฺส ปริญฺญํ ปญฺญเปนฺติ, น อตฺตวาทุ\-ปาทานสฺส ปริญฺญํ ปญฺญเปนฺติ. ตํ กิสฺส เหตุ, อิมญฺหิ เต โภนฺโต สมณพฺราหฺมณา เอกํ ฐานํ ยถาภูตํ นปฺปชานนฺติ, ตสฺมา เต โภนฺโต สมณพฺราหฺมณา สพฺพุ\-ปาทาน\-ปริญฺญา\-วาทา ปฏิชานมานา, เต น สมฺมา\csromansharedfootnote{3fa21c6d6bf0ef92}{ปฏิชานมานา น สมฺมา (?)} สพฺพุ\-ปาทาน\-ปริญฺญํ ปญฺญเปนฺติ, กามุปาทานสฺส ปริญฺญํ ปญฺญเปนฺติ, ทิฏฺฐุ\-ปาทานสฺส ปริญฺญํ ปญฺญเปนฺติ, สีลพฺพตุ\-ปาทานสฺส ปริญฺญํ ปญฺญเปนฺติ, น อตฺตวาทุ\-ปาทานสฺส ปริญฺญํ ปญฺญเปนฺติ.}

\prose{\csromanfolio{96}เอวรูเป โข ภิกฺขเว ธมฺมวินเย โย สตฺถริ ปสาโท, โส น สมฺมคฺคโต อกฺขายติ. โย ธมฺเม ปสาโท, โส น สมฺมคฺคโต อกฺขายติ. ยา สีเลสุ ปริปูรการิตา, สา น สมฺมคฺคตา อกฺขายติ. ยา สหธมฺมิเกสุ ปิยมนาปตา, สา น สมฺมคฺคตา อกฺขายติ. ตํ กิสฺส เหตุ, เอวเญฺหตํ ภิกฺขเว โหติ, ยถา ตํ ทุรกฺขาเต ธมฺมวินเย ทุปฺปเวทิเต อนิยฺยานิเก อนุปสมสํวตฺตนิเก อสมฺมาสมฺพุทฺธ\-ปฺปเวทิเต.}

\proseitem{144}{ตถาคโต จ โข ภิกฺขเว อรหํ สมฺมาสมฺพุทฺโธ สพฺพุ\-ปาทาน\-ปริญฺญา\-วาโท ปฏิชานมาโน สมฺมา สพฺพุ\-ปาทาน\-ปริญฺญํ ปญฺญเปติ, กามุปาทานสฺส ปริญฺญํ ปญฺญเปติ, ทิฏฺฐุ\-ปาทานสฺส ปริญฺญํ ปญฺญเปติ, สีลพฺพตุ\-ปาทานสฺส ปริญฺญํ ปญฺญเปติ, อตฺตวาทุ\-ปาทานสฺส ปริญฺญํ ปญฺญเปติ. เอวรูเป โข ภิกฺขเว ธมฺมวินเย โย สตฺถริ ปสาโท, โส สมฺมคฺคโต อกฺขายติ. โย ธมฺเม ปสาโท, โส สมฺมคฺคโต อกฺขายติ. ยา สีเลสุ ปริปูรการิตา, สา สมฺมคฺคตา อกฺขายติ. ยา สหธมฺมิเกสุ ปิยมนาปตา, สา สมฺมคฺคตา อกฺขายติ. ตํ กิสฺส เหตุ, เอวเญฺหตํ ภิกฺขเว โหติ, ยถา ตํ สฺวากฺขาเต ธมฺมวินเย สุปฺปเวทิเต นิยฺยานิเก อุปสม\-สํวตฺตนิเก สมฺมาสมฺพุทฺธ\-ปฺปเวทิเต.}

\proseitem{145}{อิเม จ ภิกฺขเว จตฺตาโร อุปาทานา กิํนิทานา กิํสมุทยา กิํชาติกา กิํปภวา. อิเม จตฺตาโร อุปาทานา ตณฺหานิทานา ตณฺหาสมุทยา ตณฺหาชาติกา ตณฺหาปภวา. ตณฺหา จายํ ภิกฺขเว กิํนิทานา กิํสมุทยา กิํชาติกา กิํปภวา. ตณฺหา เวทนานิทานา เวทนาสมุทยา เวทนาชาติกา เวทนาปภวา. เวทนา จายํ ภิกฺขเว กิํนิทานา กิํสมุทยา กิํชาติกา กิํปภวา. เวทนา ผสฺสนิทานา ผสฺสสมุทยา ผสฺสชาติกา ผสฺสปภวา. ผสฺโส จายํ ภิกฺขเว กิํนิทาโน กิํสมุทโย กิํชาติโก กิํปภโว. ผสฺโส สฬายต\-นนิทาโน สฬายตน\-สมุทโย สฬายตน\-ชาติโก สฬายตน\-ปภโว. สฬายตน\-ญฺจิทํ ภิกฺขเว กิํนิทานํ กิํสมุทยํ กิํชาติกํ กิํปภวํ. สฬายตนํ นามรูป\-นิทานํ นามรูป\-สมุทยํ นามรูป\-ชาติกํ นามรูป\-ปภวํ. นามรูปญฺจิทํ ภิกฺขเว กิํนิทานํ กิํสมุทยํ กิํชาติกํ กิํปภวํ. นามรูปํ วิญฺญาณนิทานํ วิญฺญาณ\-สมุทยํ วิญฺญาณชาติกํ วิญฺญาณ\-ปภวํ. วิญฺญาณญฺจิทํ ภิกฺขเว กิํนิทานํ กิํสมุทยํ \csromanfolio{97}กิํชาติกํ กิํปภวํ. วิญฺญาณํ สงฺขาร\-นิทานํ สงฺขาร\-สมุทยํ สงฺขาร\-ชาติกํ สงฺขาร\-ปภวํ. สงฺขารา จิเม ภิกฺขเว กิํนิทานา กิํสมุทยา กิํชาติกา กิํปภวา. สงฺขารา อวิชฺชานิทานา อวิชฺชาสมุทยา อวิชฺชาชาติกา อวิชฺชาปภวา.}

\prose{ยโต จ โข ภิกฺขเว ภิกฺขุโน อวิชฺชา ปหีนา โหติ, วิชฺชา อุปฺปนฺนา. โส อวิชฺชาวิราคา วิชฺชุปฺปาทา เนว กามุปาทานํ อุปาทิยติ, น ทิฏฺฐุปาทานํ อุปาทิยติ, น สีลพฺพตุ\-ปาทานํ อุปาทิยติ, น อตฺตวาทุ\-ปาทานํ อุปาทิยติ, อนุปาทิยํ น ปริตสฺสติ, อปริตสฺสํ ปจฺจตฺตญฺเญว ปรินิพฺพายติ, “ขีณา ชาติ, วุสิตํ พฺรหฺมจริยํ, กตํ กรณียํ, นาปรํ อิตฺถตฺตายา”ติ ปชานาตีติ.}

\prosewithcloserruled{อิทมโวจ ภควา. อตฺตมนา เต ภิกฺขู ภควโต ภาสิตํ อภินนฺทุนฺติ.}{จูฬสีหนาท\-สุตฺตํ นิฏฺฐิตํ ปฐมํ.}
\csromanmatikaanchor{mtk.40}
\csromantocmarkref{subsection}{2. มหาสีหนาทสุตฺต}{mtk.40}
\bookmark[dest={mtk.40},level=2]{2. มหาสีหนาทสุตฺต}
\csromanheader{2}{2. มหา\-สีหนาท\-สุตฺต}
\proseitem{146}{เอวํ เม สุตํ\csromandash{}เอกํ สมยํ ภควา เวสาลิยํ วิหรติ พหินคเร อปรปุเร วนสณฺเฑ. เตน โข ปน สมเยน สุนกฺขตฺโต ลิจฺฉวิปุตฺโต อจิรปกฺกนฺโต โหติ อิมสฺมา ธมฺมวินยา, โส เวสาลิยํ ปริสติ\csromansharedfootnote{33039af806ccf922}{ปริสติํ (สี, อิ)} เอวํ\csromansharedfootnote{1bcec884a9134fe2}{เอตํ (อิ, ก)} วาจํ ภาสติ “นตฺถิ สมณสฺส โคตมสฺส อุตฺตริ\csromansharedfootnote{37076556bb5c167a}{อุตฺตริํ (อิ)} มนุสฺสธมฺมา อลม\-ริย\-ญาณ\-ทสฺสน\-วิเสโส, ตกฺก\-ปริยาหตํ สมโณ โคตโม ธมฺมํ เทเสติ วีมํสา\-นุจริตํ สยํปฏิภานํ. ยสฺส จ ขฺวาสฺส อตฺถาย ธมฺโม เทสิโต, โส นิยฺยาติ ตกฺกรสฺส สมฺมา ทุกฺขกฺขยายา”ติ.}

\prose{อถ โข อายสฺมา สาริปุตฺโต ปุพฺพณฺหสมยํ นิวาเสตฺวา ปตฺต\-จีวรมา\-ทาย เวสาลิํ ปิณฺฑาย ปาวิสิ. อสฺโสสิ โข อายสฺมา สาริปุตฺโต สุนกฺขตฺตสฺส ลิจฺฉวิ\-ปุตฺตสฺส เวสาลิยํ ปริสติ เอวํ วาจํ ภาสมานสฺส “นตฺถิ สมณสฺส โคตมสฺส อุตฺตริ มนุสฺสธมฺมา อลม\-ริย\-ญาณ\-ทสฺสน\-วิเสโส, ตกฺก\-ปริยาหตํ สมโณ โคตโม ธมฺมํ \csromanfolio{98}เทเสติ วีมํสา\-นุจริตํ สยํปฏิภานํ. ยสฺส จ ขฺวาสฺส อตฺถาย ธมฺโม เทสิโต, โส นิยฺยาติ ตกฺกรสฺส สมฺมา ทุกฺขกฺขยายา”ติ.}

\prose{อถ โข อายสฺมา สาริปุตฺโต เวสาลิยํ ปิณฺฑาย จริตฺวา ปจฺฉาภตฺตํ ปิณฺฑปาต\-ปฏิกฺกนฺโต เยน ภควา เตนุปสงฺกมิ, อุปสงฺกมิตฺวา ภควนฺตํ อภิวาเทตฺวา เอกมนฺตํ นิสีทิ, เอกมนฺตํ นิสินฺโน โข อายสฺมา สาริปุตฺโต ภควนฺตํ เอตทโวจ “สุนกฺขตฺโต ภนฺเต ลิจฺฉวิปุตฺโต อจิรปกฺกนฺโต อิมสฺมา ธมฺมวินยา, โส เวสาลิยํ ปริสติ เอวํ วาจํ ภาสติ ‘นตฺถิ สมณสฺส โคตมสฺส อุตฺตริ มนุสฺสธมฺมา อลม\-ริย\-ญาณ\-ทสฺสน\-วิเสโส, ตกฺก\-ปริยาหตํ สมโณ โคตโม ธมฺมํ เทเสติ วิมํสา\-นุจริตํ สยํปฏิภานํ. ยสฺส จ ขฺวาสฺส อตฺถาย ธมฺโม เทสิโต, โส นิยฺยาติ ตกฺกรสฺส สมฺมา ทุกฺขกฺขยายา”ติ.}

\proseitem{147}{โกธโน เหโส สาริปุตฺต สุนกฺขตฺโต โมฆปุริโส, โกธา จ ปนสฺส เอสา วาจา ภาสิตา. “อวณฺณํ ภาสิสฺสามี”ติ โข สาริปุตฺต สุนกฺขตฺโต โมฆปุริโส วณฺณํเยว ตถาคตสฺส ภาสติ. วณฺโณ เหโส สาริปุตฺต ตถาคตสฺส, โย เอวํ วเทยฺย “ยสฺส จ ขฺวาสฺส อตฺถาย ธมฺโม เทสิโต, โส นิยฺยาติ ตกฺกรสฺส สมฺมา ทุกฺขกฺขยายา”ติ.}

\prose{อยมฺปิ หิ นาม สาริปุตฺต สุนกฺขตฺตสฺส โมฆปุริสสฺส มยิ ธมฺมนฺวโย น ภวิสฺสติ “อิติปิ โส ภควา อรหํ สมฺมาสมฺพุทฺโธ วิชฺชา\-จรณ\-สมฺปนฺโน สุคโต โลกวิทู อนุตฺตโร ปุริส\-ทมฺม\-สารถิ สตฺถา เทวมนุสฺสานํ พุทฺโธ ภควา”ติ.}

\prose{อยมฺปิ หิ นาม สาริปุตฺต สุนกฺขตฺตสฺส โมฆปุริสสฺส มยิ ธมฺมนฺวโย น ภวิสฺสติ “อิติปิ โส ภควา อเนกวิหิตํ อิทฺธิวิธํ ปจฺจนุโภติ, เอโกปิ หุตฺวา พหุธา โหติ, พหุธาปิ หุตฺวา เอโก โหติ, อาวิภาวํ ติโรภาวํ ติโรกุฏฺฏํ ติโรปาการํ ติโรปพฺพตํ อสชฺชมาโน คจฺฉติ เสยฺยถาปิ อากาเส. ปถวิยาปิ อุมฺมุชฺช\-นิมุชฺชํ กโรติ เสยฺยถาปิ อุทเก. อุทเกปิ อภิชฺชมาเน คจฺฉติ เสยฺยถาปิ ปถวิยํ. อากาเสวิ ปลฺลงฺเกน กมติ เสยฺยถาปิ ปกฺขี สกุโณ. อิเมปิ จนฺทิมสูริเย เอวํมหิทฺธิเก เอวํ\-มหา\-นุภาเว ปาณินา ปริมสติ ปริมชฺชติ, ยาวพฺรหฺมโลกาปิ กาเยน วสํ วตฺเตตี”ติ.}

\prose{\csromanfolio{99}อยมฺปิ หิ นาม สาริปุตฺต สุนกฺขตฺตสฺส โมฆปุริสสฺส มยิ ธมฺมนฺวโย น ภวิสฺสติ “อิติปิ โส ภควา ทิพฺพาย โสตธาตุยา วิสุทฺธาย อติกฺกนฺตมานุสิกาย อุโภ สทฺเท สุณาติ ทิพฺเพ จ มานุเส จ เย ทูเร สนฺติเก จา”ติ.}

\prose{อยมฺปิ หิ นาม สาริปุตฺต สุนกฺขตฺตสฺส โมฆปุริสสฺส มยิ ธมฺมนฺวโย น ภวิสฺสติ “อิติปิ โส ภควา ปรสตฺตานํ ปรปุคฺคลานํ เจตสา เจโต ปริจฺจ ปชานาติ. สราคํ วา จิตฺตํ สราคํ จิตฺตนฺติ ปชานาติ, วีตราคํ วา จิตฺตํ วีตราคํ จิตฺตนฺติ ปชานาติ. สโทสํ วา จิตฺตํ สโทสํ จิตฺตนฺติ ปชานาติ, วีตโทสํ วา จิตฺตํ วีตโทสํ จิตฺตนฺติ ปชานาติ. สโมหํ วา จิตฺตํ สโมหํ จิตฺตนฺติ ปชานาติ, วีตโมหํ วา จิตฺตํ วีตโมหํ จิตฺตนฺติ ปชานาติ. สํขิตฺตํ วา จิตฺตํ สํขิตฺตํ จิตฺตนฺติ ปชานาติ, วิกฺขิตฺตํ วา จิตฺตํ วิกฺขิตฺตํ จิตฺตนฺติ ปชานาติ. มหคฺคตํ วา จิตฺตํ มหคฺคตํ จิตฺตนฺติ ปชานาติ, อมหคฺคตํ วา จิตฺตํ อมหคฺคตํ จิตฺตนฺติ ปชานาติ. สอุตฺตรํ วา จิตฺตํ สอุตฺตรํ จิตฺตนฺติ ปชานาติ, อนุตฺตรํ วา จิตฺตํ อนุตฺตรํ จิตฺตนฺติ ปชานาติ. สมาหิตํ วา จิตฺตํ สมาหิตํ จิตฺตนฺติ ปชานาติ, อสมาหิตํ วา จิตฺตํ อสมาหิตํ จิตฺตนฺติ ปชานาติ. วิมุตฺตํ วา จิตฺตํ วิมุตฺตํ จิตฺตนฺติ ปชานาติ, อวิมุตฺตํ วา จิตฺตํ อวิมุตฺตํ จิตฺตนฺติ ปชานาตี”ติ.}

\proseitem{148}{ทส โข ปนิมานิ สาริปุตฺต ตถาคตสฺส ตถาคต\-พลานิ, เยหิ พเลหิ สมนฺนาคโต ตถาคโต อาสภํ ฐานํ ปฏิชานาติ, ปริสาสุ สีหนาทํ นทติ, พฺรหฺมจกฺกํ ปวตฺเตติ. กตมานิ ทส.}

\prose{อิธ สาริปุตฺต ตถาคโต ฐานญฺจ ฐานโต อฏฺฐานญฺจ อฏฺฐานโต ยถาภูตํ ปชานาติ. ยมฺปิ สาริปุตฺต ตถาคโต ฐานญฺจ ฐานโต อฏฺฐานญฺจ อฏฺฐานโต ยถาภูตํ ปชานาติ, อิทมฺปิ สาริปุตฺต ตถาคตสฺส ตถาคตพลํ โหติ, ยํ พลํ อาคมฺม ตถาคโต อาสภํ ฐานํ ปฏิชานาติ, ปริสาสุ สีหนาทํ นทติ, พฺรหฺมจกฺกํ ปวตฺเตติ. (1)}

\prose{ปุน จปรํ สาริปุตฺต ตถาคโต อตีตา\-นาคต\-ปจฺจุปฺปนฺนานํ กมฺม\-สมาทานานํ ฐานโส เหตุโส วิปากํ ยถาภูตํ ปชานาติ. ยมฺปิ สาริปุตฺต ตถาคโต อตีตา\-นาคต\-ปจฺจุปฺปนฺนานํ กมฺม\-สมาทานานํ ฐานโส เหตุโส วิปากํ ยถาภูตํ ปชานาติ, อิทมฺปิ สาริปุตฺต ตถาคตสฺส ตถาคตพลํ โหติ, ยํ พลํ อาคมฺม ตถาคโต \csromanfolio{100}อาสภํ ฐานํ ปฏิชานาติ, ปริสาสุ สีหนาทํ นทติ, พฺรหฺมจกฺกํ ปวตฺเตติ. (2)}

\prose{ปุน จปรํ สาริปุตฺต ตถาคโต สพฺพตฺถ\-คามินิํ ปฏิปทํ ยถาภูตํ ปชานาติ. ยมฺปิ สาริปุตฺต ตถาคโต สพฺพตฺถ\-คามินิํ ปฏิปทํ ยถาภูตํ ปชานาติ, อิทมฺปิ สาริปุตฺต ตถาคตสฺส ตถาคตพลํ โหติ, ยํ พลํ อาคมฺม ตถาคโต อาสภํ ฐานํ ปฏิชานาติ, ปริสาสุ สีหนาทํ นทติ, พฺรหฺมจกฺกํ ปวตฺเตติ. (3)}

\prose{ปุน จปรํ สาริปุตฺต ตถาคโต อเนกธาตุ\-นานาธาตุ\-โลกํ ยถาภูตํ ปชานาติ. ยมฺปิ สาริปุตฺต ตถาคโต อเนกธาตุ\-นานาธาตุ\-โลกํ ยถาภูตํ ปชานาติ, อิทมฺปิ สาริปุตฺต ตถาคตสฺส ตถาคตพลํ โหติ, ยํ พลํ อาคมฺม ตถาคโต อาสภํ ฐานํ ปฏิชานาติ, ปริสาสุ สีหนาทํ นทติ, พฺรหฺมจกฺกํ ปวตฺเตติ. (4)}

\prose{ปุน จปรํ สาริปุตฺต ตถาคโต สตฺตานํ นานา\-ธิมุตฺติกตํ ยถาภูตํ ปชานาติ. ยมฺปิ สาริปุตฺต ตถาคโต สตฺตานํ นานา\-ธิมุตฺติกตํ ยถาภูตํ ปชานาติ, อิทมฺปิ สาริปุตฺต ตถาคตสฺส ตถาคตพลํ โหติ, ยํ พลํ อาคมฺม ตถาคโต อาสภํ ฐานํ ปฏิชานาติ, ปริสาสุ สีหนาทํ นทติ, พฺรหฺมจกฺกํ ปวตฺเตติ. (5)}

\prose{ปุน จปรํ สาริปุตฺต ตถาคโต ปรสตฺตานํ ปรปุคฺคลานํ อินฺทฺริย\-ปโรปริยตฺตํ ยถาภูตํ ปชานาติ. ยมฺปิ สาริปุตฺต ตถาคโต ปรสตฺตานํ ปรปุคฺคลานํ อินฺทฺริย\-ปโรปริยตฺตํ ยถาภูตํ ปชานาติ, อิทมฺปิ สาริปุตฺต ตถาคตสฺส ตถาคตพลํ โหติ, ยํ พลํ อาคมฺม ตถาคโต อาสภํ ฐานํ ปฏิชานาติ, ปริสาสุ สีหนาทํ นทติ, พฺรหฺมจกฺกํ ปวตฺเตติ. (6)}

\prose{ปุน จปรํ สาริปุตฺต ตถาคโต ฌาน\-วิโมกฺข\-สมาธิ\-สมาปตฺตีนํ สํกิเลสํ โวทานํ วุฏฺฐานํ ยถาภูตํ ปชานาติ. ยมฺปิ สาริปุตฺต ตถาคโต ฌาน\-วิโมกฺข\-สมาธิ\-สมาปตฺตีนํ สํกิเลสํ โวทานํ วุฏฺฐานํ ยถาภูตํ ปชานาติ, อิทมฺปิ สาริปุตฺต ตถาคตสฺส ตถาคตพลํ โหติ, ยํ พลํ อาคมฺม ตถาคโต อาสภํ ฐานํ ปฏิชานาติ, ปริสาสุ สีหนาทํ นทติ, พฺรหฺมจกฺขํ ปวตฺเตติ. (7)}

\prose{\csromanfolio{101}ปุน จปรํ สาริปุตฺต ตถาคโต อเนกวิหิตํ ปุพฺเพนิวาสํ อนุสฺสรติ. เสยฺยถิทํ, เอกมฺปิ ชาติํ เทฺวปิ ชาติโย ติสฺโสปิ ชาติโย จตสฺโสปิ ชาติโย ปญฺจปิ ชาติโย ทสปิ ชาติโย วีสมฺปิ ชาติโย ติํสมฺปิ ชาติโย จตฺตาลีสมฺปิ ชาติโย ปญฺญาสมฺปิ ชาติโย ชาติสตมฺปิ ชาติสหสฺสมฺปิ ชาติ\-สต\-สหสฺสมฺปิ อเนเกปิ สํวฏฺฏกปฺเป อเนเกปิ วิวฏฺฏกปฺเป อเนเกปิ สํวฏฺฏ\-วิวฏฺฏ\-กปฺเป “อมุตฺราสิํ เอวํนาโม เอวํโคตฺโต เอวํวณฺโณ เอวมาหาโร เอวํ\-สุขทุกฺข\-ปฺปฏิสํเวที เอวมา\-ยุปริยนฺโต, โส ตโต จุโต อมุตฺร อุทปาทิํ, ตตฺราปาสิํ เอวํนาโม เอวํโคตฺโต เอวํวณฺโณ เอวมาหาโร เอวํ\-สุขทุกฺข\-ปฺปฏิสํเวที เอวมา\-ยุปริยนฺโต, โส ตโต จุโต อิธูปปนฺโน”ติ, อิติ สาการํ สอุทฺเทสํ อเนกวิหิตํ ปุพฺเพนิวาสํ อนุสฺสรติ. ยมฺปิ สาริปุตฺต ตถาคโต อเนกวิหิตํ ปุพฺเพนิวาสํ อนุสฺสรติ. เสยฺยถิทํ, เอกมฺปิ ชาติํ เทฺวปิ ชาติโย ฯเปฯ อิติ สาการํ ส- อุทฺเทสํ อเนกวิหิตํ ปุพฺเพนิวาสํ อนุสฺสรติ, อิทมฺปิ สาริปุตฺต ตถาคตสฺส ตถาคตพลํ โหติ, ยํ พลํ อาคมฺม ตถาคโต อาสภํ ฐานํ ปฏิชานาติ, ปริสาสุ สีหนาทํ นทติ, พฺรหฺมจกฺกํ ปวตฺเตติ. (8)}

\prose{ปุน จปรํ สาริปุตฺต ตถาคโต ทิพฺเพน จกฺขุนา วิสุทฺเธน อติกฺกนฺต\-มานุสเกน สตฺเต ปสฺสติ จวมาเน อุปปชฺชมาเน หีเน ปณีเต สุวณฺเณ ทุพฺพณฺเณ สุคเต ทุคฺคเต, ยถากมฺมูปเค สตฺเต ปชานาติ “อิเม วต โภนฺโต สตฺตา กาย\-ทุจฺจริเตน สมนฺนาคตา วจีทุจฺจริเตน สมนฺนาคตา มโน\-ทุจฺจริเตน สมนฺนาคตา อริยานํ อุปวาทกา มิจฺฉาทิฏฺฐิกา มิจฺฉาทิฏฺฐิ\-กมฺม\-สมาทานา, เต กายสฺส เภทา ปรํ มรณา อปายํ ทุคฺคติํ วินิปาตํ นิรยํ อุปปนฺนา. อิเม วา ปน โภนฺโต สตฺตา กายสุจริเตน สมนฺนาคตา วจีสุจริเตน สมนฺนาคตา มโนสุจริเตน สมนฺนาคตา อริยานํ อนุปวาทกา สมฺมาทิฏฺฐิกา สมฺมา\-ทิฏฺฐิ\-กมฺม\-สมาทานา, เต กายสฺส เภทา ปรํ มรณา สุคติํ สคฺคํ โลกํ อุปปนฺนา”ติ, อิติ ทิพฺเพน จกฺขุนา วิสุทฺเธน อติกฺกนฺต\-มานุสเกน สตฺเต ปสฺสติ จวมาเน อุปปชฺชมาเน หีเน ปณีเต สุวณฺเณ ทุพฺพณฺเณ สุคเต ทุคฺคเต ยถากมฺมูปเค สตฺเต ปชานาติ. ยมฺปิ สาริปุตฺต ตถาคโต ทิพฺเพน จกฺขุนา วิสุทฺเธน อติกฺกนฺต\-มานุสเกน สตฺเต ปสฺสติ จวมาเน อุปปชฺชมาเน}

\prose{\csromanfolio{102}หีเน ปณีเต สุวณฺเณ ทุพฺพณฺเณ สุคเต ทุคฺคเต, ยถากมฺมูปเค สตฺเต ปชานาติ “อิเม วต โภนฺโต สตฺตา กาย\-ทุจฺจริเตน สมนฺนาคตา วจีทุจฺจริเตน สมนฺนาคตา มโน\-ทุจฺจริเตน สมนฺนาคตา อริยานํ อุปวาทกา มิจฺฉาทิฏฺฐิกา มิจฺฉาทิฏฺฐิ\-กมฺม\-สมาทานา, เต กายสฺส เภทา ปรํ มรณา อปายํ ทุคฺคติํ วินิปาตํ นิรยํ อุปปนฺนา. อิเม วา ปน โภนฺโต สตฺตา กายสุจริเตน สมนฺนาคตา วจีสุจริเตน สมนฺนาคตา มโนสุจริเตน สมนฺนาคตา อริยานํ อนุปวาทกา สมฺมาทิฏฺฐิกา สมฺมา\-ทิฏฺฐิ\-กมฺม\-สมาทานา, เต กายสฺส เภทา ปรํ มรณา สุคติํ สคฺคํ โลกํ อุปปนฺนา”ติ, อิติ ทิพฺเพน จกฺขุนา วิสุทฺเธน อติกฺกนฺต\-มานุสเกน สตฺเต ปสฺสติ จวมาเน อุปปชฺชมาเน หีเน ปณีเต สุวณฺเณ ทุพฺพณฺเณ สุคเต ทุคฺคเต ยถากมฺมูปเค สตฺเต ปชานาติ, อิทมฺปิ สาริปุตฺต ตถาคตสฺส ตถาคตพลํ โหติ, ยํ พลํ อาคมฺม ตถาคโต อาสภํ ฐานํ ปฏิชานาติ, ปริสาสุ สีหนาทํ นทติ, พฺรหฺมจกฺกํ ปวตฺเตติ. (9)}

\prose{ปุน จปรํ สาริปุตฺต ตถาคโต อาสวานํ ขยา อนาสวํ เจโตวิมุตฺติํ ปญฺญาวิมุตฺติํ ทิฏฺเฐว ธมฺเม สยํ อภิญฺญา สจฺฉิกตฺวา อุปสมฺปชฺช วิหรติ. ยมฺปิ สาริปุตฺต ตถาคโต อาสวานํ ขยา อนาสวํ เจโตวิมุตฺติํ ปญฺญาวิมุตฺติํ ทิฏฺเฐว ธมฺเม สยํ อภิญฺญา สจฺฉิกตฺวา อุปสมฺปชฺช วิหรติ, อิทมฺปิ สาริปุตฺต ตถาคตสฺส ตถาคตพลํ โหติ, ยํ พลํ อาคมฺม ตถาคโต อาสภํ ฐานํ ปฏิชานาติ, ปริสาสุ สีหนาทํ นทติ, พฺรหฺมจกฺกํ ปวตฺเตติ. (10)}

\prose{อิมานิ โข สาริปุตฺต ทส ตถาคตสฺส ตถาคต\-พลานิ, เยหิ พเลหิ สมนฺนาคโต ตถาคโต อาสภํ ฐานํ ปฏิชานาติ, ปริสาสุ สีหนาทํ นทติ, พฺรหฺมจกฺกํ ปวตฺเตติ.}

\proseitem{149}{โย โข มํ สาริปุตฺต เอวํ ชานนฺตํ เอวํ ปสฺสนฺตํ เอวํ วเทยฺย “นตฺถิ สมณสฺส โคตมสฺส อุตฺตริ มนุสฺสธมฺมา อลม\-ริย\-ญาณ\-ทสฺสน\-วิเสโส, ตกฺก\-ปริยาหตํ สมโณ โคตโม ธมฺมํ เทเสติ วีมํสา\-นุจริตํ สยํปฏิภานนฺ”ติ. ตํ สาริปุตฺต วาจํ อปฺปหาย ตํ จิตฺตํ อปฺปหาย ตํ ทิฏฺฐิํ อปฺปฏินิสฺสชฺชิตฺวา ยถาภตํ นิกฺขิตฺโต เอวํ นิรเย. เสยฺยถาปิ สาริปุตฺต ภิกฺขุ สีลสมฺปนฺโน สมาธิ\-สมฺปนฺโน ปญฺญาสมฺปนฺโน ทิฏฺเฐว \csromanfolio{103}ธมฺเม อญฺญํ อาราเธยฺย, เอวํ สมฺปทมิทํ สาริปุตฺต วทามิ. ตํ วาจํ อปฺปหาย ตํ จิตฺตํ อปฺปหาย ตํ ทิฏฺฐิํ อปฺปฏินิสฺสชฺชิตฺวา ยถาภตํ นิกฺขิตฺโต เอวํ นิรเย.}

\proseitem{150}{จตฺตาริมานิ สาริปุตฺต ตถาคตสฺส เวสารชฺชานิ, เยหิ เวสารชฺเชหิ สมนฺนาคโต ตถาคโต อาสภํ ฐานํ ปฏิชานาติ, ปริสาสุ สีหนาทํ นทติ, พฺรหฺมจกฺกํ ปวตฺเตติ. กตมานิ จตฺตาริ.}

\prose{“สมฺมา\-สมฺพุทฺธสฺส เต ปฏิชานโต อิเม ธมฺมา อนภิสมฺพุทฺธา”ติ ตตฺร วต มํ สมโณ วา พฺราหฺมโณ วา เทโว วา มาโร วา พฺรหฺมา วา โกจิ วา โลกสฺมิํ สหธมฺเมน ปฏิโจเทสฺสตีติ นิมิตฺตเมตํ สาริปุตฺต น สมนุปสฺสามิ, เอตมหํ\csromansharedfootnote{f2c37dd41d04dba5}{เอตมฺปหํ (สี, อิ)} สาริปุตฺต นิมิตฺตํ อสมนุปสฺสนฺโต เขมปฺปตฺโต อภยปฺปตฺโต เวสารชฺช\-ปฺปตฺโต วิหรามิ. (1)}

\prose{“ขีณาสวสฺส เต ปฏิชานโต อิเม อาสวา อปริกฺขีณา”ติ ตตฺร วต มํ สมโณ วา พฺราหฺมโณ วา เทโว วา มาโร วา พฺรหฺมา วา โกจิ วา โลกสฺมิํ สหธมฺเมน ปฏิโจเทสฺสตีติ นิมิตฺตเมตํ สาริปุตฺต น สมนุปสฺสามิ, เอตมหํ สาริปุตฺต นิมิตฺตํ อสมนุปสฺสนฺโต เขมปฺปตฺโต อภยปฺปตฺโต เวสารชฺช\-ปฺปตฺโต วิหรามิ. (2)}

\prose{“เย โข ปน เต อนฺตรายิกา ธมฺมา วุตฺตา, เต ปฏิเสวโต นาลํ อนฺตรายายา”ติ ตตฺร วต มํ สมโณ วา พฺราหฺมโณ วา เทโว วา มาโร วา พฺรหฺมา วา โกจิ วา โลกสฺมิํ สหธมฺเมน ปฏิโจเทสฺสตีติ นิมิตฺตเมตํ สาริปุตฺต น สมนุปสฺสามิ, เอตมหํ สาริปุตฺต นิมิตฺตํ อสมนุปสฺสนฺโต เขมปฺปตฺโต อภยปฺปตฺโต เวสารชฺช\-ปฺปตฺโต วิหรามิ. (3)}

\prose{“ยสฺส โข ปน เต อตฺถาย ธมฺโม เทสิโต, โส น นิยฺยาติ ตกฺกรสฺส สมฺมา ทุกฺขกฺขยายา”ติ ตตฺร วต มํ สมโณ วา พฺราหฺมโณ วา เทโว วา มาโร วา พฺรหฺมา วา โกจิ วา โลกสฺมิํ สหธมฺเมน ปฏิโจเทสฺสตีติ นิมิตฺตเมตํ สาริปุตฺต น สมนุปสฺสามิ. เอตมหํ สาริปุตฺต นิมิตฺตํ อสมนุปสฺสนฺโต เขมปฺปตฺโต อภยปฺปตฺโต เวสารชฺช\-ปฺปตฺโต วิหรามิ. (4)}

\prose{\csromanfolio{104}อิมานิ โข สาริปุตฺต จตฺตาริ ตถาคตสฺส เวสารชฺชานิ, เยหิ เวสารชฺเชหิ สมนฺนาคโต ตถาคโต อาสภํ ฐานํ ปฏิชานาติ, ปริสาสุ สีหนาทํ นทติ, พฺรหฺมจกฺกํ ปวตฺเตติ.}

\prose{โย โข มํ สาริปุตฺต เอวํ ชานนฺตํ เอวํ ปสฺสนฺตํ เอวํ วเทยฺย “นตฺถิ สมณสฺส โคตมสฺส อุตฺตริ มนุสฺสธมฺมา อลม\-ริย\-ญาณ\-ทสฺสน\-วิเสโส, ตกฺก\-ปริยาหตํ สมโณ โคตโม ธมฺมํ เทเสติ วีมํสา\-นุจริตํ สยํปฏิภานนฺ”ติ. ตํ สาริปุตฺต วาจํ อปฺปหาย ตํ จิตฺตํ อปฺปหาย ตํ ทิฏฺฐิํ อปฺปฏินิสฺสชฺชิตฺวา ยถาภตํ นิกฺขิตฺโต เอวํ นิรเย. เสยฺยถาปิ สาริปุตฺต ภิกฺขุ สีลสมฺปนฺโน สมาธิ\-สมฺปนฺโน ปญฺญาสมฺปนฺโน ทิฏฺเฐว ธมฺเม อญฺญํ อาราเธยฺย, เอวํ สมฺปทมิทํ สาริปุตฺต วทามิ. ตํ วาจํ อปฺปหาย ตํ จิตฺตํ อปฺปหาย ตํ ทิฏฺฐิํ อปฺปฏินิสฺสชฺชิตฺวา ยถาภตํ นิกฺขิตฺโต เอวํ นิรเย.}

\proseitem{151}{อฏฺฐ โข อิมา สาริปุตฺต ปริสา. กตมา อฏฺฐ. ขตฺติยปริสา พฺราหฺมณปริสา คหปติปริสา สมณปริสา จาตุมหาราชิก\-ปริสา\csromansharedfootnote{52275cb0f339dc14}{จาตุมฺมหาราชิกา (สี, สฺยา, อิ)} ตาวติํส\-ปริสา มารปริสา พฺรหฺมปริสา. อิมา โข สาริปุตฺต อฏฺฐ ปริสา. อิเมหิ โข สาริปุตฺต จตูหิ เวสารชฺเชหิ สมนฺนาคโต ตถาคโต อิมา อฏฺฐ ปริสา อุปสงฺกมติ อชฺโฌคาหติ. อภิชานามิ โข ปนาหํ สาริปุตฺต อเนกสตํ ขตฺติย\-ปริสํ อุปสงฺกมิตา, ตตฺรปิ มยา สนฺนิสินฺนปุพฺพ\-ญฺเจว สลฺลปิตปุพฺพญฺจ สากจฺฉา จ สมาปชฺชิต\-ปุพฺพา. ตตฺร วต มํ ภยํ วา สารชฺชํ วา โอกฺกมิสฺสตีติ นิมิตฺตเมตํ สาริปุตฺต น สมนุปสฺสามิ, เอตมหํ สาริปุตฺต นิมิตฺตํ อสมนุปสฺสนฺโต เขมปฺปตฺโต อภยปฺปตฺโต เวสารชฺช\-ปฺปตฺโต วิหรามิ.}

\prose{อภิชานามิ โข ปนาหํ สาริปุตฺต อเนกสตํ พฺราหฺมณ\-ปริสํ ฯเปฯ คหปติ\-ปริสํ. สมณปริสํ. จาตุมหาราชิก\-ปริสํ. ตาวติํส\-ปริสํ. มารปริสํ. พฺรหฺมปริสํ อุปสงฺกมิตา, ตตฺรปิ มยา สนฺนิสินฺนปุพฺพ\-ญฺเจว สลฺลปิตปุพฺพญฺจ สากจฺฉา จ สมาปชฺชิต\-ปุพฺพา. ตตฺร วต มํ ภยํ วา สารชฺชํ วา โอกฺกมิสฺสตีติ นิมิตฺตเมตํ สาริปุตฺต น สมนุปสฺสามิ, เอตมหํ สาริปุตฺต นิมิตฺตํ อสมนุปสฺสนฺโต เขมปฺปตฺโต อภยปฺปตฺโต เวสารชฺช\-ปฺปตฺโต วิหรามิ.}

\prose{\csromanfolio{105}โย โข มํ สาริปุตฺต เอวํ ชานนฺตํ เอวํ ปสฺสนฺตํ เอวํ วเทยฺย “นตฺถิ สมณสฺส โคตมสฺส อุตฺตริ มนุสฺสธมฺมา อลม\-ริย\-ญาณ\-ทสฺสน\-วิเสโส, ตกฺก\-ปริยาหตํ สมโณ โคตโม ธมฺมํ เทเสติ วีมํสา\-นุจริตํ สยํปฏิภานนฺ”ติ. ตํ สาริปุตฺต วาจํ อปฺปหาย ตํ จิตฺตํ อปฺปหาย ตํ ทิฏฺฐิํ อปฺปฏินิสฺสชฺชิตฺวา ยถาภตํ นิกฺขิตฺโต เอวํ นิรเย. เสยฺยถาปิ สาริปุตฺต ภิกฺขุ สีลสมฺปนฺโน สมาธิ\-สมฺปนฺโน ปญฺญาสมฺปนฺโน ทิฏฺเฐว ธมฺเม อญฺญํ อาราเธยฺย, เอวํ สมฺปทมิทํ สาริปุตฺต วทามิ. ตํ วาจํ อปฺปหาย ตํ จิตฺตํ อปฺปหาย ตํ ทิฏฺฐิํ อปฺปฏินิสฺสชฺชิตฺวา ยถาภตํ นิกฺขิตฺโต เอวํ นิรเย.}

\proseitem{152}{จตสฺโส โข อิมา สาริปุตฺต โยนิโย. กตมา จตสฺโส. อณฺฑชา โยนิ ชลาพุชา โยนิ สํเสทชา โยนิ โอปปาติกา โยนิ. กตมา จ สาริปุตฺต อณฺฑชา โยนิ. เย โข เต สาริปุตฺต สตฺตา อณฺฑโกสํ อภินิพฺภิชฺช ชายนฺติ. อยํ วุจฺจติ สาริปุตฺต อณฺฑชา โยนิ. กตมา จ สาริปุตฺต ชลาพุชา โยนิ. เย โข เต สาริปุตฺต สตฺตา วตฺถิโกสํ อภินิพฺภิชฺช ชายนฺติ. อยํ วุจฺจติ สาริปุตฺต ชลาพุชา โยนิ. กตมา จ สาริปุตฺต สํเสทชา โยนิ. เย โข เต สาริปุตฺต สตฺตา ปูติมจฺเฉ วา ชายนฺติ ปูติกุณเป วา ปูติกุมฺมาเส วา จนฺทนิกาย วา โอฬิคลฺเล วา ชายนฺติ. อยํ วุจฺจติ สาริปุตฺต สํเสทชา โยนิ. กตมา จ สาริปุตฺต โอปปาติกา โยนิ. เทวา เนรยิกา เอกจฺเจ จ มนุสฺสา เอกจฺเจ จ วินิปาติกา. อยํ วุจฺจติ สาริปุตฺต โอปปาติกา โยนิ. อิมา โข สาริปุตฺต จตสฺโส โยนิโย.}

\prose{โย โข มํ สาริปุตฺต เอวํ ชานนฺตํ เอวํ ปสฺสนฺตํ เอวํ วเทยฺย “นตฺถิ สมณสฺส โคตมสฺส อุตฺตริ มนุสฺสธมฺมา อลม\-ริย\-ญาณ\-ทสฺสน\-วิเสโส, ตกฺก\-ปริยาหตํ สมโณ โคตโม ธมฺมํ เทเสติ วีมํสา\-นุจริตํ สยํปฏิภานนฺ”ติ. ตํ สาริปุตฺต วาจํ อปฺปหาย ตํ จิตฺตํ อปฺปหาย ตํ ทิฏฺฐิํ อปฺปฏินิสฺสชฺชิตฺวา ยถาภตํ นิกฺขิตฺโต เอวํ นิรเย. เสยฺยถาปิ สาริปุตฺต ภิกฺขุ สีลสมฺปนฺโน สมาธิ\-สมฺปนฺโน ปญฺญาสมฺปนฺโน ทิฏฺเฐว ธมฺเม อญฺญํ อาราเธยฺย, เอวํ สมฺปทมิทํ สาริปุตฺต วทามิ. ตํ วาจํ อปฺปหาย ตํ จิตฺตํ อปฺปหาย ตํ ทิฏฺฐิํ อปฺปฏินิสฺสชฺชิตฺวา ยถาภตํ นิกฺขิตฺโต เอวํ นิรเย.}

\proseitem{153}{\csromanfolio{106}ปญฺจ โข อิมา สาริปุตฺต คติโย. กตมา ปญฺจ. นิรโย ติรจฺฉานโยนิ เปตฺติวิสโย มนุสฺสา เทวา. นิรยญฺจาหํ สาริปุตฺต ปชานามิ นิรยคามิญฺจ มคฺคํ นิรยคามินิญฺจ ปฏิปทํ, ยถา ปฏิปนฺโน จ กายสฺส เภทา ปรํ มรณา อปายํ ทุคฺคติํ วินิปาตํ นิรยํ อุปปชฺชติ ตญฺจ ปชานามิ. ติรจฺฉานโยนิ\-ญฺจาหํ สาริปุตฺต ปชานามิ ติรจฺฉานโยนิคามิญฺจ มคฺคํ ติรจฺฉานโยนิคามินิญฺจ ปฏิปทํ, ยถา ปฏิปนฺโน จ กายสฺส เภทา ปรํ มรณา ติรจฺฉาน\-โยนิํ อุปปชฺชติ ตญฺจ ปชานามิ. เปตฺติวิสยํ จาหํ สาริปุตฺต ปชานามิ เปตฺติวิสยคามิญฺจ มคฺคํ เปตฺติวิสยคามินิญฺจ ปฏิปทํ, ยถา ปฏิปนฺโน จ กายสฺส เภทา ปรํ มรณา เปตฺติวิสยํ อุปปชฺชติ ตญฺจ ปชานามิ. มนุสฺเส จาหํ สาริปุตฺต ปชานามิ มนุสฺสโลกคามิญฺจ มคฺคํ มนุสฺสโลกคามินิญฺจ ปฏิปทํ, ยถา ปฏิปนฺโน จ กายสฺส เภทา ปรํ มรณา มนุสฺเสสุ อุปปชฺชติ ตญฺจ ปชานามิ. เทเว จาหํ สาริปุตฺต ปชานามิ เทวโลกคามิญฺจ มคฺคํ เทวโลกคามินิญฺจ ปฏิปทํ, ยถา ปฏิปนฺโน จ กายสฺส เภทา ปรํ มรณา สุคติํ สคฺคํ โลกํ อุปปชฺชติ ตญฺจ ปชานามิ. นิพฺพานญฺจาหํ สาริปุตฺต ปชานามิ นิพฺพานคามิญฺจ มคฺคํ นิพฺพานคามินิญฺจ ปฏิปทํ, ยถา ปฏิปนฺโน จ อาสวานํ ขยา อนาสวํ เจโตวิมุตฺติํ ปญฺญาวิมุตฺติํ ทิฏฺเฐว ธมฺเม สยํ อภิญฺญา สจฺฉิกตฺวา อุปสมฺปชฺช วิหรติ ตญฺจ ปชานามิ.}

\proseitem{154}{อิธาหํ สาริปุตฺต เอกจฺจํ ปุคฺคลํ เอวํ เจตสา เจโต ปริจฺจ ปชานามิ “ตถายํ ปุคฺคโล ปฏิปนฺโน, ตถา จ อิริยติ, ตญฺจ มคฺคํ สมารูโฬฺห, ยถา กายสฺส เภทา ปรํ มรณา อปายํ ทุคฺคติํ วินิปาตํ นิรยํ อุปปชฺชิสฺสตี”ติ. ตเมนํ ปสฺสามิ อปเรน สมเยน ทิพฺเพน จกฺขุนา วิสุทฺเธน อติกฺกนฺต\-มานุสเกน กายสฺส เภทา ปรํ มรณา อปายํ ทุคฺคติํ วินิปาตํ นิรยํ อุปปนฺนํ เอกนฺตทุกฺขา ติพฺพา กฏุกา เวทนา เวทยมานํ. เสยฺยถาปิ สาริปุตฺต องฺคารกาสุ สาธิกโปริสา ปูรา องฺคารานํ วีตจฺจิกานํ วีตธูมานํ. อถ ปุริโส อาคจฺเฉยฺย ฆมฺมาภิตตฺโต ฆมฺมปเรโต กิลนฺโต ตสิโต ปิปาสิโต เอกายเนน มคฺเคน ตเมว องฺคารกาสุํ ปณิธาย. ตเมนํ จกฺขุมา ปุริโส ทิสฺวา เอวํ วเทยฺย “ตถายํ ภวํ ปุริโส ปฏิปนฺโน, ตถา จ อิริยติ, ตญฺจ มคฺคํ สมารูโฬฺห, ยถา อิมํเยว องฺคารกาสุํ อาคมิสฺสตี”ติ. ตเมนํ ปสฺเสยฺย อปเรน สมเยน ตสฺสา \csromanfolio{107}องฺคารกาสุยา ปติตํ เอกนฺตทุกฺขา ติพฺพา กฏุกา เวทนา เวทยมานํ. เอวเมว โข อหํ สาริปุตฺต อิเธกจฺจํ ปุคฺคลํ เอวํ เจตสา เจโต ปริจฺจ ปชานามิ “ตถายํ ปุคฺคโล ปฏิปนฺโน, ตถา จ อิริยติ, ตญฺจ มคฺคํ สมารูโฬฺห, ยถา กายสฺส เภทา ปรํ มรณา อปายํ ทุคฺคติํ วินิปาตํ นิรยํ อุปปชฺชิสฺสตี”ติ. ตเมนํ ปสฺสามิ อปเรน สมเยน ทิพฺเพน จกฺขุนา วิสุทฺเธน อติกฺกนฺต\-มานุสเกน กายสฺส เภทา ปรํ มรณา อปายํ ทุคฺคติํ วินิปาตํ นิรยํ อุปปนฺนํ เอกนฺตทุกฺขา ติพฺพา กฏุกา เวทนา เวทยมานํ. (1)}

\prose{อิธ ปนาหํ สาริปุตฺต เอกจฺจํ ปุคฺคลํ เอวํ เจตสา เจโต ปริจฺจ ปชานามิ “ตถายํ ปุคฺคโล ปฏิปนฺโน, ตถา จ อิริยติ, ตญฺจ มคฺคํ สมารูโฬฺห, ยถา กายสฺส เภทา ปรํ มรณา ติรจฺฉาน\-โยนิํ อุปปชฺชิสฺสตี”ติ. ตเมนํ ปสฺสามิ อปเรน สมเยน ทิพฺเพน จกฺขุนา วิสุทฺเธน อติกฺกนฺต\-มานุสเกน กายสฺส เภทา ปรํ มรณา ติรจฺฉาน\-โยนิํ อุปปนฺนํ ทุกฺขา ติพฺพา กฏุกา เวทนา เวทยมานํ. เสยฺยถาปิ สาริปุตฺต คูถกูโป สาธิกโปริโส ปูโร คูถสฺส. อถ ปุริโส อาคจฺเฉยฺย ฆมฺมาภิตตฺโต ฆมฺมปเรโต กิลนฺโต ตสิโต ปิปาสิโต เอกายเนน มคฺเคน ตเมว คูถกูปํ ปณิธาย. ตเมนํ จกฺขุมา ปุริโส ทิสฺวา เอวํ วเทยฺย “ตถายํ ภวํ ปุริโส ปฏิปนฺโน, ตถา จ อิริยติ, ตญฺจ มคฺคํ สมารูโฬฺห, ยถา อิมํเยว คูถกูปํ อาคมิสฺสตี”ติ. ตเมนํ ปสฺเสยฺย อปเรน สมเยน ตสฺมิํ คูถกูเป ปติตํ ทุกฺขา ติพฺพา กฏุกา เวทนา เวทยมานํ. เอวเมว โข อหํ สาริปุตฺต อิเธกจฺจํ ปุคฺคลํ เอวํ เจตสา เจโต ปริจฺจ ปชานามิ “ตถายํ ปุคฺคโล ปฏิปนฺโน, ตถา จ อิริยติ, ตญฺจ มคฺคํ สมารูโฬฺห, ยถา กายสฺส เภทา ปรํ มรณา ติรจฺฉาน\-โยนิํ อุปปชฺชิสฺสตี”ติ. ตเมนํ ปสฺสามิ อปเรน สมเยน ทิพฺเพน จกฺขุนา วิสุทฺเธน อติกฺกนฺต\-มานุสเกน กายสฺส เภทา ปรํ มรณา ติรจฺฉาน\-โยนิํ อุปปนฺนํ ทุกฺขา ติพฺพา กฏุกา เวทนา เวทยมานํ. (2)}

\prose{อิธ ปนาหํ สาริปุตฺต เอกจฺจํ ปุคฺคลํ เอวํ เจตสา เจโต ปริจฺจ ปชานามิ “ตถายํ ปุคฺคโล ปฏิปนฺโน, ตถา จ อิริยติ, ตญฺจ มคฺคํ สมารูโฬฺห, ยถา กายสฺส เภทา ปรํ มรณา เปตฺติวิสยํ อุปปชฺชิสฺสตี”ติ. \csromanfolio{108}ตเมนํ ปสฺสามิ อปเรน สมเยน ทิพฺเพน จกฺขุนา วิสุทฺเธน อติกฺกนฺต\-มานุสเกน กายสฺส เภทา ปรํ มรณา เปตฺติวิสยํ อุปปนฺนํ ทุกฺขพหุลา เวทนา เวทยมานํ. เสยฺยถาปิ สาริปุตฺต รุกฺโข วิสเม ภูมิภาเค ชาโต ตนุปตฺต\-ปลาโส กพรจฺฉาโย. อถ ปุริโส อาคจฺเฉยฺย ฆมฺมาภิตตฺโต ฆมฺมปเรโต กิลนฺโต ตสิโต ปิปาสิโต เอกายเนน มคฺเคน ตเมว รุกฺขํ ปณิธาย. ตเมนํ จกฺขุมา ปุริโส ทิสฺวา เอวํ วเทยฺย “ตถายํ ภวํ ปุริโส ปฏิปนฺโน, ตถา จ อิริยติ, ตญฺจ มคฺคํ สมารูโฬฺห, ยถา อิมํเยว รุกฺขํ อาคมิสฺสตี”ติ. ตเมนํ ปสฺเสยฺย อปเรน สมเยน ตสฺส รุกฺขสฺส ฉายาย นิสินฺนํ วา นิปนฺนํ วา ทุกฺขพหุลา เวทนา เวทยมานํ. เอวเมว โข อหํ สาริปุตฺต อิเธกจฺจํ ปุคฺคลํ เอวํ เจตสา เจโต ปริจฺจ ปชานามิ “ตถายํ ปุคฺคโล ปฏิปนฺโน, ตถา จ อิริยติ, ตญฺจ มคฺคํ สมารูโฬฺห, ยถา กายสฺส เภทา ปรํ มรณา เปตฺติวิสยํ อุปปชฺชิสฺสตี”ติ. ตเมนํ ปสฺสามิ อปเรน สมเยน ทิพฺเพน จกฺขุนา วิสุทฺเธน อติกฺกนฺต\-มานุสเกน กายสฺส เภทา ปรํ มรณา เปตฺติวิสยํ อุปปนฺนํ ทุกฺขพหุลา เวทนา เวทยมานํ. (3)}

\proseitem{2}{อิธ ปนาหํ สาริปุตฺต เอกจฺจํ ปุคฺคลํ เอวํ เจตสา เจโต ปริจฺจ ปชานามิ “ตถายํ ปุคฺคโล ปฏิปนฺโน, ตถา จ อิริยติ, ตญฺจ มคฺคํ สมารูโฬฺห, ยถา กายสฺส เภทา ปรํ มรณา มนุสฺเสสุ อุปปชฺชิสฺสตี”ติ. ตเมนํ ปสฺสามิ อปเรน สมเยน ทิพฺเพน จกฺขุนา วิสุทฺเธน อติกฺกนฺต\-มานุสเกน กายสฺส เภทา ปรํ มรณา มนุสฺเสสุ อุปปนฺนํ สุขพหุลา เวทนา เวทยมานํ, เสยฺยถาปิ สาริปุตฺต รุกฺโข สเม ภูมิภาเค ชาโต พหล\-ปตฺต\-ปลาโส สนฺทจฺฉาโย\csromansharedfootnote{cf429f4f1c982337}{สณฺฑจฺฉาโย (สฺยา), สนฺตจฺฉาโย (ก)}, อถ ปุริโส อาคจฺเฉยฺย ฆมฺมาภิตตฺโต ฆมฺมปเรโต กิลนฺโต ตสิโต ปิปาสิโต เอกายเนน มคฺเคน ตเมว รุกฺขํ ปณิธาย, ตเมนํ จกฺขุมา ปุริโส ทิสฺวา เอวํ วเทยฺย “ตถายํ ภวํ ปุริโส ปฏิปนฺโน, ตถา จ อิริยติ, ตญฺจ มคฺคํ สมารูโฬฺห, ยถา อิมเมว รุกฺขํ อาคมิสฺสตี”ติ. ตเมนํ ปสฺเสยฺย อปเรน สมเยน ตสฺส รุกฺขสฺส ฉายาย นิสินฺนํ วา นิปนฺนํ วา สุขพหุลา เวทนา เวทยมานํ, เอวเมว โข อหํ สาริปุตฺต อิเธกจฺจํ ปุคฺคลํ เอวํ เจตสา เจโต ปริจฺจ ปชานามิ “ตถายํ ปุคฺคโล ปฏิปนฺโน, ตถา จ อิริยติ, ตญฺจ มคฺคํ สมารูโฬฺห, ยถา กายสฺส เภทา ปรํ \csromanfolio{109}มรณา มนุสฺเสสุ อุปปชฺชิสฺสตี”ติ. ตเมนํ ปสฺสามิ อปเรน สมเยน ทิพฺเพน จกฺขุนา วิสุทฺเธน อติกฺกนฺต\-มานุสเกน กายสฺส เภทา ปรํ มรณา มนุสฺเสสุ อุปปนฺนํ สุขพหุลา เวทนา เวทยมานํ. (4)}

\prose{อิธ ปนาหํ สาริปุตฺต เอกจฺจํ ปุคฺคลํ เอวํ เจตสา เจโต ปริจฺจ ปชานามิ “ตถายํ ปุคฺคโล ปฏิปนฺโน, ตถา จ อิริยติ, ตญฺจ มคฺคํ สมารูโฬฺห, ยถา กายสฺส เภทา ปรํ มรณา สุคติํ สคฺคํ โลกํ อุปปชฺชิสฺสตี”ติ. ตเมนํ ปสฺสามิ อปเรน สมเยน ทิพฺเพน จกฺขุนา วิสุทฺเธน อติกฺกนฺต\-มานุสเกน กายสฺส เภทา ปรํ มรณา สุคติํ สคฺคํ โลกํ อุปปนฺนํ เอกนฺตสุขา เวทนา เวทยมานํ. เสยฺยถาปิ สาริปุตฺต ปาสาโท ตตฺราสฺส กูฏาคารํ อุลฺลิตฺตา\-วลิตฺตํ นิวาตํ ผุสิตคฺคฬํ ปิหิต\-วาตปานํ ตตฺราสฺส ปลฺลงฺโก โคนกตฺถโต ปฏิกตฺถโต ปฏลิกตฺถโต กทลิ\-มิค\-ปวร\-ปจฺจตฺถรโณ สอุตฺตรจฺฉโท อุภโต\-โลหิตกู\-ปธาโน, อถ ปุริโส อาคจฺเฉยฺย ฆมฺมาภิตตฺโต ฆมฺมปเรโต กิลนฺโต ตสิโต ปิปาสิโต เอกายเนน มคฺเคน ตเมว ปาสาทํ ปณิธาย. ตเมนํ จกฺขุมา ปุริโส ทิสฺวา เอวํ วเทยฺย “ตถายํ ภวํ ปุริโส ปฏิปนฺโน, ตถา จ อิริยติ, ตญฺจ มคฺคํ สมารูโฬฺห, ยถา อิมํเยว ปาสาทํ อาคมิสฺสตี”ติ. ตเมนํ ปสฺเสยฺย อปเรน สมเยน ตสฺมิํ ปาสาเท ตสฺมิํ กูฏาคาเร ตสฺมิํ ปลฺลงฺเก นิสินฺนํ วา นิปนฺนํ วา เอกนฺตสุขา เวทนา เวทยมานํ, เอวเมว โข อหํ สาริปุตฺต อิเธกจฺจํ ปุคฺคลํ เอวํ เจตสา เจโต ปริจฺจ ปชานามิ “ตถายํ ปุคฺคโล ปฏิปนฺโน, ตถา จ อิริยติ, ตญฺจ มคฺคํ สมรูโฬฺห, ยถา กายสฺส เภทา ปรํ มรณา สุคติํ สคฺคํ โลกํ อุปปชฺชิสฺสตี”ติ. ตเมนํ ปสฺสามิ อปเรน สมเยน ทิพฺเพน จกฺขุนา วิสุทฺเธน อติกฺกนฺต\-มานุสเกน กายสฺส เภทา ปรํ มรณา สุคติํ สคฺคํ โลกํ อุปปนฺนํ เอกนฺตสุขา เวทนา เวทยมานํ. (5)}

\prose{อิธ ปนาหํ สาริปุตฺต เอกจฺจํ ปุคฺคลํ เจตสา เจโต ปริจฺจ ปชานามิ “ตถายํ ปุคฺคโล ปฏิปนฺโน, ตถา จ อิริยติ, ตญฺจ มคฺคํ สมรูโฬฺห, ยถา อาสวานํ ขยา อนาสวํ เจโตวิมุตฺติํ ปญฺญาวิมุตฺติํ ทิฏฺเฐว ธมฺเม สยํ อภิญฺญา สจฺฉิกตฺวา อุปสมฺปชฺช วิหริสฺสตี”ติ. ตเมนํ ปสฺสามิ อปเรน สมเยน อาสวานํ ขยา อนาสวํ เจโตวิมุตฺติํ ปญฺญาวิมุตฺติํ ทิฏฺเฐว ธมฺเม สยํ อภิญฺญา สจฺฉิกตฺวา อุปสมฺปชฺช วิหรนฺตํ เอกนฺตสุขา เวทนา เวทยมานํ. เสยฺยถาปิ สาริปุตฺต โปกฺขรณี \csromanfolio{110}อจฺโฉทกา สาโตทกา สีโตทกา เสตกา สุปติตฺถา รมณียา, อวิทูเร จสฺสา ติพฺโพ วนสณฺโฑ. อถ ปุริโส อาคจฺเฉยฺย ฆมฺมาภิตตฺโต ฆมฺมปเรโต กิลนฺโต ตสิโต ปิปาสิโต เอกายเนน มคฺเคน ตเมว โปกฺขรณิํ ปณิธาย. ตเมนํ จกฺขุมา ปุริโส ทิสฺวา เอวํ วเทยฺย “ตถา ภวํ ปุริโส ปฏิปนฺโน, ตถา จ อิริยติ, ตญฺจ มคฺคํ สมารูโฬฺห, ยถา อิมํเยว โปกฺขรณิํ อาคมิสฺสตี”ติ. ตเมนํ ปสฺเสยฺย อปเรน สมเยน ตํ โปกฺขรณิํ โอคาเหตฺวา นฺหายิตฺวา จ ปิวิตฺวา จ สพฺพ\-ทรถกิลมถปริฬาหํ ปฏิปฺปสฺสมฺเภตฺวา ปจฺจุตฺตริตฺวา ตสฺมิํ วนสณฺเฑ นิสินฺนํ วา นิปนฺนํ วา เอกนฺตสุขา เวทนา เวทยมานํ, เอวเมว โข อหํ สาริปุตฺต อิเธกจฺจํ ปุคฺคลํ เอวํ เจตสา เจโต ปริจฺจ ปชานามิ “ตถายํ ปุคฺคโล ปฏิปนฺโน ตถา จ อิริยติ ตญฺจ มคฺคํ สมารูโฬฺห, ยถา อาสวานํ ขยา อนาสวํ เจโตวิมุตฺติํ ปญฺญาวิมุตฺติํ ทิฏฺเฐว ธมฺเม สยํ อภิญฺญา สจฺฉิกตฺวา อุปสมฺปชฺช วิหริสฺสตี”ติ. ตเมนํ ปสฺสามิ อปเรน สมเยน อาสวานํ ขยา อนาสวํ เจโตวิมุตฺติํ ปญฺญาวิมุตฺติํ ทิฏฺเฐว ธมฺเม สยํ อภิญฺญา สจฺฉิกตฺวา อุปสมฺปชฺช วิหรนฺตํ เอกนฺตสุขา เวทนา เวทยมานํ. อิมา โข สาริปุตฺต ปญฺจ คติโย.}

\prose{โย โข มํ สาริปุตฺต เอวํ ชานนฺตํ เอวํ ปสฺสนฺตํ เอวํ วเทยฺย “นตฺถิ สมณสฺส โคตมสฺส อุตฺตริ\-มนุสฺสธมฺมา อลม\-ริย\-ญาณ\-ทสฺสน\-วิเสโส, ตกฺก\-ปริยาหตํ สมโณ โคตโม ธมฺมํ เทเสติ วีมํสา\-นุจริตํ สยํปฏิภานนฺ”ติ. ตํ สาริปุตฺต วาจํ อปฺปหาย ตํ จิตฺตํ อปฺปหาย ตํ ทิฏฺฐิํ อปฺปฏินิสฺสชฺชิตฺวา ยถาภตํ นิกฺขิตฺโต เอวํ นิรเย. เสยฺยถาปิ สาริปุตฺต ภิกฺขุ สีลสมฺปนฺโน สมาธิ\-สมฺปนฺโน ปญฺญาสมฺปนฺโน ทิฏฺเฐว ธมฺเม อญฺญํ อาราเธยฺย, เอวํ สมฺปทมิทํ สาริปุตฺต วทามิ. ตํ วาจํ อปฺปหาย ตํ จิตฺตํ อปฺปหาย ตํ ทิฏฺฐิํ อปฺปฏินิสฺสชฺชิตฺวา ยถาภตํ นิกฺขิตฺโต เอวํ นิรเย.}

\proseitem{155}{อภิชานามิ โข ปนาหํ สาริปุตฺต จตุรงฺคสมนฺนาคตํ พฺรหฺมจริยํ จริตา\csromansharedfootnote{964ae718f4d07f9d}{จริตฺวา (ก)}, ตปสฺสี สุทํ โหมิ ปรมตปสฺสี, ลูโข สุทํ\csromansharedfootnote{00170ca83e089c55}{ลูขสฺสุทํ (สี, อิ)} โหมิ ปรมลูโข, เชคุจฺฉี สุทํ โหมิ ปรมเชคุจฺฉี, ปวิวิตฺโต สุทํ\csromansharedfootnote{a1bbe1c7107df2dc}{ปวิวิตฺตสฺสุทํ (สี, อิ)} โหมิ \csromanfolio{111}ปรม\-ปวิวิตฺโต. ตตฺราสฺสุ เม อิทํ สาริปุตฺต ตปสฺสิตาย โหติ. อเจลโก โหมิ มุตฺตาจาโร, หตฺถา\-ปเลขโน\csromansharedfootnote{429ddba78d8d1da6}{หตฺถาวเลขโน (สฺยา)}, เนอ\-หิ\-ภทฺทนฺติโก, นติฏฺฐ\-ภทฺทนฺติโก, นาภิหฏํ, น อุทฺทิสฺสกตํ, น นิมนฺตนํ สาทิยามิ, โส น กุมฺภิมุขา ปฏิคฺคณฺหามิ, น กโฬปิมุขา ปฏิคฺคณฺหามิ, น เอฬกมนฺตรํ, น ทณฺฑมนฺตรํ, น มุสลมนฺตรํ, น ทฺวินฺนํ ภุญฺชมานานํ, น คพฺภินิยา, น ปายมานาย\csromansharedfootnote{2cd20be0c687b1d0}{ปายนฺติยา (ก)}, น ปุริส\-นฺตร\-คตาย, น สงฺกิตฺตีสุ, น ยตฺถ สา อุปฏฺฐิโต โหติ, น ยตฺถ มกฺขิกา สณฺฑ\-สณฺฑ\-จารินี, น มจฺฉํ, น มํสํ, น สุรํ, น เมรยํ, น ถุโสทกํ ปิวามิ. โส เอกาคาริโก วา โหมิ เอกาโลปิโก, ทฺวาคาริโก วา โหมิ ทฺวาโลปิโก ฯเปฯ สตฺตาคาริโก วา โหมิ สตฺตาโลปิโก. เอกิสฺสาปิ ทตฺติยา ยาเปมิ, ทฺวีหิปิ ทตฺตีหิ ยาเปมิ ฯเปฯ สตฺตหิปิ ทตฺตีหิ ยาเปมิ. เอกาหิกมฺปิ อาหารํ อาหาเรมิ, ทฺวีหิกมฺปิ อาหารํ อาหาเรมิ ฯเปฯ สตฺตาหิกมฺปิ อาหารํ อาหาเรมิ. อิติ เอวรูปํ อทฺธมาสิกมฺปิ ปริยาย\-ภตฺตโภชนานุโยคมนุยุตฺโต วิหรามิ.}

\prose{โส สากภกฺโข วา โหมิ, สามากภกฺโข วา โหมิ, นีวารภกฺโข วา โหมิ, ททฺทุลภกฺโข วา โหมิ, หฏภกฺโข วา โหมิ, กณภกฺโข วา โหมิ, อาจามภกฺโข วา โหมิ, ปิญฺญากภกฺโข วา โหมิ, ติณภกฺโข วา โหมิ, โคมยภกฺโข วา โหมิ, วนมูล\-ผลา\-หาโร ยาเปมิ ปวตฺต\-ผล\-โภชี.}

\prose{โส สาณานิปิ ธาเรมิ, มสาณานิปิ ธาเรมิ, ฉวทุสฺสานิปิ ธาเรมิ, ปํสุกูลานิปิ ธาเรมิ, ติรีฏานิปิ ธาเรมิ, อชินมฺปิ ธาเรมิ, อชินกฺขิปมฺปิ ธาเรมิ, กุสจีรมฺปิ ธาเรมิ, วากจีรมฺปิ ธาเรมิ, ผลกจีรมฺปิ ธาเรมิ, เกสกมฺพลมฺปิ ธาเรมิ, วาฬกมฺพลมฺปิ ธาเรมิ, อุลูกปกฺขมฺปิ ธาเรมิ, เกสมสฺสุ\-โลจโกปิ โหมิ เกสมสฺสุ\-โลจนา\-นุโยคม\-นุยุตฺโต, อุพฺภฏฺฐโกปิ โหมิ อาสน\-ปฏิกฺขิตฺโต, อุกฺกุฏิโกปิ โหมิ อุกฺกุฏิก\-ปฺปธานมนุยุตฺโต, กณฺฏกา\-ปสฺสยิ\-โกปิ โหมิ กณฺฏกาปสฺสเย เสยฺยํ กปฺเปมิ\csromansharedfootnote{c3fc08fa56327442}{อิมสฺสานนฺตเร อญฺโญปิ โกจิ ปาฐปเทโส อญฺเญสุ อาชีวกวตทีปกสุตฺเตสุ ทิสฺสติ.}, สาย\-ตติยกมฺปิ อุทโก\-โรหนา\-นุโยคม\-นุยุตฺโต วิหรามิ. อิติ เอวรูปํ อเนกวิหิตํ \csromanfolio{112}กายสฺส อาตาปน\-ปริตาปนา\-นุโยคม\-นุยุตฺโต วิหรามิ. อิทํสุ เม สาริปุตฺต ตปสฺสิตาย โหติ.}

\proseitem{156}{ตตฺราสฺสุ เม อิทํ สาริปุตฺต ลูขสฺมิํ โหติ. เนก\-วสฺส\-คณิกํ รโชชลฺลํ กาเย สนฺนิจิตํ โหติ ปปฏิกชาตํ. เสยฺยถาปิ สาริปุตฺต ตินฺทุกขาณุ เนก\-วสฺส\-คณิโก สนฺนิจิโต โหติ ปปฏิกชาโต, เอวเมวาสฺสุ เม สาริปุตฺต เนก\-วสฺส\-คณิกํ รโชชลฺลํ กาเย สนฺนิจิตํ โหติ ปปฏิกชาตํ. ตสฺส มยฺหํ สาริปุตฺต น เอวํ โหติ. อโห วตาหํ อิมํ รโชชาลฺลํ ปาณินา ปริมชฺเชยฺยํ. อญฺเญ วา ปน เม อิมํ รโชชลฺลํ ปาณินา ปริมชฺเชยฺยุนฺติ. เอวํปิ เม สาริปุตฺต น โหติ. อิทํสุ เม สาริปุตฺต ลูขสฺมิํ โหติ. (1)}

\prose{ตตฺราสฺสุ เม อิทํ สาริปุตฺต เชคุจฺฉิสฺมิํ โหติ. โส โข อหํ สาริปุตฺต สโตว อภิกฺกมามิ สโตว ปฏิกฺกมามิ. ยาว อุทก\-พินฺทุมฺหิปิ เม ทยา ปจฺจุปฏฺฐิตา โหติ “มาหํ ขุทฺทเก ปาเณ วิสมคเต สงฺฆาตํ อาปาเทสินฺ”ติ. อิทํสุ เม สาริปุตฺต เชคุจฺฉิสฺมิํ โหติ. (2)}

\prose{ตตฺราสฺสุ เม อิทํ สาริปุตฺต ปวิวิตฺตสฺมิํ โหติ. โส โข อหํ สาริปุตฺต อญฺญตรํ อรญฺญายตนํ อชฺโฌคาเหตฺวา วิหรามิ. ยทา ปสฺสามิ โคปาลกํ วา ปสุปาลกํ วา ติณหารกํ วา กฏฺฐหารกํ วา วนกมฺมิกํ วา, วเนน วนํ คหเนน คหนํ นินฺเนน นินฺนํ ถเลน ถลํ สํปตามิ\csromansharedfootnote{da44f98678e949ce}{ปปตามิ (สี, สฺยา, อิ)}. ตํ กิสฺส เหตุ, มา มํ เต อทฺทสํสุ, อหญฺจ มา เต อทฺทสนฺติ. เสยฺยถาปิ สาริปุตฺต อารญฺญโก มโค มนุสฺเส ทิสฺวา วเนน วนํ คหเนน คหนํ นินฺเนน นินฺนํ ถเลน ถลํ สํปตติ, เอวเมว โข อหํ สาริปุตฺต ยทา ปสฺสามิ โคปาลกํ วา ปสุปาลกํ วา ติณหารกํ วา กฏฺฐหารกํ วา วนกมฺมิกํ วา, วเนน วนํ คหเนน คหนํ นินฺเนน นินฺนํ ถเลน ถลํ สํปตามิ. ตํ กิสฺส เหตุ, มา มํ เต อทฺทสํสุ, อหญฺจ มา เต อทฺทสนฺติ. อิทํสุ เม สาริปุตฺต ปวิวิตฺตสฺมิํ โหติ. (3)}

\prose{โส โข อหํ สาริปุตฺต เย เต โคฏฺฐา ปฏฺฐิตคาโว อปคต\-โคปาลกา. ตตฺถ จตุกฺกุณฺฑิโก อุปสงฺกมิตฺวา ยานิ ตานิ วจฺฉกานํ ตรุณกานํ เธนุปกานํ โคมยานิ, ตานิ สุทํ อาหาเรมิ. ยาวกีวญฺจ \csromanfolio{113}เม สาริปุตฺต สกํ มุตฺตกรีสํ อปริยาทินฺนํ โหติ. สกํเยว สุทํ มุตฺตกรีสํ อาหาเรมิ. อิทํสุ เม สาริปุตฺต มหา\-วิกฏ\-โภชนสฺมิํ โหติ. (4)}

\proseitem{157}{โส โข อหํ สาริปุตฺต อญฺญตรํ ภิํสนกํ วนสณฺฑํ อชฺโฌคาเหตฺวา วิหรามิ. ตตฺราสฺสุทํ สาริปุตฺต ภิํสนกสฺส วนสณฺฑสฺส ภิํสนกตสฺมิํ โหติ. โย โกจิ อวีตราโค ตํ วนสณฺฑํ ปวิสติ, เยภุยฺเยน โลมานิ หํสนฺติ. โส โข อหํ สาริปุตฺต ยา ตา รตฺติโย สีตา เหมนฺติกา อนฺตรฏฺฐกา หิมปาตสมยา\csromansharedfootnote{6aef122cb4180d0b}{อนฺตรฏฺฐเก หิมปาตสมเย (สี, อิ)}, ตถารูปาสุ รตฺตีสุ รตฺติํ อพฺโภกาเส วิหรามิ, ทิวา วนสณฺเฑ. คิมฺหานํ ปจฺฉิเม มาเส ทิวา อพฺโภกาเส วิหรามิ, รตฺติํ วนสณฺเฑ. อปิสฺสุ มํ สาริปุตฺต อยํ อนจฺฉริย\-คาถา ปฏิภาสิ ปุพฺเพ อสฺสุตปุพฺพา.}

\setlength{\csromangathapagewidth}{0pt}
\setlength{\csromangathawakwidth}{0pt}
\csromangathameasuregroup{“โสตตฺโต โสสินฺโน\textsuperscript{1} เจว, \\ นคฺโค น จคฺคิมาสีโน,}{\csromangathabat{“โสตตฺโต โสสินฺโน\textsuperscript{1} เจว,}{เอโก ภิํสนเก วเน.} \\ \csromangathabat{นคฺโค น จคฺคิมาสีโน,}{เอสนาปสุโต มุนี”ติ.}}
\csromangathasetleft{“โสตตฺโต โสสินฺโน\textsuperscript{1} เจว, \\ นคฺโค น จคฺคิมาสีโน,}
\csromangathagroup{\csromangathastanza{\csromangathabat{“โสตตฺโต โสสินฺโน\csromansharedfootnote{82deeebd64ddcfb7}{โสสีโน (สี, อิ, ก), โสสิโน (สฺยา), โสสินฺโท (สทฺทนีติ)} เจว,}{เอโก ภิํสนเก วเน.} \\ \csromangathabat{นคฺโค น จคฺคิมาสีโน,}{เอสนาปสุโต มุนี”ติ.}}}

\prose{โส โข อหํ สาริปุตฺต สุสาเน เสยฺยํ กปฺเปมิ ฉวฏฺฐิกานิ อุปธาย. อปิสฺสุ มํ สาริปุตฺต คามณฺฑลา\csromansharedfootnote{2a8f821742d9eef6}{โคมณฺฑลา (พหูสุ) จริยาปิฏกอฏฺฐกถา โอโลเกตพฺพา.} อุปสงฺกมิตฺวา โอฏฺฐุภนฺติปิ โอมุตฺเตนฺติปิ ปํสุเกนปิ โอกิรนฺติ กณฺณโสเตสุปิ สลากํ ปเวเสนฺติ. น โข ปนาหํ สาริปุตฺต อภิชานามิ เตสุ ปาปกํ จิตฺตํ อุปฺปาเทตา. อิทํสุ เม สาริปุตฺต อุเปกฺขา\-วิหารสฺมิํ โหติ.}

\proseitem{158}{สนฺติ โข ปน สาริปุตฺต เอเก สมณพฺราหฺมณา เอวํวาทิโน เอวํทิฏฺฐิโน อาหาเรน สุทฺธีติ. เต เอวมาหํสุ, โกเลหิ ยาเปมาติ เต โกลํปิ ขาทนฺติ, โกลจุณฺณํปิ ขาทนฺติ, โกโลทกํปิ ปิวนฺติ, อเนกวิหิตํปิ โกลวิกติํ ปริภุญฺชนฺติ. อภิชานามิ โข ปนาหํ สาริปุตฺต เอกํเยว โกลํ อาหารํ อาหาริตา. สิยา โข ปน เต สาริปุตฺต เอวมสฺส “มหา นูน เตน สมเยน โกโล อโหสี”ติ. น โข ปเนตํ สาริปุตฺต เอวํ ทฏฺฐพฺพํ, ตทาปิ เอตปรโมเยว โกโล อโหสิ, เสยฺยถาปิ เอตรหิ. ตสฺส มยฺหํ สาริปุตฺต \csromanfolio{114}เอกํเยว โกลํ อาหารํ อาหารยโต อธิมตฺต\-กสิมานํ ปตฺโต กาโย โหติ. เสยฺยถาปิ นาม อาสีติกปพฺพานิ วา กาฬปพฺพานิ วา, เอวเมวสฺสุ เม องฺคปจฺจงฺคานิ ภวนฺติ ตาเย\-ว\-ปฺปาหารตาย. เสยฺยถาปิ นาม โอฏฺฐปทํ, เอวเมวสฺสุ เม อานิสทํ โหติ ตาเย\-ว\-ปฺปาหารตาย. เสยฺยถาปิ นาม วฏฺฏนาวฬี, เอวเมวสฺสุ เม ปิฏฺฐิกณฺฏโก อุนฺนตาวนโต โหติ ตาเย\-ว\-ปฺปาหารตาย. เสยฺยถาปิ นาม ชรสาลาย โคปานสิโย โอลุคฺควิลุคฺคา ภวนฺติ, เอวเมวสฺสุ เม ผาสุฬิโย โอลุคฺควิลุคฺคา ภวนฺติ ตาเย\-ว\-ปฺปาหารตาย. เสยฺยถาปิ นาม คมฺภีเร อุทปาเน อุทกตารกา คมฺภีรคตา โอกฺขายิกา ทิสฺสนฺติ. เอวเมวสฺสุ เม อกฺขิกูเปสุ อกฺขิตารกา คมฺภีรคตา โอกฺขายิกา ทิสฺสนฺติ ตาเย\-ว\-ปฺปาหารตาย. เสยฺยถาปิ นาม ติตฺตกาลาพุ อามกจฺฉินฺโน วาตาตเปน สํผุฏิโต\csromansharedfootnote{41bf77644c8941f9}{สมฺผุสิโต (สฺยา), สํปุฏิโต (อิ, ก) เอตฺถ สํผุฏิโตติ สงฺกุจิโตติ อตฺโถ.} โหติ สมฺมิลาโต, เอวเมวสฺสุ เม สีสจฺฉวิ สํผุฏิตา โหติ สมฺมิลาตา ตาเย\-ว\-ปฺปาหารตาย. โส โข อหํ สาริปุตฺต อุทรจฺฉวิํ ปริมสิสฺสามีติ ปิฏฺฐิ\-กณฺฏกํเยว ปริคฺคณฺหามิ. ปิฏฺฐิกณฺฏกํ ปริมสิสฺสามีติ อุทรจฺฉวิํ เยว ปริคฺคณฺหามิ. ยาวสฺสุ เม สาริปุตฺต อุทรจฺฉวิ ปิฏฺฐิกณฺฏกํ อลฺลีนา โหติ ตาเย\-ว\-ปฺปาหารตาย. โส โข อหํ สาริปุตฺต วจฺจํ วา มุตฺตํ วา กริสฺสามีติ ตตฺเถว อวกุชฺโช ปปตามิ ตาเย\-ว\-ปฺปาหารตาย. โส โข อหํ สาริปุตฺต ตเมว กายํ อสฺสาเสนฺโต ปาณินา คตฺตานิ อโนมชฺชามิ. ตสฺส มยฺหํ สาริปุตฺต ปาณินา คตฺตานิ อโนมชฺชโต ปูติมูลานิ โลมานิ กายสฺมา ปตนฺติ ตาเย\-ว\-ปฺปาหารตาย.}

\proseitem{159}{สนฺติ โข ปน สาริปุตฺต เอเก สมณพฺราหฺมณา เอวํวาทิโน เอวํทิฏฺฐิโน อาหาเรน สุทฺธีติ. เต เอวมาหํสุ, มุคฺเคหิ ยาเปม ฯเปฯ ติเลหิ ยาเปม ฯเปฯ ตณฺฑุเลหิ ยาเปมาติ เต ตณฺฑุลํปิ ขาทนฺติ, ตณฺฑุลจุณฺณํปิ ขาทนฺติ, ตณฺฑุโลทกํปิ ปิวนฺติ, อเนกวิหิตํปิ ตณฺฑุลวิกติํ ปริภุญฺชนฺติ. อภิชานามิ โข ปนาหํ สาริปุตฺต เอกํเยว ตณฺฑุลํ อาหารํ อาหาริตา. สิยา โข ปน เต สาริปุตฺต เอวมสฺส “มหา นูน เตน สมเยน ตณฺฑุโล อโหสี”ติ. น โข ปเนตํ สาริปุตฺต เอวํ ทฏฺฐพฺพํ, ตทาปิ เอตปรโมเยว ตณฺฑุโล อโหสิ, \csromanfolio{115}เสยฺยถาปิ เอตรหิ. ตสฺส มยฺหํ สาริปุตฺต เอกํเยว ตณฺฑุลํ อาหารํ อาหารยโต อธิมตฺต\-กสิมานํ ปตฺโต กาโย โหติ. เสยฺยถาปิ นาม อาสีติกปพฺพานิ วา กาฬปพฺพานิ วา, เอวเมวสฺสุ เม องฺคปจฺจงฺคานิ ภวนฺติ ตาเย\-ว\-ปฺปาหารตาย. เสยฺยถาปิ นาม โอฏฺฐปทํ, เอวเมวสฺสุ เม อานิสทํ โหติ ตาเย\-ว\-ปฺปาหารตาย. เสยฺยถาปิ นาม วฏฺฏนาวฬี, เอวเมวสฺสุ เม ปิฏฺฐิกณฺฏโก อุนฺนตาวนโต โหติ ตาเย\-ว\-ปฺปาหารตาย. เสยฺยถาปิ นาม ชรสาลาย โคปานสิโย โอลุคฺควิลุคฺคา ภวนฺติ, เอวเมวสฺสุ เม ผาสุฬิโย โอลุคฺควิลุคฺคา ภวนฺติ ตาเย\-ว\-ปฺปาหารตาย. เสยฺยถาปิ นาม คมฺภีเร อุทปาเน อุทกตารกา คมฺภีรคตา โอกฺขายิกา ทิสฺสนฺติ, เอวเมวสฺสุ เม อกฺขิกูเปสุ อกฺขิตารกา คมฺภีรคตา โอกฺขายิกา ทิสฺสนฺติ ตาเย\-ว\-ปฺปาหารตาย. เสยฺยถาปิ นาม ติตฺตกาลาพุ อามกจฺฉินฺโน วาตาตเปน สํผุฏิโต โหติ สมฺมิลาโต, เอวเมวสฺสุ เม สีสจฺฉวิ สํผุฏิตา โหติ สมฺมิลาตา ตาเย\-ว\-ปฺปาหารตาย. โส โข อหํ สาริปุตฺต อุทรจฺฉวิํ ปริมสิสฺสามีติ ปิฏฺฐิ\-กณฺฏกํเยว ปริคฺคณฺหามิ. ปิฏฺฐิกณฺฏกํ ปริมสิสฺสามีติ อุทร\-จฺฉวิํเยว ปริคฺคณฺหามิ. ยาวสฺสุ เม สาริปุตฺต อุทรจฺฉวิ ปิฏฺฐิกณฺฏกํ อลฺลีนา โหติ ตาเย\-ว\-ปฺปาหารตาย. โส โข อหํ สาริปุตฺต วจฺจํ วา มุตฺตํ วา กริสฺสามีติ ตตฺเถว อวกุชฺโช ปปตามิ ตาเย\-ว\-ปฺปาหารตาย. โส โข อหํ สาริปุตฺต ตเมว กายํ อสฺสาเสนฺโต ปาณินา คตฺตานิ อโนมชฺชามิ. ตสฺส มยฺหํ สาริปุตฺต ปาณินา คตฺตานิ อโนมชฺชโต ปูติมูลานิ โลมานิ กายสฺมา ปตนฺติ ตาเย\-ว\-ปฺปาหารตาย.}

\prose{ตายปิ โข อหํ สาริปุตฺต อิริยาย ตาย ปฏิปทาย ตาย ทุกฺกร\-การิกาย นาชฺฌคมํ อุตฺตริํ มนุสฺสธมฺมา อลม\-ริย\-ญาณ\-ทสฺสน\-วิเสสํ. ตํ กิสฺส เหตุ, อิมิสฺสาเยว อริยาย ปญฺญาย อนธิคมา, ยายํ อริยา ปญฺญา อธิคตา อริยา นิยฺยานิกา นิยฺยาติ ตกฺกรสฺส สมฺมา ทุกฺขกฺขยาย.}

\proseitem{160}{สนฺติ โข ปน สาริปุตฺต เอเก สมณพฺราหฺมณา เอวํวาทิโน เอวํทิฏฺฐิโน “สํสาเรน สุทฺธี”ติ. น โข ปน โส\csromansharedfootnote{433c5a2df9a1f514}{น โข ปเนโส (สี, สฺยา)} สาริปุตฺต สํสาโร \csromanfolio{116}สุลภรูโป, โย มยา อสํสริต\-ปุพฺโพ อิมินา ทีเฆน อทฺธุนา อญฺญตฺร สุทฺธาวาเสหิ เทเวหิ. สุทฺธาวาเส จาหํ สาริปุตฺต เทเว สํสเรยฺยํ. นยิมํ โลกํ ปุนราคจฺเฉยฺยํ.}

\prose{สนฺติ โข ปน สาริปุตฺต เอเก สมณพฺราหฺมณา เอวํวาทิโน เอวํทิฏฺฐิโน “อุปปตฺติยา สุทฺธี”ติ. น โข ปน สา สาริปุตฺต อุปปตฺติ สุลภรูปา, ยา มยา อนุปปนฺน\-ปุพฺพา อิมินา ทีเฆน อทฺธุนา อญฺญตฺร สุทฺธาวาเสหิ เทเวหิ. สุทฺธาวาเส จาหํ สาริปุตฺต เทเว อุปปชฺเชยฺยํ. นยิมํ โลกํ ปุนราคจฺเฉยฺยํ.}

\prose{สนฺติ โข ปน สาริปุตฺต เอเก สมณพฺราหฺมณา เอวํวาทิโน เอวํทิฏฺฐิโน “อาวาเสน สุทฺธี”ติ. น โข ปน โส สาริปุตฺต อาวาโส สุลภรูโป, โย มยา อนาวุฏฺฐปุพฺโพ\csromansharedfootnote{506d8989e280e559}{อนาวุตฺถปุพฺโพ (สี, อิ)} อิมินา ทีเฆน อทฺธุนา อญฺญตฺร สุทฺธาวาเสหิ เทเวหิ. สุทฺธาวาเส จาหํ สาริปุตฺต เทเว อาวเสยฺยํ. นยิมํ โลกํ ปุนราคจฺเฉยฺยํ.}

\prose{สนฺติ โข ปน สาริปุตฺต เอเก สมณพฺราหฺมณา เอวํวาทิโน เอวํทิฏฺฐิโน “ยญฺเญน สุทฺธี”ติ. น โข ปน โส สาริปุตฺต ยญฺโญ สุลภรูโป, โย มยา อยิฏฺฐปุพฺโพ อิมินา ทีเฆน อทฺธุนา. ตญฺจ โข รญฺญา วา สตา ขตฺติเยน มุทฺธา\-วสิตฺเตน พฺราหฺมเณน วา มหาสาเลน.}

\prose{สนฺติ โข ปน สาริปุตฺต เอเก สมณพฺราหฺมณา เอวํวาทิโน เอวํทิฏฺฐิโน “อคฺคิ\-ปริจริยาย สุทฺธี”ติ. น โข ปน โส สาริปุตฺต อคฺคิ สุลภรูโป, โย มยา อปริจิณฺณ\-ปุพฺโพ อิมินา ทีเฆน อทฺธุนา. ตญฺจ โข รญฺญา วา สตา ขตฺติเยน มุทฺธา\-วสิตฺเตน พฺราหฺมเณน วา มหาสาเลน.}

\proseitem{161}{สนฺติ โข ปน สาริปุตฺต เอเก สมณพฺราหฺมณา เอวํวาทิโน เอวํทิฏฺฐิโน “ยาวเทวายํ ภวํ ปุริโส ทหโร โหติ. ยุวา สุสุกาฬเกโส ภเทฺรน โยพฺพเนน สมนฺนาคโต ปฐเมน วยสา, ตาวเทว ปรเมน ปญฺญา\-เวยฺยตฺติเยน สมนฺนาคโต โหติ. ยโต จ โข อยํ ภวํ ปุริโส ชิณฺโณ โหติ วุทฺโธ มหลฺลโก อทฺธคโต วโยอนุปฺปตฺโต อาสีติโก วา นาวุติโก วา วสฺสสติโก วา ชาติยา, อถ ตมฺหา ปญฺญา\-เวยฺยตฺติยา ปริหายตี”ติ. น โข ปเนตํ \csromanfolio{117}สาริปุตฺต เอวํ ทฏฺฐพฺพํ. อหํ โข ปน สาริปุตฺต เอตรหิ ชิณฺโณ วุทฺโธ มหลฺลโก อทฺธคโต วโยอนุปฺปตฺโต อาสีติโก เม วโย วตฺตติ, อิธ เม อสฺสุ สาริปุตฺต จตฺตาโร สาวกา วสฺสสตายุกา วสฺส\-สต\-ชีวิโน ปรมาย สติยา จ คติยา จ ธิติยา จ สมนฺนาคตา ปรเมน จ ปญฺญา\-เวยฺยตฺติเยน. เสยฺยถาปิ สาริปุตฺต ทฬฺหธมฺมา\csromansharedfootnote{a681c4a569337fdf}{ทฬฺหธมฺโม (พหูสุ) ฏีกา จ โมคฺคลฺลานพฺยากรณํ จ โอโลเกตพฺพํ.} ธนุคฺคโห สิกฺขิโต กตหตฺโถ กตูปาสโน ลหุเกน อสเนน อปฺปกสิเรเนว ติริยํ ตาลจฺฉายํ อติปาเตยฺย, เอวํ อธิมตฺต\-สติมนฺโต เอวํ อธิมตฺต\-คติมนฺโต เอวํ อธิมตฺต\-ธิติมนฺโต เอวํ ปรเมน ปญฺญา\-เวยฺยตฺติเยน สมนฺนาคตา. เต มํ จตุนฺนํ สติปฏฺฐานานํ อุปาทายุปาทาย ปญฺหํ ปุจฺเฉยฺยุํ, ปุฏฺโฐ ปุฏฺโฐ จาหํ เตสํ พฺยากเรยฺยํ, พฺยากตญฺจ เม พฺยากตโต ธาเรยฺยุํ, น จ มํ ทุติยกํ อุตฺตริ ปฏิปุจฺเฉยฺยุํ, อญฺญตฺร อสิต\-ปีต\-ขายิต\-สายิตา อญฺญตฺร อุจฺจาร\-ปสฺสาว\-กมฺมา อญฺญตฺร นิทฺทา\-กิลมถ\-ปฏิวิโนทนา. อปริยาทินฺนาเย\-วสฺส สาริปุตฺต ตถาคตสฺส ธมฺมเทสนา, อปริยาทินฺนํเยวสฺส ตถาคตสฺส ธมฺม\-ปทพฺยญฺชนํ, อปริยาทินฺนํเยวสฺส ตถาคตสฺส ปญฺห\-ปฏิภานํ. อถ เม เต จตฺตาโร สาวกา วสฺสสตายุกา วสฺส\-สต\-ชีวิโน วสฺสสตสฺส อจฺจเยน กาลํ กเรยฺยุํ. มญฺจเกน เจปิ มํ สาริปุตฺต ปริหริสฺสถ, เนวตฺถิ ตถาคตสฺส ปญฺญา\-เวยฺยตฺติยสฺส อญฺญถตฺตํ. ยํ โข ตํ\csromansharedfootnote{568d02e9531711a0}{ยํ โข ปเนตํ (สี)} สาริปุตฺต สมฺมา วทมาโน วเทยฺย “อสมฺโมหธมฺโม สตฺโต โลเก อุปฺปนฺโน พหุชนหิตาย พหุชน\-สุขาย โลกานุกมฺปาย อตฺถาย หิตาย สุขาย เทวมนุสฺสานนฺ”ติ, มเมว ตํ สมฺมา วทมาโน วเทยฺย “อสมฺโมหธมฺโม สตฺโต โลเก อุปฺปนฺโน พหุชนหิตาย พหุชน\-สุขาย โลกานุกมฺปาย อตฺถาย หิตาย สุขาย เทวมนุสฺสานนฺ”ติ.}

\proseitem{162}{เตน โข ปน สมเยน อายสฺมา นาคสมาโล ภควโต ปิฏฺฐิโต ฐิโต โหติ ภควนฺตํ พีชยมาโน, อถ โข อายสฺมา นาคสมาโล ภควนฺตํ เอตทโวจ “อจฺฉริยํ ภนฺเต อพฺภุตํ ภนฺเต, อปิ หิ เม ภนฺเต อิมํ ธมฺม\-ปริยายํ สุตฺวา โลมานิ หฏฺฐานิ, โกนาโม อยํ ภนฺเต ธมฺมปริยาโย”ติ. ตสฺมาติห ตฺวํ นาคสมาล อิมํ ธมฺม\-ปริยายํ “โลมหํสน\-ปริยาโย”เตฺวว นํ ธาเรหีติ.}

\prosewithcloserruled{\csromanfolio{118}อิทมโวจ ภควา. อตฺตมโน อายสฺมา นาคสมาโล ภควโต ภาสิตํ อภินนฺทีติ.}{มหา\-สีหนาท\-สุตฺตํ นิฏฺฐิตํ ทุติยํ.}
\csromanmatikaanchor{mtk.41}
\csromantocmarkref{subsection}{3. มหาทุกฺขกฺขนฺธสุตฺต}{mtk.41}
\bookmark[dest={mtk.41},level=2]{3. มหาทุกฺขกฺขนฺธสุตฺต}
\csromanheader{2}{3. มหา\-ทุกฺขกฺขนฺธ\-สุตฺต}
\proseitem{163}{เอวํ เม สุตํ\csromandash{}เอกํ สมยํ ภควา สาวตฺถิยํ วิหรติ เชตวเน อนาถ\-ปิณฺฑิกสฺส อาราเม. อถ โข สมฺพหุลา ภิกฺขู ปุพฺพณฺหสมยํ นิวาเสตฺวา ปตฺต\-จีวรมา\-ทาย สาวตฺถิํ ปิณฺฑาย ปาวิสิํสุ. อถ โข เตสํ ภิกฺขูนํ เอตทโหสิ “อติปฺปโค โข ตาว สาวตฺถิยํ ปิณฺฑาย จริตุํ. ยํ นูน มยํ เยน อญฺญ\-ติตฺถิยานํ ปริพฺพาชกานํ อาราโม เตนุปสงฺกเมยฺยามา”ติ. อถ โข เต ภิกฺขู เยน อญฺญ\-ติตฺถิยานํ ปริพฺพาชกานํ อาราโม เตนุ\-ปสงฺกมิํสุ, อุปสงฺกมิตฺวา เตหิ อญฺญติตฺถิเยหิ ปริพฺพาชเกหิ สทฺธิํ สมฺโมทิํสุ, สมฺโมทนียํ กถํ สารณียํ วีติสาเรตฺวา เอกมนฺตํ นิสีทิํสุ. เอกมนฺตํ นิสินฺเน โข เต ภิกฺขู เต อญฺญติตฺถิยา ปริพฺพาชกา เอตทโวจุํ “สมโณ อาวุโส โคตโม กามานํ ปริญฺญํ ปญฺญเปติ, มยมฺปิ กามานํ ปริญฺญํ ปญฺญเปม. สมโณ อาวุโส โคตโม รูปานํ ปริญฺญํ ปญฺญเปติ, มยมฺปิ รูปานํ ปริญฺญํ ปญฺญเปม. สมโณ อาวุโส โคตโม เวทนานํ ปริญฺญํ ปญฺญเปติ, มยมฺปิ เวทนานํ ปริญฺญํ ปญฺญเปม. อิธ โน อาวุโส โก วิเสโส, โก อธิปฺปยาโส, กิํ นานากรณํ สมณสฺส วา โคตมสฺส อมฺหากํ วา, ยทิทํ ธมฺมเทสนาย วา ธมฺมเทสนํ อนุสาสนิยา วา อนุสาสนินฺ”ติ. อถ โข เต ภิกฺขู เตสํ อญฺญ\-ติตฺถิยานํ ปริพฺพาชกานํ ภาสิตํ เนว อภินนฺทิํสุ นปฺปฏิกฺโกสิํสุ, อนภินนฺทิตฺวา อปฺปฏิกฺโกสิตฺวา อุฏฺฐายาสนา ปกฺกมิํสุ “ภควโต สนฺติเก เอตสฺส ภาสิตสฺส อตฺถํ อาชานิสฺสามา”ติ.}

\proseitem{164}{อถ โข เต ภิกฺขู สาวตฺถิยํ ปิณฺฑาย จริตฺวา ปจฺฉาภตฺตํ ปิณฺฑปาต\-ปฏิกฺกนฺตา เยน ภควา เตนุ\-ปสงฺกมิํสุ, อุปสงฺกมิตฺวา ภควนฺตํ อภิวาเทตฺวา เอกมนฺตํ นิสีทิํสุ, เอกมนฺตํ นิสินฺนา โข เต ภิกฺขู ภควนฺตํ เอตทโวจุํ “อิธ มยํ ภนฺเต ปุพฺพณฺหสมยํ นิวาเสตฺวา ปตฺต\-จีวรมา\-ทาย \csromanfolio{119}สาวตฺถิํ ปิณฺฑาย ปาวิสิมฺห, เตสํ โน ภนฺเต อมฺหากํ เอตทโหสิ ‘อติปฺปโค โข ตาว สาวตฺถิยํ ปิณฺฑาย จริตุํ, ยํ นูน มยํ เยน อญฺญ\-ติตฺถิยานํ ปริพฺพาชกานํ อาราโม เตนุปสงฺกเมยฺยามา’ติ. อถ โข มยํ ภนฺเต เยน อญฺญ\-ติตฺถิยานํ ปริพฺพาชกานํ อาราโม เตนุ\-ปสงฺกมิมฺห, อุปสงฺกมิตฺวา เตหิ อญฺญติตฺถิเยหิ ปริพฺพาชเกหิ สทฺธิํ สมฺโมทิมฺห, สมฺโมทนียํ กถํ สารณียํ วีติสาเรตฺวา เอกมนฺตํ นิสีทิมฺห, เอกมนฺตํ นิสินฺเน โข อเมฺห ภนฺเต เต อญฺญติตฺถิยา ปริพฺพาชกา เอตทโวจุํ ‘สมโณ อาวุโส โคตโม กามานํ ปริญฺญํ ปญฺญเปติ, มยมฺปิ กามานํ ปริญฺญํ ปญฺญเปม. สมโณ อาวุโส โคตโม รูปานํ ปริญฺญํ ปญฺญเปติ, มยมฺปิ รูปานํ ปริญฺญํ ปญฺญเปม. สมโณ อาวุโส โคตโม เวทนานํ ปริญฺญํ ปญฺญเปติ, มยมฺปิ เวทนานํ ปริญฺญํ ปญฺญเปม. อิธ โน อาวุโส โก วิเสโส, โก อธิปฺปยาโส, กิํ นานากรณํ สมณสฺส วา โคตมสฺส อมฺหากํ วา, ยทิทํ ธมฺมเทสนาย วา ธมฺมเทสนํ อนุสาสนิยา วา อนุสาสนินฺ’ติ. อถ โข มยํ ภนฺเต เตสํ อญฺญ\-ติตฺถิยานํ ปริพฺพาชกานํ ภาสิตํ เนว อภินนฺทิมฺห นปฺปฏิกฺโกสิมฺห, อนภินนฺทิตฺวา อปฺปฏิกฺโกสิตฺวา อุฏฺฐายาสนา ปกฺกมิมฺห, ภควโต สนฺติเก เอตสฺส ภาสิตสฺส อตฺถํ อาชานิสฺสามา”ติ.}

\proseitem{165}{เอวํวาทิโน ภิกฺขเว อญฺญติตฺถิยา ปริพฺพาชกา เอวมสฺสุ วจนียา “โก ปนาวุโส กามานํ อสฺสาโท โก อาทีนโว กิํ นิสฺสรณํ, โก รูปานํ อสฺสาโท โก อาทีนโว กิํนิสฺสรณํ, โก เวทนานํ อสฺสาโท โก อาทีนโว กิํ นิสฺสรณนฺ”ติ. เอวํ ปุฏฺฐา ภิกฺขเว อญฺญติตฺถิยา ปริพฺพาชกา น เจว สมฺปายิสฺสนฺติ, อุตฺตริญฺจ วิฆาตํ อาปชฺชิสฺสนฺติ. ตํ กิสฺส เหตุ, ยถา ตํ ภิกฺขเว อวิสยสฺมิํ. นาหํ ตํ ภิกฺขเว ปสฺสามิ สเทวเก โลเก สมารเก สพฺรหฺมเก สสฺสมณ\-พฺราหฺมณิยา ปชาย สเทว\-มนุสฺสาย, โย อิเมสํ ปญฺหานํ เวยฺยากรเณน จิตฺตํ อาราเธยฺย อญฺญตฺร ตถาคเตน วา ตถาคต\-สาวเกน วา อิโต วา ปน สุตฺวา.}

\proseitem{166}{โก จ ภิกฺขเว กามานํ อสฺสาโท. ปญฺจิเม ภิกฺขเว กามคุณา. กตเม ปญฺจ. จกฺขุวิญฺเญยฺยา รูปา อิฏฺฐา กนฺตา มนาปา ปิยรูปา กามูปสํหิตา รชนียา. โสตวิญฺเญยฺยา สทฺทา ฯเปฯ. ฆานวิญฺเญยฺยา \csromanfolio{120}คนฺธา. ชิวฺหาวิญฺเญยฺยา รสา. กายวิญฺเญยฺยา โผฏฺฐพฺพา อิฏฺฐา กนฺตา มนาปา ปิยรูปา กามูปสํหิตา รชนียา. อิเม โข ภิกฺขเว ปญฺจ กามคุณา. ยํ โข ภิกฺขเว อิเม ปญฺจ กามคุเณ ปฏิจฺจ อุปฺปชฺชติ สุขํ โสมนสฺสํ. อยํ กามานํ อสฺสาโท.}

\proseitem{167}{โก จ ภิกฺขเว กามานํ อาทีนโว. อิธ ภิกฺขเว กุลปุตฺโต เยน สิปฺปฏฺฐาเนน ชีวิกํ กปฺเปติ, ยทิ มุทฺทาย ยทิ คณนาย ยทิ สงฺขาเนน\csromansharedfootnote{e798638017a5c596}{สงฺขาย (ก)} ยทิ กสิยา ยทิ วณิชฺชาย ยทิ โครกฺเขน ยทิ อิสฺสตฺเถน ยทิ ราชโปริเสน ยทิ สิปฺปญฺญตเรน, สีตสฺส ปุรกฺขโต อุณฺหสฺส ปุรกฺขโต ฑํสมกส\-วาตา\-ตป\-สรีสป\-สมฺผสฺเสหิ ริสฺสมาโน\csromansharedfootnote{58b349d73078d0f5}{¢รยมาโน (ก), สมฺผสฺสมาโน (ขุ 8. 262)} ขุปฺปิปาสาย มียมาโน. อยมฺปิ ภิกฺขเว กามานํ อาทีนโว สนฺทิฏฺฐิโก ทุกฺขกฺขนฺโธ กามเหตุ กามนิทานํ กามาธิกรณํ กามานเมว เหตุ.}

\prose{ตสฺส เจ ภิกฺขเว กุลปุตฺตสฺส เอวํ อุฏฺฐหโต ฆฏโต วายมโต เต โภคา นาภิ\-นิ\-ปฺผชฺชนฺติ, โส โสจติ กิลมติ ปริเทวติ อุรตฺตาฬิํ กนฺทติ สมฺโมหํ อาปชฺชติ “โมฆํ วต เม อุฏฺฐานํ, อผโล วต เม วายาโม”ติ. อยมฺปิ ภิกฺขเว กามานํ อาทีนโว สนฺทิฏฺฐิโก ทุกฺขกฺขนฺโธ กามเหตุ กามนิทานํ กามาธิกรณํ กามานเมว เหตุ.}

\prose{ตสฺส เจ ภิกฺขเว กุลปุตฺตสฺส เอวํ อุฏฺฐหโต ฆฏโต วายมโต เต โภคา อภินิ\-ปฺผชฺชนฺติ, โส เตสํ โภคานํ อารกฺขาธิกรณํ ทุกฺขํ โทมนสฺสํ ปฏิสํเวเทติ “กินฺติ เม โภเค เนว ราชาโน หเรยฺยุํ, น โจรา หเรยฺยุํ, น อคฺคิ ทเหยฺย, น อุทกํ วเหยฺย\csromansharedfootnote{4490df644bc9633e}{วาเหยฺย (ก)}, น อปฺปิยา ทายาทา หเรยฺยุนฺ”ติ. ตสฺส เอวํ อารกฺขโต โคปยโต เต โภเคราชาโน วา หรนฺติ, โจรา วา หรนฺติ, อคฺคิ วา ทหติ, อุทกํ วา วหติ, อปฺปิยา วา ทายาทา หรนฺติ, โส โสจติ กิลมติ ปริเทวติ อุรตฺตาฬิํ กนฺทติ สมฺโมหํ อาปชฺชติ “ยมฺปิ เม อโหสิ, ตมฺปิ โน นตฺถี”ติ. อยมฺปิ ภิกฺขเว กามานํ อาทีนโว สนฺทิฏฺฐิโก ทุกฺขกฺขนฺโธ กามเหตุ กามนิทานํ กามาธิกรณํ กามานเมว เหตุ.}

\proseitem{168}{ปุน จปรํ ภิกฺขเว กามเหตุ กามนิทานํ กามาธิกรณํ กามานเมว เหตุ ราชาโนปิ ราชูหิ วิวทนฺติ, ขตฺติยาปิ ขตฺติเยหิ \csromanfolio{121}วิวทนฺติ, พฺราหฺมณาปิ พฺราหฺมเณหิ วิวทนฺติ, คหปตีปิ คหปตีหิ วิวทนฺติ, มาตาปิ ปุตฺเตน วิวทติ, ปุตฺโตปิ มาตรา วิวทติ, ปิตาปิ ปุตฺเตน วิวทติ, ปุตฺโตปิ ปิตรา วิวทติ, ภาตาปิ ภาตรา วิวทติ, ภาตาปิ ภคินิยา วิวทติ, ภคินีปิ ภาตรา วิวทติ, สหาโยปิ สหาเยน วิวทติ. เต ตตฺถ กลห\-วิคฺคห\-วิวาทา\-ปนฺนา อญฺญมญฺญํ ปาณีหิปิ อุปกฺกมนฺติ, เลฑฺฑูหิปิ อุปกฺกมนฺติ, ทณฺเฑหิปิ อุปกฺกมนฺติ, สตฺเถหิปิ อุปกฺกมนฺติ. เต ตตฺถ มรณํปิ นิคจฺฉนฺติ มรณมตฺตํปิ ทุกฺขํ. อยมฺปิ ภิกฺขเว กามานํ อาทีนโว สนฺทิฏฺฐิโก ทุกฺขกฺขนฺโธ กามเหตุ กามนิทานํ กามาธิกรณํ กามานเมว เหตุ.}

\prose{ปุน จปรํ ภิกฺขเว กามเหตุ กามนิทานํ กามาธิกรณํ กามานเมว เหตุ อสิจมฺมํ คเหตฺวา ธนุกลาปํ สนฺนยฺหิตฺวา อุภโตพฺยูฬฺหํ สงฺคามํ ปกฺขนฺทนฺติ อุสูสุปิ ขิปฺปมาเนสุ สตฺตีสุปิ ขิปฺปมานาสุ อสีสุปิ วิชฺโชตลนฺเตสุ, เต ตตฺถ อุสูหิปิ วิชฺฌนฺติ, สตฺติยาปิ วิชฺฌนฺติ, อสินาปิ สีสํ ฉินฺทนฺติ. เต ตตฺถ มรณํปิ นิคจฺฉนฺติ มรณมตฺตํปิ ทุกฺขํ. อยมฺปิ ภิกฺขเว กามานํ อาทีนโว สนฺทิฏฺฐิโก ทุกฺขกฺขนฺโธ กามเหตุ กามนิทานํ กามาธิกรณํ กามานเมว เหตุ.}

\prose{ปุน จปรํ ภิกฺขเว กามเหตุ กามนิทานํ กามาธิกรณํ กามานเมว เหตุ อสิจมฺมํ คเหตฺวา ธนุกลาปํ สนฺนยฺหิตฺวา อทฺทาวเลปนา\csromansharedfootnote{8df5c9b8ff0f0132}{อฏฺฏาวเลปนา (สฺยา, ก)} อุปการิโย ปกฺขนฺทนฺติ อุสูสุปิ ขิปฺปมาเนสุ สตฺตีสุปิ ขิปฺปมานาสุ อสีสุปิ วิชฺโชตลนฺเตสุ. เต ตตฺถ อุสูหิปิ วิชฺฌนฺติ, สตฺติยาปิ วิชฺฌนฺติ, ฉกณกายปิ\csromansharedfootnote{b4d1f370f4de1b34}{ปกฏฺฐิยาปิ (สี)} โอสิญฺจนฺติ, อภิวคฺเคนปิ โอมทฺทนฺติ, อสินาปิ สีสํ ฉินฺทนฺติ. เต ตตฺถ มรณํปิ นิคจฺฉนฺติ มรณมตฺตํปิ ทุกฺขํ. อยมฺปิ ภิกฺขเว กามานํ อาทีนโว สนฺทิฏฺฐิโก ทุกฺขกฺขนฺโธ กามเหตุ กามนิทานํ กามาธิกรณํ กามานเมว เหตุ.}

\proseitem{169}{ปุน จปรํ ภิกฺขเว กามเหตุ กามนิทานํ กามาธิกรณํ กามานเมว เหตุ สนฺธิํปิ ฉินฺทนฺติ, นิลฺโลปํปิ หรนฺติ, เอกาคาริกมฺปิ กโรนฺติ, ปริปนฺเถปิ ติฏฺฐนฺติ, ปรทารํปิ คจฺฉนฺติ. ตเมนํ ราชาโน คเหตฺวา วิวิธา กมฺมการณา กาเรนฺติ, กสาหิปิ ตาเฬนฺติ, เวตฺเตหิปิ ตาเฬนฺติ, อฑฺฒ\-ทณฺฑเกหิปิ ตาเฬนฺติ, หตฺถํปิ ฉินฺทนฺติ, ปาทํปิ ฉินฺทนฺติ, หตฺถปาทํปิ ฉินฺทนฺติ, กณฺณํปิ ฉินฺทนฺติ, นาสํปิ ฉินฺทนฺติ, กณฺณนาสํปิ ฉินฺทนฺติ, พิลงฺค\-ถาลิกํปิ \csromanfolio{122}กโรนฺติ, สงฺข\-มุณฺฑิกํปิ กโรนฺติ, ราหุมุขํปิ กโรนฺติ, โชติมาลิกํปิ กโรนฺติ, หตฺถปชฺโชติกํปิ กโรนฺติ, เอรกวตฺติกํปิ กโรนฺติ, จีรก\-วาสิกํปิ กโรนฺติ, เอเณยฺยกํปิ กโรนฺติ, พฬิส\-มํสิกํปิ กโรนฺติ, กหาปณิกํปิ กโรนฺติ, ขาราปตจฺฉิกํปิ กโรนฺติ, ปลิฆปริวตฺติกํปิ กโรนฺติ, ปลาล\-ปีฐกํปิ กโรนฺติ, ตตฺเตนปิ เตเลน โอสิญฺจนฺติ, สุนเขหิปิ ขาทาเปนฺติ, ชีวนฺตํปิ สูเล อุตฺตาเสนฺติ, อสินาปิ สีสํ ฉินฺทนฺติ, เต ตตฺถ มรณํปิ นิคจฺฉนฺติ มรณมตฺตํปิ ทุกฺขํ. อยมฺปิ ภิกฺขเว กามานํ อาทีนโว สนฺทิฏฺฐิโก ทุกฺขกฺขนฺโธ กามเหตุ กามนิทานํ กามาธิกรณํ กามานเมว เหตุ.}

\prose{ปุน จปรํ ภิกฺขเว กามเหตุ กามนิทานํ กามาธิกรณํ กามานเมว เหตุ กาเยน ทุจฺจริตํ จรนฺติ, วาจาย ทุจฺจริตํ จรนฺติ, มนสา ทุจฺจริตํ จรนฺติ. เต กาเยน ทุจฺจริตํ จริตฺวา วาจาย ทุจฺจริตํ จริตฺวา มนสา ทุจฺจริตํ จริตฺวา กายสฺส เภทา ปรํ มรณา อปายํ ทุคฺคติํ วินิปาตํ นิรยํ อุปปชฺชนฺติ. อยมฺปิ ภิกฺขเว กามานํ อาทีนโว สมฺปรายิโก ทุกฺขกฺขนฺโธ กามเหตุ กามนิทานํ กามาธิกรณํ กามานเมว เหตุ.}

\proseitem{170}{กิญฺจ ภิกฺขเว กามานํ นิสฺสรณํ. โย โข ภิกฺขเว กาเมสุ ฉนฺทราค\-วินโย ฉนฺท\-ราค\-ปฺปหานํ. อิทํ กามานํ นิสฺสรณํ.}

\prose{เย หิ เกจิ ภิกฺขเว สมณา วา พฺราหฺมณา วา เอวํ กามานํ อสฺสาทญฺจ อสฺสาทโต อาทีนวญฺจ อาทีนวโต นิสฺสรณญฺจ นิสฺสรณโต ยถาภูตํ นปฺปชานนฺติ, เต วต สามํ วา กาเม ปริชานิสฺสนฺติ, ปรํ วา ตถตฺตาย สมาทเปสฺสนฺติ, ยถา ปฏิปนฺโน กาเม ปริชานิสฺสตีติ เนตํ ฐานํ วิชฺชติ. เย จ โข เกจิ ภิกฺขเว สมณา วา พฺราหฺมณา วา เอวํ กามานํ อสฺสาทญฺจ อสฺสาทโต อาทีนวญฺจ อาทีนวโต นิสฺสรณญฺจ นิสฺสรณโต ยถาภูตํ ปชานนฺติ, เต วต สามํ วา กาเม ปริชานิสฺสนฺติ, ปรํ วา ตถตฺตาย สมาทเปสฺสนฺติ, ยถา ปฏิปนฺโน กาเม ปริชานิสฺสตีติ ฐานเมตํ วิชฺชติ.}

\proseitem{171}{โก จ ภิกฺขเว รูปานํ อสฺสาโท. เสยฺยถาปิ ภิกฺขเว ขตฺติยกญฺญา วา พฺราหฺมณกญฺญา วา คหปติกญฺญา วา ปนฺนรส\-วสฺสุทฺเทสิกา วา โสฬส\-วสฺสุทฺเทสิกา วา นาติทีฆา นาติรสฺสา นาติกิสา นาติถูลา นาติกาฬี นาจฺโจทาตา, ปรมา สา ภิกฺขเว ตสฺมิํ สมเย \csromanfolio{123}สุภา วณฺณนิภาติ. เอวํ ภนฺเต. ยํ โข ภิกฺขเว สุภํ วณฺณนิภํ ปฏิจฺจ อุปฺปชฺชติ สุขํ โสมนสฺสํ. อยํ รูปานํ อสฺสาโท.}

\prose{โก จ ภิกฺขเว รูปานํ อาทีนโว. อิธ ภิกฺขเว ตเมว ภคินิํ ปสฺเสยฺย อปเรน สมเยน อาสีติกํ วา นาวุติกํ วา วสฺสสติกํ วา ชาติยา ชิณฺณํ โคปานสิวงฺกํ โภคฺคํ ทณฺฑปรายนํ ปเวธมานํ คจฺฉนฺติํ อาตุรํ คตโยพฺพนํ ขณฺฑทนฺตํ\csromansharedfootnote{b9a094a94a12f0af}{ขณฺฑทนฺติํ (สี, อิ)} ปลิตเกสํ\csromansharedfootnote{293086333d8fa0e9}{ปลิตเกสิํ,} วิลูนํ ขลิตสิรํ วลินํ ติลกา\-หต\-คตฺตํ\csromansharedfootnote{cffbe1e9df11549b}{ติลกาหตคตฺติํ (พหูสุ) อฏฺฐกถา ฏีกา โอโลเกตพฺพา.}. ตํ กิํ มญฺญถ ภิกฺขเว, ยา ปุริมา สุภา วณฺณนิภา, สา อนฺตรหิตา, อาทีนโว ปาตุภูโตติ. เอวํ ภนฺเต. อยมฺปิ ภิกฺขเว รูปานํ อาทีนโว.}

\prose{ปุน จปรํ ภิกฺขเว ตเมว ภคินิํ ปสฺเสยฺย อาพาธิกํ ทุกฺขิตํ พาฬฺหคิลานํ สเก มุตฺตกรีเส ปลิปนฺนํ เสมานํ\csromansharedfootnote{501db83a81a3c69d}{เสยฺยมานํ (ก)} อญฺเญหิ วุฏฺฐาปิยมานํ อญฺเญหิ สํเวสิยมานํ. ตํ กิํ มญฺญถ ภิกฺขเว, ยา ปุริมา สุภา วณฺณนิภา, สา อนฺตรหิตา, อาทีนโว ปาตุภูโตติ. เอวํ ภนฺเต. อยมฺปิ ภิกฺขเว รูปานํ อาทีนโว.}

\proseitem{172}{ปุน จปรํ ภิกฺขเว ตเมว ภคินิํ ปสฺเสยฺย สรีรํ สิวถิกาย ฉฑฺฑิตํ เอกาหมตํ วา ทฺวีหมตํ วา ตีหมตํ วา อุทฺธุมาตกํ วินีลกํ วิปุพฺพกชาตํ. ตํ กิํ มญฺญถ ภิกฺขเว, ยา ปุริมา สุภา วณฺณนิภา, สา อนฺตรหิตา, อาทีนโว ปาตุภูโตติ. เอวํ ภนฺเต. อยมฺปิ ภิกฺขเว รูปานํ อาทีนโว.}

\prose{ปุน จปรํ ภิกฺขเว ตเมว ภคินิํ ปสฺเสยฺย สรีรํ สิวถิกาย ฉฑฺฑิตํ กาเกหิ วา ขชฺชมานํ กุลเลหิ วา ขชฺชมานํ คิชฺเฌหิ วา ขชฺชมานํ กงฺเกหิ วา ขชฺชมานํ สุนเขหิ วา ขชฺชมานํ พฺยคฺเฆหิ วา ขชฺชมานํ ทีปีหิ วา ขชฺชมานํ สิงฺคาเลหิ วา ขชฺชมานํ วิวิเธหิ วา ปาณกชาเตหิ ขชฺชมานํ. ตํ กิํ มญฺญถ ภิกฺขเว, ยา ปุริมา สุภา วณฺณนิภา, สา อนฺตรหิตา, อาทีนโว ปาตุภูโตติ. เอวํ ภนฺเต. อยมฺปิ ภิกฺขเว รูปานํ อาทีนโว.}

\prose{ปุน จปรํ ภิกฺขเว ตเมว ภคินิํ ปสฺเสยฺย สรีรํ สิวถิกาย ฉฑฺฑิตํ อฏฺฐิก\-สงฺขลิกํ สมํสโลหิตํ นฺหารุ\-สมฺพนฺธํ. อฏฺฐิก\-สงฺขลิกํ \csromanfolio{124}นิมํส\-โลหิต\-มกฺขิตํ นฺหารุ\-สมฺพนฺธํ. อฏฺฐิก\-สงฺขลิกํ อปคต\-มํส\-โลหิตํ นฺหารุ\-สมฺพนฺธํ. อฏฺฐิกานิ อปคต\-สมฺพนฺธานิ ทิสา\-วิทิสา\-วิกฺขิตฺตานิ อญฺเญน หตฺถฏฺฐิกํ อญฺเญน ปาทฏฺฐิกํ อญฺเญน โคปฺผก\-ฏฺฐิกํ อญฺเญน ชงฺฆฏฺฐิกํ อญฺเญน อูรุฏฺฐิกํ อญฺเญน กฏิฏฺฐิกํ อญฺเญน ผาสุกฏฺฐิกํ อญฺเญน ปิฏฺฐิฏฺฐิกํ อญฺเญน ขนฺธฏฺฐิกํ อญฺเญน คีวฏฺฐิกํ อญฺเญน หนุกฏฺฐิกํ อญฺเญน ทนฺตฏฺฐิกํ อญฺเญน สีสกฏาหํ. ตํ กิํ มญฺญถ ภิกฺขเว, ยา ปุริมา สุภา วณฺณนิภา, สา อนฺตรหิตา, อาทีนโว ปาตุภูโตติ. เอวํ ภนฺเต. อยมฺปิ ภิกฺขเว รูปานํ อาทีนโว.}

\prose{ปุน จปรํ ภิกฺขเว ตเมว ภคินิํ ปสฺเสยฺย สรีรํ สิวถิกาย ฉฑฺฑิตํ อฏฺฐิกานิ เสตานิ สงฺข\-วณฺณ\-ปฏิภาคานิ. อฏฺฐิกานิ ปุญฺชกิตานิ เตโรวสฺสิกานิ. อฏฺฐิกานิ ปูตีนิ จุณฺณกชาตานิ. ตํ กิํ มญฺญถ ภิกฺขเว, ยา ปุริมา สุภา วณฺณนิภา, สา อนฺตรหิตา, อาทีนโว ปาตุภูโตติ. เอวํ ภนฺเต. อยมฺปิ ภิกฺขเว รูปานํ อาทีนโว.}

\prose{กิญฺจ ภิกฺขเว รูปานํ นิสฺสรณํ. โย ภิกฺขเว รูเปสุ ฉนฺทราค\-วินโย ฉนฺท\-ราค\-ปฺปหานํ. อิทํ รูปานํ นิสฺสรณํ.}

\prose{เย หิ เกจิ ภิกฺขเว สมณา วา พฺราหฺมณา วา เอวํ รูปานํ อสฺสาทญฺจ อสฺสาทโต อาทีนวญฺจ อาทีนวโต นิสฺสรณญฺจ นิสฺสรณโต ยถาภูตํ นปฺปชานนฺติ, เต วต สามํ วา รูเป ปริชานิสฺสนฺติ, ปรํ วา ตถตฺตาย สมาทเปสฺสนฺติ, ยถา ปฏิปนฺโน รูเป ปริชานิสฺสตีติ เนตํ ฐานํ วิชฺชติ. เย จ โข เกจิ ภิกฺขเว สมณา วา พฺราหฺมณา วา เอวํ รูปานํ อสฺสาทญฺจ อสฺสาทโต อาทีนวญฺจ อาทีนวโต นิสฺสรณญฺจ นิสฺสรณโต ยถาภูตํ ปชานนฺติ, เต วต สามํ วา รูเป ปริชานิสฺสนฺติ, ปรํ วา ตถตฺตาย สมาทเปสฺสนฺติ, ยถา ปฏิปนฺโน รูเป ปริชานิสฺสตีติ ฐานเมตํ วิชฺชติ.}

\proseitem{173}{โก จ ภิกฺขเว เวทนานํ อสฺสาโท. อิธ ภิกฺขเว ภิกฺขุ วิวิจฺเจว กาเมหิ วิวิจฺจ อกุสเลหิ ธมฺเมหิ สวิตกฺกํ สวิจารํ วิเวกชํ ปีติสุขํ ปฐมํ ฌานํ อุปสมฺปชฺช วิหรติ. ยสฺมิํ สมเย ภิกฺขเว ภิกฺขุ วิวิจฺเจว กาเมหิ วิวิจฺจ อกุสเลหิ ธมฺเมหิ สวิตกฺกํ สวิจารํ วิเวกชํ ปีติสุขํ ปฐมํ ฌานํ อุปสมฺปชฺช วิหรติ, เนว ตสฺมิํ สมเย อตฺตพฺยาพาธายปิ เจเตติ, น ปรพฺยาพาธายปิ เจเตติ, น อุภย\-พฺยาพาธายปิ \csromanfolio{125}เจเตติ, อพฺยาพชฺฌํเยว ตสฺมิํ สมเย เวทนํ เวเทติ. อพฺยาพชฺฌปรมาหํ ภิกฺขเว เวทนานํ อสฺสาทํ วทามิ.}

\prose{ปุน จปรํ ภิกฺขเว ภิกฺขุ วิตกฺก\-วิจารานํ วูปสมา อชฺฌตฺตํ สมฺปสาทนํ เจตโส เอโกทิภาวํ อวิตกฺกํ อวิจารํ สมาธิชํ ปีติสุขํ ทุติยํ ฌานํ อุปสมฺปชฺช วิหรติ. ฯเปฯ. ยสฺมิํ สมเย ภิกฺขเว ภิกฺขุ ปีติยา จ วิราคา อุเปกฺขโก จ วิหรติ สโต จ สมฺปชาโน, สุขญฺจ กาเยน ปฏิสํเวเทติ, ยํ ตํ อริยา อาจิกฺขนฺติ “อุเปกฺขโก สติมา สุขวิหารี”ติ ตติยํ ฌานํ อุปสมฺปชฺช วิหรติ ฯเปฯ. ยสฺมิํ สมเย ภิกฺขเว ภิกฺขุ สุขสฺส จ ปหานา ทุกฺขสฺส จ ปหานา ปุพฺเพว โสมนสฺส\-โทมนสฺสานํ อตฺถงฺคมา อทุกฺขมสุขํ อุเปกฺขา\-สติ\-ปาริสุทฺธิํ จตุตฺถํ ฌานํ อุปสมฺปชฺช วิหรติ, เนว ตสฺมิํ สมเย อตฺตพฺยาพาธายปิ เจเตติ, น ปรพฺยาพาธายปิ เจเตติ, น อุภย\-พฺยาพาธายปิ เจเตติ, อพฺยาพชฺฌํเยว ตสฺมิํ สมเย เวทนํ เวเทติ. อพฺยาพชฺฌปรมาหํ ภิกฺขเว เวทนานํ อสฺสาทํ วทามิ.}

\proseitem{174}{โก จ ภิกฺขเว เวทนานํ อาทีนโว. ยํ ภิกฺขเว เวทนา อนิจฺจา ทุกฺขา วิปริณาม\-ธมฺมา. อยํ เวทนานํ อาทีนโว.}

\prose{กิญฺจ ภิกฺขเว เวทนานํ นิสฺสรณํ. โย ภิกฺขเว เวทนาสุ ฉนฺทราค\-วินโย ฉนฺท\-ราค\-ปฺปหานํ. อิทํ เวทนานํ นิสฺสรณํ.}

\prose{เย หิ เกจิ ภิกฺขเว สมณา วา พฺราหฺมณา วา เอวํ เวทนานํ อสฺสาทญฺจ อสฺสาทโต อาทีนวญฺจ อาทีนวโต นิสฺสรณญฺจ นิสฺสรณโต ยถาภูตํ นปฺปชานนฺติ, เต วต สามํ วา เวทนํ ปริชานิสฺสนฺติ, ปรํ วา ตถตฺตาย สมาทเปสฺสนฺติ, ยถา ปฏิปนฺโน เวทนํ ปริชานิสฺสตีติ เนตํ ฐานํ วิชฺชติ. เย จ โข เกจิ ภิกฺขเว สมณา วา พฺราหฺมณา วา เอวํ เวทนานํ อสฺสาทญฺจ อสฺสาทโต อาทีนวญฺจ อาทีนวโต นิสฺสรณญฺจ นิสฺสรณโต ยถาภูตํ ปชานนฺติ, เต วต สามํ วา เวทนํ ปริชานิสฺสนฺติ, ปรํ วา ตถตฺตาย สมาทเปสฺสนฺติ, ยถา ปฏิปนฺโน เวทนํ ปริชานิสฺสตีติ ฐานเมตํ วิชฺชตีติ.}

\prosewithcloser{อิทมโวจ ภควา. อตฺตมนา เต ภิกฺขู ภควโต ภาสิตํ อภินนฺทุนฺติ.}{มหา\-ทุกฺขกฺขนฺธ\-สุตฺตํ นิฏฺฐิตํ ตติยํ.}
\csromanmatikaanchor{mtk.42}
\csromantocmarkref{subsection}{4. จูฬทุกฺขกฺขนฺธสุตฺต}{mtk.42}
\bookmark[dest={mtk.42},level=2]{4. จูฬทุกฺขกฺขนฺธสุตฺต}
\csromanheader{2}{4. จูฬทุกฺขกฺขนฺธ\-สุตฺต}
\proseitem{175}{\csromanfolio{126}เอวํ เม สุตํ\csromandash{}เอกํ สมยํ ภควา สกฺเกสุ วิหรติ กปิล\-วตฺถุสฺมิํ นิโคฺรธาราเม. อถ โข มหานาโม สกฺโก เยน ภควา เตนุปสงฺกมิ, อุปสงฺกมิตฺวา ภควนฺตํ อภิวาเทตฺวา เอกมนฺตํ นิสีทิ, เอกมนฺตํ นิสินฺโน โข มหานาโม สกฺโก ภควนฺตํ เอตทโวจ “ทีฆรตฺตาหํ ภนฺเต ภควตา เอวํ ธมฺมํ เทสิตํ อาชานามิ ‘โลโภ จิตฺตสฺส อุปกฺกิเลโส, โทโส จิตฺตสฺส อุปกฺกิเลโส, โมโห จิตฺตสฺส อุปกฺกิเลโส’ติ. เอวญฺจาหํ\csromansharedfootnote{803ed78fb5f89d98}{เอวํปาหํ (ก)} ภนฺเต ภควตา ธมฺมํ เทสิตํ อาชานามิ ‘โลโภ จิตฺตสฺส อุปกฺกิเลโส, โทโส จิตฺตสฺส อุปกฺกิเลโส, โมโห จิตฺตสฺส อุปกฺกิเลโส’ติ. อถ จ ปน เม เอกทา โลภธมฺมาปิ จิตฺตํ ปริยาทาย ติฏฺฐนฺติ, โทสธมฺมาปิ จิตฺตํ ปริยาทาย ติฏฺฐนฺติ, โมหธมฺมาปิ จิตฺตํ ปริยาทาย ติฏฺฐนฺติ. ตสฺส มยฺหํ ภนฺเต เอวํ โหติ, โกสุ นาม เม ธมฺโม อชฺฌตฺตํ อปฺปหีโน, เยน เม เอกทา โลภธมฺมาปิ จิตฺตํ ปริยาทาย ติฏฺฐนฺติ, โทสธมฺมาปิ จิตฺตํ ปริยาทาย ติฏฺฐนฺติ, โมหธมฺมาปิ จิตฺตํ ปริยาทาย ติฏฺฐนฺตี”ติ.}

\proseitem{176}{โส เอว โข เต มหานาม ธมฺโม อชฺฌตฺตํ อปฺปหีโน, เยน เต เอกทา โลภธมฺมาปิ จิตฺตํ ปริยาทาย ติฏฺฐนฺติ, โทสธมฺมาปิ จิตฺตํ ปริยาทาย ติฏฺฐนฺติ, โมหธมฺมาปิ จิตฺตํ ปริยาทาย ติฏฺฐนฺติ. โส จ หิ เต มหานาม ธมฺโม อชฺฌตฺตํ ปหีโน อภวิสฺส, น ตฺวํ อคารํ อชฺฌาวเสยฺยาสิ น กาเม ปริภุญฺเชยฺยาสิ. ยสฺมา จ โข เต มหานาม โส เอว ธมฺโม อชฺฌตฺตํ อปฺปหีโน, ตสฺมา ตฺวํ อคารํ อชฺฌาวสสิ กาเม ปริภุญฺชสิ.}

\proseitem{177}{“อปฺปสฺสาทา กามา พหุทุกฺขา พหุปายาสา\csromansharedfootnote{b329ce4887be7047}{พหูปายาสา (สี, สฺยา, อิ)} อาทีนโว เอตฺถ ภิยฺโย”ติ อิติ เจปิ มหานาม อริยสาวกสฺส ยถาภูตํ สมฺมปฺปญฺญาย สุทิฏฺฐํ โหติ, โส จ\csromansharedfootnote{36c249451f4c678a}{โสว (ก)} อญฺญเตฺรว กาเมหิ อญฺญตฺร อกุสเลหิ ธมฺเมหิ ปีติสุขํ นาธิคจฺฉติ อญฺญํ วา ตโต สนฺตตรํ. อถ โข โส เนว ตาว อนาวฏฺฏี กาเมสุ โหติ. ยโต จ โข มหานาม อริยสาวกสฺส “อปฺปสฺสาทา กามา พหุทุกฺขา พหุปายาสา \csromanfolio{127}อาทีนโว เอตฺถ ภิยฺโย”ติ เอวเมตํ ยถาภูตํ สมฺมปฺปญฺญาย สุทิฏฺฐํ โหติ, โส จ\csromansharedfootnote{36c249451f4c678a}{โสว (ก)} อญฺญเตฺรว กาเมหิ อญฺญตฺร อกุสเลหิ ธมฺเมหิ ปีติสุขํ อธิคจฺฉติ อญฺญํ วา ตโต สนฺตตรํ. อถ โข โส อนาวฏฺฏี กาเมสุ โหติ.}

\prose{มยฺหํปิ โข มหานาม ปุพฺเพว สมฺโพธา อนภิสมฺพุทฺธสฺส โพธิสตฺตสฺเสว สโต “อปฺปสฺสาทา กามา พหุทุกฺขา พหุปายาสา อาทีนโว เอตฺถ ภิยฺโย”ติ เอวเมตํ ยถาภูตํ สมฺมปฺปญฺญาย สุทิฏฺฐํ โหติ, โส จ\csromansharedfootnote{36c249451f4c678a}{โสว (ก)} อญฺญเตฺรว กาเมหิ อญฺญตฺร อกุสเลหิ ธมฺเมหิ ปีติสุขํ นาชฺฌคมํ อญฺญํ วา ตโต สนฺตตรํ. อถ ขฺวาหํ เนว ตาว อนาวฏฺฏี กาเมสุ ปจฺจญฺญาสิํ. ยโต จ โข เม มหานาม “อปฺปสฺสาทา กามา พหุทุกฺขา พหุปายาสา อาทีนโว เอตฺถ ภิยฺโย”ติ เอวเมตํ ยถาภูตํ สมฺมปฺปญฺญาย สุทิฏฺฐํ อโหสิ, โส จ\csromansharedfootnote{36c249451f4c678a}{โสว (ก)} อญฺญเตฺรว กาเมหิ อญฺญตฺร อกุสเลหิ ธมฺเมหิ ปีติสุขํ อชฺฌคมํ อญฺญํ วา ตโต สนฺตตรํ. อถาหํ อนาวฏฺฏี กาเมสุ ปจฺจญฺญาสิํ.}

\proseitem{178}{โก จ มหานาม กามานํ อสฺสาโท. ปญฺจิเม มหานาม กามคุณา. กตเม ปญฺจ. จกฺขุวิญฺเญยฺยา รูปา อิฏฺฐา กนฺตา มนาปา ปิยรูปา กามูปสํหิตา รชนียา. โสตวิญฺเญยฺยา สทฺทา ฯเปฯ. ฆานวิญฺเญยฺยา คนฺธา. ชิวฺหาวิญฺเญยฺยา รสา. กายวิญฺเญยฺยา โผฏฺฐพฺพา อิฏฺฐา กนฺตา มนาปา ปิยรูปา กามูปสํหิตา รชนียา. อิเม โข มหานาม ปญฺจ กามคุณา. ยํ โข มหานาม อิเม ปญฺจ กามคุเณ ปฏิจฺจ อุปฺปชฺชติ สุขํ โสมนสฺสํ. อยํ กามานํ อสฺสาโท.}

\prose{โก จ มหานาม กามานํ อาทีนโว. อิธ มหานาม กุลปุตฺโต เยน สิปฺปฏฺฐาเนน ชีวิกํ กปฺเปติ, ยทิ มุทฺทาย ยทิ คณนาย ยทิ สงฺขาเนน ยทิ กสิยา ยทิ วณิชฺชาย ยทิ โครกฺเขน ยทิ อิสฺสตฺเถน ยทิ ราชโปริเสน ยทิ สิปฺปญฺญตเรน, สีตสฺส ปุรกฺขโต อุณฺหสฺส ปุรกฺขโต ฑํสมกส\-วาตา\-ตป\-สรีสป\-สมฺผสฺเสหิ ริสฺสมาโน ขุปฺปิปาสาย มียมาโน. อยมฺปิ มหานาม กามานํ อาทีนโว สนฺทิฏฺฐิโก ทุกฺขกฺขนฺโธ กามเหตุ กามนิทานํ กามาธิกรณํ กามานเมว เหตุ.}

\prose{\csromanfolio{128}ตสฺส เจ มหานาม กุลปุตฺตสฺส เอวํ อุฏฺฐหโต ฆฏโต วายมโต เต โภคา นาภิ\-นิ\-ปฺผชฺชนฺติ, โส โสจติ กิลมติ ปริเทวติ อุรตฺตาฬิํ กนฺทติ สมฺโมหํ อาปชฺชติ “โมฆํ วต เม อุฏฺฐานํ, อผโล วต เม วายาโม”ติ. อยมฺปิ มหานาม กามานํ อาทีนโว สนฺทิฏฺฐิโก ทุกฺขกฺขนฺโธ กามเหตุ กามนิทานํ กามาธิกรณํ กามานเมว เหตุ.}

\prose{ตสฺส เจ มหานาม กุลปุตฺตสฺส เอวํ อุฏฺฐหโต ฆฏโต วายมโต เต โภคา อภินิ\-ปฺผชฺชนฺติ, โส เตสํ โภคานํ อารกฺขาธิกรณํ ทุกฺขํ โทมนสฺสํ ปฏิสํเวเทติ. กินฺติ เม โภเค เนว ราชาโน หเรยฺยุํ, น โจรา หเรยฺยุํ, น อคฺคิ ทเหยฺย, น อุทกํ วเหยฺย, น อปฺปิยา วา ทายาทา หเรยฺยุนฺติ. ตสฺส เอวํ อารกฺขโต โคปยโต เต โภเค ราชาโน วา หรนฺติ, โจรา วา หรนฺติ, อคฺคิ วา ทหติ, อุทกํ วา วหติ, อปฺปิยา วา ทายาทา หรนฺติ, โส โสจติ กิลมติ ปริเทวติ อุรตฺตาฬิํ กนฺทติ สมฺโมหํ อาปชฺชติ “ยมฺปิ เม อโหสิ, ตมฺปิ โน นตฺถี”ติ. อยมฺปิ มหานาม กามานํ อาทีนโว สนฺทิฏฺฐิโก ทุกฺขกฺขนฺโธ กามเหตุ กามนิทานํ กามาธิกรณํ กามานเมว เหตุ.}

\prose{ปุน จปรํ มหานาม กามเหตุ กามนิทานํ กามาธิกรณํ กามานเมว เหตุ ราชาโนปิ ราชูหิ วิวทนฺติ, ขตฺติยาปิ ขตฺติเยหิ วิวทนฺติ, พฺราหฺมณาปิ พฺราหฺมเณหิ วิวทนฺติ, คหปตีปิ คหปตีหิ วิวทนฺติ, มาตาปิ ปุตฺเตน วิวทติ, ปุตฺโตปิ มาตรา วิวทติ, ปิตาปิ ปุตฺเตน วิวทติ, ปุตฺโตปิ ปิตรา วิวทติ, ภาตาปิ ภาตรา วิวทติ, ภาตาปิ ภคินิยา วิวทติ, ภคินีปิ ภาตรา วิวทติ, สหาโยปิ สหาเยน วิวทติ. เต ตตฺถ กลห\-วิคฺคห\-วิวาทา\-ปนฺนา อญฺญมญฺญํ ปาณีหิปิ อุปกฺกมนฺติ, เลฑฺฑูหิปิ อุปกฺกมนฺติ, ทณฺเฑหิปิ อุปกฺกมนฺติ, สตฺเถหิปิ อุปกฺกมนฺติ. เต ตตฺถ มรณํปิ นิคจฺฉนฺติ มรณมตฺตํปิ ทุกฺขํ. อยมฺปิ มหานาม กามานํ อาทีนโว สนฺทิฏฺฐิโก ทุกฺขกฺขนฺโธ กามเหตุ กามนิทานํ กามาธิกรณํ กามานเมว เหตุ.}

\prose{ปุน จปรํ มหานาม กามเหตุ กามนิทานํ กามาธิกรณํ กามานเมว เหตุ อสิจมฺมํ คเหตฺวา ธนุกลาปํ สนฺนยฺหิตฺวา อุภโตพฺยูฬฺหํ สงฺคามํ ปกฺขนฺทนฺติ อุสูสุปิ ขิปฺปมาเนสุ สตฺตีสุปิ ขิปฺปมานาสุ อสีสุปิ วิชฺโชตลนฺเตสุ. เต ตตฺถ อุสูหิปิ วิชฺฌนฺติ, สตฺติยาปิ วิชฺฌนฺติ, อสินาปิ \csromanfolio{129}สีสํ ฉินฺทนฺติ. เต ตตฺถ มรณํปิ นิคจฺฉนฺติ มรณมตฺตํปิ ทุกฺขํ. อยมฺปิ มหานาม กามานํ อาทีนโว สนฺทิฏฺฐิโก ทุกฺขกฺขนฺโธ กามเหตุ กามนิทานํ กามาธิกรณํ กามานเมว เหตุ.}

\prose{ปุน จปรํ มหานาม กามเหตุ กามนิทานํ กามาธิกรณํ กามานเมว เหตุ อสิจมฺมํ คเหตฺวา ธนุกลาปํ สนฺนยฺหิตฺวา อทฺทาวเลปนา อุปการิโย ปกฺขนฺทนฺติ อุสูสุปิ ขิปฺปมาเนสุ สตฺตีสุปิ ขิปฺปมานาสุ อสีสุปิ วิชฺโชตลนฺเตสุ. เต ตตฺถ อุสูหิปิ วิชฺฌนฺติ, สตฺติยาปิ วิชฺฌนฺติ, ฉกณกายปิ โอสิญฺจนฺติ, อภิวคฺเคนปิ โอมทฺทนฺติ, อสินาปิ สีสํ ฉินฺทนฺติ. เต ตตฺถ มรณํปิ นิคจฺฉนฺติ มรณมตฺตํปิ ทุกฺขํ. อยมฺปิ มหานาม กามานํ อาทีนโว สนฺทิฏฺฐิโก ทุกฺขกฺขนฺโธ กามเหตุ กามนิทานํ กามาธิกรณํ กามานเมว เหตุ.}

\prose{ปุน จปรํ มหานาม กามเหตุ กามนิทานํ กามาธิกรณํ กามานเมว เหตุ สนฺธิํปิ ฉินฺทนฺติ, นิลฺโลปํปิ หรนฺติ, เอกาคาริกํปิ กโรนฺติ, ปริปนฺเถปิ ติฏฺฐนฺติ, ปรทารํปิ คจฺฉนฺติ. ตเมนํ ราชาโน คเหตฺวา วิวิธา กมฺมการณา กาเรนฺติ, กสาหิปิ ตาเฬนฺติ, เวตฺเตหิปิ ตาเฬนฺติ, อฑฺฒ\-ทณฺฑเกหิปิ ตาเฬนฺติ, หตฺถํปิ ฉินฺทนฺติ, ปาทํปิ ฉินฺทนฺติ, หตฺถปาทํปิ ฉินฺทนฺติ, กณฺณํปิ ฉินฺทนฺติ, นาสํปิ ฉินฺทนฺติ, กณฺณนาสํปิ ฉินฺทนฺติ, พิลงฺค\-ถาลิกํปิ กโรนฺติ, สงฺข\-มุณฺฑิกํปิ กโรนฺติ, ราหุมุขํปิ กโรนฺติ, โชติมาลิกํปิ กโรนฺติ, หตฺถปชฺโชติกํปิ กโรนฺติ, เอรกวตฺติกํปิ กโรนฺติ, จีรก\-วาสิกํปิ กโรนฺติ, เอเณยฺยกํปิ กโรนฺติ, พฬิส\-มํสิกํปิ กโรนฺติ, กหาปณิกํปิ กโรนฺติ, ขาราปตจฺฉิกํปิ กโรนฺติ, ปลิฆปริวตฺติกํปิ กโรนฺติ, ปลาล\-ปีฐกํปิ กโรนฺติ, ตตฺเตนปิ เตเลน โอสิญฺจนฺติ, สุนเขหิปิ ขาทาเปนฺติ, ชีวนฺตํปิ สูเล อุตฺตาเสนฺติ, อสินาปิ สีสํ ฉินฺทนฺติ. เต ตตฺถ มรณํปิ นิคจฺฉนฺติ มรณมตฺตํปิ ทุกฺขํ. อยมฺปิ มหานาม กามานํ อาทีนโว สนฺทิฏฺฐิโก ทุกฺขกฺขนฺโธ กามเหตุ กามนิทานํ กามาธิกรณํ กามานเมว เหตุ.}

\prose{ปุน จปรํ มหานาม กามเหตุ กามนิทานํ กามาธิกรณํ กามานเมว เหตุ กาเยน ทุจฺจริตํ จรนฺติ, วาจาย ทุจฺจริตํ จรนฺติ, มนสา ทุจฺจริตํ จรนฺติ. เต กาเยน ทุจฺจริตํ จริตฺวา วาจาย ทุจฺจริตํ จริตฺวา มนสา ทุจฺจริตํ จริตฺวา กายสฺส เภทา ปรํ มรณา อปายํ ทุคฺคติํ วินิปาตํ นิรยํ อุปปชฺชนฺติ. อยมฺปิ มหานาม กามานํ อาทีนโว \csromanfolio{130}สมฺปรายิโก ทุกฺขกฺขนฺโธ กามเหตุ กามนิทานํ กามาธิกรณํ กามานเมว เหตุ.}

\proseitem{179}{เอกมิทาหํ มหานาม สมยํ ราชคเห วิหรามิ คิชฺฌกูเฏ ปพฺพเต. เตน โข ปน สมเยน สมฺพหุลา นิคณฺฐา\csromansharedfootnote{7c8c47b0c617fe53}{นิคนฺถา (สฺยา, ก)} อิสิคิลิปสฺเส กาฬสิลายํ อุพฺภฏฺฐกา โหนฺติ อาสน\-ปฏิกฺขิตฺตา, โอปกฺกมิกา ทุกฺขา ติพฺพา ขรา กฏุกา เวทนา เวทยนฺติ. อถ ขฺวาหํ มหานาม สายนฺหสมยํ ปฏิสลฺลานา วุฏฺฐิโต เยน อิสิคิลิปสฺเส กาฬสิลา, เยน เต นิคณฺฐา เตนุปสงฺกมิํ, อุปสงฺกมิตฺวา เต นิคณฺเฐ เอตทโวจํ “กินฺนุ ตุเมฺห อาวุโส นิคณฺฐา อุพฺภฏฺฐกา อาสน\-ปฏิกฺขิตฺตา โอปกฺกมิกา ทุกฺขา ติพฺพา ขรา กฏุกา เวทนา เวทยถา”ติ. เอวํ วุตฺเต มหานาม เต นิคณฺฐา มํ เอตทโวจุํ “นิคณฺโฐ อาวุโส นาฏปุตฺโต\csromansharedfootnote{be303b1eabc46dd2}{นาถปุตฺโต (สี, อิ)} สพฺพญฺญู สพฺพทสฺสาวี อปริเสสํ ญาณทสฺสนํ ปฏิชานาติ ‘จรโต จ เม ติฏฺฐโต จ สุตฺตสฺส จ ชาครสฺส จ สตตํ สมิตํ ญาณทสฺสนํ ปจฺจุปฏฺฐิตนฺ’ติ. โส เอวมาห ‘อตฺถิ โข โว\csromansharedfootnote{790f0afb10ae6f23}{อตฺถิ โข โภ (สฺยา, ก)} นิคณฺฐา ปุพฺเพ ปาปกมฺมํ กตํ, ตํ อิมาย กฏุกาย ทุกฺกร\-การิกาย นิชฺชีเรถ\csromansharedfootnote{9ccf3a4dc0ba99da}{นิชฺชเรถ (สี, สฺยา, อิ)}, ยํ ปเนตฺถ\csromansharedfootnote{a04069d73487fa0b}{มยํ ปเนตฺถ (ก)} เอตรหิ กาเยน สํวุตา วาจาย สํวุตา มนสา สํวุตา ตํ อายติํ ปาปสฺส กมฺมสฺส อกรณํ, อิติ ปุราณานํ กมฺมานํ ตปสา พฺยนฺติภาวา นวานํ กมฺมานํ อกรณา อายติํ อนวสฺสโว, อายติํ อนวสฺสวา กมฺมกฺขโย, กมฺมกฺขยา ทุกฺขกฺขโย, ทุกฺขกฺขยา เวทนากฺขโย, เวทนากฺขยา สพฺพํ ทุกฺขํ นิชฺชิณฺณํ ภวิสฺสตี’ติ. ตญฺจ ปนมฺหากํ รุจฺจติ เจว ขมติ จ, เตน จมฺห อตฺตมนา”ติ.}

\proseitem{180}{เอวํ วุตฺเต อหํ มหานาม เต นิคณฺเฐ เอตทโวจํ “กิํ ปน ตุเมฺห อาวุโส นิคณฺฐา ชานาถ ‘อหุวเมฺหว มยํ ปุพฺเพ, น นาหุวมฺหา’ติ”. โน หิทํ อาวุโส. “กิํ ปน ตุเมฺห อาวุโส นิคณฺฐา ชานาถ ‘อกรเมฺหว มยํ ปุพฺเพ ปาปกมฺมํ, น นากรมฺหา’ติ”. โน หิทํ อาวุโส. “กิํ ปน ตุเมฺห อาวุโส นิคณฺฐา ชานาถ ‘เอวรูปํ วา เอวรูปํ วา ปาปกมฺมํ อกรมฺหา’ติ”. โน หิทํ อาวุโส. “กิํ ปน ตุเมฺห อาวุโส นิคณฺฐา ชานาถ ‘เอตฺตกํ วา ทุกฺขํ นิชฺชิณฺณํ, เอตฺตกํ วา ทุกฺขํ \csromanfolio{131}นิชฺชีเรตพฺพํ, เอตฺตกมฺหิ วา ทุกฺเข นิชฺชิณฺเณ สพฺพํ ทุกฺขํ นิชฺชิณฺณํ ภวิสฺสตี’ติ”. โน หิทํ อาวุโส. “กิํ ปน ตุเมฺห อาวุโส นิคณฺฐา ชานาถ ทิฏฺเฐว ธมฺเม อกุสลานํ ธมฺมานํ ปหานํ กุสลานํ ธมฺมานํ อุปสมฺปทนฺ”ติ. โน หิทํ อาวุโส.}

\prose{“อิติ กิร ตุเมฺห อาวุโส นิคณฺฐา น ชานาถ ‘อหุวเมฺหว มยํ ปุพฺเพ, น นาหุวมฺหา’ติ. น ชานาถ ‘อกรเมฺหว มยํ ปุพฺเพ ปาปกมฺมํ, น นากรมฺหา’ติ. น ชานาถ ‘เอวรูปํ วา เอวรูปํ วา ปาปกมฺมํ อกรมฺหา’ติ. น ชานาถ ‘เอตฺตกํ วา ทุกฺขํ นิชฺชิณฺณํ, เอตฺตกํ วา ทุกฺขํ นิชฺชีเรตพฺพํ, เอตฺตกมฺหิ วา ทุกฺเข นิชฺชิณฺเณ สพฺพํ ทุกฺขํ นิชฺชิณฺณํ ภวิสฺสตี’ติ. น ชานาถ ทิฏฺเฐว ธมฺเม อกุสลานํ ธมฺมานํ ปหานํ กุสลานํ ธมฺมานํ อุปสมฺปทํ. เอวํ สนฺเต อาวุโส นิคณฺฐา เย โลเก ลุทฺทา โลหิตปาณิโน กุรูรกมฺมนฺตา มนุสฺเสสุ ปจฺจาชาตา, เต นิคณฺเฐสุ ปพฺพชนฺตี”ติ. น โข อาวุโส โคตม สุเขน สุขํ อธิคนฺตพฺพํ, ทุกฺเขน โข สุขํ อธิคนฺตพฺพํ. สุเขน จาวุโส โคตม สุขํ อธิคนฺตพฺพํ อภวิสฺส, ราชา มาคโธ เสนิโย พิมฺพิสาโร สุขํ อธิคจฺเฉยฺย, ราชา มาคโธ เสนิโย พิมฺพิสาโร สุขวิหาริ\-ตโร อายสฺมตา โคตเมนาติ.}

\prose{“อทฺธา\-ยสฺมนฺเตหิ นิคณฺเฐหิ สหสา อปฺปฏิสงฺขา วาจา ภาสิตา ‘น โข อาวุโส โคตม สุเขน สุขํ อธิคนฺตพฺพํ, ทุกฺเขน โข สุขํ อธิคนฺตพฺพํ. สุเขน จาวุโส โคตม สุขํ อธิคนฺตพฺพํ อภวิสฺส, ราชา มาคโธ เสนิโย พิมฺพิสาโร สุขํ อธิคจฺเฉยฺย, ราชา มาคโธ เสนิโย พิมฺพิสาโร สุขวิหาริ\-ตโร อายสฺมตา โคตเมนา’ติ. อปิ จ อหเมว ตตฺถ ปฏิปุจฺฉิตพฺโพ ‘โก นุ โข อายสฺมนฺตานํ สุขวิหาริ\-ตโร ราชา วา มาคโธ เสนิโย พิมฺพิสาโร อายสฺมา วา โคตโม’ติ”. อทฺธาวุโส โคตม อเมฺหหิ สหสา อปฺปฏิสงฺขา วาจา ภาสิตา “น โข อาวุโส โคตม สุเขน สุขํ อธิคนฺตพฺพํ, ทุกฺเขน โข สุขํ อธิคนฺตพฺพํ. สุเขน จาวุโส โคตม สุขํ อธิคนฺตพฺพํ อภวิสฺส, ราชา มาคโธ เสนิโย พิมฺพิสาโร สุขํ อธิคจฺเฉยฺย, ราชา มาคโธ เสนิโย พิมฺพิสาโร สุขวิหาริ\-ตโร อายสฺมตา โคตเมนา”ติ. อปิ จ ติฏฺฐเตตํ. อิทานิปิ มยํ อายสฺมนฺตํ โคตมํ ปุจฺฉาม “โก นุ โข อายสฺมนฺตานํ สุขวิหาริ\-ตโร ราชา วา มาคโธ เสนิโย พิมฺพิสาโร อายสฺมา วา โคตโม”ติ.}

\proseitem{4}{\csromanfolio{132}“เตน หาวุโส นิคณฺฐา ตุเมฺหว ตตฺถ ปฏิปุจฺฉิสฺสามิ, ยถา โว ขเมยฺย, ตถา นํ พฺยากเรยฺยาถ. ตํ กิํมญฺญถาวุโส นิคณฺฐา, ปโหติ ราชา มาคโธ เสนิโย พิมฺพิสาโร อนิญฺชมาโน กาเยน อภาสมาโน วาจํ สตฺต รตฺตินฺทิวานิ เอกนฺตสุขํ ปฏิสํเวที วิหริตุนฺ”ติ. โน หิทํ อาวุโส.}

\prose{“ตํ กิํมญฺญถาวุโส นิคณฺฐา, ปโหติ ราชา มาคโธ เสนิโย พิมฺพิสาโร อนิญฺชมาโน กาเยน อภาสมาโน วาจํ ฉ รตฺตินฺทิวานิ ฯเปฯ ปญฺจ รตฺตินฺทิวานิ. จตฺตาริ รตฺตินฺทิวานิ. ตีณิ รตฺตินฺทิวานิ. เทฺว รตฺตินฺทิวานิ. เอกํ รตฺตินฺทิวํ เอกนฺตสุขํ ปฏิสํเวที วิหริตุนฺ”ติ. โน หิทํ อาวุโส.}

\prose{“อหํ โข อาวุโส นิคณฺฐา ปโหมิ อนิญฺชมาโน กาเยน อภาสมาโน วาจํ เอกํ รตฺตินฺทิวํ เอกนฺตสุขํ ปฏิสํเวที วิหริตุํ. อหํ โข อาวุโส นิคณฺฐา ปโหมิ อนิญฺชมาโน กาเยน อภาสมาโน วาจํ เทฺว รตฺตินฺทิวานิ. ตีณิ รตฺตินฺทิวานิ. จตฺตาริ รตฺตินฺทิวานิ. ปญฺจ รตฺตินฺทิวานิ. ฉ รตฺตินฺทิวานิ. สตฺต รตฺตินฺทิวานิ เอกนฺตสุขํ ปฏิสํเวที วิหริตุํ. ตํ กิํมญฺญถาวุโส นิคณฺฐา เอวํ สนฺเต โก สุขวิหาริ\-ตโร ราชา วา มาคโธ เสนิโย พิมฺพิสาโร อหํ วา”ติ. เอวํ สนฺเต อายสฺมาว โคตโม สุขวิหาริ\-ตโร รญฺญา มาคเธน เสนิเยน พิมฺพิสาเรนาติ.}

\prosewithcloserruled{อิทมโวจ ภควา. อตฺตมโน มหานาโม สกฺโก ภควโต ภาสิตํ อภินนฺทีติ.}{จูฬทุกฺขกฺขนฺธ\-สุตฺตํ นิฏฺฐิตํ จตุตฺถํ.}
\csromanmatikaanchor{mtk.43}
\csromantocmarkref{subsection}{5. อนุมานสุตฺต}{mtk.43}
\bookmark[dest={mtk.43},level=2]{5. อนุมานสุตฺต}
\csromanheader{2}{5. \textbf{อนุมานสุตฺต}}
\proseitem{181}{เอวํ เม สุตํ\csromandash{}เอกํ สมยํ อายสฺมา มหาโมคฺคลฺลาโน ภคฺเคสุ วิหรติ สุสุมารคิเร\csromansharedfootnote{6d1e3853d4a77e37}{สุํสุมารคิเร (สี, สฺยา, อิ)} เภสกฬาวเน มิคทาเย. ตตฺร โข อายสฺมา มหาโมคฺคลฺลาโน ภิกฺขู อามนฺเตสิ “อาวุโส ภิกฺขโว”ติ. “อาวุโส”ติ โข เต ภิกฺขู อายสฺมโต มหา\-โมคฺคลฺลานสฺส ปจฺจสฺโสสุํ, อายสฺมา มหาโมคฺคลฺลาโน เอตทโวจ\csromandash{}}

\prose{\csromanfolio{133}ปวาเรติ เจปิ อาวุโส ภิกฺขุ “วทนฺตุ มํ อายสฺมนฺโต, วจนีโยมฺหิ อายสฺมนฺเตหี”ติ, โส จ โหติ ทุพฺพโจ โทวจสฺส\-กรเณหิ ธมฺเมหิ สมนฺนาคโต อกฺขโม อปฺปทกฺขิณคฺคาหี อนุสาสนิํ. อถ โข นํ สพฺรหฺมจารี น เจว วตฺตพฺพํ มญฺญนฺติ, น จ อนุสาสิตพฺพํ มญฺญนฺติ, น จ ตสฺมิํ ปุคฺคเล วิสฺสาสํ อาปชฺชิตพฺพํ มญฺญนฺติ.}

\prose{กตเม จาวุโส โทวจสฺส\-กรณา ธมฺมา. อิธาวุโส ภิกฺขุ ปาปิจฺโฉ โหติ ปาปิกานํ อิจฺฉานํ วสํ คโต. ยํปาวุโส ภิกฺขุ ปาปิจฺโฉ โหติ ปาปิกานํ อิจฺฉานํ วสํ คโต, อยมฺปิ ธมฺโม โทวจสฺส\-กรโณ.}

\prose{ปุน จปรํ อาวุโส ภิกฺขุ อตฺตุกฺกํสโก โหติ ปรวมฺภี. ยํปาวุโส ภิกฺขุ อตฺตุกฺกํสโก โหติ ปรวมฺภี, อยมฺปิ ธมฺโม โทวจสฺส\-กรโณ.}

\prose{ปุน จปรํ อาวุโส ภิกฺขุ โกธโน โหติ โกธาภิภูโต. ยํปาวุโส ภิกฺขุ โกธโน โหติ โกธาภิภูโต, อยมฺปิ ธมฺโม โทวจสฺส\-กรโณ.}

\prose{ปุน จปรํ อาวุโส ภิกฺขุ โกธโน โหติ โกธเหตุ อุปนาหี. ยํปาวุโส ภิกฺขุ โกธโน โหติ โกธเหตุ อุปนาหี, อยมฺปิ ธมฺโม โทวจสฺส\-กรโณ.}

\prose{ปุน จปรํ อาวุโส ภิกฺขุ โกธโน โหติ โกธเหตุ อภิสงฺคี. ยํปาวุโส ภิกฺขุ โกธโน โหติ โกธเหตุ อภิสงฺคี, อยมฺปิ ธมฺโม โทวจสฺส\-กรโณ.}

\prose{ปุน จปรํ อาวุโส ภิกฺขุ โกธโน โหติ โกธสามนฺตา\csromansharedfootnote{dbe59350e33307fd}{โกธสามนฺตํ (สฺยา, อิ, ก)} วาจํ นิจฺฉาเรตา. ยํปาวุโส ภิกฺขุ โกธโน โหติ โกธสามนฺตา วาจํ นิจฺฉาเรตา, อยมฺปิ ธมฺโม โทวจสฺส\-กรโณ.}

\prose{ปุน จปรํ อาวุโส ภิกฺขุ โจทิโต\csromansharedfootnote{16fe989c02f529fb}{จุทิโต (สี, สฺยา, อิ)} โจทเกน โจทกํ ปฏิปฺผรติ. ยํปาวุโส ภิกฺขุ โจทิโต โจทเกน โจทกํ ปฏิปฺผรติ, อยมฺปิ ธมฺโม โทวจสฺส\-กรโณ.}

\prose{\csromanfolio{134}ปุน จปรํ อาวุโส ภิกฺขุ โจทิโต โจทเกน โจทกํ อปสาเทติ. ยํปาวุโส ภิกฺขุ โจทิโต โจทเกน โจทกํ อปสาเทติ, อยมฺปิ ธมฺโม โทวจสฺส\-กรโณ.}

\prose{ปุน จปรํ อาวุโส ภิกฺขุ โจทิโต โจทเกน โจทกสฺส ปจฺจาโรเปติ. ยํปาวุโส ภิกฺขุ โจทิโต โจทเกน โจทกสฺส ปจฺจาโรเปติ, อยมฺปิ ธมฺโม โทวจสฺส\-กรโณ.}

\prose{ปุน จปรํ อาวุโส ภิกฺขุ โจทิโต โจทเกน อญฺเญนญฺญํ ปฏิจรติ, พหิทฺธา กถํ อปนาเมติ, โกปญฺจ โทสญฺจ อปฺปจฺจยญฺจ ปาตุกโรติ. ยํปาวุโส ภิกฺขุ โจทิโต โจทเกน อญฺเญนญฺญํ ปฏิจรติ, พหิทฺธา กถํ อปนาเมติ, โกปญฺจ โทสญฺจ อปฺปจฺจยญฺจ ปาตุกโรติ, อยมฺปิ ธมฺโม โทวจสฺส\-กรโณ.}

\prose{ปุน จปรํ อาวุโส ภิกฺขุ โจทิโต โจทเกน อปทาเน น สมฺปายติ. ยํปาวุโส ภิกฺขุ โจทิโต โจทเกน อปทาเน น สมฺปายติ, อยมฺปิ ธมฺโม โทวจสฺส\-กรโณ.}

\prose{ปุน จปรํ อาวุโส ภิกฺขุ มกฺขี โหติ ปฬาสี. ยํปาวุโส ภิกฺขุ มกฺขี โหติ ปฬาสี, อยมฺปิ ธมฺโม โทวจสฺส\-กรโณ.}

\prose{ปุน จปรํ อาวุโส ภิกฺขุ อิสฺสุกี โหติ มจฺฉรี. ยํปาวุโส ภิกฺขุ อิสฺสุกี โหติ มจฺฉรี, อยมฺปิ ธมฺโม โทวจสฺส\-กรโณ.}

\prose{ปุน จปรํ อาวุโส ภิกฺขุ สโฐ โหติ มายาวี. ยํปาวุโส ภิกฺขุ สโฐ โหติ มายาวี, อยมฺปิ ธมฺโม โทวจสฺส\-กรโณ.}

\prose{ปุน จปรํ อาวุโส ภิกฺขุ ถทฺโธ โหติ อติมานี. ยํปาวุโส ภิกฺขุ ถทฺโธ โหติ อติมานี, อยมฺปิ ธมฺโม โทวจสฺส\-กรโณ.}

\prose{ปุน จปรํ อาวุโส ภิกฺขุ สนฺทิฏฺฐิ\-ปรามาสี โหติ อาธานคฺคาหี ทุปฺปฏินิสฺสคฺคี. ยํปาวุโส ภิกฺขุ สนฺทิฏฺฐิ\-ปรามาสี โหติ อาธานคฺคาหี ทุปฺปฏินิสฺสคฺคี, อยมฺปิ ธมฺโม โทวจสฺส\-กรโณ. อิเม วุจฺจนฺตาวุโส โทวจสฺส\-กรณา ธมฺมา.}

\proseitem{182}{\csromanfolio{135}โน เจปิ อาวุโส ภิกฺขุ ปวาเรติ “วทนฺตุ มํ อายสฺมนฺโต, วจนีโยมฺหิ อายสฺมนฺเตหี”ติ, โส จ โหติ สุวโจ โสวจสฺส\-กรเณหิ ธมฺเมหิ สมนฺนาคโต ขโม ปทกฺขิณ\-คฺคาหี อนุสาสนิํ. อถ โข นํ สพฺรหฺมจารี วตฺตพฺพญฺเจว มญฺญนฺติ อนุสาสิตพฺพญฺจ มญฺญนฺติ. ตสฺมิํ จ ปุคฺคเล วิสฺสาสํ อาปชฺชิตพฺพํ มญฺญนฺติ.}

\prose{กตเม จาวุโส โสวจสฺส\-กรณา ธมฺมา. อิธาวุโส ภิกฺขุ น ปาปิจฺโฉ โหติ น ปาปิกานํ อิจฺฉานํ วสํ คโต. ยํปาวุโส ภิกฺขุ น ปาปิจฺโฉ โหติ น ปาปิกานํ อิจฺฉานํ วสํ คโต, อยมฺปิ ธมฺโม โสวจสฺส\-กรโณ.}

\prose{ปุน จปรํ อาวุโส ภิกฺขุ อนตฺตุกฺกํสโก โหติ อปรวมฺภี. ยํปาวุโส ภิกฺขุ อนตฺตุกฺกํสโก โหติ อปรวมฺภี, อยมฺปิ ธมฺโม โสวจสฺส\-กรโณ.}

\prose{ปุน จปรํ อาวุโส ภิกฺขุ น โกธโน โหติ น โกธาภิภูโต. ยํปาวุโส ภิกฺขุ น โกธโน โหติ น โกธาภิภูโต, อยมฺปิ ธมฺโม โสวจสฺส\-กรโณ.}

\prose{ปุน จปรํ อาวุโส ภิกฺขุ น โกธโน โหติ น โกธเหตุ อุปนาหี. ยํปาวุโส ภิกฺขุ น โกธโน โหติ น โกธเหตุ อุปนาหี, อยมฺปิ ธมฺโม โสวจสฺส\-กรโณ.}

\prose{ปุน จปรํ อาวุโส ภิกฺขุ น โกธโน โหติ น โกธเหตุ อภิสงฺคี. ยํปาวุโส ภิกฺขุ น โกธโน โหติ น โกธเหตุ อภิสงฺคี, อยมฺปิ ธมฺโม โสวจสฺส\-กรโณ.}

\prose{ปุน จปรํ อาวุโส ภิกฺขุ น โกธโน โหติ น โกธสามนฺตา วาจํ นิจฺฉาเรตา. ยํปาวุโส ภิกฺขุ น โกธโน โหติ น โกธสามนฺตา วาจํ นิจฺฉาเรตา, อยมฺปิ ธมฺโม โสวจสฺส\-กรโณ.}

\prose{ปุน จปรํ อาวุโส ภิกฺขุ โจทิโต โจทเกน โจทกํ นปฺปฏิปฺผรติ. ยํปาวุโส ภิกฺขุ โจทิโต โจทเกน โจทกํ นปฺปฏิปฺผรติ, อยมฺปิ ธมฺโม โสวจสฺส\-กรโณ.}

\prose{\csromanfolio{136}ปุน จปรํ อาวุโส ภิกฺขุ โจทิโต โจทเกน โจทกํ น อปสาเทติ. ยํปาวุโส ภิกฺขุ โจทิโต โจทเกน โจทกํ น อปสาเทติ, อยมฺปิ ธมฺโม โสวจสฺส\-กรโณ.}

\prose{ปุน จปรํ อาวุโส ภิกฺขุ โจทิโต โจทเกน โจทกสฺส น ปจฺจาโรเปติ. ยํปาวุโส ภิกฺขุ โจทิโต โจทเกน โจทกสฺส น ปจฺจาโรเปติ, อยมฺปิ ธมฺโม โสวจสฺส\-กรโณ.}

\prose{ปุน จปรํ อาวุโส ภิกฺขุ โจทิโต โจทเกน น อญฺเญนญฺญํ ปฏิจรติ, น พหิทฺธา กถํ อปนาเมติ, น โกปญฺจ โทสญฺจ อปฺปจฺจยญฺจ ปาตุกโรติ. ยํปาวุโส ภิกฺขุ โจทิโต โจทเกน น อญฺเญนญฺญํ ปฏิจรติ, น พหิทฺธา กถํ อปนาเมติ, น โกปญฺจ โทสญฺจ อปฺปจฺจยญฺจ ปาตุกโรติ, อยมฺปิ ธมฺโม โสวจสฺส\-กรโณ.}

\prose{ปุน จปรํ อาวุโส ภิกฺขุ โจทิโต โจทเกน อปทาเน สมฺปายติ. ยํปาวุโส ภิกฺขุ โจทิโต โจทเกน อปทาเน สมฺปายติ, อยมฺปิ ธมฺโม โสวจสฺส\-กรโณ.}

\prose{ปุน จปรํ อาวุโส ภิกฺขุ อมกฺขี โหติ อปฬาสี. ยํปาวุโส ภิกฺขุ อมกฺขี โหติ อปฬาสี, อยมฺปิ ธมฺโม โสวจสฺส\-กรโณ.}

\prose{ปุน จปรํ อาวุโส ภิกฺขุ อนิสฺสุกี โหติ อมจฺฉรี. ยํปาวุโส ภิกฺขุ อนิสฺสุกี โหติ อมจฺฉรี, อยมฺปิ ธมฺโม โสวจสฺส\-กรโณ.}

\prose{ปุน จปรํ อาวุโส ภิกฺขุ อสโฐ โหติ อมายาวี. ยํปาวุโส ภิกฺขุ อสโฐ โหติ อมายาวี, อยมฺปิ ธมฺโม โสวจสฺส\-กรโณ.}

\prose{ปุน จปรํ อาวุโส ภิกฺขุ อตฺถทฺโธ โหติ อนติมานี. ยํปาวุโส ภิกฺขุ อตฺถทฺโธ โหติ อนติมานี, อยมฺปิ ธมฺโม โสวจสฺส\-กรโณ.}

\prose{ปุน จปรํ อาวุโส ภิกฺขุ อสนฺทิฏฺฐิปรามาสี โหติ อนาธานคฺคาหี สุปฺปฏินิสฺสคฺคี. ยํปาวุโส ภิกฺขุ อสนฺทิฏฺฐิปรามาสี โหติ อนาธานคฺคาหี \csromanfolio{137}สุปฺปฏินิสฺสคฺคี, อยมฺปิ ธมฺโม โสวจสฺส\-กรโณ. อิเม วุจฺจนฺตาวุโส โสวจสฺส\-กรณา ธมฺมา.}

\proseitem{183}{ตตฺราวุโส ภิกฺขุนา อตฺตนาว อตฺตานํ เอวํ อนุมินิตพฺพํ\csromansharedfootnote{ba8f5bb0db8e185f}{อนุมานิตพฺพํ (สี)}\csromandash{} “โย ขฺวายํ ปุคฺคโล ปาปิจฺโฉ ปาปิกานํ อิจฺฉานํ วสํ คโต, อยํ เม ปุคฺคโล อปฺปิโย อมนาโป. อหญฺเจว โข ปนสฺสํ ปาปิจฺโฉ ปาปิกานํ อิจฺฉานํ วสํ คโต, อหํปาสฺสํ ปเรสํ อปฺปิโย อมนาโป”ติ, เอวํ ชานนฺเตนาวุโส ภิกฺขุนา “น ปาปิจฺโฉ ภวิสฺสามิ น ปาปิกานํ อิจฺฉานํ วสํ คโต”ติ จิตฺตํ อุปฺปาเทตพฺพํ.}

\prose{“โย ขฺวายํ ปุคฺคโล อตฺตุกฺกํสโก ปรวมฺภี, อยํ เม ปุคฺคโล อปฺปิโย อมนาโป. อหญฺเจว โข ปนสฺสํ อตฺตุกฺกํสโก ปรวมฺภี, อหํปาสฺสํ ปเรสํ อปฺปิโย อมนาโป”ติ, เอวํ ชานนฺเตนาวุโส ภิกฺขุนา “อนตฺตุกฺกํสโก ภวิสฺสามิ อปรวมฺภี”ติ จิตฺตํ อุปฺปาเทตพฺพํ.}

\prose{“โย ขฺวายํ ปุคฺคโล โกธโน โกธาภิภูโต, อยํ เม ปุคฺคโล อปฺปิโย อมนาโป. อหญฺเจว โข ปนสฺสํ โกธโน โกธาภิภูโต, อหํปาสฺสํ ปเรสํ อปฺปิโย อมนาโป”ติ, เอวํ ชานนฺเตนาวุโส ภิกฺขุนา “น โกธโน ภวิสฺสามิ น โกธาภิภูโต”ติ จิตฺตํ อุปฺปาเทตพฺพํ.}

\prose{“โย ขฺวายํ ปุคฺคโล โกธโน โกธเหตุ อุปนาหี, อยํ เม ปุคฺคโล อปฺปิโย อมนาโป. อหญฺเจว โข ปนสฺสํ โกธโน โกธเหตุ อุปนาหี, อหํปาสฺสํ ปเรสํ อปฺปิโย อมนาโป”ติ, เอวํ ชานนฺเตนาวุโส ภิกฺขุนา “น โกธโน ภวิสฺสามิ น โกธเหตุ อุปนาหี”ติ จิตฺตํ อุปฺปาเทตพฺพํ.}

\prose{“โย ขฺวายํ ปุคฺคโล โกธโน โกธเหตุ อภิสงฺคี, อยํ เม ปุคฺคโล อปฺปิโย อมนาโป. อหญฺเจว โข ปนสฺสํ โกธโน โกธเหตุ อภิสงฺคี, อหํปาสฺสํ ปเรสํ อปฺปิโย อมนาโป”ติ, เอวํ ชานนฺเตนาวุโส ภิกฺขุนา “น โกธโน ภวิสฺสามิ น โกธเหตุ อภิสงฺคี”ติ จิตฺตํ อุปฺปาเทตพฺพํ.}

\prose{\csromanfolio{138}“โย ขฺวายํ ปุคฺคโล โกธโน โกธสามนฺตา วาจํ นิจฺฉาเรตา, อยํ เม ปุคฺคโล อปฺปิโย อมนาโป. อหญฺเจว โข ปนสฺสํ โกธโน โกธสามนฺตา วาจํ นิจฺฉาเรตา, อหํปาสฺสํ ปเรสํ อปฺปิโย อมนาโป”ติ, เอวํ ชานนฺเตนาวุโส ภิกฺขุนา “น โกธโน ภวิสฺสามิ น โกธสามนฺตา วาจํ นิจฺฉาเรสฺสามี”ติ จิตฺตํ อุปฺปาเทตพฺพํ.}

\prose{“โย ขฺวายํ ปุคฺคโล โจทิโต โจทเกน โจทกํ ปฏิปฺผรติ, อยํ เม ปุคฺคโล อปฺปิโย อมนาโป. อหญฺเจว โข ปน โจทิโต โจทเกน โจทกํ ปฏิปฺผเรยฺยํ, อหํปาสฺสํ ปเรสํ อปฺปิโย อมนาโป”ติ, เอวํ ชานนฺเตนาวุโส ภิกฺขุนา “โจทิโต โจทเกน โจทกํ นปฺปฏิปฺผริสฺสามี”ติ จิตฺตํ อุปฺปาเทตพฺพํ.}

\prose{“โย ขฺวายํ ปุคฺคโล โจทิโต โจทเกน โจทกํ อปสาเทติ, อยํ เม ปุคฺคโล อปฺปิโย อมนาโป. อหญฺเจว โข ปน โจทิโต โจทเกน โจทกํ อปสาเทยฺยํ, อหํปาสฺสํ ปเรสํ อปฺปิโย อมนาโป”ติ, เอวํ ชานนฺเตนาวุโส ภิกฺขุนา “โจทิโต โจทเกน โจทกํ น อปสาเทสฺสามี”ติ จิตฺตํ อุปฺปาเทตพฺพํ.}

\prose{“โย ขฺวายํ ปุคฺคโล โจทิโต โจทเกน โจทกสฺส ปจฺจาโรเปติ, อยํ เม ปุคฺคโล อปฺปิโย อมนาโป. อหญฺเจว โข ปน โจทิโต โจทเกน โจทกสฺส ปจฺจาโรเปยฺยํ, อหํปาสฺสํ ปเรสํ อปฺปิโย อมนาโป”ติ, เอวํ ชานนฺเตนาวุโส ภิกฺขุนา “โจทิโต โจทเกน โจทกสฺส น ปจฺจาโรเปสฺสามี”ติ จิตฺตํ อุปฺปาเทตพฺพํ.}

\prose{“โย ขฺวายํ ปุคฺคโล โจทิโต โจทเกน อญฺเญนญฺญํ ปฏิจรติ, พหิทฺธา กถํ อปนาเมติ, โกปญฺจ โทสญฺจ อปฺปจฺจยญฺจ ปาตุกโรติ, อยํ เม ปุคฺคโล อปฺปิโย อมนาโป. อหญฺเจว โข ปน โจทิโต โจทเกน อญฺเญนญฺญํ ปฏิจเรยฺยํ, พหิทฺธา กถํ อปนาเมยฺยํ, โกปญฺจ โทสญฺจ อปฺปจฺจยญฺจ ปาตุกเรยฺยํ, อหํปาสฺสํ ปเรสํ อปฺปิโย อมนาโป”ติ, เอวํ ชานนฺเตนาวุโส ภิกฺขุนา “โจทิโต โจทเกน น อญฺเญนญฺญํ ปฏิจริสฺสามิ, น พหิทฺธา กถํ อปนาเมสฺสามิ, น โกปญฺจ โทสญฺจ อปฺปจฺจยญฺจ ปาตุกริสฺสามี”ติ จิตฺตํ อุปฺปาเทตพฺพํ.}

\prose{\csromanfolio{139}“โย ขฺวายํ ปุคฺคโล โจทิโต โจทเกน อปทาเน น สมฺปายติ, อยํ เม ปุคฺคโล อปฺปิโย อมนาโป. อหญฺเจว โข ปน โจทิโต โจทเกน อปทาเน น สมฺปาเยยฺยํ, อหํปาสฺสํ ปเรสํ อปฺปิโย อมนาโป”ติ, เอวํ ชานนฺเตนาวุโส ภิกฺขุนา “โจทิโต โจทเกน อปทาเน สมฺปายิสฺสามี”ติ จิตฺตํ อุปฺปาเทตพฺพํ.}

\prose{“โย ขฺวายํ ปุคฺคโล มกฺขี ปฬาสี, อยํ เม ปุคฺคโล อปฺปิโย อมนาโป. อหญฺเจว โข ปนสฺสํ มกฺขี ปฬาสี, อหํปาสฺสํ ปเรสํ อปฺปิโย อมนาโป”ติ, เอวํ ชานนฺเตนาวุโส ภิกฺขุนา “อมกฺขี ภวิสฺสามิ อปฬาสี”ติ จิตฺตํ อุปฺปาเทตพฺพํ.}

\prose{“โย ขฺวายํ ปุคฺคโล อิสฺสุกี มจฺฉรี, อยํ เม ปุคฺคโล อปฺปิโย อมนาโป. อหญฺเจว โข ปนสฺสํ อิสฺสุกี มจฺฉรี, อหํปาสฺสํ ปเรสํ อปฺปิโย อมนาโป”ติ, เอวํ ชานนฺเตนาวุโส ภิกฺขุนา “อนิสฺสุกี ภวิสฺสามิ อมจฺฉรี”ติ จิตฺตํ อุปฺปาเทตพฺพํ.}

\prose{“โย ขฺวายํ ปุคฺคโล สโฐ มายาวี, อยํ เม ปุคฺคโล อปฺปิโย อมนาโป. อหญฺเจว โข ปนสฺสํ สโฐ มายาวี, อหํปาสฺสํ ปเรสํ อปฺปิโย อมนาโป”ติ, เอวํ ชานนฺเตนาวุโส ภิกฺขุนา “อสโฐ ภวิสฺสามิ อมายาวี”ติ จิตฺตํ อุปฺปาเทตพฺพํ.}

\prose{“โย ขฺวายํ ปุคฺคโล ถทฺโธ อติมานี, อยํ เม ปุคฺคโล อปฺปิโย อมนาโป. อหญฺเจว โข ปนสฺสํ ถทฺโธ อติมานี, อหํปาสฺสํ ปเรสํ อปฺปิโย อมนาโป”ติ, เอวํ ชานนฺเตนาวุโส ภิกฺขุนา “อตฺถทฺโธ ภวิสฺสามิ อนติมานี”ติ จิตฺตํ อุปฺปาเทตพฺพํ.}

\prose{“โย ขฺวายํ ปุคฺคโล สนฺทิฏฺฐิ\-ปรามาสี อาธานคฺคาหี ทุปฺปฏินิสฺสคฺคี, อยํ เม ปุคฺคโล อปฺปิโย อมนาโป. อหญฺเจว โข ปนสฺสํ สนฺทิฏฺฐิ\-ปรามาสี อาธานคฺคาหี ทุปฺปฏินิสฺสคฺคี, อหํปาสฺสํ ปเรสํ อปฺปิโย อมนาโป”ติ, เอวํ ชานนฺเตนาวุโส ภิกฺขุนา “อสนฺทิฏฺฐิ\-ปรามาสี ภวิสฺสามิ อนาธานคฺคาหี สุปฺปฏินิสฺสคฺคี”ติ จิตฺตํ อุปฺปาเทตพฺพํ.}

\proseitem{184}{ตตฺราวุโส ภิกฺขุนา อตฺตนาว อตฺถานํ เอวํ ปจฺจเวกฺขิตพฺพํ “กิํ นุ โขมฺหิ ปาปิจฺโฉ ปาปิกานํ อิจฺฉานํ วสํ คโต”ติ. สเจ \csromanfolio{140}อาวุโส ภิกฺขุ ปจฺจเวกฺขมาโน เอวํ ชานาติ “ปาปิจฺโฉ โขมฺหิ ปาปิกานํ อิจฺฉานํ วสํ คโต”ติ, เตนาวุโส ภิกฺขุนา เตสํเยว ปาปกานํ อกุสลานํ ธมฺมานํ ปหานาย วายมิตพฺพํ. สเจ ปนาวุโส ภิกฺขุ ปจฺจเวกฺขมาโน เอวํ ชานาติ “น โขมฺหิ ปาปิจฺโฉ น ปาปิกานํ อิจฺฉานํ วสํ คโต”ติ, เตนาวุโส ภิกฺขุนา เตเนว ปีติปาโมชฺเชน วิหาตพฺพํ อโห\-รตฺตา\-นุสิกฺขินา กุสเลสุ ธมฺเมสุ.}

\prose{ปุน จปรํ อาวุโส ภิกฺขุนา อตฺตนาว อตฺตานํ เอวํ ปจฺจเวกฺขิตพฺพํ “กิํ นุ โขมฺหิ อตฺตุกฺกํสโก ปรวมฺภี”ติ. สเจ อาวุโส ภิกฺขู ปจฺจเวกฺขมาโน เอวํ ชานาติ “อตฺตุกฺกํสโก โขมฺหิ ปรวมฺภี”ติ, เตนาวุโส ภิกฺขุนา เตสํเยว ปาปกานํ อกุสลานํ ธมฺมานํ ปหานาย วายมิตพฺพํ. สเจ ปนาวุโส ภิกฺขุ ปจฺจเวกฺขมาโน เอวํ ชานาติ “อนตฺตุกฺกํสโก โขมฺหิ อปรวมฺภี”ติ, เตนาวุโส ภิกฺขุนา เตเนว ปีติปาโมชฺเชน วิหาตพฺพํ อโห\-รตฺตา\-นุสิกฺขินา กุสเลสุ ธมฺเมสุ.}

\prose{ปุน จปรํ อาวุโส ภิกฺขุนา อตฺตนาว อตฺตานํ เอวํ ปจฺจเวกฺขิตพฺพํ “กิํ นุ โขมฺหิ โกธโน โกธาภิภูโต”ติ. สเจ อาวุโส ภิกฺขุ ปจฺจเวกฺขมาโน เอวํ ชานาติ “โกธโน โขมฺหิ โกธาภิภูโต”ติ, เตนาวุโส ภิกฺขุนา เตสํเยว ปาปกานํ อกุสลานํ ธมฺมานํ ปหานาย วายมิตพฺพํ. สเจ ปนาวุโส ภิกฺขุ ปจฺจเวกฺขมาโน เอวํ ชานาติ “น โขมฺหิ โกธโน โกธาภิภูโต”ติ, เตนาวุโส ภิกฺขุนา เตเนว ปีติปาโมชฺเชน วิหาตพฺพํ อโห\-รตฺตา\-นุสิกฺขินา กุสเลสุ ธมฺเมสุ.}

\prose{ปุน จปรํ อาวุโส ภิกฺขุนา อตฺตนาว อตฺตานํ เอวํ ปจฺจเวกฺขิตพฺพํ “กิํ นุ โขมฺหิ โกธโน โกธเหตุ อุปนาหี”ติ. สเจ อาวุโส ภิกฺขุ ปจฺจเวกฺขมาโน เอวํ ชานาติ “โกธโน โขมฺหิ โกธเหตุ อุปนาหี”ติ, เตนาวุโส ภิกฺขุนา เตสํเยว ปาปกานํ อกุสลานํ ธมฺมานํ ปหานาย วายมิตพฺพํ. สเจ ปนาวุโส ภิกฺขุ ปจฺจเวกฺขมาโน เอวํ ชานาติ “น โขมฺหิ โกธโน โกธเหตุ อุปนาหี”ติ, เตนาวุโส ภิกฺขุนา เตเนว ปีติปาโมชฺเชน วิหาตพฺพํ อโห\-รตฺตา\-นุสิกฺขินา กุสเลสุ ธมฺเมสุ.}

\prose{\csromanfolio{141}ปุน จปรํ อาวุโส ภิกฺขุนา อตฺตนาว อตฺตานํ เอวํ ปจฺจเวกฺขิตพฺพํ “กิํ นุ โขมฺหิ โกธโน โกธเหตุ อภิสงฺคี”ติ. สเจ อาวุโส ภิกฺขุ ปจฺจเวกฺขมาโน เอวํ ชานาติ “โกธโน โขมฺหิ โกธเหตุ อภิสงฺคี”ติ, เตนาวุโส ภิกฺขุนา เตสํเยว ปาปกานํ อกุสลานํ ธมฺมานํ ปหานาย วายมิตพฺพํ. สเจ ปนาวุโส ภิกฺขุ ปจฺจเวกฺขมาโน เอวํ ชานาติ “น โขมฺหิ โกธโน โกธเหตุ อภิสงฺคี”ติ, เตนาวุโส ภิกฺขุนา เตเนว ปีติปาโมชฺเชน วิหาตพฺพํ อโห\-รตฺตา\-นุสิกฺขินา กุสเลสุ ธมฺเมสุ.}

\prose{ปุน จปรํ อาวุโส ภิกฺขุนา อตฺตนาว อตฺตานํ เอวํ ปจฺจเวกฺขิตพฺพํ “กิํ นุ โขมฺหิ โกธโน โกธสามนฺตา วาจํ นิจฺฉาเรตา”ติ. สเจ อาวุโส ภิกฺขุ ปจฺจเวกฺขมาโน เอวํ ชานาติ “โกธโน โขมฺหิ โกธสามนฺตาวาจํ นิจฺฉาเรตา”ติ, เตนาวุโส ภิกฺขุนา เตสํเยว ปาปกานํ อกุสลานํ ธมฺมานํ ปหานาย วายมิตพฺพํ. สเจ ปนาวุโส ภิกฺขุ ปจฺจเวกฺขมาโน เอวํ ชานาติ “น โขมฺหิ โกธโน โกธสามนฺตา วาจํ นิจฺฉาเรตา”ติ, เตนาวุโส ภิกฺขุนา เตเนว ปีติปาโมชฺเชน วิหาตพฺพํ อโห\-รตฺตา\-นุสิกฺขินา กุสเลสุ ธมฺเมสุ.}

\prose{ปุน จปรํ อาวุโส ภิกฺขุนา อตฺตนาว อตฺตานํ เอวํ ปจฺจเวกฺขิตพฺพํ “กิํ นุ โขมฺหิ โจทิโต โจทเกน โจทกํ ปฏิปฺผรามี”ติ. สเจ อาวุโส ภิกฺขุ ปจฺจเวกฺขมาโน เอวํ ชานาติ “โจทิโต โขมฺหิ โจทเกน โจทกํ ปฏิปฺผรามี”ติ, เตนาวุโส ภิกฺขุนา เตสํเยว ปาปกานํ อกุสลานํ ธมฺมานํ ปหานาย วายมิตพฺพํ. สเจ ปนาวุโส ภิกฺขุ ปจฺจเวกฺขมาโน เอวํ ชานาติ “โจทิโต โขมฺหิ โจทเกน โจทกํ นปฺปฏิปฺผรามี”ติ, เตนาวุโส ภิกฺขุนา เตเนว ปีติปาโมชฺเชน วิหาตพฺพํ อโห\-รตฺตา\-นุสิกฺขินา กุสเลสุ ธมฺเมสุ.}

\prose{ปุน จปรํ อาวุโส ภิกฺขุนา อตฺตนาว อตฺตานํ เอวํ ปจฺจเวกฺขิตพฺพํ “กิํ นุ โขมฺหิ โจทิโต โจทเกน โจทกํ อปสาเทมี”ติ. สเจ อาวุโส ภิกฺขุ ปจฺจเวกฺขมาโน เอวํ ชานาติ “โจทิโต โขมฺหิ โจทเกน โจทกํ อปสาเทมี”ติ, เตนาวุโส ภิกฺขุนา เตสํเยว ปาปกานํ อกุสลานํ ธมฺมานํ ปหานาย วายมิตพฺพํ. สเจ ปนาวุโส ภิกฺขุ ปจฺจเวกฺขมาโน เอวํ ชานาติ “โจทิโต โขมฺหิ โจทเกน}

\prose{\csromanfolio{142}โจทกํ น อปสาเทมี”ติ. เตนาวุโส ภิกฺขุนา เตเนว ปีติปาโมชฺเชน วิหาตพฺพํ อโห\-รตฺตา\-นุสิกฺขินา กุสเลสุ ธมฺเมสุ.}

\prose{ปุน จปรํ อาวุโส ภิกฺขุนา อตฺตนาว อตฺตานํ เอวํ ปจฺจเวกฺขิตพฺพํ “กิํ นุ โขมฺหิ โจทิโต โจทเกน โจทกสฺส ปจฺจาโรเปมี”ติ. สเจ อาวุโส ภิกฺขุ ปจฺจเวกฺขมาโน เอวํ ชานาติ “โจทิโต โขมฺหิ โจทเกน โจทกสฺส ปจฺจาโรเปมี”ติ, เตนาวุโส ภิกฺขุนา เตสํเยว ปาปกานํ อกุสลานํ ธมฺมานํ ปหานาย วายมิตพฺพํ. สเจ ปนาวุโส ภิกฺขุ ปจฺจเวกฺขมาโน เอวํ ชานาติ “โจทิโต โขมฺหิ โจทเกน โจทกสฺส น ปจฺจาโรเปมี”ติ, เตนาวุโส ภิกฺขุนา เตเนว ปีติปาโมชฺเชน วิหาตพฺพํ อโห\-รตฺตา\-นุสิกฺขินา กุสเลสุ ธมฺเมสุ.}

\prose{ปุน จปรํ อาวุโส ภิกฺขุนา อตฺตนาว อตฺตานํ เอวํ ปจฺจเวกฺขิตพฺพํ “กิํ นุ โขมฺหิ โจทิโต โจทเกน อญฺเญนญฺญํ ปฏิจรามิ, พหิทฺธา กถํ อปนาเมมิ, โกปญฺจ โทสญฺจ อปฺปจฺจยญฺจ ปาตุกโรมี”ติ. สเจ อาวุโส ภิกฺขุ ปจฺจเวกฺขมาโน เอวํ ชานาติ “โจทิโต โขมฺหิ โจทเกน อญฺเญนญฺญํ ปฏิจรามิ, พหิทฺธา กถํ อปนาเมมิ, โกปญฺจ โทสญฺจ อปฺปจฺจยญฺจ ปาตุกโรมี”ติ, เตนาวุโส ภิกฺขุนา เตสํเยว ปาปกานํ อกุสลานํ ธมฺมานํ ปหานาย วายมิตพฺพํ. สเจ ปนาวุโส ภิกฺขุ ปจฺจเวกฺขมาโน เอวํ ชานาติ “โจทิโต โขมฺหิ โจทเกน น อญฺเญนญฺญํ ปฏิจรามิ, น พหิทฺธา กถํ อปนาเมมิ, น โกปญฺจ โทสญฺจ อปฺปจฺจยญฺจ ปาตุกโรมี”ติ, เตนาวุโส ภิกฺขุนา เตเนว ปีติปาโมชฺเชน วิหาตพฺพํ อโห\-รตฺตา\-นุสิกฺขินา กุสเลสุ ธมฺเมสุ.}

\prose{ปุน จปรํ อาวุโส ภิกฺขุนา อตฺตนาว อตฺตานํ เอวํ ปจฺจเวกฺขิตพฺพํ “กิํ นุ โขมฺหิ โจทิโต โจทเกน อปทาเน น สมฺปายามี”ติ. สเจ อาวุโส ภิกฺขุ ปจฺจเวกฺขมาโน เอวํ ชานาติ “โจทิโต โขมฺหิ โจทเกน อปทาเน น สมฺปายามี”ติ, เตนาวุโส ภิกฺขุนา เตสํเยว ปาปกานํ อกุสลานํ ธมฺมานํ ปหานาย วายมิตพฺพํ. สเจ ปนาวุโส ภิกฺขุ ปจฺจเวกฺขมาโน เอวํ ชานาติ “โจทิโต โขมฺหิ โจทเกน อปทาเน สมฺปายามี”ติ, เตนาวุโส ภิกฺขุนา เตเนว ปีติปาโมชฺเชน วิหาตพฺพํ อโห\-รตฺตา\-นุสิกฺขินา กุสเลสุ ธมฺเมสุ.}

\prose{\csromanfolio{143}ปุน จปรํ อาวุโส ภิกฺขุนา อตฺตนาว อตฺตานํ เอวํ ปจฺจเวกฺขิตพฺพํ “กิํ นุ โขมฺหิ มกฺขี ปฬาสี”ติ. สเจ อาวุโส ภิกฺขุ ปจฺจเวกฺขมาโน เอวํ ชานาติ “มกฺขี โขมฺหิ ปฬาสี”ติ, เตนาวุโส ภิกฺขุนา เตสํเยว ปาปกานํ อกุสลานํ ธมฺมานํ ปหานาย วายมิตพฺพํ. สเจ ปนาวุโส ภิกฺขุ ปจฺจเวกฺขมาโน เอวํ ชานาติ “อมกฺขี โขมฺหิ อปฬาสี”ติ, เตนาวุโส ภิกฺขุนา เตเนว ปีติปาโมชฺเชน วิหาตพฺพํ อโห\-รตฺตา\-นุสิกฺขินา กุสเลสุ ธมฺเมสุ.}

\prose{ปุน จปรํ อาวุโส ภิกฺขุนา อตฺตนาว อตฺตานํ เอวํ ปจฺจเวกฺขิตพฺพํ “กิํ นุ โขมฺหิ อิสฺสุกี มจฺฉรี”ติ. สเจ อาวุโส ภิกฺขุ ปจฺจเวกฺขมาโน เอวํ ชานาติ “อิสฺสุกี โขมฺหิ มจฺฉรี”ติ, เตนาวุโส ภิกฺขุนา เตสํเยว ปาปกานํ อกุสลานํ ธมฺมานํ ปหานาย วายมิตพฺพํ. สเจ ปนาวุโส ภิกฺขุ ปจฺจเวกฺขมาโน เอวํ ชานาติ “อนิสฺสุกี โขมฺหิ อมจฺฉรี”ติ, เตนาวุโส ภิกฺขุนา เตเนว ปีติปาโมชฺเชน วิหาตพฺพํ อโห\-รตฺตา\-นุสิกฺขินา กุสเลสุ ธมฺเมสุ.}

\prose{ปุน จปรํ อาวุโส ภิกฺขุนา อตฺตนาว อตฺตานํ เอวํ ปจฺจเวกฺขิตพฺพํ “กิํ นุ โขมฺหิ สโฐ มายาวี”ติ, สเจ อาวุโส ภิกฺขุ ปจฺจเวกฺขมาโน เอวํ ชานาติ “สโฐ โขมฺหิ มายาวี”ติ, เตนาวุโส ภิกฺขุนา เตสํเยว ปาปกานํ อกุสลานํ ธมฺมานํ ปหานาย วายมิตพฺพํ. สเจ ปนาวุโส ภิกฺขุ ปจฺจเวกฺขมาโน เอวํ ชานาติ “อสโฐ โขมฺหิ อมายาวี”ติ, เตนาวุโส ภิกฺขุนา เตเนว ปีติปาโมชฺเชน วิหาตพฺพํ อโห\-รตฺตา\-นุสิกฺขินา กุสเลสุ ธมฺเมสุ.}

\prose{ปุน จปรํ อาวุโส ภิกฺขุนา อตฺตนาว อตฺตานํ เอวํ ปจฺจเวกฺขิตพฺพํ “กิํ นุ โขมฺหิ ถทฺโธ อติมานี”ติ. สเจ อาวุโส ภิกฺขุ ปจฺจเวกฺขมาโน เอวํ ชานาติ “ถทฺโธ โขมฺหิ อติมานี”ติ, เตนาวุโส ภิกฺขุนา เตสํเยว ปาปกานํ อกุสลานํ ธมฺมานํ ปหานาย วายมิตพฺพํ. สเจ ปนาวุโส ภิกฺขุ ปจฺจเวกฺขมาโน เอวํ ชานาติ “อตฺถทฺโธ โขมฺหิ อนติมานี”ติ, เตนาวุโส ภิกฺขุนา เตเนว ปีติปาโมชฺเชน วิหาตพฺพํ อโห\-รตฺตา\-นุสิกฺขินา กุสเลสุ ธมฺเมสุ.}

\prose{ปุน จปรํ อาวุโส ภิกฺขุนา อตฺตนาว อตฺตานํ เอวํ ปจฺจเวกฺขิตพฺพํ “กิํ นุ โขมฺหิ สนฺทิฏฺฐิ\-ปรามาสี อาธานคฺคาหี ทุปฺปฏินิสฺสคฺคี”ติ. สเจ}

\prose{\csromanfolio{144}อาวุโส ภิกฺขุ ปจฺจเวกฺขมาโน เอวํ ชานาติ “สนฺทิฏฺฐิ\-ปรามาสี โขมฺหิ อาธานคฺคาหี ทุปฺปฏินิสฺสคฺคี”ติ, เตนาวุโส ภิกฺขุนา เตสํเยว ปาปกานํ อกุสลานํ ธมฺมานํ ปหานาย วายมิตพฺพํ. สเจ ปนาวุโส ภิกฺขุ ปจฺจเวกฺขมาโน เอวํ ชานาติ “อสนฺทิฏฺฐิ\-ปรามาสี โขมฺหิ อนาธานคฺคาหี สุปฺปฏินิสฺสคฺคี”ติ, เตนาวุโส ภิกฺขุนา เตเนว ปีติปาโมชฺเชน วิหาตพฺพํ อโห\-รตฺตา\-นุสิกฺขินา กุสเลสุ ธมฺเมสุ.}

\prose{สเจ อาวุโส ภิกฺขุ ปจฺจเวกฺขมาโน สพฺเพปิเม ปาปเก อกุสเล ธมฺเม อปฺปหีเน อตฺตนิ สมนุปสฺสติ, เตนาวุโส ภิกฺขุนา สพฺเพสํเยว อิเมสํ ปาปกานํ อกุสลานํ ธมฺมานํ ปหานาย วายมิตพฺพํ. สเจ ปนาวุโส ภิกฺขุ ปจฺจเวกฺขมาโน สพฺเพปิเม ปาปเก อกุสเล ธมฺเม ปหีเน อตฺตนิ สมนุปสฺสติ, เตนาวุโส ภิกฺขุนา เตเนว ปีติปาโมชฺเชน วิหาตพฺพํ อโห\-รตฺตา\-นุสิกฺขินา กุสเลสุ ธมฺเมสุ.}

\prose{เสยฺยถาปิ อาวุโส อิตฺถี วา ปุริโส วา ทหโร ยุวา มณฺฑนชาติโก อาทาเส วา ปริสุทฺเธ ปริโยทาเต อจฺเฉ วา อุทกปตฺเต สกํ มุขนิมิตฺตํ ปจฺจเวกฺขมาโน สเจ ตตฺถ ปสฺสติ รชํ วา องฺคณํ วา, ตสฺเสว รชสฺส วา องฺคณสฺส วา ปหานาย วายมติ. โน เจ ตตฺถ ปสฺสติ รชํ วา องฺคณํ วา, เตเนว อตฺตมโน โหติ “ลาภา วต เม, ปริสุทฺธํ วต เม”ติ. เอวเมว โข อาวุโส สเจ ภิกฺขุ ปจฺจเวกฺขมาโน สพฺเพปิเม ปาปเก อกุสเล ธมฺเม อปฺปหีเน อตฺตนิ สมนุปสฺสติ, เตนาวุโส ภิกฺขุนา สพฺเพสํเยว อิเมสํ ปาปกานํ อกุสลานํ ธมฺมานํ ปหานาย วายมิตพฺพํ. สเจ ปนาวุโส ภิกฺขุ ปจฺจเวกฺขมาโน สพฺเพปิเม ปาปเก อกุสเล ธมฺเม ปหีเน อตฺตนิ สมนุปสฺสติ, เตนาวุโส ภิกฺขุนา เตเนว ปีติปาโมชฺเชน วิหาตพฺพํ อโห\-รตฺตา\-นุสิกฺขินา กุสเลสุ ธมฺเมสูติ.}

\prosewithcloser{อิทมโวจายสฺมา มหาโมคฺคลฺลาโน. อตฺตมนา เต ภิกฺขู อายสฺมโต มหา\-โมคฺคลฺลานสฺส ภาสิตํ อภินนฺทุนฺติ.}{อนุมานสุตฺตํ นิฏฺฐิตํ ปญฺจมํ.}
\csromanmatikaanchor{mtk.44}
\csromantocmarkref{subsection}{6. เจโตขิลสุตฺต}{mtk.44}
\bookmark[dest={mtk.44},level=2]{6. เจโตขิลสุตฺต}
\csromanmarkh{6. เจโตขิลสุตฺต (16)}\csromanmarkhii{6. เจโตขิลสุตฺต}\csromanheadernomark{2}{6. เจโตขิลสุตฺต (16) \textbf{6. เจโตขิลสุตฺต}}
\proseitem{185}{\csromanfolio{145}เอวํ เม สุตํ\csromandash{}เอกํ สมยํ ภควา สาวตฺถิยํ วิหรติ เชตวเน อนาถ\-ปิณฺฑิกสฺส อาราเม. ตตฺร โข ภควา ภิกฺขู อามนฺเตสิ “ภิกฺขโว”ติ. “ภทนฺเต”ติ เต ภิกฺขู ภควโต ปจฺจสฺโสสุํ. ภควา เอตทโวจ\csromandash{}}

\prose{ยสฺส กสฺสจิ ภิกฺขเว ภิกฺขุโน ปญฺจ เจโตขิลา อปฺปหีนา, ปญฺจ เจตโส\-วินิพนฺธา\csromansharedfootnote{e620c2f21332a9aa}{เจตโสวินิพทฺธา (สี), เจโตวินิพทฺธา (สารตฺถทีปนีฏีกา)} อสมุจฺฉินฺนา, โส วติมสฺมิํ ธมฺมวินเย วุทฺธิํ วิรูฬฺหิํ เวปุลฺลํ อาปชฺชิสฺสตีติ เนตํ ฐานํ วิชฺชติ.}

\prose{กตมาสฺส ปญฺจ เจโตขิลา อปฺปหีนา โหนฺติ. อิธ ภิกฺขเว ภิกฺขุ สตฺถริ กงฺขติ วิจิกิจฺฉติ นาธิมุจฺจติ น สมฺปสีทติ. โย โส ภิกฺขเว ภิกฺขุ สตฺถริ กงฺขติ วิจิกิจฺฉติ นาธิมุจฺจติ น สมฺปสีทติ, ตสฺส จิตฺตํ น นมติ อาตปฺปาย อนุโยคาย สาตจฺจาย ปธานาย. ยสฺส จิตฺตํ น นมติ อาตปฺปาย อนุโยคาย สาตจฺจาย ปธานาย, เอวมสฺสายํ ปฐโม เจโตขิโล อปฺปหีโน โหติ.}

\prose{ปุน จปรํ ภิกฺขเว ภิกฺขุ ธมฺเม กงฺขติ วิจิกิจฺฉติ นาธิมุจฺจติ น สมฺปสีทติ ฯเปฯ เอวมสฺสายํ ทุติโย เจโตขิโล อปฺปหีโน โหติ.}

\prose{ปุน จปรํ ภิกฺขเว ภิกฺขุ สํเฆ กงฺขติ วิจิกิจฺฉติ นาธิมุจฺจติ น สมฺปสีทติ ฯเปฯ เอวมสฺสายํ ตติโย เจโตขิโล อปฺปหีโน โหติ.}

\prose{ปุน จปรํ ภิกฺขเว ภิกฺขุ สิกฺขาย กงฺขติ วิจิกิจฺฉติ นาธิมุจฺจติ น สมฺปสีทติ. โย โส ภิกฺขเว ภิกฺขุ สิกฺขาย กงฺขติ วิจิกิจฺฉติ นาธิมุจฺจติ น สมฺปสีทติ, ตสฺส จิตฺตํ น นมติ อาตปฺปาย อนุโยคาย สาตจฺจาย ปธานาย. ยสฺส จิตฺตํ น นมติ อาตปฺปาย อนุโยคาย สาตจฺจาย ปธานาย, เอวมสฺสายํ จตุตฺโถ เจโตขิโล อปฺปหีโน โหติ.}

\prose{ปุน จปรํ ภิกฺขเว ภิกฺขุ สพฺรหฺมจารีสุ กุปิโต โหติ อนตฺตมโน อาหตจิตฺโต ขิลชาโต. โย โส ภิกฺขเว ภิกฺขุ สพฺรหฺมจารีสุ กุปิโต โหติ อนตฺตมโน อาหตจิตฺโต ขิลชาโต, ตสฺส จิตฺตํ น นมติ อาตปฺปาย อนุโยคาย สาตจฺจาย ปธานาย. ยสฺส จิตฺตํ น \csromanfolio{146}นมติ อาตปฺปาย อนุโยคาย สาตจฺจาย ปธานาย, เอวมสฺสายํ ปญฺจโม เจโตขิโล อปฺปหีโน โหติ. อิมาสฺส ปญฺจ เจโตขิลา อปฺปหีนา โหนฺติ.}

\proseitem{186}{กตมาสฺส ปญฺจ เจตโส\-วินิพนฺธา อสมุจฺฉินฺนา โหนฺติ. อิธ ภิกฺขเว ภิกฺขุ กาเม อวีตราโค\csromansharedfootnote{ea10b88209f3a7a4}{อวิคตราโค (กตฺถจิ)} โหติ อวิคตจฺฉนฺโท อวิคตเปโม อวิคตปาโส อวิคตปริฬาโห อวิคตตโณฺห. โย โส ภิกฺขเว ภิกฺขุ กาเม อวีตราโค โหติ อวิคตจฺฉนฺโท อวิคตเปโม อวิคตปิปาโส อวิคตปริฬาโห อวิคตตโณฺห, ตสฺส จิตฺตํ น นมติ อาตปฺปาย อนุโยคาย สาตจฺจาย ปธานาย. ยสฺส จิตฺตํ น นมติ อาตปฺปาย อนุโยคาย สาตจฺจาย ปธานาย, เอวมสฺสายํ ปฐโม เจตโส\-วินิพนฺโธ อสมุจฺฉินฺโน โหติ.}

\prose{ปุน จปรํ ภิกฺขเว ภิกฺขุ กาเย อวีตราโค โหติ ฯเปฯ เอวมสฺสายํ ทุติโย เจตโส\-วินิพนฺโธ อสมุจฺฉินฺโน โหติ.}

\prose{ปุน จปรํ ภิกฺขเว ภิกฺขุ รูเป อวีตราโค โหติ ฯเปฯ เอวมสฺสายํ ตติโย เจตโส\-วินิพนฺโธ อสมุจฺฉินฺโน โหติ.}

\prose{ปุน จปรํ ภิกฺขเว ภิกฺขุ ยาวทตฺถํ อุทราวเทหกํ ภุญฺชิตฺวา เสยฺยสุขํ ปสฺสสุขํ มิทฺธสุขํ อนุยุตฺโต วิหรติ. โย โส ภิกฺขเว ภิกฺขุ ยาวทตฺถํ อุทราวเทหกํ ภุญฺชิตฺวา เสยฺยสุขํ ปสฺสสุขํ มิทฺธสุขํ อนุยุตฺโต วิหรติ, ตสฺส จิตฺตํ น นมติ อาตปฺปาย อนุโยคาย สาตจฺจาย ปธานาย. ยสฺส จิตฺตํ น นมติ อาตปฺปาย อนุโยคาย สาตจฺจาย ปธานาย, เอวมสฺสายํ จตุตฺโถ เจตโส\-วินิพนฺโธ อสมุจฺฉินฺโน โหติ.}

\prose{ปุน จปรํ ภิกฺขเว ภิกฺขุ อญฺญตรํ เทวนิกายํ ปณิธาย พฺรหฺมจริยํ จรติ “อิมินาหํ สีเลน วา วเตน วา ตเปน วา พฺรหฺมจริเยน วา เทโว วา ภวิสฺสามิ เทวญฺญตโร วา”ติ. โย โส ภิกฺขเว ภิกฺขุ อญฺญตรํ เทวนิกายํ ปณิธาย พฺรหฺมจริยํ จรติ “อิมินาหํ สีเลน วา วเตน วา ตเปน วา พฺรหฺมจริเยน วา เทโว วา ภวิสฺสามิ เทวญฺญตโร วา”ติ, ตสฺส จิตฺตํ น นมติ อาตปฺปาย อนุโยคาย สาตจฺจาย \csromanfolio{147}ปธานาย. ยสฺส จิตฺตํ น นมติ อาตปฺปาย อนุโยคาย สาตจฺจาย ปธานาย, เอวมสฺสายํ ปญฺจโม เจตโส\-วินิพนฺโธ อสมุจฺฉินฺโน โหติ. อิมาสฺส ปญฺจ เจตโส\-วินิพนฺธา อสมุจฺฉินฺนา โหนฺติ.}

\prose{ยสฺส กสฺสจิ ภิกฺขเว ภิกฺขุโน อิเม ปญฺจ เจโตขิลา อปฺปหีนา, อิเม ปญฺจ เจตโส\-วินิพนฺธา อสมุจฺฉินฺนา, โส วติมสฺมิํ ธมฺมวินเย วุทฺธิํ วิรูฬฺหิํ เวปุลฺลํ อาปชฺชิสฺสตีติ เนตํ ฐานํ วิชฺชติ.}

\proseitem{187}{ยสฺส กสฺสจิ ภิกฺขเว ภิกฺขุโน ปญฺจ เจโตขิลา ปหีนา, ปญฺจ เจตโส\-วินิพนฺธา สุสมุจฺฉินฺนา, โส วติมสฺมิํ ธมฺมวินเย วุทฺธิํ วิรูฬฺหิํ เวปุลฺลํ อาปชฺชิสฺสตีติ ฐานเมตํ วิชฺชติ.}

\prose{กตมาสฺส ปญฺจ เจโตขิลา ปหีนา โหนฺติ. อิธ ภิกฺขเว ภิกฺขุ สตฺถริ น กงฺขติ น วิจิกิจฺฉติ, อธิมุจฺจติ สมฺปสีทติ. โย โส ภิกฺขเว ภิกฺขุ สตฺถริ น กงฺขติ น วิจิกิจฺฉติ, อธิมุจฺจติ สมฺปสีทติ, ตสฺส จิตฺตํ นมติ อาตปฺปาย อนุโยคาย สาตจฺจาย ปธานาย. ยสฺส จิตฺตํ นมติ อาตปฺปาย อนุโยคาย สาตจฺจาย ปธานาย, เอวมสฺสายํ ปฐโม เจโตขิโล ปหีโน โหติ.}

\prose{ปุน จปรํ ภิกฺขเว ภิกฺขุ ธมฺเม น กงฺขติ น วิจิกิจฺฉติ, อธิมุจฺจติ สมฺปสีทติ ฯเปฯ เอวมสฺสายํ ทุติโย เจโตขิโล ปหีโน โหติ.}

\prose{ปุน จปรํ ภิกฺขเว ภิกฺขุ สํเฆ น กงฺขติ น วิจิกิจฺฉติ, อธิมุจฺจติ สมฺปสีทติ ฯเปฯ เอวมสฺสายํ ตติโย เจโตขิโล ปหีโน โหติ.}

\prose{ปุน จปรํ ภิกฺขเว ภิกฺขุ สิกฺขาย น กงฺขติ น วิจิกิจฺฉติ, อธิมุจฺจติ สมฺปสีทติ ฯเปฯ เอวมสฺสายํ จตุตฺโถ เจโตขิโล ปหีโน โหติ.}

\prose{ปุน จปรํ ภิกฺขเว ภิกฺขุ สพฺรหฺมจารีสุ น กุปิโต โหติ น อนตฺตมโน\csromansharedfootnote{f7b79e1f7e9a1479}{อตฺตมโน (สี, อิ)} อนาหตจิตฺโต อขิลชาโต. โย โส ภิกฺขเว ภิกฺขุ สพฺรหฺมจารีสุ น กุปิโต โหติ น อนตฺตมโน อนาหตจิตฺโต อขิลชาโต, ตสฺส จิตฺตํ นมติ อาตปฺปาย อนุโยคาย สาตจฺจาย ปธานาย. ยสฺส จิตฺตํ นมติ อาตปฺปาย อนุโยคาย สาตจฺจาย ปธานาย, เอวมสฺสายํ ปญฺจโม เจโตขิโล ปหีโน โหติ. อิมาสฺส ปญฺจ เจโตขิลา ปหีนา โหนฺติ.}

\proseitem{188}{\csromanfolio{148}กตมาสฺส ปญฺจ เจตโส\-วินิพนฺธา สุสมุจฺฉินฺนา โหนฺติ. อิธ ภิกฺขเว ภิกฺขุ กาเม วีตราโค โหติ วิคตจฺฉนฺโท วิคตเปโม วิคตปิปาโส วิคตปริฬาโห วิคตตโณฺห. โย โส ภิกฺขเว ภิกฺขุ กาเม วีตราโค โหติ วิคตจฺฉนฺโท วิคตเปโม วิคตปิปาโส วิคตปริฬาโห วิคตตโณฺห, ตสฺส จิตฺตํ นมติ อาตปฺปาย อนุโยคาย สาตจฺจาย ปธานาย. ยสฺส จิตฺตํ นมติ อาตปฺปาย อนุโยคาย สาตจฺจาย ปธานาย, เอวมสฺสายํ ปฐโม เจตโส\-วินิพนฺโธ สุสมุจฺฉินฺโน โหติ.}

\prose{ปุน จปรํ ภิกฺขเว ภิกฺขุ กาเย วีตราโค โหติ ฯเปฯ รูเป วีตราโค โหติ ฯเปฯ น ยาวทตฺถํ อุทราวเทหกํ ภุญฺชิตฺวา เสยฺยสุขํ ปสฺสสุขํ มิทฺธสุขํ อนุยุตฺโต วิหรติ. โย โส ภิกฺขเว ภิกฺขุ น ยาวทตฺถํ อุทราวเทหกํ ภุญฺชิตฺวา เสยฺยสุขํ ปสฺสสุขํ มิทฺธสุขํ อนุยุตฺโต วิหรติ, ตสฺส จิตฺตํ นมติ อาตปฺปาย อนุโยคาย สาตจฺจาย ปธานาย. ยสฺส จิตฺตํ นมติ อาตปฺปาย อนุโยคาย สาตจฺจาย ปธานาย, เอวมสฺสายํ จตุตฺโถ เจตโส\-วินิพนฺโธ สุสมุจฺฉินฺโน โหติ.}

\prose{ปุน จปรํ ภิกฺขเว ภิกฺขุ น อญฺญตรํ เทวนิกายํ ปณิธาย พฺรหฺมจริยํ จรติ “อิมินาหํ สีเลน วา วเตน วา ตเปน วา พฺรหฺมจริเยน วา เทโว วา ภวิสฺสามิ เทวญฺญตโร วา”ติ. โย โส ภิกฺขเว ภิกฺขุ น อญฺญตรํ เทวนิกายํ ปณิธาย พฺรหฺมจริยํ จรติ “อิมินาหํ สีเลน วา วเตน วา ตเปน วา พฺรหฺมจริเยน วา เทโว วา ภวิสฺสามิ เทวญฺญตโร วา”ติ, ตสฺส จิตฺตํ นมติ อาตปฺปาย อนุโยคาย สาตจฺจาย ปธานาย. ยสฺส จิตฺตํ นมติ อาตปฺปาย อนุโยคาย สาตจฺจาย ปธานาย, เอวมสฺสายํ ปญฺจโม เจตโส\-วินิพนฺโธ สุสมุจฺฉินฺโน โหติ. อิมาสฺส ปญฺจ เจตโส\-วินิพนฺธา สุสมุจฺฉินฺนา โหนฺติ.}

\prose{ยสฺส กสฺสจิ ภิกฺขเว ภิกฺขุโน อิเม ปญฺจ เจโตขิลา ปหีนา, อิเม ปญฺจ เจตโส\-วินิพนฺธา สุสมุจฺฉินฺนา, โส วติมสฺมิํ ธมฺมวินเย วุทฺธิํ วิรูฬฺหิํ เวปุลฺลํ อาปชฺชิสฺสตีติ ฐานเมตํ วิชฺชติ.}

\proseitem{189}{โส ฉนฺท\-สมาธิ\-ปธาน\-สงฺขาร\-สมนฺนาคตํ อิทฺธิปาทํ ภาเวติ. วีริย\-สมาธิ\-ปธาน\-สงฺขาร\-สมนฺนาคตํ อิทฺธิปาทํ ภาเวติ. จิตฺตสมาธิ- \csromanfolio{149}ปธานสงฺขาร\-สมนฺนาคตํ อิทฺธิปาทํ ภาเวติ. วีมํสา\-สมาธิ\-ปธาน\-สงฺขาร\-สมนฺนาคตํ อิทฺธิปาทํ ภาเวติ. อุสฺโสฬฺหีเยว ปญฺจมี. ส โข โส ภิกฺขเว เอวํ อุสฺโสฬฺหิ\-ปนฺนรส\-งฺค\-สมนฺนาคโต ภิกฺขุ ภพฺโพ อภินิพฺพิทาย, ภพฺโพ สมฺโพธาย, ภพฺโพ อนุตฺตรสฺส โยคกฺเขมสฺส อธิคมาย. เสยฺยถาปิ ภิกฺขเว กุกฺกุฏิยา อณฺฑานิ อฏฺฐ วา ทส วา ทฺวาทส วา, ตานสฺสุ กุกฺกุฏิยา สมฺมา อธิสยิตานิ สมฺมา ปริเสทิตานิ สมฺมา ปริภาวิตานิ. กิญฺจาปิ ตสฺสา กุกฺกุฏิยา น เอวํ อิจฺฉา อุปฺปชฺเชยฺย “อโห วติเม กุกฺกุฏโปตกา ปาท\-นข\-สิขาย วา มุขตุณฺฑเกน วา อณฺฑโกสํ ปทาเลตฺวา โสตฺถินา อภินิพฺภิ ชฺเชยฺยุนฺ”ติ. อถ โข ภพฺพาว เต กุกฺกุฏโปตกา ปาท\-นข\-สิขาย วา มุขตุณฺฑเกน วา อณฺฑโกสํ ปทาเลตฺวา โสตฺถินา อภินิ\-พฺภิชฺชิตุํ. เอวเมว โข ภิกฺขเว เอวํ อุสฺโสฬฺหิ\-ปนฺนรส\-งฺค\-สมนฺนาคโต ภิกฺขุ ภพฺโพ อภินิพฺพิทาย, ภพฺโพ สมฺโพธาย, ภพฺโพ อนุตฺตรสฺส โยคกฺเขมสฺส อธิคมายาติ.}

\prosewithcloserruled{อิทมโวจ ภควา. อตฺตมนา เต ภิกฺขู ภควโต ภาสิตํ อภินนฺทุนฺติ.}{เจโตขิล\-สุตฺตํ นิฏฺฐิตํ ฉฏฺฐํ.}
\csromanmatikaanchor{mtk.45}
\csromantocmarkref{subsection}{7. วนปตฺถสุตฺต}{mtk.45}
\bookmark[dest={mtk.45},level=2]{7. วนปตฺถสุตฺต}
\csromanheader{2}{7. วนปตฺถ\-สุตฺต}
\proseitem{190}{เอวํ เม สุตํ\csromandash{}เอกํ สมยํ ภควา สาวตฺถิยํ วิหรติ เชตวเน อนาถ\-ปิณฺฑิกสฺส อาราเม. ตตฺร โข ภควา ภิกฺขู อามนฺเตสิ “ภิกฺขโว”ติ. “ภทนฺเต”ติ เต ภิกฺขู ภควโต ปจฺจสฺโสสุํ. ภควา เอตทโวจ “วนปตฺถ\-ปริยายํ โว ภิกฺขเว เทเสสฺสามิ, ตํ สุณาถ สาธุกํ มนสิกโรถ ภาสิสฺสามี”ติ. “เอวํ ภนฺเต”ติ โข เต ภิกฺขู ภควโต ปจฺจสฺโสสุํ. ภควา เอตทโวจ\csromandash{}}

\proseitem{191}{อิธ ภิกฺขเว ภิกฺขุ อญฺญตรํ วนปตฺถํ อุปนิสฺสาย วิหรติ, ตสฺส ตํ วนปตฺถํ อุปนิสฺสาย วิหรโต อนุปฏฺฐิตา เจว สติ น อุปฏฺฐาติ, \csromanfolio{150}อสมาหิตญฺจ จิตฺตํ น สมาธิยติ, อปริกฺขีณา จ อาสวา น ปริกฺขยํ คจฺฉนฺติ, อนนุปฺปตฺตญฺจ อนุตฺตรํ โยคกฺเขมํ นานุปาปุณาติ. เย จ โข อิเม ปพฺพชิเตน ชีวิต\-ปริกฺขารา สมุทาเนตพฺพา จีวร\-ปิณฺฑปาต\-เสนาสน\-คิลานปฺปจฺจย\-เภสชฺช\-ปริกฺขารา, เต กสิเรน สมุทา\-ค\-จฺฉนฺติ. เตน ภิกฺขเว ภิกฺขุนา อิติ ปฏิสญฺจิกฺขิตพฺพํ “อหํ โข อิมํ วนปตฺถํ อุปนิสฺสาย วิหรามิ, ตสฺส เม อิมํ วนปตฺถํ อุปนิสฺสาย วิหรโต อนุปฏฺฐิตา เจว สติ น อุปฏฺฐาติ, อสมาหิตญฺจ จิตฺตํ น สมาธิยติ, อปริกฺขีณา จ อาสวา น ปริกฺขยํ คจฺฉนฺติ, อนนุปฺปตฺตญฺจ อนุตฺตรํ โยคกฺเขมํ นานุปาปุณามิ. เย จ โข อิเม ปพฺพชิเตน ชีวิต\-ปริกฺขารา สมุทาเนตพฺพา จีวร\-ปิณฺฑปาต\-เสนาสน คิลานปฺปจฺจย\-เภสชฺช\-ปริกฺขารา, เต กสิเรน สมุทาคจฺฉนฺตี”ติ. เตน ภิกฺขเว ภิกฺขุนา รตฺติภาคํ วา ทิวสภาคํ วา ตมฺหา วนปตฺถา ปกฺกมิตพฺพํ, น วตฺถพฺพํ.}

\proseitem{192}{อิธ ปน ภิกฺขเว ภิกฺขุ อญฺญตรํ วนปตฺถํ อุปนิสฺสาย วิหรติ, ตสฺส ตํ วนปตฺถํ อุปนิสฺสาย วิหรโต อนุปฏฺฐิตา เจว สติ น อุปฏฺฐาติ, อสมาหิตญฺจ จิตฺตํ น สมาธิยติ, อปริกฺขีณา จ อาสวา น ปริกฺขยํ คจฺฉนฺติ, อนนุปฺปตฺตญฺจ อนุตฺตรํ โยคกฺเขมํ นานุปาปุณาติ. เย จ โข อิเม ปพฺพชิเตน ชีวิต\-ปริกฺขารา สมุทาเนตพฺพา จีวร\-ปิณฺฑปาต\-เสนาสน\-คิลานปฺปจฺจย\-เภสชฺช\-ปริกฺขารา, เต อปฺปกสิเรน สมุทา\-ค\-จฺฉนฺติ. เตน ภิกฺขเว ภิกฺขุนา อิติ ปฏิสญฺจิกฺขิตพฺพํ “อหํ โข อิมํ วนปตฺถํ อุปนิสฺสาย วิหรามิ, ตสฺสา เม อิมํ วนปตฺถํ อุปนิสฺสาย วิหรโต อนุปฏฺฐิตา เจว สติ น อุปฏฺฐาติ, อสมาหิตญฺจ จิตฺตํ น สมาธิยติ, อปริกฺขีณา จ อาสวา น ปริกฺขยํ คจฺฉนฺติ, อนนุปฺปตฺตญฺจ อนุตฺตรํ โยคกฺเขมํ นานุปาปุณามิ. เย จ โข อิเม ปพฺพชิเตน ชีวิต\-ปริกฺขารา สมุทาเนตพฺพา จีวร\-ปิณฺฑปาต\-เสนาสน\-คิลานปฺปจฺจย\-เภสชฺช\-ปริกฺขารา, เต อปฺปกสิเรน สมุทา\-ค\-จฺฉนฺติ. น โข ปนาหํ จีวรเหตุ อคารสฺมา อนคาริยํ ปพฺพชิโต. น ปิณฺฑปาตเหตุ ฯเปฯ. น เสนาสนเหตุ ฯเปฯ. น คิลานปฺปจฺจย\-เภสชฺช\-ปริกฺขาร\-เหตุ อคารสฺมา อนคาริยํ ปพฺพชิโต. อถ จ ปน เม อิมํ วนปตฺถํ อุปนิสฺสาย วิหรโต อนุปฏฺฐิตา เจว สติ น อุปฏฺฐาติ, อสมาหิตญฺจ จิตฺตํ น สมาธิยติ, อปริกฺขีณา จ อาสวา น ปริกฺขยํ คจฺฉนฺติ, อนนุปฺปตฺตญฺจ อนุตฺตรํ โยคกฺเขมํ นานุปาปุณามี”ติ. เตน ภิกฺขเว ภิกฺขุนา สงฺขาปิ ตมฺหา วนปตฺถา ปกฺกมิตพฺพํ, น วตฺถพฺพํ.}

\proseitem{193}{\csromanfolio{151}อิธ ปน ภิกฺขเว ภิกฺขุ อญฺญตรํ วนปตฺถํ อุปนิสฺสาย วิหรติ, ตสฺส ตํ วนปตฺถํ อุปนิสฺสาย วิหรโต อนุปฏฺฐิตา เจว สติ อุปฏฺฐาติ, อสมาหิตญฺจ จิตฺตํ สมาธิยติ, อปริกฺขีณา จ อาสวา ปริกฺขยํ คจฺฉนฺติ, อนนุปฺปตฺตญฺจ อนุตฺตรํ โยคกฺเขมํ อนุปาปุณาติ. เย จ โข อิเม ปพฺพชิเตน ชีวิต\-ปริกฺขารา สมุทาเนตพฺพา จีวร\-ปิณฺฑปาต\-เสนาสน\-คิลานปฺปจฺจย\-เภสชฺช\-ปริกฺขารา, เต กสิเรน สมุทา\-ค\-จฺฉนฺติ. เตน ภิกฺขเว ภิกฺขุนา อิติ ปฏิสญฺจิกฺขิตพฺพํ “อหํ โข อิมํ วนปตฺถํ อุปนิสฺสาย วิหรามิ, ตสฺส เม อิมํ วนปตฺถํ อุปนิสฺสาย วิหรโต อนุปฏฺฐิตา เจว สติ อุปฏฺฐาติ, อสมาหิตญฺจ จิตฺตํ สมาธิยติ, อปริกฺขีณา จ อาสวา ปริกฺขยํ คจฺฉนฺติ, อนนุปฺปตฺตญฺจ อนุตฺตรํ โยคกฺเขมํ อนุปาปุณามิ. เย จ โข อิเม ปพฺพชิเตน ชีวิต\-ปริกฺขารา สมุทาเนตพฺพา จีวร\-ปิณฺฑปาต\-เสนาสน\-คิลานปฺปจฺจย\-เภสชฺช\-ปริกฺขารา, เต กสิเรน สมุทา\-ค\-จฺฉนฺติ. น โข ปนาหํ จีวรเหตุ อคารสฺมา อนคาริยํ ปพฺพชิโต. น ปิณฺฑปาตเหตุ ฯเปฯ. น เสนาสนเหตุ ฯเปฯ. น คิลานปฺปจฺจย\-เภสชฺช\-ปริกฺขาร\-เหตุ อคารสฺมา อนคาริยํ ปพฺพชิโต. อถ จ ปน เม อิมํ วนปตฺถํ อุปนิสฺสาย วิหรโต อนุปฏฺฐิตา เจว สติ อุปฏฺฐาติ, อสมาหิตญฺจ จิตฺตํ สมาธิยติ, อปริกฺขีณา จ อาสวา ปริกฺขยํ คจฺฉนฺติ, อนนุปฺปตฺตญฺจ อนุตฺตรํ โยคกฺเขมํ อนุปาปุณามี”ติ. เตน ภิกฺขเว ภิกฺขุนา สงฺขาปิ ตสฺมิํ วนปตฺเถ วตฺถพฺพํ, น ปกฺกมิตพฺพํ.}

\proseitem{194}{อิธ ปน ภิกฺขเว ภิกฺขุ อญฺญตรํ วนปตฺถํ อุปนิสฺสาย วิหรติ, ตสฺส ตํ วนปตฺถํ อุปนิสฺสาย วิหรโต อนุปฏฺฐิตา เจว สติ อุปฏฺฐาติ, อสมาหิตญฺจ จิตฺตํ สมาธิยติ, อปริกฺขีณา จ อาสวา ปริกฺขยํ คจฺฉนฺติ, อนนุปฺปตฺตญฺจ อนุตฺตรํ โยคกฺเขมํ อนุปาปุณาติ. เย จ โข อิเม ปพฺพชิเตน ชีวิต\-ปริกฺขารา สมุทาเนตพฺพา จีวร\-ปิณฺฑปาต\-เสนาสน\-คิลานปฺปจฺจย\-เภสชฺช\-ปริกฺขารา, เต อปฺปกสิเรน สมุทา\-ค\-จฺฉนฺติ. เตน ภิกฺขเว ภิกฺขุนา อิติ ปฏิสญฺจิกฺขิตพฺพํ “อหํ โข อิมํ วนปตฺถํ อุปนิสฺสาย วิหรามิ, ตสฺส เม อิมํ วนปตฺถํ อุปนิสฺสาย วิหรโต อนุปฏฺฐิตา เจว สติ อุปฏฺฐาติ, อสมาหิตญฺจ จิตฺตํ สมาธิยติ, อปริกฺขีณา จ อาสวา ปริกฺขยํ คจฺฉนฺติ, อนนุปฺปตฺตญฺจ อนุตฺตรํ โยคกฺเขมํ อนุปาปุณามิ. เย จ โข อิเม ปพฺพชิเตน ชีวิต\-ปริกฺขารา สมุทาเนตพฺพา จีวร\-ปิณฺฑปาต\-เสนาสน\-คิลานปฺปจฺจย\-เภสชฺช\-ปริกฺขารา, เต อปฺปกสิเรน \csromanfolio{152}สมุทาคจฺฉนฺตี”ติ. เตน ภิกฺขเว ภิกฺขุนา ยาวชีวมฺปิ ตสฺมิํ วนปตฺเถ วตฺถพฺพํ, น ปกฺกมิตพฺพํ.}

\proseitem{195}{อิธ ภิกฺขเว ภิกฺขุ อญฺญตรํ คามํ อุปนิสฺสาย วิหรติ ฯเปฯ อญฺญตรํ นิคมํ อุปนิสฺสาย วิหรติ ฯเปฯ อญฺญตรํ นครํ อุปนิสฺสาย วิหรติ ฯเปฯ อญฺญตรํ ชนปทํ อุปนิสฺสาย วิหรติ ฯเปฯ อญฺญตรํ ปุคฺคลํ อุปนิสฺสาย วิหรติ, ตสฺส ตํ ปุคฺคลํ อุปนิสฺสาย วิหรโต อนุปฏฺฐิตา เจว สติ น อุปฏฺฐาติ, อสมาหิตญฺจ จิตฺตํ น สมาธิยติ, อปริกฺขีณา จ อาสวา น ปริกฺขยํ คจฺฉนฺติ, อนนุปฺปตฺตญฺจ อนุตฺตรํ โยคกฺเขมํ นานุปาปุณาติ. เย จ โข อิเม ปพฺพชิเตน ชีวิต\-ปริกฺขารา สมุทาเนตพฺพา จีวร\-ปิณฺฑปาต\-เสนาสน\-คิลานปฺปจฺจย\-เภสชฺช\-ปริกฺขารา, เต กสิเรน สมุทา\-ค\-จฺฉนฺติ. เตน ภิกฺขเว ภิกฺขุนา อิติ ปฏิสญฺจิกฺขิตพฺพํ “อหํ โข อิมํ ปุคฺคลํ อุปนิสฺสาย วิหรามิ, ตสฺส เม อิมํ ปุคฺคลํ อุปนิสฺสาย วิหรโต อนุปฏฺฐิตา เจว สติ น อุปฏฺฐาติ, อสมาหิตญฺจ จิตฺตํ น สมาธิยติ, อปริกฺขีณา จ อาสวา น ปริกฺขยํ คจฺฉนฺติ, อนนุปฺปตฺตญฺจ อนุตฺตรํ โยคกฺเขมํ นานุปาปุณามิ. เย จ โข อิเม ปพฺพชิเตน ชีวิต\-ปริกฺขารา สมุทาเนตพฺพา จีวร\-ปิณฺฑปาต\-เสนาสน\-คิลานปฺปจฺจย\-เภสชฺช\-ปริกฺขารา, เต กสิเรน สมุทาคจฺฉนฺตี”ติ. เตน ภิกฺขเว ภิกฺขุนา รตฺติภาคํ วา ทิวสภาคํ วา โส ปุคฺคโล อนาปุจฺฉา ปกฺกมิตพฺพํ, นานุพนฺธิตพฺโพ.}

\proseitem{196}{อิธ ปน ภิกฺขเว ภิกฺขุ อญฺญตรํ ปุคฺคลํ อุปนิสฺสาย วิหรติ, ตสฺส ตํ ปุคฺคลํ อุปนิสฺสาย วิหรโต อนุปฏฺฐิตา เจว สติ น อุปฏฺฐาติ, อสมาหิตญฺจ จิตฺตํ น สมาธิยติ, อปริกฺขีณา จ อาสวา น ปริกฺขยํ คจฺฉนฺติ, อนนุปฺปตฺตญฺจ อนุตฺตรํ โยคกฺเขมํ นานุปาปุณาติ. เย จ โข อิเม ปพฺพชิเตน ชีวิต\-ปริกฺขารา สมุทาเนตพฺพา จีวร\-ปิณฺฑปาต\-เสนาสน\-คิลานปฺปจฺจย\-เภสชฺช\-ปริกฺขารา, เต อปฺปกสิเรน สมุทา\-ค\-จฺฉนฺติ. เตน ภิกฺขเว ภิกฺขุนา อิติ ปฏิสญฺจิกฺขิตพฺพํ “อหํ โข อิมํ ปุคฺคลํ อุปนิสฺสาย วิหรามิ, ตสฺส เม อิมํ ปุคฺคลํ อุปนิสฺสาย วิหรโต อนุปฏฺฐิตา เจว สติ น อุปฏฺฐาติ, อสมาหิตญฺจ จิตฺตํ น สมาธิยติ, อปริกฺขีณา จ อาสวา น ปริกฺขยํ คจฺฉนฺติ, อนนุปฺปตฺตญฺจ อนุตฺตรํ โยคกฺเขมํ นานุปาปุณามิ. เย จ โข อิเม ปพฺพชิเตน ชีวิต\-ปริกฺขารา สมุทาเนตพฺพา \csromanfolio{153}จีวร\-ปิณฺฑปาต\-เสนาสน\-คิลานปฺปจฺจย\-เภสชฺช\-ปริกฺขารา, เต อปฺปกสิเรน สมุทา\-ค\-จฺฉนฺติ. น โข ปนาหํ จีวรเหตุ อคารสฺมา อนคาริยํ ปพฺพชิโต. น ปิณฺฑปาตเหตุ ฯเปฯ. น เสนาสนเหตุ ฯเปฯ. น คิลานปฺปจฺจย\-เภสชฺช\-ปริกฺขาร\-เหตุ อคารสฺมา อนคาริยํ ปพฺพชิโต. อถ จ ปน เม อิมํ ปุคฺคลํ อุปนิสฺสาย วิหรโต อนุปฏฺฐิตา เจว สติ น อุปฏฺฐาติ, อสมาหิตญฺจ จิตฺตํ น สมาธิยติ, อปริกฺขีณา จ อาสวา น ปริกฺขยํ คจฺฉนฺติ, อนนุปฺปตฺตญฺจ อนุตฺตรํ โยคกฺเขมํ นานุปาปุณามี”ติ. เตน ภิกฺขเว ภิกฺขุนา สงฺขาปิ โส ปุคฺคโล อาปุจฺฉา ปกฺกมิตพฺพํ, นานุพนฺธิตพฺโพ.}

\proseitem{197}{อิธ ปน ภิกฺขเว ภิกฺขุ อญฺญตรํ ปุคฺคลํ อุปนิสฺสาย วิหรติ, ตสฺส ตํ ปุคฺคลํ อุปนิสฺสาย วิหรโต อนุปฏฺฐิตา เจว สติ อุปฏฺฐาติ, อสมาหิตญฺจ จิตฺตํ สมาธิยติ, อปริกฺขีณา จ อาสวา ปริกฺขยํ คจฺฉนฺติ, อนนุปฺปตฺตญฺจ อนุตฺตรํ โยคกฺเขมํ อนุปาปุณาติ. เย จ โข อิเม ปพฺพชิเตน ชีวิต\-ปริกฺขารา สมุทาเนตพฺพา จีวร\-ปิณฺฑปาต\-เสนาสน\-คิลานปฺปจฺจย\-เภสชฺช\-ปริกฺขารา, เต กสิเรน สมุทา\-ค\-จฺฉนฺติ. เตน ภิกฺขเว ภิกฺขุนา อิติ ปฏิสญฺจิกฺขิตพฺพํ “อหํ โข อิมํ ปุคฺคลํ อุปนิสฺสาย วิหรามิ, ตสฺส เม อิมํ ปุคฺคลํ อุปนิสฺสาย วิหรโต อนุปฏฺฐิตา เจว สติ อุปฏฺฐาติ, อสมาหิตญฺจ จิตฺตํ สมาธิยติ, อปริกฺขีณา จ อาสวา ปริกฺขยํ คจฺฉนฺติ, อนนุปฺปตฺตญฺจ อนุตฺตรํ โยคกฺเขมํ อนุปาปุณามิ. เย จ โข อิเม ปพฺพชิเตน ชีวิต\-ปริกฺขารา สมุทาเนตพฺพา จีวร\-ปิณฺฑปาต\-เสนาสน\-คิลานปฺปจฺจย\-เภสชฺช\-ปริกฺขารา, เต กสิเรน สมุทา\-ค\-จฺฉนฺติ. น โข ปนาหํ จีวรเหตุ อคารสฺมา อนคาริยํ ปพฺพชิโต. น ปิณฺฑปาตเหตุ ฯเปฯ. น เสนาสนเหตุ ฯเปฯ. น คิลานปฺปจฺจย\-เภสชฺช\-ปริกฺขาร\-เหตุ อคารสฺมา อนคาริยํ ปพฺพชิโต. อถ จ ปน เม อิมํ ปุคฺคลํ อุปนิสฺสาย วิหรโต อนุปฏฺฐิตา เจว สติ อุปฏฺฐาติ, อสมาหิตญฺจ จิตฺตํ สมาธิยติ, อปริกฺขีณา จ อาสวา ปริกฺขยํ คจฺฉนฺติ, อนนุปฺปตฺตญฺจ อนุตฺตรํ โยคกฺเขมํ อนุปาปุณามี”ติ. เตน ภิกฺขเว ภิกฺขุนา สงฺขาปิ โส ปุคฺคโล อนุพนฺธิตพฺโพ, น ปกฺกมิตพฺพํ.}

\proseitem{198}{อิธ ปน ภิกฺขเว ภิกฺขุ อญฺญตรํ ปุคฺคลํ อุปนิสฺสาย วิหรติ, ตสฺส ตํ ปุคฺคลํ อุปนิสฺสาย วิหรโต อนุปฏฺฐิตา เจว สติ อุปฏฺฐาติ, อสมาหิตญฺจ จิตฺตํ สมาธิยติ, อปริกฺขีณา จ อาสวา ปริกฺขยํ \csromanfolio{154}คจฺฉนฺติ, อนนุปฺปตฺตญฺจ อนุตฺตรํ โยคกฺเขมํ อนุปาปุณาติ. เย จ โข อิเม ปพฺพชิเตน ชีวิต\-ปริกฺขารา สมุทาเนตพฺพา จีวร\-ปิณฺฑปาต\-เสนาสน\-คิลานปฺปจฺจย\-เภสชฺช\-ปริกฺขารา, เต อปฺปกสิเรน สมุทา\-ค\-จฺฉนฺติ. เตน ภิกฺขเว ภิกฺขุนา อิติ ปฏิสญฺจิกฺขิตพฺพํ “อหํ โข อิมํ ปุคฺคลํ อุปนิสฺสาย วิหรามิ, ตสฺส เม อิมํ ปุคฺคลํ อุปนิสฺสาย วิหรโต อนุปฏฺฐิตา เจว สติ อุปฏฺฐาติ, อสมาหิตญฺจ จิตฺตํ สมาธิยติ, อปริกฺขีณา จ อาสวา ปริกฺขยํ คจฺฉนฺติ, อนนุปฺปตฺตญฺจ อนุตฺตรํ โยคกฺเขมํ อนุปาปุณามิ. เย จ โข อิเม ปพฺพชิเตน ชีวิต\-ปริกฺขารา สมุทาเนตพฺพา จีวร\-ปิณฺฑปาต\-เสนาสนคิลานปฺปจฺจย- เภสชฺช\-ปริกฺขารา, เต อปฺปกสิเรน สมุทาคจฺฉนฺตี”ติ. เตน ภิกฺขเว ภิกฺขุนา ยาวชีวมฺปิ โส ปุคฺคโล อนุพนฺธิตพฺโพ, น ปกฺกมิตพฺพํ อปิ ปนุชฺชมาเนนปีติ\csromansharedfootnote{5b69084a457c11de}{อปิ ปณุชฺชมาเนนาติ (?)}.}

\prosewithcloserruled{อิทมโวจ ภควา. อตฺตมนา เต ภิกฺขู ภควโต ภาสิตํ อภินนฺทุนฺติ.}{วนปตฺถ\-สุตฺตํ นิฏฺฐิตํ สตฺตมํ.}
\csromanmatikaanchor{mtk.46}
\csromantocmarkref{subsection}{8. มธุปิณฺฑิกสุตฺต}{mtk.46}
\bookmark[dest={mtk.46},level=2]{8. มธุปิณฺฑิกสุตฺต}
\csromanheader{2}{8. มธุปิณฺฑิก\-สุตฺต}
\proseitem{199}{เอวํ เม สุตํ\csromandash{}เอกํ สมยํ ภควา สกฺเกสุ วิหรติ กปิล\-วตฺถุสฺมิํ นิโคฺรธาราเม. อถ โข ภควา ปุพฺพณฺหสมยํ นิวาเสตฺวา ปตฺต\-จีวรมา\-ทาย กปิลวตฺถุํ ปิณฺฑาย ปาวิสิ. กปิล\-วตฺถุสฺมิํ ปิณฺฑาย จริตฺวา ปจฺฉาภตฺตํ ปิณฺฑปาต\-ปฏิกฺกนฺโต เยน มหาวนํ เตนุปสงฺกมิ ทิวาวิหาราย. มหาวนํ อชฺโฌคาเหตฺวา เพลุว\-ลฏฺฐิกาย มูเล ทิวาวิหารํ นิสีทิ. ทณฺฑปาณิปิ โข สกฺโก ชงฺฆาวิหารํ\csromansharedfootnote{8a5855202d1617f9}{ชงฺฆวิหารํ (ก)} อนุจงฺกมมาโน อนุวิจรมาโน เยน มหาวนํ เตนุปสงฺกมิ, มหาวนํ อชฺโฌคาเหตฺวา เยน เพลุวลฏฺฐิกา เยน ภควา เตนุปสงฺกมิ, อุปสงฺกมิตฺวา ภควตา สทฺธิํ สมฺโมทิ, สมฺโมทนียํ กถํ สารณียํ วีติสาเรตฺวา ทณฺฑโมลุพฺภ เอกมนฺตํ อฏฺฐาสิ, เอกมนฺตํ ฐิโต โข ทณฺฑปาณิ สกฺโก ภควนฺตํ เอตทโวจ “กิํวาที สมโณ กิมกฺขายี”ติ. ยถาวาที โข อาวุโส สเทวเก โลเก สมารเก สพฺรหฺมเก สสฺสมณ\-พฺราหฺมณิยา \csromanfolio{155}ปชาย สเทว\-มนุสฺสาย น เกนจิ โลเก วิคฺคยฺห ติฏฺฐติ, ยถา จ ปน กาเมหิ วิสํยุตฺตํ วิหรนฺตํ ตํ พฺราหฺมณํ อกถํกถิํ ฉินฺน\-กุกฺกุจฺจํ ภวาภเว วีตตณฺหํ สญฺญา นานุเสนฺติ, เอวํวาที โข อหํ อาวุโส เอวมกฺขายี”ติ. เอวํ วุตฺเต ทณฺฑปาณิ สกฺโก สีสํ โอกมฺเปตฺวา ชิวฺหํ นิลฺลาเฬตฺวา ติวิสาขํ นลาฏิกํ นลาเฏ วุฏฺฐาเปตฺวา ทณฺฑโมลุพฺภ ปกฺกามิ.}

\proseitem{200}{อถ โข ภควา สายนฺหสมยํ ปฏิสลฺลานา วุฏฺฐิโต เยน นิโคฺรธาราโม เตนุปสงฺกมิ, อุปสงฺกมิตฺวา ปญฺญตฺเต อาสเน นิสีทิ, นิสชฺช โข ภควา ภิกฺขู อามนฺเตสิ “อิธาหํ ภิกฺขเว ปุพฺพณฺหสมยํ นิวาเสตฺวา ปตฺต\-จีวรมา\-ทาย กปิลวตฺถุํ ปิณฺฑาย ปาวิสิํ. กปิล\-วตฺถุสฺมิํ ปิณฺฑาย จริตฺวา ปจฺฉาภตฺตํ ปิณฺฑปาต\-ปฏิกฺกนฺโต เยน มหาวนํ เตนุปสงฺกมิํ ทิวาวิหาราย. มหาวนํ อชฺโฌคาเหตฺวา เพลุว\-ลฏฺฐิกาย มูเล ทิวาวิหารํ นิสีทิํ. ทณฺฑปาณิปิ โข ภิกฺขเว สกฺโก ชงฺฆาวิหารํ อนุจงฺกมมาโน อนุวิจรมาโน เยน มหาวนํ เตนุปสงฺกมิ. มหาวนํ อชฺโฌคาเหตฺวา เยน เพลุวลฏฺฐิกา เยนาหํ เตนุปสงฺกมิ, อุปสงฺกมิตฺวา มยา สทฺธิํ สมฺโมทิ, สมฺโมทนียํ กถํ สารณียํ วีติสาเรตฺวา ทณฺฑโมลุพฺภ เอกมนฺตํ อฏฺฐาสิ, เอกมนฺตํ ฐิโต โข ภิกฺขเว ทณฺฑปาณิ สกฺโก มํ เอตทโวจ ‘กิํวาที สมโณ กิมกฺขายี’ติ. เอวํ วุตฺเต อหํ ภิกฺขเว ทณฺฑปาณิํ สกฺกํ เอตทโวจํ ‘ยถาวาที โข อาวุโส สเทวเก โลเก สมารเก สพฺรหฺมเก สสฺสมณ\-พฺราหฺมณิยา ปชาย สเทว\-มนุสฺสาย น เกนจิ โลเก วิคฺคยฺห ติฏฺฐติ, ยถา จ ปน กาเมหิ วิสํยุตฺตํ วิหรนฺตํ ตํ พฺราหฺมณํ อกถํกถิํ ฉินฺน\-กุกฺกุจฺจํ ภวาภเว วีตตณฺหํ สญฺญา นานุเสนฺติ, เอวํวาที โข อหํ อาวุโส เอวมกฺขายี”ติ. เอวํ วุตฺเต ภิกฺขเว ทณฺฑปาณิ สกฺโก สีสํ โอกมฺเปตฺวา ชิวฺหํ นิลฺลาเฬตฺวา ติวิสาขํ นลาฏิกํ นลาเฏ วุฏฺฐาเปตฺวา ทณฺฑโมลุพฺภ ปกฺกามี”ติ.}

\proseitem{201}{เอวํ วุตฺเต อญฺญตโร ภิกฺขุ ภควนฺตํ เอตทโวจ “กิํวาที ปน ภนฺเต ภควา สเทวเก โลเก สมารเก สพฺรหฺมเก สสฺสมณ\-พฺราหฺมณิยา ปชาย สเทว\-มนุสฺสาย น เกนจิ โลเก วิคฺคยฺห ติฏฺฐติ. กถญฺจ ปน ภนฺเต ภควนฺตํ กาเมหิ วิสํยุตฺตํ วิหรนฺตํ \csromanfolio{156}ตํ พฺราหฺมณํ อกถํกถิํ ฉินฺน\-กุกฺกุจฺจํ ภวาภเว วีตตณฺหํ สญฺญา นานุเสนฺตี”ติ. ยโตนิทานํ ภิกฺขุ ปุริสํ ปปญฺจ\-สญฺญา\-สงฺขา สมุทาจรนฺติ, เอตฺถ เจ นตฺถิ อภินนฺทิตพฺพํ อภิวทิตพฺพํ อชฺโฌสิตพฺพํ, เอเสวนฺโต ราคานุสยานํ, เอเสวนฺโต ปฏิฆา\-นุสยานํ, เอเสวนฺโต ทิฏฺฐา\-นุสยานํ, เอเสวนฺโต วิจิกิจฺฉา\-นุสยานํ, เอเสวนฺโต มานานุสยานํ, เอเสวนฺโต ภวราคา\-นุสยานํ, เอเสวนฺโต อวิชฺชา\-นุสยานํ, เอเสวนฺโต ทณฺฑาทาน\-สตฺถาทาน กลห วิคฺคห วิวาท ตุวํตุวํ เปสุญฺญ มุสาวาทานํ. เอตฺเถเต ปาปกา อกุสลา ธมฺมา อปริเสสา นิรุชฺฌนฺตีติ. อิทมโวจ ภควา, อิทํ วตฺวาน สุคโต อุฏฺฐายาสนา วิหารํ ปาวิสิ.}

\proseitem{202}{อถ โข เตสํ ภิกฺขูนํ อจิร\-ปกฺกนฺตสฺส ภควโต เอตทโหสิ “อิทํ โข โน อาวุโส ภควา สํขิตฺเตน อุทฺเทสํ อุทฺทิสิตฺวา วิตฺถาเรน อตฺถํ อวิภชิตฺวา อุฏฺฐายาสนา วิหารํ ปวิฏฺโฐ, ‘ยโตนิทานํ ภิกฺขุ ปุริสํ ปปญฺจ\-สญฺญา\-สงฺขา สมุทาจรนฺติ. เอตฺถ เจ นตฺถิ อภินนฺทิตพฺพํ อภิวทิตพฺพํ อชฺโฌสิตพฺพํ, เอเสวนฺโต ราคานุสยานํ ฯเปฯ. เอตฺเถเต ปาปกา อกุสลา ธมฺมา อปริเสสา นิรุชฺฌนฺตี’ติ. โก นุ โข อิมสฺส ภควตา สํขิตฺเตน อุทฺเทสสฺส อุทฺทิฏฺฐสฺส วิตฺถาเรน อตฺถํ อวิภตฺตสฺส วิตฺถาเรน อตฺถํ วิภเชยฺยา”ติ. อถ โข เตสํ ภิกฺขูนํ เอตทโหสิ “อยํ โข อายสฺมา มหากจฺจาโน สตฺถุ เจว สํวณฺณิโต, สมฺภาวิโต จ วิญฺญูนํ สพฺรหฺมจารีนํ, ปโหติ จายสฺมา มหากจฺจาโน อิมสฺส ภควตา สํขิตฺเตน อุทฺเทสสฺส อุทฺทิฏฺฐสฺส วิตฺถาเรน อตฺถํ อวิภตฺตสฺส วิตฺถาเรน อตฺถํ วิภชิตุํ. ยํนูน มยํ เยนายสฺมา มหากจฺจาโน เตนุปสงฺกเมยฺยาม, อุปสงฺกมิตฺวา อายสฺมนฺตํ มหากจฺจานํ เอตมตฺถํ ปฏิปุจฺเฉยฺยามา”ติ.}

\prose{อถ โข เต ภิกฺขู เยนายสฺมา มหากจฺจาโน เตนุ\-ปสงฺกมิํสุ, อุปสงฺกมิตฺวา อายสฺมตา มหากจฺจาเนน สทฺธิํ สมฺโมทิํสุ, สมฺโมทนียํ กถํ สารณียํ วีติสาเรตฺวา เอกมนฺตํ นิสีทิํสุ, เอกมนฺตํ นิสินฺนา โข เต ภิกฺขู อายสฺมนฺตํ มหากจฺจานํ เอตทโวจุํ\csromandash{}อิทํ โข โน อาวุโส กจฺจาน ภควา สํขิตฺเตน อุทฺเทสํ อุทฺทิสิตฺวา วิตฺถาเรน อตฺถํ อวิภชิตฺวา อุฏฺฐายาสนา วิหารํ ปวิฏฺโฐ, “ยโตนิทานํ ภิกฺขุ ปุริสํ ปปญฺจ\-สญฺญา\-สงฺขา สมุทาจรนฺติ. เอตฺถ เจ นตฺถิ อภินนฺทิตพฺพํ อภิวทิตพฺพํ อชฺโฌสิตพฺพํ, \csromanfolio{157}เอเสวนฺโต ราคานุสยานํ ฯเปฯ. เอตฺเถเต ปาปกา อกุสลา ธมฺมา อปริเสสานิรุชฺฌนฺตี”ติ. เตสํ โน อาวุโส กจฺจาน อมฺหากํ อจิร\-ปกฺกนฺตสฺส ภควโต เอตทโหสิ “อิทํ โข โน อาวุโส ภควา สํขิตฺเตน อุทฺเทสํ อุทฺทิสิตฺวา วิตฺถาเรน อตฺถํ อวิภชิตฺวา อุฏฺฐายาสนา วิหารํ ปวิฏฺโฐ, ‘ยโตนิทานํ ภิกฺขุ ปุริสํ ปปญฺจ\-สญฺญา\-สงฺขา สมุทาจรนฺติ. เอตฺถ เจ นตฺถิ อภินนฺทิตพฺพํ อภิวทิตพฺพํ อชฺโฌสิตพฺพํ, เอเสวนฺโต ราคานุสยานํ ฯเปฯ. เอตฺเถเต ปาปกา อกุสลา ธมฺมา อปริเสสา นิรุชฺฌนฺตี’ติ. โก นุ โข อิมสฺส ภควตา สํขิตฺเตน อุทฺเทสสฺส อุทฺทิฏฺฐสฺส วิตฺถาเรน อตฺถํ อวิภตฺตสฺส วิตฺถาเรน อตฺถํ วิภเชยฺยา”ติ. เตสํ โน อาวุโส กจฺจาน อมฺหากํ เอตทโหสิ “อยํ โข อายสฺมา มหากจฺจาโน สตฺถุ เจว สํวณฺณิโต, สมฺภาวิโต จ วิญฺญูนํ สพฺรหฺมจารีนํ, ปโหติ จายสฺมา มหากจฺจาโน อิมสฺส ภควตา สํขิตฺเตน อุทฺเทสสฺส อุทฺทิฏฺฐสฺส วิตฺถาเรน อตฺถํ อวิภตฺตสฺส วิตฺถาเรน อตฺถํ วิภชิตุํ. ยํนูน มยํ เยนายสฺมา มหากจฺจาโน เตนุปสงฺกเมยฺยาม, อุปสงฺกมิตฺวา อายสฺมนฺตํ มหากจฺจานํ เอตมตฺถํ ปฏิปุจฺเฉยฺยามา”ติ. วิภชตายสฺมา มหากจฺจาโนติ.}

\proseitem{203}{เสยฺยถาปิ อาวุโส ปุริโส สารตฺถิโก สารคเวสี สารปริเยสนํ จรมาโน มหโต รุกฺขสฺส ติฏฺฐโต สารวโต อติกฺกมฺเมว มูลํ อติกฺกมฺม ขนฺธํ สาขาปลาเส สารํ ปริเยสิตพฺพํ มญฺเญยฺย, เอวํ\-สมฺปทมิทํ อายสฺมนฺตานํ, สตฺถริ สมฺมุขีภูเต ตํ ภควนฺตํ อติสิตฺวา อเมฺห เอตมตฺถํ ปฏิปุจฺฉิตพฺพํ มญฺญถ. โส หาวุโส ภควา ชานํ ชานาติ ปสฺสํ ปสฺสติ จกฺขุภูโต ญาณภูโต ธมฺมภูโต พฺรหฺมภูโต วตฺตา ปวตฺตา อตฺถสฺส นินฺเนตา อมตสฺส ทาตา ธมฺมสฺสามี ตถาคโต. โส เจว ปเนตสฺส กาโล อโหสิ, ยํ ภควนฺตํเยว เอตมตฺถํ ปฏิปุจฺเฉยฺยาถ. ยถา โว ภควา พฺยากเรยฺย, ตถา นํ ธาเรยฺยาถาติ. อทฺธาวุโส กจฺจาน ภควา ชานํ ชานาติ ปสฺสํ ปสฺสติ จกฺขุภูโต ญาณภูโต ธมฺมภูโต พฺรหฺมภูโต วตฺตา ปวตฺตา อตฺถสฺส นินฺเนตา อมตสฺส ทาตา ธมฺมสฺสามี ตถาคโต. โส เจว ปเนตสฺส กาโล อโหสิ, ยํ ภควนฺตํเยว เอตมตฺถํ ปฏิปุจฺเฉยฺยาม. ยถา โน ภควา พฺยากเรยฺย, ตถา นํ ธาเรยฺยาม. อปิ จายสฺมา มหากจฺจาโน สตฺถุ เจว สํวณฺณิโต, สมฺภาวิโต จ วิญฺญูนํ สพฺรหฺมจารีนํ, \csromanfolio{158}ปโหติ จายสฺมา มหากจฺจาโน อิมสฺส ภควตา สํขิตฺเตน อุทฺเทสสฺส อุทฺทิฏฺฐสฺส วิตฺถาเรน อตฺถํ อวิภตฺตสฺส วิตฺถาเรน อตฺถํ วิภชิตุํ. วิภชตายสฺมา มหากจฺจาโน อครุํ กตฺวา\csromansharedfootnote{c8a32d3d24ad7f14}{อครุกตฺวา (สี), อครุกริตฺวา (สฺยา, อิ)} ติ. เตน หาวุโส สุณาถ สาธุกํ มนสิกโรถ ภาสิสฺสามีติ. “เอวมาวุโส”ติ โข เต ภิกฺขู อายสฺมโต มหากจฺจานสฺส ปจฺจสฺโสสุํ. อายสฺมา มหากจฺจาโน เอตทโวจ\csromandash{}}

\proseitem{204}{ยํ โข โน อาวุโส ภควา สํขิตฺเตน อุทฺเทสํ อุทฺทิสิตฺวา วิตฺถาเรน อตฺถํ อวิภชิตฺวา อุฏฺฐายาสนา วิหารํ ปวิฏฺโฐ, “ยโตนิทานํ ภิกฺขุ ปุริสํ ปปญฺจ\-สญฺญา\-สงฺขา สมุทาจรนฺติ. เอตฺถ เจ นตฺถิ อภินนฺทิตพฺพํ อภิวทิตพฺพํ อชฺโฌสิตพฺพํ, เอเสวนฺโต ราคานุสยานํ ฯเปฯ. เอตฺเถเต ปาปกา อกุสลา ธมฺมา อปริเสสา นิรุชฺฌนฺตี”ติ. อิมสฺส โข อหํ อาวุโส ภควตา สํขิตฺเตน อุทฺเทสสฺส อุทฺทิฏฺฐสฺส วิตฺถาเรน อตฺถํ อวิภตฺตสฺส เอวํ วิตฺถาเรน อตฺถํ อาชานามิ.}

\prose{จกฺขุญฺจาวุโส ปฏิจฺจ รูเป จ อุปฺปชฺชติ จกฺขุวิญฺญาณํ, ติณฺณํ สงฺคติ ผสฺโส, ผสฺสปจฺจยา เวทนา, ยํ เวเทติ ตํ สญฺชานาติ, ยํ สญฺชานาติ ตํ วิตกฺเกติ, ยํ วิตกฺเกติ ตํ ปปญฺเจติ, ยํ ปปญฺเจติ ตโตนิทานํ ปุริสํ ปปญฺจ\-สญฺญา\-สงฺขา สมุทาจรนฺติ อตีตา\-นาคต\-ปจฺจุปฺปนฺเนสุ จกฺขุ\-วิญฺเญยฺเยสุ รูเปสุ. โสตญฺจาวุโส ปฏิจฺจ สทฺเท จ อุปฺปชฺชติ โสตวิญฺญาณํ ฯเปฯ. ฆานญฺจาวุโส ปฏิจฺจ คนฺเธ จ อุปฺปชฺชติ ฆานวิญฺญาณํ ฯเปฯ. ชิวฺหญฺจาวุโส ปฏิจฺจ รเส จ อุปฺปชฺชติ ชิวฺหาวิญฺญาณํ ฯเปฯ. กายญฺจาวุโส ปฏิจฺจ โผฏฺฐพฺเพ จ อุปฺปชฺชติ กายวิญฺญาณํ ฯเปฯ. มนญฺจาวุโส ปฏิจฺจ ธมฺเม จ อุปฺปชฺชติ มโนวิญฺญาณํ, ติณฺณํ สงฺคติ ผสฺโส, ผสฺสปจฺจยา เวทนา, ยํ เวเทติ ตํ สญฺชานาติ, ยํ สญฺชานาติ ตํ วิตกฺเกติ, ยํ วิตกฺเกติ ตํ ปปญฺเจติ, ยํ ปปญฺเจติ ตโตนิทานํ ปุริสํ ปปญฺจ\-สญฺญา\-สงฺขา สมุทาจรนฺติ อตีตา\-นาคต\-ปจฺจุปฺปนฺเนสุ มโนวิญฺเญยฺเยสุ ธมฺเมสุ.}

\prose{โส วตาวุโส จกฺขุสฺมิํ สติ รูเป สติ จกฺขุวิญฺญาเณ สติ ผสฺส\-ปญฺญตฺติํ ปญฺญาเปสฺสตีติ ฐานเมตํ วิชฺชติ, ผสฺส\-ปญฺญตฺติยา สติ เวทนา\-ปญฺญตฺติํ ปญฺญาเปสฺสตีติ ฐานเมตํ วิชฺชติ, เวทนา\-ปญฺญตฺติยา สติ สญฺญาปญฺญตฺติํ ปญฺญาเปสฺสตีติ ฐานเมตํ วิชฺชติ, สญฺญา\-ปญฺญตฺติยา \csromanfolio{159}สติ วิตกฺก\-ปญฺญตฺติํ ปญฺญาเปสฺสตีติ ฐานเมตํ วิชฺชติ, วิตกฺก\-ปญฺญตฺติยา สติ ปปญฺจ\-สญฺญา\-สงฺขา\-สมุทาจรณ\-ปญฺญตฺติํ ปญฺญาเปสฺสตีติ ฐานเมตํ วิชฺชติ. โส วตาวุโส โสตสฺมิํ สติ สทฺเท สติ ฯเปฯ ฆานสฺมิํ สติ คนฺเธ สติ ฯเปฯ ชิวฺหาย สติ รเส สติ ฯเปฯ กายสฺมิํ สติ โผฏฺฐพฺเพ สติ ฯเปฯ มนสฺมิํ สติ ธมฺเม สติ มโนวิญฺญาเณ สติ ผสฺส\-ปญฺญตฺติํ ปญฺญาเปสฺสตีติ ฐานเมตํ วิชฺชติ. ผสฺส\-ปญฺญตฺติยา สติ เวทนา\-ปญฺญตฺติํ ปญฺญาเปสฺสตีติ ฐานเมตํ วิชฺชติ. เวทนา\-ปญฺญตฺติยา สติ สญฺญาปญฺญตฺติํ ปญฺญาเปสฺสตีติ ฐานเมตํ วิชฺชติ. สญฺญา\-ปญฺญตฺติยา สติ วิตกฺก\-ปญฺญตฺติํ ปญฺญาเปสฺสตีติ ฐานเมตํ วิชฺชติ. วิตกฺก\-ปญฺญตฺติยา สติ ปปญฺจ\-สญฺญา\-สงฺขา\-สมุทาจรณ\-ปญฺญตฺติํ ปญฺญาเปสฺสตีติ ฐานเมตํ วิชฺชติ.}

\prose{โส วตาวุโส จกฺขุสฺมิํ อสติ รูเป อสติ จกฺขุวิญฺญาเณ อสติ ผสฺส\-ปญฺญตฺติํ ปญฺญาเปสฺสตีติ เนตํ ฐานํ วิชฺชติ. ผสฺส\-ปญฺญตฺติยา อสติ เวทนา\-ปญฺญตฺติํ ปญฺญาเปสฺสตีติ เนตํ ฐานํ วิชฺชติ. เวทนา\-ปญฺญตฺติยา อสติ สญฺญาปญฺญตฺติํ ปญฺญาเปสฺสตีติ เนตํ ฐานํ วิชฺชติ. สญฺญา\-ปญฺญตฺติยา อสติ วิตกฺก\-ปญฺญตฺติํ ปญฺญาเปสฺสตีติ เนตํ ฐานํ วิชฺชติ. วิตกฺก\-ปญฺญตฺติยา อสติ ปปญฺจ\-สญฺญา\-สงฺขา\-สมุทาจรณ\-ปญฺญตฺติํ ปญฺญาเปสฺสตีติ เนตํ ฐานํ วิชฺชติ. โส วตาวุโส โสตสฺมิํ อสติ สทฺเท อสติ ฯเปฯ ฆานสฺมิํ อสติ คนฺเธ อสติ ฯเปฯ ชิวฺหาย อสติ รเส อสติ ฯเปฯ กายสฺมิํ อสติ โผฏฺฐพฺเพ อสติ ฯเปฯ มนสฺมิํ อสติ ธมฺเม อสติ มโนวิญฺญาเณ อสติ ผสฺส\-ปญฺญตฺติํ ปญฺญาเปสฺสตีติ เนตํ ฐานํ วิชฺชติ. ผสฺส\-ปญฺญตฺติยา อสติ เวทนา\-ปญฺญตฺติํ ปญฺญาเปสฺสตีติ เนตํ ฐานํ วิชฺชติ. เวทนา\-ปญฺญตฺติยา อสติ สญฺญาปญฺญตฺติํ ปญฺญาเปสฺสตีติ เนตํ ฐานํ วิชฺชติ. สญฺญา\-ปญฺญตฺติยา อสติ วิตกฺก\-ปญฺญตฺติํ ปญฺญาเปสฺสตีติ เนตํ ฐานํ วิชฺชติ. วิตกฺก\-ปญฺญตฺติยา อสติ ปปญฺจ\-สญฺญา\-สงฺขา\-สมุทาจรณ\-ปญฺญตฺติํ ปญฺญาเปสฺสตีติ เนตํ ฐานํ วิชฺชติ.}

\prose{ยํ โข โน อาวุโส ภควา สํขิตฺเตน อุทฺเทสํ อุทฺทิสิตฺวา วิตฺถาเรน อตฺถํ อวิภชิตฺวา อุฏฺฐายาสนา วิหารํ ปวิฏฺโฐ, “ยโตนิทานํ ภิกฺขุ ปุริสํ ปปญฺจ\-สญฺญา\-สงฺขา สมุทาจรนฺติ. เอตฺถ เจ นตฺถิ อภินนฺทิตพฺพํ อภิวทิตพฺพํ อชฺโฌสิตพฺพํ, เอเสวนฺโต ราคานุสยานํ ฯเปฯ. เอตฺเถเต ปาปกา อกุสลา ธมฺมา อปริเสสา นิรุชฺฌนฺตี”ติ. อิมสฺส โข อหํ อาวุโส \csromanfolio{160}ภควตา สํขิตฺเตน อุทฺเทสสฺส อุทฺทิฏฺฐสฺส วิตฺถาเรน อตฺถํ อวิภตฺตสฺส เอวํ วิตฺถาเรน อตฺถํ อาชานามิ. อากงฺขมานา จ ปน ตุเมฺห อายสฺมนฺโต ภควนฺตํเยว อุปสงฺกมิตฺวา เอตมตฺถํ ปฏิปุจฺเฉยฺยาถ. ยถา โว ภควา พฺยากโรติ, ตถา นํ ธาเรยฺยาถาติ.}

\proseitem{205}{อถ โข เต ภิกฺขู อายสฺมโต มหากจฺจานสฺส ภาสิตํ อภินนฺทิตฺวา อนุโมทิตฺวา อุฏฺฐายาสนา เยน ภควา เตนุ\-ปสงฺกมิํสุ, อุปสงฺกมิตฺวา ภควนฺตํ อภิวาเทตฺวา เอกมนฺตํ นิสีทิํสุ, เอกมนฺตํ นิสินฺนา โข เต ภิกฺขู ภควนฺตํ เอตทโวจุํ\csromandash{}ยํ โข โน ภนฺเต ภควา สํขิตฺเตน อุทฺเทสํ อุทฺทิสิตฺวา วิตฺถาเรน อตฺถํ อวิภชิตฺวา อุฏฺฐายาสนา วิหารํ ปวิฏฺโฐ, “ยโตนิทานํ ภิกฺขุ ปุริสํ ปปญฺจ\-สญฺญา\-สงฺขา สมุทาจรนฺติ. เอตฺถ เจ นตฺถิ อภินนฺทิตพฺพํ อภิวทิตพฺพํ อชฺโฌสิตพฺพํ, เอเสวนฺโต ราคานุสยานํ ฯเปฯ. เอตฺเถเต ปาปกา อกุสลา ธมฺมา อปริเสสา นิรุชฺฌนฺตี”ติ. เตสํ โน ภนฺเต อมฺหากํ อจิร\-ปกฺกนฺตสฺส ภควโต เอตทโหสิ “อิทํ โข โน อาวุโส ภควา สํขิตฺเตน อุทฺเทสํ อุทฺทิสิตฺวา วิตฺถาเรน อตฺถํ อวิภชิตฺวา อุฏฺฐายาสนา วิหารํ ปวิฏฺโฐ, ‘ยโตนิทานํ ภิกฺขุ ปุริสํ ปปญฺจ\-สญฺญา\-สงฺขา สมุทาจรนฺติ. เอตฺถ เจ นตฺถิ อภินนฺทิตพฺพํ อภิวทิตพฺพํ อชฺโฌสิตพฺพํ, เอเสวนฺโต ราคานุสยานํ, เอเสวนฺโต ปฏิฆา\-นุสยานํ, เอเสวนฺโต ทิฏฺฐา\-นุสยานํ, เอเสวนฺโต วิจิกิจฺฉา\-นุสยานํ, เอเสวนฺโต มานานุสยานํ, เอเสวนฺโต ภวราคา\-นุสยานํ, เอเสวนฺโต อวิชฺชา\-นุสยานํ, เอเสวนฺโต ทณฺฑาทาน สตฺถาทาน กลห วิคฺคห วิวาท ตุวํตุวํ เปสุญฺญ มุสาวาทานํ. เอตฺเถเต ปาปกา อกุสลา ธมฺมา อปริเสสา นิรุชฺฌนฺตี’ติ. โก นุ โข อิมสฺส ภควตา สํขิตฺเตน อุทฺเทสสฺส อุทฺทิฏฺฐสฺส วิตฺถาเรน อตฺถํ อวิภตฺตสฺส วิตฺถาเรน อตฺถํ วิภเชยฺยา”ติ. เตสํ โน ภนฺเต อมฺหากํ เอตทโหสิ “อยํ โข อายสฺมา มหากจฺจาโน สตฺถุ เจว สํวณฺณิโต, สมฺภาวิโต จ วิญฺญูนํ สพฺรหฺมจารีนํ, ปโหติ จายสฺมา มหากจฺจาโน อิมสฺส ภควตา สํขิตฺเตน อุทฺเทสสฺส อุทฺทิฏฺฐสฺส วิตฺถาเรน อตฺถํ อวิภตฺตสฺส วิตฺถาเรน อตฺถํ วิภชิตุํ, ยํนูน มยํ เยนายสฺมา มหากจฺจาโน เตนุปสงฺกเมยฺยาม, อุปสงฺกมิตฺวา อายสฺมนฺตํ มหากจฺจานํ เอตมตฺถํ ปฏิปุจฺเฉยฺยามา”ติ. อถ โข มยํ ภนฺเต เยนายสฺมา มหากจฺจาโน เตนุ\-ปสงฺกมิมฺห, อุปสงฺกมิตฺวา \csromanfolio{161}อายสฺมนฺตํ มหากจฺจานํ เอตมตฺถํ ปฏิปุจฺฉิมฺห, เตสํ โน ภนฺเต อายสฺมตา มหากจฺจาเนน อิเมหิ อากาเรหิ อิเมหิ ปเทหิ อิเมหิ พฺยญฺชเนหิ อตฺโถ วิภตฺโตติ. ปณฺฑิโต ภิกฺขเว มหากจฺจาโน, มหาปญฺโญ ภิกฺขเว มหากจฺจาโน. มํ เจปิ ตุเมฺห ภิกฺขเว เอตมตฺถํ ปฏิปุจฺเฉยฺยาถ, อหมฺปิ ตํ เอวเมวํ พฺยากเรยฺยํ, ยถา ตํ มหากจฺจาเนน พฺยากตํ, เอโส เจเวตสฺส อตฺโถ, เอวญฺจ\csromansharedfootnote{f3acba62154622b7}{เอวเมว จ (ก)} นํ ธาเรถาติ.}

\prose{เอวํ วุตฺเต อายสฺมา อานนฺโท ภควนฺตํ เอตทโวจ “เสยฺยถาปิ ภนฺเต ปุริโส ชิฆจฺฉา\-ทุพฺพลฺย\-ปเรโต มธุปิณฺฑิกํ อธิคจฺเฉยฺย, โส ยโต ยโต สาเยยฺย, ลเภเถว สาทุรสํ อเสจนกํ. เอวเมว โข ภนฺเต เจตโส ภิกฺขุ ทพฺพชาติโก ยโต ยโต อิมสฺส ธมฺม\-ปริยายสฺส ปญฺญาย อตฺถํ อุปปริกฺเขยฺย, ลเภเถว อตฺตมนตํ, ลเภเถว เจตโส ปสาทํ. โกนาโม อยํ\csromansharedfootnote{eb199ab8280164f1}{โก นามายํ (สฺยา)} ภนฺเต ธมฺมปริยาโย”ติ. ตสฺมาติห ตฺวํ อานนฺท อิมํ ธมฺม\-ปริยายํ “มธุปิณฺฑิก\-ปริยาโย”เตฺวว นํ ธาเรหีติ.}

\prosewithcloserruled{อิทมโวจ ภควา. อตฺตมโน อายสฺมา อานนฺโท ภควโต ภาสิตํ อภินนฺทีติ.}{มธุปิณฺฑิก\-สุตฺตํ นิฏฺฐิตํ อฏฺฐมํ.}
\csromanmatikaanchor{mtk.47}
\csromantocmarkref{subsection}{9. เทฺวธาวิตกฺกสุตฺต}{mtk.47}
\bookmark[dest={mtk.47},level=2]{9. เทฺวธาวิตกฺกสุตฺต}
\csromanheader{2}{9. เทฺวธา\-วิตกฺก\-สุตฺต}
\proseitem{206}{เอวํ เม สุตํ\csromandash{}เอกํ สมยํ ภควา สาวตฺถิยํ วิหรติ เชตวเน อนาถ\-ปิณฺฑิกสฺส อาราเม. ตตฺร โข ภควา ภิกฺขู อามนฺเตสิ “ภิกฺขโว”ติ. “ภทนฺเต”ติ เต ภิกฺขู ภควโต ปจฺจสฺโสสุํ. ภควา เอตทโวจ\csromandash{}}

\prose{ปุพฺเพว เม ภิกฺขเว สมฺโพธา อนภิสมฺพุทฺธสฺส โพธิสตฺตสฺเสว สโต เอตทโหสิ “ยํนูนาหํ ทฺวิธา กตฺวา ทฺวิธา กตฺวา วิตกฺเก วิหเรยฺยนฺ”ติ. โส โข อหํ ภิกฺขเว โย จายํ กามวิตกฺโก โย จ พฺยาปาทวิตกฺโก โย จ วิหิํสาวิตกฺโก, อิมํ เอกํ \csromanfolio{162}ภาคมกาสิํ. โย จายํ เนกฺขมฺม\-วิตกฺโก โย จ อพฺยาปาทวิตกฺโก โย จ อวิหิํสา\-วิตกฺโก, อิมํ ทุติยํ ภาคมกาสิํ.}

\proseitem{207}{ตสฺส มยฺหํ ภิกฺขเว เอวํ อปฺปมตฺตสฺส อาตาปิโน ปหิตตฺตสฺส วิหรโต อุปฺปชฺชติ กามวิตกฺโก. โส เอวํ ปชานามิ “อุปฺปนฺโน โข เม อยํ กามวิตกฺโก, โส จ โข อตฺตพฺยาพาธายปิ สํวตฺตติ, ปรพฺยาพาธายปิ สํวตฺตติ, อุภย\-พฺยาพาธายปิ สํวตฺตติ, ปญฺญานิโรธิโก วิฆาต\-ปกฺขิโก อนิพฺพานสํวตฺตนิโก”\csromansharedfootnote{ad8ddbf3d76339fb}{อนิพฺพานสํวตฺตนิโก”ติ (?) 2-3. อตีตกาลิกกิริยาปทานิเยว.}. ‘อตฺตพฺยาพาธาย สํวตฺตตี’ติปิ เม ภิกฺขเว ปฏิสญฺจิกฺขโต อพฺภตฺถํ คจฺฉติ, ‘ปรพฺยาพาธาย สํวตฺตตี’ติปิ เม ภิกฺขเว ปฏิสญฺจิกฺขโต อพฺภตฺถํ คจฺฉติ, ‘อุภย\-พฺยาพาธาย สํวตฺตตี’ติปิ เม ภิกฺขเว ปฏิสญฺจิกฺขโต อพฺภตฺถํ คจฺฉติ, ‘ปญฺญานิโรธิโก วิฆาต\-ปกฺขิโก อนิพฺพานสํวตฺตนิโก’ติปิ เม ภิกฺขเว ปฏิสญฺจิกฺขโต อพฺภตฺถํ คจฺฉติ. โส โข อหํ ภิกฺขเว อุปฺปนฺนุปฺปนฺนํ กามวิตกฺกํ ปชหเมว2 วิโนทเมว3, พฺยนฺตเมว\csromansharedfootnote{88f580b942aa60eb}{พฺยนฺเตว (สี, สฺยา, อิ)} นํ อกาสิํ.}

\proseitem{208}{ตสฺส มยฺหํ ภิกฺขเว เอวํ อปฺปมตฺตสฺส อาตาปิโน ปหิตตฺตสฺส วิหรโต อุปฺปชฺชติ พฺยาปาทวิตกฺโก ฯเปฯ อุปฺปชฺชติ วิหิํสาวิตกฺโก. โส เอวํ ปชานามิ “อุปฺปนฺโน โข เม อยํ วิหิํสาวิตกฺโก, โส จ โข อตฺตพฺยาพาธายปิ สํวตฺตติ, ปรพฺยาพาธายปิ สํวตฺตติ, อุภย\-พฺยาพาธายปิ สํวตฺตติ, ปญฺญานิโรธิโก วิฆาต\-ปกฺขิโก อนิพฺพานสํวตฺตนิโก”. ‘อตฺตพฺยาพาธาย สํวตฺตตี’ติปิ เม ภิกฺขเว ปฏิสญฺจิกฺขโต อพฺภตฺถํ คจฺฉติ, ‘ปรพฺยาพาธาย สํวตฺตตี’ติปิ เม ภิกฺขเว ปฏิสญฺจิกฺขโต อพฺภตฺถํ คจฺฉติ, ‘อุภย\-พฺยาพาธาย สํวตฺตตี’ติปิ เม ภิกฺขเว ปฏิสญฺจิกฺขโต อพฺภตฺถํ คจฺฉติ, ‘ปญฺญานิโรธิโก วิฆาต\-ปกฺขิโก อนิพฺพานสํวตฺตนิโก’ติปิ เม ภิกฺขเว ปฏิสญฺจิกฺขโต อพฺภตฺถํ คจฺฉติ. โส โข อหํ ภิกฺขเว อุปฺปนฺนุปฺปนฺนํ วิหิํสา\-วิตกฺกํ ปชหเมว วิโนทเมว, พฺยนฺตเมว นํ อกาสิํ.}

\prose{ยญฺญเทว ภิกฺขเว ภิกฺขุ พหุลม\-นุวิตกฺเกติ อนุวิจาเรติ, ตถา ตถา นติ โหติ เจตโส. กามวิตกฺกํ เจ ภิกฺขเว ภิกฺขุ \csromanfolio{163}พหุลม\-นุวิตกฺเกติ อนุวิจาเรติ. ปหาสิ เนกฺขมฺม\-วิตกฺกํ, กามวิตกฺกํ พหุลมกาสิ. ตสฺส ตํ กามวิตกฺกาย จิตฺตํ นมติ. พฺยาปาท\-วิตกฺกํ เจ ภิกฺขเว ฯเปฯ. วิหิํสา\-วิตกฺกํ เจ ภิกฺขเว ภิกฺขุ พหุลม\-นุวิตกฺเกติ อนุวิจาเรติ. ปหาสิ อวิหิํสา\-วิตกฺกํ, วิหิํสา\-วิตกฺกํ พหุลมกาสิ. ตสฺส ตํ วิหิํสา\-วิตกฺกาย จิตฺตํ นมติ. เสยฺยถาปิ ภิกฺขเว วสฺสานํ ปจฺฉิเม มาเส สรทสมเย กิฏฺฐสมฺพาเธ โคปาลโก คาโว รกฺเขยฺย, โส ตา คาโว ตโต ตโต ทณฺเฑน อาโกเฏยฺย ปฏิโกเฏยฺย สนฺนิรุนฺเธยฺย สนฺนิวาเรยฺย. ตํ กิสฺส เหตุ, ปสฺสติ หิ โส ภิกฺขเว โคปาลโก ตโตนิทานํ วธํ วา พนฺธนํ วา ชานิํ วา ครหํ วา. เอวเมว โข อหํ ภิกฺขเว อทฺทสํ อกุสลานํ ธมฺมานํ อาทีนวํ โอการํ สํกิเลสํ กุสลานํ ธมฺมานํ เนกฺขมฺเม อานิสํสํ โวทานปกฺขํ.}

\proseitem{209}{ตสฺส มยฺหํ ภิกฺขเว เอวํ อปฺปมตฺตสฺส อาตาปิโน ปหิตตฺตสฺส วิหรโต อุปฺปชฺชติ เนกฺขมฺม\-วิตกฺโก. โส เอวํ ปชานามิ “อุปฺปนฺโน โข เม อยํ เนกฺขมฺม\-วิตกฺโก, โส จ โข เนว\-ตฺตพฺยาพาธาย สํวตฺตติ, น ปรพฺยาพาธาย สํวตฺตติ, น อุภย\-พฺยาพาธาย สํวตฺตติ, ปญฺญาวุทฺธิโก อวิฆาต\-ปกฺขิโก นิพฺพาน\-สํวตฺตนิโก. รตฺติํ เจปิ นํ ภิกฺขเว อนุวิตกฺเกยฺยํ อนุวิจาเรยฺยํ, เนว ตโตนิทานํ ภยํ สมนุปสฺสามิ. ทิวสํ เจปิ นํ ภิกฺขเว อนุวิตกฺเกยฺยํ อนุวิจาเรยฺยํ, เนว ตโตนิทานํ ภยํ สมนุปสฺสามิ. รตฺตินฺทิวํ เจปิ นํ ภิกฺขเว อนุวิตกฺเกยฺยํ อนุวิจาเรยฺยํ, เนว ตโตนิทานํ ภยํ สมนุปสฺสามิ. อปิ จ โข เม อติจิรํ อนุวิตกฺกยโต อนุวิจารยโต กาโย กิลเมยฺย, กาเย กิลนฺเต\csromansharedfootnote{9ea31b0546c54af7}{กิลมนฺเต (ก)} จิตฺตํ อูหญฺเญยฺย, อูหเต จิตฺเต อารา จิตฺตํ สมาธิมฺหา”ติ. โส โข อหํ ภิกฺขเว อชฺฌตฺตเมว จิตฺตํ สณฺฐเปมิ สนฺนิสาเทมิ เอโกทิํ กโรมิ\csromansharedfootnote{ae4109bf89c38c6d}{เอโกทิ กโรมิ (อิ)} สมาทหามิ. ตํ กิสฺส เหตุ, มา เม จิตฺตํ อูหญฺญี\csromansharedfootnote{1230cc17c5de727d}{อุคฺฆาฏีติ (สฺยา, ก), อูหนีติ (อิ)} ติ.}

\proseitem{210}{ตสฺส มยฺหํ ภิกฺขเว เอวํ อปฺปมตฺตสฺส อาตาปิโน ปหิตตฺตสฺส วิหรโต อุปฺปชฺชติ อพฺยาปาทวิตกฺโก ฯเปฯ อุปฺปชฺชติ อวิหิํสา\-วิตกฺโก. โส เอวํ ปชานามิ “อุปฺปนฺโน โข เม อยํ อวิหิํสา\-วิตกฺโก, โส จ โข เนว\-ตฺตพฺยาพาธาย สํวตฺตติ, น ปรพฺยาพาธาย สํวตฺตติ, น \csromanfolio{164}อุภย\-พฺยาพาธาย สํวตฺตติ, ปญฺญาวุทฺธิโก อวิฆาต\-ปกฺขิโก นิพฺพาน\-สํวตฺตนิโก. รตฺติํ เจปิ นํ ภิกฺขเว อนุวิตกฺเกยฺยํ อนุวิจาเรยฺยํ, เนว ตโตนิทานํ ภยํ สมนุปสฺสามิ. ทิวสํ เจปิ นํ ภิกฺขเว อนุวิตกฺเกยฺยํ อนุวิจาเรยฺยํ, เนว ตโตนิทานํ ภยํ สมนุปสฺสามิ. รตฺตินฺทิวํ เจปิ นํ ภิกฺขเว อนุวิตกฺเกยฺยํ อนุวิจาเรยฺยํ, เนว ตโตนิทานํ ภยํ สมนุปสฺสามิ. อปิ จ โข เม อติจิรํ อนุวิตกฺกยโต อนุวิจารยโต กาโย กิลเมยฺย, กาเย กิลนฺเต จิตฺตํ อูหญฺเญยฺย, อูหเต จิตฺเต อารา จิตฺตํ สมาธิมฺหา”ติ. โส โข อหํ ภิกฺขเว อชฺฌตฺตเมว จิตฺตํ สณฺฐเปมิ สนฺนิสาเทมิ เอโกทิํ กโรมิ สมาทหามิ. ตํ กิสฺส เหตุ, มา เม จิตฺตํ อูหญฺญีติ.}

\prose{ยญฺญเทว ภิกฺขเว ภิกฺขุ พหุลม\-นุวิตกฺเกติ อนุวิจาเรติ, ตถา ตถา นติ โหติ เจตโส. เนกฺขมฺมวิตกฺก\-ญฺเจ ภิกฺขเว ภิกฺขุ พหุลม\-นุวิตกฺเกติ อนุวิจาเรติ. ปหาสิ กามวิตกฺกํ, เนกฺขมฺม\-วิตกฺกํ พหุลมกาสิ. ตสฺส ตํ เนกฺขมฺม\-วิตกฺกาย จิตฺตํ นมติ. อพฺยาปาทวิตกฺก\-ญฺเจ ภิกฺขเว ฯเปฯ. อวิหิํสาวิตกฺก\-ญฺเจ ภิกฺขเว ภิกฺขุ พหุลม\-นุวิตกฺเกติ อนุวิจาเรติ. ปหาสิ วิหิํสา\-วิตกฺกํ, อวิหิํสา\-วิตกฺกํ พหุลมกาสิ. ตสฺส ตํ อวิหิํสา\-วิตกฺกาย จิตฺตํ นมติ. เสยฺยถาปิ ภิกฺขเว คิมฺหานํ ปจฺฉิเม มาเส สพฺพสสฺเสสุ คามนฺต\-สมฺภเตสุ โคปาลโก คาโว รกฺเขยฺย, ตสฺส รุกฺขมูล\-คตสฺส วา อพฺโภกาส\-คตสฺส วา สติกรณียเมว โหติ เอตา\csromansharedfootnote{4b6b1ea4ef8ab7c5}{เอเต (ก)} คาโวติ. เอวเมวํ โข ภิกฺขเว สติกรณียเมว อโหสิ เอเต ธมฺมาติ.}

\proseitem{211}{อารทฺธํ โข ปน เม ภิกฺขเว วีริยํ อโหสิ อสลฺลีนํ, อุปฏฺฐิตา สติ อสมฺมุฏฺฐา, ปสฺสทฺโธ กาโย อสารทฺโธ, สมาหิตํ จิตฺตํ เอกคฺคํ. โส โข อหํ ภิกฺขเว วิวิจฺเจว กาเมหิ วิวิจฺจ อกุสเลหิ ธมฺเมหิ สวิตกฺกํ สวิจารํ วิเวกชํ ปีติสุขํ ปฐมํ ฌานํ อุปสมฺปชฺช วิหาสิํ. วิตกฺก\-วิจารานํ วูปสมา อชฺฌตฺตํ สมฺปสาทนํ เจตโส เอโกทิภาวํ อวิตกฺกํ อวิจารํ สมาธิชํ ปีติสุขํ ทุติยํ ฌานํ อุปสมฺปชฺช วิหาสิํ. ปีติยา จ วิราคา อุเปกฺขโก จ วิหาสิํ สโต จ สมฺปชาโน, สุขญฺจ กาเยน ปฏิสํเวเทสิํ, ยํ ตํ อริยา อาจิกฺขนฺติ “อุเปกฺขโก \csromanfolio{165}สติมา สุขวิหารี”ติ, ตติยํ ฌานํ อุปสมฺปชฺช วิหาสิํ. สุขสฺส จ ปหานา ทุกฺขสฺส จ ปหานา ปุพฺเพว โสมนสฺส\-โทมนสฺสานํ อตฺถงฺคมา อทุกฺขมสุขํ อุเปกฺขา\-สติ\-ปาริสุทฺธิํ จตุตฺถํ ฌานํ อุปสมฺปชฺช วิหาสิํ.}

\proseitem{212}{โส เอวํ สมาหิเต จิตฺเต ปริสุทฺเธ ปริโยทาเต อนงฺคเณ วิคตู\-ปกฺกิเลเส มุทุภูเต กมฺมนิเย ฐิเต อาเนญฺชปฺปตฺเต ปุพฺเพ\-นิวาสา\-นุสฺสติ\-ญาณาย จิตฺตํ อภินินฺนาเมสิํ. โส อเนกวิหิตํ ปุพฺเพนิวาสํ อนุสฺสรามิ. เสยฺยถิทํ, เอกมฺปิ ชาติํ ฯเปฯ อิติ สาการํ ส- อุทฺเทสํ อเนกวิหิตํ ปุพฺเพนิวาสํ อนุสฺสรามิ. อยํ โข เม ภิกฺขเว รตฺติยา ปฐเม ยาเม ปฐมา วิชฺชา อธิคตา, อวิชฺชา วิหตา, วิชฺชา อุปฺปนฺนา, ตโม วิหโต, อาโลโก อุปฺปนฺโน, ยถา ตํ อปฺปมตฺตสฺส อาตาปิโน ปหิตตฺตสฺส วิหรโต.}

\proseitem{213}{โส เอวํ สมาหิเต จิตฺเต ปริสุทฺเธ ปริโยทาเต อนงฺคเณ วิคตู\-ปกฺกิเลเส มุทุภูเต กมฺมนิเย ฐิเต อาเนญฺชปฺปตฺเต สตฺตานํ จุตูปปาต\-ญาณาย จิตฺตํ อภินินฺนาเมสิํ. โส ทิพฺเพน จกฺขุนา วิสุทฺเธน อติกฺกนฺต\-มานุสเกน สตฺเต ปสฺสามิ จวมาเน อุปปชฺชมาเน ฯเปฯ อิเม วต โภนฺโต สตฺตา กาย\-ทุจฺจริเตน สมนฺนาคตา ฯเปฯ อิติ ทิพฺเพน จกฺขุนา วิสุทฺเธน อติกฺกนฺต\-มานุสเกน สตฺเต ปสฺสามิ จวมาเน อุปปชฺชมาเน หีเน ปณีเต สุวณฺเณ ทุพฺพณฺเณ สุคเต ทุคฺคเต, ยถากมฺมูปเค สตฺเต ปชานามิ. อยํ โข เม ภิกฺขเว รตฺติยา มชฺฌิเม ยาเม ทุติยา วิชฺชา อธิคตา, อวิชฺชา วิหตา, วิชฺชา อุปฺปนฺนา, ตโม วิหโต, อาโลโก อุปฺปนฺโน, ยถา ตํ อปฺปมตฺตสฺส อาตาปิโน ปหิตตฺตสฺส วิหรโต.}

\proseitem{214}{โส เอวํ สมาหิเต จิตฺเต ปริสุทฺเธ ปริโยทาเต อนงฺคเณ วิคตู\-ปกฺกิเลเส มุทุภูเต กมฺมนิเย ฐิเต อาเนญฺชปฺปตฺเต อาสวานํ ขยญาณาย จิตฺตํ อภินินฺนาเมสิํ. โส “อิทํ ทุกฺขนฺ”ติ ยถาภูตํ อพฺภญฺญาสิํ, “อยํ ทุกฺขสมุทโย”ติ ยถาภูตํ อพฺภญฺญาสิํ, “อยํ ทุกฺขนิโรโธ”ติ ยถาภูตํ อพฺภญฺญาสิํ, “อยํ ทุกฺขนิโรธ\-คามินี ปฏิปทา”ติ ยถาภูตํ อพฺภญฺญาสิํ. “อิเม อาสวา”ติ ยถาภูตํ อพฺภญฺญาสิํ, “อยํ อาสวสมุทโย”ติ ยถาภูตํ อพฺภญฺญาสิํ, \csromanfolio{166}“อยํ อาสวนิโรโธ”ติ ยถาภูตํ อพฺภญฺญาสิํ, “อยํ อาสว\-นิโรธ\-คามินี ปฏิปทา”ติ ยถาภูตํ อพฺภญฺญาสิํ. ตสฺส เม เอวํ ชานโต เอวํ ปสฺสโต กามาสวาปิ จิตฺตํ วิมุจฺจิตฺถ, ภวาสวาปิ จิตฺตํ วิมุจฺจิตฺถ, อวิชฺชาสวาปิ จิตฺตํ วิมุจฺจิตฺถ, วิมุตฺตสฺมิํ “วิมุตฺตมฺ”อิติ ญาณํ อโหสิ, “ขีณา ชาติ, วุสิตํ พฺรหฺมจริยํ, กตํ กรณียํ, นาปรํ อิตฺถตฺตายา”ติ อพฺภญฺญาสิํ. อยํ โข เม ภิกฺขเว รตฺติยา ปจฺฉิเม ยาเม ตติยา วิชฺชา อธิคตา, อวิชฺชา วิหตา, วิชฺชา อุปฺปนฺนา, ตโม วิหโต, อาโลโก อุปฺปนฺโน, ยถา ตํ อปฺปมตฺตสฺส อาตาปิโน ปหิตตฺตสฺส วิหรโต.}

\proseitem{215}{เสยฺยถาปิ ภิกฺขเว อรญฺเญ ปวเน มหนฺตํ นินฺนํ ปลฺลลํ, ตเมนํ มหามิคสงฺโฆ อุปนิสฺสาย วิหเรยฺย. ตสฺส โกจิเทว ปุริโส อุปฺปชฺเชยฺย อนตฺถกาโม อหิตกาโม อโยคกฺเขม\-กาโม. โส ยฺวาสฺส มคฺโค เขโม โสวตฺถิโก ปีติคมนีโย, ตํ มคฺคํ ปิทเหยฺย, วิวเรยฺย กุมฺมคฺคํ, โอทเหยฺย โอกจรํ, ฐเปยฺย โอกจาริกํ. เอวํ หิ โส ภิกฺขเว มหามิคสงฺโฆ อปเรน สมเยน อนยพฺยสนํ\csromansharedfootnote{76e98fd26b598ec7}{อนยพฺยสนํ ตนุตฺตํ (สี, สฺยา, อิ)} อาปชฺเชยฺย. ตสฺเสว โข ปน ภิกฺขเว มหโต มิคสงฺฆสฺส โกจิเทว ปุริโส อุปฺปชฺเชยฺย อตฺถกาโม หิตกาโม โยคกฺเขมกาโม. โส ยฺวาสฺส มคฺโค เขโม โสวตฺถิโก ปีติคมนีโย, ตํ มคฺคํ วิวเรยฺย, ปิทเหยฺย กุมฺมคฺคํ, อูหเนยฺย โอกจรํ, นาเสยฺย โอกจาริกํ. เอวํ หิ โส ภิกฺขเว มหามิคสงฺโฆ อปเรน สมเยน วุทฺธิํ วิรูฬฺหิํ เวปุลฺลํ อาปชฺเชยฺย.}

\prose{อุปมา โข เม อยํ ภิกฺขเว กตา อตฺถสฺส วิญฺญาปนาย. อยํ เจ เวตฺถ อตฺโถ. มหนฺตํ นินฺนํ ปลฺลลนฺติ โข ภิกฺขเว กามานเมตํ อธิวจนํ. มหา\-มิคสงฺโฆติ โข ภิกฺขเว สตฺตานเมตํ อธิวจนํ. ปุริโส อนตฺถกาโม อหิตกาโม อโยคกฺเขม\-กาโมติ โข ภิกฺขเว มารสฺเสตํ ปาปิมโต อธิวจนํ. กุมฺมคฺโคติ โข ภิกฺขเว อฏฺฐงฺคิกสฺเสตํ มิจฺฉามคฺคสฺส อธิวจนํ. เสยฺยถิทํ, มิจฺฉาทิฏฺฐิยา มิจฺฉา\-สงฺกปฺปสฺส มิจฺฉาวาจาย มิจฺฉา\-กมฺมนฺตสฺส มิจฺฉาอาชีวสฺส มิจฺฉา\-วายามสฺส มิจฺฉาสติยา มิจฺฉา\-สมาธิสฺส. โอกจโรติ โข ภิกฺขเว นนฺทีราคสฺเสตํ อธิวจนํ.}

\prose{\csromanfolio{167}โอกจาริกาติ โข ภิกฺขเว อวิชฺชาเยตํ อธิวจนํ. ปุริโส อตฺถกาโม หิตกาโม โยคกฺเขม\-กาโมติ โข ภิกฺขเว ตถาคตสฺเสตํ อธิวจนํ อรหโต สมฺมา\-สมฺพุทฺธสฺส. เขโม มคฺโค โสวตฺถิโก ปีติคมนีโยติ โข ภิกฺขเว อริยสฺเสตํ อฏฺฐงฺคิกสฺส มคฺคสฺส อธิวจนํ. เสยฺยถิทํ, สมฺมาทิฏฺฐิยา สมฺมา\-สงฺกปฺปสฺส สมฺมาวาจาย สมฺมา\-กมฺมนฺตสฺส สมฺมาอาชีวสฺส สมฺมาวายามสฺส สมฺมาสติยา สมฺมา\-สมาธิสฺส.}

\prose{อิติ โข ภิกฺขเว วิวโฏ มยา เขโม มคฺโค โสวตฺถิโก ปีติคมนีโย, ปิหิโต กุมฺมคฺโค, อูหโต โอกจโร, นาสิตา โอกจาริกา. ยํ ภิกฺขเว สตฺถารา กรณียํ สาวกานํ หิเตสินา อนุกมฺปเกน อนุกมฺปํ อุปาทาย, กตํ โว ตํ มยา. เอตานิ ภิกฺขเว รูกฺขมูลานิ, เอตานิ สุญฺญาคารานิ. ฌายถ ภิกฺขเว มา ปมาทตฺถ, มา ปจฺฉา วิปฺปฏิสาริโน อหุวตฺถ. อยํ โว อมฺหากํ อนุสาสนีติ.}

\prosewithcloserruled{อิทมโวจ ภควา. อตฺตมนา เต ภิกฺขู ภควโต ภาสิตํ อภินนฺทุนฺติ.}{เทฺวธา\-วิตกฺก\-สุตฺตํ นิฏฺฐิตํ นวมํ.}
\csromanmatikaanchor{mtk.48}
\csromantocmarkref{subsection}{10. วิตกฺกสณฺฐานสุตฺต}{mtk.48}
\bookmark[dest={mtk.48},level=2]{10. วิตกฺกสณฺฐานสุตฺต}
\csromanheader{2}{10. วิตกฺก\-สณฺฐาน\-สุตฺต}
\proseitem{216}{เอวํ เม สุตํ\csromandash{}เอกํ สมยํ ภควา สาวตฺถิยํ วิหรติ เชตวเน อนาถ\-ปิณฺฑิกสฺส อาราเม. ตตฺร โข ภควา ภิกฺขู อามนฺเตสิ “ภิกฺขโว”ติ. “ภทนฺเต”ติ เต ภิกฺขู ภควโต ปจฺจสฺโสสุํ. ภควา เอตทโวจ\csromandash{}}

\prose{อธิจิตฺตม\-นุยุตฺเตน ภิกฺขเว ภิกฺขุนา ปญฺจ นิมิตฺตานิ กาเลน กาลํ มนสิ กาตพฺพานิ. กตมานิ ปญฺจ. อิธ ภิกฺขเว ภิกฺขุโน ยํ นิมิตฺตํ อาคมฺม ยํ นิมิตฺตํ มนสิกโรโต อุปฺปชฺชนฺติ ปาปกา อกุสลา วิตกฺกา ฉนฺทูปสํหิตาปิ โทสูปสํหิตาปิ โมหูปสํหิตาปิ. เตน ภิกฺขเว ภิกฺขุนา ตมฺหา นิมิตฺตา อญฺญํ นิมิตฺตํ มนสิ กาตพฺพํ กุสลู\-ปสํหิตํ. ตสฺส ตมฺหา นิมิตฺตา อญฺญํ นิมิตฺตํ มนสิกโรโต กุสลู\-ปสํหิตํ เย ปาปกา \csromanfolio{168}อกุสลา วิตกฺกา ฉนฺทูปสํหิตาปิ โทสูปสํหิตาปิ โมหูปสํหิตาปิ, เต ปหียนฺติ เต อพฺภตฺถํ คจฺฉนฺติ, เตสํ ปหานา อชฺฌตฺตเมว จิตฺตํ สนฺติฏฺฐติ สนฺนิสีทติ เอโกทิ โหติ\csromansharedfootnote{b2e219c0c9a2bb77}{เอโกทิโภติ (สฺยา, ก)} สมาธิยติ. เสยฺยถาปิ ภิกฺขเว ทกฺโข ปลคณฺโฑ วา ปลคณฺฑ\-นฺเตวาสี วา สุขุมาย อาณิยา โอฬาริกํ อาณิํ อภินิหเนยฺย อภินีหเรยฺย อภินิวตฺเตยฺย\csromansharedfootnote{776db8003c21b318}{อภินิวชฺเชยฺย (สี, อิ)}. เอวเมว โข ภิกฺขเว ภิกฺขุโน ยํ นิมิตฺตํ อาคมฺม ยํ นิมิตฺตํ มนสิกโรโต อุปฺปชฺชนฺติ ปาปกา อกุสลา วิตกฺกา ฉนฺทูปสํหิตาปิ โทสูปสํหิตาปิ โมหูปสํหิตาปิ. เตน ภิกฺขเว ภิกฺขุนา ตมฺหา นิมิตฺตา อญฺญํ นิมิตฺตํ มนสิ กาตพฺพํ กุสลู\-ปสํหิตํ. ตสฺส ตมฺหา นิมิตฺตา อญฺญํ นิมิตฺตํ มนสิกโรโต กุสลู\-ปสํหิตํ เย ปาปกา อกุสลา วิตกฺกา ฉนฺทูปสํหิตาปิ โทสูปสํหิตาปิ โมหูปสํหิตาปิ, เต ปหียนฺติ เต อพฺภตฺถํ คจฺฉนฺติ, เตสํ ปหานา อชฺฌตฺตเมว จิตฺตํ สนฺติฏฺฐติ สนฺนิสีทติ เอโกทิ โหติ สมาธิยติ. (1)}

\proseitem{217}{ตสฺส เจ ภิกฺขเว ภิกฺขุโน ตมฺหา นิมิตฺตา อญฺญํ นิมิตฺตํ มนสิกโรโต กุสลู\-ปสํหิตํ อุปฺปชฺชนฺเตว ปาปกา อกุสลา วิตกฺกา ฉนฺทูปสํหิตาปิ โทสูปสํหิตาปิ โมหูปสํหิตาปิ. เตน ภิกฺขเว ภิกฺขุนา เตสํ วิตกฺกานํ อาทีนโว อุปปริ\-กฺขิ\-ตพฺโพ “อิติปิเม วิตกฺกา อกุสลา, อิติปิเม วิตกฺกา สาวชฺชา, อิติปิเม วิตกฺกา ทุกฺขวิปากา”ติ. ตสฺส เตสํ วิตกฺกานํ อาทีนวํ อุปปริกฺขโต เย ปาปกา อกุสลา วิตกฺกา ฉนฺทูปสํหิตาปิ โทสูปสํหิตาปิ โมหูปสํหิตาปิ, เต ปหียนฺติ เต อพฺภตฺถํ คจฺฉนฺติ, เตสํ ปหานา อชฺฌตฺตเมว จิตฺตํ สนฺติฏฺฐติ สนฺนิสีทติ เอโกทิ โหติ สมาธิยติ. เสยฺยถาปิ ภิกฺขเว อิตฺถี วา ปุริโส วา ทหโร ยุวา มณฺฑนก\-ชาติโก อหิกุณเปน วา กุกฺกุร\-กุณเปน วา มนุสฺส\-กุณเปน วา กณฺเฐ อาสตฺเตน อฏฺฏิเยยฺย หราเยยฺย ชิคุจฺเฉยฺย. เอวเมว โข ภิกฺขเว ตสฺส เจ ภิกฺขุโน ตมฺหาปิ นิมิตฺตา อญฺญํ นิมิตฺตํ มนสิกโรโต กุสลู\-ปสํหิตํ อุปฺปชฺชนฺเตว ปาปกา อกุสลา วิตกฺกา ฉนฺทูปสํหิตาปิ โทสูปสํหิตาปิ โมหูปสํหิตาปิ. เตน ภิกฺขเว ภิกฺขุนา เตสํ วิตกฺกานํ อาทีนโว อุปปริ\-กฺขิ\-ตพฺโพ “อิติปิเม วิตกฺกา อกุสลา, \csromanfolio{169}อิติปิเม วิตกฺกา สาวชฺชา, อิติปิเม วิตกฺกา ทุกฺขวิปากา”ติ. ตสฺส เตสํ วิตกฺกานํ อาทีนวํ อุปปริกฺขโต เย ปาปกา อกุสลา วิตกฺกา ฉนฺทูปสํหิตาปิ โทสูปสํหิตาปิ โมหูปสํหิตาปิ, เต ปหียนฺติ เต อพฺภตฺถํ คจฺฉนฺติ, เตสํ ปหานา อชฺฌตฺตเมว จิตฺตํ สนฺติฏฺฐติ สนฺนิสีทติ เอโกทิ โหติ สมาธิยติ. (2)}

\proseitem{218}{ตสฺส เจ ภิกฺขเว ภิกฺขุโน เตสมฺปิ วิตกฺกานํ อาทีนวํ อุปปริกฺขโต อุปฺปชฺชนฺเตว ปาปกา อกุสลา วิตกฺกา ฉนฺทูปสํหิตาปิ โทสูปสํหิตาปิ โมหูปสํหิตาปิ. เตน ภิกฺขเว ภิกฺขุนา เตสํ วิตกฺกานํ อสติอมนสิกาโร อาปชฺชิตพฺโพ. ตสฺส เตสํ วิตกฺกานํ อสติ- อมนสิการํ อาปชฺชโต เย ปาปกา อกุสลา วิตกฺกา ฉนฺทูปสํหิตาปิ โทสูปสํหิตาปิ โมหูปสํหิตาปิ, เต ปหียนฺติ เต อพฺภตฺถํ คจฺฉนฺติ, เตสํ ปหานา อชฺฌตฺตเมว จิตฺตํ สนฺติฏฺฐติ สนฺนิสีทติ เอโกทิ โหติ สมาธิยติ. เสยฺยถาปิ ภิกฺขเว จกฺขุมา ปุริโส อาปาถคตานํ รูปานํ อทสฺสนกาโม อสฺส, โส นิมีเลยฺย วา อญฺเญน วา อปโลเกยฺย. เอวเมว โข ภิกฺขเว ตสฺส เจ ภิกฺขุโน เตสมฺปิ วิตกฺกานํ อาทีนวํ อุปปริกฺขโต อุปฺปชฺชนฺเตว ปาปกา อกุสลา วิตกฺกา ฉนฺทูปสํหิตาปิ โทสูปสํหิตาปิ โมหูปสํหิตาปิ, เต ปหียนฺติ เต อพฺภตฺถํ คจฺฉนฺติ, เตสํ ปหานา อชฺฌตฺตเมว จิตฺตํ สนฺติฏฺฐติ สนฺนิสีทติ เอโกทิ โหติ สมาธิยติ. (3)}

\proseitem{219}{ตสฺส เจ ภิกฺขเว ภิกฺขุโน เตสมฺปิ วิตกฺกานํ อสติ- อมนสิการํ อาปชฺชโต อุปฺปชฺชนฺเตว ปาปกา อกุสลา วิตกฺกา ฉนฺทูปสํหิตาปิ โทสูปสํหิตาปิ โมหูปสํหิตาปิ. เตน ภิกฺขเว ภิกฺขุนา เตสํ วิตกฺกานํ วิตกฺก\-สงฺขาร\-สณฺฐานํ มนสิ กาตพฺพํ. ตสฺส เตสํ วิตกฺกานํ วิตกฺก\-สงฺขาร\-สณฺฐานํ มนสิกโรโต เย ปาปกา อกุสลา วิตกฺกา ฉนฺทูปสํหิตาปิ โทสูปสํหิตาปิ โมหูปสํหิตาปิ, เต ปหียนฺติ เต อพฺภตฺถํ คจฺฉนฺติ, เตสํ ปหานา อชฺฌตฺตเมว จิตฺตํ สนฺติฏฺฐติ สนฺนิสีทติ เอโกทิ โหติ สมาธิยติ. เสยฺยถาปิ ภิกฺขเว ปุริโส สีฆํ คจฺเฉยฺย, ตสฺส เอวมสฺส “กิํ นุ โข อหํ สีฆํ คจฺฉามิ, ยํนูนาหํ สณิกํ คจฺเฉยฺยนฺ”ติ, โส สณิกํ คจฺเฉยฺย. ตสฺส เอวมสฺส “กิํ นุ โข อหํ สณิกํ คจฺฉามิ, ยํนูนาหํ ติฏฺเฐยฺยนฺ”ติ, โส \csromanfolio{170}ติฏฺเฐยฺย. ตสฺส เอวมสฺส “กิํ นุ โข อหํ ฐิโต, ยํนูนาหํ นิสีเทยฺยนฺ”ติ, โส นิสีเทยฺย. ตสฺส เอวมสฺส “กิํ นุ โข อหํ นิสินฺโน, ยํนูนาหํ นิปชฺเชยฺยนฺ”ติ, โส นิปชฺเชยฺย. เอวํ หิ โส ภิกฺขเว ปุริโส โอฬาริกํ โอฬาริกํ อิริยาปถํ อภินิวชฺเชตฺวา\csromansharedfootnote{d9d8523ec61d7eb7}{อภินิสฺสชฺเชตฺวา (สฺยา) 2. ทนฺเต + อภิทนฺตํ + อาธายาติ ฏีกายํ ปทจฺเฉโท. ทนฺเตภีติ ปเนตฺถ กรณตฺโถ ยุตฺโต วิย ทิสฺสติ.} สุขุมํ สุขุมํ อิริยาปถํ กปฺเปยฺย. เอวเมว โข ภิกฺขเว ตสฺส เจ ภิกฺขุโน เตสมฺปิ วิตกฺกานํ อสติอมนสิการํ อาปชฺชโต อุปฺปชฺชนฺเตว ปาปกา อกุสลา วิตกฺกา ฉนฺทูปสํหิตาปิ โทสูปสํหิตาปิ โมหูปสํหิตาปิ, เต ปหียนฺติ เต อพฺภตฺถํ คจฺฉนฺติ, เตสํ ปหานา อชฺฌตฺตเมว จิตฺตํ สนฺติฏฺฐติ สนฺนิสีทติ เอโกทิ โหติ สมาธิยติ. (4)}

\proseitem{220}{ตสฺส เจ ภิกฺขเว ภิกฺขุโน เตสมฺปิ วิตกฺกานํ วิตกฺก\-สงฺขาร\-สณฺฐานํ มนสิกโรโต อุปฺปชฺชนฺเตว ปาปกา อกุสลา วิตกฺกา ฉนฺทูปสํหิตาปิ โทสูปสํหิตาปิ โมหูปสํหิตาปิ. เตน ภิกฺขเว ภิกฺขุนา ทนฺเตภิ\-ทนฺตมา\-ธาย\csromansharedfootnote{c0be0242d8ec34ba}{ทนฺเต + อภิทนฺตํ + อาธายาติ ฏีกายํ ปทจฺเฉโท. ทนฺเตภีติ ปเนตฺถ กรณตฺโถ ยุตฺโต วิย ทิสฺสติ.} ชิวฺหาย ตาลุํ อาหจฺจ เจตสา จิตฺตํ อภินิ\-คฺคณฺหิ\-ตพฺพํ อภินิ\-ปฺปีเฬ\-ตพฺพํ อภิส\-นฺตาเป\-ตพฺพํ. ตสฺส ทนฺเตภิ\-ทนฺตมา\-ธาย ชิวฺหาย ตาลุํ อาหจฺจ เจตสา จิตฺตํ อภินิคฺคณฺหโต อภินิปฺปีฬยโต อภิสนฺตาปยโต เย ปาปกา อกุสลา วิตกฺกา ฉนฺทูปสํหิตาปิ โทสูปสํหิตาปิ โมหูปสํหิตาปิ, เต ปหียนฺติ เต อพฺภตฺถํ คจฺฉนฺติ, เตสํ ปหานา อชฺฌตฺตเมว จิตฺตํ สนฺติฏฺฐติ สนฺนิสีทติ เอโกทิ โหติ สมาธิยติ. เสยฺยถาปิ ภิกฺขเว พลวา ปุริโส ทุพฺพลตรํ ปุริสํ สีเส วา คเล วา ขนฺเธ วา คเหตฺวา อภินิคฺคเณฺหยฺย อภินิปฺปีเฬยฺย อภิสนฺตาเปยฺย. เอวเมว โข ภิกฺขเว ตสฺส เจ ภิกฺขุโน เตสมฺปิ วิตกฺกานํ วิตกฺก\-สงฺขาร\-สณฺฐานํ มนสิกโรโต อุปฺปชฺชนฺเตว ปาปกา อกุสลา วิตกฺกา ฉนฺทูปสํหิตาปิ โทสูปสํหิตาปิ โมหูปสํหิตาปิ. เตน ภิกฺขเว ภิกฺขุนา ทนฺเตภิ\-ทนฺตมา\-ธาย ชิวฺหาย ตาลุํ อาหจฺจ เจตสา จิตฺตํ อภินิ\-คฺคณฺหิ\-ตพฺพํ อภินิ\-ปฺปีเฬ\-ตพฺพํ อภิส\-นฺตาเป\-ตพฺพํ. ตสฺส ทนฺเตภิ\-ทนฺตมา\-ธาย ชิวฺหาย ตาลุํ อาหจฺจ เจตสา จิตฺตํ อภินิคฺคณฺหโต อภินิปฺปีฬยโต อภิสนฺตาปยโต เย ปาปกา อกุสลา วิตกฺกา ฉนฺทูปสํหิตาปิ โทสูปสํหิตาปิ \csromanfolio{171}โมหูปสํหิตาปิ, เต ปหียนฺติ เต อพฺภตฺถํ คจฺฉนฺติ, เตสํ ปหานา อชฺฌตฺตเมว จิตฺตํ สนฺติฏฺฐติ สนฺนิสีทติ เอโกทิ โหติ สมาธิยติ.}

\proseitem{221}{ยโต โข\csromansharedfootnote{260272ef66b278ff}{ยโต จ โข (สฺยา, ก)} ภิกฺขเว ภิกฺขุโน ยํ นิมิตฺตํ อาคมฺม ยํ นิมิตฺตํ มนสิกโรโต อุปฺปชฺชนฺติ ปาปกา อกุสลา วิตกฺกา ฉนฺทูปสํหิตาปิ โทสูปสํหิตาปิ โมหูปสํหิตาปิ. ตสฺส ตมฺหา นิมิตฺตา อญฺญํ นิมิตฺตํ มนสิกโรโต กุสลู\-ปสํหิตํ เย ปาปกา อกุสลา วิตกฺกา ฉนฺทูปสํหิตาปิ โทสูปสํหิตาปิ โมหูปสํหิตาปิ, เต ปหียนฺติ เต อพฺภตฺถํ คจฺฉนฺติ, เตสํ ปหานา อชฺฌตฺตเมว จิตฺตํ สนฺติฏฺฐติ สนฺนิสีทติ เอโกทิ โหติ สมาธิยติ. เตสมฺปิ วิตกฺกานํ อาทีนวํ อุปปริกฺขโต เย ปาปกา อกุสลา วิตกฺกา ฉนฺทูปสํหิตาปิ โทสูปสํหิตาปิ โมหูปสํหิตาปิ, เต ปหียนฺติ เต อพฺภตฺถํ คจฺฉนฺติ, เตสํ ปหานา อชฺฌตฺตเมว จิตฺตํ สนฺติฏฺฐติ สนฺนิสีทติ เอโกทิ โหติ สมาธิยติ. เตสมฺปิ วิตกฺกานํ อสติอมนสิการํ อาปชฺชโต เย ปาปกา อกุสลา วิตกฺกา ฉนฺทูปสํหิตาปิ โทสูปสํหิตาปิ โมหูปสํหิตาปิ, เต ปหียนฺติ เต อพฺภตฺถํ คจฺฉนฺติ, เตสํ ปหานา อชฺฌตฺตเมว จิตฺตํ สนฺติฏฺฐติ สนฺนิสีทติ เอโกทิ โหติ สมาธิยติ. เตสมฺปิ วิตกฺกานํ วิตกฺก\-สงฺขาร\-สณฺฐานํ มนสิกโรโต เย ปาปกา อกุสลา วิตกฺกา ฉนฺทูปสํหิตาปิ โทสูปสํหิตาปิ โมหูปสํหิตาปิ, เต ปหียนฺติ เต อพฺภตฺถํ คจฺฉนฺติ, เตสํ ปหานา อชฺฌตฺตเมว จิตฺตํ สนฺติฏฺฐติ สนฺนิสีทติ เอโกทิ โหติ สมาธิยติ. ทนฺเตภิ\-ทนฺตมา\-ธาย ชิวฺหาย ตาลุํ อาหจฺจ เจตสา จิตฺตํ อภินิคฺคณฺหโต อภินิปฺปีฬยโต อภิสนฺตาปยโต เย ปาปกา อกุสลา วิตกฺกา ฉนฺทูปสํหิตาปิ โทสูปสํหิตาปิ โมหูปสํหิตาปิ, เต ปหียนฺติ เต อพฺภตฺถํ คจฺฉนฺติ, เตสํ ปหานา อชฺฌตฺตเมว จิตฺตํ สนฺติฏฺฐติ สนฺนิสีทติ เอโกทิ โหติ สมาธิยติ. อยํ วุจฺจติ ภิกฺขเว ภิกฺขุ วสี วิตกฺก\-ปริยาย\-ปเถสุ. ยํ วิตกฺกํ อากงฺขิสฺสติ, ตํ วิตกฺกํ วิตกฺเกสฺสติ. ยํ วิตกฺกํ นากงฺขิสฺสติ, น ตํ วิตกฺกํ วิตกฺเกสฺสติ. อจฺเฉจฺฉิ ตณฺหํ, วิวตฺตยิ\csromansharedfootnote{ad9184755c407031}{วาวตฺตยิ (สี, อิ)} สญฺโญชนํ, สมฺมา มานาภิสมยา อนฺตมกาสิ ทุกฺขสฺสาติ.}

\csromanmatikaanchor{mtk.49}
\csromantocmarkref{subsection}{อุทฺทานคาถา}{mtk.49}
\bookmark[dest={mtk.49},level=2]{อุทฺทานคาถา}
\prosewithcloser{\csromanfolio{172}อิทมโวจ ภควา. อตฺตมนา เต ภิกฺขู ภควโต ภาสิตํ อภินนฺทุนฺติ.}{วิตกฺก\-สณฺฐาน\-สุตฺตํ นิฏฺฐิตํ ทสมํ.}
\nitthitammid{\textbf{สีหนาทวคฺโค นิฏฺฐิโต ทุติโย.}}
\csromancenter{\textbf{ตสฺสุทฺทานํ}}
\setlength{\csromangathapagewidth}{0pt}
\setlength{\csromangathawakwidth}{0pt}
\csromangathameasuregroup{}{จูฬสีหนาทโลมหํสวโร, \\ มหาจูฬทุกฺขกฺขนฺธอนุมานิกสุตฺตํ. \\ ขิลปตฺถมธุปิณฺฑิกทฺวิธาวิตกฺก, \\ ปญฺจนิมิตฺตกถา ปุน วคฺโค.}
\csromangathameasurewak{จูฬสีหนาทโลมหํสวโร, \\ มหาจูฬทุกฺขกฺขนฺธอนุมานิกสุตฺตํ. \\ ขิลปตฺถมธุปิณฺฑิกทฺวิธาวิตกฺก, \\ ปญฺจนิมิตฺตกถา ปุน วคฺโค.}
\setlength{\csromangathaleftcol}{0pt}
\csromangathagroup{\csromangathastanza{\csromangathawak{จูฬสีหนาทโลมหํสวโร, \\ มหาจูฬทุกฺขกฺขนฺธอนุมานิกสุตฺตํ. \\ ขิลปตฺถมธุปิณฺฑิกทฺวิธาวิตกฺก, \\ ปญฺจนิมิตฺตกถา ปุน วคฺโค.}}}

\csromanmatikaanchor{mtk.50}
\csromantocmarknopage{chapter}{3. โอปมฺมวคฺค}{mtk.50}
\bookmark[dest={mtk.50},level=0]{3. โอปมฺมวคฺค}
\chapterheadpageread{3. \textbf{โอปมฺมวคฺค}}
\csromanmatikaanchor{mtk.51}
\csromantocmarkref{subsection}{1. กกจูปมสุตฺต}{mtk.51}
\bookmark[dest={mtk.51},level=2]{1. กกจูปมสุตฺต}
\csromanheader{2}{1. กกจูปม\-สุตฺต}
\proseitem{222}{\csromanfolio{173}เอวํ เม สุตํ\csromandash{}เอกํ สมยํ ภควา สาวตฺถิยํ วิหรติ เชตวเน อนาถ\-ปิณฺฑิกสฺส อาราเม. เตน โข ปน สมเยน อายสฺมา โมฬิยผคฺคุโน ภิกฺขุนีหิ สทฺธิํ อติเวลํ สํสฏฺโฐ วิหรติ. เอวํ สํสฏฺโฐ อายสฺมา โมฬิยผคฺคุโน ภิกฺขุนีหิ สทฺธิํ วิหรติ. สเจ โกจิ ภิกฺขุ อายสฺมโต โมฬิย\-ผคฺคุนสฺส สมฺมุขา ตาสํ ภิกฺขุนีนํ อวณฺณํ ภาสติ, เตนายสฺมา โมฬิยผคฺคุโน กุปิโต อนตฺตมโน อธิกรณมฺปิ กโรติ. สเจ ปน โกจิ ภิกฺขุ ตาสํ ภิกฺขุนีนํ สมฺมุขา อายสฺมโต โมฬิย\-ผคฺคุนสฺส อวณฺณํ ภาสติ, เตน ตา ภิกฺขุนิโย กุปิตา อนตฺตมนา อธิกรณมฺปิ กโรนฺติ. เอวํ สํสฏฺโฐ อายสฺมา โมฬิยผคฺคุโน ภิกฺขุนีหิ สทฺธิํ วิหรติ. อถ โข อญฺญตโร ภิกฺขุ เยน ภควา เตนุปสงฺกมิ, อุปสงฺกมิตฺวา ภควนฺตํ อภิวาเทตฺวา เอกมนฺตํ นิสีทิ, เอกมนฺตํ นิสินฺโน โข โส ภิกฺขุ ภควนฺตํ เอตทโวจ “อายสฺมา ภนฺเต โมฬิยผคฺคุโน ภิกฺขุนีหิ สทฺธิํ อติเวลํ สํสฏฺโฐ วิหรติ. เอวํ สํสฏฺโฐ ภนฺเต อายสฺมา โมฬิยผคฺคุโน ภิกฺขุนีหิ สทฺธิํ วิหรติ. สเจ โกจิ ภิกฺขุ อายสฺมโต โมฬิย\-ผคฺคุนสฺส สมฺมุขา ตาสํ ภิกฺขุนีนํ อวณฺณํ ภาสติ, เตนายสฺมา โมฬิยผคฺคุโน กุปิโต อนตฺตมโน อธิกรณมฺปิ กโรติ. สเจ ปน โกจิ ภิกฺขุ ตาสํ ภิกฺขุนีนํ สมฺมุขา อายสฺมโต โมฬิย\-ผคฺคุนสฺส อวณฺณํ ภาสติ, เตน ตา ภิกฺขุนิโย กุปิตา อนตฺตมนา อธิกรณมฺปิ กโรนฺติ. เอวํ สํสฏฺโฐ ภนฺเต อายสฺมา โมฬิยผคฺคุโน ภิกฺขุนีหิ สทฺธิํ วิหรตี”ติ.}

\proseitem{223}{อถ โข ภควา อญฺญตรํ ภิกฺขุํ อามนฺเตสิ “เอหิ ตฺวํ ภิกฺขุ มม วจเนน โมฬิย\-ผคฺคุนํ ภิกฺขุํ อามนฺเตหิ ‘สตฺถา ตํ อาวุโส ผคฺคุน อามนฺเตตี’ติ”. เอวํ ภนฺเตติ โข โส ภิกฺขุ ภควโต ปฏิสฺสุตฺวา เยนายสฺมา โมฬิยผคฺคุโน เตนุปสงฺกมิ, อุปสงฺกมิตฺวา อายสฺมนฺตํ โมฬิย\-ผคฺคุนํ เอตทโวจ “สตฺถา ตํ อาวุโส ผคฺคุน อามนฺเตตี”ติ. เอวมาวุโสติ โข อายสฺมา โมฬิยผคฺคุโน \csromanfolio{174}ตสฺส ภิกฺขุโน ปฏิสฺสุตฺวา เยน ภควา เตนุปสงฺกมิ, อุปสงฺกมิตฺวา ภควนฺตํ อภิวาเทตฺวา เอกมนฺตํ นิสีทิ, เอกมนฺตํ นิสินฺนํ โข อายสฺมนฺตํ โมฬิย\-ผคฺคุนํ ภควา เอตทโวจ\csromandash{}}

\prose{สจฺจํ กิร ตฺวํ ผคฺคุน ภิกฺขุนีหิ สทฺธิํ อติเวลํ สํสฏฺโฐ วิหรสิ. เอวํ สํสฏฺโฐ กิร ตฺวํ ผคฺคุน ภิกฺขุนีหิ สทฺธิํ วิหรสิ. สเจ โกจิ ภิกฺขุ ตุยฺหํ สมฺมุขา ตาสํ ภิกฺขุนีนํ อวณฺณํ ภาสติ, เตน ตฺวํ กุปิโต อนตฺตมโน อธิกรณมฺปิ กโรสิ. สเจ ปน โกจิ ภิกฺขุ ตาสํ ภิกฺขุนีนํ สมฺมุขา ตุยฺหํ อวณฺณํ ภาสติ, เตน ตา ภิกฺขุนิโย กุปิตา อนตฺตมนา อธิกรณมฺปิ กโรนฺติ. เอวํ สํสฏฺโฐ กิร ตฺวํ ผคฺคุน ภิกฺขุนีหิ สทฺธิํ วิหรสีติ. เอวํ ภนฺเตติ. นนุ ตฺวํ ผคฺคุน กุลปุตฺโต สทฺธา อคารสฺมา อนคาริยํ ปพฺพชิโตติ. เอวํ ภนฺเตติ.}

\proseitem{224}{น โข เต เอตํ ผคฺคุน ปติรูปํ กุลปุตฺตสฺส สทฺธา อคารสฺมา อนคาริยํ ปพฺพชิตสฺส, ยํ ตฺวํ ภิกฺขุนีหิ สทฺธิํ อติเวลํ สํสฏฺโฐ วิหเรยฺยาสิ. ตสฺมาติห ผคฺคุน ตว เจปิ โกจิ สมฺมุขา ตาสํ ภิกฺขุนีนํ อวณฺณํ ภาเสยฺย. ตตฺราปิ ตฺวํ ผคฺคุน เย เคหสิตา\csromansharedfootnote{8f464376fb7e4b2a}{เคหสฺสิตา (?)} ฉนฺทา เย เคหสิตา วิตกฺกา, เต ปชเหยฺยาสิ. ตตฺราปิ เต ผคฺคุน เอวํ สิกฺขิตพฺพํ “น เจว เม จิตฺตํ วิปริณตํ ภวิสฺสติ, น จ ปาปิกํ วาจํ นิจฺฉาเรสฺสามิ, หิตานุกมฺปี จ วิหริสฺสามิ เมตฺตจิตฺโต น โทสนฺตโร”ติ. เอวํ หิ เต ผคฺคุน สิกฺขิตพฺพํ.}

\prose{ตสฺมาติห ผคฺคุน ตว เจปิ โกจิ สมฺมุขา ตาสํ ภิกฺขุนีนํ ปาณินา ปหารํ ทเทยฺย, เลฑฺฑุนา ปหารํ ทเทยฺย, ทณฺเฑน ปหารํ ทเทยฺย, สตฺเถน ปหารํ ทเทยฺย. ตตฺราปิ ตฺวํ ผคฺคุน เย เคหสิตา ฉนฺทา เย เคหสิตา วิตกฺกา, เต ปชเหยฺยาสิ. ตตฺราปิ เต ผคฺคุน เอวํ สิกฺขิตพฺพํ “น เจว เม จิตฺตํ วิปริณตํ ภวิสฺสติ, น จ ปาปิกํ วาจํ นิจฺฉาเรสฺสามิ, หิตานุกมฺปี จ วิหริสฺสามิ เมตฺตจิตฺโต น โทสนฺตโร”ติ. เอวํ หิ เต ผคฺคุน สิกฺขิตพฺพํ.}

\prose{ตสฺมาติห ผคฺคุน ตว เจปิ โกจิ สมฺมุขา อวณฺณํ ภาเสยฺย. ตตฺราปิ ตฺวํ ผคฺคุน เย เคหสิตา ฉนฺทา เย เคหสิตา วิตกฺกา,}

\prose{\csromanfolio{175}เต ปชเหยฺยาสิ. ตตฺราปิ เต ผคฺคุน เอวํ สิกฺขิตพฺพํ “น เจว เม จิตฺตํ วิปริณตํ ภวิสฺสติ, น จ ปาปิกํ วาจํ นิจฺฉาเรสฺสามิ, หิตานุกมฺปี จ วิหริสฺสามิ เมตฺตจิตฺโต น โทสนฺตโร”ติ. เอวํ หิ เต ผคฺคุน สิกฺขิตพฺพํ.}

\prose{ตสฺมาติห ผคฺคุน ตว เจปิ โกจิ ปาณินา ปหารํ ทเทยฺย, เลฑฺฑุนา ปหารํ ทเทยฺย, ทณฺเฑน ปหารํ ทเทยฺย, สตฺเถน ปหารํ ทเทยฺย. ตตฺราปิ ตฺวํ ผคฺคุน เย เคหสิตา ฉนฺทา เย เคหสิตา วิตกฺกา, เต ปชเหยฺยาสิ. ตตฺราปิ เต ผคฺคุน เอวํ สิกฺขิตพฺพํ “น เจว เม จิตฺตํ วิปริณตํ ภวิสฺสติ, น จ ปาปิกํ วาจํ นิจฺฉาเรสฺสามิ, หิตานุกมฺปี จ วิหริสฺสามิ เมตฺตจิตฺโต น โทสนฺตโร”ติ. เอวํ หิ เต ผคฺคุน สิกฺขิตพฺพนฺติ.}

\proseitem{225}{อถ โข ภควา ภิกฺขู อามนฺเตสิ\csromandash{}อาราธยิํสุ วต เม ภิกฺขเว ภิกฺขู เอกํ สมยํ จิตฺตํ. อิธาหํ ภิกฺขเว ภิกฺขู อามนฺเตสิํ ‘อหํ โข ภิกฺขเว เอกาสน\-โภชนํ ภุญฺชามิ, เอกาสน\-โภชนํ โข อหํ ภิกฺขเว ภุญฺชมาโน อปฺปาพาธตญฺจ สญฺชานามิ อปฺปาตงฺกตญฺจ ลหุฏฺฐานญฺจ พลญฺจ ผาสุวิหารญฺจ, เอถ ตุเมฺหปิ ภิกฺขเว เอกาสน\-โภชนํ ภุญฺชถ, เอกาสน\-โภชนํ โข ภิกฺขเว ตุเมฺหปิ ภุญฺชมานา อปฺปาพาธตญฺจ สญฺชานิสฺสถ อปฺปาตงฺกตญฺจ ลหุฏฺฐานญฺจ พลญฺจ ผาสุวิหารญฺจา’ติ. น เม ภิกฺขเว เตสุ ภิกฺขูสุ อนุสาสนี กรณียา อโหสิ, สตุปฺปาทกรณียเมว เม ภิกฺขเว เตสุ ภิกฺขูสุ อโหสิ.}

\prose{เสยฺยถาปิ ภิกฺขเว สุภูมิยํ จตุมหาปเถ อาชญฺญรโถ ยุตฺโต อสฺส ฐิโต โอธสฺตปโตโท. ตเมนํ ทกฺโข โยคฺคาจริโย อสฺส\-ทมฺม\-สารถิ อภิรุหิตฺวา วาเมน หตฺเถน รสฺมิโย คเหตฺวา ทกฺขิเณน หตฺเถน ปโตทํ คเหตฺวา เยนิจฺฉกํ ยทิจฺฉกํ สาเรยฺยปิ ปจฺจาสาเรยฺยปิ. เอวเมว โข ภิกฺขเว น เม เตสุ ภิกฺขูสุ อนุสาสนี กรณียา อโหสิ, สตุปฺปาทกรณียเมว เม ภิกฺขเว เตสุ ภิกฺขูสุ อโหสิ. ตสฺมาติห ภิกฺขเว ตุเมฺหปิ อกุสลํ ปชหถ, กุสเลสุ ธมฺเมสุ อาโยคํ กโรถ. เอวํ หิ ตุเมฺหปิ อิมสฺมิํ ธมฺมวินเย วุทฺธิํ วิรูฬฺหิํ เวปุลฺลํ อาปชฺชิสฺสถ.}

\prose{\csromanfolio{176}เสยฺยถาปิ ภิกฺขเว คามสฺส วา นิคมสฺส วา อวิทูเร มหนฺตํ สาลวนํ, ตํ จสฺส เอฬณฺเฑหิ สญฺฉนฺนํ. ตสฺส โกจิเทว ปุริโส อุปฺปชฺเชยฺย อตฺถกาโม หิตกาโม โยคกฺเขมกาโม. โส ยา ตา สาลลฏฺฐิโย กุฏิลา โอชาปหรณิโย\csromansharedfootnote{de3ff804b4206153}{โอชหรณิโย (ก)}, ตา เฉตฺวา\csromansharedfootnote{309c97b908806fe5}{ตจฺเฉตฺวา (สี, สฺยา, อิ)} พหิทฺธา นีหเรยฺย. อนฺโตวนํ สุวิโสธิตํ วิโสเธยฺย, ยา ปน ตา สาลลฏฺฐิโย อุชุกา สุชาตา, ตา สมฺมา ปริหเรยฺย. เอวํ เหตํ ภิกฺขเว สาลวนํ อปเรน สมเยน วุทฺธิํ วิรูฬฺหิํ เวปุลฺลํ อาปชฺเชยฺย. เอวเมว โข ภิกฺขเว ตุเมฺหปิ อกุสลํ ปชหถ, กุสเลสุ ธมฺเมสุ อาโยคํ กโรถ. เอวํ หิ ตุเมฺหปิ อิมสฺมิํ ธมฺมวินเย วุทฺธิํ วิรูฬฺหิํ เวปุลฺลํ อาปชฺชิสฺสถ.}

\proseitem{226}{ภูตปุพฺพํ ภิกฺขเว อิมิสฺสาเยว สาวตฺถิยา เวเทหิกา นาม คหปตานี อโหสิ. เวเทหิกาย ภิกฺขเว คหปตานิยา เอวํ กลฺยาโณ กิตฺติสทฺโท อพฺภุคฺคโต “โสรตา เวเทหิกา คหปตานี, นิวาตา เวเทหิกา คหปตานี, อุปสนฺตา เวเทหิกา คหปตานี”ติ. เวเทหิกาย โข ปน ภิกฺขเว คหปตานิยา กาฬี นาม ทาสี อโหสิ ทกฺขา อนลสา สุสํวิหิต\-กมฺมนฺตา.}

\prose{อถ โข ภิกฺขเว กาฬิยา ทาสิยา เอตทโหสิ “มยฺหํ โข อยฺยาย เอวํ กลฺยาโณ กิตฺติสทฺโท อพฺภุคฺคโต ‘โสรตา เวเทหิกา คหปตานี, นิวาตา เวเทหิกา คหปตานี, อุปสนฺตา เวเทหิกา คหปตานี’ติ. กิํ นุ โข เม อยฺยา สนฺตํเยว นุ โข อชฺฌตฺตํ โกปํ น ปาตุกโรติ อุทาหุ อสนฺตํ, อุทาหุ มยฺหเมเวเต\csromansharedfootnote{72ab47d6e98f49e0}{มเยฺหเวเต (สี, อิ)} กมฺมนฺตา สุสํวิหิตา, เยน เม อยฺยา สนฺตํเยว อชฺฌตฺตํ โกปํ น ปาตุ กโรติ โน อสนฺตํ. ยํนูนาหํ อยฺยํ วีมํเสยฺยนฺ”ติ. อถ โข ภิกฺขเว กา ทาสี ทิวา อุฏฺฐาสิ. อถ โข ภิกฺขเว เวเทหิกา คหปตานี กาฬิํ ทาสิํ เอตทโวจ “เห เช กาฬี”ติ. กิํ อยฺเยติ. กิํ เช ทิวา อุฏฺฐาสีติ. น ขฺวยฺเย\csromansharedfootnote{f0ca801d51afdb46}{น โข อยฺเย (สี, อิ)} กิญฺจีติ. “โน วต เร กิญฺจิ, ปาปิ ทาสิ\csromansharedfootnote{2d272e6f373d5f75}{ปาปทาสิ (สฺยา, ก)} ทิวา อุฏฺฐาสี”ติ กุปิตา อนตฺตมนา ภากุฏิํ\csromansharedfootnote{47f676f6783fdbe6}{ภูกุฏิํ (สี, อิ) ภกุฏิํ (สฺยา)} อกาสิ. อถ โข ภิกฺขเว กาฬิยา}

\prose{\csromanfolio{177}ทาสิยา เอตทโหสิ “สนฺตํเยว โข เม อยฺยา อชฺฌตฺตํ โกปํ น ปาตุกโรติ โน อสนฺตํ. มยฺหเมเวเต กมฺมนฺตา สุสํวิหิตา, เยน เม อยฺยา สนฺตํเยว อชฺฌตฺตํ โกปํ น ปาตุกโรติ โน อสนฺตํ. ยํนูนาหํ ภิยฺโยโส มตฺตาย อยฺยํ วีมํเสยฺยนฺ”ติ.}

\prose{อถ โข ภิกฺขเว กาฬี ทาสี ทิวาตรํเยว อุฏฺฐาสิ. อถ โข ภิกฺขเว เวเทหิกา คหปตานี กาฬิํ ทาสิํ เอตทโวจ “เห เช กาฬี”ติ. กิํ อยฺเยติ. กิํ เช ทิวาตรํ อุฏฺฐาสีติ. น ขฺวยฺเย กิญฺจีติ. “โน วต เร กิญฺจิ, ปาปิ ทาสิ ทิวาตรํ อุฏฺฐาสี”ติ กุปิตา อนตฺตมนา อนตฺตมน\-วาจํ นิจฺฉาเรสิ. อถ โข ภิกฺขเว กาฬิยา ทาสิยา เอตทโหสิ “สนฺตํเยว โข เม อยฺยา อชฺฌตฺตํ โกปํ น ปาตุกโรติ โน อสนฺตํ. มยฺหเมเวเต กมฺมนฺตา สุสํวิหิตา, เยน เม อยฺยา สนฺตํเยว อชฺฌตฺตํ โกปํ น ปาตุกโรติ โน อสนฺตํ. ยํนูนาหํ ภิยฺโยโส มตฺตาย อยฺยํ วีมํเสยฺยนฺ”ติ.}

\prose{อถ โข ภิกฺขเว กาฬี ทาสี ทิวาตรํเยว อุฏฺฐาสิ. อถ โข ภิกฺขเว เวเทหิกา คหปตานี กาฬิํ ทาสิํ เอตทโวจ “เห เช กาฬี”ติ. กิํ อยฺเยติ. กิํ เช ทิวา อุฏฺฐาสีติ. น ขฺวยฺเย กิญฺจีติ. “โน วต เร กิญฺจิ, ปาปิ ทาสิ ทิวา อุฏฺฐาสี”ติ กุปิตา อนตฺตมนา อคฺคฬสูจิํ คเหตฺวา สีเส ปหารํ อทาสิ, สีสํ โวภินฺทิ\csromansharedfootnote{f5fc3d7233f97b4e}{วิ + อว + ภินฺทิ = โวภินฺทิ}. อถ โข ภิกฺขเว กาฬี ทาสี ภินฺเนน สีเสน โลหิเตน คลนฺเตน ปฏิวิสฺสกานํ อุชฺฌาเปสิ “ปสฺสถยฺเย โสรตาย กมฺมํ, ปสฺสถยฺเย นิวาตาย กมฺมํ, ปสฺสถยฺเย อุปสนฺตาย กมฺมํ. กถํ หิ นาม เอกทาสิกาย ‘ทิวา อุฏฺฐาสี’ติ กุปิตา อนตฺตมนา อคฺคฬสูจิํ คเหตฺวา สีเส ปหารํ ทสฺสติ, สีสํ โวภินฺทิสฺสตี”ติ.}

\prose{อถ โข ภิกฺขเว เวเทหิกาย คหปตานิยา อปเรน สมเยน เอวํ ปาปโก กิตฺติสทฺโท อพฺภุคฺคจฺฉิ “จณฺฑี เวเทหิกา คหปตานี, อนิวาตา เวเทหิกา คหปตานี, อนุปสนฺตา เวเทหิกา คหปตานี”ติ.}

\prose{เอวเมว โข ภิกฺขเว อิเธกจฺโจ ภิกฺขุ ตาวเทว โสรตโสรโต โหติ, นิวาตนิวาโต โหติ, อุปสนฺตู\-ปสนฺโต โหติ, ยาว น อมนาปา วจนปถา ผุสนฺติ. ยโต จ ภิกฺขเว}

\prose{\csromanfolio{178}ภิกฺขุํ อมนาปา วจนปถา ผุสนฺติ, อถ ภิกฺขุ โสรโตติ เวทิตพฺโพ, นิวาโตติ เวทิตพฺโพ, อุปสนฺโตติ เวทิตพฺโพ. นาหํ ตํ ภิกฺขเว ภิกฺขุํ สุวโจติ วทามิ, โย จีวร\-ปิณฺฑปาต\-เสนาสน\-คิลานปฺปจฺจย\-เภสชฺช\-ปริกฺขาร\-เหตุ สุวโจ โหติ, โสวจสฺสตํ อาปชฺชติ. ตํ กิสฺส เหตุ, ตํ หิ โส ภิกฺขเว ภิกฺขุ จีวร\-ปิณฺฑปาต\-เสนาสน\-คิลานปฺปจฺจย\-เภสชฺช\-ปริกฺขารํ อลภมาโน น สุวโจ โหติ, น โสวจสฺสตํ อาปชฺชติ. โย จ โข ภิกฺขเว ภิกฺขุ ธมฺมํเยว สกฺกโรนฺโต ธมฺมํ ครุํ กโรนฺโต ธมฺมํ มาเนนฺโต ธมฺมํ ปูเชนฺโต ธมฺมํ อปจายมาโน\csromansharedfootnote{a2779b7f3732aefe}{ธมฺมํ เยว สกฺกโรนฺโต ธมฺมํ ครุกโรนฺโต ธมฺมํ อปจายมาโน(สี, สฺยา, อิ)} สุวโจ โหติ, โสวจสฺสตํ อาปชฺชติ. ตมหํ สุวโจติ วทามิ. ตสฺมา ติห ภิกฺขเว “ธมฺมํเยว สกฺกโรนฺตา ธมฺมํ ครุํ กโรนฺตา ธมฺมํ มาเนนฺตา ธมฺมํ ปูเชนฺตา ธมฺมํ อปจายมานา สุวจา ภวิสฺสาม, โสวจสฺสตํ อาปชฺชิสฺสามา”ติ เอวํ หิ โว ภิกฺขเว สิกฺขิตพฺพํ.}

\proseitem{227}{ปญฺจิเม ภิกฺขเว วจนปถา, เยหิ โว ปเร วทมานา วเทยฺยุํ\csromandash{}กาเลน วา อกาเลน วา, ภูเตน วา อภูเตน วา, สเณฺหน วา ผรุเสน วา, อตฺถสํหิเตน วา อนตฺถ\-สํหิเตน วา, เมตฺตจิตฺตา วา โทสนฺตรา วา. กาเลน วา ภิกฺขเว ปเร วทมานา วเทยฺยุํ อกาเลน วา, ภูเตน วา ภิกฺขเว ปเร วทมานา วเทยฺยุํ อภูเตน วา, สเณฺหน วา ภิกฺขเว ปเร วทมานา วเทยฺยุํ ผรุเสน วา, อตฺถสํหิเตน วา ภิกฺขเว ปเร วทมานา วเทยฺยุํ อนตฺถ\-สํหิเตน วา, เมตฺตจิตฺตา วา ภิกฺขเว ปเร วทมานา วเทยฺยุํ โทสนฺตรา วา. ตตฺราปิ โว ภิกฺขเว เอวํ สิกฺขิตพฺพํ “น เจว โน จิตฺตํ วิปริณตํ ภวิสฺสติ, น จ ปาปิกํ วาจํ นิจฺฉาเรสฺสาม, หิตานุกมฺปี จ วิหริสฺสาม เมตฺตจิตฺตา น โทสนฺตรา. ตญฺจ ปุคฺคลํ เมตฺตา\-สหคเตน เจตสา ผริตฺวา วิหริสฺสาม, ตทารมฺมณํ จ สพฺพาวนฺตํ โลกํ เมตฺตา\-สหคเตน จิตฺเตน วิปุเลน มหคฺคเตน อปฺปมาเณน อเวเรน อพฺยาพชฺเฌน\csromansharedfootnote{f4c302ccad8d3e34}{อพฺยาปชฺเฌน (สี, สฺยา, อิ), อพฺยาปชฺเชน (ก) องฺคุตฺตรติกนิปาตฏีกา โอโลเกตพฺพา.} ผริตฺวา วิหริสฺสามา”ติ. เอวํ หิ โว ภิกฺขเว สิกฺขิตพฺพํ.}

\proseitem{228}{เสยฺยถาปิ ภิกฺขเว ปุริโส อาคจฺเฉยฺย กุทาลปิฏกํ\csromansharedfootnote{4f54287fe5b55b9c}{กุทฺทาลปิฏกํ (สี, สฺยา, อิ)} อาทาย, โส เอวํ วเทยฺย “อหํ อิมํ มหาปถวิํ อปถวิํ}

\prose{\csromanfolio{179}กริสฺสามี”ติ. โส ตตฺร ตตฺร วิขเณยฺย\csromansharedfootnote{b3605cbf068ac5c6}{ขเณยฺย (สี, สฺยา, อิ)}, ตตฺร ตตฺร วิกิเรยฺย, ตตฺร ตตฺร โอฏฺฐุเภยฺย, ตตฺร ตตฺร โอมุตฺเตยฺย, “อปถวี ภวสิ อปถวี ภวสี”ติ. ตํ กิํ มญฺญถ ภิกฺขเว, อปิ นุ โส ปุริโส อิมํ มหาปถวิํ อปถวิํ กเรยฺยาติ. โน เหตํ ภนฺเต. ตํ กิสฺส เหตุ, อยํ หิ ภนฺเต มหาปถวี คมฺภีรา อปฺปเมยฺยา, สา น สุกรา อปถวี กาตุํ. ยาวเทว จ ปน โส ปุริโส กิลมถสฺส วิฆาตสฺส ภาคี อสฺสาติ. เอวเมว โข ภิกฺขเว ปญฺจิเม วจนปถา, เยหิ โว ปเร วทมานา วเทยฺยุํ\csromandash{}กาเลน วา อกาเลน วา, ภูเตน วา อภูเตน วา, สเณฺหน วา ผรุเสน วา, อตฺถสํหิเตน วา อนตฺถ\-สํหิเตน วา, เมตฺตจิตฺตา วา โทสนฺตรา วา. กาเลน วา ภิกฺขเว ปเร วทมานา วเทยฺยุํ อกาเลน วา, ภูเตน วา ภิกฺขเว ปเร วทมานา วเทยฺยุํ อภูเตน วา, สเณฺหน วา ภิกฺขเว ปเร วทมานา วเทยฺยุํ ผรุเสน วา, อตฺถสํหิเตน วา ภิกฺขเว ปเร วทมานา วเทยฺยุํ อนตฺถ\-สํหิเตน วา, เมตฺตจิตฺตา วา ภิกฺขเว ปเร วทมานา วเทยฺยุํ โทสนฺตรา วา. ตตฺราปิ โว ภิกฺขเว เอวํ สิกฺขิตพฺพํ “น เจว โน จิตฺตํ วิปริณตํ ภวิสฺสติ, น จ ปาปิกํ วาจํ นิจฺฉาเรสฺสาม, หิตานุกมฺปี จ วิหริสฺสาม เมตฺตจิตฺตา น โทสนฺตรา. ตญฺจ ปุคฺคลํ เมตฺตา\-สหคเตน เจตสา ผริตฺวา วิหริสฺสาม, ตทารมฺมณํ จ สพฺพาวนฺตํ โลกํ ปถวิสเมน เจตสา วิปุเลน มหคฺคเตน อปฺปมาเณน อเวเรน อพฺยาพชฺเฌน ผริตฺวา วิหริสฺสามา”ติ. เอวํ หิ โว ภิกฺขเว สิกฺขิตพฺพํ.}

\proseitem{229}{เสยฺยถาปิ ภิกฺขเว ปุริโส อาคจฺเฉยฺย ลาขํ วา หลิทฺทิํ วา นีลํ วา มญฺชิฏฺฐํ วา อาทาย, โส เอวํ วเทยฺย “อหํ อิมสฺมิํ อากาเส รูปํ ลิขิสฺสามิ, รูปปาตุภาวํ กริสฺสามี”ติ. ตํ กิํ มญฺญถ ภิกฺขเว, อปิ นุ โส ปุริโส อิมสฺมิํ อากาเส รูปํ ลิเขยฺย, รูปปาตุภาวํ กเรยฺยาติ. โน เหตํ ภนฺเต. ตํ กิสฺส เหตุ, อยํ หิ ภนฺเต อากาโส อรูปี อนิทสฺสโน, ตตฺถ น สุกรํ รูปํ ลิขิตุํ, รูปปาตุภาวํ กาตุํ. ยาวเทว จ ปน โส ปุริโส กิลมถสฺส วิฆาตสฺส ภาคี อสฺสาติ. เอวเมว โข ภิกฺขเว ปญฺจิเม วจนปถา, เยหิ โว ปเร วทมานา วเทยฺยุํ\csromandash{}กาเลน วา อกาเลน วา ฯเปฯ \csromanfolio{180}ตทารมฺมณํ จ สพฺพาวนฺตํ โลกํ อากาสสเมน เจตสา วิปุเลน มหคฺคเตน อปฺปมาเณน อเวเรน อพฺยาพชฺเฌน ผริตฺวา วิหริสฺสามา”ติ. เอวํ หิ โว ภิกฺขเว สิกฺขิตพฺพํ.}

\proseitem{230}{เสยฺยถาปิ ภิกฺขเว ปุริโส อาคจฺเฉยฺย อาทิตฺตํ ติณุกฺกํ อาทาย, โส เอวํ วเทยฺย “อหํ อิมาย อาทิตฺตาย ติณุกฺกาย คงฺคํ นทิํ สนฺตาเปสฺสามิ สํปริตาเปสฺสามี”ติ. ตํ กิํ มญฺญถ ภิกฺขเว, อปิ นุ โส ปุริโส อาทิตฺตาย ติณุกฺกาย คงฺคํ นทิํ สนฺตาเปยฺย สํปริตาเปยฺยาติ. โน เหตํ ภนฺเต. ตํ กิสฺส เหตุ, คงฺคา หิ ภนฺเต นที คมฺภีรา อปฺปเมยฺยา, สา น สุกรา อาทิตฺตาย ติณุกฺกาย สนฺตาเปตุํ สํปริตาเปตุํ. ยาวเทว จ ปน โส ปุริโส กิลมถสฺส วิฆาตสฺส ภาคี อสฺสาติ. เอวเมว โข ภิกฺขเว ปญฺจิเม วจนปถา, เยหิ โว ปเร วทมานา วเทยฺยุํ\csromandash{}กาเลน วา อกาเลน วา ฯเปฯ ตทารมฺมณํ จ สพฺพาวนฺตํ โลกํ คงฺคาสเมน เจตสา วิปุเลน มหคฺคเตน อปฺปมาเณน อเวเรน อพฺยาพชฺเฌน ผริตฺวา วิหริสฺสามา”ติ. เอวํ หิ โว ภิกฺขเว สิกฺขิตพฺพํ.}

\proseitem{231}{เสยฺยถาปิ ภิกฺขเว พิฬารภสฺตา มทฺทิตา สุมทฺทิตา สุปริมทฺทิตา มุทุกา ตูลินี ฉินฺนสสฺสรา ฉินฺน\-ภพฺภรา. อถ ปุริโส อาคจฺเฉยฺย กฏฺฐํ วา กถลํ\csromansharedfootnote{6da2d13966f4c009}{กฐลํ (สี, สฺยา, อิ)} วา อาทาย, โส เอวํ วเทยฺย “อหํ อิมํ พิฬารภสฺตํ มทฺทิตํ สุมทฺทิตํ สุปริมทฺทิตํ มุทุกํ ตูลินิํ ฉินฺนสสฺสรํ ฉินฺน\-ภพฺภรํ กฏฺเฐน วา กถเลน วา สรสรํ กริสฺสามิ ภรภรํ กริสฺสามี”ติ. ตํ กิํ มญฺญถ ภิกฺขเว, อปิ นุ โส ปุริโส อมุํ พิฬารภสฺตํ มทฺทิตํ สุมทฺทิตํ สุปริมทฺทิตํ มุทุกํ ตูลินิํ ฉินฺนสสฺสรํ ฉินฺน\-ภพฺภรํ กฏฺเฐน วา กถเลน วา สรสรํ กเรยฺย, ภรภรํ กเรยฺยาติ. โน เหตํ ภนฺเต. ตํ กิสฺส เหตุ, อมุ หิ ภนฺเต พิฬารภสฺตา มทฺทิตา สุมทฺทิตา สุปริมทฺทิตา มุทุกา ตูลินี ฉินฺนสสฺสรา ฉินฺน\-ภพฺภรา, สา น สุกรา กฏฺเฐน วา กถเลน วา สรสรํ กาตุํ, ภรภรํ กาตุํ. ยาวเทว จ ปน โส ปุริโส กิลมถสฺส วิฆาตสฺส ภาคี อสฺสาติ. เอวเมว โข ภิกฺขเว ปญฺจิเม วจนปถา, เยหิ โว ปเร วทมานา วเทยฺยุํ\csromandash{}กาเลน วา อกาเลน วา, ภูเตน วา อภูเตน วา,}

\prose{\csromanfolio{181}สเณฺหน วา ผรุเสน วา, อตฺถสํหิเตน วา อนตฺถ\-สํหิเตน วา, เมตฺตจิตฺตา วา โทสนฺตรา วา. กาเลน วา ภิกฺขเว ปเร วทมานา วเทยฺยุํ อกาเลน วา, ภูเตน วา ภิกฺขเว ปเร วทมานา วเทยฺยุํ อภูเตน วา, สเณฺหน วา ภิกฺขเว ปเร วทมานา วเทยฺยุํ ผรุเสน วา, อตฺถสํหิเตน วา ภิกฺขเว ปเร วทมานา วเทยฺยุํ อนตฺถ\-สํหิเตน วา, เมตฺตจิตฺตา วา ภิกฺขเว ปเร วทมานา วเทยฺยุํ โทสนฺตรา วา. ตตฺราปิ โว ภิกฺขเว เอวํ สิกฺขิตพฺพํ “น เจว โน จิตฺตํ วิปริณตํ ภวิสฺสติ, น จ ปาปิกํ วาจํ นิจฺฉาเรสฺสาม, หิตานุกมฺปี จ วิหริสฺสาม เมตฺตจิตฺตา น โทสนฺตรา. ตญฺจ ปุคฺคลํ เมตฺตา\-สหคเตน เจตสา ผริตฺวา วิหริสฺสาม, ตทารมฺมณํ จ สพฺพาวนฺตํ โลกํ พิฬารภสฺตา\-สเมน เจตสา วิปุเลน มหคฺคเตน อปฺปมาเณน อเวเรน อพฺยาพชฺเฌน ผริตฺวา วิหริสฺสามา”ติ. เอวํ หิ โว ภิกฺขเว สิกฺขิตพฺพํ.}

\proseitem{232}{อุภโต\-ทณฺฑ\-เกน เจปิ ภิกฺขเว กกเจน โจรา โอจรกา องฺคมงฺคานิ โอกนฺเตยฺยุํ, ตตฺราปิ โย มโน ปทูเสยฺย, น เม โส เตน สาสนกโร. ตตฺราปิ โว ภิกฺขเว เอวํ สิกฺขิตพฺพํ “น เจว โน จิตฺตํ วิปริณตํ ภวิสฺสติ, น จ ปาปิกํ วาจํ นิจฺฉาเรสฺสาม, หิตานุกมฺปี จ วิหริสฺสาม เมตฺตจิตฺตา น โทสนฺตรา. ตญฺจ ปุคฺคลํ เมตฺตา\-สหคเตน เจตสา ผริตฺวา วิหริสฺสาม, ตทารมฺมณํ จ สพฺพาวนฺตํ โลกํ เมตฺตา\-สหคเตน เจตสา วิปุเลน มหคฺคเตน อปฺปมาเณน อเวเรน อพฺยาพชฺเฌน ผริตฺวา วิหริสฺสามา”ติ. เอวํ หิ โว ภิกฺขเว สิกฺขิตพฺพํ.}

\proseitem{233}{อิมญฺจ\csromansharedfootnote{4e5a611f43d7a9c9}{อิมญฺเจ (?)} ตุเมฺห ภิกฺขเว กกจูปมํ โอวาทํ อภิกฺขณํ มนสิ กเรยฺยาถ, ปสฺสถ โน ตุเมฺห ภิกฺขเว ตํ วจนปถํ อณุํ วา ถูลํ วา, ยํ ตุเมฺห นาธิวาเสยฺยาถาติ. โน เหตํ ภนฺเต. ตสฺมา ติห ภิกฺขเว อิมํ กกจูปมํ โอวาทํ อภิกฺขณํ มนสิ กโรถ. ตํ โว ภวิสฺสติ ทีฆรตฺตํ หิตาย สุขายาติ.}

\prosewithcloser{อิทมโวจ ภควา. อตฺตมนา เต ภิกฺขู ภควโต ภาสิตํ อภินนฺทุนฺติ.}{กกจูปม\-สุตฺตํ นิฏฺฐิตํ ปฐมํ.}
\csromanmatikaanchor{mtk.52}
\csromantocmarkref{subsection}{2. อลคทฺทูปมสุตฺต}{mtk.52}
\bookmark[dest={mtk.52},level=2]{2. อลคทฺทูปมสุตฺต}
\csromanheader{2}{2. อลคทฺทู\-ปม\-สุตฺต}
\proseitem{234}{\csromanfolio{182}เอวํ เม สุตํ\csromandash{}เอกํ สมยํ ภควา สาวตฺถิยํ วิหรติ เชตวเน อนาถ\-ปิณฺฑิกสฺส อาราเม. เตน โข ปน สมเยน อริฏฺฐสฺส นาม ภิกฺขุโน คทฺธพาธิ\-ปุพฺพสฺส\csromansharedfootnote{ce519d550967bef5}{คนฺธพาธิปุพฺพสฺส (ก)} เอวรูปํ ปาปกํ ทิฏฺฐิคตํ อุปฺปนฺนํ โหติ “ตถาหํ ภควตา ธมฺมํ เทสิตํ อาชานามิ, ยถา เยเม อนฺตรายิกา ธมฺมา วุตฺตา ภควตา, เต ปฏิเสวโต นาลํ อนฺตรายายา”ติ. อสฺโสสุํ โข สมฺพหุลา ภิกฺขู อริฏฺฐสฺส กิร นาม ภิกฺขุโน คทฺธพาธิ\-ปุพฺพสฺส เอวรูปํ ปาปกํ ทิฏฺฐิคตํ อุปฺปนฺนํ “ตถาหํ ภควตา ธมฺมํ เทสิตํ อาชานามิ, ยถา เยเม อนฺตรายิกา ธมฺมา วุตฺตา ภควตา, เต ปฏิเสวโต นาลํ อนฺตรายายา”ติ. อถ โข เต ภิกฺขู เยน อริฏฺโฐ ภิกฺขุ คทฺธพาธิ\-ปุพฺโพ เตนุ\-ปสงฺกมิํสุ, อุปสงฺกมิตฺวา อริฏฺฐํ ภิกฺขุํ คทฺธพาธิ\-ปุพฺพํ เอตทโวจุํ “สจฺจํ กิร เต อาวุโส อริฏฺฐ เอวรูปํ ปาปกํ ทิฏฺฐิคตํ อุปฺปนฺนํ, ‘ตถาหํ ภควตา ธมฺมํ เทสิตํ อาชานามิ, ยถา เยเม อนฺตรายิกา ธมฺมา วุตฺตา ภควตา, เต ปฏิเสวโต นาลํ อนฺตรายายา’ติ”. เอวํพฺยาโข\csromansharedfootnote{b5f30e44930952cc}{เอวํ โข (?) ภควโต สมฺมุขาเยวสฺส “เอวํพฺยาโข”ติ.} อหํ อาวุโส ภควตา ธมฺมํ เทสิตํ อาชานามิ, ยถา เยเม อนฺตรายิกา ธมฺมา วุตฺตา ภควตา, เต ปฏิเสวโต นาลํ อนฺตรายายาติ. อถ โข เตปิ ภิกฺขู อริฏฺฐํ ภิกฺขุํ คทฺธพาธิ\-ปุพฺพํ เอตสฺมา ปาปกา ทิฏฺฐิคตา วิเวเจตุกามา สมนุยุญฺชนฺติ สมนุคาหนฺติ\csromansharedfootnote{ef10e6cdd2d87b32}{สมนุคฺคาหนฺติ (สฺยา)} สมนุภาสนฺติ “มา เหวํ อาวุโส อริฏฺฐ อวจ, มา ภควนฺตํ อพฺภาจิกฺขิ, น หิ สาธุ ภควโต อพฺภกฺขานํ\csromansharedfootnote{792e546ef2a1bad2}{อพฺภาจิกฺขนํ (ก)}, น หิ ภควา เอวํ วเทยฺย, อเนกปริยาเยนา\-วุโส อริฏฺฐ อนฺตรายิกา ธมฺมา อนฺตรายิกา วุตฺตา ภควตา, อลญฺจ ปน เต ปฏิเสวโต อนฺตรายาย. อปฺปสฺสาทา กามา วุตฺตา ภควตา พหุทุกฺขา พหุปายาสา, อาทีนโว เอตฺถ ภิยฺโย. อฏฺฐิ\-กงฺกลู\-ปมา กามา วุตฺตา ภควตา. มํสเปสูปมา กามา วุตฺตา ภควตา. ติณุกฺกูปมา กามา วุตฺตา ภควตา. องฺคารกาสูปมา กามา วุตฺตา ภควตา. สุปินกูปมา กามา วุตฺตา ภควตา. ยาจิตกูปมา กามา วุตฺตา ภควตา. รุกฺข\-ผลู\-ปมา กามา วุตฺตา ภควตา. อสิสูนูปมา กามา วุตฺตา ภควตา. สตฺติสูลูปมา กามา วุตฺตา ภควตา. สปฺปสิรูปมา กามา \csromanfolio{183}วุตฺตา ภควตา พหุทุกฺขา พหุปายาสา, อาทีนโว เอตฺถ ภิยฺโย”ติ. เอวํปิ โข อริฏฺโฐ ภิกฺขุ คทฺธพาธิ\-ปุพฺโพ เตหิ ภิกฺขูหิ สมนุ\-ยุญฺชิย\-มาโน สมนุ\-คาหิย\-มาโน\csromansharedfootnote{ca645fbf2b9ec61e}{สมนุคฺคาหิยมาโน (สฺยา, วินเยปิ)} สมนุ\-ภาสิย\-มาโน ตเทว\csromansharedfootnote{234d551a18128427}{ตเถว ตํ (วินเย)} ปาปกํ ทิฏฺฐิคตํ ถามสา ปรามาสา อภินิวิสฺส โวหรติ “เอวํพฺยาโข อหํ อาวุโส ภควตา ธมฺมํ เทสิตํ อาชานามิ, ยถา เยเม อนฺตรายิกา ธมฺมา วุตฺตา ภควตา, เต ปฏิเสวโต นาลํ อนฺตรายายา”ติ.}

\proseitem{235}{ยโต โข เต ภิกฺขุ นาสกฺขิํสุ อริฏฺฐํ ภิกฺขุํ คทฺธพาธิ\-ปุพฺพํ เอตสฺมา ปาปกา ทิฏฺฐิคตา วิเวเจตุํ. อถ โข เต ภิกฺขู เยน ภควา เตนุ\-ปสงฺกมิํสุ, อุปสงฺกมิตฺวา ภควนฺตํ อภิวาเทตฺวา เอกมนฺตํ นิสีทิํสุ, เอกมนฺตํ นิสินฺนา โข เต ภิกฺขู ภควนฺตํ เอตทโวจุํ\csromandash{}อริฏฺฐสฺส นาม ภนฺเต ภิกฺขุโน คทฺธพาธิ\-ปุพฺพสฺส เอวรูปํ ปาปกํ ทิฏฺฐิคตํ อุปฺปนฺนํ “ตถาหํ ภควตา ธมฺมํ เทสิตํ อาชานามิ, ยถา เยเม อนฺตรายิกา ธมฺมา วุตฺตา ภควตา, เต ปฏิเสวโต นาลํ อนฺตรายายา”ติ. อสฺสุมฺห โข มยํ ภนฺเต อริฏฺฐสฺส กิร นาม ภิกฺขุโน คทฺธพาธิ\-ปุพฺพสฺส เอวรูปํ ปาปกํ ทิฏฺฐิคตํ อุปฺปนฺนํ “ตถาหํ ภควตา ธมฺมํ เทสิตํ อาชานามิ, ยถา เยเม อนฺตรายิกา ธมฺมา วุตฺตา ภควตา, เต ปฏิเสวโต นาลํ อนฺตรายายา”ติ. อถ โข มยํ ภนฺเต เยน อริฏฺโฐ ภิกฺขุ คทฺธพาธิ\-ปุพฺโพ เตนุ\-ปสงฺกมิมฺห, อุปสงฺกมิตฺวา อริฏฺฐํ ภิกฺขุํ คทฺธพาธิ\-ปุพฺพํ เอตทโวจุมฺห “สจฺจํ กิร เต อาวุโส อริฏฺฐ เอวรูปํ ปาปกํ ทิฏฺฐิคตํ อุปฺปนฺนํ, ‘ตถาหํ ภควตา ธมฺมํ เทสิตํ อาชานามิ, ยถา เยเม อนฺตรายิกา ธมฺมา วุตฺตา ภควตา, เต ปฏิเสวโต นาลํ อนฺตรายายา’ติ”. เอวํ วุตฺเต ภนฺเต อริฏฺโฐ ภิกฺขุ คทฺธพาธิ\-ปุพฺโพ อเมฺห เอตทโวจ “เอวํพฺยาโข อหํ อาวุโส ภควตา ธมฺมํ เทสิตํ อาชานามิ, ยถา เยเม อนฺตรายิกา ธมฺมา วุตฺตา ภควตา, เต ปฏิเสวโต นาลํ อนฺตรายายา”ติ. อถ โข มยํ ภนฺเต อริฏฺฐํ ภิกฺขุํ คทฺธพาธิ\-ปุพฺพํ เอตสฺมา ปาปกา ทิฏฺฐิคตา วิเวเจตุกามา สมนุยุญฺชิมฺห สมนุคาหิมฺห สมนุภาสิมฺห “มา เหวํ อาวุโส อริฏฺฐ อวจ, มา ภควนฺตํ อพฺภาจิกฺขิ, น หิ สาธุ ภควโต \csromanfolio{184}อพฺภกฺขานํ, น หิ ภควา เอวํ วเทยฺย, อเนกปริยาเยนา\-วุโส อริฏฺฐ อนฺตรายิกา ธมฺมา อนฺตรายิกา วุตฺตา ภควตา, อลญฺจ ปน เต ปฏิเสวโต อนฺตรายาย. อปฺปสฺสาทา กามา วุตฺตา ภควตา พหุทุกฺขา พหุปายาสา, อาทีนโว เอตฺถภิยฺโย. อฏฺฐิ\-กงฺกลู\-ปมา กามา วุตฺตา ภควตา ฯเปฯ. สปฺปสิรูปมา กามา วุตฺตา ภควตา พหุทุกฺขา, พหุปายาสา อาทีนโว เอตฺถ ภิยฺโย”ติ. เอวํปิ โข ภนฺเต อริฏฺโฐ ภิกฺขุ คทฺธพาธิ\-ปุพฺโพ อเมฺหหิ สมนุ\-ยุญฺชิย\-มาโน สมนุ\-คาหิย\-มาโน สมนุ\-ภาสิย\-มาโน ตเทว ปาปกํ ทิฏฺฐิคตํ ถามสา ปรามาสา อภินิวิสฺส โวหรติ “เอวํพฺยาโข อหํ อาวุโส ภควตา ธมฺมํ เทสิตํ อาชานามิ, ยถา เยเม อนฺตรายิกา ธมฺมา วุตฺตา ภควตา, เต ปฏิเสวโต นาลํ อนฺตรายายา”ติ. ยโต โข มยํ ภนฺเต นาสกฺขิมฺห อริฏฺฐํ ภิกฺขุํ คทฺธพาธิ\-ปุพฺพํ เอตสฺมา ปาปกา ทิฏฺฐิคตา วิเวเจตุํ, อถ มยํ เอตมตฺถํ ภควโต อาโรเจมาติ.}

\proseitem{236}{อถ โข ภควา อญฺญตรํ ภิกฺขุํ อามนฺเตสิ “เอหิ ตฺวํ ภิกฺขุ มม วจเนน อริฏฺฐํ ภิกฺขุํ คทฺธพาธิ\-ปุพฺพํ อามนฺเตหิ ‘สตฺถา ตํ อาวุโส อริฏฺฐ อามนฺเตตี’ติ”. “เอวํ ภนฺเต”ติ โข โส ภิกฺขุ ภควโต ปฏิสฺสุตฺวา เยน อริฏฺโฐ ภิกฺขุ คทฺธพาธิ\-ปุพฺโพ เตนุปสงฺกมิ, อุปสงฺกมิตฺวา อริฏฺฐํ ภิกฺขุํ คทฺธพาธิ\-ปุพฺพํ เอตทโวจ “สตฺถา ตํ อาวุโส อริฏฺฐ อามนฺเตตี”ติ. “เอวมาวุโส”ติ โข อริฏฺโฐ ภิกฺขุ คทฺธพาธิ\-ปุพฺโพ ตสฺส ภิกฺขุโน ปฏิสฺสุตฺวา เยน ภควา เตนุปสงฺกมิ, อุปสงฺกมิตฺวา ภควนฺตํ อภิวาเทตฺวา เอกมนฺตํ นิสีทิ, เอกมนฺตํ นิสินฺนํ โข อริฏฺฐํ ภิกฺขุํ คทฺธพาธิ\-ปุพฺพํ ภควา เอตทโวจ “สจฺจํ กิร เต อริฏฺฐ เอวรูปํ ปาปกํ ทิฏฺฐิคตํ อุปฺปนฺนํ ‘ตถาหํ ภควตา ธมฺมํ เทสิตํ อาชานามิ, ยถา เยเม อนฺตรายิกา ธมฺมา วุตฺตา ภควตา, เต ปฏิเสวโต นาลํ อนฺตรายายา’ติ”. เอวํพฺยาโข อหํ ภนฺเต ภควตา ธมฺมํ เทสิตํ อาชานามิ, ยถา เยเม อนฺตรายิกา ธมฺมา วุตฺตา ภควตา, เต ปฏิเสวโต นาลํ อนฺตรายายาติ. กสฺส โข นาม ตฺวํ โมฆปุริส มยา เอวํ ธมฺมํ เทสิตํ อาชานาสิ, นนุ มยา โมฆปุริส อเนก\-ปริยาเยน อนฺตรายิกา ธมฺมา อนฺตรายิกา วุตฺตา, อลญฺจ ปน เต ปฏิเสวโต อนฺตรายาย. อปฺปสฺสาทา กามา วุตฺตา มยา พหุทุกฺขา พหุปายาสา, อาทีนโว เอตฺถ ภิยฺโย. อฏฺฐิ\-กงฺกลู\-ปมา กามา วุตฺตา มยา. \csromanfolio{185}มํสเปสูปมา กามา วุตฺตา มยา. ติณุกฺกูปมา กามา วุตฺตา มยา. องฺคารกาสูปมา กามา วุตฺตา มยา. สุปินกูปมา กามา วุตฺตา มยา. ยาจิตกูปมา กามา วุตฺตา มยา. รุกฺข\-ผลู\-ปมา กามา วุตฺตา มยา. อสิสูนุปมา กามา วุตฺตา มยา. สตฺติสูลูปมา กามา วุตฺตา มยา. สปฺปสิรูปมา กามา วุตฺตา มยา พหุทุกฺขา พหุปายาสา, อาทีนโว เอตฺถ ภิยฺโย. อถ จ ปน ตฺวํ โมฆปุริส อตฺตนา ทุคฺคหิเตน อเมฺห เจว อพฺภาจิกฺขสิ, อตฺตานญฺจ ขนสิ, พหุญฺจ อปุญฺญํ ปสวสิ. ตํ หิ เต โมฆปุริส ภวิสฺสติ ทีฆรตฺตํ อหิตาย ทุกฺขายาติ.}

\prose{อถ โข ภควา ภิกฺขู อามนฺเตสิ. ตํ กิํ มญฺญถ ภิกฺขเว, อปิ นายํ อริฏฺโฐ ภิกฺขุ คทฺธพาธิ\-ปุพฺโพ อุสฺมีกโตปิ อิมสฺมิํ ธมฺมวินเยติ. กิํ หิ\csromansharedfootnote{beb302dd60e9d5ef}{กิํติ (ก)} สิยา ภนฺเต, โน เหตํ ภนฺเตติ. เอวํ วุตฺเต อริฏฺโฐ ภิกฺขุ คทฺธพาธิ\-ปุพฺโพ ตุณฺหีภูโต มงฺกุภูโต ปตฺตกฺขนฺโธ อโธมุโข ปชฺฌายนฺโต อปฺปฏิภาโน นิสีทิ. อถ โข ภควา อริฏฺฐํ ภิกฺขุํ คทฺธพาธิ\-ปุพฺพํ ตุณฺหีภูตํ มงฺกุภูตํ ปตฺตกฺขนฺธํ อโธมุขํ ปชฺฌายนฺตํ อปฺปฏิภานํ วิทิตฺวา อริฏฺฐํ ภิกฺขุํ คทฺธพาธิ\-ปุพฺพํ เอตทโวจ “ปญฺญายิสฺสสิ โข ตฺวํ โมฆปุริส เอเตน สเกน ปาปเกน ทิฏฺฐิคเตน, อิธาหํ ภิกฺขู ปฏิปุจฺฉิสฺสามี”ติ.}

\proseitem{237}{อถ โข ภควา ภิกฺขู อามนฺเตสิ “ตุเมฺหปิ เม ภิกฺขเว เอวํ ธมฺมํ เทสิตํ อาชานาถ, ยถายํ อริฏฺโฐ ภิกฺขุ คทฺธพาธิ\-ปุพฺโพ อตฺตนา ทุคฺคหิเตน อเมฺห เจว อพฺภาจิกฺขติ, อตฺตานญฺจ ขนติ, พหุญฺจ อปุญฺญํ ปสวตี”ติ. โนเหตํ ภนฺเต. อเนก\-ปริยาเยน หิ โน ภนฺเต อนฺตรายิกา ธมฺมา อนฺตรายิกา วุตฺตา ภควตา, อลญฺจ ปน เต ปฏิเสวโต อนฺตรายาย. อปฺปสฺสาทา กามา วุตฺตา ภควตา พหุทุกฺขา พหุปายาสา, อาทีนโว เอตฺถ ภิยฺโย. อฏฺฐิ\-กงฺกลู\-ปมา กามา วุตฺตา ภควตา ฯเปฯ. สปฺปสิรูปมา กามา วุตฺตา ภควตา พหุทุกฺขา พหุปายาสา, อาทีนโว เอตฺถ ภิยฺโยติ. สาธุ สาธุ ภิกฺขเว, สาธุ โข เม ตุเมฺห ภิกฺขเว เอวํ ธมฺมํ เทสิตํ อาชานาถ. อเนก\-ปริยาเยน หิ โข ภิกฺขเว อนฺตรายิกา ธมฺมา วุตฺตา มยา, อลญฺจ ปน เต ปฏิเสวโต อนฺตรายาย. อปฺปสฺสาทา กามา วุตฺตา \csromanfolio{186}มยา พหุทุกฺขา พหุปายาสา, อาทีนโว เอตฺถ ภิยฺโย. อฏฺฐิ\-กงฺกลู\-ปมา กามา วุตฺตา มยา ฯเปฯ. สปฺปสิรูปมา กามา วุตฺตา มยา พหุทุกฺขา พหุปายาสา, อาทีนโว เอตฺถ ภิยฺโย. อถ จ ปนายํ อริฏฺโฐ ภิกฺขุ คทฺธพาธิ\-ปุพฺโพ อตฺตนา ทุคฺคหิเตน อเมฺห เจว อพฺภาจิกฺขติ, อตฺตานญฺจ ขนติ, พหุญฺจ อปุญฺญํ ปสวติ. ตํ หิ ตสฺส โมฆปุริสสฺส ภวิสฺสติ ทีฆรตฺตํ อหิตาย ทุกฺขาย. โส วต ภิกฺขเว อญฺญเตฺรว กาเมหิ อญฺญตฺร กามสญฺญาย อญฺญตฺร กามวิตกฺเกหิ กาเม ปฏิเสวิสฺสตีติ เนตํ ฐานํ วิชฺชติ.}

\proseitem{238}{อิธ ภิกฺขเว เอกจฺเจ โมฆปุริสา ธมฺมํ ปริยาปุณนฺติ สุตฺตํ เคยฺยํ เวยฺยากรณํ คาถํ อุทานํ อิติวุตฺตกํ ชาตกํ อพฺภุตธมฺมํ เวทลฺลํ. เต ตํ ธมฺมํ ปริยาปุณิตฺวา เตสํ ธมฺมานํ ปญฺญาย อตฺถํ น อุปปริกฺขนฺติ, เตสํ เต ธมฺมา ปญฺญาย อตฺถํ อนุปปริกฺขตํ น นิชฺฌานํ ขมนฺติ, เต อุปารมฺภา\-นิสํสา เจว ธมฺมํ ปริยาปุณนฺติ อิติวาท\-ปฺปโมกฺขา\-นิสํสา จ, ยสฺส จตฺถาย ธมฺมํ ปริยาปุณนฺติ, ตญฺจสฺส อตฺถํ นานุโภนฺติ, เตสํ เต ธมฺมา ทุคฺคหิตา ทีฆรตฺตํ อหิตาย ทุกฺขาย สํวตฺตนฺติ. ตํ กิสฺส เหตุ, ทุคฺคหิตตฺตา ภิกฺขเว ธมฺมานํ.}

\prose{เสยฺยถาปิ ภิกฺขเว ปุริโส อลคทฺทตฺถิโก อลคทฺทคเวสี อลคทฺทปริเยสนํ จรมาโน, โส ปสฺสยฺย มหนฺตํ อลคทฺทํ, ตเมนํ โภเค วา นงฺคุฏฺเฐ วา คเณฺหยฺย, ตสฺส โส อลคทฺโท ปฏิปริวตฺติตฺวา\csromansharedfootnote{b0c643fe94b4fa6f}{ปฏินิวตฺติตฺวา (สฺยา, ก)} หตฺเถ วา พาหาย วา อญฺญตรสฺมิํ วา องฺคปจฺจงฺเค ฑํเสยฺย\csromansharedfootnote{a688fb25f9dc656f}{qอเสยฺย (สี, อิ)}, โส ตโตนิทานํ มรณํ วา นิคจฺเฉยฺย มรณมตฺตํ วา ทุกฺขํ. ตํ กิสฺส เหตุ, ทุคฺคหิตตฺตา ภิกฺขเว อลคทฺทสฺส. เอวเมว โข ภิกฺขเว อิเธกจฺเจ โมฆปุริสา ธมฺมํ ปริยาปุณนฺติ สุตฺตํ เคยฺยํ เวยฺยากรณํ คาถํ อุทานํ อิติวุตฺตกํ ชาตกํ อพฺภุตธมฺมํ เวทลฺลํ. เต ตํ ธมฺมํ ปริยาปุณิตฺวา เตสํ ธมฺมานํ ปญฺญาย อตฺถํ น อุปปริกฺขนฺติ, เตสํ เต ธมฺมา ปญฺญาย อตฺถํ อนุปปริกฺขตํ น นิชฺฌานํ ขมนฺติ, เต อุปารมฺภา\-นิสํสา เจว ธมฺมํ ปริยาปุณนฺติ อิติวาท\-ปฺปโมกฺขา\-นิสํสา จ, ยสฺส จตฺถาย ธมฺมํ ปริยาปุณนฺติ, ตญฺจสฺส อตฺถํ นานุโภนฺติ, เตสํ เต ธมฺมา ทุคฺคหิตา ทีฆรตฺตํ อหิตาย ทุกฺขาย สํวตฺตนฺติ. ตํ กิสฺส เหตุ, ทุคฺคหิตตฺตา ภิกฺขเว ธมฺมานํ.}

\proseitem{239}{\csromanfolio{187}อิธ ปน ภิกฺขเว เอกจฺเจ กุลปุตฺตา ธมฺมํ ปริยาปุณนฺติ สุตฺตํ เคยฺยํ เวยฺยากรณํ คาถํ อุทานํ อิติวุตฺตกํ ชาตกํ อพฺภุตธมฺมํ เวทลฺลํ. เต ตํ ธมฺมํ ปริยาปุณิตฺวา เตสํ ธมฺมานํ ปญฺญาย อตฺถํ อุปปริกฺขนฺติ, เตสํ เต ธมฺมา ปญฺญาย อตฺถํ อุปปริกฺขตํ นิชฺฌานํ ขมนฺติ, เต น เจว อุปารมฺภา\-นิสํสา ธมฺมํ ปริยาปุณนฺติ, น อิติวาท\-ปฺปโมกฺขา\-นิสํสา จ\csromansharedfootnote{e84d8f18c8f6386e}{น จ อิติวาทปฺปโมกฺขานิสํสา (?)}, ยสฺส จตฺถาย ธมฺมํ ปริยาปุณนฺติ, ตญฺจสฺส อตฺถํ อนุโภนฺติ, เตสํ เต ธมฺมา สุคฺคหิตา ทีฆรตฺตํ หิตาย สุขาย สํวตฺตนฺติ. ตํ กิสฺส เหตุ, สุคฺคหิตตฺตา ภิกฺขเว ธมฺมานํ.}

\prose{เสยฺยถาปิ ภิกฺขเว ปุริโส อลคทฺทตฺถิโก อลคทฺทคเวสี อลคทฺทปริเยสนํ จรมาโน, โส ปสฺเสยฺย มหนฺตํ อลคทฺทํ, ตเมนํ อชปเทน ทณฺเฑน สุนิคฺคหิตํ นิคฺคเณฺหยฺย, อชปเทน ทณฺเฑน สุนิคฺคหิตํ นิคฺคหิตฺวา คีวาย สุคฺคหิตํ คเณฺหยฺย. กิญฺจาปิ โส ภิกฺขเว อลคทฺโท ตสฺส ปุริสสฺส หตฺตํ วา พาหํ วา อญฺญตรํ วา องฺคปจฺจงฺคํ โภเคหิ ปลิเวเฐยฺย. อถ โข โส เนว ตโตนิทานํ มรณํ วา นิคจฺเฉยฺย มรณมตฺตํ วา ทุกฺขํ. ตํ กิสฺส เหตุ สุคฺคหิตตฺตา ภิกฺขเว อลคทฺทสฺส. เอวเมว โข ภิกฺขเว อิเธกจฺเจ กุลปุตฺตา ธมฺมํ ปริยาปุณนฺติ สุตฺตํ เคยฺยํ เวยฺยากรณํ คาถํ อุทานํ อิติวุตฺตกํ ชาตกํ อพฺภุตธมฺมํ เวทลฺลํ, เต ตํ ธมฺมํ ปริยาปุณิตฺวา เตสํ ธมฺมานํ ปญฺญาย อตฺถํ อุปปริกฺขนฺติ, เตสํ เต ธมฺมา ปญฺญาย อตฺถํ อุปปริกฺขตํ นิชฺฌานํ ขมนฺติ, เต น เจว อุปารมฺภา\-นิสํสา ธมฺมํ ปริยาปุณนฺติ, น อิติวาท\-ปฺปโมกฺขา\-นิสํสา จ, ยสฺส จตฺถาย ธมฺมํ ปริยาปุณนฺติ, ตญฺจสฺส อตฺถํ อนุโภนฺติ. เตสํ เต ธมฺมา สุคฺคหิตา ทีฆรตฺตํ อตฺถาย หิตาย สุขาย สํวตฺตนฺติ. ตํ กิสฺส เหตุ, สุคฺคหิตตฺตา ภิกฺขเว ธมฺมานํ. ตสฺมาติห ภิกฺขเว ยสฺส เม ภาสิตสฺส อตฺถํ อาชาเนยฺยาถ, ตถา นํ ธาเรยฺยาถ. ยสฺส จ ปน เม ภาสิตสฺส อตฺถํ น อาชาเนยฺยาถ, อหํ โว ตตฺถ ปฏิปุจฺฉิตพฺโพ. เย วา ปนาสฺสุ วิยตฺตา ภิกฺขู.}

\proseitem{240}{กุลฺลูปมํ โว ภิกฺขเว ธมฺมํ เทเสสฺสามิ นิตฺถรณ\-ตฺถาย, โน คหณตฺถาย. ตํ สุณาถ สาธุกํ มนสิ กโรถ ภาสิสฺสามีติ. “เอวํ ภนฺเต”ติ โข เต ภิกฺขู ภควโต ปจฺจสฺโสสุํ. ภควา เอตทโวจ.}

\prose{\csromanfolio{188}เสยฺยถาปิ ภิกฺขเว ปุริโส อทฺธาน\-มคฺค\-ปฺปฏิปนฺโน, โส ปสฺเสยฺย มหนฺตํ อุทกณฺณวํ โอริมํ ตีรํ สาสงฺกํ สปฺปฏิภยํ, ปาริมํ ตีรํ เขมํ อปฺปฏิภยํ. น จสฺส นาวา สนฺตารณี, อุตฺตรเสตุ วา อปารา ปารํ คมนาย. ตสฺส เอวมสฺส “อยํ โข มหาอุทกณฺณโว โอริมํ ตีรํ สาสงฺกํ สปฺปฏิภยํ, ปาริมํ ตีรํ เขมํ อปฺปฏิภยํ. นตฺถิ จ นาวา สนฺตารณี, อุตฺตรเสตุ วา อปารา ปารํ คมนาย, ยํนูนาหํ ติณกฏฺฐ\-สาขา\-ปลาสํ สํกฑฺฒิตฺวา กุลฺลํ พนฺธิตฺวา ตํ กุลฺลํ นิสฺสาย หตฺเถหิ จ ปาเทหิ จ วายมมาโน โสตฺถินา ปารํ อุตฺตเรยฺยนฺ”ติ. อถ โข โส ภิกฺขเว ปุริโส ติณกฏฺฐ\-สาขา\-ปลาสํ สํกฑฺฒิตฺวา กุลฺลํ พนฺธิตฺวา ตํ กุลฺลํ นิสฺสาย หตฺเถหิ จ ปาเทหิ จ วายมมาโน โสตฺถินา ปรํ อุตฺตเรยฺย. ตสฺส ปุริสสฺส อุตฺติณฺณสฺส\csromansharedfootnote{c03ba782c61c3f1e}{ติณฺณสฺส (อิ, ก)} ปารงฺคตสฺส เอวมสฺส “พหุกาโร โข เม อยํ กุลฺโล, อิมาหํ กุลฺลํ นิสฺสาย หตฺเถหิ จ ปาเทหิ จ วายมมาโน โสตฺถินา ปารํ อุตฺติณฺโณ, ยํนูนาหํ อิมํ กุลฺลํ สีเส วา อาโรเปตฺวา ขนฺเธ วา อุจฺจาเรตฺวา\csromansharedfootnote{b8175d17bbd26851}{อุจฺโจเปตฺวา (ก)} เยน กามํ ปกฺกเมยฺยนฺ”ติ. ตํ กิํมญฺญถ ภิกฺขเว, อปิ นุ โส ปุริโส เอวํการี ตสฺมิํ กุลฺเล กิจฺจการี อสฺสาติ. โนเหตํ ภนฺเต. กถํการี จ โส ภิกฺขเว ปุริโส ตสฺมิํ กุลฺเล กิจฺจการี อสฺส. อิธ ภิกฺขเว ตสฺส ปุริสสฺส อุตฺติณฺณสฺส\csromansharedfootnote{d79208a91ffd77c2}{อุสฺสาเรตฺวา (ก)} ปารงฺคตสฺส เอวมสฺส, พหุกาโร โข เม อยํ กุลฺโล, อิมาหํ กุลฺลํ นิสฺสาย หตฺเถหิ จ ปาเทหิ จ วายมมาโน โสตฺถินา ปารํ อุตฺติณฺโณ, ยํนูนาหํ อิมํ กุลฺลํ ถเล วา อุสฺสาเทตฺวา\csromansharedfootnote{c03ba782c61c3f1e}{ติณฺณสฺส (อิ, ก)} อุทเก วา โอปิลาเปตฺวา เยน กามํ ปกฺกเมยฺยนฺติ. เอวํการี โข โส ภิกฺขเว ปุริโส ตสฺมิํ กุลฺเล กิจฺจการี อสฺส. เอวเมว โข ภิกฺขเว กุลฺลูปโม มยา ธมฺโม เทสิโต นิตฺถรณ\-ตฺถาย, โน คหณตฺถาย. กุลฺลูปมํ โว ภิกฺขเว ธมฺมํ เทสิตํ อาชานนฺเตหิ ธมฺมาปิ โว ปหาตพฺพา, ปเคว อธมฺมา.}

\proseitem{241}{ฉยิมานิ ภิกฺขเว ทิฏฺฐิฏฺฐานานิ. กตมานิ ฉ. อิธ ภิกฺขเว อสฺสุตวา ปุถุชฺชโน อริยานํ อทสฺสาวี อริยธมฺมสฺส อโกวิโท อริยธมฺเม อวินีโต สปฺปุริสานํ อทสฺสาวี สปฺปุริส\-ธมฺมสฺส อโกวิโท สปฺปุริส\-ธมฺเม อวินีโต รูปํ “เอตํ มม, เอโสหมสฺมิ, เอโส เม \csromanfolio{189}อตฺตา”ติ สมนุปสฺสติ. เวทนํ “เอตํ มม, เอโสหมสฺมิ, เอโส เม อตฺตา”ติ สมนุปสฺสติ. สญฺญํ “เอตํ มม เอโสหมสฺมิ, เอโส เม อตฺตา”ติ สมนุปสฺสติ. สงฺขาเร “เอตํ มม, เอโสหมสฺมิ, เอโส เม อตฺตา”ติ สมนุปสฺสติ. ยํปิ ตํ ทิฏฺฐํ สุตํ มุตํ วิญฺญาตํ ปตฺตํ ปริเยสิตํ อนุวิจริตํ มนสา, ตํปิ “เอตํ มม, เอโสหมสฺมิ, เอโส เม อตฺตา”ติ สมนุปสฺสติ. ยํปิ ตํ ทิฏฺฐิฏฺฐานํ ‘โส โลโก, โส อตฺตา, โส เปจฺจ ภวิสฺสามิ นิจฺโจ ธุโว สสฺสโต อวิปริณาม\-ธมฺโม, สสฺสติสมํ ตเถว ฐสฺสามี’ติ, ตํปิ “เอตํ มม, เอโสหมสฺมิ, เอโส เม อตฺตา”ติ สมนุปสฺสติ. สุตวา จ โข ภิกฺขเว อริยสาวโก อริยานํ ทสฺสาวี อริยธมฺมสฺส โกวิโท อริยธมฺเม สุวินีโต สปฺปุริสานํ ทสฺสาวี สปฺปุริส\-ธมฺมสฺส โกวิโท สปฺปุริส\-ธมฺเม สุวินีโต รูปํ “เนตํ มม, เนโสหมสฺมิ, น เมโส อตฺตา”ติ สมนุปสฺสติ. เวทนํ “เนตํ มม, เนโสหมสฺมิ, น เมโส อตฺตา”ติ สมนุปสฺสติ. สญฺญํ “เนตํ มม, เนโสหมสฺมิ, น เมโส อตฺตา”ติ สมนุปสฺสติ. สงฺขาเร “เนตํ มม, เนโสหมสฺมิ, น เมโส อตฺตา”ติ สมนุปสฺสติ. ยํปิ ตํ ทิฏฺฐํ สุตํ มุตํ วิญฺญาตํ ปตฺตํ ปริเยสิตํ อนุวิจริตํ มนสา, ตํปิ “เนตํ มม, เนโสหมสฺมิ, น เมโส อตฺตา”ติ สมนุปสฺสติ. ยํปิ ตํ ทิฏฺฐิฏฺฐานํ ‘โส โลโก, โส อตฺตา, โส เปจฺจ ภวิสฺสามิ นิจฺโจ ธุโว สสฺสโต อวิปริณาม\-ธมฺโม, สสฺสติสมํ ตเถว ฐสฺสามี’ติ, ตํปิ “เนตํ มม, เนโสหมสฺมิ, น เมโส อตฺตา”ติ สมนุปสฺสติ. โส เอวํ สมนุปสฺสนฺโต อสติ น ปริตสฺสตีติ.}

\proseitem{242}{เอวํ วุตฺเต อญฺญตโร ภิกฺขุ ภควนฺตํ เอตทโวจ “สิยา นุ โข ภนฺเต พหิทฺธา อสติ ปริตสฺสนา”ติ. สิยา ภิกฺขูติ ภควา อโวจ, อิธ ภิกฺขุ เอกจฺจสฺส เอวํ โหติ “อหุ วต เม, ตํ วต เม นตฺถิ, สิยา วต เม, ตํ วตาหํ น ลภามี”ติ. โส โสจติ กิลมติ ปริเทวติ อุรตฺตาฬิํ กนฺทติ สมฺโมหํ อาปชฺชติ. เอวํ โข ภิกฺขุ พหิทฺธา อสติ ปริตสฺสนา โหตีติ.}

\prose{สิยา ปน ภนฺเต พหิทฺธา อสติ อปริตสฺสนาติ. สิยา ภิกฺขูติ ภควา อโวจ, อิธ ภิกฺขุ เอกจฺจสฺส น เอวํ โหติ “อหุ วต เม, ตํ วต เม นตฺถิ, สิยา วต เม, ตํ วตาหํ น ลภามี”ติ. โส น \csromanfolio{190}โสจติ น กิลมติ น ปริเทวติ น อุรตฺตาฬิํ กนฺทติ น สมฺโมหํ อาปชฺชติ. เอวํ โข ภิกฺขุ พหิทฺธา อสติ อปริตสฺสนา โหตีติ.}

\prose{สิยา นุ โข ภนฺเต อชฺฌตฺตํ อสติ ปริตสฺสนาติ. สิยา ภิกฺขูติ ภควา อโวจ, อิธ ภิกฺขุ เอกจฺจสฺส เอวํ ทิฏฺฐิ โหติ, “โส โลโก, โส อตฺตา โส เปจฺจ ภวิสฺสามิ นิจฺโจ ธุโว สสฺสโต อวิปริณาม\-ธมฺโม, สสฺสติสมํ ตเถว ฐสฺสามี”ติ. โส สุณาติ ตถาคตสฺส วา ตถาคต\-สาวกสฺส วา สพฺเพสํ ทิฏฺฐิ\-ฏฺฐานา\-ธิฏฺฐาน\-ปริยุฏฺฐานา\-ภินิเวสา\-นุสยานํ สมุคฺฆาตาย สพฺพ\-สงฺขาร\-สมถาย สพฺพู\-ปธิ\-ปฏินิสฺสคฺคาย ตณฺหากฺขยาย วิราคาย นิโรธาย นิพฺพานาย ธมฺมํ เทเสนฺตสฺส, ตสฺส เอวํ โหติ “อุจฺฉิชฺชิสฺสามิ นามสฺสุ, วินสฺสิสฺสามิ นามสฺสุ, นสฺสุ นาม ภวิสฺสามี”ติ. โส โสจติ กิลมติ ปริเทวติ อุรตฺตาฬิํ กนฺทติ สมฺโมหํ อาปชฺชติ. เอวํ โข ภิกฺขุ อชฺฌตฺตํ อสติ ปริตสฺสนา โหตีติ.}

\prose{สิยา ปน ภนฺเต อชฺฌตฺตํ อสติ อปริตสฺสนาติ. สิยา ภิกฺขูติ ภควา อโวจ, อิธ ภิกฺขุ เอกจฺจสฺส น เอวํ ทิฏฺฐิ โหติ, “โส โลโก, โส อตฺตา, โส เปจฺจ ภวิสฺสามิ นิจฺโจ ธุโว สสฺสโต อวิปริณาม\-ธมฺโม, สสฺสติสมํ ตเถว ฐสฺสามี”ติ. โส สุณาติ ตถาคตสฺส วา ตถาคต\-สาวกสฺส วา สพฺเพสํ ทิฏฺฐิ\-ฏฺฐานา\-ธิฏฺฐาน\-ปริยุฏฺฐานา\-ภินิเวสา\-นุสยานํ สมุคฺฆาตาย สพฺพ\-สงฺขาร\-สมถาย สพฺพู\-ปธิ\-ปฏินิสฺสคฺคาย ตณฺหากฺขยาย วิราคาย นิโรธาย นิพฺพานาย ธมฺมํ เทเสนฺตสฺส, ตสฺส น เอวํ โหติ “อุจฺฉิชฺชิสฺสามิ นามสฺสุ, วินสฺสิสฺสามิ นามสฺสุ, นสฺสุ นาม ภวิสฺสามี”ติ. โส น โสจติ น กิลมติ น ปริเทวติ น อุรตฺตาฬิํ กนฺทติ น สมฺโมหํ อาปชฺชติ. เอวํ โข ภิกฺขุ อชฺฌตฺตํ อสติ อปริตสฺสนา โหติ.}

\proseitem{243}{ตํ\csromansharedfootnote{3ca69fbc8410b24c}{ตญฺจ (ก)} ภิกฺขเว ปริคฺคหํ ปริคฺคเณฺหยฺยาถ, ยฺวาสฺส\csromansharedfootnote{36dd81f870e1fd22}{ยฺวาสฺสุ (ก)} ปริคฺคโห นิจฺโจ ธุโว สสฺสโต อวิปริณาม\-ธมฺโม, สสฺสติสมํ ตเถว ติฏฺเฐยฺย. ปสฺสถ โน ตุเมฺห ภิกฺขเว ตํ ปริคฺคหํ, ยฺวาสฺส ปริคฺคโห นิจฺโจ ธุโว สสฺสโต อวิปริณาม\-ธมฺโม, สสฺสติสมํ ตเถว ติฏฺเฐยฺยาติ. โน เหตํ ภนฺเต. สาธุ ภิกฺขเว, อหํปิ โข ตํ ภิกฺขเว ปริคฺคหํ น \csromanfolio{191}สมนุปสฺสามิ, ยฺวาสฺส ปริคฺคโห นิจฺโจ ธุโว สสฺสโต อวิปริณาม\-ธมฺโม, สสฺสติสมํ ตเถว ติฏฺเฐยฺย.}

\prose{ตํ ภิกฺขเว อตฺตวาทุ\-ปาทานํ อุปาทิเยถ, ยํ’ส\csromansharedfootnote{b642d250ea1c3ca2}{ยสฺส (สฺยา, ก)} อตฺตวาทุ\-ปาทานํ อุปาทิยโต น อุปฺปชฺเชยฺยุํ โสก\-ปริเทว\-ทุกฺข\-โทมนสฺสุ\-ปายาสา. ปสฺสถ โน ตุเมฺห ภิกฺขเว ตํ อตฺตวาทุ\-ปาทานํ, ยํ’ส อตฺตวาทุ\-ปาทานํ อุปาทิยโต น อุปฺปชฺเชยฺยุํ โสก\-ปริเทว\-ทุกฺข\-โทมนสฺสุ\-ปายาสาติ. โน เหตํ ภนฺเต. สาธุ ภิกฺขเว, อหํปิ โข ตํ ภิกฺขเว อตฺตวาทุ\-ปาทานํ น สมนุปสฺสามิ, ยํ’ส อตฺตวาทุ\-ปาทานํ อุปาทิยโต น อุปฺปชฺเชยฺยุํ โสก\-ปริเทว\-ทุกฺข\-โทมนสฺสุ\-ปายาสา.}

\prose{ตํ ภิกฺขเว ทิฏฺฐินิสฺสยํ นิสฺสเยถ, ยํ’ส ทิฏฺฐินิสฺสยํ นิสฺสยโต น อุปฺปชฺเชยฺยุํ โสก\-ปริเทว\-ทุกฺข\-โทมนสฺสุ\-ปายาสา. ปสฺสถ โน ตุเมฺห ภิกฺขเว ตํ ทิฏฺฐินิสฺสยํ, ยํ’ส ทิฏฺฐินิสฺสยํ นิสฺสยโต น อุปฺปชฺเชยฺยุํ โสก\-ปริเทว\-ทุกฺข\-โทมนสฺสุ\-ปายาสาติ. โน เหตํ ภนฺเต. สาธุ ภิกฺขเว, อหํปิ โข ตํ ภิกฺขเว ทิฏฺฐินิสฺสยํ น สมนุปสฺสามิ, ยํ’ส ทิฏฺฐินิสฺสยํ นิสฺสยโต น อุปฺปชฺเชยฺยุํ โสก\-ปริเทว\-ทุกฺข\-โทมนสฺสุ\-ปายาสา.}

\proseitem{244}{อตฺตนิ วา ภิกฺขเว สติ “อตฺตนิยํ เม”ติ อสฺสาติ? เอวํ ภนฺเต. อตฺตนิเย วา ภิกฺขเว สติ “อตฺตา เม”ติ อสฺสาติ? เอวํ ภนฺเต. อตฺตนิ จ ภิกฺขเว อตฺตนิเย จ สจฺจโต เถตโต อนุปลพฺภมาเน ยมฺปิ ตํ ทิฏฺฐิฏฺฐานํ “โส โลโก โส อตฺตา, โส เปจฺจ ภวิสฺสามิ นิจฺโจ ธุโว สสฺสโต อวิปริณาม\-ธมฺโม, สสฺสติสมํ ตเถว ฐสฺสามี”ติ. นนายํ\csromansharedfootnote{9a390e472b004a40}{น จ โขยํ (ก)} ภิกฺขเว เกวโล ปริปูโร พาลธมฺโมติ? กิํ หิ โน สิยา ภนฺเต, เกวโล หิ ภนฺเต ปริปูโร\csromansharedfootnote{d54b673c980fd7d1}{เกวโล ปริปูโร (สี, อิ)} พาลธมฺโมติ.}

\prose{ตํ กิํ มญฺญถ ภิกฺขเว, รูปํ นิจฺจํ วา อนิจฺจํ วาติ? อนิจฺจํ ภนฺเต. ยํ ปนานิจฺจํ, ทุกฺขํ วา ตํ สุขํ วาติ? ทุกฺขํ ภนฺเต. ยํ ปนานิจฺจํ ทุกฺขํ วิปริณาม\-ธมฺมํ, กลฺลํ นุ ตํ สมนุปสฺสิตุํ “เอตํ มม, เอโสหมสฺมิ, เอโส เม อตฺตา”ติ? โน เหตํ ภนฺเต. ตํ กิํ มญฺญถ ภิกฺขเว, เวทนา ฯเปฯ สญฺญา. สงฺขารา. วิญฺญาณํ นิจฺจํ วา อนิจฺจํ วาติ? อนิจฺจํ ภนฺเต. ยํ ปนานิจฺจํ, \csromanfolio{192}ทุกฺขํ วา ตํ สุขํ วาติ? ทุกฺขํ ภนฺเต. ยํ ปนานิจฺจํ ทุกฺขํ วิปริณาม\-ธมฺมํ, กลฺลํ นุ ตํ สมนุปสฺสิตุํ “เอตํ มม, เอโสหมสฺมิ, เอโส เม อตฺตา”ติ? โน เหตํ ภนฺเต. ตสฺมาติห ภิกฺขเว ยํ กิญฺจิ รูปํ อตีตา\-นาคต\-ปจฺจุปฺปนฺนํ อชฺฌตฺตํ วา พหิทฺธา วา โอฬาริกํ วา สุขุมํ วา หีนํ วา ปณีตํ วา ยํ ทูเร สนฺติเก วา, สพฺพํ รูปํ “เนตํ มม, เนโสหมสฺมิ, น เมโส อตฺตา”ติ เอวเมตํ ยถาภูตํ สมฺมปฺปญฺญาย ทฏฺฐพฺพํ. ยา กาจิ เวทนา ฯเปฯ. ยา กาจิ สญฺญา. เย เกจิ สงฺขารา. ยํ กิญฺจิ วิญฺญาณํ อตีตา\-นาคต\-ปจฺจุปฺปนฺนํ อชฺฌตฺตํ วา พหิทฺธา วา โอฬาริกํ วา สุขุมํ วา หีนํ วา ปณีตํ วา ยํ ทูเร สนฺติเก วา, สพฺพํ วิญฺญาณํ “เนตํ มม, เนโสหมสฺมิ, น เมโส อตฺตา”ติ เอวเมตํ ยถาภูตํ สมฺมปฺปญฺญาย ทฏฺฐพฺพํ.}

\proseitem{245}{เอวํ ปสฺสํ ภิกฺขเว สุตวา อริยสาวโก รูปสฺมิํ นิพฺพินฺทติ, เวทนาย นิพฺพินฺทติ, สญฺญาย นิพฺพินฺทติ, สงฺขาเรสุ นิพฺพินฺทติ, วิญฺญาณสฺมิํ นิพฺพินฺทติ, นิพฺพิทา วิรชฺชติ\csromansharedfootnote{442b6d2595859068}{นิพฺพินฺทํ วิรชฺชติ (สี, สฺยา, อิ)}, วิราคา วิมุจฺจติ, วิมุตฺตสฺมิํ “วิมุตฺตมฺ”อิติ ญาณํ โหติ, “ขีณา ชาติ, วุสิตํ พฺรหฺมจริยํ, กตํ กรณียํ, นาปรํ อิตฺถตฺตายา”ติ ปชานาติ. อยํ วุจฺจติ ภิกฺขเว ภิกฺขุ อุกฺขิตฺต\-ปลิโฆ อิติปิ, สํกิณฺณ\-ปริกฺโข อิติปิ, อพฺพูเฬฺหสิโก อิติปิ, นิรคฺคโฬ อิติปิ, อริโย ปนฺนทฺธโช ปนฺนภาโร วิสํยุตฺโต อิติปิ. กถญฺจ ภิกฺขเว ภิกฺขุ อุกฺขิตฺต\-ปลิโฆ โหติ. อิธ ภิกฺขเว ภิกฺขุโน อวิชฺชา ปหีนา โหติ อุจฺฉินฺนมูลา ตาลาวตฺถุกตา อนภาวํกตา อายติํ อนุปฺปาทธมฺมา. เอวํ โข ภิกฺขเว ภิกฺขุ อุกฺขิตฺต\-ปลิโฆ โหติ. กถญฺจ ภิกฺขเว ภิกฺขุ สํกิณฺณ\-ปริกฺโข โหติ. อิธ ภิกฺขเว ภิกฺขุโน โปโนพฺภวิโก ชาติสํสาโร ปหีโน โหติ อุจฺฉินฺนมูโล ตาลาวตฺถุกโต อนภาวํกโต อายติํ อนุปฺปาทธมฺโม. เอวํ โข ภิกฺขเว ภิกฺขุ สํกิณฺณ\-ปริกฺโข โหติ. กถญฺจ ภิกฺขเว ภิกฺขุ อพฺพูเฬฺหสิโก โหติ. อิธ ภิกฺขเว ภิกฺขุโน ตณฺหา ปหีนา โหติ อุจฺฉินฺนมูลา ตาลาวตฺถุกตา อนภาวํกตา อายติํ อนุปฺปาทธมฺมา. เอวํ โข ภิกฺขเว ภิกฺขุ อพฺพูเฬฺหสิโก โหติ. กถญฺจ ภิกฺขเว ภิกฺขุ นิรคฺคโฬ โหติ. อิธ ภิกฺขเว ภิกฺขุโน ปญฺจ โอรมฺภาคิยานิ สํโยชนานิ ปหีนานิ โหนฺติ อุจฺฉินฺนมูลานิ ตาลา\-วตฺถุ\-กตานิ อนภาวํ\-กตานิ อายติํ \csromanfolio{193}อนุปฺปาท\-ธมฺมานิ. เอวํ โข ภิกฺขเว ภิกฺขุ นิรคฺคโฬ โหติ. กถญฺจ ภิกฺขเว ภิกฺขุ อริโย ปนฺนทฺธโช ปนฺนภาโร วิสํยุตฺโต โหติ. อิธ ภิกฺขเว ภิกฺขุโน อสฺมิมาโน ปหีโน โหติ อุจฺฉินฺนมูโล ตาลาวตฺถุกโต อนภาวํกโต อายติํ อนุปฺปาทธมฺโม. เอวํ โข ภิกฺขเว ภิกฺขุ อริโย ปนฺนทฺธโช ปนฺนภาโร วิสํยุตฺโต โหติ.}

\proseitem{246}{เอวํ วิมุตฺตจิตฺตํ โข ภิกฺขเว ภิกฺขุํ สอินฺทา เทวา สพฺรหฺมกา สปชาปติกา อเนฺวสํ นาธิคจฺฉนฺติ “อิทํ นิสฺสิตํ ตถาคตสฺส วิญฺญาณนฺ”ติ. ตํ กิสฺส เหตุ, ทิฏฺเฐวาหํ ภิกฺขเว ธมฺเม ตถาคตํ อนนุวิชฺโชติ วทามิ. เอวํวาทิํ โข มํ ภิกฺขเว เอวมกฺขายิํ เอเก สมณพฺราหฺมณา อสตา ตุจฺฉา มุสา อภูเตน อพฺภาจิกฺขนฺติ “เวนยิโก สมโณ โคตโม สโต สตฺตสฺส อุจฺเฉทํ วินาสํ วิภวํ ปญฺญาเปตี”ติ. ยถา จาหํ น ภิกฺขเว\csromansharedfootnote{9a2378115d750537}{ภิกฺขเว (สี, สฺยา, อิ)}, ยถา จาหํ น วทามิ, ตถา มํ เต โภนฺโต สมณพฺราหฺมณา อสตา ตุจฺฉา มุสา อภูเตน อพฺภาจิกฺขนฺติ “เวนยิโก สมโณ โคตโม สโต สตฺตสฺส อุจฺเฉทํ วินาสํ วิภวํ ปญฺญาเปตี”ติ. ปุพฺเพ จาหํ ภิกฺขเว เอตรหิ จ ทุกฺขํ เจว ปญฺญาเปมิ ทุกฺขสฺส จ นิโรธํ. ตตฺร เจ ภิกฺขเว ปเร ตถาคตํ อกฺโกสนฺติ ปริภาสนฺติ โรเสนฺติ วิเหเสนฺติ, ตตฺร ภิกฺขเว ตถาคตสฺส น โหติ อาฆาโต น อปฺปจฺจโย น เจตโส อนภิรทฺธิ. ตตฺร เจ ภิกฺขเว ปเร ตถาคตํ สกฺกโรนฺติ ครุํ กโรนฺติ มาเนนฺติ ปูเชนฺติ, ตตฺร ภิกฺขเว ตถาคตสฺส น โหติ อานนฺโท น โสมนสฺสํ น เจตโส อุปฺปิลาวิตตฺตํ. ตตฺร เจ ภิกฺขเว ปเร วา ตถาคตํ สกฺกโรนฺติ ครุํ กโรนฺติ มาเนนฺติ ปูเชนฺติ, ตตฺร ภิกฺขเว ตถาคตสฺส เอวํ โหติ “ยํ โข อิทํ ปุพฺเพ ปริญฺญาตํ, ตตฺถ’เม เอวรูปา การา\csromansharedfootnote{2cd1200cc5b58596}{สกฺการา (ก)} กรียนฺตี”ติ. ตสฺมาติห ภิกฺขเว ตุเมฺห เจปิ ปเร อกฺโกเสยฺยุํ ปริภาเสยฺยุํ โรเสยฺยุํ วิเหเสยฺยุํ, ตตฺร ตุเมฺห หิ น อาฆาโต น อปฺปจฺจโย น เจตโส อนภิรทฺธิ กรณียา. ตสฺมาติห ภิกฺขเว ตุเมฺห เจปิ ปเร สกฺกเรยฺยุํ ครุํ กเรยฺยุํ มาเนยฺยุํ ปูเชยฺยุํ, ตตฺร ตุเมฺหหิ น อานนฺโท น โสมนสฺสํ น เจตโส อุปฺปิลาวิตตฺตํ กรณียํ. ตสฺมาติห ภิกฺขเว ตุเมฺห เจปิ ปเร สกฺกเรยฺยุํ \csromanfolio{194}ครุํ กเรยฺยุํ มาเนยฺยุํ ปูเชยฺยุํ. ตตฺร ตุมฺหากํ เอวมสฺส “ยํ โข อิทํ ปุพฺเพ ปริญฺญาตํ, ตตฺถ’เม\csromansharedfootnote{b21de98cdc402100}{ตตฺถ โน (ก) ตตฺถ + อิเมติ ปทจฺเฉโท.} เอวรูปา การา กรียนฺตี”ติ.}

\proseitem{247}{ตสฺมาติห ภิกฺขเว ยํ น ตุมฺหากํ ตํ ปชหถ. ตํ โว ปหีนํ ทีฆรตฺตํ หิตาย สุขาย ภวิสฺสติ. กิญฺจ ภิกฺขเว น ตุมฺหากํ, รูปํ ภิกฺขเว น ตุมฺหากํ, ตํ ปชหถ. ตํ โว ปหีนํ ทีฆรตฺตํ หิตาย สุขาย ภวิสฺสติ. เวทนา ภิกฺขเว น ตุมฺหากํ, ตํ ปชหถ. สา โว ปหีนา ทีฆรตฺตํ หิตาย สุขาย ภวิสฺสติ. สญฺญา ภิกฺขเว น ตุมฺหากํ, ตํ ปชหถ. สา โว ปหีนา ทีฆรตฺตํ หิตาย สุขาย ภวิสฺสติ. สงฺขารา ภิกฺขเว น ตุมฺหากํ, เต ปชหถ. เต โว ปหีนา ทีฆรตฺตํ หิตาย สุขาย ภวิสฺสนฺติ. วิญฺญาณํ ภิกฺขเว น ตุมฺหากํ, ตํ ปชหถ. ตํ โว ปหีนํ ทีฆรตฺตํ หิตาย สุขาย ภวิสฺสติ. ตํ กิํ มญฺญถ ภิกฺขเว, ยํ อิมสฺมิํ เชตวเน ติณกฏฺฐ\-สาขา\-ปลาสํ, ตํ ชโน หเรยฺย วา ทเหยฺย วา ยถาปจฺจยํ วา กเรยฺย, อปิ นุ ตุมฺหากํ เอวมสฺส “อเมฺห ชโน หรติ วา ทหติ วา ยถาปจฺจยํ วา กโรตี”ติ. โน เหตํ ภนฺเต. ตํ กิสฺส เหตุ. น หิ โน เอตํ ภนฺเต อตฺตา วา อตฺตนิยํ วาติ. เอวเมว โข ภิกฺขเว ยํ น ตุมฺหากํ, ตํ ปชหถ. ตํ โว ปหีนํ ทีฆรตฺตํ หิตาย สุขาย ภวิสฺสติ. กิญฺจ ภิกฺขเว น ตุมฺหากํ, รูปํ ภิกฺขเว น ตุมฺหากํ, ตํ ปชหถ. ตํ โว ปหีนํ ทีฆรตฺตํ หิตาย สุขาย ภวิสฺสติ. เวทนา ภิกฺขเว ฯเปฯ สญฺญา ภิกฺขเว. สงฺขารา ภิกฺขเว. วิญฺญาณํ ภิกฺขเว น ตุมฺหากํ, ตํ ปชหถ. ตํ โว ปหีนํ ทีฆรตฺตํ หิตาย สุขาย ภวิสฺสติ.}

\proseitem{248}{เอวํ สฺวากฺขาโต ภิกฺขเว มยา ธมฺโม อุตฺตาโน วิวโฏ ปกาสิโต ฉินฺนปิโลติโก. เอวํ สฺวากฺขาเต ภิกฺขเว มยา ธมฺเม อุตฺตาเน วิวเฏ ปกาสิเต ฉินฺนปิโลติเก เย เต ภิกฺขู อรหนฺโต ขีณาสวา วุสิตวนฺโต กตกรณียา โอหิตภารา อนุปฺปตฺต\-สทตฺถา ปริกฺขีณ\-ภวสํโยชนา สมฺมทญฺญา วิมุตฺตา, วฏฺฏํ เตสํ นตฺถิ ปญฺญาปนาย. เอวํ สฺวากฺขาโต ภิกฺขเว มยา ธมฺโม อุตฺตาโน วิวโฏ ปกาสิโต ฉินฺนปิโลติโก. เอวํ สฺวากฺขาเต ภิกฺขเว มยา ธมฺเม อุตฺตาเน วิวเฏ ปกาสิเต ฉินฺนปิโลติเก เยสํ ภิกฺขูนํ ปญฺโจ\-รมฺภาคิยานิ \csromanfolio{195}สํโยชนานิ ปหีนานิ, สพฺเพ เต โอปปาติกา ตตฺถ ปรินิพฺพายิโน อนาวตฺติธมฺมา ตสฺมา โลกา. เอวํ สฺวากฺขาโต ภิกฺขเว มยา ธมฺโม อุตฺตาโน วิวโฏ ปกาสิโต ฉินฺนปิโลติโก. เอวํ สฺวากฺขาเต ภิกฺขเว มยา ธมฺเม อุตฺตาเน วิวเฏ ปกาสิเต ฉินฺนปิโลติเก เยสํ ภิกฺขูนํ ตีณิ สํโยชนานิ ปหีนานิ ราคโทสโมหา ตนุภูตา, สพฺเพ เต สกทาคามิโน สกิเทว อิมํ โลกํ อาคนฺตฺวา ทุกฺขสฺสนฺตํ กริสฺสนฺติ. เอวํ สฺวากฺขาโต ภิกฺขเว มยา ธมฺโม อุตฺตาโน วิวโฏ ปกาสิโต ฉินฺนปิโลติโก. เอวํ สฺวากฺขาเต ภิกฺขเว มยา ธมฺเม อุตฺตาเน วิวเฏ ปกาสิเต ฉินฺนปิโลติเก เยสํ ภิกฺขูนํ ตีณิ สํโยชนานิ ปหีนานิ, สพฺเพ เต โสตาปนฺนา อวินิปาต\-ธมฺมา นิยตา สมฺโพธิ\-ปรายนา. เอวํ สฺวากฺขาโต ภิกฺขเว มยา ธมฺโม อุตฺตาโน วิวโฏ ปกาสิโต ฉินฺนปิโลติโก. เอวํ สฺวากฺขาเต ภิกฺขเว มยา ธมฺเม อุตฺตาเน วิวเฏ ปกาสิเต ฉินฺนปิโลติเก เย เต ภิกฺขู ธมฺมานุสาริโน สทฺธานุสาริโน, สพฺเพ เต สมฺโพธิ\-ปรายนา. เอวํ สฺวากฺขาโต ภิกฺขเว มยา ธมฺโม อุตฺตาโน วิวโฏ ปกาสิโต ฉินฺนปิโลติโก. เอวํ สฺวากฺขาเต ภิกฺขเว มยา ธมฺเม อุตฺตาเน วิวเฏ ปกาสิเต ฉินฺนปิโลติเก เยสํ มยิ สทฺธามตฺตํ เปมมตฺตํ, สพฺเพ เต สคฺคปรายนาติ.}

\prosewithcloserruled{อิทมโวจ ภควา. อตฺตมนา เต ภิกฺขู ภควโต ภาสิตํ อภินนฺทุนฺติ.}{อลคทฺทู\-ปม\-สุตฺตํ นิฏฺฐิตํ ทุติยํ.}
\csromanmatikaanchor{mtk.53}
\csromantocmarkref{subsection}{3. วมฺมิกสุตฺต}{mtk.53}
\bookmark[dest={mtk.53},level=2]{3. วมฺมิกสุตฺต}
\csromanheader{2}{3. วมฺมิกสุตฺต}
\proseitem{249}{เอวํ เม สุตํ\csromandash{}เอกํ สมยํ ภควา สาวตฺถิยํ วิหรติ เชตวเน อนาถ\-ปิณฺฑิกสฺส อาราเม. เตน โข ปน สมเยน อายสฺมา กุมารกสฺสโป อนฺธวเน วิหรติ. อถ โข อญฺญตรา เทวตา อภิกฺกนฺตาย รตฺติยา อภิกฺกนฺต\-วณฺณา เกวลกปฺปํ อนฺธวนํ โอภาเสตฺวา เยนายสฺมา กุมารกสฺสโป เตนุปสงฺกมิ, อุปสงฺกมิตฺวา เอกมนฺตํ อฏฺฐาสิ. เอกมนฺตํ ฐิตา โข สา เทวตา อายสฺมนฺตํ กุมารกสฺสปํ เอตทโวจ\csromandash{}}

\prose{\csromanfolio{196}ภิกฺขุ ภิกฺขุ อยํ วมฺมิโก\csromansharedfootnote{65b35789d913e7ec}{วมฺมีโก (กตฺถจิ) สกฺกตานุรูปํ.} รตฺติํ ธูมายติ, ทิวา ปชฺชลติ. พฺราหฺมโณ เอวมาห “อภิกฺขณ สุเมธ สตฺถํ อาทายา”ติ. อภิกฺขณนฺโต สุเมโธ สตฺถํ อาทาย อทฺทส ลงฺคิํ. ลงฺคี ภทนฺเตติ. พฺราหฺมโณ เอวมาห “อุกฺขิป ลงฺคิํ อภิกฺขณ สุเมธ สตฺถํ อาทายา”ติ. อภิกฺขณนฺโต สุเมโธ สตฺถํ อาทาย อทฺทส อุทฺธุมายิกํ. อุทฺธุมายิกา ภทนฺเตติ. พฺราหฺมโณ เอวมาห “อุกฺขิป อุทฺธุมายิกํ อภิกฺขณ สุเมธ สตฺถํ อาทายา”ติ. อภิกฺขณนฺโต สุเมโธ สตฺถํ อาทาย อทฺทส ทฺวิธาปถํ. ทฺวิธาปโถ ภทนฺเตติ. พฺราหฺมโณ เอวมาห “อุกฺขิป ทฺวิธาปถํ อภิกฺขณ สุเมธ สตฺถํ อาทายา”ติ. อภิกฺขณนฺโต สุเมโธ สตฺถํ อาทาย อทฺทส จงฺควารํ\csromansharedfootnote{f85d0e24e1c2c95d}{ปงฺกวารํ (สฺยา), จงฺกวารํ (ก)}. จงฺควาโร ภทนฺเตติ. พฺราหฺมโณ เอวมาห “อุกฺขิป จงฺควารํ อภิกฺขณ สุเมธ สตฺถํ อาทายา”ติ. อภิกฺขณนฺโต สุเมโธ สตฺถํ อาทาย อทฺทส กุมฺมํ. กุมฺโม ภทนฺเตติ. พฺราหฺมโณ เอวมาห “อุกฺขิป กุมฺมํ อภิกฺขณ สุเมธ สตฺถํ อาทายา”ติ. อภิกฺขณนฺโต สุเมโธ สตฺถํ อาทาย อทฺทส อสิสูนํ. อสิสูนา ภทนฺเตติ. พฺราหฺมโณ เอวมาห “อุกฺขิป อสิสูนํ อภิกฺขณ สุเมธ สตฺถํ อาทายา”ติ. อภิกฺขณนฺโต สุเมโธ สตฺถํ อาทาย อทฺทส มํสเปสิํ. มํสเปสิ ภทนฺเตติ. พฺราหฺมโณ เอวมาห “อุกฺขิป มํสเปสิํ อภิกฺขณ สุเมธ สตฺถํ อาทายา”ติ. อภิกฺขณนฺโต สุเมโธ สตฺถํ อาทาย อทฺทส นาคํ. นาโค ภทนฺเตติ. พฺราหฺมโณ เอวมาห “ติฏฺฐตุ นาโค, มา นาคํ ฆฏฺเฏสิ, นโม กโรหิ นาคสฺสา”ติ. อิเม โข ตฺวํ ภิกฺขุ ปเญฺห ภควนฺตํ อุปสงฺกมิตฺวา ปุจฺโฉยฺยาสิ. ยถา จ เต ภควา พฺยากโรติ, ตถา นํ ธาเรยฺยาสิ. นาหํ ตํ ภิกฺขุ ปสฺสามิ สเทวเก โลเก สมารเก สพฺรหฺมเก สสฺสมณ\-พฺราหฺมณิยา ปชาย สเทว\-มนุสฺสาย, โย อิเมสํ ปญฺหานํ เวยฺยากรเณน จิตฺตํ อาราเธยฺย อญฺญตฺร ตถาคเตน วา ตถาคต\-สาวเกน วา อิโต วา ปน สุตฺวาติ. อิทมโวจ สา เทวตา อิทํ วตฺวา ตตฺเถ\-ว\-นฺตรธายิ.}

\proseitem{250}{อถ โข อายสฺมา กุมารกสฺสโป ตสฺสา รตฺติยา อจฺจเยน เยน ภควาเตนุปสงฺกมิ, อุปสงฺกมิตฺวา ภควนฺตํ อภิวาเทตฺวา เอกมนฺตํ นิสีทิ. เอกมนฺตํ นิสินฺโน โข อายสฺมา กุมารกสฺสโป \csromanfolio{197}ภควนฺตํ เอตทโวจ “อิมํ ภนฺเต รตฺติํ อญฺญตรา เทวตา อภิกฺกนฺตาย รตฺติยา อภิกฺกนฺต\-วณฺณา เกวลกปฺปํ อนฺธวนํ โอภาเสตฺวา เยนาหํ เตนุปสงฺกมิ, อุปสงฺกมิตฺวา เอกมนฺตํ อฏฺฐาสิ. เอกมนฺตํ ฐิตา โข ภนฺเต สา เทวตา มํ เอตทโวจ ‘ภิกฺขุ ภิกฺขุ อยํ วมฺมิโก รตฺติํ ธูมายติ, ทิวา ปชฺชลติ. พฺราหฺมโณ เอวมาห อภิกฺขณ สุเมธ สตฺถํ อาทายาติ. อภิกฺขณนฺโต สุเมโธ สตฺถํ อาทาย ฯเปฯ อิโต วา ปน สุตฺวา’ติ. อิทมโวจ ภนฺเต สา เทวตา อิทํ วตฺวา ตตฺเถ\-ว\-นฺตรธายิ. โก นุ โข ภนฺเต วมฺมิโก, กา รตฺติํ ธูมายนา, กา ทิวา ปชฺชลนา, โก พฺราหฺมโณ, โก สุเมโธ, กิํ สตฺถํ, กิํ อภิกฺขณํ, กา ลงฺคี, กา อุทฺธุมายิกา, โก ทฺวิธาปโถ, กิํ จงฺควารํ, โก กุมฺโม, กา อสิสูนา, กา มํสเปสิ, โก นาโค”ติ.}

\proseitem{251}{“วมฺมิโก”ติ โข ภิกฺขุ อิมสฺเสตํ จาตุมหาภูติกสฺส\csromansharedfootnote{561989e8be60e04f}{จาตุมฺมหาภูติกสฺส (สี, สฺยา, อิ)} กายสฺส อธิวจนํ มาตา\-เปตฺติก\-สมฺภวสฺส โอทน\-กุมฺมาสู\-ปจยสฺส อนิจฺจุจฺฉาทน ปริมทฺทน เภทน วิทฺธํสน ธมฺมสฺส. (1)}

\prose{ยํ โข ภิกฺขุ ทิวา กมฺมนฺเต\csromansharedfootnote{50f12445da1922d7}{กมฺมนฺตํ(ก)} อารพฺภ รตฺติํ อนุวิตกฺเกติ อนุวิจาเรติ, อยํ รตฺติํ ธูมายนา. ยํ โข ภิกฺขุ รตฺติํ อนุวิตกฺเกตฺวา อนุวิจาเรตฺวา ทิวา กมฺมนฺเต ปโยเชติ กาเยน วาจาย (มนสา)\csromansharedfootnote{862afa2d474204c6}{( ) นตฺถิ (สี, สฺยา)}. อยํ ทิวา ปชฺชลนา. (\csromansharedfootnote{50f12445da1922d7}{กมฺมนฺตํ(ก)}-3)}

\prose{“พฺราหฺมโณ”ติ โข ภิกฺขุ ตถาคตสฺเสตํ อธิวจนํ อรหโต สมฺมา\-สมฺพุทฺธสฺส. “สุเมโธ”ติ โข ภิกฺขุ เสกฺขสฺเสตํ ภิกฺขุโน อธิวจนํ. \mbox{(4-5)}}

\prose{“สตฺถนฺ”ติ โข ภิกฺขุ อริยาเยตํ ปญฺญาย อธิวจนํ. “อภิกฺขณนฺ”ติ โข ภิกฺขุ วีริยารมฺภ สฺเสตํ อธิวจนํ. \mbox{(6-7)}}

\prose{“ลงฺคี”ติ โข ภิกฺขุ อวิชฺชาเยตํ อธิวจนํ. อุกฺขิป ลงฺคิํ ปชห อวิชฺชํ, อภิกฺขณ สุเมธ สตฺถํ อาทายาติ อยเมตสฺส อตฺโถ. (8)}

\prose{“อุทฺธุมายิกา”ติ โข ภิกฺขุ โกธูปายาสสฺเสตํ อธิวจนํ. อุกฺขิป อุทฺธุมายิกํ ปชห โกธูปายาสํ, อภิกฺขณ สุเมธ สตฺถํ อาทายาติ อยเมตสฺส อตฺโถ. (9)}

\prose{\csromanfolio{198}“ทฺวิธาปโถ”ติ โข ภิกฺขุ วิจิกิจฺฉาเยตํ อธิวจนํ. อุกฺขิป ทฺวิธาปถํ ปชห วิจิกิจฺฉํ, อภิกฺขณ สุเมธ สตฺถํ อาทายาติ อยเมตสฺส อตฺโถ. (10)}

\prose{“จงฺควารนฺ”ติ โข ภิกฺขุ ปญฺจนฺเนตํ นีวรณานํ อธิวจนํ. เสยฺยถิทํ, กามจฺฉนฺท\-นีวรณสฺส พฺยาปาท\-นีวรณสฺส ถินมิทฺธ\-นีวรณสฺส อุทฺธจฺจกุกฺกุจฺจ\-นีวรณสฺส วิจิกิจฺฉา\-นีวรณสฺส. อุกฺขิป จงฺควารํ ปชห ปญฺจ นีวรเณ, อภิกฺขณ สุเมธ สตฺถํ อาทายาติ อยเมตสฺส อตฺโถ. (11)}

\prose{“กุมฺโม”ติ โข ภิกฺขุ ปญฺจนฺเนตํ อุปาทาน\-กฺขนฺธานํ อธิวจนํ. เสยฺยถิทํ, รูปุปาทาน\-กฺขนฺธสฺส เวทนุ\-ปาทาน\-กฺขนฺธสฺส สญฺญุ\-ปาทาน\-กฺขนฺธสฺส สงฺขารุ\-ปาทาน\-กฺขนฺธสฺส วิญฺญาณุ\-ปาทาน\-กฺขนฺธสฺส. อุกฺขิป กุมฺมํ ปชห ปญฺจุ\-ปาทานกฺขนฺเธ, อภิกฺขณ สุเมธ สตฺถํ อาทายาติ อยเมตสฺส อตฺโถ. (12)}

\prose{“อสิสูนา”ติ โข ภิกฺขุ ปญฺจนฺเนตํ กามคุณานํ อธิวจนํ. จกฺขุ\-วิญฺเญยฺยานํ รูปานํ อิฏฺฐานํ กนฺตานํ มนาปานํ ปิยรูปานํ กามู\-ปสํหิตานํ รชนียานํ. โสต\-วิญฺเญยฺยานํ สทฺทานํ ฯเปฯ. ฆาน\-วิญฺเญยฺยานํ คนฺธานํ ฯเปฯ. ชิวฺหา\-วิญฺเญยฺยานํ รสานํ ฯเปฯ. กาย\-วิญฺเญยฺยานํ โผฏฺฐพฺพานํ อิฏฺฐานํ กนฺตานํ มนาปานํ ปิยรูปานํ กามู\-ปสํหิตานํ รชนียานํ. อุกฺขิป อสิสูนํ ปชห ปญฺจ กามคุเณ, อภิกฺขณ สุเมธ สตฺถํ อาทายาติ อยเมตสฺส อตฺโถ. (13)}

\prose{“มํสเปสี”ติ โข ภิกฺขุ นนฺทีราคสฺเสตํ อธิวจนํ. อุกฺขิป มํสเปสิํ ปชห นนฺทีราคํ, อภิกฺขณ สุเมธ สตฺถํ อาทายาติ อยเมตสฺส อตฺโถ. (14)}

\prose{“นาโค”ติ โข ภิกฺขุ ขีณาสวสฺเสตํ ภิกฺขุโน อธิวจนํ. ติฏฺฐตุ นาโค, มา นาคํ ฆฏฺเฏสิ, นโม กโรหิ นาคสฺสาติ อยเมตสฺส อตฺโถติ. (15)}

\prosewithcloser{อิทมโวจ ภควา. อตฺตมโน อายสฺมา กุมารกสฺสโป ภควโต ภาสิตํ อภินนฺทีติ.}{วมฺมิกสุตฺตํ นิฏฺฐิตํ ตติยํ.}
\csromanmatikaanchor{mtk.54}
\csromantocmarkref{subsection}{4. รถวินีตสุตฺต}{mtk.54}
\bookmark[dest={mtk.54},level=2]{4. รถวินีตสุตฺต}
\csromanmarkh{4. รถวินีต\-สุตฺต (24)}\csromanmarkhii{4. รถวินีต\-สุตฺต}\csromanheadernomark{2}{4. รถวินีต\-สุตฺต (24) 4. รถวินีต\-สุตฺต}
\proseitem{252}{\csromanfolio{199}เอวํ เม สุตํ\csromandash{}เอกํ สมยํ ภควา ราชคเห วิหรติ เวฬุวเน กลนฺทก\-นิวาเป. อถ โข สมฺพหุลา ชาติภูมกา ภิกฺขู ชาติภูมิยํ วสฺสํวุฏฺฐา เยน ภควา เตนุ\-ปสงฺกมิํสุ, อุปสงฺกมิตฺวา ภควนฺตํ อภิวาเทตฺวา เอกมนฺตํ นิสีทิํสุ. เอกมนฺตํ นิสินฺเน โข เต ภิกฺขู ภควา เอตทโวจ\csromandash{}}

\prose{โก นุ โข ภิกฺขเว ชาติภูมิยํ ชาติภูมกานํ ภิกฺขูนํ สพฺรหฺมจารีนํ เอวํ สมฺภาวิโต “อตฺตนา จ อปฺปิจฺโฉ อปฺปิ\-จฺฉ\-กถญฺจ ภิกฺขูนํ กตฺตา, อตฺตนา จ สนฺตุฏฺโฐ สนฺตุฏฺฐิกถญฺจ ภิกฺขูนํ กตฺตา, อตฺตนา จ ปวิวิตฺโต ปวิเวกกถญฺจ ภิกฺขูนํ กตฺตา, อตฺตนา จ อสํสฏฺโฐ อสํสคฺคกถญฺจ ภิกฺขูนํ กตฺตา, อตฺตนา จ อารทฺธวีริโย วีริยา\-รมฺภ\-กถญฺจ ภิกฺขูนํ กตฺตา, อตฺตนา จ สีลสมฺปนฺโน สีลสมฺปทา\-กถญฺจ ภิกฺขูนํ กตฺตา, อตฺตนา จ สมาธิ\-สมฺปนฺโน สมาธิ\-สมฺปทา\-กถญฺจ ภิกฺขูนํ กตฺตา, อตฺตนา จ ปญฺญาสมฺปนฺโน ปญฺญา\-สมฺปทา\-กถญฺจ ภิกฺขูนํ กตฺตา, อตฺตนา จ วิมุตฺติ\-สมฺปนฺโน วิมุตฺติ\-สมฺปทา\-กถญฺจ ภิกฺขูนํ กตฺตา, อตฺตนา จ วิมุตฺติ\-ญาณ\-ทสฺสน\-สมฺปนฺโน วิมุตฺติ\-ญาณ\-ทสฺสน\-สมฺปทา\-กถญฺจ ภิกฺขูนํ กตฺตา, โอวาทโก วิญฺญาปโก สนฺทสฺสโก สมาทปโก สมุตฺเตชโก สมฺปหํสโก สพฺรหฺมจารีนนฺ”ติ. ปุณฺโณ นาม ภนฺเต อายสฺมา มนฺตาณิปุตฺโต ชาติภูมิยํ ชาติภูมกานํ ภิกฺขูนํ สพฺรหฺมจารีนํ เอวํ สมฺภาวิโต “อตฺตนา จ อปฺปิจฺโฉ อปฺปิ\-จฺฉ\-กถญฺจ ภิกฺขูนํ กตฺตา, อตฺตนา จ สนฺตุฏฺโฐ ฯเปฯ โอวาทโก วิญฺญาปโก สนฺทสฺสโก สมาทปโก สมุตฺเตชโก สมฺปหํสโก สพฺรหฺมจารีนนฺ”ติ.}

\proseitem{253}{เตน โข ปน สมเยน อายสฺมา สาริปุตฺโต ภควโต อวิทูเร นิสินฺโน โหติ. อถ โข อายสฺมโต สาริปุตฺตสฺส เอตทโหสิ “ลาภา อายสฺมโต ปุณฺณสฺส มนฺตาณิ\-ปุตฺตสฺส, สุลทฺธลาภา อายสฺมโต ปุณฺณสฺส มนฺตาณิ\-ปุตฺตสฺส, ยสฺส วิญฺญู สพฺรหฺมจารี สตฺถุ สมฺมุขา อนุมสฺส อนุมสฺส วณฺณํ ภาสนฺติ. ตญฺจ สตฺถา อพฺภนุโมทติ. อปฺเปว นาม มยํปิ กทาจิ กรหจิ อายสฺมตา ปุณฺเณน มนฺตาณิปุตฺเตน สทฺธิํ สมาคจฺเฉยฺยาม\csromansharedfootnote{3d78714d2879445c}{สมาคมํ คจฺเฉยฺย (ก)}, อปฺเปว นาม สิยา โกจิเทว กถาสลฺลาโป”ติ.}

\proseitem{254}{\csromanfolio{200}อถ โข ภควา ราชคเห ยถาภิรนฺตํ วิหริตฺวา เยน สาวตฺถิ เตน จาริกํ ปกฺกามิ. อนุปุพฺเพน จาริกํ จรมาโน เยน สาวตฺถิ ตทวสริ, ตตฺร สุทํ ภควา สาวตฺถิยํ วิหรติ เชตวเน อนาถ\-ปิณฺฑิกสฺส อาราเม. อสฺโสสิ โข อายสฺมา ปุณฺโณ มนฺตาณิปุตฺโต “ภควา กิร สาวตฺถิํ อนุปฺปตฺโต, สาวตฺถิยํ วิหรติ เชตวเน อนาถ\-ปิณฺฑิกสฺส อาราเม”ติ.}

\proseitem{255}{อถ โข อายสฺมา ปุณฺโณ มนฺตาณิปุตฺโต เสนาสนํ สํสาเมตฺวา ปตฺต\-จีวรมา\-ทาย เยน สาวตฺถิ เตน จาริกํ ปกฺกามิ. อนุปุพฺเพน จาริกํ จรมาโน เยน สาวตฺถิ เชตวนํ อนาถ\-ปิณฺฑิกสฺส อาราโม, เยน ภควา เตนุปสงฺกมิ, อุปสงฺกมิตฺวา ภควนฺตํ อภิวาเทตฺวา เอกมนฺตํ นิสีทิ. เอกมนฺตํ นิสินฺนํ โข อายสฺมนฺตํ ปุณฺณํ มนฺตาณิปุตฺตํ ภควา ธมฺมิยา กถาย สนฺทสฺเสสิ สมาทเปสิ สมุตฺเตเชสิ สมฺปหํเสสิ. อถ โข อายสฺมา ปุณฺโณ มนฺตาณิปุตฺโต ภควตา ธมฺมิยา กถาย สนฺทาสฺสิโต สมาทปิโต สมุตฺเตชิโต สมฺปหํสิโต ภควโต ภาสิตํ อภินนฺทิตฺวา อนุโมทิตฺวา อุฏฺฐายาสนา ภควนฺตํ อภิวาเทตฺวา ปทกฺขิณํ กตฺวา เยน อนฺธวนํ เตนุปสงฺกมิ ทิวาวิหาราย.}

\proseitem{256}{อถ โข อญฺญตโร ภิกฺขุ เยนายสฺมา สาริปุตฺโต เตนุปสงฺกมิ, อุปสงฺกมิตฺวา อายสฺมนฺตํ สาริปุตฺตํ เอตทโวจ “ยสฺส โข ตฺวํ อาวุโส สาริปุตฺต ปุณฺณสฺส นาม ภิกฺขุโน มนฺตาณิ\-ปุตฺตสฺส อภิณฺหํ กิตฺตยมาโน อโหสิ. โส ภควตา ธมฺมิยา กถาย สนฺทสฺสิโต สมาทปิโต สมุตฺเตชิโต สมฺปหํสิโต ภควโต ภาสิตํ อภินนฺทิตฺวา อนุโมทิตฺวา อุฏฺฐายาสนา ภควนฺตํ อภิวาเทตฺวา ปทกฺขิณํ กตฺวา เยน อนฺธวนํ เตน ปกฺกนฺโต ทิวาวิหารายา”ติ.}

\prose{อถ โข อายสฺมา สาริปุตฺโต ตรมานรูโป นิสีทนํ อาทาย อายสฺมนฺตํ ปุณฺณํ มนฺตาณิปุตฺตํ ปิฏฺฐิโต ปิฏฺฐิโต อนุพนฺธิ สีสานุโลกี. อถ โข อายสฺมา ปุณฺโณ มนฺตาณิปุตฺโต อนฺธวนํ อชฺโฌคาเหตฺวา อญฺญตรสฺมิํ รุกฺขมูเล ทิวาวิหารํ นิสีทิ. อายสฺมาปิ โข สาริปุตฺโต อนฺธวนํ อชฺโฌคาเหตฺวา อญฺญตรสฺมิํ รุกฺขมูเล ทิวาวิหารํ นิสีทิ. \csromanfolio{201}อถ โข อายสฺมา สาริปุตฺโต สายนฺหสมยํ ปฏิสลฺลานา วุฏฺฐิโต เยนายสฺมา ปุณฺโณ มนฺตาณิปุตฺโต เตนุปสงฺกิมิ, อุปสงฺกมิตฺวา อายสฺมตา ปุณฺเณน มนฺตาณิปุตฺเตน สทฺธิํ สมฺโมทิ, สมฺโมทนียํ กถํ สารณิยํ วีติสาเรตฺวา เอกมนฺตํ นิสีทิ. เอกมนฺตํ นิสินฺโน โข อายสฺมา สาริปุตฺโต อายสฺมนฺตํ ปุณฺณํ มนฺตาณิปุตฺตํ เอตทโวจ\csromandash{}}

\proseitem{257}{“ภควติ โน อาวุโส พฺรหฺมจริยํ วุสฺสตี”ติ. เอวมาวุโสติ. กิํ นุ โข อาวุโส สีลวิสุทฺธตฺถํ ภควติ พฺรหฺมจริยํ วุสฺสตีติ. โน หิทํ อาวุโส. กิํ ปนาวุโส จิตฺต\-วิสุทฺธตฺถํ ภควติ พฺรหฺมจริยํ วุสฺสตีติ. โน หิทํ อาวุโส. กิํ นุ โข อาวุโส ทิฏฺฐิ\-วิสุทฺธตฺถํ ภควติ พฺรหฺมจริยํ วุสฺสตีติ. โน หิทํ อาวุโส. กิํ ปนาวุโส กงฺขา\-วิตรณ\-วิสุทฺธตฺถํ ภควติ พฺรหฺมจริยํ วุสฺสตีติ. โน หิทํ อาวุโส. กิํ นุ โข อาวุโส มคฺคา\-มคฺค\-ญาณ\-ทสฺสน\-วิสุทฺธตฺถํ ภควติ พฺรหฺมจริยํ วุสฺสตีติ. โน หิทํ อาวุโส. กิํ ปนาวุโส ปฏิปทา\-ญาณ\-ทสฺสน\-วิสุทฺธตฺถํ ภควติ พฺรหฺมจริยํ วุสฺสตีติ. โน หิทํ อาวุโส. กิํ นุ โข อาวุโส ญาณ\-ทสฺสน\-วิสุทฺธตฺถํ ภควติ พฺรหฺมจริยํ วุสฺสตีติ. โน หิทํ อาวุโส.}

\prose{“กิํ นุ โข อาวุโส สีลวิสุทฺธตฺถํ ภควติ พฺรหฺมจริยํ วุสฺสตี”ติ อิติ ปุฏฺโฐ สมาโน “โน หิทํ อาวุโส”ติ วเทสิ. “กิํ ปนาวุโส จิตฺต\-วิสุทฺธตฺถํ ภควติ พฺรหฺมจริยํ วุสฺสตี”ติ อิติ ปุฏฺโฐ สมาโน “โน หิทํ อาวุโส”ติ วเทสิ. “กิํ นุ โข อาวุโส ทิฏฺฐิ\-วิสุทฺธตฺถํ ฯเปฯ กงฺขา\-วิตรณ\-วิสุทฺธตฺถํ. มคฺคา\-มคฺค\-ญาณ\-ทสฺสน\-วิสุทฺธตฺถํ. ปฏิปทา\-ญาณ\-ทสฺสน\-วิสุทฺธตฺถํ ฯเปฯ “กิํ นุ โข อาวุโส ญณทสฺสนวิสุทฺธตฺถํ ภควติ พฺรหฺมจริยํ วุสฺสตี”ติ อิติ ปุฏฺโฐ สมาโน “โน หิทํ อาวุโส”ติ วเทสิ. กิมตฺถํ จรหาวุโส ภควติ พฺรหฺมจริยํ วุสฺสตีติ. อนุปาทา\-ปรินิพฺพานตฺถํ โข อาวุโส ภควติ พฺรหฺมจริยํ วุสฺสตีติ.}

\prose{กิํ นุ โข อาวุโส สีลวิสุทฺธิ อนุปาทา\-ปรินิพฺพานนฺติ. โน หิทํ อาวุโส. กิํ ปนาวุโส จิตฺตวิสุทฺธิ อนุปาทา\-ปรินิพฺพานนฺติ. โน หิทํ อาวุโส. กิํ นุ โข อาวุโส ทิฏฺฐิวิสุทฺธิ อนุปาทา\-ปรินิพฺพานนฺติ. โน หิทํ อาวุโส. กิํ ปนาวุโส กงฺขา\-วิตรณ\-วิสุทฺธิ \csromanfolio{202}อนุปาทา\-ปรินิพฺพานนฺติ. โน หิทํ อาวุโส. กิํ นุ โข อาวุโส มคฺคา\-มคฺค\-ญาณ\-ทสฺสน\-วิสุทฺธิ อนุปาทา\-ปรินิพฺพานนฺติ. โน หิทํ อาวุโส. กิํ ปนาวุโส ปฏิปทา\-ญาณ\-ทสฺสน\-วิสุทฺธิ อนุปาทา\-ปรินิพฺพานนฺติ. โน หิทํ อาวุโส. กิํ นุ โข อาวุโส ญาณ\-ทสฺสน\-วิสุทฺธิ อนุปาทา\-ปรินิพฺพานนฺติ. โน หิทํ อาวุโส. กิํ ปนาวุโส อญฺญตฺร อิเมหิ ธมฺเมหิ อนุปาทา\-ปรินิพฺพานนฺติ. โน หิทํ อาวุโส.}

\prose{“กิํ นุ โข อาวุโส สีลวิสุทฺธิ อนุปาทาปรินิพฺพานนฺ”ติ อิติ ปุฏฺโฐ สมาโน “โน หิทํ อาวุโส”ติ วเทสิ. “กิํ ปนาวุโส จิตฺตวิสุทฺธิ อนุปาทาปรินิพฺพานนฺ”ติ อิติ ปุฏฺโฐ สมาโน “โน หิทํ อาวุโส”ติ วเทสิ. “กิํ นุ โข อาวุโส ทิฏฺฐิวิสุทฺธิ อนุปาทาปรินิพฺพานนฺ”ติ ฯเปฯ กงฺขา\-วิตรณ\-วิสุทฺธิ. มคฺคา\-มคฺค\-ญาณ\-ทสฺสน\-วิสุทฺธิ. ปฏิปทา\-ญาณ\-ทสฺสน\-วิสุทฺธิ. “กิํ นุ โข อาวุโส ญาณ\-ทสฺสน\-วิสุทฺธิ อนุปาทาปรินิพฺพานนฺ”ติ อิติ ปุฏฺโฐ สมาโน “โน หิทํ อาวุโส”ติ วเทสิ. “กิํ ปนาวุโส อญฺญตฺร อิเมหิ ธมฺเมหิ อนุปาทาปรินิพฺพานนฺ”ติ อิติ ปุฏฺโฐ สมาโน “โน หิทํ อาวุโส”ติ วเทสิ. ยถากถํ ปนาวุโส อิมสฺส ภาสิตสฺส อตฺโถ ทฏฺฐพฺโพติ.}

\proseitem{258}{สีลวิสุทฺธิํ เจ อาวุโส ภควา อนุปาทา\-ปรินิพฺพานํ ปญฺญเปยฺย\csromansharedfootnote{62063648f198c541}{ปญฺญาเปสฺส (สี, สฺยา) เอวมญฺญตฺถปิ.}, สฺเอาปาทานํเยว สมานํ อนุปาทา\-ปรินิพฺพานํ ปญฺญเปยฺย\csromansharedfootnote{62063648f198c541}{ปญฺญาเปสฺส (สี, สฺยา) เอวมญฺญตฺถปิ.}. จิตฺตวิสุทฺธิํ เจ อาวุโส ภควา อนุปาทา\-ปรินิพฺพานํ ปญฺญเปยฺย, สฺเอาปาทานํเยว สมานํ อนุปาทา\-ปรินิพฺพานํ ปญฺญเปยฺย. ทิฏฺฐิ\-วิสุทฺธิํ เจ อาวุโส ภควา อนุปาทา\-ปรินิพฺพานํ ปญฺญเปยฺย, ส- อุปาทานํเยว สมานํ อนุปาทา\-ปรินิพฺพานํ ปญฺญเปยฺย. กงฺขา\-วิตรณ\-วิสุทฺธิํ เจ อาวุโส ภควา อนุปาทา\-ปรินิพฺพานํ ปญฺญเปยฺย, สฺเอาปาทานํเยว สมานํ อนุปาทา\-ปรินิพฺพานํ ปญฺญเปยฺย. มคฺคา\-มคฺค\-ญาณ\-ทสฺสน\-วิสุทฺธิํ เจ อาวุโส ภควา อนุปาทา\-ปรินิพฺพานํ ปญฺญเปยฺย, สฺเอาปาทานํเยว สมานํ อนุปาทา\-ปรินิพฺพานํ ปญฺญเปยฺย. ปฏิปทา\-ญาณ\-ทสฺสน\-วิสุทฺธิํ เจ อาวุโส ภควา อนุปาทา\-ปรินิพฺพานํ ปญฺญเปยฺย, สฺเอาปาทานํเยว สมานํ อนุปาทา\-ปรินิพฺพานํ ปญฺญเปยฺย. ญาณ\-ทสฺสน\-วิสุทฺธิํ เจ อาวุโส ภควา อนุปาทา\-ปรินิพฺพานํ ปญฺญเปยฺย,}

\prose{\csromanfolio{203}สฺเอาปาทานํเยว สมานํ อนุปาทา\-ปรินิพฺพานํ ปญฺญเปยฺย. อญฺญตฺร เจ อาวุโส อิเมหิ ธมฺเมหิ อนุปาทา\-ปรินิพฺพานํ อภวิสฺส, ปุถุชฺชโน ปรินิพฺพาเยยฺย. ปุถุชฺชโน หิ อาวุโส อญฺญตฺร อิเมหิ ธมฺเมหิ. เตน หาวุโส อุปมํ เต กริสฺสามิ, อุปมาย\-ปิ\-เธ\-กจฺเจ วิญฺญู ปุริสา ภาสิตสฺส อตฺถํ อาชานนฺติ.}

\proseitem{259}{เสยฺยถาปิ อาวุโส รญฺโญ ปเสนทิสฺส โกสลสฺส สาวตฺถิยํ ปฏิวสนฺตสฺส สาเกเต กิญฺจิเทว อจฺจายิกํ กรณียํ อุปฺปชฺเชยฺย. ตสฺส อนฺตรา จ สาวตฺถิํ อนฺตรา จ สาเกตํ สตฺต รถวินีตานิ อุปฏฺฐเปยฺยุํ. อถ โข อาวุโส ราชา ปเสนทิ โกสโล สาวตฺถิยา นิกฺขมิตฺวา อนฺเตปุรทฺวารา ปฐมํ รถวินีตํ อภิรุเหยฺย, ปฐเมน รถวินีเตน ทุติยํ รถวินีตํ ปาปุเณยฺย. ปฐมํ รถวินีตํ วิสฺสชฺเชยฺย, ทุติยํ รถวินีตํ อภิรุเหยฺย, ทุติเยน รถวินีเตน ตติยํ รถวินีตํ ปาปุเณยฺย. ทุติยํ รถวินีตํ วิสฺสชฺเชยฺย, ตติยํ รถวินีตํ อภิรุเหยฺย, ตติเยน รถวินีเตน จตุตฺถํ รถวินีตํ ปาปุเณยฺย. ตติยํ รถวินีตํ วิสฺสชฺเชยฺย, จตุตฺถํ รถวินีตํ อภิรุเหยฺย, จตุตฺเถน รถวินีเตน ปญฺจมํ รถวินีตํ ปาปุเณยฺย. จตุตฺถํ รถวินีตํ วิสฺสชฺเชยฺย, ปญฺจมํ รถวินีตํ อภิรุเหยฺย, ปญฺจเมน รถวินีเตน ฉฏฺฐํ รถวินีตํ ปาปุเณยฺย. ปญฺจมํ รถวินีตํ วิสฺสชฺเชยฺย, ฉฏฺฐํ รถวินีตํ อภิรุเหยฺย, ฉฏฺเฐน รถวินีเตน สตฺตมํ รถวินีตํ ปาปุเณยฺย. ฉฏฺฐํ รถวินีตํ วิสฺสชฺเชยฺย, สตฺตมํ รถวินีตํ อภิรุเหยฺย, สตฺตเมน รถวินีเตน สาเกตํ อนุปาปุเณยฺย อนฺเตปุรทฺวารํ. ตเมนํ อนฺเตปุร\-ทฺวารคตํ สมานํ มิตฺตามจฺจา ญาติสาโลหิตา เอวํ ปุจฺเฉยฺยุํ “อิมินา ตฺวํ มหาราช รถวินีเตน สาวตฺถิยา สาเกตํ อนุปฺปตฺโต อนฺเตปุรทฺวารนฺ”ติ. กถํ พฺยากรมาโน นุ โข อาวุโส ราชา ปเสนทิ โกสโล สมฺมา พฺยากรมาโน พฺยากเรยฺยาติ.}

\prose{เอวํ พฺยากรมาโน โข อาวุโส ราชา ปเสนทิ โกสโล สมฺมา พฺยากรมาโน พฺยากเรยฺย\csromandash{}“อิธ เม สาวตฺถิยํ ปฏิวสนฺตสฺส สาเกเต กิญฺจิเทว อจฺจายิกํ กรณียํ อุปฺปชฺชิ\csromansharedfootnote{8a72d0c4fbf55665}{อุปฺปชฺชติ (ก)}. ตสฺส เม อนฺตรา จ สาวตฺถิํ อนฺตรา จ สาเกตํ สตฺต รถวินีตานิ อุปฏฺฐเปสุํ. อถ ขฺวาหํ \csromanfolio{204}สาวตฺถิยา นิกฺขมิตฺวา อนฺเตปุรทฺวารา ปฐมํ รถวินีตํ อภิรุหิํ, ปฐเมน รถวินีเตน ทุติยํ รถวินีตํ ปาปุณิํ. ปฐมํ รถวินีตํ วิสฺสชฺชิํ, ทุติยํ รถวินีตํ อภิรุหิํ, ทุติเยน รถวินีเตน ตติยํ รถวินีตํ ปาปุณิํ. ทุติยํ รถวินีตํ วิสฺสชฺชิํ, ตติยํ รถวินีตํ อภิรุหิํ, ตติเยน รถวินีเตน จตุตฺถํ รถวินีตํ ปาปุณิํ. ตติยํ รถวินีตํ วิสฺสชฺชิํ, จตุตฺถํ รถวินีตํ อภิรุหิํ, จตุตฺเถน รถวินีเตน ปญฺจมํ รถวินีตํ ปาปุณิํ. จตุตฺถํ รถวินีตํ วิสฺสชฺชิํ, ปญฺจมํ รถวินีตํ อภิรุหิํ, ปญฺจเมน รถวินีเตน ฉฏฺฐํ รถวินีตํ ปาปุณิํ. ปญฺจมํ รถวินีตํ วิสฺสชฺชิํ, ฉฏฺฐํ รถวินีตํ อภิรุหิํ, ฉฏฺเฐน รถวินีเตน สตฺตมํ รถวินีตํ ปาปุณิํ. ฉฏฺฐํ รถวินีตํ วิสฺสชฺชิํ, สตฺตมํ รถวินีตํ อภิรุหิํ, สตฺตเมน รถวินีเตน สาเกตํ อนุปฺปตฺโต อนฺเตปุรทฺวารนฺ”ติ. เอวํ พฺยากรมาโน โข อาวุโส ราชา ปเสนทิ โกสโล สมฺมา พฺยากรมาโน พฺยากเรยฺยาติ.}

\prose{เอวเมว โข อาวุโส สีลวิสุทฺธิ ยาวเทว จิตฺต\-วิสุทฺธ\-ตฺถา. จิตฺตวิสุทฺธิ ยาวเทว ทิฏฺฐิ\-วิสุทฺธ\-ตฺถา. ทิฏฺฐิวิสุทฺธิ ยาวเทว กงฺขา\-วิตรณ\-วิสุทฺธ\-ตฺถา. กงฺขา\-วิตรณ\-วิสุทฺธิ ยาวเทว มคฺคามคฺคญฺจณทสฺสนวิสุทฺธตฺถา. มคฺคา\-มคฺค\-ญาณ\-ทสฺสน\-วิสุทฺธิ ยาวเทว ปฏิปทา\-ญาณ\-ทสฺสน\-วิสุทฺธ\-ตฺถา. ปฏิปทา\-ญาณ\-ทสฺสน\-วิสุทฺธิ ยาวเทว ญาณ\-ทสฺสน\-วิสุทฺธ\-ตฺถา. ญาณ\-ทสฺสน\-วิสุทฺธิ ยาวเทว อนุปาทา\-ปรินิพฺพาน\-ตฺถา. อนุปาทา\-ปรินิพฺพานตฺถํ โข อาวุโส ภควติ พฺรหฺมจริยํ วุสฺสตีติ.}

\proseitem{260}{เอวํ วุตฺเต อายสฺมา สาริปุตฺโต อายสฺมนฺตํ ปุณฺณํ มนฺตาณิปุตฺตํ เอตทโวจ “โกนาโม อายสฺมา, กถญฺจ ปนายสฺมนฺตํ สพฺรหฺมจารี ชานนฺตี”ติ. “ปุณฺโณ”ติ โข เม อาวุโส นามํ, “มนฺตาณิปุตฺโต”ติ จ ปน มํ สพฺรหฺมจารี ชานนฺตีติ. อจฺฉริยํ อาวุโส, อพฺภุตํ อาวุโส. ยถา ตํ สุตวตา สาวเกน สมฺมเทว สตฺถุสาสนํ อาชานนฺเตน, เอวเมว อายสฺมตา ปุณฺเณน มนฺตาณิปุตฺเตน คมฺภีรา คมฺภีรปญฺหา อนุมสฺส อนุมสฺส พฺยากตา. ลาภา สพฺรหฺมจารีนํ สุลทฺธลาภา สพฺรหฺมจารีนํ, เย อายสฺมนฺตํ ปุณฺณํ มนฺตาณิปุตฺตํ ลภนฺติ ทสฺสนาย, ลภนฺติ ปยิรูปาสนาย. เจลณฺฑุเกน\csromansharedfootnote{61bc1505f9b5c19d}{เจลณฺฑเกน (ก), เจลณฺฑุปเกน (?)} เจปิ สพฺรหฺมจารี อายสฺมนฺตํ ปุณฺณํ \csromanfolio{205}มนฺตาณิปุตฺตํ มุทฺธนา ปริหรนฺตา ลเภยฺยุํ ทสฺสนาย, ลเภยฺยุํ ปยิรูปาสนาย, เตสํปิ ลาภา เตสํปิ สุลทฺธํ. อมฺหากํปิ ลาภา อมฺหากํปิ สุลทฺธํ, เย มยํ อายสฺมนฺตํ ปุณฺณํ มนฺตาณิปุตฺตํ ลภาม ทสฺสนาย, ลภาม ปยิรูปาสนายาติ.}

\prose{เอวํ วุตฺเต อายสฺมา ปุณฺโณ มนฺตาณิปุตฺโต อายสฺมนฺตํ สาริปุตฺตํ เอตทโวจ “โกนาโม อายสฺมา, กถญฺจ ปนายสฺมนฺตํ สพฺรหฺมจารี ชานนฺตี”ติ. “อุปติสฺโส”ติ โข เม อาวุโส นามํ, “สาริปุตฺโต”ติ จ ปน มํ สพฺรหฺมจารี ชานนฺตีติ. สตฺถุกปฺเปน วต กิร โภ\csromansharedfootnote{0f71631349a151a1}{โข (ก)} สาวเกน สทฺธิํ มนฺตยมานา น ชานิมฺห “อายสฺมา สาริปุตฺโต”ติ. สเจ หิ มยํ ชาเนยฺยาม “อายสฺมา สาริปุตฺโต”ติ, เอตฺตกํปิ โน นปฺปฏิภาเสยฺย\csromansharedfootnote{984d5067d1c2aa43}{นปฺปฏิเภยฺย (?)}. อจฺฉริยํ อาวุโส, อพฺภุตํ อาวุโส. ยถา ตํ สุตวตา สาวเกน สมฺมเทว สตฺถุสาสนํ อาชานนฺเตน, เอวเมว อายสฺมตา สาริปุตฺเตน คมฺภีรา คมฺภีรปญฺหา อนุมสฺส อนุมสฺส ปุจฺฉิตา. ลาภา สพฺรหฺมจารีนํ สุลทฺธลาภา สพฺรหฺมจารีนํ, เย อายสฺมนฺตํ สาริปุตฺตํ ลภนฺติ ทสฺสนาย, ลภนฺติ ปยิรูปาสนาย. เจลณฺฑุเกน เจปิ สพฺรหฺมจารี อายสฺมนฺตํ สาริปุตฺตํ มุทฺธนา ปริหรนฺตา ลเภยฺยุํ ทสฺสนาย, ลเภยฺยุํ ปยิรูปาสนาย, เตสํปิ ลาภา เตสํปิ สุลทฺธํ. อมฺหากํปิ ลาภา อมฺหากํปิ สุลทฺธํ, เย มยํ อายสฺมนฺตํ สาริปุตฺตํ ลภาม ทสฺสนาย, ลภาม ปยิรูปาสนายาติ.}

\prosewithcloserruled{อิติห เต อุโภปิ มหานาคา อญฺญมญฺญสฺส สุภาสิตํ สมนุ\-โมทิํสูติ.}{รถวินีต\-สุตฺตํ นิฏฺฐิตํ จตุตฺถํ.}
\csromanmatikaanchor{mtk.55}
\csromantocmarkref{subsection}{5. นิวาปสุตฺต}{mtk.55}
\bookmark[dest={mtk.55},level=2]{5. นิวาปสุตฺต}
\csromanheader{2}{5. นิวาปสุตฺต}
\proseitem{261}{เอวํ เม สุตํ\csromandash{}เอกํ สมยํ ภควา สาวตฺถิยํ วิหรติ เชตวเน อนาถ\-ปิณฺฑิกสฺส อาราเม. ตตฺร โข ภควา ภิกฺขู อามนฺเตสิ “ภิกฺขโว”ติ. “ภทนฺเต”ติ เต ภิกฺขู ภควโต ปจฺจสฺโสสุํ. ภควา เอตทโวจ\csromandash{}}

\prose{\csromanfolio{206}น ภิกฺขเว เนวาปิโก นิวาปํ นิวปติ มิคชาตานํ “อิมํ เม นิวาปํ นิวุตฺตํ มิคชาตา ปริภุญฺชนฺตา ทีฆายุกา วณฺณวนฺโต จิรํ ทีฆมทฺธานํ ยาเปนฺตู”ติ. เอวญฺจ โข ภิกฺขเว เนวาปิโก นิวาปํ นิวปติ มิคชาตานํ “อิมํ เม นิวาปํ นิวุตฺตํ มิคชาตา อนุปขชฺช มุจฺฉิตา โภชนานิ ภุญฺชิสฺสนฺติ, อนุปขชฺช มุจฺฉิตา โภชนานิ ภุญฺชมานา มทํ อาปชฺชิสฺสนฺติ, มตฺตา สมานา ปมาทํ อาปชฺชิสฺสนฺติ, ปมตฺตา สมานา ยถา\-กาม\-กรณียา ภวิสฺสนฺติ อิมสฺมิํ นิวาเป”ติ.}

\proseitem{262}{ตตฺร ภิกฺขเว ปฐมา มิคชาตา อมุํ นิวาปํ นิวุตฺตํ เนวาปิกสฺส อนุปขชฺช มุจฺฉิตา โภชนานิ ภุญฺชิํสุ. เต ตตฺถ อนุปขชฺช มุจฺฉิตา โภชนานิ ภุญฺชมานา มทํ อาปชฺชิํสุ, มตฺตา สมานา ปมาทํ อาปชฺชิํสุ, ปมตฺตา สมานา ยถา\-กาม\-กรณียา อเหสุํ เนวาปิกสฺส อมุสฺมิํ นิวาเป. เอวํ หิ เต ภิกฺขเว ปฐมา มิคชาตา น ปริมุจฺจิํสุ เนวาปิกสฺส อิทฺธานุภาวา.}

\proseitem{263}{ตตฺร ภิกฺขเว ทุติยา มิคชาตา เอวํ สมจินฺเตสุํ “เย โข เต ปฐมา มิคชาตา อมุํ นิวาปํ นิวุตฺตํ เนวาปิกสฺส อนุปขชฺช มุจฺฉิตา โภชนานิ ภุญฺชิํสุ, เต ตตฺถ อนุปขชฺช มุจฺฉิตา โภชนานิ ภุญฺชมานา มทํ อาปชฺชิํสุ, มตฺตา สมานา ปมาทํ อาปชฺชิํสุ, ปมตฺตา สมานา ยถา\-กาม\-กรณียา อเหสุํ เนวาปิกสฺส อมุสฺมิํ นิวาเป. เอวํ หิ เต ปฐมา มิคชาตา น ปริมุจฺจิํสุ เนวาปิกสฺส อิทฺธานุภาวา. ยํนูน มยํ สพฺพโส นิวาปโภชนา ปฏิวิรเมยฺยาม, ภยโภคา ปฏิวิรตา อรญฺญายตนานิ อชฺโฌคาเหตฺวา วิหเรยฺยามา”ติ. เต สพฺพโส นิวาปโภชนา ปฏิวิรมิํสุ, ภยโภคา ปฏิวิรตา อรญฺญายตนานิ อชฺโฌคาเหตฺวา วิหริํสุ. เตสํ คิมฺหานํ ปจฺฉิเม มาเส ติโณทก\-สงฺขเย อธิมตฺต\-กสิมานํ ปตฺโต กาโย โหติ. เตสํ อธิมตฺต\-กสิมานํ ปตฺตกายานํ พลวีริยํ ปริหายิ. พลวีริเย ปริหีเน ตเมว นิวาปํ นิวุตฺตํ เนวาปิกสฺส ปจฺจาคมิํสุ. เต ตตฺถ อนุปขชฺช มุจฺฉิตา โภชนานิ ภุญฺชิํสุ. เต ตตฺถ อนุปขชฺช มุจฺฉิตา โภชนานิ ภุญฺชมานา มทํ อาปชฺชิํสุ, มตฺตา สมานา ปมาทํ อาปชฺชิํสุ, ปมตฺตา สมานา ยถา\-กาม\-กรณียา อเหสุํ เนวาปิกสฺส อมุสฺมิํ นิวาเป. เอวํ หิ เต ภิกฺขเว ทุติยาปิ มิคชาตา น ปริมุจฺจิํสุ เนวาปิกสฺส อิทฺธานุภาวา.}

\proseitem{264}{\csromanfolio{207}ตตฺร ภิกฺขเว ตติยา มิคชาตา เอวํ สมจินฺเตสุํ “เย โข เต ปฐมา มิคชาตา อมุํ นิวาปํ นิวุตฺตํ เนวาปิกสฺส ฯเปฯ. เอวํ หิ เต ปฐมา มิคชาตา น ปริมุจฺจิํสุ เนวาปิกสฺส อิทฺธานุภาวา. เยปิ เต ทุติยา มิคชาตา เอวํ สมจินฺเตสุํ ‘เย โข เต ปฐมา มิคชาตา อมุํ นิวาปํ นิวุตฺตํ เนวาปิกสฺส ฯเปฯ. เอวํ หิ เต ปฐมา มิคชาตา น ปริมุจฺจิํสุ เนวาปิกสฺส อิทฺธานุภาวา. ยํนูน มยํ สพฺพโส นิวาปโภชนา ปฏิวิรเมยฺยาม, ภยโภคา ปฏิวิรตา อรญฺญายตนานิ อชฺโฌคาเหตฺวา วิหเรยฺยามา’ติ. เต สพฺพโส นิวาปโภชนา ปฏิวิรมิํสุ, ภยโภคา ปฏิวิรตา อรญฺญายตนานิ อชฺโฌคาเหตฺวา วิหริํสุ. เตสํ คิมฺหานํ ปจฺฉิเม มาเส ติโณทก\-สงฺขเย อธิมตฺต\-กสิมานํ ปตฺโต กาโย โหติ. เตสํ อธิมตฺต\-กสิมานํ ปตฺตกายานํ พลวีริยํ ปริหายิ. พลวีริเย ปริหีเน ตเมว นิวาปํ นิวุตฺตํ เนวาปิกสฺส ปจฺจาคมิํสุ. เต ตตฺถ อนุปขชฺช มุจฺฉิตา โภชนานิ ภุญฺชิํสุ. เต ตตฺถ อนุปขชฺช มุจฺฉิตา โภชนานิ ภุญฺชมานา มทํ อาปชฺชิํสุ, มตฺตา สมานา ปมาทํ อาปชฺชิํสุ, ปมตฺตา สมานา ยถา\-กาม\-กรณียา อเหสุํ เนวาปิกสฺส อมุสฺมิํ นิวาเป. เอวํ หิ เต ทุติยาปิ มิคชาตา น ปริมุจฺจิํสุ เนวาปิกสฺส อิทฺธานุภาวา. ยํนูน มยํ อมุํ นิวาปํ นิวุตฺตํ เนวาปิกสฺส อุปนิสฺสาย อาสยํ กปฺเปยฺยาม, ตตฺราสยํ กปฺเปตฺวา อมุํ นิวาปํ นิวุตฺตํ เนวาปิกสฺส อนนุปขชฺช อมุจฺฉิตา โภชนานิ ภุญฺชิสฺสาม, อนนุปขชฺช อมุจฺฉิตา โภชนานิ ภุญฺชมานา น มทํ อาปชฺชิสฺสาม. อมตฺตา สมานา น ปมาทํ อาปชฺชิสฺสาม, อปฺปมตฺตา สมานา น ยถา\-กาม\-กรณียา ภวิสฺสาม เนวาปิกสฺส อมุสฺมิํ นิวาเป”ติ. เต อมุํ นิวาปํ นิวุตฺตํ เนวาปิกสฺส อุปนิสฺสาย อาสยํ กปฺปยิํสุ. ตตฺราสยํ กปฺเปตฺวา อมุํ นิวาปํ นิวุตฺตํ เนวาปิกสฺส อนนุปขชฺช อมุจฺฉิตา โภชนานิ ภุญฺชิํสุ. เต ตตฺถ อนนุปขชฺช อมุจฺฉิตา โภชนานิ ภุญฺชมานา น มทํ อาปชฺชิํสุ, อมตฺตา สมานา น ปมาทํ อาปชฺชิํสุ, อปฺปมตฺตา สมานา น ยถา\-กาม\-กรณียา อเหสุํ เนวาปิกสฺส อมุสฺมิํ นิวาเป.}

\prose{ตตฺร ภิกฺขเว เนวาปิกสฺส จ เนวาปิก\-ปริสาย จ เอตทโหสิ “สฐาสฺสุนามิเม ตติยา มิคชาตา เกตพิโน, อิทฺธิมนฺตา\-สฺสุ\-นามิเม ตติยา มิคชาตา ปรชนา. อิมญฺจ นาม นิวาปํ นิวุตฺตํ ปริภุญฺชนฺติ, น จ \csromanfolio{208}เนสํ ชานาม อาคติํ วา คติํ วา. ยํนูน มยํ อิมํ นิวาปํ นิวุตฺตํ มหตีหิ ทณฺฑวากราหิ\csromansharedfootnote{879de0c3eb2920b4}{ทณฺฑวาคุราหิ (สฺยา)} สมนฺตา สปฺปเทสํ อนุปริวาเรยฺยาม อปฺเปว นาม ตติยานํ มิคชาตานํ อาสยํ ปสฺเสยฺยาม, ยตฺถ เต คาหํ คจฺเฉยฺยุนฺ”ติ, เต อมุํ นิวาปํ นิวุตฺตํ มหตีหิ ทณฺฑวากราหิ สมนฺตา สปฺปเทสํ อนุปริวาเรสุํ. อทฺทสํสุ โข ภิกฺขเว เนวาปิโก จ เนวาปิกปริสา จ ตติยานํ มิคชาตานํ อาสยํ, ยตฺถ เต คาหํ อคมํสุ. เอวํ หิ เต ภิกฺขเว ตติยาปิ มิคชาตา น ปริมุจฺจิํสุ เนวาปิกสฺส อิทฺธานุภาวา.}

\proseitem{265}{ตตฺร ภิกฺขเว จตุตฺถา มิคชาตา เอวํ สมจินฺเตสุํ\csromandash{}เย โข เต ปฐมา มิคชาตา ฯเปฯ. เอวํ หิ เต ปฐามา มิคชาตา น ปริมุจฺจิํสุ เนวาปิกสฺส อิทฺธานุภาวา. เยปิ เต ทุติยา มิคชาตา เอวํ สมจินฺเตสุํ “เย โข เต ปฐมา มิคชาตา ฯเปฯ. เอวํ หิ เต ปฐมา มิคชาตา น ปริมุจฺจิํสุ เนวาปิกสฺส อิทฺธานุภาวา. ยํนูน มยํ สพฺพโส นิวาปโภชนา ปฏิวิรเมยฺยาม, ภยโภคา ปฏิวิรตา อรญฺญายตนานิ อชฺโฌคาเหตฺวา วิหเรยฺยามา”ติ. เต สพฺพโส นิวาปโภชนา ปฏิวิรมิํสุ ฯเปฯ. เอวํ หิ เต ทุติยาปิ มิคชาตา น ปริมุจฺจิํสุ เนวาปิกสฺส อิทฺธานุภาวา. เยปิ เต ตติยา มิคชาตา เอวํ สมจินฺเตสุํ “เย โข เต ปฐมา มิคชาตา ฯเปฯ. เอวํ หิ เต ปฐมา มิคชาตา น ปริมุจฺจิํสุ เนวาปิกสฺส อิทฺธานุภาวา. เยปิ เต ทุติยา มิคชาตา เอวํ สมจินฺเตสุํ ‘เย โข เต ปฐมา มิคชาตา ฯเปฯ. เอวํ หิ เต ปฐมา มิคชาตา น ปริมุจฺจิํสุ เนวาปิกสฺส อิทฺธานุภาวา. ยํนูน มยํ สพฺพโส นิวาปโภชนา ปฏิวิรเมยฺยาม, ภยโภคา ปฏิวิรตา อรญฺญายตนานิ อชฺโฌคาเหตฺวา วิหเรยฺยามา’ติ. เต สพฺพโส นิวาปโภชนา ปฏิวิรมิํสุ ฯเปฯ. เอวํ หิ เต ทุติยาปิ มิคชาตา น ปริมุจฺจิํสุ เนวาปิกสฺส อิทฺธานุภาวา. ยํนูน มยํ อมุํ นิวาปํ นิวุตฺตํ เนวาปิกสฺส อุปนิสฺสาย อาสยํ กปฺเปยฺยาม, ตตฺราสยํ กปฺเปตฺวา อมุํ นิวาปํ นิวุตฺตํ เนวาปิกสฺส อนนุปขชฺช อมุจฺฉิตา โภชนานิ ภุญฺชิสฺสาม, อนนุปขชฺช อมุจฺฉิตา โภชนานิ ภุญฺชมานา น มทํ อาปชฺชิสฺสาม, อมตฺตา สมานา น ปมาทํ อาปชฺชิสฺสาม, อปฺปมตฺตา สมานา น ยถา\-กาม\-กรณียา ภวิสฺสาม \csromanfolio{209}เนวาปิกสฺส อมุสฺมิํ นิวาเป”ติ. เต อมุํ นิวาปํ นิวุตฺตํ เนวาปิกสฺส อุปนิสฺสาย อาสยํ กปฺปยิํสุ, ตตฺราสยํ กปฺเปตฺวา อมุํ นิวาปํ นิวุตฺตํ เนวาปิกสฺส อนนุปขชฺช อมุจฺฉิตา โภชนานิ ภุญฺชิํสุ, เต ตตฺถ อนนุปขชฺช อมุจฺฉิตา โภชนานิ ภุญฺชมานา น มทํ อาปชฺชิํสุ, อมตฺตา สมานา น ปมาทํ อาปชฺชิํสุ, อปฺปมตฺตา สมานา น ยถา\-กาม\-กรณียา อเหสุํ เนวาปิกสฺส อมุสฺมิํ นิวาเป. ตตฺร เนวาปิกสฺส จ เนวาปิก\-ปริสาย จ เอตทโหสิ “สฐาสฺสุนามิเม ตติยา มิคชาตา เกตพิโน, อิทฺธิมนฺตา\-สฺสุ\-นามิเม ตติยา มิคชาตา ปรชนา, อิมญฺจ นาม นิวาปํ นิวุตฺตํ ปริภุญฺชนฺติ. น จ เนสํ ชานาม อาคติํ วา คติํ วา. ยํนูน มยํ อิมํ นิวาปํ นิวุตฺตํ มหตีหิ ทณฺฑวากราหิ สมนฺตา สปฺปเทสํ อนุปริวาเรยฺยาม, อปฺเปว นาม ตติยานํ มิคชาตานํ อาสยํ ปสฺเสยฺยาม, ยตฺถ เต คาหํ คจฺเฉยฺยุนฺ”ติ. เต อมุํ นิวาปํ นิวุตฺตํ มหตีหิ ทณฺฑวากราหิ สมนฺตา สปฺปเทสํ อนุปริวาเรสุํ. อทฺทสํสุ โข เนวาปิโก จ เนวาปิกปริสา จ ตติยานํ มิคชาตานํ อาสยํ, ยตฺถ เต คาหํ อคมํสุ. เอวํ หิ เต ตติยาปิ มิคชาตา น ปริมุจฺจิํสุ เนวาปิกสฺส อิทฺธานุภาวา. ยํนูน มยํ ยตฺถ อคติ เนวาปิกสฺส จ เนวาปิก\-ปริสาย จ, ตตฺราสยํ กปฺเปยฺยาม, ตตฺราสยํ กปฺเปตฺวา อมุํ นิวาปํ นิวุตฺตํ เนวาปิกสฺส อนนุปขชฺช อมุจฺฉิตา โภชนานิ ภุญฺชิสฺสาม, อนนุปขชฺช อมุจฺฉิตา โภชนานิ ภุญฺชมานา น มทํ อาปชฺชิสฺสาม, อมตฺตา สมานา น ปมาทํ อาปชฺชิสฺสาม, อปฺปมตฺตา สมานา น ยถา\-กาม\-กรณียา ภวิสฺสาม เนวาปิกสฺส อมุสฺมิํ นิวาเปติ. เต ยตฺถ อคติ เนวาปิกสฺส จ เนวาปิก\-ปริสาย จ, ตตฺราสยํ กปฺปยิํสุ, ตตฺราสยํ กปฺเปตฺวา อมุํ นิวาปํ นิวุตฺตํ เนวาปิกสฺส อนนุปขชฺช อมุจฺฉิตา โภชนานิ ภุญฺชิํสุ, เต ตตฺถ อนนุปขชฺช อมุจฺฉิตา โภชนานิ ภุญฺชมานา น มทํ อาปชฺชิํสุ, อมตฺตา สมานา น ปมาทํ อาปชฺชิํสุ, อปฺปมตฺตา สมานา น ยถา\-กาม\-กรณียา อเหสุํ เนวาปิกสฺส อมุสฺมิํ นิวาเป.}

\prose{ตตฺร ภิกฺขเว เนวาปิกสฺส จ เนวาปิก\-ปริสาย จ เอตทโหสิ “สฐาสฺสุนามิเม จตุตฺถา มิคชาตา เกตพิโน, อิทฺธิมนฺตา\-สฺสุ\-นามิเม จตุตฺถา มิคชาตา ปรชนา, อิมญฺจ นาม นิวาปํ นิวุตฺตํ ปริภุญฺชนฺติ. น จ เนสํ ชานาม อาคติํ วา คติํวา. ยํนูน มยํ อิมํ นิวาปํ นิวุตฺตํ \csromanfolio{210}มหตีหิ ทณฺฑวากราหิ สมนฺตา สปฺปเทสํ อนุปริวาเรยฺยาม, อปฺเปว นาม จตุตฺถานํ มิคชาตานํ อาสยํ ปสฺเสยฺยาม, ยตฺถ เต คาหํ คจฺเฉยฺยุนฺ”ติ. เต อมุํ นิวาปํ นิวุตฺตํ มหตีหิ ทณฺฑวากราหิ สมนฺตา สปฺปเทสํ อนุปริวาเรสุํ. เนว โข ภิกฺขเว อทฺทสํสุ เนวาปิโก จ เนวาปิกปริสา จ จตุตฺถานํ มิคชาตานํ อาสยํ, ยตฺถ เต คาหํ คจฺเฉยฺยุํ. ตตฺร ภิกฺขเว เนวาปิกสฺส จ เนวาปิก\-ปริสาย จ เอตทโหสิ “สเจ โข มยํ จตุตฺเถ มิคชาเต ฆฏฺเฏสฺสาม. เต ฆฏฺฏิตา อญฺเญ ฆฏฺฏิสฺสนฺติ. เต ฆฏฺฏิตา อญฺเญ ฆฏฺฏิสฺสนฺติ. เอวํ อิมํ นิวาปํ นิวุตฺตํ สพฺพโส มิคชาตา ปริมุญฺจิสฺสนฺติ. ยํนูน มยํ จตุตฺเถ มิคชาเต อชฺฌุเปกฺเขยฺยามา”ติ. อชฺฌุเปกฺขิํสุ โข ภิกฺขเว เนวาปิโก จ เนวาปิกปริสา จ จตุตฺเถ มิคชาเต. เอวํ หิ เต ภิกฺขเว จตุตฺถา มิคชาตา ปริมุจฺจิํสุ เนวาปิกสฺส อิทฺธานุภาวา.}

\proseitem{266}{อุปมา โข เม อยํ ภิกฺขเว กตา อตฺถสฺส วิญฺญาปนาย. อยํ เจเวตฺถ อตฺโถ. “นิวาโป”ติ โข ภิกฺขเว ปญฺจนฺเนตํ กามคุณานํ อธิวจนํ. “เนวาปิโก”ติ โข ภิกฺขเว มารสฺเสตํ ปาปิมโต อธิวจนํ. “เนวาปิกปริสา”ติ โข ภิกฺขเว มารปริสาเยตํ อธิวจนํ. “มิคชาตา”ติ โข ภิกฺขเว สมณพฺราหฺมณานเมตํ อธิวจนํ.}

\proseitem{267}{ตตฺร ภิกฺขเว ปฐมา สมณพฺราหฺมณา อมุํ นิวาปํ นิวุตฺตํ มารสฺส อมูนิ จ โลกามิสานิ อนุปขชฺช มุจฺฉิตา โภชนานิ ภุญฺชิํสุ, เต ตตฺถ อนุปขชฺช มุจฺฉิตา โภชนานิ ภุญฺชมานา มทํ อาปชฺชิํสุ, มตฺตา สมานา ปมาทํ อาปชฺชิํสุ, ปมตฺตา สมานา ยถา\-กาม\-กรณียา อเหสุํ มารสฺส อมุสฺมิํ นิวาเป, อมุสฺมิํ จ โลกามิเส. เอวํ หิ เต ภิกฺขเว ปฐมา สมณพฺราหฺมณา น ปริมุจฺจิํสุ มารสฺส อิทฺธานุภาวา. เสยฺยถาปิ เต ภิกฺขเว ปฐมา มิคชาตา, ตถูปเม อหํ อิเม ปฐเม สมณพฺราหฺมเณ วทามิ.}

\proseitem{268}{ตตฺร ภิกฺขเว ทุติยา สมณพฺราหฺมณา เอวํ สมจินฺเตสุํ “เย โข เต ปฐมา สมณพฺราหฺมณา อมุํ นิวาปํ นิวุตฺตํ มารสฺส อมูนิ จ โลกามิสานิ อนุปขชฺช มุจฺฉิตา โภชนานิ ภุญฺชิํสุ. เต ตตฺถ อนุปขชฺช มุจฺฉิตา โภชนานิ ภุญฺชมานา มทํ อาปชฺชิํสุ, มตฺตา สมานา \csromanfolio{211}ปมาทํ อาปชฺชิํสุ, ปมตฺตา สมานา ยถา\-กาม\-กรณียา อเหสุํ มารสฺส อมุสฺมิํ นิวาเป, อมุสฺมิํ จ โลกามิเส. เอวํ หิ เต ปฐมา สมณพฺราหฺมณา น ปริมุจฺจิํสุ มารสฺส อิทฺธานุภาวา. ยํนูน มยํ สพฺพโส นิวาปโภชนา โลกามิสา ปฏิวิรเมยฺยาม, ภยโภคา ปฏิวิรตา อรญฺญยตนานิ อชฺโฌคาเหตฺวา วิหเรยฺยามา”ติ. เต สพฺพโส นิวาปโภชนา โลกามิสา ปฏิวิรมิํสุ, ภยโภคา ปฏิวิรตา อรญฺญายตนานิ อชฺโฌคาเหตฺวา วิหริํสุ. เต ตตฺถ สากภกฺขาปิ อเหสุํ, สามากภกฺขาปิ อเหสุํ, นีวารภกฺขาปิ อเหสุํ, ททฺทุลภกฺขาปิ อเหสุํ, หฏภกฺขาปิ อเหสุํ, กณภกฺขาปิ อเหสุํ, อาจามภกฺขาปิ อเหสุํ, ปิญฺญากภกฺขาปิ อเหสุํ, ติณภกฺขาปิ อเหสุํ, โคมยภกฺขาปิ อเหสุํ, วนมูล\-ผลา\-หารา ยาเปสุํ ปวตฺต\-ผล\-โภชี. เตสํ คิมฺหานํ ปจฺฉิเม มาเส ติโณทก\-สงฺขเย อธิมตฺต\-กสิมานํ ปตฺโต กาโย โหติ, เตสํ อธิมตฺต\-กสิมานํ ปตฺตกายานํ พลวีริยํ ปริหายิ, พลวีริเย ปริหีเน เจโตวิมุตฺติ ปริหายิ, เจโตวิมุตฺติยา ปริหีนาย ตเมว นิวาปํ นิวุตฺตํ มารสฺส ปจฺจาคมิํสุ ตานิ จ โลกามิสานิ. เต ตตฺถ อนุปขชฺช มุจฺฉิตา โภชนานิ ภุญฺชิํสุ, เต ตตฺถ อนุปขชฺช มุจฺฉิตา โภชนานิ ภุญฺชมานา มทํ อาปชฺชิํสุ, มตฺตา สมานา ปมาทํ อาปชฺชิํสุ, ปมตฺตา สมานา ยถา\-กาม\-กรณียา อเหสุํ มารสฺส อมุสฺมิํ นิวาเป, อมุสฺมิํ จ โลกามิเส. เอวํ หิ เต ภิกฺขเว ทุติยาปิ สมณพฺราหฺมณา น ปริมุจฺจิํสุ มารสฺส อิทฺธานุภาวา. เสยฺยถาปิ เต ภิกฺขเว ทุติยา มิคชาตา, ตถูปเม อหํ อิเม ทุติเย สมณพฺราหฺมเณ}

\proseitem{269}{ตตฺร ภิกฺขเว ตติยา สมณพฺราหฺมณา เอวํ สมจินฺเตสุํ “เย โข เต ปฐมา สมณพฺราหฺมณา อมุํ นิวาปํ นิวาตฺตํ มารสฺส อมูนิ จ โลกามิสานิ ฯเปฯ. เอวํ หิ เต ปฐมา สมณพฺราหฺมณา น ปริมุจฺจิํสุ มารสฺส อิทฺธานุภาวา. เยปิ เต ทุติยา สมณพฺราหฺมณา เอวํ สมจินฺเตสุํ ‘เย โข เต ปฐมา สมณพฺราหฺมณา อมุํ นิวาปํ นิวุตฺตํ มารสฺส อมูนิ จ โลกามิสานิ ฯเปฯ. เอวํ หิ เต ปฐมา สมณพฺราหฺมณา น ปริมุจฺจิํสุ มารสฺส อิทฺธานุภาวา. ยํนูน มยํ สพฺพโส นิวาปโภชนา โลกามิสา ปฏิวิรเมยฺยาม, ภยโภคา ปฏิวิรตา อรญฺญายตนานิ อชฺโฌคาเหตฺวา วิหเรยฺยามา’ติ. เต สพฺพโส นิวาปโภชนา \csromanfolio{212}โลกามิสา ปฏิวิรมิํสุ, ภยโภคา ปฏิวิรตา อรญฺญายตนานิ อชฺโฌคาเหตฺวา วิหริํสุ. เต ตตฺถ สากภกฺขาปิ อเหสุํ ฯเปฯ ปวตฺต\-ผล\-โภชี. เตสํ คิมฺหานํ ปจฺฉิเม มาเส ติโณทก\-สงฺขเย อธิมตฺต\-กสิมานํ ปตฺโต กาโย โหติ, เตสํ อธิมตฺต\-กสิมานํ ปตฺตกายานํ พลวีริยํ ปริหายิ, พลวีริเย ปริหีเน เจโตวิมุตฺติ ปริหายิ, เจโตวิมุตฺติยา ปริหีนาย ตเมว นิวาปํ นิวุตฺตํ มารสฺส ปจฺจาคมิํสุ ตานิ จ โลกามิสานิ. เต ตตฺถ อนุปขชฺช มุจฺฉิตา โภชนานิ ภุญฺชิํสุ, เต ตตฺถ อนุปขชฺช มุจฺฉิตา โภชนานิ ภุญฺชมานา มทํ อาปชฺชิํสุ, มตฺตา สมานา ปมาทํ อาปชฺชิํสุ, ปมตฺตา สมานา ยถา\-กาม\-กรณียา อเหสุํ มารสฺส อมุสฺมิํ นิวาเป อมุสฺมิํ จ โลกามิเส. เอวํ หิ เต ทุติยาปิ สมณพฺราหฺมณา น ปริมุจฺจิํสุ มารสฺส อิทฺธานุภาวา. ยํนูน มยํ อมุํ นิวาปํ นิวุตฺตํ มารสฺส อมูนิ จ โลกามิสานิ อุปนิสฺสาย อาสยํ กปฺเปยฺยาม, ตตฺราสยํ กปฺเปตฺวา อมุํ นิวาปํ นิวุตฺตํ มารสฺส อมูนิ จ โลกามิสานิ อนนุปขชฺช อมุจฺฉิตา โภชนานิ ภุญฺชิสฺสาม, อนนุปขชฺช อมุจฺฉิตา โภชนานิ ภุญฺชมานา น มทํ อาปชฺชิสฺสาม, อมตฺตา สมานา น ปมาทํ อาปชฺชิสฺสาม, อปฺปมตฺตา สมานา น ยถา\-กาม\-กรณียา ภวิสฺสาม มารสฺส อมุสฺมิํ นิวาเป อมุสฺมิํ จ โลกามิเส”ติ.}

\prose{เต อมุํ นิวาปํ นิวุตฺตํ มารสฺส อมูนิ จ โลกามิสานิ อุปนิสฺสาย อาสยํ กปฺปยิํสุ, ตตฺราสยํ กปฺเปตฺวา อมุํ นิวาปํ นิวุตฺตํ มารสฺส อมูนิ จ โลกามิสานิ อนนุปขชฺช อมุจฺฉิตา โภชนานิ ภุญฺชิํสุ, เต ตตฺถ อนนุปขชฺช อมุจฺฉิตา โภชนานิ ภุญฺชมานา น มทํ อาปชฺชิํสุ, อมตฺตา สมานา น ปมาทํ อาปชฺชิํสุ, อปฺปมตฺตา สมานา น ยถา\-กาม\-กรณียา อเหสุํ มารสฺส อมุสฺมิํ นิวาเป อมุสฺมิํ จ โลกามิเส. อปิ จ โข เอวํทิฏฺฐิกา อเหสุํ “สสฺสโต โลโก” อิติปิ, “อสสฺสโต โลโก” อิติปิ, “อนฺตวา โลโก” อิติปิ, “อนนฺตวา โลโก” อิติปิ, “ตํ ชีวํ ตํ สรีรํ” อิติปิ, “อญฺญํ ชีวํ อญฺญํ สรีรํ” อิติปิ, “โหติ ตถาคโต ปรํ มรณา” อิติปิ, “น โหติ ตถาคโต ปรํ มรณา” อิติปิ, “โหติ จ น จ โหติ ตถาคโต ปรํ มรณา” อิติปิ, “เนว โหติ น น โหติ ตถาคโต ปรํ มรณา” อิติปิ. เอวํ หิ เต ภิกฺขเว ตติยาปิ \csromanfolio{213}สมณพฺราหฺมณา น ปริมุจฺจิํสุ มารสฺส อิทฺธานุภาวา. เสยฺยถาปิ เต ภิกฺขเว ตติยา มิคชาตา, ตถูปเม อหํ อิเม ตติเย สมณพฺราหฺมเณ}

\proseitem{270}{ตตฺร ภิกฺขเว จตุตฺถา สมณพฺราหฺมณา เอวํ สมจินฺเตสุํ\csromandash{} เย โข เต ปฐมา สมณพฺราหฺมณา อมุํ นิวาปํ นิวุตฺตํ มารสฺส ฯเปฯ. เอวํ หิ เต ปฐมา สมณพฺราหฺมณา น ปริมุจฺจิํสุ มารสฺส อิทฺธานุภาวา. เยปิ เต ทุติยา สมณพฺราหฺมณา เอวํ สมจินฺเตสุํ “เย โข เต ปฐมา สมณพฺราหฺมณา ฯเปฯ. เอวํ หิ เต ปฐมา สมณพฺราหฺมณา น ปริมุจฺจิํสุ มารสฺส อิทฺธานุภาวา. ยํนูน มยํ สพฺพโส นิวาปโภชนา โลกามิสา ปฏิวิรเมยฺยาม, ภยโภคา ปฏิวิรตา อรญฺญายตนานิ อชฺโฌคาเหตฺวา วิหเรยฺยามา”ติ. เต สพฺพโส นิวาปโภชนา โลกามิสา ปฏิวิรมิํสุ ฯเปฯ. เอวํ หิ เต ทุติยาปิ สมณพฺราหฺมณา น ปริมุจฺจิํสุ มารสฺส อิทฺธานุภาวา. เยปิ เต ตติยา สมณพฺราหฺมณา เอวํ สมจินฺเตสุํ “เย โข เต ปฐมา สมณพฺราหฺมณา ฯเปฯ. เอวํ หิ เต ปฐมา สมณพฺราหฺมณา น ปริมุจฺจิํสุ มารสฺส อิทฺธานุภาวา. เยปิ เต ทุติยา สมณพฺราหฺมณา เอวํ สมจินฺเตสุํ ‘เย โข เต ปฐมา สมณพฺราหฺมณา ฯเปฯ. เอวํ หิ เต ปฐมา สมณพฺราหฺมณา น ปริมุจฺจิํสุ มารสฺส อิทฺธานุภาวา. ยํนูน มยํ สพฺพโส นิวาปโภชนา โลกามิสา ปฏิวิรเมยฺยาม, ภยโภคา ปฏิวิรตา อรญฺญายตนานิ อชฺโฌคาเหตฺวา วิหเรยฺยามา’ติ. เต สพฺพโส นิวาปโภชนา โลกามิสา ปฏิวิรมิํสุ ฯเปฯ. เอวํ หิ เต ทุติยาปิ สมณพฺราหฺมณา น ปริมุจฺจิํสุ มารสฺส อิทฺธานุภาวา. ยํนูน มยํ อมุํ นิวาปํ นิวุตฺตํ มารสฺส อมูนิ จ โลกามิสานิ อุปนิสฺสาย อาสยํ กปฺเปยฺยาม, ตตฺราสยํ กปฺเปตฺวา อมุํ นิวาปํ นิวุตฺตํ มารสฺส อมูนิ จ โลกามิสานิ อนนุปขชฺช อมุจฺฉิตา โภชนานิ ภุญฺชิสฺสาม, อนนุปขชฺช อมุจฺฉิตา โภชนานิ ภุญฺชมานา น มทํ อาปชฺชิสฺสาม, อมตฺตา สมานา น ปมาทํ อาปชฺชิสฺสาม, อปฺปมตฺตา สมานา น ยถา\-กาม\-กรณียา ภวิสฺสาม มารสฺส อมุสฺมิํ นิวาเป อมุสฺมิํ จ โลกามิเส”ติ. เต อมุํ นิวาปํ นิวุตฺตํ มารสฺส อมูนิ จ โลกามิสานิ อุปนิสฺสาย อาสยํ กปฺปยิํสุ, ตตฺราสยํ กปฺเปตฺวา อมุํ นิวาปํ นิวุตฺตํ มารสฺส อมูนิ จ โลกามิสานิ อนนุปขชฺช อมุจฺฉิตา โภชนานิ ภุญฺชิํสุ, เต ตตฺถ อนนุปขชฺช อมุจฺฉิตา โภชนานิ ภุญฺชมานา น มทํ อาปชฺชิํสุ, อมตฺตา \csromanfolio{214}สมานา น ปมาทํ อาปชฺชิํสุ, อปฺปมตฺตา สมานา น ยถา\-กาม\-กรณียา อเหสุํ มารสฺส อมุสฺมิํ นิวาเป อมุสฺมิํ จ โลกามิเส. อปิ จ โข เอวํทิฏฺฐิกา อเหสุํ “สสฺสโต โลโก” อิติปิ ฯเปฯ “เนว โหติ น น โหติ ตถาคโต ปรํ มรณา” อิติปิ. เอวํ หิ เต ตติยาปิ สมณพฺราหฺมณา น ปริมุจฺจิํสุ มารสฺส อิทฺธานุภาวา. ยํนูน มยํ ยตฺถ อคติ มารสฺส จ มารปริสาย จ, ตตฺราสยํ กปฺเปยฺยาม, ตตฺราสยํ กปฺเปตฺวา อมุํ นิวาปํ นิวุตฺตํ มารสฺส อมูนิ จ โลกามิสานิ อนนุปขชฺช อมุจฺฉิตา โภชนานิ ภุญฺชิสฺสาม, อนนุปขชฺช อมุจฺฉิตา โภชนานิ ภุญฺชมานา น มทํ อาปชฺชิสฺสาม, อมตฺตา สมานา น ปมาทํ อาปชฺชิสฺสาม, อปฺปมตฺตา สมานา น ยถา\-กาม\-กรณียา ภวิสฺสาม มารสฺส อมุสฺมิํ นิวาเป อมุสฺมิํ จ โลกามิเสติ.}

\prose{เต ยตฺถ อคติ มารสฺส จ มารปริสาย จ, ตตฺราสยํ กปฺปยิํสุ, ตตฺราสยํ กปฺเปตฺวา อมุํ นิวาปํ นิวุตฺตํ มารสฺส อมูนิ จ โลกามิสานิ อนนุปขชฺช อมุจฺฉิตา โภชนานิ ภุญฺชิํสุ. เต ตตฺถ อนนุปขชฺช อมุจฺฉิตา โภชนานิ ภุญฺชมานา น มทํ อาปชฺชิํสุ, อมตฺตา สมานา น ปมาทํ อาปชฺชิํสุ, อปฺปมตฺตา สมานา น ยถา\-กาม\-กรณียา อเหสุํ มารสฺส อมุสฺมิํ นิวาเป อมุสฺมิํ จ โลกามิเส. เอวํ หิ เต ภิกฺขเว จตุตฺถา สมณพฺราหฺมณา ปริมุจฺจิํสุ มารสฺส อิทฺธานุภาวา. เสยฺยถาปิ เต ภิกฺขเว จตุตฺถา มิคชาตา, ตถูปเม อหํ อิเม จตุตฺเถ สมณพฺราหฺมเณ วทามิ.}

\proseitem{271}{กถญฺจ ภิกฺขเว อคติ มารสฺส จ มารปริสาย จ. อิธ ภิกฺขเว ภิกฺขุ วิวิจฺเจว กาเมหิ วิวิจฺจ อกุสเลหิ ธมฺเมหิ สวิตกฺกํ สวิจารํ วิเวกชํ ปีติสุขํ ปฐมํ ฌานํ อุปสมฺปชฺช วิหรติ. อยํ วุจฺจติ ภิกฺขเว ภิกฺขุ อนฺธมกาสิ มารํ, อปทํ วธิตฺวา มารจกฺขุํ อทสฺสนํ คโต ปาปิมโต.}

\prose{ปุน จปรํ ภิกฺขเว ภิกฺขุ วิตกฺก\-วิจารานํ วูปสมา อชฺฌตฺตํ สมฺปสาทนํ เจตโส เอโกทิภาวํ อวิตกฺกํ อวิจารํ สมาธิชํ ปีติสุขํ ทุติยํ ฌานํ อุปสมฺปชฺช วิหรติ. อยํ วุจฺจติ ภิกฺขเว ฯเปฯ ปาปิมโต.}

\prose{ปุน จปรํ ภิกฺขเว ภิกฺขุ ปีติยา จ วิราคา อุเปกฺขโก จ วิหรติ สโต จ สมฺปชาโน, สุขญฺจ กาเยน ปฏิสํเวเทติ, ยํ ตํ อริยา \csromanfolio{215}อาจิกฺขนฺติ “อุเปกฺขโก สติมา สุขวิหารี”ติ ตติยํ ฌานํ อุปสมฺปชฺช วิหรติ. อยํ วุจฺจติ ภิกฺขเว ฯเปฯ ปาปิมโต.}

\prose{ปุน จปรํ ภิกฺขเว ภิกฺขุ สุขสฺส จ ปหานา ทุกฺขสฺส จ ปหานา ปุพฺเพว โสมนสฺส\-โทมนสฺสานํ อตฺถงฺคมา อทุกฺขมสุขํ อุเปกฺขา\-สติ\-ปาริสุทฺธิํ จตุตฺถํ ฌานํ อุปสมฺปชฺช วิหรติ. อยํ วุจฺจติ ภิกฺขเว ฯเปฯ ปาปิมโต.}

\prose{ปุน จปรํ ภิกฺขเว ภิกฺขุ สพฺพโส รูปสญฺญานํ สมติกฺกมา ปฏิฆ\-สญฺญานํ อตฺถงฺคมา นานตฺต\-สญฺญานํ อมนสิการา “อนนฺโต อากาโส”ติ อากาสา\-นญฺจา\-ยตนํ อุปสมฺปชฺช วิหรติ. อยํ วุจฺจติ ภิกฺขเว ฯเปฯ ปาปิมโต.}

\prose{ปุน จปรํ ภิกฺขเว ภิกฺขุ สพฺพโส อากาสา\-นญฺจา\-ยตนํ สมติกฺกมฺม “อนนฺตํ วิญฺญาณนฺ”ติ วิญฺญาณญฺจา\-ยตนํ อุปสมฺปชฺช วิหรติ. อยํ วุจฺจติ ภิกฺขเว ฯเปฯ ปาปิมโต.}

\prose{ปุน จปรํ ภิกฺขเว ภิกฺขุ สพฺพโส วิญฺญาณญฺจา\-ยตนํ สมติกฺกมฺม “นตฺถิ กิญฺจี”ติ อากิญฺจญฺญา\-ยตนํ อุปสมฺปชฺช วิหรติ. อยํ วุจฺจติ ภิกฺขเว ฯเปฯ ปาปิมโต.}

\prose{ปุน จปรํ ภิกฺขเว ภิกฺขุ สพฺพโส อากิญฺจญฺญา\-ยตนํ สมติกฺกมฺม เนว\-สญฺญา\-นา\-สญฺญา\-ยตนํ อุปสมฺปชฺช วิหรติ. อยํ วุจฺจติ ภิกฺขเว ฯเปฯ ปาปิมโต.}

\prose{ปุน จปรํ ภิกฺขเว ภิกฺขุ สพฺพโส เนว\-สญฺญา\-นา\-สญฺญา\-ยตนํ สมติกฺกมฺม สญฺญาเวทยิต\-นิโรธํ อุปสมฺปชฺช วิหรติ, ปญฺญาย จสฺส ทิสฺวา อาสวา ปริกฺขีณา โหนฺติ. อยํ วุจฺจติ ภิกฺขเว ภิกฺขุ อนฺธมกาสิ มารํ, อปทํ วธิตฺวา มารจกฺขุํ อทสฺสนํ คโต ปาปิมโต, ติณฺโณ โลเก วิสตฺติกนฺติ.}

\prosewithcloser{อิทมโวจ ภควา. อตฺตมนา เต ภิกฺขู ภควโต ภาสิตํ อภินนฺทุนฺติ.}{นิวาปสุตฺตํ นิฏฺฐิตํ ปญฺจมํ.}
\csromanmatikaanchor{mtk.56}
\csromantocmarkref{subsection}{6. ปาสราสิสุตฺต}{mtk.56}
\bookmark[dest={mtk.56},level=2]{6. ปาสราสิสุตฺต}
\csromanheader{2}{6. \textbf{ปาสราสิสุตฺต}}
\proseitem{272}{\csromanfolio{216}เอวํ เม สุตํ\csromandash{}เอกํ สมยํ ภควา สาวตฺถิยํ วิหรติ เชตวเน อนาถ\-ปิณฺฑิกสฺส อาราเม. อถ โข ภควา ปุพฺพณฺหสมยํ นิวาเสตฺวา ปตฺต\-จีวรมา\-ทาย สาวตฺถิํ ปิณฺฑาย ปาวิสิ. อถ โข สมฺพหุลา ภิกฺขู เยนายสฺมา อานนฺโท เตนุ\-ปสงฺกมิํสุ, อุปสงฺกมิตฺวา อายสฺมนฺตํ อานนฺทํ เอตทโวจุํ “จิรสฺสุตา โน อาวุโส อานนฺท ภควโต สมฺมุขา ธมฺมี กถา, สาธุ มยํ อาวุโส อานนฺท ลเภยฺยาม ภควโต สมฺมุขา ธมฺมิํ กถํ สวนายา”ติ. เตน หายสฺมนฺโต เยน รมฺมกสฺส พฺราหฺมณสฺส อสฺสโม เตนุ\-ปสงฺกมถ, อปฺเปว นาม ลเภยฺยาถ ภควโต สมฺมุขา ธมฺมิํ กถํ สวนายาติ. เอวมาวุโสติ โข เต ภิกฺขู อายสฺมโต อานนฺทสฺส ปจฺจสฺโสสุํ.}

\prose{อถ โข ภควา สาวตฺถิยํ ปิณฺฑาย จริตฺวา ปจฺฉาภตฺตํ ปิณฺฑปาต\-ปฏิกฺกนฺโต อายสฺมานฺตํ อานนฺทํ อามนฺเตสิ “อายามานนฺท เยน ปุพฺพาราโม มิคาร\-มาตุ\-ปาสาโท เตนุปสงฺกมิสฺสาม ทิวาวิหารายา”ติ. “เอวํ ภนฺเต”ติ โข อายสฺมา อานนฺโท ภควโต ปจฺจสฺโสสิ. อถ โข ภควา อายสฺมตา อานนฺเทน สทฺธิํ เยน ปุพฺพาราโม มิคาร\-มาตุ\-ปาสาโท เตนุปสงฺกมิ ทิวาวิหาราย. อถ โข ภควา สายนฺหสมยํ ปฏิสลฺลานา วุฏฺฐิโต อายสฺมนฺตํ อานนฺทํ อามนฺเตสิ “อายามานนฺท เยน ปุพฺพโกฏฺฐโก เตนุปสงฺกมิสฺสาม คตฺตานิ ปริสิญฺจิตุนฺ”ติ. “เอวํ ภนฺเต”ติ โข อายสฺมา อานนฺโท ภควโต ปจฺจสฺโสสิ.}

\proseitem{273}{อถ โข ภควา อายสฺมตา อานนฺเทน สทฺธิํ เยน ปุพฺพโกฏฺฐโก เตนุปสงฺกมิ คตฺตานิ ปริสิญฺจิตุํ, ปุพฺพโกฏฺฐเก คตฺตานิ ปริสิญฺจิตฺวา ปจฺจุตฺตริตฺวา เอกจีวโร อฏฺฐาสิ คตฺตานิ ปุพฺพาปยมาโน. อถ โข อายสฺมา อานนฺโท ภควนฺตํ เอตทโวจ “อยํ ภนฺเต รมฺมกสฺส พฺราหฺมณสฺส อสฺสโม อวิทูเร, รมณีโย ภนฺเต รมฺมกสฺส พฺราหฺมณสฺส อสฺสโม, ปาสาทิโก ภนฺเต รมฺมกสฺส พฺราหฺมณสฺส อสฺสโม. สาธุ ภนฺเต ภควา เยน รมฺมกสฺส พฺราหฺมณสฺส อสฺสโม, เตนุ\-ปสงฺกมตุ อนุกมฺปํ อุปาทายา”ติ. อธิวาเสสิ ภควา ตุณฺหีภาเวน.}

\prose{\csromanfolio{217}อถ โข ภควา เยน รมฺมกสฺส พฺราหฺมณสฺส อสฺสโม เตนุปสงฺกมิ. เตน โข ปน สมเยน สมฺพหุลา ภิกฺขู รมฺมกสฺส พฺราหฺมณสฺส อสฺสเม ธมฺมิยา กถาย สนฺนิสินฺนา โหนฺติ. อถ โข ภควา พหิทฺวาร\-โกฏฺฐเก อฏฺฐาสิ กถา\-ปริโยสานํ อาคมยมาโน. อถ โข ภควา กถา\-ปริโยสานํ วิทิตฺวา อุกฺกาสิตฺวา อคฺคฬํ อาโกเฏสิ. วิวริํสุ โข เต ภิกฺขู ภควโต ทฺวารํ. อถ โข ภควา รมฺมกสฺส พฺราหฺมณสฺส อสฺสมํ ปวิสิตฺวา ปญฺญตฺเต อาสเน นิสีทิ, นิสชฺช โข ภควา ภิกฺขู อามนฺเตสิ “กายนุตฺถ ภิกฺขเว เอตรหิ กถาย สนฺนิสินฺนา, กา จ ปน โว อนฺตรากถา วิปฺปกตา”ติ. ภควนฺตเมว โข โน ภนฺเต อารพฺภ ธมฺมี กถา วิปฺปกตา, อถ ภควา อนุปฺปตฺโตติ. สาธุ ภิกฺขเว, เอตํ โข ภิกฺขเว ตุมฺหากํ ปติรูปํ กุลปุตฺตานํ สทฺธา อคารสฺมา อนคาริยํ ปพฺพชิตานํ, ยํ ตุเมฺห ธมฺมิยา กถาย สนฺนิสีเทยฺยาถ. สนฺนิปติตานํ โว ภิกฺขเว ทฺวยํ กรณียํ ธมฺมี วา กถา อริโย วา ตุณฺหีภาโว.}

\proseitem{274}{เทฺวมา ภิกฺขเว ปริเยสนา อริยา จ ปริเยสนา อนริยา จ ปริเยสนา. กตมา จ ภิกฺขเว อนริยา ปริเยสนา. อิธ ภิกฺขเว เอกจฺโจ อตฺตนา ชาติธมฺโม สมาโน ชาติ\-ธมฺมํเยว ปริเยสติ, อตฺตนา ชราธมฺโม สมาโน ชรา\-ธมฺมํเยว ปริเยสติ, อตฺตนา พฺยาธิธมฺโม สมาโน พฺยาธิ\-ธมฺมํเยว ปริเยสติ, อตฺตนา มรณธมฺโม สมาโน มรณ\-ธมฺมํเยว ปริเยสติ, อตฺตนา โสกธมฺโม สมาโน โสก\-ธมฺมํเยว ปริเยสติ, อตฺตนา สํกิเลส\-ธมฺโม สมาโน สํกิเลส\-ธมฺมํเยว ปริเยสติ.}

\prose{กิญฺจ ภิกฺขเว ชาติธมฺมํ วเทถ. ปุตฺตภริยํ ภิกฺขเว ชาติธมฺมํ, ทาสิทาสํ ชาติธมฺมํ, อเชฬกํ ชาติธมฺมํ, กุกฺกุฏสูกรํ ชาติธมฺมํ, หตฺถิควาสฺสวฬวํ ชาติธมฺมํ, ชาตรูป\-รชตํ ชาติธมฺมํ. ชาติธมฺมา เหเต ภิกฺขเว อุปธโย. เอตฺถายํ คถิโต\csromansharedfootnote{42627ad098165961}{คธิโต (สฺยา, ก)} มุจฺฉิโต อชฺฌาปนฺโน อตฺตนา ชาติธมฺโม สมาโน ชาติ\-ธมฺมํเยว ปริเยสติ.}

\prose{กิญฺจ ภิกฺขเว ชราธมฺมํ วเทถ. ปุตฺตภริยํ ภิกฺขเว ชราธมฺมํ, ทาสิทาสํ ชราธมฺมํ, อเชฬกํ ชราธมฺมํ, กุกฺกุฏสูกรํ ชราธมฺมํ, หตฺถิควาสฺสวฬวํ \csromanfolio{218}ชราธมฺมํ, ชาตรูป\-รชตํ ชราธมฺมํ. ชราธมฺมา เหเต ภิกฺขเว อุปธโย. เอตฺถายํ คถิโต มุจฺฉิโต อชฺฌาปนฺโน อตฺตนา ชราธมฺโม สมาโน ชาราธมฺมํเยว ปริเยสติ.}

\prose{กิญฺจ ภิกฺขเว พฺยาธิธมฺมํ วเทถ. ปุตฺตภริยํ ภิกฺขเว พฺยาธิธมฺมํ, ทาสิทาสํ พฺยาธิธมฺมํ, อเชฬกํ พฺยาธิธมฺมํ, กุกฺกุฏสูกรํ พฺยาธิธมฺมํ, หตฺถิควาสฺสวฬวํ พฺยาธิธมฺมํ. พฺยาธิธมฺมา เหเต ภิกฺขเว อุปธโย. เอตฺถายํ คถิโต มุจฺฉิโต อชฺฌาปนฺโน อตฺตนา พฺยาธิธมฺโม สมาโน พฺยาธิ\-ธมฺมํเยว ปริเยสติ.}

\prose{กิญฺจ ภิกฺขเว มรณธมฺมํ วเทถ. ปุตฺตภริยํ ภิกฺขเว มรณธมฺมํ, ทาสิทาสํ มรณธมฺมํ, อเชฬกํ มรณธมฺมํ, กุกฺกุฏสูกรํ มรณธมฺมํ, หตฺถิควาสฺสวฬวํ มรณธมฺมํ. มรณธมฺมา เหเต ภิกฺขเว อุปธโย. เอตฺถายํ คถิโต มุจฺฉิโต อชฺฌาปนฺโน อตฺตนา มรณธมฺโม สมาโน มรณ\-ธมฺมํเยว ปริเยสติ.}

\prose{กิญฺจ ภิกฺขเว โสกธมฺมํ วเทถ. ปุตฺตภริยํ ภิกฺขเว โสกธมฺมํ, ทาสิทาสํ โสกธมฺมํ, อเชฬกํ โสกธมฺมํ, กุกฺกุฏสูกรํ โสกธมฺมํ, หตฺถิควาสฺสวฬวํ โสกธมฺมํ. โสกธมฺมา เหเต ภิกฺขเว อุปธโย. เอตฺถายํ คถิโต มุจฺฉิโต อชฺฌาปนฺโน อตฺตนา โสกธมฺโม สมาโน โสก\-ธมฺมํเยว ปริเยสติ.}

\prose{กิญฺจ ภิกฺขเว สํกิเลส\-ธมฺมํ วเทถ. ปุตฺตภริยํ ภิกฺขเว สํกิเลส\-ธมฺมํ, ทาสิทาสํ สํกิเลส\-ธมฺมํ, อเชฬกํ สํกิเลส\-ธมฺมํ, กุกฺกุฏสูกรํ สํกิเลส\-ธมฺมํ, หตฺถิควาสฺสวฬวํ สํกิเลส\-ธมฺมํ, ชาตรูป\-รชตํ สํกิเลส\-ธมฺมํ. สํกิเลส\-ธมฺมา เหเต ภิกฺขเว อุปธโย. เอตฺถายํ คถิโต มุจฺฉิโต อชฺฌาปนฺโน อตฺตนา สํกิเลส\-ธมฺโม สมาโน สํกิเลส\-ธมฺมํเยว ปริเยสติ. อยํ ภิกฺขเว อนริยา ปริเยสนา.}

\proseitem{275}{กตมา จ ภิกฺขเว อริยา ปริเยสนา. อิธ ภิกฺขเว เอกจฺโจ อตฺตนา ชาติธมฺโม สมาโน ชาติธมฺเม อาทีนวํ วิทิตฺวา อชาตํ อนุตฺตรํ โยคกฺเขมํ นิพฺพานํ ปริเยสติ, อตฺตนา ชราธมฺโม สมาโน ชราธมฺเม อาทีนวํ วิทิตฺวา อชรํ อนุตฺตรํ โยคกฺเขมํ นิพฺพานํ ปริเยสติ, อตฺตนา พฺยาธิธมฺโม สมาโน พฺยาธิธมฺเม อาทีนวํ วิทิตฺวา อพฺยาธิํ อนุตฺตรํ โยคกฺเขมํ นิพฺพานํ ปริเยสติ, อตฺตนา มรณธมฺโม สมาโน \csromanfolio{219}มรณธมฺเม อาทีนวํ วิทิตฺวา อมตํ อนุตฺตรํ โยคกฺเขมํ นิพฺพานํ ปริเยสติ, อตฺตนา โสกธมฺโม สมาโน โสกธมฺเม อาทีนวํ วิทิตฺวา อโสกํ อนุตฺตรํ โยคกฺเขมํ นิพฺพานํ ปริเยสติ, อตฺตนา สํกิเลส\-ธมฺโม สมาโน สํกิเลส\-ธมฺเม อาทีนวํ วิทิตฺวา อสํกิลิฏฺฐํ อนุตฺตรํ โยคกฺเขมํ นิพฺพานํ ปริเยสติ. อยํ ภิกฺขเว อริยา ปริเยสนา.}

\proseitem{276}{อหมฺปิ สุทํ ภิกฺขเว ปุพฺเพว สมฺโพธา อนภิสมฺพุทฺโธ โพธิสตฺโตว สมาโน อตฺตนา ชาติธมฺโม สมาโน ชาติ\-ธมฺมํเยว ปริเยสามิ, อตฺตนา ชราธมฺโม สมาโน ชรา\-ธมฺมํเยว ปริเยสามิ, อตฺตนา พฺยาธิธมฺโม สมาโน พฺยาธิ\-ธมฺมํเยว ปริเยสามิ, อตฺตนา มรณธมฺโม สมาโน มรณ\-ธมฺมํเยว ปริเยสามิ, อตฺตนา โสกธมฺโม สมาโน โสก\-ธมฺมํเยว ปริเยสามิ, อตฺตนา สํกิเลส\-ธมฺโม สมาโน สํกิเลส\-ธมฺมํเยว ปริเยสามิ. ตสฺส มยฺหํ ภิกฺขเว เอตทโหสิ “กิํ นุ โข อหํ อตฺตนา ชาติธมฺโม สมาโน ชาติ\-ธมฺมํเยว ปริเยสามิ, อตฺตนา ชราธมฺโม สมาโน, พฺยาธิธมฺโม สมาโน, มรณธมฺโม สมาโน, โสกธมฺโม สมาโน, อตฺตนา สํกิเลส\-ธมฺโม สมาโน สํกิเลส\-ธมฺมํเยว ปริเยสามิ. ยํนูนาหํ อตฺตนา ชาติธมฺโม สมาโน ชาติธมฺเม อาทีนวํ วิทิตฺวา อชาตํ อนุตฺตรํ โยคกฺเขมํ นิพฺพานํ ปริเยเสยฺยํ, อตฺตนา ชราธมฺโม สมาโน ชราธมฺเม อาทีนวํ วิทิตฺวา อชรํ อนุตฺตรํ โยคกฺเขมํ นิพฺพานํ ปริเยเสยฺยํ, อตฺตนา พฺยาธิธมฺโม สมาโน พฺยาธิธมฺเม อาทีนวํ วิทิตฺวา อพฺยาธิํ อนุตฺตรํ โยคกฺเขมํ นิพฺพานํ ปริเยเสยฺยํ, อตฺตนา มรณธมฺโม สมาโน มรณธมฺเม อาทีนวํ วิทิตฺวา อมตํ อนุตฺตรํ โยคกฺเขมํ นิพฺพานํ ปริเยเสยฺยํ, อตฺตนา โสกธมฺโม สมาโน โสกธมฺเม อาทีนวํ วิทิตฺวา อโสกํ อนุตฺตรํ โยคกฺเขมํ นิพฺพานํ ปริเยเสยฺยํ, อตฺตนา สํกิเลส\-ธมฺโม สมาโน สํกิเลส\-ธมฺเม อาทีนวํ วิทิตฺวา อสํกิลิฏฺฐํ อนุตฺตรํ โยคกฺเขมํ นิพฺพานํ ปริเยเสยฺยนฺ”ติ.}

\proseitem{277}{โส โข อหํ ภิกฺขเว อปเรน สมเยน ทหโรว สมาโน สุสุกาฬเกโส ภเทฺรน โยพฺพเนน สมนฺนาคโต ปฐเมน วยสา อกามกานํ มาตาปิตูนํ อสฺสุมุขานํ รุทนฺตานํ เกสมสฺสุํ โอหาเรตฺวา \csromanfolio{220}กาสายานิ วตฺถานิ อจฺฉาเทตฺวา อคารสฺมา อนคาริยํ ปพฺพชิํ. โส เอวํ ปพฺพชิโต สมาโน กิํ กุสลคเวสี\csromansharedfootnote{5feeee011ae4e743}{กิํกุสลํคเวสี (ก)} อนุตฺตรํ สนฺติวรปทํ ปริเยสมาโน เยน อาฬาโร กาลาโม เตนุปสงฺกมิํ, อุปสงฺกมิตฺวา อาฬารํ กาลามํ เอตทโวจํ “อิจฺฉามหํ อาวุโส กาลาม อิมสฺมิํ ธมฺมวินเย พฺรหฺมจริยํ จริตุนฺ”ติ. เอวํ วุตฺเต ภิกฺขเว อาฬาโร กาลาโม มํ เอตทโวจ “วิหรตายสฺมา ตาทิโส อยํ ธมฺโม, ยตฺถ วิญฺญู ปุริโส นจิรสฺเสว สกํ อาจริยกํ สยํ อภิญฺญา สจฺฉิกตฺวา อุปสมฺปชฺช วิหเรยฺยา”ติ. โส โข อหํ ภิกฺขเว นจิรสฺเสว ขิปฺปเมว ตํ ธมฺมํ ปริยาปุณิํ. โส โข อหํ ภิกฺขเว ตาวตเกเนว โอฏฺฐ\-ปหต\-มตฺเตน ลปิต\-ลาปน\-มตฺเตน ญาณวาทญฺจ วทามิ เถรวาทญฺจ, ชานามิ ปสฺสามีติ จ ปฏิชานามิ อหญฺเจว อญฺเญ จ. ตสฺส มยฺหํ ภิกฺขเว เอตทโหสิ “น โข อาฬาโร กาลาโม อิมํ ธมฺมํ เกวลํ สทฺธา\-มตฺตเกน ‘สยํ อภิญฺญา สจฺฉิกตฺวา อุปสมฺปชฺช วิหรามี’ติ ปเวเทติ, อทฺธา อาฬาโร กาลาโม อิมํ ธมฺมํ ชานํ ปสฺสํ วิหรตี”ติ.}

\prose{อถ ขฺวาหํ ภิกฺขเว เยน อาฬาโร กาลาโม เตนุปสงฺกมิํ, อุปสงฺกมิตฺวา อาฬารํ กาลามํ เอตทโวจํ “กิตฺตาวตา โน อาวุโส กาลาม อิมํ ธมฺมํ สยํ อภิญฺญา สจฺฉิกตฺวา อุปสมฺปชฺช วิหรามีติ ปเวเทสี”ติ\csromansharedfootnote{0087500493e88127}{อุปสมฺปชฺช ปเวเทสีติ (สี, สฺยา, อิ)}. เอวํ วุตฺเต ภิกฺขเว อาฬาโร กาลาโม อากิญฺจญฺญา\-ยตนํ ปเวเทสิ. ตสฺส มยฺหํ ภิกฺขเว เอตทโหสิ “น โข อาฬารสฺเสว กาลามสฺส อตฺถิ สทฺธา, มยฺหํปตฺถิ สทฺธา. น โข อาฬารสฺเสว กาลามสฺส อตฺถิ วีริยํ, มยฺหํปตฺถิ วีริยํ. น โข อาฬารสฺสว กาลามสฺส อตฺถิ สติ, มยฺหํปตฺถิ สติ. น โข อาฬารสฺเสว กาลามสฺส อตฺถิ สมาธิ, มยฺหํปตฺถิ สมาธิ. น โข อาฬรสฺเสว กาลามสฺส อตฺถิ ปญฺญา, มยฺหํปตฺถิ ปญฺญา. ยํนูนาหํ ยํ ธมฺมํ อาฬาโร กาลาโม สยํ อภิญฺญา สจฺฉิกตฺวา อุปสมฺปชฺช วิหรามีติ ปเวเทติ. ตสฺส ธมฺมสฺส สจฺฉิกิริยาย ปทเหยฺยนฺ”ติ. โส โข อหํ ภิกฺขเว นจิรสฺเสว ขิปฺปเมว ตํ ธมฺมํ สยํ อภิญฺญา สจฺฉิกตฺวา อุปสมฺปชฺช วิหาสิํ.}

\prose{อถ ขฺวาหํ ภิกฺขเว เยน อาฬาโร กาลาโม เตนุปสงฺกมิํ, อุปสงฺกมิตฺวา อาฬารํ กาลามํ เอตทโวจํ “เอตฺตาวตา โน อาวุโส กาลาม อิมํ ธมฺมํ สยํ อภิญฺญา สจฺฉิกตฺวา อุปสมฺปชฺช ปเวเทสี”ติ. \csromanfolio{221}เอตฺตาวตา โข อหํ อาวุโส อิมํ ธมฺมํ สยํ อภิญฺญา สจฺฉิกตฺวา อุปสมฺปชฺช ปเวเทมีติ. อหมฺปิ โข อาวุโส เอตฺตาวตา อิมํ ธมฺมํ สยํ อภิญฺญา สจฺฉิกตฺวา อุปสมฺปชฺช วิหรามีติ. ลาภา โน อาวุโส, สุลทฺธํ โน อาวุโส. เย มยํ อายสฺมนฺตํ ตาทิสํ สพฺรหฺมจาริํ ปสฺสาม. อิติ ยาหํ ธมฺมํ สยํ อภิญฺญา สจฺฉิกตฺวา อุปสมฺปชฺช ปเวเทมิ, ตํ ตฺวํ ธมฺมํ สยํ อภิญฺญา สจฺฉิกตฺวา อุปสมฺปชฺช วิหรสิ. ยํ ตฺวํ ธมฺมํ สยํ อภิญฺญา สจฺฉิกตฺวา อุปสมฺปชฺช วิหรสิ, ตมหํ ธมฺมํ สยํ อภิญฺญา สจฺฉิกตฺวา อุปสมฺปชฺช ปเวเทมิ. อิติ ยาหํ ธมฺมํ ชานามิ, ตํ ตฺวํ ธมฺมํ ชานาสิ. ยํ ตฺวํ ธมฺมํ ชานาสิ, ตมหํ ธมฺมํ ชานามิ. อิติ ยาทิโส อหํ, ตาทิโส ตุวํ. ยาทิโส ตุวํ, ตาทิโส อหํ. เอหิ ทานิ อาวุโส อุโภว สนฺตา อิมํ คณํ ปริหรามาติ. อิติ โข ภิกฺขเว อาฬาโร กาลาโม อาจริโย เม สมาโน (อตฺตโน)\csromansharedfootnote{7fdfe9259591cd52}{( ) นตฺถิ (สี, สฺยา, อิ)} อนฺเตวาสิํมํ สมานํ อตฺตนา\csromansharedfootnote{f820c9fe09286ecf}{อตฺตโน (สี, อิ)} สมสมํ ฐเปสิ, อุฬาราย จ มํ ปูชาย ปูเชสิ. ตสฺส มยฺหํ ภิกฺขเว เอตทโหสิ “นายํ ธมฺโม นิพฺพิทาย น วิราคาย น นิโรธาย น อุปสมาย น อภิญฺญาย น สมฺโพธาย น นิพฺพานาย สํวตฺตติ, ยาวเทว อากิญฺจญฺญายตนู\-ปปตฺติยา”ติ. โส โข อหํ ภิกฺขเว ตํ ธมฺมํ อนลงฺกริตฺวา ตสฺมา ธมฺมา นิพฺพิชฺช อปกฺกมิํ.}

\proseitem{278}{โส โข อหํ ภิกฺขเว กิํ กุสลคเวสี อนุตฺตรํ สนฺติวรปทํ ปริเยสมาโน เยน อุทโก\csromansharedfootnote{3a36af8b81a3f92f}{อุทฺทโก (สี, สฺยา, อิ)} รามปุตฺโต เตนุปสงฺกมิํ, อุปสงฺกมิตฺวา อุทกํ รามปุตฺตํ เอตทโวจํ “อิจฺฉามหํ อาวุโส\csromansharedfootnote{93cd63326c1513b5}{อาวุโส ราม (สี, สฺยา, ก) มหาสตฺโต รามปุตฺตเมว อโวจ, น รามํ, ราโม หิ ตตฺถ คณาจริโย ภเวยฺย, ตทา จ กาลงฺกโต อสนฺโต. เตเนเวตฺถ รามายตฺตานิ กฺริยปทานิ อตีตกาลวเสน อาคตานิ, อุทโก จ รามปุตฺโต มหาสตฺตสฺส สพฺรหฺมจารีเตฺวว วุตฺโต, น อาจริโยติ. ฏีกายํ จ “ปาฬิยํ รามสฺเสว สมาปตฺติลาภิตา อาคตา น อุทกสฺสา”ติอาทิ ปจฺฉาภาเค ปกาสิตา.} อิมสฺมิํ ธมฺมวินเย พฺรหฺมจริยํ จริตุนฺ”ติ. เอวํ วุตฺเต ภิกฺขเว อุทโก\csromansharedfootnote{3a36af8b81a3f92f}{อุทฺทโก (สี, สฺยา, อิ)} รามปุตฺโต มํ เอตทโวจ “วิหรตายสฺมา, ตาทิโส อยํ ธมฺโม, ยตฺถ วิญฺญู ปุริโส นจิรสฺเสว สกํ อาจริยกํ สยํ อภิญฺญา สจฺฉิกตฺวา อุปสมฺปชฺช วิหเรยฺยา”ติ. โส โข อหํ ภิกฺขเว นจิรสฺเสว ขิปฺปเมว ตํ ธมฺมํ ปริยาปุณิํ. โส \csromanfolio{222}โข อหํ ภิกฺขเว ตาวตเกเนว โอฏฺฐ\-ปหต\-มตฺเตน ลปิต\-ลาปน\-มตฺเตน ญาณวาทญฺจ วทามิ เถรวาทญฺจ, ชานามิ ปสฺสามีติ จ ปฏิชานามิ อหํ เจว อญฺเญ จ. ตสฺส มยฺหํ ภิกฺขเว เอตทโหสิ “น โข ราโม อิมํ ธมฺมํ เกวลํ สทฺธา\-มตฺตเกน ‘สยํ อภิญฺญา สจฺฉิกตฺวา อุปสมฺปชฺช วิหรามี’ติ ปเวเทสิ, อทฺธา ราโม อิมํ ธมฺมํ ชานํ ปสฺสํ วิหาสี”ติ. อถ ขฺวาหํ ภิกฺขเว เยน อุทโก รามปุตฺโต เตนุปสงฺกมิํ, อุปสงฺกมิตฺวา อุทกํ รามปุตฺตํ เอตทโวจํ “กิตฺตาวตา โน อาวุโส ราโม อิมํ ธมฺมํ สยํ อภิญฺญา สจฺฉิกตฺวา อุปสมฺปชฺช วิหรามีติ ปเวเทสี”ติ. เอวํ วุตฺเต ภิกฺขเว อุทโก รามปุตฺโต เนว\-สญฺญา\-นา\-สญฺญา\-ยตนํ ปเวเทสิ. ตสฺส มยฺหํ ภิกฺขเว เอตทโหสิ “น โข รามสฺเสว อโหสิ สทฺธา, มยฺหํปตฺถิ สทฺธา. น โข รามสฺเสว อโหสิ วีริยํ, มยฺหํปตฺถิ วีริยํ. น โข รามสฺเสว อโหสิ สติ, มยฺหํปตฺถิ สติ. น โข รามสฺเสว อโหสิ สมาธิ, มยฺหํปตฺถิ สมาธิ. น โข รามสฺเสว อโหสิ ปญฺญา, มยฺหํปตฺถิ ปญฺญา. ยํนูนาหํ ยํ ธมฺมํ ราโม สยํ อภิญฺญา สจฺฉิกตฺวา อุปสมฺปชฺช วิหรามีติ ปเวเทสิ. ตสฺส ธมฺมสฺส สจฺฉิกิริยาย ปทเหยฺยนฺ”ติ. โส โข อหํ ภิกฺขเว นจิรสฺเสว ขิปฺปเมว ตํ ธมฺมํ สยํ อภิญฺญา สจฺฉิกตฺวา อุปสมฺปชฺช วิหาสิํ.}

\prose{อถ ขฺวาหํ ภิกฺขเว เยน อุทโก รามปุตฺโต เตนุปสงฺกมิํ, อุปสงฺกมิตฺวา อุทกํ รามปุตฺตํ เอตทโวจํ “เอตฺตาวตา โน อาวุโส ราโม อิมํ ธมฺมํ สยํ อภิญฺญา สจฺฉิกตฺวา อุปสมฺปชฺช ปเวเทสี”ติ. เอตฺตาวตา โข อาวุโส ราโม อิมํ ธมฺมํ สยํ อภิญฺญา สจฺฉิกตฺวา อุปสมฺปชฺช ปเวเทสีติ. อหมฺปิ โข อาวุโส เอตฺตาวตา อิมํ ธมฺมํ สยํ อภิญฺญา สจฺฉิกตฺวา อุปสมฺปชฺช วิหรามีติ. ลาภา โน อาวุโส, สุลทฺธํ โน อาวุโส. เย มยํ อายสฺมนฺตํ ตาทิสํ สพฺรหฺมจาริํ ปสฺสาม. อิติ ยํ ธมฺมํ ราโม สยํ อภิญฺญา สจฺฉิกตฺวา อุปสมฺปชฺช ปเวเทสิ, ตํ ตฺวํ ธมฺมํ สยํ อภิญฺญา สจฺฉิกตฺวา อุปสมฺปชฺช วิหรสิ. ยํ ตฺวํ ธมฺมํ สยํ อภิญฺญา สจฺฉิกตฺวา อุปสมฺปชฺช วิหรสิ, ตํ ธมฺมํ ราโม สยํ อภิญฺญา สจฺฉิกตฺวา อุปสมฺปชฺช ปเวเทสิ. อิติ ยํ ธมฺมํ ราโม อภิญฺญาสิ, ตํ ตฺวํ ธมฺมํ ชานาสิ. ยํ ตฺวํ ธมฺมํ ชานาสิ, ตํ ธมฺมํ ราโม อภิญฺญาสิ. อิติ ยาทิโส ราโม อโหสิ, ตาทิโส ตุวํ. ยาทิโส ตุวํ, ตาทิโส ราโม อโหสิ. เอหิ ทานิ อาวุโส ตุวํ อิมํ คณํ \csromanfolio{223}ปริหราติ. อิติ โข ภิกฺขเว อุทโก รามปุตฺโต สพฺรหฺมจารี เม สมาโน อาจริยฏฺฐาเน มํ ฐเปสิ, อุฬาราย จ มํ ปูชาย ปูเชสิ. ตสฺส มยฺหํ ภิกฺขเว เอตทโหสิ “นายํ ธมฺโม นิพฺพิทาย น วิราคาย น นิโรธาย น อุปสมาย น อภิญฺญาย น สมฺโพธาย น นิพฺพานาย สํวตฺตติ, ยาวเทว เนวสญฺญานาสญฺญายตนู\-ปปตฺติยา”ติ. โส โข อหํ ภิกฺขเว ตํ ธมฺมํ อนลงฺกริตฺวา ตสฺมา ธมฺมา นิพฺพิชฺช อปกฺกมิํ.}

\proseitem{279}{โส โข อหํ ภิกฺขเว กิํกุสล\-คเวสี อนุตฺตรํ สนฺติวรปทํ ปริเยสมาโน มคเธสุ อนุปุพฺเพน จาริกํ จรมาโน เยน อุรุเวลา เสนานิคโม ตทวสริํ. ตตฺถทฺทสํ รมณียํ ภูมิภาคํ ปาสาทิกญฺจ วนสณฺฑํ, นทิญฺจ สนฺทนฺติํ เสตกํ สุปติตฺถํ รมณียํ, สมนฺตา\csromansharedfootnote{c08883c29a7882e2}{สามนฺตา (?)} จ โคจรคามํ. ตสฺส มยฺหํ ภิกฺขเว เอตทโหสิ “รมณีโย วต โภ ภูมิภาโค, ปาสาทิโก จ วนสณฺโฑ, นที จ สนฺทติ เสตกา สุปติตฺถา รมณียา, สมนฺตา\csromansharedfootnote{c08883c29a7882e2}{สามนฺตา (?)} จ โคจรคาโม. อลํ วติทํ กุลปุตฺตสฺส ปธาน\-ตฺถิกสฺส ปธานายา”ติ. โส โข อหํ ภิกฺขเว ตตฺเถว นิสีทิํ “อลมิทํ ปธานายา”ติ.}

\proseitem{280}{โส โข อหํ ภิกฺขเว อตฺตนา ชาติธมฺโม สมาโน ชาติธมฺเม อาทีนวํ วิทิตฺวา อชาตํ อนุตฺตรํ โยคกฺเขมํ นิพฺพานํ ปริเยสมาโน อชาตํ อนุตฺตรํ โยคกฺเขมํ นิพฺพานํ อชฺฌคมํ, อตฺตนา ชราธมฺโม สมาโน ชราธมฺเม อาทีนวํ วิทิตฺวา อชรํ อนุตฺตรํ โยคกฺเขมํ นิพฺพานํ ปริเยสมาโน อชรํ อนุตฺตรํ โยคกฺเขมํ นิพฺพานํ อชฺฌคมํ, อตฺตนา พฺยาธิธมฺโม สมาโน พฺยาธิธมฺเม อาทีนวํ วิทิตฺวา อพฺยาธิํ อนุตฺตรํ โยคกฺเขมํ นิพฺพานํ ปริเยสมาโน อพฺยาธิํ อนุตฺตรํ โยคกฺเขมํ นิพฺพานํ อชฺฌคมํ, อตฺตนา มรณธมฺโม สมาโน มรณธมฺเม อาทีนวํ วิทิตฺวา อมตํ อนุตฺตรํ โยคกฺเขมํ นิพฺพานํ อชฺฌคมํ, อตฺตนา โสกธมฺโม สมาโน โสกธมฺเม อาทีนวํ วิทิตฺวา อโสกํ อนุตฺตรํ โยคกฺเขมํ นิพฺพานํ อชฺฌคมํ, อตฺตนา สํกิเลส\-ธมฺโม สมาโน สํกิเลส\-ธมฺเม อาทีนวํ วิทิตฺวา อสํกิลิฏฺฐํ อนุตฺตรํ โยคกฺเขมํ นิพฺพานํ ปริเยสมาโน อสํกิลิฏฺฐํ อนุตฺตรํ โยคกฺขมํ นิพฺพานํ อชฺฌคมํ. ญาณญฺจ ปน เม ทสฺสนํ อุทปาทิ, อกุปฺปา เม วิมุตฺติ, อยมนฺติมา ชาติ, นตฺถิ ทานิ ปุนพฺภโวติ.}

\proseitem{281}{\csromanfolio{224}ตสฺส มยฺหํ ภิกฺขเว เอตทโหสิ “อธิคโต โข มฺยายํ ธมฺโม คมฺภีโร ทุทฺทโส ทุรนุโพโธ สนฺโต ปณีโต อตกฺกาวจโร นิปุโณ ปณฺฑิต\-เวทนีโย. อาลยรามา โข ปนายํ ปชา อาลยรตา อาลยสมฺมุทิตา. อาลยรามาย โข ปน ปชาย อาลยรตาย อาลย\-สมฺมุทิตาย ทุทฺทสํ อิทํ ฐานํ, ยทิทํ อิทปฺปจฺจยตา ปฏิจฺจ\-สมุปฺปาโท. อิทมฺปิ โข ฐานํ ทุทฺทสํ, ยทิทํ สพฺพ\-สงฺขาร\-สมโถ สพฺพู\-ปธิ\-ปฏินิสฺสคฺโค ตณฺหากฺขโย วิราโค นิโรโธ นิพฺพานํ. อหํ เจว โข ปน ธมฺมํ เทเสยฺยํ, ปเร จ เม น อาชาเนยฺยุํ “โส มมสฺส กิลมโถ, สา มมสฺส วิเหสา”ติ. อปิสฺสุ มํ ภิกฺขเว อิมา อนจฺฉริยา คาถาโย ปฏิภํสุ ปุพฺเพ อสฺสุตปุพฺพา.}

\setlength{\csromangathapagewidth}{0pt}
\setlength{\csromangathawakwidth}{0pt}
\csromangathameasuregroup{“กิจฺเฉน เม อธิคตํ, \\ ราคโทสปเรเตหิ, \\ ปฏิโสตคามิํ นิปุณํ, \\ ราครตฺตา น ทกฺขนฺติ,}{\csromangathabat{“กิจฺเฉน เม อธิคตํ,}{หลํ ทานิ ปกาสิตุํ.} \\ \csromangathabat{ราคโทสปเรเตหิ,}{นายํ ธมฺโม สุสมฺพุโธ.} \\ \csromangathabat{ปฏิโสตคามิํ นิปุณํ,}{คมฺภีรํ ทุทฺทสํ อณุํ.} \\ \csromangathabat{ราครตฺตา น ทกฺขนฺติ,}{ตโมขนฺเธน อาวุฏา”ติ\textsuperscript{1}.}}
\csromangathasetleft{“กิจฺเฉน เม อธิคตํ, \\ ราคโทสปเรเตหิ, \\ ปฏิโสตคามิํ นิปุณํ, \\ ราครตฺตา น ทกฺขนฺติ,}
\csromangathagroup{\csromangathastanza{\csromangathabat{“กิจฺเฉน เม อธิคตํ,}{หลํ ทานิ ปกาสิตุํ.} \\ \csromangathabat{ราคโทสปเรเตหิ,}{นายํ ธมฺโม สุสมฺพุโธ.}}\csromangathastanza{\csromangathabat{ปฏิโสตคามิํ นิปุณํ,}{คมฺภีรํ ทุทฺทสํ อณุํ.} \\ \csromangathabat{ราครตฺตา น ทกฺขนฺติ,}{ตโมขนฺเธน อาวุฏา”ติ\csromansharedfootnote{7bed54f7a7a1d6b2}{อาวฏาติ (สี), อาวุตา (สฺยา)}.}}}

\proseitem{282}{อิติห เม ภิกฺขเว ปฏิสญฺจิกฺขโต อปฺโปสฺสุกฺกตาย จิตฺตํ นมติ, โน ธมฺมเทสนาย. อถ โข ภิกฺขเว พฺรหฺมุโน สหมฺปติสฺส มม เจตสา เจโต\-ปริวิตกฺกม\-ญฺญาย เอตทโหสิ “นสฺสติ วต โภ โลโก, วินสฺสติ วต โภ โลโก. ยตฺร หิ นาม ตถาคตสฺส อรหโต สมฺมา\-สมฺพุทฺธสฺส อปฺโปสฺสุกฺกตาย จิตฺตํ นมติ\csromansharedfootnote{7bed54f7a7a1d6b2}{อาวฏาติ (สี), อาวุตา (สฺยา)}, โน ธมฺมเทสนายา”ติ. อถ โข ภิกฺขเว พฺรหฺมา สหมฺปติ เสยฺยถาปิ นาม พลวา ปุริโส สมิญฺชิตํ วา พาหํ ปสาเรยฺย, ปสาริตํ วา พาหํ สมิญฺเชยฺย. เอวเมว พฺรหฺมโลเก อนฺตรหิโต มม ปุรโต ปาตุรโหสิ. อถ โข ภิกฺขเว พฺรหฺมา สหมฺปติ เอกํสํ อุตฺตราสงฺคํ กริตฺวา เยนาหํ เตนญฺชลิํ ปณาเมตฺวา มํ เอตทโวจ “เทเสตุ ภนฺเต ภควา ธมฺมํ, เทเสตุ สุคโต ธมฺมํ. สนฺติ สตฺตา อปฺปรชกฺข\-ชาติกา อสฺสวนตา ธมฺมสฺส ปริหายนฺติ, ภวิสฺสนฺติ ธมฺมสฺส อญฺญาตาโร”ติ. อิทมโวจ ภิกฺขเว พฺรหฺมา สหมฺปติ, อิทํ วตฺวา อถาปรํ เอตทโวจ\csromandash{}}

\setlength{\csromangathapagewidth}{0pt}
\setlength{\csromangathawakwidth}{0pt}
\csromangathameasuregroup{}{ปาตุรโหสิ มคเธสุ ปุพฺเพ, \\ ธมฺโม อสุทฺโธ สมเลหิ จินฺติโต. \\ อปาปุเรตํ\textsuperscript{1} อมตสฺส ทฺวารํ, \\ สุณนฺตุ ธมฺมํ วิมเลนานุพุทฺธํ.}
\csromangathameasurewak{ปาตุรโหสิ มคเธสุ ปุพฺเพ, \\ ธมฺโม อสุทฺโธ สมเลหิ จินฺติโต. \\ อปาปุเรตํ\textsuperscript{1} อมตสฺส ทฺวารํ, \\ สุณนฺตุ ธมฺมํ วิมเลนานุพุทฺธํ.}
\setlength{\csromangathaleftcol}{0pt}
\csromangathagroupwithcloser{\csromangathastanza{\csromangathawak{\csromanfolio{225}ปาตุรโหสิ มคเธสุ ปุพฺเพ, \\ ธมฺโม อสุทฺโธ สมเลหิ จินฺติโต. \\ อปาปุเรตํ\csromansharedfootnote{08efd866b43f7eb5}{อวาปุเรตํ (สี)} อมตสฺส ทฺวารํ, \\ สุณนฺตุ ธมฺมํ วิมเลนา\-นุพุทฺธํ.}}}{เสเล ยถา ปพฺพตมุทฺธนิ\-ฏฺฐิโต,}
\prose{ยถาปิ ปสฺเส ชนตํ สมนฺตโต.}

\setlength{\csromangathapagewidth}{0pt}
\setlength{\csromangathawakwidth}{0pt}
\csromangathameasuregroup{}{ตถูปมํ ธมฺมมยํ สุเมธ, \\ ปาสาทมารุยฺห สมนฺตจกฺขุ. \\ โสกาวติณฺณํ\textsuperscript{1} ชนตมเปตโสโก, \\ อเวกฺขสฺสุ ชาติชราภิภูตํ. \\ อุฏฺเฐหิ วีร วิชิตสงฺคาม, \\ สตฺถวาห อณณ วิจร โลเก. \\ เทสสฺสุ\textsuperscript{1} ภควา ธมฺมํ, \\ อญฺญาตาโร ภวิสฺสนฺตีติ.}
\csromangathameasurewak{ตถูปมํ ธมฺมมยํ สุเมธ, \\ ปาสาทมารุยฺห สมนฺตจกฺขุ. \\ โสกาวติณฺณํ\textsuperscript{1} ชนตมเปตโสโก, \\ อเวกฺขสฺสุ ชาติชราภิภูตํ. \\ อุฏฺเฐหิ วีร วิชิตสงฺคาม, \\ สตฺถวาห อณณ วิจร โลเก. \\ เทสสฺสุ\textsuperscript{1} ภควา ธมฺมํ, \\ อญฺญาตาโร ภวิสฺสนฺตีติ.}
\setlength{\csromangathaleftcol}{0pt}
\csromangathagroup{\csromangathastanza{\csromangathawak{ตถูปมํ ธมฺมมยํ สุเมธ, \\ ปาสาทมารุยฺห สมนฺตจกฺขุ. \\ โสกาวติณฺณํ\csromansharedfootnote{1dc936e39dd63f21}{โสกาวกิณฺณํ (สฺยา)} ชนตม\-เปตโสโก, \\ อเวกฺขสฺสุ ชาติชรา\-ภิภูตํ.}}\csromangathastanza{\csromangathawak{อุฏฺเฐหิ วีร วิชิตสงฺคาม, \\ สตฺถวาห อณณ วิจร โลเก. \\ เทสสฺสุ\csromansharedfootnote{ba5bcb50601b6207}{เทเสตุ (สฺยา, ก)} ภควา ธมฺมํ, \\ อญฺญาตาโร ภวิสฺสนฺตีติ.}}}

\proseitem{283}{อถ โข อหํ ภิกฺขเว พฺรหฺมุโน จ อชฺเฌสนํ วิทิตฺวา สตฺเตสุ จ การุญฺญตํ ปฏิจฺจ พุทฺธจกฺขุนา โลกํ โวโลเกสิํ. อทฺทสํ โข อหํ ภิกฺขเว พุทฺธจกฺขุนา โลกํ โวโลเกนฺโต สตฺเต อปฺปรชกฺเข มหารชกฺเข ติกฺขินฺทฺริเย มุทินฺทฺริเย สฺวากาเร ทฺวากาเร สุวิญฺญาปเย ทุวิญฺญาปเย อปฺเปกจฺเจ ปรโลก\-วชฺช\-ภยทสฺสาวิเน\csromansharedfootnote{37216ce63e126ce5}{ทสฺสาวิโน (สฺยา, กํ, ก)} วิหรนฺเต อปฺเปกจฺเจ น ปรโลก\-วชฺช\-ภยทสฺสาวิเน\csromansharedfootnote{0aa167af097b82d1}{ติฏฺฐนฺติ (สี, สฺยา, อิ)} วิหรนฺเต. เสยฺยถาปิ นาม อุปฺปลินิยํ วา ปทุมินิยํ วา ปุณฺฑรีกินิยํ วา อปฺเปกจฺจานิ อุปฺปลานิ วา ปทุมานิ วา ปุณฺฑรีกานิ วา อุทเก ชาตานิ อุทเก สํวฑฺฒานิ อุทกานุคฺคตานิ อนฺโต\-นิมุคฺค\-โปสีนิ, อปฺเปกจฺจานิ อุปฺปลานิ วา ปทุมานิ วา ปุณฺฑรีกานิ วา อุทเก ชาตานิ อุทเก สํวฑฺฒานิ อุทกานุคฺคตานิ สโมทกํ ฐิตานิ, อปฺเปกจฺจานิ อุปฺปลานิ วา ปทุมานิ วา ปุณฺฑรีกานิ วา อุทเก ชาตานิ อุทเก สํวฑฺฒานิ อุทกํ อจฺจุคฺคมฺม ฐิตานิ\csromansharedfootnote{37216ce63e126ce5}{ทสฺสาวิโน (สฺยา, กํ, ก)} อนุปลิตฺตานิ อุทเกน. เอวเมว โข อหํ ภิกฺขเว \csromanfolio{226}พุทฺธจกฺขุนา โลกํ โวโลเกนฺโต อทฺทสํ สตฺเต อปฺปรชกฺเข มหารชกฺเข ติกฺขินฺทฺริเย มุทินฺทฺริเย สฺวากาเร ทฺวากาเร สุวิญฺญาปเย ทุวิญฺญาปเย อปฺเปกจฺเจ ปรโลก\-วชฺช\-ภยทสฺสาวิเน วิหรนฺเต อปฺเปกจฺเจ น ปรโลก\-วชฺช\-ภยทสฺสาวิเน วิหรนฺเต. อถ ขฺวาหํ ภิกฺขเว พฺรหฺมานํ สหมฺปติํ คาถาย ปจฺจภาสิํ.}

\setlength{\csromangathapagewidth}{0pt}
\setlength{\csromangathawakwidth}{0pt}
\csromangathameasuregroup{}{อปารุตา เตสํ อมตสฺส ทฺวารา, \\ เย โสตวนฺโต ปมุญฺจนฺตุ สทฺธํ. \\ วิหิํสสญฺญี ปคุณํ น ภาสิํ, \\ ธมฺมํ ปณีตํ มนุเชสุ พฺรหฺเมติ.}
\csromangathameasurewak{อปารุตา เตสํ อมตสฺส ทฺวารา, \\ เย โสตวนฺโต ปมุญฺจนฺตุ สทฺธํ. \\ วิหิํสสญฺญี ปคุณํ น ภาสิํ, \\ ธมฺมํ ปณีตํ มนุเชสุ พฺรหฺเมติ.}
\setlength{\csromangathaleftcol}{0pt}
\csromangathagroup{\csromangathastanza{\csromangathawak{อปารุตา เตสํ อมตสฺส ทฺวารา, \\ เย โสตวนฺโต ปมุญฺจนฺตุ สทฺธํ. \\ วิหิํสสญฺญี ปคุณํ น ภาสิํ, \\ ธมฺมํ ปณีตํ มนุเชสุ พฺรหฺเมติ.}}}

\prose{อถ โข ภิกฺขเว พฺรหฺมา สหมฺปติ “กตาวกาโส โขมฺหิ ภควตา ธมฺมเทสนายา”ติ มํ อภิวาเทตฺวา ปทกฺขิณํ กตฺวา ตตฺเถ\-ว\-นฺตรธายิ.}

\proseitem{284}{ตสฺส มยฺหํ ภิกฺขเว เอตทโหสิ “กสฺส นุ โข อหํ ปฐมํ ธมฺมํ เทเสยฺยํ, โก อิมํ ธมฺมํ ขิปฺปเมว อาชานิสฺสตี”ติ. ตสฺส มยฺหํ ภิกฺขเว เอตทโหสิ “อยํ โข อาฬาโร กาลาโม ปณฺฑิโต วิยตฺโต เมธาวี ทีฆรตฺตํ อปฺปรชกฺข\-ชาติโก, ยํนูนาหํ อาฬารสฺส กาลามสฺส ปฐมํ ธมฺมํ เทเสยฺยํ, โส อิมํ ธมฺมํ ขิปฺปเมว อาชานิสฺสตี”ติ. อถ โข มํ ภิกฺขเว เทวตา อุปสงฺกมิตฺวา เอตทโวจ “สตฺตาห\-กาลงฺกโต ภนฺเต อาฬาโร กาลาโมติ. ญาณญฺจ ปน เม ทสฺสนํ อุทปาทิ “สตฺตาห\-กาลงฺกโต อาฬาโร กาลาโม”ติ. ตสฺส มยฺหํ ภิกฺขเว เอตทโหสิ “มหาชานิโย โข อาฬาโร กาลาโม. สเจ หิ โส อิมํ ธมฺมํ สุเณยฺย, ขิปฺปเมว อาชาเนยฺยา”ติ. ตสฺส มยฺหํ ภิกฺขเว เอตทโหสิ “กสฺส นุ โข อหํ ปฐมํ ธมฺมํ เทเสยฺยํ, โก อิมํ ธมฺมํ ขิปฺปเมว อาชานิสฺสตี”ติ. ตสฺส มยฺหํ ภิกฺขเว เอตทโหสิ “อยํ โข อุทโก รามปุตฺโต ปณฺฑิโต วิยตฺโต เมธาวี ทีฆรตฺตํ อปฺปรชกฺข\-ชาติโก, ยํนูนาหํ อุทกสฺส รามปุตฺตสฺส ปฐมํ ธมฺมํ เทเสยฺยํ, โส อิมํ ธมฺมํ ขิปฺปเมว อาชานิสฺสตี”ติ. อถ โข มํ ภิกฺขเว เทวตา อุปสงฺกมิตฺวา เอตทโวจ “อภิโทส\-กาลงฺกโต ภนฺเต อุทโก รามปุตฺโต”ติ. ญาณญฺจ ปน เม ทสฺสนํ อุทปาทิ “อภิโทส\-กาลงฺกโต อุทโก รามปุตฺโต”ติ. ตสฺส มยฺหํ ภิกฺขเว เอตทโหสิ “มหาชานิโย โข อุทโก รามปุตฺโต, สเจ หิ โส อิมํ ธมฺมํ \csromanfolio{227}สุเณยฺย, ขิปฺปเมว อาชาเนยฺยา”ติ. ตสฺส มยฺหํ ภิกฺขเว เอตทโหสิ “กสฺส นุ โข อหํ ปฐมํ ธมฺมํ เทเสยฺยํ, โก อิมํ ธมฺมํ ขิปฺปเมว อาชานิสฺสตี”ติ. ตสฺส มยฺหํ ภิกฺขเว เอตทโหสิ “พหุการา โข เม ปญฺจวคฺคิยา ภิกฺขู, เย มํ ปธาน\-ปหิตตฺตํ อุปฏฺฐหิํสุ, ยํนูนาหํ ปญฺจ\-วคฺคิยานํ ภิกฺขูนํ ปฐมํ ธมฺมํ เทเสยฺยนฺ”ติ. ตสฺส มยฺหํ ภิกฺขเว เอตทโหสิ “กหํ นุ โข เอตรหิ ปญฺจวคฺคิยา ภิกฺขู วิหรนฺตี”ติ. อทฺทสํ โข อหํ ภิกฺขเว ทิพฺเพน จกฺขุนา วิสุทฺเธน อติกฺกนฺต\-มานุสเกน ปญฺจวคฺคิเย ภิกฺขู พาราณสิยํ วิหรนฺเต อิสิปตเน มิคทาเย. อถ ขฺวาหํ ภิกฺขเว อุรุเวลายํ ยถาภิรนฺตํ วิหริตฺวา เยน พาราณสี เตน จาริกํ ปกฺกมิํ\csromansharedfootnote{7d47a3b736abbad2}{ปกฺกามิํ (สฺยา, อิ, ก)}.}

\proseitem{285}{อทฺทสา โข มํ ภิกฺขเว อุปโก อาชีวโก\csromansharedfootnote{c7feccd5d6dc4970}{อาชีวิโก (สี, อิ, ก)} อนฺตรา จ คยํ อนฺตรา จ โพธิํ อทฺธาน\-มคฺค\-ปฺปฏิปนฺนํ, ทิสฺวาน มํ เอตทโวจ “วิปฺปสนฺนานิ โข เต อาวุโส อินฺทฺริยานิ, ปริสุทฺโธ ฉวิวณฺโณ ปริโยทาโต. กํสิ ตฺวํ อาวุโส อุทฺทิสฺส ปพฺพชิโต, โก วา เต สตฺถา, กสฺส วา ตฺวํ ธมฺมํ โรเจสี”ติ. เอวํ วุตฺเต อหํ ภิกฺขเว อุปกํ อาชีวกํ คาถาหิ อชฺฌภาสิํ.}

\prose{สพฺพาภิภู สพฺพวิทูหมสฺมิ, สพฺเพสุ ธมฺเมสุ อนูปลิตฺโต.}

\prose{สพฺพญฺชโห ตณฺหากฺขเย วิมุตฺโต, สยํ อภิญฺญาย กมุทฺทิเสยฺยํ.}

\prose{น เม อาจริโย อตฺถิ, สทิโส เม น วิชฺชติ.}

\prose{สเทวกสฺมิํ โลกสฺมิํ, นตฺถิ เม ปฏิปุคฺคโล.}

\prose{อหํ หิ อรหา โลเก, อหํ สตฺถา อนุตฺตโร.}

\prose{เอโกมฺหิ สมฺมาสมฺพุทฺโธ, สีติภูโตสฺมิ นิพฺพุโต.}

\prose{ธมฺมจกฺกํ ปวตฺเตตุํ, คจฺฉามิ กาสินํ ปุรํ.}

\prose{อนฺธีภูตสฺมิํ\csromansharedfootnote{f86501c9ec2c5023}{อนฺธภูตสฺมิํ (สี, สฺยา, อิ)} โลกสฺมิํ, อาหญฺฉํ อมต\-ทุนฺทุภินฺติ.}

\prose{ยถา โข ตฺวํ อาวุโส ปฏิชานาสิ “อรหสิ อนนฺตชิโน”ติ.}

\prose{มาทิสา เว ชินา โหนฺติ, เย ปตฺตา อาสวกฺขยํ.}

\prose{ชิตา เม ปาปกา ธมฺมา, ตสฺมาหมุปก ชิโนติ.}

\prose{\csromanfolio{228}เอวํ วุตฺเต ภิกฺขเว อุปโก อาชีวโก หุเปยฺยปาวุโสติ\csromansharedfootnote{d7649d7f558d1aa5}{หุเวยฺยปาวุโส (สี, อิ), หุเวยฺยาวุโส (สฺยา)} วตฺวา สีสํ โอกมฺเปตฺวา อุมฺมคฺคํ คเหตฺวา ปกฺกามิ.}

\proseitem{286}{อถ ขฺวาหํ ภิกฺขเว อนุปุพฺเพน จาริกํ จรมาโน เยน พาราณสี อิสิปตนํ มิคทาโย, เยน ปญฺจวคฺคิยา ภิกฺขู เตนุปสงฺกมิํ, อทฺทสํสุ โข มํ ภิกฺขเว ปญฺจวคฺคิยา ภิกฺขู ทูรโต อาคจฺฉนฺตํ, ทิสฺวาน อญฺญมญฺญํ สณฺฐเปสฺสุํ\csromansharedfootnote{5eae52b321c2888a}{อญฺญมญฺญํ กติกํ สณฺฐเปสุํ (วิมหาวคฺเค)} “อยํ โข อาวุโส สมโณ โคตโม อาคจฺฉติ พาหุลฺลิโก\csromansharedfootnote{21115a42ea2f5651}{พาหุลิโก (สี, อิ) สารตฺถทีปนีฏีกาย สเมติ.} ปธาน\-วิพฺภนฺโต อาวตฺโต พาหุลฺลาย, โส เนว อภิวาเทตพฺโพ, น ปจฺจุฏฺฐาตพฺโพ, นาสฺส ปตฺตจีวรํ ปฏิคฺคเหตพฺพํ, อปิ จ โข อาสนํ ฐเปตพฺพํ, สเจ อากงฺขิสฺสติ, นิสีทิสฺสตี”ติ. ยถา ยถา โข อหํ ภิกฺขเว อุปสงฺกมิํ, ตถา ตถา ปญฺจวคฺคิยา ภิกฺขู นาสกฺขิํสุ สกาย กติกาย สณฺฐาตุํ, อปฺเปกจฺเจ มํ ปจฺจุคฺคนฺตฺวา ปตฺตจีวรํ ปฏิคฺคเหสุํ, อปฺเปกจฺเจ อาสนํ ปญฺญเปสุํ, อปฺเปกจฺเจ ปาโททกํ อุปฏฺฐเปสุํ. อปิ จ โข มํ นาเมน จ อาวุโสวาเทน จ สมุทาจรนฺติ.}

\prose{เอวํ วุตฺเต อหํ ภิกฺขเว ปญฺจวคฺคิเย ภิกฺขู เอตทโวจํ “มา ภิกฺขเว ตถาคตํ นาเมน จ อาวุโสวาเทน จ สมุทาจรถ\csromansharedfootnote{2c9a6310319186f1}{สมุทาจริตฺถ (สี, สฺยา, อิ)}, อรหํ ภิกฺขเว ตถาคโต สมฺมาสมฺพุทฺโธ, โอทหถ ภิกฺขเว โสตํ, อมตม\-ธิคตํ อหมนุสาสามิ, อหํ ธมฺมํ เทเสมิ, ยถานุสิฏฺฐํ ตถา ปฏิปชฺชมานา นจิรสฺเสว, ยสฺสตฺถาย กุลปุตฺตา สมฺมเทว อคารสฺมา อนคาริยํ ปพฺพชนฺติ, ตทนุตฺตรํ พฺรหฺมจริย\-ปริโยสานํ ทิฏฺเฐว ธมฺเม สยํ อภิญฺญา สจฺฉิกตฺวา อุปสมฺปชฺช วิหริสฺสถา”ติ. เอวํ วุตฺเต ภิกฺขเว ปญฺจวคฺคิยา ภิกฺขู มํ เอตทโวจุํ “ตายปิ โข ตฺวํ อาวุโส โคตม อิริยาย ตาย ปฏิปทาย ตาย ทุกฺกร\-การิกาย นาชฺฌคมา อุตฺตริ มนุสฺสธมฺมา อลม\-ริย\-ญาณ\-ทสฺสน\-วิเสสํ, กิํ ปน ตฺวํ เอตรหิ พาหุลฺลิโก ปธาน\-วิพฺภนฺโต อาวตฺโต พาหุลฺลาย อธิคมิสฺสสิ อุตฺตริ มนุสฺสธมฺมา อลมริย\-ญาณทสฺสนวิเสสนฺ”ติ. เอวํ วุตฺเต อหํ ภิกฺขเว ปญฺจวคฺคิเย ภิกฺขู เอตทโวจํ “น ภิกฺขเว ตถาคโต พาหุลฺลิโก น ปธาน\-วิพฺภนฺโต น อาวตฺโต \csromanfolio{229}พาหุลฺลาย, อรหํ ภิกฺขเว ตถาคโต สมฺมาสมฺพุทฺโธ, โอทหถ ภิกฺขเว โสตํ, อมตม\-ธิคตํ อหมนุสาสามิ, อหํ ธมฺมํ เทเสมิ, ยถานุสิฏฺฐํ ตถา ปฏิปชฺชมานา นจิรสฺเสว, ยสฺสตฺถาย กุลปุตฺตา สมฺมเทว อคารสฺมา อนคาริยํ ปพฺพชนฺติ, ตทนุตฺตรํ พฺรหฺมจริย\-ปริโยสานํ ทิฏฺเฐว ธมฺเม สยํ อภิญฺญา สจฺฉิกตฺวา อุปสมฺปชฺช วิหริสฺสถา”ติ. ทุติยมฺปิ โข ภิกฺขเว ปญฺจวคฺคิยา ภิกฺขู มํ เอตทโวจุํ “ตายปิ โข ตฺวํ อาวุโส โคตม อิริยาย ตาย ปฏิปทาย ตาย ทุกฺกร\-การิกาย นาชฺฌคมา อุตฺตริ มนุสฺสธมฺมา อลม\-ริย\-ญาณ\-ทสฺสน\-วิเสสํ, กิํ ปน ตฺวํ เอตรหิ พาหุลฺลิโก ปธาน\-วิพฺภนฺโต อาวตฺโต พาหุลฺลาย อธิคมิสฺสสิ อุตฺตริ มนุสฺสธมฺมา อลมริย\-ญาณทสฺสนวิเสสนฺ”ติ. ทุติยมฺปิ โข อหํ ภิกฺขเว ปญฺจวคฺคิเย ภิกฺขู เอตทโวจํ “น ภิกฺขเว ตถาคโต พาหุลฺลิโก ฯเปฯ อุปสมฺปชฺช วิหริสฺสถา”ติ. ตติยมฺปิ โข ภิกฺขเว ปญฺจวคฺคิยา ภิกฺขู มํ เอตทโวจุํ “ตายปิ โข ตฺวํ อาวุโส โคตม อิริยาย ตาย ปฏิปทาย ตาย ทุกฺกร\-การิกาย นาชฺฌคมา อุตฺตริ มนุสฺสธมฺมา อลม\-ริย\-ญาณ\-ทสฺสน\-วิเสสํ, กิํ ปน ตฺวํ เอตรหิ พาหุลฺลิโก ปธาน\-วิพฺภนฺโต อาวตฺโต พาหุลฺลาย อธิคมิสฺสสิ อุตฺตริ มนุสฺสธมฺมา อลมริย\-ญาณทสฺสนวิเสสนฺ”ติ.}

\prose{เอวํ วุตฺเต อหํ ภิกฺขเว ปญฺจวคฺคิเย ภิกฺขู เอตทโวจํ “อภิชานาถ เม โน ตุเมฺห ภิกฺขเว อิโต ปุพฺเพ เอวรูปํ ปภาวิตเมตนฺ”ติ\csromansharedfootnote{12621c97c633c54c}{ภาสิตเมตนฺติ (สี, สฺยา, วินเยปิ)}. โน เหตํ ภนฺเต. อรหํ ภิกฺขเว ตถาคโต สมฺมาสมฺพุทฺโธ, โอทหถ ภิกฺขเว โสตํ, อมตม\-ธิคตํ อหมนุสาสามิ, อหํ ธมฺมํ เทเสมิ, ยถานุสิฏฺฐํ ตถา ปฏิปชฺชมานา นจิรสฺเสว, ยสฺสตฺถาย กุลปุตฺตา สมฺมเทว อคารสฺมา อนคาริยํ ปพฺพชนฺติ, ตทนุตฺตรํ พฺรหฺมจริย\-ปริโยสานํ ทิฏฺเฐว ธมฺเม สยํ อภิญฺญา สจฺฉิกตฺวา อุปสมฺปชฺช วิหริสฺสถาติ. อสกฺขิํ โข อหํ ภิกฺขเว ปญฺจวคฺคิเย ภิกฺขู สญฺญาเปตุํ. เทฺวปิ สุทํ ภิกฺขเว ภิกฺขู โอวทามิ, ตโย ภิกฺขู ปิณฺฑาย จรนฺติ, ยํ ตโย ภิกฺขู ปิณฺฑาย จริตฺวา อาหรนฺติ, เตน ฉพฺพคฺคิยา\csromansharedfootnote{8a056aa3db9cae67}{ฉพฺพคฺคา (สี, สฺยา)} ยาเปม. ตโยปิ สุทํ ภิกฺขเว ภิกฺขู โอวทามิ, เทฺว ภิกฺขู ปิณฺฑาย จรนฺติ, ยํ เทฺว ภิกฺขู ปิณฺฑาย จริตฺวา อาหรนฺติ, เตน ฉพฺพคฺคิยา ยาเปม. อถ โข ภิกฺขเว \csromanfolio{230}ปญฺจวคฺคิยา ภิกฺขู มยา เอวํ โอวทิยมานา เอวํ อนุสาสิยมานา อตฺตนา ชาติธมฺมา สมานา ชาติธมฺเม อาทีนวํ วิทิตฺวา อชาตํ อนุตฺตรํ โยคกฺเขมํ นิพฺพานํ ปริเยสมานา อชาตํ อนุตฺตรํ โยคกฺเขมํ นิพฺพานํ อชฺฌคมํสุ, อตฺตนา ชราธมฺมา สมานา ชราธมฺเม อาทีนวํ วิทิตฺวา อชรํ อนุตฺตรํ โยคกฺเขมํ นิพฺพานํ ปริเยสมานา อชรํ อนุตฺตรํ โยคกฺเขมํ นิพฺพานํ อชฺฌคมํสุ, อตฺตนา พฺยาธิธมฺมา สมานา ฯเปฯ อตฺตนา มรณธมฺมา สมานา ฯเปฯ อตฺตนา โสกธมฺมา สมานา ฯเปฯ อตฺตนา สํกิเลส\-ธมฺมา สมานา สํกิเลส\-ธมฺเม อาทีนวํ วิทิตฺวา อสํกิลิฏฺฐํ อนุตฺตรํ โยคกฺเขมํ นิพฺพานํ ปริเยสมานา อสํกิลิฏฺฐํ อนุตฺตรํ โยคกฺเขมํ นิพฺพานํ อชฺฌคมํสุ. ญาณญฺจ ปน เนสํ ทสฺสนํ อุทปาทิ “อกุปฺปา โน วิมุตฺติ\csromansharedfootnote{0fa99111a5494e6a}{อกุปฺปา เนสํ วิมุตฺติ (ก)}, อยมนฺติมา ชาติ, นตฺถิ ทานิ ปุนพฺภโว”ติ.}

\proseitem{287}{ปญฺจิเม ภิกฺขเว กามคุณา. กตเม ปญฺจ? จกฺขุวิญฺเญยฺยา รูปา อิฏฺฐา กนฺตา มนาปา ปิยรูปา กามูปสํหิตา รชนียา. โสตวิญฺเญยฺยา สทฺทา ฯเปฯ. ฆานวิญฺเญยฺยา คนฺธา. ชิวฺหาวิญฺเญยฺยา รสา. กายวิญฺเญยฺยา โผฏฺฐพฺพา อิฏฺฐา กนฺตา มนาปา ปิยรูปา กามูปสํหิตา รชนียา. อิเม โข ภิกฺขเว ปญฺจ กามคุณา. เย หิ เกจิ ภิกฺขเว สมณา วา พฺราหฺมณา วา อิเม ปญฺจ กามคุเณ คถิตา มุจฺฉิตา อชฺโฌปนฺนา อนาทีนวทสฺสาวิโน อนิสฺสรณปญฺญา ปริภุญฺชนฺติ, เต เอวมสฺสุ เวทิตพฺพา, “อนยมาปนฺนา พฺยสนมาปนฺนา ยถา\-กาม\-กรณียา ปาปิมโต”\csromansharedfootnote{e8383497e4b35829}{ปาปิมโต”ติ (?)}. เสยฺยถาปิ ภิกฺขเว อารญฺญโก มโค พทฺโธ ปาสราสิํ อธิสเยยฺย, โส เอวมสฺส เวทิตพฺโพ, “อนยมาปนฺโน พฺยสนมาปนฺโน ยถา\-กาม\-กรณีโย ลุทฺทสฺส, อาคจฺฉนฺเต จ ปน ลุทฺเท เยน กามํ น ปกฺกมิสฺสตี”ติ, เอวเมว โข ภิกฺขเว เย หิ เกจิ สมณา วา พฺราหฺมณา วา อิเม ปญฺจ กามคุเณ คถิตา มุจฺฉิตา อชฺโฌปนฺนา อนาทีนวทสฺสาวิโน อนิสฺสรณปญฺญา ปริภุญฺชนฺติ, เต เอวมสฺสุ เวทิตพฺพา, “อนยมาปนฺนา พฺยสนมาปนฺนา ยถา\-กาม\-กรณียา ปาปิมโต”. เย จ โข เกจิ ภิกฺขเว สมณา วา พฺราหฺมณา วา อิเม ปญฺจ กามคุเณ อคถิตา อมุจฺฉิตา อนชฺโฌปนฺนา อาทีนว\-ทสฺสาวิโน นิสฺสรณปญฺญา ปริภุญฺชนฺติ, เต เอวมสฺสุ เวทิตพฺพา, “น อนยมาปนฺนา น พฺยสนมาปนฺนา \csromanfolio{231}น ยถา\-กาม\-กรณียา ปาปิมโต”. เสยฺยถาปิ ภิกฺขเว อารญฺญโก มโค อพทฺโธ ปาสราสิํ อธิสเยยฺย, โส เอวมสฺส เวทิตพฺโพ, “น อนยมาปนฺโน, น พฺยสนมาปนฺโน, น ยถา\-กาม\-กรณีโย ลุทฺทสฺส, อาคจฺฉนฺเต จ ปน ลุทฺเท เยน กามํ ปกฺกมิสฺสตี”ติ, เอวเมว โข ภิกฺขเว เย หิ เกจิ สมณา วา พฺราหฺมณา วา อิเม ปญฺจ กามคุเณ อคถิตา อมุจฺฉิตา อนชฺโฌปนฺนา อาทีนว\-ทสฺสาวิโน นิสฺสรณปญฺญา ปริภุญฺชนฺติ, เต เอวมสฺสุ เวทิตพฺพา, “น อนยมาปนฺนา, น พฺยสนมาปนฺนา, น ยถา\-กาม\-กรณียา ปาปิมโต.}

\prose{เสยฺยถาปิ ภิกฺขเว อารญฺญโก มโค อรญฺเญ ปวเน จรมาโน วิสฺสตฺโถ คจฺฉติ, วิสฺสตฺโถ ติฏฺฐติ, วิสฺสตฺโถ นิสีทติ, วิสฺสตฺโถ เสยฺยํ กปฺเปติ. ตํ กิสฺส เหตุ, อนาปาถคโต ภิกฺขเว ลุทฺทสฺส. เอวเมว โข ภิกฺขเว ภิกฺขุ วิวิจฺเจว กาเมหิ วิวิจฺจ อกุสเลหิ ธมฺเมหิ สวิตกฺกํ สวิจารํ วิเวกชํ ปีติสุขํ ปฐมํ ฌานํ อุปสมฺปชฺช วิหรติ. อยํ วุจฺจติ ภิกฺขเว ภิกฺขุ อนฺธมกาสิ มารํ อปทํ, วธิตฺวา มารจกฺขุํ อทสฺสนํ คโต ปาปิมโต.}

\prose{ปุน จปรํ ภิกฺขเว ภิกฺขุ วิตกฺก\-วิจารานํ วูปสมา อชฺฌตฺตํ สมฺปสาทนํ เจตโส เอโกทิภาวํ อวิตกฺกํ อวิจารํ สมาธิชํ ปีติสุขํ ทุติยํ ฌานํ อุปสมฺปชฺช วิหรติ. อยํ วุจฺจติ ภิกฺขเว ฯเปฯ ปาปิมโต.}

\prose{ปุน จปรํ ภิกฺขเว ภิกฺขุ ปีติยา จ วิราคา อุเปกฺขโก จ วิหรติ สโต จ สมฺปชาโน, สุขญฺจ กาเยน ปฏิสํเวเทติ, ยํ ตํ อริยา อาจิกฺขนฺติ “อุเปกฺขโก สติมา สุขวิหารี”ติ ตติยํ ฌานํ อุปสมฺปชฺช วิหรติ. อยํ วุจฺจติ ภิกฺขเว ฯเปฯ ปาปิมโต.}

\prose{ปุน จปรํ ภิกฺขเว ภิกฺขุ สุขสฺส จ ปหานา ทุกฺขสฺส จ ปหานา ปุพฺเพว โสมนสฺส\-โทมนสฺสานํ อตฺถงฺคมา อทุกฺขมสุขํ อุเปกฺขา\-สติ\-ปาริสุทฺธิํ จตุตฺถํ ฌานํ อุปสมฺปชฺช วิหรติ. อยํ วุจฺจติ ภิกฺขเว ฯเปฯ ปาปิมโต.}

\prose{ปุน จปรํ ภิกฺขเว ภิกฺขุ สพฺพโส รูปสญฺญานํ สมติกฺกมา ปฏิฆ\-สญฺญานํ อตฺถงฺคมา นานตฺต\-สญฺญานํ อมนสิการา “อนนฺโต อากาโส”ติ อากาสา\-นญฺจา\-ยตนํ อุปสมฺปชฺช วิหรติ. อยํ วุจฺจติ ภิกฺขเว ฯเปฯ ปาปิมโต.}

\proseitem{6}{\csromanfolio{232}ปุน จปรํ ภิกฺขเว ภิกฺขุ สพฺพโส อากาสา\-นญฺจา\-ยตนํ สมติกฺกมฺม “อนนฺตํ วิญฺญาณนฺ”ติ วิญฺญาณญฺจา\-ยตนํ อุปสมฺปชฺช วิหรติ. อยํ วุจฺจติ ภิกฺขเว ฯเปฯ ปาปิมโต.}

\prose{ปุน จปรํ ภิกฺขเว ภิกฺขุ สพฺพโส วิญฺญาณญฺจา\-ยตนํ สมติกฺกมฺม “นตฺถิ กิญฺจี”ติ อากิญฺจญฺญา\-ยตนํ อุปสมฺปชฺช วิหรติ. อยํ วุจฺจติ ภิกฺขเว ฯเปฯ ปาปิมโต.}

\prose{ปุน จปรํ ภิกฺขเว ภิกฺขุ สพฺพโส อากิญฺจญฺญา\-ยตนํ สมติกฺกมฺม เนว\-สญฺญา\-นา\-สญฺญา\-ยตนํ อุปสมฺปชฺช วิหรติ. อยํ วุจฺจติ ภิกฺขเว ฯเปฯ ปาปิมโต.}

\prose{ปุน จปรํ ภิกฺขเว ภิกฺขุ สพฺพโส เนว\-สญฺญา\-นา\-สญฺญา\-ยตนํ สมติกฺกมฺม สญฺญาเวทยิต\-นิโรธํ อุปสมฺปชฺช วิหรติ. ปญฺญาย จสฺส ทิสฺวา อาสวา ปริกฺขีณา โหนฺติ. อยํ วุจฺจติ ภิกฺขเว ภิกฺขุ อนฺธมกาสิ มารํ, อปทํ วธิตฺวา มารจกฺขุํ อทสฺสนํ คโต ปาปิมโต. ติณฺโณ โลเก วิสตฺติกํ. วิสฺสตฺโถ คจฺฉติ, วิสฺสตฺโถ ติฏฺฐติ, วิสฺสตฺโถ นิสีทติ, วิสฺสตฺโถ เสยฺยํ กปฺเปติ. ตํ กิสฺส เหตุ, อนาปาถคโต ภิกฺขเว ปาปิมโตติ.}

\prosewithcloserruled{อิทมโวจ ภควา. อตฺตมนา เต ภิกฺขู ภควโต ภาสิตํ อภินนฺทุนฺติ.}{ปาสราสิสุตฺตํ นิฏฺฐิตํ ฉฏฺฐํ.}
\csromanmatikaanchor{mtk.57}
\csromantocmarkref{subsection}{7. จูฬหตฺถิปโทปมสุตฺต}{mtk.57}
\bookmark[dest={mtk.57},level=2]{7. จูฬหตฺถิปโทปมสุตฺต}
\csromanheader{2}{7. จูฬหตฺถิปโท\-ปม\-สุตฺต}
\proseitem{288}{เอวํ เม สุตํ\csromandash{}เอกํ สมยํ ภควา สาวตฺถิยํ วิหรติ เชตวเน อนาถ\-ปิณฺฑิกสฺส อาราเม. เตน โข ปน สมเยน ชาณุสฺโสณิ พฺราหฺมโณ สพฺพเสเตน วฬวาภิ\-รเถน\csromansharedfootnote{ef9d4c954eb63519}{วฬภีรเถน (สี, อิ)} สาวตฺถิยา นิยฺยาติ ทิวาทิวสฺส. อทฺทสา โข ชาณุสฺโสณิ พฺราหฺมโณ ปิโลติกํ ปริพฺพาชกํ ทูรโตว อาคจฺฉนฺตํ, ทิสฺวาน ปิโลติกํ ปริพฺพาชกํ เอตทโวจ “หนฺท กุโต นุ ภวํ วจฺฉายโน อาคจฺฉติ ทิวาทิวสฺสา”ติ. อิโต หิ โข \csromanfolio{233}อหํ โภ อาคจฺฉามิ สมณสฺส โคตมสฺส สนฺติกาติ. “ตํ กิํ มญฺญติ ภวํ วจฺฉายโน สมณสฺส โคตมสฺส ปญฺญา\-เวยฺยตฺติยํ ปณฺฑิโต มญฺเญ”ติ. โก จาหํ โภ, โก จ สมณสฺส โคตมสฺส ปญฺญา\-เวยฺยตฺติยํ ชานิสฺสามิ, โสปิ นูนสฺส ตาทิโสว, โย สมณสฺส โคตมสฺส ปญฺญา\-เวยฺยตฺติยํ ชาเนยฺยาติ. “อุฬาราย ขลุ ภวํ วจฺฉายโน สมณํ โคตมํ ปสํสาย ปสํสตี”ติ. โก จาหํ โภ, โก จ สมณํ โคตมํ ปสํสิสฺสามิ, ปสตฺถ\-ปสตฺโถว โส ภวํ โคตโม เสฏฺโฐ เทว\-มนุสฺสา\-นนฺติ. “กํ ปน ภวํ วจฺฉายโน อตฺถวสํ สมฺปสฺสมาโน สมเณ โคตเม เอวํ อภิปฺปสนฺโน”ติ\csromansharedfootnote{95f2fd48cd5bce49}{อภิปฺปสนฺโน โหตีติ (สฺยา)}. เสยฺยถาปิ โภ กุสโล นาควนิโก นาควนํ ปวิเสยฺย, โส ปสฺเสยฺย นาควเน มหนฺตํ หตฺถิปทํ ทีฆโต จ อายตํ ติริยญฺจ วิตฺถตํ, โส นิฏฺฐํ คจฺเฉยฺย “มหา วต โภ นาโค”ติ. เอวเมว โข อหํ โภ ยโต อทฺทสํ สมเณ โคตเม จตฺตาริ ปทานิ, อถาหํ นิฏฺฐมคมํ “สมฺมาสมฺพุทฺโธ ภควา, สฺวากฺขาโต ภควตา ธมฺโม, สุปฺปฏิปนฺโน ภควโต สาวกสํโฆ”ติ.}

\proseitem{289}{กตมานิ จตฺตาริ. อิธาหํ โภ ปสฺสามิ เอกจฺเจ ขตฺติย\-ปณฺฑิเต นิปุเณ กตปรปฺปวาเท วาลเวธิรูเป, เต ภินฺทนฺตา\csromansharedfootnote{54d077a2fcaae81e}{โวภินฺทนฺตา (สี, อิ) วิ + อว + ภินฺทนฺตา} มญฺเญ จรนฺติ ปญฺญาคเตน ทิฏฺฐิคตานิ, เต สุณนฺติ “สมโณ ขลุ โภ โคตโม อมุกํ นาม คามํ วา นิคมํ วา โอสริสฺสตี”ติ. เต ปญฺหํ อภิส\-งฺขโรนฺติ “อิมํ มยํ ปญฺหํ สมณํ โคตมํ อุปสงฺกมิตฺวา ปุจฺฉิสฺสาม, เอวํ เจ โน ปุฏฺโฐ เอวํ พฺยากริสฺสติ, เอวมสฺส มยํ วาทํ อาโรเปสฺสาม, เอวํ เจปิ โน ปุฏฺโฐ เอวํ พฺยากริสฺสติ, เอวํปิสฺส มยํ วาทํ อาโรเปสฺสามา”ติ. เต สุณนฺติ “สมโณ ขลุ โภ โคตโม อมุกํ นาม คามํ วา นิคมํ วา โอสโฏ”ติ. เต เยน สมโณ โคตโม เตนุปสงฺกมนฺติ. เต สมโณ โคตโม ธมฺมิยา กถาย สนฺทสฺเสติ สมาทเปติ สมุตฺเตเชติ สมฺปหํเสติ. เต สมเณน โคตเมน ธมฺมิยา กถาย สนฺทสฺสิตา สมาทปิตา สมุตฺเตชิตา สมฺปหํสิตา น เจว สมณํ โคตมํ ปญฺหํ ปุจฺฉนฺติ, กุโตสฺส\csromansharedfootnote{c9c173295dafb0e8}{กุตสฺส (สี, สฺยา, อิ)} วาทํ อาโรเปสฺสนฺติ, \csromanfolio{234}อญฺญทตฺถุ สมณสฺเสว โคตมสฺส สาวกา สมฺปชฺชนฺติ. ยทาหํ โภ สมเณ โคตเม อิมํ ปฐมํ ปทํ อทฺทสํ, อถาหํ นิฏฺฐมคมํ “สมฺมาสมฺพุทฺโธ ภควา, สฺวากฺขาโต ภควตา ธมฺโม, สุปฺปฏิปนฺโน ภควโต สาวกสํโฆ”ติ.}

\prose{ปุน จปราหํ โภ ปสฺสามิ. อิเธกจฺเจ พฺราหฺมณ\-ปณฺฑิเต ฯเปฯ คหปติ\-ปณฺฑิเต ฯเปฯ สมณปณฺฑิเต นิปุเณ กตปรปฺปวาเท วาลเวธิรูเป, เต ภินฺทนฺตา มญฺเญ จรนฺติ ปญฺญาคเตน ทิฏฺฐิคตานิ, เต สุณนฺติ “สมโณ ขลุ โภ โคตโม อมุกํ นาม คามํ วา นิคมํ วา โอสริสฺสตี”ติ. เต ปญฺหํ อภิส\-งฺขโรนฺติ “อิมํ มยํ ปญฺหํ สมณํ โคตมํ อุปสงฺกมิตฺวา ปุจฺฉิสฺสาม, เอวํ เจ โน ปุฏฺโฐ เอวํ พฺยากริสฺสติ, เอวมสฺส มยํ วาทํ อาโรเปสฺสาม, เอวํ เจปิ โน ปุฏฺโฐ เอวํ พฺยากริสฺสติ, เอวํปิสฺส มยํ วาทํ อาโรเปสฺสามา”ติ. เต สุณนฺติ “สมโณ ขลุ โภ โคตโม อมุกํ นาม คามํ วา นิคมํ วา โอสโฏ”ติ. เต เยน สมโณ โคตโม เตนุปสงฺกมนฺติ. เต สมโณ โคตโม ธมฺมิยา กถาย สนฺทสฺเสติ สมาทเปติ สมุตฺเตเชติ สมฺปหํเสติ. เต สมเณน โคตเมน ธมฺมิยา กถาย สนฺทสฺสิตา สมาทปิตา สมุตฺเตชิตา สมฺปหํสิตา น เจว สมณํ โคตมํ ปญฺหํ ปุจฺฉนฺติ, กุโตสฺส วาทํ อาโรเปสฺสนฺติ, อญฺญทตฺถุ สมณํเยว โคตมํ โอกาสํ ยาจนฺติ อคารสฺมา อนคาริยํ ปพฺพชฺชาย. เต สมโณ โคตโม ปพฺพาเชติ\csromansharedfootnote{8c2877e312dbd6f0}{ปพฺพาเชติ อุปสมฺปาเทติ (สี)}. เต ตตฺถ ปพฺพชิตา สมานา วูปกฏฺฐา อปฺปมตฺตา อาตาปิโน ปหิตตฺตา วิหรนฺตา นจิรสฺเสว; ยสฺสตฺถาย กุลปุตฺตา สมฺมเทว อคารสฺมา อนคาริยํ ปพฺพชนฺติ, ตทนุตฺตรํ พฺรหฺมจริย\-ปริโยสานํ ทิฏฺเฐว ธมฺเม สยํ อภิญฺญา สจฺฉิกตฺวา อุปสมฺปชฺช วิหรนฺติ. เต เอวมาหํสุ “มนํ วต โภ อนสฺสาม, มนํ วต โภ ปนสฺสาม, มยํ หิ ปุพฺเพ อสฺสมณาว สมานา สมณมฺหาติ ปฏิชานิมฺห, อพฺราหฺมณาว สมานา พฺราหฺมณมฺหาติ ปฏิชานิมฺห, อนรหนฺโตว สมานา อรหนฺตมฺหาติ ปฏิชานิมฺห, อิทานิ โขมฺห สมณา, อิทานิ โขมฺห พฺราหฺมณา, อิทานิ โขมฺห อรหนฺโต”ติ, ยทาหํ โภ สมเณ โคตเม อิมํ จตุตฺถํ ปทํ อทฺทสํ, อถาหํ นิฏฺฐมคมํ “สมฺมาสมฺพุทฺโธ ภควา, สฺวากฺขาโต ภควตา ธมฺโม, สุปฺปฏิปนฺโน ภควโต สาวกสํโฆ”ติ.}

\prose{\csromanfolio{235}ยโต โข อหํ โภ สมเณ โคตเม อิมานิ จตฺตาริ ปทานิ อทฺทสํ, อถาหํ นิฏฺฐมคมํ “สมฺมาสมฺพุทฺโธ ภควา, สฺวากฺขาโต ภควตา ธมฺโม, สุปฺปฏิปนฺโน ภควโต สาวกสํโฆ”ติ.}

\proseitem{290}{เอวํ วุตฺเต ชาณุสฺโสณิ พฺราหฺมโณ สพฺพเสตา วฬวาภิรถา โอโรหิตฺวา เอกํสํ อุตฺตราสงฺคํ กริตฺวา เยน ภควา เตนญฺชลิํ ปณาเมตฺวา ติกฺขตฺตุํ อุทานํ อุทาเนสิ “นโม ตสฺส ภควโต อรหโต สมฺมา\-สมฺพุทฺธสฺส, นโม ตสฺส ภควโต อรหโต สมฺมา\-สมฺพุทฺธสฺส, นโม ตสฺส ภควโต อรหโต สมฺมา\-สมฺพุทฺธสฺส, อปฺเปว นาม มยํปิ กทาจิ กรหจิ เตน โภตา โคตเมน สทฺธิํ สมาคจฺเฉยฺยาม, อปฺเปว นาม สิยา โกจิเทว กถาสลฺลาโป”ติ. อถ โข ชาณุสฺโสณิ พฺราหฺมโณ เยน ภควา เตนุปสงฺกมิ, อุปสงฺกมิตฺวา ภควตา สทฺธิํ สมฺโมทิ, สมฺโมทนียํ กถํ สารณียํ วีติสาเรตฺวา เอกมนฺตํ นิสีทิ, เอกมนฺตํ นิสินฺโน โข ชาณุสฺโสณิ พฺราหฺมโณ ยาวตโก อโหสิ ปิโลติเกน ปริพฺพาชเกน สทฺธิํ กถาสลฺลาโป, ตํ สพฺพํ ภควโต อาโรเจสิ. เอวํ วุตฺเต ภควา ชาณุสฺโสณิํ พฺราหฺมณํ เอตทโวจ “น โข พฺราหฺมณ เอตฺตาวตา หตฺถิปโทปโม วิตฺถาเรน ปริปูโร โหติ, อปิ จ พฺราหฺมณ ยถา หตฺถิปโทปโม วิตฺถาเรน ปริปูโร โหติ, ตํ สุณาหิ สาธุกํ มนสิ กโรหิ ภาสิสฺสามี”ติ. “เอวํ โภ”ติ โข ชาณุสฺโสณิ พฺราหฺมโณ ภควโต ปจฺจสฺโสสิ. ภควา เอตทโวจ.}

\proseitem{291}{เสยฺยถาปิ พฺราหฺมณ นาควนิโก นาควนํ ปวิเสยฺย, โส ปสฺเสยฺย นาควเน มหนฺตํ หตฺถิปทํ ทีฆโต จ อายตํ ติริยญฺจ วิตฺถตํ, โย โหติ กุสโล นาควนิโก, เนว ตาว นิฏฺฐํ คจฺฉติ “มหา วต โภ นาโค”ติ. ตํ กิสฺส เหตุ, สนฺติ หิ พฺราหฺมณ นาควเน วามนิกา นาม หตฺถินิโย มหาปทา, ตาสํเปตํ ปทํ อสฺสาติ.}

\prose{โส ตมนุคจฺฉติ, ตมนุคจฺฉนฺโต ปสฺสติ นาควเน มหนฺตํ หตฺถิปทํ ทีฆโต จ อายตํ ติริยญฺจ วิตฺถตํ อุจฺจา จ นิเสวิตํ, โย โหติ กุสโล นาควนิโก, เนว ตาว นิฏฺฐํ คจฺฉติ “มหา วต โภ นาโค”ติ. ตํ กิสฺส เหตุ, สนฺติ หิ พฺราหฺมณ นาควเน อุจฺจา กาฬาริกา นาม หตฺถินิโย มหาปทา, ตาสํเปตํ ปทํ อสฺสาติ.}

\prose{\csromanfolio{236}โส ตมนุคจฺฉติ, ตมนุคจฺฉนฺโต ปสฺสติ นาควเน มหนฺตํ หตฺถิปทํ ทีฆโต จ อายตํ ติริยญฺจ วิตฺถตํ อุจฺจา จ นิเสวิตํ อุจฺจา จ ทนฺเตหิ อารญฺชิตานิ. โย โหติ กุสโล นาควนิโก, เนว ตาว นิฏฺฐํ คจฺฉติ “มหา วต โภ นาโค”ติ. ตํ กิสฺส เหตุ, สนฺติ หิ พฺราหฺมณ นาควเน อุจฺจา กเณรุกา นาม หตฺถินิโย มหาปทา, ตาสํเปตํ ปทํ อสฺสาติ.}

\prose{โส ตมนุคจฺฉติ, ตมนุคจฺฉนฺโต ปสฺสติ นาควเน มหนฺตํ หตฺถิปทํ ทีฆโต จ อายตํ ติริยญฺจ วิตฺถตํ อุจฺจา จ นิเสวิตํ อุจฺจา จ ทนฺเตหิ อารญฺชิตานิ อุจฺจา จ สาขาภงฺคํ, ตญฺจ นาคํ ปสฺสติ รุกฺขมูลคตํ วา อพฺโภกาสคตํ วา คจฺฉนฺตํ วา ติฏฺฐนฺตํ วา นิสินฺนํ วา นิปนฺนํ วา, โส นิฏฺฐํ คจฺฉติ “อยเมว โส มหานาโค”ติ.}

\prose{เอวเมว โข พฺราหฺมณ อิธ ตถาคโต โลเก อุปฺปชฺชติ อรหํ สมฺมาสมฺพุทฺโธ วิชฺชา\-จรณ\-สมฺปนฺโน สุคโต โลกวิทู อนุตฺตโร ปุริส\-ทมฺม\-สารถิ สตฺถา เทวมนุสฺสานํ พุทฺโธ ภควา, โส อิมํ โลกํ สเทวกํ สมารกํ สพฺรหฺมกํ สสฺสมณ\-พฺราหฺมณิํ ปชํ สเทวมนุสฺสํ สยํ อภิญฺญา สจฺฉิกตฺวา ปเวเทติ, โส ธมฺมํ เทเสติ อาทิกลฺยาณํ มชฺเฌกลฺยาณํ ปริโยสาน\-กลฺยาณํ สาตฺถํ สพฺยญฺชนํ เกวล\-ปริปุณฺณํ ปริสุทฺธํ พฺรหฺมจริยํ ปกาเสติ, ตํ ธมฺมํ สุณาติ คหปติ วา คหปติปุตฺโต วา อญฺญตรสฺมิํ วา กุเล ปจฺจาชาโต, โส ตํ ธมฺมํ สุตฺวา ตถาคเต สทฺธํ ปฏิลภติ, โส เตน สทฺธา\-ปฏิลาเภน สมนฺนาคโต อิติ ปฏิสญฺจิกฺขติ “สมฺพาโธ ฆราวาโส รโชปโถ, อพฺโภกาโส ปพฺพชฺชา, นยิทํ สุกรํ อคารํ อชฺฌาวสตา เอกนฺต\-ปริปุณฺณํ เอกนฺต\-ปริสุทฺธํ สงฺขลิขิตํ พฺรหฺมจริยํ จริตุํ, ยํนูนาหํ เกสมสฺสุํ โอหาเรตฺวา กาสายานิ วตฺถานิ อจฺฉาเทตฺวา อคารสฺมา อนคาริยํ ปพฺพเชยฺยนฺ”ติ. โส อปเรน สมเยน อปฺปํ วา โภคกฺขนฺธํ ปหาย มหนฺตํ วา โภคกฺขนฺธํ ปหาย อปฺปํ วา ญาติปริวฏฺฏํ ปหาย มหนฺตํ วา ญาติปริวฏฺฏํ ปหาย เกสมสฺสุํ โอหาเรตฺวา กาสายานิ วตฺถานิ อจฺฉาเทตฺวา อคารสฺมา อนคาริยํ ปพฺพชติ.}

\proseitem{292}{โส เอวํ ปพฺพชิโต สมาโน ภิกฺขูนํ สิกฺขา\-สาชีว\-สมาปนฺโน ปาณาติปาตํ ปหาย ปาณาติปาตา ปฏิวิรโต โหติ \csromanfolio{237}นิหิตทณฺโฑ นิหิตสตฺโถ ลชฺชี ทยาปนฺโน สพฺพปาณภูต\-หิตา\-นุกมฺปี วิหรติ.}

\prose{อทินฺนาทานํ ปหาย อทินฺนาทานา ปฏิวิรโต โหติ ทินฺนาทายี ทินฺน\-ปาฏิกงฺขี, อเถเนน สุจิภูเตน อตฺตนา วิหรติ.}

\prose{อพฺรหฺมจริยํ ปหาย พฺรหฺมจารี โหติ อาราจารี วิรโต เมถุนา คามธมฺมา.}

\prose{มุสาวาทํ ปหาย มุสาวาทา ปฏิวิรโต โหติ สจฺจวาที สจฺจสนฺโธ เถโต\csromansharedfootnote{cf26b13d5a3f351a}{เฐโต (สฺยา, กํ)} ปจฺจยิโก อวิสํวาทโก โลกสฺส.}

\prose{ปิสุณํ วาจํ ปหาย ปิสุณาย วาจาย ปฏิวิรโต โหติ, อิโต สุตฺวา น อมุตฺร อกฺขาตา อิเมสํ เภทาย, อมุตฺร วา สุตฺวา น อิเมสํ อกฺขาตา อมูสํ เภทาย, อิติ ภินฺนานํ วา สนฺธาตา, สหิตานํ วา อนุปฺปทาตา, สมคฺคาราโม สมคฺครโต สมคฺคนนฺที สมคฺคกรณิํ วาจํ ภาสิตา โหติ.}

\prose{ผรุสํ วาจํ ปหาย ผรุสาย วาจาย ปฏิวิรโต โหติ, ยา สา วาจา เนลา กณฺณสุขา เปมนียา หทยงฺคมา โปรี พหุชนกนฺตา พหุชนมนาปา, ตถารูปิํ วาจํ ภาสิตา โหติ.}

\prose{สมฺผปฺปลาปํ ปหาย สมฺผปฺปลาปา ปฏิวิรโต โหติ กาลวาที ภูตวาที อตฺถวาที ธมฺมวาที วินยวาที, นิธานวติํ วาจํ ภาสิตา กาเลน สาปเทสํ ปริยนฺตวติํ อตฺถสํหิตํ.}

\proseitem{293}{โส พีชคาม\-ภูตคาม\-สมารมฺภา ปฏิวิรโต โหติ. เอกภตฺติโก โหติ รตฺตูปรโต วิรโต วิกาลโภชนา. นจฺจ\-คีต\-วาทิต\-วิสูก\-ทสฺสนา ปฏิวิรโต โหติ. มาลา\-คนฺธ\-วิเลปน\-ธารณ\-มณฺฑน\-วิภูสน\-ฏฺฐานา ปฏิวิรโต โหติ. อุจฺจาสยน\-มหาสยนา ปฏิวิรโต โหติ. ชาตรูป\-รชต\-ปฏิคฺคหณา ปฏิวิรโต โหติ. อามก\-ธญฺญ\-ปฏิคฺคหณา ปฏิวิรโต โหติ. อามก\-มํส\-ปฏิคฺคหณา ปฏิวิรโต โหติ. อิตฺถิ\-กุมาริก\-ปฏิคฺคหณา ปฏิวิรโต โหติ. ทาสิ\-ทาส\-ปฏิคฺคหณา ปฏิวิรโต โหติ. อเช\-ฬก\-ปฏิคฺคหณา ปฏิวิรโต โหติ. กุกฺกุฏ\-สูกร\-ปฏิคฺคหณา ปฏิวิรโต โหติ. \csromanfolio{238}หตฺถิ\-ควา\-สฺส\-วฬว\-ปฏิคฺคหณา ปฏิวิรโต โหติ. เขตฺตวตฺถุ\-ปฏิคฺคหณา ปฏิวิรโต โหติ. ทูเตยฺย\-ปหิณ\-คมนา\-นุโยคา ปฏิวิรโต โหติ. กยวิกฺกยา ปฏิวิรโต โหติ. ตุลา\-กูฏ\-กํส\-กูฏ\-มาน\-กูฏา ปฏิวิรโต โหติ. อุกฺโกฏน\-วญฺจน\-นิกติ\-สาจิโยคา ปฏิวิรโต โหติ. เฉทน วธ พนฺธน วิปราโมส อาโลป สหสาการา\csromansharedfootnote{3511a3ba078f619b}{สาหสาการา (ก)} ปฏิวิรโต โหติ\csromansharedfootnote{0ee5bcab7cbd132a}{อิมสฺส อนนฺตรํ “โส อิมินา อริเยน สีลกฺขนฺเธน สมนฺนาคโต อชฺฌตฺตํ อนวชฺชสุขํ ปฏิสํเวเทตี”ติ วจนํ ทีฆนิกาเย อาคตํ, ตํ อิธ สนฺโตสกถาวสาเน อาคตํ, สา จ สนฺโตสกถา ตตฺถ สติสมฺปชญฺญานนฺตรเมว อาคตา.}.}

\proseitem{294}{โส สนฺตุฏฺโฐ โหติ กาย\-ปริหาริเกน จีวเรน กุจฺฉิ\-ปริหาริเกน ปิณฺฑปาเตน, โส เยน เยเนว ปกฺกมติ, สมาทาเยว ปกฺกมติ. เสยฺยถาปิ นาม ปกฺขี สกุโณ เยน เยเนว เฑติ, สปตฺตภาโรว เฑติ. เอวเมว ภิกฺขุ สนฺตุฏฺโฐ โหติ กาย\-ปริหาริเกน จีวเรน กุจฺฉิ\-ปริหาริเกน ปิณฺฑปาเตน, โส เยน เยเนว ปกฺกมติ, สมาทาเยว ปกฺกมติ. โส อิมินา อริเยน สีลกฺขนฺเธน สมนฺนาคโต อชฺฌตฺตํ อนวชฺชสุขํ ปฏิสํเวเทติ.}

\proseitem{295}{โส จกฺขุนา รูปํ ทิสฺวา น นิมิตฺตคฺคาหี โหติ นานุพฺยญฺชน\-คฺคาหี, ยตฺวาธิกรณเมนํ จกฺขุนฺทฺริยํ อสํวุตํ วิหรนฺตํ อภิชฺฌา\-โทมนสฺสา ปาปกา อกุสลา ธมฺมา อนฺวาสฺสเวยฺยุํ, ตสฺส สํวราย ปฏิปชฺชติ, รกฺขติ จกฺขุนฺทฺริยํ, จกฺขุนฺทฺริเย สํวรํ อาปชฺชติ. โสเตน สทฺทํ สุตฺวา ฯเปฯ. ฆาเนน คนฺธํ ฆายิตฺวา. ชิวฺหาย รสํ สายิตฺวา. กาเยน โผฏฺฐพฺพํ ผุสิตฺวา. มนสา ธมฺมํ วิญฺญาย น นิมิตฺตคฺคาหี โหติ นานุพฺยญฺชน\-คฺคาหี, ยตฺวาธิกรณเมนํ มนินฺทฺริยํ อสํวุตํ วิหรนฺตํ อภิชฺฌา\-โทมนสฺสา ปาปกา อกุสลา ธมฺมา อนฺวาสฺสเวยฺยุํ, ตสฺส สํวราย ปฏิปชฺชติ, รกฺขติ มนินฺทฺริยํ, มนินฺทฺริเย สํวรํ อาปชฺชติ. โส อิมินา อริเยน อินฺทฺริย\-สํวเรน สมนฺนาคโต อชฺฌตฺตํ อพฺยาเสกสุขํ ปฏิสํเวเทติ.}

\prose{โส อภิกฺกนฺเต ปฏิกฺกนฺเต สมฺปชานการี โหติ, อาโลกิเต วิโลกิเต สมฺปชานการี โหติ, สมิญฺชิเต ปสาริเต สมฺปชานการี โหติ, สํฆาฏิ\-ปตฺต\-จีวร\-ธารเณ สมฺปชานการี โหติ, อสิเต \csromanfolio{239}ปีเต ขายิเต สายิเต สมฺปชานการี โหติ, อุจฺจาร\-ปสฺสาว\-กมฺเม สมฺปชานการี โหติ, คเต ฐิเต นิสินฺเน สุตฺเต ชาคริเต ภาสิเต ตุณฺหีภาเว สมฺปชานการี โหติ.}

\proseitem{296}{โส อิมินา จ อริเยน สีลกฺขนฺเธน สมนฺนาคโต, (อิมาย จ อริยาย สนฺตุฏฺฐิยา สมนฺนาคโต,)\csromansharedfootnote{265d0b3f2e625f48}{( ) เอตฺถนฺตเร ปาโฐ อิธ นทิสฺสติ, จตุกฺกงฺคุตฺตเร ปน อิมสฺมิํ ฐาเน ทิสฺสติ, อฏฺฐกถาฏีกาสุ จ ตทตฺโถ ปกาสิโต. ตสฺมา โส เอตฺถ ปฏิปูริโต.} อิมินา จ อริเยน อินฺทฺริย\-สํวเรน สมนฺนาคโต, อิมินา จ อริเยน สติสมฺปชญฺเญน สมนฺนาคโต วิวิตฺตํ เสนาสนํ ภชติ อรญฺญํ รุกฺขมูลํ ปพฺพตํ กนฺทรํ คิริคุหํ สุสานํ วนปตฺถํ อพฺโภกาสํ ปลาลปุญฺชํ, โส ปจฺฉาภตฺตํ ปิณฺฑปาต\-ปฏิกฺกนฺโต นิสีทติ ปลฺลงฺกํ อาภุชิตฺวา อุชุํ กายํ ปณิธาย ปริมุขํ สติํ อุปฏฺฐเปตฺวา, โส อภิชฺฌํ โลเก ปหาย วิคตา\-ภิชฺเฌน เจตสา วิหรติ, อภิชฺฌาย จิตฺตํ ปริโสเธติ. พฺยาปาท\-ปฺปโทสํ ปหาย อพฺยาปนฺนจิตฺโต วิหรติ สพฺพปาณภูต\-หิตา\-นุกมฺปี, พฺยาปาทปฺปโทสา จิตฺตํ ปริโสเธติ. ถินมิทฺธํ ปหาย วิคต\-ถินมิทฺโธ วิหรติ อาโลกสญฺญี สโต สมฺปชาโน, ถินมิทฺธา จิตฺตํ ปริโสเธติ. อุทฺธจฺจ\-กุกฺกุจฺจํ ปหาย อนุทฺธโต วิหรติ อชฺฌตฺตํ วูปสนฺตจิตฺโต, อุทฺธจฺจ\-กุกฺกุจฺจา จิตฺตํ ปริโสเธติ. วิจิกิจฺฉํ ปหาย ติณฺณ\-วิจิกิจฺโฉ วิหรติ อกถํกถี กุสเลสุ ธมฺเมสุ, วิจิกิจฺฉาย จิตฺตํ ปริโสเธติ.}

\proseitem{297}{โส อิเม ปญฺจ นีวรเณ ปหาย เจตโส อุปกฺกิเลเส ปญฺญาย ทุพฺพลีกรเณ วิวิจฺเจว กาเมหิ วิวิจฺจ อกุสเลหิ ธมฺเมหิ สวิตกฺกํ สวิจารํ วิเวกชํ ปีติสุขํ ปฐมํ ฌานํ อุปสมฺปชฺช วิหรติ. อิทมฺปิ วุจฺจติ พฺราหฺมณ ตถาคตปทํ อิติปิ ตถาคต\-นิเสวิตํ อิติปิ ตถาคตา\-รญฺชิตํ อิติปิ. น เตฺวว ตาว อริยสาวโก นิฏฺฐํ คจฺฉติ “สมฺมาสมฺพุทฺโธ ภควา, สฺวากฺขาโต ภควตา ธมฺโม, สุปฺปฏิปนฺโน ภควโต สาวกสํโฆ”ติ.}

\prose{ปุน จปรํ พฺราหฺมณ ภิกฺขุ วิตกฺก\-วิจารานํ วูปสมา อชฺฌตฺตํ สมฺปสาทนํ เจตโส เอโกทิภาวํ อวิตกฺกํ อวิจารํ สมาธิชํ ปีติสุขํ ทุติยํ ฌานํ อุปสมฺปชฺช วิหรติ. อิทมฺปิ วุจฺจติ พฺราหฺมณ ฯเปฯ สุปฺปฏิปนฺโน ภควโต สาวกสํโฆติ.}

\prose{\csromanfolio{240}ปุน จปรํ พฺราหฺมณ ภิกฺขุ ปีติยา จ วิราคา อุเปกฺขโก จ วิหรติ สโต จ สมฺปชาโน, สุขญฺจ กาเยน ปฏิสํเวเทติ, ยํ ตํ อริยา อาจิกฺขนฺติ “อุเปกฺขโก สติมา สุขวิหารี”ติ ตติยํ ฌานํ อุปสมฺปชฺช วิหรติ. อิทมฺปิ วุจฺจติ พฺราหฺมณ ฯเปฯ สุปฺปฏิปนฺโน ภควโต สาวกสํโฆติ.}

\prose{ปุน จปรํ พฺราหฺมณ ภิกฺขุ สุขสฺส จ ปหานา ทุกฺขสฺส จ ปหานา ปุพฺเพว โสมนสฺส\-โทมนสฺสานํ อตฺตงฺคมา อทุกฺขมสุขํ อุเปกฺขา\-สติ\-ปาริสุทฺธิํ จตุตฺถํ ฌานํ อุปสมฺปชฺช วิหรติ. อิทมฺปิ วุจฺจติ พฺราหฺมณ ตถาคตปทํ อิติปิ ตถาคต\-นิเสวิตํ อิติปิ ตถาคตา\-รญฺชิตํ อิติปิ. น เตฺวว ตาว อริยสาวโก นิฏฺฐํ คจฺฉติ “สมฺมาสมฺพุทฺโธ ภควา สฺวากฺขาโต ภควตา ธมฺโม, สุปฺปฏิปนฺโน ภควโต สาวกสํโฆ”ติ.}

\proseitem{298}{โส เอวํ สมาหิเต จิตฺเต ปริสุทฺเธ ปริโยทาเต อนงฺคเณ วิคตู\-ปกฺกิเลเส มุทุภูเต กมฺมนิเย ฐิเต อาเนญฺชปฺปตฺเต ปุพฺเพ\-นิวาสา\-นุสฺสติ\-ญาณาย จิตฺตํ อภินินฺนาเมติ. โส อเนกวิหิตํ ปุพฺเพนิวาสํ อนุสฺสรติ, เสยฺยถิทํ, เอกมฺปิ ชาติํ เทฺวปิ ชาติโย ฯเปฯ อิติ สาการํ สอุทฺเทสํ อเนกวิหิตํ ปุพฺเพนิวาสํ อนุสฺสรติ. อิทมฺปิ วุจฺจติ พฺราหฺมณ ตถาคตปทํ อิติปิ ตถาคต\-นิเสวิตํ อิติปิ ตถาคตา\-รญฺชิตํ อิติปิ. น เตฺวว ตาว อริยสาวโก นิฏฺฐํ คจฺฉติ “สมฺมาสมฺพุทฺโธ ภควา, สฺวากฺขาโต ภควตา ธมฺโม, สุปฺปฏิปนฺโน ภควโต สาวกสํโฆ”ติ.}

\prose{โส เอวํ สมาหิเต จิตฺเต ปริสุทฺเธ ปริโยทาเต อนงฺคเณ วิคตู\-ปกฺกิเลเส มุทุภูเต กมฺมนิเย ฐิเต อาเนญฺชปฺปตฺเต สตฺตานํ จุตูปปาต\-ญาณาย จิตฺตํ อภินินฺนาเมติ. โส ทิพฺเพน จกฺขุนา วิสุทฺเธน อติกฺกนฺต\-มานุสเกน ฯเปฯ ยถากมฺมูปเค สตฺเต ปชานาติ. อิทมฺปิ วุจฺจติ พฺราหฺมณ ตถาคตปทํ อิติปิ ตถาคต\-นิเสวิตํ อิติปิ ตถาคตา\-รญฺชิตํ อิติปิ. น เตฺวว ตาว อริยสาวโก นิฏฺฐํ คจฺฉติ “สมฺมาสมฺพุทฺโธ ภควา, สฺวากฺขาโต ภควตา ธมฺโม, สุปฺปฏิปนฺโน ภควโต สาวกสํโฆ”ติ.}

\proseitem{299}{โส เอวํ สมาหิเต จิตฺเต ปริสุทฺเธ ปริโยทาเต อนงฺคเณ วิคตู\-ปกฺกิเลเส มุทุภูเต กมฺมนิเย ฐิเต อาเนญฺชปฺปตฺเต อาสวานํ \csromanfolio{241}ขยญาณาย จิตฺตํ อภินินฺนาเมติ. โส อิทํ ทุกฺขนฺติ ยถาภูตํ ปชานาติ, อยํ ทุกฺข\-สมุทโยติ ยถาภูตํ ปชานาติ, อยํ ทุกฺข\-นิโรโธติ ยถาภูตํ ปชานาติ, อยํ ทุกฺขนิโรธ\-คามินี ปฏิปทาติ ยถาภูตํ ปชานาติ, อิเม อาสวาติ ยถาภูตํ ปชานาติ, อยํ อาสว\-สมุทโยติ ยถาภูตํ ปชานาติ, อยํ อาสวนิโรโธติ ยถาภูตํ ปชานาติ, อยํ อาสว\-นิโรธ\-คามินี ปฏิปทาติ ยถาภูตํ ปชานาติ. อิทมฺปิ วุจฺจติ พฺราหฺมณ ตถาคตปทํ อิติปิ ตถาคต\-นิเสวิตํ อิติปิ ตถาคตา\-รญฺชิตํ อิติปิ. น เตฺวว ตาว อริยสาวโก นิฏฺฐํ คโต โหติ. อปิ จ โข นิฏฺฐํ คจฺฉติ “สมฺมาสมฺพุทฺโธ ภควา, สฺวากฺขาโต ภควตา ธมฺโม, สุปฺปฏิปนฺโน ภควโต สาวกสํโฆ”ติ.}

\prose{ตสฺส เอวํ ชานโต เอวํ ปสฺสโต กามาสวาปิ จิตฺตํ วิมุจฺจติ, ภวาสวาปิ จิตฺตํ วิมุจฺจติ, อวิชฺชาสวาปิ จิตฺตํ วิมุจฺจติ, วิมุตฺตสฺมิํ “วิมุตฺตมฺ”อิติ ญาณํ โหติ, “ขีณา ชาติ, วุสิตํ พฺรหฺมจริยํ, กตํ กรณียํ, นาปรํ อิตฺถตฺตายา”ติ ปชานาติ. อิทมฺปิ วุจฺจติ พฺราหฺมณ ตถาคตปทํ อิติปิ ตถาคต\-นิเสวิตํ อิติปิ ตถาคตา\-รญฺชิตํ อิติปิ. เอตฺตาวตา โข พฺราหฺมณ อริยสาวโก นิฏฺฐํ คโต โหติ “สมฺมาสมฺพุทฺโธ ภควา, สฺวากฺขาโต ภควตา ธมฺโม, สุปฺปฏิปนฺโน ภควโต สาวกสํโฆ”ติ. เอตฺตาวตา โข พฺราหฺมณ หตฺถิปโทปโม วิตฺถาเรน ปริปูโร โหตีติ.}

\prosewithcloser{เอวํ วุตฺเต ชาณุสฺโสณิ พฺราหฺมโณ ภควนฺตํ เอตทโวจ “อภิกฺกนฺตํ โภ โคตม, อภิกฺกนฺตํ โภ โคตม, เสยฺยถาปิ โภ โคตม นิกฺกุชฺชิตํ วา อุกฺกุชฺเชยฺย, ปฏิจฺฉนฺนํ วา วิวเรยฺย, มูฬฺหสฺส วา มคฺคํ อาจิกฺเขยฺย, อนฺธกาเร วา เตลปชฺโชตํ ธาเรยฺย ‘จกฺขุมนฺโต รูปานิ ทกฺขนฺตี’ติ, เอวเมวํ โภตา โคตเมน อเนก\-ปริยาเยน ธมฺโม ปกาสิโต. เอสาหํ ภวนฺตํ โคตมํ สรณํ คจฺฉามิ ธมฺมญฺจ ภิกฺขุสํฆญฺจ, อุปาสกํ มํ ภวํ โคตโม ธาเรตุ อชฺชตคฺเค ปาณุเปตํ สรณํ คตนฺ”ติ.}{จูฬหตฺถิปโท\-ปม\-สุตฺตํ นิฏฺฐิตํ สตฺตมํ.}
\csromanmatikaanchor{mtk.58}
\csromantocmarkref{subsection}{8. มหาหตฺถิปโทปมสุตฺต}{mtk.58}
\bookmark[dest={mtk.58},level=2]{8. มหาหตฺถิปโทปมสุตฺต}
\csromanheader{2}{8. มหา\-หตฺถิปโท\-ปม\-สุตฺต}
\proseitem{300}{\csromanfolio{242}เอวํ เม สุตํ\csromandash{}เอกํ สมยํ ภควา สาวตฺถิยํ วิหรติ เชตวเน อนาถ\-ปิณฺฑิกสฺส อาราเม. ตตฺร โข อายสฺมา สาริปุตฺโต ภิกฺขู อามนฺเตสิ “อาวุโส ภิกฺขเว”ติ. “อาวุโส”ติ โข เต ภิกฺขู อายสฺมโต สาริปุตฺตสฺส ปจฺจสฺโสสุํ. อายสฺมา สาริปุตฺโต เอตทโวจ\csromandash{}เสยฺยถาปิ อาวุโส ยานิ กานิจิ ชงฺคลานํ ปาณานํ ปทชาตานิ, สพฺพานิ ตานิ หตฺถิปเท สโมธานํ คจฺฉนฺติ, หตฺถิปทํ เตสํ อคฺคมกฺขายติ ยทิทํ มหนฺตตฺเตน. เอวเมว โข อาวุโส เยเกจิ กุสลา ธมฺมา, สพฺเพเต จตูสุ อริยสจฺเจสุ สงฺคหํ คจฺฉนฺติ. กตเมสุ จตูสุ. ทุกฺเข อริยสจฺเจ ทุกฺขสมุทเย อริยสจฺเจ ทุกฺขนิโรเธ อริยสจฺเจ ทุกฺขนิโรธ\-คามินิยา ปฏิปทาย อริยสจฺเจ.}

\proseitem{301}{กตมญฺจาวุโส ทุกฺขํ อริยสจฺจํ. ชาติปิ ทุกฺขา, ชราปิ ทุกฺขา, มรณํปิ ทุกฺขํ, โสก\-ปริเทว\-ทุกฺข\-โทมนสฺสุ\-ปายาสาปิ ทุกฺขา, ยมฺปิจฺฉํ น ลภติ ตมฺปิ ทุกฺขํ, สํขิตฺเตน ปญฺจุ\-ปาทาน\-กฺขนฺธา ทุกฺขา. กตเม จาวุโส ปญฺจุ\-ปาทาน\-กฺขนฺธา. เสยฺยถิทํ, รูปุปาทาน\-กฺขนฺโธ, เวทนุ\-ปาทาน\-กฺขนฺโธ, สญฺญุ\-ปาทาน\-กฺขนฺโธ, สงฺขารุ\-ปาทาน\-กฺขนฺโธ, วิญฺญาณุ\-ปาทาน\-กฺขนฺโธ. กตโม จาวุโส รูปุปาทาน\-กฺขนฺโธ. จตฺตาริ จ มหาภูตานิ จตุนฺนญฺจ มหาภูตานํ อุปาทาย รูปํ. กตมา จาวุโส จตฺตาโร มหาภูตา. ปถวีธาตุ อาโปธาตุ เตโชธาตุ วาโยธาตุ.}

\proseitem{302}{กตมา จาวุโส ปถวีธาตุ. ปถวีธาตุ สิยา อชฺฌตฺติกา สิยา พาหิรา. กตมา จาวุโส อชฺฌตฺติกา ปถวีธาตุ. ยํ อชฺฌตฺตํ ปจฺจตฺตํ กกฺขฬํ ขริคตํ อุปาทินฺนํ, เสยฺยถิทํ, เกสา โลมา นขา ทนฺตา ตโจ มํสํ นฺหารุ อฏฺฐิ อฏฺฐิมิญฺชํ วกฺกํ หทยํ ยกนํ กิโลมกํ ปิหกํ ปปฺผาสํ อนฺตํ อนฺตคุณํ อุทริยํ กรีสํ, ยํ วา ปนญฺญํปิ กิญฺจิ อชฺฌตฺตํ ปจฺจตฺตํ กกฺขฬํ ขริคตํ อุปาทินฺนํ. อยํ วุจฺจตาวุโส อชฺฌตฺติกา ปถวีธาตุ. ยา เจว โข ปน อชฺฌตฺติกา ปถวีธาตุ, ยา จ พาหิรา ปถวีธาตุ, ปถวีธาตุเรเวสา. ตํ “เนตํ มม, เนโสหมสฺมิ, น เมโส อตฺตา”ติ เอวเมตํ ยถาภูตํ สมฺมปฺปญฺญาย ทฏฺฐพฺพํ, เอวเมตํ ยถาภูตํ สมฺมปฺปญฺญาย ทิสฺวา ปถวีธาตุยา นิพฺพินฺทติ, ปถวีธาตุยา จิตฺตํ วิราเชติ.}

\prose{\csromanfolio{243}โหติ โข โส อาวุโส สมโย, ยํ พาหิรา อาโปธาตุ ปกุปฺปติ\csromansharedfootnote{21b208c4cf95115a}{ปถวีธาตุ ปกุปฺปติ (ก)}, อนฺตรหิตา ตสฺมิํ สมเย พาหิรา ปถวีธาตุ โหติ. ตสฺสา หิ นาม อาวุโส พาหิราย ปถวีธาตุยา ตาว มหลฺลิกาย อนิจฺจตา ปญฺญายิสฺสติ, ขยธมฺมตา ปญฺญายิสฺสติ, วยธมฺมตา ปญฺญายิสฺสติ, วิปริณาม\-ธมฺมตา ปญฺญายิสฺสติ, กิํ ปนิมสฺส มตฺตฏฺฐกสฺส กายสฺส ตณฺหุ\-ปาทินฺนสฺส “อหนฺติ วา มมนฺติ วา อสฺมีติวา”. อถ ขฺวาสฺส โนเตเวตฺถ โหติ.}

\prose{ตญฺเจ อาวุโส ภิกฺขุํ ปเร อกฺโกสนฺติ ปริภาสนฺติ โรเสนฺติ วิเหเสนฺติ. โส เอวํ ปชานาติ, อุปฺปนฺนา โข เม อยํ โสต\-สมฺผสฺสชา ทุกฺขเวทนา, สา จ โข ปฏิจฺจ, โน อปฏิจฺจ, กิํ ปฏิจฺจ, ผสฺสํ ปฏิจฺจ. โส\csromansharedfootnote{804ebbd58aebf426}{โสปิโข (สฺยา), โสปิ (ก)} “ผสฺโส อนิจฺโจ”ติ ปสฺสติ, “เวทนา อนิจฺจา”ติ ปสฺสติ, “สญฺญา อนิจฺจา”ติ ปสฺสติ, “สงฺขารา อนิจฺจา”ติ ปสฺสติ, “วิญฺญาณํ อนิจฺจนฺ”ติ ปสฺสติ. ตสฺส ธาตารมฺมณเมว จิตฺตํ ปกฺขนฺทติ ปสีทติ สนฺติฏฺฐติ อธิมุจฺจติ.}

\prose{ตญฺเจ อาวุโส ภิกฺขุํ ปเร อนิฏฺเฐหิ อกนฺเตหิ อมนาเปหิ สมุทาจรนฺติ, ปาณิสมฺผสฺเสนปิ เลฑฺฑุสมฺผสฺเสนปิ ทณฺฑสมฺผสฺเสนปิ สตฺถสมฺผสฺเสนปิ. โส เอวํ ปชานาติ “ตถาภูโต โข อยํ กาโย ยถาภูตสฺมิํ กาเย ปาณิสมฺผสฺสาปิ กมนฺติ, เลฑฺฑุสมฺผสฺสาปิ กมนฺติ, ทณฺฑสมฺผสฺสาปิ กมนฺติ, สตฺถสมฺผสฺสาปิ กมนฺติ. วุตฺตํ โข ปเนตํ ภควตา กกจูปโมวาเท ‘อุภโต\-ทณฺฑ\-เกน เจปิ ภิกฺขเว กกเจน โจรา โอจรกา องฺคมงฺคานิ โอกนฺเตยฺยุํ, ตตฺราปิ โย มโน ปทูเสยฺย, น เม โส เตน สาสนกโร’ติ. อารทฺธํ โข ปน เม วีริยํ ภวิสฺสติ อสลฺลีนํ, อุปฏฺฐิตา สติ อสมฺมุฏฺฐา, ปสฺสทฺโธ กาโย อสารทฺโธ, สมาหิตํ จิตฺตํ เอกคฺคํ. กามํ ทานิ อิมสฺมิํ กาเย ปาณิสมฺผสฺสาปิ กมนฺตุ, เลฑฺฑุสมฺผสฺสาปิ กมนฺตุ, ทณฺฑสมฺผสฺสาปิ กมนฺตุ, สตฺถสมฺผสฺสาปิ กมนฺตุ. กรียติ หิทํ พุทฺธานํ สาสนนฺ”ติ.}

\prose{ตสฺส เจ อาวุโส ภิกฺขุโน เอวํ พุทฺธํ อนุสฺสรโต เอวํ ธมฺมํ อนุสฺสรโต เอวํ สํฆํ อนุสฺสรโต อุเปกฺขา กุสลนิสฺสิตา น สณฺฐาติ. โส เตน สํวิชฺชติ, สํเวคํ อาปชฺชติ “อลาภา วต เม, น}

\prose{\csromanfolio{244}วต เม ลาภา, ทุลฺลทฺธํ วต เม, น วต เม สุลทฺธํ, ยสฺส เม เอวํ พุทฺธํ อนุสฺสรโต เอวํ ธมฺมํ อนุสฺสรโต เอวํ สํฆํ อนุสฺสรโต อุเปกฺขา กุสลนิสฺสิตา น สณฺฐาตี”ติ. เสยฺยถาปิ อาวุโส สุณิสา สสุรํ ทิสฺวา สํวิชฺชติ, สํเวคํ อาปชฺชติ, เอวเมว โข อาวุโส ตสฺส เจ ภิกฺขุโน เอวํ พุทฺธํ อนุสฺสรโต เอวํ ธมฺมํ อนุสฺสรโต เอวํ สํฆํ อนุสฺสรโต อุเปกฺขา กุสลนิสฺสิตา น สณฺฐาติ. โส เตน สํวิชฺชติ, สํเวคํ อาปชฺชติ “อลาภา วต เม, น วต เม ลาภา, ทุลฺลทฺธํ วต เม, น วต เม สุลทฺธํ, ยสฺส เม เอวํ พุทฺธํ อนุสฺสรโต เอวํ ธมฺมํ อนุสฺสรโต เอวํ สํฆํ อนุสฺสรโต อุเปกฺขา กุสลนิสฺสิตา น สณฺฐาตี”ติ. ตสฺส เจ อาวุโส ภิกฺขุโน เอวํ พุทฺธํ อนุสฺสรโต เอวํ ธมฺมํ อนุสฺสรโต เอวํ สํฆํ อนุสฺสรโต อุเปกฺขา กุสลนิสฺสิตา สณฺฐาติ, โส เตน อตฺตมโน โหติ. เอตฺตาวตาปิ โข อาวุโส ภิกฺขุโน พหุกตํ โหติ.}

\proseitem{303}{กตมา จาวุโส อาโปธาตุ. อาโปธาตุ สิยา อชฺฌตฺติกา สิยา พาหิรา. กตมา จาวุโส อชฺฌตฺติกา อาโปธาตุ. ยํ อชฺฌตฺตํ ปจฺจตฺตํ อาโป อาโปคตํ อุปาทินฺนํ, เสยฺยถิทํ, ปิตฺตํ เสมฺหํ ปุพฺโพ โลหิตํ เสโท เมโท อสฺสุ วสา เขโฬ สิงฺฆาณิกา ลสิกา มุตฺตํ, ยํ วา ปนญฺญํปิ กิญฺจิ อชฺฌตฺตํ ปจฺจตฺตํ อาโป อาโปคตํ อุปาทินฺนํ. อยํ วุจฺจตาวุโส อชฺฌตฺติกา อาโปธาตุ. ยา เจว โข ปน อชฺฌตฺติกา อาโปธาตุ, ยา จ พาหิรา อาโปธาตุ, อาโปธาตุเรเวสา. ตํ “เนตํ มม, เนโสหมสฺมิ, น เมโส อตฺตา”ติ เอวเมตํ ยถาภูตํ สมฺมปฺปญฺญาย ทฏฺฐพฺพํ, เอวเมตํ ยถาภูตํ สมฺมปฺปญฺญาย ทิสฺวา อาโปธาตุยา นิพฺพินฺทติ, อาโปธาตุยา จิตฺตํ วิราเชติ.}

\prose{โหติ โข โส อาวุโส สมโย, ยํ พาหิรา อาโปธาตุ ปกุปฺปติ, สา คามํปิ วหติ, นิคมํปิ วหติ, นครํปิ วหติ, ชนปทํปิ วหติ, ชนปทปเทสํปิ วหติ. โหติ โข โส อาวุโส สมโย, ยํ มหาสมุทฺเท โยชน\-สติ\-กานิปิ อุทกานิ โอคจฺฉนฺติ, ทฺวิโยชน\-สติกานิปิ อุทตานิ โอคจฺฉนฺติ, ติโยชน\-สติกานิปิ อุทกานิ โอคจฺฉนฺติ, จตุโยชน\-สติกานิปิ อุทกานิ โอคจฺฉนฺติ, ปญฺจ\-โยชน\-สติกานิปิ อุทกานิ โอคจฺฉนฺติ, ฉโยชน\-สติกานิปิ อุทกานิ โอคจฺฉนฺติ, สตฺต\-โยชน\-สติกานิปิ อุทตานิ โอคจฺฉนฺติ. โหติ โข โส อาวุโส สมโย, ยํ \csromanfolio{245}มหาสมุทฺเท สตฺตตาลํปิ อุทกํ สณฺฐาติ, ฉตฺตาลํปิ อุทกํ สณฺฐาติ, ปญฺจตาลํปิ อุทกํ สณฺฐาติ, จตุตฺตาลํปิ อุทกํ สณฺฐาติ, ติตาลํปิ อุทกํ สณฺฐาติ, ทฺวิตาลํปิ อุทกํ สณฺฐาติ, ตาลมตฺตํปิ\csromansharedfootnote{f6efd6a4a74fd158}{ตาลํปิ (สี)} อุทกํ สณฺฐาติ. โหติ โข โส อาวุโส สมโย, ยํ มหาสมุทฺเท สตฺตโปริสํปิ อุทกํ สณฺฐาติ, ฉปฺโปริสํปิ อุทกํ สณฺฐาติ, ปญฺจโปริสํปิ อุทกํ สณฺฐาติ, จตุปฺโปริสํปิ อุทกํ สณฺฐาติ, ติโปริสํปิ อุทกํ สณฺฐาติ, ทฺวิโปริสํปิ อุทกํ สณฺฐาติ, โปริสมตฺตํปิ\csromansharedfootnote{0dc9af6453cc99e7}{โปริสํปิ (สี)} อุทกํ สณฺฐาติ. โหติ โข โส อาวุโส สมโย, ยํ มหาสมุทฺเท อฑฺฒโปริสํปิ อุทกํ สณฺฐาติ, กฏิมตฺตํปิ อุทกํ สณฺฐาติ, ชาณุกมตฺตํปิ อุทกํ สณฺฐาติ, โคปฺผกมตฺตํปิ อุทกํ สณฺฐาติ. โหติ โข โส อาวุโส สมโย, ยํ มหาสมุทฺเท องฺคุลิปพฺพ เตมนมตฺตํปิ อุทกํ น โหติ. ตสฺสา หิ นาม อาวุโส พาหิราย อาโปธาตุยา ตาว มหลฺลิกาย อนิจฺจตา ปญฺญายิสฺสติ, ขยธมฺมตา ปญฺญายิสฺสติ, วยธมฺมตา ปญฺญายิสฺสติ, วิปริณาม\-ธมฺมตา ปญฺญายิสฺสติ. กิํปนิมสฺส มตฺตฏฺฐกสฺส กายสฺส ตณฺหุ\-ปาทินฺนสฺส “อหนฺติ วา มมนฺติ วา อสฺมีติ วา”. อถ ขฺวาสฺส โนเตเวตฺถ โหติ ฯเปฯ. ตสฺส เจ อาวุโส ภิกฺขุโน เอวํ พุทฺธํ อนุสฺสรโต เอวํ ธมฺมํ อนุสฺสรโต เอวํ สํฆํ อนุสฺสรโต อุเปกฺขา กุสลนิสฺสิตา สณฺฐาติ, โส เตน อตฺตมโน โหติ. เอตฺตาวตาปิ โข อาวุโส ภิกฺขุโน พหุกตํ โหติ.}

\proseitem{304}{กตมา จาวุโส เตโชธาตุ. เตโชธาตุ สิยา อชฺฌตฺติกา สิยา พาหิรา. กตมา จาวุโส อชฺฌตฺติกา เตโชธาตุ, ยํ อชฺฌตฺตํ ปจฺจตฺตํ เตโช เตโชคตํ อุปาทินฺนํ, เสยฺยถิทํ, เยน จ สนฺตปฺปติ, เยน จ ชีรียติ, เยน จ ปริฑยฺหติ, เยน จ อสิต\-ปีต\-ขายิต\-สายิตํ สมฺมา ปริณามํ คจฺฉติ, ยํ วา ปนญฺญํปิ กิญฺจิ อชฺฌตฺตํ ปจฺจตฺตํ เตโช เตโชคตํ อุปาทินฺนํ. อยํ วุจฺจตาวุโส อชฺฌตฺติกา เตโชธาตุ. ยา เจว โข ปน อชฺฌตฺติกา เตโชธาตุ, ยา จ พาหิรา เตโชธาตุ, เตโชธาตุเรเวสา. ตํ “เนตํ มม, เนโสหมสฺมิ, น เม โส อตฺตา”ติ เอวเมตํ ยถาภูตํ สมฺมปฺปญฺญาย ทฏฺฐพฺพํ. เอวเมตํ ยถาภูตํ สมฺมปฺปญฺญาย ทิสฺวา เตโชธาตุยา นิพฺพินฺทติ, เตโชธาตุยา จิตฺตํ วิราเชติ.}

\prose{\csromanfolio{246}โหติ โข โส อาวุโส สมโย, ยํ พาหิรา เตโชธาตุ ปกุปฺปติ, สา คามํปิ ทหติ, นิคมํปิ ทหติ, นครํปิ ทหติ, ชนปทํปิ ทหติ, ชนปทปเทสํปิ ทหติ. สา หริตนฺตํ วา ปนฺถนฺตํ วา เสลนฺตํ วา อุทกนฺตํ วา รมณียํ วา ภูมิภาคํ อาคมฺม อนาหารา นิพฺพายติ. โหติ โข โส อาวุโส สมโย, ยํ กุกฺกุฏปตฺเตนปิ นฺหารุททฺทุเลนปิ อคฺคิํ คเวสนฺติ. ตสฺสา หิ นาม อาวุโส พาหิราย เตโชธาตุยา ตาว มหลฺลิกาย อนิจฺจตา ปญฺญายิสฺสติ, ขยธมฺมตา ปญฺญายิสฺสติ, วยธมฺมตา ปญฺญายิสฺสติ, วิปริณาม\-ธมฺมตา ปญฺญายิสฺสติ. กิํ ปนิมสฺส มตฺตฏฺฐกสฺส กายสฺส ตณฺหุ\-ปาทินฺนสฺส “อหนฺติ วา มมนฺติ วา อสฺมีติ วา”. อถ ขฺวาสฺส โนเตเวตฺถ โหติ ฯเปฯ. ตสฺส เจ อาวุโส ภิกฺขุโน เอวํ พุทฺธํ อนุสฺสรโต เอวํ ธมฺมํ อนุสฺสรโต เอวํ สํฆํ อนุสฺสรโต อุเปกฺขา กุสลนิสฺสิตา สณฺฐาติ, โส เตน อตฺตมโน โหติ. เอตฺตาวตาปิ โข อาวุโส ภิกฺขุโน พหุกตํ โหติ.}

\proseitem{305}{กตมา จาวุโส วาโยธาตุ. วาโยธาตุ สิยา อชฺฌตฺติกา สิยา พาหิรา. กตมา จาวุโส อชฺฌตฺติกา วาโยธาตุ. ยํ อชฺฌตฺตํ ปจฺจตฺตํ วาโย วาโยคตํ อุปาทินฺนํ, เสยฺยถิทํ, อุทฺธงฺคมา วาตา, อโธคมา วาตา, กุจฺฉิสยา วาตา, โกฏฺฐาสยา\csromansharedfootnote{7559646c64e44815}{โกฏฺฐสยา (สี, อิ)} วาตา, องฺคมงฺคา\-นุสาริโน วาตา, อสฺสาโส ปสฺสาโส อิติ, ยํ วา ปนญฺญํปิ กิญฺจิ อชฺฌตฺตํ ปจฺจตฺตํ วาโย วาโยคตํ อุปาทินฺนํ. อยํ วุจฺจตาวุโส อชฺฌตฺติกา วาโยธาตุ. ยา เจว โข ปน อชฺฌตฺติกา วาโยธาตุ, ยา จ พาหิรา วาโยธาตุ, วาโยธาตุเรเวสา. ตํ “เนตํ มม, เนโสหมสฺมิ, น เม โส อตฺตา”ติ เอวเมตํ ยถาภูตํ สมฺมปฺปญฺญาย ทฏฺฐพฺพํ. เอวเมตํ ยถาภูตํ สมฺมปฺปญฺญาย ทิสฺวา วาโยธาตุยา นิพฺพินฺทติ, วาโยธาตุยา จิตฺตํ วิราเชติ.}

\prose{โหติ โข โส อาวุโส สมโย, ยํ พาหิรา วาโยธาตุ ปกุปฺปติ, สา คามํปิ วหติ, นิคมํปิ วหติ, นครํปิ วหติ, ชนปทํปิ วหติ, ชนปทปเทสํปิ วหติ. โหติ โข โส อาวุโส สมโย, ยํ คิมฺหานํ ปจฺฉิเม มาเส ตาลวณฺเฏนปิ วิธูปเนนปิ วาตํ ปริเยสนฺติ. \csromanfolio{247}โอสฺสวเนปิ ติณานิ น อิจฺฉนฺติ. ตสฺสา หิ นาม อาวุโส พาหิราย วาโยธาตุยา ตาว มหลฺลิกาย อนิจฺจตา ปญฺญายิสฺสติ, ขยธมฺมตา ปญฺญายิสฺสติ, วยธมฺมตา ปญฺญายิสฺสติ, วิปริณาม\-ธมฺมตา ปญฺญายิสฺสติ. กิํ ปนิมสฺส มตฺตฏฺฐกสฺส กายสฺส ตณฺหุ\-ปาทินฺนสฺส “อหนฺติ วา มมนฺติ วา อสฺมีติ วา”. อถ ขฺวาสฺส โนเตเวตฺถ โหติ.}

\prose{ตญฺเจ อาวุโส ภิกฺขุํ ปเร อกฺโกสนฺติ ปริภาสนฺติ โรเสนฺติ วิเหเสนฺติ. โส เอวํ ปชานาติ, อุปฺปนฺนา โข เม อยํ โสต\-สมฺผสฺสชา ทุกฺขา เวทนา, สา จ โข ปฏิจฺจ, โน อปฏิจฺจ, กิํ ปฏิจฺจ, ผสฺสํ ปฏิจฺจ. โสปิ ผสฺโส อนิจฺโจติ ปสฺสติ, เวทนา อนิจฺจาติ ปสฺสติ, สญฺญา อนิจฺจาติ ปสฺสติ, สงฺขารา อนิจฺจาติ ปสฺสติ, วิญฺญาณํ อนิจฺจนฺติ ปสฺสติ. ตสฺส ธาตารมฺมณเมว จิตฺตํ ปกฺขนฺทติ ปสีทติ สนฺติฏฺฐติ อธิมุจฺจติ.}

\prose{ตญฺเจ อาวุโส ภิกฺขุํ ปเร อนิฏฺเฐหิ อกนฺเตหิ อมนาเปหิ สมุทาจรนฺติ, ปาณิสมฺผสฺเสนปิ เลฑฺฑุสมฺผสฺเสนปิ ทณฺฑสมฺผสฺเสนปิ สตฺถสมฺผสฺเสนปิ. โส เอวํ ปชานาติ “ตถาภูโต โข อยํ กาโย ยถาภูตสฺมิํ กาเย ปาณิสมฺผสฺสาปิ กมนฺติ, เลฑฺฑุสมฺผสฺสาปิ กมนฺติ, ทณฺฑสมฺผสฺสาปิ กมนฺติ, สตฺถสมฺผสฺสาปิ กมนฺติ. วุตฺตํ โข ปเนตํ ภควตา กกจูปโมวาเท ‘อุภโต\-ทณฺฑ\-เกน เจปิ ภิกฺขเว กกเจน โจรา โอจรกา องฺคมงฺคานิ โอกนฺเตยฺยุํ, ตตฺราปิ โย มโน ปทูเสยฺย, น เม โส เตน สาสนกโร’ติ. อารทฺธํ โข ปน เม วีริยํ ภวิสฺสติ อสลฺลีนํ, อุปฏฺฐิตา สติ อสมฺมุฏฺฐา, ปสฺสทฺโธ กาโย อสารทฺโธ, สมาหิตํ จิตฺตํ เอกคฺคํ. กามํ ทานิ อิมสฺมิํ กาเย ปาณิสมฺผสฺสาปิ กมนฺตุ, เลฑฺฑุสมฺผสฺสาปิ กมนฺตุ, ทณฺฑสมฺผสฺสาปิ กมนฺตุ, สตฺถสมฺผสฺสาปิ กมนฺตุ. กรียติ หิทํ พุทฺธานํ สาสนนฺ”ติ.}

\prose{ตสฺส เจ อาวุโส ภิกฺขุโน เอวํ พุทฺธํ อนุสฺสรโต เอวํ ธมฺมํ อนุสฺสรโต เอวํ สํฆํ อนุสฺสรโต อุเปกฺขา กุสลนิสฺสิตา น สณฺฐาติ. โส เตน สํวิชฺชติ สํเวคํ อาปชฺชติ “อลาภา วต เม, น วต เม ลาภา, ทุลฺลทฺธํ วต เม, น วต เม สุลทฺธํ. ยสฺส เม เอวํ พุทฺธํ อนุสฺสรโต เอวํ ธมฺมํ อนุสฺสรโต เอวํ สํฆํ อนุสฺสรโต อุเปกฺขา กุสลนิสฺสิตา น สณฺฐาตี”ติ. เสยฺยถาปิ อาวุโส สุณิสา สสุรํ ทิสฺวา สํวิชฺชติ สํเวคํ อาปชฺชติ. เอวเมว โข อาวุโส ตสฺส \csromanfolio{248}เจ ภิกฺขุโน เอวํ พุทฺธํ อนุสฺสรโต เอวํ ธมฺมํ อนุสฺสรโต เอวํ สํฆํ อนุสฺสรโต อุเปกฺขา กุสลนิสฺสิตา น สณฺฐาติ. โส เตน สํวิชฺชติ สํเวคํ อาปชฺชติ “อลาภา วต เม, น วต เม ลาภา, ทุลฺลทฺธํ วต เม, น วต เม สุลทฺธํ. ยสฺส เม เอวํ พุทฺธํ อนุสฺสรโต เอวํ ธมฺมํ อนุสฺสรโต เอวํ สํฆํ อนุสฺสรโต อุเปกฺขา กุสลนิสฺสิตา น สณฺฐาตี”ติ. ตสฺส เจ อาวุโส ภิกฺขุโน เอวํ พุทฺธํ อนุสฺสรโต เอวํ ธมฺมํ อนุสฺสรโต เอวํ สํฆํ อนุสฺสรโต อุเปกฺขา กุสลนิสฺสิตา สณฺฐาติ, โส เตน อตฺตมโน โหติ. เอตฺตาวตาปิ โข อาวุโส ภิกฺขุโน พหุกตํ โหติ.}

\proseitem{306}{เสยฺยถาปิ อาวุโส กฏฺฐญฺจ ปฏิจฺจ วลฺลิญฺจ ปฏิจฺจ ติณญฺจ ปฏิจฺจ มตฺติกญฺจ ปฏิจฺจ อากาโส ปริวาริโต อคารํเตฺวว สงฺขํ คจฺฉติ. เอวเมว โข อาวุโส อฏฺฐิญฺจ ปฏิจฺจ นฺหารุญฺจ ปฏิจฺจ มํสญฺจ ปฏิจฺจ จมฺมญฺจ ปฏิจฺจ อากาโส ปริวาริโต รูปํเตฺวว สงฺขํ คจฺฉติ. อชฺฌตฺติก\-ญฺเจว\csromansharedfootnote{14358106d1f45051}{อชฺฌตฺติกญฺเจ (สี, สฺยา, อิ), อชฺฌตฺติกญฺเจปิ (?)} อาวุโส จกฺขุํ อปริภินฺนํ โหติ, พาหิรา จ รูปา น อาปาถํ อาคจฺฉนฺติ, โน จ ตชฺโช สมนฺนาหาโร โหติ. เนว ตาว ตชฺชสฺส วิญฺญาณ\-ภาคสฺส ปาตุภาโว โหติ. อชฺฌตฺติก\-ญฺเจว\csromansharedfootnote{14358106d1f45051}{อชฺฌตฺติกญฺเจ (สี, สฺยา, อิ), อชฺฌตฺติกญฺเจปิ (?)} อาวุโส จกฺขุํ อปริภินฺนํ โหติ, พาหิรา จ รูปา อาปาถํ อาคจฺฉนฺติ, โน จ ตชฺโช สมนฺนาหาโร โหติ. เนว ตาว ตชฺชสฺส วิญฺญาณ\-ภาคสฺส ปาตุภาโว โหติ. ยโต จ โข อาวุโส อชฺฌตฺติก\-ญฺเจว จกฺขุํ อปริภินฺนํ โหติ, พาหิรา จ รูปา อาปาถํ อาคจฺฉนฺติ, ตชฺโช จ สมนฺนาหาโร โหติ. เอวํ ตชฺชสฺส วิญฺญาณ\-ภาคสฺส ปาตุภาโว โหติ. ยํ ตถาภูตสฺส รูปํ, ตํ รูปุปาทาน\-กฺขนฺเธ สงฺคหํ คจฺฉติ. ยา ตถาภูตสฺส เวทนา, สา เวทนุ\-ปาทาน\-กฺขนฺเธ สงฺคหํ คจฺฉติ. ยา ตถาภูตสฺส สญฺญา, สา สญฺญุ\-ปาทาน\-กฺขนฺเธ สงฺคหํ คจฺฉนฺติ. เย ตถาภูตสฺส สงฺขารา, เต สงฺขารุ\-ปาทาน\-กฺขนฺเธ สงฺคหํ คจฺฉนฺติ. ยํ ตถาภูตสฺส วิญฺญาณํ, ตํ วิญฺญาณุ\-ปาทาน\-กฺขนฺเธ สงฺคหํ คจฺฉติ. โส เอวํ ปชานาติ “เอวํ หิ กิร อิเมสํ ปญฺจนฺนํ อุปาทาน\-กฺขนฺธานํ สงฺคโห สนฺนิปาโต สมวาโย โหติ. วุตฺตํ โข ปเนตํ ภควตา ‘โย ปฏิจฺจ\-สมุปฺปาทํ ปสฺสติ, โส ธมฺมํ ปสฺสติ. โย ธมฺมํ ปสฺสติ, โส ปฏิจฺจ\-สมุปฺปาทํ ปสฺสตี’ติ. ปฏิจฺจ\-สมุปฺปนฺนา โข ปนิเม ยทิทํ ปญฺจุ\-ปาทาน\-กฺขนฺธา. โย อิเมสุ \csromanfolio{249}ปญฺจสุ อุปาทาน\-กฺขนฺเธสุ ฉนฺโท อาลโย อนุนโย อชฺโฌสานํ, โส ทุกฺขสมุทโย. โย อิเมสุ ปญฺจสุ อุปาทาน\-กฺขนฺเธสุ ฉนฺทราค\-วินโย ฉนฺท\-ราค\-ปฺปหานํ, โส ทุกฺขนิโรโธ”ติ. เอตฺตาวตาปิ โข อาวุโส ภิกฺขุโน พหุกตํ โหติ.}

\prose{อชฺฌตฺติก\-ญฺเจว อาวุโส โสตํ อปริภินฺนํ โหติ ฯเปฯ ฆานํ อปริภินฺนํ โหติ, ชิวฺหา อปริภินฺนา โหติ, กาโย อปริภินฺโน โหติ, มโน อปริภินฺโน โหติ, พาหิรา จ ธมฺมา น อาปาถํ อาคจฺฉนฺติ, โน จ ตชฺโช สมนฺนาหาโร โหติ. เนว ตาว ตชฺชสฺส วิญฺญาณ\-ภาคสฺส ปาตุภาโว โหติ. อชฺฌตฺติโก เจว อาวุโส มโน อปริภินฺโน โหติ, พาหิรา จ ธมฺมา อาปาถํ อาคจฺฉนฺติ, โน จ ตชฺโช สมนฺนาหาโร โหติ. เนว ตาว ตชฺชสฺส วิญฺญาณ\-ภาคสฺส ปาตุภาโว โหติ. ยโต จ โข อาวุโส อชฺฌตฺติโก เจว มโน อปริภินฺโน โหติ, พาหิรา จ ธมฺมา อาปาถํ อาคจฺฉนฺติ, ตชฺโช จ สมนฺนาหาโร โหติ. เอวํ ตชฺชสฺส วิญฺญาณ\-ภาคสฺส ปาตุภาโว โหติ. ยํ ตถาภูตสฺส รูปํ, ตํ รูปุปาทาน\-กฺขนฺเธ สงฺคหํ คจฺฉติ. ยา ตถาภูตสฺส เวทนา, สา เวทนุ\-ปาทาน\-กฺขนฺเธ สงฺคหํ คจฺฉติ. ยา ตถาภูตสฺส สญฺญา, สา สญฺญุ\-ปาทาน\-กฺขนฺเธ สงฺคหํ คจฺฉนฺติ. เย ตถาภูตสฺส สงฺขารา, เต สงฺขารุ\-ปาทาน\-กฺขนฺเธ สงฺคหํ คจฺฉนฺติ. ยํ ตถาภูตสฺส วิญฺญาณํ, ตํ วิญฺญาณุ\-ปาทาน\-กฺขนฺเธ สงฺคหํ คจฺฉติ, โส เอวํ ปชานาติ “เอวํ หิ กิร อิเมสํ ปญฺจนฺนํ อุปาทาน\-กฺขนฺธานํ สงฺคโห สนฺนิปาโต สมวาโย โหติ. วุตฺตํ โข ปเนตํ ภควตา ‘โย ปฏิจฺจ\-สมุปฺปาทํ ปสฺสติ, โส ธมฺมํ ปสฺสติ. โย ธมฺมํ ปสฺสติ, โส ปฏิจฺจ\-สมุปฺปาทํ ปสฺสตี’ติ. ปฏิจฺจ\-สมุปฺปนฺนา โข ปนิเม ยทิทํ ปญฺจุ\-ปาทาน\-กฺขนฺธา. โย อิเมสุ ปญฺจสุ อุปาทาน\-กฺขนฺเธสุ ฉนฺโท อาลโย อนุนโย อชฺโฌสานํ, โส ทุกฺขสมุทโย. โย อิเมสุ ปญฺจสุ อุปาทาน\-กฺขนฺเธสุ ฉนฺทราค\-วินโย ฉนฺท\-ราค\-ปฺปหานํ, โส ทุกฺขนิโรโธ”ติ. เอตฺตาวตาปิ โข อาวุโส ภิกฺขุโน พหุกตํ โหตีติ.}

\prosewithcloser{อิทมโวจ อายสฺมา สาริปุตฺโต. อตฺตมนา เต ภิกฺขู อายสฺมโต สาริปุตฺตสฺส ภาสิตํ อภินนฺทุนฺติ.}{มหา\-หตฺถิปโท\-ปม\-สุตฺตํ นิฏฺฐิตํ อฏฺฐมํ.}
\csromanmatikaanchor{mtk.59}
\csromantocmarkref{subsection}{9. มหาสาโรปมสุตฺต}{mtk.59}
\bookmark[dest={mtk.59},level=2]{9. มหาสาโรปมสุตฺต}
\csromanheader{2}{9. มหา\-สาโร\-ปม\-สุตฺต}
\proseitem{307}{\csromanfolio{250}เอวํ เม สุตํ\csromandash{}เอกํ สมยํ ภควา ราชคเห วิหรติ คิชฺฌกูเฏ ปพฺพเต อจิรปกฺกนฺเต เทวทตฺเต. ตตฺร โข ภควา เทวทตฺตํ อารพฺภ ภิกฺขู อามนฺเตสิ\csromandash{}}

\prose{อิธ ภิกฺขเว เอกจฺโจ กุลปุตฺโต สทฺธา อคารสฺมา อนคาริยํ ปพฺพชิโต โหติ “โอติณฺโณมฺหิ ชาติยา ชราย มรเณน โสเกหิ ปริเทเวหิ ทุกฺเขหิ โทมนสฺเสหิ อุปายาเสหิ, ทุกฺโขติณฺโณ ทุกฺขปเรโต. อปฺเปว นาม อิมสฺส เกวลสฺส ทุกฺข\-กฺขนฺธสฺส อนฺตกิริยา ปญฺญาเยถา”ติ. โส เอวํ ปพฺพชิโต สมาโน ลาภ\-สกฺการ\-สิโลกํ อภินิพฺพตฺเตติ. โส เตน ลาภ\-สกฺการ\-สิโลเกน อตฺตมโน โหติ ปริปุณฺณ\-สงฺกปฺโป. โส เตน ลาภ\-สกฺการ\-สิโลเกน อตฺตานุกฺกํเสติ, ปรํ วมฺเภติ “อหมสฺมิ ลาภสกฺการสิโลกวา\csromansharedfootnote{a95d744cf913e764}{ลาภี สิโลกวา (สี, อิ), ลาภี สกฺการ สิโลกวา (สฺยา)}, อิเม ปนญฺเญ ภิกฺขู อปฺปญฺญาตา อปฺเปสกฺขา”ติ. โส เตน ลาภ\-สกฺการ\-สิโลเกน มชฺชติ ปมชฺชติ ปมาทํ อาปชฺชติ, ปมตฺโต สมาโน ทุกฺขํ วิหรติ.}

\prose{เสยฺยถาปิ ภิกฺขเว ปุริโส สารตฺถิโก สารคเวสี สารปริเยสนํ จรมาโน มหโต รุกฺขสฺส ติฏฺฐโต สารวโต อติกฺกมฺเมว สารํ อติกฺกมฺม เผคฺคุํ อติกฺกมฺม ตจํ อติกฺกมฺม ปปฏิกํ สาขาปลาสํ เฉตฺวา อาทาย ปกฺกเมยฺย สารนฺติ มญฺญมาโน. ตเมนํ จกฺขุมา ปุริโส ทิสฺวา เอวํ วเทยฺย “น วตายํ ภวํ ปุริโส อญฺญาสิ สารํ, น อญฺญาสิ เผคฺคุํ, น อญฺญาสิ ตจํ, น อญฺญาสิ ปปฏิกํ, น อญฺญาสิ สาขาปลาสํ. ตถา ห’ยํ\csromansharedfootnote{b12b5215e45627a4}{ตถาปายํ (ก)} ภวํ ปุริโส สารตฺถิโก สารคเวสี สารปริเยสนํ จรมาโน มหโต รุกฺขสฺส ติฏฺฐโต สารวโต อติกฺกมฺเมว สารํ อติกฺกมฺม เผคฺคุํ อติกฺกมฺม ตจํ อติกฺกมฺม ปปฏิกํ สาขาปลาสํ เฉตฺวา อาทาย ปกฺกนฺโต สารนฺติ มญฺญมาโน. ยญฺจสฺส สาเรน สารกรณียํ, ตญฺจสฺส อตฺถํ นานุภวิสฺสตี”ติ. เอวเมว โข ภิกฺขเว อิเธกจฺโจ กุลปุตฺโต สทฺธา อคารสฺมา อนคาริยํ ปพฺพชิโต โหติ “โอติณฺโณมฺหิ ชาติยา ชราย มรเณน โสเกหิ ปริเทเวหิ ทุกฺเขหิ โทมนสฺเสหิ อุปายาเสหิ, ทุกฺโขติณฺโณ \csromanfolio{251}ทุกฺขปเรโต. อปฺเปว นาม อิมสฺส เกวลสฺส ทุกฺข\-กฺขนฺธสฺส อนฺตกิริยา ปญฺญาเยถา”ติ. โส เอวํ ปพฺพชิโต สมาโน ลาภ\-สกฺการ\-สิโลกํ อภินิพฺพตฺเตติ. โส เตน ลาภ\-สกฺการ\-สิโลเกน อตฺตมโน โหติ ปริปุณฺณ\-สงฺกปฺโป. โส เตน ลาภ\-สกฺการ\-สิโลเกน อตฺตานุกฺกํเสติ, ปรํ วมฺเภติ “อหมสฺมิ ลาภสกฺการสิโลกวา, อิเม ปนญฺเญ ภิกฺขู อปฺปญฺญาตา อปฺเปสกฺขา”ติ. โส เตน ลาภ\-สกฺการ\-สิโลเกน มชฺชติ ปมชฺชติ ปมาทํ อาปชฺชติ, ปมตฺโต สมาโน ทุกฺขํ วิหรติ. อยํ วุจฺจติ ภิกฺขเว ภิกฺขุ สาขาปลาสํ อคฺคเหสิ พฺรหฺม\-จริยสฺส, เตน จ โวสานํ อาปาทิ.}

\proseitem{308}{อิธ ปน ภิกฺขเว เอกจฺโจ กุลปุตฺโต สทฺธา อคารสฺมา อนคาริยํ ปพฺพชิโต โหติ “โอติณฺโณมฺหิ ชาติยา ชราย มรเณน โสเกหิ ปริเทเวหิ ทุกฺเขหิ โทมนสฺเสหิ อุปายาเสหิ, ทุกฺโขติณฺโณ ทุกฺขปเรโต. อปฺเปว นาม อิมสฺส เกวลสฺส ทุกฺข\-กฺขนฺธสฺส อนฺตกิริยา ปญฺญาเยถา”ติ. โส เอวํ ปพฺพชิโต สมาโน ลาภ\-สกฺการ\-สิโลกํ อภินิพฺพตฺเตติ. โส เตน ลาภ\-สกฺการ\-สิโลเกน น อตฺตมโน โหติ น ปริปุณฺณ\-สงฺกปฺโป. โส เตน ลาภ\-สกฺการ\-สิโลเกน น อตฺตานุกฺกํเสติ, น ปรํ วมฺเภติ. โส เตน ลาภ\-สกฺการ\-สิโลเกน น มชฺชติ นปฺปมชฺชติ น ปมาทํ อาปชฺชติ, อปฺปมตฺโต สมาโน สีลสมฺปทํ อาราเธติ. โส ตาย สีลสมฺปทาย อตฺตมโน โหติ ปริปุณฺณ\-สงฺกปฺโป. โส ตาย สีลสมฺปทาย อตฺตานุกฺกํเสติ, ปรํ วมฺเภติ “อหมสฺมิ สีลวา กลฺยาณธมฺโม, อิเม ปนญฺเญ ภิกฺขู ทุสฺสีลา ปาปธมฺมา”ติ. โส ตาย สีลสมฺปทาย มชฺชติ ปมชฺชติ ปมาทํ อาปชฺชติ, ปมตฺโต สมาโน ทุกฺขํ วิหรติ.}

\prose{เสยฺยถาปิ ภิกฺขเว ปุริโส สารตฺถิโก สารคเวสี สารปริเยสนํ จรมาโน มหโต รุกฺขสฺส ติฏฺฐโต สารวโต อติกฺกมฺเมว สารํ อติกฺกมฺม เผคฺคุํ อติกฺกมฺม ตจํ ปปฏิกํ เฉตฺวา อาทาย ปกฺกเมยฺย สารนฺติ มญฺญมาโน. ตเมนํ จกฺขุมา ปุริโส ทิสฺวา เอวํ วเทยฺย “น วตายํ ภวํ ปุริโส อญฺญาสิ สารํ, น อญฺญาสิ เผคฺคุํ, น อญฺญาสิ ตจํ, น อญฺญาสิ ปปฏิกํ, น อญฺญาสิ สาขาปลาสํ. ตถา ห’ยํ ภวํ ปุริโส สารตฺถิโก สารคเวสี สารปริเยสนํ จรมาโน มหโต รุกฺขสฺส ติฏฺฐโต สารวโต อติกฺกมฺเมว \csromanfolio{252}สารํ อติกฺกมฺม เผคฺคุํ อติกฺกมฺม ตจํ ปปฏิกํ เฉตฺวา อาทาย ปกฺกนฺโต สารนฺติ มญฺญมาโน. ยญฺจสฺส สาเรน สารกรณียํ, ตญฺจสฺส อตฺถํ นานุภวิสฺสตี”ติ. เอวเมว โข ภิกฺขเว อิเธกจฺโจ กุลปุตฺโต สทฺธา อคารสฺมา อนคาริยํ ปพฺพชิโต โหติ “โอติณฺโณมฺหิ ชาติยา ชราย มรเณน โสเกหิ ปริเทเวหิ ทุกฺเขหิ โทมนสฺเสหิ อุปายาเสหิ, ทุกฺโขติณฺโณ ทุกฺขปเรโต. อปฺเปว นาม อิมสฺส เกวลสฺส ทุกฺข\-กฺขนฺธสฺส อนฺตกิริยา ปญฺญาเยถา”ติ. โส เอวํ ปพฺพชิโต สมาโน ลาภ\-สกฺการ\-สิโลกํ อภินิพฺพตฺเตติ. โส เตน ลาภ\-สกฺการ\-สิโลเกน น อตฺตมโน โหติ น ปริปุณฺณ\-สงฺกปฺโป. โส เตน ลาภ\-สกฺการ\-สิโลเกน น อตฺตานุกฺกํเสติ, น ปรํ วมฺเภติ. โส เตน ลาภ\-สกฺการ\-สิโลเกน น มชฺชติ นปฺปมชฺชติ น ปมาทํ อาปชฺชติ, อปฺปมตฺโต สมาโน สีลสมฺปทํ อาราเธติ. โส ตาย สีลสมฺปทาย อตฺตมโน โหติ ปริปุณฺณ\-สงฺกปฺโป. โส ตาย สีลสมฺปทาย อตฺตานุกฺกํเสติ, ปรํ วมฺเภติ “อหมสฺมิ สีลวา กลฺยาณธมฺโม, อิเม ปนญฺเญ ภิกฺขู ทุสฺสีลา ปาปธมฺมา”ติ. โส ตาย สีลสมฺปทาย มชฺชติ ปมชฺชติ ปมาทํ อาปชฺชติ, ปมตฺโต สมาโน ทุกฺขํ วิหรติ. อยํ วุจฺจติ ภิกฺขเว ภิกฺขุ ปปฏิกํ อคฺคเหสิ พฺรหฺม\-จริยสฺส, เตน จ โวสานํ อาปาทิ.}

\proseitem{309}{อิธ ปน ภิกฺขเว เอกจฺโจ กุลปุตฺโต สทฺธา อคารสฺมา อนคาริยํ ปพฺพชิโต โหติ “โอติณฺโณมฺหิ ชาติยา ชราย มรเณน โสเกหิ ปริเทเวหิ ทุกฺเขหิ โทมนสฺเสหิ อุปายาเสหิ, ทุกฺโขติณฺโณ ทุกฺขปเรโต. อปฺเปว นาม อิมสฺส เกวลสฺส ทุกฺข\-กฺขนฺธสฺส อนฺตกิริยา ปญฺญาเยถา”ติ. โส เอวํ ปพฺพชิโต สมาโน ลาภ\-สกฺการ\-สิโลกํ อภินิพฺพตฺเตติ. โส เตน ลาภ\-สกฺการ\-สิโลเกน น อตฺตมโน โหติ น ปริปุณฺณ\-สงฺกปฺโป. โส เตน ลาภ\-สกฺการ\-สิโลเกน น อตฺตานุกฺกํเสติ, น ปรํ วมฺเภติ. โส เตน ลาภ\-สกฺการ\-สิโลเกน น มชฺชติ นปฺปมชฺชติ น ปมาทํ อาปชฺชติ, อปฺปมตฺโต สมาโน สีลสมฺปทํ อาราเธติ. โส ตาย สีลสมฺปทาย อตฺตมโน โหติ, โน จ โข ปริปุณฺณ\-สงฺกปฺโป. โส ตาย สีลสมฺปทาย น อตฺตานุกฺกํเสติ, น ปรํ วมฺเภติ. โส ตาย สีลสมฺปทาย น มชฺชติ นปฺปมชฺชติ น ปมาทํ อาปชฺชติ, อปฺปมตฺโต สมาโน สมาธิ\-สมฺปทํ \csromanfolio{253}อาราเธติ. โส ตาย สมาธิ\-สมฺปทาย อตฺตมโน โหติ ปริปุณฺณ\-สงฺกปฺโป. โส ตาย สมาธิ\-สมฺปทาย อตฺตานุกฺกํเสติ, ปรํ วมฺเภติ “อหมสฺมิ สมาหิโต เอกคฺคจิตฺโต, อิเม ปนญฺเญ ภิกฺขู อสมาหิตา วิพฺภนฺตจิตฺตา”ติ. โส ตาย สมาธิ\-สมฺปทาย มชฺชติ ปมชฺชติ ปมาทํ อาปชฺชติ, ปมตฺโต สมาโน ทุกฺขํ วิหรติ.}

\prose{เสยฺยถาปิ ภิกฺขเว ปุริโส สารตฺถิโก สารคเวสี สารปริเยสนํ จรมาโน มหโต รุกฺขสฺส ติฏฺฐโต สารวโต อติกฺกมฺเมว สารํ อติกฺกมฺม เผคฺคุํ ตจํ เฉตฺวา อาทาย ปกฺกเมยฺย สารนฺติ มญฺญมาโน. ตเมนํ จกฺขุมา ปุริโส ทิสฺวา เอวํ วเทยฺย “น วตายํ ภวํ ปุริโส อญฺญาสิ สารํ, น อญฺญาสิ เผคฺคุํ, น อญฺญาสิ ตจํ, น อญฺญาสิ ปปฏิกํ, น อญฺญาสิ สาขาปลาสํ. ตถา ห’ยํ ภวํ ปุริโส สารตฺถิโก สารคเวสี สารปริเยสนํ จรมาโน มหโต รุกฺขสฺส ติฏฺฐโต สารวโต อติกฺกมฺเมว สารํ อติกฺกมฺม เผคฺคุํ ตจํ เฉตฺวา อาทาย ปกฺกนฺโต สารนฺติ มญฺญมาโน. ยญฺจสฺส สาเรน สารกรณียํ, ตญฺจสฺส อตฺถํ นานุภวิสฺสตี”ติ. เอวเมว โข ภิกฺขเว อิเธกจฺโจ กุลปุตฺโต สทฺธา อคารสฺมา อนคาริยํ ปพฺพชิโต โหติ “โอติณฺโณมฺหิ ชาติยา ชราย มรเณน โสเกหิ ปริเทเวหิ ทุกฺเขหิ โทมนสฺเสหิ อุปายาเสหิ, ทุกฺโขติณฺโณ ทุกฺขปเรโต. อปฺเปว นาม อิมสฺส เกวลสฺส ทุกฺข\-กฺขนฺธสฺส อนฺตกิริยา ปญฺญาเยถา”ติ. โส เอวํ ปพฺพชิโต สมาโน ลาภ\-สกฺการ\-สิโลกํ อภินิพฺพตฺเตติ. โส เตน ลาภ\-สกฺการ\-สิโลเกน น อตฺตมโน โหติ น ปริปุณฺณ\-สงฺกปฺโป. โส เตน ลาภ\-สกฺการ\-สิโลเกน น อตฺตานุกฺกํเสติ, น ปรํ วมฺเภติ. โส เตน ลาภ\-สกฺการ\-สิโลเกน น มชฺชติ นปฺปมชฺชติ น ปมาทํ อาปชฺชติ, อปฺปมตฺโต สมาโน สีลสมฺปทํ อาราเธติ. โส ตาย สีลสมฺปทาย อตฺตมโน โหติ, โน จ โข ปริปุณฺณ\-สงฺกปฺโป. โส ตาย สีลสมฺปทาย น อตฺตานุกฺกํเสติ, น ปรํ วมฺเภติ. โส ตาย สีลสมฺปทาย น มชฺชติ นปฺปมชฺชติ น ปมาทํ อาปชฺชติ, อปฺปมตฺโต สมาโน สมาธิ\-สมฺปทํ อาราเธติ. โส ตาย สมาธิ\-สมฺปทาย อตฺตมโน โหติ ปริปุณฺณ\-สงฺกปฺโป. โส ตาย สมาธิ\-สมฺปทาย อตฺตานุกฺกํเสติ, ปรํ วมฺเภติ “อหมสฺมิ สมาหิโต เอกคฺคจิตฺโต, อิเม ปนญฺเญ ภิกฺขู อสมาหิตา วิพฺภนฺตจิตฺตา”ติ. โส \csromanfolio{254}ตาย สมาธิ\-สมฺปทาย มชฺชติ ปมชฺชติ ปมาทํ อาปชฺชติ, ปมตฺโต สมาโน ทุกฺขํ วิหรติ. อยํ วุจฺจติ ภิกฺขเว ภิกฺขุ ตจํ อคฺคเหสิ พฺรหฺม\-จริยสฺส, เตน จ โวสานํ อาปาทิ.}

\proseitem{310}{อิธ ปน ภิกฺขเว เอกจฺโจ กุลปุตฺโต สทฺธา อคารสฺมา อนคาริยํ ปพฺพชิโต โหติ “โอติณฺโณมฺหิ ชาติยา ชราย มรเณน โสเกหิ ปริเทเวหิ ทุกฺเขหิ โทมนสฺเสหิ อุปายาเสหิ, ทุกฺโขติณฺโณ ทุกฺขปเรโต. อปฺเปว นาม อิมสฺส เกวลสฺส ทุกฺข\-กฺขนฺธสฺส อนฺตกิริยา ปญฺญาเยถา”ติ. โส เอวํ ปพฺพชิโต สมาโน ลาภ\-สกฺการ\-สิโลกํ อภินิพฺพตฺเตติ. โส เตน ลาภ\-สกฺการ\-สิโลเกน น อตฺตมโน โหติ น ปริปุณฺณ\-สงฺกปฺโป. โส เตน ลาภ\-สกฺการ\-สิโลเกน น อตฺตานุกฺกํเสติ, น ปรํ วมฺเภติ. โส เตน ลาภ\-สกฺการ\-สิโลเกน น มชฺชติ นปฺปมชฺชติ น ปมาทํ อาปชฺชติ, อปฺปมตฺโต สมาโน สีลสมฺปทํ อาราเธติ. โส ตาย สีลสมฺปทาย อตฺตมโน โหติ, โน จ โข ปริปุณฺณ\-สงฺกปฺโป. โส ตาย สีลสมฺปทาย น อตฺตานุกฺกํเสติ, น ปรํ วมฺเภติ. โส ตาย สีลสมฺปทาย น มชฺชติ นปฺปมชฺชติ น ปมาทํ อาปชฺชติ, อปฺปมตฺโต สมาโน สมาธิ\-สมฺปทํ อาราเธติ. โส ตาย สมาธิ\-สมฺปทาย อตฺตมโน โหติ, โน จ โข ปริปุณฺณ\-สงฺกปฺโป. โส ตาย สมาธิ\-สมฺปทาย น อตฺตานุกฺกํเสติ, น ปรํ วมฺเภติ. โส ตาย สมาธิ\-สมฺปทาย น มชฺชติ นปฺปมชฺชติ น ปมาทํ อาปชฺชติ, อปฺปมตฺโต สมาโน ญาณทสฺสนํ อาราเธติ. โส เตน ญาณทสฺสเนน อตฺตมโน โหติ ปริปุณฺณ\-สงฺกปฺโป. โส เตน ญาณทสฺสเนน อตฺตานุกฺกํเสติ, ปรํ วมฺเภติ “อหมสฺมิ ชานํ ปสฺสํ วิหรามิ, อิเม ปนญฺเญ ภิกฺขู อชานํ อปสฺสํ วิหรนฺตี”ติ. โส เตน ญาณทสฺสเนน มชฺชติ ปมชฺชติ ปมาทํ อาปชฺชติ, ปมตฺโต สมาโน ทุกฺขํ วิหรติ.}

\prose{เสยฺยถาปิ ภิกฺขเว ปุริโส สารตฺถิโก สารคเวสี สารปริเยสนํ จรมาโน มหโต รุกฺขสฺส ติฏฺฐโต สารวโต อติกฺกมฺเมว สารํ เผคฺคุํ เฉตฺวา อาทาย ปกฺกเมยฺย สารนฺติ มญฺญมาโน. ตเมนํ จกฺขุมา ปุริโส ทิสฺวา เอวํ วเทยฺย “น วตายํ ภวํ ปุริโส อญฺญาสิ สารํ, น อญฺญาสิ เผคฺคุํ, น อญฺญาสิ ตจํ, น อญฺญาสิ ปปฏิกํ, น อญฺญาสิ สาขาปลาสํ. ตถา ห’ยํ ภวํ ปุริโส สารตฺถิโก สารคเวสี สารปริเยสนํ จรมาโน มหโต \csromanfolio{255}รุกฺขสฺส ติฏฺฐโต สารวโต อติกฺกมฺเมว สารํ เผคฺคุํ เฉตฺวา อาทาย ปกฺกนฺโต สารนฺติ มญฺญมาโน. ยญฺจสฺส สาเรน สารกรณียํ, ตญฺจสฺส อตฺถํ นานุภวิสฺสตี”ติ. เอวเมว โข ภิกฺขเว อิเธกจฺโจ กุลปุตฺโต สทฺธา อคารสฺมา อนคาริยํ ปพฺพชิโต โหติ “โอติณฺโณมฺหิ ชาติยา ชราย มรเณน โสเกหิ ปริเทเวหิ ทุกฺเขหิ โทมนสฺเสหิ อุปายาเสหิ, ทุกฺโขติณฺโณ ทุกฺขปเรโต. อปฺเปว นาม อิมสฺส เกวลสฺส ทุกฺข\-กฺขนฺธสฺส อนฺตกิริยา ปญฺญาเยถา”ติ. โส เอวํ ปพฺพชิโต สมาโน ลาภ\-สกฺการ\-สิโลกํ อภินิพฺพตฺเตติ. โส เตน ลาภ\-สกฺการ\-สิโลเกน น อตฺตมโน โหติ น ปริปุณฺณ\-สงฺกปฺโป. โส เตน ลาภ\-สกฺการ\-สิโลเกน น อตฺตานุกฺกํเสติ, น ปรํ วมฺเภติ. โส เตน ลาภ\-สกฺการ\-สิโลเกน น มชฺชติ นปฺปมชฺชติ น ปมาทํ อาปชฺชติ, อปฺปมตฺโต สมาโน สีลสมฺปทํ อาราเธติ. โส ตาย สีลสมฺปทาย อตฺตมโน โหติ, โน จ โข ปริปุณฺณ\-สงฺกปฺโป. โส ตาย สีลสมฺปทาย น อตฺตานุกฺกํเสติ, น ปรํ วมฺเภติ. โส ตาย สีลสมฺปทาย น มชฺชติ นปฺปมชฺชติ น ปมาทํ อาปชฺชติ, อปฺปมตฺโต สมาโน สมาธิ\-สมฺปทํ อาราเธติ. โส ตาย สมาธิ\-สมฺปทาย อตฺตมโน โหติ, โน จ โข ปริปุณฺณ\-สงฺกปฺโป. โส ตาย สมาธิ\-สมฺปทาย น อตฺตานุกฺกํเสติ, น ปรํ วมฺเภติ. โส ตาย สมาธิ\-สมฺปทาย น มชฺชติ นปฺปมชฺชติ น ปมาทํ อาปชฺชติ, อปฺปมตฺโต สมาโน ญาณทสฺสนํ อาราเธติ. โส เตน ญาณทสฺสเนน อตฺตมโน โหติ ปริปุณฺณ\-สงฺกปฺโป. โส เตน ญาณทสฺสเนน อตฺตานุกฺกํเสติ, ปรํ วมฺเภติ “อหมสฺมิ ชานํ ปสฺสํ วิหรามิ, อิเม ปนญฺเญ ภิกฺขู อชานํ อปสฺสํ วิหรนฺตี”ติ. โส เตน ญาณทสฺสเนน มชฺชติ ปมชฺชติ ปมาทํ อาปชฺชติ, ปมตฺโต สมาโน ทุกฺขํ วิหรติ. อยํ วุจฺจติ ภิกฺขเว ภิกฺขุ เผคฺคุํ อคฺคเหสิ พฺรหฺม\-จริยสฺส, เตน จ โวสานํ อาปาทิ.}

\proseitem{311}{อิธ ปน ภิกฺขเว เอกจฺโจ กุลปุตฺโต สทฺธา อคารสฺมา อนคาริยํ ปพฺพชิโต โหติ “โอติณฺโณมฺหิ ชาติยา ชราย มรเณน โสเกหิ ปริเทเวหิ ทุกฺเขหิ โทมนสฺเสหิ อุปายาเสหิ, ทุกฺโขติณฺโณ ทุกฺขปเรโต. อปฺเปว นาม อิมสฺส เกวลสฺส ทุกฺข\-กฺขนฺธสฺส อนฺตกิริยา ปญฺญาเยถา”ติ. โส เอวํ ปพฺพชิโต สมาโน ลาภ\-สกฺการ\-สิโลกํ อภินิพฺพตฺเตติ. โส เตน ลาภ\-สกฺการ\-สิโลเกน น อตฺตมโน \csromanfolio{256}โหติ น ปริปุณฺณ\-สงฺกปฺโป. โส เตน ลาภ\-สกฺการ\-สิโลเกน น อตฺตานุกฺกํเสติ, น ปรํ วมฺเภติ. โส เตน ลาภ\-สกฺการ\-สิโลเกน น มชฺชติ นปฺปมชฺชติ น ปมาทํ อาปชฺชติ, อปฺปมตฺโต สมาโน สีลสมฺปทํ อาราเธติ. โส ตาย สีลสมฺปทาย อตฺตมโน โหติ, โน จ โข ปริปุณฺณ\-สงฺกปฺโป. โส ตาย สีลสมฺปทาย น อตฺตานุกฺกํเสติ, น ปรํ วมฺเภติ. โส ตาย สีลสมฺปทาย น มชฺชติ นปฺปมชฺชติ น ปมาทํ อาปชฺชติ, อปฺปมตฺโต สมาโน สมาธิ\-สมฺปทํ อาราเธติ. โส ตาย สมาธิ\-สมฺปทาย อตฺตมโน โหติ, โน จ โข ปริปุณฺณ\-สงฺกปฺโป. โส ตาย สมาธิ\-สมฺปทาย น อตฺตานุกฺกํเสติ, น ปรํ วมฺเภติ. โส ตาย สมาธิ\-สมฺปทาย น มชฺชติ นปฺปมชฺชติ น ปมาทํ อาปชฺชติ, อปฺปมตฺโต สมาโน ญาณทสฺสนํ อาราเธติ. โส เตน ญาณทสฺสเนน อตฺตมโน โหติ, โน จ โข ปริปุณฺณ\-สงฺกปฺโป. โส เตน ญาณทสฺสเนน น อตฺตานุกฺกํเสติ, น ปรํ วมฺเภติ. โส เตน ญาณทสฺสเนน น มชฺชติ นปฺปมชฺชติ น ปมาทํ อาปชฺชติ, อปฺปมตฺโต สมาโน อสมย\-วิโมกฺขํ อาราเธติ. อฏฺฐานเมตํ\csromansharedfootnote{7bff170d54889856}{อฏฺฐานํ โข ปเนตํ (ก)} ภิกฺขเว อนวกาโส, ยํ โส ภิกฺขุ ตาย อสมย\-วิมุตฺติยา ปริหาเยถ.}

\prose{เสยฺยถาปิ ภิกฺขเว ปุริโส สารตฺถิโก สารคเวสี สารปริเยสนํ จรมาโน มหโต รุกฺขสฺส ติฏฺฐโต สารวโต สารญฺเญว เฉตฺวา อาทาย ปกฺกเมยฺย สารนฺติ ชานมาโน. ตเมนํ จกฺขุมา ปุริโส ทิสฺวา เอวํ วเทยฺย “อญฺญาสิ วตายํ ภวํ ปุริโส สารํ, อญฺญาสิ เผคฺคุํ, อญฺญาสิ ตจํ, อญฺญาสิ ปปฏิกํ, อญฺญาสิ สาขาปลาสํ. ตถา ห’ยํ ภวํ ปุริโส สารตฺถิโก สารคเวสี สารปริเยสนํ จรมาโน มหโต รุกฺขสฺส ติฏฺฐโต สารวโต สารญฺเญว เฉตฺวา อาทาย ปกฺกนฺโต สารนฺติ ชานมาโน. ยญฺจสฺส สาเรน สารกรณียํ, ตญฺจสฺส อตฺถํ อนุภวิสฺสตี”ติ. เอวเมว โข ภิกฺขเว อิเธกจฺโจ กุลปุตฺโต สทฺธา อคารสฺมา อนคาริยํ ปพฺพชิโต โหติ “โอติณฺโณมฺหิ ชาติยา ชราย มรเณน โสเกหิ ปริเทเวหิ ทุกฺเขหิ โทมนสฺเสหิ อุปายาเสหิ, ทุกฺโขติณฺโณ ทุกฺขปเรโต. อปฺเปว นาม อิมสฺส เกวลสฺส ทุกฺข\-กฺขนฺธสฺส อนฺตกิริยา ปญฺญาเยถา”ติ. \csromanfolio{257}โส เอวํ ปพฺพชิโต สมาโน ลาภ\-สกฺการ\-สิโลกํ อภินิพฺพตฺเตติ. โส เตน ลาภ\-สกฺการ\-สิโลเกน น อตฺตมโน โหติ น ปริปุณฺณ\-สงฺกปฺโป. โส เตน ลาภ\-สกฺการ\-สิโลเกน น อตฺตานุกฺกํเสติ, น ปรํ วมฺเภติ. โส เตน ลาภ\-สกฺการ\-สิโลเกน น มชฺชติ นปฺปมชฺชติ น ปมาทํ อาปชฺชติ, อปฺปมตฺโต สมาโน สีลสมฺปทํ อาราเธติ. โส ตาย สีลสมฺปทาย อตฺตมโน โหติ, โน จ โข ปริปุณฺณ\-สงฺกปฺโป. โส ตาย สีลสมฺปทาย น อตฺตานุกฺกํเสติ, น ปรํ วมฺเภติ. โส ตาย สีลสมฺปทาย น มชฺชติ นปฺปมชฺชติ น ปมาทํ อาปชฺชติ, อปฺปมตฺโต สมาโน สมาธิ\-สมฺปทํ อาราเธติ. โส ตาย สมาธิ\-สมฺปทาย อตฺตมโน โหติ, โน จ โข ปริปุณฺณ\-สงฺกปฺโป. โส ตาย สมาธิ\-สมฺปทาย น อตฺตานุกฺกํเสติ, น ปรํ วมฺเภติ. โส ตาย สมาธิ\-สมฺปทาย น มชฺชติ นปฺปมชฺชติ น ปมาทํ อาปชฺชติ, อปฺปมตฺโต สมาโน ญาณทสฺสนํ อาราเธติ. โส เตน ญาณทสฺสเนน อตฺตมโน โหติ, โน จ โข ปริปุณฺณ\-สงฺกปฺโป. โส เตน ญาณทสฺสเนน น อตฺตานุกฺกํเสติ, น ปรํ วมฺเภติ. โส เตน ญาณทสฺสเนน น มชฺชติ นปฺปมชฺชติ น ปมาทํ อาปชฺชติ, อปฺปมตฺโต สมาโน อสมย\-วิโมกฺขํ อาราเธติ. อฏฺฐานเมตํ ภิกฺขเว อนวกาโส, ยํ โส ภิกฺขุ ตาย อสมย\-วิมุตฺติยา ปริหาเยถ.}

\prose{อิติ โข ภิกฺขเว นยิทํ พฺรหฺมจริยํ ลาภสกฺการสิโลกา\-นิสํสํ, น สีลสมฺปทา\-นิสํสํ, น สมาธิสมฺปทา\-นิสํสํ, น ญาณทสฺสนา\-นิสํสํ. ยา จ โข อยํ ภิกฺขเว อกุปฺปา เจโตวิมุตฺติ, เอตทตฺถมิทํ ภิกฺขเว พฺรหฺมจริยํ, เอตํ สารํ, เอตํ ปริโยสานนฺติ.}

\prosewithcloserruled{อิทมโวจ ภควา. อตฺตมนา เต ภิกฺขู ภควโต ภาสิตํ อภินนฺทุนฺติ.}{มหา\-สาโร\-ปม\-สุตฺตํ นิฏฺฐิตํ นวมํ.}
\csromanmatikaanchor{mtk.60}
\csromantocmarkref{subsection}{10. จูฬสาโรปมสุตฺต}{mtk.60}
\bookmark[dest={mtk.60},level=2]{10. จูฬสาโรปมสุตฺต}
\csromanheader{2}{10. จูฬสาโร\-ปม\-สุตฺต}
\proseitem{312}{เอวํ เม สุตํ\csromandash{}เอกํ สมยํ ภควา สาวตฺถิยํ วิหรติ เชตวเน อนาถ\-ปิณฺฑิกสฺส อาราเม. อถ โข ปิงฺคลโกจฺโฉ พฺราหฺมโณ \csromanfolio{258}เยน ภควา เตนุปสงฺกมิ, อุปสงฺกมิตฺวา ภควตา สทฺธิํ สมฺโมทิ. สมฺโมทนียํ กถํ สารณียํ วีติสาเรตฺวา เอกมนฺตํ นิสีทิ, เอกมนฺตํ นิสินฺโน โข ปิงฺคลโกจฺโฉ พฺราหฺมโณ ภควนฺตํ เอตทโวจ “เยเม โภ โคตม สมณพฺราหฺมณา สงฺฆิโน คณิโน คณาจริยา ญาตา ยสสฺสิโน ติตฺถกรา สาธุสมฺมตา พหุชนสฺส. เสยฺยถิทํ, ปูรโณ กสฺสโป มกฺขลิ โคสาโล อชิโต เกสกมฺพโล ปกุโธ กจฺจายโน สญฺจโย\csromansharedfootnote{6bd794066d568ae0}{สญฺชโย (สี, สฺยา, อิ, ก)} เพลฏฺฐปุตฺโต นิคณฺโฐ นาฏปุตฺโต. สพฺเพเต สกาย ปฏิญฺญาย อพฺภญฺญํสุ, สพฺเพว นาพฺภญฺญํสุ, อุทาหุ เอกจฺเจ อพฺภญฺญํสุ, เอกจฺเจ นาพฺภญฺญํสู”ติ. อลํ พฺราหฺมณ ติฏฺฐเตตํ “สพฺเพเต สกาย ปฏิญฺญาย อพฺภญฺญํสุ, สพฺเพว นาพฺภญฺญํสุ, อุทาหุ เอกจฺเจ อพฺภญฺญํสุ, เอกจฺเจ นาพฺภญฺญํสู”ติ. ธมฺมํ เต พฺราหฺมณ เทเสสฺสามิ, ตํ สุณาหิ สาธุกํ มนสิกโรหิ ภาสิสฺสามีติ. “เอวํ โภ”ติ โข ปิงฺคลโกจฺโฉ พฺราหฺมโณ ภควโต ปจฺจสฺโสสิ. ภควา เอตทโวจ\csromandash{}}

\proseitem{313}{เสยฺยถาปิ พฺราหฺมณ ปุริโส สารตฺถิโก สารคเวสี สารปริเยสนํ จรมาโน มหโต รุกฺขสฺส ติฏฺฐโต สารวโต อติกฺกมฺเมว สารํ อติกฺกมฺม เผคฺคุํ อติกฺกมฺม ตจํ อติกฺกมฺม ปปฏิกํ สาขาปลาสํ เฉตฺวา อาทาย ปกฺกเมยฺย สารนฺติ มญฺญมาโน. ตเมนํ จกฺขุมา ปุริโส ทิสฺวา เอวํ วเทยฺย “น วตายํ ภวํ ปุริโส อญฺญาสิ สารํ, น อญฺญาสิ เผคฺคุํ, น อญฺญาสิ ตจํ, น อญฺญาสิ ปปฏิกํ, น อญฺญาสิ สาขาปลาสํ. ตถา ห’ยํ ภวํ ปุริโส สารตฺถิโก สารคเวสี สารปริเยสนํ จรมาโน มหโต รุกฺขสฺส ติฏฺฐโต สารวโต อติกฺกมฺเมว สารํ อติกฺกมฺม เผคฺคุํ อติกฺกมฺม ตจํ อติกฺกมฺม ปปฏิกํ สาขาปลาสํ เฉตฺวา อาทาย ปกฺกนฺโต สารนฺติ มญฺญมาโน. ยญฺจสฺส สาเรน สารกรณียํ, ตญฺจสฺส อตฺถํ นานุภวิสฺสตี”ติ.}

\proseitem{314}{เสยฺยถาปิ วา ปน พฺราหฺมณ ปุริโส สารตฺถิโก สารคเวสี สารปริเยสนํ จรมาโน มหโต รุกฺขสฺส ติฏฺฐโต สารวโต อติกฺกมฺเมว สารํ อติกฺกมฺม เผคฺคุํ อติกฺกมฺม ตจํ ปปฏิกํ เฉตฺวา \csromanfolio{259}อาทาย ปกฺกเมยฺย สารนฺติ มญฺญมาโน. ตเมนํ จกฺขุมา ปุริโส ทิสฺวา เอวํ วเทยฺย “น วตายํ ภวํ ปุริโส อญฺญาสิ สารํ, น อญฺญาสิ เผคฺคุํ, น อญฺญาสิ ตจํ, น อญฺญาสิ ปปฏิกํ, น อญฺญาสิ สาขาปลาสํ. ตถา ห’ยํ ภวํ ปุริโส สารตฺถิโก สารคเวสี สารปริเยสนํ จรมาโน มหโต รุกฺขสฺส ติฏฺฐโต สารวโต อติกฺกมฺเมว สารํ อติกฺกมฺม เผคฺคุํ อติกฺกมฺม ตจํ ปปฏิกํ เฉตฺวา อาทาย ปกฺกนฺโต สารนฺติ มญฺญมาโน. ยญฺจสฺส สาเรน สารกรณียํ, ตญฺจสฺส อตฺถํ นานุภวิสฺสตี”ติ.}

\proseitem{315}{เสยฺยถาปิ วา ปน พฺราหฺมณ ปุริโส สารตฺถิโก สารคเวสี สารปริเยสนํ จรมาโน มหโต รุกฺขสฺส ติฏฺฐโต สารวโต อติกฺกมฺเมว สารํ อติกฺกมฺม เผคฺคุํ ตจํ เฉตฺวา อาทาย ปกฺกเมยฺย สารนฺติ มญฺญมาโน. ตเมนํ จกฺขุมา ปุริโส ทิสฺวา เอวํ วเทยฺย “น วตายํ ภวํ ปุริโส อญฺญาสิ สารํ, น อญฺญาสิ เผคฺคุํ, น อญฺญาสิ ตจํ, น อญฺญาสิ ปปฏิกํ, น อญฺญาสิ สาขาปลาสํ. ตถา ห’ยํ ภวํ ปุริโส สารตฺถิโก สารคเวสี สารปริเยสนํ จรมาโน มหโต รุกฺขสฺส ติฏฺฐโต สารวโต อติกฺกมฺเมว สารํ อติกฺกมฺม เผคฺคุํ ตจํ เฉตฺวา อาทาย ปกฺกนฺโต สารนฺติ มญฺญมาโน. ยญฺจสฺส สาเรน สารกรณียํ, ตญฺจสฺส อตฺถํ นานุภวิสฺสตี”ติ.}

\proseitem{316}{เสยฺยถาปิ วา ปน พฺราหฺมณ ปุริโส สารตฺถิโก สารคเวสี สารปริเยสนํ จรมาโน มหโต รุกฺขสฺส ติฏฺฐโต สารวโต อติกฺกมฺเมว สารํ เผคฺคุํ เฉตฺวา อาทาย ปกฺกเมยฺย สารนฺติ มญฺญมาโน. ตเมนํ จกฺขุมา ปุริโส ทิสฺวา เอวํ วเทยฺย “น วตายํ ภวํ ปุริโส อญฺญาสิ สารํ, น อญฺญาสิ เผคฺคุํ, น อญฺญาสิ ตจํ, น อญฺญาสิ ปปฏิกํ, น อญฺญาสิ สาขาปลาสํ. ตถา ห’ยํ ภวํ ปุริโส สารตฺถิโก สารคเวสี สารปริเยสนํ จรมาโน มหโต รุกฺขสฺส ติฏฺฐโต สารวโต อติกฺกมฺเมว สารํ เผคฺคุํ เฉตฺวา อาทาย ปกฺกนฺโต สารนฺติ มญฺญมาโน. ยญฺจสฺส สาเรน สารกรณียํ, ตญฺจสฺส อตฺถํ นานุภวิสฺสตี”ติ.}

\proseitem{317}{เสยฺยถาปิ วา ปน พฺราหฺมณ ปุริโส สารตฺถิโก สารคเวสี สารปริเยสนํ จรมาโน มหโต รุกฺขสฺส ติฏฺฐโต สารวโต \csromanfolio{260}สารญฺเญว เฉตฺวา อาทาย ปกฺกเมยฺย สารนฺติ ชานมาโน. ตเมนํ จกฺขุมา ปุริโส ทิสฺวา เอวํ วเทยฺย “อญฺญาสิ วตายํ ภวํ ปุริโส สารํ, อญฺญาสิ เผคฺคุํ, อญฺญาสิ ตจํ, อญฺญาสิ ปปฏิกํ, อญฺญาสิ สาขาปลาสํ. ตถา ห’ยํ ภวํ ปุริโส สารตฺถิโก สารคเวสี สารปริเยสนํ จรมาโน มหโต รุกฺขสฺส ติฏฺฐโต สารวโต สารญฺเญว เฉตฺวา อาทาย ปกฺกนฺโต สารนฺติ ชานมาโน. ยญฺจสฺส สาเรน สารกรณียํ, ตญฺจสฺส อตฺถํ อนุภวิสฺสตี”ติ.}

\proseitem{318}{เอวเมว โข พฺราหฺมณ อิเธกจฺโจ ปุคฺคโล สทฺธา อคารสฺมา อนคาริยํ ปพฺพชิโต โหติ “โอติณฺโณมฺหิ ชาติยา ชราย มรเณน โสเกหิ ปริเทเวหิ ทุกฺเขหิ โทมนสฺเสหิ อุปายาเสหิ, ทุกฺโขติณฺโณ ทุกฺขปเรโต. อปฺเปว นาม อิมสฺส เกวลสฺส ทุกฺข\-กฺขนฺธสฺส อนฺตกิริยา ปญฺญาเยถา”ติ. โส เอวํ ปพฺพชิโต สมาโน ลาภ\-สกฺการ\-สิโลกํ อภินิพฺพตฺเตติ. โส เตน ลาภ\-สกฺการ\-สิโลเกน อตฺตมโน โหติ ปริปุณฺณ\-สงฺกปฺโป. โส เตน ลาภ\-สกฺการ\-สิโลเกน อตฺตานุกฺกํเสติ, ปรํ วมฺเภติ “อหมสฺมิ ลาภสกฺการสิโลกวา, อิเม ปนญฺเญ ภิกฺขู อปฺปญฺญาตา อปฺเปสกฺขา”ติ. ลาภ\-สกฺการ\-สิโลเกน จ เย อญฺเญ ธมฺมา อุตฺตริตรา จ ปณีตตรา จ, เตสํ ธมฺมานํ สจฺฉิกิริยาย น ฉนฺทํ ชเนติ, น วายมติ, โอลีนวุตฺติโก จ โหติ สาถลิโก. เสยฺยถาปิ โส พฺราหฺมณ ปุริโส สารตฺถิโก สารคเวสี สารปริเยสนํ จรมาโน มหโต รุกฺขสฺส ติฏฺฐโต สารวโต อติกฺกมฺเมว สารํ อติกฺกมฺม เผคฺคุํ อติกฺกมฺม ตจํ อติกฺกมฺม ปปฏิกํ สาขาปลาสํ เฉตฺวา อาทาย ปกฺกนฺโต สารนฺติ มญฺญมาโน. ยญฺจสฺส สาเรน สารกรณียํ, ตญฺจสฺส อตฺถํ นานุภวิสฺสติ. ตถูปมาหํ พฺราหฺมณ อิมํ ปุคฺคลํ วทามิ.}

\proseitem{319}{อิธ ปน พฺราหฺมณ เอกจฺโจ ปุคฺคโล สทฺธา อคารสฺมา อนคาริยํ ปพฺพชิโต โหติ “โอติณฺโณมฺหิ ชาติยา ชราย มรเณน โสเกหิ ปริเทเวหิ ทุกฺเขหิ โทมนสฺเสหิ อุปายาเสหิ, ทุกฺโขติณฺโณ ทุกฺขปเรโต. อปฺเปว นาม อิมสฺส เกวลสฺส ทุกฺข\-กฺขนฺธสฺส อนฺตกิริยา ปญฺญาเยถา”ติ. โส เอวํ ปพฺพชิโต สมาโน ลาภ\-สกฺการ\-สิโลกํ อภินิพฺพตฺเตติ. โส เตน ลาภ\-สกฺการ\-สิโลเกน \csromanfolio{261}น อตฺตมโน โหติ น ปริปุณฺณ\-สงฺกปฺโป. โส เตน ลาภ\-สกฺการ\-สิโลเกน น อตฺตานุกฺกํเสติ, น ปรํ วมฺเภติ. ลาภ\-สกฺการ\-สิโลเกน จ เย อญฺเญ ธมฺมา อุตฺตริตรา จ ปณีตตรา จ, เตสํ ธมฺมานํ สจฺฉิกิริยาย ฉนฺทํ ชเนติ วายมติ, อโนลีนวุตฺติโก จ โหติ อสาถลิโก. โส สีลสมฺปทํ อาราเธติ. โส ตาย สีลสมฺปทาย อตฺตมโน โหติ ปริปุณฺณ\-สงฺกปฺโป. โส ตาย สีลสมฺปทาย อตฺตานุกฺกํเสติ, ปรํ วมฺเภติ “อหมสฺมิ สีลวา กลฺยาณธมฺโม, อิเม ปนญฺเญ ภิกฺขู ทุสฺสีลา ปาปธมฺมา”ติ. สีลสมฺปทาย จ เย อญฺเญ ธมฺมา อุตฺตริตรา จ ปณีตตรา จ, เตสํ ธมฺมานํ สจฺฉิกิริยาย น ฉนฺทํ ชเนติ น วายมติ, โอลีนวุตฺติโก จ โหติ สาถลิโก. เสยฺยถาปิ โส พฺราหฺมณ ปุริโส สารตฺถิโก สารคเวสี สารปริเยสนํ จรมาโน มหโต รุกฺขสฺส ติฏฺฐโต สารวโต อติกฺกมฺเมว สารํ อติกฺกมฺม เผคฺคุํ อติกฺกมฺม ตจํ ปปฏิกํ เฉตฺวา อาทาย ปกฺกนฺโต สารนฺติ มญฺญมาโน. ยญฺจสฺส สาเรน สารกรณียํ, ตญฺจสฺส อตฺถํ นานุภวิสฺสติ. ตถูปมาหํ พฺราหฺมณ อิมํ ปุคฺคลํ}

\proseitem{320}{อิธ ปน พฺราหฺมณ เอกจฺโจ ปุคฺคโล สทฺธา อคารสฺมา อนคาริยํ ปพฺพชิโต โหติ “โอติณฺโณมฺหิ ชาติยา ชราย มรเณน โสเกหิ ปริเทเวหิ ทุกฺเขหิ โทมนสฺเสหิ อุปายาเสหิ, ทุกฺโขติณฺโณ ทุกฺขปเรโต. อปฺเปว นาม อิมสฺส เกวลสฺส ทุกฺข\-กฺขนฺธสฺส อนฺตกิริยา ปญฺญาเยถา”ติ. โส เอวํ ปพฺพชิโต สมาโน ลาภ\-สกฺการ\-สิโลกํ อภินิพฺพตฺเตติ. โส เตน ลาภ\-สกฺการ\-สิโลเกน น อตฺตมโน โหติ น ปริปุณฺณ\-สงฺกปฺโป. โส เตน ลาภ\-สกฺการ\-สิโลเกน น อตฺตานุกฺกํเสติ, น ปรํ วมฺเภติ. ลาภ\-สกฺการ\-สิโลเกน จ เย อญฺเญ ธมฺมา อุตฺตริตรา จ ปณีตตรา จ, เตสํ ธมฺมานํ สจฺฉิกิริยาย ฉนฺทํ ชเนติ วายมติ, อโนลีนวุตฺติโก จ โหติ อสาถลิโก. โส สีลสมฺปทํ อาราเธติ. โส ตาย สีลสมฺปทาย อตฺตมโน โหติ, โน จ โข ปริปุณฺณ\-สงฺกปฺโป. โส ตาย สีลสมฺปทาย น อตฺตานุกฺกํเสติ, น ปรํ วมฺเภติ. สีลสมฺปทาย จ เย อญฺเญ ธมฺมา อุตฺตริตรา จ ปณีตตรา จ, เตสํ ธมฺมานํ สจฺฉิกิริยาย ฉนฺทํ ชเนติ วายมติ, อโนลีนวุตฺติโก จ โหติ \csromanfolio{262}อสาถลิโก. โส สมาธิ\-สมฺปทํ อาราเธติ. โส ตาย สมาธิ\-สมฺปทาย อตฺตมโน โหติ ปริปุณฺณ\-สงฺกปฺโป. โส ตาย สมาธิ\-สมฺปทาย อตฺตานุกฺกํเสติ, ปรํ วมฺเภติ “อหมสฺมิ สมาหิโต เอกคฺคจิตฺโต, อิเม ปนญฺเญ ภิกฺขู อสมาหิตา วิพฺภนฺตจิตฺตา”ติ. สมาธิ\-สมฺปทาย จ เย อญฺเญ ธมฺมา อุตฺตริตรา จ ปณีตตรา จ, เตสํ ธมฺมานํ สจฺฉิกิริยาย น ฉนฺทํ ชเนติ น วายมติ, โอลีนวุตฺติโก จ โหติ สาถลิโก. เสยฺยถาปิ โส พฺราหฺมณ ปุริโส สารตฺถิโก สารคเวสี สารปริเยสนํ จรมาโน มหโต รุกฺขสฺส ติฏฺฐโต สารวโต อติกฺกมฺเมว สารํ อติกฺกมฺม เผคฺคุํ ตจํ เฉตฺวา อาทาย ปกฺกนฺโต สารนฺติ มญฺญมาโน. ยญฺจสฺส สาเรน สารกรณียํ, ตญฺจสฺส อตฺถํ นานุภวิสฺสติ, ตถูปมาหํ พฺราหฺมณ อิมํ ปุคฺคลํ}

\proseitem{321}{อิธ ปน พฺราหฺมณ เอกจฺโจ ปุคฺคโล สทฺธา อคารสฺมา อนคาริยํ ปพฺพชิโต โหติ “โอติณฺโณมฺหิ ชาติยา ชราย มรเณน ฯเปฯ อนฺตกิริยา ปญฺญาเยถา”ติ. โส เอวํ ปพฺพชิโต สมาโน ลาภ\-สกฺการ\-สิโลกํ อภินิพฺพตฺเตติ. โส เตน ลาภ\-สกฺการ\-สิโลเกน น อตฺตมโน โหติ, น ปริปุณฺณ\-สงฺกปฺโป. โส เตน ลาภ\-สกฺการ\-สิโลเกน น อตฺตานุกฺกํเสติ, น ปรํ วมฺเภติ. ลาภ\-สกฺการ\-สิโลเกน จ เย อญฺเญ ธมฺมา อุตฺตริตรา จ ปณีตตรา จ, เตสํ ธมฺมานํ สจฺฉิกิริยาย ฉนฺทํ ชเนติ วายมติ, อโนลีนวุตฺติโก จ โหติ อสาถลิโก. โส สีลสมฺปทํ อาราเธติ. โส ตาย สีลสมฺปทาย อตฺตมโน โหติ, โน จ โข ปริปุณฺณ\-สงฺกปฺโป. โส ตาย สีลสมฺปทาย น อตฺตานุกฺกํเสติ, น ปรํ วมฺเภติ. สีลสมฺปทาย จ เย อญฺเญ ธมฺมา อุตฺตริตรา จ ปณีตตรา จ, เตสํ ธมฺมานํ สจฺฉิกิริยาย ฉนฺทํ ชเนติ วายมติ, อโนลีนวุตฺติโก จ โหติ อสาถลิโก. โส สมาธิ\-สมฺปทํ อาราเธติ. โส ตาย สมาธิ\-สมฺปทาย อตฺตมโน โหติ, โน จ โข ปริปุณฺณ\-สงฺกปฺโป. โส ตาย สมาธิ\-สมฺปทาย น อตฺตานุกฺกํเสติ, น ปรํ วมฺเภติ. สมาธิ\-สมฺปทาย จ เย อญฺเญ ธมฺมา อุตฺตริตรา จ ปณีตตรา จ, เตสํ ธมฺมานํ สจฺฉิกิริยาย ฉนฺทํ ชเนติ วายมติ, อโนลีนวุตฺติโก จ โหติ อสาถลิโก. โส ญาณทสฺสนํ อาราเธติ. โส เตน ญาณทสฺสเนน อตฺตมโน โหติ, \csromanfolio{263}ปริปุณฺณ\-สงฺกปฺโป. โส เตน ญาณทสฺสเนน อตฺตานุกฺกํเสติ, ปรํ วมฺเภติ “อหมสฺมิ ชานํ ปสฺสํ วิหรามิ, อิเม ปนญฺเญ ภิกฺขุ อชานํ อปสฺสํ วิหรนฺตี”ติ. ญาณทสฺสเนน จ เย อญฺเญ ธมฺมา อุตฺตริตรา จ ปณีตตรา จ, เตสํ ธมฺมานํ สจฺฉิกิริยาย น ฉนฺทํ ชเนติ น วายมติ, โอลีนวุตฺติโก จ โหติ สาถลิโก. เสยฺยถาปิ โส พฺราหฺมณ ปุริโส สารตฺถิโก สารคเวสี สารปริเยสนํ จรมาโน มหโต รุกฺขสฺส ติฏฺฐโต สารวโต อติกฺกมฺเมว สารํ เผคฺคุํ เฉตฺวา อาทาย ปกฺกนฺโต สารนฺติ มญฺญมาโน. ยญฺจสฺส สาเรน สารกรณียํ, ตญฺจสฺส อตฺถํ นานุภวิสฺสติ, ตถูปมาหํ พฺราหฺมณ อิมํ ปุคฺคลํ วทามิ.}

\proseitem{322}{อิธ ปน พฺราหฺมณ เอกจฺโจ ปุคฺคโล สทฺธา อคารสฺมา อนคาริยํ ปพฺพชิโต โหติ “โอติณฺโณมฺหิ ชาติยา ชราย มรเณน โสเกหิ ปริเทเวหิ ทุกฺเขหิ โทมนสฺเสหิ อุปายาเสหิ, ทุกฺโขติณฺโณ ทุกฺขปเรโต. อปฺเปว นาม อิมสฺส เกวลสฺส ทุกฺข\-กฺขนฺธสฺส อนฺตกิริยา ปญฺญาเยถา”ติ. โส เอวํ ปพฺพชิโต สมาโน ลาภ\-สกฺการ\-สิโลกํ อภินิพฺพตฺเตติ. โส เตน ลาภ\-สกฺการ\-สิโลเกน น อตฺตมโน โหติ, น ปริปุณฺณ\-สงฺกปฺโป. โส เตน ลาภ\-สกฺการ\-สิโลเกน น อตฺตานุกฺกํเสติ, น ปรํ วมฺเภติ. ลาภ\-สกฺการ\-สิโลเกน จ เย อญฺเญ ธมฺมา อุตฺตริตรา จ ปณีตตรา จ, เตสํ ธมฺมานํ สจฺฉิกิริยาย ฉนฺทํ ชเนติ วายมติ, อโนลีนวุตฺติโก จ โหติ อสาถลิโก. โส สีลสมฺปทํ อาราเธติ. โส ตาย สีลสมฺปทาย อตฺตมโน โหติ, โน จ โข ปริปุณฺณ\-สงฺกปฺโป. โส ตาย สีลสมฺปทาย น อตฺตานุกฺกํเสติ, น ปรํ วมฺเภติ. สีลสมฺปทาย จ เย อญฺเญ ธมฺมา อุตฺตริตรา จ ปณีตตรา จ, เตสํ ธมฺมานํ สจฺฉิกิริยาย ฉนฺทํ ชเนติ วายมติ, อโนลีนวุตฺติโก จ โหติ อสาถลิโก. โส สมาธิ\-สมฺปทํ อาราเธติ. โส ตาย สมาธิ\-สมฺปทาย อตฺตมโน โหติ, โน จ โข ปริปุณฺณ\-สงฺกปฺโป. โส ตาย สมาธิ\-สมฺปทาย น อตฺตานุกฺกํเสติ, น ปรํ วมฺเภติ. สมาธิ\-สมฺปทาย จ เย อญฺเญ ธมฺมา อุตฺตริตรา จ ปณีตตรา จ, เตสํ ธมฺมานํ สจฺฉิกิริยาย ฉนฺทํ ชเนติ วายมติ, อโนลีนวุตฺติโก จ โหติ อสาถลิโก. โส ญาณทสฺสนํ อาราเธติ. โส เตน ญาณทสฺสเนน อตฺตมโน โหติ, โน จ โข ปริปุณฺณ\-สงฺกปฺโป. โส เตน ญาณทสฺสเนน น อตฺตานุกฺกํเสติ, \csromanfolio{264}น ปรํ วมฺเภติ. ญาณทสฺสเนน จ เย อญฺเญ ธมฺมา อุตฺตริตรา จ ปณีตตรา จ, เตสํ ธมฺมานํ สจฺฉิกิริยาย ฉนฺทํ ชเนติ วายมติ, อโนลีนวุตฺติโก จ โหติ อสาถลิโก.}

\proseitem{323}{กตเม จ พฺราหฺมณ ธมฺมา ญาณทสฺสเนน อุตฺตริตรา จ ปณีตตรา จ. อิธ พฺราหฺมณ ภิกฺขุ วิวิจฺเจว กาเมหิ วิวิจฺจ อกุสเลหิ ธมฺเมหิ สวิตกฺกํ สวิจารํ วิเวกชํ ปีติสุขํ ปฐมํ ฌานํ อุปสมฺปชฺช วิหรติ. อยมฺปิ โข พฺราหฺมณ ธมฺโม ญาณทสฺสเนน อุตฺตริตโร จ ปณีตตโร จ.}

\prose{ปุน จปรํ พฺราหฺมณ ภิกฺขุ วิตกฺก\-วิจารานํ วูปสมา อชฺฌตฺตํ สมฺปสาทนํ เจตโส เอโกทิภาวํ อวิตกฺกํ อวิจารํ สมาธิชํ ปีติสุขํ ทุติยํ ฌานํ อุปสมฺปชฺช วิหรติ. อยมฺปิ โข พฺราหฺมณ ธมฺโม ญาณทสฺสเนน อุตฺตริตโร จ ปณีตตโร จ.}

\prose{ปุน จปรํ พฺราหฺมณ ภิกฺขุ ปีติยา จ วิราคา อุเปกฺขโก จ วิหรติ สโต จ สมฺปชาโน, สุขญฺจ กาเยน ปฏิสํเวเทติ, ยํ ตํ อริยา อาจิกฺขนฺติ “อุเปกฺขโก สติมา สุขวิหารี”ติ ตติยํ ฌานํ อุปสมฺปชฺช วิหรติ. อยมฺปิ โข พฺราหฺมณ ธมฺโม ญาณทสฺสเนน อุตฺตริตโร จ ปณีตตโร จ.}

\prose{ปุน จปรํ พฺราหฺมณ ภิกฺขุ สุขสฺส จ ปหานา ทุกฺขสฺส จ ปหานา ปุพฺเพว โสมนสฺส\-โทมนสฺสานํ อตฺถงฺคมา อทุกฺขมสุขํ อุเปกฺขา\-สติ\-ปาริสุทฺธิํ จตุตฺถํ ฌานํ อุปสมฺปชฺช วิหรติ. อยมฺปิ โข พฺราหฺมณ ธมฺโม ญาณทสฺสเนน อุตฺตริตโร จ ปณีตตโร จ.}

\prose{ปุน จปรํ พฺราหฺมณ ภิกฺขุ สพฺพโส รูปสญฺญานํ สมติกฺกมา ปฏิฆ\-สญฺญานํ อตฺถงฺคมา นานตฺต\-สญฺญานํ อมนสิการา “อนนฺโต อากาโส”ติ อากาสา\-นญฺจา\-ยตนํ อุปสมฺปชฺช วิหรติ. อยมฺปิ โข พฺราหฺมณ ธมฺโม ญาณทสฺสเนน อุตฺตริตโร จ ปณีตตโร จ.}

\prose{ปุน จปรํ พฺราหฺมณ ภิกฺขุ สพฺพโส อากาสา\-นญฺจา\-ยตนํ สมติกฺกมฺม “อนนฺตํ วิญฺญาณนฺ”ติ วิญฺญาณญฺจา\-ยตนํ อุปสมฺปชฺช วิหรติ. อยมฺปิ โข พฺราหฺมณ ธมฺโม ญาณทสฺสเนน อุตฺตริตโร จ ปณีตตโร จ.}

\prose{ปุน จปรํ พฺราหฺมณ ภิกฺขุ สพฺพโส วิญฺญาณญฺจา\-ยตนํ สมติกฺกมฺม “นตฺถิ กิญฺจี”ติ อากิญฺจญฺญา\-ยตนํ อุปสมฺปชฺช วิหรติ. อยมฺปิ โข พฺราหฺมณ ธมฺโม ญาณทสฺสเนน อุตฺตริตโร จ ปณีตตโร จ.}

\csromanmatikaanchor{mtk.61}
\csromantocmarkref{subsection}{อุทฺทานคาถา}{mtk.61}
\bookmark[dest={mtk.61},level=2]{อุทฺทานคาถา}
\prose{\csromanfolio{265}ปุน จปรํ พฺราหฺมณ ภิกฺขุ สพฺพโส อากิญฺจญฺญา\-ยตนํ สมติกฺกมฺม เนว\-สญฺญา\-นา\-สญฺญา\-ยตนํ อุปสมฺปชฺช วิหรติ. อยมฺปิ โข พฺราหฺมณ ธมฺโม ญาณทสฺสเนน อุตฺตริตโร จ ปณีตตโร จ.}

\prose{ปุน จปรํ พฺราหฺมณ ภิกฺขุ สพฺพโส เนว\-สญฺญา\-นา\-สญฺญา\-ยตนํ สมติกฺกมฺม สญฺญาเวทยิต\-นิโรธํ อุปสมฺปชฺช วิหรติ. ปญฺญาย จสฺส ทิสฺวา อาสวา ปริกฺขีณา โหนฺติ. อยมฺปิ โข พฺราหฺมณ ธมฺโม ญาณทสฺสเนน อุตฺตริตโร จ ปณีตตโร จ. อิเม โข พฺราหฺมณ ธมฺมา ญาณทสฺสเนน อุตฺตริตรา จ ปณีตตรา จ.}

\proseitem{324}{เสยฺยถาปิ โส พฺราหฺมณ ปุริโส สารตฺถิโก สารคเวสี สารปริเยสนํ จรมาโน มหโต รุกฺขสฺส ติฏฺฐโต สารวโต สารํเยว เฉตฺวา อาทาย ปกฺกนฺโต สารนฺติ ชานมาโน. ยญฺจสฺส สาเรน สารกรณียํ, ตญฺจสฺส อตฺถํ อนุภวิสฺสติ, ตถูปมาหํ พฺราหฺมณ อิมํ ปุคฺคลํ}

\prose{อิติ โข พฺราหฺมณ นยิทํ พฺรหฺมจริยํ ลาภสกฺการสิโลกา\-นิสํสํ, น สีลสมฺปทา\-นิสํสํ, น สมาธิสมฺปทา\-นิสํสํ, น ญาณทสฺสนา\-นิสํสํ. ยา จ โข อยํ พฺราหฺมณ อกุปฺปา เจโตวิมุตฺติ, เอตทตฺถมิทํ พฺราหฺมณ พฺรหฺมจริยํ, เอตํ สารํ, เอตํ ปริโยสานนฺติ.}

\prosewithcloser{เอวํ วุตฺเต ปิงฺคลโกจฺโฉ พฺราหฺมโณ ภควนฺตํ เอตทโวจ “อภิกฺกนฺตํ โภ โคตม, อภิกฺกนฺตํ โภ โคตม ฯเปฯ อุปาสกํ มํ ภวํ โคตโม ธาเรตุ อชฺชตคฺเค ปาณุเปตํ สรณํ คตนฺ”ติ.}{จูฬสาโร\-ปม\-สุตฺตํ นิฏฺฐิตํ ทสมํ.}
\nitthitammid{\textbf{โอปมฺมวคฺโค นิฏฺฐิโต ตติโย.}}
\csromancenter{\textbf{ตสฺสุทฺทานํ}}
\setlength{\csromangathapagewidth}{0pt}
\setlength{\csromangathawakwidth}{0pt}
\csromangathameasuregroup{โมฬิยผคฺคุนริฏฺฐํ จ นาโม, \\ ราสิกเณรุมหาคชนาโม,}{\csromangathabat{โมฬิยผคฺคุนริฏฺฐํ จ นาโม,}{อนฺธวเน กถิปุณฺณํ นิวาโป.} \\ \csromangathabat{ราสิกเณรุมหาคชนาโม,}{สารูปโม\textsuperscript{1} ปุน ปิงฺคลโกจฺโฉ.}}
\csromangathasetleft{โมฬิยผคฺคุนริฏฺฐํ จ นาโม, \\ ราสิกเณรุมหาคชนาโม,}
\csromangathagroup{\csromangathastanza{\csromangathabat{โมฬิยผคฺคุนริฏฺฐํ จ นาโม,}{อนฺธวเน กถิปุณฺณํ นิวาโป.} \\ \csromangathabat{ราสิกเณรุมหาคชนาโม,}{สารูปโม\csromansharedfootnote{3e46bebd4c9f46bd}{สารวโร (สฺยา), สารวโน (ก)} ปุน ปิงฺคลโกจฺโฉ.}}}

\csromanmatikaanchor{mtk.62}
\csromantocmarknopage{chapter}{4. มหายมกวคฺค}{mtk.62}
\bookmark[dest={mtk.62},level=0]{4. มหายมกวคฺค}
\chapterheadpageread{4. มหา\-ยมก\-วคฺค}
\csromanmatikaanchor{mtk.63}
\csromantocmarkref{subsection}{1. จูฬโคสิงฺคสุตฺต}{mtk.63}
\bookmark[dest={mtk.63},level=2]{1. จูฬโคสิงฺคสุตฺต}
\csromanheader{2}{1. จูฬโคสิงฺค\-สุตฺต}
\proseitem{325}{\csromanfolio{266}เอวํ เม สุตํ\csromandash{}เอกํ สมยํ ภควา นาติเก\csromansharedfootnote{0909fef6a8b1a95d}{นาทิเก (สี, สฺยา, อิ), ญาติเก (ก)} วิหรติ คิญฺชกาวสเถ. เตน โข ปน สมเยน อายสฺมา จ อนุรุทฺโธ อายสฺมา จ นนฺทิโย อายสฺมา จ กิมิโล\csromansharedfootnote{84cab62150801a02}{กิมฺพิโล (สี, อิ, ก)} โคสิงฺค\-สาลวน\-ทาเย วิหรนฺติ. อถ โข ภควา สายนฺหสมยํ ปฏิสลฺลานา วุฏฺฐิโต เยน โคสิงฺค\-สาลวน\-ทาโย เตนุปสงฺกมิ. อทฺทสา โข ทายปาโล ภควนฺตํ ทูรโตว อาคจฺฉนฺตํ, ทิสฺวาน ภควนฺตํ เอตทโวจ “มา สมณ เอตํ ทายํ ปาวิสิ, สนฺเตตฺถ ตโย กุลปุตฺตา อตฺตกามรูปา วิหรนฺติ. มา เตสํ อผาสุมกาสี”ติ.}

\prose{อสฺโสสิ โข อายสฺมา อนุรุทฺโธ ทายปาลสฺส ภควตา สทฺธิํ มนฺตยมานสฺส, สุตฺวาน ทายปาลํ เอตทโวจ “มา อาวุโส ทายปาล ภควนฺตํ วาเรสิ, สตฺถา โน ภควา อนุปฺปตฺโต”ติ. อถ โข อายสฺมา อนุรุทฺโธ เยนายสฺมา จ นนฺทิโย อายสฺมา จ กิมิโล เตนุปสงฺกมิ, อุปสงฺกมิตฺวา อายสฺมนฺตํ จ นนฺทิยํ อายสฺมนฺตํ จ กิมิลํ เอตทโวจ “อภิกฺกมถา\-ยสฺมนฺโต อภิกฺกมถา\-ยสฺมนฺโต, สตฺถา โน ภควา อนุปฺปตฺโต”ติ. อถ โข อายสฺมา จ อนุรุทฺโธ อายสฺมา จ นนฺทิโย อายสฺมา จ กิมิโล ภควนฺตํ ปจฺจุคฺคนฺตฺวา เอโก ภควโต ปตฺตจีวรํ ปฏิคฺคเหสิ, เอโก อาสนํ ปญฺญเปสิ, เอโก ปาโททกํ อุปฏฺฐาเปสิ. นิสีทิ ภควา ปญฺญตฺเต อาสเน, นิสชฺช โข ภควา ปาเท ปกฺขาเลสิ. เตปิ โข อายสฺมนฺโต ภควนฺตํ อภิวาเทตฺวา เอกมนฺตํ นิสีทิํสุ. เอกมนฺตํ นิสินฺนํ โข อายสฺมนฺตํ อนุรุทฺธํ ภควา เอตทโวจ\csromandash{}}

\proseitem{326}{กจฺจิ โว อนุรุทฺธา ขมนียํ กจฺจิ ยาปนียํ, กจฺจิ ปิณฺฑเกน น กิลมถาติ. ขมนียํ ภควา ยาปนียํ ภควา, น จ มยํ ภนฺเต ปิณฺฑเกน กิลมามาติ. กจฺจิ ปน โว อนุรุทฺธา สมคฺคา สมฺโมทมานา อวิวทมานา ขีโรทกีภูตา อญฺญมญฺญํ ปิยจกฺขูหิ สมฺปสฺสนฺตา วิหรถาติ. ตคฺฆ \csromanfolio{267}มยํ ภนฺเต สมคฺคา สมฺโมทมานา อวิวทมานา ขีโรทกีภูตา อญฺญมญฺญํ ปิยจกฺขูหิ สมฺปสฺสนฺตา วิหรามาติ. ยถา กถํ ปน ตุเมฺห อนุรุทฺธา สมคฺคา สมฺโมทมานา อวิวทมานา ขีโรทกีภูตา อญฺญมญฺญํ ปิยจกฺขูหิ สมฺปสฺสนฺตา วิหรถาติ. อิธ มยฺหํ ภนฺเต เอวํ โหติ “ลาภา วต เม, สุลทฺธํ วต เม, โยหํ เอวรูเปหิ สพฺรหฺมจารีหิ สทฺธิํ วิหรามี”ติ. ตสฺส มยฺหํ ภนฺเต อิเมสุ อายสฺมนฺเตสุ เมตฺตํ กายกมฺมํ ปจฺจุปฏฺฐิตํ อาวิ เจว รโห จ, เมตฺตํ วจีกมฺมํ ปจฺจุปฏฺฐิตํ อาวิ เจว รโห จ, เมตฺตํ มโนกมฺมํ ปจฺจุปฏฺฐิตํ อาวิ เจว รโห จ. ตสฺส มยฺหํ ภนฺเต เอวํ โหติ “ยํนูนาหํ สกํ จิตฺตํ นิกฺขิปิตฺวา อิเมสํเยว อายสฺมนฺตานํ จิตฺตสฺส วเสน วตฺเตยฺยนฺ”ติ. โส โข อหํ ภนฺเต สกํ จิตฺตํ นิกฺขิปิตฺวา อิเมสํเยว อายสฺมนฺตานํ จิตฺตสฺส วเสน วตฺตามิ. นานา หิ โข โน ภนฺเต กายา, เอกํ จ ปน มญฺเญ จิตฺตนฺติ.}

\prose{อายสฺมาปิ โข นนฺทิโย ฯเปฯ. อายสฺมาปิ โข กิมิโล ภควนฺตํ เอตทโวจ “มยฺหมฺปิ ภนฺเต เอวํ โหติ ‘ลาภา วต เม, สุลทฺธํ วต เม, โยหํ เอวรูเปหิ สพฺรหฺมจารีหิ สทฺธิํ วิหรามี’ติ. ตสฺส มยฺหํ ภนฺเต อิเมสุ อายสฺมนฺเตสุ เมตฺตํ กายกมฺมํ ปจฺจุปฏฺฐิตํ อาวิ เจว รโห จ, เมตฺตํ วจีกมฺมํ ปจฺจุปฏฺฐิตํ อาวิ เจว รโห จ, เมตฺตํ มโนกมฺมํ ปจฺจุปฏฺฐิตํ อาวิ เจว รโห จ. ตสฺส มยฺหํ ภนฺเต เอวํ โหติ ‘ยํนูนาหํ สกํ จิตฺตํ นิกฺขิปิตฺวา อิเมสํเยว อายสฺมนฺตานํ จิตฺตสฺส วเสน วตฺเตยฺยนฺ’ติ. โส โข อหํ ภนฺเต สกํ จิตฺตํ นิกฺขิปิตฺวา อิเมสํเยว อายสฺมนฺตานํ จิตฺตสฺส วเสน วตฺตามิ. นานา หิ โข โน ภนฺเต กายา, เอกํ จ ปน มญฺเญ จิตฺตนฺติ.}

\prose{เอวํ โข มยํ ภนฺเต สมคฺคา สมฺโมทมานา อวิวทมานา ขีโรทกีภูตา อญฺญมญฺญํ ปิยจกฺขูหิ สมฺปสฺสนฺตา วิหรามา”ติ.}

\proseitem{327}{สาธุ สาธุ อนุรุทฺธา, กจฺจิ ปน โว อนุรุทฺธา อปฺปมตฺตา อาตาปิโน ปหิตตฺตา วิหรถาติ. ตคฺฆ มยํ ภนฺเต อปฺปมตฺตา อาตาปิโน ปหิตตฺตา วิหรามาติ. ยถา กถํ ปน ตุเมฺห อนุรุทฺธา อปฺปมตฺตา อาตาปิโน ปหิตตฺตา วิหรถาติ. อิธ ภนฺเต อมฺหากํ โย ปฐมํ คามโต ปิณฺฑาย ปฏิกฺกมติ, โส อาสนานิ ปญฺญเปติ, ปานียํ ปริโภชนียํ อุปฏฺฐาเปติ, อวกฺการปาติํ อุปฏฺฐาเปติ. โย ปจฺฉา \csromanfolio{268}คามโต ปิณฺฑาย ปฏิกฺกมติ. สเจ โหติ ภุตฺตาวเสโส, สเจ อากงฺขติ ภุญฺชติ. โน เจ อากงฺขติ, อปฺปหริเต วา ฉฑฺเฑติ, อปฺปาณเก วา อุทเก โอปิลาเปติ, โส อาสนานิ ปฏิสาเมติ, ปานียํ ปริโภชนียํ ปฏิสาเมติ, อวกฺการปาติํ ปฏิสาเมติ, ภตฺตคฺคํ สมฺมชฺชติ. โย ปสฺสติ ปานียฆฏํ วา ปริโภชนียฆฏํ วา วจฺจฆฏํ วา ริตฺตํ ตุจฺฉํ, โส อุปฏฺฐาเปติ. สจสฺส โหติ อวิสยฺหํ, หตฺถวิกาเรน ทุติยํ อามนฺเตตฺวา หตฺถ\-วิลงฺฆเกน อุปฏฺฐาเปม, น เตฺวว มยํ ภนฺเต ตปฺปจฺจยา วาจํ ภินฺทาม. ปญฺจาหิกํ โข ปน มยํ ภนฺเต สพฺพรตฺติกํ ธมฺมิยา กถาย สนฺนิสีทาม. เอวํ โข มยํ ภนฺเต อปฺปมตฺตา อาตาปิโน ปหิตตฺตา วิหรามาติ.}

\proseitem{328}{สาธุ สาธุ อนุรุทฺธา, อตฺถิ ปน โว อนุรุทฺธา เอวํ อปฺปมตฺตานํ อาตาปีนํ ปหิตตฺตานํ วิหรนฺตานํ อุตฺตริ มนุสฺสธมฺมา อลม\-ริย\-ญาณ\-ทสฺสน\-วิเสโส อธิคโต ผาสุวิหาโรติ. กิํ หิ โน สิยา ภนฺเต. อิธ มยํ ภนฺเต ยาวเทว อากงฺขาม, วิวิจฺเจว กาเมหิ วิวิจฺจ อกุสเลหิ ธมฺเมหิ สวิตกฺกํ สวิจารํ วิเวกชํ ปีติสุขํ ปฐมํ ฌานํ อุปสมฺปชฺช วิหราม. อยํ โข โน ภนฺเต อมฺหากํ อปฺปมตฺตานํ อาตาปีนํ ปหิตตฺตานํ วิหรนฺตานํ อุตฺตริ มนุสฺสธมฺมา อลม\-ริย\-ญาณ\-ทสฺสน\-วิเสโส อธิคโต ผาสุวิหาโรติ.}

\prose{สาธุ สาธุ อนุรุทฺธา, เอตสฺส ปน โว อนุรุทฺธา วิหารสฺส สมติกฺกมาย เอตสฺส วิหารสฺส ปฏิปฺปสฺสทฺธิยา อตฺถญฺโญ อุตฺตริ มนุสฺสธมฺมา อลม\-ริย\-ญาณ\-ทสฺสน\-วิเสโส อธิคโต ผาสุวิหาโรติ. กิํ หิ โน สิยา ภนฺเต. อิธ มยํ ภนฺเต ยาวเทว อากงฺขาม, วิตกฺก\-วิจารานํ วูปสมา อชฺฌตฺตํ สมฺปสาทนํ เจตโส เอโกทิภาวํ อวิตกฺกํ อวิจารํ สมาธิชํ ปีติสุขํ ทุติยํ ฌานํ อุปสมฺปชฺช วิหราม. เอตสฺส ภนฺเต วิหารสฺส สมติกฺกมาย เอตสฺส วิหารสฺส ปฏิปฺปสฺสทฺธิยา อยมญฺโญ อุตฺตริ มนุสฺสธมฺมา อลม\-ริย\-ญาณ\-ทสฺสน\-วิเสโส อธิคโต ผาสุวิหาโรติ.}

\prose{สาธุ สาธุ อนุรุทฺธา, เอตสฺส ปน โว อนุรุทฺธา วิหารสฺส สมติกฺกมาย เอตสฺส วิหารสฺส ปฏิปฺปสฺสทฺธิยา อตฺถญฺโญ อุตฺตริ มนุสฺสธมฺมา อลม\-ริย\-ญาณ\-ทสฺสน\-วิเสโส อธิคโต ผาสุวิหาโรติ. \csromanfolio{269}กิํ หิ โน สิยา ภนฺเต. อิธ มยํ ภนฺเต ยาวเทว อากงฺขาม, ปีติยา จ วิราคา อุเปกฺขกา จ วิหราม สตา จ สมฺปชานา, สุขญฺจ กาเยน ปฏิสํเวเทม, ยํ ตํ อริยา อาจิกฺขนฺติ “อุเปกฺขโก สติมา สุขวิหารี”ติ ตติยํ ฌานํ อุปสมฺปชฺช วิหราม. เอตสฺส ภนฺเต วิหารสฺส สมติกฺกมาย เอตสฺส วิหารสฺส ปฏิปฺปสฺสทฺธิยา อยมญฺโญ อุตฺตริ มนุสฺสธมฺมา อลม\-ริย\-ญาณ\-ทสฺสน\-วิเสโส อธิคโต ผาสุวิหาโรติ.}

\prose{สาธุ สาธุ อนุรุทฺธา, เอตสฺส ปน โว อนุรุทฺธา วิหารสฺส สมติกฺกมาย เอตสฺส วิหารสฺส ปฏิปฺปสฺสทฺธิยา อตฺถญฺโญ อุตฺตริ มนุสฺสธมฺมา อลม\-ริย\-ญาณ\-ทสฺสน\-วิเสโส อธิคโต ผาสุวิหาโรติ. กิํ หิ โน สิยา ภนฺเต. อิธ มยํ ภนฺเต ยาวเทว อากงฺขาม, สุขสฺส จ ปหานา ทุกฺขสฺส จ ปหานา ปุพฺเพว โสมนสฺส\-โทมนสฺสานํ อตฺถงฺคมา อทุกฺขมสุขํ อุเปกฺขา\-สติ\-ปาริสุทฺธิํ จตุตฺถํ ฌานํ อุปสมฺปชฺช วิหราม. เอตสฺส ภนฺเต วิหารสฺส สมติกฺกมาย เอตสฺส วิหารสฺส ปฏิปฺปสฺสทฺธิยา อยมญฺโญ อุตฺตริ มนุสฺสธมฺมา อลม\-ริย\-ญาณ\-ทสฺสน\-วิเสโส อธิคโต ผาสุวิหาโรติ.}

\prose{สาธุ สาธุ อนุรุทฺธา, เอตสฺส ปน โว อนุรุทฺธา วิหารสฺส สมติกฺกมาย เอตสฺส วิหารสฺส ปฏิปฺปสฺสทฺธิยา อตฺถญฺโญ อุตฺตริ มนุสฺสธมฺมา อลม\-ริย\-ญาณ\-ทสฺสน\-วิเสโส อธิคโต ผาสุวิหาโรติ. กิํ หิ โน สิยา ภนฺเต. อิธ มยํ ภนฺเต ยาวเทว อากงฺขาม, สพฺพโส รูปสญฺญานํ สมติกฺกมา ปฏิฆ\-สญฺญานํ อตฺถงฺคมา นานตฺต\-สญฺญานํ อมนสิการา “อนนฺโต อากาโส”ติ อากาสา\-นญฺจา\-ยตนํ อุปสมฺปชฺช วิหราม. เอตสฺส ภนฺเต วิหารสฺส สมติกฺกมาย เอตสฺส วิหารสฺส ปฏิปฺปสฺสทฺธิยา อยมญฺโญ อุตฺตริ มนุสฺสธมฺมา อลม\-ริย\-ญาณ\-ทสฺสน\-วิเสโส อธิคโต ผาสุวิหาโรติ.}

\prose{สาธุ สาธุ อนุรุทฺธา, เอตสฺส ปน โว อนุรุทฺธา วิหารสฺส สมติกฺกมาย เอตสฺส วิหารสฺส ปฏิปฺปสฺสทฺธิยา อตฺถญฺโญ อุตฺตริ มนุสฺสธมฺมา อลม\-ริย\-ญาณ\-ทสฺสน\-วิเสโส อธิคโต ผาสุวิหาโรติ. กิํ หิ โน สิยา ภนฺเต. อิธ มยํ ภนฺเต ยาวเทว อากงฺขาม, สพฺพโส อากาสา\-นญฺจา\-ยตนํ สมติกฺกมฺม “อนนฺตํ วิญฺญาณนฺ”ติ วิญฺญาณญฺจา\-ยตนํ อุปสมฺปชฺช วิหราม ฯเปฯ สพฺพโส วิญฺญาณญฺจา\-ยตนํ สมติกฺกมฺม “นตฺถิ กิญฺจี”ติ อากิญฺจญฺญา\-ยตนํ อุปสมฺปชฺช วิหราม ฯเปฯ สพฺพโส อากิญฺจญฺญา\-ยตนํ \csromanfolio{270}สมติกฺกมฺม เนว\-สญฺญา\-นา\-สญฺญา\-ยตนํ อุปสมฺปชฺช วิหราม. เอตสฺส ภนฺเต วิหารสฺส สมติกฺกมาย เอตสฺส วิหารสฺส ปฏิปฺปสฺสทฺธิยา อยมญฺโญ อุตฺตริ มนุสฺสธมฺมา อลม\-ริย\-ญาณ\-ทสฺสน\-วิเสโส อธิคโต ผาสุวิหาโรติ.}

\proseitem{329}{สาธุ สาธุ อนุรุทฺธา, เอตสฺส ปน โว อนุรุทฺธา วิหารสฺส สมติกฺกมาย เอตสฺส วิหารสฺส ปฏิปฺปสฺสทฺธิยา อตฺถญฺโญ อุตฺตริ มนุสฺสธมฺมา อลม\-ริย\-ญาณ\-ทสฺสน\-วิเสโส อธิคโต ผาสุวิหาโรติ. กิํ หิ โน สิยา ภนฺเต. อิธ มยํ ภนฺเต ยาวเทว อากงฺขาม, สพฺพโส เนว\-สญฺญา\-นา\-สญฺญา\-ยตนํ สมติกฺกมฺม สญฺญาเวทยิต\-นิโรธํ อุปสมฺปชฺช วิหราม, ปญฺญาย จ โน ทิสฺวา อาสวา ปริกฺขีณา. เอตสฺส ภนฺเต วิหารสฺส สมติกฺกมาย เอตสฺส วิหารสฺส ปฏิปฺปสฺสทฺธิยา อยมญฺโญ อุตฺตริ มนุสฺสธมฺมา อลม\-ริย\-ญาณ\-ทสฺสน\-วิเสโส อธิคโต ผาสุวิหาโร. อิมมฺหา จ มยํ ภนฺเต ผาสุวิหารา อญฺญํ ผาสุวิหารํ อุตฺตริตรํ วา ปณีตตรํ วา น สมนุปสฺสามาติ. สาธุ สาธุ อนุรุทฺธา, อิมมฺหา ผาสุวิหารา อุตฺตริตโร วา ปณีตตโร วา ผาสุวิหาโร นตฺถีติ.}

\proseitem{330}{อถ โข ภควา อายสฺมนฺตญฺจ อนุรุทฺธํ อายสฺมนฺตญฺจ นนฺทิยํ อายสฺมนฺตญฺจ กิมิลํ ธมฺมิยา กถาย สนฺทสฺเสตฺวา สมาทเปตฺวา สมุตฺเตเชตฺวา สมฺปหํเสตฺวา อุฏฺฐายาสนา ปกฺกามิ. อถ โข อายสฺมา จ อนุรุทฺโธ อายสฺมา จ นนฺทิโย อายสฺมา จ กิมิโล ภควนฺตํ อนุสํยายิตฺวา\csromansharedfootnote{3b0f94f867f00577}{อนุสํสาเวตฺวา (สี), อนุสาเวตฺวา (ฏีกา)} ตโต ปฏินิวตฺติตฺวา อายสฺมา จ นนฺทิโย อายสฺมา จ กิมิโล อายสฺมนฺตํ อนุรุทฺธํ เอตทโวจุํ “กิํ นุ โข มยํ อายสฺมโต อนุรุทฺธสฺส เอวมาโรจิมฺห ‘อิมาสญฺจ อิมาสญฺจ วิหาร\-สมาปตฺตีนํ มยํ ลาภิโน’ติ. ยํ โน อายสฺมา อนุรุทฺโธ ภควโต สมฺมุขา ยาว อาสวานํ ขยา ปกาเสตี”ติ. น โข เม อายสฺมนฺโต เอวมาโรเจสุํ “อิมาสญฺจ อิมาสญฺจ วิหาร\-สมาปตฺตีนํ มยํ ลาภิโน”ติ. อปิ จ เม อายสฺมนฺตานํ เจตสา เจโต ปริจฺจ วิทิโต “อิมาสญฺจ อิมาสญฺจ วิหาร\-สมาปตฺตีนํ อิเม อายสฺมนฺโต ลาภิโน”ติ. เทวตาปิ \csromanfolio{271}เม เอตมตฺถํ อาโรเจสุํ “อิมาสญฺจ อิมาสญฺจ วิหาร\-สมาปตฺตีนํ อิเม อายสฺมนฺโต ลาภิโน”ติ. ตเมนํ ภควตา ปญฺหา\-ภิปุฏฺเฐน พฺยากตนฺติ.}

\proseitem{331}{อถ โข ทีโฆ ปรชโน ยกฺโข เยน ภควา เตนุปสงฺกมิ, อุปสงฺกมิตฺวา ภควนฺตํ อภิวาเทตฺวา เอกมนฺตํ อฏฺฐาสิ, เอกมนฺตํ ฐิโต โข ทีโฆ ปรชโน ยกฺโข ภควนฺตํ เอตทโวจ “ลาภา วต ภนฺเต วชฺชีนํ, สุลทฺธลาภา วชฺชิปชาย, ยตฺถ ตถาคโต วิหรติ อรหํ สมฺมาสมฺพุทฺโธ, อิเม จ ตโย กุลปุตฺตา อายสฺมา จ อนุรุทฺโธ อายสฺมา จ นนฺทิโย อายสฺมา จ กิมิโล”ติ. ทีฆสฺส ปรชนสฺส ยกฺขสฺส สทฺทํ สุตฺวา ภุมฺมา เทวา สทฺทม\-นุสฺสาเวสุํ “ลาภา วต โภ วชฺชีนํ, สุลทฺธลาภา วชฺชิปชาย, ยตฺถ ตถาคโต วิหรติ อรหํ สมฺมาสมฺพุทฺโธ, อิเม จ ตโย กุลปุตฺตา อายสฺมา จ อนุรุทฺโธ อายสฺมา จ นนฺทิโย อายสฺมา จ กิมิโล”ติ. ภุมฺมานํ เทวานํ สทฺทํ สุตฺวา จาตุมหาราชิตา เทวา ฯเปฯ ตาวติํสา เทวา ฯเปฯ ยามา เทวา ฯเปฯ ตุสิตา เทวา ฯเปฯ นิมฺมานรตี เทวา ฯเปฯ ปรนิมฺมิต\-วส\-วตฺตี เทวา ฯเปฯ พฺรหฺมกายิกา เทวา สทฺทม\-นุสฺสาเวสุํ “ลาภา วต โภ วชฺชีนํ, สุลทฺธลาภา วชฺชิปชาย, ยตฺถ ตถาคโต วิหรติ อรหํ สมฺมาสมฺพุทฺโธ, อิเม จ ตโย กุลปุตฺตา อายสฺมา จ อนุรุทฺโธ อายสฺมา จ นนฺทิโย อายสฺมา จ กิมิโล”ติ. อิติห เต อายสฺมนฺโต เตน ขเณน (เตน ลเยน)\csromansharedfootnote{b6e4c5acbaaffdb0}{( ) สีสฺยาอิโปตฺถเกสุ นตฺถิ.} เตน มุหุตฺเตน ยาวพฺรหฺมโลกา วิทิตา\csromansharedfootnote{05f4d2da259a8b8b}{สํวิทิตา (ก)} อเหสุํ.}

\prose{เอวเมตํ ทีฆ เอวเมตํ ทีฆ, ยสฺมาปิ ทีฆ กุลา เอเต ตโย กุลปุตฺตา อคารสฺมา อนคาริยํ ปพฺพชิตา, ตญฺเจปิ กุลํ เอเต ตโย กุลปุตฺเต ปสนฺนจิตฺตํ อนุสฺสเรยฺย, ตสฺสปาสฺส กุลสฺส ทีฆรตฺตํ หิตาย สุขาย. ยสฺมาปิ ทีฆ กุลปริวฏฺฏา เอเต ตโย กุลปุตฺตา อคารสฺมา อนคาริยํ ปพฺพชิตา, โส เจปิ กุลปริวฏฺโฏ เอเต ตโย กุลปุตฺเต ปสนฺนจิตฺโต อนุสฺสเรยฺย, ตสฺสปาสฺส กุลปริวฏฺฏสฺส ทีฆรตฺตํ หิตาย สุขาย. ยสฺมาปิ ทีฆ คามา เอเต ตโย กุลปุตฺตา อคารสฺมา อนคาริยํ ปพฺพชิตา, โส เจปิ คาโม เอเต ตโย \csromanfolio{272}กุลปุตฺเต ปสนฺนจิตฺโต อนุสฺสเรยฺย, ตสฺสปาสฺส คามสฺส ทีฆรตฺตํ หิตาย สุขาย. ยสฺมาปิ ทีฆ นิคมา เอเต ตโย กุลปุตฺตา อคารสฺมา อนคาริยํ ปพฺพชิตา, โส เจปิ นิคโม เอเต ตโย กุลปุตฺเต ปสนฺนจิตฺโต อนุสฺสเรยฺย, ตสฺสปาสฺส นิคมสฺส ทีฆรตฺตํ หิตาย สุขาย. ยสฺมาปิ ทีฆ นครา เอเต ตโย กุลปุตฺตา อคารสฺมา อนคาริยํ ปพฺพชิตา, ตญฺเจปิ นครํ เอเต ตโย กุลปุตฺเต ปสนฺนจิตฺตํ อนุสฺสเรยฺย, ตสฺสปาสฺส นครสฺส ทีฆรตฺตํ หิตาย สุขาย. ยสฺมาปิ ทีฆ ชนปทา เอเต ตโย กุลปุตฺตา อคารสฺมา อนคาริยํ ปพฺพชิตา, โส เจปิ ชนปโท เอเต ตโย กุลปุตฺเต ปสนฺนจิตฺโต อนุสฺสเรยฺย, ตสฺสปาสฺส ชนปทสฺส ทีฆรตฺตํ หิตาย สุขาย. สพฺเพ เจปิ ทีฆ ขตฺติยา เอเต ตโย กุลปุตฺเต ปสนฺนจิตฺตา อนุสฺสเรยฺยุํ, สพฺเพสานํปาสฺส ขตฺติยานํ ทีฆรตฺตํ หิตาย สุขาย. สพฺเพ เจปิ ทีฆ พฺราหฺมณา ฯเปฯ. สพฺเพ เจปิ ทีฆ เวสฺสา ฯเปฯ. สพฺเพ เจปิ ทีฆ สุทฺทา เอเต ตโย กุลปุตฺเต ปสนฺนจิตฺตา อนุสฺสเรยฺยุํ, สพฺเพสานํปาสฺส สุทฺทานํ ทีฆรตฺตํ หิตาย สุขาย. สเทวโก เจปิ ทีฆ โลโก สมารโก สพฺรหฺมโก สสฺสมณ\-พฺราหฺมณี ปชา สเทวมนุสฺสา เอเต ตโย กุลปุตฺเต ปสนฺนจิตฺตา อนุสฺสเรยฺย, สเทวกสฺส\-ปาสฺส โลกสฺส สมารกสฺส สพฺรหฺมกสฺส สสฺสมณ\-พฺราหฺมณิยา ปชาย สเทว\-มนุสฺสาย ทีฆรตฺตํ หิตาย สุขาย. ปสฺส ทีฆ ยาว เอเต ตโย กุลปุตฺตา พหุชนหิตาย ปฏิปนฺนา พหุชน\-สุขาย โลกานุกมฺปาย อตฺถาย หิตาย สุขาย เทว\-มนุสฺสา\-นนฺติ.}

\prosewithcloserruled{อิทมโวจ ภควา. อตฺตมโน ทีโฆ ปรชโน ยกฺโข ภควโต ภาสิตํ อภินนฺทีติ.}{จูฬโคสิงฺค\-สุตฺตํ นิฏฺฐิตํ ปฐมํ.}
\csromanmatikaanchor{mtk.64}
\csromantocmarkref{subsection}{2. มหาโคสิงฺคสุตฺต}{mtk.64}
\bookmark[dest={mtk.64},level=2]{2. มหาโคสิงฺคสุตฺต}
\csromanheader{2}{2. มหา\-โคสิงฺค\-สุตฺต}
\proseitem{332}{เอวํ เม สุตํ\csromandash{}เอกํ สมยํ ภควา โคสิงฺค\-สาลวน\-ทาเย วิหรติ สมฺพหุเลหิ อภิญฺญาเตหิ อภิญฺญาเตหิ เถเรหิ สาวเกหิ สทฺธิํ อายสฺมตา จ สาริปุตฺเตน อายสฺมตา จ มหา\-โมคฺคลฺลาเนน \csromanfolio{273}อายสฺมตา จ มหากสฺสเปน อายสฺมตา จ อนุรุทฺเธน อายสฺมตา จ เรวเตน อายสฺมตา จ อานนฺเทน อญฺเญหิ จ อภิญฺญาเตหิ อภิญฺญาเตหิ เถเรหิ สาวเกหิ สทฺธิํ. อถ โข อายสฺมา มหาโมคฺคลฺลาโน สายนฺหสมยํ ปฏิสลฺลานา วุฏฺฐิโต เยนายสฺมา มหากสฺสโป เตนุปสงฺกมิ, อุปสงฺกมิตฺวา อายสฺมนฺตํ มหากสฺสปํ เอตทโวจ “อายามาวุโส กสฺสป เยนายสฺมา สาริปุตฺโต เตนุปสงฺกมิสฺสาม ธมฺม\-สฺสวนายา”ติ. “เอวมาวุโส”ติ โข อายสฺมา มหากสฺสโป อายสฺมโต มหา\-โมคฺคลฺลานสฺส ปจฺจสฺโสสิ. อถ โข อายสฺมา จ มหาโมคฺคลฺลาโน อายสฺมา จ มหากสฺสโป อายสฺมา จ อนุรุทฺโธ เยนายสฺมา สาริปุตฺโต เตนุ\-ปสงฺกมิํสุ ธมฺม\-สฺสวนาย. อทฺทสา โข อายสฺมา อานนฺโท อายสฺมนฺตํ จ มหา\-โมคฺคลฺลานํ อายสฺมนฺตํ จ มหากสฺสปํ อายสฺมนฺตํ จ อนุรุทฺธํ เยนายสฺมา สาริปุตฺโต เตนุ\-ปสงฺกมนฺเต ธมฺม\-สฺสวนาย, ทิสฺวาน เยนายสฺมา เรวโต เตนุปสงฺกมิ, อุปสงฺกมิตฺวา อายสฺมนฺตํ เรวตํ เอตทโวจ “อุปสงฺกมนฺตา โข อมู อาวุโส\csromansharedfootnote{3933a2827da8deb8}{อายสฺมนฺตาวุโส (ก)} เรวต สปฺปุริสา เยนายสฺมา สาริปุตฺโต เตน ธมฺม\-สฺสวนาย. อายามาวุโส เรวต เยนายสฺมา สาริปุตฺโต เตนุปสงฺกมิสฺสาม ธมฺม\-สฺสวนายา”ติ. “เอวมาวุโส”ติ โข อายสฺมา เรวโต อายสฺมโต อานนฺทสฺส ปจฺจสฺโสสิ. อถ โข อายสฺมา จ เรวโต อายสฺมา จ อานนฺโท เยนายสฺมา สาริปุตฺโต เตนุ\-ปสงฺกมิํสุ ธมฺม\-สฺสวนาย.}

\proseitem{333}{อทฺทสา โข อายสฺมา สาริปุตฺโต อายสฺมนฺตํ จ เรวตํ อายสฺมนฺตํ จ อานนฺทํ ทูรโตว อาคจฺฉนฺเต, ทิสฺวาน อายสฺมนฺตํ อานนฺทํ เอตทโวจ “เอตุ โข อายสฺมา อานนฺโท, สฺวาคตํ อายสฺมโต อานนฺทสฺส ภควโต อุปฏฺฐากสฺส ภควโต สนฺติกา\-วจรสฺส. รมณียํ อาวุโส อานนฺท โคสิงฺค\-สาลวนํ, โทสินา รตฺติ, สพฺพ\-ผาลิผุลฺลา\csromansharedfootnote{fa0c589d6ea6f0c0}{สพฺพปาลิผุลฺลา (สี)} สาลา, ทิพฺพา มญฺเญ คนฺธา สมฺปวนฺติ. กถํรูเปน อาวุโส อานนฺท ภิกฺขุนา โคสิงฺค\-สาลวนํ โสเภยฺยา”ติ. อิธาวุโส สาริปุตฺต ภิกฺขุ พหุสฺสุโต โหติ สุตธโร สุตสนฺนิจโย. เย เต ธมฺมา อาทิกลฺยาณา มชฺเฌกลฺยาณา \csromanfolio{274}ปริโยสาน\-กลฺยาณา สาตฺถา สพฺยญฺชนา เกวล\-ปริปุณฺณํ ปริสุทฺธํ พฺรหฺมจริยํ อภิวทนฺติ, ตถารูปาสฺส ธมฺมา พหุสฺสุตา โหนฺติ ธาตา\csromansharedfootnote{acf84a73c1bbfd91}{ธตา (สี, สฺยา, กํ, อิ)} วจสา ปริจิตา มนสา\-นุเปกฺขิตา ทิฏฺฐิยา สุปฺปฏิวิทฺธา. โส จตสฺสนฺนํ ปริสานํ ธมฺมํ เทเสติ ปริมณฺฑเลหิ ปทพฺยญฺชเนหิ อนุปฺปพนฺเธหิ\csromansharedfootnote{a313d5a418d3d312}{อปฺปพทฺเธหิ (สี, อิ)} อนุสย\-สมุคฺฆาตาย. เอวรูเปน โข อาวุโส สาริปุตฺต ภิกฺขุนา โคสิงฺค\-สาลวนํ โสเภยฺยาติ.}

\proseitem{334}{เอวํ วุตฺเต อายสฺมา สาริปุตฺโต อายสฺมนฺตํ เรวตํ เอตทโวจ “พฺยากตํ โข อาวุโส เรวต อายสฺมตา อานนฺเทน ยถาสกํ ปฏิภานํ. ตตฺถ ทานิ มยํ อายสฺมนฺตํ เรวตํ ปุจฺฉาม, รมณียํ อาวุโส เรวต โคสิงฺค\-สาลวนํ, โทสินา รตฺติ, สพฺพ\-ผาลิผุลฺลา สาลา, ทิพฺพา มญฺเญ คนฺธา สมฺปวนฺติ. กถํรูเปน อาวุโส เรวต ภิกฺขุนา โคสิงฺค\-สาลวนํ โสเภยฺยา”ติ. อิธาวุโส สาริปุตฺต ภิกฺขุ ปฏิสลฺลานา\-ราโม โหติ ปฏิสลฺลาน\-รโต, อชฺฌตฺตํ เจโตสมถม\-นุยุตฺโต, อนิรากต\-ชฺฌาโน, วิปสฺสนาย สมนฺนาคโต, พฺรูเหตา สุญฺญาคารานํ. เอวรูเปน โข อาวุโส สาริปุตฺต ภิกฺขุนา โคสิงฺค\-สาลวนํ โสเภยฺยาติ.}

\proseitem{335}{เอวํ วุตฺเต อายสฺมา สาริปุตฺโต อายสฺมนฺตํ อนุรุทฺธํ เอตทโวจ “พฺยากตํ โข อาวุโส อนุรุทฺธ อายสฺมตา เรวเตน ยถาสกํ ปฏิภานํ. ตตฺถ ทานิ มยํ อายสฺมนฺตํ อนุรุทฺธํ ปุจฺฉาม, รมณียํ อาวุโส อนุรุทฺธ โคสิงฺค\-สาลวนํ, โทสินา รตฺติ, สพฺพ\-ผาลิผุลฺลา สาลา, ทิพฺพา มญฺเญ คนฺธา สมฺปวนฺติ. กถํรูเปน อาวุโส อนุรุทฺธ ภิกฺขุนา โคสิงฺค\-สาลวนํ โสเภยฺยา”ติ. อิธาวุโส สาริปุตฺต ภิกฺขุ ทิพฺเพน จกฺขุนา วิสุทฺเธน อติกฺกนฺต\-มานุสเกน สหสฺสํ โลกานํ โวโลเกติ. เสยฺยถาปิ อาวุโส สาริปุตฺต จกฺขุมา ปุริโส อุปริปาสาท\-วรคโต สหสฺสํ เนมิมณฺฑลานํ โวโลเกยฺย. เอวเมว โข อาวุโส สาริปุตฺต ภิกฺขุ ทิพฺเพน จกฺขุนา วิสุทฺเธน อติกฺกนฺต\-มานุสเกน สหสฺสํ โลกานํ โวโลเกติ. เอวรูเปน โข อาวุโส สาริปุตฺต ภิกฺขุนา โคสิงฺค\-สาลวนํ โสเภยฺยาติ.}

\proseitem{336}{\csromanfolio{275}เอวํ วุตฺเต อายสฺมา สาริปุตฺโต อายสฺมนฺตํ มหากสฺสปํ เอตทโวจ “พฺยากตํ โข อาวุโส กสฺสป อายสฺมตา อนุรุทฺเธน ยถาสกํ ปฏิภานํ. ตตฺถ ทานิ มยํ อายสฺมนฺตํ มหากสฺสปํ ปุจฺฉาม, รมณียํ อาวุโส กสฺสป โคสิงฺค\-สาลวนํ, โทสินา รตฺติ, สพฺพ\-ผาลิผุลฺลา สาลา, ทิพฺพา มญฺเญ คนฺธา สมฺปวนฺติ. กถํรูเปน อาวุโส กสฺสป ภิกฺขุนา โคสิงฺค\-สาลวนํ โสเภยฺยา”ติ. อิธาวุโส สาริปุตฺต ภิกฺขุ อตฺตนา จ อารญฺญิโก โหติ อารญฺญิกตฺตสฺส จ วณฺณวาที, อตฺตนา จ ปิณฺฑปาติโก โหติ ปิณฺฑ\-ปาติ\-ก\-ตฺตสฺส จ วณฺณวาที, อตฺตนา จ ปํสุกูลิโก โหติ ปํสุกูลิ\-ก\-ตฺตสฺส จ วณฺณวาที, อตฺตนา จ เตจีวริโก โหติ เตจีวริกตฺตสฺส จ วณฺณวาที, อตฺตนา จ อปฺปิจฺโฉ โหติ อปฺปิจฺฉตาย จ วณฺณวาที, อตฺตนา จ สนฺตุฏฺโฐ โหติ สนฺตุฏฺฐิยา จ วณฺณวาที, อตฺตนา จ ปวิวิตฺโต โหติ ปวิเวกสฺส จ วณฺณวาที, อตฺตนา จ อสํสฏฺโฐ โหติ อสํสคฺคสฺส จ วณฺณวาที, อตฺตนา จ อารทฺธวีริโย โหติ วีริยา\-รมฺภสฺส จ วณฺณวาที, อตฺตนา จ สีลสมฺปนฺโน โหติ สีลสมฺปทาย จ วณฺณวาที, อตฺตนา จ สมาธิ\-สมฺปนฺโน โหติ สมาธิ\-สมฺปทาย จ วณฺณวาที, อตฺตนา จ ปญฺญาสมฺปนฺโน โหติ ปญฺญาสมฺปทาย จ วณฺณวาที, อตฺตนา จ วิมุตฺติ\-สมฺปนฺโน โหติ วิมุตฺติ\-สมฺปทาย จ วณฺณวาที, อตฺตนา จ วิมุตฺติ\-ญาณ\-ทสฺสน\-สมฺปนฺโน โหติ วิมุตฺติ\-ญาณ\-ทสฺสน\-สมฺปทาย จ วณฺณวาที. เอวรูเปน โข อาวุโส สาริปุตฺต ภิกฺขุนา โคสิงฺค\-สาลวนํ โสเภยฺยาติ.}

\proseitem{337}{เอวํ วุตฺเต อายสฺมา สาริปุตฺโต อายสฺมนฺตํ มหา\-โมคฺคลฺลานํ เอตทโวจ “พฺยากตํ โข อาวุโส โมคฺคลฺลาน อายสฺมตา มหากสฺสเปน ยถาสกํ ปฏิภานํ. ตตฺถ ทานิ มยํ อายสฺมนฺตํ มหา\-โมคฺคลฺลานํ ปุจฺฉาม, รมณียํ อาวุโส โมคฺคลฺลาน โคสิงฺค\-สาลวนํ, โทสินา รตฺติ, สพฺพ\-ผาลิผุลฺลา สาลา, ทิพฺพา มญฺเญ คนฺธา สมฺปวนฺติ. กถํรูเปน อาวุโส โมคฺคลฺลาน ภิกฺขุนา โคสิงฺค\-สาลวนํ โสเภยฺยา”ติ. อิธาวุโส สาริปุตฺต เทฺว ภิกฺขู อภิธมฺมกถํ กเถนฺติ. เต อญฺญมญฺญํ ปญฺหํ ปุจฺฉนฺติ, อญฺญมญฺญสฺส ปญฺหํ ปุฏฺฐา วิสฺสชฺเชนฺติ, โน จ สํสาเทนฺติ\csromansharedfootnote{9f1443e099dce6aa}{สํสาเรนฺติ (ก)}, ธมฺมี จ เนสํกถา ปวตฺตินี โหติ. เอวรูเปน โข อาวุโส สาริปุตฺต ภิกฺขุนา โคสิงฺค\-สาลวนํ โสเภยฺยาติ.}

\proseitem{338}{\csromanfolio{276}อถ โข อายสฺมา มหาโมคฺคลฺลาโน อายสฺมนฺตํ สาริปุตฺตํ เอตทโวจ “พฺยากตํ โข อาวุโส สาริปุตฺต อเมฺหหิ สพฺเพเหว ยถาสกํ ปฏิภานํ. ตตฺถ ทานิ มยํ อายสฺมนฺตํ สาริปุตฺตํ ปุจฺฉาม, รมณียํ อาวุโส สาริปุตฺต โคสิงฺค\-สาลวนํ, โทสินา รตฺติ, สพฺพ\-ผาลิผุลฺลา สาลา, ทิพฺพา มญฺเญ คนฺธา สมฺปวนฺติ. กถํรูเปน อาวุโส สาริปุตฺต ภิกฺขุนา โคสิงฺค\-สาลวนํ โสเภยฺยา”ติ. อิธาวุโส โมคฺคลฺลาน ภิกฺขุ จิตฺตํ วสํ วตฺเตติ, โน จ ภิกฺขุ จิตฺตสฺส วเสน วตฺตติ. โส ยาย วิหาร\-สมาปตฺติยา อากงฺขติ ปุพฺพณฺหสมยํ วิหริตุํ, ตาย วิหาร\-สมาปตฺติยา ปุพฺพณฺหสมยํ วิหรติ. ยาย วิหาร\-สมาปตฺติยา อากงฺขติ มชฺฌนฺหิกสมยํ\csromansharedfootnote{e5aba9944229c365}{มชฺฌนฺติกสมยํ (สี, สฺยา, กํ, อิ, ก)} วิหริตุํ, ตาย วิหาร\-สมาปตฺติยา มชฺฌนฺหิกสมยํ วิหรติ. ยาย วิหาร\-สมาปตฺติยา อากงฺขติ สายนฺหสมยํ วิหริตุํ, ตาย วิหาร\-สมาปตฺติยา สายนฺหสมยํ วิหรติ. เสยฺยถาปิ อาวุโส โมคฺคลฺลาน รญฺโญ วา ราช\-มหามตฺตสฺส วา นานารตฺตานํ ทุสฺสานํ ทุสฺสกรณฺฑโก ปูโร อสฺส. โส ยญฺญเทว ทุสฺสยุคํ อากงฺเขยฺย ปุพฺพณฺหสมยํ ปารุปิตุํ, ตํ ตเทว ทุสฺสยุคํ ปุพฺพณฺหสมยํ ปารุเปยฺย. ยญฺญเทว ทุสฺสยุคํ อากงฺเขยฺย มชฺฌนฺหิกสมยํ ปารุปิตุํ, ตํ ตเทว ทุสฺสยุคํ มชฺฌนฺหิกสมยํ ปารุเปยฺย. ยญฺญเทว ทุสฺสยุคํ อากงฺเขยฺย สายนฺหสมยํ ปารุปิตุํ, ตํ ตเทว ทุสฺสยุคํ สายนฺหสมยํ ปารุเปยฺย. เอวเมว โข อาวุโส โมคฺคลฺลาโน ภิกฺขุ จิตฺตํ วสํ วตฺเตติ, น จ ภิกฺขุ จิตฺตสฺส วเสน วตฺตติ. โส ยาย วิหาร\-สมาปตฺติยา อากงฺขติ ปุพฺพณฺหสมยํ วิหริตุํ, ตาย วิหาร\-สมาปตฺติยา ปุพฺพณฺหสมยํ วิหรติ. ยาย วิหาร\-สมาปตฺติยา อากงฺขติ มชฺฌนฺหิกสมยํ วิหริตุํ, ตาย วิหาร\-สมาปตฺติยา มชฺฌนฺหิกสมยํ วิหรติ. ยาย วิหาร\-สมาปตฺติยา อากงฺขติ สายนฺหสมยํ วิหริตุํ, ตาย วิหาร\-สมาปตฺติยา สายนฺหสมยํ วิหรติ. เอวรูเปน โข อาวุโส โมคฺคลฺลาน ภิกฺขุนา โคสิงฺค\-สาลวนํ โสเภยฺยาติ.}

\proseitem{339}{อถ โข อายสฺมา สาริปุตฺโต เต อายสฺมนฺเต เอตทโวจ “พฺยากตํ โข อาวุโส อเมฺหหิ สพฺเพเหว ยถาสกํ ปฏิภานํ. อายามาวุโส เยน ภควา เตนุปสงฺกมิสฺสาม, อุปสงฺกมิตฺวา \csromanfolio{277}เอตมตฺถํ ภควโต อาโรเจสฺสาม. ยถา โน ภควา พฺยากริสฺสติ ตถา นํ ธาเรสฺสามา”ติ. เอวมาวุโสติ โข เต อายสฺมนฺโต อายสฺมโต สาริปุตฺตสฺส ปจฺจสฺโสสุํ. อถ โข เต อายสฺมนฺโต เยน ภควา เตนุ\-ปสงฺกมิํสุ, อุปสงฺกมิตฺวา ภควนฺตํ อภิวาเทตฺวา เอกมนฺตํ นิสีทิํสุ, เอกมนฺตํ นิสินฺโน โข อายสฺมา สาริปุตฺโต ภควนฺตํ เอตทโวจ “อิธ ภนฺเต อายสฺมา จ เรวโต อายสฺมา จ อาอานนฺโท เยนาหํ เตนุ\-ปสงฺกมิํสุ ธมฺม\-สฺสวนาย. อทฺทสํ โข อหํ ภนฺเต อายสฺมนฺตํ จ เรวตํ อายสฺมนฺตํ จ อานนฺทํ ทูรโตว อาคจฺฉนฺเต, ทิสฺวาน อายสฺมนฺตํ อานนฺทํ เอตทโวจํ ‘เอตุ โข อายสฺมา อานนฺโท, สฺวาคตํ อายสฺมโต อานนฺทสฺส ภควโต อุปฏฺฐากสฺส ภควโต สนฺติกา\-วจรสฺส. รมณียํ อาวุโส อานนฺท โคสิงฺค\-สาลวนํ, โทสินา รตฺติ, สพฺพ\-ผาลิผุลฺลา สาลา, ทิพฺพา มญฺเญ คนฺธา สมฺปวนฺติ. กถํรูเปน อาวุโส อานนฺท ภิกฺขุนา โคสิงฺค\-สาลวนํ โสเภยฺยา’ติ. เอวํ วุตฺเต ภนฺเต อายสฺมา อานนฺโท มํ เอตทโวจ ‘อิธาวุโส สาริปุตฺต ภิกฺขุ พหุสฺสุโต โหติ สุตธโร ฯเปฯ อนุสย\-สมุคฺฆาตาย. เอวรูเปน โข อาวุโส สาริปุตฺต ภิกฺขุนา โคสิงฺค\-สาลวนํ โสเภยฺยา”ติ. สาธุ สาธุ สาริปุตฺต, ยถา ตํ อานนฺโทว สมฺมา พฺยากรมาโน พฺยากเรยฺย. อานนฺโท หิ สาริปุตฺต พหุสฺสุโต สุตธโร สุตสนฺนิจโย. เย เต ธมฺมา อาทิกลฺยาณา มชฺเฌกลฺยาณา ปริโยสาน\-กลฺยาณา สาตฺถา สพฺยญฺชนา เกวล\-ปริปุณฺณํ ปริสุทฺธํ พฺรหฺมจริยํ อภิวทนฺติ, ตถารูปาสฺส ธมฺมา พหุสฺสุตา โหนฺติ ธาตา วจสา ปริจิตา มนสา\-นุเปกฺขิตา ทิฏฺฐิยา สุปฺปฏิวิทฺธา. โส จตสฺสนฺนํ ปริสานํ ธมฺมํ เทเสติ ปริมณฺฑเลหิ ปทพฺยญฺชเนหิ อนุปฺปพนฺเธหิ อนุสย\-สมุคฺฆาตายาติ.}

\proseitem{340}{เอวํ วุตฺเต อหํ ภนฺเต อายสฺมนฺตํ เรวตํ เอตทโวจํ ‘พฺยากตํ โข อาวุโส เรวต อายสฺมตา อานนฺเทน ยถาสกํ ปฏิภานํ. ตตฺถ ทานิ มยํ อายสฺมนฺตํ เรวตํ ปุจฺฉาม ‘รมณียํ อาวุโส เรวต โคสิงฺค\-สาลวนํ, โทสินา รตฺติ, สพฺพ\-ผาลิผุลฺลา สาลา, ทิพฺพา มญฺเญ คนฺธา สมฺปวนฺติ. กถํรูเปน อาวุโส เรวต ภิกฺขุนา โคสิงฺค\-สาลวนํ โสเภยฺยา’ติ. เอวํ วุตฺเต ภนฺเต อายสฺมา เรวโต มํ เอตทโวจ ‘อิธาวุโส สาริปุตฺต ภิกฺขุ ปฏิสลฺลานา\-ราโม โหติ \csromanfolio{278}ปฏิสลฺลาน\-รโต, อชฺฌตฺตํ เจโตสมถม\-นุยุตฺโต, อนิรากต\-ชฺฌาโน, วิปสฺสนาย สมนฺนาคโต, พฺรูเหตา สุญฺญาคารานํ. เอวรูเปน โข อาวุโส สาริปุตฺต ภิกฺขุนา โคสิงฺค\-สาลวนํ โสเภยฺยา”ติ. สาธุ สาธุ สาริปุตฺต, ยถา ตํ เรวโตว สมฺมา พฺยากรมาโน พฺยากเรยฺย. เรวโต หิ สาริปุตฺต ปฏิสลฺลานา\-ราโม ปฏิสลฺลาน\-รโต, อชฺฌตฺตํ เจโตสมถม\-นุยุตฺโต, อนิรากต\-ชฺฌาโน, วิปสฺสนาย สมนฺนาคโต, พฺรูเหตา สุญฺญาคารานนฺติ.}

\proseitem{341}{เอวํ วุตฺเต อหํ ภนฺเต อายสฺมนฺตํ อนุรุทฺธํ เอตทโวจํ ‘พฺยากตํ โข อาวุโส อนุรุทฺธ อายสฺมตา เรวเตน ฯเปฯ. กถํรูเปน อาวุโส อนุรุทฺธ ภิกฺขุนา โคสิงฺค\-สาลวนํ โสเภยฺยา’ติ. เอวํ วุตฺเต ภนฺเต อายสฺมา อนุรุทฺโธ มํ เอตทโวจ ‘อิธาวุโส สาริปุตฺต ภิกฺขุ ทิพฺเพน จกฺขุนา วิสุทฺเธน อติกฺกนฺต\-มานุสเกน สหสฺสํ โลกานํ โวโลเกติ. เสยฺยถาปิ อาวุโส สาริปุตฺต จกฺขุมา ปุริโส ฯเปฯ. เอวรูเปน โข อาวุโส สาริปุตฺต ภิกฺขุนา โคสิงฺค\-สาลวนํ โสเภยฺยาติ. สาธุ สาธุ สาริปุตฺต, ยถา ตํ อนุรุทฺโธว สมฺมา พฺยากรมาโน พฺยากเรยฺย. อนุรุทฺโธ หิ สาริปุตฺต ทิพฺเพน จกฺขุนา วิสุทฺเธน อติกฺกนฺต\-มานุสเกน สหสฺสํ โลกานํ โวโลเกตีติ.}

\proseitem{342}{เอวํ วุตฺเต อหํ ภนฺเต อายสฺมนฺตํ มหากสฺสปํ เอตทโวจํ “พฺยากตํ โข อาวุโส กสฺสป อายสฺมตา อนุรุทฺเธน ยถาสกํ ปฏิภานํ. ตตฺถ ทานิ มยํ อายสฺมนฺตํ มหากสฺสปํ ปุจฺฉาม ฯเปฯ. กถํ รูเปน โข อาวุโส กสฺสป ภิกฺขุนา โคสิงฺค\-สาลวนํ โสเภยฺยา”ติ. เอวํ วุตฺเต ภนฺเต อายสฺมา มหากสฺสโป มํ เอตทโวจ “อิธาวุโส สาริปุตฺต ภิกฺขุ อตฺตนา จ อารญฺญิโก โหติ อารญฺญิกตฺตสฺส จ วณฺณวาที, อตฺตนา จ ปิณฺฑปาติโก โหติ ฯเปฯ อตฺตนา จ ปํสุกูลิโก โหติ ฯเปฯ อตฺตนา จ เตจีวริโก โหติ ฯเปฯ อตฺตนา จ อปฺปิจฺโฉ โหติ ฯเปฯ อตฺตนา จ สนฺตุฏฺโฐ โหติ ฯเปฯ อตฺตนา จ ปวิวิตฺโต โหติ ฯเปฯ อตฺตนา จ อสํสฏฺโฐ โหติ ฯเปฯ อตฺตนา จ อารทฺธวีริโย โหติ ฯเปฯ อตฺตนา จ สีลสมฺปนฺโน โหติ ฯเปฯ อตฺตนา จ สมาธิ\-สมฺปนฺโน โหติ ฯเปฯ อตฺตนา จ ปญฺญาสมฺปนฺโน โหติ ฯเปฯ อตฺตนา จ วิมุตฺติ\-สมฺปนฺโน \csromanfolio{279}โหติ ฯเปฯ อตฺตนา จ วิมุตฺติ\-ญาณ\-ทสฺสน\-สมฺปนฺโน โหติ วิมุตฺติ\-ญาณ\-ทสฺสน\-สมฺปทาย จ วณฺณวาที. เอวรูเปน โข อาวุโส สาริปุตฺต ภิกฺขุนา โคสิงฺค\-สาลวนํ โสเภยฺยา”ติ. สาธุ สาธุ สาริปุตฺต, ยถา ตํ กสฺสโปว สมฺมา พฺยากรมาโน พฺยากเรยฺย. กสฺสโป หิ สาริปุตฺต อตฺตนา จ อารญฺญิโก อารญฺญิกตฺตสฺส จ วณฺณวาที, อตฺตนา จ ปิณฺฑปาติโก ปิณฺฑ\-ปาติ\-ก\-ตฺตสฺส จ วณฺณวาที, อตฺตนา จ ปํสุกูลิโก ปํสุกูลิ\-ก\-ตฺตสฺส จ วณฺณวาที, อตฺตนา จ เตจีวริโก เตจีวริกตฺตสฺส จ วณฺณวาที, อตฺตนา จ อปฺปิจฺโฉ อปฺปิจฺฉตาย จ วณฺณวาที, อตฺตนา จ สนฺตุฏฺโฐ สนฺตุฏฺฐิยา จ วณฺณวาที, อตฺตนา จ ปวิวิตฺโต ปวิเวกสฺส จ วณฺณวาที, อตฺตนา จ อสํสฏฺโฐ อสํสคฺคสฺส จ วณฺณวาที, อตฺตนา จ อารทฺธวีริโย วีริยา\-รมฺภสฺส จ วณฺณวาที, อตฺตนา จ สีลสมฺปนฺโน สีลสมฺปทาย จ วณฺณวาที, อตฺตนา จ สมาธิ\-สมฺปนฺโน สมาธิ\-สมฺปทาย จ วณฺณวาที, อตฺตนา จ ปญฺญาสมฺปนฺโน ปญฺญาสมฺปทาย จ วณฺณวาที, อตฺตนา จ วิมุตฺติ\-สมฺปนฺโน วิมุตฺติ\-สมฺปทาย จ วณฺณวาที, อตฺตนา จ วิมุตฺติ\-ญาณ\-ทสฺสน\-สมฺปนฺโน วิมุตฺติ\-ญาณ\-ทสฺสน\-สมฺปทาย จ วณฺณวาทีติ.}

\proseitem{343}{เอวํ วุตฺเต อหํ ภนฺเต อายสฺมนฺตํ มหา\-โมคฺคลฺลานํ เอตทโวจํ “พฺยากตํ โข อาวุโส โมคฺคลฺลาน อายสฺมตา มหากสฺสเปน ยถาสกํ ปฏิภานํ. ตตฺถ ทานิ มยํ อายสฺมนฺตํ มหา\-โมคฺคลฺลานํ ปุจฺฉาม ฯเปฯ. กถํรูเปน อาวุโส โมคฺคลฺลาน ภิกฺขุนา โคสิงฺค\-สาลวนํ โสเภยฺยา”ติ. เอวํ วุตฺเต ภนฺเต อายสฺมา มหาโมคฺคลฺลาโน มํ เอตทโวจ “อิธาวุโส สาริปุตฺต เทฺว ภิกฺขู อภิธมฺมกถํ กเถนฺติ, เต อญฺญมญฺญํ ปญฺหํ ปุจฺฉนฺติ, อญฺญมญฺญสฺส ปญฺหํ ปุฏฺฐา วิสฺสชฺเชนฺติ, โน จ สํสาเทนฺติ, ธมฺมี จ เนสํ กถา ปวตฺตินี โหติ, เอวรูเปน โข อาวุโส สาริปุตฺต ภิกฺขุนา โคสิงฺค\-สาลวนํ โสเภยฺยา”ติ. สาธุ สาธุ สาริปุตฺต, ยถา ตํ โมคฺคลฺลาโนว สมฺมา พฺยากรมาโน พฺยากเรยฺย, โมคฺคลฺลาโน หิ สาริปุตฺต ธมฺม\-กถิโกติ.}

\proseitem{344}{เอวํ วุตฺเต อายสฺมา มหาโมคฺคลฺลาโน ภควนฺตํ เอตทโวจ “อถ ขฺวาหํ ภนฺเต อายสฺมนฺตํ สาริปุตฺตํ เอตทโวจํ ‘พฺยากตํ โข อาวุโส สาริปุตฺต อเมฺหหิ สพฺเพเหว ยถาสกํ ปฏิภานํ. ตตฺถ \csromanfolio{280}ทานิ มยํ อายสฺมนฺตํ สาริปุตฺตํ ปุจฺฉาม, รมณียํ อาวุโส สาริปุตฺต โคสิงฺค\-สาลวนํ, โทสินา รตฺติ, สพฺพ\-ผาลิผุลฺลา สาลา, ทิพฺพา มญฺเญ คนฺธา สมฺปวนฺติ. กถํรูเปน อาวุโส สาริปุตฺต ภิกฺขุนา โคสิงฺค\-สาลวนํ โสเภยฺยา”ติ. เอวํ วุตฺเต ภนฺเต อายสฺมา สริปุตฺโต มํ เอตทโวจ “อิธาวุโส โมคฺคลฺลาน ภิกฺขุ จิตฺตํ วสํ วตฺเตติ, โน จ ภิกฺขุ จิตฺตสฺส วเสน วตฺตติ, โส ยาย วิหาร\-สมาปตฺติยา อากงฺขติ ปุพฺพณฺหสมยํ วิหริตุํ, ตาย วิหาร\-สมาปตฺติยา ปุพฺพณฺหสมยํ วิหรติ. ยาย วิหาร\-สมาปตฺติยา อากงฺขติ มชฺฌนฺหิกสมยํ วิหริตุํ, ตาย วิหาร\-สมาปตฺติยา มชฺฌนฺหิกสมยํ วิหรติ. ยาย วิหาร\-สมาปตฺติยา อากงฺขติ สายนฺหสมยํ วิหริตุํ, ตาย วิหาร\-สมาปตฺติยา สายนฺหสมยํ วิหรติ. เสยฺยถาปิ อาวุโส โมคฺคลฺลาน รญฺโญ วา ราช\-มหามตฺตสฺส วา นานารตฺตานํ ทุสฺสานํ ทุสฺสกรณฺฑโก ปูโร อสฺส, โส ยญฺญเทว ทุสฺสยุคํ อากงฺเขยฺย ปุพฺพณฺหสมยํ ปารุปิตุํ, ตํ ตเทว ทุสฺสยุคํ ปุพฺพณฺหสมยํ ปารุเปยฺย. ยญฺญเทว ทุสฺสยุคํ อากงฺเขยฺย มชฺฌนฺหิกสมยํ ปารุปิตุํ, ตํ ตเทว ทุสฺสยุคํ มชฺฌนฺหิกสมยํ ปารุเปยฺย. ยญฺญเทว ทุสฺสยุคํ อากงฺเขยฺย สายนฺหสมยํ ปารุปิตุํ, ตํ ตเทว ทุสฺสยุคํ สายนฺหสมยํ ปารุเปยฺย. เอวเมว โข อาวุโส โมคฺคลฺลาน ภิกฺขุ จิตฺตํ วสํ วตฺเตติ, โน จ ภิกฺขุ จิตฺตสฺส วเสน วตฺตติ, โส ยาย วิหาร\-สมาปตฺติยา อากงฺขติ ปุพฺพณฺหสมยํ วิหริตุํ, ตาย วิหาร\-สมาปตฺติยา ปุพฺพณฺหสมยํ วิหรติ. ยาย วิหาร\-สมาปตฺติยา อากงฺขติ มชฺฌนฺหิกสมยํ วิหริตุํ, ตาย วิหาร\-สมาปตฺติยา มชฺฌนฺหิกสมยํ วิหรติ. ยาย วิหาร\-สมาปตฺติยา อากงฺขติ สายนฺหสมยํ วิหริตุํ, ตาย วิหาร\-สมาปตฺติยา สายนฺหสมยํ วิหรติ. เอวรูเปน โข อาวุโส โมคฺคลฺลาน ภิกฺขุนา โคสิงฺค\-สาลวนํ โสเภยฺยา”ติ. สาธุ สาธุ โมคฺคลฺลาน, ยถา ตํ สาริปุตฺโตว สมฺมา พฺยากรมาโน พฺยากเรยฺย. สาริปุตฺโต หิ โมคฺคลฺลาน จิตฺตํ วสํ วตฺเตติ, โน จ สาริปุตฺโต จิตฺตสฺส วเสน วตฺตติ. โส ยาย วิหาร\-สมาปตฺติยา อากงฺขติ ปุพฺพณฺหสมยํ วิหริตุํ, ตาย วิหาร\-สมาปตฺติยา ปุพฺพณฺหสมยํ วิหรติ. ยาย วิหาร\-สมาปตฺติยา อากงฺขติ มชฺฌนฺหิกสมยํ วิหริตุํ, ตาย วิหาร\-สมาปตฺติยา มชฺฌนฺหิกสมยํ วิหรติ. ยาย วิหาร\-สมาปตฺติยา อากงฺขติ สายนฺหสมยํ วิหริตุํ, ตาย วิหาร\-สมาปตฺติยา สายนฺหสมยํ วิหรตีติ.}

\proseitem{345}{\csromanfolio{281}เอวํ วุตฺเต อายสฺมา สาริปุตฺโต ภควนฺตํ เอตทโวจ “กสฺส นุ โข ภนฺเต สุภาสิตนฺ”ติ. สพฺเพสํ โว สาริปุตฺต สุภาสิตํ ปริยาเยน. อปิ จ มมปิ สุณาถ, ยถารูเปน ภิกฺขุนา โคสิงฺค\-สาลวนํ โสเภยฺย, อิธ สาริปุตฺต ภิกฺขุ ปจฺฉาภตฺตํ ปิณฺฑปาต\-ปฏิกฺกนฺโต นิสีทติ ปลฺลงฺกํ อาภุชิตฺวา อุชุํ กายํ ปณิธาย ปริมุขํ สติํ อุปฏฺฐเปตฺวา “น ตาวาหํ อิมํ ปลฺลงฺกํ ภินฺทิสฺสามิ, ยาว เม นานุปาทาย อาสเวหิ จิตฺตํ วิมุจฺจิสฺสตี”ติ. เอวรูเปน โข สาริปุตฺต ภิกฺขุนา โคสิงฺค\-สาลวนํ โสเภยฺยาติ.}

\prosewithcloserruled{อิทมโวจ ภควา. อตฺตมนา เต อายสฺมนฺโต\csromansharedfootnote{cb62f0f5e397900f}{เต ภิกฺขู (ก)} ภควโต ภาสิตํ อภินนฺทุนฺติ.}{มหา\-โคสิงฺค\-สุตฺตํ นิฏฺฐิตํ ทุติยํ.}
\csromanmatikaanchor{mtk.65}
\csromantocmarkref{subsection}{3. มหาโคปาลกสุตฺต}{mtk.65}
\bookmark[dest={mtk.65},level=2]{3. มหาโคปาลกสุตฺต}
\csromanheader{2}{3. มหา\-โคปาลก\-สุตฺต}
\proseitem{346}{เอวํ เม สุตํ\csromandash{}เอกํ สมยํ ภควา สาวตฺถิยํ วิหรติ เชตวเน อนาถ\-ปิณฺฑิกสฺส อาราเม. ตตฺร โข ภควา ภิกฺขู อามนฺเตสิ “ภิกฺขโว”ติ. “ภทนฺเต”ติ เต ภิกฺขู ภควโต ปจฺจสฺโสสุํ. ภควา เอตทโวจ\csromandash{}}

\prose{เอกาทสหิ ภิกฺขเว องฺเคหิ สมนฺนาคโต โคปาลโก อภพฺโพ โคคณํ ปริหริตุํ ผาติํ กาตุํ\csromansharedfootnote{8f615a0aa712bc43}{ผาติกตฺตุํ (สี, อิ), ผาติกาตุํ (สฺยา, กํ)}. กตเมหิ เอกาทสหิ. อิธ ภิกฺขเว โคปาลโก น รูปญฺญู โหติ, น ลกฺขณกุสโล โหติ, น อาสาฏิกํ หาเรตา\csromansharedfootnote{234f184101d8feb4}{สาเฏตา (สี, สฺยา, กํ, อิ)} โหติ, น วณํ ปฏิจฺฉาเทตา โหติ, น ธูมํ กตฺตา โหติ, น ติตฺถํ ชานาติ, น ปีตํ ชานาติ, น วีถิํ ชานาติ, น โคจรกุสโล โหติ, อนวเสสโทหี จ โหติ, เย เต อุสภา โคปิตโร โคปริณายกา, เต น อติเรกปูชาย ปูเชตา โหติ. อิเมหิ โข ภิกฺขเว เอกาทสหิ องฺเคหิ สมนฺนาคโต โคปาลโก อภพฺโพ โคคณํ ปริหริตุํ ผาติํ กาตุํ. เอวเมว โข ภิกฺขเว เอกาทสหิ ธมฺเมหิ สมนฺนาคโต ภิกฺขุ อภพฺโพ อิมสฺมิํ \csromanfolio{282}ธมฺมวินเย วุทฺธิํ วิรูฬฺหิํ เวปุลฺลํ อาปชฺชิตุํ. กตเมหิ เอกาทสหิ. อิธ ภิกฺขเว ภิกฺขุ น รูปญฺญู โหติ, น ลกฺขณกุสโล โหติ, น อาสาฏิกํ หาเรตา โหติ, น วณํ ปฏิจฺฉาเทตา โหติ, น ธูมํ กตฺตา โหติ, น ติตฺถํ ชานาติ, น ปีตํ ชานาติ, น วีถิํ ชานาติ, น โคจรกุสโล โหติ, อนวเสสโทหี จ โหติ, เย เต ภิกฺขู เถรา รตฺตญฺญู จิรปพฺพชิตา สํฆปิตโร สํฆปริณายกา, เต น อติเรกปูชาย ปูเชตา โหติ.}

\proseitem{347}{กถญฺจ ภิกฺขเว ภิกฺขุ น รูปญฺญู โหติ. อิธ ภิกฺขเว ภิกฺขุ “ยํ กิญฺจิ รูปํ สพฺพํ รูปํ จตฺตาริ มหาภูตานิ จตุนฺนญฺจ มหาภูตานํ อุปาทายรูปนฺ”ติ ยถาภูตํ นปฺปชานาติ. เอวํ โข ภิกฺขเว ภิกฺขุ น รูปญฺญู โหติ. (1)}

\prose{กถญฺจ ภิกฺขเว ภิกฺขุ น ลกฺขณกุสโล โหติ. อิธ ภิกฺขเว ภิกฺขุ “กมฺมลกฺขโณ พาโล, กมฺมลกฺขโณ ปณฺฑิโต”ติ ยถาภูตํ นปฺปชานาติ. เอวํ โข ภิกฺขเว ภิกฺขุ น ลกฺขณกุสโล โหติ. (2)}

\prose{กถญฺจ ภิกฺขเว ภิกฺขุ น อาสาฏิกํ หาเรตา โหติ. อิธ ภิกฺขเว ภิกฺขุ อุปฺปนฺนํ กามวิตกฺกํ อธิวาเสติ, นปฺปชหติ, น วิโนเทติ, น พฺยนฺตี กโรติ, น อนภาวํ คเมติ. อุปฺปนฺนํ พฺยาปาท\-วิตกฺกํ ฯเปฯ. อุปฺปนฺนํ วิหิํสา\-วิตกฺกํ ฯเปฯ. อุปฺปนฺนุปฺปนฺเน ปาปเก อกุสเล ธมฺเม อธิวาเสติ, นปฺปชหติ, น วิโนเทติ, น พฺยนฺตี กโรติ, น อนภาวํ คเมติ. เอวํ โข ภิกฺขเว ภิกฺขุ น อาสาฏิกํ หาเรตา โหติ. (3)}

\prose{กถญฺจ ภิกฺขเว ภิกฺขุ น วณํ ปฏิจฺฉาเทตา โหติ. อิธ ภิกฺขเว ภิกฺขุ จกฺขุนา รูปํ ทิสฺวา นิมิตฺตคฺคาหี โหติ อนุพฺยญฺชนคฺคาหี, ยตฺวาธิกรณเมนํ จกฺขุนฺทฺริยํ อสํวุตํ วิหรนฺตํ อภิชฺฌา\-โทมนสฺสา ปาปกา อกุสลา ธมฺมา อนฺวาสฺสเวยฺยุํ, ตสฺส สํวราย น ปฏิปชฺชติ, น รกฺขติ จกฺขุนฺทฺริยํ, จกฺขุนฺทฺริเย น สํวรํ อาปชฺชติ. โสเตน สทฺทํ สุตฺวา ฯเปฯ. ฆาเนน คนฺธํ ฆายิตฺวา ฯเปฯ. ชิวฺหาย รสํ สายิตฺวา ฯเปฯ. กาเยน โผฏฺฐพฺพํ ผุสิตฺวา ฯเปฯ. มนสา ธมฺมํ วิญฺญาย นิมิตฺตคฺคาหี โหติ อนุพฺยญฺชนคฺคาหี, ยตฺวาธิกรณเมนํ มนินฺทฺริยํ อสํวุตํ วิหรนฺตํ อภิชฺฌา\-โทมนสฺสา ปาปกา อกุสลา ธมฺมา อนฺวาสฺสเวยฺยุํ, ตสฺส สํวราย น ปฏิปชฺชติ, \csromanfolio{283}น รกฺขติ มนินฺทฺริยํ, มนินฺทฺริเย น สํวรํ อาปชฺชติ. เอวํ โข ภิกฺขเว ภิกฺขุ น วณํ ปฏิจฺฉาเทตา โหติ. (4)}

\prose{กถญฺจ ภิกฺขเว ภิกฺขุ น ธูมํ กตฺตา โหติ. อิธ ภิกฺขเว ภิกฺขุ ยถาสุตํ ยถา\-ปริยตฺตํ ธมฺมํ น วิตฺถาเรน ปเรสํ เทเสตา โหติ. เอวํ โข ภิกฺขเว ภิกฺขุ น ธูมํ กตฺตา โหติ. (5)}

\prose{กถญฺจ ภิกฺขเว ภิกฺขุ น ติตฺถํ ชานาติ. อิธ ภิกฺขเว ภิกฺขุ เย เต ภิกฺขู พหุสฺสุตา อาคตาคมา ธมฺมธรา วินยธรา มาติกาธรา, เต กาเลน กาลํ อุปสงฺกมิตฺวา น ปริปุจฺฉติ, น ปริปญฺหติ “อิทํ ภนฺเต กถํ, อิมสฺส โก อตฺโถ”ติ. ตสฺส เต อายสฺมนฺโต อวิวฏญฺเจว น วิวรนฺติ, อนุตฺตานีกตญฺจ น อุตฺตานี กโรนฺติ, อเนกวิหิเตสุ จ กงฺขา\-ฐานีเยสุ ธมฺเมสุ กงฺขํ น ปฏิวิโนเทนฺติ. เอวํ โข ภิกฺขเว ภิกฺขุ น ติตฺถํ ชานาติ. (6)}

\prose{กถญฺจ ภิกฺขเว ภิกฺขุ น ปีตํ ชานาติ. อิธ ภิกฺขเว ภิกฺขุ ตถาคต\-ปฺปเวทิเต ธมฺมวินเย เทสิยมาเน น ลภติ อตฺถเวทํ, น ลภติ ธมฺมเวทํ, น ลภติ ธมฺมู\-ปสํหิตํ ปาโมชฺชํ. เอวํ โข ภิกฺขเว ภิกฺขุ น ปีตํ ชานาติ. (7)}

\prose{กถญฺจ ภิกฺขเว ภิกฺขุ น วีถิํ ชานาติ. อิธ ภิกฺขเว ภิกฺขุ อริยํ อฏฺฐงฺคิกํ มคฺคํ ยถาภูตํ นปฺปชานาติ. เอวํ โข ภิกฺขเว ภิกฺขุ น วีถิํ ชานาติ. (8)}

\prose{กถญฺจ ภิกฺขเว ภิกฺขุ น โคจรกุสโล โหติ. อิธ ภิกฺขเว ภิกฺขุ จตฺตโร สติปฏฺฐาเน ยถาภูตํ นปฺปชานาติ. เอวํ โข ภิกฺขเว ภิกฺขุ น โคจรกุสโล โหติ. (9)}

\prose{กถญฺจ ภิกฺขเว ภิกฺขุ อนวเสสโทหี โหติ. อิธ ภิกฺขเว ภิกฺขุํ สทฺธา คหปติกา อภิหฏฺฐุํ ปวาเรนฺติ จีวร\-ปิณฺฑปาต\-เสนาสน\-คิลานปฺปจฺจย\-เภสชฺช\-ปริกฺขาเรหิ, ตตฺร ภิกฺขุ มตฺตํ น ชานาติ ปฏิคฺคหณาย. เอวํ โข ภิกฺขเว ภิกฺขุ อนวเสสโทหี โหติ. (10)}

\prose{กถญฺจ ภิกฺขเว ภิกฺขุ เย เต ภิกฺขู เถรา รตฺตญฺญู จิรปพฺพชิตา สํฆปิตโร สํฆปริณายกา, เต น อติเรกปูชาย ปูเชตา \csromanfolio{284}โหติ. อิธ ภิกฺขเว ภิกฺขุ เย เต ภิกฺขู เถรา รตฺตญฺญู จิรปพฺพชิตา สํฆปิตโร สํฆปริณายกา, เตสุ น เมตฺตํ กายกมฺมํ ปจฺจุปฏฺฐาเปติ อาวิ เจว รโห จ. น เมตฺตํ วจีกมฺมํ ปจฺจุปฏฺฐาเปติ อาวิ เจว รโห จ. น เมตฺตํ มโนกมฺมํ ปจฺจุปฏฺฐาเปติ อาวิ เจว รโห จ. เอวํ โข ภิกฺขเว ภิกฺขุ เย เต ภิกฺขู เถรา รตฺตญฺญู จิรปพฺพชิตา สํฆปิตโร สํฆปริณายกา, เต น อติเรกปูชาย ปูเชตา โหติ. (11)}

\proseitem{3}{อิเมหิ โข ภิกฺขเว เอกาทสหิ ธมฺเมหิ สมนฺนาคโต ภิกฺขุ อภพฺโพ อิมสฺมิํ ธมฺมวินเย วุทฺธิํ วิรูฬฺหิํ เวปุลฺลํ อาปชฺชิตุํ.}

\proseitem{348}{เอกาทสหิ ภิกฺขเว องฺเคหิ สมนฺนาคโต โคปาลโก ภพฺโพ โคคณํ ปริหริตุํ ผาติํ กาตุํ. กตเมหิ เอกาทสหิ. อิธ ภิกฺขเว โคปาลโก รูปญฺญู โหติ, ลกฺขณกุสโล โหติ, อาสาฏิกํ หาเรตา โหติ, วณํ ปฏิจฺฉาเทตา โหติ, ธูมํ กตฺตา โหติ, ติตฺถํ ชานาติ, ปีตํ ชานาติ, วีถิํ ชานาติ, โคจรกุสโล โหติ, สาวเสสโทหี จ โหติ, เย เต อุสภา โคปิตโร โคปริณายกา, เต อติเรกปูชาย ปูเชตา โหติ. อิเมหิ โข ภิกฺขเว เอกาทสหิ องฺเคหิ สมนฺนาคโต โคปาลโก ภพฺโพ โคคณํ ปริหริตุํ ผาติํ กาตุํ. เอวเมว โข ภิกฺขเว เอกาทสหิ ธมฺเมหิ สมนฺนาคโต ภิกฺขุ ภพฺโพ อิมสฺมิํ ธมฺมวินเย วุทฺธิํ วิรูฬฺหิํ เวปุลฺลํ อาปชฺชิตุํ. กตเมหิ เอกาทสหิ. อิธ ภิกฺขเว ภิกฺขุ รูปญฺญู โหติ, ลกฺขณกุสโล โหติ, อาสาฏิกํ หาเรตา โหติ, วณํ ปฏิจฺฉาเทตา โหติ, ธูมํ กตฺตา โหติ, ติตฺถํ ชานาติ, ปีตํ ชานาติ, วีถิํ ชานาติ, โคจรกุสโล โหติ, สาวเสสโทหี จ โหติ, เย เต ภิกฺขู เถรา รตฺตญฺญู จิรปพฺพชิตา สํฆปิตโร สํฆปริณายกา, เต อติเรกปูชาย ปูเชตา โหติ.}

\proseitem{349}{กถญฺจ ภิกฺขเว ภิกฺขุ รูปญฺญู โหติ. อิธ ภิกฺขเว ภิกฺขุ “ยํ กิญฺจิ รูปํ สพฺพํ รูปํ จตฺตาริ มหาภูตานิ จตุนฺนญฺจ มหาภูตานํ อุปาทายรูปนฺ”ติ ยถาภูตํ ปชานาติ. เอวํ โข ภิกฺขเว ภิกฺขุ รูปญฺญู โหติ. (1)}

\prose{กถญฺจ ภิกฺขเว ภิกฺขุ ลกฺขณกุสโล โหติ. อิธ ภิกฺขเว ภิกฺขุ “กมฺมลกฺขโณ พาโล, กมฺมลกฺขโณ ปณฺฑิโต”ติ ยถาภูตํ ปชานาติ. เอวํ โข ภิกฺขเว ภิกฺขุ ลกฺขณกุสโล โหติ. (2)}

\prose{\csromanfolio{285}กถญฺจ ภิกฺขเว ภิกฺขุ อาสาฏิกํ หาเรตา โหติ. อิธ ภิกฺขเว ภิกฺขุ อุปฺปนฺนํ กามวิตกฺกํ นาธิวาเสติ, ปชหติ, วิโนเทติ, พฺยนฺตี กโรติ, อนภาวํ คเมติ. อุปฺปนฺนํ พฺยาปาท\-วิตกฺกํ ฯเปฯ. อุปฺปนฺนํ วิหิํสา\-วิตกฺกํ ฯเปฯ. อุปฺปนฺนุปฺปนฺเน ปาปเก อกุสเล ธมฺเม นาธิวาเสติ, ปชหติ, วิโนเทติ, พฺยนฺตี กโรติ, อนภาวํ คเมติ. เอวํ โข ภิกฺขเว ภิกฺขุ อาสาฏิกํ หาเรตา โหติ. (3)}

\prose{กถญฺจ ภิกฺขเว ภิกฺขุ วณํ ปฏิจฺฉาเทตา โหติ. อิธ ภิกฺขเว ภิกฺขุ จกฺขุนา รูปํ ทิสฺวา น นิมิตฺตคฺคาหี โหติ นานุพฺยญฺชน\-คฺคาหี, ยตฺวาธิกรณเมนํ จกฺขุนฺทฺริยํ อสํวุตํ วิหรนฺตํ อภิชฺฌา\-โทมนสฺสา ปาปกา อกุสลา ธมฺมา อนฺวาสฺสเวยฺยุํ, ตสฺส สํวราย ปฏิปชฺชติ, รกฺขติ จกฺขุนฺทฺริยํ, จกฺขุนฺทฺริเย สํวรํ อาปชฺชติ. โสเตน สทฺทํ สุตฺวา ฯเปฯ. ฆาเนน คนฺธํ ฆายิตฺวา ฯเปฯ. ชิวฺหาย รสํ สายิตฺวา ฯเปฯ. กาเยน โผฏฺฐพฺพํ ผุสิตฺวา ฯเปฯ. มนสา ธมฺมํ วิญฺญาย น นิมิตฺตคฺคาหี โหติ นานุพฺยญฺชน\-คฺคาหี, ยตฺวาธิกรณเมนํ มนินฺทฺริยํ อสํวุตํ วิหรนฺตํ อภิชฺฌา\-โทมนสฺสา ปาปกา อกุสลา ธมฺมา อนฺวาสฺสเวยฺยุํ, ตสฺส สํวราย ปฏิปชฺชติ, รกฺขติ มนินฺทฺริยํ, มนินฺทฺริเย สํวรํ อาปชฺชติ. เอวํ โข ภิกฺขเว ภิกฺขุ วณํ ปฏิจฺฉาเทตา โหติ. (4)}

\prose{กถญฺจ ภิกฺขเว ภิกฺขุ ธูมํ กตฺตา โหติ. อิธ ภิกฺขเว ภิกฺขุ ยถาสุตํ ยถา\-ปริยตฺตํ ธมฺมํ วิตฺถาเรน ปเรสํ เทเสตา โหติ. เอวํ โข ภิกฺขเว ภิกฺขุ ธูมํ กตฺตา โหติ. (5)}

\prose{กถญฺจ ภิกฺขเว ภิกฺขุ ติตฺถํ ชานาติ. อิธ ภิกฺขเว ภิกฺขุ เย เต ภิกฺขู พหุสฺสุตา อาคตาคมา ธมฺมธรา วินยธรา มาติกาธรา, เต กาเลน กาลํ อุปสงฺกมิตฺวา ปริปุจฺฉติ, ปริปญฺหติ “อิทํ ภนฺเต กถํ, อิมสฺส โก อตฺโถ”ติ. ตสฺส เต อายสฺมนฺโต อวิวฏญฺเจว วิวรนฺติ. อนุตฺตานีกตญฺจ อุตฺตานี กโรนฺติ. อเนกวิหิเตสุ จ กงฺขา\-ฐานีเยสุ ธมฺเมสุ กงฺขํ ปฏิวิโนเทนฺติ. เอวํ โข ภิกฺขเว ภิกฺขุ ติตฺถํ ชานาติ. (6)}

\prose{กถญฺจ ภิกฺขเว ภิกฺขุ ปีตํ ชานาติ. อิธ ภิกฺขเว ภิกฺขุ ตถาคต\-ปฺปเวทิเต ธมฺมวินเย เทสิยมาเน ลภติ อตฺถเวทํ, ลภติ ธมฺมเวทํ, ลภติ ธมฺมู\-ปสํหิตํ ปาโมชฺชํ. เอวํ โข ภิกฺขเว ภิกฺขุ ปีตํ ชานาติ. (7)}

\prose{\csromanfolio{286}กถญฺจ ภิกฺขเว ภิกฺขุ วีถิํ ชานาติ. อิธ ภิกฺขเว ภิกฺขุ อริยํ อฏฺฐงฺคิกํ มคฺคํ ยถาภูตํ ปชานาติ. เอวํ โข ภิกฺขเว ภิกฺขุ วีถิํ ชานาติ. (8)}

\prose{กถญฺจ ภิกฺขเว ภิกฺขุ โคจรกุสโล โหติ. อิธ ภิกฺขเว ภิกฺขุ จตฺตาโร สติปฏฺฐาเน ยถาภูตํ ปชานาติ. เอวํ โข ภิกฺขเว ภิกฺขุ โคจรกุสโล โหติ. (9)}

\prose{กถญฺจ ภิกฺขเว ภิกฺขุ สาวเสสโทหี โหติ. อิธ ภิกฺขเว ภิกฺขุํ สทฺธา คหปติกา อภิหฏฺฐุํ ปวาเรนฺติ จีวร\-ปิณฺฑปาต\-เสนาสน\-คิลานปฺปจฺจย\-เภสชฺช\-ปริกฺขาเรหิ, ตตฺร ภิกฺขุ มตฺตํ ชานาติ ปฏิคฺคหณาย. เอวํ โข ภิกฺขเว ภิกฺขุ สาวเสสโทหี โหติ. (10)}

\prose{กถญฺจ ภิกฺขเว ภิกฺขุ เย เต ภิกฺขู เถรา รตฺตญฺญู จิรปพฺพชิตา สํฆปิตโร สํฆปริณายกา, เต อติเรกปูชาย ปูเชตา โหติ. อิธ ภิกฺขเว ภิกฺขุ เย เต ภิกฺขู เถรา รตฺตญฺญู จิรปพฺพชิตา สํฆปิตโร สํฆปริณายกา, เตสุ เมตฺตํ กายกมฺมํ ปจฺจุปฏฺฐาเปติ อาวิ เจว รโห จ. เมตฺตํ วจีกมฺมํ ปจฺจุปฏฺฐาเปติ อาวิ เจว รโห จ. เมตฺตํ มโนกมฺมํ ปจฺจุปฏฺฐาเปติ อาวิ เจว รโห จ. เอวํ โข ภิกฺขเว ภิกฺขุ เย เต ภิกฺขู เถรา รตฺตญฺญู จิรปพฺพชิตา สํฆปิตโร สํฆปริณายกา, เต อติเรกปูชาย ปูเชตา โหติ. (11)}

\prose{อิเมหิ โข ภิกฺขเว เอกาทสหิ ธมฺเมหิ สมนฺนาคโต ภิกฺขุ ภพฺโพ อิมสฺมิํ ธมฺมวินเย วุทฺธิํ วิรูฬฺหิํ เวปุลฺลํ อาปชฺชิตุนฺติ.}

\prosewithcloserruled{อิทมโวจ ภควา. อตฺตมนา เต ภิกฺขู ภควโต ภาสิตํ อภินนฺทุนฺติ.}{มหา\-โคปาลก\-สุตฺตํ นิฏฺฐิตํ ตติยํ.}
\csromanmatikaanchor{mtk.66}
\csromantocmarkref{subsection}{4. จูฬโคปาลกสุตฺต}{mtk.66}
\bookmark[dest={mtk.66},level=2]{4. จูฬโคปาลกสุตฺต}
\csromanheader{2}{4. จูฬโคปาลก\-สุตฺต}
\proseitem{350}{เอวํ เม สุตํ\csromandash{}เอกํ สมยํ ภควา วชฺชีสุ วิหรติ อุกฺกเจลายํ คงฺคาย นทิยา ตีเร. ตตฺร โข ภควา ภิกฺขู อามนฺเตสิ “ภิกฺขโว”ติ. “ภทนฺเต”ติ เต ภิกฺขู ภควโต ปจฺจสฺโสสุํ. ภควา เอตทโวจ\csromandash{}}

\prose{\csromanfolio{287}ภูตปุพฺพํ ภิกฺขเว มาคธโก โคปาลโก ทุปฺปญฺญชาติโก วสฺสานํ ปจฺฉิเม มาเส สรทสมเย อสมเวกฺขิตฺวา คงฺคาย นทิยา โอริมํ ตีรํ อสมเวกฺขิตฺวา ปาริมํ ตีรํ อติตฺเถเนว คาโว ปตาเรสิ อุตฺตรํ ตีรํ สุวิเทหานํ. อถ โข ภิกฺขเว คาโว มชฺเฌคงฺคาย นทิยา โสเต อามณฺฑลิยํ กริตฺวา ตตฺเถว อนยพฺยสนํ อาปชฺชิํสุ. ตํ กิสฺส เหตุ, ตถา หิ โส ภิกฺขเว มาคธโก โคปาลโก ทุปฺปญฺญชาติโก วสฺสานํ ปจฺฉิเม มาเส สรทสมเย อสมเวกฺขิตฺวา คงฺคาย นทิยา โอริมํ ตีรํ อสมเวกฺขิตฺวา ปาริมํ ตีรํ อติตฺเถเนว คาโว ปตาเรสิ อุตฺตรํ ตีรํ สุวิเทหานํ. เอวเมว โข ภิกฺขเว เย หิ เกจิ\csromansharedfootnote{54a4e251f1bb97b8}{เย เกจิ (สฺยา, กํ)} สมณา วา พฺราหฺมณา วา อกุสลา อิมสฺส โลกสฺส, อกุสลา ปรสฺส โลกสฺส, อกุสลา มารเธยฺยสฺส, อกุสลา อมารเธยฺยสฺส, อกุสลา มจฺจุเธยฺยสฺส, อกุสลา อมจฺจุเธยฺยสฺส. เตสํ เย โสตพฺพํ สทฺทหาตพฺพํ มญฺญิสฺสนฺติ, เตสํ ตํ ภวิสฺสติ ทีฆรตฺตํ อหิตาย ทุกฺขาย.}

\proseitem{351}{ภูตปุพฺพํ ภิกฺขเว มาคธโก โคปาลโก สปฺปญฺญชาติโก วสฺสานํ ปจฺฉิเม มาเส สรทสมเย สมเวกฺขิตฺวา คงฺคาย นทิยา โอริมํ ตีรํ สมเวกฺขิตฺวา ปาริมํ ตีรํ ติตฺเถเนว คาโว ปตาเรสิ อุตฺตรํ ตีรํ สุวิเทหานํ. โส ปฐมํ ปตาเรสิ เย เต อุสภา โคปิตโร โคปริณายกา, เต ติริยํ คงฺคาย โสตํ เฉตฺวา โสตฺถินา ปารํ อคมํสุ. อถาปเร ปตาเรสิ พลวคาโว ทมฺมคาโว เตปิ ติริยํ คงฺคาย โสตํ เฉตฺวา โสตฺถินา ปารํ อคมํสุ. อถาปเร ปตาเรสิ วจฺฉตเร วจฺฉตริโย, เตปิ ติริยํ คงฺคาย โสตํ เฉตฺวา โสตฺถินา ปารํ อคมํสุ. อถาปเร ปตาเรสิ วจฺฉเก กิสาพลเก\csromansharedfootnote{5152d8e4bdc9eb17}{กิสพลเก (สี, สฺยา, อิ)}, เตปิ ติริยํ คงฺคาย โสตํ เฉตฺวา โสตฺถินา ปารํ อคมํสุ. ภูตปุพฺพํ ภิกฺขเว วจฺฉโก ตรุณโก ตาวเทว ชาตโก มาตุโครวเกน วุยฺหมาโน, โสปิ ติริยํ คงฺคาย โสตํ เฉตฺวา โสตฺถินา ปรํ อคมาสิ. ตํ กิสฺส เหตุ, ตถา หิ โส ภิกฺขเว มาคธโก โคปาลโก สปฺปญฺญชาติโก วสฺสานํ ปจฺฉิเม มาเส สรทสมเย สมเวกฺขิตฺวา คงฺคาย นทิยา \csromanfolio{288}โอริมํ ตีรํ สมเวกฺขิตฺวา ปาริมํ ตีรํ ติตฺเถเนว คาโว ปตาเรสิ อุตฺตรํ ตีรํ สุวิเทหานํ. เอวเมว โข ภิกฺขเว เย หิ เกจิ สมณา วา พฺราหฺมณา วา กุสลา อิมสฺส โลกสฺส, กุสลา ปรสฺส โลกสฺส, กุสลา มารเธยฺยสฺส, กุสลา อมารเธยฺยสฺส, กุสลา มจฺจุเธยฺยสฺส, กุสลา อมจฺจุเธยฺยสฺส. เตสํ เย โสตพฺพํ สทฺทหาตพฺพํ มญฺญิสฺสนฺติ, เตสํ ตํ ภวิสฺสติ ทีฆรตฺตํ หิตาย สุขาย.}

\proseitem{352}{เสยฺยถาปิ ภิกฺขเว เย เต อุสภา โคปิตโร โคปริณายกา, เต ติริยํ คงฺคาย โสตํ เฉตฺวา โสตฺถินา ปารํ อคมํสุ. เอวเมว โข ภิกฺขเว เย เต ภิกฺขู อรหนฺโต ขีณาสวา วุสิตวนฺโต กตกรณียา โอหิตภารา อนุปฺปตฺต\-สทตฺถา ปริกฺขีณ\-ภวสํโยชนา สมฺมทญฺญา วิมุตฺตา, เต ติริยํ มารสฺส โสตํ เฉตฺวา โสตฺถินา ปารํ คตา.}

\prose{เสยฺยถาปิ เต ภิกฺขเว พลวคาโว ทมฺมคาโว ติริยํ คงฺคาย โสตํ เฉตฺวา โสตฺถินา ปารํ อคมํสุ. เอวเมว โข ภิกฺขเว เย เต ภิกฺขู ปญฺจนฺนํ โอรมฺภาคิยานํ สํโยชนานํ ปริกฺขยา โอปปาติกา ตตฺถ ปรินิพฺพายิโน, อนาวตฺติธมฺมา ตสฺมา โลกา, เตปิ ติริยํ มารสฺส โสตํ เฉตฺวา โสตฺถินา ปรํ คมิสฺสนฺติ.}

\prose{เสยฺยถาปิ เต ภิกฺขเว วจฺฉตรา วจฺฉตริโย ติริยํ คงฺคาย โสตํ เฉตฺวา โสตฺถินา ปารํ อคมํสุ. เอวเมว โข ภิกฺขเว เย เต ภิกฺขู ติณฺณํ สํโยชนานํ ปริกฺขยา ราค\-โทส\-โมหานํ ตนุตฺตา สกทาคามิโน สกิเทว อิมํ โลกํ อาคนฺตฺวา ทุกฺขสฺสนฺตํ กริสฺสนฺติ, เตปิ ติริยํ มารสฺส โสตํ เฉตฺวา โสตฺถินา ปารํ คมิสฺสนฺติ.}

\prose{เสยฺยถาปิ เต ภิกฺขเว วจฺฉกา กิสาพลกา ติริยํ คงฺคาย โสตํ เฉตฺวา โสตฺถินา ปารํ อคมํสุ. เอวเมว โข ภิกฺขเว เย เต ภิกฺขู ติณฺณํ สํโยชนานํ ปริกฺขยา โสตาปนฺนา อวินิปาต\-ธมฺมา นิยตา สมฺโพธิ\-ปรายนา, เตปิ ติริยํ มารสฺส โสตํ เฉตฺวา โสตฺถินา ปารํ คมิสฺสนฺติ.}

\prose{เสยฺยถาปิ โส ภิกฺขเว วจฺฉโก ตรุณโก ตาวเทว ชาตโก มาตุโครวเกน วุยฺหมาโน ติริยํ คงฺคาย โสตํ เฉตฺวา \csromanfolio{289}โสตฺถินา ปารํ อคมาสิ. เอวเมว โข ภิกฺขเว เย เต ภิกฺขู ธมฺมานุสาริโน สทฺธานุสาริโน, เตปิ ติริยํ มารสฺส โสตํ เฉตฺวา โสตฺถินา ปารํ คมิสฺสนฺติ.}

\prose{อหํ โข ปน ภิกฺขเว กุสโล อิมสฺส โลกสฺส, กุสโล ปรสฺส โลกสฺส, กุสโล มารเธยฺยสฺส, กุสโล อมารเธยฺยสฺส, กุสโล มจฺจุเธยฺยสฺส, กุสโล อมจฺจุเธยฺยสฺส. ตสฺส มยฺหํ ภิกฺขเว เย โสตพฺพํ สทฺทหาตพฺพํ มญฺญิสฺสนฺติ, เตสํ ตํ ภวิสฺสติ ทีฆรตฺตํ หิตาย สุขายาติ.}

\prose{อิทมโวจ ภควา, อิทํ วตฺวา สุคโต อถาปรํ เอตทโวจ สตฺถา.}

\setlength{\csromangathapagewidth}{0pt}
\setlength{\csromangathawakwidth}{0pt}
\csromangathameasuregroup{“อยํ โลโก ปโร โลโก, \\ ยญฺจ มาเรน สมฺปตฺตํ, \\ สพฺพํ โลกํ อภิญฺญาย, \\ วิวฏํ อมตทฺวารํ, \\ ฉินฺนํ ปาปิมโต โสตํ, \\ ปาโมชฺชพหุลา โหถ,}{\csromangathabat{“อยํ โลโก ปโร โลโก,}{ชานตา สุปฺปกาสิโต.} \\ \csromangathabat{ยญฺจ มาเรน สมฺปตฺตํ,}{อปฺปตฺตํ ยญฺจ มจฺจุนา.} \\ \csromangathabat{สพฺพํ โลกํ อภิญฺญาย,}{สมฺพุทฺเธน ปชานตา.} \\ \csromangathabat{วิวฏํ อมตทฺวารํ,}{เขมํ นิพฺพานปตฺติยา.} \\ \csromangathabat{ฉินฺนํ ปาปิมโต โสตํ,}{วิทฺธสฺตํ วินฬีกตํ.} \\ \csromangathabat{ปาโมชฺชพหุลา โหถ,}{เขมํ ปตฺต’ตฺถ\textsuperscript{1} ภิกฺขโว”ติ.}}
\csromangathasetleft{“อยํ โลโก ปโร โลโก, \\ ยญฺจ มาเรน สมฺปตฺตํ, \\ สพฺพํ โลกํ อภิญฺญาย, \\ วิวฏํ อมตทฺวารํ, \\ ฉินฺนํ ปาปิมโต โสตํ, \\ ปาโมชฺชพหุลา โหถ,}
\csromangathagroupwithcloserruled{\csromangathastanza{\csromangathabat{“อยํ โลโก ปโร โลโก,}{ชานตา สุปฺปกาสิโต.} \\ \csromangathabat{ยญฺจ มาเรน สมฺปตฺตํ,}{อปฺปตฺตํ ยญฺจ มจฺจุนา.}}\csromangathastanza{\csromangathabat{สพฺพํ โลกํ อภิญฺญาย,}{สมฺพุทฺเธน ปชานตา.} \\ \csromangathabat{วิวฏํ อมตทฺวารํ,}{เขมํ นิพฺพานปตฺติยา.}}\csromangathastanza{\csromangathabat{ฉินฺนํ ปาปิมโต โสตํ,}{วิทฺธสฺตํ วินฬีกตํ.} \\ \csromangathabat{ปาโมชฺชพหุลา โหถ,}{เขมํ ปตฺต’ตฺถ\csromansharedfootnote{0f7cf750adb43a65}{ปตฺเถถ (สฺยา, กํ, ก, อฏฺฐกถายํ สํวณฺเณตพฺพปาโฐ)} ภิกฺขโว”ติ.}}}{จูฬโคปาลก\-สุตฺตํ นิฏฺฐิตํ จตุตฺถํ.}
\csromanmatikaanchor{mtk.67}
\csromantocmarkref{subsection}{5. จูฬสจฺจกสุตฺต}{mtk.67}
\bookmark[dest={mtk.67},level=2]{5. จูฬสจฺจกสุตฺต}
\csromanheader{2}{5. จูฬสจฺจก\-สุตฺต}
\proseitem{353}{เอวํ เม สุตํ\csromandash{}เอกํ สมยํ ภควา เวสาลิยํ วิหรติ มหาวเน กูฏาคาร\-สาลายํ. เตน โข ปน สมเยน สจฺจโก นิคณฺฐปุตฺโต เวสาลิยํ ปฏิวสติ ภสฺส\-ปฺปวาทโก ปณฺฑิตวาโท สาธุสมฺมโต พหุชนสฺส. โส เวสาลิยํ ปริสติ เอวํ วาจํ ภาสติ “นาหํ ตํ ปสฺสามิ สมณํ วา พฺราหฺมณํ วา สงฺฆิํ คณิํ คณาจริยํ, อปิ อรหนฺตํ สมฺมา\-สมฺพุทฺธํ ปฏิชานมานํ, โย มยา วาเทน วาทํ สมารทฺโธ น สงฺกมฺเปยฺย น สมฺปกมฺเปยฺย น สมฺปเวเธยฺย, ยสฺส น กจฺเฉหิ เสทา มุจฺเฉยฺยุํ, ถูณํ เจปาหํ อเจตนํ วาเทน สมารเภยฺยํ, สาปิ มยา วาเทน วาทํ \csromanfolio{290}สมารทฺธา สงฺกมฺเปยฺย สมฺปกมฺเปยฺย สมฺปเวเธยฺย, โก ปน วาโท มนุสฺสภูตสฺสา”ติ. อถ โข อายสฺมา อสฺสชิ ปุพฺพณฺหสมยํ นิวาเสตฺวา ปตฺต\-จีวรมา\-ทาย เวสาลิํ ปิณฺฑาย ปาวิสิ. อทฺทสา โข สจฺจโก นิคณฺฐปุตฺโต เวสาลิยํ ชงฺฆาวิหารํ อนุจงฺกมมาโน อนุวิจรมาโน อายสฺมนฺตํ อสฺสชิํ ทูรโตว อาคจฺฉนฺตํ, ทิสฺวาน เยนายสฺมา อสฺสชิ เตนุปสงฺกมิ, อุปสงฺกมิตฺวา อายสฺมตา อสฺสชินา สทฺธิํ สมฺโมทิ, สมฺโมทนียํ กถํ สารณียํ วีติสาเรตฺวา เอกมนฺตํ อฏฺฐาสิ, เอกมนฺตํ ฐิโต โข สจฺจโก นิคณฺฐปุตฺโต อายสฺมนฺตํ อสฺสชิํ เอตทโวจ “กถํ ปน โภ อสฺสชิ สมโณ โคตโม สาเวเก วิเนติ, กถํภาคา จ ปน สมณสฺส โคตมสฺส สาวเกสุ อนุสาสนี พหุลา ปวตฺตตี”ติ. เอวํ โข อคฺคิเวสฺสน ภควา สาวเก วิเนติ, เอวํภาคา จ ปน ภควโต สาวเกสุ อนุสาสนี พหุลา ปวตฺตติ “รูปํ ภิกฺขเว อนิจฺจํ, เวทนา อนิจฺจา, สญฺญา อนิจฺจา, สงฺขารา อนิจฺจา, วิญฺญาณํ อนิจฺจํ. รูปํ ภิกฺขเว อนตฺตา, เวทนา อนตฺตา, สญฺญา อนตฺตา, สงฺขารา อนตฺตา, วิญฺญาณํ อนตฺตา. สพฺเพ สงฺขารา อนิจฺจา, สพฺเพ ธมฺมา อนตฺตา”ติ. เอวํ โข อคฺคิเวสฺสน ภควา สาวเก วิเนติ, เอวํภาคา จ ปน ภควโต สาวเกสุ อนุสาสนี พหุลา ปวตฺตตีติ. ทุสฺสุตํ วต โภ อสฺสชิ อสฺสุมฺห, เย มยํ เอวํวาทิํ สมณํ โคตมํ อสฺสุมฺห, อปฺเปว นาม มยํ กทาจิ กรหจิ เตน โภตา โคตเมน สทฺธิํ สมาคจฺเฉยฺยาม, อปฺเปว นาม สิยา โกจิเทว กถาสลฺลาโป, อปฺเปว นาม ตสฺมา ปาปกา ทิฏฺฐิคตา วิเวเจยฺยามาติ.}

\proseitem{354}{เตน โข ปน สมเยน ปญฺจมตฺตานิ ลิจฺฉวิสตานิ สนฺถาคาเร\csromansharedfootnote{bd376d5b59240561}{สนฺธาคาเร (ก)} สนฺนิปติตานิ โหนฺติ เกนจิเทว กรณีเยน. อถ โข สจฺจโก นิคณฺฐปุตฺโต เยน เต ลิจฺฉวี เตนุปสงฺกมิ, อุปสงฺกมิตฺวา เต ลิจฺฉวี เอตทโวจ “อภิกฺกมนฺตุ โภนฺโต ลิจฺฉวี, อภิกฺกมนฺตุ โภนฺโต ลิจฺฉวี, อชฺช เม สมเณน โคตเมน สทฺธิํ กถาสลฺลาโป ภวิสฺสติ. สเจ เม สมโณ โคตโม ตถา ปติฏฺฐิสฺสติ, ยถา จ เม\csromansharedfootnote{4e91206fdb6aac46}{ยถาสฺส เม (สี, อิ)} ญาตญฺญตเรน สาวเกน อสฺสชินา นาม ภิกฺขุนา ปติฏฺฐิตํ. เสยฺยถาปิ นาม พลวา ปุริโส ทีฆโลมิกํ เอฬกํ โลเมสุ คเหตฺวา อากฑฺเฒยฺย ปริกฑฺเฒยฺย \csromanfolio{291}สมฺปริกฑฺเฒยฺย, เอวเมวาหํ สมณํ โคตมํ วาเทน วาทํ อากฑฺฒิสฺสามิ ปริกฑฺฒิสฺสามิ สมฺปริกฑฺฒิสฺสามิ. เสยฺยถาปิ นาม พลวา โสณฺฑิกา\-กมฺมกาโร มหนฺตํ โสณฺฑิกา\-กิฬญฺชํ คมฺภีเร อุทกรหเท ปกฺขิปิตฺวา กณฺเณ คเหตฺวา อากฑฺเฒยฺย ปริกฑฺเฒยฺย สมฺปริกฑฺเฒยฺย, เอวเมวาหํ สมณํ โคตมํ วาเทน วาทํ อากฑฺฒิสฺสามิ ปริกฑฺฒิสฺสามิ สมฺปริกฑฺฒิสฺสามิ. เสยฺยถาปิ นาม พลวา โสณฺฑิกาธุตฺโต วาลํ\csromansharedfootnote{f35df81502f8fe06}{ถาลํ (ก)} กณฺเณ คเหตฺวา โอธุเนยฺย นิทฺธุเนยฺย นิปฺโผเฏยฺย\csromansharedfootnote{eed4384fa5a97614}{นิจฺฉาเทยฺย (สี, อิ, ก), นิจฺโฉเฏยฺย (ก), นิปฺโปเฐยฺย (สฺยา, กํ)}, เอวเมวาหํ สมณํ โคตมํ วาเทน วาทํ โอธุนิสฺสามิ นิทฺธุนิสฺสามิ นิปฺโผเฏสฺสามิ. เสยฺยถาปิ นาม กุญฺชโร สฏฺฐิหายโน คมฺภีรํ โปกฺขรณิํ โอคาเหตฺวา สาณโธวิกํ นาม กีฬิตชาตํ กีฬติ, เอวเมวาหํ สมณํ โคตมํ สาณโธวิกํ มญฺเญ กีฬิตชาตํ กีฬิสฺสามิ. อภิกฺกมนฺตุ โภนฺโต ลิจฺฉวี, อภิกฺกมนฺตุ โภนฺโต ลิจฺฉวี, อชฺช เม สมเณน โคตเมน สทฺธิํ กถาสลฺลาโป ภวิสฺสตี”ติ. ตเตฺรกจฺเจ ลิจฺฉวี เอวมาหํสุ “กิํ สมโณ โคตโม สจฺจกสฺส นิคณฺฐ\-ปุตฺตสฺส วาทํ อาโรเปสฺสติ, อถ โข สจฺจโก นิคณฺฐปุตฺโต สมณสฺส โคตมสฺส วาทํ อาโรเปสฺสตี”ติ. เอกจฺเจ ลิจฺฉวี เอวมาหํสุ “กิํ โส ภวมาโน สจฺจโก นิคณฺฐปุตฺโต, โย ภควโต วาทํ อาโรเปสฺสติ, อถ โข ภควา สจฺจกสฺส นิคณฺฐ\-ปุตฺตสฺส วาทํ อาโรเปสฺสตี”ติ. อถ โข สจฺจโก นิคณฺฐปุตฺโต ปญฺจมตฺเตหิ ลิจฺฉวิสเตหิ ปริวุโต เยน มหาวนํ กูฏาคารสาลา เตนุปสงฺกมิ.}

\proseitem{355}{เตน โข ปน สมเยน สมฺพหุลา ภิกฺขู อพฺโภกาเส จงฺกมนฺติ. อถ โข สจฺจโก นิคณฺฐปุตฺโต เยน เต ภิกฺขู เตนุปสงฺกมิ, อุปสงฺกมิตฺวา เต ภิกฺขู เอตทโวจ “กหํ นุ โข โภ เอตรหิ โส ภวํ โคตโม วิหรติ, ทสฺสนกามา หิ มยํ ตํ ภวนฺตํ โคตมนฺ”ติ. เอส อคฺคิเวสฺสน ภควา มหาวนํ อชฺโฌคาเหตฺวา อญฺญตรสฺมิํ รุกฺขมูเล ทิวาวิหารํ นิสินฺโนติ. อถ โข สจฺจโก นิคณฺฐปุตฺโต มหติยา ลิจฺฉวิ\-ปริสาย สทฺธิํ มหาวนํ อชฺโฌคาเหตฺวา เยน ภควา เตนุปสงฺกมิ, อุปสงฺกมิตฺวา ภควตา สทฺธิํ สมฺโมทิ, สมฺโมทนียํ กถํ สารณียํ วีติสาเรตฺวา เอกมนฺตํ นิสีทิ. เตปิ โข ลิจฺฉวี อปฺเปกจฺเจ ภควนฺตํ อภิวาเทตฺวา เอกมนฺตํ นิสีทิํสุ. อปฺเปกจฺเจ \csromanfolio{292}ภควตา สทฺธิํ สมฺโมทิํสุ, สมฺโมทนียํกถํ สารณียํ วีติสาเรตฺวา เอกมนฺตํ นิสีทิํสุ. อปฺเปกจฺเจ เยน ภควา เตนญฺชลิํ ปณาเมตฺวา เอกมนฺตํ นิสีทิํสุ. อปฺเปกจฺเจ ภควโต สนฺติเก นามโคตฺตํ สาเวตฺวา เอกมนฺตํ นิสีทิํสุ. อปฺเปกจฺเจ ตุณฺหีภูตา เอกมนฺตํ นิสีทิํสุ.}

\proseitem{356}{เอกมนฺตํ นิสินฺโน โข สจฺจโก นิคณฺฐปุตฺโต ภควนฺตํ เอตทโวจ “ปุจฺเฉยฺยาหํ ภวนฺตํ โคตมํ กิญฺจิเทว เทสํ, สเจ เม ภวํ โคตโม โอกาสํ กโรติ ปญฺหสฺส เวยฺยากรณายา”ติ. ปุจฺฉ อคฺคิเวสฺสน ยทากงฺขสีติ. กถํ ปน ภวํ โคตโม สาวเก วิเนติ, กถํภาคา จ ปน โภโต โคตมสฺส สาวเกสุ อนุสาสนี พหุลา ปวตฺตตีติ. เอวํ โข อหํ อคฺคิเวสฺสน สาวเก วิเนมิ, เอวํภาคา จ ปน เม สาวเกสุ อนุสาสนี พหุลา ปวตฺตติ “รูปํ ภิกฺขเว อนิจฺจํ, เวทนา อนิจฺจา, สญฺญา อนิจฺจา, สงฺขารา อนิจฺจา, วิญฺญาณํ อนิจฺจํ. รูปํ ภิกฺขเว อนตฺตา, เวทนา อนตฺตา, สญฺญา อนตฺตา, สงฺขารา อนตฺตา, วิญฺญาณํ อนตฺตา. สพฺเพ สงฺขารา อนิจฺจา, สพฺเพ ธมฺมา อนตฺตา”ติ. เอวํ โข อหํ อคฺคิเวสฺสน สาวเก วิเนมิ, เอวํภาคา จ ปน เม สาวเกสุ อนุสาสนี พหุลา ปวตฺตตีติ. อุปมา มํ โภ โคตม ปฏิภาตีติ. “ปฏิภาตุ ตํ อคฺคิเวสฺสนา”ติ ภควา อโวจ.}

\prose{เสยฺยถาปิ โภ โคตม เย เกจิเม พีชคาม\-ภูตคามา วุทฺธิํ วิรูฬฺหิํ เวปุลฺลํ อาปชฺชนฺติ, สพฺเพ เต ปถวิํ นิสฺสาย ปถวิยํ ปติฏฺฐาย, เอวเมเต พีชคาม\-ภูตคามา วุทฺธิํ วิรูฬฺหิํ เวปุลฺลํ อาปชฺชนฺติ. เสยฺยถาปิ วา ปน โภ โคตม เย เกจิเม พลกรณียา กมฺมนฺตา กรียนฺติ, สพฺเพ เต ปถวิํ นิสฺสาย ปถวิยํ ปติฏฺฐาย, เอวเมเต พลกรณียา กมฺมนฺตา กรียนฺติ. เอวเมว โข โภ โคตม รูปตฺตายํ ปุริสปุคฺคโล รูเป ปติฏฺฐาย ปุญฺญํ วา อปุญฺญํ วา ปสวติ, เวทนตฺตายํ ปุริสปุคฺคโล เวทนายํ ปติฏฺฐาย ปุญฺญํ วา อปุญฺญํ วา ปสวติ, สญฺญตฺตายํ ปุริสปุคฺคโล สญฺญายํ ปติฏฺฐาย ปุญฺญํ วา อปุญฺญํ วา ปสวติ, สงฺขารตฺตายํ ปุริสปุคฺคโล สงฺขาเรสุ ปติฏฺฐาย ปุญฺญํ วา อปุญฺญํ วา ปสวติ, วิญฺญาณตฺตายํ ปุริสปุคฺคโล วิญฺญาเณ ปติฏฺฐาย ปุญฺญํ วา อปุญฺญํ วา ปสวตีติ.}

\prose{นนุ ตฺวํ อคฺคิเวสฺสน เอวํ วเทสิ, รูปํ เม อตฺตา, เวทนา เม อตฺตา, สญฺญา เม อตฺตา, สงฺขารา เม อตฺตา, วิญฺญาณํ เม อตฺตาติ. อหํ หิ โภ \csromanfolio{293}โคตม เอวํ วทามิ “รูปํ เม อตฺตา, เวทนา เม อตฺตา, สญฺญา เม อตฺตา, สงฺขารา เม อตฺตา, วิญฺญาณํ เม อตฺตา”ติ, อยญฺจ มหตี ชนตาติ.}

\prose{กิํ หิ เต อคฺคิเวสฺสน มหตี ชนตา กริสฺสติ, อิงฺฆ ตฺวํ อคฺคิเวสฺสน สกญฺเญว วาทํ นิพฺเพเฐหีติ. อหํ หิ โภ โคตม เอวํ วทามิ “รูปํ เม อตฺตา, เวทนา เม อตฺตา, สญฺญา เม อตฺตา, สงฺขารา เม อตฺตา, วิญฺญาณํ เม อตฺตา”ติ.}

\proseitem{357}{เตน หิ อคฺคิเวสฺสน ตญฺเญเวตฺถ ปฏิปุจฺฉิสฺสามิ, ยถา เต ขเมยฺย, ตถา นํ\csromansharedfootnote{52e041efbe88099f}{ตถา ตํ (ก)} พฺยากเรยฺยาสิ, ตํ กิํ มญฺญสิ อคฺคิเวสฺสน, วตฺเตยฺย รญฺโญ ขตฺติยสฺส มุทฺธา\-วสิตฺตสฺส สกสฺมิํ วิชิเต วโส ฆาเตตายํ วา ฆาเตตุํ ชาเปตายํ วา ชาเปตุํ ปพฺพาเชตายํ วา ปพฺพาเชตุํ, เสยฺยถาปิ รญฺโญ ปเสนทิสฺส โกสลสฺส, เสยฺยถาปิ วา ปน รญฺโญ มาคธสฺส อชาตสตฺตุสฺส เวเทหิปุตฺตสฺสาติ. วตฺเตยฺย โภ โคตม รญฺโญ ขตฺติยสฺส มุทฺธา\-วสิตฺตสฺส สกสฺมิํ วิชิเต วโส ฆาเตตายํ วา ฆาเตตุํ ชาเปตายํ วา ชาเปตุํ ปพฺพาเชตายํ วา ปพฺพาเชตุํ, เสยฺยถาปิ รญฺโญ ปเสนทิสฺส โกสลสฺส, เสยฺยถาปิ วา ปน รญฺโญ มาคธสฺส อชาตสตฺตุสฺส เวเทหิปุตฺตสฺส. อิเมสํ ปิ หิ โภ โคตม สํฆานํ คณานํ, เสยฺยถิทํ, วชฺชีนํ มลฺลานํ วตฺตติ สกสฺมิํ วิชิเต วโส ฆาเตตายํ วา ฆาเตตุํ ชาเปตายํ วา ชาเปตุํ ปพฺพาเชตายํ วา ปพฺพาเชตุํ, กิํ ปน รญฺโญ ขตฺติยสฺส มุทฺธา\-วสิตฺตสฺส, เสยฺยถาปิ รญฺโญ ปเสนทิสฺส โกสลสฺส, เสยฺยถาปิ วา ปน รญฺโญ มาคธสฺส อชาตสตฺตุสฺส เวเทหิปุตฺตสฺส, วตฺเตยฺย โภ โคตม วตฺติตุํ จ มรหตีติ.}

\prose{ตํ กิํ มญฺญสิ อคฺคิเวสฺสน, ยํ ตฺวํ เอวํ วเทสิ “รูปํ เม อตฺตา”ติ, วตฺตติ เต ตสฺมิํ รูเป วโส “เอวํ เม รูปํ โหตุ, เอวํ เม รูปํ มา อโหสี”ติ. เอวํ วุตฺเต สจฺจโก นิคณฺฐปุตฺโต ตุณฺหี อโหสิ. ทุติยมฺปิ โข ภควา สจฺจกํ นิคณฺฐปุตฺตํ เอตทโวจ, ตํ กิํ มญฺญสิ อคฺคิเวสฺสน, ยํ ตฺวํ เอวํ วเทสิ “รูปํ เม อตฺตา”ติ, วตฺตติ เต ตสฺมิํ รูเป วโส “เอวํ เม รูปํ โหตุ, เอวํ เม รูปํ มา อโหสี”ติ. \csromanfolio{294}ทุติยมฺปิ โข สจฺจโก นิคณฺฐปุตฺโต ตุณฺหี อโหสิ. อถ โข ภควา สจฺจกํ นิคณฺฐปุตฺตํ เอตทโวจ “พฺยากโรหิ ทานิ อคฺคิเวสฺสน น ทานิ เต ตุณฺหีภาวสฺส กาโล, โย โกจิ อคฺคิเวสฺสน ตถาคเตน ยาวตติยํ สหธมฺมิกํ ปญฺหํ ปุฏฺโฐ น พฺยากโรติ. เอตฺเถวสฺส สตฺตธา มุทฺธา ผลตี”ติ.}

\prose{เตน โข ปน สมเยน วชิรปาณิ ยกฺโข อายสํ วชิรํ อาทาย อาทิตฺตํ สมฺปชฺชลิตํ สโชติภูตํ สจฺจกสฺส นิคณฺฐ\-ปุตฺตสฺส อุปริเวหาสํ ฐิโต โหติ “สจายํ สจฺจโก นิคณฺฐปุตฺโต ภควตา ยาวตติยํ สหธมฺมิกํ ปญฺหํ ปุฏฺโฐ น พฺยากริสฺสติ, เอตฺเถวสฺส สตฺตธา มุทฺธํ ผาเลสฺสามี”ติ. ตํ โข ปน วชิรปาณิํ ยกฺขํ ภควา เจว ปสฺสติ, สจฺจโก จ นิคณฺฐปุตฺโต. อถ โข สจฺจโก นิคณฺฐปุตฺโต ภีโต สํวิคฺโค โลมหฏฺฐชาโต ภควนฺตํเยว ตาณํ คเวสี ภควนฺตํเยว เลณํ คเวสี ภควนฺตํเยว สรณํ คเวสี ภควนฺตํ เอตทโวจ “ปุจฺฉตุ มํ ภวํ โคตโม พฺยากริสฺสามี”ติ.}

\proseitem{358}{ตํ กิํ มญฺญสิ อคฺคิเวสฺสน, ยํ ตฺวํ เอวํ วเทสิ “รูปํ เม อตฺตา”ติ, วตฺตติ เต ตสฺมิํ รูเป วโส “เอวํ เม รูปํ โหตุ, เอวํ เม รูปํ มา อโหสี”ติ. โน หิทํ โภ โคตม.}

\prose{มนสิ กโรหิ อคฺคิเวสฺสน, มนสิ กริตฺวา โข อคฺคิเวสฺสน พฺยากโรหิ, น โข เต สนฺธิยติ ปุริเมน วา ปจฺฉิมํ ปจฺฉิเมน วา ปุริมํ. ตํ กิํ มญฺญสิ อคฺคิเวสฺสน, ยํ ตฺวํ เอวํ วเทสิ “เวทนา เม อตฺตา”ติ, วตฺตติ เต ติสฺสํ เวทนายํ\csromansharedfootnote{d90b3958e9d0bff0}{ตายํ เวทนายํ (สี, สฺยา)} วโส “เอวํ เม เวทนา โหตุ, เอวํ เม เวทนา มา อโหสี”ติ. โน หิทํ โภ โคตม.}

\prose{มนสิ กโรหิ อคฺคิเวสฺสน, มนสิ กริตฺวา โข อคฺคิเวสฺสน พฺยากโรหิ, น โข เต สนฺธิยติ ปุริเมน วา ปจฺฉิมํ ปจฺฉิเมน วา ปุริมํ. ตํ กิํ มญฺญสิ อคฺคิเวสฺสน, ยํ ตฺวํ เอวํ วเทสิ “สญฺญา เม อตฺตา”ติ, วตฺตติ เต ติสฺสํ สญฺญายํ วโส “เอวํ เม สญฺญา โหตุ, เอวํ เม สญฺญา มา อโหสี”ติ. โน หิทํ โภ โคตม.}

\prose{\csromanfolio{295}มนสิ กโรหิ อคฺคิเวสฺสน, มนสิ กริตฺวา โข อคฺคิเวสฺสน พฺยากโรหิ, น โข เต สนฺธิยติ ปุริเมน วา ปจฺฉิมํ ปจฺฉิเมน วา ปุริมํ. ตํ กิํ มญฺญสิ อคฺคิเวสฺสน, ยํ ตฺวํ เอวํ วเทสิ “สงฺขารา เม อตฺตา”ติ, วตฺตติ เต เตสุ สงฺขาเรสุ วโส “เอวํ เม สงฺขารา โหนฺตุ, เอวํ เม สงฺขารา มา อเหสุนฺ”ติ. โน หิทํ โภ โคตม.}

\prose{มนสิ กโรหิ อคฺคิเวสฺสน, มนสิ กริตฺวา โข อคฺคิเวสฺสน พฺยากโรหิ, น โข เต สนฺธิยติ ปุริเมน วา ปจฺฉิมํ ปจฺฉิเมน วา ปุริมํ. ตํ กิํ มญฺญสิ อคฺคิเวสฺสน, ยํ ตฺวํ เอวํ วเทสิ “วิญฺญาณํ เม อตฺตา”ติ, วตฺตติ เต ตสฺมิํ วิญฺญาเณ วโส “เอวํ เม วิญฺญาณํ โหตุ, เอวํ เม วิญฺญาณํ มา อโหสี”ติ. โน หิทํ โภ โคตม.}

\prose{มนสิ กโรหิ อคฺคิเวสฺสน, มนสิ กริตฺวา โข อคฺคิเวสฺสน พฺยากโรหิ, น โข เต สนฺธิยติ ปุริเมน วา ปจฺฉิมํ ปจฺฉิเมน วา ปุริมํ. ตํ กิํ มญฺญสิ อคฺคิเวสฺสน, รูปํ นิจฺจํ วา อนิจฺจํ วาติ. อนิจฺจํ โภ โคตม. ยํ ปนานิจฺจํ, ทุกฺขํ วา ตํ สุขํ วาติ. ทุกฺขํ โภ โคตม. ยํ ปนานิจฺจํ ทุกฺขํ วิปริณาม\-ธมฺมํ, กลฺลํ นุ ตํ สมนุปสฺสิตุํ “เอตํ มม, เอโสหมสฺมิ, เอโส เม อตฺตา”ติ. โน หิทํ โภ โคตม.}

\prose{ตํ กิํ มญฺญสิ อคฺคิเวสฺสน, เวทนา ฯเปฯ สญฺญา ฯเปฯ สงฺขารา ฯเปฯ. ตํ กิํ มญฺญสิ อคฺคิเวสฺสน, วิญฺญาณํ นิจฺจํ วา อนิจฺจํ วาติ. อนิจฺจํ โภ โคตม. ยํ ปนานิจฺจํ, ทุกฺขํ วา ตํ สุขํ วาติ. ทุกฺขํ โภ โคตม. ยํ ปนานิจฺจํ ทุกฺขํ วิปริณาม\-ธมฺมํ, กลฺลํ นุ ตํ สมนุปสฺสิตุํ “เอตํ มม, เอโสหมสฺมิ, เอโส เม อตฺตา”ติ. โน หิทํ โภ โคตม.}

\prose{ตํ กิํ มญฺญสิ อคฺคิเวสฺสน, โย นุ โข ทุกฺขํ อลฺลีโน ทุกฺขํ อุปคโต ทุกฺขํ อชฺโฌสิโต ทุกฺขํ “เอตํ มม, เอโสหมสฺมิ, เอโส เม อตฺตา”ติ สมนุปสฺสติ, อปิ นุ โข โส สามํ วา ทุกฺขํ ปริชาเนยฺย, ทุกฺขํ วา ปริกฺเขเปตฺวา วิหเรยฺยาติ. กิํ หิ สิยา โภ โคตม, โน หิทํ โภ โคตมาติ.}

\prose{ตํ กิํ มญฺญสิ อคฺคิเวสฺสน, นนุ ตฺวํ เอวํ สนฺเต ทุกฺขํ อลฺลีโน ทุกฺขํ อุปคโต ทุกฺขํ อชฺโฌสิโต, ทุกฺขํ “เอตํ มม, เอโสหมสฺมิ, เอโส เม อตฺตา”ติ สมนุปสฺสสีติ. กิํ หิ โน สิยา โภ โคตม, เอวเมตํ โภ โคตมาติ.}

\proseitem{359}{\csromanfolio{296}เสยฺยถาปิ อคฺคิเวสฺสน ปุริโส สารตฺถิโก สารคเวสี สารปริเยสนํ จรมาโน ติณฺหํ กุฐาริํ\csromansharedfootnote{a214774892a5b525}{กุธาริํ (สฺยา, กํ, ก)} อาทาย วนํ ปวิเสยฺย. โส ตตฺถ ปสฺเสยฺย มหนฺตํ กทลิ\-กฺขนฺธํ อุชุํ นวํ อกุกฺกุกชาตํ\csromansharedfootnote{e6046496017e39a8}{อกุกฺกุฏชาตํ (สฺยา, กํ)}. ตเมนํ มูเล ฉินฺเทยฺย มูเล เฉตฺวา อคฺเค ฉินฺเทยฺย อคฺเค เฉตฺวา ปตฺตวฏฺฏิํ วินิพฺภุเชยฺย\csromansharedfootnote{8ce953b2b9e8245d}{วินิพฺภุชฺเชยฺย (ก)}. โส ตตฺถ ปตฺตวฏฺฏิํ วินิพฺภุชนฺโต เผคฺคุํปิ นาธิคจฺเฉยฺย, กุโต สารํ. เอวเมว โข ตฺวํ อคฺคิเวสฺสน มยา สกสฺมิํ วาเท สมนุ\-ยุญฺชิย\-มาโน สมนุ\-คาหิย\-มาโน สมนุ\-ภาสิย\-มาโน ริตฺโต ตุจฺโฉ อปรทฺโธ. ภาสิตา โข ปน เต เอสา อคฺคิเวสฺสน เวสาลิยํ ปริสติ วาจา “นาหํ ตํ ปสฺสามิ สมณํ วา พฺราหฺมณํ วา สงฺฆิํ คณิํ คณาจริยํ, อปิ อรหนฺตํ สมฺมา\-สมฺพุทฺธํ ปฏิชานมานํ, โย มยา วาเทน วาทํ สมารทฺโธ น สงฺกมฺเปยฺย น สมฺปกมฺเปยฺย น สมฺปเวเธยฺย, ยสฺส น กจฺเฉหิ เสทา มุจฺเจยฺยุํ. ถูณํ เจปาหํ อเจตนํ วาเทน วาทํ สมารเภยฺยํ, สาปิ มยา วาเทน วาทํ สมารทฺธา สงฺกมฺเปยฺย สมฺปกมฺเปยฺย สมฺปเวเธยฺย. โก ปน วาโท มนุสฺสภูตสฺสา”ติ. ตุยฺหํ โข ปน อคฺคิเวสฺสน อปฺเปกจฺจานิ เสทผุสิตานิ นลาฏา มุตฺตานิ อุตฺตราสงฺคํ วินิภินฺทิตฺวา ภูมิยํ ปติฏฺฐิตานิ, มยฺหํ โข ปน อคฺคิเวสฺสน นตฺถิ เอตรหิ กายสฺมิํ เสโทติ. อิติ ภควา ตสฺมิํ\csromansharedfootnote{f4f4354d27d1851e}{ตสฺสํ (?)} ปริสติ สุวณฺณวณฺณํ กายํ วิวริ. เอวํ วุตฺเต สจฺจโก นิคณฺฐปุตฺโต ตุณฺหีภูโต มงฺกุภูโต ปตฺตกฺขนฺโธ อโธมุโข ปชฺฌายนฺโต อปฺปฏิภาโน นิสีทิ.}

\proseitem{360}{อถ โข ทุมฺมุโข ลิจฺฉวิปุตฺโต สจฺจกํ นิคณฺฐปุตฺตํ ตุณฺหีภูตํ มงฺกุภูตํ ปตฺตกฺขนฺธํ อโธมุขํ ปชฺฌายนฺตํ อปฺปฏิภานํ วิทิตฺวา ภควนฺตํ เอตทโวจ “อุปมา มํ ภควา ปฏิภาตี”ติ. “ปฏิภาตุ ตํ ทุมฺมุขา”ติ ภควา อโวจ. เสยฺยถาปิ ภนฺเต คามสฺส วา นิคมสฺส วา อวิทูเร โปกฺขรณี, ตตฺราสฺส กกฺกฏโก. อถ โข ภนฺเต สมฺพหุลา กุมารกา วา กุมาริกา วา ตมฺหา คามา วา นิคมา วา นิกฺขมิตฺวา เยน สา โปกฺขรณี เตนุปสงฺกเมยฺยุํ, อุปสงฺกมิตฺวา ตํ โปกฺขรณิํ โอคาเหตฺวา ตํ กกฺกฏกํ อุทกา อุทฺธริตฺวา ถเล ปติฏฺฐาเปยฺยุํ. ยญฺญเทว หิ โส ภนฺเต กกฺกฏโก อฬํ อภินินฺนาเมยฺย. ตํ ตเทว เต กุมารกา วา กุมาริกา วา กฏฺเฐน วา กถเลน วา สญฺฉินฺเทยฺยุํ \csromanfolio{297}สมฺภญฺเชยฺยุํ สมฺปลิภญฺเชยฺยุํ. เอวํ หิ โส ภนฺเต กกฺกฏโก สพฺเพหิ อเฬหิ สญฺฉินฺเนหิ สมฺภคฺเคหิ สมฺปลิภคฺเคหิ อภพฺโพ ตํ โปกฺขรณิํ ปุน โอตริตุํ เสยฺยถาปิ ปุพฺเพ. เอวเมว โข ภนฺเต ยานิ สจฺจกสฺส นิคณฺฐ\-ปุตฺตสฺส วิสูกายิตานิ วิเสวิตานิ วิปฺผนฺทิตานิ, ตานิปิ สพฺพานิ\csromansharedfootnote{54869f1c2eb9d7b0}{วิปฺผนฺทิตานิ กานิจิ กานิจิ ตานิ (สี, สฺยา, กํ, อิ)} ภควตา สญฺฉินฺนานิ สมฺภคฺคานิ สมฺปลิภคฺคานิ. อภพฺโพ จ ทานิ ภนฺเต สจฺจโก นิคณฺฐปุตฺโต ปุน ภควนฺตํ อุปสงฺกมิตุํ ยทิทํ วาทาธิปฺปาโยติ. เอวํ วุตฺเต สจฺจโก นิคณฺฐปุตฺโต ทุมฺมุขํ ลิจฺฉวิปุตฺตํ เอตทโวจ, อาคเมหิ ตฺวํ ทุมฺมุข, อาคเมหิ ตฺวํ ทุมฺมุข, ( )\csromansharedfootnote{c8211f4c7f1b7308}{(มุขโรสิ ตฺวํ ทุมฺมุข) (สฺยา, กํ)} น มยํ ตยา สทฺธิํ มนฺเตม, อิธ มยํ โภตา โคตเมน สทฺธิํ มนฺเตม.}

\proseitem{361}{ติฏฺฐเตสา โภ โคตม อมฺหากญฺเจว อญฺเญสญฺจ ปุถุสมณ\-พฺราหฺมณานํ วาจา, วิลาปํ วิลปิตํ มญฺเญ. กิตฺตาวตา จ นุ โข โภโต โคตมสฺส สาวโก สาสนกโร โหติ โอวาทปติกโร, ติณฺณ\-วิจิกิจฺโฉ วิคต\-กถํกโถ เวสารชฺช\-ปฺปตฺโต อปรปฺปจฺจโย สตฺถุสาสเน วิหรตีติ. อิธ อคฺคิเวสฺสน มม สาวโก ยํ กิญฺจิ รูปํ อตีตา\-นาคต\-ปจฺจุปฺปนฺนํ อชฺฌตฺตํ วา พหิทฺธา วา โอฬาริกํ วา สุขุมํ วา หีนํ วา ปณีตํ วา ยํ ทูเร สนฺติเก วา, สพฺพํ รูปํ “เนตํ มม, เนโสหมสฺมิ, น เมโส อตฺตา”ติ เอวเมตํ ยถาภูตํ สมฺมปฺปญฺญาย ปสฺสติ. ยา กาจิ เวทนา ฯเปฯ. ยา กาจิ สญฺญา ฯเปฯ. เย เกจิ สงฺขารา ฯเปฯ. ยํ กิญฺจิ วิญฺญาณํ อตีตา\-นาคต\-ปจฺจุปฺปนฺนํ อชฺฌตฺตํ วา พหิทฺธา วา โอฬาริกํ วา สุขุมํ วา หีนํ วา ปณีตํ วา ยํ ทูเร สนฺติเก วา, สพฺพํ วิญฺญาณํ “เนตํ มม, เนโสหมสฺมิ, น เมโส อตฺตา”ติ เอวเมตํ ยถาภูตํ สมฺมปฺปญฺญาย ปสฺสติ. เอตฺตาวตา โข อคฺคิเวสฺสน มม สาวโก สาสนกโร โหติ โอวาทปติกโร, ติณฺณ\-วิจิกิจฺโฉ วิคต\-กถํกโถ เวสารชฺช\-ปฺปตฺโต อปรปฺปจฺจโย สตฺถุสาสเน วิหรตีติ.}

\prose{กิตฺตาวตา ปน โภ โคตม ภิกฺขุ อรหํ โหติ ขีณาสโว วุสิตวา กตกรณีโย โอหิตภาโร อนุปฺปตฺต\-สทตฺโถ ปริกฺขีณ\-ภวสํโยชโน สมฺมทญฺญา วิมุตฺโตติ. อิธ อคฺคิเวสฺสน ภิกฺขุ ยํ กิญฺจิ รูปํ อตีตา\-นาคต\-ปจฺจุปฺปนฺนํ อชฺฌตฺตํ วา พหิทฺธา วา โอฬาริกํ วา สุขุมํ วา หีนํ วา ปณีตํ วา ยํ ทูเร สนฺติเก วา, สพฺพํ รูปํ “เนตํ มม, เนโสหมสฺมิ, \csromanfolio{298}น เมโส อตฺตา”ติ เอวเมตํ ยถาภูตํ สมฺมปฺปญฺญาย ทิสฺวา อนุปาทา วิมุตฺโต โหติ. ยา กาจิ เวทนา ฯเปฯ. ยา กาจิ สญฺญา ฯเปฯ. เย เกจิ สงฺขารา ฯเปฯ. ยํ กิญฺจิ วิญฺญาณํ อตีตา\-นาคต\-ปจฺจุปฺปนฺนํ อชฺฌตฺตํ วา พหิทฺธา วา โอฬาริกํ วา สุขุมํ วา หีนํ วา ปณีตํ วา ยํ ทูเร สนฺติเก วา, สพฺพํ วิญฺญาณํ “เนตํ มม, เนโสหมสฺมิ, น เมโส อตฺตา”ติ เอวเมตํ ยถาภูตํ สมฺมปฺปญฺญาย ทิสฺวา อนุปาทา วิมุตฺโต โหติ. เอตฺตาวตา โข อคฺคิเวสฺสน ภิกฺขุ อรหํ โหติ ขีณาสโว วุสิตวา กตกรณีโย โอหิตภาโร อนุปฺปตฺต\-สทตฺโถ ปริกฺขีณ\-ภวสํโยชโน สมฺมทญฺญา วิมุตฺโต. เอวํ วิมุตฺตจิตฺโต โข อคฺคิเวสฺสน ภิกฺขุ ตีหิ อนุตฺตริเยหิ สมนฺนาคโต โหติ ทสฺสนา\-นุตฺตริเยน ปฏิปทา\-นุตฺตริเยน วิมุตฺตา\-นุตฺตริเยน. เอวํ วิมุตฺตจิตฺโต โข อคฺคิเวสฺสน ภิกฺขุ ตถาคตญฺเญว สกฺกโรติ ครุํ กโรติ มาเนติ ปูเชติ “พุทฺโธ โส ภควา โพธาย ธมฺมํ เทเสติ, ทนฺโต โส ภควา ทมถาย ธมฺมํ เทเสติ, สนฺโต โส ภควา สมถาย ธมฺมํ เทเสติ, ติณฺโณ โส ภควา ตรณาย ธมฺมํ เทเสติ, ปรินิพฺพุโต โส ภควา ปรินิพฺพานาย ธมฺมํ เทเสตี”ติ.}

\proseitem{362}{เอวํ วุตฺเต สจฺจโก นิคณฺฐปุตฺโต ภควนฺตํ เอตทโวจ “มยเมว โภ โคตม ธํสี, มยํ ปคพฺพา, เย มยํ ภวนฺตํ โคตมํ วาเทน วาทํ อาสาเทตพฺพํ อมญฺญิมฺห. สิยา หิ โภ โคตม หตฺถิํ ปภินฺนํ อาสชฺช ปุริสสฺส โสตฺถิภาโว, น เตฺวว ภวนฺตํ โคตมํ อาสชฺช สิยา ปุริสสฺส โสตฺถิภาโว. สิยา หิ โภ โคตม ปชฺชลิตํ\csromansharedfootnote{19f85a9b99e0de51}{ชลนฺตํ (สี, อิ)} อคฺคิกฺขนฺธํ อาสชฺช ปุริสสฺส โสตฺถิภาโว, น เตฺวว ภวนฺตํ โคตมํ อาสชฺช สิยา ปุริสสฺส โสตฺถิภาโว. สิยา หิ โภ โคตม อาสีวิสํ โฆรวิสํ อาสชฺช ปุริสสฺส โสตฺถิภาโว, น เตฺวว ภวนฺตํ โคตมํ อาสชฺช สิยา ปุริสสฺส โสตฺถิภาโว. มยเมว โภ โคตม ธํสี, มยํ ปคพฺพา, เย มยํ ภวนฺตํ โคตมํ วาเทน วาทํ อาสาเทตพฺพํ อมญฺญิมฺห. อธิวาเสตุ\csromansharedfootnote{cf9e9d704c617278}{อธิวาเสตุ จ (อิ, ก)} เม ภวํ โคตโม สฺวาตนาย ภตฺตํ สทฺธิํ ภิกฺขุสํเฆนา”ติ. อธิวาเสสิ ภควา ตุณฺหีภาเวน.}

\proseitemwithcloserruled{363}{อถ โข สจฺจโก นิคณฺฐปุตฺโต ภควโต อธิวาสนํ วิทิตฺวา เต ลิจฺฉวี อามนฺเตสิ “สุณนฺตุ เม โภนฺโต ลิจฺฉวี, \csromanfolio{299}สมโณ เม โคตโม นิมนฺติโต สฺวาตนาย สทฺธิํ ภิกฺขุ\-สํเฆน, เตน เม อภิหเรยฺยาถ ยมสฺส ปติรูปํ มญฺเญยฺยาถา”ติ. อถ โข เต ลิจฺฉวี ตสฺสา รตฺติยา อจฺจเยน สจฺจกสฺส นิคณฺฐ\-ปุตฺตสฺส ปญฺจมตฺตานิ ถาลิปาก\-สตานิ ภตฺตาภิหารํ อภิหริํสุ. อถ โข นิคณฺฐปุตฺโต สเก อาราเม ปณีตํ ขาทนียํ โภชนียํ ปฏิยาทาเปตฺวา ภควโต กาลํ อาโรจาเปสิ “กาโล โภ โคตม นิฏฺฐิตํ ภตฺตนฺ”ติ. อถ โข ภควา ปุพฺพณฺหสมยํ นิวาเสตฺวา ปตฺต\-จีวรมา\-ทาย เยน สจฺจกสฺส นิคณฺฐ\-ปุตฺตสฺส อาราโม เตนุปสงฺกมิ, อุปสงฺกมิตฺวา ปญฺญตฺเต อาสเน นิสีทิ สทฺธิํ ภิกฺขุ\-สํเฆน. อถ โข สจฺจโก นิคณฺฐปุตฺโต พุทฺธ\-ปฺปมุขํ ภิกฺขุสํฆํ ปณีเตน ขาทนีเยน โภชนีเยน สหตฺถา สนฺตปฺเปสิ สมฺปวาเรสิ. อถ โข สจฺจโก นิคณฺฐปุตฺโต ภควนฺตํ ภุตฺตาวิํ โอนีต\-ปตฺต\-ปาณิํ อญฺญตรํ นีจํ อาสนํ คเหตฺวา เอกมนฺตํ นิสีทิ, เอกมนฺตํ นิสินฺโน โข สจฺจโก นิคณฺฐปุตฺโต ภควนฺตํ เอตทโวจ “ยมิทํ โภ โคตม ทาเน ปุญฺญญฺจ ปุญฺญมหี จ, ตํ ทายกานํ สุขาย โหตู”ติ. ยํ โข อคฺคิเวสฺสน ตาทิสํ ทกฺขิเณยฺยํ อาคมฺม อวีตราคํ อวีตโทสํ อวีตโมหํ, ตํ ทายกานํ ภวิสฺสติ. ยํ โข อคฺคิเวสฺสน มาทิสํ ทกฺขิเณยฺยํ อาคมฺม วีตราคํ วีตโทสํ วีตโมหํ, ตํ ตุยฺหํ ภวิสฺสตีติ.}{จูฬสจฺจก\-สุตฺตํ นิฏฺฐิตํ ปญฺจมํ.}
\csromanmatikaanchor{mtk.68}
\csromantocmarkref{subsection}{6. มหาสจฺจกสุตฺต}{mtk.68}
\bookmark[dest={mtk.68},level=2]{6. มหาสจฺจกสุตฺต}
\csromanheader{2}{6. มหา\-สจฺจก\-สุตฺต}
\proseitem{364}{เอวํ เม สุตํ\csromandash{}เอกํ สมยํ ภควา เวสาลิยํ วิหรติ มหาวเน กูฏาคาร\-สาลายํ. เตน โข ปน สมเยน ภควา ปุพฺพณฺหสมยํ สุนิวตฺโถ โหติ ปตฺต\-จีวรมา\-ทาย เวสาลิํ ปิณฺฑาย ปวิสิตุกาโม\csromansharedfootnote{4073d11db399f221}{ปุพฺพณฺหสมยํ นิวาเสตฺวา ปตฺตจีวรมาทาย ... ปวิสิตุกาโม โหติ (สี)}. อถ โข สจฺจโก นิคณฺฐปุตฺโต ชงฺฆาวิหารํ อนุจงฺกมมาโน อนุวิจรมาโน เยน มหาวนํ กูฏาคารสาลา เตนุปสงฺกมิ. อทฺทสา โข อายสฺมา อานนฺโท สจฺจกํ นิคณฺฐปุตฺตํ ทูรโตว อาคจฺฉนฺตํ, ทิสฺวาน ภควนฺตํ เอตทโวจ “อยํ ภนฺเต \csromanfolio{300}สจฺจโก นิคณฺฐปุตฺโต อาคจฺฉติ ภสฺส\-ปฺปวาทโก ปณฺฑิตวาโท สาธุสมฺมโต พหุชนสฺส. เอโส โข ภนฺเต อวณฺณกาโม พุทฺธสฺส, อวณฺณกาโม ธมฺมสฺส, อวณฺณกาโม สํฆสฺส. สาธุ ภนฺเต ภควา มุหุตฺตํ นิสีทตุ อนุกมฺปํ อุปาทายา”ติ. นิสีทิ ภควา ปญฺญตฺเต อาสเน. อถ โข สจฺจโก นิคณฺฐปุตฺโต เยน ภควา เตนุปสงฺกมิ, อุปสงฺกมิตฺวา ภควตา สทฺธิํ สมฺโมทิ. สมฺโมทนียํ กถํ สารณียํ วีติสาเรตฺวา เอกมนฺตํ นิสีทิ. เอกมนฺตํ นิสินฺโน โข สจฺจโก นิคณฺฐปุตฺโต ภควนฺตํ เอตทโวจ\csromandash{}}

\proseitem{365}{สนฺติ โภ โคตม เอเก สมณพฺราหฺมณา กาย\-ภาวนา\-นุโยคม\-นุยุตฺตา วิหรนฺติ, โน จิตฺตภาวนํ. ผุสนฺติ หิ เต โภ โคตม สารีริกํ ทุกฺขํ เวทนํ. ภูตปุพฺพํ โภ โคตม สารีริกาย ทุกฺขาย เวทนาย ผุฏฺฐสฺส สโต อูรุกฺขมฺโภปิ นาม ภวิสฺสติ, หทยํปิ นาม ผลิสฺสติ, อุณฺหํปิ โลหิตํ มุขโต อุคฺคมิสฺสติ, อุมฺมาทํปิ ปาปุณิสฺสติ\csromansharedfootnote{7351245f5a224af2}{ปาปุณิสฺสนฺติ (สฺยา, กํ)} จิตฺตกฺเขปํ. ตสฺส โข เอตํ โภ โคตม กายนฺวยํ จิตฺตํ โหติ, กายสฺส วเสน วตฺตติ. ตํ กิสฺส เหตุ, อภาวิตตฺตา จิตฺตสฺส. สนฺติ ปน โภ โคตม เอเก สมณพฺราหฺมณา จิตฺต\-ภาวนา\-นุโยคม\-นุยุตฺตา วิหรนฺติ, โน กายภาวนํ. ผุสนฺติ หิ เต โภ โคตม เจตสิกํ ทุกฺขํ เวทนํ. ภูตปุพฺพํ โภ โคตม เจตสิกาย ทุกฺขาย เวทนาย ผุฏฺฐสฺส สโต อูรุกฺขมฺโภปิ นาม ภวิสฺสติ, หทยํปิ นาม ผลิสฺสติ, อุณฺหํ โลหิตํ มุขโต อุคฺคมิสฺสติ, อุมฺมาทํปิ ปาปุณิสฺสติ จิตฺตกฺเขปํ. ตสฺส โข เอโส โภ โคตม จิตฺตนฺวโย กาโย โหติ, จิตฺตสฺส วเสน วตฺตติ. ตํ กิสฺส เหตุ, อภาวิตตฺตา กายสฺส. ตสฺส มยฺหํ โภ โคตม เอวํ โหติ “อทฺธา โภโต โคตมสฺส สาวกา จิตฺต\-ภาวนา\-นุโยคม\-นุยุตฺตา วิหรนฺติ, โน กายภาวนนฺ”ติ.}

\proseitem{366}{กินฺติ ปน เต อคฺคิเวสฺสน กายภาวนา สุตาติ. เสยฺยถิทํ, นนฺโท วจฺโฉ, กิโส สํกิจฺโจ, มกฺขลิ โคสาโล. เอเตหิ โภ โคตม อเจลกา มุตฺตาจารา หตฺถา\-ปเลขนา, นเอหิภทฺทนฺติกา, นติฏฺฐ\-ภทฺทนฺติกา\csromansharedfootnote{6671718475145780}{นเอหิภทนฺติกา, นติฏฺฐภทนฺติกา (สี, สฺยา, กํ, อิ, ก)}, น อภิหฏํ น อุทฺทิสฺสกตํ น นิมนฺตนํ สาทิยนฺติ, เต น \csromanfolio{301}กุมฺภิมุขา ปฏิคฺคณฺหนฺติ, น กโฬปิมุขา ปฏิคฺคณฺหนฺติ, น เอฬกมนฺตรํ, น ทณฺฑมนฺตรํ, น มุสลมนฺตรํ, น ทฺวินฺนํ ภุญฺชมานานํ, น คพฺภินิยา, น ปายมานาย, น ปุริส\-นฺตร\-คตาย, น สงฺกิตฺตีสุ, น ยตฺถ สา อุปฏฺฐิโต โหติ, น ยตฺถ มกฺขิกา สณฺฑ\-สณฺฑ\-จารินี, น มจฺฉํ, น มํสํ, น สุรํ, น เมรยํ, น ถุโสทกํ ปิวนฺติ, เต เอกาคาริกา วา โหนฺติ เอกาโลปิกา, ทฺวาคาริกา วา โหนฺติ ทฺวาโลปิกา ฯเปฯ สตฺตาคาริกา วา โหนฺติ สตฺตาโลปิกา, เอกิสฺสาปิ ทตฺติยา ยาเปนฺติ, ทฺวีหิปิ ทตฺตีหิ ยาเปนฺติ ฯเปฯ สตฺตหิปิ ทตฺตีหิ ยาเปนฺติ, เอกาหิกมฺปิ อาหารํ อาหาเรนฺติ, ทฺวีหิกมฺปิ อาหารํ อาหาเรนฺติ ฯเปฯ สตฺตาหิกมฺปิ อาหารํ อาหาเรนฺติ. อิติ เอวรูปํ อทฺธมาสิกมฺปิ ปริยายภตฺตโภชนานุโยคมนุยุตฺตา วิหรนฺตีติ.}

\prose{กิํ ปน เต อคฺคิเวสฺสน ตาวตเกเนว ยาเปนฺตีติ. โน หิทํ โภ โคตม. อปฺเปกทา โภ โคตม อุฬารานิ อุฬารานิ ขาทนียานิ ขาทนฺติ, อุฬารานิ อุฬารานิ โภชนานิ ภุญฺชนฺติ, อุฬารานิ อุฬารานิ สายนียานิ สายนฺติ, อุฬารานิ อุฬารานิ ปานานิ ปิวนฺติ. เต อิมํ กายํ พลํ คาเหนฺติ นาม พฺรูเหนฺติ นาม เมเทนฺติ นามาติ.}

\prose{ยํ โข เต อคฺคิเวสฺสน ปุริมํ ปหาย ปจฺฉา อุปจินนฺติ. เอวํ อิมสฺส กายสฺส อาจยาปจโย โหติ. กินฺติ ปน เต อคฺคิเวสฺสน จิตฺตภาวนา สุตาติ. จิตฺตภาวนาย โข สจฺจโก นิคณฺฐปุตฺโต ภควตา ปุฏฺโฐ สมาโน น สมฺปายาสิ.}

\proseitem{367}{อถ โข ภควา สจฺจกํ นิคณฺฐปุตฺตํ เอตทโวจ “ยาปิ โข เต เอสา อคฺคิเวสฺสน ปุริมา กายภาวนา ภาสิตา, สาปิ อริยสฺส วินเย โน ธมฺมิกา กายภาวนา. กายภาวนมฺปิ\csromansharedfootnote{bfeb3f6146c9b532}{กายภาวนํ หิ (สี, อิ, ก)} โข ตฺวํ อคฺคิเวสฺสน น อญฺญาสิ, กุโต ปน ตฺวํ จิตฺตภาวนํ ชานิสฺสสิ. อปิ จ อคฺคิเวสฺสน ยถา อภาวิตกาโย จ โหติ อภาวิตจิตฺโต จ. ภาวิตกาโย จ โหติ ภาวิตจิตฺโต จ. ตํ สุณาหิ, สาธุกํ มนสิ กโรหิ, ภาสิสฺสามี”ติ. “เอวํ โภ”ติ โข สจฺจโก นิคณฺฐปุตฺโต ภควโต ปจฺจสฺโสสิ. ภควา เอตทโวจ\csromandash{}}

\proseitem{368}{\csromanfolio{302}กถญฺจ อคฺคิเวสฺสน อภาวิตกาโย จ โหติ อภาวิตจิตฺโต จ. อิธ อคฺคิเวสฺสน อสฺสุตวโต ปุถุชฺชนสฺส อุปฺปชฺชติ สุขา เวทนา. โส สุขาย เวทนาย ผุฏฺโฐ สมาโน สุขสาราคี จ โหติ, สุขสาราคิตญฺจ อาปชฺชติ. ตสฺส สา สุขา เวทนา นิรุชฺฌติ. สุขาย เวทนาย นิโรธา อุปฺปชฺชติ ทุกฺขา เวทนา. โส ทุกฺขาย เวทนาย ผุฏฺโฐ สมาโน โสจติ กิลมติ ปริเทวติ อุรตฺตาฬิํ กนฺทติ สมฺโมหํ อาปชฺชติ. ตสฺส โข เอสา อคฺคิเวสฺสน อุปฺปนฺนาปิ สุขา เวทนา จิตฺตํ ปริยาทาย ติฏฺฐติ, อภาวิตตฺตา กายสฺส. อุปฺปนฺนาปิ ทุกฺขา เวทนา จิตฺตํ ปริยาทาย ติฏฺฐติ, อภาวิตตฺตา จิตฺตสฺส. ยสฺส กสฺสจิ อคฺคิเวสฺสน เอวํ อุภโตปกฺขํ อุปฺปนฺนาปิ สุขา เวทนา จิตฺตํ ปริยาทาย ติฏฺฐติ, อภาวิตตฺตา กายสฺส. อุปฺปนฺนาปิ ทุกฺขา เวทนา จิตฺตํ ปริยาทาย ติฏฺฐติ, อภาวิตตฺตา จิตฺตสฺส. เอวํ โข อคฺคิเวสฺสน อภาวิตกาโย จ โหติ อภาวิตจิตฺโต จ.}

\proseitem{369}{กถญฺจ อคฺคิเวสฺสน ภาวิตกาโย จ โหติ ภาวิตจิตฺโต จ. อิธ อคฺคิเวสฺสน สุตวโต อริยสาวกสฺส อุปฺปชฺชติ สุขา เวทนา. โส สุขาย เวทนาย ผุฏฺโฐ สมาโน น สุขสาราคี จ โหติ, น สุขสาราคิตญฺจ อาปชฺชติ. ตสฺส สา สุขา เวทนา นิรุชฺฌติ. สุขาย เวทนาย นิโรธา อุปฺปชฺชติ ทุกฺขา เวทนา. โส ทุกฺขาย เวทนาย ผุฏฺโฐ สมาโน น โสจติ น กิลมติ น ปริเทวติ น อุรตฺตาฬิํ กนฺทติ น สมฺโมหํ อาปชฺชติ. ตสฺส โข เอสา อคฺคิเวสฺสน อุปฺปนฺนาปิ สุขา เวทนา จิตฺตํ น ปริยาทาย ติฏฺฐติ, ภาวิตตฺตา กายสฺส. อุปฺปนฺนาปิ ทุกฺขา เวทนา จิตฺตํ น ปริยาทาย ติฏฺฐติ, ภาวิตตฺตา จิตฺตสฺส. ยสฺส กสฺสจิ อคฺคิเวสฺสน เอวํ อุภโตปกฺขํ อุปฺปนฺนาปิ สุขา เวทนา จิตฺตํ น ปริยาทาย ติฏฺฐติ, ภาวิตตฺตา กายสฺส. อุปฺปนฺนาปิ ทุกฺขา เวทนา จิตฺตํ น ปริยาทาย ติฏฺฐติ, ภาวิตตฺตา จิตฺตสฺส. เอวํ โข อคฺคิเวสฺสน ภาวิตกาโย จ โหติ ภาวิตจิตฺโต จาติ.}

\proseitem{370}{เอวํ ปสนฺโน อหํ โภโต โคตมสฺส. ภวํ หิ โคตโม ภาวิตกาโย จ โหติ ภาวิตจิตฺโต จาติ. อทฺธา โข เต อยํ อคฺคิเวสฺสน อาสชฺช อุปนีย วาจา ภาสิตา. อปิ จ เต อหํ \csromanfolio{303}พฺยากริสฺสามิ. ยโต โข อหํ อคฺคิเวสฺสน เกสมสฺสุํ โอหาเรตฺวา กาสายานิ วตฺถานิ อจฺฉาเทตฺวา อคารสฺมา อนคาริยํ ปพฺพชิโต. ตํ วต เม อุปฺปนฺนา วา สุขา เวทนา จิตฺตํ ปริยาทาย ฐสฺสติ, อุปฺปนฺนา วา ทุกฺขา เวทนา จิตฺตํ ปริยาทาย ฐสฺสตีติ เนตํ ฐานํ\csromansharedfootnote{341ba0a44f0add92}{เนตํ โข ฐานํ (สี, อิ)} วิชฺชตีติ.}

\prose{น หิ นูน\csromansharedfootnote{664d4c83cbde1cd3}{น หนูน (สี, สฺยา, กํ, อิ)} โภโต โคตมสฺส อุปฺปชฺชติ ตถารูปา สุขา เวทนา, ยถารูปา อุปฺปนฺนา สุขา เวทนา จิตฺตํ ปริยาทาย ติฏฺเฐยฺย. น หิ นูน โภโต โคตมสฺส อุปฺปชฺชติ ตถารูปา ทุกฺขา เวทนา, ยถารูปา อุปฺปนฺนา ทุกฺขา เวทนา จิตฺตํ ปริยาทาย ติฏฺเฐยฺยาติ.}

\proseitem{371}{กิํ หิ โน สิยา อคฺคิเวสฺสน. อิธ เม อคฺคิเวสฺสน ปุพฺเพว สมฺโพธา อนภิสมฺพุทฺธสฺส โพธิสตฺตสฺเสว สโต เอตทโหสิ “สมฺพาโธ ฆราวาโส รชาปโถ, อพฺโภกาโส ปพฺพชฺชา. นยิทํ สุกรํ อคารํ อชฺฌาวสตา เอกนฺต\-ปริปุณฺณํ เอกนฺต\-ปริสุทฺธํ สงฺขลิขิตํ พฺรหฺมจริยํ จริตุํ. ยํนูนาหํ เกสมสฺสุํ โอหาเรตฺวา กาสายานิ วตฺถานิ อจฺฉาเทตฺวา อคารสฺมา อนคาริยํ ปพฺพเชยฺยนฺ”ติ. โส โข อหํ อคฺคิเวสฺสน อปเรน สมเยน ทหโรว สมาโน สุสุกาฬเกโส ภเทฺรน โยพฺพเนน สมนฺนาคโต ปฐเมน วยสา อกามกานํ มาตาปิตูนํ อสฺสุมุขานํ รุทนฺตานํ เกสมสฺสุํ โอหาเรตฺวา กาสายานิ วตฺถานิ อจฺฉาเทตฺวา อคารสฺมา อนคาริยํ ปพฺพชิํ. โส เอวํ ปพฺพชิโต สมาโน กิํกุสล\-คเวสี อนุตฺตรํ สนฺติวรปทํ ปริเยสมาโน เยน อาฬาโร กาลาโม เตนุปสงฺกมิํ, อุปสงฺกมิตฺวา อาฬารํ กาลามํ เอตทโวจํ “อิจฺฉามหํ อาวุโส กาลาม อิมสฺมิํ ธมฺมวินเย พฺรหฺมจริยํ จริตุนฺ”ติ. เอวํ วุตฺเต อคฺคิเวสฺสน อาฬาโร กาลาโม มํ เอตทโวจ “วิหรตายสฺมา, ตาทิโส อยํ ธมฺโม, ยตฺถ วิญฺญู ปุริโส นจิรสฺเสว สกํ อาจริยกํ สยํ อภิญฺญา สจฺฉิกตฺวา อุปสมฺปชฺช วิหเรยฺยา”ติ. โส โข อหํ อคฺคิเวสฺสน นจิรสฺเสว ขิปฺปเมว ตํ ธมฺมํ ปริยาปุณิํ. โส โข อหํ อคฺคิเวสฺสน ตาวตเกเนว โอฏฺฐ\-ปหต\-มตฺเตน ลปิตลาปนมตฺตน ญาณวาทญฺจ วทามิ เถรวาทญฺจ. ชานามิ ปสฺสามีติ จ ปฏิชานามิ อหญฺเจว อญฺเญ จ. ตสฺส มยฺหํ อคฺคิเวสฺสน เอตทโหสิ “น โข อาฬาโร กาลาโม อิมํ ธมฺมํ เกวลํ สทฺธา\-มตฺตเกน ‘สยํ \csromanfolio{304}อภิญฺญา สจฺฉิกตฺวา อุปสมฺปชฺช วิหรามี’ติ ปเวเทติ, อทฺธา อาฬาโร กาลาโม อิมํ ธมฺมํ ชานํ ปสฺสํ วิหรตี”ติ.}

\prose{อถ ขฺวาหํ อคฺคิเวสฺสน เยน อาฬาโร กาลาโม เตนุปสงฺกมิํ, อุปสงฺกมิตฺวา อาฬารํ กาลามํ เอตทโวจํ “กิตฺตาวตา โน อาวุโส กาลาม อิมํ ธมฺมํ สยํ อภิญฺญา สจฺฉิกตฺวา อุปสมฺปชฺช วิหรามีติ ปเวเทสี”ติ. เอวํ วุตฺเต อคฺคิเวสฺสน อาฬาโร กาลาโม อากิญฺจญฺญา\-ยตนํ ปเวเทสิ. ตสฺส มยฺหํ อคฺคิเวสฺสน เอตทโหสิ “น โข อาฬารสฺเสว กาลามสฺส อตฺถิ สทฺธา, มยฺหํปตฺถิ สทฺธา. น โข อาฬารสฺเสว กาลามสฺส อตฺถิ วีริยํ, มยฺหํปตฺถิ วีริยํ. น โข อาฬารสฺเสว กาลามสฺส อตฺถิ สติ, มยฺหํปตฺถิ สติ. น โข อาฬารสฺเสว กาลามสฺส อตฺถิ สมาธิ, มยฺหํปตฺถิ สมาธิ. น โข อาฬารสฺเสว กาลามสฺส อตฺถิ ปญฺญา, มยฺหํปตฺถิ ปญฺญา. ยํนูนาหํ ‘ยํ ธมฺมํ อาฬาโร กาลาโม สยํ อภิญฺญา สจฺฉิกตฺวา อุปสมฺปชฺช วิหรามี’ติ ปเวเทติ, ตสฺส ธมฺมสฺส สจฺฉิกิริยาย ปทเหยฺยนฺ”ติ. โส โข อหํ อคฺคิเวสฺสน นจิรสฺเสว ขิปฺปเมว ตํ ธมฺมํ สยํ อภิญฺญา สจฺฉิกตฺวา อุปสมฺปชฺช วิหาสิํ.}

\prose{อถ ขฺวาหํ อคฺคิเวสฺสน เยน อาฬาโร กาลาโม เตนุปสงฺกมิํ, อุปสงฺกมิตฺวา อาฬารํ กาลามํ เอตทโวจํ “เอตฺตาวตา โน อาวุโส กาลาม อิมํ ธมฺมํ สยํ อภิญฺญา สจฺฉิกตฺวา อุปสมฺปชฺช ปเวเทสี”ติ. เอตฺตาวตา โข อหํ อาวุโส อิมํ ธมฺมํ สยํ อภิญฺญา สจฺฉิกตฺวา อุปสมฺปชฺช ปเวเทมีติ. อหมฺปิ โข อาวุโส เอตฺตาวตา อิมํ ธมฺมํ สยํ อภิญฺญา สจฺฉิกตฺวา อุปสมฺปชฺช วิหรามีติ. ลาภา โน อาวุโส, สุลทฺธํ โน อาวุโส, เย มยํ อายสฺมนฺตํ ตาทิสํ สพฺรหฺมจาริํ ปสฺสาม. อิติ ยาหํ ธมฺมํ สยํ อภิญฺญา สจฺฉิกตฺวา อุปสมฺปชฺช ปเวเทมิ, ตํ ตฺวํ ธมฺมํ สยํ อภิญฺญา สจฺฉิกตฺวา อุปสมฺปชฺช วิหรสิ. ยํ ตฺวํ ธมฺมํ สยํ อภิญฺญา สจฺฉิกตฺวา อุปสมฺปชฺช วิหรสิ, ตมหํ ธมฺมํ สยํ อภิญฺญา สจฺฉิกตฺวา อุปสมฺปชฺช ปเวเทมิ. อิติ ยาหํ ธมฺมํ ชานามิ, ตํ ตฺวํ ธมฺมํ ชานาสิ. ยํ ตฺวํ ธมฺมํ ชานาสิ, ตมหํ ธมฺมํ ชานามิ. อิติ ยาทิโส อหํ, ตาทิโส ตุวํ. ยาทิโส ตุวํ, ตาทิโส อหํ. เอหิ ทานิ อาวุโส อุโภว สนฺตา อิมํ คณํ ปริหรามาติ. \csromanfolio{305}อิติ โข อคฺคิเวสฺสน อาฬาโร กาลาโม อาจริโย เม สมาโน (อตฺตโน)\csromansharedfootnote{987ccb21c4be8712}{( ) นตฺถิ (สี, อิ)} อนฺเตวาสิํ มํ สมานํ อตฺตนา สมสมํ ฐเปสิ, อุฬาราย จ มํ ปูชาย ปูเชสิ. ตสฺส มยฺหํ อคฺคิเวสฺสน เอตทโหสิ “นายํ ธมฺโม นิพฺพิทาย น วิราคาย น นิโรธาย น อุปสมาย น อภิญฺญาย น สมฺโพธาย น นิพฺพานาย สํวตฺตติ. ยาวเทว อากิญฺจญฺญายตนู\-ปปตฺติยา”ติ. โส โข อหํ อคฺคิเวสฺสน ตํ ธมฺมํ อนลงฺกริตฺวา ตสฺมา ธมฺมา นิพฺพิชฺช อปกฺกมิํ.}

\proseitem{372}{โส โข อหํ อคฺคิเวสฺสน กิํกุสล\-คเวสี อนุตฺตรํ สนฺติวรปทํ ปริเยสมาโน เยน อุทโก รามปุตฺโต เตนุปสงฺกมิํ, อุปสงฺกมิตฺวา อุทกํ รามปุตฺตํ เอตทโวจํ “อิจฺฉามหํ อาวุโส\csromansharedfootnote{04a7f2081132ce39}{ปสฺส ปาสราสิสุตฺเต (221) ปิฏฺเฐ.} อิมสฺมิํ ธมฺมวินเย พฺรหฺมจริยํ จริตุนฺ”ติ. เอวํ วุตฺเต อคฺคิเวสฺสน อุทโก รามปุตฺโต มํ เอตทโวจ “วิหรตายสฺมา, ตาทิโส อยํ ธมฺโม, ยตฺถ วิญฺญู ปุริโส นจิรสฺเสว สกํ อาจริยกํ สยํ อภิญฺญา สจฺฉิกตฺวา อุปสมฺปชฺช วิหเรยฺยา”ติ. โส โข อหํ อคฺคิเวสฺสน นจิรสฺเสว ขิปฺปเมว ตํ ธมฺมํ ปริยาปุณิํ. โส โข อหํ อคฺคิเวสฺสน ตาวตเกเนว โอฏฺฐ\-ปหต\-มตฺเตน ลปิต\-ลาปน\-มตฺเตน ญาณวาทญฺจ วทามิ เถรวาทญฺจ, ชานามิปสฺสามีติ จ ปฏิชานามิ อหญฺเจว อญฺเญ จ. ตสฺส มยฺหํ อคฺคิเวสฺสน เอตทโหสิ “น โข ราโม อิมํ ธมฺมํ เกวลํ สทฺธา\-มตฺตเกน ‘สยํ อภิญฺญา สจฺฉิกตฺวา อุปสมฺปชฺช วิหรามี’ติ ปเวเทสิ, อทฺธา ราโม อิมํ ธมฺมํ ชานํ ปสฺสํ วิหาสี”ติ. อถ ขฺวาหํ อคฺคิเวสฺสน เยน อุทโก รามปุตฺโต เตนุปสงฺกมิํ, อุปสงฺกมิตฺวา อุทกํ รามปุตฺตํ เอตทโวจํ “กิตฺตาวตา โน อาวุโส ราโม อิมํ ธมฺมํ สยํ อภิญฺญา สจฺฉิกตฺวา อุปสมฺปชฺช วิหรามีติ ปเวเทสี”ติ. เอวํ วุตฺเต อคฺคิเวสฺสน อุทโก รามปุตฺโต เนว\-สญฺญา\-นา\-สญฺญา\-ยตนํ ปเวเทสิ. ตสฺส มยฺหํ อคฺคิเวสฺสน เอตทโหสิ “น โข รามสฺเสว อโหสิ สทฺธา, มยฺหํปตฺถิ สทฺธา. น โข รามสฺเสว อโหสิ วีริยํ, มยฺหํปตฺถิ วีริยํ. น โข รามสฺเสว อโหสิ สติ, มยฺหํปตฺถิ สติ. น โข รามสฺเสว อโหสิ สมาธิ, มยฺหํปตฺถิ สมาธิ. น โข รามสฺเสว อโหสิ ปญฺญา, มยฺหํปตฺถิ ปญฺญา. ยํนูนาหํ ‘ยํ ธมฺมํ ราโม สยํ อภิญฺญา สจฺฉิกตฺวา อุปสมฺปชฺช วิหรามี’ติ ปเวเทสิ. ตสฺส ธมฺมสฺส \csromanfolio{306}สจฺฉิกิริยาย ปทเหยฺยนฺ”ติ. โส โข อหํ อคฺคิเวสฺสน นจิรสฺเสว ขิปฺปเมว ตํ ธมฺมํ สยํ อภิญฺญา สจฺฉิกตฺวา อุปสมฺปชฺช วิหาสิํ.}

\prose{อถ ขฺวาหํ อคฺคิเวสฺสน เยน อุทโก รามปุตฺโต เตนุปสงฺกมิํ, อุปสงฺกมิตฺวา อุทกํ รามปุตฺตํ เอตทโวจํ “เอตฺตาวตา โน อาวุโส ราโม อิมํ ธมฺมํ สยํ อภิญฺญา สจฺฉิกตฺวา อุปสมฺปชฺช ปเวเทสี”ติ. เอตฺตาวตา โข อาวุโส ราโม อิมํ ธมฺมํ สยํ อภิญฺญา สจฺฉิกตฺวา อุปสมฺปชฺช ปเวเทสีติ. อหมฺปิ โข อาวุโส เอตฺตาวตา อิมํ ธมฺมํ สยํ อภิญฺญา สจฺฉิกตฺวา อุปสมฺปชฺช วิหรามีติ. ลาภา โน อาวุโส, สุลทฺธํ โน อาวุโส. เย มยํ อายสฺมนฺตํ ตาทิสํ สพฺรหฺมจาริํ ปสฺสาม. อิติ ยํ ธมฺมํ ราโม สยํ อภิญฺญา สจฺฉิกตฺวา อุปสมฺปชฺช ปเวเทสิ, ตํ ตฺวํ ธมฺมํ สยํ อภิญฺญา สจฺฉิกตฺวา อุปสมฺปชฺช วิหรสิ. ยํ ตฺวํ ธมฺมํ สยํ อภิญฺญา สจฺฉิกตฺวา อุปสมฺปชฺช วิหรสิ, ตํ ธมฺมํ ราโม สยํ อภิญฺญา สจฺฉิกตฺวา อุปสมฺปชฺช ปเวเทสิ. อิติ ยํ ธมฺมํ ราโม อภิญฺญาสิ, ตํ ตฺวํ ธมฺมํ ชานาสิ. ยํ ตฺวํ ธมฺมํ ชานาสิ, ตํ ธมฺมํ ราโม อภิญฺญาสิ. อิติ ยาทิโส ราโม อโหสิ, ตาทิโส ตุวํ. ยาทิโส ตุวํ, ตาทิโส ราโม อโหสิ. เอหิ ทานิ อาวุโส ตุวํ อิมํ คณํ ปริหราติ. อิติ โข อคฺคิเวสฺสน อุทโก รามปุตฺโต สพฺรหฺมจารี เม สมาโน อาจริยฏฺฐาเน จ มํ ฐเปสิ, อุฬาราย จ มํ ปูชาย ปูเชสิ. ตสฺส มยฺหํ อคฺคิเวสฺสน เอตทโหสิ “นายํ ธมฺโม นิพฺพิทาย น วิราคาย น นิโรธาย น อุปสมาย น อภิญฺญาย น สมฺโพธาย น นิพฺพานาย สํวตฺตติ. ยาวเทว เนวสญฺญานาสญฺญายตนู\-ปปตฺติยา”ติ. โส โข อหํ อคฺคิเวสฺสน ตํ ธมฺมํ อนลงฺกริตฺวา ตสฺมา ธมฺมา นิพฺพิชฺช อปกฺกมิํ.}

\proseitem{373}{โส โข อหํ อคฺคิเวสฺสน กิํกุสล\-คเวสี อนุตฺตรํ สนฺติวรปทํ ปริเยสมาโน มคเธสุ อนุปุพฺเพน จาริกํ จรมาโน เยน อุรุเวลา เสนานิคโม ตทวสริํ. ตตฺถทฺทสํ รมณียํ ภูมิภาคํ ปาสาทิกญฺจ วนสณฺฑํ นทิญฺจ สนฺทนฺติํ เสตกํ สุปติตฺถํ รมณียํ, สมนฺตา จ โคจรคามํ. ตสฺส มยฺหํ อคฺคิเวสฺสน เอตทโหสิ “รมณีโย วต โภ ภูมิภาโค, ปาสาทิโก จ วนสณฺโฑ, นที จ สนฺทติ เสตกา สุปติตฺถา รมณียา, สมนฺตา จ โคจรคาโม, อลํ วติทํ กุลปุตฺตสฺส \csromanfolio{307}ปธาน\-ตฺถิกสฺส ปธานายาติ. โส โข อหํ อคฺคิเวสฺสน ตตฺเถว นิสีทิํ อลมิทํ ปธานายา”ติ.}

\proseitem{374}{อปิสฺสุมํ อคฺคิเวสฺสน ติสฺโส อุปมา ปฏิภํสุ อนจฺฉริยา ปุพฺเพ อสฺสุตปุพฺพา. เสยฺยถาปิ อคฺคิเวสฺสน อลฺลํ กฏฺฐํ สสฺเนหํ อุทเก นิกฺขิตฺตํ. อถ ปุริโส อาคจฺเฉยฺย อุตฺตรารณิํ อาทาย “อคฺคิํ อภินิพฺพตฺเตสฺสามิ เตโช ปาตุกริสฺสามี”ติ. ตํ กิํ มญฺญสิ อคฺคิเวสฺสน, อปิ นุ โส ปุริโส อมุํ อลฺลํ กฏฺฐํ สสฺเนหํ อุทเก นิกฺขิตฺตํ อุตฺตรารณิํ อาทาย อภิมนฺเถนฺโต อคฺคิํ อภินิพฺพตฺเตยฺย เตโช ปาตุกเรยฺยาติ. โน หิทํ โภ โคตม. ตํ กิสฺส เหตุ, อทุํ หิ โภ โคตม อลฺลํ กฏฺฐํ สสฺเนหํ, ตญฺจ ปน อุทเก นิกฺขิตฺตํ. ยาวเทว จ ปน โส ปุริโส กิลมถสฺส วิฆาตสฺส ภาคี อสฺสาติ. เอวเมว โข อคฺคิเวสฺสน เย หิ เกจิ สมณา วา พฺราหฺมณา วา กาเยน เจว จิตฺเตน จ กาเมหิ อวูปกฏฺฐา วิหรนฺติ, โย จ เนสํ กาเมสุ กามจฺฉนฺโท กามสฺเนโห กามมุจฺฉา กามปิปาสา กามปริฬาโห, โส จ อชฺฌตฺตํ น สุปฺปหีโน โหติ น สุปฺปฏิปฺปสฺสทฺโธ. โอปกฺกมิกา เจปิ เต โภนฺโต สมณพฺราหฺมณา ทุกฺขา ติพฺพา ขรา กฏุกา เวทนา เวทยนฺติ, อภพฺพาว เต ญาณาย ทสฺสนาย อนุตฺตราย สมฺโพธาย. โน เจปิ เต โภนฺโต สมณพฺราหฺมณา โอปกฺกมิกา ทุกฺขา ติพฺพา ขรา กฏุกา เวทนา เวทยนฺติ, อภพฺพาว เต ญาณาย ทสฺสนาย อนุตฺตราย สมฺโพธาย. อยํ โข มํ อคฺคิเวสฺสน ปฐมา อุปมา ปฏิภาสิ อนจฺฉริยา ปุพฺเพ อสฺสุตปุพฺพา.}

\proseitem{375}{อปราปิ โข มํ อคฺคิเวสฺสน ทุติยา อุปมา ปฏิภาสิ อนจฺฉริยา ปุพฺเพ อสฺสุตปุพฺพา. เสยฺยถาปิ อคฺคิเวสฺสน อลฺลํ กฏฺฐํ สสฺเนหํ อารกา อุทกา ถเล นิกฺขิตฺตํ. อถ ปุริโส อาคจฺเฉยฺย อุตฺตรารณิํ อาทาย “อคฺคิํ อภินิพฺพตฺเตสฺสามิ เตโช ปาตุกริสฺสามี”ติ. ตํ กิํ มญฺญสิ อคฺคิเวสฺสน, อปิ นุ โส ปุริโส อมุํ อลฺลํ กฏฺฐํ สสฺเนหํ อารกา อุทกา ถเล นิกฺขิตฺตํ อุตฺตรารณิํ อาทาย อภิมนฺเถนฺโต อคฺคิํ อภินิพฺพตฺเตยฺย เตโช ปาตุกเรยฺยาติ. โน หิทํ โภ โคตม. ตํ กิสฺส เหตุ, อทุํ หิ โภ โคตม อลฺลํ กฏฺฐํ สสฺเนหํ กิญฺจาปิ อารกา อุทกา ถเล นิกฺขิตฺตํ. ยาวเทว จ ปน โส ปุริโส กิลมถสฺส \csromanfolio{308}วิฆาตสฺส ภาคี อสฺสาติ. เอวเมว โข อคฺคิเวสฺสน เย หิ เกจิ สมณา วา พฺราหฺมณา วา กาเยน เจว จิตฺเตน จ กาเมหิ วูปกฏฺฐา วิหรนฺติ, โย จ เนสํ กาเมสุ กามจฺฉนฺโท กามสฺเนโห กามมุจฺฉา กามปิปาสา กามปริฬาโห, โส จ อชฺฌตฺตํ น สุปฺปหีโน โหติ น สุปฺปฏิปฺปสฺสทฺโธ. โอปกฺกมิกา เจปิ เต โภนฺโต สมณพฺราหฺมณา ทุกฺขา ติพฺพา ขรา กฏุกา เวทนา เวทยนฺติ, อภพฺพาว เต ญาณาย ทสฺสนาย อนุตฺตราย สมฺโพธาย. โน เจปิ เต โภนฺโต สมณพฺราหฺมณา โอปกฺกมิกา ทุกฺขา ติพฺพา ขรา กฏุกา เวทนา เวทยนฺติ, อภพฺพาว เต ญาณาย ทสฺสนาย อนุตฺตราย สมฺโพธาย. อยํ โข มํ อคฺคิเวสฺสน ทุติยา อุปมา ปฏิภาสิ อนจฺฉริยา ปุพฺเพ อสฺสุตปุพฺพา.}

\proseitem{376}{อปราปิ โข มํ อคฺคิเวสฺสน ตติยา อุปมา ปฏิภาสิ อนจฺฉริยา ปุพฺเพ อสฺสุตปุพฺพา. เสยฺยถาปิ อคฺคิเวสฺสน สุกฺขํ กฏฺฐํ โกฬาปํ อารกา อุทกา ถเล นิกฺขิตฺตํ. อถ ปุริโส อาคจฺเฉยฺย อุตฺตรารณิํ อาทาย “อคฺคิํ อภินิพฺพตฺเตสฺสามิ เตโช ปาตุกริสฺสามี”ติ. ตํ กิํ มญฺญสิ อคฺคิเวสฺสน, อปิ นุ โส ปุริโส อมุํ สุกฺขํ กฏฺฐํ โกฬาปํ อารกา อุทกา ถเล นิกฺขิตฺตํ อุตฺตรารณิํ อาทาย อภิมนฺเถนฺโต อคฺคิํ อภินิพฺพตฺเตยฺย เตโช ปาตุกเรยฺยาติ. เอวํ โภ โคตม. ตํ กิสฺส เหตุ, อทุํ หิ โภ โคตม สุกฺขํ กฏฺฐํ โกฬาปํ, ตญฺจ ปน อารกา อุทกา ถเล นิกฺขิตฺตนฺติ. เอวเมว โข อคฺคิเวสฺสน เย หิ เกจิ สมณา วา พฺราหฺมณา วา กาเยน เจว จิตฺเตน จ กาเมหิ วูปกฏฺฐา วิหรนฺติ, โย จ เนสํ กาเมสุ กามจฺฉนฺโท กามสฺเนโห กามมุจฺฉา กามปิปาสา กามปริฬาโห, โส จ อชฺฌตฺตํ สุปฺปหีโน โหติ สุปฺปฏิปฺปสฺสทฺโธ. โอปกฺกมิกา เจปิ เต โภนฺโต สมณพฺราหฺมณา ทุกฺขา ติพฺพา ขรา กฏุกา เวทนา เวทยนฺติ, ภพฺพาว เต ญาณาย ทสฺสนาย อนุตฺตราย สมฺโพธาย. โน เจปิ เต โภนฺโต สมณพฺราหฺมณา โอปกฺกมิกา ทุกฺขา ติพฺพา ขรา กฏุกา เวทนา เวทยนฺติ, ภพฺพาว เต ญาณาย ทสฺสนาย อนุตฺตราย สมฺโพธาย. อยํ โข มํ อคฺคิเวสฺสน ตติยา อุปมา ปฏิภาสิ อนจฺฉริยา ปุพฺเพ อสฺสุตปุพฺพา. อิมา โข มํ อคฺคิเวสฺสน ติสฺโส อุปมา ปฏิภํสุ อนจฺฉริยา ปุพฺเพ อสฺสุตปุพฺพา.}

\proseitem{377}{\csromanfolio{309}ตสฺส มยฺหํ อคฺคิเวสฺสน เอตทโหสิ “ยํนูนาหํ ทนฺเตภิ\-ทนฺตมา\-ธาย\csromansharedfootnote{31d331dc88ab235a}{ปสฺส วิตกฺกสณฺฐานสุตฺเต (171) ปิฏฺเฐ.} ชิวฺหาย ตาลุํ อาหจฺจ เจตสา จิตฺตํ อภินิคฺคเณฺหยฺยํ อภินิปฺปีเฬยฺยํ อภิสนฺตาเปยฺยนฺ”ติ. โส โข อหํ อคฺคิเวสฺสน ทนฺเตภิ\-ทนฺตมา\-ธาย ชิวฺหาย ตาลุํ อาหจฺจ เจตสา จิตฺตํ อภินิคฺคณฺหามิ อภินิปฺปีเฬมิ อภิสนฺตาเปมิ. ตสฺส มยฺหํ อคฺคิเวสฺสน ทนฺเตภิ\-ทนฺตมา\-ธาย ชิวฺหาย ตาลุํ อาหจฺจ เจตสา จิตฺตํ อภินิคฺคณฺหโต อภินิปฺปีฬยโต อภิสนฺตาปยโต กจฺเฉหิ เสทา มุจฺจนฺติ. เสยฺยถาปิ อคฺคิเวสฺสน พลวา ปุริโส ทุพฺพลตรํ ปุริสํ สีเส วา คเหตฺวา ขนฺเธ วา คเหตฺวา อภินิคฺคเณฺหยฺย อภินิปฺปีเฬยฺย อภิสนฺตาเปยฺย. เอวเมว โข เม อคฺคิเวสฺสน ทนฺเตภิ\-ทนฺตมา\-ธาย ชิวฺหาย ตาลุํ อาหจฺจ เจตสา จิตฺตํ อภินิคฺคณฺหโต อภินิปฺปีฬยโต อภิสนฺตาปยโต กจฺเฉหิ เสทา มุจฺจนฺติ. อารทฺธํ โข ปน เม อคฺคิเวสฺสน วีริยํ โหติ อสลฺลีนํ, อุปฏฺฐิตา สติ อสมฺมุฏฺฐา. สารทฺโธ จ ปน เม กาโย โหติ อปฺปฏิปฺปสฺสทฺโธ เตเนว ทุกฺข\-ปฺปธาเนน ปธานา\-ภิตุนฺนสฺส สโต. เอวรูปาปิ โข เม อคฺคิเวสฺสน อุปฺปนฺนา ทุกฺขา เวทนา จิตฺตํ น ปริยาทาย ติฏฺฐติ.}

\proseitem{378}{ตสฺส มยฺหํ อคฺคิเวสฺสน เอตทโหสิ “ยํนูนาหํ อปฺปาณกํเยว ฌานํ ฌาเยยฺยนฺ”ติ. โส โข อหํ อคฺคิเวสฺสน มุขโต จ นาสโต จ อสฺสาสปสฺสาเส อุปรุนฺธิํ. ตสฺส มยฺหํ อคฺคิเวสฺสน มุขโต จ นาสโต จ อสฺสาส\-ปสฺสาเสสุ อุปรุทฺเธสุ กณฺณโสเตหิ วาตานํ นิกฺขมนฺตานํ อธิมตฺโต สทฺโท โหติ. เสยฺยถาปิ นาม กมฺมาร\-คคฺคริยา ธมมานาย อธิมตฺโต สทฺโท โหติ. เอวเมว โข เม อคฺคิเวสฺสน มุขโต จ นาสโต จ อสฺสาส\-ปสฺสาเสสุ อุปรุทฺเธสุ กณฺณโสเตหิ วาตานํ นิกฺขมนฺตานํ อธิมตฺโต สทฺโท โหติ. อารทฺธํ โข ปน เม อคฺคิเวสฺสน วีริยํ โหติ อสลฺลีนํ, อุปฏฺฐิตา สติ อสมฺมุฏฺฐา. สารทฺโธ จ ปน เม กาโย โหติ อปฺปฏิปฺปสฺสทฺโธ เตเนว ทุกฺข\-ปฺปธาเนน ปธานา\-ภิตุนฺนสฺส สโต. เอวรูปาปิ โข เม อคฺคิเวสฺสน อุปฺปนฺนา ทุกฺขา เวทนา จิตฺตํ น ปริยาทาย ติฏฺฐติ.}

\prose{ตสฺส มยฺหํ อคฺคิเวสฺสน เอตทโหสิ “ยํนูนาหํ อปฺปาณกํเยว ฌานํ ฌาเยยฺยนฺ”ติ. โส โข อหํ อคฺคิเวสฺสน มุขโต จ นาสโต จ \csromanfolio{310}กณฺณโต จ อสฺสาสปสฺสาเส อุปรุนฺธิํ. ตสฺส มยฺหํ อคฺคิเวสฺสน มุขโต จ นาสโต จ กณฺณโต จ อสฺสาส\-ปสฺสาเสสุ อุปรุทฺเธสุ อธิมตฺตา วาตา มุทฺธนิ อูหนนฺติ\csromansharedfootnote{46dffcde156e29c1}{อูหนฺติ (สี), โอหนนฺติ (สฺยา, กํ), อุหนนฺติ (ก)}. เสยฺยถาปิ อคฺคิเวสฺสน พลวา ปุริโส ติเณฺหน สิขเรน มุทฺธนิ อภิมตฺเถยฺย\csromansharedfootnote{4c01165ef7b5185a}{มุทฺธานํ อภิมนฺเถยฺย (สี, อิ), มุทฺธานํ อภิมตฺเถยฺย (สฺยา, กํ)}. เอวเมว โข เม อคฺคิเวสฺสน มุขโต จ นาสโต จ กณฺณโต จ อสฺสาส\-ปสฺสาเสสุ อุปรุทฺเธสุ อธิมตฺตา วาตา มุทฺธนิ อูหนนฺติ. อารทฺธํ โข ปน เม อคฺคิเวสฺสน วีริยํ โหติ อสลฺลีนํ, อุปฏฺฐิตา สติ อสมฺมุฏฺฐา. สารทฺโธ จ ปน เม กาโย โหติ อปฺปฏิปฺปสฺสทฺโธ เตเนว ทุกฺข\-ปฺปธาเนน ปธานา\-ภิตุนฺนสฺส สโต. เอวรูปาปิ โข เม อคฺคิเวสฺสน อุปฺปนฺนา ทุกฺขา เวทนา จิตฺตํ น ปริยาทาย ติฏฺฐติ.}

\prose{ตสฺส มยฺหํ อคฺคิเวสฺสน เอตทโหสิ “ยํนูนาหํ อปฺปาณกํเยว ฌานํ ฌาเยยฺยนฺ”ติ. โส โข อหํ อคฺคิเวสฺสน มุขโต จ นาสโต จ กณฺณโต จ อสฺสาสปสฺสาเส อุปรุนฺธิํ. ตสฺส มยฺหํ อคฺคิเวสฺสน มุขโต จ นาสโต จ กณฺณโต จ อสฺสาส\-ปสฺสาเสสุ อุปรุทฺเธสุ อธิมตฺตา สีเส สีสเวทนา โหนฺติ. เสยฺยถาปิ อคฺคิเวสฺสน พลวา ปุริโส ทเฬฺหน วรตฺต\-กฺขณฺเฑน\csromansharedfootnote{e54a1a83680c8582}{วรตฺตกพนฺธเนน (สี)} สีเส สีสเวฐํ ทเทยฺย. เอวเมว โข เม อคฺคิเวสฺสน มุขโต จ นาสโต จ กณฺณโต จ อสฺสาส\-ปสฺสาเสสุ อุปรุทฺเธสุ อธิมตฺตา สีเส สีสเวทนา โหนฺติ. อารทฺธํ โข ปน เม อคฺคิเวสฺสน วีริยํ โหติ อสลฺลีนํ, อุปฏฺฐิตา สติ อสมฺมุฏฺฐา. สารทฺโธ จ ปน เม กาโย โหติ อปฺปฏิปฺปสฺสทฺโธ เตเนว ทุกฺข\-ปฺปธาเนน ปธานา\-ภิตุนฺนสฺส สโต. เอวรูปาปิ โข เม อคฺคิเวสฺสน อุปฺปนฺนา ทุกฺขา เวทนา จิตฺตํ น ปริยาทาย ติฏฺฐติ.}

\prose{ตสฺส มยฺหํ อคฺคิเวสฺสน เอตทโหสิ “ยํนูนาหํ อปฺปาณกํเยว ฌานํ ฌาเยยฺยนฺ”ติ. โส โข อหํ อคฺคิเวสฺสน มุขโต จ นาสโต จ กณฺณโต จ อสฺสาสปสฺสาเส อุปรุนฺธิํ. ตสฺส มยฺหํ อคฺคิเวสฺสน มุขโต จ นาสโต จ กณฺณโต จ อสฺสาส\-ปสฺสาเสสุ อุปรุทฺเธสุ อธิมตฺตา วาตา กุจฺฉิํ ปริกนฺตนฺติ. เสยฺยถาปิ อคฺคิเวสฺสน ทกฺโข โคฆาตโก วา โคฆาตก\-นฺเตวาสี วา \csromanfolio{311}ติเณฺหน โควิกนฺตเนน กุจฺฉิํ ปริกนฺเตยฺย. เอวเมว โข เม อคฺคิเวสฺสน มุขโต จ นาสโต จ กณฺณโต จ อสฺสาส\-ปสฺสาเสสุ อุปรุทฺเธสุ อธิมตฺตา วาตา กุจฺฉิํ ปริกนฺตนฺติ. อารทฺธํ โข ปน เม อคฺคิเวสฺสน วีริยํ โหติ อสลฺลีนํ, อุปฏฺฐิตา สติ อสมฺมุฏฺฐา. สารทฺโธ จ ปน เม กาโย โหติ อปฺปฏิปฺปสฺสทฺโธ เตเนว ทุกฺข\-ปฺปธาเนน ปธานา\-ภิตุนฺนสฺส สโต. เอวรูปาปิ โข เม อคฺคิเวสฺสน อุปฺปนฺนา ทุกฺขา เวทนา จิตฺตํ น ปริยาทาย ติฏฺฐติ.}

\prose{ตสฺส มยฺหํ อคฺคิเวสฺสน เอตทโหสิ “ยํนูนาหํ อปฺปาณกํเยว ฌานํ ฌาเยยฺยนฺ”ติ. โส โข อหํ อคฺคิเวสฺสน มุขโต จ นาสโต จ กณฺณโต จ อสฺสาสปสฺสาเส อุปรุนฺธิํ. ตสฺส มยฺหํ อคฺคิเวสฺสน มุขโต จ นาสโต จ กณฺณโต จ อสฺสาส\-ปสฺสาเสสุ อุปรุทฺเธสุ อธิมตฺโต กายสฺมิํ ฑาโห โหติ. เสยฺยถาปิ อคฺคิเวสฺสน เทฺว พลวนฺโต ปุริสา ทุพฺพลตรํ ปุริสํ นานาพาหาสุ คเหตฺวา องฺคารกาสุยา สนฺตาเปยฺยุํ สมฺปริตาเปยฺยุํ. เอวเมว โข เม อคฺคิเวสฺสน มุขโต จ นาสโต จ กณฺณโต จ อสฺสาส\-ปสฺสาเสสุ อุปรุทฺเธสุ อธิมตฺโต กายสฺมิํ ฑาโห โหติ. อารทฺธํ โข ปน เม อคฺคิเวสฺสน วีริยํ โหติ อสลฺลีนํ, อุปฏฺฐิตา สติ อสมฺมุฏฺฐา. สารทฺโธ จ ปน เม กาโย โหติ อปฺปฏิปฺปสฺสทฺโธ เตเนว ทุกฺข\-ปฺปธาเนน ปธานา\-ภิตุนฺนสฺส สโต. เอวรูปาปิ โข เม อคฺคิเวสฺสน อุปฺปนฺนา ทุกฺขา เวทนา จิตฺตํ น ปริยาทาย ติฏฺฐติ. อปิสฺสุ มํ อคฺคิเวสฺสน เทวตา ทิสฺวา เอวมาหํสุ “กาลงฺกโต สมโณ โคตโม”ติ. เอกจฺจา เทวตา เอวมาหํสุ “น กาลงฺกโต สมโณ โคตโม, อปิ จ กาลงฺกโรตี”ติ. เอกจฺจา เทวตา เอวมาหํสุ “น กาลงฺกโต สมโณ โคตโม, นปิ กาลงฺกโรติ. อรหํ สมโณ โคตโม, วิหาโรเตฺวว โส\csromansharedfootnote{60d11e19a5379e76}{วิหาโรเตฺวเวโส (สี)} อรหโต เอวรูโป โหตี”ติ\csromansharedfootnote{eaf0c06433966913}{วิหาโรเตฺวเวโส อรหโต”ติ (?)}.}

\proseitem{379}{ตสฺส มยฺหํ อคฺคิเวสฺสน เอตทโหสิ “ยํนูนาหํ สพฺพโส อาหารุ\-ปจฺเฉทาย ปฏิปชฺเชยฺยนฺ”ติ. อถ โข มํ อคฺคิเวสฺสน เทวตา \csromanfolio{312}อุปสงฺกมิตฺวา เอตทโวจุํ “มา โข ตฺวํ มาริส สพฺพโส อาหารุ\-ปจฺเฉทาย ปฏิปชฺชิ. สเจ โข ตฺวํ มาริส สพฺพโส อาหารุ\-ปจฺเฉทาย ปฏิปชฺชิสฺสสิ, ตสฺส เต มยํ ทิพฺพํ โอชํ โลมกูเปหิ อชฺโฌหาเรสฺสาม\csromansharedfootnote{447dbcb3e8f9b303}{อชฺโฌหริสฺสาม (สฺยา, กํ, อิ, ก)}, ตาย ตฺวํ ยาเปสฺสสี”ติ. ตสฺส มยฺหํ อคฺคิเวสฺสน เอตทโหสิ “อหญฺเจว โข ปน สพฺพโส อชชฺชิตํ\csromansharedfootnote{d53fb5ce3bf3e6a2}{อชทฺธุกํ (สี, อิ), ชทฺธุกํ (สฺยา, กํ)} ปฏิชาเนยฺยํ. อิมา จ เม เทวตา ทิพฺพํ โอชํ โลมกูเปหิ อชฺโฌหาเรยฺยุํ\csromansharedfootnote{9be773cd5db336f5}{อชฺโฌหเรยฺยุํ (สฺยา, กํ, อิ, ก)}, ตาย จาหํ ยาเปยฺยํ, ตํ มมสฺส มุสา”ติ. โส โข อหํ อคฺคิเวสฺสน ตา เทวตา ปจฺจาจิกฺขามิ “หลนฺ”ติ วทามิ.}

\proseitem{380}{ตสฺส มยฺหํ อคฺคิเวสฺสน เอตทโหสิ “ยํนูนาหํ โถกํ โถกํ อาหารํ อาหาเรยฺยํ ปสตํ ปสตํ, ยทิ วา มุคฺคยูสํ ยทิ วา กุลตฺถยูสํ ยทิ วา กฬายยูสํ ยทิ วา หเรณุกยูสนฺ”ติ. โส โข อหํ อคฺคิเวสฺสน โถกํ โถกํ อาหารํ อาหาเรสิํ ปสตํ ปสตํ, ยทิ วา มุคฺคยูสํ ยทิ วา กุลตฺถยูสํ ยทิ วา กฬายยูสํ ยทิ วา หเรณุกยูสํ. ตสฺส มยฺหํ อคฺคิเวสฺสน โถกํ โถกํ อาหารํ อาหารยโต ปสตํ ปสตํ, ยทิ วา มุคฺคยูสํ ยทิ วา กุลตฺถยูสํ ยทิ วา กฬายยูสํ ยทิ วา หเรณุกยูสํ, อธิมตฺต\-กสิมานํ ปตฺโต กาโย โหติ. เสยฺยถาปิ นาม อาสีติกปพฺพานิ วา กาฬปพฺพานิ วา, เอวเมวสฺสุ เม องฺคปจฺจงฺคานิ ภวนฺติ ตาเย\-ว\-ปฺปาหารตาย. เสยฺยถาปิ นาม โอฏฺฐปทํ, เอวเมวสฺสุ เม อานิสทํ โหติ ตาเย\-ว\-ปฺปาหารตาย. เสยฺยถาปิ นาม วฏฺฏนาวฬี, เอวเมวสฺสุ เม ปิฏฺฐิกณฺฏโก อุณฺณตาวนโต โหติ ตาเย\-ว\-ปฺปาหารตาย. เสยฺยถาปิ นาม ชรสาลาย โคปาณสิโย โอลุคฺควิลุคฺคา ภวนฺติ, เอวเมวสฺสุ เม ผาสุฬิโย โอลุคฺควิลุคฺคา ภวนฺติ ตาเย\-ว\-ปฺปาหารตาย. เสยฺยถาปิ นาม คมฺภีเร อุทปาเน อุทกตารกา คมฺภีรคตา โอกฺขายิกา ทิสฺสนฺติ, เอวเมวสฺสุ เม อกฺขิกูเปสุ อกฺขิตารกา คมฺภีรคตา โอกฺขายิกา ทิสฺสนฺติ ตาเย\-ว\-ปฺปาหารตาย. เสยฺยถาปิ นาม ติตฺตกาลาพุ อามกจฺฉินฺโน วาตาตเปน สํผุฏิโต โหติ สมฺมิลาโต, เอวเมวสฺสุ เม สีสจฺฉวิ สํผุฏิตา โหติ \csromanfolio{313}สมฺมิลาตา ตาเย\-ว\-ปฺปาหารตาย. โส โข อหํ อคฺคิเวสฺสน อุทรจฺฉวิํ ปริมสิสฺสามีติ ปิฏฺฐิ\-กณฺฏกํเยว ปริคฺคณฺหามิ, ปิฏฺฐิกณฺฏกํ ปริมสิสฺสามีติ อุทร\-จฺฉวิํเยว ปริคฺคณฺหามิ. ยาวสฺสุ เม อคฺคิเวสฺสน อุทรจฺฉวิ ปิฏฺฐิกณฺฏกํ อลฺลีนา โหติ ตาเย\-ว\-ปฺปาหารตาย. โส โข อหํ อคฺคิเวสฺสน วจฺจํ วา มุตฺตํ วา กริสฺสามีติ ตตฺเถว อวกุชฺโช ปปตามิ ตาเย\-ว\-ปฺปาหารตาย. โส โข อหํ อคฺคิเวสฺสน อิมเมว กายํ อสฺสาเสนฺโต ปาณินา คตฺตานิ อนุมชฺชามิ. ตสฺส มยฺหํ อคฺคิเวสฺสน ปาณินา คตฺตานิ อนุมชฺชโต ปูติมูลานิ โลมานิ กายสฺมา ปปตนฺติ ตาเย\-ว\-ปฺปาหารตาย. อปิสฺสุ มํ อคฺคิเวสฺสน มนุสฺสา ทิสฺวา เอวมาหํสุ “กาโฬ สมโณ โคตโม”ติ. เอกจฺเจ มนุสฺสา เอวมาหํสุ “น กาโฬ สมโณ โคตโม, สาโม สมโณ โคตโม”ติ. เอกจฺเจ มนุสฺสา เอวมาหํสุ “น กาโฬ สมโณ โคตโม นปิ สาโม, มงฺคุรจฺฉวิ สมโณ โคตโม”ติ. ยาวสฺสุ เม อคฺคิเวสฺสน ตาว ปริสุทฺโธ ฉวิวณฺโณ ปริโยทาโต อุปหโต โหติ ตาเย\-ว\-ปฺปาหารตาย.}

\proseitem{381}{ตสฺส มยฺหํ อคฺคิเวสฺสน เอตทโหสิ “เย โข เกจิ อตีตมทฺธานํ สมณา วา พฺราหฺมณา วา โอปกฺกมิกา ทุกฺขา ติพฺพา ขรา กฏุกา เวทนา เวทยิํสุ, เอตาวปรมํ นยิโต ภิยฺโย. เยปิ หิ เกจิ อนาคตมทฺธานํ สมณา วา พฺราหฺมณา วา โอปกฺกมิกา ทุกฺกา ติพฺพา ขรา กฏุกา เวทนา เวทยิสฺสนฺติ, เอตาวปรมํ นยิโต ภิยฺโย. เยปิ หิ เกจิ เอตรหิ สมณา วา พฺราหฺมณา วา โอปกฺกมิกา ทุกฺขา ติพฺพา ขรา กฏุกา เวทนา เวทยนฺติ, เอตาวปรมํ นยิโต ภิยฺโย. น โข ปนาหํ อิมาย กฏุกาย ทุกฺกร\-การิกาย อธิคจฺฉามิ อุตฺตริ มนุสฺสธมฺมา อลม\-ริย\-ญาณ\-ทสฺสน\-วิเสสํ. สิยา นุ โข อญฺโญ มคฺโค โพธายา”ติ. ตสฺส มยฺหํ อคฺคิเวสฺสน เอตทโหสิ “อภิชานามิ โข ปนาหํ ปิตุ สกฺกสฺส กมฺมนฺเต สีตาย ชมฺพุจฺฉายาย นิสินฺโน วิวิจฺเจว กาเมหิ วิวิจฺจ อกุสเลหิ ธมฺเมหิ สวิตกฺกํ สวิจารํ วิเวกชํ ปีติสุขํ ปฐมํ ฌานํ อุปสมฺปชฺช วิหริตา. สิยา นุ โข เอโส มคฺโค โพธายา”ติ. ตสฺส มยฺหํ อคฺคิเวสฺสน สตานุสาริ วิญฺญาณํ อโหสิ “เอเสว มคฺโค โพธายา”ติ. ตสฺส มยฺหํ อคฺคิเวสฺสน เอตทโหสิ “กิํ นุ โข อหํ ตสฺส สุขสฺส ภายามิ, ยํ ตํ สุขํ อญฺญเตฺรว กาเมหิ \csromanfolio{314}อญฺญตฺร อกุสเลหิ ธมฺเมหี”ติ. ตสฺส มยฺหํ อคฺคิเวสฺสน เอตทโหสิ “น โข อหํ ตสฺส สุขสฺส ภายามิ, ยํ ตํ สุขํ อญฺญเตฺรว กาเมหิ อญฺญตฺร อกุสเลหิ ธมฺเมหี”ติ.}

\proseitem{382}{ตสฺส มยฺหํ อคฺคิเวสฺสน เอตทโหสิ “น โข ตํ สุกรํ สุขํ อธิคนฺตุํ เอวํ อธิมตฺต\-กสิมานํ ปตฺตกาเยน. ยํนูนาหํ โอฬาริกํ อาหารํ อาหาเรยฺยํ โอทนกุมฺมาสนฺ”ติ. โส โข อหํ อคฺคิเวสฺสน โอฬาริกํ อาหารํ อาหาเรสิํ โอทนกุมฺมาสํ. เตน โข ปน มํ อคฺคิเวสฺสน สมเยน ปญฺจ\csromansharedfootnote{f63682c6c2edf8cf}{ปญฺจวคฺคิยา (อญฺญสุตฺเตสุ)} ภิกฺขู ปจฺจุปฏฺฐิตา โหนฺติ “ยํ โข สมโณ โคตโม ธมฺมํ อธิคมิสฺสติ, ตํ โน อาโรเจสฺสตี”ติ. ยโต โข อหํ อคฺคิเวสฺสน โอฬาริกํ อาหารํ อาหาเรสิํ โอทนกุมฺมาสํ, อถ เม เต ปญฺจ\csromansharedfootnote{e10cfdaac2749903}{พาหุลิโก (สี, อิ) สํฆเภทสิกฺขาปทฏีกาย สเมติ.} ภิกฺขู นิพฺพิชฺช ปกฺกมิํสุ “พาหุลฺลิโก\csromansharedfootnote{f63682c6c2edf8cf}{ปญฺจวคฺคิยา (อญฺญสุตฺเตสุ)} สมโณ โคตโม ปธาน\-วิพฺภนฺโต อาวตฺโต พาหุลฺลายา”ติ.}

\proseitem{383}{โส โข อหํ อคฺคิเวสฺสน โอฬาริกํ อาหารํ อาหาเรตฺวา พลํ คเหตฺวา วิวิจฺเจว กาเมหิ วิวิจฺจ อกุสเลหิ ธมฺเมหิ สวิตกฺกํ สวิจารํ วิเวกชํ ปีติสุขํ ปฐมํ ฌานํ อุปสมฺปชฺช วิหาสิํ. เอวรูปาปิ โข เม อคฺคิเวสฺสน อุปฺปนฺนา สุขา เวทนา จิตฺตํ น ปริยาทาย ติฏฺฐติ. วิตกฺก\-วิจารานํ วูปสมา อชฺฌตฺตํ สมฺปสาทนํ เจตโส เอโกทิภาวํ อวิตกฺกํ อวิจารํ สมาธิชํ ปีติสุขํ ทุติยํ ฌานํ อุปสมฺปชฺช วิหาสิํ. เอวรูปาปิ โข เม อคฺคิเวสฺสน อุปฺปนฺนา สุขา เวทนา จิตฺตํ น ปริยาทาย ติฏฺฐติ. ปีติยา จ วิราคา อุเปกฺขโก จ วิหาสิํ สโต จ สมฺปชาโน, สุขญฺจ กาเยน ปฏิสํเวเทสิํ, ยํ ตํ อริยา อาจิกฺขนฺติ “อุเปกฺขโก สติมา สุขวิหารี”ติ ตติยํ ฌานํ อุปสมฺปชฺช วิหาสิํ. เอวรูปาปิ โข เม อคฺคิเวสฺสน อุปฺปนฺนา สุขา เวทนา จิตฺตํ น ปริยาทาย ติฏฺฐติ. สุขสฺส จ ปหานา ทุกฺขสฺส จ ปหานา ปุพฺเพว โสมนสฺส\-โทมนสฺสานํ อตฺถงฺคมา อทุกฺขมสุขํ อุเปกฺขา\-สติ\-ปาริสุทฺธิํ จตุตฺถํ ฌานํ อุปสมฺปชฺช วิหาสิํ. เอวรูปาปิ โข เม อคฺคิเวสฺสน อุปฺปนฺนา สุขา เวทนา จิตฺตํ น ปริยาทาย ติฏฺฐติ.}

\proseitem{384}{โส เอวํ สมาหิเต จิตฺเต ปริสุทฺเธ ปริโยทาเต อนงฺคเณ วิคตู\-ปกฺกิเลเส มุทุภูเต กมฺมนิเย ฐิเต อาเนญฺชปฺปตฺเต ปุพฺเพ\-นิวาสา\-นุสฺสติ\-ญาณาย จิตฺตํ อภินินฺนาเมสิํ. โส อเนกวิหิตํ ปุพฺเพนิวาสํ \csromanfolio{315}อนุสฺสรามิ. เสยฺยถิทํ, เอกมฺปิ ชาติํ ฯเปฯ อิติ สาการํ สฺเอาทฺเทสํ อเนกวิหิตํ ปุพฺเพนิวาสํ อนุสฺสรามิ. อยํ โข เม อคฺคิเวสฺสน รตฺติยา ปฐเม ยาเม ปฐมา วิชฺชา อธิคตา, อวิชฺชา วิหตา, วิชฺชา อุปฺปนฺนา, ตโม วิหโต, อาโลโก อุปฺปนฺโน, ยถา ตํ อปฺปมตฺตสฺส อาตาปิโน ปหิตตฺตสฺส วิหรโต. เอวรูปาปิ โข เม อคฺคิเวสฺสน อุปฺปนฺนา สุขา เวทนา จิตฺตํ น ปริยาทาย ติฏฺฐติ.}

\proseitem{385}{โส เอวํ สมาหิเต จิตฺเต ปริสุทฺเธ ปริโยทาเต อนงฺคเณ วิคตู\-ปกฺกิเลเส มุทุภูเต กมฺมนิเย ฐิเต อาเนญฺชปฺปตฺเต สตฺตานํ จุตูปปาต\-ญาณาย จิตฺตํ อภินินฺนาเมสิํ. โส ทิพฺเพน จกฺขุนา วิสุทฺเธน อติกฺกนฺต\-มานุสเกน สตฺเต ปสฺสามิ จวมาเน อุปปชฺชมาเน หีเน ปณีเต สุวณฺเณ ทุพฺพณฺเณ สุคเต ทุคฺคเต ยถากมฺมูปเค สตฺเต ปชานามิ ฯเปฯ. อยํ โข เม อคฺคิเวสฺสน รตฺติยา มชฺฌิเม ยาเม ทุติยา วิชฺชา อธิคตา, อวิชฺชา วิหตา, วิชฺชา อุปฺปนฺนา, ตโม วิหโต, อาโลโก อุปฺปนฺโน, ยถา ตํ อปฺปมตฺตสฺส อาตาปิโน ปหิตตฺตสฺส วิหรโต. เอวรูปาปิ โข เม อคฺคิเวสฺสน อุปฺปนฺนา สุขา เวทนา จิตฺตํ น ปริยาทาย ติฏฺฐติ.}

\proseitem{386}{โส เอวํ สมาหิเต จิตฺเต ปริสุทฺเธ ปริโยทาเต อนงฺคเณ วิคตู\-ปกฺกิเลเส มุทุภูเต กมฺมนิเย ฐิเต อาเนญฺชปฺปตฺเต อาสวานํ ขยญาณาย จิตฺตํ อภินินฺนาเมสิํ. โส อิทํ ทุกฺขนฺติ ยถาภูตํ อพฺภญฺญาสิํ, อยํ ทุกฺข\-สมุทโยติ ยถาภูตํ อพฺภญฺญาสิํ, อยํ ทุกฺข\-นิโรโธติ ยถาภูตํ อพฺภญฺญาสิํ, อยํ ทุกฺขนิโรธ\-คามินี ปฏิปทาติ ยถาภูตํ อพฺภญฺญาสิํ, อิเม อาสวาติ ยถาภูตํ อพฺภญฺญาสิํ, อยํ อาสว\-สมุทโยติ ยถาภูตํ อพฺภญฺญาสิํ, อยํ อาสวนิโรโธติ ยถาภูตํ อพฺภญฺญาสิํ, อยํ อาสว\-นิโรธ\-คามินี ปฏิปทาติ ยถาภูตํ อพฺภญฺญาสิํ. ตสฺส เม เอวํ ชานโต เอวํ ปสฺสโต กามาสวาปิ จิตฺตํ วิมุจฺจิตฺถ, ภวาสวาปิ จิตฺตํ วิมุจฺจิตฺถ, อวิชฺชาสวาปิ จิตฺตํ วิมุจฺจิตฺถ, วิมุตฺตสฺมิํ “วิมุตฺตมฺ”อิติ ญาณํ อโหสิ, “ขีณา ชาติ, วุสิตํ พฺรหฺมจริยํ, กตํ กรณียํ, นาปรํ อิตฺถตฺตายา”ติ อพฺภญฺญาสิํ. อยํ โข เม อคฺคิเวสฺสน รตฺติยา ปจฺฉิเม ยาเม ตติยา วิชฺชา อธิคตา, อวิชฺชา วิหตา, วิชฺชา อุปฺปนฺนา, ตโม วิหโต, อาโลโก อุปฺปนฺโน, \csromanfolio{316}ยถา ตํ อปฺปมตฺตสฺส อาตาปิโน ปหิตตฺตสฺส วิหรโต. เอวรูปาปิ โข เม อคฺคิเวสฺสน อุปฺปนฺนา สุขา เวทนา จิตฺตํ น ปริยาทาย ติฏฺฐติ.}

\proseitem{387}{อภิชานามิ โข ปนาหํ อคฺคิเวสฺสน อเนกสตาย ปริสาย ธมฺมํ เทเสตา, อปิสฺสุ มํ เอกเมโก เอวํ มญฺญติ “มเมวารพฺภ สมโณ โคตโม ธมฺมํ เทเสตี”ติ. น โข ปเนตํ อคฺคิเวสฺสน เอวํ ทฏฺฐพฺพํ. ยาวเทว วิญฺญาปน\-ตฺถาย ตถาคโต ปเรสํ ธมฺมํ เทเสติ. โส โข อหํ อคฺคิเวสฺสน ตสฺสาเยว กถาย ปริโยสาเน ตสฺมิํเยว ปุริมสฺมิํ สมาธินิมิตฺเต อชฺฌตฺตเมว จิตฺตํ สณฺฐเปมิ สนฺนิสาเทมิ, เอโกทิํ กโรมิ, สมาทหามิ “เยน สุทํ นิจฺจกปฺปํ วิหรามี”ติ.}

\prose{โอกปฺปนิยเมตํ โภโต โคตมสฺส. ยถา ตํ อรหโต สมฺมา\-สมฺพุทฺธสฺส. อภิชานาติ โข ปน ภวํ โคตโม “ทิวา สุปิตา”ติ. อภิชานามหํ อคฺคิเวสฺสน คิมฺหานํ ปจฺฉิเม มาเส ปจฺฉาภตฺตํ ปิณฺฑปาต\-ปฏิกฺกนฺโต จตุคฺคุณํ สํฆาฏิํ ปญฺญเปตฺวา ทกฺขิเณน ปสฺเสน สโต สมฺปชาโน นิทฺทํ โอกฺกมิตาติ. เอตํ โข โภ โคตม เอเก สมณพฺราหฺมณา สมฺโมห\-วิหารสฺมิํ วทนฺตีติ. น โข อคฺคิเวสฺสน เอตฺตาวตา สมฺมูโฬฺห วา โหติ อสมฺมูโฬฺห วา, อปิ จ อคฺคิเวสฺสน ยถา สมฺมูโฬฺห จ โหติ อสมฺมูโฬฺห จ, ตํ สุณาหิ สาธุกํ มนสิ กโรหิ ภาสิสฺสามีติ. เอวํ โภติ โข สจฺจโก นิคณฺฐปุตฺโต ภควโต ปจฺจสฺโสสิ. ภควา เอตทโวจ\csromandash{}}

\proseitem{388}{ยสฺส กสฺสจิ อคฺคิเวสฺสน เย อาสวา สํกิเลสิกา โปโนพฺภวิกา สทรา ทุกฺขวิปากา อายติํ ชาติ\-ชรา\-มรณิยา อปฺปหีนา, ตมหํ สมฺมูโฬฺหติ วทามิ, อาสวานํ หิ อคฺคิเวสฺสน อปฺปหานา สมฺมูโฬฺห โหติ. ยสฺส กสฺสจิ อคฺคิเวสฺสน เย อาสวา สํกิเลสิกา โปโนพฺภวิกา สทรา ทุกฺขวิปากา อายติํ ชาติชรามรณียา ปหีนา, ตมหํ อสมฺมูโฬฺหติ วทามิ, อาสวานํ หิ อคฺคิเวสฺสน ปหานา อสมฺมูโฬฺห โหติ.}

\prose{ตถาคตสฺส โข อคฺคิเวสฺสน เย อาสวา สํกิเลสิกา โปโนพฺภวิกา สทรา ทุกฺขวิปากา อายติํ ชาติ\-ชรา\-มรณิยา ปหีนา อุจฺฉินฺนมูลา ตาลาวตฺถุกตา อนภาวํกตา อายติํ \csromanfolio{317}อนุปฺปาทธมฺมา. เสยฺยถาปิ อคฺคิเวสฺสน ตาโล มตฺถก\-จฺฉินฺโน อภพฺโพ ปุน วิรูฬฺหิยา. เอวเมว โข อคฺคิเวสฺสน ตถาคตสฺส เย อาสวา สํกิเลสิกา โปโนพฺภวิกา สทรา ทุกฺขวิปากา อายติํ ชาติ\-ชรา\-มรณิยา ปหีนา อุจฺฉินฺนมูลา ตาลาวตฺถุกตา อนภาวํกตา อายติํ อนุปฺปาทธมฺมาติ.}

\proseitemwithcloser{389}{เอวํ วุตฺเต สจฺจโก นิคณฺฐปุตฺโต ภควนฺตํ เอตทโวจ “อจฺฉริยํ โภ โคตม, อพฺภุตํ โภ โคตม, ยาวญฺจิทํ โภโต โคตมสฺส เอวํ อาสชฺช อาสชฺช วุจฺจมานสฺส อุปนีเตหิ วจนปฺปเถหิ สมุทา\-จ\-ริย\-มานสฺส ฉวิวณฺโณ เจว ปริโยทายติ, มุขวณฺโณ จ วิปฺปสีทติ, ยถา ตํ อรหโต สมฺมา\-สมฺพุทฺธสฺส. อภิชานามหํ โภ โคตม ปูรณํ กสฺสปํ วาเทน วาทํ สมารภิตา, โสปิ มยา วาเทน วาทํ สมารทฺโธ อญฺเญนญฺญํ ปฏิจริ, พหิทฺธา กถํ อปนาเมสิ, โกปญฺจ โทสญฺจ อปฺปจฺจยญฺจ ปาตฺวากาสิ. โภโต ปน\csromansharedfootnote{bfe12379ac10e719}{โภโต โข ปน (สี)} โคตมสฺส เอวํ อาสชฺช อาสชฺช วุจฺจมานสฺส อุปนีเตหิ วจนปฺปเถหิ สมุทา\-จ\-ริย\-มานสฺส ฉวิวณฺโณ เจว ปริโยทายติ, มุขวณฺโณ จ วิปฺปสีทติ, ยถา ตํ อรหโต สมฺมา\-สมฺพุทฺธสฺส. อภิชานามหํ โภ โคตม มกฺขลิํ โคสาลํ ฯเปฯ อชิตํ เกสกมฺพลํ. ปกุธํ กจฺจายนํ. สญฺชยํ เพลฏฺฐปุตฺตํ. นิคณฺฐํ นาฏปุตฺตํ วาเทน วาทํ สมารภิตา, โสปิ มยา วาเทน วาทํ สมารทฺโธ อญฺเญนญฺญํ ปฏิจริ, พหิทฺธา กถํ อปนาเมสิ, โกปญฺจ โทสญฺจ อปฺปจฺจยญฺจ ปาตฺวากาสิ. โภโต ปน โคตมสฺส เอวํ อาสชฺช อาสชฺช วุจฺจมานสฺส อุปนีเตหิ วจนปฺปเถหิ สมุทา\-จ\-ริย\-มานสฺส ฉวิวณฺโณ เจว ปริโยทายติ, มุขวณฺโณ จ วิปฺปสีทติ, ยถา ตํ อรหโต สมฺมา\-สมฺพุทฺธสฺส. หนฺท จ ทานิ มยํ โภ โคตม คจฺฉาม, พหุกิจฺจา มยํ พหุกรณียา”ติ. ยสฺสทานิ ตฺวํ อคฺคิเวสฺสน กาลํ มญฺญสีติ. อถ โข สจฺจโก นิคณฺฐปุตฺโต ภควโต ภาสิตํ อภินนฺทิตฺวา อนุโมทิตฺวา อุฏฺฐายาสนา ปกฺกามีติ.}{มหา\-สจฺจก\-สุตฺตํ นิฏฺฐิตํ ฉฏฺฐํ.}
\csromanmatikaanchor{mtk.69}
\csromantocmarkref{subsection}{7. จูฬตณฺหาสงฺขยสุตฺต}{mtk.69}
\bookmark[dest={mtk.69},level=2]{7. จูฬตณฺหาสงฺขยสุตฺต}
\csromanheader{2}{7. จูฬตณฺหาสงฺขย\-สุตฺต}
\proseitem{390}{\csromanfolio{318}เอวํ เม สุตํ\csromandash{}เอกํ สมยํ ภควา สาวตฺถิยํ วิหรติ ปุพฺพาราเม มิคาร\-มาตุ\-ปาสาเท. อถ โข สกฺโก เทวานมินฺโท เยน ภควา เตนุปสงฺกมิ, อุปสงฺกมิตฺวา ภควนฺตํ อภิวาเทตฺวา เอกมนฺตํ อฏฺฐาสิ, เอกมนฺตํ ฐิโต โข สกฺโก เทวานมินฺโท ภควนฺตํ เอตทโวจ “กิตฺตาวตา นุ โข ภนฺเต ภิกฺขุ สํขิตฺเตน ตณฺหาสงฺขย\-วิมุตฺโต โหติ อจฺจนฺตนิฏฺโฐ อจฺจนฺตโย คกฺเขมี อจฺจนฺต\-พฺรหฺมจารี อจฺจนฺต\-ปริโยสาโน เสฏฺโฐ เทวมนุสฺสานนฺ”ติ.}

\prose{อิธ เทวานมินฺท ภิกฺขุโน สุตํ โหติ “สพฺเพ ธมฺมา นาลํ อภินิเวสายา”ติ, เอวํ เจตํ เทวานมินฺท ภิกฺขุโน สุตํ โหติ “สพฺเพ ธมฺมา นาลํ อภินิเวสายา”ติ. โส สพฺพํ ธมฺมํ อภิชานาติ, สพฺพํ ธมฺมํ อภิญฺญาย สพฺพํ ธมฺมํ ปริชานาติ, สพฺพํ ธมฺมํ ปริญฺญาย ยํ กิญฺจิ เวทนํ เวเทติ สุขํ วา ทุกฺขํ วา อทุกฺขมสุขํ วา. โส ตาสุ เวทนาสุ อนิจฺจานุปสฺสี วิหรติ, วิราคานุปสฺสี วิหรติ, นิโรธานุปสฺสี วิหรติ, ปฏินิสฺสคฺคา\-นุปสฺสี วิหรติ. โส ตาสุ เวทนาสุ อนิจฺจานุปสฺสี วิหรนฺโต วิราคานุปสฺสี วิหรนฺโต นิโรธานุปสฺสี วิหรนฺโต ปฏินิสฺสคฺคา\-นุปสฺสี วิหรนฺโต น กิญฺจิ โลเก อุปาทิยติ, อนุปาทิยํ น ปริตสฺสติ, อปริตสฺสํ ปจฺจตฺตญฺเญว ปรินิพฺพายติ, “ขีณา ชาติ, วุสิตํ พฺรหฺมจริยํ, กตํ กรณียํ, นาปรํ อิตฺถตฺตายา”ติ ปชานาติ. เอตฺตาวตา โข เทวานมินฺท ภิกฺขุ สํขิตฺเตน ตณฺหาสงฺขย\-วิมุตฺโต โหติ อจฺจนฺตนิฏฺโฐ อจฺจนฺต\-โยคกฺเขมี อจฺจนฺต\-พฺรหฺมจารี อจฺจนฺต\-ปริโยสาโน เสฏฺโฐ เทว\-มนุสฺสา\-นนฺติ.}

\prose{อถ โข สกฺโก เทวานมินฺโท ภควโต ภาสิตํ อภินนฺทิตฺวา อนุโมทิตฺวา ภควนฺตํ อภิวาเทตฺวา ปทกฺขิณํ กตฺวา ตตฺเถ\-ว\-นฺตรธายิ.}

\proseitem{391}{เตน โข ปน สมเยน อายสฺมา มหาโมคฺคลฺลาโน ภควโต อวิทูเร นิสินฺโน โหติ. อถ โข อายสฺมโต มหา\-โมคฺคลฺลานสฺส เอตทโหสิ “กิํ นุ โข โส ยกฺโข ภควโต ภาสิตํ อภิสเมจฺจ อนุโมทิ, อุทาหุ โน. ยํนูนาหํ ตํ ยกฺขํ ชาเนยฺยํ ยทิ วา โส ยกฺโข ภควโต ภาสิตํ อภิสเมจฺจ \csromanfolio{319}อนุโมทิ, ยทิ วา โน”ติ. อถ โข อายสฺมา มหาโมคฺคลฺลาโน เสยฺยถาปิ นาม พลวา ปุริโส สมิญฺชิตํ วา พาหํ ปสาเรยฺย, ปสาริตํ วา พาหํ สมิญฺเชยฺย. เอวเมว ปุพฺพาราเม มิคาร\-มาตุ\-ปาสาเท อนฺตรหิโต เทเวสุ ตาวติํเสสุ ปาตุรโหสิ. เตน โข ปน สมเยน สกฺโก เทวานมินฺโท เอกปุณฺฑรีเก อุยฺยาเน ทิพฺเพหิ ปญฺจหิ ตูริยสเตหิ\csromansharedfootnote{419239d9907aebbf}{ตุริยสเตหิ (สี, สฺยา, กํ, อิ)} สมปฺปิโต สมงฺคีภูโต ปริจาเรติ. อทฺทสา โข สกฺโก เทวานมินฺโท อายสฺมนฺตํ มหา\-โมคฺคลฺลานํ ทูรโตว อาคจฺฉนฺตํ, ทิสฺวาน ตานิ ทิพฺพานิ ปญฺจ ตูริยสตานิ ปฏิปฺปณาเมตฺวา เยนายสฺมา มหาโมคฺคลฺลาโน เตนุปสงฺกมิ, อุปสงฺกมิตฺวา อายสฺมนฺตํ มหา\-โมคฺคลฺลานํ เอตทโวจ “เอหิ โข มาริส โมคฺคลฺลาน, สฺวาคตํ มาริส โมคฺคลฺลาน, จิรสฺสํ โข มาริส โมคฺคลฺลาน อิมํ ปริยายํ อกาสิ, ยทิทํ อิธาคมนาย. นิสีท มาริส โมคฺคลฺลาน, อิทมาสนํ ปญฺญตฺตนฺ”ติ. นิสีทิ โข อายสฺมา มหาโมคฺคลฺลาโน ปญฺญตฺเต อาสเน. สกฺโกปิ โข เทวานมินฺโท อญฺญตรํ นีจํ อาสนํ คเหตฺวา เอกมนฺตํ นิสีทิ. เอกมนฺตํ นิสินฺนํ โข สกฺกํ เทวานมินฺทํ อายสฺมา มหาโมคฺคลฺลาโน เอตทโวจ “ยถา กถํ ปน โข โกสิย ภควา สํขิตฺเตน ตณฺหาสงฺขย\-วิมุตฺติํ อภาสิ, สาธุ มยมฺปิ เอติสฺสา กถาย ภาคิโน อสฺสาม สวนายา”ติ.}

\proseitem{392}{มยํ โข มาริส โมคฺคลฺลาน พหุกิจฺจา พหุกรณียา อปฺเปว สเกน กรณีเยน, อปิ จ เทวานํเยว ตาวติํสานํ กรณีเยน. อปิ จ มาริส โมคฺคลฺลาน สุสฺสุตํเยว โหติ สุคฺคหิตํ สุมนสิกตํ สูปธาริตํ, ยํ โน ขิปฺปเมว อนฺตรธายติ. ภูตปุพฺพํ มาริส โมคฺคลฺลาน เทวา\-สุร\-สงฺคาโม สมุปพฺยูโฬฺห\csromansharedfootnote{90e981cf930b81dd}{สมูปพฺยุโฬฺห (สฺยา, กํ), สมูปพฺพูโฬฺห (สี)} อโหสิ. ตสฺมิํ โข ปน มาริส โมคฺคลฺลาน สงฺคาเม เทวา ชินิํสุ อสุรา ปราชินิํสุ. โส โข อหํ มาริส โมคฺคลฺลาน ตํ สงฺคามํ อภิวิชินิตฺวา วิชิตสงฺคาโม, ตโต ปฏินิวตฺติตฺวา เวชยนฺตํ นาม ปาสาทํ มาเปสิํ. เวชยนฺตสฺส โข มาริส โมคฺคลฺลาน ปาสาทสฺส เอกสตํ นิยฺยูหํ, เอเกกสฺมิํ นิยฺยูเห สตฺตสตฺต กูฏาคารสตานิ, เอกเมกสฺมิํ กูฏาคาเร สตฺตสตฺต อจฺฉราโย, เอกเมกิสฺสา อจฺฉราย สตฺตสตฺต ปริจาริกาโย. อิจฺเฉยฺยาสิ โน \csromanfolio{320}ตฺวํ มาริส โมคฺคลฺลาน เวชยนฺตสฺส ปาสาทสฺส รามเณยฺยกํ ทฏฺฐุนฺติ. อธิวาเสสิ โข อายสฺมา มหาโมคฺคลฺลาโน ตุณฺหีภาเวน.}

\proseitem{393}{อถ โข สกฺโก จ เทวานมินฺโท เวสฺสวโณ จ มหาราชา อายสฺมนฺตํ มหา\-โมคฺคลฺลานํ ปุรกฺขตฺวา เยน เวชยนฺโต ปาสาโท เตนุ\-ปสงฺกมิํสุ. อทฺทสํสุ โข สกฺกสฺส เทวานมินฺทสฺส ปริจาริกาโย อายสฺมนฺตํ มหา\-โมคฺคลฺลานํ ทูรโตว อาคจฺฉนฺตํ, ทิสฺวา โอตฺตปฺปมานา หิรียมานา สกํ สกํ โอวรกํ ปวิสิํสุ, เสยฺยถาปิ นาม สุณิสา สสุรํ ทิสฺวา โอตฺตปฺปติ หิรียติ. เอวเมว สกฺกสฺส เทวานมินฺทสฺส ปริจาริกาโย อายสฺมนฺตํ มหา\-โมคฺคลฺลานํ ทิสฺวา โอตฺตปฺปมานา หิรียมานา สกํ สกํ โอวรกํ ปวิสิํสุ. อถ โข สกฺโก จ เทวานมินฺโท เวสฺสวโณ จ มหาราชา อายสฺมนฺตํ มหา\-โมคฺคลฺลานํ เวชยนฺเต ปาสาเท อนุจงฺกมาเปนฺติ อนุวิจราเปนฺติ. อิทมฺปิ มาริส โมคฺคลฺลาน ปสฺส เวชยนฺตสฺส ปาสาทสฺส รามเณยฺยกํ, อิทมฺปิ มาริส โมคฺคลฺลาน ปสฺส เวชยนฺตสฺส ปาสาทสฺส รามเณยฺยกนฺติ. โสภติ อิทํ อายสฺมโต โกสิยสฺส, ยถา ตํ ปุพฺเพ กตปุญฺญสฺส. มนุสฺสาปิ กิญฺจิเทว รามเณยฺยกํ ทิสฺวา\csromansharedfootnote{8aacddaeae27f9eb}{ทิฏฺฐา (สี, อิ, ก)} เอวมาหํสุ “โสภติ วต โภ ยถา เทวานํ ตาวติํสานนฺ”ติ. ตยิทํ อายสฺมโต โกสิยสฺส โสภติ, ยถา ตํ ปุพฺเพ กตปุญฺญสฺสาติ. อถ โข อายสฺมโต มหา\-โมคฺคลฺลานสฺส เอตทโหสิ “อติพาฬฺหํ โข อยํ ยกฺโข ปมตฺโต วิหรติ. ยํนูนาหํ อิมํ ยกฺขํ สํเวเชยฺยนฺ”ติ. อถ โข อายสฺมา มหาโมคฺคลฺลาโน ตถารูปํ อิทฺธา\-ภิสงฺขารํ อภิสงฺขาสิ\csromansharedfootnote{055a7ac747f41cab}{อภิสงฺขาเรสิ (ก), อภิสงฺขาเรติ (สฺยา, กํ)}. ยถา เวชยนฺตํ ปาสาทํ ปาท\-งฺคุฏฺฐเกน สงฺกมฺเปสิ สมฺปกมฺเปสิ สมฺปเวเธสิ. อถ โข สกฺโก จ เทวานมินฺโท เวสฺสวโณ จ มหาราชา เทวา จ ตาวติํสา อจฺฉริย\-พฺภุต\-จิตฺต\-ชาตา อเหสุํ “อจฺฉริยํ วต โภ, อพฺภุตํ วต โภ สมณสฺส มหิทฺธิกตา มหานุภาวตา. ยตฺร หิ นาม ทิพฺพภวนํ ปาท\-งฺคุฏฺฐเกน สงฺกมฺเปสฺสติ สมฺปกมฺเปสฺสติ สมฺปเวเธสฺสตี”ติ. อถ โข อายสฺมา มหาโมคฺคลฺลาโน สกฺกํ เทวานมินฺทํ สํวิคฺคํ โลมหฏฺฐ\-ชาตํ วิทิตฺวา สกฺกํ เทวานมินฺทํ เอตทโวจ “ยถา กถํ ปน โข โกสิย ภควา สํขิตฺเตน ตณฺหาสงฺขย\-วิมุตฺติํ อภาสิ, สาธุ มยมฺปิ เอติสฺสา กถาย ภาคิโน อสฺสาม สวนายา”ติ.}

\proseitem{394}{\csromanfolio{321}อิธาหํ มาริส โมคฺคลฺลาน เยน ภควา เตนุปสงฺกมิํ, อุปสงฺกมิตฺวา ภควนฺตํ อภิวาเทตฺวา เอกมนฺตํ อฏฺฐาสิํ, เอกมนฺตํ ฐิโต โข อหํ มาริส โมคฺคลฺลาน ภควนฺตํ เอตทโวจํ “กิตฺตาวตา นุ โข ภนฺเต ภิกฺขุ สํขิตฺเตน ตณฺหาสงฺขย\-วิมุตฺโต โหติ อจฺจนฺตนิฏฺโฐ อจฺจนฺต\-โยคกฺเขมี อจฺจนฺต\-พฺรหฺมจารี อจฺจนฺต\-ปริโยสาโน เสฏฺโฐ เทวมนุสฺสานนฺ”ติ. เอวํ วุตฺเต มาริส โมคฺคลฺลาน ภควา มํ เอตทโวจ “อิธ เทวานมินฺท ภิกฺขุโน สุตํ โหติ ‘สพฺเพ ธมฺมา นาลํ อภินิเวสายา’ติ. เอวํ เจตํ เทวานมินฺท ภิกฺขุโน สุตํ โหติ ‘สพฺเพ ธมฺมา นาลํ อภินิเวสายา’ติ, โส สพฺพํ ธมฺมํ อภิชานาติ, สพฺพํ ธมฺมํ อภิญฺญาย สพฺพํ ธมฺมํ ปริชานาติ, สพฺพํ ธมฺมํ ปริญฺญาย ยํ กิญฺจิ เวทนํ เวเทติ สุขํ วา ทุกฺขํ วา อทุกฺขมสุขํ วา. โส ตาสุ เวทนาสุ อนิจฺจานุปสฺสี วิหรติ, วิราคานุปสฺสี วิหรติ, นิโรธานุปสฺสี วิหรติ, ปฏินิสฺสคฺคา\-นุปสฺสี วิหรติ. โส ตาสุ เวทนาสุ อนิจฺจานุปสฺสี วิหรนฺโต วิราคานุปสฺสี วิหรนฺโต นิโรธานุปสฺสี วิหรนฺโต ปฏินิสฺสคฺคา\-นุปสฺสี วิหรนฺโต น กิญฺจิ โลเก อุปาทิยติ, อนุปาทิยํ น ปริตสฺสติ, อปริตสฺสํ ปจฺจตฺตญฺเญว ปรินิพฺพายติ, ‘ขีณา ชาติ, วุสิตํ พฺรหฺมจริยํ, กตํ กรณียํ, นาปรํ อิตฺถตฺตายา’ติ ปชานาติ. เอตฺตาวตา โข เทวานมินฺท ภิกฺขุ สํขิตฺเตน ตณฺหาสงฺขย\-วิมุตฺโต โหติ อจฺจนฺตนิฏฺโฐ อจฺจนฺต\-โยคกฺเขมี อจฺจนฺต\-พฺรหฺมจารี อจฺจนฺต\-ปริโยสาโน เสฏฺโฐ เทวมนุสฺสานนฺ”ติ. เอวํ โข เม มาริส โมคฺคลฺลาน ภควา สํขิตฺเตน ตณฺหาสงฺขย\-วิมุตฺติํ อภาสีติ.}

\prose{อถ โข อายสฺมา มหาโมคฺคลฺลาโน สกฺกสฺส เทวานมินฺทสฺส ภาสิตํ อภินนฺทิตฺวา อนุโมทิตฺวา เสยฺยถาปิ นาม พลวา ปุริโส สมิญฺชิตํ วา พาหํ ปสาเรยฺย, ปสาริตํ วา พาหํ สมิญฺเชยฺย. เอวเมว เทเวสุ ตาวติํเสสุ อนฺตรหิโต ปุพฺพาราเม มิคาร\-มาตุ\-ปาสาเท ปาตุรโหสิ. อถ โข สกฺกสฺส เทวานมินฺทสฺส ปริจาริกาโย อจิรปกฺกนฺเต อายสฺมนฺเต มหาโมคฺคลฺลาเน สกฺกํ เทวานมินฺทํ เอตทโวจุํ “เอโส นุ เต มาริส โส ภควา สตฺถา”ติ. น โข เม มาริส โส ภควา สตฺถา, สพฺรหฺมจารี เม เอโส อายสฺมา มหา\-โมคฺคลฺลาโนติ. ลาภา เต มาริส, (สุลทฺธํ เต มาริส)\csromansharedfootnote{987ccb21c4be8712}{( ) นตฺถิ (สี, อิ)}, ยสฺส เต สพฺรหฺมจาริ เอวํมหิทฺธิโก เอวํ\-มหา\-นุภาโว, อโห นูน เต โส ภควา สตฺถาติ.}

\proseitem{395}{\csromanfolio{322}อถ โข อายสฺมา มหาโมคฺคลฺลาโน เยน ภควา เตนุปสงฺกมิ, อุปสงฺกมิตฺวา ภควนฺตํ อภิวาเทตฺวา เอกมนฺตํ นิสีทิ, เอกมนฺตํ นิสินฺโน โข อายสฺมา มหาโมคฺคลฺลาโน ภควนฺตํ เอตทโวจ “อภิชานาติ โน ภนฺเต ภควา อหุ\csromansharedfootnote{c0c04961c7289aa5}{อหุนญฺเญว (สี, สฺยา, กํ)} ญาตญฺญตรสฺส มเหสกฺขสฺส ยกฺขสฺส สํขิตฺเตน ตณฺหาสงฺขย\-วิมุตฺติํ ภาสิตา”ติ\csromansharedfootnote{9db2310600433271}{อภาสิตฺถาติ (ก)}. อภิชานามหํ โมคฺคลฺลาน, อิธ สกฺโก เทวานมินฺโท เยนาหํ เตนุปสงฺกมิ, อุปสงฺกมิตฺวา มํ อภิวาเทตฺวา เอกมนฺตํ อฏฺฐาสิ, เอกมนฺตํ ฐิโต โข โมคฺคลฺลาน สกฺโก เทวานมินฺโท มํ เอตทโวจ “กิตฺตาวตา นุ โข ภนฺเต ภิกฺขุ สํขิตฺเตน ตณฺหาสงฺขย\-วิมุตฺโต โหติ อจฺจนฺตนิฏฺโฐ อจฺจนฺต\-โยคกฺเขมี อจฺจนฺต\-พฺรหฺมจารี อจฺจนฺต\-ปริโยสาโน เสฏฺโฐ เทวมนุสฺสานนฺ”ติ. เอวํ วุตฺเต อหํ โมคฺคลฺลาน สกฺกํ เทวานมินฺทํ เอตทโวจํ “อิธ เทวานมินฺท ภิกฺขุโน สุตํ โหติ ‘สพฺเพ ธมฺมา นาลํ อภินิเวสายา’ติ. เอวํ เจตํ เทวานมินฺท ภิกฺขุโน สุตํ โหติ ‘สพฺเพ ธมฺมา นาลํ อภินิเวสายา’ติ. โส สพฺพํ ธมฺมํ อภิชานาติ, สพฺพํ ธมฺมํ อภิญฺญาย สพฺพํ ธมฺมํ ปริชานาติ, สพฺพํ ธมฺมํ ปริญฺญาย ยํ กิญฺจิ เวทนํ เวเทติ สุขํ วา ทุกฺขํ วา อทุกฺขมสุขํ วา. โส ตาสุ เวทนาสุ อนิจฺจานุปสฺสี วิหรติ, วิราคานุปสฺสี วิหรติ, นิโรธานุปสฺสี วิหรติ, ปฏินิสฺสคฺคา\-นุปสฺสี วิหรติ. โส ตาสุ เวทนาสุ อนิจฺจานุปสฺสี วิหรนฺโต วิราคานุปสฺสี วิหรนฺโต นิโรธานุปสฺสี วิหรนฺโต ปฏินิสฺสคฺคา\-นุปสฺสี วิหรนฺโต น กิญฺจิ โลเก อุปาทิยติ, อนุปาทิยํ น ปริตสฺสติ, อปริตสฺสํ ปจฺจตฺตญฺเญว ปรินิพฺพายติ, ‘ขีณา ชาติ, วุสิตํ พฺรหฺมจริยํ, กตํ กรณียํ, นาปรํ อิตฺถตฺตายา’ติ ปชานาติ. เอตฺตาวตา โข เทวานมินฺท ภิกฺขุ สํขิตฺเตน ตณฺหาสงฺขย\-วิมุตฺโต โหติ อจฺจนฺตนิฏฺโฐ อจฺจนฺต\-โยคกฺเขมี อจฺจนฺต\-พฺรหฺมจารี อจฺจนฺต\-ปริโยสาโน เสฏฺโฐ เทวมนุสฺสานนฺ”ติ. เอวํ โข อหํ โมคฺคลฺลาน อภิชานามิ “สกฺกสฺส เทวานมินฺทสฺส สํขิตฺเตน ตณฺหาสงฺขย\-วิมุตฺติํ ภาสิตา”ติ\csromansharedfootnote{9db2310600433271}{อภาสิตฺถาติ (ก)}.}

\prosewithcloser{อิทมโวจ ภควา. อตฺตมโน อายสฺมา มหาโมคฺคลฺลาโน ภควโต ภาสิตํ อภินนฺทีติ.}{จูฬตณฺหาสงฺขย\-สุตฺตํ นิฏฺฐิตํ สตฺตมํ.}
\csromanmatikaanchor{mtk.70}
\csromantocmarkref{subsection}{8. มหาตณฺหาสงฺขยสุตฺต}{mtk.70}
\bookmark[dest={mtk.70},level=2]{8. มหาตณฺหาสงฺขยสุตฺต}
\csromanmarkh{8. มหา\-ตณฺหาสงฺขย\-สุตฺต (38)}\csromanmarkhii{8. มหา\-ตณฺหาสงฺขย\-สุตฺต}\csromanheadernomark{2}{8. มหา\-ตณฺหาสงฺขย\-สุตฺต (38) 8. มหา\-ตณฺหาสงฺขย\-สุตฺต}
\proseitem{396}{\csromanfolio{323}เอวํ เม สุตํ\csromandash{}เอกํ สมยํ ภควา สาวตฺถิยํ วิหรติ เชตวเน อนาถ\-ปิณฺฑิกสฺส อาราเม. เตน โข ปน สมเยน สาติสฺส นาม ภิกฺขุโน เกวฏฺฏ\-ปุตฺตสฺส เอวรูปํ ปาปกํ ทิฏฺฐิคตํ อุปฺปนฺนํ โหติ “ตถาหํ ภควตา ธมฺมํ เทสิตํ อาชานามิ, ยถา ตเทวิทํ วิญฺญาณํ สนฺธาวติ สํสรติ อนญฺญนฺ”ติ. อสฺโสสุํ โข สมฺพหุลา ภิกฺขู สาติสฺส กิร นาม ภิกฺขุโน เกวฏฺฏ\-ปุตฺตสฺส เอวรูปํ ปาปกํ ทิฏฺฐิคตํ อุปฺปนฺนํ “ตถาหํ ภควตา ธมฺมํ เทสิตํ อาชานามิ, ยถา ตเทวิทํ วิญฺญาณํ สนฺธาวติ สํสรติ อนญฺญนฺ”ติ. อถ โข เต ภิกฺขู เยน สาติ ภิกฺขุ เกวฏฺฏปุตฺโต เตนุ\-ปสงฺกมิํสุ, อุปสงฺกมิตฺวา สาติํ ภิกฺขุํ เกวฏฺฏปุตฺตํ เอตทโวจุํ “สจฺจํ กิร เต อาวุโส สาติ เอวรูปํ ปาปกํ ทิฏฺฐิคตํ อุปฺปนฺนํ ‘ตถาหํ ภควตา ธมฺมํ เทสิตํ อาชานามิ, ยถา ตเทวิทํ วิญฺญาณํ สนฺธาวติ สํสรติ อนญฺญนฺ’ติ”. เอวํพฺยา โข อหํ อาวุโส ภควตา ธมฺมํ เทสิตํ อาชานามิ, ยถา ตเทวิทํ วิญฺญาณํ สนฺธาวติ สํสรติ อนญฺญนฺติ. อถ โข เต ภิกฺขู สาติํ ภิกฺขุํ เกวฏฺฏปุตฺตํ เอตสฺมา ปาปกา ทิฏฺฐิคตา วิเวเจตุกามา สมนุยุญฺชนฺติ สมนุคาหนฺติ สมนุภาสนฺติ “มา เอวํ อาวุโส สาติ อวจ, มา ภควนฺตํ อพฺภาจิกฺขิ, น หิ สาธุ ภควโต อพฺภกฺขานํ, น หิ ภควา เอวํ วเทยฺย, อเนกปริยาเยนา\-วุโส สาติ ปฏิจฺจ\-สมุปฺปนฺนํ วิญฺญาณํ วุตฺตํ ภควตา ‘อญฺญตฺร ปจฺจยา นตฺถิ วิญฺญาณสฺส สมฺภโว’ติ”. เอวมฺปิ โข สาติ ภิกฺขุ เกวฏฺฏปุตฺโต เตหิ ภิกฺขูหิ สมนุ\-ยุญฺชิย\-มาโน สมนุ\-คาหิย\-มาโน สมนุ\-ภาสิย\-มาโน ตเทว ปาปกํ ทิฏฺฐิคตํ ถามสา ปรามาสา อภินิวิสฺส โวหรติ “เอวํพฺยา โข อหํ อาวุโส ภควตา ธมฺมํ เทสิตํ อาชานามิ, ยถา ตเทวิทํ วิญฺญาณํ สนฺธาวติ สํสรติ อนญฺญนฺ”ติ.}

\proseitem{397}{ยโต โข เต ภิกฺขู นาสกฺขิํสุ สาติํ ภิกฺขุํ เกวฏฺฏปุตฺตํ เอตสฺมา ปาปกา ทิฏฺฐิคตา วิเวเจตุํ. อถ โข เต ภิกฺขู เยน ภควา เตนุ\-ปสงฺกมิํสุ, อุปสงฺกมิตฺวา ภควนฺตํ อภิวาเทตฺวา เอกมนฺตํ นิสีทิํสุ, เอกมนฺตํ นิสินฺนา โข เต ภิกฺขู ภควนฺตํ เอตทโวจุํ\csromandash{} สาติสฺส นาม ภนฺเต ภิกฺขุโน เกวฏฺฏ\-ปุตฺตสฺส เอวรูปํ ปาปกํ ทิฏฺฐิคตํ \csromanfolio{324}อุปฺปนฺนํ “ตถาหํ ภควตา ธมฺมํ เทสิตํ อาชานามิ, ยถา ตเทวิทํ วิญฺญาณํ สนฺธาวติ สํสรติ อนญฺญนฺ”ติ. อสฺสุมฺห โข มยํ ภนฺเต สาติสฺส กิร นาม ภิกฺขุโน เกวฏฺฏ\-ปุตฺตสฺส เอวรูปํ ปาปกํ ทิฏฺฐิคตํ อุปฺปนฺนํ “ตถาหํ ภควตา ธมฺมํ เทสิตํ อาชานามิ, ยถา ตเทวิทํ วิญฺญาณํ สนฺธาวติ สํสรติ อนญฺญนฺ”ติ. อถ โข มยํ ภนฺเต เยน สาติ ภิกฺขุ เกวฏฺฏปุตฺโต เตนุ\-ปสงฺกมิมฺห, อุปสงฺกมิตฺวา สาติํ ภิกฺขุํ เกวฏฺฏปุตฺตํ เอตทโวจุมฺห “สจฺจํ กิร เต อาวุโส สาติ เอวรูปํ ปาปกํ ทิฏฺฐิคตํ อุปฺปนฺนํ ‘ตถาหํ ภควตา ธมฺมํ เทสิตํ อาชานามิ, ยถา ตเทวิทํ วิญฺญาณํ สนฺธาวติ สํสรติ อนญฺญนฺ’ติ”. เอวํ วุตฺเต ภนฺเต สาติ ภิกฺขุ เกวฏฺฏปุตฺโต อเมฺห เอตทโวจ “เอวํพฺยา โข อหํ อาวุโส ภควตา ธมฺมํ เทสิตํ อาชานามิ, ยถา ตเทวิทํ วิญฺญาณํ สนฺธาวติ สํสรติ อนญฺญนฺ”ติ. อถ โข มยํ ภนฺเต สาติํ ภิกฺขุํ เกวฏฺฏปุตฺตํ เอตสฺมา ปาปกา ทิฏฺฐิคตา วิเวเจตุกามา สมนุยุญฺชิมฺห สมนุคาหิมฺห สมนุภาสิมฺห “มา เอวํ อาวุโส สาติ อวจ, มา ภควนฺตํ อพฺภาจิกฺขิ, น หิ สาธุ ภควโต อพฺภกฺขานํ, น หิ ภควา เอวํ วเทยฺย, อเนกปริยาเยนา\-วุโส สาติ ปฏิจฺจ\-สมุปฺปนฺนํ วิญฺญาณํ วุตฺตํ ภควตา ‘อญฺญตฺร ปจฺจยา นตฺถิ วิญฺญาณสฺส สมฺภโว’ติ”. เอวมฺปิ โข ภนฺเต สาติ ภิกฺขุ เกวฏฺฏปุตฺโต อเมฺหหิ สมนุ\-ยุญฺชิย\-มาโน สมนุ\-คาหิย\-มาโน สมนุ\-ภาสิย\-มาโน ตเทว ปาปกํ ทิฏฺฐิคตํ ถามสา ปรามสา อภินิวิสฺส โวหรติ “เอวํพฺยา โข อหํ อาวุโส ภควตา ธมฺมํ เทสิตํ อาชานามิ, ยถา ตเทวิทํ วิญฺญาณํ สนฺธาวติ สํสรติ อนญฺญนฺ”ติ. ยโต โข มยํ ภนฺเต นาสกฺขิมฺห สาติํ ภิกฺขุํ เกวฏฺฏปุตฺตํ เอตสฺมา ปาปกา ทิฏฺฐิคตา วิเวเจตุํ. อถ มยํ เอตมตฺถํ ภควโต อาโรเจมาติ.}

\proseitem{398}{อถ โข ภควา อญฺญตรํ ภิกฺขุํ อามนฺเตสิ “เอหิ ตฺวํ ภิกฺขุ มม วจเนน สาติํ ภิกฺขุํ เกวฏฺฏปุตฺตํ อามนฺเตหิ ‘สตฺถา ตํ อาวุโส สาติ อามนฺเตตี’ติ”. “เอวํ ภนฺเต”ติ โข โส ภิกฺขุ ภควโต ปฏิสฺสุตฺวา เยน สาติ ภิกฺขุ เกวฏฺฏปุตฺโต เตนุปสงฺกมิ, อุปสงฺกมิตฺวา สาติํ ภิกฺขุํ เกวฏฺฏปุตฺตํ เอตทโวจ “สตฺถา ตํ อาวุโส สาติ อามนฺเตตี”ติ. “เอวมาวุโส”ติ โข สาติ ภิกฺขุ เกวฏฺฏปุตฺโต ตสฺส ภิกฺขุโน ปฏิสฺสุตฺวา เยน ภควา เตนุปสงฺกมิ, อุปสงฺกมิตฺวา \csromanfolio{325}ภควนฺตํ อภิวาเทตฺวา เอกมนฺตํ นิสีทิ, เอกมนฺตํ นิสินฺนํ โข สาติํ ภิกฺขุํ เกวฏฺฏปุตฺตํ ภควา เอตทโวจ “สจฺจํ กิร เต สาติ เอวรูปํ ปาปกํ ทิฏฺฐิคตํ อุปฺปนฺนํ ‘ตถาหํ ภควตา ธมฺมํ เทสิตํ อาชานามิ, ยถา ตเทวิทํ วิญฺญาณํ สนฺธาวติ สํสรติ อนญฺญนฺ’ติ”. เอวํพฺยา โข อหํ ภนฺเต ภควตา ธมฺมํ เทสิตํ อาชานามิ, ยถา ตเทวิทํ วิญฺญาณํ สนฺธาวติ สํสรติ อนญฺญนฺติ. กตมํ ตํ สาติ วิญฺญาณนฺติ. ยฺวายํ ภนฺเต วโท เวเทยฺโย, ตตฺร ตตฺร กลฺยาณ\-ปาปกานํ กมฺมานํ วิปากํ ปฏิสํเวเทตีติ. กสฺส นุ โข นาม ตฺวํ โมฆปุริส มยา เอวํ ธมฺมํ เทสิตํ อาชานาสิ, นนุ มยา โมฆปุริส อเนก\-ปริยาเยน ปฏิจฺจ\-สมุปฺปนฺนํ วิญฺญาณํ วุตฺตํ “อญฺญตฺร ปจฺจยา นตฺถิ วิญฺญาณสฺส สมฺภโว”ติ. อถ จ ปน ตฺวํ โมฆปุริส อตฺตนา ทุคฺคหิเตน อเมฺห เจว อพฺภาจิกฺขสิ, อตฺตานญฺจ ขณสิ, พหุญฺจ อปุญฺญํ ปสวสิ. ตํ หิ เต โมฆปุริส ภวิสฺสติ ทีฆรตฺตํ อหิตาย ทุกฺขายาติ.}

\proseitem{399}{อถ โข ภควา ภิกฺขู อามนฺเตสิ “ตํ กิํ มญฺญถ ภิกฺขเว, อปิ นายํ สาติ ภิกฺขุ เกวฏฺฏปุตฺโต อุสฺมีกโตปิ อิมสฺมิํ ธมฺมวินเย”ติ. กิํ หิ สิยา ภนฺเต, โน เหตํ ภนฺเตติ. เอวํ วุตฺเต สาติ ภิกฺขุ เกวฏฺฏปุตฺโต ตุณฺหีภูโต มงฺกุภูโต ปตฺตกฺขนฺโธ อโธมุโข ปชฺฌายนฺโต อปฺปฏิภาโน นิสีทิ. อถ โข ภควา สาติํ ภิกฺขุํ เกวฏฺฏปุตฺตํ ตุณฺหีภูตํ มงฺกุภูตํ ปตฺตกฺขนฺธํ อโธมุขํ ปชฺฌายนฺตํ อปฺปฏิภานํ วิทิตฺวา สาติํ ภิกฺขุํ เกวฏฺฏปุตฺตํ เอตทโวจ “ปญฺญายิสฺสสิ โข ตฺวํ โมฆปุริส เอเตน สเกน ปาปเกน ทิฏฺฐิคเตน. อิธาหํ ภิกฺขู ปฏิปุจฺฉิสฺสามี”ติ. อถ โข ภควา ภิกฺขู อามนฺเตสิ “ตุเมฺหปิ เม ภิกฺขเว เอวํ ธมฺมํ เทสิตํ อาชานาถ, ยถายํ สาติ ภิกฺขุ เกวฏฺฏปุตฺโต อตฺตนา ทุคฺคหิเตน อเมฺห เจว อพฺภาจิกฺขติ, อตฺตานญฺจ ขณติ, พหุญฺจ อปุญฺญํ ปสวตี”ติ. โน เหตํ ภนฺเต. อเนก\-ปริยาเยน หิ โน ภนฺเต ปฏิจฺจ\-สมุปฺปนฺนํ วิญฺญาณํ วุตฺตํ ภควตา “อญฺญตฺร ปจฺจยา นตฺถิ วิญฺญาณสฺส สมฺภโว”ติ, สาธุ สาธุ ภิกฺขเว, สาธุ โข เม ตุเมฺห ภิกฺขเว เอวํ ธมฺมํ เทสิตํ อาชานาถ, อเนก\-ปริยาเยน หิ โว ภิกฺขเว ปฏิจฺจ\-สมุปฺปนฺนํ วิญฺญาณํ วุตฺตํ มยา “อญฺญตฺร ปจฺจยา นตฺถิ วิญฺญาณสฺส สมฺภโว”ติ. อถ จ ปนายํ สาติ ภิกฺขุ เกวฏฺฏปุตฺโต อตฺตนา ทุคฺคหิเตน อเมฺห เจว อพฺภาจิกฺขติ, อตฺตานญฺจ ขณติ, พหุญฺจ \csromanfolio{326}อปุญฺญํ ปสวติ. ตํ หิ ตสฺส โมฆปุริสสฺส ภวิสฺสติ ทีฆรตฺตํ อหิตาย ทุกฺขาย.}

\proseitem{400}{ยํ ยเทว ภิกฺขเว ปจฺจยํ ปฏิจฺจ อุปฺปชฺชติ วิญฺญาณํ, เตน เตเนว วิญฺญาณํเตฺวว สงฺขฺยํ คจฺฉติ\csromansharedfootnote{cc6a29783d8db2c6}{สงฺขํ คจฺฉติ (สี, อิ)}. จกฺขุญฺจ ปฏิจฺจ รูเป จ อุปฺปชฺชติ วิญฺญาณํ, จกฺขุวิญฺญาณํเตฺวว สงฺขฺยํ คจฺฉติ. โสตญฺจ ปฏิจฺจ สทฺเท จ อุปฺปชฺชติ วิญฺญาณํ, โสตวิญฺญาณํเตฺวว สงฺขฺยํ คจฺฉติ. ฆานญฺจ ปฏิจฺจ คนฺเธ จ อุปฺปชฺชติ วิญฺญาณํ, ฆานวิญฺญาณํเตฺวว สงฺขฺยํ คจฺฉติ. ชิวฺหญฺจ ปฏิจฺจ รเส จ อุปฺปชฺชติ วิญฺญาณํ, ชิวฺหาวิญฺญาณํเตฺวว สงฺขฺยํ คจฺฉติ. กายญฺจ ปฏิจฺจ โผฏฺฐพฺเพ จ อุปฺปชฺชติ วิญฺญาณํ, กายวิญฺญาณํเตฺวว สงฺขฺยํ คจฺฉติ. มนญฺจ ปฏิจฺจ ธมฺเม จ อุปฺปชฺชติ วิญฺญาณํ, มโนวิญฺญาณํเตฺวว สงฺขฺยํ คจฺฉติ.}

\prose{เสยฺยถาปิ ภิกฺขเว ยํ ยเทว ปจฺจยํ ปฏิจฺจ อคฺคิ ชลติ, เตน เตเนว สงฺขฺยํ คจฺฉติ. กฏฺฐญฺจ ปฏิจฺจ อคฺคิ ชลติ, กฏฺฐคฺคิเตฺวว สงฺขฺยํ คจฺฉติ. สกลิกญฺจ ปฏิจฺจ อคฺคิ ชลติ, สกลิกคฺคิเตฺวว สงฺขฺยํ คจฺฉติ. ติณญฺจ ปฏิจฺจ อคฺคิ ชลติ, ติณคฺคิเตฺวว สงฺขํ คจฺฉติ. โคมยญฺจ ปฏิจฺจ อคฺคิ ชลติ, โคมยคฺคิเตฺวว สงฺขฺยํ คจฺฉติ. ถุสญฺจ ปฏิจฺจ อคฺคิ ชลติ, ถุสคฺคิเตฺวว สงฺขฺยํ คจฺฉติ. สงฺการญฺจ ปฏิจฺจ อคฺคิ ชลติ, สงฺการ\-คฺคิเตฺวว สงฺขํ คจฺฉติ. เอวเมว โข ภิกฺขเว ยํ ยเทว ปจฺจยํ ปฏิจฺจ อุปฺปชฺชติ วิญฺญาณํ, เตน เตเนว สงฺขฺยํ คจฺฉติ. จกฺขุญฺจ ปฏิจฺจ รูเป จ อุปฺปชฺชติ วิญฺญาณํ, จกฺขุวิญฺญาณํเตฺวว สงฺขฺยํ คจฺฉติ. โสตญฺจ ปฏิจฺจ สทฺเท จ อุปฺปชฺชติ วิญฺญาณํ, โสตวิญฺญาณํเตฺวว สงฺขฺยํ คจฺฉติ. ฆานญฺจ ปฏิจฺจ คนฺเธ จ อุปฺปชฺชติ วิญฺญาณํ, ฆานวิญฺญาณํเตฺวว สงฺขฺยํ คจฺฉติ. ชิวฺหญฺจ ปฏิจฺจ รเส จ อุปฺปชฺชติ วิญฺญาณํ, ชิวฺหาวิญฺญาณํเตฺวว สงฺขฺยํ คจฺฉติ. กายญฺจ ปฏิจฺจ โผฏฺฐพฺเพ จ อุปฺปชฺชติ วิญฺญาณํ, กายวิญฺญาณํเตฺวว สงฺขํ คจฺฉติ. มนญฺจ ปฏิจฺจ ธมฺเม จ อุปฺปชฺชติ วิญฺญาณํ, มโนวิญฺญาณํเตฺวว สงฺขฺยํ คจฺฉติ.}

\proseitem{401}{ภูตมิทนฺติ ภิกฺขเว ปสฺสถาติ. เอวํ ภนฺเต. ตทา\-หาร\-สมฺภวนฺติ ภิกฺขเว ปสฺสถาติ. เอวํ ภนฺเต. ตทา\-หาร\-นิโรธา ยํ ภูตํ, ตํ นิโรธ\-ธมฺมนฺติ ภิกฺขเว ปสฺสถาติ. เอวํ ภนฺเต. ภูตมิทํ โนสฺสูติ ภิกฺขเว กงฺขโต อุปฺปชฺชติ วิจิกิจฺฉาติ. เอวํ ภนฺเต. ตทา\-หาร\-สมฺภวํ โนสฺสูติ \csromanfolio{327}ภิกฺขเว กงฺขโต อุปฺปชฺชติ วิจิกิจฺฉาติ. เอวํ ภนฺเต. ตทา\-หาร\-นิโรธา ยํ ภูตํ, ตํ นิโรธธมฺมํ โนสฺสูติ ภิกฺขเว กงฺขโต อุปฺปชฺชติ วิจิกิจฺฉาติ. เอวํ ภนฺเต. ภูตมิทนฺติ ภิกฺขเว ยถาภูตํ สมฺมปฺปญฺญาย ปสฺสโต ยา วิจิกิจฺฉา, สา ปหียตีติ. เอวํ ภนฺเต. ตทา\-หาร\-สมฺภวนฺติ ภิกฺขเว ยถาภูตํ สมฺมปฺปญฺญาย ปสฺสโต ยา วิจิกิจฺฉา, สา ปหียตีติ. เอวํ ภนฺเต. ตทา\-หาร\-นิโรธา ยํ ภูตํ, ตํ นิโรธ\-ธมฺมนฺติ ภิกฺขเว ยถาภูตํ สมฺมปฺปญฺญาย ปสฺสโต ยา วิจิกิจฺฉา, สา ปหียตีติ. เอวํ ภนฺเต. ภูตมิทนฺติ ภิกฺขเว อิติปิ โว เอตฺถ นิพฺพิจิกิจฺฉาติ. เอวํ ภนฺเต. ตทา\-หาร\-สมฺภวนฺติ ภิกฺขเว อิติปิ โว เอตฺถ นิพฺพิจิกิจฺฉาติ. เอวํ ภนฺเต. ตทา\-หาร\-นิโรธา ยํ ภูตํ, ตํ นิโรธ\-ธมฺมนฺติ ภิกฺขเว อิติปิ โว เอตฺถ นิพฺพิจิกิจฺฉาติ. เอวํ ภนฺเต. ภูตมิทนฺติ ภิกฺขเว ยถาภูตํ สมฺมปฺปญฺญาย สุทิฏฺฐนฺติ. เอวํ ภนฺเต. ตทา\-หาร\-สมฺภวนฺติ ภิกฺขเว ยถาภูตํ สมฺมปฺปญฺญาย สุทิฏฺฐนฺติ. เอวํ ภนฺเต. ตทา\-หาร\-นิโรธา ยํ ภูตํ, ตํ นิโรธ\-ธมฺมนฺติ ภิกฺขเว ยถาภูตํ สมฺมปฺปญฺญาย สุทิฏฺฐนฺติ. เอวํ ภนฺเต. อิมํ เจ ตุเมฺห ภิกฺขเว ทิฏฺฐิํ เอวํ ปริสุทฺธํ เอวํ ปริโยทาตํ อลฺลีเยถ เกลาเยถ ธนาเยถ มมาเยถ, อปิ นุ เม ตุเมฺห ภิกฺขเว กุลฺลูปมํ ธมฺมํ เทสิตํ อาชาเนยฺยาถ นิตฺถรณ\-ตฺถาย, โน คหณตฺถายาติ. โน เหตํ ภนฺเต. อิมํ เจ ตุเมฺห ภิกฺขเว ทิฏฺฐิํ เอวํ ปริสุทฺธํ เอวํ ปริโยทาตํ น อลฺลีเยถ น เกลาเยถ น ธนาเยถ น มมาเยถ, อปิ นุ เม ตุเมฺห ภิกฺขเว กุลฺลูปมํ ธมฺมํ เทสิตํ อาชาเนยฺยาถ นิตฺถรณ\-ตฺถาย, โน คหณตฺถายาติ. เอวํ ภนฺเต.}

\proseitem{402}{จตฺตาโรเม ภิกฺขเว อาหารา ภูตานํ วา สตฺตานํ ฐิติยา สมฺภเวสีนํ วา อนุคฺคหาย. กตเม จตฺตาโร. กพฬีกาโร อาหาโร โอฬาริโก วา สุขุโม วา, ผสฺโส ทุติโย, มโนสญฺเจตนา ตติยา, วิญฺญาณํ จตุตฺถํ. อิเม จ ภิกฺขเว จตฺตาโร อาหารา กิํนิทานา กิํสมุทยา กิํชาติกา กิํปภวา. อิเม จตฺตาโร อาหารา ตณฺหานิทานา ตณฺหาสมุทยา ตณฺหาชาติกา ตณฺหาปภวา. ตณฺหา จายํ ภิกฺขเว กิํนิทานา กิํสมุทยา กิํชาติกา กิํปภวา. ตณฺหา เวทนานิทานา เวทนาสมุทยา วิทนาชาติกา เวทนาปภวา. เวทนา จายํ ภิกฺขเว กิํนิทานา กิํสมุทยา กิํชาติกา กิํปภวา. เวทนา ผสฺสนิทานา ผสฺสสมุทยา ผสฺสชาติกา \csromanfolio{328}ผสฺสปภวา. ผสฺโส จายํ ภิกฺขเว กิํนิทาโน กิํสมุทโย กิํชาติโก กิํปภโว. ผสฺโส สฬายต\-นนิทาโน สฬายตน\-สมุทโย สฬายตน\-ชาติโก สฬายตน\-ปภโว. สฬายตนํ จิทํ ภิกฺขเว กิํนิทานํ กิํสมุทยํ กิํชาติกํ กิํปภวํ. สฬายตนํ นามรูป\-นิทานํ นามรูป\-สมุทยํ นามรูป\-ชาติกํ นามรูป\-ปภวํ. นามรูปํ จิทํ ภิกฺขเว กิํนิทานํ กิํสมุทยํ กิํชาติกํ กิํปภวํ. นามรูปํ วิญฺญาณนิทานํ วิญฺญาณ\-สมุทยํ วิญฺญาณชาติกํ วิญฺญาณ\-ปภวํ. วิญฺญาณํ จิทํ ภิกฺขเว กิํนิทานํ กิํสมุทยํ กิํชาติกํ กิํปภวํ. วิญฺญาณํ สงฺขาร\-นิทานํ สงฺขาร\-สมุทยํ สงฺขาร\-ชาติกํ สงฺขาร\-ปภวํ. สงฺขารา จิเม ภิกฺขเว กิํนิทานา กิํสมุทยา กิํชาติกา กิํปภวา. สงฺขารา อวิชฺชานิทานา อวิชฺชาสมุทยา อวิชฺชาชาติกา อวิชฺชาปภวา. อิติ โข ภิกฺขเว อวิชฺชาปจฺจยา สงฺขารา, สงฺขาร\-ปจฺจยา วิญฺญาณํ, วิญฺญาณปจฺจยา นามรูปํ, นามรูป\-ปจฺจยา สฬายตนํ, สฬายตน\-ปจฺจยา ผสฺโส, ผสฺสปจฺจยา เวทนา, เวทนาปจฺจยา ตณฺหา, ตณฺหาปจฺจยา อุปาทานํ, อุปาทานปจฺจยา ภโว, ภวปจฺจยา ชาติ, ชาติปจฺจยา ชรามรณํ โสก\-ปริเทว\-ทุกฺข\-โทมนสฺสุ\-ปายาสา สมฺภวนฺติ, เอวเมตสฺส เกวลสฺส ทุกฺข\-กฺขนฺธสฺส สมุทโย โหติ.}

\proseitem{403}{“ชาติปจฺจยา ชรามรณนฺ”ติ อิติ โข ปเนตํ วุตฺตํ, ชาติปจฺจยา นุ โข ภิกฺขเว ชรามรณํ โน วา, กถํ วา เอตฺถ\csromansharedfootnote{31c39d6fe6bf7cbe}{กถํ วา โว เอตฺถ (?)} โหตีติ. ชาติปจฺจยา ภนฺเต ชรามรณํ, เอวํ โน เอตฺถ โหติ\csromansharedfootnote{e07b95fe32e66b2b}{เอวํ โน เอตฺถ โหตีติ (ก)} “ชาติปจฺจยา ชรามรณนฺ”ติ. “ภวปจฺจยา ชาตี”ติ อิติ โข ปเนตํ วุตฺตํ, ภวปจฺจยา นุ โข ภิกฺขเว ชาติ โน วา, กถํ วา เอตฺถ โหตีติ. ภวปจฺจยา ภนฺเต ชาติ, เอวํ โน เอตฺถ โหติ “ภวปจฺจยา ชาตี”ติ. “อุปาทานปจฺจยา ภโว”ติ อิติ โข ปเนตํ วุตฺตํ, อุปาทานปจฺจยา นุ โข ภิกฺขเว ภโว โน วา, กถํ วา เอตฺถ โหตีติ. อุปาทานปจฺจยา ภนฺเต ภโว, เอวํ โน เอตฺถ โหติ “อุปาทานปจฺจยา ภโว”ติ. “ตณฺหาปจฺจยา อุปาทานนฺ”ติ อิติ โข ปเนตํ วุตฺตํ, ตณฺหาปจฺจยา นุ โข ภิกฺขเว อุปาทานํ โน วา, กถํ วา เอตฺถ โหตีติ. ตณฺหาปจฺจยา ภนฺเต อุปาทานํ, เอวํ โน เอตฺถ โหติ “ตณฺหาปจฺจยา อุปาทานนฺ”ติ. \csromanfolio{329}“เวทนาปจฺจยา ตณฺหา”ติ อิติ โข ปเนตํ วุตฺตํ, เวทนาปจฺจยา นุ โข ภิกฺขเว ตณฺหา โน วา, กถํ วา เอตฺถ โหตีติ. เวทนาปจฺจยา ภนฺเต ตณฺหา, เอวํ โน เอตฺถ โหติ “เวทนาปจฺจยา ตณฺหา”ติ. “ผสฺสปจฺจยา เวทนา”ติ อิติ โข ปเนตํ วุตฺตํ, ผสฺสปจฺจยา นุ โข ภิกฺขเว เวทนา โน วา, กถํ วา เอตฺถ โหตีติ. ผสฺสปจฺจยา ภนฺเต เวทนา, เอวํ โน เอตฺถ โหติ “ผสฺสปจฺจยา เวทนา”ติ. “สฬายตน\-ปจฺจยา ผสฺโส”ติ อิติ โข ปเนตํ วุตฺตํ, สฬายตน\-ปจฺจยา นุ โข ภิกฺขเว ผสฺโส โน วา, กถํ วา เอตฺถ โหตีติ. สฬายตน\-ปจฺจยา ภนฺเต ผสฺโส, เอวํ โน เอตฺถ โหติ “สฬายตน\-ปจฺจยา ผสฺโส”ติ. “นามรูป\-ปจฺจยา สฬายตนนฺ”ติ อิติ โข ปเนตํ วุตฺตํ, นามรูป\-ปจฺจยา นุ โข ภิกฺขเว สฬายตนํ โน วา, กถํ วา เอตฺถ โหตีติ. นามรูป\-ปจฺจยา ภนฺเต สฬายตนํ, เอวํ โน เอตฺถ โหติ “นามรูป\-ปจฺจยา สฬายตนนฺ”ติ. “วิญฺญาณปจฺจยา นามรูปนฺ”ติ อิติ โข ปเนตํ วุตฺตํ, วิญฺญาณปจฺจยา นุ โข ภิกฺขเว นามรูปํ โน วา, กถํ วา เอตฺถ โหตีติ. วิญฺญาณปจฺจยา ภนฺเต นามรูปํ, เอวํ โน เอตฺถ โหติ “วิญฺญาณปจฺจยา นามรูปนฺ”ติ. “สงฺขาร\-ปจฺจยา วิญฺญาณนฺ”ติ อิติ โข ปเนตํ วุตฺตํ, สงฺขาร\-ปจฺจยา นุ โข ภิกฺขเว วิญฺญาณํ, โน วา, กถํ วา เอตฺถ โหตีติ. สงฺขาร\-ปจฺจยา ภนฺเต วิญฺญาณํ, เอวํ โน เอตฺถ โหติ “สงฺขาร\-ปจฺจยา วิญฺญาณนฺ”ติ. “อวิชฺชาปจฺจยา สงฺขารา”ติ อิติ โข ปเนตํ วุตฺตํ, อวิชฺชาปจฺจยา นุ โข ภิกฺขเว สงฺขารา โน วา, กถํ วา เอตฺถ โหตีติ. อวิชฺชาปจฺจยา ภนฺเต สงฺขารา, เอวํ โน เอตฺถ โหติ “อวิชฺชาปจฺจยา สงฺขารา”ติ.}

\proseitem{404}{สาธุ ภิกฺขเว. อิติ โข ภิกฺขเว ตุเมฺหปิ เอวํ วเทถ, อหมฺปิ เอวํ วทามิ, อิมสฺมิํ สติ อิทํ โหติ, อิมสฺสุปฺปาทา อิทํ อุปฺปชฺชติ. ยทิทํ อวิชฺชาปจฺจยา สงฺขารา, สงฺขาร\-ปจฺจยา วิญฺญาณํ, วิญฺญาณปจฺจยา นามรูปํ, นามรูป\-ปจฺจยา สฬายตนํ, สฬายตน\-ปจฺจยา ผสฺโส, ผสฺสปจฺจยา เวทนา, เวทนาปจฺจยา ตณฺหา, ตณฺหาปจฺจยา อุปาทานํ, อุปาทานปจฺจยา ภโว, ภวปจฺจยา ชาติ, ชาติปจฺจยา ชรามรณํ โสกปริเทว ทุกฺขโทมนสฺสุ\-ปายาสา สมฺภวนฺติ, เอวเมตสฺส เกวลสฺส ทุกฺข\-กฺขนฺธสฺส สมุทโย โหติ.}

\prose{\csromanfolio{330}อวิชฺชายเตฺวว อเสส\-วิราค\-นิโรธา สงฺขาร\-นิโรโธ, สงฺขาร\-นิโรธา วิญฺญาณนิโรโธ, วิญฺญาณนิโรธา นามรูป\-นิโรโธ, นามรูป\-นิโรธา สฬายตน\-นิโรโธ, สฬายตน\-นิโรธา ผสฺสนิโรโธ, ผสฺสนิโรธา เวทนานิโรโธ, เวทนานิโรธา ตณฺหานิโรโธ, ตณฺหานิโรธา อุปาทานนิโรโธ, อุปาทานนิโรธา ภวนิโรโธ, ภวนิโรธา ชาตินิโรโธ, ชาตินิโรธา ชรามรณํ โสก\-ปริเทว\-ทุกฺข\-โทมนสฺสุ\-ปายาสา นิรุชฺฌนฺติ, เอวเมตสฺส เกวลสฺส ทุกฺข\-กฺขนฺธสฺส นิโรโธ โหติ.}

\proseitem{405}{“ชาตินิโรธา ชรามรณ\-นิโรโธ”ติ อิติ โข ปเนตํ วุตฺตํ, ชาตินิโรธา นุ โข ภิกฺขเว ชรามรณ\-นิโรโธ โน วา, กถํ วา เอตฺถ โหตีติ. ชาตินิโรธา ภนฺเต ชรามรณ\-นิโรโธ, เอวํ โน เอตฺถ โหติ “ชาตินิโรธา ชรามรณ\-นิโรโธ”ติ. “ภวนิโรธา ชาตินิโรโธ”ติ อิติ โข ปเนตํ วุตฺตํ, ภวนิโรธา นุ โข ภิกฺขเว ชาตินิโรโธ โน วา, กถํ วา เอตฺถ โหตีติ. ภวนิโรธา ภนฺเต ชาตินิโรโธ, เอวํ โน เอตฺถ โหติ “ภวนิโรธา ชาตินิโรโธ”ติ. “อุปาทานนิโรธา ภวนิโรโธ”ติ อิติ โข ปเนตํ วุตฺตํ, อุปาทานนิโรธา นุ โข ภิกฺขเว ภวนิโรโธ โน วา, กถํ วา เอตฺถ โหตีติ. อุปาทานนิโรธา ภนฺเต ภวนิโรโธ, เอวํ โน เอตฺถ โหติ “อุปาทานนิโรธา ภวนิโรโธ”ติ. “ตณฺหานิโรธา อุปาทานนิโรโธ”ติ อิติ โข ปเนตํ วุตฺตํ, ตณฺหานิโรธา นุ โข ภิกฺขเว อุปาทานนิโรโธ โน วา, กถํ วา เอตฺถ โหตีติ. ตณฺหานิโรธา ภนฺเต อุปาทานนิโรโธ, เอวํ โน เอตฺถ โหติ “ตณฺหานิโรธา อุปาทานนิโรโธ”ติ. “เวทนานิโรธา ตณฺหานิโรโธ”ติ อิติ โข ปเนตํ วุตฺตํ, เวทนานิโรธา นุ โข ภิกฺขเว ตณฺหานิโรโธ โน วา, กถํ วา เอตฺถ โหตีติ. เวทนานิโรธา ภนฺเต ตณฺหานิโรโธ, เอวํ โน เอตฺถ โหติ “เวทนานิโรธา ตณฺหานิโรโธ”ติ. “ผสฺสนิโรธา เวทนานิโรโธ”ติ อิติ โข ปเนตํ วุตฺตํ, ผสฺสนิโรธา นุ โข ภิกฺขเว เวทนานิโรโธ โน วา กถํ วา เอตฺถ โหตีติ. ผสฺสนิโรธา ภนฺเต เวทนานิโรโธ, เอวํ โน เอตฺถ โหติ “ผสฺสนิโรธา เวทนานิโรโธ”ติ. “สฬายตน\-นิโรธา ผสฺสนิโรโธ”ติ อิติ โข ปเนตํ วุตฺตํ, สฬายตน\-นิโรธา นุ โข ภิกฺขเว ผสฺสนิโรโธ โน \csromanfolio{331}วา, กถํ วา เอตฺถ โหตีติ. สฬายตน\-นิโรธา ภนฺเต ผสฺสนิโรโธ, เอวํ โน เอตฺถ โหติ “สฬายตน\-นิโรธา ผสฺสนิโรโธ”ติ. “นามรูป\-นิโรธา สฬายตน\-นิโรโธ”ติ อิติ โข ปเนตํ วุตฺตํ, นามรูป\-นิโรธา นุ โข ภิกฺขเว สฬายตน\-นิโรโธ โน วา, กถํ วา เอตฺถ โหตีติ. นามรูป\-นิโรธา ภนฺเต สฬายตน\-นิโรโธ, เอวํ โน เอตฺถ โหติ “นามรูป\-นิโรธา สฬายตน\-นิโรโธ”ติ. “วิญฺญาณนิโรธา นามรูป\-นิโรโธ”ติ อิติ โข ปเนตํ วุตฺตํ. วิญฺญาณนิโรธา นุ โข ภิกฺขเว นามรูป\-นิโรโธ โน วา, กถํ วา เอตฺถ โหตีติ. วิญฺญาณนิโรธา ภนฺเต นามรูป\-นิโรโธ, เอวํ โน เอตฺถ โหติ “วิญฺญาณนิโรธา นามรูป\-นิโรโธ”ติ. “สงฺขาร\-นิโรธา วิญฺญาณนิโรโธ”ติ อิติ โข ปเนตํ วุตฺตํ, สงฺขาร\-นิโรธา นุ โข ภิกฺขเว วิญฺญาณนิโรโธ โน วา, กถํ วา เอตฺถ โหตีติ. สงฺขาร\-นิโรธา ภนฺเต วิญฺญาณนิโรโธ, เอวํ โน เอตฺถ โหติ “สงฺขาร\-นิโรธา วิญฺญาณนิโรโธ”ติ. “อวิชฺชานิโรธา สงฺขาร\-นิโรโธ”ติ อิติ โข ปเนตํ วุตฺตํ, อวิชฺชานิโรธา นุ โข ภิกฺขเว สงฺขาร\-นิโรโธ โน วา, กถํ วา เอตฺถ โหตีติ. อวิชฺชานิโรธา ภนฺเต สงฺขาร\-นิโรโธ, เอวํ โน เอตฺถ โหติ “อวิชฺชานิโรธา สงฺขาร\-นิโรโธ”ติ.}

\proseitem{406}{สาธุ ภิกฺขเว. อิติ โข ภิกฺขเว ตุเมฺหปิ เอวํ วเทถ, อหมฺปิ เอวํ วทามิ, อิมสฺมิํ อสติ อิทํ น โหติ, อิมสฺส นิโรธา อิทํ นิรุชฺฌติ. ยทิทํ อวิชฺชานิโรธา สงฺขาร\-นิโรโธ, สงฺขาร\-นิโรธา วิญฺญาณนิโรโธ, วิญฺญาณนิโรธา นามรูป\-นิโรโธ, นามรูป\-นิโรธา สฬายตน\-นิโรโธ, สฬายตน\-นิโรธา ผสฺสนิโรโธ, ผสฺสนิโรธา เวทนานิโรโธ, เวทนานิโรธา ตณฺหานิโรโธ, ตณฺหานิโรธา อุปาทานนิโรโธ, อุปาทานนิโรธา ภวนิโรโธ, ภวนิโรธา ชาตินิโรโธ, ชาตินิโรธา ชรามรณํ โสก\-ปริเทว\-ทุกฺข\-โทมนสฺสุ\-ปายาสา นิรุชฺฌนฺติ, เอวเมตสฺส เกวลสฺส ทุกฺข\-กฺขนฺธสฺส นิโรโธ โหติ.}

\proseitem{407}{อปิ นุ ตุเมฺห ภิกฺขเว เอวํ ชานนฺตา เอวํ ปสฺสนฺตา ปุพฺพนฺตํ วา ปฏิธาเวยฺยาถ, “อเหสุมฺห นุ โข มยํ อตีตมทฺธานํ, นนุ โข อเหสุมฺห อตีตมทฺธานํ, กิํ นุ โข อเหสุมฺห อตีตมทฺธานํ, กถํ นุ โข อเหสุมฺห อตีตมทฺธานํ, กิํ หุตฺวา กิํ อเหสุมฺห นุ โข มยํ อตีตมทฺธานนฺ”ติ. โน เหตํ ภนฺเต. อปิ นุ ตุเมฺห ภิกฺขเว เอวํ ชานนฺตา เอวํ ปสฺสนฺตา \csromanfolio{332}อปรนฺตํ วา ปฏิธาเวยฺยาถ, “ภวิสฺสาม นุ โข มยํ อนาคตมทฺธานํ, นนุ โข ภวิสฺสาม อนาคตมทฺธานํ, กิํ นุ โข ภวิสฺสาม อนาคตมทฺธานํ, กถํ นุ โข ภวิสฺสาม อนาคตมทฺธานํ, กิํหุตฺวา กิํ ภวิสฺสาม นุ โข มยํ อนาคตมทฺธานนฺ”ติ. โน เหตํ ภนฺเต. อปิ นุ ตุเมฺห ภิกฺขเว เอวํ ชานนฺตา เอวํ ปสฺสนฺตา เอตรหิ วา ปจฺจุปฺปนฺน\-ม\-ทฺธานํ อชฺฌตฺตํ กถํกถี อสฺสถ, “อหํ นุ โขสฺมิ, โน นุ โขสฺมิ, กิํ นุ โขสฺมิ, กถํ นุ โขสฺมิ, อยํ นุ โข สตฺโต กุโต อาคโต, โส กุหิํ คามี ภวิสฺสตี”ติ. โน เหตํ ภนฺเต. อปิ นุ ตุเมฺห ภิกฺขเว เอวํ ชานนฺตา เอวํ ปสฺสนฺตา เอวํ วเทยฺยาถ “สตฺถา โน ครุ, สตฺถุคารเวน จ มยํ เอวํ วเทมา”ติ. โน เหตํ ภนฺเต. อปิ นุ ตุเมฺห ภิกฺขเว เอวํ ชานนฺตา เอวํ ปสฺสนฺตา เอวํ วเทยฺยาถ “สมโณ เอวมาห, สมณา จ นาม มยํ เอวํ วเทมา”ติ. โน เหตํ ภนฺเต. อปิ นุ ตุเมฺห ภิกฺขเว เอวํ ชานนฺตา เอวํ ปสฺสนฺตา อญฺญํ สตฺถารํ อุทฺทิเสยฺยาถาติ. โน เหตํ ภนฺเต. อปิ นุ ตุเมฺห ภิกฺขเว เอวํ ชานนฺตา เอวํ ปสฺสนฺตา ยานิ ตานิ ปุถุสมณ\-พฺราหฺมณานํ วต โกตูหล\-มงฺคลานิ, ตานิ สารโต ปจฺจาคจฺเฉยฺยาถาติ. โน เหตํ ภนฺเต. นนุ ภิกฺขเว ยเทว ตุมฺหากํ สามํ ญาตํ สามํ ทิฏฺฐํ สามํ วิทิตํ, ตเทว ตุเมฺห วเทถาติ. เอวํ ภนฺเต. สาธุ ภิกฺขเว, อุปนีตา โข เม ตุเมฺห ภิกฺขเว อิมินา สนฺทิฏฺฐิเกน ธมฺเมน อกาลิเกน เอหิปสฺสิเกน โอปเนยฺยิเกน ปจฺจตฺตํ เวทิตพฺเพน วิญฺญูหิ. “สนฺทิฏฺฐิโก อยํ ภิกฺขเว ธมฺโม อกาลิโก เอหิปสฺสิโก โอปเนยฺยิโก ปจฺจตฺตํ เวทิตพฺโพ วิญฺญูหิ” อิติ ยนฺตํ วุตฺตํ, อิทเมตํ ปฏิจฺจ วุตฺตนฺติ.}

\proseitem{408}{ติณฺณํ โข ปน ภิกฺขเว สนฺนิปาตา คพฺภสฺสา\-วกฺกนฺติ โหติ. อิธ มาตาปิตโร จ สนฺนิปติตา โหนฺติ, มาตา จ น อุตุนี โหติ, คนฺธพฺโพ จ น ปจฺจุปฏฺฐิโต โหติ, เนว ตาว คพฺภสฺสา\-วกฺกนฺติ โหติ. อิธ มาตาปิตโร จ สนฺนิปติตา โหนฺติ, มาตา จ อุตุนี โหติ, คนฺธพฺโพ จ น ปจฺจุปฏฺฐิโต โหติ, เนว ตาว คพฺภสฺสา\-วกฺกนฺติ โหติ. ยโต จ โข ภิกฺขเว มาตาปิตโร จ สนฺนิปติตา โหนฺติ, มาตา จ อุตุนี โหติ, คนฺธพฺโพ จ ปจฺจุปฏฺฐิโต โหติ. เอวํ ติณฺณํ สนฺนิปาตา คพฺภสฺสา\-วกฺกนฺติ โหติ. ตเมนํ ภิกฺขเว มาตา นว วา ทส วา มาเส คพฺภํ กุจฺฉินา ปริหรติ มหตา สํสเยน ครุภารํ\csromansharedfootnote{ee3dd3073117e1a7}{ครุมฺภารํ (สี, อิ)}, ตเมนํ ภิกฺขเว \csromanfolio{333}มาตา นวนฺนํ วา ทสนฺนํ วา มาสานํ อจฺจเยน วิชายติ มหตา สํสเยน ครุภารํ. ตเมนํ ชาตํ สมานํ สเกน โลหิเตน โปเสติ. โลหิตเญฺหตํ ภิกฺขเว อริยสฺส วินเย, ยทิทํ มาตุถญฺญํ. ส โข โส ภิกฺขเว กุมาโร วุทฺธิมนฺวาย อินฺทฺริยานํ ปริปากม\-นฺวาย ยานิ ตานิ กุมารกานํ กีฬาปนกานิ, เตหิ กีฬติ, เสยฺยถิทํ, วงฺกกํ ฆฏิกํ โมกฺขจิกํ จิงฺคุลกํ ปตฺตาฬฺหกํ รถกํ ธนุกํ. ส โข โส ภิกฺขเว กุมาโร วุทฺธิมนฺวาย อินฺทฺริยานํ ปริปากม\-นฺวาย ปญฺจหิ กามคุเณหิ สมปฺปิโต สมงฺคีภูโต ปริจาเรติ. จกฺขุ\-วิญฺเญยฺเยหิ รูเปหิ อิฏฺเฐหิ กนฺเตหิ มนาเปหิ ปิยรูเปหิ กามู\-ปสํหิเตหิ รชนีเยหิ. โสตวิญฺเญยฺเยหิ สทฺเทหิ. ฆาน\-วิญฺเญยฺเยหิ คนฺเธหิ. ชิวฺหา\-วิญฺเญยฺเยหิ รเสหิ. กายวิญฺเญยฺเยหิ โผฏฺฐพฺเพหิ อิฏฺเฐหิ กนฺเตหิ มนาเปหิ ปิยรูเปหิ กามู\-ปสํหิเตหิ รชนีเยหิ.}

\proseitem{409}{โส จกฺขุนา รูปํ ทิสฺวา ปิยรูเป รูเป สารชฺชติ, อปฺปิยรูเป รูเป พฺยาปชฺชติ อนุปฏฺฐิต\-กายสติ จ วิหรติ ปริตฺตเจตโส, ตญฺจ เจโตวิมุตฺติํ ปญฺญาวิมุตฺติํ ยถาภูตํ นปฺปชานาติ, ยตฺถสฺส เต ปาปกา อกุสลา ธมฺมา อปริเสสา นิรุชฺฌนฺติ. โส เอวํ อนุโรธ\-วิโรธํ สมาปนฺโน ยํ กิญฺจิ เวทนํ เวเทติ สุขํ วา ทุกฺขํ วา อทุกฺขมสุขํ วา, โส ตํ เวทนํ อภินนฺทติ อภิวทติ, อชฺโฌสาย ติฏฺฐติ. ตสฺส ตํ เวทนํ อภินนฺทโต อภิวทโต อชฺโฌสาย ติฏฺฐโต อุปฺปชฺชติ นนฺที, ยา เวทนาสุ นนฺที, ตทุปาทานํ. ตสฺสุ\-ปาทานปจฺจยา ภโว, ภวปจฺจยา ชาติ, ชาติปจฺจยา ชรามรณํ โสก\-ปริเทว\-ทุกฺข\-โทมนสฺสุ\-ปายาสา สมฺภวนฺติ, เอวเมตสฺส เกวลสฺส ทุกฺข\-กฺขนฺธสฺส สมุทโย โหติ. โสเตน สทฺทํ สุตฺวา ฯเปฯ. ฆาเนน คนฺธํ ฆายิตฺวา ฯเปฯ. ชิวฺหาย รสํ สายิตฺวา ฯเปฯ. กาเยน โผฏฺฐพฺพํ ผุสิตฺวา ฯเปฯ. มนสา ธมฺมํ วิญฺญาย ปิยรูเป ธมฺเม สารชฺชติ, อปฺปิยรูเป ธมฺเม พฺยาปชฺชติ, อนุปฏฺฐิต\-กายสติ จ วิหรติ ปริตฺตเจตโส, ตญฺจ เจโตวิมุตฺติํ ปญฺญาวิมุตฺติํ ยถาภูตํ นปฺปชานาติ, ยตฺถสฺส เต ปาปกา อกุสลา ธมฺมา อปริเสสา นิรุชฺฌนฺติ. โส เอวํ อนุโรธ\-วิโรธํ สมาปนฺโน ยํ กิญฺจิ เวทนํ เวเทติ สุขํ วา ทุกฺขํ วา อทุกฺขมสุขํ วา, โส ตํ เวทนํ อภินนฺทติ อภิวทติ, อชฺโฌสาย ติฏฺฐติ. ตสฺส ตํ เวทนํ อภินนฺทโต อภิวทโต อชฺโฌสาย ติฏฺฐโต อุปฺปชฺชติ นนฺที, ยา เวทนาสุ นนฺที, ตทุปาทานํ. \csromanfolio{334}ตสฺสุ\-ปาทานปจฺจยา ภโว, ภวปจฺจยา ชาติ, ชาติปจฺจยา ชรามรณํ โสก\-ปริเทว\-ทุกฺข\-โทมนสฺสุ\-ปายาสา สมฺภวนฺติ, เอวเมตสฺส เกวลสฺส ทุกฺข\-กฺขนฺธสฺส สมุทโย โหติ.}

\proseitem{410}{อิธ ภิกฺขเว ตถาคโต โลเก อุปฺปชฺชติ อรหํ สมฺมาสมฺพุทฺโธ วิชฺชา\-จรณ\-สมฺปนฺโน สุคโต โลกวิทู อนุตฺตโร ปุริส\-ทมฺม\-สารถิ สตฺถา เทวมนุสฺสานํ พุทฺโธ ภควา. โส อิมํ โลกํ สเทวกํ สมารกํ สพฺรหฺมกํ สสฺสมณ\-พฺราหฺมณิํ ปชํ สเทวมนุสฺสํ สยํ อภิญฺญา สจฺฉิกตฺวา ปเวเทติ, โส ธมฺมํ เทเสติ อาทิกลฺยาณํ มชฺเฌกลฺยาณํ ปริโยสาน\-กลฺยาณํ สาตฺถํ สพฺยญฺชนํ เกวล\-ปริปุณฺณํ ปริสุทฺธํ พฺรหฺมจริยํ ปกาเสติ. ตํ ธมฺมํ สุณาติ คหปติ วา คหปติปุตฺโต วา อญฺญตรสฺมิํ วา กุเล ปจฺจาชาโต, โส ตํ ธมฺมํ สุตฺวา ตถาคเต สทฺธํ ปฏิลภติ, โส เตน สทฺธา\-ปฏิลาเภน สมนฺนาคโต อิติ ปฏิสญฺจิกฺขติ “สมฺพาโธ ฆราวาโส รชาปโถ, อพฺโภกาโส ปพฺพชฺชา, นยิทํ สุกรํ อคารํ อชฺฌาวสตา เอกนฺต\-ปริปุณฺณํ เอกนฺต\-ปริสุทฺธํ สงฺขลิขิตํ พฺรหฺมจริยํ จริตุํ, ยํนูนาหํ เกสมสฺสุํ โอหาเรตฺวา กาสายานิ วตฺถานิ อจฺฉาเทตฺวา อคารสฺมา อนคาริยํ ปพฺพเชยฺยนฺ”ติ. โส อปเรน สมเยน อปฺปํ วา โภคกฺขนฺธํ ปหาย มหนฺตํ วา โภคกฺขนฺทํ ปหาย อปฺปํ วา ญาติปริวฏฺฏํ ปหาย มหนฺตํ วา ญาติปริวฏฺฏํ ปหาย เกสมสฺสุํ โอหาเรตฺวา กาสายานิ วตฺถานิ อจฺฉาเทตฺวา อคารสฺมา อนคาริยํ ปพฺพชติ.}

\proseitem{411}{โส เอวํ ปพฺพชิโต สมาโน ภิกฺขูนํ สิกฺขา\-สาชีว\-สมาปนฺโน ปาณาติปาตํ ปหาย ปาณาติปาตา ปฏิวิรโต โหติ นิหิตทณฺโฑ นิหิตสตฺโถ ลชฺชี ทยาปนฺโน, สพฺพปาณภูต\-หิตา\-นุกมฺปี วิหรติ.}

\prose{อทินฺนาทานํ ปหาย อทินฺนาทานา ปฏิวิรโต โหติ ทินฺนาทายี ทินฺน\-ปาฏิกงฺขี, อเถเนน สุจิภูเตน อตฺตนา วิหรติ.}

\prose{อพฺรหฺมจริยํ ปหาย พฺรหฺมจารี โหติ อาราจารี วิรโต เมถุนา คามธมฺมา. \csromanfolio{335}มุสาวาทํ ปหาย มุสาวาทา ปฏิวิรโต โหติ สจฺจวาที สจฺจสนฺโธ เถโต ปจฺจยิโก อวิสํวาทโก โลกสฺส.}

\prose{ปิสุณํ วาจํ ปหาย ปิสุณาย วาจาย ปฏิวิรโต โหติ, อิโต สุตฺวา น อมุตฺร อกฺขาตา อิเมสํ เภทาย, อมุตฺร วา สุตฺวา น อิเมสํ อกฺขาตา อมูสํ เภทาย, อิติ ภินฺนานํ วา สนฺธาตา, สหิตานํ วา อนุปฺปทาตา, สมคฺคาราโม สมคฺครโต สมคฺคนนฺที สมคฺคกรณิํ วาจํ ภาสิตา โหติ.}

\prose{ผรุสํ วาจํ ปหาย ผรุสาย วาจาย ปฏิวิรโต โหติ, ยา สา วาจา เนลา กณฺณสุขา เปมนียา หทยงฺคมา โปรี พหุชนกนฺตา พหุชนมนาปา, ตถารูปิํ วาจํ ภาสิตา โหติ.}

\prose{สมฺผปฺปลาปํ ปหาย สมฺผปฺปลาปา ปฏิวิรโต โหติ กาลวาที ภูตวาที อตฺถวาที ธมฺมวาที วินยวาที, นิธานวติํ วาจํ ภาสิตา กาเลน สาปเทสํ ปริยนฺตวติํ อตฺถสํหิตํ.}

\prose{โส พีชคาม\-ภูตคาม\-สมารมฺภา ปฏิวิรโต โหติ. เอกภตฺติโก โหติ รตฺตูปรโต วิรโต วิกาลโภชนา. นจฺจ\-คีต\-วาทิต\-วิสูก\-ทสฺสนา ปฏิวิรโต โหติ. มาลา\-คนฺธ\-วิเลปน\-ธารณ\-มณฺฑน\-วิภูสน\-ฏฺฐานา ปฏิวิรโต โหติ. อุจฺจาสยน\-มหาสยนา ปฏิวิรโต โหติ. ชาตรูป\-รชต\-ปฏิคฺคหณา ปฏิวิรโต โหติ. อามก\-ธญฺญ\-ปฏิคฺคหณา ปฏิวิรโต โหติ. อามก\-มํส\-ปฏิคฺคหณา ปฏิวิรโต โหติ. อิตฺถิ\-กุมาริก\-ปฏิคฺคหณา ปฏิวิรโต โหติ. ทาสิ\-ทาส\-ปฏิคฺคหณา ปฏิวิรโต โหติ. อเช\-ฬก\-ปฏิคฺคหณา ปฏิวิรโต โหติ. กุกฺกุฏ\-สูกร\-ปฏิคฺคหณา ปฏิวิรโต โหติ. หตฺถิ\-ควา\-สฺส\-วฬว\-ปฏิคฺคหณา ปฏิวิรโต โหติ. เขตฺตวตฺถุ\-ปฏิคฺคหณา ปฏิวิรโต โหติ. ทูเตยฺย\-ปหิณ\-คมนา\-นุโยคา ปฏิวิรโต โหติ. กยวิกฺกยา ปฏิวิรโต โหติ. ตุลา\-กูฏ\-กํส\-กูฏ\-มาน\-กูฏา ปฏิวิรโต โหติ. อุกฺโกฏน วญฺจน นิกติ สาจิโยคา ปฏิวิรโต โหติ. เฉทน วธพนฺธน\-วิปราโมส อาโลป สหสาการา ปฏิวิรโต โหติ\csromansharedfootnote{0f9a95619017dc70}{ปสฺส จูฬหตฺถิปโทปเม (238) ปิฏฺเฐ.}.}

\prose{โส สนฺตุฏฺโฐ โหติ กาย\-ปริหาริเกน จีวเรน กุจฺฉิ\-ปริหาริเกน ปิณฺฑปาเตน, โส เยน เยเนว ปกฺกมติ, สมาทาเยว \csromanfolio{336}ปกฺกมติ. เสยฺยถาปิ นาม ปกฺขี สกุโณ เยน เยเนว เฑติ, สปตฺตภาโรว เฑติ. เอวเมว ภิกฺขุ สนฺตุฏฺโฐ โหติ กาย\-ปริหาริเกน จีวเรน กุจฺฉิ\-ปริหาริเกน ปิณฺฑปาเตน, โส เยน เยเนว ปกฺกมติ, สมาทาเยว ปกฺกมติ. โส อิมินา อริเยน สีลกฺขนฺเธน สมนฺนาคโต อชฺฌตฺตํ อนวชฺชสุขํ ปฏิสํเวเทติ.}

\proseitem{8}{โส จกฺขุนา รูปํ ทิสฺวา น นิมิตฺตคฺคาหี โหติ นานุพฺยญฺชน\-คฺคาหี, ยตฺวาธิกรณเมนํ จกฺขุนฺทฺริยํ อสํวุตํ วิหรนฺตํ อภิชฺฌา\-โทมนสฺสา ปาปกา อกุสลา ธมฺมา อนฺวาสฺสเวยฺยุํ, ตสฺส สํวราย ปฏิปชฺชติ, รกฺขติ จกฺขุนฺทฺริยํ, จกฺขุนฺทฺริเย สํวรํ อาปชฺชติ. โสเตน สทฺทํ สุตฺวา ฯเปฯ. ฆาเนน คนฺธํ ฆายิตฺวา ฯเปฯ. ชิวฺหาย รสํ สายิตฺวา ฯเปฯ. กาเยน โผฏฺฐพฺพํ ผุสิตฺวา ฯเปฯ. มนสา ธมฺมํ วิญฺญาย น นิมิตฺตคฺคาหี โหติ นานุพฺยญฺชน\-คฺคาหี, ยตฺวาธิกรณเมนํ มนินฺทฺริยํ อสํวุตํ วิหรนฺตํ อภิชฺฌา\-โทมนสฺสา ปาปกา อกุสลา ธมฺมา อนฺวาสฺสเวยฺยุํ, ตสฺส สํวราย ปฏิปชฺชติ, รกฺขติ มนินฺทฺริยํ, มนินฺทฺริเย สํวรํ อาปชฺชติ. โส อิมินา อริเยน อินฺทฺริย\-สํวเรน สมนฺนาคโต อชฺฌตฺตํ อพฺยาเสกสุขํ ปฏิสํเวเทติ.}

\prose{โส อภิกฺกนฺเต ปฏิกฺกนฺเต สมฺปชานการี โหติ, อาโลกิเต วิโลกิเต สมฺปชานการี โหติ, สมิญฺชิเต ปสาริเต สมฺปชานการี โหติ, สํฆาฏิ\-ปตฺต\-จีวร\-ธารเณ สมฺปชานการี โหติ, อสิเต ปีเต ขายิเต สายิเต สมฺปชานการี โหติ, อุจฺจาร\-ปสฺสาว\-กมฺเม สมฺปชานการี โหติ, คเต ฐิเต นิสินฺเน สุตฺเต ชาคริเต ภาสิเต ตุณฺหีภาเว สมฺปชานการี โหติ.}

\proseitem{412}{โส อิมินา จ อริเยน สีลกฺขนฺเธน สมนฺนาคโต, (อิมาย จ อริยาย สนฺตุฏฺฐิยา สมนฺนาคโต)\csromansharedfootnote{f6113e777bd388d1}{ปสฺส จูฬหตฺถิปโทปเม (239) ปิฏฺเฐ.} อิมินา จ อริเยน อินฺทฺริย\-สํวเรน สมนฺนาคโต, อิมินา จ อริเยน สติสมฺปชญฺเญน สมนฺนาคโต วิวิตฺตํ เสนาสนํ ภชติ อรญฺญํ รุกฺขมูลํ ปพฺพตํ กนฺทรํ คิริคุหํ สุสานํ วนปตฺถํ อพฺโภกาสํ ปลาลปุญฺชํ. โส ปจฺฉาภตฺตํ ปิณฺฑปาต\-ปฏิกฺกนฺโต นิสีทติ ปลฺลงฺกํ อาภุชิตฺวา อุชุํ กายํ ปณิธาย ปริมุขํ สติํ อุปฏฺฐเปตฺวา. โส อภิชฺฌํ โลเก ปหาย วิคตา\-ภิชฺเฌน เจตสา วิหรติ, อภิชฺฌาย จิตฺตํ ปริโสเธติ. พฺยาปาทปโทสํ \csromanfolio{337}ปหาย อพฺยาปนฺนจิตฺโต วิหรติ สพฺพปาณภูต\-หิตา\-นุกมฺปี, พฺยาปาทปโทสา จิตฺตํ ปริโสเธติ. ถินมิทฺธํ ปหาย วิคต\-ถินมิทฺโธ วิหรติ อาโลกสญฺญี สโต สมฺปชาโน, ถินมิทฺธา จิตฺตํ ปริโสเธติ. อุทฺธจฺจ\-กุกฺกุจฺจํ ปหาย อนุทฺธโต วิหรติ อชฺฌตฺตํ วูปสนฺตจิตฺโต, อุทฺธจฺจ\-กุกฺกุจฺจา จิตฺตํ ปริโสเธติ. วิจิกิจฺฉํ ปหาย ติณฺณ\-วิจิกิจฺโฉ วิหรติ อกถํกถี กุสเลสุ ธมฺเมสุ, วิจิกิจฺฉาย จิตฺตํ ปริโสเธติ.}

\proseitem{413}{โส อิเม ปญฺจ นีวรเณ ปหาย เจตโส อุปกฺกิเลเส ปญฺญาย ทุพฺพลีกรเณ วิวิจฺเจว กาเมหิ วิวิจฺจ อกุสเลหิ ธมฺเมหิ สวิตกฺกํ สวิจารํ วิเวกชํ ปีติสุขํ ปฐมํ ฌานํ อุปสมฺปชฺช วิหรติ. ปุน จปรํ ภิกฺขเว ภิกฺขุ วิตกฺก\-วิจารานํ วูปสมา อชฺฌตฺตํ สมฺปสาทนํ เจตโส เอโกทิภาวํ อวิตกฺกํ อวิจารํ สมาธิชํ ปีติสุขํ ทุติยํ ฌานํ ฯเปฯ ตติยํ ฌานํ ฯเปฯ จตุตฺถํ ฌานํ อุปสมฺปชฺช วิหรติ.}

\csromanmatikaanchor{mtk.71}
\csromantocmarkref{subsection}{9. มหาอสฺสปุรสุตฺต}{mtk.71}
\bookmark[dest={mtk.71},level=2]{9. มหาอสฺสปุรสุตฺต}
\proseitem{414}{โส จกฺขุนา รูปํ ทิสฺวา ปิยรูเป รูเป น สารชฺชติ, อปฺปิยรูเป รูเป น พฺยาปชฺชติ, อุปฏฺฐิต\-กายสติ จ วิหรติ อปฺปมาณเจตโส, ตญฺจ เจโตวิมุตฺติํ ปญฺญาวิมุตฺติํ ยถาภูตํ ปชานาติ, ยตฺถสฺส เต ปาปกา อกุสลา ธมฺมา อปริเสสา นิรุชฺฌนฺติ. โส เอวํ อนุโรธ\-วิโรธ\-วิปฺปหีโน ยํ กิญฺจิ เวทนํ เวเทติ สุขํ วา ทุกฺขํ วา อทุกฺขมสุขํ วา, โส ตํ เวทนํ นาภินนฺทติ นาภิวทติ, นาชฺโฌสาย ติฏฺฐติ. ตสฺส ตํ เวทนํ อนภินนฺทโต อนภิวทโต อนชฺโฌสาย ติฏฺฐโต ยา เวทนาสุ นนฺที, สา นิรุชฺฌติ, ตสฺส นนฺทีนิโรธา อุปาทานนิโรโธ, อุปาทานนิโรธา ภวนิโรโธ, ภวนิโรธา ชาตินิโรโธ, ชาตินิโรธา ชรามรณํ โสก\-ปริเทว\-ทุกฺข\-โทมนสฺสุ\-ปายาสา นิรุชฺฌนฺติ, เอวเมตสฺส เกวลสฺส ทุกฺข\-กฺขนฺธสฺส นิโรโธ โหติ. โสเตน สทฺทํ สุตฺวา ฯเปฯ. ฆาเนน คนฺธํ ฆายิตฺวา ฯเปฯ. ชิวฺหาย รสํ สายิตฺวา ฯเปฯ. กาเยน โผฏฺฐพฺพํ ผุสิตฺวา ฯเปฯ. มนสา ธมฺมํ วิญฺญาย ปิยรูเป ธมฺเม น สารชฺชติ, อปฺปิยรูเป ธมฺเม น พฺยาปชฺชติ, อุปฏฺฐิต\-กายสติ จ วิหรติ อปฺปมาณเจตโส, ตญฺจ เจโตวิมุตฺติํ ปญฺญาวิมุตฺติํ ยถาภูตํ ปชานาติ, ยตฺถสฺส เต ปาปกา อกุสลา ธมฺมา อปริเสสา นิรุชฺฌนฺติ. โส เอวํ อนุโรธ\-วิโรธ\-วิปฺปหีโน ยํ กิญฺจิ เวทนํ เวเทติ สุขํ วา ทุกฺขํ วา อทุกฺขมสุขํ วา, โส ตํ เวทนํ นาภินนฺทติ \csromanfolio{338}นาภิวทติ, นาชฺโฌสาย ติฏฺฐติ. ตสฺส ตํ เวทนํ อนภินนฺทโต อนภิวทโต อนชฺโฌสาย ติฏฺฐโต ยา เวทนาสุ นนฺที, สา นิรุชฺฌติ. ตสฺส นนฺทีนิโรธา อุปาทานนิโรโธ, อุปาทานนิโรธา ภวนิโรโธ, ภวนิโรธา ชาตินิโรโธ, ชาตินิโรธา ชรามรณํ โสก\-ปริเทว\-ทุกฺข\-โทมนสฺสุ\-ปายาสา นิรุชฺฌนฺติ, เอวเมตสฺส เกวลสฺส ทุกฺข\-กฺขนฺธสฺส นิโรโธ โหติ. อิมํ โข เม ตุเมฺห ภิกฺขเว สํขิตฺเตน ตณฺหาสงฺขย\-วิมุตฺติํ ธาเรถ. สาติํ ปน ภิกฺขุํ เกวฏฺฏปุตฺตํ มหาตณฺหาชาล- ตณฺหาสงฺฆาฏปฺปฏิมุกฺกนฺติ.}

\prosewithcloserruled{อิทมโวจ ภควา. อตฺตมนา เต ภิกฺขู ภควโต ภาสิตํ อภินนฺทุนฺติ.}{มหา\-ตณฺหาสงฺขย\-สุตฺตํ นิฏฺฐิตํ อฏฺฐมํ.}
\prose{9. มหา\-อสฺสปุร\-สุตฺต}

\proseitem{415}{เอวํ เม สุตํ\csromandash{}เอกํ สมยํ ภควา องฺเคสุ วิหรติ อสฺสปุรํ นาม องฺคานํ นิคโม. ตตฺร โข ภควา ภิกฺขู อามนฺเตสิ “ภิกฺขโว”ติ. “ภทนฺเต”ติ เต ภิกฺขู ภควโต ปจฺจสฺโสสุํ. ภควา เอตทโวจ\csromandash{}}

\prose{“สมณา สมณา”ติ โว ภิกฺขเว ชโน สญฺชานาติ, ตุเมฺห จ ปน “เก ตุเมฺห”ติ ปุฏฺฐา สมานา “สมณามฺหา”ติ ปฏิชานาถ, เตสํ โว ภิกฺขเว เอวํสมญฺญานํ สตํ เอวํปฏิญฺญานํ สตํ “เย ธมฺมา สมณกรณา จ พฺราหฺมณกรณา จ, เต ธมฺเม สมาทาย วตฺติสฺสาม, เอวํ โน อยํ อมฺหากํ สมญฺญา จ สจฺจา ภวิสฺสติ ปฏิญฺญา จ ภูตา. เยสญฺจ มยํ จีวรปิณฺฑปาตเสนาสน\-คิลานปฺปจฺจย\-เภสชฺช- ปริกฺขารํ ปริภุญฺชาม, เตสํ เต การา อเมฺหสุ มหปฺผลา ภวิสฺสนฺติ มหานิสํสา, อมฺหากญฺเจวายํ ปพฺพชฺชา อวญฺฌา ภวิสฺสติ สผลา ส- อุทฺรยา”ติ เอวํ หิ โว ภิกฺขเว สิกฺขิตพฺพํ.}

\proseitem{416}{กตเม จ ภิกฺขเว ธมฺมา สมณกรณา จ พฺราหฺมณกรณา จ, “หิโรตฺตปฺเปน สมนฺนาคตา ภวิสฺสามา”ติ เอวํ หิ โว ภิกฺขเว สิกฺขิตพฺพํ. สิยา โข ปน ภิกฺขเว ตุมฺหากํ เอวมสฺส, “หิโรตฺตปฺเปนมฺห \csromanfolio{339}สมนฺนาคตา, อลเมตฺตาวตา, กตเมตฺตาวตา, อนุปฺปตฺโต โน สามญฺญตฺโถ, นตฺถิ โน กิญฺจิ อุตฺตริํ กรณียนฺ”ติ ตาวตเกเนว ตุฏฺฐิํ อาปชฺเชยฺยาถ. อาโรจยามิ โว ภิกฺขเว, ปฏิเวทยามิ โว ภิกฺขเว, มา โว สามญฺญ\-ตฺถิกานํ สตํ สามญฺญตฺโถ ปริหายิ สติ อุตฺตริํ กรณีเย.}

\proseitem{417}{กิญฺจ ภิกฺขเว อุตฺตริํ กรณียํ, “ปริสุทฺโธ โน กายสมาจาโร ภวิสฺสติ อุตฺตาโน วิวโฏ, น จ ฉิทฺทวา, สํวุโต จ, ตาย จ ปน ปริสุทฺธ\-กาย\-สมาจาร\-ตาย เนว\-ตฺตานุกฺกํเสสฺสาม น ปรํ วมฺเภสฺสามา”ติ\csromansharedfootnote{8f4a7df92a550e53}{เนวตฺตานุกฺกํสิสฺสาม น ปรํ วมฺภิสฺสามาติ (สพฺพตฺถ)} เอวํ หิ โว ภิกฺขเว สิกฺขิตพฺพํ. สิยา โข ปน ภิกฺขเว ตุมฺหากํ เอวมสฺส, “หิโรตฺตปฺเปนมฺห สมนฺนาคตา, ปริสุทฺโธ โน กายสมาจาโร. อลเมตฺตาวตา, กตเมตฺตาวตา, อนุปฺปตฺโต โน สามญฺญตฺโถ, นตฺถิ โน กิญฺจิ อุตฺตริํ กรณียนฺ”ติ ตาวตเกเนว ตุฏฺฐิํ อาปชฺเชยฺยาถ. อาโรจยามิ โว ภิกฺขเว, ปฏิเวทยามิ โว ภิกฺขเว, มา โว สามญฺญ\-ตฺถิกานํ สตํ สามญฺญตฺโถ ปริหายิ สติ อุตฺตริํ กรณีเย.}

\proseitem{418}{กิญฺจ ภิกฺขเว อุตฺตริํ กรณียํ, “ปริสุทฺโธ โน วจีสมาจาโร ภวิสฺสติ อุตฺตาโน วิวโฏ, น จ ฉิทฺทวา, สํวุโต จ, ตาย จ ปน ปริสุทฺธ\-วจี\-สมาจาร\-ตาย เนว\-ตฺตานุกฺกํเสสฺสาม น ปรํ วมฺเภสฺสามา”ติ เอวํ หิ โว ภิกฺขเว สิกฺขิตพฺพํ. สิยา โข ปน ภิกฺขเว ตุมฺหากํ เอวมสฺส, “หิโรตฺตปฺเปนมฺห สมนฺนาคตา, ปริสุทฺโธ โน กายสมาจาโร, ปริสุทฺโธ วจีสมาจาโร. อลเมตฺตาวตา, กตเมตฺตาวตา, อนุปฺปตฺโต โน สามญฺญตฺโถ, นตฺถิ โน กิญฺจิ อุตฺตริํ กรณียนฺ”ติ ตาวตเกเนว ตุฏฺฐิํ อาปชฺเชยฺยาถ. อาโรจยามิ โว ภิกฺขเว, ปฏิเวทยามิ โว ภิกฺขเว, มา โว สามญฺญ\-ตฺถิกานํ สตํ สามญฺญตฺโถ ปริหายิ สติ อุตฺตริํ กรณีเย.}

\proseitem{419}{กิญฺจ ภิกฺขเว อุตฺตริํ กรณียํ, “ปริสุทฺโธ โน มโนสมาจาโร ภวิสฺสติ อุตฺตาโน วิวโฏ, น จ ฉิทฺทวา, สํวุโต จ, ตาย จ \csromanfolio{340}ปน ปริสุทฺธ\-มโน\-สมาจาร\-ตาย เนว\-ตฺตานุกฺกํเสสฺสาม น ปรํ วมฺเภสฺสามา”ติ เอวํ หิ โว ภิกฺขเว สิกฺขิตพฺพํ. สิยา โข ปน ภิกฺขเว ตุมฺหากํ เอวมสฺส, “หิโรตฺตปฺเปนมฺห สมนฺนาคตา, ปริสุทฺโธ โน กายสมาจาโร, ปริสุทฺโธ วจีสมาจาโร, ปริสุทฺโธ มโนสมาจาโร. อลเมตฺตาวตา, กตเมตฺตาวตา, อนุปฺปตฺโต โน สามญฺญตฺโถ, นตฺถิ โน กิญฺจิ อุตฺตริํ กรณียนฺ”ติ ตาวตเกเนว ตุฏฺฐิํ อาปชฺเชยฺยาถ. อาโรจยามิ โว ภิกฺขเว, ปฏิเวทยามิ โว ภิกฺขเว, มา โว สามญฺญ\-ตฺถิกานํ สตํ สามญฺญตฺโถ ปริหายิ สติ อุตฺตริํ กรณีเย.}

\proseitem{420}{กิญฺจ ภิกฺขเว อุตฺตริํ กรณียํ, “ปริสุทฺโธ โน อาชีโว ภวิสฺสติ อุตฺตาโน วิวโฏ, น จ ฉิทฺทวา, สํวุโต จ, ตาย จ ปน ปริสุทฺธา\-ชีว\-ตาย เนว\-ตฺตานุกฺกํเสสฺสาม น ปรํ วมฺเภสฺสามา”ติ เอวํ หิ โว ภิกฺขเว สิกฺขิตพฺพํ. สิยา โข ปน ภิกฺขเว ตุมฺหากํ เอวมสฺส, “หิโรตฺตปฺเปนมฺห สมนฺนาคตา, ปริสุทฺโธ โน กายสมาจาโร, ปริสุทฺโธ วจีสมาจาโร, ปริสุทฺโธ มโนสมาจาโร, ปริสุทฺโธ อาชีโว. อลเมตฺตาวตา, กตเมตฺตาวตา, อนุปฺปตฺโต โน สามญฺญตฺโถ, นตฺถิ โน กิญฺจิ อุตฺตริํ กรณียนฺ”ติ ตาวตเกเนว ตุฏฺฐิํ อาปชฺเชยฺยาถ. อาโรจยามิ โว ภิกฺขเว, ปฏิเวทยามิ โว ภิกฺขเว, มา โว สามญฺญ\-ตฺถิกานํ สตํ สามญฺญตฺโถ ปริหายิ สติ อุตฺตริํ กรณีเย.}

\proseitem{421}{กิญฺจ ภิกฺขเว อุตฺตริํ กรณียํ, “อินฺทฺริเยสุ คุตฺตทฺวารา ภวิสฺสาม, จกฺขุนา รูปํ ทิสฺวา น นิมิตฺตคฺคาหี นานุพฺยญฺชน\-คฺคาหี, ยตฺวาธิกรณเมนํ จกฺขุนฺทฺริยํ อสํวุตํ วิหรนฺตํ อภิชฺฌา\-โทมนสฺสา ปาปกา อกุสลา ธมฺมา อนฺวาสฺสเวยฺยุํ, ตสฺส สํวราย ปฏิปชฺชิสฺสาม, รกฺขิสฺสาม จกฺขุนฺทฺริยํ, จกฺขุนฺทฺริเย สํวรํ อาปชฺชิสฺสาม. โสเตน สทฺทํ สุตฺวา ฯเปฯ. ฆาเนน คนฺธํ ฆายิตฺวา ฯเปฯ. ชิวฺหาย รสํ สายิตฺวา ฯเปฯ. กาเยน โผฏฺฐพฺพํ ผุสิตฺวา ฯเปฯ. มนสา ธมฺมํ วิญฺญาย น นิมิตฺตคฺคาหี นานุพฺยาญฺชนคฺคาหี, ยตฺวาธิกรณเมนํ มนินฺทฺริยํ อสํวุตํ วิหรนฺตํ อภิชฺฌา\-โทมนสฺสา ปาปกา อกุสลา ธมฺมา อนฺวาสฺสเวยฺยุํ, ตสฺส สํวราย ปฏิปชฺชิสฺสาม, รกฺขิสฺสาม มนินฺทฺริยํ, มนินฺทฺริเย สํวรํ อาปชฺชิสฺสามา”ติ เอวํ หิ โว ภิกฺขเว สิกฺขิตพฺพํ. สิยา โข ปน ภิกฺขเว ตุมฺหากํ เอวมสฺส, \csromanfolio{341}“หิโรตฺตปฺเปนมฺห สมนฺนาคตา, ปริสุทฺโธ โน กายสมาจาโร, ปริสุทฺโธ วจีสมาจาโร, ปริสุทฺโธ มโนสมาจาโร, ปริสุทฺโธ อาชีโว, อินฺทฺริเยสุมฺห คุตฺตทฺวารา. อลเมตฺตาวตา, กตเมตฺตาวตา, อนุปฺปตฺโต โน สามญฺญตฺโถ, นตฺถิ โน กิญฺจิ อุตฺตริํ กรณียนฺ”ติ ตาวตเกเนว ตุฏฺฐิํ อาปชฺเชยฺยาถ. อาโรจยามิ โว ภิกฺขเว, ปฏิเวทยามิ โว ภิกฺขเว, มา โว สามญฺญ\-ตฺถิกานํ สตํ สามญฺญตฺโถ ปริหายิ สติ อุตฺตริํ กรณีเย.}

\proseitem{422}{กิญฺจ ภิกฺขเว อุตฺตริํ กรณียํ, “โภชเน มตฺตญฺญุโน ภวิสฺสาม, ปฏิสงฺขา โยนิโส อาหารํ อาหริสฺสาม, เนว ทวาย น มทาย น มณฺฑนาย น วิภูสนาย, ยาว เทว อิมสฺส กายสฺส ฐิติยา ยาปนาย, วิหิํสู\-ปรติยา พฺรหฺมจริยา\-นุคฺคหาย, อิติ ปุราณญฺจ เวทนํ ปฏิหงฺขาม, นวญฺจ เวทนํ น อุปฺปาเทสฺสาม, ยาตฺรา จ โน ภวิสฺสติ อนวชฺชตา จ ผาสุ วิหาโร จา”ติ เอวํ หิ โว ภิกฺขเว สิกฺขิตพฺพํ. สิยา โข ปน ภิกฺขเว ตุมฺหากํ เอวมสฺส, “หิโรตฺตปฺเปนมฺห สมนฺนาคตา, ปริสุทฺโธ โน กายสมาจาโร, ปริสุทฺโธ วจีสมาจาโร, ปริสุทฺโธ มโนสมาจาโร, ปริสุทฺโธ อาชีโว, อินฺทฺริเยสุมฺห คุตฺตทฺวารา, โภชเน มตฺตญฺญุโน. อลเมตฺตาวตา, กตเมตฺตาวตา, อนุปฺปตฺโต โน สามญฺญตฺโถ, นตฺถิ โน กิญฺจิ อุตฺตริํ กรณียนฺ”ติ ตาวตเกเนว ตุฏฺฐิํ อาปชฺเชยฺยาถ. อาโรจยามิ โว ภิกฺขเว, ปฏิเวทยามิ โว ภิกฺขเว, มา โว สามญฺญ\-ตฺถิกานํ สตํ สามญฺญตฺโถ ปริหายิ สติ อุตฺตริํ กรณีเย.}

\proseitem{423}{กิญฺจ ภิกฺขเว อุตฺตริํ กรณียํ, “ชาคริยํ อนุยุตฺตา ภวิสฺสาม, ทิวสํ จงฺกเมน นิสชฺชาย อาวรณีเยหิ ธมฺเมหิ จิตฺตํ ปริโสเธสฺสาม, รตฺติยา ปฐมํ ยามํ จงฺกเมน นิสชฺชาย อาวรณีเยหิ ธมฺเมหิ จิตฺตํ ปริโสเธสฺสาม, รตฺติยา มชฺฌิมํ ยามํ ทกฺขิเณน ปสฺเสน สีหเสยฺยํ กปฺเปสฺสาม ปาเท ปาทํ อจฺจาธาย สโต สมฺปชาโน อุฏฺฐานสญฺญํ มนสิ กริตฺวา, รตฺติยา ปจฺฉิมํ ยามํ ปจฺจุฏฺฐาย จงฺกเมน นิสชฺชาย อาวรณีเยหิ ธมฺเมหิ จิตฺตํ ปริโสเธสฺสามาติ เอวํ หิ โว ภิกฺขเว สิกฺขิตพฺพํ. สิยา โข ปน ภิกฺขเว ตุมฺหากํ เอวมสฺส, “หิโรตฺตปฺเปนมฺห สมนฺนาคตา, ปริสุทฺโธ โน กายสมาจาโร, ปริสุทฺโธ วจีสมาจาโร, \csromanfolio{342}ปริสุทฺโธ มโนสมาจาโร, ปริสุทฺโธ อาชีโว, อินฺทฺริเยสุมฺห คุตฺตทฺวารา, โภชเน มตฺตญฺญุโน, ชาคริยํ อนุยุตฺตา. อลเมตฺตาวตา, กตเมตฺตาวตา, อนุปฺปตฺโต โน สามญฺญตฺโถ, นตฺถิ โน กิญฺจิ อุตฺตริํ กรณียนฺ”ติ. ตาวตเกเนว ตุฏฺฐิํ อาปชฺเชยฺยาถ. อาโรจยามิ โว ภิกฺขเว, ปฏิเวทยามิ โว ภิกฺขเว, มา โว สามญฺญ\-ตฺถิกานํ สตํ สามญฺญตฺโถ ปริหายิ สติ อุตฺตริํ กรณีเย.}

\proseitem{424}{กิญฺจ ภิกฺขเว อุตฺตริํ กรณียํ, “สติสมฺปชญฺเญน สมนฺนาคตา ภวิสฺสาม, อภิกฺกนฺเต ปฏิกฺกนฺเต สมฺปชานการี, อาโลกิเต วิโลกิเต สมฺปชานการี, สมิญฺชิเต ปสาริเต สมฺปชานการี, สํฆาฏิ\-ปตฺต\-จีวร\-ธารเณ สมฺปชานการี, อสิเต ปีเต ขายิเต สายิเต สมฺปชานการี, อุจฺจาร\-ปสฺสาว\-กมฺเม สมฺปชานการี, คเต ฐิเต นิสินฺเน สุตฺเต ชาคริเต ภาสิเต ตุณฺหีภาเว สมฺปชานการี”ติ เอวํ หิ โว ภิกฺขเว สิกฺขิตพฺพํ. สิยา โข ปน ภิกฺขเว ตุมฺหากํ เอวมสฺส, “หิโรตฺตปฺเปนมฺห สมนฺนาคตา, ปริสุทฺโธ โน กายสมาจาโร, ปริสุทฺโธ วจีสมาจาโร, ปริสุทฺโธ มโนสมาจาโร, ปริสุทฺโธ อาชีโว, อินฺทฺริเยสุมฺห คุตฺตทฺวารา, โภชเน มตฺตญฺญุโน, ชาคริยํ อนุยุตฺตา, สติสมฺปชญฺเญน สมนฺนาคตา. อลเมตฺตาวตา, กตเมตฺตาวตา, อนุปฺปตฺโต โน สามญฺญตฺโถ, นตฺถิ โน กิญฺจิ อุตฺตริํ กรณียนฺ”ติ ตาวตเกเนว ตุฏฺฐิํ อาปชฺเชยฺยาถ. อาโรจยามิ โว ภิกฺขเว, ปฏิเวทยามิ โว ภิกฺขเว, มา โว สามญฺญ\-ตฺถิกานํ สตํ สามญฺญตฺโถ ปริหายิ สติ อุตฺตริํ กรณีเย.}

\proseitem{425}{กิญฺจ ภิกฺขเว อุตฺตริํ กรณียํ, อิธ ภิกฺขเว ภิกฺขุ วิวิตฺตํ เสนาสนํ ภชติ อรญฺญํ รุกฺขมูลํ ปพฺพตํ กนฺทรํ คิริคุหํ สุสานํ วนปฺปตฺถํ อพฺโภกาสํ ปลาลปุญฺชํ. โส ปจฺฉาภตฺตํ ปิณฺฑปาต\-ปฏิกฺกนฺโต นิสีทติ ปลฺลงฺกํ อาภุชิตฺวา อุชุํ กายํ ปณิธาย ปริมุขํ สติํ อุปฏฺฐเปตฺวา, โส อภิชฺฌํ โลเก ปหาย วิคตา\-ภิชฺเฌน เจตสา วิหรติ, อภิชฺฌาย จิตฺตํ ปริโสเธติ. พฺยาปาทปโทสํ ปหาย อพฺยาปนฺนจิตฺโต วิหรติ สพฺพปาณภูต\-หิตา\-นุกมฺปี, พฺยาปาทปโทสา จิตฺตํ ปริโสเธติ. ถินมิทฺธํ ปหาย วิคต\-ถินมิทฺโธ วิหรติ อาโลกสญฺญี สโต สมฺปชาโน, ถินมิทฺธา จิตฺตํ ปริโสเธติ. อุทฺธจฺจ\-กุกฺกุจฺจํ ปหาย อนุทฺธโต \csromanfolio{343}วิหรติ อชฺฌตฺตํ วูปสนฺตจิตฺโต, อุทฺธจฺจ\-กุกฺกุจฺจา จิตฺตํ ปริโสเธติ. วิจิกิจฺฉํ ปหาย ติณฺณ\-วิจิกิจฺโฉ วิหรติ อกถํกถี กุสเลสุ ธมฺเมสุ, วิจิกิจฺฉาย จิตฺตํ ปริโสเธติ.}

\proseitem{426}{เสยฺยถาปิ ภิกฺขเว ปุริโส อิณํ อาทาย กมฺมนฺเต ปโยเชยฺย, ตสฺส เต กมฺมนฺตา สมิชฺเฌยฺยุํ\csromansharedfootnote{1c152a7ec1612dd2}{สมฺปชฺเชยฺยุํ (สฺยา, กํ, ก)}, โส ยานิ จ โปราณานิ อิณมูลานิ, ตานิ จ พฺยนฺตี\csromansharedfootnote{a71a148becd7de83}{พฺยนฺติํ (ก), พฺยนฺติ (อิ)} กเรยฺย, สิยา จสฺส อุตฺตริํ อวสิฏฺฐํ ทารภรณาย. ตสฺส เอวมสฺส “อหํ โข ปุพฺเพ อิณํ อาทาย กมฺมนฺเต ปโยเชสิํ, ตสฺส เม เต กมฺมนฺตา สมิชฺฌิํสุ, โสหํ ยานิ จ โปราณานิ อิณมูลานิ, ตานิ จ พฺยนฺตี อกาสิํ, อตฺถิ จ เม อุตฺตริํ อวสิฏฺฐํ ทารภรณายา”ติ. โส ตโตนิทานํ ลเภถ ปาโมชฺชํ, อธิคจฺเฉยฺย โสมนสฺสํ.}

\prose{เสยฺยถาปิ ภิกฺขเว ปุริโส อาพาธิโก อสฺส ทุกฺขิโต พาฬฺหคิลาโน, ภตฺตญฺจสฺส นจฺฉาเทยฺย, น จสฺส กาเย พลมตฺตา, โส อปเรน สมเยน ตมฺหา อาพาธา มุจฺเจยฺย, ภตฺตญฺจสฺส ฉาเทยฺย, สิยา จสฺส กาเย พลมตฺตา. ตสฺส เอวมสฺส “อหํ โข ปุพฺเพ อาพาธิโก อโหสิํ ทุกฺขิโต พาฬฺหคิลาโน, ภตฺตํ จ เม นจฺฉาเทสิ, น จ เม อาสิ กาเย พลมตฺตา, โสมฺหิ เอตรหิ ตมฺหา อาพาธา มุตฺโต, ภตฺตํ จ เม ฉาเทติ, อตฺถิ จ เม กาเย พลมตฺตา”ติ. โส ตโตนิทานํ ลเภถ ปาโมชฺชํ, อธิคจฺเฉยฺย โสมนสฺสํ.}

\prose{เสยฺยถาปิ ภิกฺขเว ปุริโส พนฺธนาคาเร พทฺโธ อสฺส, โส อปเรน สมเยน ตมฺหา พนฺธนา มุจฺเจยฺย โสตฺถินา อพฺภเยน\csromansharedfootnote{86a572cd834564d7}{อพฺยเยน (สี, อิ)}, น จสฺส กิญฺจิ โภคานํ วโย. ตสฺส เอวมสฺส “อหํ โข ปุพฺเพ พนฺธนาคาเร พทฺโธ อโหสิํ, โสมฺหิ เอตรหิ ตมฺหา พนฺธนา มุตฺโต โสตฺถินา อพฺภเยน, นตฺถิ จ เม กิญฺจิ โภคานํ วโย”ติ. โส ตโตนิทานํ ลเภถ ปาโมชฺชํ, อธิคจฺเฉยฺย โสมนสฺสํ. เสยฺยถาปิ ภิกฺขเว ปุริโส ทาโส อสฺส อนตฺตาธีโน ปราธีโน น เยนกามํคโม, โส อปเรน สมเยน ตมฺหา ทาสพฺยา มุจฺเจยฺย อตฺตาธีโน อปราธีโน ภุชิสฺโส เยนกามํคโม. ตสฺส เอวมสฺส “อหํ โข ปุพฺเพ ทาโส \csromanfolio{344}อโหสิํ อนตฺตาธีโน ปราธีโน น เยนกามํคโม, โสมฺหิ เอตรหิ ตมฺหา ทาสพฺยา มุตฺโต อตฺตาธีโน อปราธีโน ภุชิสฺโส เยนกามํคโม”ติ. โส ตโตนิทานํ ลเภถ ปาโมชฺชํ, อธิคจฺเฉยฺย โสมนสฺสํ.}

\prose{เสยฺยถาปิ ภิกฺขเว ปุริโส สธโน สโภโค กนฺตาร\-ทฺธาน\-มคฺคํ ปฏิปชฺเชยฺย\csromansharedfootnote{1b607b28ebe739b8}{สีลกฺขนฺธวคฺคปาฬิยา กิญฺจิ วิสทิสํ.}, โส อปเรน สมเยน ตมฺหา กนฺตารา นิตฺถเรยฺย โสตฺถินา อพฺภเยน, น จสฺส กิญฺจิ โภคานํ วโย. ตสฺส เอวมสฺส “อหํ โข ปุพฺเพ สธโน สโภโค กนฺตาร\-ทฺธาน\-มคฺคํ ปฏิปชฺชิํ, โสมฺหิ เอตรหิ ตมฺหา กนฺตารา นิตฺถิณฺโณ โสตฺถินา อพฺภเยน, นตฺถิ จ เม กิญฺจิ โภคานํ วโย”ติ. โส ตโตนิทานํ ลเภถ ปาโมชฺชํ, อธิคจฺเฉยฺย โสมนสฺสํ.}

\prose{เอวเมว โข ภิกฺขเว ภิกฺขุ ยถา อิณํ ยถา โรคํ ยถา พนฺธนาคารํ ยถา ทาสพฺยํ ยถา กนฺตาร\-ทฺธาน\-มคฺคํ อิเม ปญฺจ นีวรเณ อปฺปหีเน อตฺตนิ สมนุปสฺสติ. เสยฺยถาปิ ภิกฺขเว อาณณฺยํ ยถา อาโรคฺยํ ยถา พนฺธนา\-โมกฺขํ ยถา ภุชิสฺสํ ยถา เขมนฺตภูมิํ. เอวเมว ภิกฺขุ อิเม ปญฺจ นีวรเณ ปหีเน อตฺตนิ สมนุปสฺสติ.}

\proseitem{427}{โส อิเม ปญฺจ นีวรเณ ปหาย เจตโส อุปกฺกิเลเส ปญฺญาย ทุพฺพลีกรเณ วิวิจฺเจว กาเมหิ วิวิจฺจ อกุสเลหิ ธมฺเมหิ สวิตกฺกํ สวิจารํ วิเวกชํ ปีติสุขํ ปฐมํ ฌานํ อุปสมฺปชฺช วิหรติ, โส อิมเมว กายํ วิเวกเชน ปีติสุเขน อภิสนฺเทติ ปริสนฺเทติ ปริปูเรติ ปริปฺผรติ, นาสฺส กิญฺจิ สพฺพาวโต กายสฺส วิเวกเชน ปีติสุเขน อปฺผุฏํ โหติ. เสยฺยถาปิ ภิกฺขเว ทกฺโข นฺหาปโก\csromansharedfootnote{af4513fcb08f01ed}{นหาปโก (สี, สฺยา, กํ, อิ)} วา นฺหาปกนฺเตวาสี วา กํสถาเล นฺหานียจุณฺณานิ\csromansharedfootnote{afd8b3ff1722a88b}{นหานียจุณฺณานิ (สี, สฺยา, กํ, อิ)} อากิริตฺวา อุทเกน ปริปฺโผสกํ ปริปฺโผสกํ สนฺเนยฺย, สายํ นฺหานียปิณฺฑิ สฺเนหานุคตา สฺเนหปเรตา สนฺตรพาหิรา ผุฏา สฺเนเหน น จ ปคฺฆริณี. เอวเมว โข ภิกฺขเว ภิกฺขุ อิมเมว กายํ วิเวกเชน ปีติสุเขน อภิสนฺเทติ ปริสนฺเทติ ปริปูเรติ ปริปฺผรติ, นาสฺส กิญฺจิ สพฺพาวโต กายสฺส วิเวกเชน ปีติสุเขน อปฺผุฏํ โหติ.}

\proseitem{428}{\csromanfolio{345}ปุน จปรํ ภิกฺขเว ภิกฺขุ วิตกฺก\-วิจารานํ วูปสมา อชฺฌตฺตํ สมฺปสาทนํ เจตโส เอโกทิภาวํ อวิตกฺกํ อวิจารํ สมาธิชํ ปีติสุขํ ทุติยํ ฌานํ อุปสมฺปชฺช วิหรติ, โส อิมเมว กายํ สมาธิเชน ปีติสุเขน อภิสนฺเทติ ปริสนฺเทติ ปริปูเรติ ปริปฺผรติ, นาสฺส กิญฺจิ สพฺพาวโต กายสฺส สมาธิเชน ปีติสุเขน อปฺผุฏํ โหติ. เสยฺยถาปิ ภิกฺขเว อุทกรหโท อุพฺภิโททโก\csromansharedfootnote{cc7b4b46377afb13}{อุพฺภิโตทโก (ก)}. ตสฺส เนวสฺส ปุรตฺถิมาย ทิสาย อุทกสฺส อายมุขํ, น ปจฺฉิมาย ทิสาย อุทกสฺส อายมุขํ, น อุตฺตราย ทิสาย อุทกสฺส อายมุขํ, น ทกฺขิณาย ทิสาย อุทกสฺส อายมุขํ, เทโว จ น กาเลนกาลํ สมฺมาธารํ อนุปฺปเวจฺเฉยฺย. อถ โข ตมฺหาว อุทกรหทา สีตา วาริธารา อุพฺภิชฺชิตฺวา ตเมว อุทกรหทํ สีเตน วารินา อภิสนฺเทยฺย ปริสนฺเทยฺย ปริปูเรยฺย, ปริปฺผเรยฺย, นาสฺส กิญฺจิ สพฺพาวโต อุทกรหทสฺส สีเตน วารินา อปฺผุฏํ อสฺส. เอวเมว โข ภิกฺขเว ภิกฺขุ อิมเมว กายํ สมาธิเชน ปีติสุเขน อภิสนฺเทติ ปริสนฺเทติ ปริปูเรติ ปริปฺผรติ, นาสฺส กิญฺจิ สพฺพาวโต กายสฺส สมาธิเชน ปีติสุเขน อปฺผุฏํ โหติ.}

\proseitem{429}{ปุน จปรํ ภิกฺขเว ภิกฺขุ ปีติยา จ วิราคา อุเปกฺขโก จ วิหรติ สโต จ สมฺปชาโน, สุขญฺจ กาเยน ปฏิสํเวเทติ, ยํ ตํ อริยา อาจิกฺขนฺติ “อุเปกฺขโก สติมา สุขวิหารี”ติ, ตติยํ ฌานํ อุปสมฺปชฺช วิหรติ, โส อิมเมว กายํ นิปฺปีติเกน สุเขน อภิสนฺเทติ ปริสนฺเทติ ปริปูเรติ ปริปฺผรติ, นาสฺส กิญฺจิ สพฺพาวโต กายสฺส นิปฺปีติเกน สุเขน อปฺผุฏํ โหติ. เสยฺยถาปิ ภิกฺขเว อุปฺปลินิยํ วา ปทุมินิยํ วา ปุณฺฑรีกินิยํ วา อปฺเปกจฺจานิ อุปฺปลานิ วา ปทุมานิ วา ปุณฺฑรีกานิ วา อุทเก ชาตานิ อุทเก สํวฑฺฒานิ อุทกานุคฺคตานิ อนฺโต\-นิมุคฺค\-โปสีนิ, ตานิ ยาว จคฺคา ยาว จ มูลา สีเตน วารินา อภิสนฺนานิ ปริสนฺนานิ ปริปูรานิ ปริปฺผุฏานิ, นาสฺส\csromansharedfootnote{cf04144ba339ae39}{น เนสํ (สี)} กิญฺจิ สพฺพาวตํ อุปฺปลานํ วา ปทุมานํ วา ปุณฺฑรีกานํ วา สีเตน วารินา อปฺผุฏํ อสฺส. เอวเมว โข ภิกฺขเว ภิกฺขุ อิมเมว กายํ นิปฺปีติเกน สุเขน อภิสนฺเทติ ปริสนฺเทติ ปริปูเรติ ปริปฺผรติ, นาสฺส กิญฺจิ สพฺพาวโต กายสฺส นิปฺปีติเกน สุเขน อปฺผุฏํ โหติ.}

\proseitem{430}{\csromanfolio{346}ปุน จปรํ ภิกฺขเว ภิกฺขุ สุขสฺส จ ปหานา ทุกฺขสฺส จ ปหานา ปุพฺเพว โสมนสฺส\-โทมนสฺสานํ อตฺถงฺคมา อทุกฺขมสุขํ อุเปกฺขา\-สติ\-ปาริสุทฺธิํ จตุตฺถํ ฌานํ อุปสมฺปชฺช วิหรติ, โส อิมเมว กายํ ปริสุทฺเธน เจตสา ปริโยทาเตน ผริตฺวา นิสินฺโน โหติ, นาสฺส กิญฺจิ สพฺพาวโต กายสฺส ปริสุทฺเธน เจตสา ปริโยทาเตน อปฺผุฏํ โหติ. เสยฺยถาปิ ภิกฺขเว ปุริโส โอทาเตน วตฺเถน สสีสํ ปารุเปตฺวา นิสินฺโน อสฺส, นาสฺส กิญฺจิ สพฺพาวโต กายสฺส โอทาเตน วตฺเถน อปฺผุฏํ อสฺส. เอวเมว โข ภิกฺขเว ภิกฺขุ อิมเมว กายํ ปริสุทฺเธน เจตสา ปริโยทาเตน ผริตฺวา นิสินฺโน โหติ. นาสฺส กิญฺจิ สพฺพาวโต กายสฺส ปริสุทฺเธน เจตสา ปริโยทาเตน อปฺผุฏํ โหติ.}

\proseitem{431}{โส เอวํ สมาหิเต จิตฺเต ปริสุทฺเธ ปริโยทาเต อนงฺคเณ วิคตู\-ปกฺกิเลเส มุทุภูเต กมฺมนิเย ฐิเต อาเนญฺชปฺปตฺเต ปุพฺเพ\-นิวาสา\-นุสฺสติ\-ญาณาย จิตฺตํ อภินินฺนาเมติ. โส อเนกวิหิตํ ปุพฺเพนิวาสํ อนุสฺสรติ. เสยฺยถิทํ, เอกมฺปิ ชาติํ เทฺวปิ ชาติโย ฯเปฯ อิติสาการํ สอุทฺเทสํ อเนกวิหิตํ ปุพฺเพนิวาสํ อนุสฺสรติ. เสยฺยถาปิ ภิกฺขเว ปุริโส สกมฺหา คามา อญฺญํ คามํ คจฺเฉยฺย, ตมฺหาปิ คามา อญฺญํ คามํ คจฺเฉยฺย, โส ตมฺหา คามา สกํเยว คามํ ปจฺจาคจฺเฉยฺย. ตสฺส เอวมสฺส “อหํ โข สกมฺหา คามา อมุํ คามํ อคจฺฉิํ\csromansharedfootnote{3f48df4acff37ffe}{อคญฺฉิํ (สี, สฺยา, กํ, อิ)}. ตตฺรปิ เอวํ อฏฺฐาสิํ, เอวํ นิสีทิํ, เอวํ อภาสิํ, เอวํ ตุณฺหี อโหสิํ. ตมฺหาปิ คามา อมุํ คามํ อคจฺฉิํ. ตตฺรปิ เอวํ อฏฺฐาสิํ, เอวํ นิสีทิํ, เอวํ อภาสิํ, เอวํ ตุณฺหี อโหสิํ. โสมฺหิ ตมฺหา คามา สกํเยว คามํ ปจฺจาคโต”ติ. เอวเมว โข ภิกฺขเว ภิกฺขุ อเนกวิหิตํ ปุพฺเพนิวาสํ อนุสฺสรติ. เสยฺยถิทํ, เอกมฺปิ ชาติํ เทฺวปิ ชาติโย ฯเปฯ อิติ สาการํ สอุทฺเทสํ อเนกวิหิตํ ปุพฺเพนิวาสํ อนุสฺสรติ.}

\proseitem{432}{โส เอวํ สมาหิเต จิตฺเต ปริสุทฺเธ ปริโยทาเต อนงฺคเณ วิคตู\-ปกฺกิเลเส มุทุภูเต กมฺมนิเย ฐิเต อาเนญฺชปฺปตฺเต สตฺตานํ จุตูปปาต\-ญาณาย จิตฺตํ อภินินฺนาเมติ. โส ทิพฺเพน จกฺขุนา วิสุทฺเธน อติกฺกนฺต\-มานุสเกน สตฺเต ปสฺสติ จวมาเน อุปปชฺชมาเน หีเน \csromanfolio{347}ปณีเต สุวณฺเณ ทุพฺพณฺเณ สุคเต ทุคฺคเต, ยถากมฺมูปเค สตฺเต ปชานาติ ฯเปฯ. เสยฺยถาปิ ภิกฺขเว เทฺว อคารา สทฺวารา\csromansharedfootnote{808501aeda6edde2}{สนฺนทฺวารา (ก)}. ตตฺถ จกฺขุมา ปุริโส มชฺเฌ ฐิโต ปสฺเสยฺย มนุสฺเส เคหํ ปวิสนฺเตปิ นิกฺขมนฺเตปิ อนุจงฺกม\-นฺเตปิ อนุวิจรนฺเตปิ. เอวเมว โข ภิกฺขเว ภิกฺขุ ทิพฺเพน จกฺขุนา วิสุทฺเธน อติกฺกนฺต\-มานุสเกน สตฺเต ปสฺสติ จวมาเน อุปปชฺชมาเน หีเน ปณีเต สุวณฺเณ ทุพฺพณฺเณ สุคเต ทุคฺคเต, ยถากมฺมูปเค สตฺเต ปชานาติ ฯเปฯ.}

\proseitem{433}{โส เอวํ สมาหิเต จิตฺเต ปริสุทฺเธ ปริโยทาเต อนงฺคเณ วิคตู\-ปกฺกิเลเส มุทุภูเต กมฺมนิเย ฐิเต อาเนญฺชปฺปตฺเต อาสวานํ ขยญาณาย จิตฺตํ อภินินฺนาเมติ. โส อิทํ ทุกฺขนฺติ ยถาภูตํ ปชานาติ, อยํ ทุกฺข\-สมุทโยติ ยถาภูตํ ปชานาติ, อยํ ทุกฺข\-นิโรโธติ ยถาภูตํ ปชานาติ, อยํ ทุกฺขนิโรธ\-คามินี ปฏิปทาติ ยถาภูตํ ปชานาติ. อิเม อาสวาติ ยถาภูตํ ปชานาติ, อยํ อาสว\-สมุทโยติ ยถาภูตํ ปชานาติ, อยํ อาสวนิโรโธติ ยถาภูตํ ปชานาติ, อยํ อาสว\-นิโรธ\-คามินี ปฏิปทาติ ยถาภูตํ ปชานาติ. ตสฺส เอวํ ชานโต เอวํ ปสฺสโต กามาสวาปิ จิตฺตํ วิมุจฺจติ, ภวาสวาปิ จิตฺตํ วิมุจฺจติ, อวิชฺชาสวาปิ จิตฺตํ วิมุจฺจติ, วิมุตฺตสฺมิํ “วิมุตฺตมฺ”อิติ ญาณํ โหติ, “ขีณา ชาติ วุสิตํ พฺรหฺมจริยํ กตํ กรณียํ นาปรํ อิตฺถตฺตายา”ติ ปชานาติ. เสยฺยถาปิ ภิกฺขเว ปพฺพต\-สงฺเขเป อุทกรหโท อจฺโฉ วิปฺปสนฺโน อนาวิโล, ตตฺถ จกฺขุมา ปุริโส ตีเร ฐิโต ปสฺเสยฺย สิปฺปิ\-สมฺพุกมฺปิ\csromansharedfootnote{7598763466180c37}{สิปฺปิกสมฺพุกมฺปิ (สฺยา, กํ, ก)} สกฺขร\-กถลมฺปิ มจฺฉคุมฺพมฺปิ จรนฺตมฺปิ ติฏฺฐนฺตมฺปิ. ตสฺส เอวมสฺส “อยํ โข อุทกรหโท อจฺโฉ วิปฺปสนฺโน อนาวิโล, ตตฺริเม สิปฺปิสมฺพุกาปิ สกฺขรกถลาปิ มจฺฉคุมฺพาปิ จรนฺติปิ ติฏฺฐนฺติปี”ติ. เอวเมว โข ภิกฺขเว ภิกฺขุ อิทํ ทุกฺขนฺติ ยถาภูตํ ปชานาติ ฯเปฯ นาปรํ อิตฺถตฺตายาติ ปชานาติ.}

\proseitem{434}{อยํ วุจฺจติ ภิกฺขเว ภิกฺขุ สมโณ อิติปิ พฺราหฺมโณ อิติปิ นฺหาตโก อิติปิ เวทคู อิติปิ โสตฺติโย อิติปิ อริโย อิติปิ อรหํ อิติปิ. กถญฺจ ภิกฺขเว ภิกฺขุ สมโณ โหติ. สมิตาสฺส โหนฺติ ปาปกา อกุสลา ธมฺมา สํกิเลสิกา โปโนพฺภวิกา สทรา \csromanfolio{348}ทุกฺขวิปากา อายติํ ชาติ\-ชรา\-มรณิยา. เอวํ โข ภิกฺขเว ภิกฺขุ สมโณ โหติ.}

\prose{กถญฺจ ภิกฺขเว ภิกฺขุ พฺราหฺมโณ โหติ. พาหิตาสฺส โหนฺติ ปาปกา อกุสลา ธมฺมา สํกิเลสิกา โปโนพฺภวิกา สทรา ทุกฺขวิปากา อายติํ ชาติ\-ชรา\-มรณิยา. เอวํ โข ภิกฺขเว ภิกฺขุ พฺราหฺมโณ โหติ.}

\prose{กถญฺจ ภิกฺขเว ภิกฺขุ นฺหาตโก\csromansharedfootnote{ee82a07bfaa3ffad}{นหาตโก (สี, สฺยา, กํ, อิ)} โหติ. นฺหาตาสฺส โหนฺติ ปาปกา อกุสลา ธมฺมา สํกิเลสิกา โปโนพฺภวิกา สทรา ทุกฺขวิปากา อายติํ ชาติ\-ชรา\-มรณิยา. เอวํ โข ภิกฺขเว ภิกฺขุ นฺหาตโก โหติ.}

\prose{กถญฺจ ภิกฺขเว ภิกฺขุ เวทคู โหติ. วิทิตาสฺส โหนฺติ ปาปกา อกุสลา ธมฺมา สํกิเลสิกา โปโนพฺภวิกา สทรา ทุกฺขวิปากา อายติํ ชาติ\-ชรา\-มรณิยา. เอวํ โข ภิกฺขเว ภิกฺขุ เวทคู โหติ.}

\prose{กถญฺจ ภิกฺขเว ภิกฺขุ โสตฺติโย โหติ. นิสฺสุตาสฺส โหนฺติ ปาปกา อกุสลา ธมฺมา สํกิเลสิกา โปโนพฺภวิกา สทรา ทุกฺขวิปากา อายติํ ชาติ\-ชรา\-มรณิยา. เอวํ โข ภิกฺขเว ภิกฺขุ โสตฺติโย โหติ.}

\prose{กถญฺจ ภิกฺขเว ภิกฺขุ อริโย โหติ. อารกาสฺส โหนฺติ ปาปกา อกุสลา ธมฺมา สํกิเลสิกา โปโนพฺภวิกา สทรา ทุกฺขวิปากา อายติํ ชาติ\-ชรา\-มรณิยา. เอวํ โข ภิกฺขเว ภิกฺขุ อริโย โหติ.}

\prose{กถญฺจ ภิกฺขเว ภิกฺขุ อรหํ โหติ. อารกาสฺส โหนฺติ ปาปกา อกุสลา ธมฺมา สํกิเลสิกา โปโนพฺภวิกา สทรา ทุกฺขวิปากา อายติํ ชาติ\-ชรา\-มรณิยา. เอวํ โข ภิกฺขเว ภิกฺขุ อรหํ โหตีติ.}

\prosewithcloser{อิทมโวจ ภควา. อตฺตมนา เต ภิกฺขู ภควโต ภาสิตํ อภินนฺทุนฺติ.}{มหา\-อสฺสปุร\-สุตฺตํ นิฏฺฐิตํ นวมํ.}
\csromanmatikaanchor{mtk.72}
\csromantocmarkref{subsection}{10. จูฬอสฺสปุรสุตฺต}{mtk.72}
\bookmark[dest={mtk.72},level=2]{10. จูฬอสฺสปุรสุตฺต}
\csromanmarkh{10. จูฬอสฺสปุร\-สุตฺต (40)}\csromanmarkhii{10. จูฬอสฺสปุร\-สุตฺต}\csromanheadernomark{2}{10. จูฬอสฺสปุร\-สุตฺต (40) 10. จูฬอสฺสปุร\-สุตฺต}
\proseitem{435}{\csromanfolio{349}เอวํ เม สุตํ\csromandash{}เอกํ สมยํ ภควา องฺเคสุ วิหรติ อสฺสปุรํ นาม องฺคานํ นิคโม. ตตฺร โข ภควา ภิกฺขู อามนฺเตสิ “ภิกฺขโว”ติ. “ภทนฺเต”ติ เต ภิกฺขู ภควโต ปจฺจสฺโสสุํ. ภควา เอตทโวจ\csromandash{}“สมณา สมณา”ติ โว ภิกฺขเว ชโน สญฺชานาติ, ตุเมฺห จ ปน “เก ตุเมฺห”ติ ปุฏฺฐา สมานา “สมณามฺหา”ติ ปฏิชานาถ, เตสํ โว ภิกฺขเว เอวํสมญฺญานํ สตํ เอวํปฏิญฺญานํ สตํ “ยา สมณ\-สามีจิ\-ปฺปฏิปทา, ตํ ปฏิปชฺชิสฺสาม, เอวํ โน อยํ อมฺหากํ สมญฺญา จ สจฺจา ภวิสฺสติ ปฏิญฺญา จ ภูตา. เยสญฺจ มยํ จีวร\-ปิณฺฑปาต\-เสนาสน\-คิลานปฺปจฺจย\-เภสชฺช\-ปริกฺขารํ ปริภุญฺชาม, เตสํ เต การา อเมฺหสุ มหปฺผลา ภวิสฺสนฺติ มหานิสํสา, อมฺหากญฺเจวายํ ปพฺพชฺชา อวญฺฌา ภวิสฺสติ สผลา สอุทฺรยา”ติ เอวํ หิ โว ภิกฺขเว สิกฺขิตพฺพํ.}

\proseitem{436}{กถญฺจ ภิกฺขเว ภิกฺขุ น สมณ\-สามีจิ\-ปฺปฏิปทํ ปฏิปนฺโน โหติ. ยสฺส กสฺสจิ ภิกฺขเว ภิกฺขุโน อภิชฺฌาลุสฺส อภิชฺฌา อปฺปหีนา โหติ, พฺยาปนฺน\-จิตฺตสฺส พฺยาปาโท อปฺปหีโน โหติ, โกธนสฺส โกโธ อปฺปหีโน โหติ, อุปนาหิสฺส อุปนาโห อปฺปหีโน โหติ, มกฺขิสฺส มกฺโข อปฺปหีโน โหติ, ปฬาสิสฺส ปฬาโส อปฺปหีโน โหติ, อิสฺสุกิสฺส อิสฺสา อปฺปหีนา โหติ, มจฺฉริสฺส มจฺฉริยํ อปฺปหีนํ โหติ, สฐสฺส สาเฐยฺยํ อปฺปหีนํ โหติ, มายาวิสฺส มายา อปฺปหีนา โหติ, ปาปิจฺฉสฺส ปาปิกา อิจฺฉา อปฺปหีนา โหติ, มิจฺฉา\-ทิฏฺฐิกสฺส มิจฺฉาทิฏฺฐิ อปฺปหีนา โหติ. อิเมสํ โข อหํ ภิกฺขเว สมณมลานํ สมณโทสานํ สมณ\-กสฏานํ อาปายิกานํ ฐานานํ ทุคฺคติ\-เวทนิยานํ อปฺปหานา “น สมณ\-สามีจิ\-ปฺปฏิปทํ ปฏิปนฺโน”ติ วทามิ. เสยฺยถาปิ ภิกฺขเว มตชํ นาม อาวุธชาตํ อุภโตธารํ ปีตนิสิตํ, ตทสฺส สํฆาฏิยา สมฺปารุตํ สมฺปลิ- เวฐิตํ. ตถูปมาหํ ภิกฺขเว อิมสฺส ภิกฺขุโน ปพฺพชฺชํ วทามิ.}

\proseitem{437}{นาหํ ภิกฺขเว สํฆาฏิกสฺส สํฆาฏิธารณมตฺเตน สามญฺญํ วทามิ. นาหํ ภิกฺขเว อเจลกสฺส อเจลกมตฺเตน สามญฺญํ วทามิ. นาหํ ภิกฺขเว รโชชลฺลิกสฺส รโชชลฺลิก\-มตฺเตน สามญฺญํ วทามิ. นาหํ \csromanfolio{350}ภิกฺขเว อุทโกโรหกสฺส อุทโกโรหณ\-มตฺเตน\csromansharedfootnote{51eea462f968d5e9}{อุทโกโรหกมตฺเตน (สี, อิ)} สามญฺญํ วทามิ. นาหํ ภิกฺขเว รุกฺข\-มูลิ\-กสฺส รูกฺขมูลิกมตฺเตน สามญฺญํ วทามิ. นาหํ ภิกฺขเว อพฺโภกาสิกสฺส อพฺโภกาสิก\-มตฺเตน สามญฺญํ วทามิ. นาหํ ภิกฺขเว อุพฺภฏฺฐกสฺส อุพฺภฏฺฐก\-มตฺเตน สามญฺญํ วทามิ. นาหํ ภิกฺขเว ปริยาย\-ภตฺติ\-กสฺส ปริยายภตฺติก\-มตฺเตน สามญฺญํ วทามิ. นาหํ ภิกฺขเว มนฺต\-ชฺฌายกสฺส มนฺตชฺฌายก\-มตฺเตน สามญฺญํ วทามิ. นาหํ ภิกฺขเว ชฏิลกสฺส ชฏาธารณ\-มตฺเตน สามญฺญํ วทามิ.}

\prose{สํฆาฏิกสฺส เจ ภิกฺขเว สํฆาฏิธารณมตฺเตน อภิชฺฌาลุสฺส อภิชฺฌา ปหีเยถ, พฺยาปนฺน\-จิตฺตสฺส พฺยาปาโท ปหีเยถ, โกธนสฺส โกโธ ปหีเยถ, อุปนาหิสฺส อุปนาโห ปหีเยถ, มกฺขิสฺส มกฺโข ปหีเยถ, ปฬาสิสฺส ปฬาโส ปหีเยถ, อิสฺสุกิสฺส อิสฺสา ปหีเยถ, มจฺฉริสฺส มจฺฉริยํ ปหีเยถ, สฐสฺส สาเฐยฺยํ ปหีเยถ, มายาวิสฺส มายา ปหีเยถ, ปาปิจฺฉสฺส ปาปิกา อิจฺฉา ปหีเยถ, มิจฺฉา\-ทิฏฺฐิกสฺส มิจฺฉาทิฏฺฐิ ปหีเยถ. ตเมนํ มิตฺตามจฺจา ญาติสาโลหิตา ชาตเมว นํ สํฆาฏิกํ กเรยฺยุํ, สํฆาฏิกตฺตเมว\csromansharedfootnote{9fbee858d49a4391}{สํฆาฏิกตฺเต เจว (ก)} สมาทเปยฺยุํ, “เอหิ ตฺวํ ภทฺรมุข สํฆาฏิโก โหติ, สํฆาฏิกสฺส เต สโต สํฆาฏิธารณมตฺเตน อภิชฺฌาลุสฺส อภิชฺฌา ปหียิสฺสติ, พฺยาปนฺน\-จิตฺตสฺส พฺยาปาโท ปหียิสฺสติ, โกธนสฺส โกโธ ปหียิสฺสติ, อุปนาหิสฺส อุปนาโห ปหียิสฺสติ, มกฺขิสฺส มกฺโข ปหียิสฺสติ, ปฬาสิสฺส ปฬาโส ปหียิสฺสติ, อิสฺสุกิสฺส อิสฺสา ปหียิสฺสติ, มจฺฉริสฺส มจฺฉริยํ ปหียิสฺสติ, สฐสฺส สาเฐยฺยํ ปหียิสฺสติ, มายาวิสฺส มายา ปหียิสฺสติ, ปาปิจฺฉสฺส ปาปิกา อิจฺฉา ปหียิสฺสติ, มิจฺฉา\-ทิฏฺฐิกสฺส มิจฺฉาทิฏฺฐิ ปหียิสฺสตี”ติ. ยสฺมา จ โข อหํ ภิกฺขเว สํฆาฏิกมฺปิ อิเธกจฺจํ ปสฺสามิ อภิชฺฌาลุํ พฺยาปนฺนจิตฺตํ โกธนํ อุปนาหิํ มกฺขิํ ปฬาสิํ อิสฺสุกิํ มจฺฉริํ สฐํ มายาวิํ ปาปิจฺฉํ มิจฺฉา\-ทิฏฺฐิกํ. ตสฺมา น สํฆาฏิกสฺส สํฆาฏิธารณมตฺเตน สามญฺญํ วทามิ.}

\prose{อเจลกสฺส เจ ภิกฺขเว ฯเปฯ. รโชชลฺลิกสฺส เจ ภิกฺขเว ฯเปฯ. อุทโกโรหกสฺส เจ ภิกฺขเว ฯเปฯ. รุกฺข\-มูลิ\-กสฺส เจ ภิกฺขเว ฯเปฯ. อพฺโภกาสิกสฺส เจ ภิกฺขเว ฯเปฯ. อุพฺภฏฺฐกสฺส เจ ภิกฺขเว ฯเปฯ. ปริยาย\-ภตฺติ\-กสฺส เจ ภิกฺขเว ฯเปฯ. มนฺต\-ชฺฌายกสฺส เจ ภิกฺขเว \csromanfolio{351}ชฏิลกสฺส เจ ภิกฺขเว ชฏาธารณ\-มตฺเตน อภิชฺฌาลุสฺส อภิชฺฌา ปหีเยถ, พฺยาปนฺน\-จิตฺตสฺส พฺยาปาโท ปหีเยถ, โกธนสฺส โกโธ ปหีเยถ, อุปนาหิสฺส อุปนาโห ปหีเยถ, มกฺขิสฺส มกฺโข ปหีเยถ, ปฬาสิสฺส ปฬาโส ปหีเยถ, อิสฺสุกิสฺส อิสฺสา ปหีเยถ, มจฺฉริสฺส มจฺฉริยํ ปหีเยถ, สฐสฺส สาเฐยฺยํ ปหีเยถ, มายาวิสฺส มายา ปหีเยถ, ปาปิจฺฉสฺส ปาปิกา อิจฺฉา ปหีเยถ, มิจฺฉา\-ทิฏฺฐิกสฺส มิจฺฉาทิฏฺฐิ ปหีเยถ. ตเมนํ มิตฺตามจฺจา ญาติสาโลหิตา ชาตเมว นํ ชฏิลกํ กเรยฺยํ, ชฏิลกตฺตเมว\csromansharedfootnote{4bb1eaca2a42f726}{ชฏิลกตฺเต เจว (ก)} สมาทเปยฺยุํ, “เอหิ ตฺวํ ภทฺรมุข ชฏิลโก โหหิ, ชฏิลกสฺส เต สโต ชฏาธารณ\-มตฺเตน อภิชฺฌาลุสฺส อภิชฺฌา ปหียิสฺสติ, พฺยาปนฺน\-จิตฺตสฺส พฺยาปาโท ปหียิสฺสติ, โกธนสฺส โกโธ ปหียิสฺสติ ฯเปฯ ปาปิจฺฉสฺส ปาปิกา อิจฺฉา ปหียิสฺสติ, มิจฺฉา\-ทิฏฺฐิกสฺส มิจฺฉาทิฏฺฐิ ปหียิสฺสตี”ติ. ยสฺมา จ โข อหํ ภิกฺขเว ชฏิลกมฺปิ อิเธกจฺจํ ปสฺสามิ อภิชฺฌาลุํ พฺยาปนฺนจิตฺตํ โกธนํ อุปนาหิํ มกฺขิํ ปลาสิํ อิสฺสุกิํ มจฺฉริํ สฐํ มายาวิํ ปาปิจฺฉํ มิจฺฉาทิฏฺฐิํ. ตสฺมา น ชฏิลกสฺส ชฏาธารณ\-มตฺเตน สามญฺญํ วทามิ.}

\proseitem{438}{กถญฺจ ภิกฺขเว ภิกฺขุ สมณ\-สามีจิ\-ปฺปฏิปทํ ปฏิปนฺโน โหติ. ยสฺส กสฺสจิ ภิกฺขเว ภิกฺขุโน อภิชฺฌาลุสฺส อภิชฺฌา ปหีนา โหติ, พฺยาปนฺน\-จิตฺตสฺส พฺยาปาโท ปหีโน โหติ, โกธนสฺส โกโธ ปหีโน โหติ, อุปนาหิสฺส อุปนาโห ปหีโน โหติ, มกฺขิสฺส มกฺโข ปหีโน โหติ, ปฬาสิสฺส ปฬาโส ปหีโน โหติ, อิสฺสุกิสฺส อิสฺสา ปหีนา โหติ, มจฺฉริสฺส มจฺฉริยํ ปหีนํ โหติ, สฐสฺส สาเฐยฺยํ ปหีนํ โหติ, มายาวิสฺส มายา ปหีนา โหติ, ปาปิจฺฉสฺส ปาปิกา อิจฺฉา ปหีนา โหติ, มิจฺฉา\-ทิฏฺฐิกสฺส มิจฺฉาทิฏฺฐิ ปหีนา โหติ. อิเมสํ โข อหํ ภิกฺขเว สมณมลานํ สมณโทสานํ สมณ\-กสฏานํ อาปายิกานํ ฐานานํ ทุคฺคติ\-เวทนิยานํ ปหานา “สมณ\-สามีจิ\-ปฺปฏิปทํ ปฏิปนฺโน”ติ วทามิ. โส สพฺเพหิ อิเมหิ ปาปเกหิ อกุสเลหิ ธมฺเมหิ วิสุทฺธม\-ตฺตานํ สมนุปสฺสติ ( )\csromansharedfootnote{3d0464bbd2c549c2}{(วิมุตฺตมตฺตานํ สมนุปสฺสติ) (สี, สฺยา, กํ, อิ)}. ตสฺส สพฺเพหิ อิเมหิ ปาปเกหิ อกุสเลหิ ธมฺเมหิ วิสุทฺธม\-ตฺตานํ}

\prose{\csromanfolio{352}สมนุปสฺสโต ( )\csromansharedfootnote{bb9b2cb7b6a67685}{(วิมุตฺตมตฺตานํ สมนุปสฺสโต) (สี, สฺยา, กํ, อิ) 2-2. ตมหํ (ก)} ปาโมชฺชํ ชายติ, ปมุทิตสฺส ปีติ ชายติ, ปีติมนสฺส กาโย ปสฺสมฺภติ, ปสฺสทฺธกาโย สุขํ เวเทติ, สุขิโน จิตฺตํ สมาธิยติ.}

\prose{โส เมตฺตา\-สหคเตน เจตสา เอกํ ทิสํ ผริตฺวา วิหรติ. ตถา ทุติยํ. ตถา ตติยํ. ตถา จตุตฺถํ. อิติ อุทฺธมโธ ติริยํ สพฺพธิ สพฺพตฺตตาย สพฺพาวนฺตํ โลกํ เมตฺตา\-สหคเตน เจตสา วิปุเลน มหคฺคเตน อปฺปมาเณน อเวเรน อพฺยาพชฺเฌน ผริตฺวา วิหรติ. กรุณา\-สหคเตน เจตสา ฯเปฯ. มุทิตา\-สหคเตน เจตสา ฯเปฯ. อุเปกฺขา\-สหคเตน เจตสา เอกํ ทิสํ ผริตฺวา วิหรติ. ตถา ทุติยํ. ตถา ตติยํ. ตถา จตุตฺถํ. อิติ อุทฺธมโธ ติริยํ สพฺพธิ สพฺพตฺตตาย สพฺพาวนฺตํ โลกํ อุเปกฺขา\-สหคเตน เจตสา วิปุเลน มหคฺคเตน อปฺปมาเณน อเวเรน อพฺยาพชฺเฌน ผริตฺวา วิหรติ. เสยฺยถาปิ ภิกฺขเว โปกฺขรณี อจฺโฉทกา สาโตทกา สีโตทกา เสตกา สุปติตฺถา รมณียา, ปุรตฺถิมาย เจปิ ทิสาย ปุริโส อาคจฺเฉยฺย ฆมฺมาภิตตฺโต ฆมฺมปเรโต กิลนฺโต ตสิโต ปิปาสิโต, โส ตํ โปกฺขรณิํ อาคมฺม วิเนยฺย อุทกปิปาสํ, วิเนยฺย ฆมฺม\-ปริฬาหํ. ปจฺฉิมาย เจปิ ทิสาย ปุริโส อาคจฺเฉยฺย. อุตฺตราย เจปิ ทิสาย ปุริโส อาคจฺเฉยฺย. ทกฺขิณาย เจปิ ทิสาย ปุริโส อาคจฺเฉยฺย. ยโต กุโต เจปิ นํ ปุริโส อาคจฺเฉยฺย ฆมฺมาภิตตฺโต ฆมฺมปเรโต กิลนฺโต ตสิโต ปิปาสิโต, โส ตํ โปกฺขรณิํ อาคมฺม วิเนยฺย อุทกปิปาสํ, วิเนยฺย ฆมฺม\-ปริฬาหํ. เอวเมว โข ภิกฺขเว ขตฺติยกุลา เจปิ อคารสฺมา อนคาริยํ ปพฺพชิโต โหติ, โส จ ตถาคต\-ปฺปเวทิตํ ธมฺมวินยํ อาคมฺม เอวํ เมตฺตํ กรุณํ มุทิตํ อุเปกฺขํ ภาเวตฺวา ลภติ อชฺฌตฺตํ วูปสมํ,2 อชฺฌตฺตํ วูปสมา2 “สมณ\-สามีจิ\-ปฺปฏิปทํ ปฏิปนฺโน”ติ วทามิ. พฺราหฺมณกุลา เจปิ ฯเปฯ. เวสฺสกุลา เจปิ ฯเปฯ. สุทฺทกุลา เจปิ ฯเปฯ. ยสฺมา กสฺมา เจปิ กุลา อคารสฺมา อนคาริยํ ปพฺพชิโต โหติ, โส จ ตถาคต\-ปฺปเวทิตํ ธมฺมวินยํ อาคมฺม เอวํ เมตฺตํ กรุณํ มุทิตํ อุเปกฺขํ ภาเวตฺวา ลภติ อชฺฌตฺตํ วูปสมํ, อชฺฌตฺตํ วูปสมา “สมณ\-สามีจิ\-ปฺปฏิปทํ ปฏิปนฺโน”ติ วทามิ.}

\csromanmatikaanchor{mtk.73}
\csromantocmarkref{subsection}{อุทฺทานคาถา}{mtk.73}
\bookmark[dest={mtk.73},level=2]{อุทฺทานคาถา}
\prose{\csromanfolio{353}ขตฺติยกุลา เจปิ อคารสฺมา อนคาริยํ ปพฺพชิโต โหติ, โส จ อาสวานํ ขยา อนาสวํ เจโตวิมุตฺติํ ปญฺญาวิมุตฺติํ ทิฏฺเฐว ธมฺเม สยํ อภิญฺญา สจฺฉิกตฺวา อุปสมฺปชฺช วิหรติ, อาสวานํ ขยา สมโณ โหติ. พฺราหฺมณกุลา เจปิ ฯเปฯ. เวสฺสกุลา เจปิ ฯเปฯ. สุทฺทกุลา เจปิ ฯเปฯ. ยสฺมา กสฺมา เจปิ กุลา อคารสฺมา อนคาริยํ ปพฺพชิโต โหติ, โส จ อาสวานํ ขยา อนาสวํ เจโตวิมุตฺติํ ปญฺญาวิมุตฺติํ ทิฏฺเฐว ธมฺเม สยํ อภิญฺญา สจฺฉิกตฺวา อุปสมฺปชฺช วิหรติ, อาสวานํ ขยา สมโณ โหตีติ.}

\prosewithcloser{อิทมโวจ ภควา. อตฺตมนา เต ภิกฺขู ภควโต ภาสิตํ อภินนฺทุนฺติ.}{จูฬอสฺสปุร\-สุตฺตํ นิฏฺฐิตํ ทสมํ.}
\nitthitammid{มหา\-ยมก\-วคฺโค นิฏฺฐิโต จตุตฺโถ.}
\csromancenter{\textbf{ตสฺสุทฺทานํ}}
\setlength{\csromangathapagewidth}{0pt}
\setlength{\csromangathawakwidth}{0pt}
\csromangathameasuregroup{คิญฺชกสาลวนํ ปริหริตุํ, \\ มุขวณฺณปสีทนตาปินฺโท,}{\csromangathabat{คิญฺชกสาลวนํ ปริหริตุํ,}{ปญฺญวโต ปุน สจฺจกนิเสโธ.} \\ \csromangathabat{มุขวณฺณปสีทนตาปินฺโท,}{เกวฏฺฏอสฺสปุรชฏิเลน.}}
\csromangathasetleft{คิญฺชกสาลวนํ ปริหริตุํ, \\ มุขวณฺณปสีทนตาปินฺโท,}
\csromangathagroup{\csromangathastanza{\csromangathabat{คิญฺชกสาลวนํ ปริหริตุํ,}{ปญฺญวโต ปุน สจฺจกนิเสโธ.} \\ \csromangathabat{มุขวณฺณปสีทนตาปินฺโท,}{เกวฏฺฏอสฺสปุรชฏิเลน.}}}

\csromanmatikaanchor{mtk.74}
\csromantocmarknopage{chapter}{5. จูฬยมกวคฺค}{mtk.74}
\bookmark[dest={mtk.74},level=0]{5. จูฬยมกวคฺค}
\chapterheadpageread{5. จูฬยมก\-วคฺค}
\csromanmatikaanchor{mtk.75}
\csromantocmarkref{subsection}{1. สาเลยฺยกสุตฺต}{mtk.75}
\bookmark[dest={mtk.75},level=2]{1. สาเลยฺยกสุตฺต}
\csromanheader{2}{1. \textbf{สาเลยฺยกสุตฺต}}
\proseitem{439}{\csromanfolio{354}เอวํ เม สุตํ\csromandash{}เอกํ สมยํ ภควา โกสเลสุ จาริกํ จรมาโน มหตา ภิกฺขุ\-สํเฆน สทฺธิํ เยน สาลา นาม โกสลานํ พฺราหฺมณคาโม ตทวสริ. อสฺโสสุํ โข สาเลยฺยกา พฺราหฺมณ\-คหปติกา “สมโณ ขลุ โภ โคตโม สกฺยปุตฺโต สกฺยกุลา ปพฺพชิโต โกสเลสุ จาริกํ จรมาโน มหตา ภิกฺขุ\-สํเฆน สทฺธิํ สาลํ อนุปฺปตฺโต. ตํ โข ปน ภวนฺตํ โคตมํ เอวํ กลฺยาโณ กิตฺติสทฺโท อพฺภุคฺคโต, อิติปิ โส ภควา อรหํ สมฺมาสมฺพุทฺโธ วิชฺชา\-จรณ\-สมฺปนฺโน สุคโต โลกวิทู อนุตฺตโร ปุริส\-ทมฺม\-สารถิ สตฺถา เทวมนุสฺสานํ พุทฺโธ ภควา. โส อิมํ โลกํ สเทวกํ สมารกํ สพฺรหฺมกํ สสฺสมณ\-พฺราหฺมณิํ ปชํ สเทวมนุสฺสํ สยํ อภิญฺญา สจฺฉิกตฺวา ปเวเทติ. โส ธมฺมํ เทเสติ อาทิกลฺยาณํ มชฺเฌกลฺยาณํ ปริโยสาน\-กลฺยาณํ สาตฺถํ สพฺยญฺชนํ เกวล\-ปริปุณฺณํ ปริสุทฺธํ พฺรหฺมจริยํ ปกาเสติ. สาธุ โข ปน ตถารูปานํ อรหตํ ทสฺสนํ โหตี”ติ.}

\prose{อถ โข สาเลยฺยกา พฺราหฺมณ\-คหปติกา เยน ภควา เตนุ\-ปสงฺกมิํสุ, อุปสงฺกมิตฺวา อปฺเปกจฺเจ ภควนฺตํ อภิวาเทตฺวา เอกมนฺตํ นิสีทิํสุ, อปฺเปกจฺเจ ภควตา สทฺธิํ สมฺโมทิํสุ, สมฺโมทนียํ กถํ สารณียํ วีติสาเรตฺวา เอกมนฺตํ นิสีทิํสุ, อปฺเปกจฺเจ เยน ภควา เตนญฺชลิํ ปณาเมตฺวา เอกมนฺตํ นิสีทิํสุ, อปฺเปกจฺเจ ภควโต สนฺติเก นามโคตฺตํ สาเวตฺวา เอกมนฺตํ นิสีทิํสุ, อปฺเปกจฺเจ ตุณฺหีภูตา เอกมนฺตํ นิสีทิํสุ. เอกมนฺตํ นิสินฺนา โข สาเลยฺยกา พฺราหฺมณ\-คหปติกา ภควนฺตํ เอตทโวจุํ “โก นุ โข โภ โคตม เหตุ โก ปจฺจโย, เยน มิเธกจฺเจ สตฺตา กายสฺส เภทา ปรํ มรณา อปายํ ทุคฺคติํ วินิปาตํ นิรยํ อุปปชฺชนฺติ. โก ปน โภ โคตม เหตุ โก ปจฺจโย, เยน มิเธกจฺเจ สตฺตา กายสฺส เภทา ปรํ มรณา สุคติํ สคฺคํ โลกํ อุปปชฺชนฺตี”ติ.}

\prose{อธมฺมจริยา\-วิสมจริยา\-เหตุ โข คหปตโย เอวมิเธกจฺเจ สตฺตา กายสฺส เภทา ปรํ มรณา อปายํ ทุคฺคติํ วินิปาตํ นิรยํ อุปปชฺชนฺติ. \csromanfolio{355}ธมฺมจริยา\-สมจริยา\-เหตุ โข คหปตโย เอวมิเธกจฺเจ สตฺตา กายสฺส เภทา ปรํ มรณา สุคติํ สคฺคํ โลกํ อุปปชฺชนฺตีติ.}

\prose{น โข มยํ อิมสฺส โภโต โคตมสฺส สํขิตฺเตน ภาสิตสฺส วิตฺถาเรน อตฺถํ อวิภตฺตสฺส วิตฺถาเรน อตฺถํ อาชานาม. สาธุ โน ภวํ โคตโม ตถา ธมฺมํ เทเสตุ, ยถา มยํ อิมสฺส โภโต โคตมสฺส สํขิตฺเตน ภาสิตสฺส วิตฺถาเรน อตฺถํ อวิภตฺตสฺส วิตฺถาเรน อตฺถํ อาชาเนยฺยามาติ. เตน หิ คหปตโย สุณาถ สาธุกํ มนสิ กโรถ ภาสิสฺสามีติ. “เอวํ โภ”ติ โข สาเลยฺยกา พฺราหฺมณ\-คหปติกา ภควโต ปจฺจสฺโสสุํ. ภควา เอตทโวจ\csromandash{}}

\proseitem{440}{ติวิธํ โข คหปตโย กาเยน อธมฺมจริยา\-วิสมจริยา โหติ, จตุพฺพิธํ วาจาย อธมฺมจริยา\-วิสมจริยา โหติ, ติวิธํ มนสา อธมฺมจริยา\-วิสมจริยา โหติ.}

\prose{กถญฺจ คหปตโย ติวิธํ กาเยน อธมฺมจริยา\-วิสมจริยา โหติ. อิธ คหปตโย เอกจฺโจ ปาณาติปาตี โหติ, ลุทฺโท\csromansharedfootnote{4f5f6465f81d6a32}{ลุทฺโท ทารุโณ (ก) ฏีกา โอโลเกตพฺพา.} โลหิตปาณิ หตปฺปหเต นิวิฏฺโฐ อทยาปนฺโน ปาณภูเตสุ\csromansharedfootnote{9c7c48ea0b825be2}{สพฺพปาณภูเตสุ (สฺยา, กํ, ก)}.}

\prose{อทินฺนาทายี โข ปน โหติ, ยํ ตํ ปรสฺส ปรวิตฺตู\-ปกรณํ คามคตํ วา อรญฺญคตํ วา, ตํ อทินฺนํ เถยฺยสงฺขาตํ อาทาตา โหติ.}

\prose{กาเมสุ\-มิจฺฉาจารี โข ปน โหติ, ยา ตา มาตุรกฺขิตา ปิตุรกฺขิตา มาตาปิตุ\-รกฺขิตา ภาตุรกฺขิตา ภคินิ\-รกฺขิตา ญาติรกฺขิตา โคตฺตรกฺขิตา ธมฺมรกฺขิตา สสฺสามิกา สปริทณฺฑา อนฺตมโส มาลาคุฬปริกฺขิตฺตาปิ, ตถารูปาสุ จาริตฺตํ อาปชฺชิตา โหติ. เอวํ โข คหปตโย ติวิธํ กาเยน อธมฺมจริยา\-วิสมจริยา โหติ.}

\prose{กถญฺจ คหปตโย จตุพฺพิธํ วาจาย อธมฺมจริยา\-วิสมจริยา โหติ. อิธ คหปตโย เอกจฺโจ มุสาวาที โหติ, สภาคโต วา ปริสาคโต วา ญาติมชฺฌคโต วา ปูคมชฺฌคโต วา ราชกุล\-มชฺฌคโต วา อภินีโต สกฺขิปุฏฺโฐ “เอหมฺโภ ปุริส ยํ ชานาสิ, ตํ \csromanfolio{356}วเทหี”ติ โส อชานํ วา อาห “ชานามี”ติ, ชานํ วา อาห “น ชานามี”ติ, อปสฺสํ วา อาห “ปสฺสามี”ติ, ปสฺสํ วา อาห “น ปสฺสามี”ติ\csromansharedfootnote{cdadf7b82a467bd3}{โส อาห อชานํ วา อหํ ชานามีติ ชานํ วา อหํ น ชานามีติ อปสฺสํ วา อหํ ปสฺสามีติ ปสฺสํ วา อหํ น ปสฺสามีติ (ก)}. อิติ อตฺตเหตุ วา ปรเหตุ วา อามิส\-กิญฺจิกฺข\-เหตุ วา สมฺปชานมุสา ภาสิตา โหติ.}

\prose{ปิสุณวาโจ โข ปน โหติ, อิโต สุตฺวา อมุตฺร อกฺขาตา อิเมสํ เภทาย, อมุตฺร วา สุตฺวา อิเมสํ อกฺขาตา อมูสํ เภทาย. อิติ สมคฺคานํ วา เภตฺตา\csromansharedfootnote{14c78f5279890028}{เภทกา (ก), เภเทตา (สฺยา, กํ), ตทฏฺฐกถายํ ปน เภตฺตาติ ทิสฺสติ.} ภินฺนานํ วา อนุปฺปทาตา, วคฺคาราโม วคฺครโต วคฺคนนฺที วคฺคกรณิํ วาจํ ภาสิตา โหติ.}

\prose{ผรุสวาโจ โข ปน โหติ. ยา สา วาจา อณฺฑกา\csromansharedfootnote{e25bf23b695957dd}{กณฺฑกา (ก)} กกฺกสา ปรกฏุกา ปราภิสชฺชนี โกธสามนฺตา อสมาธิสํวตฺตนิกา, ตถารูปิํ วาจํ ภาสิตา โหติ.}

\prose{สมฺผปฺปลาปี โข ปน โหติ. อกาลวาที อภูตวาที อนตฺถวาที อธมฺมวาที อวินยวาที, อนิธานวติํ วาจํ ภาสิตา โหติ อกาเลน อนปเทสํ อปริยนฺตวติํ อนตฺถ\-สํหิตํ. เอวํ โข คหปตโย จตุพฺพิธํ วาจาย อธมฺมจริยา\-วิสมจริยา โหติ.}

\prose{กถญฺจ คหปตโย ติวิธํ มนสา อธมฺมจริยา\-วิสมจริยา โหติ. อิธ คหปตโย เอกจฺโจ อภิชฺฌาลุ โหติ. ยํ ตํ ปรสฺส ปรวิตฺตู\-ปกรณํ, ตํ อภิชฺฌาตา โหติ “อโห วต ยํ ปรสฺส, ตํ มมสฺสา”ติ.}

\prose{พฺยาปนฺนจิตฺโต โข ปน โหติ, ปทุฏฺฐ\-มนสงฺกปฺโป “อิเม สตฺตา หญฺญนฺตุ วา วชฺฌนฺตุ วา อุจฺฉิชฺชนฺตุ วา วินสฺสนฺตุ วา มา วา อเหสุนฺ”ติ\csromansharedfootnote{cbd67f2c14d3d0c8}{มา วา อเหสุํ อิติ วาติ (สี, อิ, ก)}.}

\prose{มิจฺฉาทิฏฺฐิโก โข ปน โหติ วิปรีต\-ทสฺสโน “นตฺถิ ทินฺนํ, นตฺถิ ยิฏฺฐํ, นตฺถิ หุตํ, นตฺถิ สุกต\-ทุกฺกฏานํ กมฺมานํ ผลํ วิปาโก, นตฺถิ อยํ โลโก, นตฺถิ ปโร โลโก, นตฺถิ มาตา, นตฺถิ ปิตา, นตฺถิ สตฺตา \csromanfolio{357}โอปปาติกา, นตฺถิ โลเก สมณพฺราหฺมณา สมฺมคฺคตา สมฺมาปฏิปนฺนา, เย อิมญฺจ โลกํ ปรญฺจ โลกํ สยํ อภิญฺญา สจฺฉิกตฺวา ปเวเทนฺตี”ติ. เอวํ โข คหปตโย ติวิธํ มนสา อธมฺมจริยา\-วิสมจริยา โหติ.}

\prose{เอวํ อธมฺมจริยา\-วิสมจริยา\-เหตุ โข คหปตโย เอวมิเธกจฺเจ สตฺตา กายสฺส เภทา ปรํ มรณา อปายํ ทุคฺคติํ วินิปาตํ นิรยํ อุปปชฺชนฺติ.}

\proseitem{441}{ติวิธํ โข คหปตโย กาเยน ธมฺมจริยา\-สมจริยา โหติ. จตุพฺพิธํ วาจาย ธมฺมจริยา\-สมจริยา โหติ. ติวิธํ มนสา ธมฺมจริยา\-สมจริยา โหติ.}

\prose{กถญฺจ คหปตโย ติวิธํ กาเยน ธมฺมจริยา\-สมจริยา โหติ. อิธ คหปตโย เอกจฺโจ ปาณาติปาตํ ปหาย ปาณาติปาตา ปฏิวิรโต โหติ นิหิตทณฺโฑ นิหิตสตฺโถ ลชฺชี ทยาปนฺโน สพฺพปาณภูต\-หิตา\-นุกมฺปี วิหรติ.}

\prose{อทินฺนาทานํ ปหาย อทินฺนาทานา ปฏิวิรโต โหติ, ยํ ตํ ปรสฺส ปรวิตฺตู\-ปกรณํ คามคตํ วา อรญฺญคตํ วา, ตํ นาทินฺนํ เถยฺยสงฺขาตํ อาทาตา โหติ.}

\prose{กาเมสุ\-มิจฺฉาจารํ ปหาย กาเมสุ\-มิจฺฉาจารา ปฏิวิรโต โหติ, ยา ตา มาตุรกฺขิตา ปิตุรกฺขิตา มาตาปิตุ\-รกฺขิตา ภาตุรกฺขิตา ภคินิ\-รกฺขิตา ญาติรกฺขิตา โคตฺตรกฺขิตา ธมฺมรกฺขิตา สสฺสามิกา สปริทณฺฑา อนฺตมโส มาลาคุฬปริกฺขิตฺตาปิ, ตถารูปาสุ น จาริตฺตํ อาปชฺชิตา โหติ. เอวํ โข คหปตโย ติวิธํ กาเยน ธมฺมจริยา\-สมจริยา โหติ.}

\prose{กถญฺจ คหปตโย จตุพฺพิธํ วาจาย ธมฺมจริยา\-สมจริยา โหติ. อิธ คหปตโย เอกจฺโจ มุสาวาทํ ปหาย มุสาวาทา ปฏิวิรโต โหติ, สภาคโต วา ปริสาคโต วา ญาติมชฺฌคโต วา ปูคมชฺฌคโต วา ราชกุล\-มชฺฌคโต วา อภินีโต สกฺขิปุฏฺโฐ “เอหมฺโภ ปุริส ยํ ชานาสิ, ตํ วเทหี”ติ โส อชานํ วา อาห “น ชานามี”ติ, ชานํ วา อาห “ชานามี”ติ, อปสฺสํ วา อาห “น ปสฺสามี”ติ, ปสฺสํ วา \csromanfolio{358}อาห “ปสฺสามี”ติ. อิติ อตฺตเหตุ วา ปรเหตุ วา อามิส\-กิญฺจิกฺข\-เหตุ วา น สมฺปชานมุสา ภาสิตา โหติ.}

\prose{ปิสุณํ วาจํ ปหาย ปิสุณาย วาจาย ปฏิวิรโต โหติ, อิโต สุตฺวา น อมุตฺร อกฺขาตา อิเมสํ เภทาย, อมุตฺร วา สุตฺวา น อิเมสํ อกฺขาตา อมูสํ เภทาย, อิติ ภินฺนานํ วา สนฺธาตา สหิตานํ วา อนุปฺปทาตา, สมคฺคาราโม สมคฺครโต สมคฺคนนฺที สมคฺคกรณิํ วาจํ ภาสิตา โหติ.}

\prose{ผรุสํ วาจํ ปหาย ผรุสาย วาจาย ปฏิวิรโต โหติ, ยา สา วาจา เนลา กณฺณสุขา เปมนียา หทยงฺคมา โปรี พหุชนกนฺตา พหุชนมนาปา, ตถารูปิํ วาจํ ภาสิตา โหติ.}

\prose{สมฺผปฺปลาปํ ปหาย สมฺผปฺปลาปา ปฏิวิรโต โหติ กาลวาที ภูตวาที อตฺถวาที ธมฺมวาที วินยวาที, นิธานวติํ วาจํ ภาสิตา โหติ กาเลน สาปเทสํ ปริยนฺตวติํ อตฺถสํหิตํ. เอวํ โข คหปตโย จตุพฺพิธํ วาจาย ธมฺมจริยา\-สมจริยา โหติ.}

\prose{กถญฺจ คหปตโย ติวิธํ มนสา ธมฺมจริยา\-สมจริยา โหติ. อิธ คหปตโย เอกจฺโจ อนภิชฺฌาลุ โหติ, ยํ ตํ ปรสฺส ปรวิตฺตู\-ปกรณํ, ตํ นาภิชฺฌาตา โหติ “อโห วต ยํ ปรสฺส, ตํ มมสฺสา”ติ.}

\prose{อพฺยาปนฺนจิตฺโต โข ปน โหติ อปฺปทุฏฺฐมนสงฺกปฺโป “อิเม สตฺตา อเวรา อพฺยาพชฺฌา อนีฆา สุขี อตฺตานํ ปริหรนฺตู”ติ.}

\prose{สมฺมาทิฏฺฐิโก โข ปน โหติ อวิปรีต\-ทสฺสโน “อตฺถิ ทินฺนํ, อตฺถิ ยิฏฺฐํ, อตฺถิ หุตํ, อตฺถิ สุกต\-ทุกฺกฏานํ กมฺมานํ ผลํ วิปาโก, อตฺถิ อยํ โลโก, อตฺถิ ปโร โลโก, อตฺถิ มาตา, อตฺถิ ปิตา, อตฺถิ สตฺตา โอปปาติกา, อตฺถิ โลเก สมณพฺราหฺมณา สมฺมคฺคตา สมฺมาปฏิปนฺนา, เย อิมญฺจ โลกํ ปรญฺจ โลกํ สยํ อภิญฺญา สจฺฉิกตฺวา ปเวเทนฺตี”ติ. เอวํ โข คหปตโย ติวิธํ มนสา ธมฺมจริยา\-สมจริยา โหติ.}

\prose{เอวํ ธมฺมจริยา\-สมจริยา\-เหตุ โข คหปตโย เอวมิเธกจฺเจ สตฺตา กายสฺส เภทา ปรํ มรณา สุคติํ สคฺคํ โลกํ อุปปชฺชนฺติ.}

\proseitem{442}{\csromanfolio{359}อากงฺเขยฺย เจ คหปตโย ธมฺมจารี สมจารี “อโห วตาหํ กายสฺส เภทา ปรํ มรณา ขตฺติย\-มหาสาลานํ สหพฺยตํ อุปปชฺเชยฺยนฺ”ติ. ฐานํ โข ปเนตํ วิชฺชติ, ยํ โส กายสฺส เภทา ปรํ มรณา ขตฺติย\-มหาสาลานํ สหพฺยตํ อุปปชฺเชยฺย. ตํ กิสฺส เหตุ, ตถา หิ โส ธมฺมจารี สมจารี.}

\prose{อากงฺเขยฺย เจ คหปตโย ธมฺมจารี สมจารี “อโห วตาหํ กายสฺส เภทา ปรํ มรณา พฺราหฺมณ\-มหาสาลานํ ฯเปฯ คหปติ\-มหาสาลานํ สหพฺยตํ อุปปชฺเชยฺยนฺ”ติ. ฐานํ โข ปเนตํ วิชฺชติ, ยํ โส กายสฺส เภทา ปรํ มรณา คหปติ\-มหาสาลานํ สหพฺยตํ อุปปชฺเชยฺย. ตํ กิสฺส เหตุ, ตถา หิ โส ธมฺมจารี สมจารี.}

\prose{อากงฺเขยฺย เจ คหปตโย ธมฺมจารี สมจารี “อโห วตาหํ กายสฺส เภทา ปรํ มรณา จาตุมหาราชิกานํ เทวานํ สหพฺยตํ อุปปชฺเชยฺยนฺ”ติ. ฐานํ โข ปเนตํ วิชฺชติ, ยํ โส กายสฺส เภทา ปรํ มรณา จาตุมหาราชิกานํ เทวานํ สหพฺยตํ อุปปชฺเชยฺย. ตํ กิสฺส เหตุ, ตถา หิ โส ธมฺมจารี สมจารี.}

\prose{อากงฺเขยฺย เจ คหปตโย ธมฺมจารี สมจารี “อโห วตาหํ กายสฺส เภทา ปรํ มรณา ตาวติํสานํ เทวานํ ฯเปฯ ยามานํ เทวานํ ฯเปฯ ตุสิตานํ เทวานํ ฯเปฯ นิมฺมานรตีนํ เทวานํ ฯเปฯ ปรนิมฺมิต\-วส\-วตฺตีนํ เทวานํ ฯเปฯ พฺรหฺม\-กายิกานํ เทวานํ สหพฺยตํ อุปปชฺเชยฺยนฺ”ติ. ฐานํ โข ปเนตํ วิชฺชติ, ยํ โส กายสฺส เภทา ปรํ มรณา พฺรหฺม\-กายิกานํ เทวานํ สหพฺยตํ อุปปชฺเชยฺย. ตํ กิสฺส เหตุ, ตถา หิ โส ธมฺมจารี สมจารี.}

\prose{อากงฺเขยฺย เจ คหปตโย ธมฺมจารี สมจารี “อโห วตาหํ กายสฺส เภทา ปรํ มรณา อาภานํ เทวานํ สหพฺยตํ อุปปชฺเชยฺยนฺ”ติ. ฐานํ โข ปเนตํ วิชฺชติ, ยํ โส กายสฺส เภทา ปรํ มรณา อาภานํ เทวานํ สหพฺยตํ อุปปชฺเชยฺย. ตํ กิสฺส เหตุ, ตถา หิ โส ธมฺมจารี สมจารี.}

\prose{อากงฺเขยฺย เจ คหปตโย ธมฺมจารี สมจารี “อโห วตาหํ กายสฺส เภทา ปรํ มรณา ปริตฺตาภานํ เทวานํ ฯเปฯ อปฺปมาณาภานํ เทวานํ ฯเปฯ อาภสฺสรานํ เทวานํ ฯเปฯ ปริตฺต\-สุภานํ เทวานํ ฯเปฯ อปฺปมาณ\-สุภานํ เทวานํ ฯเปฯ สุภกิณฺหานํ เทวานํ ฯเปฯ เวหปฺผลานํ เทวานํ ฯเปฯ \csromanfolio{360}อวิหานํ เทวานํ ฯเปฯ อตปฺปานํ เทวานํ ฯเปฯ สุทสฺสานํ เทวานํ ฯเปฯ สุทสฺสีนํ เทวานํ ฯเปฯ อกนิฏฺฐานํ เทวานํ ฯเปฯ อากาสานญฺจายตนู\-ปคานํ เทวานํ ฯเปฯ วิญฺญาณญฺจายตนู\-ปคานํ เทวานํ ฯเปฯ อากิญฺจญฺญายตนู\-ปคานํ เทวานํ ฯเปฯ เนว\-สญฺญา\-นา\-สญฺญา\-ยตนู\-ปคานํ เทวานํ สหพฺยตํ อุปปชฺเชยฺยนฺ”ติ. ฐานํ โข ปเนตํ วิชฺชติ, ยํ โส กายสฺส เภทา ปรํ มรณา เนว\-สญฺญา\-นา\-สญฺญา\-ยตนู\-ปคานํ เทวานํ สหพฺยตํ อุปปชฺเชยฺย. ตํ กิสฺส เหตุ, ตถา หิ โส ธมฺมจารี สมจารี.}

\prose{อากงฺเขยฺย เจ คหปตโย ธมฺมจารี สมจารี “อโห วตาหํ อาสวานํ ขยา อนาสวํ เจโตวิมุตฺติํ ปญฺญาวิมุตฺติํ ทิฏฺเฐว ธมฺเม สยํ อภิญฺญา สจฺฉิกตฺวา อุปสมฺปชฺช วิหเรยฺยนฺ”ติ. ฐานํ โข ปเนตํ วิชฺชติ, ยํ โส อาสวานํ ขยา อนาสวํ เจโตวิมุตฺติํ ปญฺญาวิมุตฺติํ ทิฏฺเฐว ธมฺเม สยํ อภิญฺญา สจฺฉิกตฺวา อุปสมฺปชฺช วิหเรยฺย. ตํ กิสฺส เหตุ, ตถา หิ โส ธมฺมจารี สมจารีติ.}

\proseitemwithcloserruled{443}{เอวํ วุตฺเต สาเลยฺยกา พฺราหฺมณ\-คหปติกา ภควนฺตํ เอตทโวจุํ “อภิกฺกนฺตํ โภ โคตม, อภิกฺกนฺตํ โภ โคตม, เสยฺยถาปิ โภ โคตม นิกฺกุชฺชิตํ วา อุกฺกุชฺเชยฺย, ปฏิจฺฉนฺนํ วา วิวเรยฺย, มูฬฺหสฺส วา มคฺคํ อาจิกฺเขยฺย, อนฺธกาเร วา เตลปชฺโชตํ ธาเรยฺย ‘จกฺขุมนฺโต รูปานิ ทกฺขนฺตี’ติ. เอวเมวํ โภตา โคตเมน อเนก\-ปริยาเยน ธมฺโม ปกาสิโต. เอเต มยํ ภวนฺตํ โคตมํ สรณํ คจฺฉาม ธมฺมญฺจ ภิกฺขุสํฆญฺจ. อุปาสเก โน ภวํ โคตโม ธาเรตุ อชฺชตคฺเค ปาณุเปเต\csromansharedfootnote{ebd455903f695188}{ปาณุเปตํ (ก)} สรณํ คเต”ติ.}{สาเลยฺยก\-สุตฺตํ นิฏฺฐิตํ ปฐมํ.}
\csromanmatikaanchor{mtk.76}
\csromantocmarkref{subsection}{2. เวรญฺชกสุตฺต}{mtk.76}
\bookmark[dest={mtk.76},level=2]{2. เวรญฺชกสุตฺต}
\csromanheader{2}{2. \textbf{เวรญฺชกสุตฺต}}
\proseitem{444}{เอวํ เม สุตํ\csromandash{}เอกํ สมยํ ภควา สาวตฺถิยํ วิหรติ เชตวเน อนาถ\-ปิณฺฑิกสฺส อาราเม. เตน โข ปน สมเยน เวรญฺชกา พฺราหฺมณ\-คหปติกา สาวตฺถิยํ ปฏิวสนฺติ เกนจิเทว กรณีเยน. \csromanfolio{361}อสฺโสสุํ โข เวรญฺชกา พฺราหฺมณ\-คหปติกา “สมโณ ขลุ โภ โคตโม สกฺยปุตฺโต สกฺยกุลา ปพฺพชิโต สาวตฺถิยํ วิหรติ เชตวเน อนาถ\-ปิณฺฑิกสฺส อาราเม. ตํ โข ปน ภวนฺตํ โคตมํ เอวํ กลฺยาโณ กิตฺติสทฺโท อพฺภุคฺคโต, อิติปิ โส ภควา อรหํ สมฺมาสมฺพุทฺโธ วิชฺชา\-จรณ\-สมฺปนฺโน สุคโต โลกวิทู อนุตฺตโร ปุริส\-ทมฺม\-สารถิ สตฺถา เทวมนุสฺสานํ พุทฺโธ ภควา. โส อิมํ โลกํ สเทวกํ สมารกํ สพฺรหฺมกํ สสฺสมณ\-พฺราหฺมณิํ ปชํ สเทวมนุสฺสํ สยํ อภิญฺญา สจฺฉิกตฺวา ปเวเทติ. โส ธมฺมํ เทเสติ อาทิกลฺยาณํ มชฺเฌกลฺยาณํ ปริโยสาน\-กลฺยาณํ สาตฺถํ สพฺยญฺชนํ เกวล\-ปริปุณฺณํ ปริสุทฺธํ พฺรหฺมจริยํ ปกาเสติ. สาธุ โข ปน ตถารูปานํ อรหตํ ทสฺสนํ โหตี”ติ.}

\prose{อถ โข เวรญฺชกา พฺราหฺมณ\-คหปติกา เยน ภควา เตนุ\-ปสงฺกมิํสุ, อุปสงฺกมิตฺวา อปฺเปกจฺเจ ภควนฺตํ อภิวาเทตฺวา เอกมนฺตํ นิสีทิํสุ. อปฺเปกจฺเจ ภควตา สทฺธิํ สมฺโมทิํสุ, สมฺโมทนียํ กถํ สารณียํ วีติสาเรตฺวา เอกมนฺตํ นิสีทิํสุ, อปฺเปกจฺเจ เยน ภควา เตนญฺชลิํ ปณาเมตฺวา เอกมนฺตํ นิสีทิํสุ, อปฺเปกจฺเจ ภควโต สนฺติเก นามโคตฺตํ สาเวตฺวา เอกมนฺตํ นิสีทิํสุ, อปฺเปกจฺเจ ตุณฺหีภูตา เอกมนฺตํ นิสีทิํสุ. เอกมนฺตํ นิสินฺนา โข เวรญฺชกา พฺราหฺมณ\-คหปติกา ภควนฺตํ เอตทโวจุํ “โก นุ โข โภ โคตม เหตุ โก ปจฺจโย, เยน มิเธกจฺเจ สตฺตา กายสฺส เภทา ปรํ มรณา อปายํ ทุคฺคติํ วินิปาตํ นิรยํ อุปปชฺชนฺติ. โก ปน โภ โคตม เหตุ โก ปจฺจโย, เยน มิเธกจฺเจ สตฺตา กายสฺส เภทา ปรํ มรณา สุคติํ สคฺคํ โลกํ อุปปชฺชนฺตี”ติ.}

\prose{อธมฺมจริยา\-วิสมจริยา\-เหตุ โข คหปตโย เอวมิเธกจฺเจ สตฺตา กายสฺส เภทา ปรํ มรณา อปายํ ทุคฺคติํ วินิปาตํ นิรยํ อุปปชฺชนฺติ. ธมฺมจริยา\-สมจริยา\-เหตุ โข คหปตโย เอวมิเธกจฺเจ สตฺตา กายสฺส เภทา ปรํ มรณา สุคติํ สคฺคํ โลกํ อุปปชฺชนฺตีติ.}

\prose{น โข มยํ อิมสฺส โภโต โคตมสฺส สํขิตฺเตน ภาสิตสฺส วิตฺถาเรน อตฺถํ อวิภตฺตสฺส วิตฺถาเรน อตฺถํ อาชานาม. สาธุ โน ภวํ โคตโม ตถา ธมฺมํ เทเสตุ, ยถา มยํ อิมสฺส โภโต โคตมสฺส สํขิตฺเตน ภาสิตสฺส วิตฺถาเรน อตฺถํ อวิภตฺตสฺส วิตฺถาเรน \csromanfolio{362}อตฺถํ อาชาเนยฺยามาติ. เตน หิ คหปตโย สุณาถ สาธุกํ มนสิ กโรถ ภาสิสฺสามีติ. “เอวํ โภ”ติ โข เวรญฺชกา พฺราหฺมณ\-คหปติกา ภควโต ปจฺจสฺโสสุํ. ภควา เอตทโวจ\csromandash{}}

\proseitem{445}{ติวิธํ โข คหปตโย กาเยน อธมฺมจารี วิสมจารี โหติ, จตุพฺพิธํ วาจาย อธมฺมจารี วิสมจารี โหติ, ติวิธํ มนสา อธมฺมจารี วิสมจารี โหติ.}

\prose{กถญฺจ คหปตโย ติวิธํ กาเยน อธมฺมจารี วิสมจารี โหติ. อิธ คหปตโย เอกจฺโจ ปาณาติปาตี โหติ ลุทฺโท โลหิตปาณิ หตปฺปหเต นิวิฏฺโฐ อทยาปนฺโน ปาณภูเตสุ. อทินฺนาทายี โข ปน โหติ, ยํ ตํ ปรสฺส ปรวิตฺตู\-ปกรณํ ฯเปฯ ตํ อทินฺนํ เถยฺยสงฺขาตํ อาทาตา โหติ. กาเมสุ\-มิจฺฉาจารี โข ปน โหติ, ยา ตา มาตุรกฺขิตา ฯเปฯ ตถารูปาสุ จาริตฺตํ อาปชฺชิตา โหติ. เอวํ โข คหปตโย ติวิธํ กาเยน อธมฺมจารี วิสมจารี โหติ.}

\prose{กถญฺจ คหปตโย จตุพฺพิธํ วาจาย อธมฺมจารี วิสมจารี โหติ. อิธ คหปตโย เอกจฺโจ มุสาวาที โหติ, สภาคโต วา ฯเปฯ สมฺปชานมุสา ภาสิตา โหติ. ปิสุณวาโจ โข ปน โหติ, อิโต สุตฺวา อมุตฺร อกฺขาตา ฯเปฯ วคฺคกรณิํ วาจํ ภาสิตา โหติ. ผรุสวาโจ โข ปน โหติ, ยา สา วาจา อณฺฑกา กกฺกสา ฯเปฯ ตถารูปิํ วาจํ ภาสิตา โหติ. สมฺผปฺปลาปี โข ปน โหติ, อกาลวาที ฯเปฯ อปริยนฺตวติํ อนตฺถ\-สํหิตํ. เอวํ โข คหปตโย จตุพฺพิธํ วาจาย อธมฺมจารี วิสมจารี โหติ.}

\prose{กถญฺจ คหปตโย ติวิธํ มนสา อธมฺมจารี วิสมจารี โหติ. อิธ คหปตโย เอกจฺโจ อภิชฺฌาลุ โหติ ฯเปฯ ตํ มมสฺสา”ติ. พฺยาปนฺนจิตฺโต โข ปน โหติ ปทุฏฺฐ\-มนสงฺกปฺโป “อิเม สตฺตา หญฺญนฺตุ วา ฯเปฯ มา วา อเหสุนฺ”ติ. มิจฺฉาทิฏฺฐิโก โข ปน โหติ วิปรีต\-ทสฺสโน “นตฺถิ ทินฺนํ, นตฺถิ ยิฏฺฐํ ฯเปฯ สจฺฉิกตฺวา ปเวเทนฺตี”ติ. เอวํ โข คหปตโย ติวิธํ มนสา อธมฺมจารี วิสมจารี โหติ.}

\prose{เอวํ อธมฺมจริยา\-วิสมจริยา\-เหตุ โข คหปตโย เอวมิเธกจฺเจ สตฺตา กายสฺส เภทา ปรํ มรณา อปายํ ทุคฺคติํ วินิปาตํ นิรยํ อุปปชฺชนฺติ.}

\proseitem{446}{\csromanfolio{363}ติวิธํ โข คหปตโย กาเยน ธมฺมจารี สมจารี โหติ, จตุพฺพิธํ วาจาย ธมฺมจารี สมจารี โหติ, ติวิธํ มนสา ธมฺมจารี สมจารี โหติ.}

\prose{กถญฺจ คหปตโย ติวิธํ กาเยน ธมฺมจารี สมจารี โหติ. อิธ คหปตโย เอกจฺโจ ปาณาติปาตํ ปหาย ปาณาติปาตา ปฏิวิรโต โหติ นิหิตทณฺโฑ นิหิตสตฺโถ ลชฺชี ทยาปนฺโน สพฺพปาณภูต\-หิตา\-นุกมฺปี วิหรติ. อทินฺนาทานํ ปหาย อทินฺนาทานา ปฏิวิรโต โหติ, ยํ ตํ ปรสฺส ฯเปฯ ตํ นาทินฺนํ เถยฺยสงฺขาตํ อาทาตา โหติ. กาเมสุ\-มิจฺฉาจารํ ปหาย ฯเปฯ ตถารูปาสุ น จาริตฺตํ อาปชฺชิตา โหติ. เอวํ โข คหปตโย ติวิธํ กาเยน ธมฺมจารี สมจารี โหติ.}

\prose{กถญฺจ คหปตโย จตุพฺพิธํ วาจาย ธมฺมจารี สมจารี โหติ. อิธ คหปตโย เอกจฺโจ มุสาวาทํ ปหาย มุสาวาทา ปฏิวิรโต โหติ, สภาคโต วา ฯเปฯ น สมฺปชานมุสา ภาสิตา โหติ. ปิสุณํ วาจํ ปหาย ฯเปฯ สมคฺคกรณิํ วาจํ ภาสิตา โหติ. ผรุสํ วาจํ ปหาย ฯเปฯ ตถารูปํ วาจํ ภาสิตา โหติ. สมฺผปฺปลาปํ ปหาย ฯเปฯ กาเลน สาปเทสํ ปริยนฺตวติํ อตฺถสํหิตํ. เอวํ โข คหปตโย จตุพฺพิธํ วาจาย ธมฺมจารี สมจารี โหติ.}

\prose{กถญฺจ คหปตโย ติวิธํ มนสา ธมฺมจารี สมจารี โหติ. อิธ คหปตโย เอกจฺโจ อนภิชฺฌาลุ โหติ, ยํ ตํ ปรสฺส ปรวิตฺตู\-ปกรณํ, ตํ นาภิชฺฌาตา โหติ “อโห วต ยํ ปรสฺส, ตํ มมสฺสา”ติ. อพฺยาปนฺนจิตฺโต โข ปน โหติ อปฺปทุฏฺฐมนสงฺกปฺโป “อิเม สตฺตา อเวรา อพฺยาพชฺฌา อนีฆา สุขี อตฺตานํ ปริหรนฺตู”ติ. สมฺมาทิฏฺฐิโก โข ปน โหติ อวิปรีต\-ทสฺสโน “อตฺถิ ทินฺนํ, อตฺถิ ยิฏฺฐํ ฯเปฯ สยํ อภิญฺญา สจฺฉิกตฺวา ปเวเทนฺตี”ติ. เอวํ โข คหปตโย ติวิธํ มนสา ธมฺมจารี สมจารี โหติ.}

\prose{เอวํ ธมฺมจริยา\-สมจริยา\-เหตุ โข คหปตโย เอวมิเธกจฺเจ สตฺตา กายสฺส เภทา ปรํ มรณา สุคติํ สคฺคํ โลกํ อุปปชฺชนฺติ.}

\proseitem{447}{อากงฺเขยฺย เจ คหปตโย ธมฺมจารี สมาจารี “อโห วตาหํ กายสฺส เภทา ปรํ มรณา ขตฺติย\-มหาสาลานํ สหพฺยตํ \csromanfolio{364}อุปปชฺเชยฺยนฺ”ติ. ฐานํ โข ปเนตํ วิชฺชติ, ยํ โส กายสฺส เภทา ปรํ มรณา ขตฺติย\-มหาสาลานํ สหพฺยตํ อุปปชฺเชยฺย. ตํ กิสฺส เหตุ, ตถา หิ โส ธมฺมจารี สมจารี.}

\prose{อากงฺเขยฺย เจ คหปตโย ธมฺมจารี สมจารี “อโห วตาหํ กายสฺส เภทา ปรํ มรณา พฺราหฺมณ\-มหาสาลานํ ฯเปฯ คหปติ\-มหาสาลานํ สหพฺยตํ อุปปชฺเชยฺยนฺ”ติ. ฐานํ โข ปเนตํ วิชฺชติ, ยํ โส กายสฺส เภทา ปรํ มรณา คหปติ\-มหาสาลานํ สหพฺยตํ อุปปชฺเชยฺย. ตํ กิสฺส เหตุ, ตถา หิ โส ธมฺมจารี สมจารี.}

\prose{อากงฺเขยฺย เจ คหปตโย ธมฺมจารี สมจารี “อโห วตาหํ กายสฺส เภทา ปรํ มรณา จาตุมหาราชิกานํ เทวานํ สหพฺยตํ อุปปชฺเชยฺยนฺ”ติ. ฐานํ โข ปเนตํ วิชฺชติ, ยํ โส กายสฺส เภทา ปรํ มรณา จาตุมหาราชิกานํ เทวานํ สหพฺยตํ อุปปชฺเชยฺย. ตํ กิสฺส เหตุ, ตถา หิ โส ธมฺมจารี สมจารี.}

\prose{อากงฺเขยฺย เจ คหปตโย ธมฺมจารี สมจารี “อโห วตาหํ กายสฺส เภทา ปรํ มรณา ตาวติํสานํ เทวานํ ฯเปฯ ยามานํ เทวานํ ฯเปฯ ตุสิตานํ เทวานํ ฯเปฯ นิมฺมานรตีนํ เทวานํ ฯเปฯ ปรนิมฺมิต\-วส\-วตฺตีนํ เทวานํ ฯเปฯ พฺรหฺม\-กายิกานํ เทวานํ สหพฺยตํ อุปปชฺเชยฺยนฺ”ติ. ฐานํ โข ปเนตํ วิชฺชติ, ยํ โส กายสฺส เภทา ปรํ มรณา พฺรหฺม\-กายิกานํ เทวานํ สหพฺยตํ อุปปชฺเชยฺย. ตํ กิสฺส เหตุ, ตถา หิ โส ธมฺมจารี สมจารี.}

\prose{อากงฺเขยฺย เจ คหปตโย ธมฺมจารี สมจารี “อโห วตาหํ กายสฺส เภทา ปรํ มรณา อาภานํ เทวานํ สหพฺยตํ อุปปชฺเชยฺยนฺ”ติ. ฐานํ โข ปเนตํ วิชฺชติ, ยํ โส กายสฺส เภทา ปรํ มรณา อาภานํ เทวานํ สหพฺยตํ อุปปชฺเชยฺย. ตํ กิสฺส เหตุ, ตถา หิ โส ธมฺมจารี สมจารี.}

\prose{อากงฺเขยฺย เจ คหปตโย ธมฺมจารี สมจารี “อโห วตาหํ กายสฺส เภทา ปรํ มรณา ปริตฺตาภานํ เทวานํ ฯเปฯ อปฺปมาณาภานํ เทวานํ ฯเปฯ อาภสฺสรานํ เทวานํ ฯเปฯ ปริตฺต\-สุภานํ เทวานํ. ฯเปฯ อปฺปมาณ\-สุภานํ เทวานํ ฯเปฯ สุภกิณฺหานํ เทวานํ ฯเปฯ เวหปฺผลานํ เทวานํ ฯเปฯ อวิหานํ เทวานํ ฯเปฯ อตปฺปานํ เทวานํ ฯเปฯ สุทสฺสานํ เทวานํ ฯเปฯ สุทสฺสีนํ เทวานํ ฯเปฯ อกนิฏฺฐานํ เทวานํ ฯเปฯ อากาสานญฺจายตนู\-ปคานํ เทวานํ ฯเปฯ \csromanfolio{365}วิญฺญาณญฺจายตนู\-ปคานํ เทวานํ ฯเปฯ อากิญฺจญฺญายตนู\-ปคานํ เทวานํ ฯเปฯ เนว\-สญฺญา\-นา\-สญฺญา\-ยตนู\-ปคานํ เทวานํ สหพฺยตํ อุปปชฺเชยฺยนฺ”ติ. ฐานํ โข ปเนตํ วิชฺชติ, ยํ โส กายสฺส เภทา ปรํ มรณา เนว\-สญฺญา\-นา\-สญฺญา\-ยตนู\-ปคานํ เทวานํ สหพฺยตํ อุปปชฺเชยฺย. ตํ กิสฺส เหตุ, ตถา หิ โส ธมฺมจารี สมจารี.}

\prose{อากงฺเขยฺย เจ คหปตโย ธมฺมจารี สมจารี “อโห วตาหํ อาสวานํ ขยา อนาสวํ เจโตวิมุตฺติํ ปญฺญาวิมุตฺติํ ทิฏฺเฐว ธมฺเม สยํ อภิญฺญา สจฺฉิกตฺวา อุปสมฺปชฺช วิหเรยฺยนฺ”ติ. ฐานํ โข ปเนตํ วิชฺชติ, ยํ โส อาสวานํ ขยา อนาสวํ เจโตวิมุตฺติํ ปญฺญาวิมุตฺติํ ทิฏฺเฐว ธมฺเม สยํ อภิญฺญา สจฺฉิกตฺวา อุปสมฺปชฺช วิหเรยฺย. ตํ กิสฺส เหตุ, ตถา หิ โส ธมฺมจารี สมจารีติ.}

\proseitemwithcloserruled{448}{เอวํ วุตฺเต เวรญฺชกา พฺราหฺมณ\-คหปติกา ภควนฺตํ เอตทโวจุํ “อภิกฺกนฺตํ โภ โคตม, อภิกฺกนฺตํ โภ โคตม. เสยฺยถาปิ โภ โคตม นิกฺกุชฺชิตํ วา อุกฺกุชฺเชยฺย, ปฏิจฺฉนฺนํ วา วิวเรยฺย, มูฬฺหสฺส วา มคฺคํ อาจิกฺเขยฺย, อนฺธกาเร วา เตลปชฺโชตํ ธาเรยฺย ‘จกฺขุมนฺโต รูปานิ ทกฺขนฺตี’ติ. เอวเมวํ โภตา โคตเมน อเนก\-ปริยาเยน ธมฺโม ปกาสิโต. เอเต มยํ ภวนฺตํ โคตมํ สรณํ คจฺฉาม ธมฺมญฺจ ภิกฺขุสํฆญฺจ, อุปาสเก โน ภวํ โคตโม ธาเรตุ อชฺชตคฺเค ปาณุเปเต สรณํ คเต”ติ.}{เวรญฺชก\-สุตฺตํ นิฏฺฐิตํ ทุติยํ.}
\csromanmatikaanchor{mtk.77}
\csromantocmarkref{subsection}{3. มหาเวทลฺลสุตฺต}{mtk.77}
\bookmark[dest={mtk.77},level=2]{3. มหาเวทลฺลสุตฺต}
\csromanheader{2}{3. มหาเวทลฺล\-สุตฺต}
\proseitem{449}{เอวํ เม สุตํ\csromandash{}เอกํ สมยํ ภควา สาวตฺถิยํ วิหรติ เชตวเน อนาถ\-ปิณฺฑิกสฺส อาราเม. อถ โข อายสฺมา มหาโกฏฺฐิโก สายนฺหสมยํ ปฏิสลฺลานา วุฏฺฐิโต เยนายสฺมา สาริปุตฺโต เตนุปสงฺกมิ, อุปสงฺกมิตฺวา อายสฺมตา สาริปุตฺเตน สทฺธิํ สมฺโมทิ, สมฺโมทนียํ กถํ สารณียํ วีติสาเรตฺวา เอกมนฺตํ นิสีทิ, เอกมนฺตํ นิสินฺโน โข อายสฺมา มหาโกฏฺฐิโก อายสฺมนฺตํ สาริปุตฺตํ \csromanfolio{366}เอตทโวจ “ทุปฺปญฺโญ ทุปฺปญฺโญติ อาวุโส วุจฺจติ, กิตฺตาวตา นุ โข อาวุโส ‘ทุปฺปญฺโญ’ติ วุจฺจตี”ติ. “นปฺปชานาติ นปฺปชานาตี”ติ โข อาวุโส ตสฺมา “ทุปฺปญฺโญ”ติ วุจฺจติ. กิญฺจ นปฺปชานาติ. อิทํ ทุกฺขนฺติ นปฺปชานาติ, อยํ ทุกฺข\-สมุทโยติ นปฺปชานาติ, อยํ ทุกฺข\-นิโรโธติ นปฺปชานาติ, อยํ ทุกฺขนิโรธ\-คามินี ปฏิปทาติ นปฺปชานาติ. “นปฺปชานาติ นปฺปชานาตี”ติ โข อาวุโส ตสฺมา “ทุปฺปญฺโญ”ติ วุจฺจตีติ.}

\prose{“สาธาวุโส”ติ โข อายสฺมา มหาโกฏฺฐิโก อายสฺมโต สาริปุตฺตสฺส ภาสิตํ อภินนฺทิตฺวา อนุโมทิตฺวา อายสฺมนฺตํ สาริปุตฺตํ อุตฺตริํ ปญฺหํ อปุจฺฉิ “ปญฺญวา ปญฺญวาติ อาวุโส วุจฺจติ, กิตฺตาวตา นุ โข อาวุโส ‘ปญฺญวา’ติ วุจฺจตี”ติ. “ปชานาติ ปชานาตี”ติ โข อาวุโส ตสฺมา “ปญฺญวา”ติ วุจฺจติ. กิญฺจ ปชานาติ. อิทํ ทุกฺขนฺติ ปชานาติ, อยํ ทุกฺข\-สมุทโยติ ปชานาติ, อยํ ทุกฺข\-นิโรโธติ ปชานาติ, อยํ ทุกฺขนิโรธ\-คามินี ปฏิปทาติ ปชานาติ. “ปชานาติ ปชานาตี”ติ โข อาวุโส ตสฺมา “ปญฺญวา”ติ วุจฺจตีติ.}

\prose{“วิญฺญาณํ วิญฺญาณนฺ”ติ อาวุโส วุจฺจติ, กิตฺตาวตา นุ โข อาวุโส “วิญฺญาณนฺ”ติ วุจฺจตีติ. “วิชานาติ วิชานาตี”ติ โข อาวุโส ตสฺมา “วิญฺญาณนฺ”ติ วุจฺจติ. กิญฺจ วิชานาติ. สุขนฺติปิ วิชานาติ, ทุกฺขนฺติปิ วิชานาติ, อทุกฺขมสุขนฺติปิ วิชานาติ. “วิชานาติ วิชานาติ”ติ โข อาวุโส ตสฺมา “วิญฺญาณนฺ”ติ วุจฺจตีติ.}

\prose{ยา จาวุโส ปญฺญา ยญฺจ วิญฺญาณํ, อิเม ธมฺมา สํสฏฺฐา, อุทาหุ วิสํสฏฺฐา, ลพฺภา จ ปนิเมสํ ธมฺมานํ วินิพฺภุชิตฺวา วินิพฺภุชิตฺวา\csromansharedfootnote{a4d814a2cecaf7cf}{วินิพฺภุชฺชิตฺวา วินิพฺภุชฺชิตฺวา (ก)} นานากรณํ ปญฺญาเปตุนฺติ? ยา จาวุโส ปญฺญา ยญฺจ วิญฺญาณํ, อิเม ธมฺมา สํสฏฺฐา, โน วิสํสฏฺฐา, น จ ลพฺภา อิเมสํ ธมฺมานํ วินิพฺภุชิตฺวา วินิพฺภุชิตฺวา นานากรณํ ปญฺญาเปตุํ. ยํ หาวุโส\csromansharedfootnote{b1a23f6c0ca9f687}{ยญฺจาวุโส (สฺยา, กํ, ก)} ปชานาติ, ตํ วิชานาติ. ยํ วิชานาติ, ตํ ปชานาติ. ตสฺมา อิเม ธมฺมา สํสฏฺฐา, โน วิสํสฏฺฐา, น จ ลพฺภา อิเมสํ ธมฺมานํ วินิพฺภุชิตฺวา วินิพฺภุชิตฺวา นานากรณํ ปญฺญาเปตุนฺติ.}

\prose{ยา จาวุโส ปญฺญา ยญฺจ วิญฺญาณํ, อิเมสํ ธมฺมานํ สํสฏฺฐานํ โน วิสํสฏฺฐานํ กิํ นานากรณนฺติ. ยา จาวุโส ปญฺญา ยญฺจ วิญฺญาณํ, \csromanfolio{367}อิเมสํ ธมฺมานํ สํสฏฺฐานํ โน วิสํสฏฺฐานํ ปญฺญา ภาเวตพฺพา, วิญฺญาณํ ปริญฺเญยฺยํ. อิทํ เนสํ นานากรณนฺติ.}

\proseitem{450}{“เวทนา เวทนา”ติ อาวุโส วุจฺจติ, กิตฺตาวตา นุ โข อาวุโส “เวทนา”ติ วุจฺจตีติ. “เวเทติ เวเทตี”ติ โข อาวุโส ตสฺมา “เวทนา”ติ วุจฺจติ. กิญฺจ เวเทติ. สุขมฺปิ เวเทติ, ทุกฺขมฺปิ เวเทติ, อทุกฺขม\-สุขมฺปิ เวเทติ. “เวเทติ เวเทตี”ติ โข อาวุโส ตสฺมา “เวทนา”ติ วุจฺจตีติ.}

\prose{“สญฺญา สญฺญา”ติ อาวุโส วุจฺจติ, กิตฺตาวตา นุ โข อาวุโส “สญฺญา”ติ วุจฺจตีติ. “สญฺชานาติ สญฺชานาตี”ติ โข อาวุโส ตสฺมา “สญฺญา”ติ วุจฺจติ. กิญฺจ สญฺชานาติ. นีลกมฺปิ สญฺชานาติ, ปีตกมฺปิ สญฺชานาติ, โลหิตกมฺปิ สญฺชานาติ, โอทาตมฺปิ สญฺชานาติ. “สญฺชานาติ สญฺชานาตี”ติ โข อาวุโส ตสฺมา “สญฺญา”ติ วุจฺจตีติ.}

\prose{ยา จาวุโส เวทนา ยา จ สญฺญา ยญฺจ วิญฺญาณํ, อิเม ธมฺมา สํสฏฺฐา, อุทาหุ วิสํสฏฺฐา, ลพฺภา จ ปนิเมสํ ธมฺมานํ วินิพฺภุชิตฺวา วินิพฺภุชิตฺวา นานากรณํ ปญฺญาเปตุนฺติ? ยา จาวุโส เวทนา ยา จ สญฺญา ยญฺจ วิญฺญาณํ, อิเม ธมฺมา สํสฏฺฐา, โน วิสํสฏฺฐา, น จ ลพฺภา อิเมสํ ธมฺมานํ วินิพฺภุชิตฺวา วินิพฺภุชิตฺวา นานากรณํ ปญฺญาเปตุํ. ยํ หาวุโส\csromansharedfootnote{b1a23f6c0ca9f687}{ยญฺจาวุโส (สฺยา, กํ, ก)} เวเทติ, ตํ สญฺชานาติ, ยํ สญฺชานาติ, ตํ วิชานาติ. ตสฺมา อิเม ธมฺมา สํสฏฺฐา, โน วิสํสฏฺฐา, น จ ลพฺภา อิเมสํ ธมฺมานํ วินิพฺภุชิตฺวา วินิพฺภุชิตฺวา นานากรณํ ปญฺญาเปตุนฺติ.}

\proseitem{451}{นิสฺสฏฺเฐน หาวุโส\csromansharedfootnote{a4bdb364054df3e2}{นิสฺสฏฺเฐน ปนาวุโส (?)} ปญฺจหิ อินฺทฺริเยหิ ปริสุทฺเธน มโนวิญฺญาเณน กิํ เนยฺยนฺติ. นิสฺสฏฺเฐน อาวุโส ปญฺจหิ อินฺทฺริเยหิ ปริสุทฺเธน มโนวิญฺญาเณน “อนนฺโต อากาโส”ติ อากาสา\-นญฺจา\-ยตนํ เนยฺยํ. “อนนฺตํ วิญฺญาณนฺ”ติ วิญฺญาณญฺจา\-ยตนํ เนยฺยํ. “นตฺถิ กิญฺจี”ติ อากิญฺจญฺญา\-ยตนํ เนยฺยนฺติ. เนยฺยํ ปนาวุโส ธมฺมํ เกน ปชานาตีติ. เนยฺยํ โข อาวุโส ธมฺมํ ปญฺญาจกฺขุนา ปชานาตีติ. ปญฺญา ปนาวุโส กิมตฺถิยาติ. ปญฺญา โข อาวุโส อภิญฺญตฺถา ปริญฺญตฺถา ปหานตฺถาติ.}

\proseitem{452}{\csromanfolio{368}กติ ปนาวุโส ปจฺจยา สมฺมาทิฏฺฐิยา อุปฺปาทายาติ. เทฺว โข อาวุโส ปจฺจยา สมฺมาทิฏฺฐิยา อุปฺปาทาย, ปรโต จ โฆโส โยนิโส จ มนสิกาโร, อิเม โข อาวุโส เทฺว ปจฺจยา สมฺมาทิฏฺฐิยา อุปฺปาทายาติ.}

\prose{กติหิ ปนาวุโส องฺเคหิ อนุคฺคหิตา สมฺมาทิฏฺฐิ เจโตวิมุตฺติ\-ผลา จ โหติ เจโตวิมุตฺติ\-ผลา\-นิสํสา จ, ปญฺญา\-วิมุตฺติ\-ผลา จ โหติ ปญฺญา\-วิมุตฺติ\-ผลา\-นิสํสา จาติ. ปญฺจหิ โข อาวุโส องฺเคหิ อนุคฺคหิตา สมฺมาทิฏฺฐิ เจโตวิมุตฺติ\-ผลา จ โหติ เจโตวิมุตฺติ\-ผลา\-นิสํสา จ, ปญฺญา\-วิมุตฺติ\-ผลา จ โหติ ปญฺญา\-วิมุตฺติ\-ผลา\-นิสํสา จ. อิธาวุโส สมฺมาทิฏฺฐิ สีลานุคฺคหิตา จ โหติ, สุตานุคฺคหิตา จ โหติ, สากจฺฉา\-นุคฺคหิตา จ โหติ, สมถา\-นุคฺคหิตา จ โหติ, วิปสฺสนา\-นุคฺคหิตา จ โหติ. อิเมหิ โข อาวุโส ปญฺจหงฺเคหิ อนุคฺคหิตา สมฺมาทิฏฺฐิ เจโตวิมุตฺติ\-ผลา จ โหติ เจโตวิมุตฺติ\-ผลา\-นิสํสา จ, ปญฺญา\-วิมุตฺติ\-ผลา จ โหติ ปญฺญา\-วิมุตฺติ\-ผลา\-นิสํสา จาติ.}

\proseitem{453}{กติ ปนาวุโส ภวาติ. ตโยเม อาวุโส ภวา กามภโว รูปภโว อรูปภโวติ. กถํ ปนาวุโส อายติํ ปุนพฺภวา\-ภินิพฺพตฺติ โหตีติ. อวิชฺชา\-นีวรณานํ โข อาวุโส สตฺตานํ ตณฺหา\-สํโยชนานํ ตตฺร\-ตตฺรา\-ภินนฺทนา เอวํ อายติํ ปุนพฺภวา\-ภินิพฺพตฺติ โหตีติ. กถํ ปนาวุโส อายติํ ปุนพฺภวา\-ภินิพฺพตฺติ น โหตีติ. อวิชฺชาวิราคา โข อาวุโส วิชฺชุปฺปาทา ตณฺหานิโรธา เอวํ อายติํ ปุนพฺภวา\-ภินิพฺพตฺติ น โหตีติ.}

\proseitem{454}{กตมํ ปนาวุโส ปฐมํ ฌานนฺติ. อิธาวุโส ภิกฺขุ วิวิจฺเจว กาเมหิ วิวิจฺจ อกุสเลหิ ธมฺเมหิ สวิตกฺกํ สวิจารํ วิเวกชํ ปีติสุขํ ปฐมํ ฌานํ อุปสมฺปชฺช วิหรติ. อิทํ วุจฺจติ อาวุโส ปฐมํ ฌานนฺติ. ปฐมํ ปนาวุโส ฌานํ กติองฺคิกนฺติ. ปฐมํ โข อาวุโส ฌานํ ปญฺจงฺคิกํ. อิธาวุโส ปฐมํ ฌานํ สมาปนฺนสฺส ภิกฺขุโน วิตกฺโก จ วตฺตติ วิจาโร จ ปีติ จ สุขญฺจ จิตฺเตกคฺคตา จ. ปฐมํ โข อาวุโส ฌานํ เอวํ ปญฺจงฺคิกนฺติ. ปฐมํ ปนาวุโส ฌานํ กตงฺค\-วิปฺปหีนํ, กตงฺค\-สมนฺนาคตนฺติ. ปฐมํ โข อาวุโส ฌานํ ปญฺจ\-งฺค\-วิปฺปหีนํ, ปญฺจ\-งฺค\-สมนฺนาคตํ. อิธาวุโส ปฐมํ ฌานํ สมาปนฺนสฺส ภิกฺขุโน กามจฺฉนฺโท \csromanfolio{369}ปหีโน โหติ, พฺยาปาโท ปหีโน โหติ, ถินมิทฺธํ ปหีนํ โหติ, อุทฺธจฺจ\-กุกฺกุจฺจํ ปหีนํ โหติ, วิจิกิจฺฉา ปหีนา โหติ. วิตกฺโก จ วตฺตติ วิจาโร จ ปีติ จ สุขญฺจ จิตฺเตกคฺคตา จ. ปฐมํ โข อาวุโส ฌานํ เอวํ ปญฺจ\-งฺค\-วิปฺปหีนํ ปญฺจ\-งฺค\-สมนฺนาคตนฺติ.}

\proseitem{455}{ปญฺจิมานิ อาวุโส อินฺทฺริยานิ นานาวิสยานิ นานาโคจรานิ น อญฺญมญฺญสฺส โคจรวิสยํ ปจฺจนุโภนฺติ. เสยฺยถิทํ, จกฺขุนฺทฺริยํ โสตินฺทฺริยํ ฆานินฺทฺริยํ ชิวฺหินฺทฺริยํ กายินฺทฺริยํ. อิเมสํ โข อาวุโส ปญฺจนฺนํ อินฺทฺริยานํ นานาวิสยานํ นานาโคจรานํ น อญฺญมญฺญสฺส โคจรวิสยํ ปจฺจนุโภนฺตานํ กิํ ปฏิสรณํ, โก จ เนสํ โคจรวิสยํ ปจฺจนุโภตีติ. ปญฺจิมานิ อาวุโส อินฺทฺริยานิ นานาวิสยานิ นานาโคจรานิ น อญฺญมญฺญสฺส โคจรวิสยํ ปจฺจนุโภนฺติ. เสยฺยถิทํ, จกฺขุนฺทฺริยํ โสตินฺทฺริยํ ฆานินฺทฺริยํ ชิวฺหินฺทฺริยํ กายินฺทฺริยํ. อิเมสํ โข อาวุโส ปญฺจนฺนํ อินฺทฺริยานํ นานาวิสยานํ นานาโคจรานํ น อญฺญมญฺญสฺส โคจรวิสยํ ปจฺจนุโภนฺตานํ มโน ปฏิสรณํ, มโน จ เนสํ โคจรวิสยํ ปจฺจนุโภตีติ.}

\proseitem{456}{ปญฺจิมานิ อาวุโส อินฺทฺริยานิ. เสยฺยถิทํ, จกฺขุนฺทฺริยํ โสตินฺทฺริยํ ฆานินฺทฺริยํ ชิวฺหินฺทฺริยํ กายินฺทฺริยํ. อิมานิ โข อาวุโส ปญฺจินฺทฺริยานิ กิํ ปฏิจฺจ ติฏฺฐนฺตีติ. ปญฺจิมานิ อาวุโส อินฺทฺริยานิ. เสยฺยถิทํ, จกฺขุนฺทฺริยํ โสตินฺทฺริยํ ฆานินฺทฺริยํ ชิวฺหินฺทฺริยํ กายินฺทฺริยํ. อิมานิ โข อาวุโส ปญฺจินฺทฺริยานิ อายุํ ปฏิจฺจ ติฏฺฐนฺตีติ.}

\prose{อายุ ปนาวุโส กิํ ปฏิจฺจ ติฏฺฐตีติ. อายุ อุสฺมํ ปฏิจฺจ ติฏฺฐตีติ. อุสฺมา ปนาวุโส กิํ ปฏิจฺจ ติฏฺฐตีติ. อุสฺมา อายุํ ปฏิจฺจ ติฏฺฐตีติ. อิทาเนว โข มยํ อาวุโส อายสฺมโต สาริปุตฺตสฺส ภาสิตํ เอวํ อาชานาม “อายุ อุสฺมํ ปฏิจฺจ ติฏฺฐตี”ติ. อิทาเนว ปน มยํ อาวุโส อายสฺมโต สาริปุตฺตสฺส ภาสิตํ เอวํ อาชานาม “อุสฺมา อายุํ ปฏิจฺจ ติฏฺฐตี”ติ. ยถา กถํ ปนาวุโส อิมสฺส ภาสิตสฺส อตฺโถ ทฏฺฐพฺโพติ. เตน หาวุโส อุปมํ เต กริสฺสามิ, อุปมาย\-ปิ\-เธ\-กจฺเจ วิญฺญู ปุริสา ภาสิตสฺส อตฺถํ อาชานนฺติ. เสยฺยถาปิ อาวุโส เตลปฺปทีปสฺส ฌายโต อจฺจิํ ปฏิจฺจ อาภา ปญฺญายติ, อาภํ ปฏิจฺจ อจฺจิ ปญฺญายติ. เอวเมว โข อาวุโส อายุ อุสฺมํ ปฏิจฺจ ติฏฺฐติ, อุสฺมา อายุํ ปฏิจฺจ ติฏฺฐตีติ.}

\proseitem{457}{\csromanfolio{370}เตว นุ โข อาวุโส อายุสงฺขารา, เต เวทนิยา ธมฺมา. อุทาหุ อญฺเญ อายุสงฺขารา, อญฺเญ เวทนิยา ธมฺมาติ. น โข อาวุโส เตว อายุสงฺขารา, เต เวทนิยา ธมฺมา. เต จ หาวุโส อายุสงฺขารา อภวิํสุ เต เวทนิยา ธมฺมา. น ยิทํ สญฺญาเวทยิต\-นิโรธํ สมาปนฺนสฺส ภิกฺขุโน วุฏฺฐานํ ปญฺญาเยถ. ยสฺมา จ โข อาวุโส อญฺเญ อายุสงฺขารา อญฺเญ เวทนิยา ธมฺมา, ตสฺมา สญฺญาเวทยิต\-นิโรธํ สมาปนฺนสฺส ภิกฺขุโน วุฏฺฐานํ ปญฺญายตีติ. ยทา นุ โข อาวุโส อิมํ กายํ กติ ธมฺมา ชหนฺติ, อถายํ กาโย อุชฺฌิโต อวกฺขิตฺโต เสติ ยถา กฏฺฐํ อเจตนนฺติ. ยทา โข อาวุโส อิมํ กายํ ตโย ธมฺมา ชหนฺติ อายุ อุสฺมา จ วิญฺญาณํ, อถายํ กาโย อุชฺฌิโต อวกฺขิตฺโต เสติ ยถา กฏฺฐํ อเจตนนฺติ.}

\prose{ยฺวายํ อาวุโส มโต กาลงฺกโต, โย จายํ ภิกฺขุ สญฺญาเวทยิต\-นิโรธํ สมาปนฺโน, อิเมสํ กิํ นานากรณนฺติ. ยฺวายํ อาวุโส มโต กาลงฺกโต, ตสฺส กายสงฺขารา นิรุทฺธา ปฏิปฺปสฺสทฺธา, วจีสงฺขารา นิรุทฺธา ปฏิปฺปสฺสทฺธา, จิตฺตสงฺขารา นิรุทฺธา ปฏิปฺปสฺสทฺธา, อายุ ปริกฺขีโณ, อุสฺมา วูปสนฺตา, อินฺทฺริยานิ ปริภินฺนานิ. โย จายํ ภิกฺขุ สญฺญาเวทยิต\-นิโรธํ สมาปนฺโน, ตสฺสปิ กายสงฺขารา นิรุทฺธา ปฏิปฺปสฺสทฺธา, วจีสงฺขารา นิรุทฺธา ปฏิปฺปสฺสทฺธา, จิตฺตสงฺขารา นิรุทฺธา ปฏิปฺปสฺสทฺธา, อายุ น ปริกฺขีโณ, อุสฺมา อวูปสนฺตา, อินฺทฺริยานิ วิปฺปสนฺนานิ. ยฺวายํ อาวุโส มโต กาลงฺกโต, โย จายํ ภิกฺขุ สญฺญาเวทยิต\-นิโรธํ สมาปนฺโน. อิทํ เนสํ นานากรณนฺติ.}

\proseitem{458}{กติ ปนาวุโส ปจฺจยา อทุกฺขม\-สุขาย เจโตวิมุตฺติยา สมาปตฺติยาติ. จตฺตาโร โข อาวุโส ปจฺจยา อทุกฺขม\-สุขาย เจโตวิมุตฺติยา สมาปตฺติยา. อิธาวุโส ภิกฺขุ สุขสฺส จ ปหานา ทุกฺขสฺส จ ปหานา ปุพฺเพว โสมนสฺส\-โทมนสฺสานํ อตฺถงฺคมา อทุกฺขมสุขํ อุเปกฺขา\-สติ\-ปาริสุทฺธิํ จตุตฺถํ ฌานํ อุปสมฺปชฺช วิหรติ. อิเม โข อาวุโส จตฺตาโร ปจฺจยา อทุกฺขม\-สุขาย เจโตวิมุตฺติยา สมาปตฺติยาติ.}

\prose{กติ ปนาวุโส ปจฺจยา อนิมิตฺตาย เจโตวิมุตฺติยา สมาปตฺติยาติ. เทฺว โข อาวุโส ปจฺจยา อนิมิตฺตาย เจโตวิมุตฺติยา \csromanfolio{371}สมาปตฺติยา. สพฺพนิมิตฺตานญฺจ อมนสิกาโร, อนิมิตฺตาย จ ธาตุยา มนสิกาโร. อิเม โข อาวุโส เทฺว ปจฺจยา อนิมิตฺตาย เจโตวิมุตฺติยา สมาปตฺติยาติ.}

\prose{กติ ปนาวุโส ปจฺจยา อนิมิตฺตาย เจโตวิมุตฺติยา ฐิติยาติ. ตโย โข อาวุโส ปจฺจยา อนิมิตฺตาย เจโตวิมุตฺติยา ฐิติยา. สพฺพนิมิตฺตานญฺจ อมนสิกาโร, อนิมิตฺตาย จ ธาตุยา มนสิกาโร, ปุพฺเพ จ อภิสงฺขาโร. อิเม โข อาวุโส ตโย ปจฺจยา อนิมิตฺตาย เจโตวิมุตฺติยา ฐิติยาติ.}

\prose{กติ ปนาวุโส ปจฺจยา อนิมิตฺตาย เจโตวิมุตฺติยา วุฏฺฐานายาติ. เทฺว โข อาวุโส ปจฺจยา อนิมิตฺตาย เจโตวิมุตฺติยา วุฏฺฐานาย. สพฺพนิมิตฺตานญฺจ มนสิกาโร, อนิมิตฺตาย จ ธาตุยา อมนสิกาโร. อิเม โข อาวุโส เทฺว ปจฺจยา อนิมิตฺตาย เจโตวิมุตฺติยา วุฏฺฐานายาติ.}

\proseitem{459}{ยา จายํ อาวุโส อปฺปมาณา เจโตวิมุตฺติ ยา จ อากิญฺจญฺญา เจโตวิมุตฺติ ยา จ สุญฺญตา เจโตวิมุตฺติ ยา จ อนิมิตฺตา เจโตวิมุตฺติ, อิเม ธมฺมา นานาตฺถา เจว นานาพฺยญฺชนา จ, อุทาหุ เอกตฺถา พฺยญฺชนเมว นานนฺติ? ยา จายํ อาวุโส อปฺปมาณา เจโตวิมุตฺติ ยา จ อากิญฺจญฺญา เจโตวิมุตฺติ ยา จ สุญฺญตา เจโตวิมุตฺติ ยา จ อนิมิตฺตา เจโตวิมุตฺติ, อตฺถิ โข อาวุโส ปริยาโย, ยํ ปริยายํ อาคมฺม อิเม ธมฺมา นานาตฺถา เจว นานาพฺยญฺชนา จ. อตฺถิ จ โข อาวุโส ปริยาโย, ยํ ปริยายํ อาคมฺม อิเม ธมฺมา เอกตฺถา พฺยญฺชนเมว นานํ.}

\prose{กตโม จาวุโส ปริยาโย, ยํ ปริยายํ อาคมฺม อิเม ธมฺมา นานาตฺถา เจว นานาพฺยญฺชนา จ. อิธาวุโส ภิกฺขุ เมตฺตา\-สหคเตน เจตสา เอกํ ทิสํ ผริตฺวา วิหรติ. ตถา ทุติยํ. ตถา ตติยํ. ตถา จตุตฺถํ. อิติ อุทฺธมโธ ติริยํ สพฺพธิ สพฺพตฺตตาย สพฺพาวนฺตํ โลกํ เมตฺตา\-สหคเตน เจตสา วิปุเลน มหคฺคเตน อปฺปมาเณน อเวเรน อพฺยาพชฺเฌน ผริตฺวา วิหรติ. กรุณา\-สหคเตน เจตสา ฯเปฯ. มุทิตา\-สหคเตน เจตสา ฯเปฯ. อุเปกฺขา\-สหคเตน เจตสา เอกํ ทิสํ ผริตฺวา วิหรติ. ตถา ทุติยํ. ตถา ตติยํ. ตถา จตุตฺถํ. อิติ \csromanfolio{372}อุทฺธมโธ ติริยํ สพฺพธิ สพฺพตฺตตาย สพฺพาวนฺตํ โลกํ อุเปกฺขา\-สหคเตน เจตสา วิปุเลน มหคฺคเตน อปฺปมาเณน อเวเรน อพฺยาพชฺเฌน ผริตฺวา วิหรติ. อยํ วุจฺจตาวุโส อปฺปมาณา เจโตวิมุตฺติ.}

\prose{กตมา จาวุโส อากิญฺจญฺญา เจโตวิมุตฺติ. อิธาวุโส ภิกฺขุ สพฺพโส วิญฺญาณญฺจา\-ยตนํ สมติกฺกมฺม “นตฺถิ กิญฺจี”ติ อากิญฺจญฺญา\-ยตนํ อุปสมฺปชฺช วิหรติ. อยํ วุจฺจตาวุโส อากิญฺจญฺญา เจโตวิมุตฺติ.}

\prose{กตมา จาวุโส สุญฺญตา เจโตวิมุตฺติ. อิธาวุโส ภิกฺขุ อรญฺญคโต วา รุกฺขมูลคโต วา สุญฺญาคารคโต วา อิติ ปฏิสญฺจิกฺขติ “สุญฺญมิทํ อตฺเตน วา อตฺตนิเยน วา”ติ. อยํ วุจฺจตาวุโส สุญฺญตา เจโตวิมุตฺติ.}

\prose{กตมา จาวุโส อนิมิตฺตา เจโตวิมุตฺติ. อิธาวุโส ภิกฺขุ สพฺพ\-นิมิตฺตานํ อมนสิการา อนิมิตฺตํ เจโตสมาธิํ อุปสมฺปชฺช วิหรติ. อยํ วุจฺจตาวุโส อนิมิตฺตา เจโตวิมุตฺติ. อยํ โข อาวุโส ปริยาโย, ยํ ปริยายํ อาคมฺม อิเม ธมฺมา นานาตฺถา เจว นานาพฺยญฺชนา จ.}

\prose{กตโม จาวุโส ปริยาโย, ยํ ปริยายํ อาคมฺม อิเม ธมฺมา เอกตฺถา พฺยญฺชนเมว นานํ. ราโค โข อาวุโส ปมาณกรโณ, โทโส ปมาณกรโณ, โมโห ปมาณกรโณ. เต ขีณาสวสฺส ภิกฺขุโน ปหีนา อุจฺฉินฺนมูลา ตาลาวตฺถุกตา อนภาวํกตา อายติํ อนุปฺปาทธมฺมา. ยาวตา โข อาวุโส อปฺปมาณา เจโตวิมุตฺติโย, อกุปฺปา ตาสํ เจโตวิมุตฺติ อคฺคมกฺขายติ. สา โข ปนากุปฺปา เจโตวิมุตฺติ สุญฺญา ราเคน, สุญฺญา โทเสน, สุญฺญา โมเหน. ราโค โข อาวุโส กิญฺจโน, โทโส กิญฺจโน, โมโห กิญฺจโน. เต ขีณาสวสฺส ภิกฺขุโน ปหีนา อุจฺฉินฺนมูลา ตาลาวตฺถุกตา อนภาวํกตา อายติํ อนุปฺปาทธมฺมา. ยาวตา โข อาวุโส อากิญฺจญฺญา เจโตวิมุตฺติโย, อกุปฺปา ตาสํ เจโตวิมุตฺติ อคฺคมกฺขายติ. สา โข ปนากุปฺปา เจโตวิมุตฺติ สุญฺญา ราเคน, สุญฺญา โทเสน, สุญฺญา โมเหน. ราโค โข อาวุโส นิมิตฺตกรโณ, โทโส นิมิตฺตกรโณ, โมโห นิมิตฺตกรโณ.}

\prose{\csromanfolio{373}เต ขีณาสวสฺส ภิกฺขุโน ปหีนา อุจฺฉินฺนมูลา ตาลาวตฺถุกตา อนภาวํกตา อายติํ อนุปฺปาทธมฺมา. ยาวตา โข อาวุโส อนิมิตฺตา เจโตวิมุตฺติโย, อกุปฺปา ตาสํ เจโตวิมุตฺติ อคฺคมกฺขายติ. สา โข ปนากุปฺปา เจโตวิมุตฺติ สุญฺญา ราเคน, สุญฺญา โทเสน, สุญฺญา โมเหน. อยํ โข อาวุโส ปริยาโย, ยํ ปริยายํ อาคมฺม อิเม ธมฺมา เอกตฺถา พฺยญฺชนเมว นานนฺติ.}

\prosewithcloserruled{อิทมโวจายสฺมา สาริปุตฺโต. อตฺตมโน อายสฺมา มหาโกฏฺฐิโก อายสฺมโต สาริปุตฺตสฺส ภาสิตํ อภินนฺทีติ.}{มหาเวทลฺล\-สุตฺตํ นิฏฺฐิตํ ตติยํ.}
\csromanmatikaanchor{mtk.78}
\csromantocmarkref{subsection}{4. จูฬเวทลฺลสุตฺต}{mtk.78}
\bookmark[dest={mtk.78},level=2]{4. จูฬเวทลฺลสุตฺต}
\csromanheader{2}{4. จูฬเวทลฺล\-สุตฺต}
\proseitem{460}{เอวํ เม สุตํ\csromandash{}เอกํ สมยํ ภควา ราชคเห วิหรติ เวฬุวเน กลนฺทก\-นิวาเป. อถ โข วิสาโข อุปาสโก เยน ธมฺมทินฺนา ภิกฺขุนี เตนุปสงฺกมิ, อุปสงฺกมิตฺวา ธมฺมทินฺนํ ภิกฺขุนิํ อภิวาเทตฺวา เอกมนฺตํ นิสีทิ, เอกมนฺตํ นิสินฺโน โข วิสาโข อุปาสโก ธมฺมทินฺนํ ภิกฺขุนิํ เอตทโวจ “สกฺกาโย สกฺกาโยติ อยฺเย วุจฺจติ, กตโม นุ โข อยฺเย สกฺกาโย วุตฺโต ภควตา”ติ. ปญฺจ โข อิเม อาวุโส วิสาข อุปาทาน\-กฺขนฺธา สกฺกาโย วุตฺโต ภควตา. เสยฺยถิทํ, รูปุปาทาน\-กฺขนฺโธ เวทนุ\-ปาทาน\-กฺขนฺโธ สญฺญุ\-ปาทาน\-กฺขนฺโธ สงฺขารุ\-ปาทาน\-กฺขนฺโธ วิญฺญาณุ\-ปาทาน\-กฺขนฺโธ. อิเม โข อาวุโส วิสาข ปญฺจุ\-ปาทาน\-กฺขนฺธา สกฺกาโย วุตฺโต ภควตาติ.}

\prose{สาธยฺเยติ โข วิสาโข อุปาสโก ธมฺมทินฺนาย ภิกฺขุนิยา ภาสิตํ อภินนฺทิตฺวา อนุโมทิตฺวา ธมฺมทินฺนํ ภิกฺขุนิํ อุตฺตริํ ปญฺหํ อปุจฺฉิ “สกฺกาย\-สมุทโย สกฺกาย\-สมุทโยติ อยฺเย วุจฺจติ, กตโม นุ โข อยฺเย สกฺกาย\-สมุทโย วุตฺโต ภควตา”ติ. ยายํ อาวุโส วิสาข ตณฺหา โปโนพฺภวิกา นนฺที\-ราค\-สหคตา ตตฺร\-ตตฺรา\-ภินนฺทินี. เสยฺยถิทํ, กามตณฺหา ภวตณฺหา วิภวตณฺหา. อยํ โข อาวุโส วิสาข สกฺกาย\-สมุทโย วุตฺโต ภควตาติ.}

\prose{\csromanfolio{374}“สกฺกายนิโรโธ สกฺกาย\-นิโรโธติ อยฺเย วุจฺจติ, กตโม นุ โข อยฺเย สกฺกายนิโรโธ วุตฺโต ภควตา”ติ. โย โข อาวุโส วิสาข ตสฺสาเยว ตณฺหาย อเสส\-วิราค\-นิโรโธ จาโค ปฏินิสฺสคฺโค มุตฺติ อนาลโย. อยํ โข อาวุโส วิสาข สกฺกายนิโรโธ วุตฺโต ภควตาติ.}

\prose{“สกฺกาย\-นิโรธ\-คามินี ปฏิปทา สกฺกาย\-นิโรธ\-คามินี ปฏิปทาติ อยฺเย วุจฺจติ, กตมา นุ โข อยฺเย สกฺกาย\-นิโรธ\-คามินี ปฏิปทา วุตฺตา ภควตา”ติ. อยเมว โข อาวุโส วิสาข อริโย อฏฺฐงฺคิโก มคฺโค สกฺกาย\-นิโรธ\-คามินี ปฏิปทา วุตฺตา ภควตา. เสยฺยถิทํ, สมฺมาทิฏฺฐิ สมฺมาสงฺกปฺโป สมฺมาวาจา สมฺมากมฺมนฺโต สมฺมาอาชีโว สมฺมาวายาโม สมฺมาสติ สมฺมาสมาธีติ. “ตญฺเญว นุ โข อยฺเย อุปาทานํ เต\csromansharedfootnote{1e5cac16e6d26073}{เตว (สี)} ปญฺจุ\-ปาทาน\-กฺขนฺธา, อุทาหุ อญฺญตฺร ปญฺจหุ\-ปาทานกฺขนฺเธหิ อุปาทานนฺ”ติ. น โข อาวุโส วิสาข ตญฺเญว อุปาทานํ เต ปญฺจุ\-ปาทาน\-กฺขนฺธา, นาปิ อญฺญตฺร ปญฺจหุ\-ปาทานกฺขนฺเธหิ อุปาทานํ. โย โข อาวุโส วิสาข ปญฺจสุ อุปาทาน\-กฺขนฺเธสุ ฉนฺทราโค, ตํ ตตฺถ อุปาทานนฺติ.}

\proseitem{461}{“กถํ ปนายฺเย สกฺกายทิฏฺฐิ โหตี”ติ. อิธาวุโส วิสาข อสฺสุตวา ปุถุชฺชโน อริยานํ อทสฺสาวี อริยธมฺมสฺส อโกวิโท อริยธมฺเม อวินีโต สปฺปุริสานํ อทสฺสาวี สปฺปุริส\-ธมฺมสฺส อโกวิโท สปฺปุริส\-ธมฺเม อวินีโต รูปํ อตฺตโต สมนุปสฺสติ, รูปวนฺตํ วา อตฺตานํ, อตฺตนิ วา รูปํ, รูปสฺมิํ วา อตฺตานํ. เวทนํ ฯเปฯ. สญฺญํ ฯเปฯ. สงฺขาเร ฯเปฯ. วิญฺญาณํ อตฺตโต สมนุปสฺสติ, วิญฺญาณวนฺตํ วา อตฺตานํ, อตฺตนิ วา วิญฺญาณํ, วิญฺญาณสฺมิํ วา อตฺตานํ. เอวํ โข อาวุโส วิสาข สกฺกายทิฏฺฐิ โหตีติ.}

\prose{“กถํ ปนายฺเย สกฺกายทิฏฺฐิ น โหตี”ติ. อิธาวุโส วิสาข สุตวา อริยสาวโก อริยานํ ทสฺสาวี อริยธมฺมสฺส โกวิโท อริยธมฺเม สุวินีโต สปฺปุริสานํ ทสฺสาวี สปฺปุริส\-ธมฺมสฺส โกวิโท สปฺปุริส\-ธมฺเม สุวินีโต น รูปํ อตฺตโต สมนุปสฺสติ, น รูปวนฺตํ วา อตฺตานํ, น อตฺตนิ วา รูปํ, น รูปสฺมิํ วา อตฺตานํ. น เวทนํ ฯเปฯ. น สญฺญํ ฯเปฯ. น สงฺขาเร ฯเปฯ. น วิญฺญาณํ อตฺตโต สมนุปสฺสติ, น วิญฺญาณวนฺตํ วา \csromanfolio{375}อตฺตานํ, น อตฺตนิ วา วิญฺญาณํ, น วิญฺญาณสฺมิํ วา อตฺตานํ. เอวํ โข อาวุโส วิสาข สกฺกายทิฏฺฐิ น โหตีติ.}

\proseitem{462}{“กตโม ปนายฺเย อริโย อฏฺฐงฺคิโก มคฺโค”ติ. อยเมว โข อาวุโส วิสาข อริโย อฏฺฐงฺคิโก มคฺโค. เสยฺยถิทํ, สมฺมาทิฏฺฐิ สมฺมาสงฺกปฺโป สมฺมาวาจา สมฺมากมฺมนฺโต สมฺมาอาชีโว สมฺมาวายาโม สมฺมาสติ สมฺมาสมาธีติ. “อริโย ปนายฺเย อฏฺฐงฺคิโก มคฺโค สงฺขโต, อุทาหุ อสงฺขโต”ติ. อริโย โข อาวุโส วิสาข อฏฺฐงฺคิโก มคฺโค สงฺขโตติ.}

\prose{“อริเยน นุ โข อยฺเย อฏฺฐงฺคิเกน มคฺเคน ตโย ขนฺธา สงฺคหิตา, อุทาหุ ตีหิ ขนฺเธหิ อริโย อฏฺฐงฺคิโก มคฺโค สงฺคหิโต”ติ. น โข อาวุโส วิสาข อริเยน อฏฺฐงฺคิเกน มคฺเคน ตโย ขนฺธา สงฺคหิตา. ตีหิ จ โข อาวุโส วิสาข ขนฺเธหิ อริโย อฏฺฐงฺคิโก มคฺโค สงฺคหิโต. ยา จาวุโส วิสาข สมฺมาวาจา, โย จ สมฺมากมฺมนฺโต, โย จ สมฺมาอาชีโว. อิเม ธมฺมา สีลกฺขนฺเธ สงฺคหิตา. โย จ สมฺมาวายาโม, ยา จ สมฺมาสติ, โย จ สมฺมาสมาธิ. อิเม ธมฺมา สมาธิ\-กฺขนฺเธ สงฺคหิตา. ยา จ สมฺมาทิฏฺฐิ, โย จ สมฺมาสงฺกปฺโป. อิเม ธมฺมา ปญฺญากฺขนฺเธ สงฺคหิตาติ.}

\prose{“กตโม ปนายฺเย สมาธิ, กตเม ธมฺมา สมาธินิมิตฺตา, กตเม ธมฺมา สมาธิ\-ปริกฺขารา, กตมา สมาธิภาวนา”ติ. ยา โข อาวุโส วิสาข จิตฺตสฺส เอกคฺคตา, อยํ สมาธิ. จตฺตาโร สติปฏฺฐานา สมาธินิมิตฺตา. จตฺตาโร สมฺมปฺปธานา สมาธิ\-ปริกฺขารา. ยา เตสํเยว ธมฺมานํ อาเสวนา ภาวนา พหุลีกมฺมํ, อยํ เอตฺถ สมาธิภาวนาติ.}

\proseitem{463}{“กติ ปนายฺเย สงฺขารา”ติ. ตโยเม อาวุโส วิสาข สงฺขารา กายสงฺขาโร วจีสงฺขาโร จิตฺต\-สงฺขาโรติ. “กตโม ปนายฺเย กายสงฺขาโร, กตโม วจีสงฺขาโร, กตโม จิตฺตสงฺขาโร”ติ. อสฺสาสปสฺสาสา โข อาวุโส วิสาข กายสงฺขาโร, วิตกฺกวิจารา วจีสงฺขาโร, สญฺญา จ เวทนา จ จิตฺต\-สงฺขาโรติ. “กสฺมา ปนายฺเย อสฺสาสปสฺสาสา กายสงฺขาโร, กสฺมา วิตกฺกวิจารา วจีสงฺขาโร, กสฺมา สญฺญา จ เวทนา จ จิตฺตสงฺขาโร”ติ. อสฺสาสปสฺสาสา โข \csromanfolio{376}อาวุโส วิสาข กายิกา เอเต ธมฺมา กาย\-ปฺปฏิพทฺธา, ตสฺมา อสฺสาสปสฺสาสา กายสงฺขาโร. ปุพฺเพ โข อาวุโส วิสาข วิตกฺเกตฺวา วิจาเรตฺวา ปจฺฉา วาจํ ภินฺทติ, ตสฺมา วิตกฺกวิจารา วจีสงฺขาโร. สญฺญา จ เวทนา จ เจตสิกา เอเต ธมฺมา จิตฺต\-ปฺปฏิพทฺธา, ตสฺมา สญฺญา จ เวทนา จ จิตฺต\-สงฺขาโรติ.}

\proseitem{464}{“กถํ ปนายฺเย สญฺญา\-เวทยิต\-นิโรธ\-สมาปตฺติ โหตี”ติ. น โข อาวุโส วิสาข สญฺญาเวทยิต\-นิโรธํ สมา\-ป\-ชฺช\-นฺตสฺส ภิกฺขุโน เอวํ โหติ “อหํ สญฺญาเวทยิต\-นิโรธํ สมาปชฺชิสฺสนฺติ วา, อหํ สญฺญาเวทยิต\-นิโรธํ สมาปชฺชามีติ วา, อหํ สญฺญาเวทยิต\-นิโรธํ สมาปนฺโนติ วา”. อถ ขฺวาสฺส ปุพฺเพว ตถา จิตฺตํ ภาวิตํ โหติ, ยํ ตํ ตถตฺตาย อุปเนตีติ.}

\prose{“สญฺญาเวทยิต\-นิโรธํ สมา\-ป\-ชฺช\-นฺตสฺส ปนายฺเย ภิกฺขุโน กตเม ธมฺมา ปฐมํ นิรุชฺฌนฺติ, ยทิ วา กายสงฺขาโร ยทิ วา วจีสงฺขาโร ยทิ วา จิตฺตสงฺขาโร”ติ. สญฺญาเวทยิต\-นิโรธํ สมา\-ป\-ชฺช\-นฺตสฺส โข อาวุโส วิสาข ภิกฺขุโน ปฐมํ นิรุชฺฌติ วจีสงฺขาโร, ตโต กายสงฺขาโร, ตโต จิตฺต\-สงฺขาโรติ.}

\prose{“กถํ ปนายฺเย สญฺญา\-เวทยิต\-นิโรธ\-สมาปตฺติยา วุฏฺฐานํ โหตี”ติ. น โข อาวุโส วิสาข สญฺญา\-เวทยิต\-นิโรธ\-สมาปตฺติยา วุฏฺฐหนฺตสฺส ภิกฺขุโน เอวํ โหติ “อหํ สญฺญา\-เวทยิต\-นิโรธ\-สมาปตฺติยา วุฏฺฐหิสฺสนฺติ วา, อหํ สญฺญา\-เวทยิต\-นิโรธ\-สมาปตฺติยา วุฏฺฐหามีติ วา, อหํ สญฺญา\-เวทยิต\-นิโรธ\-สมาปตฺติยา วุฏฺฐิโตติ วา”. อถ ขฺวาสฺส ปุพฺเพว ตถา จิตฺตํ ภาวิตํ โหติ, ยํ ตํ ตถตฺตาย อุปเนตีติ.}

\prose{“สญฺญา\-เวทยิต\-นิโรธ\-สมาปตฺติยา วุฏฺฐหนฺตสฺส ปนายฺเย ภิกฺขุโน กตเม ธมฺมา ปฐมํ อุปฺปชฺชนฺติ, ยทิ วา กายสงฺขาโร ยทิ วา วจีสงฺขาโร ยทิ วา จิตฺตสงฺขาโร”ติ. สญฺญา\-เวทยิต\-นิโรธ\-สมาปตฺติยา วุฏฺฐหนฺตสฺส โข อาวุโส วิสาข ภิกฺขุโน ปฐมํ อุปฺปชฺชติ จิตฺตสงฺขาโร, ตโต กายสงฺขาโร, ตโต วจีสงฺขาโรติ.}

\prose{“สญฺญา\-เวทยิต\-นิโรธ\-สมาปตฺติยา วุฏฺฐิตํ ปนายฺเย ภิกฺขุํ กติ ผสฺสา ผุสนฺตี”ติ. สญฺญา\-เวทยิต\-นิโรธ\-สมาปตฺติยา วุฏฺฐิตํ โข \csromanfolio{377}อาวุโส วิสาข ภิกฺขุํ ตโย ผสฺสา ผุสนฺติ สุญฺญโต ผสฺโส อนิมิตฺโต ผสฺโส อปฺปณิหิโต ผสฺโสติ.}

\prose{“สญฺญา\-เวทยิต\-นิโรธ\-สมาปตฺติยา วุฏฺฐิตสฺส ปนายฺเย ภิกฺขุโน กิํ นินฺนํ จิตฺตํ โหติ กิํ โปณํ กิํปพฺภารนฺ”ติ, สญฺญา\-เวทยิต\-นิโรธ\-สมาปตฺติยา วุฏฺฐิตสฺส โข อาวุโส วิสาข ภิกฺขุโน วิเวกนินฺนํ จิตฺตํ โหติ วิเวกโปณํ วิเวก\-ปพฺภารนฺติ.}

\proseitem{465}{“กติ ปนายฺเย เวทนา”ติ. ติสฺโส โข อิมา อาวุโส วิสาข เวทนา สุขา เวทนา ทุกฺขา เวทนา อทุกฺขมสุขา เวทนาติ. “กตมา ปนายฺเย สุขา เวทนา, กตมา ทุกฺขา เวทนา, กตมา อทุกฺขมสุขา เวทนา”ติ. ยํ โข อาวุโส วิสาข กายิกํ วา เจตสิกํ วา สุขํ สาตํ เวทยิตํ, อยํ สุขา เวทนา. ยํ โข อาวุโส วิสาข กายิกํ วา เจตสิกํ วา ทุกฺขํ อสาตํ เวทยิตํ, อยํ ทุกฺขา เวทนา. ยํ โข อาวุโส วิสาข กายิกํ วา เจตสิกํ วา เนว สาตํ นาสาตํ เวทยิตํ, อยํ อทุกฺขมสุขา เวทนาติ.}

\prose{“สุขา ปนายฺเย เวทนา กิํสุขา กิํทุกฺขา, ทุกฺขา เวทนา กิํสุขา กิํทุกฺขา, อทุกฺขมสุขา เวทนา กิํสุขา กิํทุกฺขา”ติ. สุขา โข อาวุโส วิสาข เวทนา ฐิติสุขา วิปริณาม\-ทุกฺขา. ทุกฺขา เวทนา ฐิติทุกฺขา วิปริณาม\-สุขา. อทุกฺขมสุขา เวทนา ญาณสุขา อญฺญาณทุกฺขาติ.}

\prose{“สุขาย ปนายฺเย เวทนาย กิํ อนุสโย อนุเสติ, ทุกฺขาย เวทนาย กิํ อนุสโย อนุเสติ, อทุกฺขม\-สุขาย เวทนาย กิํ อนุสโย อนุเสตี”ติ. สุขาย โข อาวุโส วิสาข เวทนาย ราคานุสโย อนุเสติ, ทุกฺขาย เวทนาย ปฏิฆานุสโย อนุเสติ, อทุกฺขม\-สุขาย เวทนาย อวิชฺชานุสโย อนุเสตีติ.}

\prose{“สพฺพาย นุ โข อยฺเย สุขาย เวทนาย ราคานุสโย อนุเสติ, สพฺพาย ทุกฺขาย เวทนาย ปฏิฆานุสโย อนุเสติ, สพฺพาย อทุกฺขม\-สุขาย เวทนาย อวิชฺชานุสโย อนุเสตี”ติ. น โข อาวุโส วิสาข สพฺพาย สุขาย เวทนาย ราคานุสโย อนุเสติ, น สพฺพาย ทุกฺขาย เวทนาย ปฏิฆานุสโย อนุเสติ, น สพฺพาย อทุกฺขมสุกฺขาย เวทนาย อวิชฺชานุสโย อนุเสตีติ.}

\prose{\csromanfolio{378}“สุขาย ปนายฺเย เวทนาย กิํ ปหาตพฺพํ, ทุกฺขาย เวทนาย กิํ ปหาตพฺพํ, อทุกฺขม\-สุขาย เวทนาย กิํ ปหาตพฺพนฺ”ติ. สุขาย โข อาวุโส วิสาข เวทนาย ราคานุสโย ปหาตพฺโพ, ทุกฺขาย เวทนาย ปฏิฆานุสโย ปหาตพฺโพ, อทุกฺขม\-สุขาย เวทนาย อวิชฺชานุสโย ปหาตพฺโพติ.}

\prose{“สพฺพาย นุ โข อยฺเย สุขาย เวทนาย ราคานุสโย ปหาตพฺโพ, สพฺพาย ทุกฺขาย เวทนาย ปฏิฆานุสโย ปหาตพฺโพ, สพฺพาย อทุกฺขม\-สุขาย เวทนาย อวิชฺชานุสโย ปหาตพฺโพ”ติ. น โข อาวุโส วิสาข สพฺพาย สุขาย เวทนาย ราคานุสโย ปหาตพฺโพ, น สพฺพาย ทุกฺขาย เวทนาย ปฏิฆานุสโย ปหาตพฺโพ, น สพฺพาย อทุกฺขม\-สุขาย เวทนาย อวิชฺชานุสโย ปหาตพฺโพ. อิธาวุโส วิสาข ภิกฺขุ วิวิจฺเจว กาเมหิ วิวิจฺจ อกุสเลหิ ธมฺเมหิ สวิตกฺกํ สวิจารํ วิเวกชํ ปีติสุขํ ปฐมํ ฌานํ อุปสมฺปชฺช วิหรติ. ราคํ เตน ปชหติ, น ตตฺถ ราคานุสโย อนุเสติ. อิธาวุโส วิสาข ภิกฺขุ อิติ ปฏิสญฺจิกฺขติ “กุทาสฺสุ นามาหํ ตทายตนํ อุปสมฺปชฺช วิหริสฺสามิ, ยทริยา เอตรหิ อายตนํ อุปสมฺปชฺช วิหรนฺตี”ติ. อิติ อนุตฺตเรสุ วิโมกฺเขสุ ปิหํ อุปฏฺฐาปยโต อุปฺปชฺชติ ปิหาปฺปจฺจยา โทมนสฺสํ, ปฏิฆํ เตน ปชหติ, น ตตฺถ ปฏิฆานุสโย อนุเสติ. อิธาวุโส วิสาข ภิกฺขุ สุขสฺส จ ปหานา ทุกฺขสฺส จ ปหานา ปุพฺเพว โสมนสฺส\-โทมนสฺสานํ อตฺถงฺคมา อทุกฺขมสุขํ อุเปกฺขา\-สติ\-ปาริสุทฺธิํ จตุตฺถํ ฌานํ อุปสมฺปชฺช วิหรติ. อวิชฺชํ เตน ปชหติ, น ตตฺถ อวิชฺชานุสโย อนุเสตีติ.}

\proseitem{466}{“สุขาย ปนายฺเย เวทนาย กิํ ปฏิภาโค”ติ. สุขาย โข อาวุโส วิสาข เวทนาย ทุกฺขา เวทนา ปฏิภาโคติ. “ทุกฺขาย ปนายฺเย เวทนาย กิํ ปฏิภาโค”ติ. ทุกฺขาย โข อาวุโส วิสาข เวทนาย สุขา เวทนา ปฏิภาโคติ. “อทุกฺขม\-สุขาย ปนายฺเย เวทนาย กิํ ปฏิภาโค”ติ. อทุกฺขม\-สุขาย โข อาวุโส วิสาข เวทนาย อวิชฺชา ปฏิภาโคติ. “อวิชฺชาย ปนายฺเย กิํ ปฏิภาโค”ติ. อวิชฺชาย โข อาวุโส วิสาข วิชฺชา ปฏิภาโคติ. “วิชฺชาย ปนายฺเย กิํ ปฏิภาโค”ติ. วิชฺชาย โข อาวุโส วิสาข วิมุตฺติ ปฏิภาโคติ. “วิมุตฺติยา ปนายฺเย กิํ ปฏิภาโค”ติ. วิมุตฺติยา โข อาวุโส วิสาข \csromanfolio{379}นิพฺพานํ ปฏิภาโคติ. “นิพฺพานสฺส ปนายฺเย กิํ ปฏิภาโค”ติ. อจฺจยาสิ อาวุโส\csromansharedfootnote{57d32c358bafca64}{อจฺจสราวุโส (สี, อิ), อจฺจสฺสราวุโส (สฺยา, กํ)} วิสาข ปญฺหํ, นาสกฺขิ ปญฺหานํ ปริยนฺตํ คเหตุํ. นิพฺพาโนคธํ หิ อาวุโส วิสาข พฺรหฺมจริยํ นิพฺพาน\-ปรายนํ นิพฺพาน\-ปริโยสานํ. อากงฺขมาโน จ ตฺวํ อาวุโส วิสาข ภควนฺตํ อุปสงฺกมิตฺวา เอตมตฺถํ ปุจฺเฉยฺยาสิ. ยถา จ เต ภควา พฺยากโรติ, ตถา นํ ธาเรยฺยาสีติ.}

\proseitem{467}{อถ โข วิสาโข อุปาสโก ธมฺมทินฺนาย ภิกฺขุนิยา ภาสิตํ อภินนฺทิตฺวา อนุโมทิตฺวา อุฏฺฐายาสนา ธมฺมทินฺนํ ภิกฺขุนิํ อภิวาเทตฺวา ปทกฺขิณํ กตฺวา เยน ภควา เตนุปสงฺกมิ, อุปสงฺกมิตฺวา ภควนฺตํ อภิวาเทตฺวา เอกมนฺตํ นิสีทิ, เอกมนฺตํ นิสินฺโน โข วิสาโข อุปาสโก ยาวตโก อโหสิ ธมฺมทินฺนาย ภิกฺขุนิยา สทฺธิํ กถาสลฺลาโป, ตํ สพฺพํ ภควโต อาโรเจสิ. เอวํ วุตฺเต ภควา วิสาขํ อุปาสกํ เอตทโวจ “ปณฺฑิตา วิสาข ธมฺมทินฺนา ภิกฺขุนี, มหาปญฺญา วิสาข ธมฺมทินฺนา ภิกฺขุนี. มํ เจปิ ตฺวํ วิสาข เอตมตฺถํ ปุจฺเฉยฺยาสิ, อหมฺปิ ตํ เอวเมวํ พฺยากเรยฺยํ, ยถา ตํ ธมฺมทินฺนาย ภิกฺขุนิยา พฺยากตํ, เอโส เจเวตสฺส\csromansharedfootnote{931a8162755b11c4}{เอโสเวตสฺส (สฺยา, กํ)} อตฺโถ, เอวญฺจ นํ\csromansharedfootnote{26bbab31f88a1781}{เอวเมตํ (สี, สฺยา, กํ)} ธาเรหี”ติ.}

\prosewithcloserruled{อิทมโวจ ภควา. อตฺตมโน วิสาโข อุปาสโก ภควโต ภาสิตํ อภินนฺทีติ.}{จูฬเวทลฺล\-สุตฺตํ นิฏฺฐิตํ จตุตฺถํ.}
\csromanmatikaanchor{mtk.79}
\csromantocmarkref{subsection}{5. จูฬธมฺมสมาทานสุตฺต}{mtk.79}
\bookmark[dest={mtk.79},level=2]{5. จูฬธมฺมสมาทานสุตฺต}
\csromanheader{2}{5. จูฬธมฺมสมาทาน\-สุตฺต}
\proseitem{468}{เอวํ เม สุตํ\csromandash{}เอกํ สมยํ ภควา สาวตฺถิยํ วิหรติ เชตวเน อนาถ\-ปิณฺฑิกสฺส อาราเม. ตตฺร โข ภควา ภิกฺขู อามนฺเตสิ “ภิกฺขโว”ติ. “ภทนฺเต”ติ เต ภิกฺขู ภควโต ปจฺจสฺโสสุํ. ภควา เอตทโวจ “จตฺตาริมานิ ภิกฺขเว ธมฺม\-สมาทานานิ. กตมานิ จตฺตาริ. อตฺถิ ภิกฺขเว ธมฺม\-สมาทานํ ปจฺจุปฺปนฺนสุขํ อายติํ ทุกฺขวิปากํ, อตฺถิ \csromanfolio{380}ภิกฺขเว ธมฺม\-สมาทานํ ปจฺจุปฺปนฺนทุกฺข\-ญฺเจว อายติํ จ ทุกฺขวิปากํ, อตฺถิ ภิกฺขเว ธมฺม\-สมาทานํ ปจฺจุปฺปนฺน\-ทุกฺขํ อายติํ สุขวิปากํ, อตฺถิ ภิกฺขเว ธมฺม\-สมาทานํ ปจฺจุปฺปนฺนสุข\-ญฺเจว อายติํ จ สุขวิปากํ.}

\proseitem{469}{กตมญฺจ ภิกฺขเว ธมฺม\-สมาทานํ ปจฺจุปฺปนฺนสุขํ อายติํ ทุกฺขวิปากํ. สนฺติ ภิกฺขเว เอเก สมณพฺราหฺมณา เอวํวาทิโน เอวํทิฏฺฐิโน “นตฺถิ กาเมสุ โทโส”ติ. เต กาเมสุ ปาตพฺยตํ อาปชฺชนฺติ. เต โข โมฬิพทฺธาหิ\csromansharedfootnote{0a8c72be0780d38f}{โมฬิพนฺธาหิ (สฺยา, กํ, ก)} ปริพฺพาชิกาหิ ปริจาเรนฺติ. เต เอวมาหํสุ “กิํสุ นาม เต โภนฺโต สมณพฺราหฺมณา กาเมสุ อนาคตภยํ สมฺปสฺสมานา กามานํ ปหานมาหํสุ, กามานํ ปริญฺญํ ปญฺญเปนฺติ. สุโข อิมิสฺสา ปริพฺพาชิกาย ตรุณาย มุทุกาย โลมสาย พาหาย สมฺผสฺโส”ติ. เต กาเมสุ ปาตพฺยตํ อาปชฺชนฺติ. เต กาเมสุ ปาตพฺยตํ อาปชฺชิตฺวา กายสฺส เภทา ปรํ มรณา อปายํ ทุคฺคติํ วินิปาตํ นิรยํ อุปปชฺชนฺติ. เต ตตฺถ ทุกฺขา ติพฺพา ขรา กฏุกา เวทนา เวทยนฺติ. เต เอวมาหํสุ “อิทํ โข เต โภนฺโต สมณพฺราหฺมณา กาเมสุ อนาคตภยํ สมฺปสฺสมานา กามานํ ปหานมาหํสุ, กามานํ ปริญฺญํ ปญฺญเปนฺติ, อิเม หิ มยํ กามเหตุ กามนิทานํ ทุกฺขา ติพฺพา ขรา กฏุกา เวทนา เวทยามา”ติ. เสยฺยถาปิ ภิกฺขเว คิมฺหานํ ปจฺฉิเม มาเส มาลุวาสิปาฏิกา ผเลยฺย, อถ โข ตํ ภิกฺขเว มาลุวาพีชํ อญฺญตรสฺมิํ สาลมูเล นิปเตยฺย. อถ โข ภิกฺขเว ยา ตสฺมิํ สาเล อธิวตฺถา เทวตา สา ภีตา สํวิคฺคา สนฺตาสํ อาปชฺเชยฺย. อถ โข ภิกฺขเว ตสฺมิํ สาเล อธิวตฺถาย เทวตาย มิตฺตามจฺจา ญาติสาโลหิตา อารามเทวตา วนเทวตา รุกฺขเทวตา โอสธิ\-ติณ\-วนปฺปตีสุ อธิวตฺถา เทวตา สงฺคมฺม สมาคมฺม เอวํ สมสฺสาเสยฺยุํ “มา ภวํ ภายิ มา ภวํ ภายิ, อปฺเปว นาเมตํ มาลุวาพีชํ โมโร วา คิเลยฺย\csromansharedfootnote{ca9fe125ddb1f573}{โมโร วา คิเลยฺย, โคธา วา ขาเทยฺย (ก)}, มโค วา ขาเทยฺย, ทวฑาโห\csromansharedfootnote{c410275eee33e281}{วนทาโห (ก)} วา ฑเหยฺย, วนกมฺมิกา วา อุทฺธเรยฺยุํ, อุปจิกา วา อุฏฺฐเหยฺยุํ\csromansharedfootnote{a9a4f2f79bfde578}{อุทฺรเภยฺยุํ (สี, อิ, ก)}, อพีชํ วา ปนสฺสา”ติ. อถ โข ตํ ภิกฺขเว มาลุวาพีชํ เนว โมโร คิเลยฺย, น มโค ขาเทยฺย, น ทวฑาโห ฑเหยฺย, น วนกมฺมิกา อุทฺธเรยฺยุํ, น อุปจิกา อุฏฺฐเหยฺยุํ, พีชํ จ ปนสฺส, ตํ ปาวุสฺสเกน เมเฆน \csromanfolio{381}อภิปฺปวุฏฺฐํ สมฺมเทว วิรุเหยฺย. สาสฺส มาลุวาลตา ตรุณา มุทุกา โลมสา วิลมฺพินี. สา ตํ สาลํ อุปนิเสเวยฺย. อถ โข ภิกฺขเว ตสฺมิํ สาเล อธิวตฺถาย เทวตาย เอวมสฺส “กิํสุ นาม เต โภนฺโต มิตฺตามจฺจา ญาติสาโลหิตา อารามเทวตา วนเทวตา รุกฺขเทวตา โอสธิ\-ติณ\-วนปฺปตีสุ อธิวตฺถา เทวตา มาลุวาพีเช อนาคตภยํ สมฺปสฺสมานา สงฺคมฺม สมาคมฺม เอวํ สมสฺสาเสสุํ\csromansharedfootnote{7a66549bc9c60156}{สมสฺสาเสยฺยุํ (ก)} ‘มา ภวํ ภายิ มา ภวํ ภายิ, อปฺเปว นาเมตํ มาลุวาพีชํ โมโร วา คิเลยฺย, มโค วา ขาเทยฺย, ทวฑโห วา ฑเหยฺย, วนกมฺมิกา วา อุทฺธเรยฺยุํ, อุปจิกา วา อุฏฺฐเหยฺยุํ, อพีชํ วา ปนสฺสา’ติ. สุโข อิมิสฺสา มาลุวาลตาย ตรุณาย มุทุกาย โลมสาย วิลมฺพินิยา สมฺผสฺโส”ติ. สา ตํ สาลํ อนุปริหเรยฺย, สา ตํ สาลํ อนุปริหริตฺวา อุปริ วิฏภิํ\csromansharedfootnote{1e4832ef19a7f7ff}{วิฏปํ (สฺยาฐ)} กเรยฺย, อุปริ วิฏภิํ กริตฺวา โอฆนํ ชเนยฺย, โอฆนํ ชเนตฺวา เย ตสฺส สาลสฺส มหนฺตา มหนฺตา ขนฺธา, เต ปทาเลยฺย. อถ โข ภิกฺขเว ตสฺมิํ สาเล อธิวตฺถาย เทวตาย เอวมสฺส “อิทํ โข เต โภนฺโต มิตฺตามจฺจา ญาติสาโลหิตา อารามเทวตา วนเทวตา รุกฺขเทวตา โอสธิ\-ติณ\-วนปฺปตีสุ อธิวตฺถา เทวตา มาลุวาพีเช อนาคตภยํ สมฺปสฺสมานา สงฺคมฺม สมาคมฺม เอวํ สมสฺสาเสสุํ\csromansharedfootnote{5a3b33cba8baef58}{ยํ วาหํ (ก), สฺวาหํ (สฺยา, กํ)} ‘มา ภวํ ภายิ มา ภวํ ภายิ, อปฺเปว นาเมตํ มาลุวาพีชํ โมโร วา คิเลยฺย, มโค วา ขาเทยฺย, ทวฑาโห วา ฑเหยฺย, วนกมฺมิกา วา อุทฺธเรยฺยุํ, อุปจิกา วา อุฏฺฐเหยฺยุํ, อพีชํ วา ปนสฺสา’ติ. ยญฺจาหํ\csromansharedfootnote{7a66549bc9c60156}{สมสฺสาเสยฺยุํ (ก)} มาลุวาพีชเหตุ ทุกฺขา ติพฺพา ขรา กฏุกา เวทนา เวทยามี”ติ. เอวเมว โข ภิกฺขเว สนฺติ เอเก สมณพฺราหฺมณา เอวํวาทิโน เอวํทิฏฺฐิโน “นตฺถิ กาเมสุ โทโส”ติ. เต กาเมสุ ปาตพฺยตํ อาปชฺชนฺติ, เต โมฬิพทฺธาหิ ปริพฺพาชิกาหิ ปริจาเรนฺติ. เต เอวมาหํสุ “กิํสุ นาม เต โภนฺโต สมณพฺราหฺมณา กาเมสุ อนาคตภยํ สมฺปสฺสมานา กามานํ ปหานมาหํสุ, กามานํ ปริญฺญํ ปญฺญเปนฺติ, สุโข อิมิสฺสา ปริพฺพาชิกาย ตรุณาย มุทุกาย โลมสาย พาหาย สมฺผสฺโส”ติ. เต กาเมสุ ปาตพฺยตํ อาปชฺชนฺติ, เต กาเมสุ ปาตพฺยตํ อาปชฺชิตฺวา กายสฺส เภทา ปรํ มรณา อปายํ ทุคฺคติํ วินิปาตํ นิรยํ อุปปชฺชนฺติ, เต ตตฺถ ทุกฺขา ติพฺพา \csromanfolio{382}ขรา กฏุกา เวทนา เวทยนฺติ. เต เอวมาหํสุ “อิทํ โข เต โภนฺโต สมณพฺราหฺมณา กาเมสุ อนาคตภยํ สมฺปสฺสมานา กามานํ ปหานมาหํสุ, กามานํ ปริญฺญํ ปญฺญเปนฺติ, อิเม หิ มยํ กามเหตุ กามนิทานํ ทุกฺขา ติพฺพา ขรา กฏุกา เวทนา เวทยามา”ติ. อิทํ วุจฺจติ ภิกฺขเว ธมฺม\-สมาทานํ ปจฺจุปฺปนฺนสุขํ อายติํ ทุกฺขวิปากํ.}

\proseitem{470}{กตมญฺจ ภิกฺขเว ธมฺม\-สมาทานํ ปจฺจุปฺปนฺนทุกฺข\-ญฺเจว อายติํ จ ทุกฺขวิปากํ. อิธ ภิกฺขเว เอกจฺโจ อเจลโก โหติ มุตฺตาจาโร, หตฺถา\-ปเลขโน, นเอหิภทฺทนฺติโก, นติฏฺฐ\-ภทฺทนฺติโก, นาภิหฏํ, น อุทฺทิสฺสกตํ, น นิมนฺตนํ สาทิยติ, โส น กุมฺภิมุขา ปฏิคฺคณฺหาติ, น กโฬวิมุขา ปฏิคฺคณฺหาติ, น เอฬกมนฺตรํ, น ทณฺฑมนฺตรํ, น มุสลมนฺตรํ, น ทฺวินฺนํ ภุญฺชมานานํ, น คพฺภินิยา, น ปายมานาย, น ปุริส\-นฺตร\-คตาย, น สงฺกิตฺตีสุ. น ยตฺถ สา อุปฏฺฐิโต โหติ, น ยตฺถ มกฺขิกา สณฺฑ\-สณฺฑ\-จารินี, น มจฺฉํ, น มํสํ, น สุรํ, น เมรยํ, น ถุโสทกํ ปิวติ. โส เอกาคาริโก วา โหติ เอกาโลปิโก, ทฺวาคาริโก วา โหติ ทฺวาโลปิโก ฯเปฯ สตฺตาคาริโก วา โหติ สตฺตาโลปิโก. เอกิสฺสาปิ ทตฺติยา ยาเปติ, ทฺวีหิปิ ทตฺตีหิ ยาเปติ ฯเปฯ สตฺตหิปิ ทตฺตีหิ ยาเปติ. เอกาหิกมฺปิ อาหารํ อาหาเรติ, ทฺวีหิกมฺปิ อาหารํ อาหาเรติ ฯเปฯ สตฺตาหิกมฺปิ อาหารํ อาหาเรติ, อิติ เอวรูปํ อทฺธมาสิกมฺปิ ปริยาย\-ภตฺตโภชนานุโยคมนุยุตฺโต วิหรติ. โส สากภกฺโข วา โหติ, สามากภกฺโข วา โหติ, นีวารภกฺโข วา โหติ, ททฺทุลภกฺโข วา โหติ, หฏภกฺโข วา โหติ, กณภกฺโข วา โหติ, อาจามภกฺโข วา โหติ, ปิญฺญากภกฺโข วา โหติ, ติณภกฺโข วา โหติ, โคมยภกฺโข วา โหติ, วนมูล\-ผลา\-หาโร ยาเปติ ปวตฺต\-ผล\-โภชี. โส สาณานิปิ ธาเรติ, มสาณานิปิ ธาเรติ, ฉวทุสฺสานิปิ ธาเรติ, ปํสุกูลานิปิ ธาเรติ, ติรีฏานิปิ ธาเรติ, อชินมฺปิ ธาเรติ, อชินกฺขิปมฺปิ ธาเรติ, กุสจีรมฺปิ ธาเรติ, วากจีรมฺปิ ธาเรติ, ผลกจีรมฺปิ ธาเรติ, เกสกมฺพลมฺปิ ธาเรติ, วาฬกมฺพลมฺปิ ธาเรติ, อุลูกปกฺขมฺปิ ธาเรติ, เกสมสฺสุ\-โลจโกปิ โหติ, เกสมสฺสุ\-โลจนา\-นุโยคม\-นุยุตฺโต, อุพฺภฏฺฐโกปิ โหติ, อาสน\-ปฏิกฺขิตฺโต, อุกฺกุฏิโกปิ โหติ อุกฺกุฏิก\-ปฺปธานมนุยุตฺโต, กณฺฏกา\-ปสฺสยิ\-โกปิ โหติ, กณฺฏกาปสฺสเย \csromanfolio{383}เสยฺยํ กปฺเปติ\csromansharedfootnote{62b59c6a9fdf6da9}{ปสฺส มหาสีหนาทสุตฺเต (111) ปิฏฺเฐ.}, สาย\-ตติยกมฺปิ อุทโก\-โรหนา\-นุโยคม\-นุยุตฺโต วิหรติ. อิติ เอวรูปํ อเนกวิหิตํ กายสฺส อาตาปน\-ปริตาปนา\-นุโยคม\-นุยุตฺโต วิหรติ, โส กายสฺส เภทา ปรํ มรณา อปายํ ทุคฺคติํ วินิปาตํ นิรยํ อุปปชฺชติ. อิทํ วุจฺจติ ภิกฺขเว ธมฺม\-สมาทานํ ปจฺจุปฺปนฺนทุกฺข\-ญฺเจว อายติํ จ ทุกฺขวิปากํ.}

\proseitem{471}{กตมญฺจ ภิกฺขเว ธมฺม\-สมาทานํ ปจฺจุปฺปนฺน\-ทุกฺขํ อายติํ สุขวิปากํ. อิธ ภิกฺขเว เอกจฺโจ ปกติยา ติพฺพ\-ราค\-ชาติโก โหติ, โส อภิกฺขณํ ราคชํ ทุกฺขํ โทมนสฺสํ ปฏิสํเวเทติ. ปกติยา ติพฺพ\-โทส\-ชาติโก โหติ, โส อภิกฺขณํ โทสชํ ทุกฺขํ โทมนสฺสํ ปฏิสํเวเทติ. ปกติยา ติพฺพ\-โมห\-ชาติโก โหติ, โส อภิกฺขณํ โมหชํ ทุกฺขํ โทมนสฺสํ ปฏิสํเวเทติ. โส สหาปิ ทุกฺเขน สหาปิ โทมนสฺเสน อสฺสุมุโขปิ รุทมาโน ปริปุณฺณํ ปริสุทฺธํ พฺรหฺมจริยํ จรติ, โส กายสฺส เภทา ปรํ มรณา สุคติํ สคฺคํ โลกํ อุปปชฺชติ. อิทํ วุจฺจติ ภิกฺขเว ธมฺม\-สมาทานํ ปจฺจุปฺปนฺน\-ทุกฺขํ อายติํ สุขวิปากํ.}

\proseitemwithcloserruled{472}{กตมญฺจ ภิกฺขเว ธมฺม\-สมาทานํ ปจฺจุปฺปนฺนสุข\-ญฺเจว อายติํ จ สุขวิปากํ. อิธ ภิกฺขเว เอกจฺโจ ปกติยา น ติพฺพ\-ราค\-ชาติโก โหติ, โส น อภิกฺขณํ ราคชํ ทุกฺขํ โทมนสฺสํ ปฏิสํเวเทติ. ปกติยา น ติพฺพ\-โทส\-ชาติโก โหติ, โส น อภิกฺขณํ โทสชํ ทุกฺขํ โทมนสฺสํ ปฏิสํเวเทติ. ปกติยา น ติพฺพ\-โมห\-ชาติโก โหติ, โส น อภิกฺขณํ โมหชํ ทุกฺขํ โทมนสฺสํ ปฏิสํเวเทติ. โส วิวิจฺเจว กาเมหิ วิวิจฺจ อกุสเลหิ ธมฺเมหิ สวิตกฺกํ สวิจารํ วิเวกชํ ปีติสุขํ ปฐมํ ฌานํ อุปสมฺปชฺช วิหรติ. วิตกฺก\-วิจารานํ วูปสมา อชฺฌตฺตํ สมฺปสาทนํ เจตโส เอโกทิภาวํ อวิตกฺกํ อวิจารํ สมาธิชํ ปีติสุขํ ทุติยํ ฌานํ ฯเปฯ ตติยํ ฌานํ ฯเปฯ จตุตฺถํ ฌานํ อุปสมฺปชฺช วิหรติ. โส กายสฺส เภทา ปรํ มรณา สุคติํ สคฺคํ โลกํ อุปปชฺชติ. อิทํ วุจฺจติ ภิกฺขเว ธมฺม\-สมาทานํ ปจฺจุปฺปนฺนสุข\-ญฺเจว อายติํ จ สุขวิปากํ. อิมานิ โข ภิกฺขเว จตฺตาริ ธมฺมสมาทานานีติ. \csromanfolio{384}อิทมโวจ ภควา. อตฺตมนา เต ภิกฺขู ภควโต ภาสิตํ อภินนฺทุนฺติ.}{จูฬธมฺมสมาทาน\-สุตฺตํ นิฏฺฐิตํ ปญฺจมํ.}
\csromanmatikaanchor{mtk.80}
\csromantocmarkref{subsection}{6. มหาธมฺมสมาทานสุตฺต}{mtk.80}
\bookmark[dest={mtk.80},level=2]{6. มหาธมฺมสมาทานสุตฺต}
\csromanheader{2}{6. มหา\-ธมฺมสมาทาน\-สุตฺต}
\proseitem{473}{เอวํ เม สุตํ\csromandash{}เอกํ สมยํ ภควา สาวตฺถิยํ วิหรติ เชตวเน อนาถ\-ปิณฺฑิกสฺส อาราเม. ตตฺร โข ภควา ภิกฺขู อามนฺเตสิ “ภิกฺขโว”ติ. “ภทนฺเต”ติ เต ภิกฺขู ภควโต ปจฺจสฺโสสุํ. ภควา เอตทโวจ “เยภุยฺเยน ภิกฺขเว สตฺตา เอวํกามา เอวํฉนฺทา เอวํอธิปฺปายา ‘อโห วต อนิฏฺฐา อกนฺตา อมนาปา ธมฺมา ปริหาเยยฺยุํ, อิฏฺฐา กนฺตา มนาปา ธมฺมา อภิวฑฺเฒยฺยุนฺ’ติ, เตสํ ภิกฺขเว สตฺตานํ เอวํกามานํ เอวํฉนฺทานํ เอวํ\-อธิปฺปายานํ อนิฏฺฐา อกนฺตา อมนาปา ธมฺมา อภิวฑฺฒนฺติ, อิฏฺฐา กนฺตา มนาปา ธมฺมา ปริหายนฺติ, ตตฺร ตุเมฺห ภิกฺขเว กํ เหตุํ ปจฺเจถา”ติ. ภควํมูลกา โน ภนฺเต ธมฺมา ภควํ\-เนตฺติกา ภควํ\-ปฏิสรณา, สาธุ วต ภนฺเต ภควนฺตญฺเญว ปฏิภาตุ เอตสฺส ภาสิตสฺส อตฺโถ, ภควโต สุตฺวา ภิกฺขู ธาเรสฺสนฺตีติ. เตน หิ ภิกฺขเว สุณาถ, สาธุกํ มนสิ กโรถ ภาสิสฺสามีติ. “เอวํ ภนฺเต”ติ โข เต ภิกฺขู ภควโต ปจฺจสฺโสสุํ. ภควา เอตทโวจ\csromandash{}}

\proseitem{474}{อิธ ภิกฺขเว อสฺสุตวา ปุถุชฺชโน อริยานํ อทสฺสาวี อริยธมฺมสฺส อโกวิโท อริยธมฺเม อวินีโต สปฺปุริสานํ อทสฺสาวี สปฺปุริส\-ธมฺมสฺส อโกวิโท สปฺปุริส\-ธมฺเม อวินีโต เสวิตพฺเพ ธมฺเม น ชานาติ, อเสวิตพฺเพ ธมฺเม น ชานาติ, ภชิตพฺเพ ธมฺเม น ชานาติ, อภชิตพฺเพ ธมฺเม น ชานาติ. โส เสวิตพฺเพ ธมฺเม อชานนฺโต อเสวิตพฺเพ ธมฺเม อชานนฺโต ภชิตพฺเพ ธมฺเม อชานนฺโต อภชิตพฺเพ ธมฺเม อชานนฺโต อเสวิตพฺเพ ธมฺเม เสวติ, เสวิตพฺเพ ธมฺเม น เสวติ, อภชิตพฺเพ ธมฺเม ภชติ, ภชิตพฺเพ ธมฺเม น ภชติ. ตสฺส อเสวิตพฺเพ ธมฺเม เสวโต เสวิตพฺเพ ธมฺเม อเสวโต อภชิตพฺเพ ธมฺเม ภชโต ภชิตพฺเพ ธมฺเม อภชโต อนิฏฺฐา อกนฺตา อมนาปา \csromanfolio{385}ธมฺมา อภิวฑฺฒนฺติ, อิฏฺฐา กนฺตา มนาปา ธมฺมา ปริหายนฺติ. ตํ กิสฺส เหตุ, เอวํ เหตํ ภิกฺขเว โหติ, ยถา ตํ อวิทฺทสุโน.}

\prose{สุตวา จ โข ภิกฺขเว อริยสาวโก อริยานํ ทสฺสาวี อริยธมฺมสฺส โกวิโท อริยธมฺเม สุวินีโต สปฺปุริสานํ ทสฺสาวี สปฺปุริส\-ธมฺมสฺส โกวิโท สปฺปุริส\-ธมฺเม สุวินีโต เสวิตพฺเพ ธมฺเม ชานาติ, อเสวิตพฺเพ ธมฺเม ชานาติ, ภชิตพฺเพ ธมฺเม ชานาติ, อภชิตพฺเพ ธมฺเม ชานาติ. โส เสวิตพฺเพ ธมฺเม ชานนฺโต อเสวิตพฺเพ ธมฺเม ชานนฺโต ภชิตพฺเพ ธมฺเม ชานนฺโต อภชิตพฺเพ ธมฺเม ชานนฺโต อเสวิตพฺเพ ธมฺเม น เสวติ, เสวิตพฺเพ ธมฺเม เสวติ, อภชิตพฺเพ ธมฺเม น ภชติ, ภชิตพฺเพ ธมฺเม ภชติ. ตสฺส อเสวิตพฺเพ ธมฺเม อเสวโต เสวิตพฺเพ ธมฺเม เสวโต อภชิตพฺเพ ธมฺเม อภชโต ภชิตพฺเพ ธมฺเม ภชโต อนิฏฺฐา อกนฺตา อมนาปา ธมฺมา ปริหายนฺติ, อิฏฺฐา กนฺตา มนาปา ธมฺมา อภิวฑฺฒนฺติ. ตํ กิสฺส เหตุ, เอวํ เหตํ ภิกฺขเว โหติ, ยถา ตํ วิทฺทสุโน.}

\proseitem{475}{จตฺตาริมานิ ภิกฺขเว ธมฺม\-สมาทานานิ. กตมานิ จตฺตาริ, อตฺถิ ภิกฺขเว ธมฺม\-สมาทานํ ปจฺจุปฺปนฺนทุกฺข\-ญฺเจว อายติํ จ ทุกฺขวิปากํ. อตฺถิ ภิกฺขเว ธมฺม\-สมาทานํ ปจฺจุปฺปนฺนสุขํ อายติํ ทุกฺขวิปากํ. อตฺถิ ภิกฺขเว ธมฺม\-สมาทานํ ปจฺจุปฺปนฺน\-ทุกฺขํ อายติํ สุขวิปากํ. อตฺถิ ภิกฺขเว ธมฺม\-สมาทานํ ปจฺจุปฺปนฺนสุข\-ญฺเจว อายติํ จ สุขวิปากํ.}

\proseitem{476}{ตตฺร ภิกฺขเว ยมิทํ\csromansharedfootnote{ee166d3bd0f984bc}{ยทิทํ (สี)} ธมฺม\-สมาทานํ ปจฺจุปฺปนฺนทุกฺข\-ญฺเจว อายติํ จ ทุกฺขวิปากํ. ตํ อวิทฺวา อวิชฺชาคโต ยถาภูตํ นปฺปชานาติ “อิทํ โข ธมฺม\-สมาทานํ ปจฺจุปฺปนฺนทุกฺข\-ญฺเจว อายติํ จ ทุกฺขวิปากนฺ”ติ, ตํ อวิทฺวา อวิชฺชาคโต ยถาภูตํ อปฺปชานนฺโต ตํ เสวติ, ตํ น ปริวชฺเชติ, ตสฺส ตํ เสวโต ตํ อปริวชฺชยโต อนิฏฺฐา อกนฺตา อมนาปา ธมฺมา อภิวฑฺฒนฺติ, อิฏฺฐา กนฺตา มนาปา ธมฺมา ปริหายนฺติ. ตํ กิสฺส เหตุ, เอวํ เหตํ ภิกฺขเว โหติ, ยถา ตํ อวิทฺทสุโน. (1)}

\prose{ตตฺร ภิกฺขเว ยมิทํ ธมฺม\-สมาทานํ ปจฺจุปฺปนฺนสุขํ อายติํ ทุกฺขวิปากํ, ตํ อวิทฺวา อวิชฺชาคโต ยถาภูตํ นปฺปชานาติ “อิทํ โข ธมฺม\-สมาทานํ \csromanfolio{386}ปจฺจุปฺปนฺนสุขํ อายติํ ทุกฺขวิปากนฺ”ติ, ตํ อวิทฺวา อวิชฺชาคโต ยถาภูตํ อปฺปชานนฺโต ตํ เสวติ, ตํ น ปริวชฺเชติ, ตสฺส ตํ เสวโต ตํ อปริวชฺชยโต อนิฏฺฐา อกนฺตา อมนาปา ธมฺมา อภิวฑฺฒนฺติ, อิฏฺฐา กนฺตา มนาปา ธมฺมา ปริหายนฺติ. ตํ กิสฺส เหตุ, เอวํ เหตํ ภิกฺขเว โหติ, ยถา ตํ อวิทฺทสุโน. (2)}

\prose{ตตฺร ภิกฺขเว ยมิทํ ธมฺม\-สมาทานํ ปจฺจุปฺปนฺน\-ทุกฺขํ อายติํ สุขวิปากํ, ตํ อวิทฺวา อวิชฺชาคโต ยถาภูตํ นปฺปชานาติ, “อิทํ โข ธมฺม\-สมาทานํ ปจฺจุปฺปนฺน\-ทุกฺขํ อายติํ สุขวิปากนฺ”ติ, ตํ อวิทฺวา อวิชฺชาคโต ยถาภูตํ อปฺปชานนฺโต ตํ น เสวติ, ตํ ปริวชฺเชติ, ตสฺส ตํ อเสวโต ตํ ปริวชฺชยโต อนิฏฺฐา อกนฺตา อมนาปา ธมฺมา อภิวฑฺฒนฺติ, อิฏฺฐา กนฺตา มนาปา ธมฺมา ปริหายนฺติ. ตํ กิสฺส เหตุ, เอวํ เหตํ ภิกฺขเว โหติ, ยถา ตํ อวิทฺทสุโน. (3)}

\prose{ตตฺร ภิกฺขเว ยมิทํ ธมฺม\-สมาทานํ ปจฺจุปฺปนฺนสุข\-ญฺเจว อายติํ จ สุขวิปากํ, ตํ อวิทฺวา อวิชฺชาคโต ยถาภูตํ นปฺปชานาติ “อิทํ โข ธมฺม\-สมาทานํ ปจฺจุปฺปนฺนสุข\-ญฺเจว อายติํ จ สุขวิปากนฺ”ติ, ตํ อวิทฺวา อวิชฺชาคโต ยถาภูตํ อปฺปชานนฺโต ตํ น เสวติ, ตํ ปริวชฺเชติ, ตสฺส ตํ อเสวโต ตํ ปริวชฺชยโต อนิฏฺฐา อกนฺตา อมนาปา ธมฺมา อภิวฑฺฒนฺติ, อิฏฺฐา กนฺตา มนาปา ธมฺมา ปริหายนฺติ. ตํ กิสฺส เหตุ, เอวํ เหตํ ภิกฺขเว โหติ, ยถา ตํ อวิทฺทสุโน. (4)}

\proseitem{477}{ตตฺร ภิกฺขเว ยมิทํ ธมฺม\-สมาทานํ ปจฺจุปฺปนฺนทุกฺข\-ญฺเจว อายติํ จ ทุกฺขวิปากํ. ตํ วิทฺวา วิชฺชาคโต ยถาภูตํ ปชานาติ “อิทํ โข ธมฺม\-สมาทานํ ปจฺจุปฺปนฺนทุกฺข\-ญฺเจว อายติํ จ ทุกฺขวิปากนฺ”ติ, ตํ วิทฺวา วิชฺชาคโต ยถาภูตํ ปชานนฺโต ตํ น เสวติ, ตํ ปริวชฺเชติ, ตสฺส ตํ อเสวโต ตํ ปริวชฺชยโต อนิฏฺฐา อกนฺตา อมนาปา ธมฺมา ปริหายนฺติ, อิฏฺฐา กนฺตา มนาปา ธมฺมา อภิวฑฺฒนฺติ. ตํ กิสฺส เหตุ, เอวํ เหตํ ภิกฺขเว โหติ, ยถา ตํ วิทฺทสุโน. (1)}

\prose{ตตฺร ภิกฺขเว ยมิทํ ธมฺม\-สมาทานํ ปจฺจุปฺปนฺนสุขํ อายติํ ทุกฺขวิปากํ, ตํ วิทฺวา วิชฺชาคโต ยถาภูตํ ปชานาติ “อิทํ โข ธมฺม\-สมาทานํ ปจฺจุปฺปนฺนสุขํ อายติํ ทุกฺขวิปากนฺ”ติ, ตํ วิทฺวา วิชฺชาคโต ยถาภูตํ ปชานนฺโต ตํ น เสวติ, ตํ ปริวชฺเชติ, ตสฺส ตํ อเสวโต ตํ \csromanfolio{387}ปริวชฺชยโต อนิฏฺฐา อกนฺตา อมนาปา ธมฺมา ปริหายนฺติ, อิฏฺฐา กนฺตา มนาปา ธมฺมา อภิวฑฺฒนฺติ. ตํ กิสฺส เหตุ, เอวํ เหตํ ภิกฺขเว โหติ, ยถา ตํ วิทฺทสุโน. (2)}

\prose{ตตฺร ภิกฺขเว ยมิทํ ธมฺม\-สมาทานํ ปจฺจุปฺปนฺน\-ทุกฺขํ อายติํ สุขวิปากํ, ตํ วิทฺวา วิชฺชาคโต ยถาภูตํ ปชานาติ “อิทํ โข ธมฺม\-สมาทานํ ปจฺจุปฺปนฺน\-ทุกฺขํ อายติํ สุขวิปากนฺ”ติ, ตํ วิทฺวา วิชฺชาคโต ยถาภูตํ ปชานนฺโต ตํ เสวติ, ตํ น ปริวชฺเชติ, ตสฺส ตํ เสวโต ตํ อปริวชฺชยโต อนิฏฺฐา อกนฺตา อมนาปา ธมฺมา ปริหายนฺติ, อิฏฺฐา กนฺตา มนาปา ธมฺมา อภิวฑฺฒนฺติ. ตํ กิสฺส เหตุ, เอวํ เหตํ ภิกฺขเว โหติ, ยถา ตํ วิทฺทสุโน. (3)}

\prose{ตตฺร ภิกฺขเว ยมิทํ ธมฺม\-สมาทานํ ปจฺจุปฺปนฺนสุข\-ญฺเจว อายติํ จ สุขวิปากํ, ตํ วิทฺวา วิชฺชาคโต ยถาภูตํ ปชานาติ “อิทํ โข ธมฺม\-สมาทานํ ปจฺจุปฺปนฺนสุข\-ญฺเจว อายติํ จ สุขวิปากนฺ”ติ, ตํ วิทฺวา วิชฺชาคโต ยถาภูตํ ปชานนฺโต ตํ เสวติ, ตํ น ปริวชฺเชติ, ตสฺส ตํ เสวโต ตํ อปริวชฺชยโต อนิฏฺฐา อกนฺตา อมนาปา ธมฺมา ปริหายนฺติ, อิฏฺฐา กนฺตา มนาปา ธมฺมา อภิวฑฺฒนฺติ. ตํ กิสฺส เหตุ, เอวํ เหตํ ภิกฺขเว โหติ, ยถา ตํ วิทฺทสุโน. (4)}

\proseitem{478}{กตมญฺจ ภิกฺขเว ธมฺม\-สมาทานํ ปจฺจุปฺปนฺนทุกฺข\-ญฺเจว อายติํ จ ทุกฺขวิปากํ. อิธ ภิกฺขเว เอกจฺโจ สหาปิ ทุกฺเขน สหาปิ โทมนสฺเสน ปาณาติปาตี โหติ, ปาณาติปาต\-ปจฺจยา จ ทุกฺขํ โทมนสฺสํ ปฏิสํเวเทติ. สหาปิ ทุกฺเขน สหาปิ โทมนสฺเสน อทินฺนาทายี โหติ, อทินฺนาทาน\-ปจฺจยา จ ทุกฺขํ โทมนสฺสํ ปฏิสํเวเทติ. สหาปิ ทุกฺเขน สหาปิ โทมนสฺเสน กาเมสุ มิจฺฉาจารี โหติ, กาเมสุ มิจฺฉาจาร\-ปจฺจยา จ ทุกฺขํ โทมนสฺสํ ปฏิสํเวเทติ. สหาปิ ทุกฺเขน สหาปิ โทมนสฺเสน มุสาวาที โหติ, มุสาวาท\-ปจฺจยา จ ทุกฺขํ โทมนสฺสํ ปฏิสํเวเทติ. สหาปิ ทุกฺเขน สหาปิ โทมนสฺเสน ปิสุณวาโจ โหติ, ปิสุณวาจา\-ปจฺจยา จ ทุกฺขํ โทมนสฺสํ ปฏิสํเวเทติ. สหาปิ ทุกฺเขน สหาปิ โทมนสฺเสน ผรุสวาโจ โหติ, ผรุสวาจา\-ปจฺจยา จ ทุกฺขํ โทมนสฺสํ ปฏิสํเวเทติ. สหาปิ ทุกฺเขน สหาปิ โทมนสฺเสน สมฺผปฺปลาปี โหติ, สมฺผปฺปลาป\-ปจฺจยา จ ทุกฺขํ โทมนสฺสํ ปฏิสํเวเทติ. \csromanfolio{388}สหาปิ ทุกฺเขน สหาปิ โทมนสฺเสน อภิชฺฌาลุ โหติ, อภิชฺฌา\-ปจฺจยา จ ทุกฺขํ โทมนสฺสํ ปฏิสํเวเทติ. สหาปิ ทุกฺเขน สหาปิ โทมนสฺเสน พฺยาปนฺนจิตฺโต โหติ, พฺยาปาทปจฺจยา จ ทุกฺขํ โทมนสฺสํ ปฏิสํเวเทติ. สหาปิ ทุกฺเขน สหาปิ โทมนสฺเสน มิจฺฉาทิฏฺฐิ โหติ, มิจฺฉาทิฏฺฐิ\-ปจฺจยา จ ทุกฺขํ โทมนสฺสํ ปฏิสํเวเทติ. โส กายสฺส เภทา ปรํ มรณา อปายํ ทุคฺคติํ วินิปาตํ นิรยํ อุปปชฺชติ. อิทํ วุจฺจติ ภิกฺขเว ธมฺม\-สมาทานํ ปจฺจุปฺปนฺนทุกฺข\-ญฺเจว อายติํ จ ทุกฺขวิปากํ. (1)}

\proseitem{479}{กตมญฺจ ภิกฺขเว ธมฺม\-สมาทานํ ปจฺจุปฺปนฺนสุขํ อายติํ ทุกฺขวิปากํ. อิธ ภิกฺขเว เอกจฺโจ สหาปิ สุเขน สหาปิ โสมนสฺเสน ปาณาติปาตี โหติ, ปาณาติปาต\-ปจฺจยา จ สุขํ โสมนสฺสํ ปฏิสํเวเทติ. สหาปิ สุเขน สหาปิ โสมนสฺเสน อทินฺนาทายี โหติ, อทินฺนาทาน\-ปจฺจยา จ สุขํ โสมนสฺสํ ปฏิสํเวเทติ. สหาปิ สุเขน สหาปิ โสมนสฺเสน กาเมสุ\-มิจฺฉาจารี โหติ, กาเมสุ\-มิจฺฉาจาร\-ปจฺจยา จ สุขํ โสมนสฺสํ ปฏิสํเวเทติ. สหาปิ สุเขน สหาปิ โสมนสฺเสน มุสาวาที โหติ, มุสาวาท\-ปจฺจยา จ สุขํ โสมนสฺสํ ปฏิสํเวเทติ. สหาปิ สุเขน สหาปิ โสมนสฺเสน ปิสุณวาโจ โหติ, ปิสุณวาจา\-ปจฺจยา จ สุขํ โสมนสฺสํ ปฏิสํเวเทติ. สหาปิ สุเขน สหาปิ โสมนสฺเสน ผรุสวาโจ โหติ, ผรุสวาจา\-ปจฺจยา จ สุขํ โสมนสฺสํ ปฏิสํเวเทติ. สหาปิ สุเขน สหาปิ โสมนสฺเสน สมฺผปฺปลาปี โหติ, สมฺผปฺปลาป\-ปจฺจยา จ สุขํ โสมนสฺสํ ปฏิสํเวเทติ. สหาปิ สุเขน สหาปิ โสมนสฺเสน อภิชฺฌาลุ โหติ, อภิชฺฌา\-ปจฺจยา จ สุขํ โสมนสฺสํ ปฏิสํเวเทติ. สหาปิ สุเขน สหาปิ โสมนสฺเสน พฺยาปนฺนจิตฺโต โหติ, พฺยาปาทปจฺจยา จ สุขํ โสมนสฺสํ ปฏิสํเวเทติ. สหาปิ สุเขน สหาปิ โสมนสฺเสน มิจฺฉาทิฏฺฐิ โหติ, มิจฺฉาทิฏฺฐิ\-ปจฺจยา จ สุขํ โสมนสฺสํ ปฏิสํเวเทติ. โส กายสฺส เภทา ปรํ มรณา อปายํ ทุคฺคติํ วินิปาตํ นิรยํ อุปปชฺชติ. อิทํ วุจฺจติ ภิกฺขเว ธมฺม\-สมาทานํ ปจฺจุปฺปนฺนสุขํ อายติํ ทุกฺขวิปากํ. (2)}

\proseitem{480}{กตมญฺจ ภิกฺขเว ธมฺม\-สมาทานํ ปจฺจุปฺปนฺน\-ทุกฺขํ อายติํ สุขวิปากํ. อิธ ภิกฺขเว เอกจฺโจ สหาปิ ทุกฺเขน สหาปิ โทมนสฺเสน ปาณาติปาตา ปฏิวิรโต โหติ, ปาณาติปาตา เวรมณี\-ปจฺจยา จ ทุกฺขํ โทมนสฺสํ \csromanfolio{389}ปฏิสํเวเทติ. สหาปิ ทุกฺเขน สหาปิ โทมนสฺเสน อทินฺนาทานา ปฏิวิรโต โหติ, อทินฺนาทานา เวรมณี\-ปจฺจยา จ ทุกฺขํ โทมนสฺสํ ปฏิสํเวเทติ. สหาปิ ทุกฺเขน สหาปิ โทมนสฺเสน กาเมสุ\-มิจฺฉาจารา ปฏิวิรโต โหติ, กาเมสุ\-มิจฺฉาจารา เวรมณี\-ปจฺจยา จ ทุกฺขํ โทมนสฺสํ ปฏิสํเวเทติ. สหาปิ ทุกฺเขน สหาปิ โทมนสฺเสน มุสาวาทา ปฏิวิรโต โหติ, มุสาวาทา เวรมณี\-ปจฺจยา จ ทุกฺขํ โทมนสฺสํ ปฏิสํเวเทติ. สหาปิ ทุกฺเขน สหาปิ โทมนสฺเสน ปิสุณาย วาจาย ปฏิวิรโต โหติ, ปิสุณาย วาจาย เวรมณี\-ปจฺจยา จ ทุกฺขํ โทมนสฺสํ ปฏิสํเวเทติ. สหาปิ ทุกฺเขน สหาปิ โทมนสฺเสน ผรุสาย วาจาย ปฏิวิรโต โหติ, ผรุสาย วาจาย เวรมณี\-ปจฺจยา จ ทุกฺขํ โทมนสฺสํ ปฏิสํเวเทติ. สหาปิ ทุกฺเขน สหาปิ โทมนสฺเสน สมฺผปฺปลาปา ปฏิวิรโต โหติ, สมฺผปฺปลาปา เวรมณี ปจฺจยา จ ทุกฺขํ โทมนสฺสํ ปฏิสํเวเทติ. สหาปิ ทุกฺเขน สหาปิ โทมนสฺเสน อนภิชฺฌาลุ โหติ, อนภิชฺฌา\-ปจฺจยา จ ทุกฺขํ โทมนสฺสํ ปฏิสํเวเทติ. สหาปิ ทุกฺเขน สหาปิ โทมนสฺเสน อพฺยาปนฺนจิตฺโต โหติ, อพฺยาปาท\-ปจฺจยา จ ทุกฺขํ โทมนสฺสํ ปฏิสํเวเทติ. สหาปิ ทุกฺเขน สหาปิ โทมนสฺเสน สมฺมาทิฏฺฐิ โหติ, สมฺมาทิฏฺฐิ\-ปจฺจยา จ ทุกฺขํ โทมนสฺสํ ปฏิสํเวเทติ. โส กายสฺส เภทา ปรํ มรณา สุคติํ สคฺคํ โลกํ อุปปชฺชติ. อิทํ วุจฺจติ ภิกฺขเว ธมฺม\-สมาทานํ ปจฺจุปฺปนฺน\-ทุกฺขํ อายติํ สุขวิปากํ. (3)}

\proseitem{481}{กตมญฺจ ภิกฺขเว ธมฺม\-สมาทานํ ปจฺจุปฺปนฺนสุข\-ญฺเจว อายติํ จ สุขวิปากํ. อิธ ภิกฺขเว เอกจฺโจ สหาปิ สุเขน สหาปิ โสมนสฺเสน ปาณาติปาตา ปฏิวิรโต โหติ, ปาณาติปาตา เวรมณี\-ปจฺจยา จ สุขํ โสมนสฺสํ ปฏิสํเวเทติ. สหาปิ สุเขน สหาปิ โสมนสฺเสน อทินฺนาทานา ปฏิวิรโต โหติ, อทินฺนาทานา เวรมณี\-ปจฺจยา จ สุขํ โสมนสฺสํ ปฏิสํเวเทติ. สหาปิ สุเขน สหาปิ โสมนสฺเสน กาเมสุ\-มิจฺฉาจารา ปฏิวิรโต โหติ, กาเมสุ\-มิจฺฉาจารา เวรมณี\-ปจฺจยา จ สุขํ โสมนสฺสํ ปฏิสํเวเทติ. สหาปิ สุเขน สหาปิ โสมนสฺเสน มุสาวาทา ปฏิวิรโต โหติ, มุสาวาทา เวรมณี\-ปจฺจยา จ สุขํ โสมนสฺสํ ปฏิสํเวเทติ. สหาปิ สุเขน สหาปิ โสมนสฺเสน ปิสุณาย วาจาย ปฏิวิรโต โหติ, ปิสุณาย \csromanfolio{390}วาจาย เวรมณี\-ปจฺจยา จ สุขํ โสมนสฺสํ ปฏิสํเวเทติ. สหาปิ สุเขน สหาปิ โสมนสฺเสน ผรุสาย วาจาย ปฏิวิรโต โหติ, ผรุสาย วาจาย เวรมณี\-ปจฺจยา จ สุขํ โสมนสฺสํ ปฏิสํเวเทติ. สหาปิ สุเขน สหาปิ โสมนสฺเสน สมฺผปฺปลาปา ปฏิวิรโต โหติ, สมฺผปฺปลาปา เวรมณี\-ปจฺจยา จ สุขํ โสมนสฺสํ ปฏิสํเวเทติ. สหาปิ สุเขน สหาปิ โสมนสฺเสน อนภิชฺฌาลุ โหติ, อนภิชฺฌา\-ปจฺจยา จ สุขํ โสมนสฺสํ ปฏิสํเวเทติ. สหาปิ สุเขน สหาปิ โสมนสฺเสน อพฺยาปนฺนจิตฺโต โหติ, อพฺยาปาท\-ปจฺจยา จ สุขํ โสมนสฺสํ ปฏิสํเวเทติ. สหาปิ สุเขน สหาปิ โสมนสฺเสน สมฺมาทิฏฺฐิ โหติ, สมฺมาทิฏฺฐิ\-ปจฺจยา จ สุขํ โสมนสฺสํ ปฏิสํเวเทติ. โส กายสฺส เภทา ปรํ มรณา สุคติํ สคฺคํ โลกํ อุปปชฺชติ. อิทํ วุจฺจติ ภิกฺขเว ธมฺม\-สมาทานํ ปจฺจุปฺปนฺนสุข\-ญฺเจว อายติํ จ สุขวิปากํ. อิมานิ โข ภิกฺขเว จตฺตาริ ธมฺม\-สมาทานานิ. (4)}

\proseitem{482}{เสยฺยถาปิ ภิกฺขเว ติตฺตกาลาพุ วิเสน สํสฏฺโฐ, อถ ปุริโส อาคจฺเฉยฺย ชีวิตุกาโม อมริตุกาโม สุขกาโม ทุกฺข\-ปฺปฏิกูโล, ตเมนํ เอวํ วเทยฺยุํ “อมฺโภ ปุริส อยํ ติตฺตกาลาพุ วิเสน สํสฏฺโฐ, สเจ อากงฺขสิ, ปิว\csromansharedfootnote{a55fc6d1d3689a4e}{ปิป (สี, อิ)}. ตสฺส เต ปิวโต\csromansharedfootnote{88efb13f6fcbdafb}{ปิปโต (สี, อิ)} เจว นจฺฉาเทสฺสติ วณฺเณนปิ คนฺเธนปิ รเสนปิ, ปิวิตฺวา\csromansharedfootnote{bb603970a2e6d618}{ปีตฺวา (สี)} จ ปน มรณํ วา นิคจฺฉสิ มรณมตฺตํ วา ทุกฺขนฺ”ติ. โส ตํ อปฺปฏิสงฺขาย ปิเวยฺย, นปฺปฏินิสฺสชฺเชยฺย, ตสฺส ตํ ปิวโต เจว นจฺฉาเทยฺย วณฺเณนปิ คนฺเธนปิ รเสนปิ, ปิวิตฺวา จ ปน มรณํ วา นิคจฺเฉยฺย มรณมตฺตํ วา ทุกฺขํ. ตถูปมาหํ ภิกฺขเว อิมํ ธมฺม\-สมาทานํ วทามิ, ยมิทํ ธมฺม- สมาทานํ ปจฺจุปฺปนฺนทุกฺข\-ญฺเจว อายติํ จ ทุกฺขวิปากํ. (1)}

\proseitem{483}{เสยฺยถาปิ ภิกฺขเว อาปานียกํโส วณฺณสมฺปนฺโน คนฺธสมฺปนฺโน รสสมฺปนฺโน, โส จ โข วิเสน สํสฏฺโฐ, อถ ปุริโส อาคจฺเฉยฺย ชีวิตุกาโม อมริตุกาโม สุขกาโม ทุกฺข\-ปฺปฏิกูโล, ตเมนํ เอวํ วเทยฺยุํ “อมฺโภ ปุริส อยํ อาปานียกํโส วณฺณสมฺปนฺโน คนฺธสมฺปนฺโน รสสมฺปนฺโน, โส จ โข วิเสน สํสฏฺโฐ, สเจ อากงฺขสิ, ปิว. ตสฺส เต ปิวโตหิ\csromansharedfootnote{eda87181fd3f1ff8}{ปิวโตปิ (ก)} โข ฉาเทสฺสติ วณฺเณนปิ คนฺเธนปิ รเสนปิ, ปิวิตฺวา จ ปน มรณํ วา นิคจฺฉสิ มรณมตฺตํ วา ทุกฺขนฺ”ติ. โส ตํ \csromanfolio{391}อปฺปฏิสงฺขาย ปิเวยฺย, นปฺปฏินิสฺสชฺเชยฺย, ตสฺส ตํ ปิวโตหิ โข ฉาเทยฺย วณฺเณนปิ คนฺเธนปิ รเสนปิ, ปิวิตฺวา จ ปน มรณํ วา นิคจฺเฉยฺย มรณมตฺตํ วา ทุกฺขํ. ตถูปมาหํ ภิกฺขเว อิมํ ธมฺม\-สมาทานํ วทามิ, ยมิทํ ธมฺม\-สมาทานํ ปจฺจุปฺปนฺนสุขํ อายติํ ทุกฺขวิปากํ. (2)}

\proseitem{484}{เสยฺยถาปิ ภิกฺขเว ปูติมุตฺตํ นานาเภสชฺเชหิ สํสฏฺฐํ, อถ ปุริโส อาคจฺเฉยฺย ปณฺฑุกโรคี, ตเมนํ เอวํ วเทยฺยุํ “อมฺโภ ปุริส อิทํ ปูติมุตฺตํ นานาเภสชฺเชหิ สํสฏฺฐํ, สเจ อากงฺขสิ, ปิว. ตสฺส เต ปิวโตหิ โข นจฺฉาเทสฺสติ วณฺเณนปิ คนฺเธนปิ รเสนปิ, ปิวิตฺวา จ ปน สุขี ภวิสฺสสี”ติ. โส ตํ ปฏิสงฺขาย ปิเวยฺย, นปฺปฏินิสฺสชฺเชยฺย, ตสฺส ตํ ปิวโตหิ โข นจฺฉาเทยฺย วณฺเณนปิ คนฺเธนปิ รเสนปิ, ปิวิตฺวา จ ปน สุขี อสฺส. ตถูปมาหํ ภิกฺขเว อิมํ ธมฺม\-สมาทานํ วทามิ, ยมิทํ ธมฺม\-สมาทานํ ปจฺจุปฺปนฺน\-ทุกฺขํ อายติํ สุขวิปากํ. (3)}

\proseitem{485}{เสยฺยถาปิ ภิกฺขเว ทธิ จ มธุ จ สปฺปิ จ ผาณิตํ จ เอกชฺฌํ สํสฏฺฐํ, อถ ปุริโส อาคจฺเฉยฺย โลหิต\-ปกฺขนฺทิโก, ตเมนํ เอวํ วเทยฺยุํ “อมฺโภ ปุริส อิทํ ทธิ จ มธุ จ สปฺปิ จ ผาณิตํ จ เอกชฺฌํ สํสฏฺฐํ, สเจ อากงฺขสิ, ปิว. ตสฺส เต ปิวโต เจว ฉาเทสฺสติ วณฺเณนปิ คนฺเธนปิ รเสนปิ, ปิวิตฺวา จ ปน สุขี ภวิสฺสสี”ติ. โส ตํ ปฏิสงฺขาย ปิเวยฺย, นปฺปฏินิสฺสชฺเชยฺย, ตสฺส ตํ ปิวโต เจว ฉาเทยฺย วณฺเณนปิ คนฺเธนปิ รเสนปิ, ปิวิตฺวา จ ปน สุขี อสฺส. ตถูปมาหํ ภิกฺขเว อิมํ ธมฺม\-สมาทานํ วทามิ, ยมิทํ ธมฺม\-สมาทานํ ปจฺจุปฺปนฺนสุข\-ญฺเจว อายติํ จ สุขวิปากํ. (4)}

\proseitem{486}{เสยฺยถาปิ ภิกฺขเว วสฺสานํ ปจฺฉิเม มาเส สรทสมเย วิทฺเธ วิคตวลาหเก เทเว อาทิจฺโจ นภํ อพฺภุสฺสกฺกมาโน สพฺพํ อาตาสคตํ ตมคตํ อภิวิหจฺจ ภาสเต จ ตปเต จ วิโรจเต จ. เอวเมว โข ภิกฺขเว ยมิทํ ธมฺม\-สมาทานํ ปจฺจุปฺปนฺนสุข\-ญฺเจว อายติํ จ สุขวิปากํ, ตทญฺเญ ปุถุสมณพฺราหฺมณปรปฺปวาเท อภิวิหจฺจ ภาสเต จ ตปเต จ วิโรจเต จาติ.}

\prosewithcloser{อิทมโวจ ภควา. อตฺตมนา เต ภิกฺขู ภควโต ภาสิตํ อภินนฺทุนฺติ.}{มหา\-ธมฺมสมาทาน\-สุตฺตํ นิฏฺฐิตํ ฉฏฺฐํ.}
\csromanmatikaanchor{mtk.81}
\csromantocmarkref{subsection}{7. วีมํสกสุตฺต}{mtk.81}
\bookmark[dest={mtk.81},level=2]{7. วีมํสกสุตฺต}
\csromanheader{2}{7. \textbf{วีมํสกสุตฺต}}
\proseitem{487}{\csromanfolio{392}เอวํ เม สุตํ\csromandash{}เอกํ สมยํ ภควา สาวตฺถิยํ วิหรติ เชตวเน อนาถ\-ปิณฺฑิกสฺส อาราเม. ตตฺร โข ภควา ภิกฺขู อามนฺเตสิ “ภิกฺขโว”ติ. “ภทนฺเต”ติ เต ภิกฺขู ภควโต ปจฺจสฺโสสุํ. ภควา เอตทโวจ “วีมํสเกน ภิกฺขเว ภิกฺขุนา ปรสฺส เจโตปริยายํ อชานนฺเตน\csromansharedfootnote{06b34b3a60fe4b33}{อาชานนฺเตน (อิ, ก), อชานนฺเตน กินฺติ (?)} ตถาคเต สมนฺเนสนา กาตพฺพา สมฺมาสมฺพุทฺโธ วา โน วา อิติ วิญฺญาณายา”ติ. ภควํมูลกา โน ภนฺเต ธมฺมา ภควํ\-เนตฺติกา ภควํ\-ปฏิสรณา, สาธุ วต ภนฺเต ภควนฺตํเยว ปฏิภาตุ เอตสฺส ภาสิตสฺส อตฺโถ, ภควโต สุตฺวา ภิกฺขู ธาเรสฺสนฺตีติ. เตน หิ ภิกฺขเว สุณาถ สาธุกํ มนสิ กโรถ ภาสิสฺสามีติ. “เอวํ ภนฺเต”ติ โข เต ภิกฺขู ภควโต ปจฺจสฺโสสุํ. ภควา เอตทโวจ\csromandash{}}

\proseitem{488}{วีมํสเกน ภิกฺขเว ภิกฺขุนา ปรสฺส เจโตปริยายํ อชานนฺเตน ทฺวีสุ ธมฺเมสุ ตถาคโต สมนฺเนสิตพฺโพ จกฺขุ\-โสต\-วิญฺเญยฺเยสุ ธมฺเมสุ “เย สํกิลิฏฺฐา จกฺขุ\-โสต\-วิญฺเญยฺยา ธมฺมา, สํวิชฺชนฺติ วา เต ตถาคตสฺส โน วา”ติ. ตเมนํ สมนฺเนสมาโน เอวํ ชานาติ “เย สํกิลิฏฺฐา จกฺขุ\-โสต\-วิญฺเญยฺยา ธมฺมา, น เต ตถาคตสฺส สํวิชฺชนฺตี”ติ.}

\prose{ยโต นํ สมนฺเนสมาโน เอวํ ชานาติ “เย สํกิลิฏฺฐา จกฺขุ\-โสต\-วิญฺเญยฺยา ธมฺมา, น เต ตถาคตสฺส สํวิชฺชนฺตี”ติ, ตโต นํ อุตฺตริํ สมนฺเนสติ “เย วีติมิสฺสา จกฺขุ\-โสต\-วิญฺเญยฺยา ธมฺมา, สํวิชฺชนฺติ วา เต ตถาคตสฺส โน วา”ติ. ตเมนํ สมนฺเนสมาโน เอวํ ชานาติ “เย วีติมิสฺสา จกฺขุ\-โสต\-วิญฺเญยฺยา ธมฺมา, น เต ตถาคตสฺส สํวิชฺชนฺตี”ติ.}

\prose{ยโต นํ สมนฺเนสมาโน เอวํ ชานาติ “เย วีติมิสฺสา จกฺขุ\-โสต\-วิญฺเญยฺยา ธมฺมา, น เต ตถาคตสฺส สํวิชฺชนฺตี”ติ, ตโต นํ อุตฺตริํ สมนฺเนสติ “เย โวทาตา จกฺขุ\-โสต\-วิญฺเญยฺยา ธมฺมา, สํวิชฺชนฺติ วา เต ตถาคตสฺส โน วา”ติ. ตเมนํ สมนฺเนสมาโน \csromanfolio{393}เอวํ ชานาติ “เย โวทาตา จกฺขุ\-โสต\-วิญฺเญยฺยา ธมฺมา, สํวิชฺชนฺติ เต ตถาคตสฺสา”ติ.}

\prose{ยโต นํ สมนฺเนสมาโน เอวํ ชานาติ “เย โวทาตา จกฺขุ\-โสต\-วิญฺเญยฺยา ธมฺมา, สํวิชฺชนฺติ เต ตถาคตสฺสา”ติ, ตโต นํ อุตฺตริํ สมนฺเนสติ “ทีฆรตฺตํ สมาปนฺโน อยมายสฺมา อิมํ กุสลํ ธมฺมํ, อุทาหุ อิตฺตร\-สมาปนฺโน”ติ. ตเมนํ สมนฺเนสมาโน เอวํ ชานาติ “ทีฆรตฺตํ สมาปนฺโน อยมายสฺมา อิมํ กุสลํ ธมฺมํ, นายมายสฺมา อิตฺตร\-สมาปนฺโน”ติ.}

\prose{ยโต นํ สมนฺเนสมาโน เอวํ ชานาติ “ทีฆรตฺตํ สมาปนฺโน อยมายสฺมา อิมํ กุสลํ ธมฺมํ, นายมายสฺมา อิตฺตร\-สมาปนฺโน”ติ, ตโต นํ อุตฺตริํ สมนฺเนสติ “ญตฺตชฺฌาปนฺโน อยมายสฺมา ภิกฺขุ ยสปฺปตฺโต, สํวิชฺชนฺต’สฺส อิเธกจฺเจ อาทีนวา”ติ. น ตาว ภิกฺขเว ภิกฺขุโน อิเธกจฺเจ อาทีนวา สํวิชฺชนฺติ, ยาว น ญตฺตชฺฌาปนฺโน โหติ ยสปฺปตฺโต. ยโต จ โข ภิกฺขเว ภิกฺขุ ญตฺตชฺฌาปนฺโน โหติ ยสปฺปตฺโต, อถ’สฺส อิเธกจฺเจ อาทีนวา สํวิชฺชนฺติ. ตเมนํ สมนฺเนสมาโน เอวํ ชานาติ “ญตฺตชฺฌาปนฺโน อยมายสฺมา ภิกฺขุ ยสปฺปตฺโต, นาสฺส อิเธกจฺเจ อาทีนวา สํวิชฺชนฺตี”ติ.}

\prose{ยโต นํ สมนฺเนสมาโน เอวํ ชานาติ “ญตฺตชฺฌาปนฺโน อยมายสฺมา ภิกฺขุ ยสปฺปตฺโต, นาสฺส อิเธกจฺเจ อาทีนวา สํวิชฺชนฺตี”ติ. ตโต นํ อุตฺตริํ สมนฺเนสติ “อภยูปรโต อยมายสฺมา, นายมายสฺมา ภยูปรโต, วีตราคตฺตา กาเม น เสวติ ขยา ราคสฺสา”ติ. ตเมนํ สมนฺเนสมาโน เอวํ ชานาติ “อภยูปรโต อยมายสฺมา, นายมายสฺมา ภยูปรโต, วีตราคตฺตา กาเม น เสวติ ขยา ราคสฺสา”ติ. ตญฺเจ ภิกฺขเว ภิกฺขุํ ปเร เอวํ ปุจฺเฉยฺยุํ “เก ปนายสฺมโต อาการา เก อนฺวยา, เยนายสฺมา เอวํ วเทสิ, อภยูปรโต อยมายสฺมา, นายมายสฺมา ภยูปรโต, วีตราคตฺตา กาเม น เสวติ ขยา ราคสฺสา”ติ. สมฺมา พฺยากรมาโน ภิกฺขเว ภิกฺขุ เอวํ พฺยากเรยฺย “ตถา หิ ปน อยมายสฺมา สํเฆ วา วิหรนฺโต เอโก วา วิหรนฺโต, เย จ ตตฺถ สุคตา, เย จ ตตฺถ ทุคฺคตา, เย จ ตตฺถ คณมนุสาสนฺติ, เย จ อิเธกจฺเจ อามิเสสุ สํทิสฺสนฺติ, เย จ อิเธกจฺเจ อามิเสน อนุปลิตฺตา, นายมายสฺมา ตํ เตน \csromanfolio{394}อวชานาติ. สมฺมุขา โข ปน เมตํ ภควโต สุตํ สมฺมุขา ปฏิคฺคหิตํ ‘อภยูปรโตหมสฺมิ, นาหมสฺมิ ภยูปรโต, วีตราคตฺตา กาเม น เสวามิ ขยา ราคสฺสา’ติ”.}

\proseitem{489}{ตตฺร ภิกฺขเว ตถาคโตว อุตฺตริํ ปฏิปุจฺฉิตพฺโพ “เย สํกิลิฏฺฐา จกฺขุ\-โสต\-วิญฺเญยฺยา ธมฺมา, สํวิชฺชนฺติ วา เต ตถาคตสฺส โน วา”ติ. พฺยากรมาโน ภิกฺขเว ตถาคโต เอวํ พฺยากเรยฺย “เย สํกิลิฏฺฐา จกฺขุ\-โสต\-วิญฺเญยฺยา ธมฺมา, น เต ตถาคตสฺส สํวิชฺชนฺตี”ติ.}

\prose{เย วีติมิสฺสา จกฺขุ\-โสต\-วิญฺเญยฺยา ธมฺมา, สํวิชฺชนฺติ วา เต ตถาคตสฺส โน วาติ. พฺยากรมาโน ภิกฺขเว ตถาคโต เอวํ พฺยากเรยฺย “เย วีติมิสฺสา จกฺขุ\-โสต\-วิญฺเญยฺยา ธมฺมา, น เต ตถาคตสฺส สํวิชฺชนฺตี”ติ.}

\prose{เย โวทาตา จกฺขุ\-โสต\-วิญฺเญยฺยา ธมฺมา, สํวิชฺชนฺติ วา เต ตถาคตสฺส โน วาติ. พฺยากรมาโน ภิกฺขเว ตถาคโต เอวํ พฺยากเรยฺย “เย โวทาตา จกฺขุ\-โสต\-วิญฺเญยฺยา ธมฺมา, สํวิชฺชนฺติ เต ตถาคตสฺส, เอตํปโถหมสฺมิ เอตํโคจโร\csromansharedfootnote{742f93ae537d85ec}{เอตปโถหมสฺมิ เอตโคจโร (สี, สฺยา, กํ, อิ)}, โน จ เตน ตมฺมโย”ติ.}

\prose{เอวํวาทิํ โข ภิกฺขเว สตฺถารํ อรหติ สาวโก อุปสงฺกมิตุํ ธมฺม\-สฺสวนาย, ตสฺส สตฺถา ธมฺมํ เทเสติ อุตฺตรุตฺตริํ ปณีตปณีตํ กณฺห\-สุกฺก\-สปฺปฏิภาคํ, ยถา ยถา โข ภิกฺขเว ภิกฺขุโน สตฺถา ธมฺมํ เทเสติ อุตฺตรุตฺตริํ ปณีตปณีตํ กณฺห\-สุกฺก\-สปฺปฏิภาคํ, ตถา ตถา โส ตสฺมิํ ธมฺเม อภิญฺญาย อิเธกจฺจํ ธมฺมํ ธมฺเมสุ นิฏฺฐํ คจฺฉติ, สตฺถริ ปสีทติ “สมฺมาสมฺพุทฺโธ ภควา, สฺวากฺขาโต ภควตา ธมฺโม, สุปฺปฏิปนฺโน สํโฆ”ติ. ตญฺเจ ภิกฺขเว ภิกฺขุํ ปเร เอวํ ปุจฺเฉยฺยุํ “เก ปนายสฺมโต อาการา เก อนฺวยา, เยนายสฺมา เอวํ วเทสิ ‘สมฺมาสมฺพุทฺโธ ภควา, สฺวากฺขาโต ภควตา ธมฺโม, สุปฺปฏิปนฺโน สํโฆ’ติ”. สมฺมา พฺยากรมาโน ภิกฺขเว ภิกฺขุ เอวํ พฺยากเรยฺย “อิธาหํ อาวุโส เยน ภควา เตนุปสงฺกมิํ ธมฺม\-สฺสวนาย. ตสฺส เม ภควา ธมฺมํ เทเสติ อุตฺตรุตฺตริํ ปณีตปณีตํ กณฺห\-สุกฺก\-สปฺปฏิภาคํ. ยถา ยถา เม \csromanfolio{395}อาวุโส ภควา ธมฺมํ เทเสติ อุตฺตรุตฺตริํ ปณีตปณีตํ กณฺห\-สุกฺก\-สปฺปฏิภาคํ, ตถา ตถาหํ ตสฺมิํ ธมฺเม อภิญฺญาย อิเธกจฺจํ ธมฺมํ ธมฺเมสุ นิฏฺฐมคมํ, สตฺถริ ปสีทิํ ‘สมฺมาสมฺพุทฺโธ ภควา, สฺวากฺขาโต ภควตา ธมฺโม, สุปฺปฏิปนฺโน สํโฆ’ติ”.}

\proseitem{490}{ยสฺส กสฺสจิ ภิกฺขเว อิเมหิ อากาเรหิ อิเมหิ ปเทหิ อิเมหิ พฺยญฺชเนหิ ตถาคเต สทฺธา นิวิฏฺฐา โหติ, มูลชาตา ปติฏฺฐิตา. อยํ วุจฺจติ ภิกฺขเว อาการวตี สทฺธา ทสฺสนมูลิกา ทฬฺหา อสํหาริยา สมเณน วา พฺราหฺมเณน วา เทเวน วา มาเรน วา พฺรหฺมุนา วา เกนจิ วา โลกสฺมิํ. เอวํ โข ภิกฺขเว ตถาคเต ธมฺม\-สมนฺเนสนา โหติ. เอวญฺจ ปน ตถาคโต ธมฺมตา\-สุสมนฺนิฏฺโฐ โหตีติ.}

\prosewithcloserruled{อิทมโวจ ภควา. อตฺตมนา เต ภิกฺขู ภควโต ภาสิตํ อภินนฺทุนฺติ.}{วีมํสก\-สุตฺตํ นิฏฺฐิตํ สตฺตมํ.}
\csromanmatikaanchor{mtk.82}
\csromantocmarkref{subsection}{8. โกสมฺพิยสุตฺต}{mtk.82}
\bookmark[dest={mtk.82},level=2]{8. โกสมฺพิยสุตฺต}
\csromanheader{2}{8. โกสมฺพิยสุตฺต}
\proseitem{491}{เอวํ เม สุตํ\csromandash{}เอกํ สมยํ ภควา โกสมฺพิยํ วิหรติ โฆสิตาราเม. เตน โข ปน สมเยน โกสมฺพิยํ ภิกฺขู ภณฺฑนชาตา กลหชาตา วิวาทาปนฺนา อญฺญมญฺญํ มุขสตฺตีหิ วิตุทนฺตา วิหรนฺติ. เต น เจว อญฺญมญฺญํ สญฺญาเปนฺติ, น จ สญฺญตฺติํ อุเปนฺติ, น จ อญฺญมญฺญํ นิชฺฌาเปนฺติ, น จ นิชฺฌตฺติํ อุเปนฺติ. อถ โข อญฺญตโร ภิกฺขุ เยน ภควา เตนุปสงฺกมิ, อุปสงฺกมิตฺวา ภควนฺตํ อภิวาเทตฺวา เอกมนฺตํ นิสีทิ, เอกมนฺตํ นิสินฺโน โข โส ภิกฺขุ ภควนฺตํ เอตทโวจ “อิธ ภนฺเต โกสมฺพิยํ ภิกฺขู ภณฺฑนชาตา กลหชาตา วิวาทาปนฺนา อญฺญมญฺญํ มุขสตฺตีหิ วิตุทนฺตา วิหรนฺติ, เต น เจว อญฺญมญฺญํ สญฺญาเปนฺติ, น จ สญฺญตฺติํ อุเปนฺติ, น จ อญฺญมญฺญํ นิชฺฌาเปนฺติ, น จ นิชฺฌตฺติํ อุเปนฺตี”ติ.}

\prose{อถ โข ภควา อญฺญตรํ ภิกฺขุํ อามนฺเตสิ “เอหิ ตฺวํ ภิกฺขุ มม วจเนน เต ภิกฺขู อามนฺเตหิ, สตฺถา โว อายสฺมนฺเต อามนฺเตตี”ติ. \csromanfolio{396}“เอวํ ภนฺเต”ติ โข โส ภิกฺขุ ภควโต ปฏิสฺสุตฺวา เยน เต ภิกฺขู เตนุปสงฺกมิ, อุปสงฺกมิตฺวา เต ภิกฺขู เอตทโวจ “สตฺถา อายสฺมนฺเต อามนฺเตตี”ติ. “เอวมาวุโส”ติ โข เต ภิกฺขู ตสฺส ภิกฺขุโน ปฏิสฺสุตฺวา เยน ภควา เตนุ\-ปสงฺกมิํสุ, อุปสงฺกมิตฺวา ภควนฺตํ อภิวาเทตฺวา เอกมนฺตํ นิสีทิํสุ, เอกมนฺตํ นิสินฺเน โข เต ภิกฺขู ภควา เอตทโวจ “สจฺจํ กิร ตุเมฺห ภิกฺขเว ภณฺฑนชาตา กลหชาตา วิวาทาปนฺนา อญฺญมญฺญํ มุขสตฺตีหิ วิตุทนฺตา วิหรถ, เตน เจว อญฺญมญฺญํ สญฺญาเปถ, น จ สญฺญตฺติํ อุเปถ, น จ อญฺญมญฺญํ นิชฺฌาเปถ, น จ นิชฺฌตฺติํ อุเปถา”ติ. เอวํ ภนฺเต. ตํ กิํ มญฺญถ ภิกฺขเว, ยสฺมิํ ตุเมฺห สมเย ภณฺฑนชาตา กลหชาตา วิวาทาปนฺนา อญฺญมญฺญํ มุขสตฺตีหิ วิตุทนฺตา วิหรถ, อปิ นุ ตุมฺหากํ ตสฺมิํ สมเย เมตฺตํ กายกมฺมํ ปจฺจุปฏฺฐิตํ โหติ สพฺรหฺมจารีสุ อาวิ เจว รโห จ. เมตฺตํ วจีกมฺมํ ฯเปฯ เมตฺตํ มโนกมฺมํ ปจฺจุปฏฺฐิตํ โหติ สพฺรหฺมจารีสุ อาวิ เจว รโห จาติ. โน เหตํ ภนฺเต. อิติ กิร ภิกฺขเว ยสฺมิํ ตุเมฺห สมเย ภณฺฑนชาตา กลหชาตา วิวาทาปนฺนา อญฺญมญฺญํ มุขสตฺตีหิ วิตุทนฺตา วิหรถ, เนว ตุมฺหากํ ตสฺมิํ สมเย เมตฺตํ กายกมฺมํ ปจฺจุปฏฺฐิตํ โหติ สพฺรหฺมจารีสุ อาวิ เจว รโห จ. น เมตฺตํ วจีกมฺมํ ฯเปฯ น เมตฺตํ มโนกมฺมํ ปจฺจุปฏฺฐิตํ โหติ สพฺรหฺมจารีสุ อาวิ เจว รโห จ. อถ กิญฺจรหิ ตุเมฺห โมฆปุริสา กิํ ชานนฺตา กิํ ปสฺสนฺตา ภณฺฑนชาตา กลหชาตา วิวาทาปนฺนา อญฺญมญฺญํ มุขสตฺตีหิ วิตุทนฺตา วิหรถ. เต น เจว อญฺญมญฺญํ สญฺญาเปถ, น จ สญฺญตฺติํ อุเปถ, น จ อญฺญมญฺญํ นิชฺฌาเปถ, น จ นิชฺฌตฺติํ อุเปถ. ตํ หิ ตุมฺหากํ โมฆปุริสา ภวิสฺสติ ทีฆรตฺตํ อหิตาย ทุกฺขายาติ.}

\proseitem{492}{อถ โข ภควา ภิกฺขู อามนฺเตสิ, ฉยิเม ภิกฺขเว ธมฺมา สารณียา ปิยกรณา ครุกรณา สงฺคหาย อวิวาทาย สามคฺคิยา เอกีภาวาย สํวตฺตนฺติ. กตเม ฉ, อิธ ภิกฺขเว ภิกฺขุโน เมตฺตํ กายกมฺมํ ปจฺจุปฏฺฐิตํ โหติ สพฺรหฺมจารีสุ อาวิ เจว รโห จ. อยมฺปิ ธมฺโม สารณีโย ปิยกรโณ ครุกรโณ สงฺคหาย อวิวาทาย สามคฺคิยา เอกีภาวาย สํวตฺตติ. (1)}

\prose{ปุน จปรํ ภิกฺขเว ภิกฺขุโน เมตฺตํ วจีกมฺมํ ปจฺจุปฏฺฐิตํ โหติ สพฺรหฺมจารีสุ อาวิ เจว รโห จ. อยมฺปิ ธมฺโม สารณีโย \csromanfolio{397}ปิยกรโณ ครุกรโณ สงฺคหาย อวิวาทาย สามคฺคิยา เอกีภาวาย สํวตฺตติ. (2)}

\prose{ปุน จปรํ ภิกฺขเว ภิกฺขุโน เมตฺตํ มโนกมฺมํ ปจฺจุปฏฺฐิตํ โหติ สพฺรหฺมจารีสุ อาวิ เจว รโห จ. อยมฺปิ ธมฺโม สารณีโย ปิยกรโณ ครุกรโณ สงฺคหาย อวิวาทาย สามคฺคิยา เอกีภาวาย สํวตฺตติ. (3)}

\prose{ปุน จปรํ ภิกฺขเว ภิกฺขุ เย เต ลาภา ธมฺมิกา ธมฺมลทฺธา, อนฺตมโส ปตฺต\-ปริยาปนฺน\-มตฺตํปิ, ตถารูเปหิ ลาเภหิ อปฺปฏิวิภตฺต\-โภคี โหติ สีลวนฺเตหิ สพฺรหฺมจารีหิ สาธารณโภคี. อยมฺปิ ธมฺโม สารณีโย ปิยกรโณ ครุกรโณ สงฺคหาย อวิวาทาย สามคฺคิยา เอกีภาวาย สํวตฺตติ. (4)}

\prose{ปุน จปรํ ภิกฺขเว ภิกฺขุ ยานิ ตานิ สีลานิ อขณฺฑานิ อจฺฉิทฺทานิ อสพลานิ อกมฺมาสานิ ภุชิสฺสานิ วิญฺญุ\-ปฺปสตฺถานิ อปรามฏฺฐานิ สมาธิ\-สํวตฺตนิกานิ, ตถารูเปสุ สีเลสุ สีลสามญฺญคโต วิหรติ สพฺรหฺมจารีหิ อาวิ เจว รโห จ. อยมฺปิ ธมฺโม สารณีโย ปิยกรโณ ครุกรโณ สงฺคหาย อวิวาทาย สามคฺคิยา เอกีภาวาย สํวตฺตติ. (5)}

\prose{ปุน จปรํ ภิกฺขเว ภิกฺขุ ยายํ ทิฏฺฐิ อริยา นิยฺยานิกา นิยฺยาติตกฺกรสฺส สมฺมา ทุกฺขกฺขยาย, ตถารูปาย ทิฏฺฐิยา ทิฏฺฐิ\-สามญฺญคโต วิหรติ สพฺรหฺมจารีหิ อาวิ เจว รโห จ. อยมฺปิ ธมฺโม สารณีโย ปิยกรโณ ครุกรโณ สงฺคหาย อวิวาทาย สามคฺคิยา เอกีภาวาย สํวตฺตติ. (6)}

\prose{อิเม โข ภิกฺขเว ฉ สารณียา ธมฺมา ปิยกรณา ครุกรณา สงฺคหาย อวิวาทาย สามคฺคิยา เอกีภาวาย สํวตฺตนฺติ. อิเมสํ โข ภิกฺขเว ฉนฺนํ สารณียานํ ธมฺมานํ เอตํ อคฺคํ เอตํ สงฺคาหิกํ\csromansharedfootnote{745296188da6ee10}{สงฺคาหกํ (?)} เอตํ สงฺฆาฏนิกํ ยทิทํ ยายํ ทิฏฺฐิ อริยา นิยฺยานิกา นิยฺยาติ ตกฺกรสฺส สมฺมา ทุกฺขกฺขยาย. เสยฺยถาปิ ภิกฺขเว กูฏาคารสฺส เอตํ อคฺคํ เอตํ สงฺคาหิกํ เอตํ สงฺฆาฏนิกํ ยทิทํ กูฏํ. เอวเมว โข ภิกฺขเว อิเมสํ ฉนฺนํ สารณียานํ ธมฺมานํ เอตํ อคฺคํ เอตํ สงฺคาหิกํ เอตํ \csromanfolio{398}สงฺฆาฏนิกํ ยทิทํ ยายํ ทิฏฺฐิ อริยา นิยฺยานิกา นิยฺยาติ ตกฺกรสฺส สมฺมา ทุกฺขกฺขยาย.}

\proseitem{493}{กถญฺจ ภิกฺขเว ยายํ ทิฏฺฐิ อริยา นิยฺยานิกา นิยฺยาติ ตกฺกรสฺส สมฺมา ทุกฺขกฺขยาย. อิธ ภิกฺขเว ภิกฺขุ อรญฺญคโต วา รุกฺขมูลคโต วา สุญฺญาคารคโต วา อิติ ปฏิสญฺจิกฺขติ “อตฺถิ นุ โข เม ตํ ปริยุฏฺฐานํ อชฺฌตฺตํ อปฺปหีนํ, เยนาหํ ปริยุฏฺฐาเนน ปริยุฏฺฐิต\-จิตฺโต ยถาภูตํ นปฺปชาเนยฺยํ น ปสฺเสยฺยนฺ”ติ. สเจ ภิกฺขเว ภิกฺขุ กาม\-ราค\-ปริยุฏฺฐิโต โหติ, ปริยุฏฺฐิต\-จิตฺโตว โหติ. สเจ ภิกฺขเว ภิกฺขุ พฺยาปาท\-ปริยุฏฺฐิโต โหติ, ปริยุฏฺฐิต\-จิตฺโตว โหติ. สเจ ภิกฺขเว ภิกฺขุ ถินมิทฺธ\-ปริยุฏฺฐิโต โหติ, ปริยุฏฺฐิต\-จิตฺโตว โหติ. สเจ ภิกฺขเว ภิกฺขุ อุทฺธจฺจ\-กุกฺกุจฺจ\-ปริยุฏฺฐิโต โหติ, ปริยุฏฺฐิต\-จิตฺโตว โหติ. สเจ ภิกฺขเว ภิกฺขุ วิจิกิจฺฉา\-ปริยุฏฺฐิโต โหติ, ปริยุฏฺฐิต\-จิตฺโตว โหติ. สเจ ภิกฺขเว ภิกฺขุ อิธโลก\-จินฺตาย ปสุโต โหติ, ปริยุฏฺฐิต\-จิตฺโตว โหติ. สเจ ภิกฺขเว ภิกฺขุ ปรโลก\-จินฺตาย ปสุโต โหติ, ปริยุฏฺฐิต\-จิตฺโตว โหติ. สเจ ภิกฺขเว ภิกฺขุ ภณฺฑนชาโต กลหชาโต วิวาทาปนฺโน อญฺญมญฺญํ มุขสตฺตีหิ วิตุทนฺโต วิหรติ, ปริยุฏฺฐิต\-จิตฺโตว โหติ. โส เอวํ ปชานาติ “นตฺถิ โข เม ตํ ปริยุฏฺฐานํ อชฺฌตฺตํ อปฺปหีนํ, เยนาหํ ปริยุฏฺฐาเนน ปริยุฏฺฐิต\-จิตฺโต ยถาภูตํ นปฺปชาเนยฺยํ น ปสฺเสยฺยํ, สุปฺปณิหิตํ เม มานสํ สจฺจานํ โพธายา”ติ. อิทมสฺส ปฐมํ ญาณํ อธิคตํ โหติ อริยํ โลกุตฺตรํ อสาธารณํ ปุถุชฺชเนหิ.}

\proseitem{494}{ปุน จปรํ ภิกฺขเว อริยสาวโก อิติ ปฏิสญฺจิกฺขติ “อิมํ นุ โข อหํ ทิฏฺฐิํ อาเสวนฺโต ภาเวนฺโต พหุลีกโรนฺโต ลภามิ ปจฺจตฺตํ สมถํ, ลภามิ ปจฺจตฺตํ นิพฺพุตินฺ”ติ. โส เอวํ ปชานาติ “อิมํ โข อหํ ทิฏฺฐิํ อาเสวนฺโต ภาเวนฺโต พหุลีกโรนฺโต ลภามิ ปจฺจตฺตํ สมถํ, ลภามิ ปจฺจตฺตํ นิพฺพุตินฺ”ติ. อิทมสฺส ทุติยํ ญาณํ อธิคตํ โหติ อริยํ โลกุตฺตรํ อสาธารณํ ปุถุชฺชเนหิ.}

\proseitem{495}{ปุน จปรํ ภิกฺขเว อริยสาวโก อิติ ปฏิสญฺจิกฺขติ “ยถารูปายาหํ ทิฏฺฐิยา สมนฺนาคโต, อตฺถิ นุ โข อิโต พหิทฺธา อญฺโญ \csromanfolio{399}สมโณ วา ภาหฺมโณ วา ตถารูปาย ทิฏฺฐิยา สมนฺนาคโต”ติ. โส เอวํ ปชานาติ “ยถารูปายาหํ ทิฏฺฐิยา สมนฺนาคโต, นตฺถิ อิโต พหิทฺธา อญฺโญ สมโณ วา พฺราหฺมโณ วา ตถารูปาย ทิฏฺฐิยา สมนฺนาคโต”ติ. อิทมสฺส ตติยํ ญาณํ อธิคตํ โหติ อริยํ โลกุตฺตรํ อสาธารณํ ปุถุชฺชเนหิ.}

\proseitem{496}{ปุน จปรํ ภิกฺขเว อริยสาวโก อิติ ปฏิสญฺจิกฺขติ “ยถารูปาย ธมฺมตาย ทิฏฺฐิสมฺปนฺโน ปุคฺคโล สมนฺนาคโต, อหมฺปิ ตถารูปาย ธมฺมตาย สมนฺนาคโต”ติ. กถํรูปาย จ ภิกฺขเว ธมฺมตาย ทิฏฺฐิสมฺปนฺโน ปุคฺคโล สมนฺนาคโต. ธมฺมตา เอสา ภิกฺขเว ทิฏฺฐิ\-สมฺปนฺนสฺส ปุคฺคลสฺส, กิญฺจาปิ ตถารูปิํ อาปตฺติํ อาปชฺชติ, ยถารูปาย อาปตฺติยา วุฏฺฐานํ ปญฺญายติ, อถ โข นํ ขิปฺปเมว สตฺถริ วา วิญฺญูสุ วา สพฺรหฺมจารีสุ เทเสติ วิวรติ อุตฺตานีกโรติ, เทเสตฺวา วิวริตฺวา อุตฺตานีกตฺวา อายติํ สํวรํ อาปชฺชติ. เสยฺยถาปิ ภิกฺขเว ทหโร กุมาโร มนฺโท อุตฺตานเสยฺยโก หตฺเถน วา ปาเทน วา องฺคารํ อกฺกมิตฺวา ขิปฺปเมว ปฏิสํหรติ. เอวเมว โข ภิกฺขเว ธมฺมตา เอสา ทิฏฺฐิ\-สมฺปนฺนสฺส ปุคฺคลสฺส, กิญฺจาปิ ตถารูปิํ อาปตฺติํ อาปชฺชติ, ยถารูปาย อาปตฺติยา วุฏฺฐานํ ปญฺญายติ, อถ โข นํ ขิปฺปเมว สตฺถริ วา วิญฺญูสุ วา สพฺรหฺมจารีสุ เทเสติ วิวรติ อุตฺตานีกโรติ, เทเสตฺวา วิวริตฺวา อุตฺตานีกตฺวา อายติํ สํวรํ อาปชฺชติ. โส เอวํ ปชานาติ “ยถารูปาย ธมฺมตาย ทิฏฺฐิสมฺปนฺโน ปุคฺคโล สมนฺนาคโต, อหมฺปิ ตถารูปาย ธมฺมตาย สมนฺนาคโต”ติ. อิทมสฺส จตุตฺถํ ญาณํ อธิคตํ โหติ อริยํ โลกุตฺตรํ อสาธารณํ ปุถุชฺชเนหิ.}

\proseitem{497}{ปุน จปรํ ภิกฺขเว อริยสาวโก อิติ ปฏิสญฺจิกฺขติ “ยถารูปาย ธมฺมตาย ทิฏฺฐิสมฺปนฺโน ปุคฺคโล สมนฺนาคโต, อหมฺปิ ตถารูปาย ธมฺมตาย สมนฺนาคโต”ติ. กถํรูปาย จ ภิกฺขเว ธมฺมตาย ทิฏฺฐิสมฺปนฺโน ปุคฺคโล สมนฺนาคโต. ธมฺมตา เอสา ภิกฺขเว ทิฏฺฐิ\-สมฺปนฺนสฺส ปุคฺคลสฺส, กิญฺจาปิ ยานิ ตานิ สพฺรหฺมจารีนํ อุจฺจาวจานิ กิํกรณียานิ, ตตฺถ อุสฺสุกฺกํ อาปนฺโน โหติ, อถ ขฺวาสฺส ติพฺพาเปกฺขา โหติ อธิสีล\-สิกฺขาย อธิจิตฺต\-สิกฺขาย อธิปญฺญา\-สิกฺขาย. เสยฺยถาปิ ภิกฺขเว คาวี ตรุณวจฺฉา ถมฺพญฺจ อาลุมฺปติ, วจฺฉกญฺจ อปจินติ. เอวเมว โข \csromanfolio{400}ภิกฺขเว ธมฺมตา เอสา ทิฏฺฐิ\-สมฺปนฺนสฺส ปุคฺคลสฺส, กิญฺจาปิ ยานิ ตานิ สพฺรหฺมจารีนํ อุจฺจาวจานิ กิํกรณียานิ, ตตฺถ อุสฺสุกฺกํ อาปนฺโน โหติ, อถ ขฺวาสฺส ติพฺพาเปกฺขา โหติ อธิสีล\-สิกฺขาย อธิจิตฺต\-สิกฺขาย อธิปญฺญา\-สิกฺขาย. โส เอวํ ปชานาติ “ยถารูปาย ธมฺมตาย ทิฏฺฐิสมฺปนฺโน ปุคฺคโล สมนฺนาคโต, อหมฺปิ ตถารูปาย ธมฺมตาย สมนฺนาคโต”ติ. อิทมสฺส ปญฺจมํ ญาณํ อธิคตํ โหติ อริยํ โลกุตฺตรํ อสาธารณํ ปุถุชฺชเนหิ.}

\proseitem{498}{ปุน จปรํ ภิกฺขเว อริยสาวโก อิติ ปฏิสญฺจิกฺขติ “ยถารูปาย พลตาย ทิฏฺฐิสมฺปนฺโน ปุคฺคโล สมนฺนาคโต, อหมฺปิ ตถารูปาย พลตาย สมนฺนาคโต”ติ. กถํรูปาย จ ภิกฺขเว พลตาย ทิฏฺฐิสมฺปนฺโน ปุคฺคโล สมนฺนาคโต. พลตา เอสา ภิกฺขเว ทิฏฺฐิ\-สมฺปนฺนสฺส ปุคฺคลสฺส, ยํ ตถาคต\-ปฺปเวทิเต ธมฺมวินเย เทสิยมาเน อฏฺฐิํกตฺวา มนสิกตฺวา สพฺพเจตสา\csromansharedfootnote{ccde576fb7f2d630}{สพฺพเจตโส (สี, สฺยา, กํ, อิ) สพฺพํ เจตสา (ก)} สมนฺนาหริตฺวา โอหิตโสโต ธมฺมํ สุณาติ. โส เอวํ ปชานาติ “ยถารูปาย พลตาย ทิฏฺฐิสมฺปนฺโน ปุคฺคโล สมนฺนาคโต, อหมฺปิ ตถารูปาย พลตาย สมนฺนาคโต”ติ. อิทมสฺส ฉฏฺฐํ ญาณํ อธิคตํ โหติ อริยํ โลกุตฺตรํ อสาธารณํ ปุถุชฺชเนหิ.}

\proseitem{499}{ปุน จปรํ ภิกฺขเว อริยสาวโก อิติ ปฏิสญฺจิกฺขติ “ยถารูปาย พลตาย ทิฏฺฐิสมฺปนฺโน ปุคฺคโล สมนฺนาคโต, อหมฺปิ ตถารูปาย พลตาย สมนฺนาคโต”ติ. กถํรูปาย จ ภิกฺขเว พลตาย ทิฏฺฐิสมฺปนฺโน ปุคฺคโล สมนฺนาคโต. พลตา เอสา ภิกฺขเว ทิฏฺฐิ\-สมฺปนฺนสฺส ปุคฺคลสฺส ยํ ตถาคต\-ปฺปเวทิเต ธมฺมวินเย เทสิยมาเน ลภติ อตฺถเวทํ, ลภติ ธมฺมเวทํ, ลภติ ธมฺมู\-ปสํหิตํ ปาโมชฺชํ. โส เอวํ ปชานาติ “ยถารูปาย พลตาย ทิฏฺฐิสมฺปนฺโน ปุคฺคโล สมนฺนาคโต, อหมฺปิ ตถารูปาย พลตาย สมนฺนาคโต”ติ. อิทมสฺส สตฺตมํ ญาณํ อธิคตํ โหติ อริยํ โลกุตฺตรํ อสาธารณํ ปุถุชฺชเนหิ.}

\proseitem{500}{เอวํ สตฺต\-งฺค\-สมนฺนาคตสฺส โข ภิกฺขเว อริยสาวกสฺส ธมฺมตา สุสมนฺนิฏฺฐา โหติ โสตาปตฺติผล\-สจฺฉิกิริยาย. เอวํ \csromanfolio{401}สตฺต\-งฺค\-สมนฺนาคโต โข ภิกฺขเว อริยสาวโก โสตาปตฺติผล\-สมนฺนาคโต โหตีติ.}

\prosewithcloserruled{อิทมโวจ ภควา. อตฺตมนา เต ภิกฺขู ภควโต ภาสิตํ อภินนฺทุนฺติ.}{โกสมฺพิย\-สุตฺตํ นิฏฺฐิตํ อฏฺฐมํ.}
\csromanmatikaanchor{mtk.83}
\csromantocmarkref{subsection}{9. พฺรหฺมนิมนฺตนิกสุตฺต}{mtk.83}
\bookmark[dest={mtk.83},level=2]{9. พฺรหฺมนิมนฺตนิกสุตฺต}
\csromanheader{2}{9. พฺรหฺม\-นิมนฺตนิก\-สุตฺต}
\proseitem{501}{เอวํ เม สุตํ\csromandash{}เอกํ สมยํ ภควา สาวตฺถิยํ วิหรติ เชตวเน อนาถ\-ปิณฺฑิกสฺส อาราเม. ตตฺร โข ภควา ภิกฺขู อามนฺเตสิ “ภิกฺขโว”ติ. “ภทนฺเต”ติ เต ภิกฺขู ภควโต ปจฺจสฺโสสุํ. ภควา เอตทโวจ\csromandash{}}

\prose{เอกมิทาหํ ภิกฺขเว สมยํ อุกฺกฏฺฐายํ วิหรามิ สุภควเน สาลราชมูเล. เตน โข ปน ภิกฺขเว สมเยน พกสฺส พฺรหฺมุโน เอวรูปํ ปาปกํ ทิฏฺฐิคตํ อุปฺปนฺนํ โหติ “อิทํ นิจฺจํ อิทํ ธุวํ อิทํ สสฺสตํ อิทํ เกวลํ อิทํ อจวนธมฺมํ, อิทํ หิ น ชายติ น ชียติ น มียติ น จวติ น อุปปชฺชติ, อิโต จ ปนญฺญํ อุตฺตริ นิสฺสรณํ นตฺถี”ติ. อถ ขฺวาหํ ภิกฺขเว พกสฺส พฺรหฺมุโน เจตสา เจโต\-ปริวิตกฺกม\-ญฺญาย เสยฺยถาปิ นาม พลวา ปุริโส สมิญฺชิตํ วา พาหํ ปสาเรยฺย ปสาริตํ วา พาหํ สมิญฺเชยฺย, เอวเมว อุกฺกฏฺฐายํ สุภควเน สาลราชมูเล อนฺตรหิโต ตสฺมิํ พฺรหฺมโลเก ปาตุรโหสิํ. อทฺทสา โข มํ ภิกฺขเว พโก พฺรหฺมา ทูรโตว อาคจฺฉนฺตํ, ทิสฺวาน มํ เอตทโวจ “เอหิ โข มาริส, สฺวาคตํ มาริส, จิรสฺสํ โข มาริส อิมํ ปริยายมกาสิ ยทิทํ อิธาคมนาย. อิทํ หิ มาริส นิจฺจํ อิทํ ธุวํ อิทํ สสฺสตํ อิทํ เกวลํ อิทํ อจวนธมฺมํ, อิทํ หิ น ชายติ น ชียติ น มียติ น จวติ น อุปปชฺชติ, อิโต จ ปนญฺญํ อุตฺตริ นิสฺสรณํ นตฺถี”ติ.}

\prose{เอวํ วุตฺเต อหํ ภิกฺขเว พกํ พฺรหฺมานํ เอตทโวจํ “อวิชฺชาคโต วต โภ พโก พฺรหฺมา, อวิชฺชาคโต วต โภ พโก พฺรหฺมา. ยตฺร หิ นาม อนิจฺจํเยว สมานํ ‘นิจฺจนฺ’ติ วกฺขติ, อทฺธุวํเยว สมานํ ‘ธุวนฺ’ติ วกฺขติ, อสสฺสตํเยว สมานํ ‘สสฺสตนฺ’ติ วกฺขติ, อเกวลํเยว \csromanfolio{402}สมานํ ‘เกวลนฺ’ติ วกฺขติ, จวน\-ธมฺมํเยว สมานํ ‘อจวนธมฺมนฺ’ติ วกฺขติ. ยตฺถ จ ปน ชายติ ชียติ มียติ จวติ อุปปชฺชติ, ตญฺจ วกฺขติ ‘อิทํ หิ น ชายติ น ชียติ น มียติ น จวติ น อุปปชฺชตี’ติ, สนฺตญฺจ ปนญฺญํ อุตฺตริ นิสฺสรณํ ‘นตฺถญฺญํ อุตฺตริ นิสฺสรณนฺ’ติ วกฺขตี”ติ.}

\proseitem{502}{อถ โข ภิกฺขเว มาโร ปาปิมา อญฺญตรํ พฺรหฺม\-ปาริสชฺชํ อนฺวาวิสิตฺวา มํ เอตทโวจ “ภิกฺขุ ภิกฺขุ เมตมาสโท, เมตมาสโท. เอโส หิ ภิกฺขุ พฺรหฺมา มหาพฺรหฺมา อภิภู อนภิภูโต อญฺญทตฺถุทโส วสวตฺตี อิสฺสโร กตฺตา นิมฺมาตา เสฏฺโฐ สชิตา\csromansharedfootnote{5fe431da91f42d67}{สชฺชิตา (สฺยา, กํ, ก), สญฺชิตา (สี, อิ)} วสี ปิตา ภูตภพฺยานํ. อเหสุํ โข เย ภิกฺขุ ตยา ปุพฺเพ สมณพฺราหฺมโณ โลกสฺมิํ ปถวี\-ครหกา ปถวี\-ชิคุจฺฉกา อาปครหกา อาปชิคุจฺฉกา เตชครหกา เตชชิคุจฺฉกา วายครหกา วายชิคุจฺฉกา ภูตครหกา ภูตชิคุจฺฉกา เทวครหกา เทวชิคุจฺฉกา ปชาปติ\-ครหกา ปชาปติ\-ชิคุจฺฉกา พฺรหฺมครหกา พฺรหฺม\-ชิคุจฺฉกา, เต กายสฺส เภทา ปาณุปจฺเฉทา หีเน กาเย ปติฏฺฐิตา. อเหสุํ เย ปน ภิกฺขุ ตยา ปุพฺเพ สมณพฺราหฺมณา โลกสฺมิํ ปถวี\-ปสํสกา ปถวา\-ภินนฺทิโน อาปปสํสกา อาปาภินนฺทิโน เตชปสํสกา เตชาภินนฺทิโน วายปสํสกา วายาภินนฺทิโน ภูตปสํสกา ภูตา\-ภินนฺทิโน เทวปสํสกา เทวาภินนฺทิโน ปชาปติ\-ปสํสกา ปชาปตา\-ภินนฺทิโน พฺรหฺม\-ปสํสกา พฺรหฺมา\-ภินนฺทิโน, เต กายสฺส เภทา ปาณุปจฺเฉทา ปณีเต กาเย ปติฏฺฐิตา. ตํ ตาหํ ภิกฺขุ เอวํ วทามิ\csromandash{}อิงฺฆ ตฺวํ มาริส ยเทว เต พฺรหฺมา อาห, ตเทว ตฺวํ กโรหิ. มา ตฺวํ พฺรหฺมุโน วจนํ อุปาติวตฺติตฺโถ. สเจ โข ตฺวํ ภิกฺขุ พฺรหฺมุโน วจนํ อุปาติวตฺติสฺสสิ. เสยฺยถาปิ นาม ปุริโส สิริํ อาคจฺฉนฺติํ ทณฺเฑน ปฏิปฺปณาเมยฺย, เสยฺยถาปิ วา ปน ภิกฺขุ ปุริโส นรกปฺปปาเต ปปตนฺโต หตฺเถหิ จ ปาเทหิ จ ปถวิํ วิราเธยฺย, เอวํ สมฺปทมิทํ ภิกฺขุ ตุยฺหํ ภวิสฺสติ. อิงฺฆ ตฺวํ มาริส ยเทว เต พฺรหฺมา อาห, ตเทว ตฺวํ กโรหิ. มา ตฺวํ พฺรหฺมุโน วจนํ อุปาติวตฺติตฺโถ. นนุ ตฺวํ ภิกฺขุ ปสฺสสิ พฺรหฺมปริสํ สนฺนิปติตนฺ”ติ. อิติ โข มํ ภิกฺขเว มาโร ปาปิมา พฺรหฺมปริสํ อุปเนสิ.}

\prose{\csromanfolio{403}เอวํ วุตฺเต อหํ ภิกฺขเว มารํ ปาปิมนฺตํ เอตทโวจํ “ชานามิ โข ตาหํ ปาปิม, มา ตฺวํ มญฺญิตฺโถ ‘น มํ ชานาตี’ติ, มาโร ตฺวมสิ ปาปิม. โย เจว ปาปิม พฺรหฺมา ยา จ พฺรหฺมปริสา เย จ พฺรหฺม\-ปาริสชฺชา, สพฺเพว ตว หตฺถคตา สพฺเพว ตว วสํคตา, ตุยฺหํ หิ ปาปิม เอวํ โหติ ‘เอโสปิ เม อสฺส หตฺถคโต, เอโสปิ เม อสฺส วสํคโต’ติ, อหํ โข ปน ปาปิม เนว ตว หตฺถคโต, เนว ตว วสํคโต”ติ.}

\proseitem{503}{เอวํ วุตฺเต ภิกฺขเว พโก พฺรหฺมา มํ เอตทโวจ “อหํ หิ มาริส นิจฺจํเยว สมานํ ‘นิจฺจนฺ’ติ วทามิ, ธุวํเยว สมานํ ‘ธุวนฺ’ติ วทามิ, สสฺสตํเยว สมานํ ‘สสฺสตนฺ’ติ วทามิ, เกวลํเยว สมานํ ‘เกวลนฺ’ติ วทามิ, อจวน\-ธมฺมํเยว สมานํ ‘อจวนธมฺมนฺ’ติ วทามิ. ยตฺถ จ ปน น ชายติ น ชียติ น มียติ น จวติ น อุปปชฺชติ, ตเทวาหํ วทามิ ‘อิทํ หิ น ชายติ น ชียติ น มียติ น จวติ น อุปปชฺชตี’ติ, อสนฺตญฺจ ปนญฺญํ อุตฺตริ นิสฺสรณํ ‘นตฺถญฺญํ อุตฺตริ นิสฺสรณนฺ’ติ วทามิ. อเหสุํ โข ภิกฺขุ ตยา ปุพฺเพ สมณพฺราหฺมณา โลกสฺมิํ ยาวตกํ ตุยฺหํ กสิณํ อายุ, ตาวตกํ เตสํ ตโปกมฺมเมว อโหสิ. เต โข เอวํ ชาเนยฺยุํ สนฺตญฺจ ปนญฺญํ อุตฺตริ นิสฺสรณํ ‘อตฺถญฺญํ อุตฺตริ นิสฺสรณนฺ’ติ, อสนฺตํ วา อญฺญํ อุตฺตริ นิสฺสรณํ ‘นตฺถญฺญํ อุตฺตริ นิสฺสรณนฺ’ติ. ตํ ตาหํ ภิกฺขุ เอวํ วทามิ\csromandash{}น เจวญฺญํ อุตฺตริ นิสฺสรณํ ทกฺขิสฺสสิ, ยาวเทว จ ปน กิลมถสฺส วิฆาตสฺส ภาคี ภวิสฺสสิ. สเจ โข ตฺวํ ภิกฺขุ ปถวิํ อชฺโฌสิสฺสสิ, โอปสายิโก เม ภวิสฺสสิ วตฺถุสายิโก ยถา\-กาม\-กรณีโย พาหิเตยฺโย. สเจ อาปํ. เตชํ. วายํ. ภูเต. เทเว. ปชาปติํ. พฺรหฺมํ อชฺโฌสิสฺสสิ, โอปสายิโก เม ภวิสฺสสิ วตฺถุสายิโก ยถา\-กาม\-กรณีโย พาหิเตยฺโย”ติ.}

\prose{อหมฺปิ โข เอวํ พฺรหฺเม ชานามิ “สเจ ปถวิํ อชฺโฌสิสฺสามิ, โอปสายิโก เต ภวิสฺสามิ วตฺถุสายิโก ยถา\-กาม\-กรณีโย พาหิเตยฺโย. สเจ อาปํ. เตชํ. วายํ. ภูเต. เทเว. ปชาปติํ. พฺรหฺมํ อชฺโฌสิสฺสามิ, โอปสายิโก เต ภวิสฺสามิ วตฺถุสายิโก ยถา\-กาม\-กรณีโย พาหิเตยฺโย”ติ. อปิ จ เต อหํ พฺรหฺเม คติํ จ ปชานามิ, ชุติํ จ ปชานามิ “เอวํ มหิทฺธิโก พโก พฺรหฺมา, เอวํ มหานุภาโว พโก พฺรหฺมา, เอวํ มเหสกฺโข พโก พฺรหฺมา”ติ.}

\prose{\csromanfolio{404}ยถากถํ ปน เม ตฺวํ มาริส คติํ จ ปชานาสิ, ชุติํ จ ปชานาสิ “เอวํ มหิทฺธิโก พโก พฺรหฺมา, เอวํ มหานุภาโว พโก พฺรหฺมา, เอวํ มเหสกฺโข พโก พฺรหฺมา”ติ.}

\setlength{\csromangathapagewidth}{0pt}
\setlength{\csromangathawakwidth}{0pt}
\csromangathameasuregroup{“ยาวตา จนฺทิมสูริยา, \\ ตาว สหสฺสธา โลโก, \\ ปโรปรํ จ\textsuperscript{1} ชานาสิ, \\ อิตฺถภาวญฺญถาภาวํ,}{\csromangathabat{“ยาวตา จนฺทิมสูริยา,}{ปริหรนฺติ ทิสา ภนฺติ วิโรจนา.} \\ \csromangathabat{ตาว สหสฺสธา โลโก,}{เอตฺถ เต วตฺตเต\textsuperscript{1} วโส.} \\ \csromangathabat{ปโรปรํ จ\textsuperscript{1} ชานาสิ,}{อโถ ราควิราคินํ.} \\ \csromangathabat{อิตฺถภาวญฺญถาภาวํ,}{สตฺตานํ อาคติํ คตินฺ”ติ.}}
\csromangathasetleft{“ยาวตา จนฺทิมสูริยา, \\ ตาว สหสฺสธา โลโก, \\ ปโรปรํ จ\textsuperscript{1} ชานาสิ, \\ อิตฺถภาวญฺญถาภาวํ,}
\csromangathagroup{\csromangathastanza{\csromangathabat{“ยาวตา จนฺทิมสูริยา,}{ปริหรนฺติ ทิสา ภนฺติ วิโรจนา.} \\ \csromangathabat{ตาว สหสฺสธา โลโก,}{เอตฺถ เต วตฺตเต\csromansharedfootnote{f121716c87fec8f2}{วตฺตตี (สี, สฺยา, กํ, อิ)} วโส.}}\csromangathastanza{\csromangathabat{ปโรปรํ จ\csromansharedfootnote{0a064fbdbb1d634b}{ปโรวรญฺจ (สี, อิ)} ชานาสิ,}{อโถ ราควิราคินํ.} \\ \csromangathabat{อิตฺถ\-ภาว\-ญฺญถา\-ภาวํ,}{สตฺตานํ อาคติํ คตินฺ”ติ.}}}

\prose{เอวํ โข เต อหํ พฺรหฺเม คติํ จ ปชานามิ, ชุติํ จ ปชานามิ “เอวํ มหิทฺธิโก พโก พฺรหฺมา, เอวํ มหานุภาโว พโก พฺรหฺมา, เอวํ มเหสกฺโข พโก พฺรหฺมา”ติ.}

\proseitem{504}{อตฺถิ โข พฺรหฺเม อญฺโญ กาโย, ตํ ตฺวํ น ชานาสิ น ปสฺสสิ, ตมหํ ชานามิ ปสฺสามิ. อตฺถิ โข พฺรหฺเม อาภสฺสรา นาม กาโย ยโต ตฺวํ จุโต อิธูปปนฺโน, ตสฺส เต อติจิร\-นิวาเสน สา สติ ปมุฏฺฐา. เตน ตํ ตฺวํ น ชานาสิ น ปสฺสสิ, ตมหํ ชานามิ ปสฺสามิ. เอวํปิ โข อหํ พฺรหฺเม เนว เต สมสโม อภิญฺญาย, กุโต นีเจยฺยํ, อถ โข อหเมว ตยา ภิยฺโย. อตฺถิ โข พฺรหฺเม สุภกิโณฺห นาม กาโย. เวหปฺผโล นาม กาโย. อภิภู นาม กาโย, ตํ ตฺวํ น ชานาสิ น ปสฺสสิ, ตมหํ ชานามิ ปสฺสามิ. เอวํปิ โข อหํ พฺรหฺเม เนว เต สมสโม อภิญฺญาย, กุโต นีเจยฺยํ, อถ โข อหเมว ตยา ภิยฺโย. ปถวิํ โข อหํ พฺรหฺเม ปถวิโต อภิญฺญาย ยาวตา ปถวิยา ปถวตฺเตน อนนุภูตํ, ตทภิญฺญาย ปถวิํ นาปโหสิํ, ปถวิยา นาปโหสิํ, ปถวิโต นาปโหสิํ, ปถวิํ เมติ นาปโหสิํ, ปถวิํ นาภิวทิํ. เอวํปิ โข อหํ พฺรหฺเม เนว เต สมสโม อภิญฺญาย, กุโต นีเจยฺยํ, อถ โข อหเมว ตยา ภิยฺโย. อาปํ โข อหํ พฺรหฺเม ฯเปฯ. เตชํ โข อหํ พฺรหฺเม ฯเปฯ. วายํ โข อหํ พฺรหฺเม ฯเปฯ. ภูเต โข อหํ พฺรหฺเม ฯเปฯ. เทเว โข อหํ พฺรหฺเม ฯเปฯ. ปชาปติํ โข อหํ พฺรหฺเม ฯเปฯ. พฺรหฺมํ โข อหํ พฺรหฺเม ฯเปฯ. อาภสฺสเร โข อหํ พฺรหฺเม ฯเปฯ. สุภกิเณฺห โข อหํ \csromanfolio{405}พฺรหฺเม ฯเปฯ. เวหปฺผเล โข อหํ พฺรหฺเม ฯเปฯ. อภิภุํ โข อหํ พฺรหฺเม ฯเปฯ. สพฺพํ โข อหํ พฺรหฺเม สพฺพโต อภิญฺญาย ยาวตา สพฺพสฺส สพฺพตฺเตน อนนุภูตํ, ตทภิญฺญาย สพฺพํ นาปโหสิํ, สพฺพสฺมิํ นาปโหสิํ, สพฺพโต นาปโหสิํ, สพฺพํ เมติ นาปโหสิํ, สพฺพํ นาภิวทิํ. เอวํปิ โข อหํ พฺรหฺเม เนว เต สมสโม อภิญฺญาย, กุโต นีเจยฺยํ, อถ โข อหเมว ตยา ภิยฺโยติ.}

\prose{สเจ โข มาริส สพฺพสฺส สพฺพตฺเตน อนนุภูตํ. ตทภิญฺญาย มา เหว เต ริตฺตกเมว อโหสิ, ตุจฺฉกเมว อโหสีติ.}

\prose{วิญฺญาณํ อนิทสฺสนํ อนนฺตํ สพฺพโต ปภํ, ตํ ปถวิยา ปถวตฺเตน อนนุภูตํ, อาปสฺส อาปตฺเตน อนนุภูตํ, เตชสฺส เตชตฺเตน อนนุภูตํ, วายสฺส วายตฺเตน อนนุภูตํ, ภูตานํ ภูตตฺเตน อนนุภูตํ, เทวานํ เทวตฺเตน อนนุภูตํ, ปชาปติสฺส ปชาปติตฺเตน อนนุภูตํ, พฺรหฺมานํ พฺรหฺมตฺเตน อนนุภูตํ, อาภสฺสรานํ อาภสฺสรตฺเตน อนนุภูตํ, สุภกิณฺหานํ สุภกิณฺหตฺเตน อนนุภูตํ, เวหปฺผลานํ เวหปฺผลตฺเตน อนนุภูตํ, อภิภุสฺส อภิภุตฺเตน อนนุภูตํ, สพฺพสฺส สพฺพตฺเตน อนนุภูตํ.}

\prose{หนฺท จรหิ\csromansharedfootnote{1533573682f73fc3}{หนฺท จ หิ (สี, อิ)} เต มาริส ปสฺส อนฺตรธายามีติ. หนฺท จรหิ เม ตฺวํ พฺรหฺเม อนฺตรธายสฺสุ, สเจ วิสหสีติ. อถ โข ภิกฺขเว พโก พฺรหฺมา “อนฺตรธายิสฺสามิ สมณสฺส โคตมสฺส, อนฺตรธายิสฺสามิ สมณสฺส โคตมสฺสา”ติ เนวสฺสุ เม สกฺโกติ อนฺตรธายิตุํ.}

\prose{เอวํ วุตฺเต อหํ ภิกฺขเว พกํ พฺรหฺมานํ เอตทโวจํ “หนฺท จรหิ เต พฺรหฺเม อนฺตรธายามี”ติ. หนฺท จรหิ เม ตฺวํ มาริส อนฺตรธายสฺสุ, สเจ วิสหสีติ. อถ โข อหํ ภิกฺขเว ตถารูปํ อิทฺธา\-ภิสงฺขารํ อภิสงฺขาสิํ “เอตฺตาวตา พฺรหฺมา จ พฺรหฺมปริสา จ พฺรหฺม\-ปาริสชฺชา จ สทฺทญฺจ เม โสสฺสนฺติ\csromansharedfootnote{a44d39759098a72e}{สทฺทเมว สุยฺยนฺติ (ก)}, น จ มํ ทกฺขนฺตี”ติ อนฺตรหิโต อิมํ คาถํ อภาสิํ\csromandash{}}

\setlength{\csromangathapagewidth}{0pt}
\setlength{\csromangathawakwidth}{0pt}
\csromangathameasuregroup{“ภเววาหํ ภยํ ทิสฺวา, \\ ภวํ นาภิวทิํ กิญฺจิ,}{\csromangathabat{“ภเววาหํ ภยํ ทิสฺวา,}{ภวญฺจ วิภเวสินํ.} \\ \csromangathabat{ภวํ นาภิวทิํ กิญฺจิ,}{นนฺทิญฺจ น อุปาทิยินฺ”ติ.}}
\csromangathasetleft{“ภเววาหํ ภยํ ทิสฺวา, \\ ภวํ นาภิวทิํ กิญฺจิ,}
\csromangathagroup{\csromangathastanza{\csromangathabat{“ภเววาหํ ภยํ ทิสฺวา,}{ภวญฺจ วิภเวสินํ.} \\ \csromangathabat{ภวํ นาภิวทิํ กิญฺจิ,}{นนฺทิญฺจ น อุปาทิยินฺ”ติ.}}}

\proseitem{9}{\csromanfolio{406}อถ โข ภิกฺขเว พฺรหฺมา จ พฺรหฺมปริสา จ พฺรหฺม\-ปาริสชฺชา จ อจฺฉริย\-พฺภุต\-จิตฺต\-ชาตา อเหสุํ “อจฺฉริยํ วต โภ, อพฺภุตํ วต โภ สมณสฺส โคตมสฺส มหิทฺธิกตา มหานุภาวตา, น จ วต โน อิโต ปุพฺเพ ทิฏฺโฐ วา สุโต วา อญฺโญ สมโณ วา พฺราหฺมโณ วา เอวํ มหิทฺธิโก เอวํ มหานุภาโว, ยถายํ สมโณ โคตโม สกฺยปุตฺโต สกฺยกุลา ปพฺพชิโต, ภวรามาย วต โภ ปชาย ภวรตาย ภวสมฺมุทิตาย สมูลํ ภวํ อุทพฺพหี”ติ.}

\proseitem{505}{อถ โข ภิกฺขเว มาโร ปาปิมา อญฺญตรํ พฺรหฺม\-ปาริสชฺชํ อนฺวาวิสิตฺวา มํ เอตทโวจ “สเจ โข ตฺวํ มาริส เอวํ ปชานาสิ, สเจ ตฺวํ เอวํ อนุพุทฺโธ. มา สาวเก อุปเนสิ มา ปพฺพชิเต, มา สาวกานํ ธมฺมํ เทเสสิ มา ปพฺพชิตานํ, มา สาวเกสุ เคธิมกาสิ มา ปพฺพชิเตสุ. อเหสุํ โข ภิกฺขุ ตยา ปุพฺเพ สมณพฺราหฺมณา โลกสฺมิํ อรหนฺโต สมฺมาสมฺพุทฺธา ปฏิชานมานา, เต สาวเก อุปเนสุํ ปพฺพชิเต, สาวกานํ ธมฺมํ เทเสสุํ ปพฺพชิตานํ, สาวเกสุ เคธิมกํสุ ปพฺพชิเตสุ. เต สาวเก อุปเนตฺวา ปพฺพชิเต, สาวกานํ ธมฺมํ เทเสตฺวา ปพฺพชิตานํ, สาวเกสุ เคธิตจิตฺตา ปพฺพชิเตสุ กายสฺส เภทา ปาณุปจฺเฉทา หีเน กาเย ปติฏฺฐิตา. อเหสุํ เย ปน ภิกฺขุ ตยา ปุพฺเพ สมณพฺราหฺมณา โลกสฺมิํ อรหนฺโต สมฺมาสมฺพุทฺธา ปฏิชานมานา, เต น สาวเก อุปเนสุํ น ปพฺพชิเต, น สาวกานํ ธมฺมํ เทเสสุํ น ปพฺพชิตานํ, น สาวเกสุ เคธิมกํสุ น ปพฺพชิเตสุ. เต น สาวเก อุปเนตฺวา น ปพฺพชิเต, น สาวกานํ ธมฺมํ เทเสตฺวา น ปพฺพชิตานํ, น สาวเกสุ เคธิตจิตฺตา น ปพฺพชิเตสุ กายสฺส เภทา ปาณุปจฺเฉทา ปณีเต กาเย ปติฏฺฐิตา. ตํ ตาหํ ภิกฺขุ เอวํ วทามิ อิงฺฆ ตฺวํ มาริส อปฺโปสฺสุกฺโก ทิฏฺฐ\-ธมฺม\-สุข\-วิหารม\-นุยุตฺโต วิหรสฺสุ. อนกฺขาตํ กุสลํ หิ มาริส มา ปรํ โอวทาหี”ติ.}

\csromanmatikaanchor{mtk.84}
\csromantocmarkref{subsection}{10. มูรตชฺชนียสุตฺต}{mtk.84}
\bookmark[dest={mtk.84},level=2]{10. มูรตชฺชนียสุตฺต}
\prose{เอวํ วุตฺเต อหํ ภิกฺขเว มารํ ปาปิมนฺตํ เอตทโวจํ “ชานามิ โข ตาหํ ปาปิม, มา ตฺวํ มญฺญิตฺโถ ‘น มํ ชานาตี’ติ. มาโร ตฺวมสิ ปาปิม, น มํ ตฺวํ ปาปิม หิตานุกมฺปี เอวํ วเทสิ. อหิตานุกมฺปี มํ ตฺวํ ปาปิม เอวํ วเทสิ. ตุยฺหํ หิ ปาปิม เอวํ โหติ ‘เยสํ สมโณ โคตโม ธมฺมํ เทเสสฺสติ, เต เม วิสยํ อุปาติวตฺติสฺสนฺตี’ติ. อสมฺมาสมฺพุทฺธาว ปน \csromanfolio{407}เต ปาปิม สมานา ‘สมฺมาสมฺพุทฺธามฺหา’ติ ปฏิชานิํสุ. อหํ โข ปน ปาปิม สมฺมา\-สมฺพุทฺโธว สมาโน ‘สมฺมาสมฺพุทฺโธมฺหี’ติ ปฏิชานามิ. เทเสนฺโตปิ หิ ปาปิม ตถาคโต สาวกานํ ธมฺมํ ตาทิโสว, อเทเสนฺโตปิ หิ ปาปิม ตถาคโต สาวกานํ ธมฺมํ ตาทิโสว. อุปเนนฺโตปิ หิ ปาปิม ตถาคโต สาวเก ตาทิโสว, อนุปเนนฺโตปิ หิ ปาปิม ตถาคโต สาวเก ตาทิโสว. ตํ กิสฺส เหตุ, ตถาคตสฺส ปาปิม เย อาสวา สํกิเลสิกา โปโนพฺภวิกา สทรา ทุกฺขวิปากา อายติํ ชาติ\-ชรา\-มรณิยา, เต ปหีนา อุจฺฉินฺนมูลา ตาลาวตฺถุกตา อนภาวํกตา อายติํ อนุปาทธมฺมา. เสยฺยถาปิ ปาปิม ตาโล มตฺถก\-จฺฉินฺโน อภพฺโพ ปุน วิรูฬฺหิยา. เอวเมว โข ปาปิม ตถาคตสฺส เย อาสวา สํกิเลสิกา โปโนพฺภวิกา สทรา ทุกฺขวิปากา อายติํ ชาติ\-ชรา\-มรณิยา, เต ปหีนา อุจฺฉินฺนมูลา ตาลาวตฺถุกตา อนภาวํกตา อายติํ อนุปฺปาทธมฺมา”ติ.}

\prosewithcloserruled{อิติ หิ ทํ มารสฺส จ อนาลปนตาย พฺรหฺมุโน จ อภินิมนฺตนตาย, ตสฺมา อิมสฺส เวยฺยากรณสฺส พฺรหฺมนิมนฺตนิกํเตว อธิวจนนฺติ.}{พฺรหฺม\-นิมนฺตนิก\-สุตฺตํ นิฏฺฐิตํ นวมํ.}
\csromanheader{2}{10. มาร\-ตชฺชนีย\-สุตฺต}
\proseitem{506}{เอวํ เม สุตํ\csromandash{}เอกํ สมยํ อายสฺมา มหาโมคฺคลฺลาโน ภคฺเคสุ วิหรติ สุสุมารคิเร เภสกฬาวเน มิคทาเย. เตน โข ปน สมเยน อายสฺมา มหาโมคฺคลฺลาโน อพฺโภกาเส จงฺกมติ. เตน โข ปน สมเยน มาโร ปาปิมา อายสฺมโต มหา\-โมคฺคลฺลานสฺส กุจฺฉิคโต โหติ โกฏฺฐม\-นุปวิฏฺโฐ. อถ โข อายสฺมโต มหา\-โมคฺคลฺลานสฺส เอตทโหสิ “กิํ นุ โข เม กุจฺฉิ ครุคโร วิย\csromansharedfootnote{cad715b4ac66f05e}{ครุ ครุ วิย (สี, อิ, ฏีกายํ ปาฐนฺตรํ)} มาสาจิตํ มญฺเญ”ติ. อถ โข อายสฺมา มหาโมคฺคลฺลาโน จงฺกมา โอโรหิตฺวา วิหารํ ปวิสิตฺวา ปญฺญตฺเต อาสเน นิสีทิ, นิสชฺช โข อายสฺมา มหาโมคฺคลฺลาโน ปจฺจตฺตํ โยนิโส มนสากาสิ. อทฺทสา โข อายสฺมา มหาโมคฺคลฺลาโน มารํ ปาปิมนฺตํ กุจฺฉิคตํ โกฏฺฐม\-นุปวิฏฺฐํ. \csromanfolio{408}ทิสฺวาน มารํ ปาปิมนฺตํ เอตทโวจ “นิกฺขม ปาปิม, นิกฺขม ปาปิม. มา ตถาคตํ วิเหเสสิ มา ตถาคต\-สาวกํ, มา เต อโหสิ ทีฆรตฺตํ อหิตาย ทุกฺขายา”ติ. อถ โข มารสฺส ปาปิมโต เอตทโหสิ “อชานเมว โข มํ อยํ สมโณ อปสฺสํ เอวมาห ‘นิกฺขม ปาปิม, นิกฺขม ปาปิม. มา ตถาคตํ วิเหเสสิ มา ตถาคต\-สาวกํ, มา เต อโหสิ ทีฆรตฺตํ อหิตาย ทุกฺขายา’ติ. โยปิสฺส โส สตฺถา, โสปิ มํ เนว ขิปฺปํ ชาเนยฺย, กุโต ปน\csromansharedfootnote{507853473b12f4d2}{กุโต จ ปน (สฺยา)} มํ อยํ สาวโก ชานิสฺสตี”ติ. อถ โข อายสฺมา มหาโมคฺคลฺลาโน มารํ ปาปิมนฺตํ เอตทโวจ “เอวํปิ โข ตาหํ ปาปิม ชานามิ, มา ตฺวํ มญฺญิตฺโถ ‘น มํ ชานาตี’ติ, มาโร ตฺวมสิ ปาปิม, ตุยฺหํ หิ ปาปิม เอวํ โหติ ‘อชานเมว โข มํ อยํ สมโณ อปสฺสํ เอวมาห นิกฺขม ปาปิม, นิกฺขม ปาปิม. มา ตถาคตํ วิเหเสสิ มา ตถาคต\-สาวกํ, มา เต อโหสิ ทีฆรตฺตํ อหิตาย ทุกฺขายาติ. โยปิสฺส โส สตฺถา, โสปิ มํ เนว ขิปฺปํ ชาเนยฺย, กุโต ปน มํ อยํ สาวโก ชานิสฺสตี’ติ”.}

\prose{อถ โข มารสฺส ปาปิมโต เอตทโหสิ “ชานเมว โข มํ อยํ สมโณ ปสฺสํ เอวมาห ‘นิกฺขม ปาปิม, นิกฺขม ปาปิม. มา ตถาคตํ วิเหเสสิ มา ตถาคต\-สาวกํ, มา เต อโหสิ ทีฆรตฺตํ อหิตาย ทุกฺขายา’ติ”. อถ โข มาโร ปาปิมา อายสฺมโต มหา\-โมคฺคลฺลานสฺส มุขโต อุคฺคนฺตฺวา ปจฺจคฺคเฬ อฏฺฐาสิ.}

\proseitem{507}{อทฺทสา โข อายสฺมา มหาโมคฺคลฺลาโน มารํ ปาปิมนฺตํ ปจฺจคฺคเฬ ฐิตํ, ทิสฺวาน มารํ ปาปิมนฺตํ เอตทโวจ\csromandash{}เอตฺถาปิ โข ตาหํ ปาปิม ปสฺสามิ, มา ตฺวํ มญฺญิตฺโถ “น มํ ปสฺสตี”ติ. เอโส ตฺวํ ปาปิม ปจฺจคฺคเฬ ฐิโต. ภูตปุพฺพาหํ ปาปิม ทูสี นาม มาโร อโหสิํ, ตสฺส เม กาฬี นาม ภคินี. ตสฺสา ตฺวํ ปุตฺโต โส เม ตฺวํ ภาคิเนยฺโย อโหสิ. เตน โข ปน ปาปิม สมเยน กกุสนฺโธ ภควา อรหํ สมฺมาสมฺพุทฺโธ โลเก อุปฺปนฺโน โหติ. กกุสนฺธสฺส โข ปน ปาปิม ภควโต อรหโต สมฺมา\-สมฺพุทฺธสฺส วิธุร\-สญฺชีวํ นาม สาวกยุคํ อโหสิ อคฺคํ ภทฺทยุคํ. ยาวตา โข ปน ปาปิม กกุสนฺธสฺส \csromanfolio{409}ภควโต อรหโต สมฺมา\-สมฺพุทฺธสฺส สาวกา. เตสุ น จ โกจิ อายสฺมตา วิธุเรน สมสโม โหติ ยทิทํ ธมฺมเทสนาย. อิมินา โข เอวํ\csromansharedfootnote{441d372539a2de49}{เอตํ (สี, สฺยา, อิ)} ปาปิม ปริยาเยน อายสฺมโต วิธุรสฺส วิธุโรเตว\csromansharedfootnote{ea8578f7a0819094}{วิธุรสฺส วิธุโร วิธุโรเตฺวว (สี, สฺยา, กํ, อิ)} สมญฺญา อุทปาทิ.}

\prose{อายสฺมา ปน ปาปิม สญฺชีโว อรญฺญคโตปิ รุกฺขมูล\-คโตปิ สุญฺญาคาร\-คโตปิ อปฺปกสิเรเนว สญฺญาเวทยิต\-นิโรธํ สมาปชฺชติ. ภูตปุพฺพํ ปาปิม อายสฺมา สญฺชีโว อญฺญตรสฺมิํ รุกฺขมูเล สญฺญาเวทยิต\-นิโรธํ สมาปนฺโน นิสินฺโน โหติ. อทฺทสํสุ โข ปาปิม โคปาลกา ปสุปาลกา กสฺสกา ปถาวิโน อายสฺมนฺตํ สญฺชีวํ อญฺญตรสฺมิํ รุกฺขมูเล สญฺญาเวทยิต\-นิโรธํ สมาปนฺนํ นิสินฺนํ, ทิสฺวาน เตสํ เอตทโหสิ “อจฺฉริยํ วต โภ, อพฺภุตํ วต โภ, อยํ สมโณ นิสินฺนโกว กาลงฺกโต, หนฺท นํ ทหามา”ติ. อถ โข เต ปาปิม โคปาลกา ปสุปาลกา กสฺสกา ปถาวิโน ติณญฺจ กฏฺฐญฺจ โคมยญฺจ สํกฑฺฒิตฺวา อายสฺมโต สญฺชีวสฺส กาเย อุปจินิตฺวา อคฺคิํ ทตฺวา ปกฺกมิํสุ. อถ โข ปาปิม อายสฺมา สญฺชีโว ตสฺสา รตฺติยา อจฺจเยน ตาย สมาปตฺติยา วุฏฺฐหิตฺวา จีวรานิ ปปฺโผเฏตฺวา ปุพฺพณฺหสมยํ นิวาเสตฺวา ปตฺต\-จีวรมา\-ทาย คามํ ปิณฺฑาย ปาวิสิ. อทฺทสํสุ โข เต ปาปิม โคปาลกา ปสุปาลกา กสฺสกา ปถาวิโน อายสฺมนฺตํ สญฺชีวํ ปิณฺฑาย จรนฺตํ, ทิสฺวาน เนสํ เอตทโหสิ “อจฺฉริยํ วต โภ, อพฺภุตํ วต โภ, อยํ สมโณ นิสินฺนโกว กาลงฺกโต, สฺวายํ ปฏิสญฺชีวิโต”ติ. อิมินา โข เอวํ ปาปิม ปริยาเยน อายสฺมโต สญฺชีวสฺส สญฺชีโวเตว\csromansharedfootnote{8c42029f5402b122}{สญฺชีโว สญฺชีโวเตฺวว (สี, สฺยา, กํ, อิ)} สมญฺญา อุทปาทิ.}

\proseitem{508}{อถ โข ปาปิม ทูสิสฺส มารสฺส เอตทโหสิ “อิเมสํ โข อหํ ภิกฺขูนํ สีลวนฺตานํ กลฺยาณ\-ธมฺมานํ เนว ชานามิ อาคติํ วา คติํ วา, ยํนูนาหํ พฺราหฺมณคหปติเก อนฺวาวิเสยฺยํ, เอถ ตุเมฺห ภิกฺขู สีลวนฺเต กลฺยาณธมฺเม อกฺโกสถ ปริภาสถ โรเสถ วิเหเสถ. อปฺเปว นาม ตุเมฺหหิ อกฺโกสิย\-มานานํ ปริภาสิย\-มานานํ โรสิยมานานํ วิเหสิย\-มานานํ สิยา จิตฺตสฺส อญฺญถตฺตํ, ยถา ตํ \csromanfolio{410}ทูสี มาโร ลเภถ โอตารนฺ”ติ. อถ โข เต ปาปิม ทูสี มาโร พฺราหฺมณคหปติเก อนฺวาวิสิ “เอถ ตุเมฺห ภิกฺขู สีลวนฺเต กลฺยาณธมฺเม อกฺโกสถ ปริภาสถ โรเสถ วิเหเสถ. อปฺเปว นาม ตุเมฺหหิ อกฺโกสิย\-มานานํ ปริภาสิย\-มานานํ โรสิยมานานํ วิเหสิย\-มานานํ สิยา จิตฺตสฺส อญฺญถตฺตํ, ยถา ตํ ทูสี มาโร ลเภถ โอตารนฺ”ติ.}

\prose{อถ โข เต ปาปิม พฺราหฺมณ\-คหปติกา อนฺวาวิสิฏฺฐา ทูสินา มาเรน ภิกฺขู สีลวนฺเต กลฺยาณธมฺเม อกฺโกสนฺติ ปริภาสนฺติ โรเสนฺติ วิเหเสนฺติ “อิเม ปน มุณฺฑกา สมณกา อิพฺภา กิณฺหา\csromansharedfootnote{17c156dde9deb1c8}{กณฺหา (สฺยา, กํ, ก)} พนฺธุ\-ปาทา\-ปจฺจา ‘ฌายิโนสฺมา ฌายิโนสฺมา’ติ ปตฺตกฺขนฺธา อโธมุขา มธุรกชาตา ฌายนฺติ ปชฺฌายนฺติ นิชฺฌายนฺติ อปชฺฌายนฺติ. เสยฺยถาปิ นาม อุลูโก รุกฺขสาขายํ มูสิกํ มคฺคยมาโน ฌายติ ปชฺฌายติ นิชฺฌายติ อปชฺฌายติ. เอวเมวิเม มุณฺฑกา สมณกา อิพฺภา กิณฺหา พนฺธุ\-ปาทา\-ปจฺจา ‘ฌายิโนสฺมา ฌายิโนสฺมา’ติ ปตฺตกฺขนฺธา อโธมุขา มธุรกชาตา ฌายนฺติ ปชฺฌายนฺติ นิชฺฌายนฺติ อปชฺฌายนฺติ. เสยฺยถาปิ นาม โกตฺถุ นทีตีเร มจฺเฉ มคฺคยมาโน ฌายติ ปชฺฌายติ นิชฺฌายติ อปชฺฌายติ. เอวเมวิเม มุณฺฑกา สมณกา อิพฺภา กิณฺหา พนฺธุ\-ปาทา\-ปจฺจา ‘ฌายิโนสฺมา ฌายิโนสฺมา’ติ ปตฺตกฺขนฺธา อโธมุขา มธุรกชาตา ฌายนฺติ ปชฺฌายนฺติ นิชฺฌายนฺติ อปชฺฌายนฺติ. เสยฺยถาปิ นาม พิฬาโร สนฺธิ\-สมล\-สงฺกฏีเร มูสิกํ มคฺคยมาโน ฌายติ ปชฺฌายติ นิชฺฌายติ อปชฺฌายติ. เอวเมวิเม มุณฺฑกา สมณกา อิพฺภา กิณฺหา พนฺธุ\-ปาทา\-ปจฺจา ‘ฌายิโนสฺมา ฌายิโนสฺมา’ติ ปตฺตกฺขนฺธา อโธมุขา มธุรกชาตา ฌายนฺติ ปชฺฌายนฺติ นิชฺฌายนฺติ อปชฺฌายนฺติ. เสยฺยถาปิ นาม คทฺรโภ วหจฺฉินฺโน สนฺธิ\-สมล\-สงฺกฏีเร ฌายติ ปชฺฌายติ นิชฺฌายติ อปชฺฌายติ. เอวเมวิเม มุณฺฑกา สมณกา อิพฺภา กิณฺหา พนฺธุ\-ปาทา\-ปจฺจา ‘ฌายิโนสฺมา ฌายิโนสฺมา’ติ ปตฺตกฺขนฺธา อโธมุขา มธุรกชาตา ฌายนฺติ ปชฺฌายนฺติ นิชฺฌายนฺติ อปชฺฌายนฺตี”ติ.}

\prose{เย โข ปน ปาปิม เตน สมเยน มนุสฺสา กาลงฺกโรนฺติ เยภุยฺเยน กายสฺส เภทา ปรํ มรณา อปายํ ทุคฺคติํ วินิปาตํ นิรยํ อุปปชฺชนฺติ.}

\proseitem{509}{\csromanfolio{411}อถ โข ปาปิม กกุสนฺโธ ภควา อรหํ สมฺมาสมฺพุทฺโธ ภิกฺขู อามนฺเตสิ “อนฺวาวิฏฺฐา โข ภิกฺขเว พฺราหฺมณ\-คหปติกา ทูสินา มาเรน. เอถ ตุเมฺห ภิกฺขู สีลวนฺเต กลฺยาณธมฺเม อกฺโกสถ ปริภาสถ โรเสถ วิเหเสถ. อปฺเปว นาม ตุเมฺหหิ อกฺโกสิย\-มานานํ ปริภาสิย\-มานานํ โรสิยมานานํ วิเหสิย\-มานานํ สิยา จิตฺตสฺส อญฺญถตฺตํ, ยถา ตํ ทูสี มาโร ลเภถ โอตารนฺติ. เอถ ตุเมฺห ภิกฺขเว เมตฺตา\-สหคเตน เจตสา เอกํ ทิสํ ผริตฺวา วิหรถ. ตถา ทุติยํ. ตถา ตติยํ. ตถา จตุตฺถํ. อิติ อุทฺธมโธ ติริยํ สพฺพธิ สพฺพตฺตตาย สพฺพาวนฺตํ โลกํ เมตฺตา\-สหคเตน เจตสา วิปุเลน มหคฺคเตน อปฺปมาเณน อเวเรน อพฺยาพชฺเฌน ผริตฺวา วิหรถ. กรุณา\-สหคเตน เจตสา ฯเปฯ. มุทิตา\-สหคเตน เจตสา ฯเปฯ. อุเปกฺขา\-สหคเตน เจตสา เอกํ ทิสํผริตฺวา วิหรถ. ตถา ทุติยํ. ตถา ตติยํ. ตถา จตุตฺถํ. อิติ อุทฺธมโธ ติริยํ สพฺพธิ สพฺพตฺตตาย สพฺพาวนฺตํ โลกํ อุเปกฺขา\-สหคเตน เจตสา วิปุเลน มหคฺคเตน อปฺปมาเณน อเวเรน อพฺยาพชฺเฌน ผริตฺวา วิหรถา”ติ.}

\prose{อถ โข เต ปาปิม ภิกฺขู กกุสนฺเธน ภควตา อรหตา สมฺมา\-สมฺพุทฺเธน เอวํ โอวทิยมานา เอวํ อนุสาสิยมานา อรญฺญคตาปิ รุกฺขมูลคตาปิ สุญฺญาคารคตาปิ เมตฺตา\-สหคเตน เจตสา เอกํ ทิสํ ผริตฺวา วิหริํสุ. ตถา ทุติยํ. ตถา ตติยํ. ตถา จตุตฺถํ. อิติ อุทฺธมโธ ติริยํ สพฺพธิ สพฺพตฺตตาย สพฺพาวนฺตํ โลกํ เมตฺตา\-สหคเตน เจตสา วิปุเลน มหคฺคเตน อปฺปมาเณน อเวเรน อพฺยาพชฺเฌน ผริตฺวา วิหริํสุ. กรุณา\-สหคเตน เจตสา ฯเปฯ. มุทิตา\-สหคเตน เจตสา ฯเปฯ. อุเปกฺขา\-สหคเตน เจตสา เอกํ ทิสํ ผริตฺวา วิหริํสุ. ตถา ทุติยํ. ตถา ตติยํ. ตถา จตุตฺถํ. อิติ อุทฺธมโธ ติริยํ สพฺพธิ สพฺพตฺตตาย สพฺพาวนฺตํ โลกํ อุเปกฺขา\-สหคเตน เจตสา วิปุเลน มหคฺคเตน อปฺปมาเณน อเวเรน อพฺยาพชฺเฌน ผริตฺวา วิหริํสุ.}

\proseitem{510}{อถ โข ปาปิม ทูสิสฺส มารสฺส เอตทโหสิ “เอวํปิ โข อหํ กโรนฺโต อิเมสํ ภิกฺขูนํ สีลวนฺตานํ กลฺยาณ\-ธมฺมานํ เนว ชานามิ อาคติํ วา คติํ วา, ยํนูนาหํ พฺราหฺมณคหปติเก อนฺวาวิเสยฺยํ, เอถ ตุเมฺห ภิกฺขู สีลวนฺเต กลฺยาณธมฺเม สกฺกโรถ ครุํ กโรถ \csromanfolio{412}มาเนถ ปูเชถ. อปฺเปว นาม ตุเมฺหหิ สกฺกริยมานานํ ครุกริยมานานํ มานิยมานานํ ปูชิยมานานํ สิยา จิตฺตสฺส อญฺญถตฺตํ, ยถา ตํ ทูสี มาโร ลเภถ โอตารนฺ”ติ. อถ โข เต ปาปิม ทูสี มาโร พฺราหฺมณคหปติเก อนฺวาวิสิ. เอถ ตุเมฺห ภิกฺขู สีลวนฺเต กลฺยาณธมฺเม สกฺกโรถ ครุํ กโรถ มาเนถ ปูเชถ. อปฺเปว นาม ตุเมฺหหิ สกฺกริยมานานํ ครุกริยมานานํ มานิยมานานํ ปูชิยมานานํ สิยา จิตฺตสฺส อญฺญถตฺตํ, ยถา ตํ ทูสี มาโร ลเภถ โอตารนฺติ. อถ โข เต ปาปิม พฺราหฺมณ\-คหปติกา อนฺวาวิฏฺฐา ทูสินา มาเรน ภิกฺขู สีลวนฺเต กลฺยาณธมฺเม สกฺกโรนฺติ ครุํ กโรนฺติ มาเนนฺติ ปูเชนฺติ.}

\prose{เย โข ปน ปาปิม เตน สมเยน มนุสฺสา กาลงฺกโรนฺติ. เยภุยฺเยน กายสฺส เภทา ปรํ มรณา สุคติํ สคฺคํ โลกํ อุปปชฺชนฺติ.}

\proseitem{511}{อถ โข ปาปิม กกุสนฺโธ ภควา อรหํ สมฺมาสมฺพุทฺโธ ภิกฺขู อามนฺเตสิ “อนฺวาวิฏฺฐา โข ภิกฺขเว พฺราหฺมณ\-คหปติกา ทูสินา มาเรน. เอถ ตุเมฺห ภิกฺขู สีลวนฺเต กลฺยาณธมฺเม สกฺกโรถ ครุํ กโรถ มาเนถ ปูเชถ. อปฺเปว นาม ตุเมฺหหิ สกฺกริยมานานํ ครุกริยมานานํ มานิยมานานํ ปูชิยมานานํ สิยา จิตฺตสฺส อญฺญถตฺตํ, ยถา ตํ ทูสี มาโร ลเภถ โอตารนฺติ. เอถ ตุเมฺห ภิกฺขเว อสุภา\-นุปสฺสิโน กาเย วิหรถ, อาหาเร ปฏิกูล\-สญฺญิโน สพฺพโลเก อนภิรติ\-สญฺญิโน\csromansharedfootnote{5df86dc4242cf880}{อนภิรตสญฺญิโน (สี, สฺยา, กํ, อิ)} สพฺพ\-สงฺขาเรสุ อนิจฺจา\-นุปสฺสิโน”ติ.}

\prose{อถ โข เต ปาปิม ภิกฺขู กกุสนฺเธน ภควตา อรหตา สมฺมา\-สมฺพุทฺเธน เอวํ โอวทิยมานา เอวํ อนุสาสิยมานา อรญฺญคตาปิ รุกฺขมูลคตาปิ สุญฺญาคารคตาปิ อสุภา\-นุปสฺสิโน กาเย วิหริํสุ, อาหาเร ปฏิกูล\-สญฺญิโน สพฺพโลเก อนภิรติ\-สญฺญิโน สพฺพ\-สงฺขาเรสุ อนิจฺจา\-นุปสฺสิโน.}

\proseitem{512}{อถ โข ปาปิม กกุสนฺโธ ภควา อรหํ สมฺมาสมฺพุทฺโธ ปุพฺพณฺหสมยํ นิวาเสตฺวา ปตฺต\-จีวรมา\-ทาย อายสฺมตา วิธุเรน ปจฺฉาสมเณน คามํ ปิณฺฑาย ปาวิสิ. อถ โข ปาปิม ทูสี มาโร อญฺญตรํ \csromanfolio{413}กุมารกํ\csromansharedfootnote{134e7b72a15210e4}{กุมารํ (สี, อิ)} อนฺวาวิสิตฺวา สกฺขรํ คเหตฺวา อายสฺมโต วิธุรสฺส สีเส ปหารมทาสิ, สีสํ โวภินฺทิ\csromansharedfootnote{e6b481b11774ec4b}{สีสํ เต ภินฺทิสฺสามีติ (ก)}. อถ โข ปาปิม อายสฺมา วิธุโร ภินฺเนน สีเสน โลหิเตน คฬนฺเตน กกุสนฺธํเยว ภควนฺตํ อรหนฺตํ สมฺมา\-สมฺพุทฺธํ ปิฏฺฐิโต ปิฏฺฐิโต อนุพนฺธิ. อถ โข ปาปิม กกุสนฺโธ ภควา อรหํ สมฺมาสมฺพุทฺโธ นาคาปโลกิตํ อปโลเกสิ “น วายํ ทูสี มาโร มตฺตมญฺญาสี”ติ. สหาปโลกนาย จ ปน ปาปิม ทูสี มาโร ตมฺหา จ ฐานา จวิ, มหานิรยํ จ อุปปชฺชิ.}

\prose{ตสฺส โข ปน ปาปิม มหานิรยสฺส ตโย นามเธยฺยา โหนฺติ ฉผสฺสา\-ยตนิโก อิติปิ สงฺกุสมาหโต อิติปิ ปจฺจตฺต\-เวทนิโย อิติปิ. อถ โข มํ ปาปิม นิรยปาลา อุปสงฺกมิตฺวา เอตทโวจุํ “ยทา โข เต\csromansharedfootnote{310c1f58487b3fb9}{ยโต เต (ก)} มาริส สงฺกุนา สงฺกุ หทเย สมาคจฺเฉยฺย. อถ นํ ตฺวํ ชาเนยฺยาสิ วสฺสสหสฺสํ เม นิรเย ปจฺจมานสฺสา”ติ. โส โข อหํ ปาปิม พหูนิ วสฺสานิ พหูนิ วสฺสสตานิ พหูนิ วสฺสสหสฺสานิ ตสฺมิํ มหานิรเย อปจฺจิํ. ทสวสฺสสหสฺสานิ ตสฺเสว มหานิรยสฺส อุสฺสเท อปจฺจิํ วุฏฺฐานิมํ นาม เวทนํ เวทิยมาโน. ตสฺส มยฺหํ ปาปิม เอวรูโป กาโย โหติ เสยฺยถาปิ มนุสฺสสฺส. เอวรูปํ สีสํ โหติ เสยฺยถาปิ มจฺฉสฺส.}

\csromanmatikaanchor{mtk.85}
\csromantocmarkref{subsection}{อุทฺทานคาถา}{mtk.85}
\bookmark[dest={mtk.85},level=2]{อุทฺทานคาถา}
\setlength{\csromangathapagewidth}{0pt}
\setlength{\csromangathawakwidth}{0pt}
\csromangathameasuregroup{กีทิโส นิรโย อาสิ, \\ วิธุรํ สาวกมาสชฺช, \\ สตํ อาสิ อโยสงฺกู, \\ ¢ทิโส นิรโย อาสิ, \\ วิธุรํ สาวกมาสชฺช, \\ โย เอตมภิชานาติ, \\ ตาทิสํ ภิกฺขุมาสชฺช, \\ มชฺเฌ สรสฺส ติฏฺฐนฺติ, \\ เวฬุริยวณฺณา รุจิรา, \\ อจฺฉรา ตตฺถ นจฺจนฺติ, \\ โย เอตมภิชานาติ, \\ ตาทิสํ ภิกฺขุมาสชฺช, \\ โย เว พุทฺเธน โจทิโต, \\ มิคารมาตุปาสาทํ, \\ โย เอตมภิชานาติ, \\ ตาทิสํ ภิกฺขุมาสชฺช, \\ โย เวชยนฺตํ ปาสาทํ, \\ อิทฺธิพเลนุปตฺถทฺโธ, \\ โย เอตมภิชานาติ, \\ ตาทิสํ ภิกฺขุมาสชฺช, \\ โย เวชยนฺตปาสาเท, \\ อปิ วาสว ชานาสิ, \\ ตสฺส สกฺโก วิยากาสิ, \\ โย เอตมภิชานาติ, \\ ตาทิสํ ภิกฺขุมาสชฺช, \\ โย พฺรหฺมํ ปริปุจฺฉติ, \\ อชฺชาปิ ตฺยาวุโส ทิฏฺฐิ, \\ ปสฺสสิ วีติวตฺตนฺตํ, \\ ตสฺส พฺรหฺมา วิยากาสิ, \\ น เม มาริส สา ทิฏฺฐิ, \\ ปสฺสามิ วีติวตฺตนฺตํ, \\ โสหํ อชฺช กถํ วชฺชํ, \\ โย เอตมภิชานาติ, \\ ตาทิสํ ภิกฺขุมาสชฺช, \\ โย มหาเมรุโน กูฏํ, \\ วนํ ปุพฺพวิเทหานํ, \\ โย เอตมภิชานาติ, \\ ตาทิสํ ภิกฺขุมาสชฺช, \\ น เว อคฺคิ เจตยติ\textsuperscript{1}, \\ พาโล จ ชลิตํ อคฺคิํ, \\ เอวเมว ตุวํ มาร, \\ สยํ ฑหิสฺสสิ อตฺตานํ, \\ อปุญฺญํ ปสวี มาโร, \\ กินฺนุ มญฺญสิ ปาปิม, \\ กโรโต จียติ ปาปํ, \\ มาร นิพฺพินฺท พุทฺธมฺหา, \\ อิติ มารํ อตชฺเชสิ, \\ ตโต โส ทุมฺมโน ยกฺโข,}{\csromangathabat{กีทิโส นิรโย อาสิ,}{ยตฺถ ทูสี อปจฺจถ.} \\ \csromangathabat{วิธุรํ สาวกมาสชฺช,}{กกุสนฺธญฺจ พฺราหฺมณํ.} \\ \csromangathabat{สตํ อาสิ อโยสงฺกู,}{สพฺเพ ปจฺจตฺตเวทนา.} \\ \csromangathabat{¢ทิโส นิรโย อาสิ,}{ยตฺถ ทูสี อปจฺจถ.} \\ \csromangathabat{วิธุรํ สาวกมาสชฺช,}{กกุสนฺธญฺจ พฺราหฺมณํ.} \\ \csromangathabat{โย เอตมภิชานาติ,}{ภิกฺขุ พุทฺธสฺส สาวโก.} \\ \csromangathabat{ตาทิสํ ภิกฺขุมาสชฺช,}{กณฺห ทุกฺขํ นิคจฺฉสิ.} \\ \csromangathabat{มชฺเฌ สรสฺส ติฏฺฐนฺติ,}{วิมานา กปฺปฏฺฐายิโน.} \\ \csromangathabat{เวฬุริยวณฺณา รุจิรา,}{อจฺจิมนฺโต ปภสฺสรา.} \\ \csromangathabat{อจฺฉรา ตตฺถ นจฺจนฺติ,}{ปุถุ นานตฺตวณฺณิโย.} \\ \csromangathabat{โย เอตมภิชานาติ,}{ภิกฺขุ พุทฺธสฺส สาวโก.} \\ \csromangathabat{ตาทิสํ ภิกฺขุมาสชฺช,}{กณฺห ทุกฺขํ นิคจฺฉสิ.} \\ \csromangathabat{โย เว พุทฺเธน โจทิโต,}{ภิกฺขุ สงฺฆสฺส เปกฺขโต.} \\ \csromangathabat{มิคารมาตุปาสาทํ,}{ปาทงฺคุฏฺเฐน กมฺปยิ.} \\ \csromangathabat{โย เอตมภิชานาติ,}{ภิกฺขุ พุทฺธสฺส สาวโก.} \\ \csromangathabat{ตาทิสํ ภิกฺขุมาสชฺช,}{กณฺห ทุกฺขํ นิคจฺฉสิ.} \\ \csromangathabat{โย เวชยนฺตํ ปาสาทํ,}{ปาทงฺคุฏฺเฐน กมฺปยิ.} \\ \csromangathabat{อิทฺธิพเลนุปตฺถทฺโธ,}{สํเวเชสิ จ เทวตา.} \\ \csromangathabat{โย เอตมภิชานาติ,}{ภิกฺขุ พุทฺธสฺส สาวโก.} \\ \csromangathabat{ตาทิสํ ภิกฺขุมาสชฺช,}{กณฺห ทุกฺขํ นิคจฺฉสิ.} \\ \csromangathabat{โย เวชยนฺตปาสาเท,}{สกฺกํ โส ปริปุจฺฉติ.} \\ \csromangathabat{อปิ วาสว ชานาสิ,}{ตณฺหากฺขยวิมุตฺติโย.} \\ \csromangathabat{ตสฺส สกฺโก วิยากาสิ,}{ปญฺหํ ปุฏฺโฐ ยถาตถํ.} \\ \csromangathabat{โย เอตมภิชานาติ,}{ภิกฺขุ พุทฺธสฺส สาวโก.} \\ \csromangathabat{ตาทิสํ ภิกฺขุมาสชฺช,}{กณฺห ทุกฺขํ นิคจฺฉสิ.} \\ \csromangathabat{โย พฺรหฺมํ ปริปุจฺฉติ,}{สุธมฺมายาภิโต สภํ.} \\ \csromangathabat{อชฺชาปิ ตฺยาวุโส ทิฏฺฐิ,}{ยา เต ทิฏฺฐิ ปุเร อหุ.} \\ \csromangathabat{ปสฺสสิ วีติวตฺตนฺตํ,}{พฺรหฺมโลเก ปภสฺสรํ.} \\ \csromangathabat{ตสฺส พฺรหฺมา วิยากาสิ,}{อนุปุพฺพํ ยถาตถํ.} \\ \csromangathabat{น เม มาริส สา ทิฏฺฐิ,}{ยา เม ทิฏฺฐิ ปุเร อหุ.} \\ \csromangathabat{ปสฺสามิ วีติวตฺตนฺตํ,}{พฺรหฺมโลเก ปภสฺสรํ.} \\ \csromangathabat{โสหํ อชฺช กถํ วชฺชํ,}{อหํ นิจฺโจมฺหิ สสฺสโต.} \\ \csromangathabat{โย เอตมภิชานาติ,}{ภิกฺขุ พุทฺธสฺส สาวโก.} \\ \csromangathabat{ตาทิสํ ภิกฺขุมาสชฺช,}{กณฺห ทุกฺขํ นิคจฺฉสิ.} \\ \csromangathabat{โย มหาเมรุโน กูฏํ,}{วิโมกฺเขน อผสฺสยิ.} \\ \csromangathabat{วนํ ปุพฺพวิเทหานํ,}{เย จ ภูมิสยา นรา.} \\ \csromangathabat{โย เอตมภิชานาติ,}{ภิกฺขุ พุทฺธสฺส สาวโก.} \\ \csromangathabat{ตาทิสํ ภิกฺขุมาสชฺช,}{กณฺห ทุกฺขํ นิคจฺฉสิ.} \\ \csromangathabat{น เว อคฺคิ เจตยติ\textsuperscript{1},}{อหํ พาลํ ฑหามีติ.} \\ \csromangathabat{พาโล จ ชลิตํ อคฺคิํ,}{อาสชฺช นํ ส ฑยฺหติ.} \\ \csromangathabat{เอวเมว ตุวํ มาร,}{อาสชฺช นํ ตถาคตํ.} \\ \csromangathabat{สยํ ฑหิสฺสสิ อตฺตานํ,}{พาโล อคฺคิํว สํผุสํ.} \\ \csromangathabat{อปุญฺญํ ปสวี มาโร,}{อาสชฺช นํ ตถาคตํ.} \\ \csromangathabat{กินฺนุ มญฺญสิ ปาปิม,}{น เม ปาปํ วิปจฺจติ.} \\ \csromangathabat{กโรโต จียติ ปาปํ,}{จิรรตฺตาย อนฺตก.} \\ \csromangathabat{มาร นิพฺพินฺท พุทฺธมฺหา,}{อาสํ มากาสิ ภิกฺขุสุ.} \\ \csromangathabat{อิติ มารํ อตชฺเชสิ,}{ภิกฺขุ เภสกฬาวเน.} \\ \csromangathabat{ตโต โส ทุมฺมโน ยกฺโข,}{ตตฺเถวนฺตรธายถาติ.}}
\csromangathasetleft{กีทิโส นิรโย อาสิ, \\ วิธุรํ สาวกมาสชฺช, \\ สตํ อาสิ อโยสงฺกู, \\ ¢ทิโส นิรโย อาสิ, \\ วิธุรํ สาวกมาสชฺช, \\ โย เอตมภิชานาติ, \\ ตาทิสํ ภิกฺขุมาสชฺช, \\ มชฺเฌ สรสฺส ติฏฺฐนฺติ, \\ เวฬุริยวณฺณา รุจิรา, \\ อจฺฉรา ตตฺถ นจฺจนฺติ, \\ โย เอตมภิชานาติ, \\ ตาทิสํ ภิกฺขุมาสชฺช, \\ โย เว พุทฺเธน โจทิโต, \\ มิคารมาตุปาสาทํ, \\ โย เอตมภิชานาติ, \\ ตาทิสํ ภิกฺขุมาสชฺช, \\ โย เวชยนฺตํ ปาสาทํ, \\ อิทฺธิพเลนุปตฺถทฺโธ, \\ โย เอตมภิชานาติ, \\ ตาทิสํ ภิกฺขุมาสชฺช, \\ โย เวชยนฺตปาสาเท, \\ อปิ วาสว ชานาสิ, \\ ตสฺส สกฺโก วิยากาสิ, \\ โย เอตมภิชานาติ, \\ ตาทิสํ ภิกฺขุมาสชฺช, \\ โย พฺรหฺมํ ปริปุจฺฉติ, \\ อชฺชาปิ ตฺยาวุโส ทิฏฺฐิ, \\ ปสฺสสิ วีติวตฺตนฺตํ, \\ ตสฺส พฺรหฺมา วิยากาสิ, \\ น เม มาริส สา ทิฏฺฐิ, \\ ปสฺสามิ วีติวตฺตนฺตํ, \\ โสหํ อชฺช กถํ วชฺชํ, \\ โย เอตมภิชานาติ, \\ ตาทิสํ ภิกฺขุมาสชฺช, \\ โย มหาเมรุโน กูฏํ, \\ วนํ ปุพฺพวิเทหานํ, \\ โย เอตมภิชานาติ, \\ ตาทิสํ ภิกฺขุมาสชฺช, \\ น เว อคฺคิ เจตยติ\textsuperscript{1}, \\ พาโล จ ชลิตํ อคฺคิํ, \\ เอวเมว ตุวํ มาร, \\ สยํ ฑหิสฺสสิ อตฺตานํ, \\ อปุญฺญํ ปสวี มาโร, \\ กินฺนุ มญฺญสิ ปาปิม, \\ กโรโต จียติ ปาปํ, \\ มาร นิพฺพินฺท พุทฺธมฺหา, \\ อิติ มารํ อตชฺเชสิ, \\ ตโต โส ทุมฺมโน ยกฺโข,}
\csromangathagroupwithcloser{\csromangathastanza{\csromangathaitembat{513}{กีทิโส นิรโย อาสิ,}{ยตฺถ ทูสี อปจฺจถ.} \\ \csromangathacontbat{วิธุรํ สาวกมาสชฺช,}{กกุสนฺธญฺจ พฺราหฺมณํ.} \\ \csromangathacontbat{สตํ อาสิ อโยสงฺกู,}{สพฺเพ ปจฺจตฺตเวทนา.} \\ \csromangathacontbat{¢ทิโส นิรโย อาสิ,}{ยตฺถ ทูสี อปจฺจถ.} \\ \csromangathacontbat{วิธุรํ สาวกมาสชฺช,}{กกุสนฺธญฺจ พฺราหฺมณํ.} \\ \csromangathacontbat{โย เอตมภิชานาติ,}{ภิกฺขุ พุทฺธสฺส สาวโก.} \\ \csromangathacontbat{ตาทิสํ ภิกฺขุมาสชฺช,}{กณฺห ทุกฺขํ นิคจฺฉสิ.} \\ \csromangathacontbat{มชฺเฌ สรสฺส ติฏฺฐนฺติ,}{วิมานา กปฺปฏฺฐายิโน.} \\ \csromangathacontbat{เวฬุริยวณฺณา รุจิรา,}{อจฺจิมนฺโต ปภสฺสรา.} \\ \csromangathacontbat{อจฺฉรา ตตฺถ นจฺจนฺติ,}{ปุถุ นานตฺตวณฺณิโย.}}\csromangathastanza{\csromanfolio{414}\csromangathabat{โย เอตมภิชานาติ,}{ภิกฺขุ พุทฺธสฺส สาวโก.} \\ \csromangathabat{ตาทิสํ ภิกฺขุมาสชฺช,}{กณฺห ทุกฺขํ นิคจฺฉสิ.}}\csromangathastanza{\csromangathabat{โย เว พุทฺเธน โจทิโต,}{ภิกฺขุ สงฺฆสฺส เปกฺขโต.} \\ \csromangathabat{มิคาร\-มาตุ\-ปาสาทํ,}{ปาทงฺคุฏฺเฐน กมฺปยิ.}}\csromangathastanza{\csromangathabat{โย เอตมภิชานาติ,}{ภิกฺขุ พุทฺธสฺส สาวโก.} \\ \csromangathabat{ตาทิสํ ภิกฺขุมาสชฺช,}{กณฺห ทุกฺขํ นิคจฺฉสิ.}}\csromangathastanza{\csromangathabat{โย เวชยนฺตํ ปาสาทํ,}{ปาทงฺคุฏฺเฐน กมฺปยิ.} \\ \csromangathabat{อิทฺธิพเลนุปตฺถทฺโธ,}{สํเวเชสิ จ เทวตา.}}\csromangathastanza{\csromangathabat{โย เอตมภิชานาติ,}{ภิกฺขุ พุทฺธสฺส สาวโก.} \\ \csromangathabat{ตาทิสํ ภิกฺขุมาสชฺช,}{กณฺห ทุกฺขํ นิคจฺฉสิ.}}\csromangathastanza{\csromangathabat{โย เวชยนฺต\-ปาสาเท,}{สกฺกํ โส ปริปุจฺฉติ.} \\ \csromangathabat{อปิ วาสว ชานาสิ,}{ตณฺหากฺขยวิมุตฺติโย.}}\csromangathastanza{\csromangathabat{ตสฺส สกฺโก วิยากาสิ,}{ปญฺหํ ปุฏฺโฐ ยถาตถํ.} \\ \csromangathabat{โย เอตมภิชานาติ,}{ภิกฺขุ พุทฺธสฺส สาวโก.}}\csromangathastanza{\csromangathabat{ตาทิสํ ภิกฺขุมาสชฺช,}{กณฺห ทุกฺขํ นิคจฺฉสิ.} \\ \csromangathabat{โย พฺรหฺมํ ปริปุจฺฉติ,}{สุธมฺมายาภิโต สภํ.}}\csromangathastanza{\csromangathabat{อชฺชาปิ ตฺยาวุโส ทิฏฺฐิ,}{ยา เต ทิฏฺฐิ ปุเร อหุ.} \\ \csromangathabat{ปสฺสสิ วีติวตฺตนฺตํ,}{พฺรหฺมโลเก ปภสฺสรํ.}}\csromangathastanza{\csromangathabat{ตสฺส พฺรหฺมา วิยากาสิ,}{อนุปุพฺพํ ยถาตถํ.} \\ \csromangathabat{น เม มาริส สา ทิฏฺฐิ,}{ยา เม ทิฏฺฐิ ปุเร อหุ.}}\csromangathastanza{\csromangathabat{ปสฺสามิ วีติวตฺตนฺตํ,}{พฺรหฺมโลเก ปภสฺสรํ.} \\ \csromangathabat{โสหํ อชฺช กถํ วชฺชํ,}{อหํ นิจฺโจมฺหิ สสฺสโต.}}\csromangathastanza{\csromangathabat{โย เอตมภิชานาติ,}{ภิกฺขุ พุทฺธสฺส สาวโก.} \\ \csromangathabat{ตาทิสํ ภิกฺขุมาสชฺช,}{กณฺห ทุกฺขํ นิคจฺฉสิ.}}\csromangathastanza{\csromangathabat{โย มหาเมรุโน กูฏํ,}{วิโมกฺเขน อผสฺสยิ.} \\ \csromangathabat{วนํ ปุพฺพวิเทหานํ,}{เย จ ภูมิสยา นรา.}}\csromangathastanza{\csromangathabat{โย เอตมภิชานาติ,}{ภิกฺขุ พุทฺธสฺส สาวโก.} \\ \csromangathabat{ตาทิสํ ภิกฺขุมาสชฺช,}{กณฺห ทุกฺขํ นิคจฺฉสิ.}}\csromangathastanza{\csromanfolio{415}\csromangathabat{น เว อคฺคิ เจตยติ\csromansharedfootnote{4869e704beee69aa}{เวฐยติ (สี)},}{อหํ พาลํ ฑหามีติ.} \\ \csromangathabat{พาโล จ ชลิตํ อคฺคิํ,}{อาสชฺช นํ ส ฑยฺหติ.}}\csromangathastanza{\csromangathabat{เอวเมว ตุวํ มาร,}{อาสชฺช นํ ตถาคตํ.} \\ \csromangathabat{สยํ ฑหิสฺสสิ อตฺตานํ,}{พาโล อคฺคิํว สํผุสํ.}}\csromangathastanza{\csromangathabat{อปุญฺญํ ปสวี มาโร,}{อาสชฺช นํ ตถาคตํ.} \\ \csromangathabat{กินฺนุ มญฺญสิ ปาปิม,}{น เม ปาปํ วิปจฺจติ.}}\csromangathastanza{\csromangathabat{กโรโต จียติ ปาปํ,}{จิรรตฺตาย อนฺตก.} \\ \csromangathabat{มาร นิพฺพินฺท พุทฺธมฺหา,}{อาสํ มากาสิ ภิกฺขุสุ.}}\csromangathastanza{\csromangathabat{อิติ มารํ อตชฺเชสิ,}{ภิกฺขุ เภสกฬาวเน.} \\ \csromangathabat{ตโต โส ทุมฺมโน ยกฺโข,}{ตตฺเถวนฺตรธายถาติ.}}}{มาร\-ตชฺชนีย\-สุตฺตํ นิฏฺฐิตํ ทสมํ.}
\nitthitammid{จูฬยมก\-วคฺโค นิฏฺฐิโต ปญฺจโม.}
\csromancenter{\textbf{ตสฺสุทฺทานํ}}
\setlength{\csromangathapagewidth}{0pt}
\setlength{\csromangathawakwidth}{0pt}
\csromangathameasuregroup{สาเลยฺย เวรญฺชทุเว จ ตุฏฺฐิ, \\ วีมํสกา โกสมฺพิ จ พฺราหฺมโณ,}{\csromangathabat{สาเลยฺย เวรญฺชทุเว จ ตุฏฺฐิ,}{จูฬมหาธมฺมสมาทานญฺจ.} \\ \csromangathabat{วีมํสกา โกสมฺพิ จ พฺราหฺมโณ,}{ทูสี จ มาโร ทสโม จ วคฺโค.} \\ สาเลยฺยวคฺโค นิฏฺฐิโต ปญฺจโม.}
\csromangathameasurewak{สาเลยฺยวคฺโค นิฏฺฐิโต ปญฺจโม.}
\csromangathasetleft{สาเลยฺย เวรญฺชทุเว จ ตุฏฺฐิ, \\ วีมํสกา โกสมฺพิ จ พฺราหฺมโณ,}
\csromangathagroup{\csromangathastanza{\csromangathabat{สาเลยฺย เวรญฺชทุเว จ ตุฏฺฐิ,}{จูฬมหาธมฺมสมาทานญฺจ.} \\ \csromangathabat{วีมํสกา โกสมฺพิ จ พฺราหฺมโณ,}{ทูสี จ มาโร ทสโม จ วคฺโค.}}\csromangathastanza{\csromangathawak{สาเลยฺยวคฺโค นิฏฺฐิโต ปญฺจโม.}}}

\vspace{\tipitakaparskip}
\csromancenter{อิทํ วคฺคานมุ\-ทฺทานํ}
\setlength{\csromangathapagewidth}{0pt}
\setlength{\csromangathawakwidth}{0pt}
\csromangathameasuregroup{มูลปริยาโย เจว, \\ กกโจ เจว โคสิงฺโค,}{\csromangathabat{มูลปริยาโย เจว,}{สีหนาโท จ อุตฺตโม.} \\ \csromangathabat{กกโจ เจว โคสิงฺโค,}{สาเลยฺโย จ อิเม ปญฺจ.} \\ \textbf{มูลปณฺณาสกํ สมตฺตํ.}}
\csromangathasetleft{มูลปริยาโย เจว, \\ กกโจ เจว โคสิงฺโค,}
\csromangathagroup{\csromangathastanza{\csromangathabat{มูลปริยาโย เจว,}{สีหนาโท จ อุตฺตโม.} \\ \csromangathabat{กกโจ เจว โคสิงฺโค,}{สาเลยฺโย จ อิเม ปญฺจ.} \\ \textbf{มูลปณฺณาสกํ สมตฺตํ.}}}

\orphannote{นมิสฺสติ (?)}

